%% =====================================================================
%% answers.tex　附录：章末习题 A 组答案与 B 组提示
%% 用法：在 document.tex 的 \appendix 之后 %% =====================================================================
%% answers.tex　附录：章末习题 A 组答案与 B 组提示
%% 用法：在 document.tex 的 \appendix 之后 %% =====================================================================
%% answers.tex　附录：章末习题 A 组答案与 B 组提示
%% 用法：在 document.tex 的 \appendix 之后 %% =====================================================================
%% answers.tex　附录：章末习题 A 组答案与 B 组提示
%% 用法：在 document.tex 的 \appendix 之后 \input{answers}（须在 \input{exercises} 之后，
%%      exercises.tex 中每道题都有标签 ex:<章键>:<组号>，本文件按标签引用题号）.
%% 章键：elim vs mat lmap det sol poly eig ip dec quad bil（依次为第一至十二章）.
%% =====================================================================

\providecommand{\ans}[2]{\par\smallskip\noindent\textbf{\ref{#1}.}\ #2\par}
\setlength{\emergencystretch}{3em}
\chapter{习题答案与提示}\label{app:answers}

本附录给出各章章末习题 A 组的答案与 B 组的提示. A 组的计算题给出最终结果，必要时附关键步骤；证明题只给出思路，完整的书写请读者自行完成. B 组只给出关键的一步或方向，希望读者在看提示之前先独立思考至少一刻钟. C 组一般不提供提示；正文或其他题目的答案中引用到的 C 组题例外，在该章末尾给出证明要点.

计算题的答案往往不唯一（例如基础解系、过渡矩阵的选取、特征向量的倍数），只要满足题目的全部要求，与这里的答案不同也是正确的. 检验的方法通常比答案本身更重要：代回原方程、计算迹与行列式、验证正交性.


\section{第一章　线性方程组与消元法}

\subsection*{A 组答案}

\ans{ex:elim:A1}{第 1、4 个是简化阶梯形；第 2 个是阶梯形但不是简化阶梯形（主元 $2$ 不为 $1$）；第 3、5 个不是阶梯形（第 3 个第二行主元在第一行主元左侧；第 5 个零行在非零行上方）.}

\ans{ex:elim:A2}{$\begin{pmatrix}3&0&-1&2\\0&1&4&0\\1&-1&0&5\end{pmatrix}$；反过来为 $x_1-2x_2=3$，$x_2+5x_3=-1$（未知量为 $x_1,x_2,x_3$，最后一列为常数项；$x_3$ 在第一个方程中的系数为 $0$）.}

\ans{ex:elim:A3}{$(x_1,x_2,x_3)=(1,2,3)$.}

\ans{ex:elim:A4}{$(x_1,x_2,x_3,x_4)=(1,-1,2,0)$.}

\ans{ex:elim:A5}{无解. $(r_3-r_1)-(r_2-2r_1)$ 给出 $(0,0,0\mid 2)$，即 $0=2$.}

\ans{ex:elim:A6}{主变量 $x_1,x_3$，自由变量 $x_2,x_4$. 全部解为 $x_1=4-2s-2t$，$x_2=s$，$x_3=1-t$，$x_4=t$（$s,t$ 任意）.}

\ans{ex:elim:A7}{$x_1=\frac27x_3+\frac37x_4$，$x_2=\frac57x_3+\frac47x_4$. 令 $x_3=7s$，$x_4=7t$，全部解为 $s(2,5,7,0)^\trans+t(3,4,0,7)^\trans$.}

\ans{ex:elim:A8}{$\begin{pmatrix}1&2&0&2\\0&0&1&-1\\0&0&0&0\end{pmatrix}$. 一种做法：$r_1\leftrightarrow r_2$；$r_2-2r_1$，$r_3-3r_1$；$-\frac14r_2$；$r_3+2r_2$；$r_1-r_2$.}

\ans{ex:elim:A9}{共四类：$\vv O$；$\begin{pmatrix}1&a\\0&0\end{pmatrix}$（$a$ 任意）；$\begin{pmatrix}0&1\\0&0\end{pmatrix}$；$\vv I_2$.}

\ans{ex:elim:A10}{由前两个方程得 $x=2$，$y=-1$，代入第三个得 $4-k=5$，$k=-1$.}

\ans{ex:elim:A11}{有解当且仅当 $b=3a$. 这时全部解为 $x=a-2t$，$y=t$（$t$ 任意）.}

\ans{ex:elim:A12}{$y=x^2-2x+3$.}

\ans{ex:elim:A13}{$y=1-2x+x^3$，即 $a=1$，$b=-2$，$c=0$，$d=1$.}

\ans{ex:elim:A14}{$30$ 元、$50$ 元、$80$ 元票分别为 $40$，$40$，$20$ 张.}

\ans{ex:elim:A15}{守恒方程 $3x_1=x_3$，$8x_1=2x_4$，$2x_2=2x_3+x_4$. 解为 $(x_1,x_2,x_3,x_4)=t(1,5,3,4)$，最小正整数解：$\mathrm{C_3H_8}+5\,\mathrm{O_2}\to3\,\mathrm{CO_2}+4\,\mathrm{H_2O}$.}

\ans{ex:elim:A16}{(1) 四个路口的方程为 $x_1-x_4=100$，$x_1-x_2=150$，$x_3-x_2=150$，$x_3-x_4=100$，全部解为 $x_1=t+100$，$x_2=t-50$，$x_3=t+100$，$x_4=t$. (2) 由 $x_2\ge0$ 得 $t\ge50$，故 $x_1$ 最少为 $150$. (3) 方程组无解：外部驶入总量 $1000$，驶出总量变为 $1100$，第一、三个方程之和减去第二、四个方程，得 $0=-100$. 守恒律是``进入路网的车辆总数等于离开的总数''.}

\ans{ex:elim:A17}{若 $a\ne0$，用第一个方程消去第二个方程中的 $x$，得 $(ad-bc)y=af-ce$. 解唯一当且仅当 $y$ 的系数 $ad-bc\ne0$. $a=0$ 时交换两个方程（或对 $c$ 讨论），同样得到条件 $ad-bc\ne0$.}

\ans{ex:elim:A18}{设 $\vv A\vv x_1=\vv A\vv x_2=\vv b$（这里先借用第三章的记号，也可以逐个方程写出），则 $\vv A(\vv x_1+t(\vv x_2-\vv x_1))=(1-t)\vv b+t\vv b=\vv b$. 若有两个不同的解，$t$ 取不同值给出不同的解，于是有无穷多个.}

\ans{ex:elim:A19}{第二个方程组的第一个方程等于第一个方程组两个方程之和，第二个方程相同，所以 (I) 的解都是 (II) 的解；反之 (I) 的第二个方程是 (II) 的第一个方程减去第二个方程. 两者的解集都是 $\{(t,t,3-2t)\}$.}

\ans{ex:elim:A20}{例如方程组 $x=0$，$y=0$. 同时做两步后得到 $x+y=0$，$x+y=0$，解集变大. 不矛盾：``同时''做的这两步合起来不是一次初等行变换，也不能被撤销（它把两行变成了相同的两行，丢失了信息）.}

\ans{ex:elim:A21}{有解当且仅当 $a=0$. 这时系数部分有 $2$ 个主元，$4$ 个未知量，有 $2$ 个自由变量.}

\subsection*{B 组提示}

\ans{ex:elim:B1}{用 $r_4-3r_1$，再加上 $r_2$，得 $(0,0,0,a-1\mid0)$；$r_3+r_2$ 得 $(0,0,a-1,0\mid b+1)$. $a\ne1$ 时解唯一；$a=1$，$b\ne-1$ 时无解；$a=1$，$b=-1$ 时有无穷多解：$(-1+s+t,\ 1-2s-2t,\ s,\ t)$.}

\ans{ex:elim:B2}{把第三个方程移到最前面再消元，得到的阶梯形后两行为 $(0,\lambda,-\lambda\mid3-\lambda)$ 与 $\big(0,0,-\lambda(\lambda+3)\mid-(\lambda+3)(\lambda-1)\big)$. $\lambda\ne0,-3$ 时有三个主元，解唯一；$\lambda=0$ 时倒数第二行变为 $(0,0,0\mid3)$，无解；$\lambda=-3$ 时最后一行为零行，只有两个主元，有无穷多解.}

\ans{ex:elim:B3}{齐次方程组只有零解说明简化阶梯形有 $n$ 个主元，即每一列都有主元. 而系数部分只有 $n$ 行，不可能出现 $(0,\ldots,0\mid d)$ 形的行.}

\ans{ex:elim:B4}{初等行变换得到的每一行都是原方程的线性组合. 无解时简化阶梯形中出现 $(0,\ldots,0\mid d)$，$d\ne0$，把这一行除以 $d$ 即可. 反之若组合得 $0=1$，任何解代入都会矛盾.}

\ans{ex:elim:B5}{依次做 $r_i+r_j$，$r_j-r_i$，$r_i+r_j$，最后把第 $j$ 行乘以 $-1$，逐步写出两行的变化即可.}

\ans{ex:elim:B6}{消元只用到四则运算，简化阶梯形的元素仍是有理数. 有解时令全部自由变量取 $0$，主变量由有理数算出. 解唯一时没有自由变量，唯一的解就是这个有理数解.}

\ans{ex:elim:B7}{把 $n$ 个方程相加，左边为 $0$，故必要条件是 $a_1+\cdots+a_n=0$. 它也充分：令 $x_n=t$，则 $x_k=t+a_k+a_{k+1}+\cdots+a_{n-1}$（$k<n$）.}

\ans{ex:elim:B8}{把第一个方程减第二个、加第三个……，交替相加. $n$ 为奇数时得 $2x_1=a_1-a_2+a_3-\cdots+a_n$，$x_1$ 唯一确定，其余依次确定；$n$ 为偶数时左边相消为 $0$.}

\ans{ex:elim:B9}{方程组为 $(\vv I-\vv A)\vv x=\vv d$（逐个写出即可），解得各部门总产出为 $(100,\ 200,\ 150)^\trans$.}

\ans{ex:elim:B10}{处理第 $k$ 列时，需要对下面 $n-k$ 行各做一次倍加，每次约 $n-k+1$ 次乘法，求和 $\sum_{k}(n-k)(n-k+1)\approx\frac13n^3$；回代只需约 $\frac12n^2$ 次.}

\ans{ex:elim:B11}{考察三个方程的系数部分化为阶梯形后的主元个数：$3$ 个时交于一点；$2$ 个且相容时交于一条直线；$2$ 个且出现 $0=d$ 时，三条交线方向相同（都平行于系数矩阵的解空间方向），没有公共点. 例子可以取 $z=0$，$x=0$，$x+z=1$.}

\ans{ex:elim:B12}{直接验证 $\vv A(\vv x+\vv y)=\vv A\vv x+\vv A\vv y$，$\vv A(k\vv x)=k\vv A\vv x$（或逐个方程验证）. 后半句：设 $\vv x_0$ 是一个解，任一解 $\vv x$ 满足 $\vv x-\vv x_0$ 是齐次方程组的解.}

\ans{ex:elim:B13}{开放题. 一种设定：未知量为三位律师各承办三类案件的件数，共 $9$ 个；已知条件给出 $3$ 个独立方程，故至少还缺 $6$ 条独立的信息. 关键是先写清楚未知量，再数主元个数.}

\ans{ex:elim:B14}{每行之和为零即 $\vv A(1,\ldots,1)^\trans=\zero$. 后半句由 B12 的结论.}

\ans{ex:elim:B15}{(I) 的每个解满足 (II) 的每个方程，因为后者是前者的线性组合；反之亦然. 逆命题对有解的方程组成立，证明要用到第六章的结论（齐次方程组同解当且仅当系数矩阵行向量组等价，第六章章末习题~\ref{ex:sol:B11}）.}

\ans{ex:elim:B16}{$\lambda\ne1,-2$ 时解唯一：$x=-\frac{\lambda+1}{\lambda+2}$，$y=\frac1{\lambda+2}$，$z=\frac{(\lambda+1)^2}{\lambda+2}$；$\lambda=1$ 时三个方程相同，有无穷多解；$\lambda=-2$ 时三个方程相加得 $0=3$，无解.}

\subsection*{C 组证明要点（正文引用的题目）}

\ans{ex:elim:C9}{设 $\vv R_1,\vv R_2$ 都是 $\vv M$ 的简化阶梯形. 它们行等价，以它们为系数矩阵的齐次方程组同解（定理~\ref{thm:sameSolution}）. 对列数 $n$ 归纳. $n=1$ 时，简化阶梯形是零列或 $\vv e_1$，由 $\vv M$ 是否为零唯一决定.
	
	设 $n>1$. 删去最后一列，得到的 $\vv R_1',\vv R_2'$ 仍是简化阶梯形，且仍行等价（用同一串初等行变换），由归纳假设 $\vv R_1'=\vv R_2'$. 于是两者在前 $n-1$ 列中的主元位置相同，设共有 $r$ 个主元. 再比较最后一列：
	\begin{enumerate}
		\item 若最后一列在 $\vv R_1$ 中是主元列，则 $\vv R_1$ 的第 $r+1$ 行为 $(0,\ldots,0,1)$，$\vv R_1\vv x=\zero$ 的解都有 $x_n=0$. 若它在 $\vv R_2$ 中不是主元列，令 $x_n=1$、其余自由变量取 $0$，可得 $\vv R_2\vv x=\zero$ 的一个解，其 $x_n=1$，与同解矛盾. 所以它在两者中同为主元列，都等于 $\vv e_{r+1}$.
		\item 若最后一列在两者中都不是主元列，令 $x_n=1$、其余自由变量取 $0$，两个方程组各有唯一一个这样的解，其主变量的值分别是两矩阵最后一列前 $r$ 个分量的相反数. 两方程组同解，故这两个解相同，最后一列相同.
\end{enumerate}}

\section{第二章　向量空间}

\subsection*{A 组答案}

\ans{ex:vs:A1}{(1) 是. (2) 不是：$\frac12\cdot\sqrt2=\frac{\sqrt2}2$ 的 $\sqrt2$ 系数不是整数，乘法不封闭. (3) 不是：不含 $1$，也不对除法封闭. (4) 不是：$\sqrt[3]2\cdot\sqrt[3]2=\sqrt[3]4$ 不在其中. (5) 不是：$\mathrm i\cdot\mathrm i=-1$ 不是纯虚数.}

\ans{ex:vs:A2}{(1) 是. (2) 不是：$(1,0)^\trans$ 的负向量不在其中. (3) 不是：不含零多项式，且两个 $n$ 次多项式之和可能降次. (4) 是. (5) 是. (6) 不是：不含零向量.}

\ans{ex:vs:A3}{(1) 线性相关：$(4,1,-1)^\trans=2(1,2,-1)^\trans+(2,-3,1)^\trans$. (2) 线性无关. (3) 线性相关：$\mathbb{R}[x]_3$ 是三维的，四个向量必相关. (4) 线性相关：全 $1$ 矩阵等于前两个之和.}

\ans{ex:vs:A4}{$\vv\beta=\frac54\vv\alpha_1+\frac14\vv\alpha_2-\frac14\vv\alpha_3-\frac14\vv\alpha_4$. （这四个向量两两``对应分量乘积之和''为零，每个系数都可以直接算出.）}

\ans{ex:vs:A5}{以三个向量为列作矩阵并消元：$r_2-r_1$ 得 $(0,2,-2)$，再 $r_3+\frac12r_2$ 得 $(0,0,t-1)$. $t=1$ 时出现自由变量，线性相关；$t\ne1$ 时线性无关.（学完第五章后，可用行列式 $2t-2$ 验证.）}

\ans{ex:vs:A6}{秩为 $3$，一个极大无关组为 $\vv\alpha_1,\vv\alpha_2,\vv\alpha_4$；$\vv\alpha_3=3\vv\alpha_1+\vv\alpha_2$，$\vv\alpha_5=\vv\alpha_1+\vv\alpha_2+\vv\alpha_4$.}

\ans{ex:vs:A7}{(1) 线性相关：三者之和为 $\zero$. (2) 线性无关. (3) 线性无关：设 $k_1(\vv\alpha_1+2\vv\alpha_2)+k_2(\vv\alpha_2+2\vv\alpha_3)+k_3(\vv\alpha_3+2\vv\alpha_1)=\zero$，由 $\vv\alpha_i$ 线性无关得 $k_1+2k_3=0$，$2k_1+k_2=0$，$2k_2+k_3=0$. 依次代入得 $k_1=-2k_3$，$k_2=4k_3$，$k_3=-8k_3$，故 $k_3=0$，从而 $k_1=k_2=0$.}

\ans{ex:vs:A8}{$1\cdot\zero=\zero$ 是系数不全为零的组合；$\vv\alpha-\vv\alpha=\zero$ 同理；$k\vv\alpha=\zero$ 且 $\vv\alpha\ne\zero$ 时由命题~\ref{prop:linspc}(5) 得 $k=0$.}

\ans{ex:vs:A9}{由定理~\ref{thm:relvandexp}，两个向量相关当且仅当其中一个可由另一个表出，即是另一个的倍数. 三个向量时只能说``其中一个是另两个的线性组合''，不一定是某一个的倍数，例如 $\vv e_1,\vv e_2,\vv e_1+\vv e_2$.}

\ans{ex:vs:A10}{(1) $\vv\alpha+\vv\alpha=1\vv\alpha+1\vv\alpha=(1+1)\vv\alpha$. (2) 两边乘 $k^{-1}$. (3) $-\vv\alpha$ 与 $\vv\alpha$ 之和为零，由负向量的唯一性.}

\ans{ex:vs:A11}{(1) $\dim=3$，基 $(-1,1,0,0)^\trans$，$(1,0,1,0)^\trans$，$(1,0,0,1)^\trans$. (2) $\dim=\frac{n(n+1)}2$，基 $\vv E_{ij}$（$i\le j$）. (3) $\dim=n^2-1$，基 $\vv E_{ij}$（$i\ne j$）与 $\vv E_{11}-\vv E_{ii}$（$i=2,\ldots,n$）. (4) $\dim=3$，基 $x-1$，$(x-1)^2$，$(x-1)^3$.}

\ans{ex:vs:A12}{(1) 是. (2) 不是：$(1,0,0)^\trans+(0,1,0)^\trans$ 不在其中. (3) 不是：对实数数乘不封闭. (4) 是.}

\ans{ex:vs:A13}{维数为 $3$，一组基为 $\vv\alpha_1,\vv\alpha_2,\vv\alpha_4$（$\vv\alpha_3=\vv\alpha_1+\vv\alpha_2$）.}

\ans{ex:vs:A14}{$2x^2-x+3=4+3(x-1)+2(x-1)^2$，坐标为 $(4,3,2)^\trans$.}

\ans{ex:vs:A15}{坐标为 $(-1,4,-5,4)^\trans$（从右下角元素开始依次确定）.}

\ans{ex:vs:A16}{(1) 是. (2) 不是：$\mathbb{Q}$ 对实数的数乘不封闭. (3) 作为 $\mathbb{R}$ 上的空间是 $2$ 维，作为 $\mathbb{C}$ 上的空间是 $1$ 维.}

\ans{ex:vs:A17}{例如 $a_0+a_1x+\cdots+a_{n-1}x^{n-1}\mapsto(a_0,\ldots,a_{n-1})^\trans$；$\begin{pmatrix}a&b\\c&d\end{pmatrix}\mapsto(a,b,c,d)^\trans$. 两者都是取坐标的映射，保持线性运算.}

\ans{ex:vs:A18}{$W_1+W_2$ 是直和（$(1,1,1)^\trans\notin W_1$，交为零），等于 $\mathbb{R}^3$，维数 $3$；$W_3\subset W_1$，$W_1+W_3=W_1$ 不是直和，维数 $2$.}

\ans{ex:vs:A19}{$\dim(W_1+W_2)=3$，$\dim(W_1\cap W_2)=1$：$(0,-1,1,1)^\trans=-(1,1,0,0)^\trans+(1,0,1,1)^\trans$ 同属两者. $2+2-1=3$，与维数公式一致.}

\ans{ex:vs:A20}{它们都可由一个含 $r$ 个向量的极大无关组线性表出，由基本引理任意 $r+1$ 个相关. 任意 $r$ 个无关的向量再添一个就相关，由引理~\ref{lem:addone} 的逆否形式，组中每个向量都可由它们线性表出.}

\subsection*{B 组提示}

\ans{ex:vs:B1}{设两者互不包含，取 $\vv\alpha\in W_1\setminus W_2$，$\vv\beta\in W_2\setminus W_1$，证明 $\vv\alpha+\vv\beta$ 不属于并集.}
\ans{ex:vs:B2}{取 (II) 的一个极大无关组，它含 $r$ 个向量且可由 (I) 的极大无关组表出，二者秩相同，用定理~\ref{thm:ndimbasis}(2) 在 (I) 张成的空间中论证.}
\ans{ex:vs:B3}{(III) 可由两者的极大无关组合并表出. 第二个不等式：左边的每个向量可由两组的极大无关组合并表出.}
\ans{ex:vs:B4}{新组含 $n$ 个向量，只需证明它能表出 $\vv\alpha_j$：$a_j\ne0$ 时 $\vv\alpha_j=\frac1{a_j}(\vv\beta-\sum_{i\ne j}a_i\vv\alpha_i)$；$a_j=0$ 时新组线性相关.}
\ans{ex:vs:B5}{按次数从高到低比较系数，证明线性无关；个数等于维数 $n$.}
\ans{ex:vs:B6}{$a_1,a_2$ 确定整个数列，映射 $\{a_n\}\mapsto(a_1,a_2)$ 是同构. 代入验证 $2^n$ 与 $(-1)^n$ 满足递推且线性无关. $a_1=a_2=1$ 时 $a_n=\frac13\big(2^n-(-1)^n\big)$.}
\ans{ex:vs:B7}{(1) 设组合为零，令 $x=0$ 并逐次求导，或除以 $\mathrm e^{3x}$ 后令 $x\to+\infty$. (2) $\sin(x+1)=\cos1\cdot\sin x+\sin1\cdot\cos x$. (3) 秩为 $3$：$\cos2x=2\cos^2x-1$.}
\ans{ex:vs:B8}{证明对加、减、乘封闭，再对除法用两次``分母有理化''. 线性无关可设 $a+b\sqrt2=-(c+d\sqrt2)\sqrt3$ 后平方，利用 $\sqrt2,\sqrt3,\sqrt6$ 都是无理数. 维数为 $4$.}
\ans{ex:vs:B9}{按维数分：$\{\zero\}$；过原点的直线；（$\mathbb{R}^3$ 中）过原点的平面；整个空间.}
\ans{ex:vs:B10}{用维数公式：$\dim(W_1\cap W_2)=k_1+k_2-\dim(W_1+W_2)\ge k_1+k_2-n>0$.}
\ans{ex:vs:B11}{$f=f(a)+(f-f(a))$. 交为零：常数多项式在 $a$ 处为零只能是零.}
\ans{ex:vs:B12}{$f(x)=\frac{f(x)+f(-x)}2+\frac{f(x)-f(-x)}2$；$\mathrm e^x=\cosh x+\sinh x$.}
\ans{ex:vs:B13}{对 $k$ 归纳，反复使用维数公式：$\dim(U\cap W)\ge\dim U+\dim W-n$.}
\ans{ex:vs:B14}{这 $n-1$ 个向量线性无关（用坐标验证），$\dim W=n-1$. 补空间可取 $\spn(\vv\alpha_n)$（需验证 $\vv\alpha_n\notin W$：$W$ 中向量的坐标交错和为零）.}
\ans{ex:vs:B15}{$\vv x=x_1(1,\ldots,1)^\trans+(\vv x-x_1(1,\ldots,1)^\trans)$. 两种分解都是直和；与 $V_1$ 互补的子空间有无穷多个，第九章会说明例~\ref{eg:pricedecomp} 中的 $R$ 是其中``垂直''的那一个.}
\ans{ex:vs:B16}{每个矩阵唯一地拆成``对角线及以上''与``对角线以下''两部分. $\dim W_1=\frac{n(n+1)}2$.}
\ans{ex:vs:B17}{(1) 设 $W$ 的基为 $\vv\varepsilon_1,\ldots,\vv\varepsilon_k$（$k<n$），扩充为 $\mathbb{R}^n$ 的基 $\vv\varepsilon_1,\ldots,\vv\varepsilon_n$. ``取 $\vv x$ 在这组基下的第 $n$ 个坐标''是 $x_1,\ldots,x_n$ 的一次齐次式（解一个线性方程组即可写出），它在 $W$ 上恒为零，又在 $\vv\varepsilon_n$ 处为 $1$，系数不全为零. (2) 条件是 $a_1+a_2t+\cdots+a_nt^{n-1}=0$，这是一个次数不超过 $n-1$ 的非零多项式，至多 $n-1$ 个根. (3) 例如 $t=2$：$\vv v(2)=(1,2,4)^\trans$ 不满足三个方程中的任何一个.}

\ans{ex:vs:B18}{对 $k$ 归纳. 已选好线性无关的 $\vv\varepsilon_1,\ldots,\vv\varepsilon_{k-1}$（$k\le n$）后，$U=\spn(\vv\varepsilon_1,\ldots,\vv\varepsilon_{k-1})$ 的维数 $k-1<n$，也是真子空间. 由覆盖定理取 $\vv\varepsilon_k\notin W_1\cup\cdots\cup W_s\cup U$；由引理~\ref{lem:addone}，$\vv\varepsilon_1,\ldots,\vv\varepsilon_k$ 仍线性无关. 选满 $n$ 个即得一组基.}

\section{第三章　矩阵}

\subsection*{A 组答案}

\ans{ex:mat:A1}{$\vv{AB}=\begin{pmatrix}5&9\\-2&4\end{pmatrix}$，$\vv{BA}=\begin{pmatrix}2&1&5\\-1&1&-2\\1&11&6\end{pmatrix}$，$\vv A^\trans\vv B^\trans=(\vv{BA})^\trans$. 由此也看出 $(\vv{AB})^\trans=\vv B^\trans\vv A^\trans\ne\vv A^\trans\vv B^\trans$（后者甚至与前者不同型）.}

\ans{ex:mat:A2}{$2(1,3)^\trans-(2,-1)^\trans+(0,4)^\trans=(0,11)^\trans$. $\vv e_i^\trans\vv A\vv e_j=a_{ij}$；$\vv A\vv e_j$ 是 $\vv A$ 的第 $j$ 列，$\vv e_i^\trans\vv A$ 是第 $i$ 行.}

\ans{ex:mat:A3}{$\begin{pmatrix}1&n\\0&1\end{pmatrix}$；$\begin{pmatrix}\lambda^n&n\lambda^{n-1}\\0&\lambda^n\end{pmatrix}$；$\begin{pmatrix}\cos n\theta&-\sin n\theta\\ \sin n\theta&\cos n\theta\end{pmatrix}$；第四个矩阵记为 $\vv N$，$\vv N^2=\vv E_{13}$，$n\ge3$ 时 $\vv N^n=\vv O$. 均可用归纳法证明.}

\ans{ex:mat:A4}{$(\vv A+\vv B)^2=\begin{pmatrix}10&8\\12&10\end{pmatrix}$，$\vv A^2+2\vv{AB}+\vv B^2=\begin{pmatrix}16&8\\12&4\end{pmatrix}$，不相等；$(\vv A+\vv B)(\vv A-\vv B)=\begin{pmatrix}-6&4\\-6&6\end{pmatrix}\ne\vv A^2-\vv B^2=\begin{pmatrix}0&4\\-6&0\end{pmatrix}$. 原因是 $\vv{AB}\ne\vv{BA}$.}

\ans{ex:mat:A5}{依次为
$\begin{pmatrix}1&3&-2\\-\frac32&-3&\frac52\\1&1&-1\end{pmatrix}$，
$\begin{pmatrix}1&-1&0\\0&1&-1\\0&0&1\end{pmatrix}$，
$\frac14\begin{pmatrix}3&-2&1\\-2&4&-2\\1&-2&3\end{pmatrix}$，
$\begin{pmatrix}0&0&0&\frac14\\0&0&\frac13&0\\0&\frac12&0&0\\1&0&0&0\end{pmatrix}$.}

\ans{ex:mat:A6}{(1) $\vv X=\begin{pmatrix}2&-23\\0&8\end{pmatrix}$. (2) $\vv X=\begin{pmatrix}1&-1\\-2&4\end{pmatrix}$（注意 $\vv X=\vv B\vv A^{-1}$，逆矩阵乘在右边）.}

\ans{ex:mat:A7}{$r_2-r_1$，$r_3-ar_1$，再 $r_3+r_2$，得阶梯形的后两行为 $(0,a-1,1-a)$ 与 $\big(0,0,(1-a)(a+2)\big)$. $a\ne1$ 且 $a\ne-2$ 时秩为 $3$；$a=1$ 时后两行全为零，秩为 $1$；$a=-2$ 时第二行非零、第三行为零，秩为 $2$.（学完第五章后，可用行列式 $-(a-1)^2(a+2)$ 验证.）}

\ans{ex:mat:A8}{秩为 $3$. 简化阶梯形的主元在第 $1,3,5$ 列，故第 $1,3,5$ 列构成极大无关组；第 $2$ 列 $=2\times$第 $1$ 列，第 $4$ 列 $=$ 第 $1$ 列 $+$ 第 $3$ 列.}

\ans{ex:mat:A9}{$\vv E(1,3)=\begin{pmatrix}0&0&1\\0&1&0\\1&0&0\end{pmatrix}$，逆为自身；$\vv E(2(5))=\diag(1,5,1)$，逆为 $\diag(1,\frac15,1)$；$\vv E(3,1(-2))=\begin{pmatrix}1&0&0\\0&1&0\\-2&0&1\end{pmatrix}$，逆为 $\vv E(3,1(2))$. $\vv E(3,1(-2))\vv A$ 把 $\vv A$ 第 $1$ 行的 $-2$ 倍加到第 $3$ 行；$\vv A\vv E(3,1(-2))$ 把 $\vv A$ 第 $3$ 列的 $-2$ 倍加到第 $1$ 列.}

\ans{ex:mat:A10}{$\begin{pmatrix}2&1\\4&3\end{pmatrix}=\vv E(2,1(2))\,\vv E(1,2(1))\,\vv E(1(2))=\begin{pmatrix}1&0\\2&1\end{pmatrix}\begin{pmatrix}1&1\\0&1\end{pmatrix}\begin{pmatrix}2&0\\0&1\end{pmatrix}$. 写法不唯一.}

\ans{ex:mat:A11}{$(\vv{AB})_{ij}=\sum_ka_{ik}b_{kj}$，$i>j$ 时每一项中 $k<i$ 使 $a_{ik}=0$ 或 $k\ge i>j$ 使 $b_{kj}=0$；$i=j$ 时只剩 $a_{ii}b_{ii}$. 可逆上三角矩阵的逆也是上三角矩阵：对 $(\vv A\mid\vv I)$ 作消元时只需用下面的行消去上面的行.}

\ans{ex:mat:A12}{(1) 用 $(\vv{AB})^\trans=\vv B^\trans\vv A^\trans$ 逐一验证. (2) $(\vv{AB})^\trans=\vv B^\trans\vv A^\trans=\vv{BA}$. (3) $(\vv C^\trans\vv C)^\trans=\vv C^\trans\vv C$.}

\ans{ex:mat:A13}{(1) $\vv x\vv y^\trans$ 的 $(i,i)$ 元为 $x_iy_i$，求和即得. (2) $\tr(\vv A\vv E_{ji})=\vv e_i^\trans\vv A\vv e_j=a_{ij}$，故每个 $a_{ij}=0$. (3) 两边取迹，由命题~\ref{prop:trace}，左边为 $0$，右边为 $n$.}

\ans{ex:mat:A14}{$\vv A(\vv A-\vv I)=2\vv I$，故 $\vv A^{-1}=\frac12(\vv A-\vv I)$；$(\vv A+2\vv I)(\vv A-3\vv I)=\vv A^2-\vv A-6\vv I=-4\vv I$，故 $(\vv A+2\vv I)^{-1}=-\frac14(\vv A-3\vv I)$. $\vv A+\vv I$ 不一定可逆：$\vv A=-\vv I$ 满足条件，$\vv A+\vv I=\vv O$.}

\ans{ex:mat:A15}{左乘 $\vv A^{-1}$. 反例：$\vv A=\begin{pmatrix}1&0\\0&0\end{pmatrix}$，$\vv B=\vv O$，$\vv C=\begin{pmatrix}0&0\\0&1\end{pmatrix}$.}

\ans{ex:mat:A16}{$\begin{pmatrix}\vv A^{-1}&\vv O\\ \vv O&\vv B^{-1}\end{pmatrix}$，$\begin{pmatrix}\vv O&\vv B^{-1}\\ \vv A^{-1}&\vv O\end{pmatrix}$. 所求逆矩阵为 $\begin{pmatrix}0&0&\frac12&0\\0&0&0&\frac15\\3&-2&0&0\\-1&1&0&0\end{pmatrix}$.}

\ans{ex:mat:A17}{记 $\vv B=\begin{pmatrix}2&1\\1&3\end{pmatrix}$，矩阵为 $\begin{pmatrix}\vv I&\vv B\\ \vv O&\vv I\end{pmatrix}$，其平方为 $\begin{pmatrix}\vv I&2\vv B\\ \vv O&\vv I\end{pmatrix}=\begin{pmatrix}1&0&4&2\\0&1&2&6\\0&0&1&0\\0&0&0&1\end{pmatrix}$，$n$ 次方为 $\begin{pmatrix}\vv I&n\vv B\\ \vv O&\vv I\end{pmatrix}$.}

\ans{ex:mat:A18}{$\vv Q\vv C=\begin{pmatrix}5600&8300\\6750&9800\end{pmatrix}$：第 $i$ 行是第 $i$ 个工厂每天的总成本与总销售额. 每天利润分别为 $2700$ 元与 $3050$ 元.}

\ans{ex:mat:A19}{$\vv A=\begin{pmatrix}0&1&1&0\\1&0&1&0\\1&1&0&1\\0&0&1&0\end{pmatrix}$，$\vv A^2=\begin{pmatrix}2&1&1&1\\1&2&1&1\\1&1&3&0\\1&1&0&1\end{pmatrix}$，$\vv A^3=\begin{pmatrix}2&3&4&1\\3&2&4&1\\4&4&2&3\\1&1&3&0\end{pmatrix}$. $(\vv A^2)_{ii}$ 等于城市 $i$ 的航线数；从 $1$ 到 $4$ 恰乘三次航班有 $(\vv A^3)_{14}=1$ 种走法（$1\to2\to3\to4$）.}

\ans{ex:mat:A20}{$\vv x^\trans\vv A\vv x$ 是一个数，等于它的转置 $\vv x^\trans\vv A^\trans\vv x=-\vv x^\trans\vv A\vv x$，故为 $0$. $(\vv A^2)^\trans=(\vv A^\trans)^2=(-\vv A)^2=\vv A^2$.}

\ans{ex:mat:A21}{$\vv P$ 是初等矩阵之积，$\vv{PA}$ 与 $\vv A$ 行等价，因而有相同的简化阶梯形（简化阶梯形的唯一性见第一章章末习题~\ref{ex:elim:C9}）.}

\ans{ex:mat:A22}{存在性由提示中的等式. 唯一性：若 $\vv S_1+\vv K_1=\vv S_2+\vv K_2$（$\vv S_i$ 对称，$\vv K_i$ 反对称），则 $\vv S_1-\vv S_2=\vv K_2-\vv K_1$ 既对称又反对称，只能是零矩阵.}

\subsection*{B 组提示}

\ans{ex:mat:B1}{设 $\vv A=\begin{pmatrix}a&b\\c&d\end{pmatrix}$，用 Cayley–Hamilton 型恒等式 $\vv A^2=(a+d)\vv A-(ad-bc)\vv I$（直接验证）. (1) $\vv A^2=\vv O$ 当且仅当 $a+d=0$ 且 $ad-bc=0$. (2) $\vv A=\pm\vv I$，或 $a+d=0$ 且 $ad-bc=-1$. (3) $\vv A=\vv O$，$\vv A=\vv I$，或 $a+d=1$ 且 $ad-bc=0$.}
\ans{ex:mat:B2}{比较 $\vv{AD}$ 与 $\vv{DA}$ 的 $(i,j)$ 元：$a_ix_{ij}=x_{ij}a_j$.}
\ans{ex:mat:B3}{由 $\vv{AE}_{ij}=\vv E_{ij}\vv A$ 比较元素，得 $\vv A$ 的非对角元为零，且 $a_{ii}=a_{jj}$.}
\ans{ex:mat:B4}{用二项式定理（可交换时成立）：$(\vv A+\vv B)^{p+q}$ 的每一项中 $\vv A$ 的次数 $\ge p$ 或 $\vv B$ 的次数 $\ge q$. 反例：$\vv A=\vv E_{12}$，$\vv B=\vv E_{21}$.}
\ans{ex:mat:B5}{等式：分块对角矩阵的行（列）向量组的极大无关组由两块各自的极大无关组拼成. 不等式：用初等变换把 $\vv A$、$\vv B$ 分别化为标准形，看非零子矩阵. 例子：$\vv A=\vv B=(0)$，$\vv C=(1)$.}
\ans{ex:mat:B6}{把右边与 $\vv I+\vv u\vv v^\trans$ 相乘，利用 $\vv v^\trans\vv u$ 是一个数.}
\ans{ex:mat:B7}{直接验证 $(\vv I_n+\vv{BA})\big(\vv I_n-\vv B(\vv I_m+\vv{AB})^{-1}\vv A\big)=\vv I_n$，展开后把 $\vv B\vv A\vv B$ 写成 $\vv B(\vv{AB})$.}
\ans{ex:mat:B8}{各列都是某个非零列 $\vv u$ 的倍数，倍数组成 $\vv v$. $\vv A^2=\vv u(\vv v^\trans\vv u)\vv v^\trans$，而 $\vv v^\trans\vv u=\tr\vv A$. 于是 $\vv A^k=(\tr\vv A)^{k-1}\vv A$.}
\ans{ex:mat:B9}{用相抵标准形 $\vv A=\vv P\diag(\vv I_r,\vv O)\vv Q$，把 $\vv I_r$ 拆成 $r$ 个 $\vv E_{ii}$. 后半句：秩一矩阵之和的秩不超过个数.}
\ans{ex:mat:B10}{取 $\vv B$ 为 $\vv A$ 列向量组的一个极大无关组排成的矩阵，$\vv C$ 的各列是 $\vv A$ 各列在这组向量下的坐标. 对第 A8 题，$\vv B$ 取第 $1,3,5$ 列，$\vv C$ 是简化阶梯形的前三行.}
\ans{ex:mat:B11}{全 $1$ 矩阵 $\vv J$ 满足 $\vv J^2=3\vv J$，故 $\vv J^n=3^{n-1}\vv J$. 第二个：$\vv N^3=\vv O$，$(2\vv I+\vv N)^n=2^n\vv I+n2^{n-1}\vv N+\binom n22^{n-2}\vv N^2$.}
\ans{ex:mat:B12}{对 $k$ 归纳，计算 $\vv N^k$ 的 $(i,j)$ 元 $\sum_l(\vv N^{k-1})_{il}n_{lj}$ 中哪些项可能非零.}
\ans{ex:mat:B13}{$\tr(\vv A^\trans\vv A)=\sum_{i,j}a_{ij}^2$. 复矩阵不成立：$\vv A=\begin{pmatrix}1&0\\\mathrm i&0\end{pmatrix}$ 满足 $\vv A^\trans\vv A=\vv O$；复矩阵应改用 $\bar{\vv A}^\trans\vv A$.}
\ans{ex:mat:B14}{置换矩阵就是对单位矩阵的行作排列得到的矩阵. $\vv{PA}$ 对 $\vv A$ 的行作相应排列，$\vv{AP}$ 对列作排列.}
\ans{ex:mat:B15}{秩为 $2$. 先用行变换化为阶梯形（记录 $\vv P$），再用列变换（记录 $\vv Q$）. 答案不唯一，可用 $\vv{PAQ}$ 直接检验.}
\ans{ex:mat:B16}{$(\vv I+\vv A)(\vv I-\vv A+\vv A^2-\cdots+(-1)^{k-1}\vv A^{k-1})=\vv I$. 若 $\vv A^2=2\vv A$，则 $(\vv I-\vv A)^2=\vv I$，可逆且逆为自身.}

\section{第四章　线性映射}

\subsection*{A 组答案}

\ans{ex:lmap:A1}{(1) 线性，核为 $\{\zero\}$. (2) 不线性：$\zero\mapsto(1,0)^\trans$. (3) 不线性：$(-1,0)^\trans$ 与 $(1,0)^\trans$ 之和的像不等于像之和. (4) 线性，核为 $\{0\}$. (5) 线性，核为常数多项式全体. (6) 线性，核为 $\{\vv O\}$. (7) 不线性：$\det(2\vv X)=4\det\vv X$.}

\ans{ex:lmap:A2}{第一个：核的基 $(-2,1,0)^\trans$，$(1,0,1)^\trans$；像的基 $(1,2)^\trans$；$3=2+1$. 第二个：核的基 $(-2,1,1,0)^\trans$，$(-1,-1,0,1)^\trans$；像的基 $(1,0,1)^\trans$，$(0,1,1)^\trans$；$4=2+2$.}

\ans{ex:lmap:A3}{$\begin{pmatrix}1&-2&1\\0&3&4\end{pmatrix}$.}

\ans{ex:lmap:A4}{(1) $\begin{pmatrix}0&-1\\1&0\end{pmatrix}$. (2) $\begin{pmatrix}0&1\\1&0\end{pmatrix}$，平方为 $\vv I$：反射两次回到原处. (3) $\frac12\begin{pmatrix}1&1\\1&1\end{pmatrix}$，平方等于自身：投影两次与投影一次相同. (4) $\begin{pmatrix}1&2\\0&1\end{pmatrix}$.}

\ans{ex:lmap:A5}{两组基下的矩阵都是 $\begin{pmatrix}0&1&0\\0&0&2\\0&0&0\end{pmatrix}$：$\mathscr D(x-1)=1$，$\mathscr D(x-1)^2=2(x-1)$.}

\ans{ex:lmap:A6}{$\begin{pmatrix}1&0&0&0\\0&0&1&0\\0&1&0&0\\0&0&0&1\end{pmatrix}$，它的平方是 $\vv I_4$.}

\ans{ex:lmap:A7}{过渡矩阵 $\vv P=\begin{pmatrix}-1&-1&-4\\1&0&6\\1&3&1\end{pmatrix}$. $\vv\alpha$ 在 $\vv\varepsilon$ 下的坐标为 $(-1,2,2)^\trans$，在 $\vv\eta$ 下的坐标为 $(-16,5,3)^\trans$. 可用 $\vv x=\vv P\vv x'$ 检验.}

\ans{ex:lmap:A8}{$\diag(3,1)$. 过渡矩阵 $\vv P=\begin{pmatrix}1&1\\1&-1\end{pmatrix}$，也可以直接看出 $\mathscr T\vv\eta_1=3\vv\eta_1$，$\mathscr T\vv\eta_2=\vv\eta_2$.}

\ans{ex:lmap:A9}{$\mathscr S\mathscr T(a\vv\alpha+b\vv\beta)=\mathscr S(a\mathscr T\vv\alpha+b\mathscr T\vv\beta)=a\mathscr{ST}\vv\alpha+b\mathscr{ST}\vv\beta$. 逆映射：设 $\vv\alpha=\mathscr T\vv\alpha'$，$\vv\beta=\mathscr T\vv\beta'$，则 $\mathscr T^{-1}(a\vv\alpha+b\vv\beta)=\mathscr T^{-1}\mathscr T(a\vv\alpha'+b\vv\beta')=a\mathscr T^{-1}\vv\alpha+b\mathscr T^{-1}\vv\beta$.}

\ans{ex:lmap:A10}{若 $\sum k_i\mathscr T\vv\alpha_i=\zero$，则 $\sum k_i\vv\alpha_i\in\ker\mathscr T=\{\zero\}$. 反之，若 $\mathscr T\vv\alpha=\zero$ 而 $\vv\alpha\ne\zero$，则线性无关组 $\{\vv\alpha\}$ 被映为相关组 $\{\zero\}$.}

\ans{ex:lmap:A11}{(1) 存在：前两个向量是基，$\mathscr T(x,y)^\trans=\big(\frac{x+y}2,\frac{x-y}2\big)^\trans$，它恰好满足第三个条件. (2) 不存在：$(2,0)^\trans=(1,1)^\trans+(1,-1)^\trans$，线性性要求 $\mathscr T(2,0)^\trans=(1,1)^\trans$.}

\ans{ex:lmap:A12}{(1) $5-3=2$. (2) 都不存在：像的维数 $\le3<5$；核的维数 $\ge5-3=2>0$.}

\ans{ex:lmap:A13}{$\mathscr T\vv\alpha=\zero\Rightarrow\mathscr{ST}\vv\alpha=\zero$；$\mathscr{ST}\vv\alpha=\mathscr S(\mathscr T\vv\alpha)$. 由前者 $\rk(\mathscr{ST})=\dim U-\dim\ker(\mathscr{ST})\le\dim U-\dim\ker\mathscr T=\rk\mathscr T$，由后者 $\rk(\mathscr{ST})\le\rk\mathscr S$.}

\ans{ex:lmap:A14}{$\tr$ 是满射（$\tr(c\vv E_{11})=c$），由秩—零化度定理 $\dim\ker\tr=n^2-1$.}

\ans{ex:lmap:A15}{$\tr(\vv P^{-1}\vv{AP})=\tr(\vv{AP}\vv P^{-1})=\tr\vv A$. 两矩阵迹分别为 $4$ 与 $3$，不相似.}

\ans{ex:lmap:A16}{$\vv P^{-1}\vv I\vv P=\vv I$. 由于 $\begin{pmatrix}1&1\\0&1\end{pmatrix}\ne\vv I$，它与 $\vv I$ 不相似.}

\ans{ex:lmap:A17}{矩阵为 $\begin{pmatrix}1&0&0\\1&1&1\\1&2&4\end{pmatrix}$. 依次做 $r_2-r_1$，$r_3-r_1$，$r_3-2r_2$，得 $\begin{pmatrix}1&0&0\\0&1&1\\0&0&2\end{pmatrix}$，有 $3$ 个主元，由定理~\ref{thm:invertible} 矩阵可逆，故 $\mathscr E$ 是同构.（第五章会看到，它的行列式是 Vandermonde 型的.）}

\ans{ex:lmap:A18}{$\mathscr T(x_1,x_2,x_3)^\trans=(x_1+x_3,\ x_1+x_2+2x_3,\ x_2+x_3)^\trans$. $\ker\mathscr T=\spn\big((1,1,-1)^\trans\big)$，$\imm\mathscr T=\spn\big((1,1,0)^\trans,(0,1,1)^\trans\big)$.}

\ans{ex:lmap:A19}{$\vv A_2(\vv A_1\vv x+\vv b_1)+\vv b_2=(\vv A_2\vv A_1)\vv x+(\vv A_2\vv b_1+\vv b_2)$. 仿射变换线性当且仅当把 $\zero$ 映为 $\zero$，即 $\vv b=\zero$.}

\ans{ex:lmap:A20}{第一列是 $\vv e_1$ 转 $\theta$ 后的位置 $(\cos\theta,\sin\theta)^\trans$，第二列同理. 复合``先转 $\beta$ 再转 $\alpha$''的矩阵是 $\vv R(\alpha)\vv R(\beta)$（定理~\ref{thm:composition}），比较 $(1,1)$ 元与 $(2,1)$ 元，得 $\cos(\alpha+\beta)=\cos\alpha\cos\beta-\sin\alpha\sin\beta$，$\sin(\alpha+\beta)=\sin\alpha\cos\beta+\cos\alpha\sin\beta$.}

\subsection*{B 组提示}

\ans{ex:lmap:B1}{取定基后 $\mathscr T\mapsto$ 矩阵是 $L(V,W)$ 到 $M_{m\times n}(\mathbb{F})$ 的同构.}
\ans{ex:lmap:B2}{$\mathscr T^2=\mathscr O$ 即 $\mathscr T(\mathscr T\vv\alpha)=\zero$. 取 $\imm\mathscr T$ 的基 $\vv\beta_1,\ldots,\vv\beta_r$，$\vv\beta_i=\mathscr T\vv\alpha_i$，证明 $\vv\beta_1,\vv\alpha_1,\ldots,\vv\beta_r,\vv\alpha_r$ 线性无关，再用 $\ker\mathscr T$ 中的向量补成基.}
\ans{ex:lmap:B3}{对 $\mathscr T|_{\ker(\mathscr{ST})}$ 用秩—零化度定理：它的像含于 $\ker\mathscr S$，核含于 $\ker\mathscr T$.}
\ans{ex:lmap:B4}{$\mathscr{ST}=\mathscr I$ 说明 $\mathscr T$ 是单射，有限维时由推论~\ref{cor:nullity} 它是双射，于是 $\mathscr S=\mathscr T^{-1}$.}
\ans{ex:lmap:B5}{$\mathscr{JD}(1)=0$. 不矛盾：B4 依赖有限维，$\mathbb{R}[x]$ 是无限维的.}
\ans{ex:lmap:B6}{$\tr(\vv{AX}-\vv{XA})=0$. 后半句：$\vv A^k$ 与 $\vv A$ 可交换.}
\ans{ex:lmap:B7}{逐个计算 $\vv A\vv E_{11}$ 等并取坐标. $\vv X\mapsto\vv{AX}$ 的矩阵为 $\begin{pmatrix}\vv A&\vv O\\ \vv O&\vv A\end{pmatrix}$（在给定的基序下）；$\vv X\mapsto\vv{XA}$ 的矩阵为 $\begin{pmatrix}a\vv I&c\vv I\\ b\vv I&d\vv I\end{pmatrix}$.}
\ans{ex:lmap:B8}{证明 $V=\imm\mathscr P\oplus\ker\mathscr P$：$\vv\alpha=\mathscr P\vv\alpha+(\vv\alpha-\mathscr P\vv\alpha)$，且 $\mathscr P$ 在 $\imm\mathscr P$ 上是恒等. 合并两者的基.}
\ans{ex:lmap:B9}{按提示分解，并验证两部分分别在 $\ker(\mathscr T\mp\mathscr I)$ 中；交为零是因为 $\mathscr T\vv\alpha=\vv\alpha=-\vv\alpha$ 推出 $\vv\alpha=\zero$.}
\ans{ex:lmap:B10}{对 $\mathscr T|_U$ 用秩—零化度定理.}
\ans{ex:lmap:B11}{在基 $\mathrm e^x\cos x$，$\mathrm e^x\sin x$ 下 $\mathscr D$ 的矩阵为 $\begin{pmatrix}1&1\\-1&1\end{pmatrix}$，逆为 $\frac12\begin{pmatrix}1&-1\\1&1\end{pmatrix}$，故 $\int\mathrm e^x\sin x\,\mathrm dx=\frac12\mathrm e^x(\sin x-\cos x)+C$.}
\ans{ex:lmap:B12}{平移的矩阵为 $\begin{pmatrix}1&0&a\\0&1&b\\0&0&1\end{pmatrix}$. 两种顺序的矩阵不相等：先转后移把原点送到 $(1,2)$，先移后转把原点送到 $(-2,1)$.}
\ans{ex:lmap:B13}{$\vv P^{-1}\vv A^k\vv P=(\vv P^{-1}\vv{AP})^k$；转置时用 $\vv Q=(\vv P^{-1})^\trans$.}
\ans{ex:lmap:B14}{两者秩都为 $1$，故相抵. 若相似，则 $\vv P^{-1}\vv E_{11}\vv P$ 是幂等的，而 $\vv E_{12}^2=\vv O\ne\vv E_{12}$. 相抵允许在定义域与陪域中各自换基，相似必须在同一个空间中用同一组基.}
\ans{ex:lmap:B15}{设 $\sum c_i\mathscr T^i\vv\alpha=\zero$，依次施以 $\mathscr T^{k-1},\mathscr T^{k-2},\ldots$. 所求矩阵是主对角线之上一条斜线全为 $1$、其余为 $0$ 的矩阵.}

\section{第五章　行列式}

\subsection*{A 组答案}

\ans{ex:det:A1}{$\tau(315426)=5$，奇排列；$\tau(42531)=7$，奇排列；第三个排列的逆序数为 $0+1+2+\cdots+(n-1)=\frac{n(n-1)}2$，$n\equiv0,1\pmod4$ 时为偶排列，$n\equiv2,3\pmod4$ 时为奇排列.}

\ans{ex:det:A2}{(1) $\tau(3421)=5$，带负号. (2) 按行号排好后列号排列为 $i\,3\,5\,4\,j$，$\{i,j\}=\{1,2\}$. $\tau(13542)=4$，故 $i=1$，$j=2$.}

\ans{ex:det:A3}{$-41$，$-16$，$48$. 第三个：把各列加到第一列后提出 $6$，再化为三角形得 $6\cdot2^3=48$.}

\ans{ex:det:A4}{$6$.}

\ans{ex:det:A5}{$(-1)^{\frac{n(n-1)}2}a_1a_2\cdots a_n$：只有一项非零，它的列号排列是 $n(n-1)\cdots1$.}

\ans{ex:det:A6}{第一列减去第二列、再加上第三列（或按多重线性拆成 $8$ 个行列式，其中有两列相同的都为零），得 $2\begin{vmatrix}a&b&c\\a_1&b_1&c_1\\a_2&b_2&c_2\end{vmatrix}$.}

\ans{ex:det:A7}{第一个行列式等于 $(x-1)^2(x+2)$，根为 $x=1$（二重）与 $x=-2$；第二个是 Vandermonde 行列式 $(2-1)(x-1)(x-2)$，根为 $x=1,2$.}

\ans{ex:det:A8}{$24$，$-3$，$\frac13$，$9$，$9$. 最后一个：$\vv A^*=3\vv A^{-1}$，故 $3\vv A^{-1}-2\vv A^*=-3\vv A^{-1}$，行列式为 $(-3)^3\cdot\frac13=-9$.}

\ans{ex:det:A9}{(1) $(\det\vv A)^2=\det(\vv A^\trans\vv A)=1$. (2) $(\det\vv A)^2=1$. (3) 若 $\det\vv A\ne0$，则 $\vv A$ 可逆，由 $\vv A^2=\vv A$ 得 $\vv A=\vv I$，矛盾.}

\ans{ex:det:A10}{$\begin{pmatrix}d&-b\\-c&a\end{pmatrix}$，$ad-bc\ne0$ 时逆为 $\frac1{ad-bc}\begin{pmatrix}d&-b\\-c&a\end{pmatrix}$. 第二个矩阵行列式为 $1$，伴随矩阵就是逆矩阵：$\begin{pmatrix}-24&18&5\\20&-15&-4\\-5&4&1\end{pmatrix}$.}

\ans{ex:det:A11}{$D=17$，$D_1=17$，$D_2=-17$，$D_3=34$，解为 $(1,-1,2)$.}

\ans{ex:det:A12}{面积为 $\frac12\left|\det\begin{pmatrix}3&1\\1&4\end{pmatrix}\right|=\frac{11}2$. 三点共线：$(2,3)^\trans$ 与 $(6,9)^\trans$ 张成的平行四边形面积为 $0$.}

\ans{ex:det:A13}{平行六面体体积为 $7$，四面体体积为 $\frac76$.}

\ans{ex:det:A14}{系数行列式为 $(5-\lambda)(6-\lambda)(4-\lambda)-4(4-\lambda)-4(6-\lambda)=-(\lambda-2)(\lambda-5)(\lambda-8)$，故 $\lambda=2,5,8$.}

\ans{ex:det:A15}{把第一行换成 $(1,1,1,1)$ 后按第一行展开即得所求之和，计算得 $16$.}

\ans{ex:det:A16}{点 $(x,y)$ 在直线上当且仅当第一行是后两行的``仿射组合''，即三行线性相关. 也可以按第一行展开，得到一个关于 $x,y$ 的一次方程，并验证两点满足它.}

\ans{ex:det:A17}{按定义，非零项必须在左上角 $\vv A$ 中取 $m$ 个元素、在右下角 $\vv B$ 中取 $n$ 个，符号相乘. 第二个等于 $(-1)^{mn}\det\vv A\det\vv B$：把后 $n$ 列逐一移到前面，共 $mn$ 次相邻列交换.}

\ans{ex:det:A18}{$5$（一般地，这类 $n$ 阶行列式等于 $n+1$，见 B1）.}

\ans{ex:det:A19}{$\vv A\vv A^*=(\det\vv A)\vv I$ 两边取行列式得 $\det\vv A^*=(\det\vv A)^{n-1}\ne0$；再由 $\vv A^*=(\det\vv A)\vv A^{-1}$ 得 $(\vv A^*)^{-1}=\frac1{\det\vv A}\vv A$；对 $\vv A^{-1}$ 用同一公式得 $(\vv A^{-1})^*=\det(\vv A^{-1})\vv A=\frac1{\det\vv A}\vv A$.}

\ans{ex:det:A20}{行列式为 $0$，线性相关（$(7,8,9)^\trans=2(4,5,6)^\trans-(1,2,3)^\trans$）；改成 $10$ 后行列式为 $-3$，线性无关.}

\ans{ex:det:A21}{$\det\vv A=\det\vv A^\trans=\det(-\vv A)=(-1)^n\det\vv A=-\det\vv A$，故 $\det\vv A=0$.}

\subsection*{B 组提示}

\ans{ex:det:B1}{由 $D_n=2D_{n-1}-D_{n-2}$，$D_1=2$，$D_2=3$，得 $D_n-D_{n-1}=1$，$D_n=n+1$.}
\ans{ex:det:B2}{第一个：各列加到第一列，提出 $a+(n-1)b$，各行减去第一行，得 $[a+(n-1)b](a-b)^{n-1}$. 第二个：先证递推 $D_n=(x-a)D_{n-1}+a(x+a)^{n-1}$（或利用转置后 $a$ 变 $-a$），得 $\frac12\big[(x+a)^n+(x-a)^n\big]$.}
\ans{ex:det:B3}{第 $i$ 列乘 $\frac1{a_i}$ 后从第一列减去，化为上三角形：$a_1\cdots a_n\big(a_0-\sum_{i=1}^n\frac1{a_i}\big)$.}
\ans{ex:det:B4}{计算 $\begin{pmatrix}\vv I&\zero\\ -\vv v^\trans&1\end{pmatrix}\begin{pmatrix}\vv I+\vv u\vv v^\trans&\vv u\\ \zero^\trans&1\end{pmatrix}\begin{pmatrix}\vv I&\zero\\ \vv v^\trans&1\end{pmatrix}$. B2 的第一个行列式即 $\det\big((a-b)\vv I+b\vv 1\vv 1^\trans\big)$.}
\ans{ex:det:B5}{把 $\vv{AB}$ 的两列按第三章``列是线性组合''的看法展开，用多重线性拆成若干个行列式，其中列号重复的为零. $\vv{BA}$ 的秩 $\le2<3$.}
\ans{ex:det:B6}{必要性：$\det\vv A\det\vv A^{-1}=1$，两者都是整数. 充分性：$\vv A^{-1}=\frac1{\det\vv A}\vv A^*$，$\vv A^*$ 是整数矩阵.}
\ans{ex:det:B7}{可逆时用 $\vv A^*=(\det\vv A)\vv A^{-1}$ 两次. 不可逆时 $\rk\vv A^*\le1$，$n\ge3$ 时 $(\vv A^*)^*=\vv O$；$n=2$ 时直接验证.}
\ans{ex:det:B8}{对分块矩阵先把第二块行加到第一块行，再把第一块列从第二块列减去，化为 $\begin{pmatrix}\vv A+\vv B&\vv O\\ \vv B&\vv A-\vv B\end{pmatrix}$.}
\ans{ex:det:B9}{左乘 $\begin{pmatrix}\vv I&\vv O\\ -\vv{CA}^{-1}&\vv I\end{pmatrix}$. 可交换时 $\det\vv A\det(\vv D-\vv{CA}^{-1}\vv B)=\det(\vv{AD}-\vv{ACA}^{-1}\vv B)=\det(\vv{AD}-\vv{CB})$.}
\ans{ex:det:B10}{两边取行列式：$\det\vv F^n=(\det\vv F)^n=(-1)^n$.}
\ans{ex:det:B11}{按第一行展开得 $D_n=(1+x^2)D_{n-1}-x^2D_{n-2}$，即 $D_n-D_{n-1}=x^2(D_{n-1}-D_{n-2})$，故 $D_n=1+x^2+\cdots+x^{2n}$.}
\ans{ex:det:B12}{把第二、三列分别乘 $\omega^k$，$\omega^{2k}$ 加到第一列，提出公因式；也可直接展开验证第二个等号.}
\ans{ex:det:B13}{从最后一行起，每一行减去上一行，化为对角线全为 $1$ 的上三角形.}
\ans{ex:det:B14}{圆（或直线）的方程可写成 $a(x^2+y^2)+bx+cy+d=0$，系数不全为零. 四点都满足它，即一个以 $(a,b,c,d)$ 为未知量的齐次方程组有非零解.}
\ans{ex:det:B15}{$\rk\vv A=n$ 时 $\vv A^*=(\det\vv A)\vv A^{-1}$ 可逆. $\rk\vv A\le n-2$ 时，由定理~\ref{thm:minorrank} 所有 $n-1$ 阶子式为零，$\vv A^*=\vv O$. $\rk\vv A=n-1$ 时有非零的 $n-1$ 阶子式，故 $\vv A^*\ne\vv O$；又 $\vv A\vv A^*=(\det\vv A)\vv I=\vv O$，$\vv A^*$ 的每一列都属于 $\ker\vv A$，而 $\dim\ker\vv A=n-(n-1)=1$（推论~\ref{cor:nullity}），故 $\rk\vv A^*=1$.}
\ans{ex:det:B16}{$x=1,2,3$ 时有两行相同；每行之和都是 $x+6$，把各列加到第一列可提出因子 $x+6$. $f$ 是四次多项式，首项系数为 $1$，故 $f(x)=(x+6)(x-1)(x-2)(x-3)$，根为 $1,2,3,-6$.}

\section{第六章　线性方程组解的结构}

\subsection*{A 组答案}

\ans{ex:sol:A1}{(1) 基础解系 $(2,5,7,0)^\trans$，$(3,4,0,7)^\trans$，通解为二者的任意线性组合. (2) 基础解系 $(-2,1,0,0)^\trans$，$(1,0,0,1)^\trans$.}

\ans{ex:sol:A2}{(1) $\vv x=\big(\tfrac54,-\tfrac14,0,0\big)^\trans+k_1(3,3,2,0)^\trans+k_2(-3,7,0,4)^\trans$. (2) $\vv x=\big(\tfrac12,0,0,0\big)^\trans+k_1(-1,2,0,0)^\trans+k_2(1,0,2,0)^\trans$.}

\ans{ex:sol:A3}{(1) $5-3=2$ 个. (2) 不一定，要看 $\rk(\vv A\mid\vv b)$ 是否为 $3$. (3) 一定有解（行满秩，推论~\ref{cor:sol-fullrank}），且有无穷多解.}

\ans{ex:sol:A4}{(1) $\rk\vv A=\rk(\vv A\mid\vv b)=2$，有解（且唯一）. (2) $\rk\vv A=1$，$\rk(\vv A\mid\vv b)=2$，无解.}

\ans{ex:sol:A5}{系数和为 $1$ 的 $2\vv\eta_1-\vv\eta_2$，$\frac12(\vv\eta_1+\vv\eta_2)$，$3\vv\eta_1-2\vv\eta_2$ 是 $\vv{Ax}=\vv b$ 的解；系数和为 $0$ 的 $\vv\eta_1-\vv\eta_2$ 是齐次方程组的解；$\vv\eta_1+\vv\eta_2$ 两者都不是（命题~\ref{prop:sol-affine}）.}

\ans{ex:sol:A6}{基础解系含 $3-2=1$ 个向量，$\vv\eta_1-\vv\eta_2=(0,4,4)^\trans$ 是非零解. 通解为 $(1,2,3)^\trans+k(0,1,1)^\trans$.}

\ans{ex:sol:A7}{例如 $x_1-x_2=1$，$x_1+x_3=3$. 做法：求与 $(1,1,-1)^\trans$ 垂直（即满足 $y_1+y_2-y_3=0$）的两个线性无关的系数向量，再代入特解求常数项. 答案不唯一.}

\ans{ex:sol:A8}{$\rk\vv A=2$（第三个法向量是前两个之和，且任意两个不平行），$\rk(\vv A\mid\vv b)=3$：三个平面两两相交成三条平行直线，没有公共点. 常数改为 $3$ 后 $\rk(\vv A\mid\vv b)=2$，三个平面交于一条直线.}

\ans{ex:sol:A9}{例如 $\vv A=\begin{pmatrix}1\\0\end{pmatrix}$，$\vv b=(0,1)^\trans$：$\vv{Ax}=\zero$ 只有零解，$\vv{Ax}=\vv b$ 无解. 若有解，两个解之差是齐次方程组的解，只能为零.}

\ans{ex:sol:A10}{$\vv X=\begin{pmatrix}3&2\\-2&-3\\1&3\end{pmatrix}$.}

\ans{ex:sol:A11}{$\vv X=\begin{pmatrix}a&b\\-a&-b\end{pmatrix}$（$a,b$ 任意），构成二维空间：$\vv X$ 的每一列都属于 $\ker\vv A=\spn\big((1,-1)^\trans\big)$.}

\ans{ex:sol:A12}{$\vv B$ 的非零列是 $\vv{Ax}=\zero$ 的非零解，故 $\det\vv A=7t+21=0$，$t=-3$. 这时 $\rk\vv A=2$，由 $\rk\vv A+\rk\vv B\le3$ 且 $\vv B\ne\vv O$ 得 $\rk\vv B=1$.}

\ans{ex:sol:A13}{通解 $(1,0,0)^\trans+s(-1,1,0)^\trans+t(-1,0,1)^\trans$. 取 $\vv e_1,\vv e_2,\vv e_3$：任一解 $(x_1,x_2,x_3)^\trans=x_1\vv e_1+x_2\vv e_2+x_3\vv e_3$，而 $x_1+x_2+x_3=1$.}

\ans{ex:sol:A14}{第一条直线：$(0,0,1)^\trans+t(1,1,-2)^\trans$. 第二条直线：$x_1-x_3=1$，$x_2=2$（答案不唯一）.}

\ans{ex:sol:A15}{方程组 $1+s=0$，$s=1+t$，$0=1-t$ 无解，两直线不相交（它们的方向也不平行，是异面直线）.}

\ans{ex:sol:A16}{方程 $x_1+x_3=30$，$x_1-x_2=0$，$x_2+x_3=30$（秩为 $2$）. 全部方案为 $x_1=x_2=s$，$x_3=30-s$. $x_1+x_3=30$ 与 $x_1-x_2=0$ 被唯一确定（它们是方程左边的线性组合）；$x_3$ 不被确定：它在基础解系 $(1,1,-1)^\trans$ 上的值不为零.}

\ans{ex:sol:A17}{$\vv L=\begin{pmatrix}1&0&0\\2&1&0\\1&2&1\end{pmatrix}$，$\vv U=\begin{pmatrix}2&1&0\\0&1&1\\0&0&1\end{pmatrix}$.}

\ans{ex:sol:A18}{$\rk(\vv A\mid\vv b)$ 不超过行数 $m=\rk\vv A$，故相容；解空间维数 $n-m>0$，有无穷多解.}

\ans{ex:sol:A19}{两个方向 $(1,-1,0)^\trans$，$(-1,1,2)^\trans$；求与二者都垂直的系数向量，即解齐次方程组，得 $(1,1,0)^\trans$. 平面方程为 $x+y=2$.}

\ans{ex:sol:A20}{解分别为 $(2,0)$ 与 $(1,1)$. 右端的相对变化约为 $\frac{0.0001}{2\sqrt2}\approx3.5\times10^{-5}$，解的相对变化为 $\frac{\sqrt2}2\approx0.71$，二者之比约为 $2\times10^4$.}

\subsection*{B 组提示}

\ans{ex:sol:B1}{把四个方程联立求解，得公共解 $k(-1,1,2,1)^\trans$.}
\ans{ex:sol:B2}{$\vv A^\trans\vv{Ax}=\zero\Rightarrow\vv x^\trans\vv A^\trans\vv A\vv x=0$，即 $\vv{Ax}$ 各分量的平方和为零. 两个齐次方程组同解，解空间维数相同，再用推论~\ref{cor:nullity}.}
\ans{ex:sol:B3}{$\vv A+\vv B$ 的列是 $(\vv A\mid\vv B)$ 的列的线性组合；后者的列由两个极大无关组合并表出.}
\ans{ex:sol:B4}{$(\vv A+\vv I)(\vv A-\vv I)=\vv O$ 给出 $\le n$（Sylvester 不等式的推论）；$(\vv A+\vv I)-(\vv A-\vv I)=2\vv I$ 给出 $\ge n$（B3）.}
\ans{ex:sol:B5}{由定理~\ref{thm:sylvester} 证明中的等式 $\rk\vv B=\rk(\vv{AB})+\dim(\ker\vv A\cap\imm\vv B)$.}
\ans{ex:sol:B6}{设 $k_0\vv\eta^*+\sum k_i(\vv\eta^*+\vv\xi_i)=\zero$，左乘 $\vv A$ 得 $(k_0+\sum k_i)\vv b=\zero$，再用 $\vv\xi_i$ 线性无关.}
\ans{ex:sol:B7}{$9$ 个未知数，$6$ 个方程的秩为 $5$，自由度为 $4$，至少再公开 $4$ 个单元格. 公开左上角的 $2\times2$ 块即可；但若公开的 $4$ 格中有一行的全部 $3$ 格，那一行有一格是多余的，全表仍不能确定.}
\ans{ex:sol:B8}{三行满足 $r_1+r_2=r_3$，故条件为 $b_1+b_2=b_3$.}
\ans{ex:sol:B9}{解空间是一维的，$(1,\ldots,1)^\trans$ 是非零解. 通解为 $k(1,\ldots,1)^\trans$.}
\ans{ex:sol:B10}{$\rk\vv A+\rk\vv B\le n$（B4 中的推论）且 $\ge\rk(\vv A+\vv B)=n$（B3）.}
\ans{ex:sol:B11}{同解时，$\vv B$ 的每一行在 $\vv A$ 的解空间上恒为零，由可确定性判据（定理~\ref{thm:sol-determinable} 的 (2)$\Rightarrow$(3)）它是 $\vv A$ 各行的线性组合.}
\ans{ex:sol:B12}{必要性：对 $\vv b=\vv e_1,\ldots,\vv e_m$ 分别取解，排成 $\vv C$. 充分性：$\vv{A}(\vv C\vv b)=\vv b$.}
\ans{ex:sol:B13}{$\vv A\vv A^*=(\det\vv A)\vv I=\vv O$，$\vv A^*$ 的每一列都在一维解空间中.}
\ans{ex:sol:B14}{$x_1-x_4$ 是第一、二个方程之差，被确定为 $-3$；后两个量都不被确定：找不到系数 $a,b,c$ 使 $a,b,c$ 倍的三个方程之和的左边恰为所求的量.}
\ans{ex:sol:B15}{对 $\vv B|_{\ker(\vv{AB})}$ 用秩—零化度定理：它的核是 $\ker\vv B$，像含于 $\ker\vv A$.}

\section{第七章　多项式}

\subsection*{A 组答案}

\ans{ex:poly:A1}{(1) 商 $x^3-x^2+1$，余式 $-x-2$. (2) 商 $2x^3+x^2+6x+7$，余式 $20$（$=f(2)$）.}

\ans{ex:poly:A2}{依次为 $1$，$1$，$100x-99$. 前两个用 $x^2\equiv1$，$x^2\equiv-1$；第三个设余式为 $ax+b$，由 $x^{100}-(ax+b)$ 以 $1$ 为二重根，在 $x=1$ 处函数值与导数都为零.}

\ans{ex:poly:A3}{$\omega^4=\omega$，$\omega^2=-1-\omega$，代入得 $(1-a)\omega+(b-a)=0$，而 $\omega$ 不是实数，故 $a=b=1$. 验证：$x^4+x^2+1=(x^2+x+1)(x^2-x+1)$.}

\ans{ex:poly:A4}{(1) $x+3$. (2) $x-1$（$x^3-2x^2+x=x(x-1)^2$，$x^2-1=(x-1)(x+1)$）.}

\ans{ex:poly:A5}{$f$ 与 $g$ 的公因式整除 $g-hf$；$f$ 与 $g-hf$ 的公因式整除 $(g-hf)+hf=g$. 两组公因式相同，最大公因式也相同.}

\ans{ex:poly:A6}{$x^2+1=(x-1)(x+1)+2$，故 $u=\frac12$，$v=-\frac12(x-1)$.}

\ans{ex:poly:A7}{有理根只可能是 $\pm1,\pm2,\pm\frac12$. 得 $2x^3-5x^2+x+2=(x-1)(x-2)(2x+1)$，有理根为 $1,2,-\frac12$.}

\ans{ex:poly:A8}{(1) 在 $\mathbb{Q}$ 与 $\mathbb{R}$ 上为 $(x^2+x+1)(x^2-x+1)$；在 $\mathbb{C}$ 上为 $(x-\omega)(x-\bar\omega)(x+\omega)(x+\bar\omega)$，$\omega=\mathrm e^{2\pi\mathrm i/3}$. (2) 在 $\mathbb{Q}$ 与 $\mathbb{R}$ 上为 $(x-1)(x+1)(x^2+x+1)(x^2-x+1)$；在 $\mathbb{C}$ 上为 $\prod_{k=0}^5(x-\mathrm e^{k\pi\mathrm i/3})$. (3) 在 $\mathbb{Q}$ 上不可约；在 $\mathbb{R}$ 上为 $(x-\sqrt[3]2)(x^2+\sqrt[3]2x+\sqrt[3]4)$；在 $\mathbb{C}$ 上为 $(x-\sqrt[3]2)(x-\sqrt[3]2\,\omega)(x-\sqrt[3]2\,\omega^2)$.}

\ans{ex:poly:A9}{$f(1)=f'(1)=0$，$f''(1)=20\ne0$，是二重根. 也可以连续两次用 $x-1$ 作综合除法.}

\ans{ex:poly:A10}{有重根当且仅当与导数 $3x^2-3$ 有公共根，即 $f(1)=0$ 或 $f(-1)=0$. $a=2$ 时重根为 $1$；$a=-2$ 时重根为 $-1$.}

\ans{ex:poly:A11}{$\sigma_1=2$，$\sigma_2=3$，$\sigma_3=4$. 所求依次为 $2$；$\sigma_1^2-2\sigma_2=-2$；$\sigma_2/\sigma_3=\frac34$.}

\ans{ex:poly:A12}{$(x^2-2x-1)(x^2+1)=x^4-2x^3-2x-1$. （$1+\sqrt2$ 的共轭 $1-\sqrt2$、$\mathrm i$ 的共轭 $-\mathrm i$ 也必须是根，见第七章章末习题~\ref{ex:poly:B10}.）}

\ans{ex:poly:A13}{$(x-2)(x^2-2x+2)=x^3-4x^2+6x-4$.}

\ans{ex:poly:A14}{$f(\vv A)=(\vv A-\vv I)^2=\vv O$. 于是 $\vv A^2=2\vv A-\vv I$，$\vv A^3=3\vv A-2\vv I$，$g(\vv A)=2\vv A-2\vv I=\begin{pmatrix}0&4\\0&0\end{pmatrix}$.}

\ans{ex:poly:A15}{$f(x)=2x^2-3x+1$（四个点恰好落在一条抛物线上，插值多项式的三次项系数为 $0$）.}

\ans{ex:poly:A16}{$x^n-1=(x-1)(x^{n-1}+x^{n-2}+\cdots+1)$. 一般地 $x^n-a^n=(x-a)(x^{n-1}+ax^{n-2}+\cdots+a^{n-1})$；也可由余数定理：$x=a$ 时 $x^n-a^n=0$.}

\ans{ex:poly:A17}{$f(x)=2-2(x-1)+2(x-1)^3+(x-1)^4$.}

\ans{ex:poly:A18}{设余式为 $ax+b$，由 $a+b=3$，$2a+b=5$ 得余式为 $2x+1$.}

\ans{ex:poly:A19}{$60.5v^2+55v-100=0$，即 $121v^2+110v-200=0$，正根 $v=\frac{10}{11}$，内部收益率 $r=10\%$.}

\ans{ex:poly:A20}{比较次数得 $\deg f=\deg g$，$f=hg$ 中 $h$ 为常数，且 $f\ne0$ 时 $h\ne0$.}

\ans{ex:poly:A21}{设 $h$ 是 $f_1,g_1$ 的公因式，则 $dh$ 是 $f,g$ 的公因式，故 $dh\mid d$，$h$ 是常数.}

\subsection*{B 组提示}

\ans{ex:poly:B1}{由 Bézout 等式 $uf+vg=1$，$u'f+v'h=1$，两式相乘.}
\ans{ex:poly:B2}{用提示中的等式反复降低指数，指数的变化与整数的辗转相除 $\gcd(m,n)=\gcd(m-n,n)$ 完全平行.}
\ans{ex:poly:B3}{$f/\gcd(f,f')=(x-1)(x+2)$. 一般地，$p_i^{k_i}$ 在 $f'$ 中恰好以 $p_i^{k_i-1}$ 出现（需要用到 $p_i$ 与 $p_i'$ 互素）.}
\ans{ex:poly:B4}{$p$ 不可约，由引理~\ref{lem:primediv} 即得. 反例：$p=x^2$，$f=x$.}
\ans{ex:poly:B5}{两边都是次数小于 $n$ 的多项式，在 $x_1,\ldots,x_n$ 处取值相同（$\ell_i(x_j)$ 的性质），用推论~\ref{cor:rootcount}.}
\ans{ex:poly:B6}{$f(x)-f(0)$ 以 $0,1,2,\ldots$ 为根，有无穷多个根.}
\ans{ex:poly:B7}{三次多项式可约当且仅当有一次因式，即有有理根；可能的有理根只有 $\pm1$，都不是根. 四次多项式可能分解为两个二次因式之积而没有有理根，例如 $(x^2+1)^2$.}
\ans{ex:poly:B8}{设余式为 $ax+b$，代入 $x=1,2$：$(2^n-1)x+2-2^n$.}
\ans{ex:poly:B9}{由 $f(2)=4$ 得 $c=1$，$f(x)=x^2$.}
\ans{ex:poly:B10}{$(x-a)^2-2b^2$ 是有理系数的，用它除 $f$，余式 $rx+s$ 在 $a+b\sqrt2$ 处为零，由 $\sqrt2$ 无理得 $r=s=0$.}
\ans{ex:poly:B11}{对整数 $n$，$f(n)$ 与 $f(n\bmod2)$ 的奇偶性相同（$n-m\mid f(n)-f(m)$）.}
\ans{ex:poly:B12}{依次为 $x^2+1$，$x-2$，$(x-2)^2$.}
\ans{ex:poly:B13}{$N'(v)=C_1+2C_2v+\cdots>0$. 由 $N(0)=-C_0<0$ 与 $N(1)=\sum C_i-C_0>0$，用零点定理.}
\ans{ex:poly:B14}{$\binom xk$ 在整数点取整数值；把 $f$ 写成 $\sum c_k\binom xk$，依次代入 $x=0,1,\ldots,n$ 解出 $c_k$，它们都是整数.}
\ans{ex:poly:B15}{核中的多项式有 $n$ 个不同的根而次数小于 $n$，只能是零.}
\ans{ex:poly:B16}{$\gcd(p,f)$ 整除 $p$，只能是 $1$ 或 $p$ 的常数倍.}

\section{第八章　特征值与特征向量}

\subsection*{A 组答案}

\ans{ex:eig:A1}{以下``属于 $\lambda$''后给出特征子空间的一组基，特征向量是它们的非零线性组合.
\begin{enumerate}
	\item $\lambda=3$：$(1,0)^\trans$；$\lambda=2$：$(1,-1)^\trans$.
	\item $\lambda=3$：$(1,1)^\trans$；$\lambda=-1$：$(1,-1)^\trans$.
	\item $\lambda=\mathrm i$：$(1,\mathrm i)^\trans$；$\lambda=-\mathrm i$：$(1,-\mathrm i)^\trans$（在 $\mathbb{R}$ 上没有特征值）.
	\item $\lambda=2$：$(1,1)^\trans$；$\lambda=3$：$(2,1)^\trans$.
	\item $\lambda=2$（二重）：$(1,0,0)^\trans$，$(0,1,1)^\trans$；$\lambda=0$：$(0,1,-1)^\trans$.
	\item $\lambda=5$：$(1,1,1)^\trans$；$\lambda=-1$（二重）：$(1,-1,0)^\trans$，$(1,0,-1)^\trans$.
\end{enumerate}}

\ans{ex:eig:A2}{除第 3 个（在 $\mathbb{R}$ 上没有特征值）外都可对角化，$\vv P$ 取 A1 中各特征向量为列. 例如第 4 个：$\vv P=\begin{pmatrix}1&2\\1&1\end{pmatrix}$，$\vv P^{-1}\vv{AP}=\diag(2,3)$；第 6 个：$\vv P=\begin{pmatrix}1&1&1\\1&-1&0\\1&0&-1\end{pmatrix}$，$\vv P^{-1}\vv{AP}=\diag(5,-1,-1)$.}

\ans{ex:eig:A3}{第 1 个可以（两个不同的特征值）；第 2 个不可以（$2$ 的几何重数为 $1$）；第 3 个不可以（$2$ 是三重特征值，$2\vv I-\vv A$ 的秩为 $1$，几何重数为 $2$）；第 4 个可以（$1$ 的代数重数与几何重数都是 $2$）.}

\ans{ex:eig:A4}{$\det\vv A=6$，$\tr\vv A=6$；$\vv A^{-1}$ 的特征值为 $1,\frac12,\frac13$；$\vv A^2-2\vv A+\vv I=(\vv A-\vv I)^2$ 的特征值为 $0,1,4$；$\det\vv A^*=(\det\vv A)^2=36$；$\det(\vv A^2+\vv I)=2\cdot5\cdot10=100$.}

\ans{ex:eig:A5}{特征多项式为 $\lambda^2-5\lambda+6$，特征值 $2,3$；第二种情形为 $\lambda^2-2\lambda+5$，特征值 $1\pm2\mathrm i$.}

\ans{ex:eig:A6}{$\det(\lambda\vv I-\vv A^\trans)=\det\big((\lambda\vv I-\vv A)^\trans\big)$. 特征向量一般不同：$\vv A=\begin{pmatrix}1&1\\0&2\end{pmatrix}$ 属于 $1$ 的特征向量为 $(1,0)^\trans$，$\vv A^\trans$ 属于 $1$ 的特征向量为 $(1,-1)^\trans$.}

\ans{ex:eig:A7}{$\lambda\vv I-\vv A$ 仍是三角矩阵，行列式为 $\prod(\lambda-a_{ii})$.}

\ans{ex:eig:A8}{若 $\vv{Ax}=\lambda\vv x$，$\vv x\ne\zero$，则 $\vv A^k\vv x=\lambda^k\vv x$. 由 $\vv A^k=\vv O$ 得 $\lambda^k=0$；由 $\vv A^2=\vv A$ 得 $\lambda^2=\lambda$；由 $\vv A^2=\vv I$ 得 $\lambda^2=1$.}

\ans{ex:eig:A9}{$\vv A=\vv P\diag(2,3)\vv P^{-1}$，$\vv P=\begin{pmatrix}1&2\\1&1\end{pmatrix}$，
$\vv A^n=\begin{pmatrix}2\cdot3^n-2^n&2^{n+1}-2\cdot3^n\\3^n-2^n&2^{n+1}-3^n\end{pmatrix}$.}

\ans{ex:eig:A10}{特征值 $1$ 与 $0.7$. 稳态分布 $\vv\pi=(\frac23,\frac13)^\trans$. $\vv P^n\to\begin{pmatrix}\frac23&\frac23\\ \frac13&\frac13\end{pmatrix}$：两列相同，说明无论开始时顾客在哪一家，长期看都以同样的概率分布在两家，初始状态被``遗忘''了.}

\ans{ex:eig:A11}{$\vv A\vv\xi=(-1,a+2,b+1)^\trans=\lambda(1,1,-1)^\trans$，得 $\lambda=-1$，$a=-3$，$b=0$.}

\ans{ex:eig:A12}{特征多项式 $(\lambda-1)^2(\lambda+1)$. 可对角化当且仅当 $\rk(\vv I-\vv A)=1$，即 $x=-1$.}

\ans{ex:eig:A13}{$\vv A^2=\begin{pmatrix}a^2+bc&b(a+d)\\c(a+d)&d^2+bc\end{pmatrix}$，代入直接验证.}

\ans{ex:eig:A14}{依次为 $(x-1)(x-2)$，$(x-1)^2(x-2)$，$x^3$，$x(x-3)$.}

\ans{ex:eig:A15}{相似. 取 $\vv P=\begin{pmatrix}0&1\\1&0\end{pmatrix}$（交换两个基向量的次序），$\vv P^{-1}\begin{pmatrix}1&1\\0&1\end{pmatrix}\vv P=\begin{pmatrix}1&0\\1&1\end{pmatrix}$.}

\ans{ex:eig:A16}{$\vv A\vv e_1=a_{11}\vv e_1$，$\vv A\vv e_2=a_{12}\vv e_1+a_{22}\vv e_2$，……，$\vv A\vv e_k$ 只含 $\vv e_1,\ldots,\vv e_k$.}

\ans{ex:eig:A17}{特征根 $2,-1$，$a_k=\frac13\big(2^k-(-1)^k\big)$.}

\ans{ex:eig:A18}{矩阵 $\begin{pmatrix}0&1.1\\0.6&0.5\end{pmatrix}$ 的特征值为 $1.1$ 与 $-0.6$. 长期年增长率为 $10\%$；属于 $1.1$ 的特征向量为 $(1,1)^\trans$，幼体与成年个体之比稳定在 $1:1$.}

\ans{ex:eig:A19}{特征多项式 $\lambda^2-(\tr\vv A)\lambda+\det\vv A$ 没有实根，当且仅当判别式 $(\tr\vv A)^2-4\det\vv A<0$.}

\ans{ex:eig:A20}{投影：方向 $(1,2)^\trans$ 上的向量不动（特征值 $1$），垂直方向 $(2,-1)^\trans$ 上的向量变为零（特征值 $0$），矩阵 $\frac15\begin{pmatrix}1&2\\2&4\end{pmatrix}$. 反射：$(1,2)^\trans$ 属于 $1$，$(2,-1)^\trans$ 属于 $-1$，矩阵 $\frac15\begin{pmatrix}-3&4\\4&3\end{pmatrix}$.}

\subsection*{B 组提示}

\ans{ex:eig:B1}{$\vv A\vv u=(\vv v^\trans\vv u)\vv u$；$\rk\vv A=1$，$0$ 的几何重数为 $n-1$. 特征值为 $\vv v^\trans\vv u$ 与 $0$（$n-1$ 重）. $\vv v^\trans\vv u=0$ 时 $0$ 是 $n$ 重特征值，但几何重数只有 $n-1$.}
\ans{ex:eig:B2}{$\vv A\vv B\vv p=\vv B\vv A\vv p=\lambda\vv B\vv p$，而 $V_\lambda$ 是一维的. 后半句：用 Lagrange 插值构造多项式 $f$ 使 $f(\lambda_i)$ 等于 $\vv B$ 的对应特征值.}
\ans{ex:eig:B3}{$\vv P^{-1}f(\vv A)\vv P=f(\vv D)$. 逆命题不成立：$\vv A=\begin{pmatrix}0&1\\0&0\end{pmatrix}$ 不可对角化，$\vv A^2=\vv O$ 可以.}
\ans{ex:eig:B4}{$\vv A|_W$ 的最小多项式整除 $\vv A$ 的最小多项式，用定理~\ref{thm:minpolydiag}.}
\ans{ex:eig:B5}{取 $\vv v\ne\zero$，证明 $\vv v,\vv{Av}$ 线性无关，在这组基下 $\vv A$ 的矩阵为 $\begin{pmatrix}0&-1\\1&0\end{pmatrix}$.}
\ans{ex:eig:B6}{设 $\vv{Ax}=\lambda\vv x$，取 $|x_i|$ 最大的分量，看第 $i$ 个方程 $(\lambda-a_{ii})x_i=\sum_{j\ne i}a_{ij}x_j$.}
\ans{ex:eig:B7}{$\vv A(1,\ldots,1)^\trans=s(1,\ldots,1)^\trans$. 可逆时 $s\ne0$，$\vv A^{-1}$ 每行之和为 $\frac1s$.}
\ans{ex:eig:B8}{特征值为 $a+b\omega^k+c\omega^{2k}$（$k=0,1,2$），行列式是三者之积.}
\ans{ex:eig:B9}{特征根 $2$ 是二重的，通项形如 $(c_1+c_2k)2^k$，得 $a_k=(k+1)2^k$. $k2^k$ 来自 Jordan 块的幂（第八章选读节）.}
\ans{ex:eig:B10}{对三角化（或可对角化）的情形看对角元. 特征值为 $0,\sqrt3,-\sqrt3$.}
\ans{ex:eig:B11}{$x^3-x=x(x-1)(x+1)$，最小多项式整除它.}
\ans{ex:eig:B12}{$2$ 是二重特征值，需 $\rk(2\vv I-\vv A)=1$，条件为 $a=0$（$b,c$ 任意）.}
\ans{ex:eig:B13}{转移矩阵 $\begin{pmatrix}0&0&1\\ \frac12&0&0\\ \frac12&1&0\end{pmatrix}$，稳态分布 $\frac15(2,1,2)^\trans$.}
\ans{ex:eig:B14}{特征值 $\pm1$，$\mathrm e^{t\vv A}=\begin{pmatrix}\cosh t&\sinh t\\ \sinh t&\cosh t\end{pmatrix}$. 对 $t$ 求导验证.}
\ans{ex:eig:B15}{稳态分布 $(\frac13,\frac13,\frac13)^\trans$；特征值 $1,\frac14,\frac14$，偏离部分每一步缩小为原来的 $\frac14$.}
\ans{ex:eig:B16}{几何重数 $n-\rk(\lambda\vv I-\vv A)$ 不超过代数重数 $k$.}
\ans{ex:eig:B17}{得分（按判决 $1,2,3,4$）约为：$d=0.5$ 时 $(0.169,0.212,0.265,0.354)$；$d=0.85$ 时 $(0.128,0.183,0.261,0.428)$；$d=0.95$ 时 $(0.119,0.175,0.258,0.448)$；$d=1$ 时恰为 $\frac1{35}(4,6,9,16)$. 排序不变；$d$ 越大，随机跳转越少，得分越集中到引证链末端的判决 $4$ 上.}

\subsection*{C 组证明要点（正文引用的题目）}

\ans{ex:eig:C1}{直接计算两个分块乘积：
	\[\begin{pmatrix}\vv I_m&-\vv A\\ \vv O&\lambda\vv I_n\end{pmatrix}\begin{pmatrix}\lambda\vv I_m&\vv A\\ \vv B&\vv I_n\end{pmatrix}=\begin{pmatrix}\lambda\vv I_m-\vv{AB}&\vv O\\ \lambda\vv B&\lambda\vv I_n\end{pmatrix},\]\[
	\begin{pmatrix}\vv I_m&\vv O\\ -\vv B&\lambda\vv I_n\end{pmatrix}\begin{pmatrix}\lambda\vv I_m&\vv A\\ \vv B&\vv I_n\end{pmatrix}=\begin{pmatrix}\lambda\vv I_m&\vv A\\ \vv O&\lambda\vv I_n-\vv{BA}\end{pmatrix}.\]
	记中间矩阵的行列式为 $D$. 对两式取行列式（第五章章末习题~\ref{ex:det:A17}），得 $\lambda^nD=\lambda^n\det(\lambda\vv I_m-\vv{AB})$ 与 $\lambda^nD=\lambda^m\det(\lambda\vv I_n-\vv{BA})$，因此 $\lambda^n\det(\lambda\vv I_m-\vv{AB})=\lambda^m\det(\lambda\vv I_n-\vv{BA})$. 两边作为 $\lambda$ 的多项式相等，因子 $\lambda^n$、$\lambda^m$ 只影响根 $0$ 的重数，所以两个特征多项式的非零根（连同重数）相同.}

\section{第九章　内积空间}

\subsection*{A 组答案}

\ans{ex:ip:A1}{$\|\vv\alpha\|=\|\vv\beta\|=3$，$\langle\vv\alpha,\vv\beta\rangle=4$，夹角 $\arccos\frac49\approx63.6^\circ$，$\|\vv\alpha-\vv\beta\|=\sqrt{10}$.}

\ans{ex:ip:A2}{(1) 是. (2) 不是：$\vv x=(1,-1)^\trans$ 时 $\langle\vv x,\vv x\rangle=-2$. (3) 是：$\langle\vv x,\vv x\rangle=x_1^2+(x_1+x_2)^2$. (4) 不是：$\langle\vv e_2,\vv e_2\rangle=-1$.}

\ans{ex:ip:A3}{(1) $\frac1{\sqrt3}(1,1,1)^\trans$，$\frac1{\sqrt2}(0,1,-1)^\trans$，$\frac1{\sqrt6}(-2,1,1)^\trans$. (2) $\frac1{\sqrt2}(1,1,0,0)^\trans$，$\frac1{\sqrt6}(1,-1,2,0)^\trans$，$\frac1{2\sqrt3}(-1,1,1,3)^\trans$.}

\ans{ex:ip:A4}{解 $x_1+x_2+x_3=0$，$x_1-x_3=0$，得 $\pm\frac1{\sqrt6}(1,-2,1)^\trans$.}

\ans{ex:ip:A5}{逐对验证内积；坐标为 $\langle\vv\alpha,\vv e_i\rangle$，即 $(5,4,-2)^\trans$.}

\ans{ex:ip:A6}{两个生成向量已经正交，投影为 $\frac32(1,1,0,0)^\trans+\frac72(0,0,1,1)^\trans=\big(\frac32,\frac32,\frac72,\frac72\big)^\trans$，
	
	距离为 $$\|(-\frac12,\frac12,-\frac12,\frac12)^\trans\|=1$$.}

\ans{ex:ip:A7}{正规方程 $5a+10b=15$，$10a+30b=42$，得 $y=0.6+1.2x$，残差平方和为 $1.6$.}

\ans{ex:ip:A8}{第二列须与第一、三列正交且为单位向量：$a=c$，$b=a$，$3a^2=1$. 故 $a=b=c=\frac1{\sqrt3}$ 或 $a=b=c=-\frac1{\sqrt3}$.}

\ans{ex:ip:A9}{展开 $\big\langle\sum\vv\alpha_i,\sum\vv\alpha_j\big\rangle$，交叉项都为零.}

\ans{ex:ip:A10}{由三角不等式 $\|\vv\alpha\|\le\|\vv\alpha-\vv\beta\|+\|\vv\beta\|$，再交换 $\vv\alpha,\vv\beta$.}

\ans{ex:ip:A11}{(1) 对 $(a_1,\ldots,a_n)$ 与 $(1,\ldots,1)$ 用 Cauchy–Schwarz 不等式. (2) 对 $(\sqrt{a_1},\ldots,\sqrt{a_n})$ 与 $\big(\frac1{\sqrt{a_1}},\ldots,\frac1{\sqrt{a_n}}\big)$ 用之.}

\ans{ex:ip:A12}{$\|1\|=1$，$\|x\|=\frac1{\sqrt3}$，$\langle1,x\rangle=\frac12$，夹角 $\frac\pi6$. 正交化并单位化得 $1$，$\sqrt3(2x-1)$.}

\ans{ex:ip:A13}{$\ker\vv A$ 的标准正交基：$\frac1{\sqrt3}(1,-1,1,0)^\trans$，$\frac1{\sqrt{15}}(-1,-2,-1,3)^\trans$；行空间的标准正交基：$\frac1{\sqrt3}(1,1,0,1)^\trans$，$\frac1{\sqrt{15}}(-2,1,3,1)^\trans$. 两组之间两两内积为零，维数 $2+2=4$，故互为正交补.}

\ans{ex:ip:A14}{$\vv Q=\begin{pmatrix}\frac1{\sqrt2}&\frac1{\sqrt6}\\ \frac1{\sqrt2}&-\frac1{\sqrt6}\\0&\frac2{\sqrt6}\end{pmatrix}$，$\vv R=\begin{pmatrix}\sqrt2&\frac1{\sqrt2}\\0&\sqrt{\frac32}\end{pmatrix}$.}

\ans{ex:ip:A15}{$\vv H=\frac19\begin{pmatrix}7&-4&-4\\-4&1&-8\\-4&-8&1\end{pmatrix}$；$\vv H\vv u=-\vv u$；$(2,-1,0)^\trans$ 与 $\vv u$ 正交，$\vv H(2,-1,0)^\trans=(2,-1,0)^\trans$.}

\ans{ex:ip:A16}{第一个行列式为 $1$，是旋转，旋转角 $\arctan\frac43\approx53.1^\circ$；第二个行列式为 $-1$，是关于直线 $y=\frac x2$（方向 $(2,1)^\trans$）的反射.}

\ans{ex:ip:A17}{设计矩阵只有一列 $\vv x$，正规方程为 $(\vv x^\trans\vv x)b=\vv x^\trans\vv y$.}

\ans{ex:ip:A18}{投影为 $2(1,1,1)^\trans$，距离为 $\|(-1,0,1)^\trans\|=\sqrt2$.}

\ans{ex:ip:A19}{过渡矩阵的第 $j$ 列是新基第 $j$ 个向量在旧（标准正交）基下的坐标，坐标的标准内积等于向量的内积，故各列标准正交.}

\ans{ex:ip:A20}{内积打分：$6,1,4$，排序为甲（第一条）$>$ 丙 $>$ 乙. 余弦相似度：约 $0.926$，$0.140$，$0.943$，排序为丙 $>$ 甲 $>$ 乙. 第一条法条更长，内积占了便宜.}

\ans{ex:ip:A21}{$\vv P^\trans=(\vv u^\trans)^\trans\vv u^\trans=\vv P$，$\vv P^2=\vv u(\vv u^\trans\vv u)\vv u^\trans=\vv P$. $\vv{Px}=\langle\vv x,\vv u\rangle\vv u$ 正是定理~\ref{thm:projformula} 中 $k=1$ 的公式. 一般非零 $\vv u$ 时投影矩阵为 $\dfrac{\vv u\vv u^\trans}{\vv u^\trans\vv u}$.}

\subsection*{B 组提示}

\ans{ex:ip:B1}{展开两个平方相减. 取 $\vv x=(1,0)^\trans$，$\vv y=(0,1)^\trans$ 检验平行四边形法则 $\|\vv x+\vv y\|^2+\|\vv x-\vv y\|^2=2\|\vv x\|^2+2\|\vv y\|^2$，左边为 $8$，右边为 $4$.}
\ans{ex:ip:B2}{对 $\vv\alpha-\sum\langle\vv\alpha,\vv e_i\rangle\vv e_i$ 用勾股定理.}
\ans{ex:ip:B3}{证明 $\vv x-\vv{Px}$ 与 $\imm\vv P$ 正交：$\langle\vv x-\vv{Px},\vv{Py}\rangle=\vv y^\trans\vv P^\trans(\vv I-\vv P)\vv x=0$. 只有幂等时是沿 $\ker\vv P$ 向 $\imm\vv P$ 的``斜投影''，例如 $\begin{pmatrix}1&1\\0&0\end{pmatrix}$.}
\ans{ex:ip:B4}{目标函数是 $\|\vv y-\vv A\vv\beta\|_W^2$，其中 $\langle\vv x,\vv y\rangle_W=\vv x^\trans\vv W\vv y$ 是加权内积；残差在这个内积下与各列正交.}
\ans{ex:ip:B5}{对 $\ln y=\ln c+kt$ 作最小二乘，约得 $k\approx0.30$，$c\approx2.0$. 它最小化的是对数残差的平方和，相当于对相对误差加权.}
\ans{ex:ip:B6}{$\vv G=\vv A^\trans\vv A$（$\vv A$ 以 $\vv\alpha_i$ 为列），用第十一章命题~\ref{prop:gram} 或直接 $\vv x^\trans\vv G\vv x=\|\vv A\vv x\|^2$. $k=2$ 时 $\det\vv G=\|\vv\alpha\|^2\|\vv\beta\|^2\sin^2\theta$.}
\ans{ex:ip:B7}{距离是 $\vv x_0-\vv x_1$（$\vv x_1$ 为超平面上任一点）在法向 $\vv a$ 上投影的长度.}
\ans{ex:ip:B8}{设 $\vv{Qz}=\lambda\vv z$，考察 $\bar{\vv z}^\trans\vv Q^\trans\vv Q\vv z$.}
\ans{ex:ip:B9}{用积化和差逐对计算积分. 投影为 $2\sin x-\sin2x$（偶函数的系数为零）.}
\ans{ex:ip:B10}{全部最小二乘解为 $\hat{\vv x}+\ker\vv A$，把 $\hat{\vv x}$ 分解到 $\imm\vv A^\trans\oplus\ker\vv A$，用勾股定理.}
\ans{ex:ip:B11}{$\|\vv x-\vv y\|^2=2\|\vv x\|^2-2\langle\vv x,\vv y\rangle$，直接计算 $\vv H\vv x=\vv x-2\vv u\vv u^\trans\vv x$.}
\ans{ex:ip:B12}{对中心化的数据用 A17，再代回截距；正规方程的第一个方程就是 $\bar y=a+b\bar x$.}
\ans{ex:ip:B13}{保持夹角就保持正交. 把一组标准正交基的像归一化，证明 $\|\mathscr T\vv e_i\|$ 都相等（考察 $\vv e_i\pm\vv e_j$ 的像是否正交）.}
\ans{ex:ip:B14}{截距为第一组均值 $4$，虚拟变量系数为两组均值之差 $8-4=4$.}
\ans{ex:ip:B15}{$\vv R(\alpha)\vv R(\beta)=\vv R(\alpha+\beta)$；$\vv F(\alpha)\vv F(\beta)=\vv R(\alpha-\beta)$；计算 $\vv R(\alpha)\vv F(\beta)$ 与 $\vv F(\beta)\vv R(\alpha)$ 比较.}

\section{第十章　矩阵分解}

\subsection*{A 组答案}

\ans{ex:dec:A1}{（$\vv Q$ 的列依次对应所列特征值，答案不唯一.）
\begin{enumerate}
	\item $\vv Q=\frac1{\sqrt2}\begin{pmatrix}1&1\\1&-1\end{pmatrix}$，对角矩阵 $\diag(3,-1)$.
	\item $\vv Q=\frac1{\sqrt2}\begin{pmatrix}1&1\\1&-1\end{pmatrix}$，对角矩阵 $\diag(2,4)$.
	\item 特征值 $4,0,1$，特征向量 $\frac1{\sqrt2}(1,1,0)^\trans$，$\frac1{\sqrt2}(1,-1,0)^\trans$，$(0,0,1)^\trans$.
	\item 特征值 $3,0,0$，特征向量 $\frac1{\sqrt3}(1,1,1)^\trans$，$\frac1{\sqrt2}(1,-1,0)^\trans$，$\frac1{\sqrt6}(1,1,-2)^\trans$（后两个须在特征子空间内正交化）.
\end{enumerate}}

\ans{ex:dec:A2}{$\begin{pmatrix}3&-1\\-1&3\end{pmatrix}=2\cdot\frac12\begin{pmatrix}1&1\\1&1\end{pmatrix}+4\cdot\frac12\begin{pmatrix}1&-1\\-1&1\end{pmatrix}$. 两个投影矩阵之和为 $\vv I$，乘积为 $\vv O$.}

\ans{ex:dec:A3}{$\vv A=1\cdot\frac12\begin{pmatrix}1&1\\1&1\end{pmatrix}+2\cdot\frac12\begin{pmatrix}1&-1\\-1&1\end{pmatrix}=\frac12\begin{pmatrix}3&-1\\-1&3\end{pmatrix}$.}

\ans{ex:dec:A4}{属于 $1$ 的特征子空间是 $(0,1,1)^\trans$ 的正交补，故 $\vv A=\vv I-2\vv P$，$\vv P$ 是向 $\spn\big((0,1,1)^\trans\big)$ 的投影：$\vv A=\begin{pmatrix}1&0&0\\0&0&-1\\0&-1&0\end{pmatrix}$.}

\ans{ex:dec:A5}{特征值 $4,1,0$. 最大值 $4$，在 $\pm\frac1{\sqrt2}(1,1,0)^\trans$ 处取到；最小值 $0$，在 $\pm\frac1{\sqrt2}(1,-1,0)^\trans$ 处取到.}

\ans{ex:dec:A6}{矩阵为 $\begin{pmatrix}1&2\\2&1\end{pmatrix}$，最大值 $3$（在 $\pm\frac1{\sqrt2}(1,1)$ 处），最小值 $-1$（在 $\pm\frac1{\sqrt2}(1,-1)$ 处）.}

\ans{ex:dec:A7}{(1) $\sigma_1=3$，$\sigma_2=2$；$\vv U=\begin{pmatrix}0&1\\-1&0\end{pmatrix}$，$\vv\Sigma=\diag(3,2)$，$\vv V=\begin{pmatrix}0&1\\1&0\end{pmatrix}$. (2) $\vv U=\vv V=\frac1{\sqrt2}\begin{pmatrix}1&1\\1&-1\end{pmatrix}$，$\vv\Sigma=\diag(2,0)$. (3) $\vv U=(1)$，$\vv\Sigma=(5,\ 0)$，$\vv V=\frac15\begin{pmatrix}3&-4\\4&3\end{pmatrix}$.}

\ans{ex:dec:A8}{$\vv A^\trans\vv A=\begin{pmatrix}1&1\\1&2\end{pmatrix}$，特征值 $\frac{3\pm\sqrt5}2$，奇异值为 $\frac{\sqrt5+1}2$ 与 $\frac{\sqrt5-1}2$（黄金分割比及其倒数）；特征值是 $1,1$.}

\ans{ex:dec:A9}{对称正定矩阵的奇异值分解就是谱分解. $\vv A_1=4\cdot\frac12\begin{pmatrix}1&1\\1&1\end{pmatrix}=\begin{pmatrix}2&2\\2&2\end{pmatrix}$，误差 $\|\vv A-\vv A_1\|_F=\sigma_2=2$.}

\ans{ex:dec:A10}{$\vv A=9\vv P_1+\vv P_2$，$\vv P_1=\frac12\begin{pmatrix}1&1\\1&1\end{pmatrix}$，$\vv P_2=\frac12\begin{pmatrix}1&-1\\-1&1\end{pmatrix}$. $\vv A^{10}=9^{10}\vv P_1+\vv P_2$；$\vv B=3\vv P_1+\vv P_2=\begin{pmatrix}2&1\\1&2\end{pmatrix}$.}

\ans{ex:dec:A11}{$\vv A=\vv Q\vv O\vv Q^\trans=\vv O$. 反例：$\begin{pmatrix}0&1\\0&0\end{pmatrix}$.}

\ans{ex:dec:A12}{$\vv A^2=\vv A^\trans\vv A=\vv I$，特征值满足 $\lambda^2=1$. 它是关于特征子空间 $V_1$ 的``镜面反射''：$V_1$ 中的向量不动，$V_{-1}=V_1^\perp$ 中的向量反向.}

\ans{ex:dec:A13}{由第八章章末习题~\ref{ex:eig:C1}（取 $\vv B=\vv A^\trans$），或直接：若 $\vv A^\trans\vv A\vv x=\lambda\vv x$，$\lambda\ne0$，则 $\vv{AA}^\trans(\vv{Ax})=\lambda\vv{Ax}$ 且 $\vv{Ax}\ne\zero$；两边的特征子空间维数相同.}

\ans{ex:dec:A14}{$\tr(\vv A^\trans\vv A)=\sum_{i,j}a_{ij}^2$；它也等于 $\vv A^\trans\vv A$ 的特征值之和 $\sum\sigma_i^2$.}

\ans{ex:dec:A15}{$\det\vv A=\det\vv U\det\vv\Sigma\det\vv V^\trans=\pm\sigma_1\cdots\sigma_n$；$\vv A^{-1}=\vv V\vv\Sigma^{-1}\vv U^\trans$.}

\ans{ex:dec:A16}{只有第二个可以：能被正交矩阵对角化的实矩阵必须对称（谱定理的逆命题）.}

\ans{ex:dec:A17}{$\vv S=\frac13\begin{pmatrix}10&8\\8&10\end{pmatrix}$，特征值 $6$ 与 $\frac23$. 第一主成分方向 $\frac1{\sqrt2}(1,1)^\trans$，解释方差比例 $\frac6{6+2/3}=90\%$.}

\ans{ex:dec:A18}{$\kappa\big(\diag(10,1)\big)=10$；$\begin{pmatrix}1&1\\1&1.01\end{pmatrix}$ 的特征值约为 $2.005$ 与 $0.004988$，$\kappa\approx402$. 二者行列式都等于 $0.01$ 的 $\diag(10,0.001)$ 条件数为 $10^4$，与后者不同：行列式不决定条件数.}

\ans{ex:dec:A19}{$40\times(800+600+1)=56040$ 个数，原图 $480000$ 个，约为原来的 $11.7\%$. $k\ge343$ 时存储量超过原图.}

\ans{ex:dec:A20}{由谱定理 $\vv A$ 相似于 $\diag(\lambda_1,\ldots,\lambda_n)$，秩不变. 奇异值的平方是 $\vv A^\trans\vv A$ 的非零特征值，个数为 $\rk(\vv A^\trans\vv A)=\rk\vv A$.}

\ans{ex:dec:A21}{取正交矩阵 $\vv Q$ 使 $\vv Q^\trans\vv{AQ}=\diag(\lambda_1,\ldots,\lambda_n)$，令 $\vv y=\vv Q^\trans\vv x$，则 $\vv x^\trans\vv A\vv x=\sum\lambda_iy_i^2$，且 $\vv x\ne\zero\iff\vv y\ne\zero$. 若某个 $\lambda_i\le0$，取 $\vv y=\vv e_i$；若全为正，则 $\vv y\ne\zero$ 时和为正.}

\subsection*{B 组提示}

\ans{ex:dec:B1}{$\vv A$ 的每个特征子空间都是 $\vv B$-不变的，$\vv B$ 限制在其上仍是对称的，在其中取 $\vv B$ 的标准正交特征向量.}
\ans{ex:dec:B2}{存在性：$\vv B=\vv Q\diag(\sqrt{\lambda_i})\vv Q^\trans$. 唯一性：$\vv B$ 的每个特征子空间（属于 $\mu$）都含于 $\vv A$ 的属于 $\mu^2$ 的特征子空间，而 $\mu\ge0$ 由 $\mu^2$ 唯一确定.}
\ans{ex:dec:B3}{在 $\vv q_1^\perp=\spn(\vv q_2,\ldots,\vv q_n)$ 中重复定理~\ref{thm:rayleigh} 的证明.}
\ans{ex:dec:B4}{$\|\vv{Ax}\|^2=\vv x^\trans\vv A^\trans\vv A\vv x\le\lambda_{\max}(\vv A^\trans\vv A)\|\vv x\|^2$，在 $\vv v_1$ 处取等.}
\ans{ex:dec:B5}{$\vv{Ax}=\lambda\vv x$ 时 $|\lambda|\,\|\vv x\|=\|\vv{Ax}\|$，再用 B4 及其``最小值''版本.}
\ans{ex:dec:B6}{$\vv A^\trans\vv A$ 的非零特征值为 $3,1$（$\vv A\vv A^\trans=\begin{pmatrix}2&1\\1&2\end{pmatrix}$ 更易算），$\sigma_1=\sqrt3$，$\sigma_2=1$. 先求 $\vv U$ 的列：$\frac1{\sqrt2}(1,1)^\trans$，$\frac1{\sqrt2}(1,-1)^\trans$，再用 $\vv v_i=\vv A^\trans\vv u_i/\sigma_i$，第三个 $\vv v_3=\frac1{\sqrt3}(1,1,-1)^\trans$.}
\ans{ex:dec:B7}{$\vv A^+\vv b$ 满足正规方程（在 $\vv V$、$\vv U$ 的坐标下逐分量验证），且属于行空间 $\imm\vv A^\trans$，由第九章章末习题~\ref{ex:ip:B10} 它是范数最小的一个.}
\ans{ex:dec:B8}{$\vv J$ 的特征值为 $p,0,\ldots,0$，故矩阵的特征值为 $1+(p-1)\rho$ 与 $1-\rho$（$p-1$ 重）. 第一主成分 $\frac1{\sqrt p}(1,\ldots,1)^\trans$，比例 $\frac{1+(p-1)\rho}p$，$\rho\to1$ 时趋于 $1$，$\rho\to0$ 时趋于 $\frac1p$.}
\ans{ex:dec:B9}{协方差矩阵的第一主成分几乎就是第一个指标的方向（方差 $100$ 远大于 $1$）；相关系数为 $0.6$，相关系数矩阵的第一主成分是 $\frac1{\sqrt2}(1,1)^\trans$. 量纲改变方差，从而改变主成分.}
\ans{ex:dec:B10}{$\vv A-\vv A_k=\sum_{i>k}\sigma_i\vv u_i\vv v_i^\trans$，它的最大奇异值是 $\sigma_{k+1}$.}
\ans{ex:dec:B11}{$\vv v_i$ 是 $\vv A^\trans\vv A$ 的单重特征值的单位特征向量，只能差一个符号；$\vv u_i=\vv A\vv v_i/\sigma_i$ 随之确定.}
\ans{ex:dec:B12}{设 $\vv A=\begin{pmatrix}a&b\\c&d\end{pmatrix}$，比较 $\vv A^\trans\vv A$ 与 $\vv{AA}^\trans$ 的元素，得 $b^2=c^2$ 与 $(a-d)(b-c)=0$.}
\ans{ex:dec:B13}{第一句由第八章章末习题~\ref{ex:eig:B6}（特征值为实数）. 第二句取 $\vv x=\vv e_i$ 用定理~\ref{thm:rayleigh}.}
\ans{ex:dec:B14}{$k=1$ 时每份文书只剩一个数，查询与四份文书中三份合同类文书的相似度都是 $1$，与第四份为 $0$——同 $k=2$ 时的结果，但``合同''内部的差别完全消失，信息损失更多.}
\ans{ex:dec:B15}{$\vv x^\trans(\vv A-a\vv I)\vv x\ge0$ 对一切 $\vv x$ 成立当且仅当 $\vv A-a\vv I$ 的特征值 $\lambda_i-a\ge0$.}

\section{第十一章　二次型}

\subsection*{A 组答案}

\ans{ex:quad:A1}{(1) $\begin{pmatrix}1&2&0\\2&-3&-1\\0&-1&0\end{pmatrix}$，秩 $3$. (2) 对角元为 $0$、非对角元为 $\frac12$ 的三阶矩阵，秩 $3$. (3) 元素全为 $1$ 的三阶矩阵，秩 $1$. (4) 对角元为 $1$、非对角元为 $\frac12$ 的 $n$ 阶矩阵，即 $\frac12(\vv I+\vv J)$，秩 $n$（特征值 $\frac{n+1}2$ 与 $\frac12$，都不为零）.}

\ans{ex:quad:A2}{$x_1^2+2x_2^2-2x_1x_2+6x_2x_3$.}

\ans{ex:quad:A3}{(1) $(x_1+x_2)^2+(x_2+3x_3)^2-4x_3^2$，令 $y_1=x_1+x_2$，$y_2=x_2+3x_3$，$y_3=x_3$，得 $y_1^2+y_2^2-4y_3^2$. (2) $(x_1-2x_2)^2-x_3^2$，令 $y_1=x_1-2x_2$，$y_2=x_3$，$y_3=x_2$，得 $y_1^2-y_2^2$. (3) $x_2(x_1-x_3)$，令 $y_1=\frac12(x_1+x_2-x_3)$，$y_2=\frac12(x_2-x_1+x_3)$，$y_3=x_3$，得 $y_1^2-y_2^2$. 替换都不唯一.}

\ans{ex:quad:A4}{特征值 $2,5,1$，$\vv Q=\begin{pmatrix}1&0&0\\0&\frac1{\sqrt2}&\frac1{\sqrt2}\\0&\frac1{\sqrt2}&-\frac1{\sqrt2}\end{pmatrix}$，标准形 $2y_1^2+5y_2^2+y_3^2$. 曲面是椭球面.}

\ans{ex:quad:A5}{(1) $p=2$，$q=1$，符号差 $1$，规范形 $z_1^2+z_2^2-z_3^2$. (2)(3) 都是 $p=q=1$，符号差 $0$，规范形 $z_1^2-z_2^2$.}

\ans{ex:quad:A6}{(1) 正定（顺序主子式 $2,1,3$）. (2) 不正定（$\Delta_2=-3$）. (3) 正定（顺序主子式 $5,1,1$）.}

\ans{ex:quad:A7}{矩阵 $\begin{pmatrix}2&1&0\\1&1&\frac t2\\0&\frac t2&1\end{pmatrix}$，$\Delta_3=1-\frac{t^2}2>0$，即 $-\sqrt2<t<\sqrt2$.}

\ans{ex:quad:A8}{(1) 特征值 $3,-1$，标准方程 $3u^2-v^2=1$，双曲线；$u$ 轴沿 $\frac1{\sqrt2}(1,-1)$，$v$ 轴沿 $\frac1{\sqrt2}(1,1)$. (2) 特征值 $4,2$，标准方程 $\frac{u^2}2+\frac{v^2}4=1$，椭圆；沿 $\frac1{\sqrt2}(1,1)$ 的半轴为 $\sqrt2$，沿 $\frac1{\sqrt2}(1,-1)$ 的半轴为 $2$.}

\ans{ex:quad:A9}{(1) 合同（秩都为 $3$，$p$ 都为 $2$）：取 $\vv C$ 使 $\vv C\vv e_1=\frac1{\sqrt3}\vv e_3$，$\vv C\vv e_2=\sqrt5\,\vv e_1$，$\vv C\vv e_3=\sqrt3\,\vv e_2$. (2) 不合同：$p=1$ 与 $p=2$. (3) 合同：$\vv C=\frac1{\sqrt2}\begin{pmatrix}1&1\\1&-1\end{pmatrix}$.}

\ans{ex:quad:A10}{$\vv A^*=(\det\vv A)\vv A^{-1}$ 是对称矩阵，特征值为 $\frac{\det\vv A}{\lambda_i}>0$.}

\ans{ex:quad:A11}{驻点 $(0,0)$ 与 $(1,1)$. Hessian $\begin{pmatrix}6x&-3\\-3&6y\end{pmatrix}$：在 $(0,0)$ 处不定，是鞍点；在 $(1,1)$ 处正定，取极小值 $-1$.}

\ans{ex:quad:A12}{顺序主子式 $2,3,4$，正定. $\vv L=\begin{pmatrix}\sqrt2&0&0\\-\frac1{\sqrt2}&\sqrt{\frac32}&0\\0&-\sqrt{\frac23}&\frac2{\sqrt3}\end{pmatrix}$.}

\ans{ex:quad:A13}{$\vv L=\begin{pmatrix}2&0&0\\1&3&0\\-1&1&2\end{pmatrix}$.}

\ans{ex:quad:A14}{方差 $\vv w^\trans\vv\Sigma\vv w=0.0215$，标准差约 $14.7\%$.}

\ans{ex:quad:A15}{$\vv w^*\propto\vv\Sigma^{-1}\vv 1\propto\big(\frac1{0.09},\frac1{0.04}\big)$，即 $\vv w^*=\big(\frac4{13},\frac9{13}\big)$，方差 $\frac9{325}$，标准差约 $16.6\%$.}

\ans{ex:quad:A16}{$(\vv C^\trans\vv{AC})^\trans=\vv C^\trans\vv A^\trans\vv C$；秩不变见推论~\ref{cor:rankinvariant}. $\det(\vv C^\trans\vv{AC})=(\det\vv C)^2\det\vv A$：保持符号，不保持数值.}

\ans{ex:quad:A17}{在 $\mathbb{R}$ 上为 $z_1^2+z_2^2-z_3^2$；在 $\mathbb{C}$ 上为 $z_1^2+z_2^2+z_3^2$.}

\ans{ex:quad:A18}{依次为单叶双曲面（$p=2,q=1$）、双叶双曲面（$p=1,q=2$）、圆锥面、圆柱面（秩为 $2$）.}

\ans{ex:quad:A19}{特征值 $3$ 与 $-1$. 取 $\vv x=(1,-1)^\trans$，二次型的值为 $-2$.}

\ans{ex:quad:A20}{对主子矩阵 $\vv A_I$（行、列指标集都为 $I$），把 $\vv x_I\ne\zero$ 在其余位置补零，$\vv x_I^\trans\vv A_I\vv x_I=\vv x^\trans\vv A\vv x>0$. 正定矩阵的行列式为正（特征值之积）.}

\subsection*{B 组提示}

\ans{ex:quad:B1}{可交换的对称矩阵可同时正交对角化（第十章章末习题~\ref{ex:dec:B1}），对角元对应相乘. 反例：$\vv A=\begin{pmatrix}2&1\\1&1\end{pmatrix}$，$\vv B=\diag(1,2)$.}
\ans{ex:quad:B2}{取 $\vv C$ 使 $\vv C^\trans\vv{AC}=\vv I$，$\vv C^\trans\vv{BC}=\diag(\mu_i)$，$\mu_i\ge0$，则 $\det(\vv A+\vv B)=\det\vv A\prod(1+\mu_i)$.}
\ans{ex:quad:B3}{$\vv x^\trans(\vv A+t\vv B)\vv x\ge(\lambda_{\min}(\vv A)-|t|\max_i|\lambda_i(\vv B)|)\|\vv x\|^2$，用 B5.}
\ans{ex:quad:B4}{第 $i,j$ 行列组成的二阶主子式为正.}
\ans{ex:quad:B5}{由谱定理写成 $\sum\lambda_iy_i^2$，$\|\vv y\|=\|\vv x\|$.}
\ans{ex:quad:B6}{特征值 $2,-1,-1$，$p=1$，$q=2$. 配方时先作 $x_1=y_1+y_2$，$x_2=y_1-y_2$.}
\ans{ex:quad:B7}{$\Delta_1=\lambda$，$\Delta_2=\lambda^2-1$，$\Delta_3=(\lambda-2)(\lambda+1)^2$，故 $\lambda>2$.}
\ans{ex:quad:B8}{驻点条件 $2\vv{Ax}=\mu\vv c$，代入约束定出 $\mu$. $\vv A$ 正定保证驻点是最小值点（第十一章最小二乘例题中的配方法同样适用）.}
\ans{ex:quad:B9}{驻点条件给出 $\vv w=a\vv 1+b\vv\mu$，由两个约束解出 $a,b$（关于 $m$ 线性），方差 $\sigma^2\|\vv w\|^2$ 是 $m$ 的二次函数.}
\ans{ex:quad:B10}{解 $\det(\vv B-\lambda\vv A)=3\lambda^2-1=0$，得 $\lambda=\pm\frac1{\sqrt3}$；对每个 $\lambda$ 求 $(\vv B-\lambda\vv A)\vv x=\zero$ 的解并使 $\vv x^\trans\vv A\vv x=1$，以它们为列作 $\vv C$，则 $\vv C^\trans\vv{BC}=\diag(\frac1{\sqrt3},-\frac1{\sqrt3})$.}
\ans{ex:quad:B11}{配方：$f=(\vv x+\frac12\vv A^{-1}\vv b)^\trans\vv A(\vv x+\frac12\vv A^{-1}\vv b)+c-\frac14\vv b^\trans\vv A^{-1}\vv b$. 半正定时不一定有最小值，例如 $\vv A=\vv O$，$\vv b\ne\zero$.}
\ans{ex:quad:B12}{写 $\vv A=\vv B^\trans\vv B$（$\vv B=\vv D^{1/2}\vv Q^\trans$），则 $\|\vv B\vv x\|^2=0$，$\vv{Ax}=\vv B^\trans\vv B\vv x=\zero$.}
\ans{ex:quad:B13}{用定理~\ref{thm:posdef} 证明中的分块合同变换（把 $\vv a$ 换成矩阵块）.}
\ans{ex:quad:B14}{主轴定理把椭圆化为 $\lambda_1u^2+\lambda_2v^2\le1$，面积 $\frac\pi{\sqrt{\lambda_1\lambda_2}}$，而正交变换保持面积.}
\ans{ex:quad:B15}{$\vv x^\trans(\vv A^\trans\vv A+\mu\vv I)\vv x=\|\vv{Ax}\|^2+\mu\|\vv x\|^2$. 目标函数对 $\vv x$ 的驻点条件给出正规方程的``岭''版本.}
\ans{ex:quad:B16}{用谱分解，$\vv B=\vv Q\diag(\sqrt{\lambda_i})\vv Q^\trans$；$\vv A^{-1}$ 的特征值为 $\frac1{\lambda_i}$.}

\section{第十二章　*对偶空间与双线性型}

\subsection*{A 组答案}

\ans{ex:bil:A1}{(1) 是. (2) 不是：$2\vv x\mapsto4x_1x_2$. (3) 是. (4) 不是：$f\mapsto2f$ 时值变为 $4$ 倍. (5) 是.}

\ans{ex:bil:A2}{以三个基向量为列作 $\vv P$，对偶基是 $\vv P^{-1}$ 的各行：$\vv\varepsilon_1^*(\vv x)=x_1-x_2$，$\vv\varepsilon_2^*(\vv x)=x_2-x_3$，$\vv\varepsilon_3^*(\vv x)=x_3$.}

\ans{ex:bil:A3}{$f=f(\vv\varepsilon_1)\vv\varepsilon_1^*+f(\vv\varepsilon_2)\vv\varepsilon_2^*+f(\vv\varepsilon_3)\vv\varepsilon_3^*=\vv\varepsilon_1^*+3\vv\varepsilon_2^*+6\vv\varepsilon_3^*$.}

\ans{ex:bil:A4}{$W^0$ 是一维的，基为 $\vv x\mapsto x_1-x_2+x_3$.}

\ans{ex:bil:A5}{$\mathscr T^*$ 的矩阵为 $\begin{pmatrix}1&0&3\\2&1&-1\end{pmatrix}$；$\mathscr T^*$ 把 $\vv y\mapsto y_1+y_2+y_3$ 变为 $\vv x\mapsto4x_1+2x_2$.}

\ans{ex:bil:A6}{标准基下 $\vv A=\begin{pmatrix}1&2\\-1&3\end{pmatrix}$；新基下 $\vv B=\begin{pmatrix}5&-5\\1&3\end{pmatrix}$，可直接验证 $\vv B=\vv P^\trans\vv{AP}$，$\vv P=\begin{pmatrix}1&1\\1&-1\end{pmatrix}$.}

\ans{ex:bil:A7}{对称部分的度量矩阵 $\begin{pmatrix}1&\frac12\\ \frac12&3\end{pmatrix}$，交错部分 $\begin{pmatrix}0&\frac32\\-\frac32&0\end{pmatrix}$.}

\ans{ex:bil:A8}{(1)(2) 非退化（度量矩阵行列式 $-1$）；(3) 退化（度量矩阵全为 $1$，不可逆）.}

\ans{ex:bil:A9}{对应二次型 $x_1^2+2x_1x_2+3x_2^2$. 由二次型得 $f(\vv x,\vv y)=x_1y_1-2x_1y_2-2x_2y_1+2x_2y_2$.}

\ans{ex:bil:A10}{$f(\vv x,\vv x)=0$，$f(\vv x,\vv y)=-f(\vv y,\vv x)$. 度量矩阵 $\begin{pmatrix}0&1\\-1&0\end{pmatrix}$，可逆，非退化.}

\ans{ex:bil:A11}{Lagrange 基：$\frac12(x-1)(x-2)$，$-x(x-2)$，$\frac12x(x-1)$.}

\ans{ex:bil:A12}{$\dim W=2$，$\dim W^0=2$，基为 $\vv x\mapsto x_1+x_2$ 与 $\vv x\mapsto x_3-x_4$.}

\ans{ex:bil:A13}{双线性显然，对称性由 $\tr(\vv{AB})=\tr(\vv{BA})$；非退化：$\tr(\vv A\vv E_{ji})=a_{ij}$. 不是内积：$\vv A=\begin{pmatrix}0&1\\-1&0\end{pmatrix}$ 时 $\tr(\vv A^2)=-2<0$.}

\ans{ex:bil:A14}{坐标为 $\big(\varphi(1),\varphi(x),\varphi(x^2)\big)=\big(1,\frac12,\frac13\big)$.}

\ans{ex:bil:A15}{对任意 $\vv\alpha=\sum a_i\vv e_i$，$\langle\vv\alpha,\vv\beta\rangle=\sum a_if(\vv e_i)=f(\vv\alpha)$. 唯一性：若 $\langle\vv\alpha,\vv\beta-\vv\beta'\rangle=0$ 对一切 $\vv\alpha$ 成立，取 $\vv\alpha=\vv\beta-\vv\beta'$.}

\subsection*{B 组提示}

\ans{ex:bil:B1}{取 $V$ 的一组基，映射的矩阵的第 $i$ 行恰是 $f_i$ 在对偶基下的坐标. 满射当且仅当这个 $k\times n$ 矩阵行满秩，即各行线性无关.}
\ans{ex:bil:B2}{取 $\vv\alpha_0$ 使 $f(\vv\alpha_0)=1$，则 $V=\ker f\oplus\spn(\vv\alpha_0)$，$c=g(\vv\alpha_0)$. 推广：对映射 $\vv\alpha\mapsto(f_1(\vv\alpha),\ldots,f_k(\vv\alpha))$ 用同样的思路，或用零化子的维数.}
\ans{ex:bil:B3}{记 $U=\spn(f_1,\ldots,f_k)\subset V^*$. 把 $f_i$ 写成坐标后，$\ker f_1\cap\cdots\cap\ker f_k$ 就是以这些坐标为行的齐次方程组的解空间，维数为 $\dim V-\dim U$. 方程组语言：齐次方程组只有零解，当且仅当各方程的系数行向量张成整个行向量空间.}
\ans{ex:bil:B4}{由定理~\ref{thm:dualrank}：$\rk\mathscr T^*=\rk\mathscr T$.}
\ans{ex:bil:B5}{$a=c=\frac13$，$b=\frac43$. 积分是 $\mathbb{R}[x]_3$ 上的线性函数，三个求值函数构成 $\mathbb{R}[x]_3^*$ 的基（A11）. 三次时验证 $\int_0^2(x-1)^3\,\mathrm dx=0$ 与右边相等.}
\ans{ex:bil:B6}{双线性型由度量矩阵唯一确定，反之任一矩阵给出双线性型. 后半句即第三章章末习题~\ref{ex:mat:A22}.}
\ans{ex:bil:B7}{由定理~\ref{thm:alternating}，非退化时没有零块，$\dim V=2s$.}
\ans{ex:bil:B8}{度量矩阵 $\begin{pmatrix}1&1&0\\1&1&0\\0&0&1\end{pmatrix}$ 的核为 $\spn\big((1,-1,0)^\trans\big)$.}
\ans{ex:bil:B9}{迷向向量满足 $x_1^2=x_2^2$，即两条直线 $x_2=\pm x_1$. 后半句：设 $\vv P=\begin{pmatrix}a&b\\c&d\end{pmatrix}$，由条件得 $a^2-c^2=1$，$d^2-b^2=1$，$ab=cd$，再令 $a=\cosh t$.}
\ans{ex:bil:B10}{对 $\vv\alpha\in W$，$f\in W^0$，有 $(\mathscr T^*f)(\vv\alpha)=f(\mathscr T\vv\alpha)=0$.}
（须在 %% =====================================================================
%% exercises.tex　《高等代数讲义》章末习题（A/B/C 三组）
%%
%% 用法
%%   1. 在导言区（\begin{document} 之前、全书统一记号之后）写 \input{exercises}.
	%%   2. 在每章末尾（原 \begin{problemset}...\end{problemset} 之后，或替换它）
	%%      写对应的命令：
	%%        第一章 \ExElimination      第二章 \ExVectorSpace   第三章 \ExMatrix
	%%        第四章 \ExLinearMap        第五章 \ExDeterminant   第六章 \ExSolutions
	%%        第七章 \ExPolynomial       第八章 \ExEigen         第九章 \ExInnerProduct
	%%        第十章 \ExDecomposition    第十一章 \ExQuadratic   第十二章 \ExBilinear
	%%   3. 分组标准
	%%        A 组：直接运用本章定义与定理的计算和证明；
	%%        B 组：需要综合多个结论、或有一定技巧的计算和证明；
	%%        C 组：有挑战性的题目，或超出本课程大纲、为继续深造准备的内容.
	%%   4. 题目与正文中已有的节末练习、原 problemset 不重复；个别题目注明``接正文某题''.
	%%      第六章按扩充版（ch06-solutions.tex）的内容编写.
	%%   5. 每道题都带标签 ex:<章键>:<组号>，正文、习题与附录答案一律用 \ref 引用，
	%%      渲染为``A4''``B12''这样的题号. 章键依次为
	%%        elim vs mat lmap det sol poly eig ip dec quad bil（第一至十二章）.
	%%      例：第一章章末习题~\ref{ex:elim:B12}. 新增题目请加在组末，并补上标签.
	%%      A 组答案、B 组提示在 answers.tex 中，按同一标签对应.
	%% 依赖：导言区的 \vv \zero \trans \spn \rk \imm \diag，以及 enumitem（elegantbook 已加载）.
	%% 引用的前文标签：thm:sameSolution thm:solvability prop:linspc eg:pricedecomp cor:homogeneous thm:rouche thm:sol-determinable
	%%   thm:spectral thm:posdef thm:svd thm:orthcanon
	%% =====================================================================
	
	\providecommand{\tr}{\operatorname{tr}}
	\newcommand{\exsetheading}{%
		\section*{章末习题}%
		\markright{章末习题}%
	}
\newcommand{\exgroup}[2]{%
	\par\medskip{\centering\sffamily\bfseries #1\quad #2\par}\smallskip}
	\newenvironment{exA}{\begin{enumerate}[label=\textbf{A\arabic*.},ref=A\arabic*,leftmargin=2.8em,itemsep=4pt]}{\end{enumerate}}
	\newenvironment{exB}{\begin{enumerate}[label=\textbf{B\arabic*.},ref=B\arabic*,leftmargin=2.8em,itemsep=4pt]}{\end{enumerate}}
	\newenvironment{exC}{\begin{enumerate}[label=\textbf{C\arabic*.},ref=C\arabic*,leftmargin=2.8em,itemsep=4pt]}{\end{enumerate}}
	\newcommand{\Agroup}{\exgroup{A 组}{基础巩固}}
	\newcommand{\Bgroup}{\exgroup{B 组}{综合提升}}
	\newcommand{\Cgroup}{\exgroup{C 组}{拓展探索}}
	
	% =====================================================================
	% 第一章　线性方程组与消元法
	% =====================================================================
	\newcommand{\ExElimination}{%
		\exsetheading
		\Agroup
		\begin{exA}
			\item\label{ex:elim:A1} 判断下列矩阵哪些是阶梯形，哪些是简化阶梯形：
			\[
			\begin{pmatrix}1&2&0\\0&0&1\end{pmatrix},\quad
			\begin{pmatrix}1&0&3\\0&2&1\\0&0&0\end{pmatrix},\quad
			\begin{pmatrix}0&1&0\\1&0&0\end{pmatrix},\quad
			\begin{pmatrix}1&0&2&0\\0&1&-1&0\\0&0&0&1\end{pmatrix},\quad
			\begin{pmatrix}0&0\\0&1\end{pmatrix}.
			\]
			\item\label{ex:elim:A2} 写出方程组 $3x_1-x_3=2$，$x_2+4x_3=0$，$x_1-x_2=5$ 的增广矩阵；反过来，写出以 $\begin{pmatrix}1&-2&0&3\\0&1&5&-1\end{pmatrix}$ 为增广矩阵的方程组.
			\item\label{ex:elim:A3} 用消元法解方程组 $\left\{\begin{array}{l}x_1+2x_2+3x_3=14,\\2x_1-x_2+x_3=3,\\3x_1+x_2-x_3=2.\end{array}\right.$
			\item\label{ex:elim:A4} 用消元法解方程组 $\left\{\begin{array}{l}x_1+x_2+x_3+x_4=2,\\x_1-x_2+x_3-x_4=4,\\2x_1+x_2-x_3+3x_4=-1,\\x_1+2x_2+3x_3-2x_4=5.\end{array}\right.$
			\item\label{ex:elim:A5} 判断方程组 $\left\{\begin{array}{l}x_1-x_2+x_3=1,\\2x_1+x_2-x_3=2,\\x_1+2x_2-2x_3=3\end{array}\right.$ 是否有解，并说明消元过程中哪一行``暴露''了结论.
			\item\label{ex:elim:A6} 求方程组 $\left\{\begin{array}{l}x_1+2x_2-x_3+x_4=3,\\2x_1+4x_2-x_3+3x_4=7,\\-x_1-2x_2+2x_3=-2\end{array}\right.$ 的全部解，并指出主变量与自由变量.
			\item\label{ex:elim:A7} 求齐次方程组 $\left\{\begin{array}{l}x_1+x_2-x_3-x_4=0,\\2x_1-5x_2+3x_3+2x_4=0,\\7x_1-7x_2+3x_3+x_4=0\end{array}\right.$ 的全部解.
			\item\label{ex:elim:A8} 把矩阵 $\begin{pmatrix}2&4&-2&6\\1&2&1&1\\3&6&1&5\end{pmatrix}$ 化为简化阶梯形，并写出每一步所用的初等行变换.
			\item\label{ex:elim:A9} 写出全部 $2\times2$ 简化阶梯形矩阵的``形状''（可以含参数），并说明共有几类.
			\item\label{ex:elim:A10} $k$ 取何值时，方程组 $x+y=1$，$x-y=3$，$2x+ky=5$ 有解？
			\item\label{ex:elim:A11} 设 $a,b$ 为常数. 方程组 $x+2y=a$，$3x+6y=b$ 何时有解？有解时写出全部解.
			\item\label{ex:elim:A12} 求过三点 $(1,2)$，$(2,3)$，$(3,6)$ 的抛物线 $y=ax^2+bx+c$.
			\item\label{ex:elim:A13} 求过四点 $(-1,2)$，$(0,1)$，$(1,0)$，$(2,5)$ 的三次多项式函数 $y=a+bx+cx^2+dx^3$.
			\item\label{ex:elim:A14} 某剧场三种票价分别为 $30$，$50$，$80$ 元. 某场演出共售出 $100$ 张票，票房 $4800$ 元，且 $80$ 元票的张数恰为 $30$ 元票的一半. 三种票各售出多少张？
			\item\label{ex:elim:A15} 配平化学方程式 $x_1\,\mathrm{C_3H_8}+x_2\,\mathrm{O_2}\to x_3\,\mathrm{CO_2}+x_4\,\mathrm{H_2O}$：对 C，H，O 三种原子分别写出守恒方程，解这个齐次方程组，并取最小的正整数解.
			\item\label{ex:elim:A16} 某城区四个路口 $A,B,C,D$ 由单行道依次连成环路 $A\to B\to C\to D\to A$，环路上的车流量依次记为 $x_1$（$A\to B$），$x_2$（$B\to C$），$x_3$（$C\to D$），$x_4$（$D\to A$）（单位：辆/小时）. 各路口与外部道路的进出量为：$A$ 驶入 $400$、驶出 $300$；$B$ 驶入 $200$、驶出 $350$；$C$ 驶入 $250$、驶出 $100$；$D$ 驶入 $150$、驶出 $250$. 设每个路口驶入总量等于驶出总量.
			\begin{enumerate}
				\item 列出方程组并求全部解；
				\item 若各路段车流量都不能为负，$x_1$ 最少是多少？
				\item 若把 $D$ 的驶出量改为 $350$，方程组会怎样？这反映了什么守恒律？
			\end{enumerate}
			\item\label{ex:elim:A17} 证明：二元一次方程组 $ax+by=e$，$cx+dy=f$ 对任意的 $e,f$ 都有唯一解，当且仅当 $ad-bc\ne0$.（只用消元法，不用行列式.）
			\item\label{ex:elim:A18} 证明：若线性方程组有两个不同的解 $\vv x_1,\vv x_2$，则对任意数 $t$，$\vv x_1+t(\vv x_2-\vv x_1)$ 都是解. 由此说明：线性方程组的解要么不存在，要么唯一，要么有无穷多个.
			\item\label{ex:elim:A19} 证明方程组 $\{x+y+z=3,\ x-y=0\}$ 与 $\{2x+z=3,\ x+y+z=3\}$ 同解.
			\item\label{ex:elim:A20} 举例说明：若\textbf{同时}做``第一个方程加到第二个方程上''与``第二个方程加到第一个方程上''（两步都用原来的方程），所得方程组可能与原方程组不同解. 这与定理~\ref{thm:sameSolution}（初等行变换不改变解集）矛盾吗？
			\item\label{ex:elim:A21} 已知一个方程组的增广矩阵化为阶梯形后为 $\begin{pmatrix}1&3&-1&2&0\\0&0&2&1&4\\0&0&0&0&a\end{pmatrix}$. 讨论 $a$ 取何值时方程组有解；有解时它有几个自由变量？
		\end{exA}
		
		\Bgroup
		\begin{exB}
			\item\label{ex:elim:B1} 讨论 $a,b$ 取何值时，方程组
			\[\left\{\begin{array}{l}x_1+x_2+x_3+x_4=0,\\x_2+2x_3+2x_4=1,\\-x_2+(a-3)x_3-2x_4=b,\\3x_1+2x_2+x_3+ax_4=-1\end{array}\right.\]
			有唯一解、无解、有无穷多解；有无穷多解时求出全部解.
			\item\label{ex:elim:B2} 讨论 $\lambda$ 取何值时，方程组
			$(1+\lambda)x_1+x_2+x_3=0$，$x_1+(1+\lambda)x_2+x_3=3$，$x_1+x_2+(1+\lambda)x_3=\lambda$
			有唯一解、无解、有无穷多解.
			\item\label{ex:elim:B3} 设 $n$ 个方程、$n$ 个未知量的方程组对应的齐次方程组只有零解. 证明：无论常数项取什么值，原方程组都有唯一解.（提示：考察简化阶梯形中主元的个数.）
			\item\label{ex:elim:B4} 证明：线性方程组无解，当且仅当可以把若干个方程分别乘以适当的数再相加，得到矛盾方程 $0=1$.
			\item\label{ex:elim:B5} 证明：交换两行这一初等变换，可以由``倍加''与``某行乘以 $-1$''这两类变换的组合实现. 由此说明三类初等行变换并不是``独立''的.
			\item\label{ex:elim:B6} 证明：若线性方程组的系数和常数项都是有理数，且它有实数解，则它一定有有理数解；若它的实数解唯一，则这个解就是有理数解.
			\item\label{ex:elim:B7} 讨论方程组
			\[x_1-x_2=a_1,\quad x_2-x_3=a_2,\quad\ldots,\quad x_{n-1}-x_n=a_{n-1},\quad x_n-x_1=a_n\]
			有解的充要条件，并在有解时求全部解.
			\item\label{ex:elim:B8} 讨论方程组
			\[x_1+x_2=a_1,\quad x_2+x_3=a_2,\quad\ldots,\quad x_{n-1}+x_n=a_{n-1},\quad x_n+x_1=a_n\]
			的解. 证明：$n$ 为奇数时解唯一；$n$ 为偶数时有解当且仅当 $a_1-a_2+a_3-\cdots-a_n=0$.
			\item\label{ex:elim:B9} 某经济体有三个部门，直接消耗系数矩阵与最终需求分别为
			\[\vv A=\begin{pmatrix}0.1&0.2&0.2\\0.2&0.1&0.2\\0.1&0.3&0.1\end{pmatrix},\qquad \vv d=\begin{pmatrix}20\\130\\65\end{pmatrix}.\]
			按第一章``投入产出模型''的方式列出方程组，求各部门的总产出.
			\item\label{ex:elim:B10} 估计用消元法（化为阶梯形后回代）解 $n$ 个方程、$n$ 个未知量的方程组所需的乘除法次数，证明它约为 $\frac13n^3$（即与 $\frac13n^3$ 之比在 $n\to\infty$ 时趋于 $1$）.
			\item\label{ex:elim:B11} 设三个平面 $a_ix+b_iy+c_iz=d_i$（$i=1,2,3$）两两相交. 用消元法的语言说明：它们恰好交于一点、交于一条直线、或没有公共点（三条交线互相平行），并各举一例.
			\item\label{ex:elim:B12} 证明：$n$ 元齐次方程组的解集对加法和数乘封闭；而非齐次方程组的两个解之差是对应齐次方程组的解. 由此说明：非齐次方程组的全部解，等于它的任一个解加上齐次方程组的全部解.
			\item\label{ex:elim:B13} 一家律所有甲、乙、丙三位律师，本月分别承办了若干件合同、劳动、知识产权纠纷. 已知：三类案件合计 $30$ 件；甲承办的数量 $=$ 乙 $+$ 丙；每件合同、劳动、知识产权案件分别需要 $10$，$6$，$15$ 个工时，三人本月共用 $300$ 个工时……
			这些条件够不够确定每位律师承办了多少案件？请自己补全一个合理的设定（明确未知量是什么），列出方程组，用自由变量的个数说明还缺多少条独立的信息.
			\item\label{ex:elim:B14} 设方程组的系数矩阵每一行的元素之和都等于 $0$. 证明 $x_1=x_2=\cdots=x_n=1$ 是对应齐次方程组的解；若非齐次方程组有解 $\vv x_0$，则对任意数 $t$，$\vv x_0+t(1,\ldots,1)^\trans$ 也是解.
			\item\label{ex:elim:B15} 证明：如果方程组 (I) 的每个方程都是方程组 (II) 各方程的``线性组合''（各乘一个数再相加），反之亦然，那么 (I) 与 (II) 同解. 逆命题对于有解的方程组成立吗？先在两个未知量的情形下思考.
			\item\label{ex:elim:B16} 讨论 $\lambda$ 取何值时，方程组 $\lambda x+y+z=1$，$x+\lambda y+z=\lambda$，$x+y+\lambda z=\lambda^2$ 有唯一解、无解、有无穷多解；有唯一解时求出解.（第六章例~\ref{eg:lambda} 将用秩与行列式重新组织这个讨论.）
		\end{exB}
		
		\Cgroup
		\begin{exC}
			\item\label{ex:elim:C1} （百鸡问题）《张丘建算经》（约公元 5 世纪）：``今有鸡翁一，直钱五；鸡母一，直钱三；鸡雏三，直钱一. 凡百钱买鸡百只，问鸡翁、母、雏各几何？''列出线性方程组，求它的全部解，再从中找出所有\textbf{非负整数}解. 线性代数在哪一步之后就无能为力了？
			\item\label{ex:elim:C2} （幻方）若一个 $3\times3$ 数表的每行、每列和两条对角线上的三个数之和都等于同一个数 $s$，称它为幻方（不要求数互不相同）. 把 $9$ 个元素作为未知量、$s$ 也作为未知量，写出方程组，证明：
			\begin{enumerate}
				\item 中心元素一定等于 $s/3$；
				\item 全部幻方可以用 $3$ 个自由参数表示，写出通式；
				\item 用 $1,2,\ldots,9$ 各一次填成的幻方本质上只有一个（其余由旋转、翻转得到）.
			\end{enumerate}
			\item\label{ex:elim:C3} （点灯游戏）$3\times3$ 方格中每格有一盏灯，按下某格会使该格及其上下左右相邻的格子开关状态翻转. 在``$1+1=0$''的算术（即模 $2$ 运算）中，把``按哪些格''作为未知量，把``从某一状态变为全灭''写成线性方程组. 证明：从任意初始状态出发，总能把灯全部熄灭.（模 $2$ 运算的集合 $\{0,1\}$ 不是本书意义下的数域，但消元法照样适用——想一想为什么.）
			\item\label{ex:elim:C4} （舍入误差与主元）在``只保留三位有效数字''的计算机上解方程组 $0.0001x+y=1$，$x+y=2$：
			\begin{enumerate}
				\item 直接用第一个方程消去第二个方程中的 $x$；
				\item 先交换两个方程再消元.
			\end{enumerate}
			把两种结果与精确解比较，说明数值软件为什么要``选绝对值最大的元素作主元''.
			\item\label{ex:elim:C5} 设两个都有解的方程组 (I)(II) 都有 $m$ 个方程、$n$ 个未知量，且同解. 证明：它们的增广矩阵行等价.（提示：先证明二者的简化阶梯形的非零行完全相同.）
			\item\label{ex:elim:C6} （Jacobi 迭代）对方程组 $4x+y=9$，$x+3y=5$，把它改写为 $x=(9-y)/4$，$y=(5-x)/3$，从 $(x_0,y_0)=(0,0)$ 出发反复代入. 计算前五步，观察是否收敛到真解. 再对 $x+4y=9$，$3x+y=5$ 做同样的事，会发生什么？试找出一个保证收敛的充分条件.
			\item\label{ex:elim:C7} （电路）一个电路中有三条支路电流 $I_1,I_2,I_3$. 节点 $A$ 处 $I_1=I_2+I_3$；两个回路的电压方程为 $2I_1+I_2=8$，$I_2-2I_3=0$（单位从略）. 求各电流. 再说明：若把第三个节点 $B$ 的电流方程也写进来，方程组为什么不会因此``多一个条件''？
			\item\label{ex:elim:C8} 证明：一个 $m\times n$ 矩阵（$m,n\ge1$）通过初等行变换所能得到的不同的简化阶梯形``形状''（即主元所在列的位置组合）共有 $\sum_{r=0}^{\min(m,n)}\binom nr$ 种.
			\item\label{ex:elim:C9} 证明：一个矩阵的简化阶梯形是唯一的.（提示：设 $\vv R_1,\vv R_2$ 都是 $\vv M$ 的简化阶梯形，则以它们为系数矩阵的齐次方程组同解. 对列数归纳：删去最后一列后由归纳假设前面各列相同，再比较最后一列，考察它在两者中是否为主元列.）
		\end{exC}
	}
	
	% =====================================================================
	% 第二章　向量空间
	% =====================================================================
	\newcommand{\ExVectorSpace}{%
		\exsetheading
		\Agroup
		\begin{exA}
			\item\label{ex:vs:A1} 下列集合关于通常的四则运算是否构成数域？
			(1) $\{a+b\sqrt3\mid a,b\in\mathbb{Q}\}$；(2) $\{a+b\sqrt2\mid a\in\mathbb{Q},\ b\in\mathbb{Z}\}$；(3) 全体偶数；(4) $\{a+b\sqrt[3]{2}\mid a,b\in\mathbb{Q}\}$；(5) 全体纯虚数与 $0$.
			\item\label{ex:vs:A2} 下列集合关于通常的加法与数乘是否构成实向量空间？
			(1) 全体 $n$ 阶实对称矩阵；(2) $\{(x,y)^\trans\in\mathbb{R}^2\mid x\ge0\}$；(3) 次数恰为 $n$ 的实系数多项式全体；(4) $\{f\in\mathbb{R}[x]\mid f(1)=0\}$；(5) 满足 $a_{n+2}=a_{n+1}+a_n$ 的全体实数列；(6) $\{(x,y,z)^\trans\mid x+y+z=1\}$.
			\item\label{ex:vs:A3} 判断下列向量组的线性相关性：
			\begin{enumerate}
				\item $(1,2,-1)^\trans$，$(2,-3,1)^\trans$，$(4,1,-1)^\trans$；
				\item 在 $\mathbb{R}[x]_3$ 中：$1+x$，$1-x$，$x^2$；
				\item 在 $\mathbb{R}[x]_3$ 中：$1+x+x^2$，$1-x$，$2+x^2$，$x+x^2$；
				\item 在 $M_2(\mathbb{R})$ 中：$\begin{pmatrix}1&0\\0&1\end{pmatrix}$，$\begin{pmatrix}0&1\\1&0\end{pmatrix}$，$\begin{pmatrix}1&1\\1&1\end{pmatrix}$，$\begin{pmatrix}1&0\\0&0\end{pmatrix}$.
			\end{enumerate}
			\item\label{ex:vs:A4} 把 $\vv\beta=(1,2,1,1)^\trans$ 表示为 $\vv\alpha_1=(1,1,1,1)^\trans$，$\vv\alpha_2=(1,1,-1,-1)^\trans$，$\vv\alpha_3=(1,-1,1,-1)^\trans$，$\vv\alpha_4=(1,-1,-1,1)^\trans$ 的线性组合.
			\item\label{ex:vs:A5} $t$ 取何值时，$(1,1,0)^\trans$，$(1,3,-1)^\trans$，$(5,3,t)^\trans$ 线性相关？
			\item\label{ex:vs:A6} 求向量组 $\vv\alpha_1=(1,-1,2,4)^\trans$，$\vv\alpha_2=(0,3,1,2)^\trans$，$\vv\alpha_3=(3,0,7,14)^\trans$，$\vv\alpha_4=(1,-1,2,0)^\trans$，$\vv\alpha_5=(2,1,5,6)^\trans$ 的秩和一个极大无关组，并把其余向量用这个极大无关组线性表出.
			\item\label{ex:vs:A7} 设 $\vv\alpha_1,\vv\alpha_2,\vv\alpha_3$ 线性无关. 判断下列向量组的线性相关性：
			(1) $\vv\alpha_1-\vv\alpha_2$，$\vv\alpha_2-\vv\alpha_3$，$\vv\alpha_3-\vv\alpha_1$；
			(2) $\vv\alpha_1$，$\vv\alpha_1+\vv\alpha_2$，$\vv\alpha_1+\vv\alpha_2+\vv\alpha_3$；
			(3) $\vv\alpha_1+2\vv\alpha_2$，$\vv\alpha_2+2\vv\alpha_3$，$\vv\alpha_3+2\vv\alpha_1$.
			\item\label{ex:vs:A8} 证明：含有零向量的向量组线性相关；含有两个相同向量的向量组线性相关；单个向量 $\vv\alpha$ 线性无关当且仅当 $\vv\alpha\ne\zero$.
			\item\label{ex:vs:A9} 证明：两个向量线性相关，当且仅当其中一个是另一个的倍数. 三个向量的情形能这样说吗？
			\item\label{ex:vs:A10} 只用向量空间的八条公理与命题~\ref{prop:linspc}，证明：(1) $\vv\alpha+\vv\alpha=2\vv\alpha$；(2) 若 $k\ne0$ 且 $k\vv\alpha=k\vv\beta$，则 $\vv\alpha=\vv\beta$；(3) $-(-\vv\alpha)=\vv\alpha$.
			\item\label{ex:vs:A11} 证明下列集合是相应空间的子空间，并求一组基和维数：
			\begin{enumerate}
				\item $\{\vv x\in\mathbb{R}^4\mid x_1+x_2=x_3+x_4\}$；
				\item $n$ 阶上三角矩阵全体；
				\item 迹为零的 $n$ 阶方阵全体 $\{\vv X\mid x_{11}+\cdots+x_{nn}=0\}$；
				\item $\{f\in\mathbb{R}[x]_4\mid f(1)=0\}$.
			\end{enumerate}
			\item\label{ex:vs:A12} 判断下列子集是否为子空间：(1) $\{(a,b,c)^\trans\mid a=2b\}$；(2) $\{(a,b,c)^\trans\mid ab=0\}$；(3) $\{(a,b,c)^\trans\mid a,b,c\in\mathbb{Q}\}\subset\mathbb{R}^3$；(4) $\{f\in C[0,1]\mid f(0)=2f(1)\}$.
			\item\label{ex:vs:A13} 求 $\vv\alpha_1=(1,2,1,0)^\trans$，$\vv\alpha_2=(-1,1,1,1)^\trans$，$\vv\alpha_3=(0,3,2,1)^\trans$，$\vv\alpha_4=(1,0,1,2)^\trans$ 张成的子空间的维数和一组基.
			\item\label{ex:vs:A14} 求多项式 $2x^2-x+3$ 在基 $1,\ x-1,\ (x-1)^2$ 下的坐标.
			\item\label{ex:vs:A15} 求矩阵 $\begin{pmatrix}2&3\\-1&4\end{pmatrix}$ 在 $M_2(\mathbb{R})$ 的基 $\begin{pmatrix}1&0\\0&0\end{pmatrix},\begin{pmatrix}1&1\\0&0\end{pmatrix},\begin{pmatrix}1&1\\1&0\end{pmatrix},\begin{pmatrix}1&1\\1&1\end{pmatrix}$ 下的坐标.
			\item\label{ex:vs:A16} (1) $\mathbb{R}$ 是否是 $\mathbb{Q}$ 上的向量空间？(2) $\mathbb{Q}$ 是否是 $\mathbb{R}$ 上的向量空间？(3) $\mathbb{C}$ 作为 $\mathbb{R}$ 上的向量空间和作为 $\mathbb{C}$ 上的向量空间，维数各是多少？
			\item\label{ex:vs:A17} 写出 $\mathbb{F}[x]_n$ 与 $\mathbb{F}^n$ 之间、$M_2(\mathbb{F})$ 与 $\mathbb{F}^4$ 之间的一个同构映射，并验证它保持线性运算.
			\item\label{ex:vs:A18} 设 $W_1=\spn\big((1,0,0)^\trans,(0,1,0)^\trans\big)$，$W_2=\spn\big((1,1,1)^\trans\big)$，$W_3=\spn\big((1,1,0)^\trans\big)$. 分别判断 $W_1+W_2$，$W_1+W_3$ 是否为直和，并求和空间的维数.
			\item\label{ex:vs:A19} 设 $W_1=\spn\big((1,1,0,0)^\trans,(1,0,1,1)^\trans\big)$，$W_2=\spn\big((0,-1,1,1)^\trans,(1,2,0,1)^\trans\big)$. 求 $W_1+W_2$ 与 $W_1\cap W_2$ 的维数，并验证维数公式.
			\item\label{ex:vs:A20} 设 $\vv\alpha_1,\ldots,\vv\alpha_s$ 的秩为 $r$. 证明：其中任意 $r+1$ 个向量线性相关；任意 $r$ 个线性无关的向量都构成极大无关组.
		\end{exA}
		
		\Bgroup
		\begin{exB}
			\item\label{ex:vs:B1} 证明：两个子空间的并 $W_1\cup W_2$ 是子空间，当且仅当 $W_1\subset W_2$ 或 $W_2\subset W_1$.
			\item\label{ex:vs:B2} 证明：若向量组 $\vv\beta_1,\ldots,\vv\beta_t$ 可由 $\vv\alpha_1,\ldots,\vv\alpha_s$ 线性表出，且两组的秩相等，则两组等价.
			\item\label{ex:vs:B3} 设向量组 (I) 与 (II) 的秩分别为 $r_1,r_2$，把它们合并为 (III). 证明 $\max\{r_1,r_2\}\le\rk(\mathrm{III})\le r_1+r_2$. 再证明：$\rk(\vv\alpha_1+\vv\beta_1,\ldots,\vv\alpha_s+\vv\beta_s)\le\rk(\vv\alpha_1,\ldots,\vv\alpha_s)+\rk(\vv\beta_1,\ldots,\vv\beta_s)$.
			\item\label{ex:vs:B4} 设 $\vv\alpha_1,\ldots,\vv\alpha_n$ 是 $V$ 的基，$\vv\beta=a_1\vv\alpha_1+\cdots+a_n\vv\alpha_n$. 证明：把 $\vv\alpha_j$ 换成 $\vv\beta$ 后所得的向量组仍是基，当且仅当 $a_j\ne0$.
			\item\label{ex:vs:B5} 设 $p_0,p_1,\ldots,p_{n-1}\in\mathbb{F}[x]$ 满足 $\deg p_k=k$. 证明它们构成 $\mathbb{F}[x]_n$ 的一组基.
			\item\label{ex:vs:B6} 记 $V$ 为满足 $a_{n+2}=a_{n+1}+2a_n$（$n\ge1$）的全体实数列. 证明 $V$ 是向量空间，$\dim V=2$，并说明 $\{2^n\}$ 与 $\{(-1)^n\}$ 构成 $V$ 的一组基. 由此写出 $a_1=1$，$a_2=1$ 时的通项.
			\item\label{ex:vs:B7} 在 $C(\mathbb{R})$ 中：(1) 证明 $\mathrm e^x,\mathrm e^{2x},\mathrm e^{3x}$ 线性无关；(2) 证明 $\sin x,\cos x,\sin(x+1)$ 线性相关；(3) $1,\cos x,\cos2x,\cos^2x$ 的秩是多少？
			\item\label{ex:vs:B8} 证明 $\{a+b\sqrt2+c\sqrt3+d\sqrt6\mid a,b,c,d\in\mathbb{Q}\}$ 是数域，并证明 $1,\sqrt2,\sqrt3,\sqrt6$ 在 $\mathbb{Q}$ 上线性无关. 这个数域作为 $\mathbb{Q}$ 上的向量空间是几维的？
			\item\label{ex:vs:B9} 写出 $\mathbb{R}^2$ 与 $\mathbb{R}^3$ 的全部子空间（按维数分类，并说明几何意义）.
			\item\label{ex:vs:B10} 证明：$\mathbb{R}^5$ 中任意两个三维子空间的交至少是一维的. 一般地，$n$ 维空间中维数分别为 $k_1,k_2$ 的两个子空间，若 $k_1+k_2>n$，则交不为 $\{\zero\}$.
			\item\label{ex:vs:B11} 设 $a\in\mathbb{F}$. 证明 $\mathbb{F}[x]_n=\mathbb{F}\cdot1\oplus\{f\in\mathbb{F}[x]_n\mid f(a)=0\}$，并写出 $f$ 按这个直和的分解.
			\item\label{ex:vs:B12} 设 $V=C[-1,1]$，$E$，$O$ 分别为其中的偶函数、奇函数全体. 证明 $V=E\oplus O$，并把 $\mathrm e^x$ 按此分解.
			\item\label{ex:vs:B13} 设 $\dim V=n$，$W_1,\ldots,W_k$ 是 $V$ 的子空间，且 $\dim W_1+\cdots+\dim W_k>(k-1)n$. 证明 $W_1\cap\cdots\cap W_k\ne\{\zero\}$.
			\item\label{ex:vs:B14} 设 $\vv\alpha_1,\ldots,\vv\alpha_n$ 是 $V$ 的基，$W=\spn(\vv\alpha_1+\vv\alpha_2,\ \vv\alpha_2+\vv\alpha_3,\ \ldots,\ \vv\alpha_{n-1}+\vv\alpha_n)$. 求 $\dim W$，并求 $W$ 的一个补空间.
			\item\label{ex:vs:B15} 设 $V_1=\{\vv x\in\mathbb{F}^n\mid x_1=x_2=\cdots=x_n\}$，$V_2=\{\vv x\in\mathbb{F}^n\mid x_1=0\}$. 证明 $\mathbb{F}^n=V_1\oplus V_2$. 与第二章例~\ref{eg:pricedecomp} 的分解相比，哪一种分解更``自然''？（提示：思考补空间的不唯一性.）
			\item\label{ex:vs:B16} 在 $M_n(\mathbb{F})$ 中，设 $W_1$ 为对角线以下元素全为零的矩阵（上三角矩阵）全体，$W_2$ 为对角线及以上元素全为零的矩阵全体. 证明 $M_n(\mathbb{F})=W_1\oplus W_2$，并求 $\dim W_1$.
			\item\label{ex:vs:B17} （覆盖定理的另一个证明）设 $W_1,\ldots,W_s$ 是 $\mathbb{R}^n$ 的真子空间，令 $\vv v(t)=(1,t,t^2,\ldots,t^{n-1})^\trans$.
			\begin{enumerate}
				\item 证明：每个真子空间都包含于某个超平面 $\{\vv x\mid a_1x_1+\cdots+a_nx_n=0\}$（$a_i$ 不全为零）中.（提示：取 $W$ 的基并扩充为 $\mathbb{R}^n$ 的基，看``最后一个坐标''.）
				\item 证明：对上述超平面，满足 $\vv v(t)$ 落在其中的实数 $t$ 至多有 $n-1$ 个.
				\item 由此推出：$\vv v(t)$ 落在 $W_1\cup\cdots\cup W_s$ 中的 $t$ 至多有 $s(n-1)$ 个. 这给出了覆盖定理在 $\mathbb{R}^n$ 中的一个``构造性''证明. 对 $n=3$ 以及三个子空间 $x_1=x_2$，$x_2=x_3$，$x_1+x_2+x_3=0$，写出一个不在它们之中的 $\vv v(t)$.
			\end{enumerate}
			\item\label{ex:vs:B18} 设 $\dim V=n$，$W_1,\ldots,W_s$ 是 $V$ 的真子空间. 证明：存在 $V$ 的一组基 $\vv\varepsilon_1,\ldots,\vv\varepsilon_n$，其中\textbf{每个}向量都不属于 $W_1\cup\cdots\cup W_s$.（提示：逐个选取. 选 $\vv\varepsilon_k$ 时，要躲开的除了 $W_1,\ldots,W_s$，还有哪个真子空间？）
		\end{exB}
		
		\Cgroup
		\begin{exC}
			\item\label{ex:vs:C1} （覆盖定理在有限``数系''上的失效）设 $p$ 是素数，$\mathbb{F}_p=\{0,1,\ldots,p-1\}$ 按模 $p$ 作加法和乘法（可以承认它对加、减、乘以及除以非零元素都封闭. 它不包含 $\mathbb{Q}$，所以不是本书意义下的数域；但第二章中除定理~\ref{thm:Qsmallest} 与覆盖定理以外的定义和结论，在它上面照样成立）.
			\begin{enumerate}
				\item 证明 $\mathbb{F}_p^2$ 恰有 $p+1$ 个一维子空间，并且它们的并是整个 $\mathbb{F}_p^2$.
				\item 证明 $\mathbb{F}_p^2$ 不能写成 $p$ 个真子空间的并.（提示：每个一维子空间恰含 $p-1$ 个非零向量，数一数.）
				\item 指出定理~\ref{thm:cover} 的证明在 $\mathbb{F}_p$ 上从哪一步开始失效. 由第 2 问，``个数不超过 $p$ 时盖不住''，你能把正文的证明改写成这个``定量''的版本吗？
			\end{enumerate}
			\item\label{ex:vs:C2} 证明 $\ln2,\ln3,\ln5$ 在 $\mathbb{Q}$ 上线性无关. 由此说明 $\mathbb{R}$ 作为 $\mathbb{Q}$ 上的向量空间不是有限维的.（提示：算术基本定理.）
			\item\label{ex:vs:C3} 证明：全体实数列构成的向量空间不是有限维的；而``只有有限项非零''的实数列全体是它的子空间，且 $\vv e_1,\vv e_2,\ldots$（第 $i$ 项为 $1$，其余为 $0$）是它的一组``基''——每个元素都能唯一地写成其中\textbf{有限个}向量的线性组合. 这组向量为什么不是全体实数列的基？
			\item\label{ex:vs:C4} （模律）设 $W_1,W_2,W_3$ 是 $V$ 的子空间. 证明 $(W_1\cap W_3)+(W_2\cap W_3)\subset(W_1+W_2)\cap W_3$，在 $\mathbb{R}^2$ 中举例说明等号可以不成立；再证明：若 $W_1\subset W_3$，则等号一定成立.
			\item\label{ex:vs:C5} （有限域上的向量空间）在只含 $0,1$ 两个元素、按模 $2$ 运算的``数系'' $\mathbb{F}_2$ 上考虑 $\mathbb{F}_2^3$. (1) 它有多少个向量？(2) 有多少个一维子空间、多少个二维子空间？(3) 它有多少组（有序）基？推广到 $\mathbb{F}_2^n$.
			\item\label{ex:vs:C6} （塔律）设 $\mathbb{F}\subset\mathbb{K}\subset\mathbb{L}$ 是三个数域，$\mathbb{L}$ 作为 $\mathbb{K}$ 上的向量空间是 $m$ 维的，$\mathbb{K}$ 作为 $\mathbb{F}$ 上的向量空间是 $n$ 维的. 证明 $\mathbb{L}$ 作为 $\mathbb{F}$ 上的向量空间是 $mn$ 维的. 用它重新计算 B8 中数域的维数.
			\item\label{ex:vs:C7} 证明 $\mathbb{Q}(\sqrt2+\sqrt3)=\mathbb{Q}(\sqrt2,\sqrt3)$，即包含 $\sqrt2+\sqrt3$ 的最小数域与包含 $\sqrt2,\sqrt3$ 的最小数域相同.（提示：计算 $(\sqrt2+\sqrt3)^3$.）
			\item\label{ex:vs:C8} 设 $V$ 是 $n$ 维空间，$W$ 是 $k$ 维子空间（$0<k<n$）. 证明 $W$ 有无穷多个补空间. 在 $\mathbb{F}_2^2$ 中，一维子空间有几个补空间？
			\item\label{ex:vs:C9} （Fibonacci 空间的另一组基）在 B6 的记号下，改为研究满足 $a_{n+2}=a_{n+1}+a_n$ 的数列空间. 求出两个公比 $q$ 使 $\{q^n\}$ 属于这个空间，证明它们构成一组基，并推出 Fibonacci 数列的通项（Binet 公式）. 第八章会从另一个角度再次得到它.
		\end{exC}
	}
	
	% =====================================================================
	% 第三章　矩阵
	% =====================================================================
	\newcommand{\ExMatrix}{%
		\exsetheading
		\Agroup
		\begin{exA}
			\item\label{ex:mat:A1} 设 $\vv A=\begin{pmatrix}1&-1&2\\0&3&1\end{pmatrix}$，$\vv B=\begin{pmatrix}2&1\\-1&0\\1&4\end{pmatrix}$. 计算 $\vv{AB}$，$\vv{BA}$，$\vv A^\trans\vv B^\trans$，并验证 $(\vv{AB})^\trans=\vv B^\trans\vv A^\trans$.
			\item\label{ex:mat:A2} 按``列的线性组合''计算 $\begin{pmatrix}1&2&0\\3&-1&4\end{pmatrix}\begin{pmatrix}2\\-1\\1\end{pmatrix}$；再说明 $\vv e_i^\trans\vv A\vv e_j$ 是 $\vv A$ 的哪个元素，$\vv A\vv e_j$、$\vv e_i^\trans\vv A$ 又分别是什么.
			\item\label{ex:mat:A3} 计算 $\begin{pmatrix}1&1\\0&1\end{pmatrix}^n$，$\begin{pmatrix}\lambda&1\\0&\lambda\end{pmatrix}^n$，$\begin{pmatrix}\cos\theta&-\sin\theta\\ \sin\theta&\cos\theta\end{pmatrix}^n$，$\begin{pmatrix}0&1&0\\0&0&1\\0&0&0\end{pmatrix}^n$.
			\item\label{ex:mat:A4} 设 $\vv A=\begin{pmatrix}1&2\\0&1\end{pmatrix}$，$\vv B=\begin{pmatrix}1&0\\3&1\end{pmatrix}$. 计算 $(\vv A+\vv B)^2$ 与 $\vv A^2+2\vv{AB}+\vv B^2$，二者是否相等？$(\vv A+\vv B)(\vv A-\vv B)$ 与 $\vv A^2-\vv B^2$ 呢？
			\item\label{ex:mat:A5} 用初等行变换求下列矩阵的逆：
			\[\begin{pmatrix}1&2&3\\2&2&1\\3&4&3\end{pmatrix},\qquad\begin{pmatrix}1&1&1\\0&1&1\\0&0&1\end{pmatrix},\qquad \begin{pmatrix}2&1&0\\1&2&1\\0&1&2\end{pmatrix},\qquad\begin{pmatrix}0&0&0&1\\0&0&2&0\\0&3&0&0\\4&0&0&0\end{pmatrix}.\]
			\item\label{ex:mat:A6} 解矩阵方程：(1) $\begin{pmatrix}2&5\\1&3\end{pmatrix}\vv X=\begin{pmatrix}4&-6\\2&1\end{pmatrix}$；(2) $\vv X\begin{pmatrix}2&1\\1&1\end{pmatrix}=\begin{pmatrix}1&0\\0&2\end{pmatrix}$.
			\item\label{ex:mat:A7} 求矩阵 $\begin{pmatrix}1&1&a\\1&a&1\\a&1&1\end{pmatrix}$ 的秩（对 $a$ 分情况讨论）.
			\item\label{ex:mat:A8} 求矩阵 $\begin{pmatrix}1&2&-1&0&3\\2&4&1&3&0\\-1&-2&2&1&-3\\1&2&2&3&0\end{pmatrix}$ 的秩，以及列向量组的一个极大无关组，并用它表出其余各列.
			\item\label{ex:mat:A9} 写出 $3$ 阶初等矩阵 $\vv E(1,3)$，$\vv E(2(5))$，$\vv E(3,1(-2))$，并分别写出它们的逆. 计算 $\vv E(3,1(-2))\vv A$ 与 $\vv A\vv E(3,1(-2))$，其中 $\vv A$ 是任一 $3$ 阶方阵，说明效果.
			\item\label{ex:mat:A10} 把 $\begin{pmatrix}2&1\\4&3\end{pmatrix}$ 写成若干个初等矩阵的乘积.
			\item\label{ex:mat:A11} 证明：两个上三角矩阵之积仍是上三角矩阵，且积的对角元等于两者对角元之积. 可逆上三角矩阵的逆也是上三角矩阵吗？
			\item\label{ex:mat:A12} 设 $\vv A,\vv B$ 都是 $n$ 阶对称矩阵. 证明：(1) $\vv A^k$，$\vv A+\vv B$，$\vv P^\trans\vv A\vv P$ 都对称；(2) $\vv{AB}$ 对称当且仅当 $\vv{AB}=\vv{BA}$；(3) 对任意 $m\times n$ 矩阵 $\vv C$，$\vv C^\trans\vv C$ 与 $\vv C\vv C^\trans$ 都对称.
			\item\label{ex:mat:A13} (1) 设 $\vv x,\vv y\in\mathbb{F}^n$，证明 $\tr(\vv x\vv y^\trans)=\vv y^\trans\vv x$；(2) 设 $n$ 阶方阵 $\vv A$ 满足：对一切 $n$ 阶方阵 $\vv X$ 都有 $\tr(\vv{AX})=0$. 证明 $\vv A=\vv O$（提示：取 $\vv X=\vv E_{ji}$）；(3) 证明：不存在 $n$ 阶方阵 $\vv A,\vv B$ 使 $\vv{AB}-\vv{BA}=\vv I$.
			\item\label{ex:mat:A14} 设方阵 $\vv A$ 满足 $\vv A^2-\vv A-2\vv I=\vv O$. 证明 $\vv A$ 与 $\vv A+2\vv I$ 都可逆，并用 $\vv A$ 表示它们的逆. $\vv A+\vv I$ 一定可逆吗？
			\item\label{ex:mat:A15} 证明：若 $\vv A$ 可逆且 $\vv{AB}=\vv{AC}$，则 $\vv B=\vv C$；举例说明 $\vv A$ 不可逆时结论可能不成立.
			\item\label{ex:mat:A16} 设 $\vv A,\vv B$ 可逆. 求 $\begin{pmatrix}\vv A&\vv O\\ \vv O&\vv B\end{pmatrix}^{-1}$ 与 $\begin{pmatrix}\vv O&\vv A\\ \vv B&\vv O\end{pmatrix}^{-1}$，并用后者求 $\begin{pmatrix}0&0&1&2\\0&0&1&3\\2&0&0&0\\0&5&0&0\end{pmatrix}^{-1}$.
			\item\label{ex:mat:A17} 用分块乘法计算 $\begin{pmatrix}1&0&2&1\\0&1&1&3\\0&0&1&0\\0&0&0&1\end{pmatrix}^2$，并猜出 $n$ 次方的一般形式.
			\item\label{ex:mat:A18} 某公司在两个工厂生产三种产品，各工厂每天的产量（件）与三种产品的单位成本、单位售价（元）分别为
			\[\vv Q=\begin{pmatrix}100&80&60\\120&50&90\end{pmatrix}\ (\text{行：工厂；列：产品}),\qquad \vv C=\begin{pmatrix}20&30\\15&25\\40&55\end{pmatrix}\ (\text{列：成本、售价}).\]
			计算 $\vv Q\vv C$，说明它每个元素的含义，并求两个工厂每天的利润.
			\item\label{ex:mat:A19} 四座城市之间的直达航线为：$1\leftrightarrow2$，$1\leftrightarrow3$，$2\leftrightarrow3$，$3\leftrightarrow4$（双向）. 写出邻接矩阵 $\vv A$，计算 $\vv A^2$ 与 $\vv A^3$. $(\vv A^2)_{ii}$ 与城市 $i$ 的航线数有什么关系？从城市 $1$ 到城市 $4$ 恰好乘三次航班有几种走法？
			\item\label{ex:mat:A20} 设 $\vv A$ 是实反对称矩阵. 证明：对一切 $\vv x\in\mathbb{R}^n$ 有 $\vv x^\trans\vv A\vv x=0$；$\vv A^2$ 是对称矩阵.
			\item\label{ex:mat:A21} 证明：若 $\vv A$ 是 $m\times n$ 矩阵，$\vv P$ 是 $m$ 阶可逆矩阵，则 $\vv{PA}$ 与 $\vv A$ 的简化阶梯形相同.
			\item\label{ex:mat:A22} 证明：任一 $n$ 阶方阵都可以唯一地写成一个对称矩阵与一个反对称矩阵之和.（提示：$\vv A=\frac12(\vv A+\vv A^\trans)+\frac12(\vv A-\vv A^\trans)$.）
		\end{exA}
		
		\Bgroup
		\begin{exB}
			\item\label{ex:mat:B1} 分别求出全部满足 (1) $\vv A^2=\vv O$；(2) $\vv A^2=\vv I$；(3) $\vv A^2=\vv A$ 的二阶实方阵.
			\item\label{ex:mat:B2} 设 $a_1,\ldots,a_n$ 互不相同. 证明：与 $\diag(a_1,\ldots,a_n)$ 可交换的矩阵一定是对角矩阵.
			\item\label{ex:mat:B3} 证明：与\textbf{一切} $n$ 阶方阵都可交换的矩阵必为数量矩阵 $c\vv I$.（提示：考虑与 $\vv E_{ij}$ 交换.）
			\item\label{ex:mat:B4} 设 $\vv A,\vv B$ 都是幂零矩阵且 $\vv{AB}=\vv{BA}$. 证明 $\vv A+\vv B$ 与 $\vv{AB}$ 都是幂零矩阵. 举例说明``可交换''不能去掉.
			\item\label{ex:mat:B5} 证明：$\rk\begin{pmatrix}\vv A&\vv O\\ \vv O&\vv B\end{pmatrix}=\rk\vv A+\rk\vv B$，$\rk\begin{pmatrix}\vv A&\vv C\\ \vv O&\vv B\end{pmatrix}\ge\rk\vv A+\rk\vv B$，并举例说明后者可以取严格不等号.
			\item\label{ex:mat:B6} （Sherman–Morrison 公式）设 $\vv u,\vv v\in\mathbb{F}^n$，$1+\vv v^\trans\vv u\ne0$. 证明
			\[(\vv I+\vv u\vv v^\trans)^{-1}=\vv I-\frac{\vv u\vv v^\trans}{1+\vv v^\trans\vv u}.\]
			\item\label{ex:mat:B7} 设 $\vv A$ 为 $m\times n$ 矩阵，$\vv B$ 为 $n\times m$ 矩阵. 证明：若 $\vv I_m+\vv{AB}$ 可逆，则 $\vv I_n+\vv{BA}$ 也可逆，并且 $(\vv I_n+\vv{BA})^{-1}=\vv I_n-\vv B(\vv I_m+\vv{AB})^{-1}\vv A$.
			\item\label{ex:mat:B8} 证明：秩为 $1$ 的 $m\times n$ 矩阵都可以写成 $\vv u\vv v^\trans$（$\vv u\in\mathbb{F}^m$，$\vv v\in\mathbb{F}^n$ 非零）. 对这样的矩阵证明 $\vv A^2=(\tr\vv A)\vv A$（$m=n$ 时），并由此求 $\vv A^k$.
			\item\label{ex:mat:B9} 证明：秩为 $r$ 的矩阵可以写成 $r$ 个秩为 $1$ 的矩阵之和，但不能写成少于 $r$ 个秩为 $1$ 的矩阵之和.
			\item\label{ex:mat:B10} （满秩分解）证明：秩为 $r$ 的 $m\times n$ 矩阵 $\vv A$ 可以写成 $\vv A=\vv{BC}$，其中 $\vv B$ 是列满秩的 $m\times r$ 矩阵，$\vv C$ 是行满秩的 $r\times n$ 矩阵. 对 A8 中的矩阵写出一个满秩分解.
			\item\label{ex:mat:B11} 计算 $\begin{pmatrix}1&1&1\\1&1&1\\1&1&1\end{pmatrix}^n$ 与 $\begin{pmatrix}2&1&0\\0&2&1\\0&0&2\end{pmatrix}^n$.（提示：后者写成 $2\vv I+\vv N$.）
			\item\label{ex:mat:B12} 证明：$n$ 阶严格上三角矩阵（对角线及以下全为零）$\vv N$ 满足 $\vv N^n=\vv O$. 更精确地，$\vv N^k$ 的非零元只可能出现在第 $k$ 条``上对角线''及其右上方.
			\item\label{ex:mat:B13} 设 $\vv A$ 为实矩阵. 证明：若 $\vv A^\trans\vv A=\vv O$，则 $\vv A=\vv O$；若 $\tr(\vv A^\trans\vv A)=0$，结论同样成立. 复矩阵呢？
			\item\label{ex:mat:B14} 每行、每列恰有一个 $1$、其余元素为 $0$ 的方阵称为\textbf{置换矩阵}. 证明：置换矩阵之积仍是置换矩阵，置换矩阵 $\vv P$ 满足 $\vv P^\trans\vv P=\vv I$；$\vv{PA}$ 与 $\vv{AP}$ 分别对 $\vv A$ 做了什么？
			\item\label{ex:mat:B15} 对 $\vv A=\begin{pmatrix}1&2&1\\2&4&3\\1&2&2\end{pmatrix}$，求可逆矩阵 $\vv P,\vv Q$ 使 $\vv{PAQ}$ 为相抵标准形.
			\item\label{ex:mat:B16} 设 $\vv A$ 是 $n$ 阶方阵，$\vv A^k=\vv O$. 证明 $\vv I+\vv A$ 可逆，并求其逆. 若 $\vv A^2=2\vv A$，$\vv I-\vv A$ 是否可逆？
		\end{exB}
		
		\Cgroup
		\begin{exC}
			\item\label{ex:mat:C1} 设 $f:M_n(\mathbb{F})\to\mathbb{F}$ 满足 $f(\vv A+\vv B)=f(\vv A)+f(\vv B)$，$f(k\vv A)=kf(\vv A)$，且 $f(\vv{AB})=f(\vv{BA})$ 对一切 $\vv A,\vv B$ 成立. 证明 $f=c\cdot\tr$ 对某个常数 $c$ 成立.（提示：计算 $f(\vv E_{ij})$.）
			\item\label{ex:mat:C2} 设 $n$ 阶方阵 $\vv A,\vv B$ 满足 $\vv A+\vv B=\vv{AB}$. 证明 $\vv{AB}=\vv{BA}$.（提示：考察 $(\vv I-\vv A)(\vv I-\vv B)$.）
			\item\label{ex:mat:C3} （Strassen 算法）对二阶分块矩阵的乘法，直接计算需要 $8$ 次块乘法. 验证下面的公式只需 $7$ 次：令
			$M_1=(A_{11}+A_{22})(B_{11}+B_{22})$，$M_2=(A_{21}+A_{22})B_{11}$，$M_3=A_{11}(B_{12}-B_{22})$，$M_4=A_{22}(B_{21}-B_{11})$，$M_5=(A_{11}+A_{12})B_{22}$，$M_6=(A_{21}-A_{11})(B_{11}+B_{12})$，$M_7=(A_{12}-A_{22})(B_{21}+B_{22})$，则
			$C_{11}=M_1+M_4-M_5+M_7$，$C_{12}=M_3+M_5$，$C_{21}=M_2+M_4$，$C_{22}=M_1-M_2+M_3+M_6$. 由此说明 $n=2^k$ 阶矩阵乘法可以用约 $n^{\log_27}\approx n^{2.81}$ 次数的乘法完成.
			\item\label{ex:mat:C4} 设 $\vv A$ 是幂等矩阵. 用满秩分解 $\vv A=\vv{BC}$ 证明 $\vv{CB}=\vv I_r$，进而证明 $\rk\vv A=\tr\vv A$.
			\item\label{ex:mat:C5} （Kronecker 积）对 $m\times n$ 矩阵 $\vv A$ 与 $p\times q$ 矩阵 $\vv B$，定义 $\vv A\otimes\vv B=(a_{ij}\vv B)$（$mp\times nq$ 分块矩阵）. 证明 $(\vv A\otimes\vv B)(\vv C\otimes\vv D)=(\vv{AC})\otimes(\vv{BD})$，并由此说明 $\vv A,\vv B$ 可逆时 $(\vv A\otimes\vv B)^{-1}=\vv A^{-1}\otimes\vv B^{-1}$.
			\item\label{ex:mat:C6} （Hadamard 矩阵）元素全为 $\pm1$ 且满足 $\vv H\vv H^\trans=n\vv I$ 的 $n$ 阶矩阵称为 Hadamard 矩阵. 验证 $\vv H_2=\begin{pmatrix}1&1\\1&-1\end{pmatrix}$ 是，并证明 $\vv H_2\otimes\vv H_2$，$\vv H_2\otimes\vv H_2\otimes\vv H_2$ 也是. 证明：$n>2$ 时 Hadamard 矩阵的阶数必是 $4$ 的倍数.（这类矩阵用于实验设计与纠错编码.）
			\item\label{ex:mat:C7} 证明：全体形如 $\begin{pmatrix}a&-b\\b&a\end{pmatrix}$（$a,b\in\mathbb{R}$）的矩阵关于矩阵的加法和乘法``像复数一样''运算：加法、乘法封闭，乘法可交换，非零元可逆. 写出它与 $\mathbb{C}$ 之间的对应，并用它验证 $(1+\mathrm i)^8=16$.
			\item\label{ex:mat:C8} 设 $\vv A$ 是元素非负、各列之和都小于 $1$ 的方阵. 证明 $\vv A^k$ 的每个元素都随 $k\to\infty$ 趋于 $0$，并且级数 $\vv I+\vv A+\vv A^2+\cdots$（按元素）收敛到 $(\vv I-\vv A)^{-1}$. 解释它在投入产出模型中的意义：``直接消耗 $+$ 一次间接消耗 $+$ 二次间接消耗 $+\cdots$''.
			\item\label{ex:mat:C9} 设 $\vv A$ 为 $m\times n$ 矩阵. 证明：$\vv A$ 有右逆（存在 $\vv C$ 使 $\vv{AC}=\vv I_m$）当且仅当 $\rk\vv A=m$；有左逆当且仅当 $\rk\vv A=n$. 非方阵的单侧逆唯一吗？
		\end{exC}
	}
	
	% =====================================================================
	% 第四章　线性映射
	% =====================================================================
	\newcommand{\ExLinearMap}{%
		\exsetheading
		\Agroup
		\begin{exA}
			\item\label{ex:lmap:A1} 判断下列映射是否线性，是线性映射的求出核：
			\begin{enumerate}
				\item $\mathbb{R}^2\to\mathbb{R}^2$，$(x,y)^\trans\mapsto(x+y,\,2x)^\trans$；
				\item $\mathbb{R}^2\to\mathbb{R}^2$，$(x,y)^\trans\mapsto(x+1,\,y)^\trans$；
				\item $\mathbb{R}^2\to\mathbb{R}$，$(x,y)^\trans\mapsto|x|+y$；
				\item $\mathbb{R}[x]\to\mathbb{R}[x]$，$f(x)\mapsto f(x+1)$；
				\item $\mathbb{R}[x]\to\mathbb{R}[x]$，$f(x)\mapsto xf'(x)$；
				\item $M_n(\mathbb{R})\to M_n(\mathbb{R})$，$\vv X\mapsto\vv X^\trans$；
				\item $M_2(\mathbb{R})\to\mathbb{R}$，$\vv X\mapsto x_{11}x_{22}-x_{12}x_{21}$.
			\end{enumerate}
			\item\label{ex:lmap:A2} 求下列矩阵所定义的线性映射 $\vv x\mapsto\vv{Ax}$ 的核与像（各给出一组基），并验证秩—零化度定理：
			\[\begin{pmatrix}1&2&-1\\2&4&-2\end{pmatrix},\qquad\begin{pmatrix}1&0&2&1\\0&1&-1&1\\1&1&1&2\end{pmatrix}.\]
			\item\label{ex:lmap:A3} 写出线性映射 $\mathscr T:\mathbb{R}^3\to\mathbb{R}^2$，$\mathscr T(x_1,x_2,x_3)^\trans=(x_1-2x_2+x_3,\ 3x_2+4x_3)^\trans$ 在标准基下的矩阵.
			\item\label{ex:lmap:A4} 求下列平面变换在标准基下的矩阵：(1) 逆时针旋转 $90^\circ$；(2) 关于直线 $y=x$ 的反射；(3) 向直线 $y=x$ 的正投影；(4) 沿 $x$ 方向的剪切 $(x,y)\mapsto(x+2y,y)$. 计算 (2)(3) 的矩阵的平方，并解释结果.
			\item\label{ex:lmap:A5} 求求导变换 $\mathscr D:\mathbb{R}[x]_3\to\mathbb{R}[x]_3$ 在基 $1,x,x^2$ 下的矩阵，以及在基 $1,\ x-1,\ (x-1)^2$ 下的矩阵.
			\item\label{ex:lmap:A6} 求转置变换 $\vv X\mapsto\vv X^\trans$ 在 $M_2(\mathbb{F})$ 的基 $\vv E_{11},\vv E_{12},\vv E_{21},\vv E_{22}$ 下的矩阵. 它的平方是什么？
			\item\label{ex:lmap:A7} 在 $\mathbb{R}^3$ 中，求从基 $\vv\varepsilon_1=(1,0,0)^\trans$，$\vv\varepsilon_2=(1,1,0)^\trans$，$\vv\varepsilon_3=(1,1,1)^\trans$ 到基 $\vv\eta_1=(1,2,1)^\trans$，$\vv\eta_2=(2,3,3)^\trans$，$\vv\eta_3=(3,7,1)^\trans$ 的过渡矩阵，并求 $\vv\alpha=(3,4,2)^\trans$ 在两组基下的坐标.
			\item\label{ex:lmap:A8} 设线性变换 $\mathscr T$ 在基 $\vv\varepsilon_1,\vv\varepsilon_2$ 下的矩阵为 $\begin{pmatrix}2&1\\1&2\end{pmatrix}$. 求它在基 $\vv\eta_1=\vv\varepsilon_1+\vv\varepsilon_2$，$\vv\eta_2=\vv\varepsilon_1-\vv\varepsilon_2$ 下的矩阵.
			\item\label{ex:lmap:A9} 证明：线性映射的复合仍是线性映射；可逆（双射）线性映射的逆映射仍是线性映射.
			\item\label{ex:lmap:A10} 证明：线性映射 $\mathscr T$ 是单射，当且仅当它把每个线性无关组都映为线性无关组.
			\item\label{ex:lmap:A11} 是否存在线性变换 $\mathscr T:\mathbb{R}^2\to\mathbb{R}^2$，使
			(1) $\mathscr T(1,1)^\trans=(1,0)^\trans$，$\mathscr T(1,-1)^\trans=(0,1)^\trans$，$\mathscr T(2,0)^\trans=(1,1)^\trans$？
			(2) 同上，但最后一个条件改为 $\mathscr T(2,0)^\trans=(1,2)^\trans$？
			存在时求出 $\mathscr T(x,y)^\trans$.
			\item\label{ex:lmap:A12} (1) 线性映射 $\mathscr T:\mathbb{R}^5\to\mathbb{R}^3$ 是满射，$\dim\ker\mathscr T$ 是多少？(2) 是否存在从 $\mathbb{R}^3$ 到 $\mathbb{R}^5$ 的满射线性映射？是否存在从 $\mathbb{R}^5$ 到 $\mathbb{R}^3$ 的单射线性映射？
			\item\label{ex:lmap:A13} 设 $\mathscr T:U\to V$，$\mathscr S:V\to W$ 线性. 证明 $\ker\mathscr T\subset\ker(\mathscr{ST})$，$\imm(\mathscr{ST})\subset\imm\mathscr S$，从而 $\rk(\mathscr{ST})\le\min\{\rk\mathscr S,\rk\mathscr T\}$.
			\item\label{ex:lmap:A14} 求迹映射 $\tr:M_n(\mathbb{F})\to\mathbb{F}$ 的核的维数.
			\item\label{ex:lmap:A15} 证明：相似矩阵有相同的迹（利用命题~\ref{prop:trace}(3)）. 由此判断 $\begin{pmatrix}1&2\\0&3\end{pmatrix}$ 与 $\begin{pmatrix}1&0\\0&2\end{pmatrix}$ 是否相似.
			\item\label{ex:lmap:A16} 证明：与单位矩阵相似的矩阵只有单位矩阵本身. 由此说明 $\begin{pmatrix}1&1\\0&1\end{pmatrix}$ 与 $\vv I$ 不相似.
			\item\label{ex:lmap:A17} 定义求值映射 $\mathscr E:\mathbb{R}[x]_3\to\mathbb{R}^3$，$f\mapsto(f(0),f(1),f(2))^\trans$. 求它在基 $1,x,x^2$ 与标准基下的矩阵，并证明 $\mathscr E$ 是同构.
			\item\label{ex:lmap:A18} 设 $\mathscr T$ 是 $\mathbb{R}^3$ 上的线性变换，$\mathscr T\vv e_1=(1,1,0)^\trans$，$\mathscr T\vv e_2=(0,1,1)^\trans$，$\mathscr T\vv e_3=(1,2,1)^\trans$. 求 $\mathscr T(x_1,x_2,x_3)^\trans$，$\ker\mathscr T$ 与 $\imm\mathscr T$.
			\item\label{ex:lmap:A19} 称 $\vv x\mapsto\vv{Ax}+\vv b$ 为\textbf{仿射变换}. 证明两个仿射变换的复合仍是仿射变换；仿射变换是线性变换当且仅当 $\vv b=\zero$.
			\item\label{ex:lmap:A20} 设 $\vv R(\theta)=\begin{pmatrix}\cos\theta&-\sin\theta\\ \sin\theta&\cos\theta\end{pmatrix}$ 是平面上绕原点旋转角 $\theta$ 的变换在标准基下的矩阵. 从几何上说明``先转 $\beta$ 再转 $\alpha$''就是转 $\alpha+\beta$，从而 $\vv R(\alpha)\vv R(\beta)=\vv R(\alpha+\beta)$；比较两边矩阵的元素，推出两角和的正弦、余弦公式.
		\end{exA}
		
		\Bgroup
		\begin{exB}
			\item\label{ex:lmap:B1} 记 $L(V,W)$ 为 $V$ 到 $W$ 的全体线性映射. 证明它关于映射的加法与数乘构成向量空间，且 $\dim L(V,W)=\dim V\cdot\dim W$.（提示：取定基后，线性映射与矩阵一一对应.）
			\item\label{ex:lmap:B2} 设 $\mathscr T$ 是 $n$ 维空间 $V$ 上的线性变换，$\mathscr T^2=\mathscr O$. 证明 $\imm\mathscr T\subset\ker\mathscr T$，从而 $\rk\mathscr T\le n/2$；并证明存在一组基，使 $\mathscr T$ 的矩阵只在 $(1,2),(3,4),\ldots,(2r-1,2r)$ 这 $r$ 个位置上为 $1$，其余元素全为 $0$，其中 $r=\rk\mathscr T$.
			\item\label{ex:lmap:B3} 设 $\mathscr T,\mathscr S$ 是有限维空间上的线性变换. 证明 $\dim\ker(\mathscr{ST})\le\dim\ker\mathscr S+\dim\ker\mathscr T$.
			\item\label{ex:lmap:B4} 设 $V$ 有限维，$\mathscr S,\mathscr T$ 是 $V$ 上的线性变换且 $\mathscr{ST}=\mathscr I$. 证明 $\mathscr{TS}=\mathscr I$.
			\item\label{ex:lmap:B5} 在无限维空间 $\mathbb{R}[x]$ 上，令 $\mathscr D$ 为求导，$\mathscr J$ 为 $f\mapsto\int_0^xf(t)\,\mathrm dt$. 证明 $\mathscr{DJ}=\mathscr I$，但 $\mathscr{JD}\ne\mathscr I$；$\mathscr D$ 是满射但不是单射，$\mathscr J$ 是单射但不是满射. 这与 B4 是否矛盾？
			\item\label{ex:lmap:B6} 固定 $\vv A\in M_n(\mathbb{F})$，令 $\mathscr T(\vv X)=\vv{AX}-\vv{XA}$. 证明 $\imm\mathscr T$ 含于迹为零的矩阵全体，从而 $\mathscr T$ 不是满射；并证明 $\ker\mathscr T$ 包含 $\vv I,\vv A,\vv A^2,\ldots$.
			\item\label{ex:lmap:B7} 设 $\vv A=\begin{pmatrix}a&b\\c&d\end{pmatrix}$. 求线性变换 $\vv X\mapsto\vv{AX}$ 在 $M_2(\mathbb{F})$ 的基 $\vv E_{11},\vv E_{21},\vv E_{12},\vv E_{22}$ 下的矩阵. 若换成 $\vv X\mapsto\vv{XA}$ 呢？
			\item\label{ex:lmap:B8} 设线性变换 $\mathscr P$ 满足 $\mathscr P^2=\mathscr P$. 证明存在一组基使 $\mathscr P$ 的矩阵为 $\begin{pmatrix}\vv I_r&\vv O\\ \vv O&\vv O\end{pmatrix}$，其中 $r=\rk\mathscr P$.
			\item\label{ex:lmap:B9} 设线性变换 $\mathscr T$ 满足 $\mathscr T^2=\mathscr I$. 证明 $V=\ker(\mathscr T-\mathscr I)\oplus\ker(\mathscr T+\mathscr I)$，从而存在一组基使 $\mathscr T$ 的矩阵为 $\diag(\vv I_p,-\vv I_q)$.（提示：$\vv\alpha=\frac12(\vv\alpha+\mathscr T\vv\alpha)+\frac12(\vv\alpha-\mathscr T\vv\alpha)$.）
			\item\label{ex:lmap:B10} 设 $\mathscr T:V\to W$ 线性，$U$ 是 $V$ 的子空间. 证明 $\dim\mathscr T(U)=\dim U-\dim(U\cap\ker\mathscr T)$.
			\item\label{ex:lmap:B11} 在由 $\mathrm e^x\cos x$，$\mathrm e^x\sin x$ 张成的二维函数空间中，求求导变换的矩阵，证明它可逆，并利用逆矩阵求出 $\int\mathrm e^x\sin x\,\mathrm dx$（不必用分部积分）.
			\item\label{ex:lmap:B12} （齐次坐标）平面上的平移 $(x,y)\mapsto(x+a,y+b)$ 不是线性的. 把点 $(x,y)$ 写成三维向量 $(x,y,1)^\trans$，证明平移、旋转、伸缩都可以写成 $3\times3$ 矩阵的乘法，且仿射变换的复合对应矩阵的乘积. 写出``先绕原点旋转 $90^\circ$，再平移 $(1,2)$''与``先平移再旋转''的矩阵，它们相等吗？（计算机图形学正是这样处理变换的.）
			\item\label{ex:lmap:B13} 证明：若 $\vv A$ 与 $\vv B$ 相似，则 $\vv A^k$ 与 $\vv B^k$ 相似，$\vv A^\trans$ 与 $\vv B^\trans$ 相似，$f(\vv A)$ 与 $f(\vv B)$ 相似（$f$ 为多项式）.
			\item\label{ex:lmap:B14} 证明 $\begin{pmatrix}1&0\\0&0\end{pmatrix}$ 与 $\begin{pmatrix}0&1\\0&0\end{pmatrix}$ 相抵但不相似. 从``线性映射 $V\to W$''与``线性变换 $V\to V$''的区别解释：为什么相抵只需看秩，而相似要复杂得多？
			\item\label{ex:lmap:B15} 设 $\mathscr T$ 是 $V$ 上的线性变换，$\vv\alpha\in V$ 满足 $\mathscr T^{k-1}\vv\alpha\ne\zero$，$\mathscr T^k\vv\alpha=\zero$. 证明 $\vv\alpha,\mathscr T\vv\alpha,\ldots,\mathscr T^{k-1}\vv\alpha$ 线性无关. 若 $k=n=\dim V$，求 $\mathscr T$ 在基 $\mathscr T^{n-1}\vv\alpha,\ldots,\mathscr T\vv\alpha,\vv\alpha$ 下的矩阵.
		\end{exB}
		
		\Cgroup
		\begin{exC}
			\item\label{ex:lmap:C1} （Fitting 分解）设 $\mathscr T$ 是 $n$ 维空间 $V$ 上的线性变换. 证明：
			\begin{enumerate}
				\item $\ker\mathscr T\subset\ker\mathscr T^2\subset\cdots$，$\imm\mathscr T\supset\imm\mathscr T^2\supset\cdots$，且两条链都在某一步之后不再变化；
				\item $V=\ker\mathscr T^n\oplus\imm\mathscr T^n$，两者都是 $\mathscr T$ 的不变子空间，$\mathscr T$ 在前者上幂零，在后者上可逆.
			\end{enumerate}
			\item\label{ex:lmap:C2} （商空间）设 $W$ 是 $V$ 的子空间. 对 $\vv\alpha\in V$，称集合 $\vv\alpha+W$ 为 $\vv\alpha$ 所在的陪集，全体陪集记为 $V/W$. 定义 $(\vv\alpha+W)+(\vv\beta+W)=(\vv\alpha+\vv\beta)+W$，$k(\vv\alpha+W)=k\vv\alpha+W$.
			\begin{enumerate}
				\item 证明这两个运算是良定义的，且 $V/W$ 构成向量空间；
				\item 证明 $\dim V/W=\dim V-\dim W$；
				\item （第一同构定理/同态基本定理）对线性映射 $\mathscr T:V\to U$，证明 $V/\ker\mathscr T\cong\imm\mathscr T$. 这给出秩—零化度定理的另一种理解.
			\end{enumerate}
			\item\label{ex:lmap:C3} 把 $\mathbb{C}$ 看作 $\mathbb{R}$ 上的二维空间，基为 $1,\mathrm i$. (1) 求``乘以 $a+b\mathrm i$''这个变换的矩阵；(2) 求``取共轭''这个变换的矩阵；(3) 证明 $\mathbb{R}$-线性变换 $\begin{pmatrix}p&q\\r&s\end{pmatrix}$ 是``乘以某个复数''，当且仅当 $p=s$，$q=-r$.
			\item\label{ex:lmap:C4} 证明：$n$ 维空间上的任一线性变换 $\mathscr T$，都存在次数不超过 $n^2$ 的非零多项式 $f$ 使 $f(\mathscr T)=\mathscr O$.（提示：$\mathscr I,\mathscr T,\ldots,\mathscr T^{n^2}$ 在 $L(V,V)$ 中线性相关. 第八章的 Cayley–Hamilton 定理会把 $n^2$ 改进为 $n$.）
			\item\label{ex:lmap:C5} （Haar 小波）在 $\mathbb{R}^4$ 中取基 $\vv h_1=(1,1,1,1)^\trans$，$\vv h_2=(1,1,-1,-1)^\trans$，$\vv h_3=(1,-1,0,0)^\trans$，$\vv h_4=(0,0,1,-1)^\trans$. 求信号 $\vv x=(8,6,3,4)^\trans$ 在这组基下的坐标，并解释每个坐标的含义. 若把绝对值最小的坐标改为 $0$ 再变换回来，误差是多少？这是图像与音频压缩的思想之一.
			\item\label{ex:lmap:C6} 设 $\mathscr T$ 是 $V$ 上的线性变换. 若 $\mathscr T$ 与 $V$ 上的一切线性变换都可交换，证明 $\mathscr T=c\mathscr I$. 若只要求与一切\textbf{可逆}线性变换可交换呢？
			\item\label{ex:lmap:C7} 设 $\mathscr T:V\to V$ 线性，$\dim V=n$. 证明：$\rk\mathscr T^k-\rk\mathscr T^{k+1}$ 关于 $k$ 单调不增. 由此说明：若 $\rk\mathscr T=\rk\mathscr T^2$，则对一切 $k$，$\rk\mathscr T^k=\rk\mathscr T$.
			\item\label{ex:lmap:C8} （差分算子）在全体实数列构成的空间上，定义左移 $\mathscr S(a_1,a_2,a_3,\ldots)=(a_2,a_3,\ldots)$. 证明 $\mathscr S$ 是线性的、满射但不是单射；满足 $a_{n+2}=a_{n+1}+a_n$ 的数列全体恰好是 $\ker(\mathscr S^2-\mathscr S-\mathscr I)$. 由第二章章末习题~\ref{ex:vs:C9}，这个核是二维的.
		\end{exC}
	}
	
	% =====================================================================
	% 第五章　行列式
	% =====================================================================
	\newcommand{\ExDeterminant}{%
		\exsetheading
		\Agroup
		\begin{exA}
			\item\label{ex:det:A1} 求下列排列的逆序数并判断奇偶：$315426$，$42531$，$13\cdots(2n-1)\,24\cdots(2n)$.
			\item\label{ex:det:A2} (1) 在 $4$ 阶行列式中，项 $a_{13}a_{24}a_{32}a_{41}$ 带什么符号？(2) 选择 $i,j$，使 $a_{1i}a_{23}a_{35}a_{5j}a_{44}$ 成为 $5$ 阶行列式中带正号的一项.
			\item\label{ex:det:A3} 计算：
			\[\begin{vmatrix}2&1&-1\\3&0&2\\1&4&5\end{vmatrix},\qquad
			\begin{vmatrix}1&2&0&1\\2&1&1&0\\0&1&2&1\\1&0&1&2\end{vmatrix},\qquad
			\begin{vmatrix}3&1&1&1\\1&3&1&1\\1&1&3&1\\1&1&1&3\end{vmatrix}.\]
			\item\label{ex:det:A4} 计算 $\begin{vmatrix}1&-1&2&3\\2&0&1&-1\\0&1&-1&2\\3&1&0&1\end{vmatrix}$，分别用``打洞''化三角形与按第一列展开两种方法，比较哪种更快.
			\item\label{ex:det:A5} 计算 $n$ 阶行列式 $\begin{vmatrix}0&\cdots&0&a_1\\0&\cdots&a_2&0\\ \vdots&&&\vdots\\a_n&\cdots&0&0\end{vmatrix}$（只有副对角线上有元素）.
			\item\label{ex:det:A6} 不展开，证明：
			\[\begin{vmatrix}a+b&b+c&c+a\\a_1+b_1&b_1+c_1&c_1+a_1\\a_2+b_2&b_2+c_2&c_2+a_2\end{vmatrix}=2\begin{vmatrix}a&b&c\\a_1&b_1&c_1\\a_2&b_2&c_2\end{vmatrix}.\]
			\item\label{ex:det:A7} 解方程 $\begin{vmatrix}x&1&1\\1&x&1\\1&1&x\end{vmatrix}=0$ 与 $\begin{vmatrix}1&1&1\\1&2&x\\1&4&x^2\end{vmatrix}=0$.
			\item\label{ex:det:A8} 设 $\vv A$ 为 $3$ 阶方阵，$\det\vv A=3$. 求 $\det(2\vv A)$，$\det(-\vv A)$，$\det(\vv A^{-1})$，$\det(\vv A^*)$，$\det(\vv A^\trans\vv A)$，$\det(3\vv A^{-1}-2\vv A^*)$.
			\item\label{ex:det:A9} 证明：(1) 若 $n$ 阶方阵满足 $\vv A^\trans\vv A=\vv I$，则 $\det\vv A=\pm1$；(2) 若 $\vv A^2=\vv I$，则 $\det\vv A=\pm1$；(3) 若 $\vv A^2=\vv A$ 且 $\vv A\ne\vv I$，则 $\det\vv A=0$.
			\item\label{ex:det:A10} 求下列矩阵的伴随矩阵与逆矩阵：\[\begin{pmatrix}a&b\\c&d\end{pmatrix},\qquad\begin{pmatrix}1&2&3\\0&1&4\\5&6&0\end{pmatrix}.\]
			\item\label{ex:det:A11} 用 Cramer 法则解方程组 $2x_1+x_2+x_3=3$，$x_1+3x_2+x_3=0$，$x_1+x_2+4x_3=8$.
			\item\label{ex:det:A12} 求以 $(1,1)$，$(4,2)$，$(2,5)$ 为顶点的三角形的面积；判断 $(1,2)$，$(3,5)$，$(7,11)$ 三点是否共线.
			\item\label{ex:det:A13} 求以 $(1,0,1)^\trans$，$(2,1,0)^\trans$，$(0,3,1)^\trans$ 为棱的平行六面体的体积，以及以原点和这三个向量的终点为顶点的四面体的体积.
			\item\label{ex:det:A14} $\lambda$ 取何值时，齐次方程组 $(5-\lambda)x+2y+2z=0$，$2x+(6-\lambda)y=0$，$2x+(4-\lambda)z=0$ 有非零解？
			\item\label{ex:det:A15} 设 \[\vv A=\begin{pmatrix}1&2&3&4\\2&3&4&1\\3&4&1&2\\4&1&2&3\end{pmatrix}.\] 求 $A_{11}+A_{12}+A_{13}+A_{14}$（$A_{1j}$ 是代数余子式）.（提示：它是把第一行换成全 $1$ 后的行列式.）
			\item\label{ex:det:A16} 证明：过平面上两个不同点 $(x_1,y_1)$，$(x_2,y_2)$ 的直线方程可以写成 \[\begin{vmatrix}x&y&1\\x_1&y_1&1\\x_2&y_2&1\end{vmatrix}=0.\]
			\item\label{ex:det:A17} 设 $\vv A,\vv B$ 分别是 $m$ 阶、$n$ 阶方阵，证明 $\begin{vmatrix}\vv A&\vv O\\ \vv C&\vv B\end{vmatrix}=\det\vv A\det\vv B$，并计算 $\begin{vmatrix}\vv O&\vv A\\ \vv B&\vv O\end{vmatrix}$.
			\item\label{ex:det:A18} 计算 $4$ 阶三对角行列式 $\begin{vmatrix}2&-1&0&0\\-1&2&-1&0\\0&-1&2&-1\\0&0&-1&2\end{vmatrix}$.
			\item\label{ex:det:A19} 设 $\vv A$ 是 $n$ 阶可逆矩阵，$n\ge2$. 证明 $\vv A^*$ 可逆，且 $(\vv A^*)^{-1}=\frac1{\det\vv A}\vv A=(\vv A^{-1})^*$.
			\item\label{ex:det:A20} 用行列式判断：向量 $(1,2,3)^\trans$，$(4,5,6)^\trans$，$(7,8,9)^\trans$ 是否线性无关？把 $9$ 改成 $10$ 呢？
			\item\label{ex:det:A21} 设 $\vv A$ 是 $n$ 阶反对称矩阵（$\vv A^\trans=-\vv A$），$n$ 为奇数. 证明 $\det\vv A=0$.（提示：对 $\vv A^\trans=-\vv A$ 两边取行列式.）
		\end{exA}
		
		\Bgroup
		\begin{exB}
			\item\label{ex:det:B1} 计算 $n$ 阶三对角行列式 $D_n=\begin{vmatrix}2&-1&&\\-1&2&\ddots&\\&\ddots&\ddots&-1\\&&-1&2\end{vmatrix}$.（提示：按第一行展开得 $D_n=2D_{n-1}-D_{n-2}$.）
			\item\label{ex:det:B2} 计算 $n$ 阶行列式 $\begin{vmatrix}a&b&\cdots&b\\b&a&\cdots&b\\ \vdots&&\ddots&\vdots\\b&b&\cdots&a\end{vmatrix}$ 与 $\begin{vmatrix}x&a&\cdots&a\\-a&x&\cdots&a\\ \vdots&&\ddots&\vdots\\-a&-a&\cdots&x\end{vmatrix}$.
			\item\label{ex:det:B3} （爪形行列式）设 $a_1\cdots a_n\ne0$，计算 $\begin{vmatrix}a_0&1&1&\cdots&1\\1&a_1&0&\cdots&0\\1&0&a_2&\cdots&0\\ \vdots&&&\ddots&\\1&0&0&\cdots&a_n\end{vmatrix}$.
			\item\label{ex:det:B4} （矩阵行列式引理）证明 $\det(\vv I_n+\vv u\vv v^\trans)=1+\vv v^\trans\vv u$，更一般地，$\vv A$ 可逆时 $\det(\vv A+\vv u\vv v^\trans)=(1+\vv v^\trans\vv A^{-1}\vv u)\det\vv A$. 用它重做 B2 的第一个行列式.
			\item\label{ex:det:B5} 设 $\vv A$ 为 $2\times3$ 矩阵，$\vv B$ 为 $3\times2$ 矩阵. 证明 $\det(\vv{AB})$ 等于 $\vv A$ 的三个 $2$ 阶子式与 $\vv B$ 的对应 $2$ 阶子式乘积之和（这是 Cauchy–Binet 公式最简单的情形）. 由此说明 $\vv{BA}$（$3$ 阶方阵）的行列式为零.
			\item\label{ex:det:B6} 设 $\vv A$ 是整数矩阵. 证明：$\vv A^{-1}$ 存在且也是整数矩阵，当且仅当 $\det\vv A=\pm1$.
			\item\label{ex:det:B7} 证明：$n\ge2$ 时 $(\vv A^*)^*=(\det\vv A)^{n-2}\vv A$.（先设 $\vv A$ 可逆；不可逆的情形另行讨论.）
			\item\label{ex:det:B8} 设 $\vv A,\vv B$ 都是 $n$ 阶方阵. 证明 $\begin{vmatrix}\vv A&\vv B\\ \vv B&\vv A\end{vmatrix}=\det(\vv A+\vv B)\det(\vv A-\vv B)$.（提示：分块初等变换.）
			\item\label{ex:det:B9} 设 $\vv A$ 可逆. 证明 $\begin{vmatrix}\vv A&\vv B\\ \vv C&\vv D\end{vmatrix}=\det\vv A\cdot\det(\vv D-\vv{CA}^{-1}\vv B)$；若还有 $\vv{AC}=\vv{CA}$，则它等于 $\det(\vv{AD}-\vv{CB})$.
			\item\label{ex:det:B10} 用 Fibonacci 矩阵 $\vv F=\begin{pmatrix}1&1\\1&0\end{pmatrix}$ 满足 $\vv F^n=\begin{pmatrix}F_{n+1}&F_n\\F_n&F_{n-1}\end{pmatrix}$（归纳法证明），推出 Cassini 恒等式 $F_{n+1}F_{n-1}-F_n^2=(-1)^n$.
			\item\label{ex:det:B11} 计算 $n$ 阶行列式 $D_n=\begin{vmatrix}1+x^2&x&&\\x&1+x^2&\ddots&\\&\ddots&\ddots&x\\&&x&1+x^2\end{vmatrix}$.
			\item\label{ex:det:B12} 设 $\omega=\mathrm e^{2\pi\mathrm i/3}$. 证明
			\[\begin{vmatrix}a&b&c\\c&a&b\\b&c&a\end{vmatrix}=(a+b+c)(a+\omega b+\omega^2c)(a+\omega^2b+\omega c)=a^3+b^3+c^3-3abc.\]
			\item\label{ex:det:B13} 证明：元素为 $a_{ij}=\min(i,j)$ 的 $n$ 阶行列式等于 $1$.
			\item\label{ex:det:B14} 证明：平面上四点 $(x_i,y_i)$（$i=1,2,3,4$）共圆（或共线），当且仅当 $\begin{vmatrix}x_i^2+y_i^2&x_i&y_i&1\end{vmatrix}_{4\times4}=0$（每行对应一个点）.
						\item\label{ex:det:B15} 设 $\vv A$ 为 $n$ 阶方阵（$n\ge2$）. 证明
			\[\rk\vv A^*=\begin{cases}n,&\rk\vv A=n,\\1,&\rk\vv A=n-1,\\0,&\rk\vv A\le n-2.\end{cases}\]
			（提示：$\vv A^*$ 的元素都是 $\vv A$ 的 $n-1$ 阶子式，用定理~\ref{thm:minorrank}；$\rk\vv A=n-1$ 时考察 $\vv A\vv A^*=(\det\vv A)\vv I$.）
			\item\label{ex:det:B16} 设 $f(x)=\begin{vmatrix}x&1&2&3\\1&x&2&3\\1&2&x&3\\1&2&3&x\end{vmatrix}$. 不展开，求 $f(x)=0$ 的全部根.
		\end{exB}
		
		\Cgroup
		\begin{exC}
			\item\label{ex:det:C1} （行列式的唯一性）设 $D:M_n(\mathbb{F})\to\mathbb{F}$ 对每一行都是线性的、交换两行变号，且 $D(\vv I)=1$. 证明 $D=\det$.（提示：先说明 $D$ 在初等矩阵上的取值，再对任意矩阵用初等变换.）
			\item\label{ex:det:C2} （Cauchy 行列式）证明
			\[\det\Big(\frac1{a_i+b_j}\Big)_{n\times n}=\frac{\prod_{i<j}(a_j-a_i)(b_j-b_i)}{\prod_{i,j}(a_i+b_j)}.\]
			先对 $n=2$ 直接验证，再用``各行减去最后一行''归纳. 取 $a_i=i$，$b_j=j-1$，得到 Hilbert 矩阵 $(1/(i+j-1))$ 的行列式，它为什么在数值计算中臭名昭著？
			\item\label{ex:det:C3} （排列的符号是同态）对排列 $\sigma$，令 $\vv P_\sigma$ 为把 $\vv e_j$ 映到 $\vv e_{\sigma(j)}$ 的置换矩阵. 证明 $\vv P_\sigma\vv P_\tau=\vv P_{\sigma\tau}$，$\det\vv P_\sigma=(-1)^{\tau(\sigma)}$，并由行列式的乘积公式推出 $\operatorname{sgn}(\sigma\tau)=\operatorname{sgn}\sigma\cdot\operatorname{sgn}\tau$.
			\item\label{ex:det:C4} （矩阵—树定理）对一个连通无向图，令 $\vv D$ 为各顶点度数构成的对角矩阵，$\vv A$ 为邻接矩阵，$\vv L=\vv D-\vv A$ 称为 Laplace 矩阵. 定理说：图的生成树个数等于 $\vv L$ 的任一个 $(n-1)$ 阶主子式. 对 $4$ 个顶点的完全图验证它有 $16$ 棵生成树，对一个 $4$ 个顶点的环验证它有 $4$ 棵生成树.
			\item\label{ex:det:C5} 证明：若 $n$ 阶矩阵的元素全为 $\pm1$，则它的行列式能被 $2^{n-1}$ 整除. 求 $3$ 阶此类行列式的最大值.
			\item\label{ex:det:C6} （Jacobi 公式的雏形）证明 $\det(\vv I+t\vv B)=1+t\tr\vv B+(\text{$t^2$ 及更高次项})$. 由此说明：若 $\vv A(t)$ 的元素都是 $t$ 的可导函数，且 $\vv A(0)=\vv I$，则 $\frac{\mathrm d}{\mathrm dt}\det\vv A(t)\big|_{t=0}=\tr\vv A'(0)$.
			\item\label{ex:det:C7} （Sylvester 结式）设 $f=a_0x^2+a_1x+a_2$，$g=b_0x^2+b_1x+b_2$（$a_0b_0\ne0$）. 证明 $f,g$ 有公共根当且仅当
			\[\begin{vmatrix}a_0&a_1&a_2&0\\0&a_0&a_1&a_2\\b_0&b_1&b_2&0\\0&b_0&b_1&b_2\end{vmatrix}=0.\]
			（提示：有公共根当且仅当存在次数 $\le1$ 的不全为零的 $u,v$ 使 $uf+vg=0$，与第七章的 Bézout 等式对照.）
			\item\label{ex:det:C8} （Dodgson 凝聚法）对 $3$ 阶矩阵 $\vv A$，记 $M_{ij}$ 为余子式，$c=a_{22}$. 验证 $\det\vv A\cdot c=M_{11}M_{33}-M_{13}M_{31}$. 它允许只计算 $2$ 阶行列式就得到高阶行列式. 《爱丽丝漫游奇境》的作者刘易斯·卡罗尔（即 Dodgson）于 1866 年发表了这一方法.
			\item\label{ex:det:C9} 设 $\vv A$ 是 $n$ 阶实反对称矩阵. 证明 $\det(\vv I+\vv A)\ge1$.（提示：$\det(\vv I+\vv A)=\det(\vv I-\vv A)$，考察 $\det(\vv I+t\vv A)$ 作为 $t$ 的多项式只含偶次项，且 $\vv I+t\vv A$ 对一切实数 $t$ 都可逆.）
		\end{exC}
	}
	
	% =====================================================================
	% 第六章　线性方程组解的结构
	% =====================================================================
	\newcommand{\ExSolutions}{%
		\exsetheading
		\Agroup
		\begin{exA}
			\item\label{ex:sol:A1} 求下列齐次方程组的基础解系与通解：
			\[\text{(1)}\ \left\{\begin{array}{l}x_1+x_2-x_3-x_4=0,\\2x_1-5x_2+3x_3+2x_4=0,\\7x_1-7x_2+3x_3+x_4=0;\end{array}\right.\qquad
			\text{(2)}\ \left\{\begin{array}{l}x_1+2x_2+x_3-x_4=0,\\3x_1+6x_2-x_3-3x_4=0,\\5x_1+10x_2+x_3-5x_4=0.\end{array}\right.\]
			\item\label{ex:sol:A2} 求下列方程组的通解，写成``特解 $+$ 基础解系的线性组合''：
			\[\text{(1)}\ \left\{\begin{array}{l}x_1+x_2-3x_3-x_4=1,\\3x_1-x_2-3x_3+4x_4=4,\\x_1+5x_2-9x_3-8x_4=0;\end{array}\right.\qquad
			\text{(2)}\ \left\{\begin{array}{l}2x_1+x_2-x_3+x_4=1,\\4x_1+2x_2-2x_3+x_4=2,\\2x_1+x_2-x_3-x_4=1.\end{array}\right.\]
			\item\label{ex:sol:A3} 设 $\vv A$ 是 $4\times5$ 矩阵，$\rk\vv A=3$. (1) $\vv{Ax}=\zero$ 的基础解系含几个向量？(2) $\vv{Ax}=\vv b$ 是否一定有解？(3) 若 $\rk\vv A=4$ 呢？
			\item\label{ex:sol:A4} 不解方程，用秩判断下列方程组是否有解：(1) $x_1+x_2=1$，$x_1-x_2=1$，$2x_1+3x_2=2$；(2) $x_1+x_2+x_3=1$，$2x_1+2x_2+2x_3=3$.
			\item\label{ex:sol:A5} 设 $\vv\eta_1,\vv\eta_2$ 是 $\vv{Ax}=\vv b$（$\vv b\ne\zero$）的两个解. 下列向量中哪些是 $\vv{Ax}=\vv b$ 的解，哪些是 $\vv{Ax}=\zero$ 的解？
			$\vv\eta_1+\vv\eta_2$，$2\vv\eta_1-\vv\eta_2$，$\frac12(\vv\eta_1+\vv\eta_2)$，$\vv\eta_1-\vv\eta_2$，$3\vv\eta_1-2\vv\eta_2$.
			\item\label{ex:sol:A6} 设 $3$ 元方程组 $\vv{Ax}=\vv b$ 的系数矩阵秩为 $2$，$\vv\eta_1=(1,2,3)^\trans$ 与 $\vv\eta_2=(1,-2,-1)^\trans$ 是它的两个解. 求通解.
			\item\label{ex:sol:A7} 写出一个以 $\vv x=(1,0,2)^\trans+k(1,1,-1)^\trans$（$k$ 任意）为全部解的线性方程组.
			\item\label{ex:sol:A8} 判断三个平面 $x+y+z=1$，$x+2y+3z=2$，$2x+3y+4z=4$ 的位置关系；若把最后一个常数改为 $3$ 呢？
			\item\label{ex:sol:A9} 举例说明：$\vv{Ax}=\zero$ 只有零解时，$\vv{Ax}=\vv b$ 可能无解；但若有解，解必唯一.
			\item\label{ex:sol:A10} 解矩阵方程 $\begin{pmatrix}1&2&3\\2&2&1\\3&4&3\end{pmatrix}\vv X=\begin{pmatrix}2&5\\3&1\\4&3\end{pmatrix}$.
			\item\label{ex:sol:A11} 设 $\vv A=\begin{pmatrix}1&1\\2&2\end{pmatrix}$. 求全部满足 $\vv{AX}=\vv O$ 的二阶矩阵 $\vv X$，它们构成的空间是几维的？
			\item\label{ex:sol:A12} 设 $\vv A=\begin{pmatrix}1&2&-2\\4&t&3\\3&-1&1\end{pmatrix}$，且存在三阶非零矩阵 $\vv B$ 使 $\vv{AB}=\vv O$. 求 $t$，并求 $\rk\vv B$ 的可能值.
			\item\label{ex:sol:A13} 求方程 $x_1+x_2+x_3=1$ 的通解，并写出它的三个解，使得通解中的每个解都是这三个解的仿射组合.
			\item\label{ex:sol:A14} 把空间直线 $\{x+y+z=1,\ x-y=0\}$ 写成参数式；反过来，把直线 $\vv x=(1,2,0)^\trans+t(1,0,1)^\trans$ 写成两个平面的交线.
			\item\label{ex:sol:A15} 判断空间直线 $\vv x=(1,0,0)^\trans+s(1,1,0)^\trans$ 与 $\vv x=(0,1,1)^\trans+t(0,1,-1)^\trans$ 是否相交.（列出关于 $s,t$ 的方程组.）
			\item\label{ex:sol:A16} 物流网络：三个仓库 $A,B,C$ 之间有线路 $A\to B$（$x_1$），$B\to C$（$x_2$），$A\to C$（$x_3$）. $A$ 每天净发出 $30$ 吨，$C$ 每天净收到 $30$ 吨，$B$ 不囤货. 求全部方案. 哪些量（如 $x_1+x_3$，$x_1-x_2$，$x_3$）能被方程组唯一确定？用可确定性判据说明.
			\item\label{ex:sol:A17} 求 $\vv A=\begin{pmatrix}2&1&0\\4&3&1\\2&3&3\end{pmatrix}$ 的 LU 分解.
			\item\label{ex:sol:A18} 设 $\vv A$ 是 $m\times n$ 矩阵，$\rk\vv A=m<n$. 证明 $\vv{Ax}=\vv b$ 对每个 $\vv b$ 都有无穷多解.
			\item\label{ex:sol:A19} 求经过点 $(1,1,1)$，$(2,0,1)$，$(0,2,3)$ 的平面方程（用齐次方程组求法向量）.
			\item\label{ex:sol:A20} 设 $\vv A=\begin{pmatrix}1&1\\1&1.0001\end{pmatrix}$. 分别求 $\vv{Ax}=(2,2)^\trans$ 与 $\vv{Ax}=(2,2.0001)^\trans$ 的解，计算右端的相对变化与解的相对变化之比.
		\end{exA}
		
		\Bgroup
		\begin{exB}
			\item\label{ex:sol:B1} 求两个齐次方程组 $\{x_1+x_2=0,\ x_2-x_4=0\}$ 与 $\{x_1-x_2+x_3=0,\ x_2-x_3+x_4=0\}$ 的公共非零解.
			\item\label{ex:sol:B2} 设 $\vv A$ 是实矩阵. 证明 $\vv{Ax}=\zero$ 与 $\vv A^\trans\vv{Ax}=\zero$ 同解，从而 $\rk(\vv A^\trans\vv A)=\rk\vv A$.（提示：$\vv x^\trans\vv A^\trans\vv{Ax}=\|\vv{Ax}\|^2$ 在第九章才有名字，但这里只需``实数的平方和为零则每个都为零''.）
			\item\label{ex:sol:B3} 证明 $\rk(\vv A+\vv B)\le\rk(\vv A\mid\vv B)\le\rk\vv A+\rk\vv B$.
			\item\label{ex:sol:B4} 设 $\vv A$ 为 $n$ 阶方阵，$\vv A^2=\vv I$. 证明 $\rk(\vv A+\vv I)+\rk(\vv A-\vv I)=n$.
			\item\label{ex:sol:B5} 证明 $\rk(\vv{AB})=\rk\vv B$ 当且仅当 $\ker\vv A\cap\imm\vv B=\{\zero\}$.
			\item\label{ex:sol:B6} 设 $\vv\eta^*$ 是 $\vv{Ax}=\vv b$（$\vv b\ne\zero$）的一个解，$\vv\xi_1,\ldots,\vv\xi_k$ 是对应齐次方程组的基础解系. 证明 $\vv\eta^*,\vv\eta^*+\vv\xi_1,\ldots,\vv\eta^*+\vv\xi_k$ 线性无关.
			\item\label{ex:sol:B7} 一张 $3\times3$ 的统计表只公开了三个行和与三个列和. 至少还要再公开几个单元格，才能确定全表？给出一种``公开哪几个格''的选法，并说明为什么某些同样个数的选法却不行.
			\item\label{ex:sol:B8} 对 $\vv A=\begin{pmatrix}1&2&1\\2&4&3\\3&6&4\end{pmatrix}$，求 $\vv b=(b_1,b_2,b_3)^\trans$ 应满足的条件，使 $\vv{Ax}=\vv b$ 有解.（提示：找出把各方程组合成 $0=\cdots$ 的系数.）
			\item\label{ex:sol:B9} 设 $\vv A$ 为 $n$ 阶方阵，$\rk\vv A=n-1$，且 $\vv A$ 的每行元素之和为零. 求 $\vv{Ax}=\zero$ 的通解.
			\item\label{ex:sol:B10} 设 $\vv A,\vv B$ 为 $n$ 阶方阵，$\vv{AB}=\vv O$，$\vv A+\vv B$ 可逆. 证明 $\rk\vv A+\rk\vv B=n$.
			\item\label{ex:sol:B11} 证明：两个 $m\times n$ 齐次方程组 $\vv{Ax}=\zero$ 与 $\vv{Bx}=\zero$ 同解，当且仅当 $\vv A$ 与 $\vv B$ 的行向量组等价.
			\item\label{ex:sol:B12} 设 $\vv A$ 为 $m\times n$ 矩阵. 证明：$\vv{Ax}=\vv b$ 对一切 $\vv b$ 有解，当且仅当存在 $n\times m$ 矩阵 $\vv C$ 使 $\vv{AC}=\vv I_m$；这时 $\vv C\vv b$ 就是一个解.
			\item\label{ex:sol:B13} 设 $\vv A$ 为 $n$ 阶方阵（$n\ge2$），$\rk\vv A=n-1$，$\vv\xi$ 是 $\vv{Ax}=\zero$ 的非零解. 证明 $\vv A^*$ 的每一列都是 $\vv\xi$ 的倍数.
			\item\label{ex:sol:B14} 用可确定性判据回答：方程组 $x_1+x_2+x_3=6$，$x_2+x_3+x_4=9$，$x_3+x_4+x_5=12$（例如三段相邻时间窗口内的观测总量）能确定 $x_1-x_4$ 吗？$x_1+x_2+x_3+x_4+x_5$ 呢？$x_1-x_4+x_5$ 呢？
			\item\label{ex:sol:B15} 设 $\vv A$ 为 $m\times n$ 矩阵，$\vv B$ 为 $n\times p$ 矩阵. 证明 $\dim\ker(\vv{AB})\le\dim\ker\vv A+\dim\ker\vv B$，并由此重新推出 Sylvester 不等式.
		\end{exB}
		
		\Cgroup
		\begin{exC}
			\item\label{ex:sol:C1} （Farkas 引理的二维形式）设 $\vv a_1,\vv a_2\in\mathbb{R}^2$ 线性无关，$\vv b\in\mathbb{R}^2$. 证明：下面两件事恰有一件成立：(1) 存在 $x_1,x_2\ge0$ 使 $x_1\vv a_1+x_2\vv a_2=\vv b$；(2) 存在 $\vv y$ 使 $\vv y^\trans\vv a_1\ge0$，$\vv y^\trans\vv a_2\ge0$，$\vv y^\trans\vv b<0$. 画图解释. 这是``无套利定价''与线性规划对偶理论的出发点.
			\item\label{ex:sol:C2} 设 $\vv A$ 是 $n$ 阶整数矩阵. 证明：若对任意整数向量 $\vv b$，$\vv{Ax}=\vv b$ 都有整数解，则 $\det\vv A=\pm1$.
			\item\label{ex:sol:C3} 证明：方阵 $\vv A$ 不用行交换就能做 LU 分解（$\vv L$ 单位下三角，$\vv U$ 上三角且可逆），当且仅当 $\vv A$ 的各阶顺序主子式都不为零；这时分解唯一.
			\item\label{ex:sol:C4} 设 $\vv A$ 是三对角矩阵（只有主对角线和相邻两条斜线上可能有非零元）. 证明：若它有 LU 分解，则 $\vv L,\vv U$ 都只有两条非零斜线；因此解三对角方程组只需与 $n$ 成正比的运算量. 这种方法称为``追赶法''.
			\item\label{ex:sol:C5} 用分块初等变换
			\[\begin{pmatrix}\vv I_n&\vv O\\ -\vv A&\vv I_m\end{pmatrix}\begin{pmatrix}\vv I_n&\vv B\\ \vv A&\vv O\end{pmatrix}\begin{pmatrix}\vv I_n&-\vv B\\ \vv O&\vv I_p\end{pmatrix}=\begin{pmatrix}\vv I_n&\vv O\\ \vv O&-\vv{AB}\end{pmatrix}\]
			给出 Sylvester 不等式 $\rk(\vv{AB})\ge\rk\vv A+\rk\vv B-n$ 的另一个证明.（提示：用第三章章末习题~\ref{ex:mat:B5} 的分块秩不等式.）
			\item\label{ex:sol:C6} 证明：$\mathbb{F}^n$ 中任意 $k$ 个线性流形（不必过原点）的交，要么是空集，要么仍是线性流形. 给出两个二维线性流形在 $\mathbb{R}^3$ 中交为空集的例子，并用秩的语言刻画这种情形.
			\item\label{ex:sol:C7} （列昂惕夫价格模型）设 $\vv A$ 为直接消耗系数矩阵，$\vv v$ 为各部门单位产品的增加值（工资、利润等）. 价格 $\vv p$ 满足``单位价格 $=$ 单位投入成本 $+$ 单位增加值''：$\vv p=\vv A^\trans\vv p+\vv v$. 说明它是第一章产量模型 $\vv x=\vv{Ax}+\vv d$ 的``转置''；证明在两个模型都可解时有 $\vv p^\trans\vv d=\vv v^\trans\vv x$，并解释其经济含义（国民收入的两种核算方法：支出法与收入法）.
			\item\label{ex:sol:C8} 设 $\vv A$ 为 $n$ 阶方阵. 证明存在正整数 $k$ 使 $\rk\vv A^k=\rk\vv A^{k+1}=\rk\vv A^{k+2}=\cdots$，且这样的最小的 $k$ 不超过 $n$.
		\end{exC}
	}
	
	% =====================================================================
	% 第七章　多项式
	% =====================================================================
	\newcommand{\ExPolynomial}{%
		\exsetheading
		\Agroup
		\begin{exA}
			\item\label{ex:poly:A1} 求商与余式：(1) $x^5-1$ 除以 $x^2+x+1$；(2) $2x^4-3x^3+4x^2-5x+6$ 除以 $x-2$（用综合除法）.
			\item\label{ex:poly:A2} 求 $x^{100}$ 分别除以 $x^2-1$，$x^2+1$，$(x-1)^2$ 所得的余式.
			\item\label{ex:poly:A3} 求 $a,b$，使 $x^2+x+1$ 整除 $x^4+ax^2+b$.（提示：设 $\omega$ 是 $x^2+x+1$ 的根，则 $\omega^3=1$.）
			\item\label{ex:poly:A4} 用辗转相除法求最大公因式：(1) $x^4+3x^3-x^2-4x-3$ 与 $3x^3+10x^2+2x-3$；(2) $x^3-2x^2+x$ 与 $x^2-1$.
			\item\label{ex:poly:A5} 证明：对任意多项式 $h$，$\gcd(f,g)=\gcd(f,\ g-hf)$. 这正是辗转相除法每一步的依据.
			\item\label{ex:poly:A6} 求 $u,v$ 使 $u(x)(x^2+1)+v(x)(x+1)=1$.
			\item\label{ex:poly:A7} 求 $2x^3-5x^2+x+2$ 的全部有理根，并在 $\mathbb{Q}$ 上把它分解为一次因式之积.
			\item\label{ex:poly:A8} 分别在 $\mathbb{Q}$，$\mathbb{R}$，$\mathbb{C}$ 上分解：(1) $x^4+x^2+1$；(2) $x^6-1$；(3) $x^3-2$.
			\item\label{ex:poly:A9} 判断 $1$ 是 $x^5-5x+4$ 的几重根.
			\item\label{ex:poly:A10} $a$ 取何值时，$x^3-3x+a$ 有重根？求出重根.
			\item\label{ex:poly:A11} 设 $\alpha_1,\alpha_2,\alpha_3$ 是 $x^3-2x^2+3x-4$ 的根. 不解方程，求 $\alpha_1+\alpha_2+\alpha_3$，$\alpha_1^2+\alpha_2^2+\alpha_3^2$，$\frac1{\alpha_1}+\frac1{\alpha_2}+\frac1{\alpha_3}$.
			\item\label{ex:poly:A12} 求次数最低的首一有理系数多项式，使 $1+\sqrt2$ 与 $\mathrm i$ 都是它的根.
			\item\label{ex:poly:A13} 求一个三次实系数首一多项式，以 $2$ 和 $1+\mathrm i$ 为根.
			\item\label{ex:poly:A14} 设 $f(x)=x^2-2x+1$，$\vv A=\begin{pmatrix}1&2\\0&1\end{pmatrix}$，计算 $f(\vv A)$；再设 $g(x)=x^3-x$，计算 $g(\vv A)$.
			\item\label{ex:poly:A15} 求过点 $(0,1)$，$(1,0)$，$(2,3)$，$(3,10)$ 的三次以下的插值多项式（用 Lagrange 公式）.
			\item\label{ex:poly:A16} 证明 $(x-1)\mid x^n-1$，并写出商式. 一般地，对任意 $a$，$(x-a)\mid x^n-a^n$.
			\item\label{ex:poly:A17} 把 $f(x)=x^4-2x^3+3$ 按 $x-1$ 的幂展开，即写成 $\sum c_k(x-1)^k$.（提示：反复用 $x-1$ 作综合除法，或用 $c_k=f^{(k)}(1)/k!$.）
			\item\label{ex:poly:A18} 多项式 $f$ 除以 $x-1$ 余 $3$，除以 $x-2$ 余 $5$. 求 $f$ 除以 $(x-1)(x-2)$ 的余式.
			\item\label{ex:poly:A19} 一个项目现在投入 $100$ 万元，一年后回收 $55$ 万元，两年后回收 $60.5$ 万元. 写出内部收益率满足的方程并精确求解.
			\item\label{ex:poly:A20} 证明：若 $f\mid g$ 且 $g\mid f$，则 $f=cg$（$c$ 为非零常数）.
			\item\label{ex:poly:A21} 设 $\gcd(f,g)=d$，$f=df_1$，$g=dg_1$. 证明 $\gcd(f_1,g_1)=1$.
		\end{exA}
		
		\Bgroup
		\begin{exB}
			\item\label{ex:poly:B1} 证明：若 $\gcd(f,g)=1$ 且 $\gcd(f,h)=1$，则 $\gcd(f,gh)=1$.
			\item\label{ex:poly:B2} 证明：$\gcd(x^m-1,\ x^n-1)=x^{\gcd(m,n)}-1$.（提示：模仿整数的辗转相除，先证 $x^m-1=x^{m-n}(x^n-1)+(x^{m-n}-1)$.）
			\item\label{ex:poly:B3} 设 $f=(x-1)^3(x+2)^2$. 计算 $f/\gcd(f,f')$，说明它恰好是``去掉重复因式''的结果. 一般地证明：若 $f=c\,p_1^{k_1}\cdots p_m^{k_m}$ 是标准分解，则 $f/\gcd(f,f')=c\,p_1\cdots p_m$.
			\item\label{ex:poly:B4} 证明：若 $p$ 不可约且 $p\mid f^2$，则 $p\mid f$. 举例说明``不可约''不能去掉.
			\item\label{ex:poly:B5} 设 $\ell_1,\ldots,\ell_n$ 是关于互异节点 $x_1,\ldots,x_n$ 的 Lagrange 基函数. 证明 $\sum_i\ell_i(x)=1$，$\sum_ix_i\ell_i(x)=x$，一般地 $\sum_ix_i^k\ell_i(x)=x^k$（$0\le k\le n-1$）.
			\item\label{ex:poly:B6} 证明：若复系数多项式 $f$ 满足 $f(x+1)=f(x)$ 对一切 $x$ 成立，则 $f$ 是常数.
			\item\label{ex:poly:B7} 证明 $x^3+x+1$ 在 $\mathbb{Q}$ 上不可约. 为什么对三次多项式只需检验有无有理根？四次多项式可以这样检验吗？
			\item\label{ex:poly:B8} 求 $x^n$ 除以 $x^2-3x+2$ 的余式.
			\item\label{ex:poly:B9} （Hermite 插值）求次数最低的多项式 $f$，使 $f(1)=1$，$f'(1)=2$，$f(2)=4$.（提示：设 $f=1+2(x-1)+c(x-1)^2$.）
			\item\label{ex:poly:B10} 设 $f$ 是有理系数多项式，$a,b\in\mathbb{Q}$，$\sqrt2\notin\mathbb{Q}$. 证明：若 $a+b\sqrt2$ 是 $f$ 的根，则 $a-b\sqrt2$ 也是 $f$ 的根.（提示：用 $(x-a)^2-2b^2$ 除 $f$.）
			\item\label{ex:poly:B11} 设 $f$ 是整系数多项式，$f(0)$ 与 $f(1)$ 都是奇数. 证明 $f$ 没有整数根.
			\item\label{ex:poly:B12} 求最小多项式意义下的``零化''：分别求出次数最低的首一多项式 $g$，使 $g(\vv A)=\vv O$，其中 $\vv A=\begin{pmatrix}0&-1\\1&0\end{pmatrix}$，$\vv A=\begin{pmatrix}2&0\\0&2\end{pmatrix}$，$\vv A=\begin{pmatrix}2&1\\0&2\end{pmatrix}$.
			\item\label{ex:poly:B13} 证明：现金流 $-C_0,\ C_1,\ldots,C_n$（$C_0>0$，$C_i\ge0$ 且不全为零）的净现值 $N(v)=-C_0+C_1v+\cdots+C_nv^n$ 在 $v>0$ 时严格递增，从而 $N(v)=0$ 至多一个正根；当 $C_1+\cdots+C_n>C_0$ 时这个正根落在 $(0,1)$ 中，即 IRR 唯一且为正.
			\item\label{ex:poly:B14} 设 $f$ 是 $n$ 次多项式，在 $n+1$ 个整数点 $0,1,\ldots,n$ 上都取整数值. 证明 $f$ 在一切整数点上都取整数值.（提示：用基 $\binom x0,\binom x1,\ldots,\binom xn$，其中 $\binom xk=\frac{x(x-1)\cdots(x-k+1)}{k!}$.）
			\item\label{ex:poly:B15} 证明 Lagrange 插值的``唯一性''也可以由线性代数推出：求值映射 $\mathbb{F}[x]_n\to\mathbb{F}^n$，$f\mapsto(f(x_1),\ldots,f(x_n))^\trans$ 是单射，从而是同构.
			\item\label{ex:poly:B16} 设 $p$ 是 $\mathbb{F}$ 上的不可约多项式，$f\in\mathbb{F}[x]$. 证明：$\gcd(p,f)=1$ 或 $p\mid f$，二者必居其一.
		\end{exB}
		
		\Cgroup
		\begin{exC}
			\item\label{ex:poly:C1} （Eisenstein 判别法）设 $f=a_nx^n+\cdots+a_0$ 是整系数多项式，$p$ 是素数，$p\nmid a_n$，$p\mid a_i$（$i<n$），$p^2\nmid a_0$. 证明 $f$ 在 $\mathbb{Q}$ 上不可约（可以先承认 Gauss 引理：整系数多项式若能分解为两个低次有理系数多项式之积，则也能分解为两个低次整系数多项式之积）. 由此说明 $x^n-2$ 对一切 $n$ 都不可约，所以 $\mathbb{Q}[x]$ 中有任意高次的不可约多项式.
			\item\label{ex:poly:C2} 设 $p$ 是素数. 证明分圆多项式 $\Phi_p(x)=x^{p-1}+x^{p-2}+\cdots+1$ 在 $\mathbb{Q}$ 上不可约.（提示：对 $\Phi_p(x+1)$ 用 Eisenstein 判别法.）
			\item\label{ex:poly:C3} （Newton 恒等式）记 $s_k=\alpha_1^k+\alpha_2^k+\alpha_3^k$，$\sigma_1,\sigma_2,\sigma_3$ 为初等对称多项式. 证明 $s_1=\sigma_1$，$s_2=\sigma_1s_1-2\sigma_2$，$s_3=\sigma_1s_2-\sigma_2s_1+3\sigma_3$，并求 $x^3-2x^2+3x-4$ 的根的立方和.
			\item\label{ex:poly:C4} （三次方程的判别式）设 $x^3+px+q$ 的三个复根为 $\alpha_1,\alpha_2,\alpha_3$，$\Delta=\prod_{i<j}(\alpha_i-\alpha_j)^2$. 证明 $\Delta=-4p^3-27q^2$，并说明（$p,q$ 为实数时）：$\Delta>0$ 时三个不同实根，$\Delta<0$ 时一实二虚，$\Delta=0$ 时有重根.
			\item\label{ex:poly:C5} （Chebyshev 多项式）证明存在 $n$ 次多项式 $T_n$ 使 $\cos n\theta=T_n(\cos\theta)$，且 $T_{n+1}=2xT_n-T_{n-1}$. 求 $T_n$ 的全部根. 由此证明：$\cos20^\circ$ 是 $8x^3-6x-1$ 的根，并验证这个三次多项式在 $\mathbb{Q}$ 上不可约. （结合``尺规作图只能得到经若干次开平方得到的数''，这说明 $60^\circ$ 角不能用尺规三等分.）
			\item\label{ex:poly:C6} 证明：在 $\mathbb{R}$ 上处处非负的实系数多项式，都能写成两个实系数多项式的平方和.（提示：利用实根的重数必为偶数，以及 $|a+b\mathrm i|^2=a^2+b^2$ 与 $(a^2+b^2)(c^2+d^2)=(ac-bd)^2+(ad+bc)^2$.）
			\item\label{ex:poly:C7} 设 $f\in\mathbb{Q}[x]$ 不可约，$g\in\mathbb{Q}[x]$ 与 $f$ 有一个公共复根. 证明 $f\mid g$. 由此说明：一个复数若是某个有理系数多项式的根，则以它为根的次数最低的首一有理系数多项式（称为它的\textbf{极小多项式}）唯一，且不可约.
			\item\label{ex:poly:C8} （$\mathbb{F}_2$ 上的多项式）在系数按模 $2$ 运算的多项式中，证明 $x^2+x+1$ 是唯一的二次不可约多项式，并找出全部三次不可约多项式. 这类多项式是循环冗余校验（CRC）与密码学的基础.
			\item\label{ex:poly:C9} 设 $f$ 是实系数多项式，它的系数序列的``变号次数''为 $V$（忽略零系数）. Descartes 符号法则说：$f$ 的正根个数（按重数计）不超过 $V$，且与 $V$ 的差是偶数. 用它重新解释第七章``内部收益率''建模实验中两个项目的结论.
		\end{exC}
	}
	
	% =====================================================================
	% 第八章　特征值与特征向量
	% =====================================================================
	\newcommand{\ExEigen}{%
		\exsetheading
		\Agroup
		\begin{exA}
			\item\label{ex:eig:A1} 求下列矩阵的特征值与特征向量（在 $\mathbb{C}$ 上考虑）：
			\[\begin{pmatrix}3&1\\0&2\end{pmatrix},\quad\begin{pmatrix}1&2\\2&1\end{pmatrix},\quad\begin{pmatrix}0&1\\-1&0\end{pmatrix},\quad\begin{pmatrix}4&-2\\1&1\end{pmatrix},\]\[\begin{pmatrix}2&0&0\\0&1&1\\0&1&1\end{pmatrix},\quad\begin{pmatrix}1&2&2\\2&1&2\\2&2&1\end{pmatrix}.\]
			\item\label{ex:eig:A2} 判断 A1 中各矩阵能否在 $\mathbb{R}$ 上对角化；能的，求可逆矩阵 $\vv P$ 使 $\vv P^{-1}\vv{AP}$ 为对角矩阵.
			\item\label{ex:eig:A3} 判断下列矩阵能否对角化：$\begin{pmatrix}1&1\\0&2\end{pmatrix}$，$\begin{pmatrix}2&1\\0&2\end{pmatrix}$，$\begin{pmatrix}2&0&0\\0&2&0\\1&0&2\end{pmatrix}$，$\begin{pmatrix}1&1&0\\0&2&0\\0&0&1\end{pmatrix}$.
			\item\label{ex:eig:A4} 设 $3$ 阶方阵 $\vv A$ 的特征值为 $1,2,3$. 求 $\det\vv A$，$\tr\vv A$，$\vv A^{-1}$ 的特征值，$\vv A^2-2\vv A+\vv I$ 的特征值，$\det(\vv A^*)$，$\det(\vv A^2+\vv I)$.
			\item\label{ex:eig:A5} 已知二阶方阵的迹为 $5$，行列式为 $6$，求它的特征值. 若迹为 $2$、行列式为 $5$ 呢？
			\item\label{ex:eig:A6} 证明：$\vv A$ 与 $\vv A^\trans$ 有相同的特征多项式. 它们的特征向量是否相同？举例说明.
			\item\label{ex:eig:A7} 证明：三角矩阵的特征值恰为其对角元.
			\item\label{ex:eig:A8} 证明：若 $\vv A^k=\vv O$，则 $\vv A$ 的特征值全为 $0$；若 $\vv A^2=\vv A$，则特征值只能是 $0$ 或 $1$；若 $\vv A^2=\vv I$，则特征值只能是 $\pm1$.
			\item\label{ex:eig:A9} 设 $\vv A=\begin{pmatrix}4&-2\\1&1\end{pmatrix}$. 利用对角化计算 $\vv A^n$.
			\item\label{ex:eig:A10} 设品牌转移矩阵为 $\vv P=\begin{pmatrix}0.9&0.2\\0.1&0.8\end{pmatrix}$. 求稳态分布，并利用对角化求 $\lim_{n\to\infty}\vv P^n$. 解释极限矩阵两列相同的含义.
			\item\label{ex:eig:A11} 已知 $\vv\xi=(1,1,-1)^\trans$ 是 $\vv A=\begin{pmatrix}2&-1&2\\5&a&3\\-1&b&-2\end{pmatrix}$ 的特征向量. 求 $a,b$ 与对应的特征值.
			\item\label{ex:eig:A12} 设 $\vv A=\begin{pmatrix}0&0&1\\1&1&x\\1&0&0\end{pmatrix}$. $x$ 取何值时 $\vv A$ 可对角化？
			\item\label{ex:eig:A13} 对二阶方阵直接验证 Cayley–Hamilton 定理：$\vv A^2-(\tr\vv A)\vv A+(\det\vv A)\vv I=\vv O$.
			\item\label{ex:eig:A14} 求下列矩阵的最小多项式：$\diag(1,1,2)$，$\begin{pmatrix}1&1&0\\0&1&0\\0&0&2\end{pmatrix}$，$\begin{pmatrix}0&1&0\\0&0&1\\0&0&0\end{pmatrix}$，$\begin{pmatrix}1&1&1\\1&1&1\\1&1&1\end{pmatrix}$.
			\item\label{ex:eig:A15} 判断 $\begin{pmatrix}1&1\\0&1\end{pmatrix}$ 与 $\begin{pmatrix}1&0\\1&1\end{pmatrix}$ 是否相似；若相似，求出 $\vv P$.
			\item\label{ex:eig:A16} 设 $\vv A$ 是上三角矩阵. 证明 $\spn(\vv e_1)$，$\spn(\vv e_1,\vv e_2)$，……都是 $\vv A$ 的不变子空间.
			\item\label{ex:eig:A17} 解差分方程 $a_{k+1}=a_k+2a_{k-1}$，$a_0=0$，$a_1=1$.
			\item\label{ex:eig:A18} （种群模型）某种群分为幼年与成年两组，第 $k$ 年的数量记为 $\vv x_k=(y_k,z_k)^\trans$. 每只成年个体每年平均产下 $1.1$ 只幼体，幼体有 $60\%$ 存活并长为成年，成年个体有一半活到下一年，即
			$y_{k+1}=1.1z_k$，$z_{k+1}=0.6y_k+0.5z_k$. 求矩阵的特征值与特征向量，说明种群的长期增长率和幼体与成年个体的稳定比例.
			\item\label{ex:eig:A19} 证明：实二阶方阵 $\vv A$ 没有实特征值，当且仅当 $(\tr\vv A)^2<4\det\vv A$.
			\item\label{ex:eig:A20} 求平面上向直线 $y=2x$ 的正投影、关于该直线的反射这两个变换的特征值与特征向量（不必写出矩阵，从几何上直接回答），再写出矩阵加以验证.
		\end{exA}
		
		\Bgroup
		\begin{exB}
			\item\label{ex:eig:B1} 设 $\vv u,\vv v\in\mathbb{R}^n$ 非零，$\vv A=\vv u\vv v^\trans$. 求 $\vv A$ 的全部特征值及其重数，并证明：$\vv A$ 可对角化当且仅当 $\vv v^\trans\vv u\ne0$.
			\item\label{ex:eig:B2} 设 $\vv A$ 有 $n$ 个互异的特征值，$\vv{AB}=\vv{BA}$. 证明 $\vv A$ 的特征向量都是 $\vv B$ 的特征向量，从而 $\vv B$ 也可对角化；进一步证明 $\vv B$ 是 $\vv A$ 的多项式.
			\item\label{ex:eig:B3} 证明：若 $\vv A$ 可对角化，则 $\vv A^k$ 与 $f(\vv A)$ 都可对角化. 逆命题成立吗？
			\item\label{ex:eig:B4} 设 $\vv A$ 可对角化，$W$ 是 $\vv A$ 的不变子空间. 证明 $\vv A|_W$ 也可对角化.（提示：最小多项式.）
			\item\label{ex:eig:B5} 证明：满足 $\vv A^2=-\vv I$ 的二阶实方阵都与 $\begin{pmatrix}0&-1\\1&0\end{pmatrix}$ 相似.
			\item\label{ex:eig:B6} （Gershgorin 圆盘定理）证明：方阵 $\vv A$ 的每个特征值都落在某个圆盘 $\{z\in\mathbb{C}:|z-a_{ii}|\le\sum_{j\ne i}|a_{ij}|\}$ 中. 由此证明：严格对角占优的矩阵（$|a_{ii}|>\sum_{j\ne i}|a_{ij}|$）可逆.
			\item\label{ex:eig:B7} 设 $\vv A$ 的每行元素之和都是 $s$. 证明 $s$ 是 $\vv A$ 的特征值. 若 $\vv A$ 可逆，求 $\vv A^{-1}$ 每行元素之和.
			\item\label{ex:eig:B8} 设 $\vv A=\begin{pmatrix}a&b&c\\c&a&b\\b&c&a\end{pmatrix}$（循环矩阵），$\omega=\mathrm e^{2\pi\mathrm i/3}$. 证明 $(1,\omega^k,\omega^{2k})^\trans$（$k=0,1,2$）都是 $\vv A$ 的特征向量，并求对应的特征值. 由此重新得到第五章章末习题~\ref{ex:det:B12} 的行列式公式.
			\item\label{ex:eig:B9} 解差分方程 $a_{k+1}=4a_k-4a_{k-1}$，$a_0=1$，$a_1=4$. 解释通项中为什么出现 $k\cdot2^k$.
			\item\label{ex:eig:B10} 证明 $\tr(\vv A^k)=\lambda_1^k+\cdots+\lambda_n^k$. 若三阶方阵满足 $\tr\vv A=0$，$\tr\vv A^2=6$，$\det\vv A=0$，求它的特征值.
			\item\label{ex:eig:B11} 证明：若 $\vv A^3=\vv A$，则 $\vv A$ 可对角化，且特征值在 $\{0,1,-1\}$ 中.
			\item\label{ex:eig:B12} $a,b,c$ 满足什么条件时，$\begin{pmatrix}2&0&0\\a&2&0\\b&c&1\end{pmatrix}$ 可对角化？
			\item\label{ex:eig:B13} 三个网页之间的链接为：网页 $1$ 链接到 $2,3$，网页 $2$ 链接到 $3$，网页 $3$ 链接到 $1$.（与判例不同，网页之间可以``互相引用''，形成环路.） 取 $d=1$ 写出转移矩阵，手算它的稳态分布.
			\item\label{ex:eig:B14} （矩阵指数）对可对角化矩阵 $\vv A=\vv P\vv D\vv P^{-1}$，定义 $\mathrm e^{t\vv A}=\vv P\diag(\mathrm e^{t\lambda_1},\ldots,\mathrm e^{t\lambda_n})\vv P^{-1}$. 对 $\vv A=\begin{pmatrix}0&1\\1&0\end{pmatrix}$ 计算 $\mathrm e^{t\vv A}$，并验证 $\vv x(t)=\mathrm e^{t\vv A}\vv x_0$ 满足微分方程组 $x_1'=x_2$，$x_2'=x_1$.
			\item\label{ex:eig:B15} 三状态马尔可夫链的转移矩阵为 $\begin{pmatrix}0.5&0.25&0.25\\0.25&0.5&0.25\\0.25&0.25&0.5\end{pmatrix}$. 求稳态分布与全部特征值，说明偏离稳态的部分每一步缩小为原来的多少.
			\item\label{ex:eig:B16} 证明：若 $\lambda$ 是 $\vv A$ 的 $k$ 重特征值（代数重数），则 $\rk(\lambda\vv I-\vv A)\ge n-k$.
			\item\label{ex:eig:B17} 对第三章``判例引证网络''中的四份判决，分别取阻尼系数 $d=0.5,\ 0.85,\ 0.95,\ 1$，按第八章建模实验的方法求 PageRank 得分（可用计算机），列表比较各判决的得分与排序. $d$ 越接近 $1$，得分越``集中''在哪些判决上？为什么？
		\end{exB}
		
		\Cgroup
		\begin{exC}
			\item\label{ex:eig:C1} 设 $\vv A$ 为 $m\times n$ 矩阵，$\vv B$ 为 $n\times m$ 矩阵. 证明 $\lambda^n\det(\lambda\vv I_m-\vv{AB})=\lambda^m\det(\lambda\vv I_n-\vv{BA})$. 由此说明 $\vv{AB}$ 与 $\vv{BA}$ 的非零特征值（连同重数）相同.（提示：计算 $\begin{pmatrix}\vv I_m&-\vv A\\ \vv O&\lambda\vv I_n\end{pmatrix}\begin{pmatrix}\lambda\vv I_m&\vv A\\ \vv B&\vv I_n\end{pmatrix}$ 等两种乘积的行列式.）
			\item\label{ex:eig:C2} （Schur 三角化）证明：每个复方阵都相似于一个上三角矩阵. 由此给出 Cayley–Hamilton 定理的另一个证明：对上三角矩阵 $\vv T$，证明 $(\vv T-t_{11}\vv I)\cdots(\vv T-t_{nn}\vv I)=\vv O$.
			\item\label{ex:eig:C3} 求下列矩阵的 Jordan 标准形：
			\[\begin{pmatrix}2&1&0\\0&2&0\\0&0&2\end{pmatrix},\quad\begin{pmatrix}1&1&1\\0&1&1\\0&0&1\end{pmatrix},\quad\begin{pmatrix}0&1&0&0\\0&0&0&0\\0&0&0&1\\0&0&0&0\end{pmatrix}.\]
			特征多项式为 $(\lambda-2)^4$ 的 $4$ 阶复矩阵，在相似的意义下共有几类？
			\item\label{ex:eig:C4} 证明：若一组可对角化的矩阵两两可交换，则它们可以\textbf{同时}对角化（存在同一个 $\vv P$ 使它们都化为对角矩阵）.
			\item\label{ex:eig:C5} 设 $\vv A$ 可对角化. 证明 $\vv A^k\to\vv O$（按元素）当且仅当 $\vv A$ 的所有特征值的模都小于 $1$. 用它说明：投入产出模型中，若直接消耗系数矩阵 $\vv A$ 可对角化且特征值的模都小于 $1$，则 $(\vv I-\vv A)^{-1}=\vv I+\vv A+\vv A^2+\cdots$.
			\item\label{ex:eig:C6} （友矩阵）证明 $\begin{pmatrix}0&0&\cdots&0&-a_0\\1&0&\cdots&0&-a_1\\0&1&\cdots&0&-a_2\\ \vdots&&\ddots&&\vdots\\0&0&\cdots&1&-a_{n-1}\end{pmatrix}$ 的特征多项式为 $\lambda^n+a_{n-1}\lambda^{n-1}+\cdots+a_0$. 由此说明：求多项式的根等价于求矩阵的特征值（数值软件正是这样求多项式的根的）.
			\item\label{ex:eig:C7} （赌徒破产）一位赌徒有 $k$ 元（$0\le k\le N$），每局以概率 $\frac12$ 赢或输 $1$ 元，输光或达到 $N$ 元时停止. 设 $p_k$ 为最终达到 $N$ 元的概率，证明 $p_k=\frac12(p_{k+1}+p_{k-1})$，并用特征根方法解出 $p_k=k/N$. 赢的概率改为 $p\ne\frac12$ 时呢？
			\item\label{ex:eig:C8} 证明：$n$ 阶三对角矩阵 $\vv T=\begin{pmatrix}2&-1&&\\-1&2&\ddots&\\&\ddots&\ddots&-1\\&&-1&2\end{pmatrix}$ 的特征值为 $2-2\cos\frac{k\pi}{n+1}$（$k=1,\ldots,n$），对应特征向量的第 $j$ 个分量为 $\sin\frac{jk\pi}{n+1}$. 对 $n=2,3$ 直接验证.
			\item\label{ex:eig:C9} （Perron–Frobenius 的二阶情形）设 $\vv A$ 是元素全为正的二阶矩阵. 证明 $\vv A$ 有两个互异实特征值，较大者 $r$ 为正且大于另一个的绝对值，属于 $r$ 的特征向量可取为各分量都为正. 由此解释：为什么 PageRank 引入阻尼系数后，``全部元素为正''能保证排序唯一.
			\item\label{ex:eig:C10} 证明：满足 $\vv A^2=-\vv I$ 的 $n$ 阶实方阵存在当且仅当 $n$ 为偶数，且这时 $\vv A$ 相似于 $\diag(\vv J,\ldots,\vv J)$，$\vv J=\begin{pmatrix}0&-1\\1&0\end{pmatrix}$.
		\end{exC}
	}
	
	% =====================================================================
	% 第九章　内积空间
	% =====================================================================
	\newcommand{\ExInnerProduct}{%
		\exsetheading
		\Agroup
		\begin{exA}
			\item\label{ex:ip:A1} 在 $\mathbb{R}^4$ 中，设 $\vv\alpha=(1,2,2,0)^\trans$，$\vv\beta=(2,0,1,2)^\trans$. 求 $\|\vv\alpha\|$，$\|\vv\beta\|$，$\langle\vv\alpha,\vv\beta\rangle$，它们的夹角，以及 $\|\vv\alpha-\vv\beta\|$.
			\item\label{ex:ip:A2} 判断下列二元函数是否为 $\mathbb{R}^2$ 上的内积：
			(1) $\langle\vv x,\vv y\rangle=x_1y_1+2x_2y_2$；(2) $\langle\vv x,\vv y\rangle=x_1y_2+x_2y_1$；(3) $\langle\vv x,\vv y\rangle=2x_1y_1+x_1y_2+x_2y_1+x_2y_2$；(4) $\langle\vv x,\vv y\rangle=x_1y_1-x_2y_2$.
			\item\label{ex:ip:A3} 用 Gram–Schmidt 方法把下列向量组正交化并单位化：
			(1) $(1,1,1)^\trans$，$(1,2,0)^\trans$，$(0,1,2)^\trans$；(2) $(1,1,0,0)^\trans$，$(1,0,1,0)^\trans$，$(-1,0,0,1)^\trans$.
			\item\label{ex:ip:A4} 求与 $(1,1,1)^\trans$ 和 $(1,0,-1)^\trans$ 都正交的单位向量.
			\item\label{ex:ip:A5} 已知 $\vv e_1=\frac13(1,2,2)^\trans$，$\vv e_2=\frac13(2,1,-2)^\trans$，$\vv e_3=\frac13(2,-2,1)^\trans$. 验证它们构成 $\mathbb{R}^3$ 的标准正交基，并只用内积求 $(3,6,0)^\trans$ 在这组基下的坐标.
			\item\label{ex:ip:A6} 求 $\vv\alpha=(1,2,3,4)^\trans$ 在 $W=\spn\big((1,1,0,0)^\trans,(0,0,1,1)^\trans\big)$ 上的正交投影，以及 $\vv\alpha$ 到 $W$ 的距离.
			\item\label{ex:ip:A7} 用最小二乘法求过数据点 $(0,1)$，$(1,2)$，$(2,2)$，$(3,4)$，$(4,6)$ 的最佳直线 $y=a+bx$，并计算残差平方和.
			\item\label{ex:ip:A8} 求 $a,b,c$ 使 $\begin{pmatrix}\frac1{\sqrt2}&a&\frac1{\sqrt6}\\0&b&-\frac2{\sqrt6}\\-\frac1{\sqrt2}&c&\frac1{\sqrt6}\end{pmatrix}$ 为正交矩阵.
			\item\label{ex:ip:A9} 证明：若 $\vv\alpha_1,\ldots,\vv\alpha_k$ 两两正交，则 $\|\vv\alpha_1+\cdots+\vv\alpha_k\|^2=\|\vv\alpha_1\|^2+\cdots+\|\vv\alpha_k\|^2$.
			\item\label{ex:ip:A10} 证明：$\big|\|\vv\alpha\|-\|\vv\beta\|\big|\le\|\vv\alpha-\vv\beta\|$.
			\item\label{ex:ip:A11} 用 Cauchy–Schwarz 不等式证明：(1) $(a_1+\cdots+a_n)^2\le n(a_1^2+\cdots+a_n^2)$；(2) 若 $a_i>0$，则 $(a_1+\cdots+a_n)\big(\frac1{a_1}+\cdots+\frac1{a_n}\big)\ge n^2$.
			\item\label{ex:ip:A12} 在 $C[0,1]$ 中取 $\langle f,g\rangle=\int_0^1f(x)g(x)\,\mathrm dx$. 求 $\|1\|$，$\|x\|$，$1$ 与 $x$ 的夹角，并把 $1,x$ 正交化.
			\item\label{ex:ip:A13} 设 $\vv A=\begin{pmatrix}1&1&0&1\\0&1&1&1\end{pmatrix}$. 求 $\ker\vv A$ 与行空间 $\imm\vv A^\trans$ 的标准正交基，并验证它们互为正交补.
			\item\label{ex:ip:A14} 求 $\vv A=\begin{pmatrix}1&1\\1&0\\0&1\end{pmatrix}$ 的 QR 分解.
			\item\label{ex:ip:A15} 设 $\vv u=\frac13(1,2,2)^\trans$. 写出 Householder 矩阵 $\vv H=\vv I-2\vv u\vv u^\trans$，验证它是对称的正交矩阵，并求 $\vv H\vv u$ 与 $\vv H(2,-1,0)^\trans$.
			\item\label{ex:ip:A16} 判断下列正交矩阵是旋转还是反射；若是旋转求旋转角，若是反射求反射轴：$\frac15\begin{pmatrix}3&-4\\4&3\end{pmatrix}$，$\frac15\begin{pmatrix}3&4\\4&-3\end{pmatrix}$.
			\item\label{ex:ip:A17} 对``过原点的直线''模型 $y=bx$，推导最小二乘解 $\hat b=\sum x_iy_i\big/\sum x_i^2$.
			\item\label{ex:ip:A18} 在 $\mathbb{R}^3$ 中，求点 $(1,2,3)$ 到直线 $\vv x=t(1,1,1)^\trans$ 的距离.
			\item\label{ex:ip:A19} 证明：两组标准正交基之间的过渡矩阵是正交矩阵.
			\item\label{ex:ip:A20} 三个法条的词频向量（关键词：违约、赔偿、合同、离婚）分别为 $(2,1,3,0)^\trans$，$(0,1,0,4)^\trans$，$(1,1,2,0)^\trans$，查询为 $(1,1,1,0)^\trans$. 分别用内积和余弦相似度给三个法条排序，结果有何不同？
			\item\label{ex:ip:A21} 设 $\vv u\in\mathbb{R}^n$ 是单位向量. 证明 $\vv P=\vv u\vv u^\trans$ 满足 $\vv P^\trans=\vv P$，$\vv P^2=\vv P$，且对任意 $\vv x$，$\vv P\vv x$ 是 $\vv x$ 在 $\spn(\vv u)$ 上的正交投影. 若 $\vv u$ 不是单位向量，投影矩阵应当怎样修改？
		\end{exA}
		
		\Bgroup
		\begin{exB}
			\item\label{ex:ip:B1} （极化恒等式）证明实内积空间中 $\langle\vv\alpha,\vv\beta\rangle=\frac14\big(\|\vv\alpha+\vv\beta\|^2-\|\vv\alpha-\vv\beta\|^2\big)$. 再证明：$\mathbb{R}^2$ 上的``范数'' $\|\vv x\|_1=|x_1|+|x_2|$ 不满足平行四边形法则，因而不能来自任何内积.
			\item\label{ex:ip:B2} （Bessel 不等式）设 $\vv e_1,\ldots,\vv e_k$ 是标准正交组. 证明 $\sum_{i=1}^k\langle\vv\alpha,\vv e_i\rangle^2\le\|\vv\alpha\|^2$，等号成立当且仅当 $\vv\alpha\in\spn(\vv e_1,\ldots,\vv e_k)$.
			\item\label{ex:ip:B3} 设实矩阵 $\vv P$ 满足 $\vv P^\trans=\vv P$，$\vv P^2=\vv P$. 证明 $\vv P$ 是向 $\imm\vv P$ 的正交投影的矩阵. 若只有 $\vv P^2=\vv P$，它是什么样的投影？举例说明.
			\item\label{ex:ip:B4} （加权最小二乘）若各观测值的可靠程度不同，给第 $i$ 个残差权重 $w_i>0$，最小化 $\sum w_i(y_i-a-bx_i)^2$. 证明正规方程变为 $\vv A^\trans\vv W\vv A\hat{\vv\beta}=\vv A^\trans\vv W\vv y$，$\vv W=\diag(w_1,\ldots,w_n)$. 它对应于哪个内积下的正交投影？
			\item\label{ex:ip:B5} 某商品的销售额 $y$ 随时间 $t$ 近似按指数增长 $y=c\,\mathrm e^{kt}$. 已知 $t=0,1,2,3$ 时 $y$ 约为 $2.0,2.7,3.7,5.0$. 两边取对数化为线性模型，用最小二乘估计 $c,k$. 这样做最小化的是哪个量？与直接最小化 $\sum(y_i-c\mathrm e^{kt_i})^2$ 有何不同？
			\item\label{ex:ip:B6} （Gram 行列式）设 $\vv\alpha_1,\ldots,\vv\alpha_k\in\mathbb{R}^n$，$\vv G=(\langle\vv\alpha_i,\vv\alpha_j\rangle)$. 证明 $\det\vv G\ge0$，等号成立当且仅当它们线性相关. 在 $k=2$ 时说明 $\sqrt{\det\vv G}$ 是两向量张成的平行四边形的面积.
			\item\label{ex:ip:B7} 证明：$\mathbb{R}^n$ 中点 $\vv x_0$ 到超平面 $\vv a^\trans\vv x=b$（$\vv a\ne\zero$）的距离为 $|\vv a^\trans\vv x_0-b|/\|\vv a\|$.
			\item\label{ex:ip:B8} 证明：实正交矩阵的复特征值的模都等于 $1$.
			\item\label{ex:ip:B9} （Fourier 逼近）在 $C[-\pi,\pi]$ 中取 $\langle f,g\rangle=\int_{-\pi}^\pi fg\,\mathrm dx$. 证明 $1,\cos x,\sin x,\cos2x,\sin2x$ 两两正交. 求 $f(x)=x$ 在它们张成的子空间上的正交投影.
			\item\label{ex:ip:B10} 设 $\vv A$ 为实矩阵. 证明：在 $\vv{Ax}=\vv y$ 的全部最小二乘解中，范数最小的那一个唯一，并且它属于行空间 $\imm\vv A^\trans$.
			\item\label{ex:ip:B11} 设 $\vv x,\vv y\in\mathbb{R}^n$，$\|\vv x\|=\|\vv y\|$，$\vv x\ne\vv y$. 令 $\vv u=(\vv x-\vv y)/\|\vv x-\vv y\|$. 证明 Householder 矩阵 $\vv I-2\vv u\vv u^\trans$ 把 $\vv x$ 变为 $\vv y$. 这是用 Householder 反射做 QR 分解的基础.
			\item\label{ex:ip:B12} 设两个变量的样本标准差分别为 $s_x,s_y$，相关系数为 $r$. 证明回归直线 $y=a+bx$ 的斜率为 $b=r\,s_y/s_x$，并且直线经过点 $(\bar x,\bar y)$.
			\item\label{ex:ip:B13} 证明：保持\textbf{夹角}的可逆线性变换（对一切非零向量，$\mathscr T\vv\alpha$ 与 $\mathscr T\vv\beta$ 的夹角等于 $\vv\alpha$ 与 $\vv\beta$ 的夹角）必是正交变换的非零倍数.
			\item\label{ex:ip:B14} 在回归中加入一个``虚拟变量''：设有两组数据，第一组 $y$ 值为 $3,5,4$，第二组为 $7,8,9$. 以 $\vv A=(\vv 1,\vv d)$ 为设计矩阵（$\vv d$ 在第二组为 $1$、第一组为 $0$），求最小二乘解，并解释两个系数的含义.
			\item\label{ex:ip:B15} 证明：在 $\mathbb{R}^2$ 中，任意两个旋转可交换，而一个旋转与一个反射一般不可交换；两个反射的乘积是旋转.
		\end{exB}
		
		\Cgroup
		\begin{exC}
			\item\label{ex:ip:C1} （Hadamard 不等式）用 QR 分解证明：对 $n$ 阶实方阵 $\vv A=(\vv a_1,\ldots,\vv a_n)$，$|\det\vv A|\le\|\vv a_1\|\cdots\|\vv a_n\|$，并说明它的几何意义（平行多面体的体积不超过棱长之积）.
			\item\label{ex:ip:C2} （Gram–Schmidt 的数值不稳定）设 $\varepsilon$ 很小，对 $\vv a_1=(1,\varepsilon,0,0)^\trans$，$\vv a_2=(1,0,\varepsilon,0)^\trans$，$\vv a_3=(1,0,0,\varepsilon)^\trans$ 作 Gram–Schmidt，在计算中把 $1+\varepsilon^2$ 舍入为 $1$. 证明所得的 $\vv e_2,\vv e_3$ 之间的夹角是 $60^\circ$ 而不是 $90^\circ$. 再用``修正的 Gram–Schmidt''（每得到一个 $\vv e_k$，立即从\textbf{剩余的}所有向量中减去它的分量）重做，比较结果.
			\item\label{ex:ip:C3} （Legendre 多项式）令 $P_n(x)=\frac1{2^nn!}\frac{\mathrm d^n}{\mathrm dx^n}(x^2-1)^n$. 证明 $P_0,P_1,P_2,\ldots$ 在 $\langle f,g\rangle=\int_{-1}^1fg\,\mathrm dx$ 下两两正交，并求 $\|P_n\|$.
			\item\label{ex:ip:C4} （Gauss–Markov 定理的线性代数部分）设 $\vv A$ 列满秩. 若 $\vv C\vv y$ 是参数 $\vv\beta$ 的``线性无偏估计''，即 $\vv{CA}=\vv I$. 令 $\vv C_0=(\vv A^\trans\vv A)^{-1}\vv A^\trans$（最小二乘）. 证明 $\vv C\vv C^\trans-\vv C_0\vv C_0^\trans$ 是半正定矩阵（第十一章）. 学过概率论后可以知道，这说明最小二乘估计的方差在一切线性无偏估计中最小.
			\item\label{ex:ip:C5} （复内积空间）在 $\mathbb{C}^n$ 上定义 $\langle\vv x,\vv y\rangle=\sum x_i\bar y_i$. 证明它满足共轭对称性、对第一个变量线性、正定性；证明 Cauchy–Schwarz 不等式仍成立. 满足 $\bar{\vv U}^\trans\vv U=\vv I$ 的复矩阵称为\textbf{酉矩阵}，证明它保持这个内积.
			\item\label{ex:ip:C6} 设 $f:\mathbb{R}^n\to\mathbb{R}^n$ 保持距离（$\|f(\vv x)-f(\vv y)\|=\|\vv x-\vv y\|$）且 $f(\zero)=\zero$. 证明 $f$ 保持内积，进而证明 $f$ 是线性的正交变换. 由此说明：保持距离的变换都可以分解为正交变换与平移的复合.
			\item\label{ex:ip:C7} （伪逆）设 $\vv A$ 列满秩，$\vv A^+=(\vv A^\trans\vv A)^{-1}\vv A^\trans$. 验证 Penrose 四条件：$\vv{AA}^+\vv A=\vv A$，$\vv A^+\vv{AA}^+=\vv A^+$，$(\vv{AA}^+)^\trans=\vv{AA}^+$，$(\vv A^+\vv A)^\trans=\vv A^+\vv A$. 说明 $\vv{AA}^+$ 是向哪个子空间的正交投影.
			\item\label{ex:ip:C8} 求空间中两条异面直线 $\vv x=\vv p+s\vv u$ 与 $\vv x=\vv q+t\vv v$ 之间的最短距离：把问题写成关于 $(s,t)$ 的最小二乘问题，并对 $\vv p=\zero$，$\vv u=(1,0,0)^\trans$，$\vv q=(0,1,1)^\trans$，$\vv v=(0,1,-1)^\trans$ 计算.
		\end{exC}
	}
	
	% =====================================================================
	% 第十章　矩阵分解
	% =====================================================================
	\newcommand{\ExDecomposition}{%
		\exsetheading
		\Agroup
		\begin{exA}
			\item\label{ex:dec:A1} 求正交矩阵 $\vv Q$ 使 $\vv Q^\trans\vv{AQ}$ 为对角矩阵：
			\[\begin{pmatrix}1&2\\2&1\end{pmatrix},\qquad\begin{pmatrix}3&-1\\-1&3\end{pmatrix},\]\[\begin{pmatrix}2&2&0\\2&2&0\\0&0&1\end{pmatrix},\qquad\begin{pmatrix}1&1&1\\1&1&1\\1&1&1\end{pmatrix}.\]
			\item\label{ex:dec:A2} 写出 $\begin{pmatrix}3&-1\\-1&3\end{pmatrix}$ 的谱分解 $\lambda_1\vv q_1\vv q_1^\trans+\lambda_2\vv q_2\vv q_2^\trans$，并验证两个投影矩阵之和为 $\vv I$、之积为 $\vv O$.
			\item\label{ex:dec:A3} 求二阶实对称矩阵 $\vv A$，使它的特征值为 $1,2$，且 $(1,1)^\trans$ 是属于 $1$ 的特征向量.
			\item\label{ex:dec:A4} 设三阶实对称矩阵 $\vv A$ 的特征值为 $1,1,-1$，属于 $-1$ 的特征向量为 $(0,1,1)^\trans$. 求 $\vv A$.
			\item\label{ex:dec:A5} 求 $\vv x^\trans\vv{Ax}$ 在单位球面 $\|\vv x\|=1$ 上的最大值与最小值，以及取到它们的 $\vv x$：$\vv A=\begin{pmatrix}2&2&0\\2&2&0\\0&0&1\end{pmatrix}$.
			\item\label{ex:dec:A6} 求 $x_1^2+4x_1x_2+x_2^2$ 在单位圆 $x_1^2+x_2^2=1$ 上的最大值与最小值.
			\item\label{ex:dec:A7} 求下列矩阵的奇异值分解：$\begin{pmatrix}2&0\\0&-3\end{pmatrix}$，$\begin{pmatrix}1&1\\1&1\end{pmatrix}$，$(3,4)$（$1\times2$ 矩阵）.
			\item\label{ex:dec:A8} 求 $\begin{pmatrix}1&1\\0&1\end{pmatrix}$ 的奇异值，并与它的特征值比较.（答案与黄金分割比有关.）
			\item\label{ex:dec:A9} 求 $\begin{pmatrix}3&1\\1&3\end{pmatrix}$ 的最佳秩一近似，并计算误差 $\|\vv A-\vv A_1\|_F$.
			\item\label{ex:dec:A10} 设 $\vv A=\begin{pmatrix}5&4\\4&5\end{pmatrix}$. 利用谱分解求 $\vv A^{10}$（写成两个矩阵之和的形式即可），并求对称正定矩阵 $\vv B$ 使 $\vv B^2=\vv A$.
			\item\label{ex:dec:A11} 证明：特征值全为 $0$ 的实对称矩阵是零矩阵. 对非对称矩阵举出反例.
			\item\label{ex:dec:A12} 证明：既对称又正交的实矩阵 $\vv A$ 满足 $\vv A^2=\vv I$，其特征值为 $\pm1$；并说明它在几何上是什么变换.
			\item\label{ex:dec:A13} 证明：$\vv A^\trans\vv A$ 与 $\vv{AA}^\trans$ 的非零特征值相同（重数也相同）.
			\item\label{ex:dec:A14} 证明：$\|\vv A\|_F^2=\tr(\vv A^\trans\vv A)=\sigma_1^2+\cdots+\sigma_r^2$.
			\item\label{ex:dec:A15} 设 $\vv A$ 是可逆实方阵. 证明 $|\det\vv A|=\sigma_1\sigma_2\cdots\sigma_n$，$\vv A^{-1}$ 的奇异值为 $1/\sigma_i$.
			\item\label{ex:dec:A16} 判断下列矩阵能否用正交矩阵对角化，说明理由：$\begin{pmatrix}1&2\\0&3\end{pmatrix}$，$\begin{pmatrix}1&2\\2&-2\end{pmatrix}$，$\begin{pmatrix}0&-1\\1&0\end{pmatrix}$（在实数范围内）.
			\item\label{ex:dec:A17} 四个已中心化的二维数据点为 $(2,1)$，$(-2,-1)$，$(1,2)$，$(-1,-2)$. 求样本协方差矩阵、第一主成分方向及其解释的方差比例.
			\item\label{ex:dec:A18} 计算 $\diag(10,1)$ 与 $\begin{pmatrix}1&1\\1&1.01\end{pmatrix}$ 的条件数（后者用近似值）. 后者的行列式与 $\diag(10,0.001)$ 相同，二者的条件数相同吗？
			\item\label{ex:dec:A19} 一张 $800\times600$ 像素的灰度图片，保留前 $k=40$ 个奇异值时需存储多少个数？压缩比是多少？$k$ 取多少时存储量反而超过原图？
			\item\label{ex:dec:A20} 证明：实对称矩阵的秩等于非零特征值的个数；奇异值的个数（非零）等于秩.
			\item\label{ex:dec:A21} 证明：实对称矩阵 $\vv A$ 满足``对一切 $\vv x\ne\zero$，$\vv x^\trans\vv A\vv x>0$''，当且仅当 $\vv A$ 的特征值全为正.（提示：用谱定理把 $\vv x^\trans\vv A\vv x$ 写成 $\sum\lambda_iy_i^2$.）
		\end{exA}
		
		\Bgroup
		\begin{exB}
			\item\label{ex:dec:B1} 设 $\vv A,\vv B$ 都是 $n$ 阶实对称矩阵且 $\vv{AB}=\vv{BA}$. 证明存在正交矩阵 $\vv Q$ 使 $\vv Q^\trans\vv{AQ}$ 与 $\vv Q^\trans\vv{BQ}$ 同时为对角矩阵.
			\item\label{ex:dec:B2} 证明：半正定的实对称矩阵 $\vv A$（特征值全 $\ge0$）有唯一的半正定平方根 $\vv B$（$\vv B^2=\vv A$，$\vv B$ 对称且特征值全 $\ge0$）.
			\item\label{ex:dec:B3} 设实对称矩阵 $\vv A$ 的特征值为 $\lambda_1\ge\lambda_2\ge\cdots\ge\lambda_n$，$\vv q_1$ 是属于 $\lambda_1$ 的单位特征向量. 证明 $\lambda_2=\max\{\vv x^\trans\vv{Ax}:\|\vv x\|=1,\ \vv x\perp\vv q_1\}$. 这说明第二主成分是``与第一主成分垂直的方向中方差最大的''.
			\item\label{ex:dec:B4} 证明：对任意 $\vv x$，$\|\vv{Ax}\|\le\sigma_1\|\vv x\|$，且等号可以取到. 因此 $\sigma_1=\max_{\vv x\ne\zero}\|\vv{Ax}\|/\|\vv x\|$，称为 $\vv A$ 的\textbf{谱范数}\index{pu3 fan4shu4@谱范数} $\|\vv A\|_2$.
			\item\label{ex:dec:B5} 证明：方阵 $\vv A$ 的任一特征值 $\lambda$ 满足 $\sigma_{\min}\le|\lambda|\le\sigma_{\max}$.
			\item\label{ex:dec:B6} 求 $\vv A=\begin{pmatrix}1&0&1\\0&1&1\end{pmatrix}$ 的奇异值分解.
			\item\label{ex:dec:B7} （伪逆）设 $\vv A=\vv U\vv\Sigma\vv V^\trans$，定义 $\vv A^+=\vv V\vv\Sigma^+\vv U^\trans$，其中 $\vv\Sigma^+$ 是把 $\vv\Sigma$ 转置并把非零元取倒数. 证明 $\vv A^+\vv b$ 是 $\vv{Ax}=\vv b$ 的范数最小的最小二乘解. 当 $\vv A$ 列满秩时，它与第九章的 $(\vv A^\trans\vv A)^{-1}\vv A^\trans$ 一致.
			\item\label{ex:dec:B8} 设 $p$ 个标准化指标两两相关系数都为 $\rho$（$0<\rho<1$），相关系数矩阵为 $(1-\rho)\vv I+\rho\vv J$（$\vv J$ 为全 $1$ 矩阵）. 求它的特征值与第一主成分，并求第一主成分解释的方差比例. 当 $\rho\to1$ 与 $\rho\to0$ 时比例分别趋于多少？
			\item\label{ex:dec:B9} 设两个指标的协方差矩阵为 $\begin{pmatrix}100&6\\6&1\end{pmatrix}$（第一个指标以``元''计，第二个以``万元''计）. 对协方差矩阵与相关系数矩阵分别作主成分分析，比较第一主成分. 由此说明：主成分分析对量纲敏感，为什么通常先标准化.
			\item\label{ex:dec:B10} 证明 $\|\vv A-\vv A_k\|_2=\sigma_{k+1}$.
			\item\label{ex:dec:B11} 证明：若 $\vv A$ 的奇异值互不相同且都不为零，则 $\vv A$ 的奇异值分解中的 $\vv u_i,\vv v_i$ 除了同时改变符号之外是唯一的.
			\item\label{ex:dec:B12} （正规矩阵）满足 $\vv A^\trans\vv A=\vv{AA}^\trans$ 的实方阵称为正规矩阵. 证明二阶实正规矩阵只有两类：对称矩阵，以及形如 $\begin{pmatrix}a&-b\\b&a\end{pmatrix}$ 的矩阵（旋转与伸缩的复合）.
			\item\label{ex:dec:B13} 设 $\vv A$ 是实对称矩阵. 证明 $\vv A$ 的特征值都落在 $[\min_i(a_{ii}-r_i),\ \max_i(a_{ii}+r_i)]$ 中，$r_i=\sum_{j\ne i}|a_{ij}|$. 并证明 $\lambda_{\max}\ge\max_ia_{ii}$.
			\item\label{ex:dec:B14} 对第十章``潜在语义分析''的例子，取 $k=1$，说明此时四份文书分别被表示成什么，查询``解约''与哪几份文书相似. 与 $k=2$ 比较，``压缩过头''的现象是否更严重？
			\item\label{ex:dec:B15} 证明：实对称矩阵 $\vv A$ 的特征值都在 $[a,b]$ 中，当且仅当 $\vv A-a\vv I$ 与 $b\vv I-\vv A$ 的二次型都非负.
		\end{exB}
		
		\Cgroup
		\begin{exC}
			\item\label{ex:dec:C1} （Eckart–Young 定理的证明）设 $\rk\vv B\le k$. 证明 $\ker\vv B$ 与 $\spn(\vv v_1,\ldots,\vv v_{k+1})$ 有非零公共向量 $\vv x$，由此证明 $\|\vv A-\vv B\|_2\ge\sigma_{k+1}$. 再借助 Weyl 不等式 $\sigma_{i+j-1}(\vv X+\vv Y)\le\sigma_i(\vv X)+\sigma_j(\vv Y)$（可承认）证明 Frobenius 范数的情形.
			\item\label{ex:dec:C2} （Cauchy 交错定理）设实对称矩阵 $\vv A$ 的特征值为 $\lambda_1\ge\cdots\ge\lambda_n$，删去最后一行一列所得矩阵的特征值为 $\mu_1\ge\cdots\ge\mu_{n-1}$. 证明 $\lambda_1\ge\mu_1\ge\lambda_2\ge\mu_2\ge\cdots\ge\mu_{n-1}\ge\lambda_n$.（提示：Courant–Fischer 极小极大刻画.）
			\item\label{ex:dec:C3} （复谱定理）证明：复方阵 $\vv A$ 可以被酉矩阵对角化（存在 $\vv U$ 使 $\bar{\vv U}^\trans\vv{AU}$ 为对角矩阵），当且仅当 $\vv A\bar{\vv A}^\trans=\bar{\vv A}^\trans\vv A$. 实对称矩阵、实正交矩阵、实反对称矩阵都满足这个条件.
			\item\label{ex:dec:C4} （正交 Procrustes 问题）给定 $n\times k$ 矩阵 $\vv A,\vv B$，求正交矩阵 $\vv Q$ 使 $\|\vv{QA}-\vv B\|_F$ 最小. 证明：若 $\vv{BA}^\trans=\vv U\vv\Sigma\vv V^\trans$，则最优解为 $\vv Q=\vv{UV}^\trans$. 这个问题出现在比较两组测量数据、对齐两份词向量等场合.
			\item\label{ex:dec:C5} （Ky Fan 定理）设实对称矩阵的特征值为 $\lambda_1\ge\cdots\ge\lambda_n$. 证明对任意列标准正交的 $n\times k$ 矩阵 $\vv X$，$\tr(\vv X^\trans\vv{AX})\le\lambda_1+\cdots+\lambda_k$. 由此说明前 $k$ 个主成分``一起''解释的方差也是最大的.
			\item\label{ex:dec:C6} （幂迭代）设 $\vv A$ 可对角化，特征值满足 $|\lambda_1|>|\lambda_2|\ge\cdots$. 证明：对``一般的''初始向量 $\vv x_0$，$\vv x_k=\vv A^k\vv x_0/\|\vv A^k\vv x_0\|$ 收敛到属于 $\lambda_1$ 的单位特征向量（可能差一个符号），收敛速度由 $|\lambda_2/\lambda_1|$ 决定. 对 $\begin{pmatrix}2&1\\1&2\end{pmatrix}$ 从 $(1,0)^\trans$ 出发迭代五步验证.
			\item\label{ex:dec:C7} （图的 Laplace 矩阵）设 $\vv L=\vv D-\vv A$ 是无向图的 Laplace 矩阵（见第五章章末习题~\ref{ex:det:C4}）. 证明 $\vv x^\trans\vv L\vv x=\sum_{\text{边 }ij}(x_i-x_j)^2$，因而 $\vv L$ 半正定；特征值 $0$ 的重数等于图的连通分支数. 第二小的特征值对应的特征向量（Fiedler 向量）的正负号常被用来把图``切''成两半——这就是谱聚类的思想.
			\item\label{ex:dec:C8} （正交回归）对中心化数据点 $\vv z_1,\ldots,\vv z_n\in\mathbb{R}^2$，求过原点的直线使各点到直线的\textbf{垂直}距离平方和最小. 证明最优方向就是第一主成分方向. 与普通最小二乘（只看竖直方向的残差）比较，两条直线何时相同？
			\item\label{ex:dec:C9} （扰动分析）设 $\vv A$ 可逆，$\vv{Ax}=\vv b$，$\vv A(\vv x+\delta\vv x)=\vv b+\delta\vv b$. 证明 $\dfrac{\|\delta\vv x\|}{\|\vv x\|}\le\kappa(\vv A)\dfrac{\|\delta\vv b\|}{\|\vv b\|}$，$\kappa(\vv A)=\sigma_1/\sigma_n$，并说明等号可以取到. 用它解释第六章病态方程组的例子.
		\end{exC}
	}
	
	% =====================================================================
	% 第十一章　二次型
	% =====================================================================
	\newcommand{\ExQuadratic}{%
		\exsetheading
		\Agroup
		\begin{exA}
			\item\label{ex:quad:A1} 写出下列二次型的矩阵与秩：
			(1) $x_1^2-3x_2^2+4x_1x_2-2x_2x_3$；(2) $x_1x_2+x_2x_3+x_3x_1$；(3) $(x_1+x_2+x_3)^2$；(4) $\sum_{i=1}^n x_i^2+\sum_{i<j}x_ix_j$.
			\item\label{ex:quad:A2} 写出以 $\begin{pmatrix}1&-1&0\\-1&2&3\\0&3&0\end{pmatrix}$ 为矩阵的二次型.
			\item\label{ex:quad:A3} 用配方法化为标准形，并写出所用的可逆变量替换：
			(1) $x_1^2+2x_2^2+5x_3^2+2x_1x_2+6x_2x_3$；(2) $x_1^2-4x_1x_2+4x_2^2-x_3^2$；(3) $x_1x_2-x_2x_3$.
			\item\label{ex:quad:A4} 用正交变换法把 $2x_1^2+3x_2^2+3x_3^2+4x_2x_3$ 化为标准形，写出正交矩阵. 曲面 $2x_1^2+3x_2^2+3x_3^2+4x_2x_3=1$ 是什么形状？
			\item\label{ex:quad:A5} 求 A3 中各二次型的正、负惯性指数与符号差，并写出规范形.
			\item\label{ex:quad:A6} 判断下列二次型是否正定：
			(1) $2x_1^2+x_2^2+3x_3^2+2x_1x_2$；(2) $x_1^2+x_2^2+x_3^2+4x_1x_2$；(3) $5x_1^2+x_2^2+5x_3^2+4x_1x_2-8x_1x_3-4x_2x_3$.
			\item\label{ex:quad:A7} $t$ 取何值时，$2x_1^2+x_2^2+x_3^2+2x_1x_2+tx_2x_3$ 正定？
			\item\label{ex:quad:A8} 判断曲线类型并求主轴方向：(1) $x^2-4xy+y^2=1$；(2) $3x^2+2xy+3y^2=8$.
			\item\label{ex:quad:A9} 判断下列矩阵是否合同：(1) $\diag(1,2,-3)$ 与 $\diag(-1,5,6)$；(2) $\begin{pmatrix}0&1\\1&0\end{pmatrix}$ 与 $\vv I_2$；(3) $\begin{pmatrix}0&1\\1&0\end{pmatrix}$ 与 $\diag(1,-1)$. 合同时求出 $\vv C$.
			\item\label{ex:quad:A10} 证明：若 $\vv A$ 正定，则 $\vv A^*$ 也正定.
			\item\label{ex:quad:A11} 求 $f(x,y)=x^3+y^3-3xy$ 的全部驻点，并用 Hessian 矩阵判断其类型.
			\item\label{ex:quad:A12} 判断 $\begin{pmatrix}2&-1&0\\-1&2&-1\\0&-1&2\end{pmatrix}$ 是否正定，并求它的 Cholesky 分解.
			\item\label{ex:quad:A13} 求 $\begin{pmatrix}4&2&-2\\2&10&2\\-2&2&6\end{pmatrix}$ 的 Cholesky 分解.
			\item\label{ex:quad:A14} 三种资产收益率的协方差矩阵为 $\vv\Sigma=\begin{pmatrix}0.04&0.01&0\\0.01&0.09&0\\0&0&0.01\end{pmatrix}$. 求组合 $\vv w=(0.5,0.3,0.2)^\trans$ 的方差与标准差.
			\item\label{ex:quad:A15} 两种资产收益率的标准差为 $30\%$ 与 $20\%$，互不相关. 求最小方差组合及其标准差.
			\item\label{ex:quad:A16} 证明：合同变换保持矩阵的对称性与秩. 合同变换保持行列式的符号吗？保持行列式的值吗？
			\item\label{ex:quad:A17} 分别在 $\mathbb{R}$ 上与 $\mathbb{C}$ 上写出二次型 $x_1^2-x_2^2+2x_3^2$ 的规范形.
			\item\label{ex:quad:A18} 判断下列方程表示的二次曲面类型：$x^2+y^2-z^2=1$，$x^2-y^2-z^2=1$，$x^2+y^2-z^2=0$，$x^2+y^2=1$.（提示：先看正、负惯性指数.）
			\item\label{ex:quad:A19} 说明 $\begin{pmatrix}1&2\\2&1\end{pmatrix}$ 不正定，并找出一个使二次型取负值的向量.
			\item\label{ex:quad:A20} 设 $\vv A$ 正定. 证明：$\vv A$ 的任一主子矩阵（取相同的若干行和列）都正定，从而所有主子式都为正.
		\end{exA}
		
		\Bgroup
		\begin{exB}
			\item\label{ex:quad:B1} 设 $\vv A,\vv B$ 正定且 $\vv{AB}=\vv{BA}$. 证明 $\vv{AB}$ 正定. 举例说明``可交换''不能去掉（此时 $\vv{AB}$ 甚至不对称）.
			\item\label{ex:quad:B2} 设 $\vv A$ 正定，$\vv B$ 半正定. 证明 $\det(\vv A+\vv B)\ge\det\vv A$，等号成立当且仅当 $\vv B=\vv O$.（提示：同时合同对角化.）
			\item\label{ex:quad:B3} 设 $\vv A$ 正定，$\vv B$ 对称. 证明：当 $|t|$ 充分小时 $\vv A+t\vv B$ 正定.
			\item\label{ex:quad:B4} 设 $\vv A$ 正定. 证明对 $i\ne j$ 有 $|a_{ij}|<\sqrt{a_{ii}a_{jj}}$，从而 $\vv A$ 的绝对值最大的元素在对角线上.
			\item\label{ex:quad:B5} 证明：实对称矩阵 $\vv A$ 满足 $\lambda_{\min}\|\vv x\|^2\le\vv x^\trans\vv{Ax}\le\lambda_{\max}\|\vv x\|^2$.
			\item\label{ex:quad:B6} 求 $\begin{pmatrix}0&1&1\\1&0&1\\1&1&0\end{pmatrix}$ 的正、负惯性指数（用两种方法：特征值；配方法）.
			\item\label{ex:quad:B7} $\lambda$ 取何值时，$\lambda(x_1^2+x_2^2+x_3^2)+2x_1x_2+2x_1x_3-2x_2x_3$ 正定？
			\item\label{ex:quad:B8} 求二次型 $\vv x^\trans\vv{Ax}$ 在约束 $\vv c^\trans\vv x=1$ 下的最小值（$\vv A$ 正定）：用拉格朗日乘数法证明最优解为 $\vv x^*=\dfrac{\vv A^{-1}\vv c}{\vv c^\trans\vv A^{-1}\vv c}$. 最小方差组合是它的特例.
			\item\label{ex:quad:B9} 设三种资产互不相关、方差都为 $\sigma^2$，期望收益率为 $\mu_1<\mu_2<\mu_3$. 求期望收益为 $m$ 的最小方差组合的权重，并说明有效前沿是一条抛物线（在``方差—期望收益''平面上）.
			\item\label{ex:quad:B10} 设 $\vv A=\begin{pmatrix}2&1\\1&2\end{pmatrix}$，$\vv B=\begin{pmatrix}1&0\\0&-1\end{pmatrix}$. 求可逆矩阵 $\vv C$ 使 $\vv C^\trans\vv{AC}=\vv I$ 且 $\vv C^\trans\vv{BC}$ 为对角矩阵.
			\item\label{ex:quad:B11} 设 $\vv A$ 正定. 求 $f(\vv x)=\vv x^\trans\vv{Ax}+\vv b^\trans\vv x+c$ 的最小值点与最小值. 若 $\vv A$ 只是半正定，最小值一定存在吗？
			\item\label{ex:quad:B12} 设 $\vv A$ 半正定，$\vv x^\trans\vv{Ax}=0$. 证明 $\vv{Ax}=\zero$.
			\item\label{ex:quad:B13} （Schur 补判据）证明：对称分块矩阵 $\begin{pmatrix}\vv A&\vv B\\ \vv B^\trans&\vv C\end{pmatrix}$ 正定，当且仅当 $\vv A$ 正定且 $\vv C-\vv B^\trans\vv A^{-1}\vv B$ 正定.
			\item\label{ex:quad:B14} 设 $\vv A$ 是二阶正定矩阵. 证明椭圆区域 $\{\vv x:\vv x^\trans\vv{Ax}\le1\}$ 的面积为 $\pi/\sqrt{\det\vv A}$.（提示：用主轴定理化为标准椭圆.）
			\item\label{ex:quad:B15} 设 $\vv A$ 为 $m\times n$ 实矩阵，$\mu>0$. 证明 $\vv A^\trans\vv A+\mu\vv I$ 正定，并说明岭回归 $\min\|\vv y-\vv{Ax}\|^2+\mu\|\vv x\|^2$ 的解为 $(\vv A^\trans\vv A+\mu\vv I)^{-1}\vv A^\trans\vv y$.
			\item\label{ex:quad:B16} 证明：实对称矩阵 $\vv A$ 正定，当且仅当存在正定矩阵 $\vv B$ 使 $\vv A=\vv B^2$；当且仅当 $\vv A^{-1}$ 存在且正定.
		\end{exB}
		
		\Cgroup
		\begin{exC}
			\item\label{ex:quad:C1} （Schur 乘积定理）对同阶矩阵定义逐元素乘积 $(\vv A\circ\vv B)_{ij}=a_{ij}b_{ij}$. 证明：若 $\vv A,\vv B$ 都半正定，则 $\vv A\circ\vv B$ 也半正定.（提示：把 $\vv B$ 写成 $\sum\lambda_k\vv q_k\vv q_k^\trans$.）
			\item\label{ex:quad:C2} （Minkowski 行列式不等式）设 $\vv A,\vv B$ 都是 $n$ 阶正定矩阵. 证明 $\det(\vv A+\vv B)^{1/n}\ge\det(\vv A)^{1/n}+\det(\vv B)^{1/n}$.（提示：同时合同对角化后用 $\prod(1+\mu_i)^{1/n}\ge1+\prod\mu_i^{1/n}$.）
			\item\label{ex:quad:C3} （有理数域上的二次型）证明 $x^2+y^2$ 与 $3x^2+3y^2$ 在 $\mathbb{R}$ 上合同，但在 $\mathbb{Q}$ 上不合同（即不存在有理系数的可逆替换把一个变为另一个）.（提示：若合同，则存在有理数 $a,b$ 使 $a^2+b^2=3$.）这说明惯性定理依赖于实数可以随意开平方.
			\item\label{ex:quad:C4} 函数 $f=x^2+y^4$ 与 $g=x^2-y^4$ 在原点处的 Hessian 矩阵相同且只是半正定，但前者在原点取极小值，后者不取极值. 由此说明 Hessian 半正定时二阶判别法失效.
			\item\label{ex:quad:C5} （广义特征值）设 $\vv A$ 对称，$\vv B$ 正定. 证明 $\max_{\vv x\ne\zero}\dfrac{\vv x^\trans\vv{Ax}}{\vv x^\trans\vv{Bx}}$ 等于方程 $\det(\vv A-\lambda\vv B)=0$ 的最大根. 统计学中的判别分析（Fisher LDA）正是这样求最佳投影方向的.
			\item\label{ex:quad:C6} （Mahalanobis 距离）设 $\vv\Sigma$ 正定，定义 $d(\vv x,\vv y)=\sqrt{(\vv x-\vv y)^\trans\vv\Sigma^{-1}(\vv x-\vv y)}$. 证明它满足三角不等式，且在仿射变换 $\vv x\mapsto\vv{Px}+\vv b$、$\vv\Sigma\mapsto\vv P\vv\Sigma\vv P^\trans$ 下不变. 它与欧氏距离有何不同？
			\item\label{ex:quad:C7} （Hilbert 矩阵）证明 $n$ 阶矩阵 $\vv H=\big(\frac1{i+j-1}\big)$ 正定.（提示：$\frac1{i+j-1}=\int_0^1x^{i-1}x^{j-1}\,\mathrm dx$，它是某个内积下的 Gram 矩阵.）
			\item\label{ex:quad:C8} （切点组合）设有无风险资产，收益率为 $r_f$. 证明使``夏普比率''$\dfrac{\vv\mu^\trans\vv w-r_f}{\sqrt{\vv w^\trans\vv\Sigma\vv w}}$（约束 $\vv 1^\trans\vv w=1$）最大的风险资产组合为 $\vv w\propto\vv\Sigma^{-1}(\vv\mu-r_f\vv 1)$.（提示：先去掉约束，比值关于 $\vv w$ 是零次齐次的；再用 Cauchy–Schwarz 不等式.）
			\item\label{ex:quad:C9} 证明：秩为 $r$、正惯性指数为 $p$ 的实二次型可以写成 $p$ 个线性型的平方和减去 $r-p$ 个线性型的平方和，并且不能写成少于 $r$ 个线性型的平方（带 $\pm$ 号）的组合.
		\end{exC}
	}
	
	% =====================================================================
	% 第十二章　*对偶空间与双线性型（选读）
	% =====================================================================
	\newcommand{\ExBilinear}{%
		\exsetheading
		\Agroup
		\begin{exA}
			\item\label{ex:bil:A1} 判断下列映射是否为相应空间上的线性函数：(1) $\mathbb{R}^3\to\mathbb{R}$，$\vv x\mapsto x_1-2x_3$；(2) $\vv x\mapsto x_1x_2$；(3) $\mathbb{R}[x]\to\mathbb{R}$，$f\mapsto f'(0)$；(4) $f\mapsto\int_0^1f(x)^2\,\mathrm dx$；(5) $M_n(\mathbb{R})\to\mathbb{R}$，$\vv X\mapsto\tr(\vv{AX})$（$\vv A$ 固定）.
			\item\label{ex:bil:A2} 求 $\mathbb{R}^3$ 的基 $(1,0,0)^\trans$，$(1,1,0)^\trans$，$(1,1,1)^\trans$ 的对偶基（把每个对偶基向量写成 $\vv x\mapsto a_1x_1+a_2x_2+a_3x_3$ 的形式）.
			\item\label{ex:bil:A3} 把线性函数 $f(\vv x)=x_1+2x_2+3x_3$ 用 A2 中的对偶基线性表出.
			\item\label{ex:bil:A4} 求 $W=\spn\big((1,1,0)^\trans,(0,1,1)^\trans\big)$ 的零化子 $W^0\subset(\mathbb{R}^3)^*$ 的一组基.
			\item\label{ex:bil:A5} 设线性映射 $\mathscr T:\mathbb{R}^2\to\mathbb{R}^3$ 在标准基下的矩阵为 $\begin{pmatrix}1&2\\0&1\\3&-1\end{pmatrix}$. 求对偶映射 $\mathscr T^*$ 在对偶基下的矩阵，并计算 $\mathscr T^*$ 作用在线性函数 $\vv y\mapsto y_1+y_2+y_3$ 上的结果.
			\item\label{ex:bil:A6} 求双线性型 $f(\vv x,\vv y)=x_1y_1+2x_1y_2-x_2y_1+3x_2y_2$ 在标准基下的度量矩阵，以及在基 $(1,1)^\trans$，$(1,-1)^\trans$ 下的度量矩阵，并验证 $\vv B=\vv P^\trans\vv{AP}$.
			\item\label{ex:bil:A7} 把 A6 中的 $f$ 分解为对称部分与交错部分之和.
			\item\label{ex:bil:A8} 判断下列双线性型是否非退化：(1) $x_1y_1-x_2y_2$；(2) $x_1y_2+x_2y_1$；(3) $(x_1+x_2)(y_1+y_2)$.
			\item\label{ex:bil:A9} 写出对称双线性型 $f(\vv x,\vv y)=x_1y_1+x_1y_2+x_2y_1+3x_2y_2$ 对应的二次型；反过来，由二次型 $x_1^2-4x_1x_2+2x_2^2$ 用极化恒等式求对应的对称双线性型.
			\item\label{ex:bil:A10} 验证 $\mathbb{R}^2$ 上的 $f(\vv x,\vv y)=\det(\vv x,\vv y)$ 是交错双线性型，写出它的度量矩阵. 它非退化吗？
			\item\label{ex:bil:A11} 设 $a_1,a_2,a_3$ 互异，$f_i(p)=p(a_i)$ 是 $\mathbb{R}[x]_3$ 上的线性函数. 对 $a_i=0,1,2$，求以 $f_1,f_2,f_3$ 为对偶基的 $\mathbb{R}[x]_3$ 的基.
			\item\label{ex:bil:A12} 在 $\mathbb{R}^4$ 中，设 $W$ 由 $x_1+x_2=0$，$x_3-x_4=0$ 定义. 求 $\dim W$ 与 $\dim W^0$，并写出 $W^0$ 的一组基.
			\item\label{ex:bil:A13} 证明 $\langle\vv A,\vv B\rangle=\tr(\vv{AB})$ 是 $M_n(\mathbb{R})$ 上非退化的对称双线性型. 它是内积吗？
			\item\label{ex:bil:A14} 对 $\mathbb{R}[x]_3$ 上的线性函数 $\varphi(p)=\int_0^1p(x)\,\mathrm dx$，求它在 $1,x,x^2$ 的对偶基下的坐标.
			\item\label{ex:bil:A15} （Riesz 表示定理的有限维情形）设 $V$ 是有限维实内积空间. 证明：对每个 $f\in V^*$，存在唯一的 $\vv\beta\in V$ 使 $f(\vv\alpha)=\langle\vv\alpha,\vv\beta\rangle$ 对一切 $\vv\alpha$ 成立.（提示：取标准正交基 $\vv e_1,\ldots,\vv e_n$，令 $\vv\beta=\sum f(\vv e_i)\vv e_i$.）
		\end{exA}
		
		\Bgroup
		\begin{exB}
			\item\label{ex:bil:B1} 证明：$f_1,\ldots,f_k\in V^*$ 线性无关，当且仅当映射 $V\to\mathbb{F}^k$，$\vv\alpha\mapsto(f_1(\vv\alpha),\ldots,f_k(\vv\alpha))^\trans$ 是满射.
			\item\label{ex:bil:B2} 设 $f,g\in V^*$ 且 $\ker f\subset\ker g$. 证明存在常数 $c$ 使 $g=cf$. 推广：若 $\ker f_1\cap\cdots\cap\ker f_k\subset\ker g$，则 $g$ 是 $f_1,\ldots,f_k$ 的线性组合.
			\item\label{ex:bil:B3} 证明：$\ker f_1\cap\cdots\cap\ker f_k=\{\zero\}$ 当且仅当 $f_1,\ldots,f_k$ 张成 $V^*$. 用方程组的语言解释这个结论.
			\item\label{ex:bil:B4} 证明：$\mathscr T$ 是单射当且仅当 $\mathscr T^*$ 是满射；$\mathscr T$ 是满射当且仅当 $\mathscr T^*$ 是单射.
			\item\label{ex:bil:B5} （Simpson 公式）求常数 $a,b,c$，使 $\int_0^2p(x)\,\mathrm dx=a\,p(0)+b\,p(1)+c\,p(2)$ 对一切 $p\in\mathbb{R}[x]_3$ 成立. 用对偶基的语言说明为什么这样的 $a,b,c$ 存在且唯一. 验证这个公式对三次多项式也成立.
			\item\label{ex:bil:B6} 证明：$V$ 上全体双线性型构成的空间与 $M_n(\mathbb{F})$ 同构，其中对称型、交错型分别构成维数为 $\frac{n(n+1)}2$、$\frac{n(n-1)}2$ 的子空间，且整个空间是二者的直和.
			\item\label{ex:bil:B7} 证明：若 $V$ 上存在非退化的交错双线性型，则 $\dim V$ 是偶数.
			\item\label{ex:bil:B8} 求双线性型 $f(\vv x,\vv y)=x_1y_1+x_1y_2+x_2y_1+x_2y_2+x_3y_3$ 的``根'' $\{\vv\alpha:f(\vv\alpha,\cdot)=0\}$.
			\item\label{ex:bil:B9} 在 $\mathbb{R}^2$ 上取 $f(\vv x,\vv y)=x_1y_1-x_2y_2$. 求全部``迷向向量''（$f(\vv\alpha,\vv\alpha)=0$），并说明它们构成两条直线. 再证明保持 $f$ 的线性变换 $\vv P$（$\vv P^\trans\diag(1,-1)\vv P=\diag(1,-1)$，$\det\vv P=1$，且保持 $x_1>|x_2|$ 的区域）都形如 $\begin{pmatrix}\cosh t&\sinh t\\ \sinh t&\cosh t\end{pmatrix}$.
			\item\label{ex:bil:B10} 设 $W$ 是 $V$ 的子空间，$\mathscr T:V\to V$ 线性且 $\mathscr T(W)\subset W$. 证明 $\mathscr T^*(W^0)\subset W^0$.
		\end{exB}
		
		\Cgroup
		\begin{exC}
			\item\label{ex:bil:C1} 证明：对无限维空间 $\mathbb{F}[x]$，自然映射 $\vv\alpha\mapsto\hat{\vv\alpha}$ 是单射但不是满射.（提示：$\mathbb{F}[x]^*$ 中的元素由它在 $1,x,x^2,\ldots$ 上的取值任意决定，可以``无限多个非零''.）
			\item\label{ex:bil:C2} 证明：$V$ 上的双线性型与线性映射 $V\to V^*$ 一一对应：$f\mapsto(\vv\alpha\mapsto f(\vv\alpha,\cdot))$. 用这个对应和对偶基的换基规则，重新推出度量矩阵的换基公式 $\vv B=\vv P^\trans\vv{AP}$.
			\item\label{ex:bil:C3} 证明：二阶矩阵 $\vv P$ 满足 $\vv P^\trans\vv J\vv P=\vv J$（$\vv J=\begin{pmatrix}0&1\\-1&0\end{pmatrix}$）当且仅当 $\det\vv P=1$. 四阶时，满足 $\vv P^\trans\diag(\vv J,\vv J)\vv P=\diag(\vv J,\vv J)$ 的矩阵（辛矩阵）的行列式是多少？
			\item\label{ex:bil:C4} 求 $M_n(\mathbb{R})$ 上对称双线性型 $\tr(\vv{AB})$ 的正、负惯性指数.（提示：对称矩阵与反对称矩阵在这个型下``互相正交''.）
			\item\label{ex:bil:C5} 设 $f$ 是 $V$ 上的对称双线性型，$W$ 是子空间，$W^\perp=\{\vv\alpha:f(\vv\alpha,\vv\beta)=0\ \forall\vv\beta\in W\}$. 证明 $\dim W^\perp\ge\dim V-\dim W$，$f$ 非退化时等号成立；举例说明即使 $f$ 非退化，$W\cap W^\perp$ 也可能不为零.
			\item\label{ex:bil:C6} 证明：$\mathbb{F}^n$ 上全体交错 $n$ 重线性函数构成的空间是一维的，由行列式张成. 这给出第五章章末习题~\ref{ex:det:C1} 的``对偶空间版本''.
			\item\label{ex:bil:C7} （线性规划对偶）考虑问题：在 $\vv{Ax}=\vv b$，$\vv x\ge\zero$ 下最小化 $\vv c^\trans\vv x$. 证明对任意满足 $\vv A^\trans\vv y\le\vv c$ 的 $\vv y$ 与任意可行的 $\vv x$，都有 $\vv b^\trans\vv y\le\vv c^\trans\vv x$（弱对偶）. 用``价格是数量空间上的线性函数''解释 $\vv y$ 的经济含义（影子价格）.
		\end{exC}
	} 之后，
%%      exercises.tex 中每道题都有标签 ex:<章键>:<组号>，本文件按标签引用题号）.
%% 章键：elim vs mat lmap det sol poly eig ip dec quad bil（依次为第一至十二章）.
%% =====================================================================

\providecommand{\ans}[2]{\par\smallskip\noindent\textbf{\ref{#1}.}\ #2\par}
\setlength{\emergencystretch}{3em}
\chapter{习题答案与提示}\label{app:answers}

本附录给出各章章末习题 A 组的答案与 B 组的提示. A 组的计算题给出最终结果，必要时附关键步骤；证明题只给出思路，完整的书写请读者自行完成. B 组只给出关键的一步或方向，希望读者在看提示之前先独立思考至少一刻钟. C 组一般不提供提示；正文或其他题目的答案中引用到的 C 组题例外，在该章末尾给出证明要点.

计算题的答案往往不唯一（例如基础解系、过渡矩阵的选取、特征向量的倍数），只要满足题目的全部要求，与这里的答案不同也是正确的. 检验的方法通常比答案本身更重要：代回原方程、计算迹与行列式、验证正交性.


\section{第一章　线性方程组与消元法}

\subsection*{A 组答案}

\ans{ex:elim:A1}{第 1、4 个是简化阶梯形；第 2 个是阶梯形但不是简化阶梯形（主元 $2$ 不为 $1$）；第 3、5 个不是阶梯形（第 3 个第二行主元在第一行主元左侧；第 5 个零行在非零行上方）.}

\ans{ex:elim:A2}{$\begin{pmatrix}3&0&-1&2\\0&1&4&0\\1&-1&0&5\end{pmatrix}$；反过来为 $x_1-2x_2=3$，$x_2+5x_3=-1$（未知量为 $x_1,x_2,x_3$，最后一列为常数项；$x_3$ 在第一个方程中的系数为 $0$）.}

\ans{ex:elim:A3}{$(x_1,x_2,x_3)=(1,2,3)$.}

\ans{ex:elim:A4}{$(x_1,x_2,x_3,x_4)=(1,-1,2,0)$.}

\ans{ex:elim:A5}{无解. $(r_3-r_1)-(r_2-2r_1)$ 给出 $(0,0,0\mid 2)$，即 $0=2$.}

\ans{ex:elim:A6}{主变量 $x_1,x_3$，自由变量 $x_2,x_4$. 全部解为 $x_1=4-2s-2t$，$x_2=s$，$x_3=1-t$，$x_4=t$（$s,t$ 任意）.}

\ans{ex:elim:A7}{$x_1=\frac27x_3+\frac37x_4$，$x_2=\frac57x_3+\frac47x_4$. 令 $x_3=7s$，$x_4=7t$，全部解为 $s(2,5,7,0)^\trans+t(3,4,0,7)^\trans$.}

\ans{ex:elim:A8}{$\begin{pmatrix}1&2&0&2\\0&0&1&-1\\0&0&0&0\end{pmatrix}$. 一种做法：$r_1\leftrightarrow r_2$；$r_2-2r_1$，$r_3-3r_1$；$-\frac14r_2$；$r_3+2r_2$；$r_1-r_2$.}

\ans{ex:elim:A9}{共四类：$\vv O$；$\begin{pmatrix}1&a\\0&0\end{pmatrix}$（$a$ 任意）；$\begin{pmatrix}0&1\\0&0\end{pmatrix}$；$\vv I_2$.}

\ans{ex:elim:A10}{由前两个方程得 $x=2$，$y=-1$，代入第三个得 $4-k=5$，$k=-1$.}

\ans{ex:elim:A11}{有解当且仅当 $b=3a$. 这时全部解为 $x=a-2t$，$y=t$（$t$ 任意）.}

\ans{ex:elim:A12}{$y=x^2-2x+3$.}

\ans{ex:elim:A13}{$y=1-2x+x^3$，即 $a=1$，$b=-2$，$c=0$，$d=1$.}

\ans{ex:elim:A14}{$30$ 元、$50$ 元、$80$ 元票分别为 $40$，$40$，$20$ 张.}

\ans{ex:elim:A15}{守恒方程 $3x_1=x_3$，$8x_1=2x_4$，$2x_2=2x_3+x_4$. 解为 $(x_1,x_2,x_3,x_4)=t(1,5,3,4)$，最小正整数解：$\mathrm{C_3H_8}+5\,\mathrm{O_2}\to3\,\mathrm{CO_2}+4\,\mathrm{H_2O}$.}

\ans{ex:elim:A16}{(1) 四个路口的方程为 $x_1-x_4=100$，$x_1-x_2=150$，$x_3-x_2=150$，$x_3-x_4=100$，全部解为 $x_1=t+100$，$x_2=t-50$，$x_3=t+100$，$x_4=t$. (2) 由 $x_2\ge0$ 得 $t\ge50$，故 $x_1$ 最少为 $150$. (3) 方程组无解：外部驶入总量 $1000$，驶出总量变为 $1100$，第一、三个方程之和减去第二、四个方程，得 $0=-100$. 守恒律是``进入路网的车辆总数等于离开的总数''.}

\ans{ex:elim:A17}{若 $a\ne0$，用第一个方程消去第二个方程中的 $x$，得 $(ad-bc)y=af-ce$. 解唯一当且仅当 $y$ 的系数 $ad-bc\ne0$. $a=0$ 时交换两个方程（或对 $c$ 讨论），同样得到条件 $ad-bc\ne0$.}

\ans{ex:elim:A18}{设 $\vv A\vv x_1=\vv A\vv x_2=\vv b$（这里先借用第三章的记号，也可以逐个方程写出），则 $\vv A(\vv x_1+t(\vv x_2-\vv x_1))=(1-t)\vv b+t\vv b=\vv b$. 若有两个不同的解，$t$ 取不同值给出不同的解，于是有无穷多个.}

\ans{ex:elim:A19}{第二个方程组的第一个方程等于第一个方程组两个方程之和，第二个方程相同，所以 (I) 的解都是 (II) 的解；反之 (I) 的第二个方程是 (II) 的第一个方程减去第二个方程. 两者的解集都是 $\{(t,t,3-2t)\}$.}

\ans{ex:elim:A20}{例如方程组 $x=0$，$y=0$. 同时做两步后得到 $x+y=0$，$x+y=0$，解集变大. 不矛盾：``同时''做的这两步合起来不是一次初等行变换，也不能被撤销（它把两行变成了相同的两行，丢失了信息）.}

\ans{ex:elim:A21}{有解当且仅当 $a=0$. 这时系数部分有 $2$ 个主元，$4$ 个未知量，有 $2$ 个自由变量.}

\subsection*{B 组提示}

\ans{ex:elim:B1}{用 $r_4-3r_1$，再加上 $r_2$，得 $(0,0,0,a-1\mid0)$；$r_3+r_2$ 得 $(0,0,a-1,0\mid b+1)$. $a\ne1$ 时解唯一；$a=1$，$b\ne-1$ 时无解；$a=1$，$b=-1$ 时有无穷多解：$(-1+s+t,\ 1-2s-2t,\ s,\ t)$.}

\ans{ex:elim:B2}{把第三个方程移到最前面再消元，得到的阶梯形后两行为 $(0,\lambda,-\lambda\mid3-\lambda)$ 与 $\big(0,0,-\lambda(\lambda+3)\mid-(\lambda+3)(\lambda-1)\big)$. $\lambda\ne0,-3$ 时有三个主元，解唯一；$\lambda=0$ 时倒数第二行变为 $(0,0,0\mid3)$，无解；$\lambda=-3$ 时最后一行为零行，只有两个主元，有无穷多解.}

\ans{ex:elim:B3}{齐次方程组只有零解说明简化阶梯形有 $n$ 个主元，即每一列都有主元. 而系数部分只有 $n$ 行，不可能出现 $(0,\ldots,0\mid d)$ 形的行.}

\ans{ex:elim:B4}{初等行变换得到的每一行都是原方程的线性组合. 无解时简化阶梯形中出现 $(0,\ldots,0\mid d)$，$d\ne0$，把这一行除以 $d$ 即可. 反之若组合得 $0=1$，任何解代入都会矛盾.}

\ans{ex:elim:B5}{依次做 $r_i+r_j$，$r_j-r_i$，$r_i+r_j$，最后把第 $j$ 行乘以 $-1$，逐步写出两行的变化即可.}

\ans{ex:elim:B6}{消元只用到四则运算，简化阶梯形的元素仍是有理数. 有解时令全部自由变量取 $0$，主变量由有理数算出. 解唯一时没有自由变量，唯一的解就是这个有理数解.}

\ans{ex:elim:B7}{把 $n$ 个方程相加，左边为 $0$，故必要条件是 $a_1+\cdots+a_n=0$. 它也充分：令 $x_n=t$，则 $x_k=t+a_k+a_{k+1}+\cdots+a_{n-1}$（$k<n$）.}

\ans{ex:elim:B8}{把第一个方程减第二个、加第三个……，交替相加. $n$ 为奇数时得 $2x_1=a_1-a_2+a_3-\cdots+a_n$，$x_1$ 唯一确定，其余依次确定；$n$ 为偶数时左边相消为 $0$.}

\ans{ex:elim:B9}{方程组为 $(\vv I-\vv A)\vv x=\vv d$（逐个写出即可），解得各部门总产出为 $(100,\ 200,\ 150)^\trans$.}

\ans{ex:elim:B10}{处理第 $k$ 列时，需要对下面 $n-k$ 行各做一次倍加，每次约 $n-k+1$ 次乘法，求和 $\sum_{k}(n-k)(n-k+1)\approx\frac13n^3$；回代只需约 $\frac12n^2$ 次.}

\ans{ex:elim:B11}{考察三个方程的系数部分化为阶梯形后的主元个数：$3$ 个时交于一点；$2$ 个且相容时交于一条直线；$2$ 个且出现 $0=d$ 时，三条交线方向相同（都平行于系数矩阵的解空间方向），没有公共点. 例子可以取 $z=0$，$x=0$，$x+z=1$.}

\ans{ex:elim:B12}{直接验证 $\vv A(\vv x+\vv y)=\vv A\vv x+\vv A\vv y$，$\vv A(k\vv x)=k\vv A\vv x$（或逐个方程验证）. 后半句：设 $\vv x_0$ 是一个解，任一解 $\vv x$ 满足 $\vv x-\vv x_0$ 是齐次方程组的解.}

\ans{ex:elim:B13}{开放题. 一种设定：未知量为三位律师各承办三类案件的件数，共 $9$ 个；已知条件给出 $3$ 个独立方程，故至少还缺 $6$ 条独立的信息. 关键是先写清楚未知量，再数主元个数.}

\ans{ex:elim:B14}{每行之和为零即 $\vv A(1,\ldots,1)^\trans=\zero$. 后半句由 B12 的结论.}

\ans{ex:elim:B15}{(I) 的每个解满足 (II) 的每个方程，因为后者是前者的线性组合；反之亦然. 逆命题对有解的方程组成立，证明要用到第六章的结论（齐次方程组同解当且仅当系数矩阵行向量组等价，第六章章末习题~\ref{ex:sol:B11}）.}

\ans{ex:elim:B16}{$\lambda\ne1,-2$ 时解唯一：$x=-\frac{\lambda+1}{\lambda+2}$，$y=\frac1{\lambda+2}$，$z=\frac{(\lambda+1)^2}{\lambda+2}$；$\lambda=1$ 时三个方程相同，有无穷多解；$\lambda=-2$ 时三个方程相加得 $0=3$，无解.}

\subsection*{C 组证明要点（正文引用的题目）}

\ans{ex:elim:C9}{设 $\vv R_1,\vv R_2$ 都是 $\vv M$ 的简化阶梯形. 它们行等价，以它们为系数矩阵的齐次方程组同解（定理~\ref{thm:sameSolution}）. 对列数 $n$ 归纳. $n=1$ 时，简化阶梯形是零列或 $\vv e_1$，由 $\vv M$ 是否为零唯一决定.
	
	设 $n>1$. 删去最后一列，得到的 $\vv R_1',\vv R_2'$ 仍是简化阶梯形，且仍行等价（用同一串初等行变换），由归纳假设 $\vv R_1'=\vv R_2'$. 于是两者在前 $n-1$ 列中的主元位置相同，设共有 $r$ 个主元. 再比较最后一列：
	\begin{enumerate}
		\item 若最后一列在 $\vv R_1$ 中是主元列，则 $\vv R_1$ 的第 $r+1$ 行为 $(0,\ldots,0,1)$，$\vv R_1\vv x=\zero$ 的解都有 $x_n=0$. 若它在 $\vv R_2$ 中不是主元列，令 $x_n=1$、其余自由变量取 $0$，可得 $\vv R_2\vv x=\zero$ 的一个解，其 $x_n=1$，与同解矛盾. 所以它在两者中同为主元列，都等于 $\vv e_{r+1}$.
		\item 若最后一列在两者中都不是主元列，令 $x_n=1$、其余自由变量取 $0$，两个方程组各有唯一一个这样的解，其主变量的值分别是两矩阵最后一列前 $r$ 个分量的相反数. 两方程组同解，故这两个解相同，最后一列相同.
\end{enumerate}}

\section{第二章　向量空间}

\subsection*{A 组答案}

\ans{ex:vs:A1}{(1) 是. (2) 不是：$\frac12\cdot\sqrt2=\frac{\sqrt2}2$ 的 $\sqrt2$ 系数不是整数，乘法不封闭. (3) 不是：不含 $1$，也不对除法封闭. (4) 不是：$\sqrt[3]2\cdot\sqrt[3]2=\sqrt[3]4$ 不在其中. (5) 不是：$\mathrm i\cdot\mathrm i=-1$ 不是纯虚数.}

\ans{ex:vs:A2}{(1) 是. (2) 不是：$(1,0)^\trans$ 的负向量不在其中. (3) 不是：不含零多项式，且两个 $n$ 次多项式之和可能降次. (4) 是. (5) 是. (6) 不是：不含零向量.}

\ans{ex:vs:A3}{(1) 线性相关：$(4,1,-1)^\trans=2(1,2,-1)^\trans+(2,-3,1)^\trans$. (2) 线性无关. (3) 线性相关：$\mathbb{R}[x]_3$ 是三维的，四个向量必相关. (4) 线性相关：全 $1$ 矩阵等于前两个之和.}

\ans{ex:vs:A4}{$\vv\beta=\frac54\vv\alpha_1+\frac14\vv\alpha_2-\frac14\vv\alpha_3-\frac14\vv\alpha_4$. （这四个向量两两``对应分量乘积之和''为零，每个系数都可以直接算出.）}

\ans{ex:vs:A5}{以三个向量为列作矩阵并消元：$r_2-r_1$ 得 $(0,2,-2)$，再 $r_3+\frac12r_2$ 得 $(0,0,t-1)$. $t=1$ 时出现自由变量，线性相关；$t\ne1$ 时线性无关.（学完第五章后，可用行列式 $2t-2$ 验证.）}

\ans{ex:vs:A6}{秩为 $3$，一个极大无关组为 $\vv\alpha_1,\vv\alpha_2,\vv\alpha_4$；$\vv\alpha_3=3\vv\alpha_1+\vv\alpha_2$，$\vv\alpha_5=\vv\alpha_1+\vv\alpha_2+\vv\alpha_4$.}

\ans{ex:vs:A7}{(1) 线性相关：三者之和为 $\zero$. (2) 线性无关. (3) 线性无关：设 $k_1(\vv\alpha_1+2\vv\alpha_2)+k_2(\vv\alpha_2+2\vv\alpha_3)+k_3(\vv\alpha_3+2\vv\alpha_1)=\zero$，由 $\vv\alpha_i$ 线性无关得 $k_1+2k_3=0$，$2k_1+k_2=0$，$2k_2+k_3=0$. 依次代入得 $k_1=-2k_3$，$k_2=4k_3$，$k_3=-8k_3$，故 $k_3=0$，从而 $k_1=k_2=0$.}

\ans{ex:vs:A8}{$1\cdot\zero=\zero$ 是系数不全为零的组合；$\vv\alpha-\vv\alpha=\zero$ 同理；$k\vv\alpha=\zero$ 且 $\vv\alpha\ne\zero$ 时由命题~\ref{prop:linspc}(5) 得 $k=0$.}

\ans{ex:vs:A9}{由定理~\ref{thm:relvandexp}，两个向量相关当且仅当其中一个可由另一个表出，即是另一个的倍数. 三个向量时只能说``其中一个是另两个的线性组合''，不一定是某一个的倍数，例如 $\vv e_1,\vv e_2,\vv e_1+\vv e_2$.}

\ans{ex:vs:A10}{(1) $\vv\alpha+\vv\alpha=1\vv\alpha+1\vv\alpha=(1+1)\vv\alpha$. (2) 两边乘 $k^{-1}$. (3) $-\vv\alpha$ 与 $\vv\alpha$ 之和为零，由负向量的唯一性.}

\ans{ex:vs:A11}{(1) $\dim=3$，基 $(-1,1,0,0)^\trans$，$(1,0,1,0)^\trans$，$(1,0,0,1)^\trans$. (2) $\dim=\frac{n(n+1)}2$，基 $\vv E_{ij}$（$i\le j$）. (3) $\dim=n^2-1$，基 $\vv E_{ij}$（$i\ne j$）与 $\vv E_{11}-\vv E_{ii}$（$i=2,\ldots,n$）. (4) $\dim=3$，基 $x-1$，$(x-1)^2$，$(x-1)^3$.}

\ans{ex:vs:A12}{(1) 是. (2) 不是：$(1,0,0)^\trans+(0,1,0)^\trans$ 不在其中. (3) 不是：对实数数乘不封闭. (4) 是.}

\ans{ex:vs:A13}{维数为 $3$，一组基为 $\vv\alpha_1,\vv\alpha_2,\vv\alpha_4$（$\vv\alpha_3=\vv\alpha_1+\vv\alpha_2$）.}

\ans{ex:vs:A14}{$2x^2-x+3=4+3(x-1)+2(x-1)^2$，坐标为 $(4,3,2)^\trans$.}

\ans{ex:vs:A15}{坐标为 $(-1,4,-5,4)^\trans$（从右下角元素开始依次确定）.}

\ans{ex:vs:A16}{(1) 是. (2) 不是：$\mathbb{Q}$ 对实数的数乘不封闭. (3) 作为 $\mathbb{R}$ 上的空间是 $2$ 维，作为 $\mathbb{C}$ 上的空间是 $1$ 维.}

\ans{ex:vs:A17}{例如 $a_0+a_1x+\cdots+a_{n-1}x^{n-1}\mapsto(a_0,\ldots,a_{n-1})^\trans$；$\begin{pmatrix}a&b\\c&d\end{pmatrix}\mapsto(a,b,c,d)^\trans$. 两者都是取坐标的映射，保持线性运算.}

\ans{ex:vs:A18}{$W_1+W_2$ 是直和（$(1,1,1)^\trans\notin W_1$，交为零），等于 $\mathbb{R}^3$，维数 $3$；$W_3\subset W_1$，$W_1+W_3=W_1$ 不是直和，维数 $2$.}

\ans{ex:vs:A19}{$\dim(W_1+W_2)=3$，$\dim(W_1\cap W_2)=1$：$(0,-1,1,1)^\trans=-(1,1,0,0)^\trans+(1,0,1,1)^\trans$ 同属两者. $2+2-1=3$，与维数公式一致.}

\ans{ex:vs:A20}{它们都可由一个含 $r$ 个向量的极大无关组线性表出，由基本引理任意 $r+1$ 个相关. 任意 $r$ 个无关的向量再添一个就相关，由引理~\ref{lem:addone} 的逆否形式，组中每个向量都可由它们线性表出.}

\subsection*{B 组提示}

\ans{ex:vs:B1}{设两者互不包含，取 $\vv\alpha\in W_1\setminus W_2$，$\vv\beta\in W_2\setminus W_1$，证明 $\vv\alpha+\vv\beta$ 不属于并集.}
\ans{ex:vs:B2}{取 (II) 的一个极大无关组，它含 $r$ 个向量且可由 (I) 的极大无关组表出，二者秩相同，用定理~\ref{thm:ndimbasis}(2) 在 (I) 张成的空间中论证.}
\ans{ex:vs:B3}{(III) 可由两者的极大无关组合并表出. 第二个不等式：左边的每个向量可由两组的极大无关组合并表出.}
\ans{ex:vs:B4}{新组含 $n$ 个向量，只需证明它能表出 $\vv\alpha_j$：$a_j\ne0$ 时 $\vv\alpha_j=\frac1{a_j}(\vv\beta-\sum_{i\ne j}a_i\vv\alpha_i)$；$a_j=0$ 时新组线性相关.}
\ans{ex:vs:B5}{按次数从高到低比较系数，证明线性无关；个数等于维数 $n$.}
\ans{ex:vs:B6}{$a_1,a_2$ 确定整个数列，映射 $\{a_n\}\mapsto(a_1,a_2)$ 是同构. 代入验证 $2^n$ 与 $(-1)^n$ 满足递推且线性无关. $a_1=a_2=1$ 时 $a_n=\frac13\big(2^n-(-1)^n\big)$.}
\ans{ex:vs:B7}{(1) 设组合为零，令 $x=0$ 并逐次求导，或除以 $\mathrm e^{3x}$ 后令 $x\to+\infty$. (2) $\sin(x+1)=\cos1\cdot\sin x+\sin1\cdot\cos x$. (3) 秩为 $3$：$\cos2x=2\cos^2x-1$.}
\ans{ex:vs:B8}{证明对加、减、乘封闭，再对除法用两次``分母有理化''. 线性无关可设 $a+b\sqrt2=-(c+d\sqrt2)\sqrt3$ 后平方，利用 $\sqrt2,\sqrt3,\sqrt6$ 都是无理数. 维数为 $4$.}
\ans{ex:vs:B9}{按维数分：$\{\zero\}$；过原点的直线；（$\mathbb{R}^3$ 中）过原点的平面；整个空间.}
\ans{ex:vs:B10}{用维数公式：$\dim(W_1\cap W_2)=k_1+k_2-\dim(W_1+W_2)\ge k_1+k_2-n>0$.}
\ans{ex:vs:B11}{$f=f(a)+(f-f(a))$. 交为零：常数多项式在 $a$ 处为零只能是零.}
\ans{ex:vs:B12}{$f(x)=\frac{f(x)+f(-x)}2+\frac{f(x)-f(-x)}2$；$\mathrm e^x=\cosh x+\sinh x$.}
\ans{ex:vs:B13}{对 $k$ 归纳，反复使用维数公式：$\dim(U\cap W)\ge\dim U+\dim W-n$.}
\ans{ex:vs:B14}{这 $n-1$ 个向量线性无关（用坐标验证），$\dim W=n-1$. 补空间可取 $\spn(\vv\alpha_n)$（需验证 $\vv\alpha_n\notin W$：$W$ 中向量的坐标交错和为零）.}
\ans{ex:vs:B15}{$\vv x=x_1(1,\ldots,1)^\trans+(\vv x-x_1(1,\ldots,1)^\trans)$. 两种分解都是直和；与 $V_1$ 互补的子空间有无穷多个，第九章会说明例~\ref{eg:pricedecomp} 中的 $R$ 是其中``垂直''的那一个.}
\ans{ex:vs:B16}{每个矩阵唯一地拆成``对角线及以上''与``对角线以下''两部分. $\dim W_1=\frac{n(n+1)}2$.}
\ans{ex:vs:B17}{(1) 设 $W$ 的基为 $\vv\varepsilon_1,\ldots,\vv\varepsilon_k$（$k<n$），扩充为 $\mathbb{R}^n$ 的基 $\vv\varepsilon_1,\ldots,\vv\varepsilon_n$. ``取 $\vv x$ 在这组基下的第 $n$ 个坐标''是 $x_1,\ldots,x_n$ 的一次齐次式（解一个线性方程组即可写出），它在 $W$ 上恒为零，又在 $\vv\varepsilon_n$ 处为 $1$，系数不全为零. (2) 条件是 $a_1+a_2t+\cdots+a_nt^{n-1}=0$，这是一个次数不超过 $n-1$ 的非零多项式，至多 $n-1$ 个根. (3) 例如 $t=2$：$\vv v(2)=(1,2,4)^\trans$ 不满足三个方程中的任何一个.}

\ans{ex:vs:B18}{对 $k$ 归纳. 已选好线性无关的 $\vv\varepsilon_1,\ldots,\vv\varepsilon_{k-1}$（$k\le n$）后，$U=\spn(\vv\varepsilon_1,\ldots,\vv\varepsilon_{k-1})$ 的维数 $k-1<n$，也是真子空间. 由覆盖定理取 $\vv\varepsilon_k\notin W_1\cup\cdots\cup W_s\cup U$；由引理~\ref{lem:addone}，$\vv\varepsilon_1,\ldots,\vv\varepsilon_k$ 仍线性无关. 选满 $n$ 个即得一组基.}

\section{第三章　矩阵}

\subsection*{A 组答案}

\ans{ex:mat:A1}{$\vv{AB}=\begin{pmatrix}5&9\\-2&4\end{pmatrix}$，$\vv{BA}=\begin{pmatrix}2&1&5\\-1&1&-2\\1&11&6\end{pmatrix}$，$\vv A^\trans\vv B^\trans=(\vv{BA})^\trans$. 由此也看出 $(\vv{AB})^\trans=\vv B^\trans\vv A^\trans\ne\vv A^\trans\vv B^\trans$（后者甚至与前者不同型）.}

\ans{ex:mat:A2}{$2(1,3)^\trans-(2,-1)^\trans+(0,4)^\trans=(0,11)^\trans$. $\vv e_i^\trans\vv A\vv e_j=a_{ij}$；$\vv A\vv e_j$ 是 $\vv A$ 的第 $j$ 列，$\vv e_i^\trans\vv A$ 是第 $i$ 行.}

\ans{ex:mat:A3}{$\begin{pmatrix}1&n\\0&1\end{pmatrix}$；$\begin{pmatrix}\lambda^n&n\lambda^{n-1}\\0&\lambda^n\end{pmatrix}$；$\begin{pmatrix}\cos n\theta&-\sin n\theta\\ \sin n\theta&\cos n\theta\end{pmatrix}$；第四个矩阵记为 $\vv N$，$\vv N^2=\vv E_{13}$，$n\ge3$ 时 $\vv N^n=\vv O$. 均可用归纳法证明.}

\ans{ex:mat:A4}{$(\vv A+\vv B)^2=\begin{pmatrix}10&8\\12&10\end{pmatrix}$，$\vv A^2+2\vv{AB}+\vv B^2=\begin{pmatrix}16&8\\12&4\end{pmatrix}$，不相等；$(\vv A+\vv B)(\vv A-\vv B)=\begin{pmatrix}-6&4\\-6&6\end{pmatrix}\ne\vv A^2-\vv B^2=\begin{pmatrix}0&4\\-6&0\end{pmatrix}$. 原因是 $\vv{AB}\ne\vv{BA}$.}

\ans{ex:mat:A5}{依次为
$\begin{pmatrix}1&3&-2\\-\frac32&-3&\frac52\\1&1&-1\end{pmatrix}$，
$\begin{pmatrix}1&-1&0\\0&1&-1\\0&0&1\end{pmatrix}$，
$\frac14\begin{pmatrix}3&-2&1\\-2&4&-2\\1&-2&3\end{pmatrix}$，
$\begin{pmatrix}0&0&0&\frac14\\0&0&\frac13&0\\0&\frac12&0&0\\1&0&0&0\end{pmatrix}$.}

\ans{ex:mat:A6}{(1) $\vv X=\begin{pmatrix}2&-23\\0&8\end{pmatrix}$. (2) $\vv X=\begin{pmatrix}1&-1\\-2&4\end{pmatrix}$（注意 $\vv X=\vv B\vv A^{-1}$，逆矩阵乘在右边）.}

\ans{ex:mat:A7}{$r_2-r_1$，$r_3-ar_1$，再 $r_3+r_2$，得阶梯形的后两行为 $(0,a-1,1-a)$ 与 $\big(0,0,(1-a)(a+2)\big)$. $a\ne1$ 且 $a\ne-2$ 时秩为 $3$；$a=1$ 时后两行全为零，秩为 $1$；$a=-2$ 时第二行非零、第三行为零，秩为 $2$.（学完第五章后，可用行列式 $-(a-1)^2(a+2)$ 验证.）}

\ans{ex:mat:A8}{秩为 $3$. 简化阶梯形的主元在第 $1,3,5$ 列，故第 $1,3,5$ 列构成极大无关组；第 $2$ 列 $=2\times$第 $1$ 列，第 $4$ 列 $=$ 第 $1$ 列 $+$ 第 $3$ 列.}

\ans{ex:mat:A9}{$\vv E(1,3)=\begin{pmatrix}0&0&1\\0&1&0\\1&0&0\end{pmatrix}$，逆为自身；$\vv E(2(5))=\diag(1,5,1)$，逆为 $\diag(1,\frac15,1)$；$\vv E(3,1(-2))=\begin{pmatrix}1&0&0\\0&1&0\\-2&0&1\end{pmatrix}$，逆为 $\vv E(3,1(2))$. $\vv E(3,1(-2))\vv A$ 把 $\vv A$ 第 $1$ 行的 $-2$ 倍加到第 $3$ 行；$\vv A\vv E(3,1(-2))$ 把 $\vv A$ 第 $3$ 列的 $-2$ 倍加到第 $1$ 列.}

\ans{ex:mat:A10}{$\begin{pmatrix}2&1\\4&3\end{pmatrix}=\vv E(2,1(2))\,\vv E(1,2(1))\,\vv E(1(2))=\begin{pmatrix}1&0\\2&1\end{pmatrix}\begin{pmatrix}1&1\\0&1\end{pmatrix}\begin{pmatrix}2&0\\0&1\end{pmatrix}$. 写法不唯一.}

\ans{ex:mat:A11}{$(\vv{AB})_{ij}=\sum_ka_{ik}b_{kj}$，$i>j$ 时每一项中 $k<i$ 使 $a_{ik}=0$ 或 $k\ge i>j$ 使 $b_{kj}=0$；$i=j$ 时只剩 $a_{ii}b_{ii}$. 可逆上三角矩阵的逆也是上三角矩阵：对 $(\vv A\mid\vv I)$ 作消元时只需用下面的行消去上面的行.}

\ans{ex:mat:A12}{(1) 用 $(\vv{AB})^\trans=\vv B^\trans\vv A^\trans$ 逐一验证. (2) $(\vv{AB})^\trans=\vv B^\trans\vv A^\trans=\vv{BA}$. (3) $(\vv C^\trans\vv C)^\trans=\vv C^\trans\vv C$.}

\ans{ex:mat:A13}{(1) $\vv x\vv y^\trans$ 的 $(i,i)$ 元为 $x_iy_i$，求和即得. (2) $\tr(\vv A\vv E_{ji})=\vv e_i^\trans\vv A\vv e_j=a_{ij}$，故每个 $a_{ij}=0$. (3) 两边取迹，由命题~\ref{prop:trace}，左边为 $0$，右边为 $n$.}

\ans{ex:mat:A14}{$\vv A(\vv A-\vv I)=2\vv I$，故 $\vv A^{-1}=\frac12(\vv A-\vv I)$；$(\vv A+2\vv I)(\vv A-3\vv I)=\vv A^2-\vv A-6\vv I=-4\vv I$，故 $(\vv A+2\vv I)^{-1}=-\frac14(\vv A-3\vv I)$. $\vv A+\vv I$ 不一定可逆：$\vv A=-\vv I$ 满足条件，$\vv A+\vv I=\vv O$.}

\ans{ex:mat:A15}{左乘 $\vv A^{-1}$. 反例：$\vv A=\begin{pmatrix}1&0\\0&0\end{pmatrix}$，$\vv B=\vv O$，$\vv C=\begin{pmatrix}0&0\\0&1\end{pmatrix}$.}

\ans{ex:mat:A16}{$\begin{pmatrix}\vv A^{-1}&\vv O\\ \vv O&\vv B^{-1}\end{pmatrix}$，$\begin{pmatrix}\vv O&\vv B^{-1}\\ \vv A^{-1}&\vv O\end{pmatrix}$. 所求逆矩阵为 $\begin{pmatrix}0&0&\frac12&0\\0&0&0&\frac15\\3&-2&0&0\\-1&1&0&0\end{pmatrix}$.}

\ans{ex:mat:A17}{记 $\vv B=\begin{pmatrix}2&1\\1&3\end{pmatrix}$，矩阵为 $\begin{pmatrix}\vv I&\vv B\\ \vv O&\vv I\end{pmatrix}$，其平方为 $\begin{pmatrix}\vv I&2\vv B\\ \vv O&\vv I\end{pmatrix}=\begin{pmatrix}1&0&4&2\\0&1&2&6\\0&0&1&0\\0&0&0&1\end{pmatrix}$，$n$ 次方为 $\begin{pmatrix}\vv I&n\vv B\\ \vv O&\vv I\end{pmatrix}$.}

\ans{ex:mat:A18}{$\vv Q\vv C=\begin{pmatrix}5600&8300\\6750&9800\end{pmatrix}$：第 $i$ 行是第 $i$ 个工厂每天的总成本与总销售额. 每天利润分别为 $2700$ 元与 $3050$ 元.}

\ans{ex:mat:A19}{$\vv A=\begin{pmatrix}0&1&1&0\\1&0&1&0\\1&1&0&1\\0&0&1&0\end{pmatrix}$，$\vv A^2=\begin{pmatrix}2&1&1&1\\1&2&1&1\\1&1&3&0\\1&1&0&1\end{pmatrix}$，$\vv A^3=\begin{pmatrix}2&3&4&1\\3&2&4&1\\4&4&2&3\\1&1&3&0\end{pmatrix}$. $(\vv A^2)_{ii}$ 等于城市 $i$ 的航线数；从 $1$ 到 $4$ 恰乘三次航班有 $(\vv A^3)_{14}=1$ 种走法（$1\to2\to3\to4$）.}

\ans{ex:mat:A20}{$\vv x^\trans\vv A\vv x$ 是一个数，等于它的转置 $\vv x^\trans\vv A^\trans\vv x=-\vv x^\trans\vv A\vv x$，故为 $0$. $(\vv A^2)^\trans=(\vv A^\trans)^2=(-\vv A)^2=\vv A^2$.}

\ans{ex:mat:A21}{$\vv P$ 是初等矩阵之积，$\vv{PA}$ 与 $\vv A$ 行等价，因而有相同的简化阶梯形（简化阶梯形的唯一性见第一章章末习题~\ref{ex:elim:C9}）.}

\ans{ex:mat:A22}{存在性由提示中的等式. 唯一性：若 $\vv S_1+\vv K_1=\vv S_2+\vv K_2$（$\vv S_i$ 对称，$\vv K_i$ 反对称），则 $\vv S_1-\vv S_2=\vv K_2-\vv K_1$ 既对称又反对称，只能是零矩阵.}

\subsection*{B 组提示}

\ans{ex:mat:B1}{设 $\vv A=\begin{pmatrix}a&b\\c&d\end{pmatrix}$，用 Cayley–Hamilton 型恒等式 $\vv A^2=(a+d)\vv A-(ad-bc)\vv I$（直接验证）. (1) $\vv A^2=\vv O$ 当且仅当 $a+d=0$ 且 $ad-bc=0$. (2) $\vv A=\pm\vv I$，或 $a+d=0$ 且 $ad-bc=-1$. (3) $\vv A=\vv O$，$\vv A=\vv I$，或 $a+d=1$ 且 $ad-bc=0$.}
\ans{ex:mat:B2}{比较 $\vv{AD}$ 与 $\vv{DA}$ 的 $(i,j)$ 元：$a_ix_{ij}=x_{ij}a_j$.}
\ans{ex:mat:B3}{由 $\vv{AE}_{ij}=\vv E_{ij}\vv A$ 比较元素，得 $\vv A$ 的非对角元为零，且 $a_{ii}=a_{jj}$.}
\ans{ex:mat:B4}{用二项式定理（可交换时成立）：$(\vv A+\vv B)^{p+q}$ 的每一项中 $\vv A$ 的次数 $\ge p$ 或 $\vv B$ 的次数 $\ge q$. 反例：$\vv A=\vv E_{12}$，$\vv B=\vv E_{21}$.}
\ans{ex:mat:B5}{等式：分块对角矩阵的行（列）向量组的极大无关组由两块各自的极大无关组拼成. 不等式：用初等变换把 $\vv A$、$\vv B$ 分别化为标准形，看非零子矩阵. 例子：$\vv A=\vv B=(0)$，$\vv C=(1)$.}
\ans{ex:mat:B6}{把右边与 $\vv I+\vv u\vv v^\trans$ 相乘，利用 $\vv v^\trans\vv u$ 是一个数.}
\ans{ex:mat:B7}{直接验证 $(\vv I_n+\vv{BA})\big(\vv I_n-\vv B(\vv I_m+\vv{AB})^{-1}\vv A\big)=\vv I_n$，展开后把 $\vv B\vv A\vv B$ 写成 $\vv B(\vv{AB})$.}
\ans{ex:mat:B8}{各列都是某个非零列 $\vv u$ 的倍数，倍数组成 $\vv v$. $\vv A^2=\vv u(\vv v^\trans\vv u)\vv v^\trans$，而 $\vv v^\trans\vv u=\tr\vv A$. 于是 $\vv A^k=(\tr\vv A)^{k-1}\vv A$.}
\ans{ex:mat:B9}{用相抵标准形 $\vv A=\vv P\diag(\vv I_r,\vv O)\vv Q$，把 $\vv I_r$ 拆成 $r$ 个 $\vv E_{ii}$. 后半句：秩一矩阵之和的秩不超过个数.}
\ans{ex:mat:B10}{取 $\vv B$ 为 $\vv A$ 列向量组的一个极大无关组排成的矩阵，$\vv C$ 的各列是 $\vv A$ 各列在这组向量下的坐标. 对第 A8 题，$\vv B$ 取第 $1,3,5$ 列，$\vv C$ 是简化阶梯形的前三行.}
\ans{ex:mat:B11}{全 $1$ 矩阵 $\vv J$ 满足 $\vv J^2=3\vv J$，故 $\vv J^n=3^{n-1}\vv J$. 第二个：$\vv N^3=\vv O$，$(2\vv I+\vv N)^n=2^n\vv I+n2^{n-1}\vv N+\binom n22^{n-2}\vv N^2$.}
\ans{ex:mat:B12}{对 $k$ 归纳，计算 $\vv N^k$ 的 $(i,j)$ 元 $\sum_l(\vv N^{k-1})_{il}n_{lj}$ 中哪些项可能非零.}
\ans{ex:mat:B13}{$\tr(\vv A^\trans\vv A)=\sum_{i,j}a_{ij}^2$. 复矩阵不成立：$\vv A=\begin{pmatrix}1&0\\\mathrm i&0\end{pmatrix}$ 满足 $\vv A^\trans\vv A=\vv O$；复矩阵应改用 $\bar{\vv A}^\trans\vv A$.}
\ans{ex:mat:B14}{置换矩阵就是对单位矩阵的行作排列得到的矩阵. $\vv{PA}$ 对 $\vv A$ 的行作相应排列，$\vv{AP}$ 对列作排列.}
\ans{ex:mat:B15}{秩为 $2$. 先用行变换化为阶梯形（记录 $\vv P$），再用列变换（记录 $\vv Q$）. 答案不唯一，可用 $\vv{PAQ}$ 直接检验.}
\ans{ex:mat:B16}{$(\vv I+\vv A)(\vv I-\vv A+\vv A^2-\cdots+(-1)^{k-1}\vv A^{k-1})=\vv I$. 若 $\vv A^2=2\vv A$，则 $(\vv I-\vv A)^2=\vv I$，可逆且逆为自身.}

\section{第四章　线性映射}

\subsection*{A 组答案}

\ans{ex:lmap:A1}{(1) 线性，核为 $\{\zero\}$. (2) 不线性：$\zero\mapsto(1,0)^\trans$. (3) 不线性：$(-1,0)^\trans$ 与 $(1,0)^\trans$ 之和的像不等于像之和. (4) 线性，核为 $\{0\}$. (5) 线性，核为常数多项式全体. (6) 线性，核为 $\{\vv O\}$. (7) 不线性：$\det(2\vv X)=4\det\vv X$.}

\ans{ex:lmap:A2}{第一个：核的基 $(-2,1,0)^\trans$，$(1,0,1)^\trans$；像的基 $(1,2)^\trans$；$3=2+1$. 第二个：核的基 $(-2,1,1,0)^\trans$，$(-1,-1,0,1)^\trans$；像的基 $(1,0,1)^\trans$，$(0,1,1)^\trans$；$4=2+2$.}

\ans{ex:lmap:A3}{$\begin{pmatrix}1&-2&1\\0&3&4\end{pmatrix}$.}

\ans{ex:lmap:A4}{(1) $\begin{pmatrix}0&-1\\1&0\end{pmatrix}$. (2) $\begin{pmatrix}0&1\\1&0\end{pmatrix}$，平方为 $\vv I$：反射两次回到原处. (3) $\frac12\begin{pmatrix}1&1\\1&1\end{pmatrix}$，平方等于自身：投影两次与投影一次相同. (4) $\begin{pmatrix}1&2\\0&1\end{pmatrix}$.}

\ans{ex:lmap:A5}{两组基下的矩阵都是 $\begin{pmatrix}0&1&0\\0&0&2\\0&0&0\end{pmatrix}$：$\mathscr D(x-1)=1$，$\mathscr D(x-1)^2=2(x-1)$.}

\ans{ex:lmap:A6}{$\begin{pmatrix}1&0&0&0\\0&0&1&0\\0&1&0&0\\0&0&0&1\end{pmatrix}$，它的平方是 $\vv I_4$.}

\ans{ex:lmap:A7}{过渡矩阵 $\vv P=\begin{pmatrix}-1&-1&-4\\1&0&6\\1&3&1\end{pmatrix}$. $\vv\alpha$ 在 $\vv\varepsilon$ 下的坐标为 $(-1,2,2)^\trans$，在 $\vv\eta$ 下的坐标为 $(-16,5,3)^\trans$. 可用 $\vv x=\vv P\vv x'$ 检验.}

\ans{ex:lmap:A8}{$\diag(3,1)$. 过渡矩阵 $\vv P=\begin{pmatrix}1&1\\1&-1\end{pmatrix}$，也可以直接看出 $\mathscr T\vv\eta_1=3\vv\eta_1$，$\mathscr T\vv\eta_2=\vv\eta_2$.}

\ans{ex:lmap:A9}{$\mathscr S\mathscr T(a\vv\alpha+b\vv\beta)=\mathscr S(a\mathscr T\vv\alpha+b\mathscr T\vv\beta)=a\mathscr{ST}\vv\alpha+b\mathscr{ST}\vv\beta$. 逆映射：设 $\vv\alpha=\mathscr T\vv\alpha'$，$\vv\beta=\mathscr T\vv\beta'$，则 $\mathscr T^{-1}(a\vv\alpha+b\vv\beta)=\mathscr T^{-1}\mathscr T(a\vv\alpha'+b\vv\beta')=a\mathscr T^{-1}\vv\alpha+b\mathscr T^{-1}\vv\beta$.}

\ans{ex:lmap:A10}{若 $\sum k_i\mathscr T\vv\alpha_i=\zero$，则 $\sum k_i\vv\alpha_i\in\ker\mathscr T=\{\zero\}$. 反之，若 $\mathscr T\vv\alpha=\zero$ 而 $\vv\alpha\ne\zero$，则线性无关组 $\{\vv\alpha\}$ 被映为相关组 $\{\zero\}$.}

\ans{ex:lmap:A11}{(1) 存在：前两个向量是基，$\mathscr T(x,y)^\trans=\big(\frac{x+y}2,\frac{x-y}2\big)^\trans$，它恰好满足第三个条件. (2) 不存在：$(2,0)^\trans=(1,1)^\trans+(1,-1)^\trans$，线性性要求 $\mathscr T(2,0)^\trans=(1,1)^\trans$.}

\ans{ex:lmap:A12}{(1) $5-3=2$. (2) 都不存在：像的维数 $\le3<5$；核的维数 $\ge5-3=2>0$.}

\ans{ex:lmap:A13}{$\mathscr T\vv\alpha=\zero\Rightarrow\mathscr{ST}\vv\alpha=\zero$；$\mathscr{ST}\vv\alpha=\mathscr S(\mathscr T\vv\alpha)$. 由前者 $\rk(\mathscr{ST})=\dim U-\dim\ker(\mathscr{ST})\le\dim U-\dim\ker\mathscr T=\rk\mathscr T$，由后者 $\rk(\mathscr{ST})\le\rk\mathscr S$.}

\ans{ex:lmap:A14}{$\tr$ 是满射（$\tr(c\vv E_{11})=c$），由秩—零化度定理 $\dim\ker\tr=n^2-1$.}

\ans{ex:lmap:A15}{$\tr(\vv P^{-1}\vv{AP})=\tr(\vv{AP}\vv P^{-1})=\tr\vv A$. 两矩阵迹分别为 $4$ 与 $3$，不相似.}

\ans{ex:lmap:A16}{$\vv P^{-1}\vv I\vv P=\vv I$. 由于 $\begin{pmatrix}1&1\\0&1\end{pmatrix}\ne\vv I$，它与 $\vv I$ 不相似.}

\ans{ex:lmap:A17}{矩阵为 $\begin{pmatrix}1&0&0\\1&1&1\\1&2&4\end{pmatrix}$. 依次做 $r_2-r_1$，$r_3-r_1$，$r_3-2r_2$，得 $\begin{pmatrix}1&0&0\\0&1&1\\0&0&2\end{pmatrix}$，有 $3$ 个主元，由定理~\ref{thm:invertible} 矩阵可逆，故 $\mathscr E$ 是同构.（第五章会看到，它的行列式是 Vandermonde 型的.）}

\ans{ex:lmap:A18}{$\mathscr T(x_1,x_2,x_3)^\trans=(x_1+x_3,\ x_1+x_2+2x_3,\ x_2+x_3)^\trans$. $\ker\mathscr T=\spn\big((1,1,-1)^\trans\big)$，$\imm\mathscr T=\spn\big((1,1,0)^\trans,(0,1,1)^\trans\big)$.}

\ans{ex:lmap:A19}{$\vv A_2(\vv A_1\vv x+\vv b_1)+\vv b_2=(\vv A_2\vv A_1)\vv x+(\vv A_2\vv b_1+\vv b_2)$. 仿射变换线性当且仅当把 $\zero$ 映为 $\zero$，即 $\vv b=\zero$.}

\ans{ex:lmap:A20}{第一列是 $\vv e_1$ 转 $\theta$ 后的位置 $(\cos\theta,\sin\theta)^\trans$，第二列同理. 复合``先转 $\beta$ 再转 $\alpha$''的矩阵是 $\vv R(\alpha)\vv R(\beta)$（定理~\ref{thm:composition}），比较 $(1,1)$ 元与 $(2,1)$ 元，得 $\cos(\alpha+\beta)=\cos\alpha\cos\beta-\sin\alpha\sin\beta$，$\sin(\alpha+\beta)=\sin\alpha\cos\beta+\cos\alpha\sin\beta$.}

\subsection*{B 组提示}

\ans{ex:lmap:B1}{取定基后 $\mathscr T\mapsto$ 矩阵是 $L(V,W)$ 到 $M_{m\times n}(\mathbb{F})$ 的同构.}
\ans{ex:lmap:B2}{$\mathscr T^2=\mathscr O$ 即 $\mathscr T(\mathscr T\vv\alpha)=\zero$. 取 $\imm\mathscr T$ 的基 $\vv\beta_1,\ldots,\vv\beta_r$，$\vv\beta_i=\mathscr T\vv\alpha_i$，证明 $\vv\beta_1,\vv\alpha_1,\ldots,\vv\beta_r,\vv\alpha_r$ 线性无关，再用 $\ker\mathscr T$ 中的向量补成基.}
\ans{ex:lmap:B3}{对 $\mathscr T|_{\ker(\mathscr{ST})}$ 用秩—零化度定理：它的像含于 $\ker\mathscr S$，核含于 $\ker\mathscr T$.}
\ans{ex:lmap:B4}{$\mathscr{ST}=\mathscr I$ 说明 $\mathscr T$ 是单射，有限维时由推论~\ref{cor:nullity} 它是双射，于是 $\mathscr S=\mathscr T^{-1}$.}
\ans{ex:lmap:B5}{$\mathscr{JD}(1)=0$. 不矛盾：B4 依赖有限维，$\mathbb{R}[x]$ 是无限维的.}
\ans{ex:lmap:B6}{$\tr(\vv{AX}-\vv{XA})=0$. 后半句：$\vv A^k$ 与 $\vv A$ 可交换.}
\ans{ex:lmap:B7}{逐个计算 $\vv A\vv E_{11}$ 等并取坐标. $\vv X\mapsto\vv{AX}$ 的矩阵为 $\begin{pmatrix}\vv A&\vv O\\ \vv O&\vv A\end{pmatrix}$（在给定的基序下）；$\vv X\mapsto\vv{XA}$ 的矩阵为 $\begin{pmatrix}a\vv I&c\vv I\\ b\vv I&d\vv I\end{pmatrix}$.}
\ans{ex:lmap:B8}{证明 $V=\imm\mathscr P\oplus\ker\mathscr P$：$\vv\alpha=\mathscr P\vv\alpha+(\vv\alpha-\mathscr P\vv\alpha)$，且 $\mathscr P$ 在 $\imm\mathscr P$ 上是恒等. 合并两者的基.}
\ans{ex:lmap:B9}{按提示分解，并验证两部分分别在 $\ker(\mathscr T\mp\mathscr I)$ 中；交为零是因为 $\mathscr T\vv\alpha=\vv\alpha=-\vv\alpha$ 推出 $\vv\alpha=\zero$.}
\ans{ex:lmap:B10}{对 $\mathscr T|_U$ 用秩—零化度定理.}
\ans{ex:lmap:B11}{在基 $\mathrm e^x\cos x$，$\mathrm e^x\sin x$ 下 $\mathscr D$ 的矩阵为 $\begin{pmatrix}1&1\\-1&1\end{pmatrix}$，逆为 $\frac12\begin{pmatrix}1&-1\\1&1\end{pmatrix}$，故 $\int\mathrm e^x\sin x\,\mathrm dx=\frac12\mathrm e^x(\sin x-\cos x)+C$.}
\ans{ex:lmap:B12}{平移的矩阵为 $\begin{pmatrix}1&0&a\\0&1&b\\0&0&1\end{pmatrix}$. 两种顺序的矩阵不相等：先转后移把原点送到 $(1,2)$，先移后转把原点送到 $(-2,1)$.}
\ans{ex:lmap:B13}{$\vv P^{-1}\vv A^k\vv P=(\vv P^{-1}\vv{AP})^k$；转置时用 $\vv Q=(\vv P^{-1})^\trans$.}
\ans{ex:lmap:B14}{两者秩都为 $1$，故相抵. 若相似，则 $\vv P^{-1}\vv E_{11}\vv P$ 是幂等的，而 $\vv E_{12}^2=\vv O\ne\vv E_{12}$. 相抵允许在定义域与陪域中各自换基，相似必须在同一个空间中用同一组基.}
\ans{ex:lmap:B15}{设 $\sum c_i\mathscr T^i\vv\alpha=\zero$，依次施以 $\mathscr T^{k-1},\mathscr T^{k-2},\ldots$. 所求矩阵是主对角线之上一条斜线全为 $1$、其余为 $0$ 的矩阵.}

\section{第五章　行列式}

\subsection*{A 组答案}

\ans{ex:det:A1}{$\tau(315426)=5$，奇排列；$\tau(42531)=7$，奇排列；第三个排列的逆序数为 $0+1+2+\cdots+(n-1)=\frac{n(n-1)}2$，$n\equiv0,1\pmod4$ 时为偶排列，$n\equiv2,3\pmod4$ 时为奇排列.}

\ans{ex:det:A2}{(1) $\tau(3421)=5$，带负号. (2) 按行号排好后列号排列为 $i\,3\,5\,4\,j$，$\{i,j\}=\{1,2\}$. $\tau(13542)=4$，故 $i=1$，$j=2$.}

\ans{ex:det:A3}{$-41$，$-16$，$48$. 第三个：把各列加到第一列后提出 $6$，再化为三角形得 $6\cdot2^3=48$.}

\ans{ex:det:A4}{$6$.}

\ans{ex:det:A5}{$(-1)^{\frac{n(n-1)}2}a_1a_2\cdots a_n$：只有一项非零，它的列号排列是 $n(n-1)\cdots1$.}

\ans{ex:det:A6}{第一列减去第二列、再加上第三列（或按多重线性拆成 $8$ 个行列式，其中有两列相同的都为零），得 $2\begin{vmatrix}a&b&c\\a_1&b_1&c_1\\a_2&b_2&c_2\end{vmatrix}$.}

\ans{ex:det:A7}{第一个行列式等于 $(x-1)^2(x+2)$，根为 $x=1$（二重）与 $x=-2$；第二个是 Vandermonde 行列式 $(2-1)(x-1)(x-2)$，根为 $x=1,2$.}

\ans{ex:det:A8}{$24$，$-3$，$\frac13$，$9$，$9$. 最后一个：$\vv A^*=3\vv A^{-1}$，故 $3\vv A^{-1}-2\vv A^*=-3\vv A^{-1}$，行列式为 $(-3)^3\cdot\frac13=-9$.}

\ans{ex:det:A9}{(1) $(\det\vv A)^2=\det(\vv A^\trans\vv A)=1$. (2) $(\det\vv A)^2=1$. (3) 若 $\det\vv A\ne0$，则 $\vv A$ 可逆，由 $\vv A^2=\vv A$ 得 $\vv A=\vv I$，矛盾.}

\ans{ex:det:A10}{$\begin{pmatrix}d&-b\\-c&a\end{pmatrix}$，$ad-bc\ne0$ 时逆为 $\frac1{ad-bc}\begin{pmatrix}d&-b\\-c&a\end{pmatrix}$. 第二个矩阵行列式为 $1$，伴随矩阵就是逆矩阵：$\begin{pmatrix}-24&18&5\\20&-15&-4\\-5&4&1\end{pmatrix}$.}

\ans{ex:det:A11}{$D=17$，$D_1=17$，$D_2=-17$，$D_3=34$，解为 $(1,-1,2)$.}

\ans{ex:det:A12}{面积为 $\frac12\left|\det\begin{pmatrix}3&1\\1&4\end{pmatrix}\right|=\frac{11}2$. 三点共线：$(2,3)^\trans$ 与 $(6,9)^\trans$ 张成的平行四边形面积为 $0$.}

\ans{ex:det:A13}{平行六面体体积为 $7$，四面体体积为 $\frac76$.}

\ans{ex:det:A14}{系数行列式为 $(5-\lambda)(6-\lambda)(4-\lambda)-4(4-\lambda)-4(6-\lambda)=-(\lambda-2)(\lambda-5)(\lambda-8)$，故 $\lambda=2,5,8$.}

\ans{ex:det:A15}{把第一行换成 $(1,1,1,1)$ 后按第一行展开即得所求之和，计算得 $16$.}

\ans{ex:det:A16}{点 $(x,y)$ 在直线上当且仅当第一行是后两行的``仿射组合''，即三行线性相关. 也可以按第一行展开，得到一个关于 $x,y$ 的一次方程，并验证两点满足它.}

\ans{ex:det:A17}{按定义，非零项必须在左上角 $\vv A$ 中取 $m$ 个元素、在右下角 $\vv B$ 中取 $n$ 个，符号相乘. 第二个等于 $(-1)^{mn}\det\vv A\det\vv B$：把后 $n$ 列逐一移到前面，共 $mn$ 次相邻列交换.}

\ans{ex:det:A18}{$5$（一般地，这类 $n$ 阶行列式等于 $n+1$，见 B1）.}

\ans{ex:det:A19}{$\vv A\vv A^*=(\det\vv A)\vv I$ 两边取行列式得 $\det\vv A^*=(\det\vv A)^{n-1}\ne0$；再由 $\vv A^*=(\det\vv A)\vv A^{-1}$ 得 $(\vv A^*)^{-1}=\frac1{\det\vv A}\vv A$；对 $\vv A^{-1}$ 用同一公式得 $(\vv A^{-1})^*=\det(\vv A^{-1})\vv A=\frac1{\det\vv A}\vv A$.}

\ans{ex:det:A20}{行列式为 $0$，线性相关（$(7,8,9)^\trans=2(4,5,6)^\trans-(1,2,3)^\trans$）；改成 $10$ 后行列式为 $-3$，线性无关.}

\ans{ex:det:A21}{$\det\vv A=\det\vv A^\trans=\det(-\vv A)=(-1)^n\det\vv A=-\det\vv A$，故 $\det\vv A=0$.}

\subsection*{B 组提示}

\ans{ex:det:B1}{由 $D_n=2D_{n-1}-D_{n-2}$，$D_1=2$，$D_2=3$，得 $D_n-D_{n-1}=1$，$D_n=n+1$.}
\ans{ex:det:B2}{第一个：各列加到第一列，提出 $a+(n-1)b$，各行减去第一行，得 $[a+(n-1)b](a-b)^{n-1}$. 第二个：先证递推 $D_n=(x-a)D_{n-1}+a(x+a)^{n-1}$（或利用转置后 $a$ 变 $-a$），得 $\frac12\big[(x+a)^n+(x-a)^n\big]$.}
\ans{ex:det:B3}{第 $i$ 列乘 $\frac1{a_i}$ 后从第一列减去，化为上三角形：$a_1\cdots a_n\big(a_0-\sum_{i=1}^n\frac1{a_i}\big)$.}
\ans{ex:det:B4}{计算 $\begin{pmatrix}\vv I&\zero\\ -\vv v^\trans&1\end{pmatrix}\begin{pmatrix}\vv I+\vv u\vv v^\trans&\vv u\\ \zero^\trans&1\end{pmatrix}\begin{pmatrix}\vv I&\zero\\ \vv v^\trans&1\end{pmatrix}$. B2 的第一个行列式即 $\det\big((a-b)\vv I+b\vv 1\vv 1^\trans\big)$.}
\ans{ex:det:B5}{把 $\vv{AB}$ 的两列按第三章``列是线性组合''的看法展开，用多重线性拆成若干个行列式，其中列号重复的为零. $\vv{BA}$ 的秩 $\le2<3$.}
\ans{ex:det:B6}{必要性：$\det\vv A\det\vv A^{-1}=1$，两者都是整数. 充分性：$\vv A^{-1}=\frac1{\det\vv A}\vv A^*$，$\vv A^*$ 是整数矩阵.}
\ans{ex:det:B7}{可逆时用 $\vv A^*=(\det\vv A)\vv A^{-1}$ 两次. 不可逆时 $\rk\vv A^*\le1$，$n\ge3$ 时 $(\vv A^*)^*=\vv O$；$n=2$ 时直接验证.}
\ans{ex:det:B8}{对分块矩阵先把第二块行加到第一块行，再把第一块列从第二块列减去，化为 $\begin{pmatrix}\vv A+\vv B&\vv O\\ \vv B&\vv A-\vv B\end{pmatrix}$.}
\ans{ex:det:B9}{左乘 $\begin{pmatrix}\vv I&\vv O\\ -\vv{CA}^{-1}&\vv I\end{pmatrix}$. 可交换时 $\det\vv A\det(\vv D-\vv{CA}^{-1}\vv B)=\det(\vv{AD}-\vv{ACA}^{-1}\vv B)=\det(\vv{AD}-\vv{CB})$.}
\ans{ex:det:B10}{两边取行列式：$\det\vv F^n=(\det\vv F)^n=(-1)^n$.}
\ans{ex:det:B11}{按第一行展开得 $D_n=(1+x^2)D_{n-1}-x^2D_{n-2}$，即 $D_n-D_{n-1}=x^2(D_{n-1}-D_{n-2})$，故 $D_n=1+x^2+\cdots+x^{2n}$.}
\ans{ex:det:B12}{把第二、三列分别乘 $\omega^k$，$\omega^{2k}$ 加到第一列，提出公因式；也可直接展开验证第二个等号.}
\ans{ex:det:B13}{从最后一行起，每一行减去上一行，化为对角线全为 $1$ 的上三角形.}
\ans{ex:det:B14}{圆（或直线）的方程可写成 $a(x^2+y^2)+bx+cy+d=0$，系数不全为零. 四点都满足它，即一个以 $(a,b,c,d)$ 为未知量的齐次方程组有非零解.}
\ans{ex:det:B15}{$\rk\vv A=n$ 时 $\vv A^*=(\det\vv A)\vv A^{-1}$ 可逆. $\rk\vv A\le n-2$ 时，由定理~\ref{thm:minorrank} 所有 $n-1$ 阶子式为零，$\vv A^*=\vv O$. $\rk\vv A=n-1$ 时有非零的 $n-1$ 阶子式，故 $\vv A^*\ne\vv O$；又 $\vv A\vv A^*=(\det\vv A)\vv I=\vv O$，$\vv A^*$ 的每一列都属于 $\ker\vv A$，而 $\dim\ker\vv A=n-(n-1)=1$（推论~\ref{cor:nullity}），故 $\rk\vv A^*=1$.}
\ans{ex:det:B16}{$x=1,2,3$ 时有两行相同；每行之和都是 $x+6$，把各列加到第一列可提出因子 $x+6$. $f$ 是四次多项式，首项系数为 $1$，故 $f(x)=(x+6)(x-1)(x-2)(x-3)$，根为 $1,2,3,-6$.}

\section{第六章　线性方程组解的结构}

\subsection*{A 组答案}

\ans{ex:sol:A1}{(1) 基础解系 $(2,5,7,0)^\trans$，$(3,4,0,7)^\trans$，通解为二者的任意线性组合. (2) 基础解系 $(-2,1,0,0)^\trans$，$(1,0,0,1)^\trans$.}

\ans{ex:sol:A2}{(1) $\vv x=\big(\tfrac54,-\tfrac14,0,0\big)^\trans+k_1(3,3,2,0)^\trans+k_2(-3,7,0,4)^\trans$. (2) $\vv x=\big(\tfrac12,0,0,0\big)^\trans+k_1(-1,2,0,0)^\trans+k_2(1,0,2,0)^\trans$.}

\ans{ex:sol:A3}{(1) $5-3=2$ 个. (2) 不一定，要看 $\rk(\vv A\mid\vv b)$ 是否为 $3$. (3) 一定有解（行满秩，推论~\ref{cor:sol-fullrank}），且有无穷多解.}

\ans{ex:sol:A4}{(1) $\rk\vv A=\rk(\vv A\mid\vv b)=2$，有解（且唯一）. (2) $\rk\vv A=1$，$\rk(\vv A\mid\vv b)=2$，无解.}

\ans{ex:sol:A5}{系数和为 $1$ 的 $2\vv\eta_1-\vv\eta_2$，$\frac12(\vv\eta_1+\vv\eta_2)$，$3\vv\eta_1-2\vv\eta_2$ 是 $\vv{Ax}=\vv b$ 的解；系数和为 $0$ 的 $\vv\eta_1-\vv\eta_2$ 是齐次方程组的解；$\vv\eta_1+\vv\eta_2$ 两者都不是（命题~\ref{prop:sol-affine}）.}

\ans{ex:sol:A6}{基础解系含 $3-2=1$ 个向量，$\vv\eta_1-\vv\eta_2=(0,4,4)^\trans$ 是非零解. 通解为 $(1,2,3)^\trans+k(0,1,1)^\trans$.}

\ans{ex:sol:A7}{例如 $x_1-x_2=1$，$x_1+x_3=3$. 做法：求与 $(1,1,-1)^\trans$ 垂直（即满足 $y_1+y_2-y_3=0$）的两个线性无关的系数向量，再代入特解求常数项. 答案不唯一.}

\ans{ex:sol:A8}{$\rk\vv A=2$（第三个法向量是前两个之和，且任意两个不平行），$\rk(\vv A\mid\vv b)=3$：三个平面两两相交成三条平行直线，没有公共点. 常数改为 $3$ 后 $\rk(\vv A\mid\vv b)=2$，三个平面交于一条直线.}

\ans{ex:sol:A9}{例如 $\vv A=\begin{pmatrix}1\\0\end{pmatrix}$，$\vv b=(0,1)^\trans$：$\vv{Ax}=\zero$ 只有零解，$\vv{Ax}=\vv b$ 无解. 若有解，两个解之差是齐次方程组的解，只能为零.}

\ans{ex:sol:A10}{$\vv X=\begin{pmatrix}3&2\\-2&-3\\1&3\end{pmatrix}$.}

\ans{ex:sol:A11}{$\vv X=\begin{pmatrix}a&b\\-a&-b\end{pmatrix}$（$a,b$ 任意），构成二维空间：$\vv X$ 的每一列都属于 $\ker\vv A=\spn\big((1,-1)^\trans\big)$.}

\ans{ex:sol:A12}{$\vv B$ 的非零列是 $\vv{Ax}=\zero$ 的非零解，故 $\det\vv A=7t+21=0$，$t=-3$. 这时 $\rk\vv A=2$，由 $\rk\vv A+\rk\vv B\le3$ 且 $\vv B\ne\vv O$ 得 $\rk\vv B=1$.}

\ans{ex:sol:A13}{通解 $(1,0,0)^\trans+s(-1,1,0)^\trans+t(-1,0,1)^\trans$. 取 $\vv e_1,\vv e_2,\vv e_3$：任一解 $(x_1,x_2,x_3)^\trans=x_1\vv e_1+x_2\vv e_2+x_3\vv e_3$，而 $x_1+x_2+x_3=1$.}

\ans{ex:sol:A14}{第一条直线：$(0,0,1)^\trans+t(1,1,-2)^\trans$. 第二条直线：$x_1-x_3=1$，$x_2=2$（答案不唯一）.}

\ans{ex:sol:A15}{方程组 $1+s=0$，$s=1+t$，$0=1-t$ 无解，两直线不相交（它们的方向也不平行，是异面直线）.}

\ans{ex:sol:A16}{方程 $x_1+x_3=30$，$x_1-x_2=0$，$x_2+x_3=30$（秩为 $2$）. 全部方案为 $x_1=x_2=s$，$x_3=30-s$. $x_1+x_3=30$ 与 $x_1-x_2=0$ 被唯一确定（它们是方程左边的线性组合）；$x_3$ 不被确定：它在基础解系 $(1,1,-1)^\trans$ 上的值不为零.}

\ans{ex:sol:A17}{$\vv L=\begin{pmatrix}1&0&0\\2&1&0\\1&2&1\end{pmatrix}$，$\vv U=\begin{pmatrix}2&1&0\\0&1&1\\0&0&1\end{pmatrix}$.}

\ans{ex:sol:A18}{$\rk(\vv A\mid\vv b)$ 不超过行数 $m=\rk\vv A$，故相容；解空间维数 $n-m>0$，有无穷多解.}

\ans{ex:sol:A19}{两个方向 $(1,-1,0)^\trans$，$(-1,1,2)^\trans$；求与二者都垂直的系数向量，即解齐次方程组，得 $(1,1,0)^\trans$. 平面方程为 $x+y=2$.}

\ans{ex:sol:A20}{解分别为 $(2,0)$ 与 $(1,1)$. 右端的相对变化约为 $\frac{0.0001}{2\sqrt2}\approx3.5\times10^{-5}$，解的相对变化为 $\frac{\sqrt2}2\approx0.71$，二者之比约为 $2\times10^4$.}

\subsection*{B 组提示}

\ans{ex:sol:B1}{把四个方程联立求解，得公共解 $k(-1,1,2,1)^\trans$.}
\ans{ex:sol:B2}{$\vv A^\trans\vv{Ax}=\zero\Rightarrow\vv x^\trans\vv A^\trans\vv A\vv x=0$，即 $\vv{Ax}$ 各分量的平方和为零. 两个齐次方程组同解，解空间维数相同，再用推论~\ref{cor:nullity}.}
\ans{ex:sol:B3}{$\vv A+\vv B$ 的列是 $(\vv A\mid\vv B)$ 的列的线性组合；后者的列由两个极大无关组合并表出.}
\ans{ex:sol:B4}{$(\vv A+\vv I)(\vv A-\vv I)=\vv O$ 给出 $\le n$（Sylvester 不等式的推论）；$(\vv A+\vv I)-(\vv A-\vv I)=2\vv I$ 给出 $\ge n$（B3）.}
\ans{ex:sol:B5}{由定理~\ref{thm:sylvester} 证明中的等式 $\rk\vv B=\rk(\vv{AB})+\dim(\ker\vv A\cap\imm\vv B)$.}
\ans{ex:sol:B6}{设 $k_0\vv\eta^*+\sum k_i(\vv\eta^*+\vv\xi_i)=\zero$，左乘 $\vv A$ 得 $(k_0+\sum k_i)\vv b=\zero$，再用 $\vv\xi_i$ 线性无关.}
\ans{ex:sol:B7}{$9$ 个未知数，$6$ 个方程的秩为 $5$，自由度为 $4$，至少再公开 $4$ 个单元格. 公开左上角的 $2\times2$ 块即可；但若公开的 $4$ 格中有一行的全部 $3$ 格，那一行有一格是多余的，全表仍不能确定.}
\ans{ex:sol:B8}{三行满足 $r_1+r_2=r_3$，故条件为 $b_1+b_2=b_3$.}
\ans{ex:sol:B9}{解空间是一维的，$(1,\ldots,1)^\trans$ 是非零解. 通解为 $k(1,\ldots,1)^\trans$.}
\ans{ex:sol:B10}{$\rk\vv A+\rk\vv B\le n$（B4 中的推论）且 $\ge\rk(\vv A+\vv B)=n$（B3）.}
\ans{ex:sol:B11}{同解时，$\vv B$ 的每一行在 $\vv A$ 的解空间上恒为零，由可确定性判据（定理~\ref{thm:sol-determinable} 的 (2)$\Rightarrow$(3)）它是 $\vv A$ 各行的线性组合.}
\ans{ex:sol:B12}{必要性：对 $\vv b=\vv e_1,\ldots,\vv e_m$ 分别取解，排成 $\vv C$. 充分性：$\vv{A}(\vv C\vv b)=\vv b$.}
\ans{ex:sol:B13}{$\vv A\vv A^*=(\det\vv A)\vv I=\vv O$，$\vv A^*$ 的每一列都在一维解空间中.}
\ans{ex:sol:B14}{$x_1-x_4$ 是第一、二个方程之差，被确定为 $-3$；后两个量都不被确定：找不到系数 $a,b,c$ 使 $a,b,c$ 倍的三个方程之和的左边恰为所求的量.}
\ans{ex:sol:B15}{对 $\vv B|_{\ker(\vv{AB})}$ 用秩—零化度定理：它的核是 $\ker\vv B$，像含于 $\ker\vv A$.}

\section{第七章　多项式}

\subsection*{A 组答案}

\ans{ex:poly:A1}{(1) 商 $x^3-x^2+1$，余式 $-x-2$. (2) 商 $2x^3+x^2+6x+7$，余式 $20$（$=f(2)$）.}

\ans{ex:poly:A2}{依次为 $1$，$1$，$100x-99$. 前两个用 $x^2\equiv1$，$x^2\equiv-1$；第三个设余式为 $ax+b$，由 $x^{100}-(ax+b)$ 以 $1$ 为二重根，在 $x=1$ 处函数值与导数都为零.}

\ans{ex:poly:A3}{$\omega^4=\omega$，$\omega^2=-1-\omega$，代入得 $(1-a)\omega+(b-a)=0$，而 $\omega$ 不是实数，故 $a=b=1$. 验证：$x^4+x^2+1=(x^2+x+1)(x^2-x+1)$.}

\ans{ex:poly:A4}{(1) $x+3$. (2) $x-1$（$x^3-2x^2+x=x(x-1)^2$，$x^2-1=(x-1)(x+1)$）.}

\ans{ex:poly:A5}{$f$ 与 $g$ 的公因式整除 $g-hf$；$f$ 与 $g-hf$ 的公因式整除 $(g-hf)+hf=g$. 两组公因式相同，最大公因式也相同.}

\ans{ex:poly:A6}{$x^2+1=(x-1)(x+1)+2$，故 $u=\frac12$，$v=-\frac12(x-1)$.}

\ans{ex:poly:A7}{有理根只可能是 $\pm1,\pm2,\pm\frac12$. 得 $2x^3-5x^2+x+2=(x-1)(x-2)(2x+1)$，有理根为 $1,2,-\frac12$.}

\ans{ex:poly:A8}{(1) 在 $\mathbb{Q}$ 与 $\mathbb{R}$ 上为 $(x^2+x+1)(x^2-x+1)$；在 $\mathbb{C}$ 上为 $(x-\omega)(x-\bar\omega)(x+\omega)(x+\bar\omega)$，$\omega=\mathrm e^{2\pi\mathrm i/3}$. (2) 在 $\mathbb{Q}$ 与 $\mathbb{R}$ 上为 $(x-1)(x+1)(x^2+x+1)(x^2-x+1)$；在 $\mathbb{C}$ 上为 $\prod_{k=0}^5(x-\mathrm e^{k\pi\mathrm i/3})$. (3) 在 $\mathbb{Q}$ 上不可约；在 $\mathbb{R}$ 上为 $(x-\sqrt[3]2)(x^2+\sqrt[3]2x+\sqrt[3]4)$；在 $\mathbb{C}$ 上为 $(x-\sqrt[3]2)(x-\sqrt[3]2\,\omega)(x-\sqrt[3]2\,\omega^2)$.}

\ans{ex:poly:A9}{$f(1)=f'(1)=0$，$f''(1)=20\ne0$，是二重根. 也可以连续两次用 $x-1$ 作综合除法.}

\ans{ex:poly:A10}{有重根当且仅当与导数 $3x^2-3$ 有公共根，即 $f(1)=0$ 或 $f(-1)=0$. $a=2$ 时重根为 $1$；$a=-2$ 时重根为 $-1$.}

\ans{ex:poly:A11}{$\sigma_1=2$，$\sigma_2=3$，$\sigma_3=4$. 所求依次为 $2$；$\sigma_1^2-2\sigma_2=-2$；$\sigma_2/\sigma_3=\frac34$.}

\ans{ex:poly:A12}{$(x^2-2x-1)(x^2+1)=x^4-2x^3-2x-1$. （$1+\sqrt2$ 的共轭 $1-\sqrt2$、$\mathrm i$ 的共轭 $-\mathrm i$ 也必须是根，见第七章章末习题~\ref{ex:poly:B10}.）}

\ans{ex:poly:A13}{$(x-2)(x^2-2x+2)=x^3-4x^2+6x-4$.}

\ans{ex:poly:A14}{$f(\vv A)=(\vv A-\vv I)^2=\vv O$. 于是 $\vv A^2=2\vv A-\vv I$，$\vv A^3=3\vv A-2\vv I$，$g(\vv A)=2\vv A-2\vv I=\begin{pmatrix}0&4\\0&0\end{pmatrix}$.}

\ans{ex:poly:A15}{$f(x)=2x^2-3x+1$（四个点恰好落在一条抛物线上，插值多项式的三次项系数为 $0$）.}

\ans{ex:poly:A16}{$x^n-1=(x-1)(x^{n-1}+x^{n-2}+\cdots+1)$. 一般地 $x^n-a^n=(x-a)(x^{n-1}+ax^{n-2}+\cdots+a^{n-1})$；也可由余数定理：$x=a$ 时 $x^n-a^n=0$.}

\ans{ex:poly:A17}{$f(x)=2-2(x-1)+2(x-1)^3+(x-1)^4$.}

\ans{ex:poly:A18}{设余式为 $ax+b$，由 $a+b=3$，$2a+b=5$ 得余式为 $2x+1$.}

\ans{ex:poly:A19}{$60.5v^2+55v-100=0$，即 $121v^2+110v-200=0$，正根 $v=\frac{10}{11}$，内部收益率 $r=10\%$.}

\ans{ex:poly:A20}{比较次数得 $\deg f=\deg g$，$f=hg$ 中 $h$ 为常数，且 $f\ne0$ 时 $h\ne0$.}

\ans{ex:poly:A21}{设 $h$ 是 $f_1,g_1$ 的公因式，则 $dh$ 是 $f,g$ 的公因式，故 $dh\mid d$，$h$ 是常数.}

\subsection*{B 组提示}

\ans{ex:poly:B1}{由 Bézout 等式 $uf+vg=1$，$u'f+v'h=1$，两式相乘.}
\ans{ex:poly:B2}{用提示中的等式反复降低指数，指数的变化与整数的辗转相除 $\gcd(m,n)=\gcd(m-n,n)$ 完全平行.}
\ans{ex:poly:B3}{$f/\gcd(f,f')=(x-1)(x+2)$. 一般地，$p_i^{k_i}$ 在 $f'$ 中恰好以 $p_i^{k_i-1}$ 出现（需要用到 $p_i$ 与 $p_i'$ 互素）.}
\ans{ex:poly:B4}{$p$ 不可约，由引理~\ref{lem:primediv} 即得. 反例：$p=x^2$，$f=x$.}
\ans{ex:poly:B5}{两边都是次数小于 $n$ 的多项式，在 $x_1,\ldots,x_n$ 处取值相同（$\ell_i(x_j)$ 的性质），用推论~\ref{cor:rootcount}.}
\ans{ex:poly:B6}{$f(x)-f(0)$ 以 $0,1,2,\ldots$ 为根，有无穷多个根.}
\ans{ex:poly:B7}{三次多项式可约当且仅当有一次因式，即有有理根；可能的有理根只有 $\pm1$，都不是根. 四次多项式可能分解为两个二次因式之积而没有有理根，例如 $(x^2+1)^2$.}
\ans{ex:poly:B8}{设余式为 $ax+b$，代入 $x=1,2$：$(2^n-1)x+2-2^n$.}
\ans{ex:poly:B9}{由 $f(2)=4$ 得 $c=1$，$f(x)=x^2$.}
\ans{ex:poly:B10}{$(x-a)^2-2b^2$ 是有理系数的，用它除 $f$，余式 $rx+s$ 在 $a+b\sqrt2$ 处为零，由 $\sqrt2$ 无理得 $r=s=0$.}
\ans{ex:poly:B11}{对整数 $n$，$f(n)$ 与 $f(n\bmod2)$ 的奇偶性相同（$n-m\mid f(n)-f(m)$）.}
\ans{ex:poly:B12}{依次为 $x^2+1$，$x-2$，$(x-2)^2$.}
\ans{ex:poly:B13}{$N'(v)=C_1+2C_2v+\cdots>0$. 由 $N(0)=-C_0<0$ 与 $N(1)=\sum C_i-C_0>0$，用零点定理.}
\ans{ex:poly:B14}{$\binom xk$ 在整数点取整数值；把 $f$ 写成 $\sum c_k\binom xk$，依次代入 $x=0,1,\ldots,n$ 解出 $c_k$，它们都是整数.}
\ans{ex:poly:B15}{核中的多项式有 $n$ 个不同的根而次数小于 $n$，只能是零.}
\ans{ex:poly:B16}{$\gcd(p,f)$ 整除 $p$，只能是 $1$ 或 $p$ 的常数倍.}

\section{第八章　特征值与特征向量}

\subsection*{A 组答案}

\ans{ex:eig:A1}{以下``属于 $\lambda$''后给出特征子空间的一组基，特征向量是它们的非零线性组合.
\begin{enumerate}
	\item $\lambda=3$：$(1,0)^\trans$；$\lambda=2$：$(1,-1)^\trans$.
	\item $\lambda=3$：$(1,1)^\trans$；$\lambda=-1$：$(1,-1)^\trans$.
	\item $\lambda=\mathrm i$：$(1,\mathrm i)^\trans$；$\lambda=-\mathrm i$：$(1,-\mathrm i)^\trans$（在 $\mathbb{R}$ 上没有特征值）.
	\item $\lambda=2$：$(1,1)^\trans$；$\lambda=3$：$(2,1)^\trans$.
	\item $\lambda=2$（二重）：$(1,0,0)^\trans$，$(0,1,1)^\trans$；$\lambda=0$：$(0,1,-1)^\trans$.
	\item $\lambda=5$：$(1,1,1)^\trans$；$\lambda=-1$（二重）：$(1,-1,0)^\trans$，$(1,0,-1)^\trans$.
\end{enumerate}}

\ans{ex:eig:A2}{除第 3 个（在 $\mathbb{R}$ 上没有特征值）外都可对角化，$\vv P$ 取 A1 中各特征向量为列. 例如第 4 个：$\vv P=\begin{pmatrix}1&2\\1&1\end{pmatrix}$，$\vv P^{-1}\vv{AP}=\diag(2,3)$；第 6 个：$\vv P=\begin{pmatrix}1&1&1\\1&-1&0\\1&0&-1\end{pmatrix}$，$\vv P^{-1}\vv{AP}=\diag(5,-1,-1)$.}

\ans{ex:eig:A3}{第 1 个可以（两个不同的特征值）；第 2 个不可以（$2$ 的几何重数为 $1$）；第 3 个不可以（$2$ 是三重特征值，$2\vv I-\vv A$ 的秩为 $1$，几何重数为 $2$）；第 4 个可以（$1$ 的代数重数与几何重数都是 $2$）.}

\ans{ex:eig:A4}{$\det\vv A=6$，$\tr\vv A=6$；$\vv A^{-1}$ 的特征值为 $1,\frac12,\frac13$；$\vv A^2-2\vv A+\vv I=(\vv A-\vv I)^2$ 的特征值为 $0,1,4$；$\det\vv A^*=(\det\vv A)^2=36$；$\det(\vv A^2+\vv I)=2\cdot5\cdot10=100$.}

\ans{ex:eig:A5}{特征多项式为 $\lambda^2-5\lambda+6$，特征值 $2,3$；第二种情形为 $\lambda^2-2\lambda+5$，特征值 $1\pm2\mathrm i$.}

\ans{ex:eig:A6}{$\det(\lambda\vv I-\vv A^\trans)=\det\big((\lambda\vv I-\vv A)^\trans\big)$. 特征向量一般不同：$\vv A=\begin{pmatrix}1&1\\0&2\end{pmatrix}$ 属于 $1$ 的特征向量为 $(1,0)^\trans$，$\vv A^\trans$ 属于 $1$ 的特征向量为 $(1,-1)^\trans$.}

\ans{ex:eig:A7}{$\lambda\vv I-\vv A$ 仍是三角矩阵，行列式为 $\prod(\lambda-a_{ii})$.}

\ans{ex:eig:A8}{若 $\vv{Ax}=\lambda\vv x$，$\vv x\ne\zero$，则 $\vv A^k\vv x=\lambda^k\vv x$. 由 $\vv A^k=\vv O$ 得 $\lambda^k=0$；由 $\vv A^2=\vv A$ 得 $\lambda^2=\lambda$；由 $\vv A^2=\vv I$ 得 $\lambda^2=1$.}

\ans{ex:eig:A9}{$\vv A=\vv P\diag(2,3)\vv P^{-1}$，$\vv P=\begin{pmatrix}1&2\\1&1\end{pmatrix}$，
$\vv A^n=\begin{pmatrix}2\cdot3^n-2^n&2^{n+1}-2\cdot3^n\\3^n-2^n&2^{n+1}-3^n\end{pmatrix}$.}

\ans{ex:eig:A10}{特征值 $1$ 与 $0.7$. 稳态分布 $\vv\pi=(\frac23,\frac13)^\trans$. $\vv P^n\to\begin{pmatrix}\frac23&\frac23\\ \frac13&\frac13\end{pmatrix}$：两列相同，说明无论开始时顾客在哪一家，长期看都以同样的概率分布在两家，初始状态被``遗忘''了.}

\ans{ex:eig:A11}{$\vv A\vv\xi=(-1,a+2,b+1)^\trans=\lambda(1,1,-1)^\trans$，得 $\lambda=-1$，$a=-3$，$b=0$.}

\ans{ex:eig:A12}{特征多项式 $(\lambda-1)^2(\lambda+1)$. 可对角化当且仅当 $\rk(\vv I-\vv A)=1$，即 $x=-1$.}

\ans{ex:eig:A13}{$\vv A^2=\begin{pmatrix}a^2+bc&b(a+d)\\c(a+d)&d^2+bc\end{pmatrix}$，代入直接验证.}

\ans{ex:eig:A14}{依次为 $(x-1)(x-2)$，$(x-1)^2(x-2)$，$x^3$，$x(x-3)$.}

\ans{ex:eig:A15}{相似. 取 $\vv P=\begin{pmatrix}0&1\\1&0\end{pmatrix}$（交换两个基向量的次序），$\vv P^{-1}\begin{pmatrix}1&1\\0&1\end{pmatrix}\vv P=\begin{pmatrix}1&0\\1&1\end{pmatrix}$.}

\ans{ex:eig:A16}{$\vv A\vv e_1=a_{11}\vv e_1$，$\vv A\vv e_2=a_{12}\vv e_1+a_{22}\vv e_2$，……，$\vv A\vv e_k$ 只含 $\vv e_1,\ldots,\vv e_k$.}

\ans{ex:eig:A17}{特征根 $2,-1$，$a_k=\frac13\big(2^k-(-1)^k\big)$.}

\ans{ex:eig:A18}{矩阵 $\begin{pmatrix}0&1.1\\0.6&0.5\end{pmatrix}$ 的特征值为 $1.1$ 与 $-0.6$. 长期年增长率为 $10\%$；属于 $1.1$ 的特征向量为 $(1,1)^\trans$，幼体与成年个体之比稳定在 $1:1$.}

\ans{ex:eig:A19}{特征多项式 $\lambda^2-(\tr\vv A)\lambda+\det\vv A$ 没有实根，当且仅当判别式 $(\tr\vv A)^2-4\det\vv A<0$.}

\ans{ex:eig:A20}{投影：方向 $(1,2)^\trans$ 上的向量不动（特征值 $1$），垂直方向 $(2,-1)^\trans$ 上的向量变为零（特征值 $0$），矩阵 $\frac15\begin{pmatrix}1&2\\2&4\end{pmatrix}$. 反射：$(1,2)^\trans$ 属于 $1$，$(2,-1)^\trans$ 属于 $-1$，矩阵 $\frac15\begin{pmatrix}-3&4\\4&3\end{pmatrix}$.}

\subsection*{B 组提示}

\ans{ex:eig:B1}{$\vv A\vv u=(\vv v^\trans\vv u)\vv u$；$\rk\vv A=1$，$0$ 的几何重数为 $n-1$. 特征值为 $\vv v^\trans\vv u$ 与 $0$（$n-1$ 重）. $\vv v^\trans\vv u=0$ 时 $0$ 是 $n$ 重特征值，但几何重数只有 $n-1$.}
\ans{ex:eig:B2}{$\vv A\vv B\vv p=\vv B\vv A\vv p=\lambda\vv B\vv p$，而 $V_\lambda$ 是一维的. 后半句：用 Lagrange 插值构造多项式 $f$ 使 $f(\lambda_i)$ 等于 $\vv B$ 的对应特征值.}
\ans{ex:eig:B3}{$\vv P^{-1}f(\vv A)\vv P=f(\vv D)$. 逆命题不成立：$\vv A=\begin{pmatrix}0&1\\0&0\end{pmatrix}$ 不可对角化，$\vv A^2=\vv O$ 可以.}
\ans{ex:eig:B4}{$\vv A|_W$ 的最小多项式整除 $\vv A$ 的最小多项式，用定理~\ref{thm:minpolydiag}.}
\ans{ex:eig:B5}{取 $\vv v\ne\zero$，证明 $\vv v,\vv{Av}$ 线性无关，在这组基下 $\vv A$ 的矩阵为 $\begin{pmatrix}0&-1\\1&0\end{pmatrix}$.}
\ans{ex:eig:B6}{设 $\vv{Ax}=\lambda\vv x$，取 $|x_i|$ 最大的分量，看第 $i$ 个方程 $(\lambda-a_{ii})x_i=\sum_{j\ne i}a_{ij}x_j$.}
\ans{ex:eig:B7}{$\vv A(1,\ldots,1)^\trans=s(1,\ldots,1)^\trans$. 可逆时 $s\ne0$，$\vv A^{-1}$ 每行之和为 $\frac1s$.}
\ans{ex:eig:B8}{特征值为 $a+b\omega^k+c\omega^{2k}$（$k=0,1,2$），行列式是三者之积.}
\ans{ex:eig:B9}{特征根 $2$ 是二重的，通项形如 $(c_1+c_2k)2^k$，得 $a_k=(k+1)2^k$. $k2^k$ 来自 Jordan 块的幂（第八章选读节）.}
\ans{ex:eig:B10}{对三角化（或可对角化）的情形看对角元. 特征值为 $0,\sqrt3,-\sqrt3$.}
\ans{ex:eig:B11}{$x^3-x=x(x-1)(x+1)$，最小多项式整除它.}
\ans{ex:eig:B12}{$2$ 是二重特征值，需 $\rk(2\vv I-\vv A)=1$，条件为 $a=0$（$b,c$ 任意）.}
\ans{ex:eig:B13}{转移矩阵 $\begin{pmatrix}0&0&1\\ \frac12&0&0\\ \frac12&1&0\end{pmatrix}$，稳态分布 $\frac15(2,1,2)^\trans$.}
\ans{ex:eig:B14}{特征值 $\pm1$，$\mathrm e^{t\vv A}=\begin{pmatrix}\cosh t&\sinh t\\ \sinh t&\cosh t\end{pmatrix}$. 对 $t$ 求导验证.}
\ans{ex:eig:B15}{稳态分布 $(\frac13,\frac13,\frac13)^\trans$；特征值 $1,\frac14,\frac14$，偏离部分每一步缩小为原来的 $\frac14$.}
\ans{ex:eig:B16}{几何重数 $n-\rk(\lambda\vv I-\vv A)$ 不超过代数重数 $k$.}
\ans{ex:eig:B17}{得分（按判决 $1,2,3,4$）约为：$d=0.5$ 时 $(0.169,0.212,0.265,0.354)$；$d=0.85$ 时 $(0.128,0.183,0.261,0.428)$；$d=0.95$ 时 $(0.119,0.175,0.258,0.448)$；$d=1$ 时恰为 $\frac1{35}(4,6,9,16)$. 排序不变；$d$ 越大，随机跳转越少，得分越集中到引证链末端的判决 $4$ 上.}

\subsection*{C 组证明要点（正文引用的题目）}

\ans{ex:eig:C1}{直接计算两个分块乘积：
	\[\begin{pmatrix}\vv I_m&-\vv A\\ \vv O&\lambda\vv I_n\end{pmatrix}\begin{pmatrix}\lambda\vv I_m&\vv A\\ \vv B&\vv I_n\end{pmatrix}=\begin{pmatrix}\lambda\vv I_m-\vv{AB}&\vv O\\ \lambda\vv B&\lambda\vv I_n\end{pmatrix},\]\[
	\begin{pmatrix}\vv I_m&\vv O\\ -\vv B&\lambda\vv I_n\end{pmatrix}\begin{pmatrix}\lambda\vv I_m&\vv A\\ \vv B&\vv I_n\end{pmatrix}=\begin{pmatrix}\lambda\vv I_m&\vv A\\ \vv O&\lambda\vv I_n-\vv{BA}\end{pmatrix}.\]
	记中间矩阵的行列式为 $D$. 对两式取行列式（第五章章末习题~\ref{ex:det:A17}），得 $\lambda^nD=\lambda^n\det(\lambda\vv I_m-\vv{AB})$ 与 $\lambda^nD=\lambda^m\det(\lambda\vv I_n-\vv{BA})$，因此 $\lambda^n\det(\lambda\vv I_m-\vv{AB})=\lambda^m\det(\lambda\vv I_n-\vv{BA})$. 两边作为 $\lambda$ 的多项式相等，因子 $\lambda^n$、$\lambda^m$ 只影响根 $0$ 的重数，所以两个特征多项式的非零根（连同重数）相同.}

\section{第九章　内积空间}

\subsection*{A 组答案}

\ans{ex:ip:A1}{$\|\vv\alpha\|=\|\vv\beta\|=3$，$\langle\vv\alpha,\vv\beta\rangle=4$，夹角 $\arccos\frac49\approx63.6^\circ$，$\|\vv\alpha-\vv\beta\|=\sqrt{10}$.}

\ans{ex:ip:A2}{(1) 是. (2) 不是：$\vv x=(1,-1)^\trans$ 时 $\langle\vv x,\vv x\rangle=-2$. (3) 是：$\langle\vv x,\vv x\rangle=x_1^2+(x_1+x_2)^2$. (4) 不是：$\langle\vv e_2,\vv e_2\rangle=-1$.}

\ans{ex:ip:A3}{(1) $\frac1{\sqrt3}(1,1,1)^\trans$，$\frac1{\sqrt2}(0,1,-1)^\trans$，$\frac1{\sqrt6}(-2,1,1)^\trans$. (2) $\frac1{\sqrt2}(1,1,0,0)^\trans$，$\frac1{\sqrt6}(1,-1,2,0)^\trans$，$\frac1{2\sqrt3}(-1,1,1,3)^\trans$.}

\ans{ex:ip:A4}{解 $x_1+x_2+x_3=0$，$x_1-x_3=0$，得 $\pm\frac1{\sqrt6}(1,-2,1)^\trans$.}

\ans{ex:ip:A5}{逐对验证内积；坐标为 $\langle\vv\alpha,\vv e_i\rangle$，即 $(5,4,-2)^\trans$.}

\ans{ex:ip:A6}{两个生成向量已经正交，投影为 $\frac32(1,1,0,0)^\trans+\frac72(0,0,1,1)^\trans=\big(\frac32,\frac32,\frac72,\frac72\big)^\trans$，
	
	距离为 $$\|(-\frac12,\frac12,-\frac12,\frac12)^\trans\|=1$$.}

\ans{ex:ip:A7}{正规方程 $5a+10b=15$，$10a+30b=42$，得 $y=0.6+1.2x$，残差平方和为 $1.6$.}

\ans{ex:ip:A8}{第二列须与第一、三列正交且为单位向量：$a=c$，$b=a$，$3a^2=1$. 故 $a=b=c=\frac1{\sqrt3}$ 或 $a=b=c=-\frac1{\sqrt3}$.}

\ans{ex:ip:A9}{展开 $\big\langle\sum\vv\alpha_i,\sum\vv\alpha_j\big\rangle$，交叉项都为零.}

\ans{ex:ip:A10}{由三角不等式 $\|\vv\alpha\|\le\|\vv\alpha-\vv\beta\|+\|\vv\beta\|$，再交换 $\vv\alpha,\vv\beta$.}

\ans{ex:ip:A11}{(1) 对 $(a_1,\ldots,a_n)$ 与 $(1,\ldots,1)$ 用 Cauchy–Schwarz 不等式. (2) 对 $(\sqrt{a_1},\ldots,\sqrt{a_n})$ 与 $\big(\frac1{\sqrt{a_1}},\ldots,\frac1{\sqrt{a_n}}\big)$ 用之.}

\ans{ex:ip:A12}{$\|1\|=1$，$\|x\|=\frac1{\sqrt3}$，$\langle1,x\rangle=\frac12$，夹角 $\frac\pi6$. 正交化并单位化得 $1$，$\sqrt3(2x-1)$.}

\ans{ex:ip:A13}{$\ker\vv A$ 的标准正交基：$\frac1{\sqrt3}(1,-1,1,0)^\trans$，$\frac1{\sqrt{15}}(-1,-2,-1,3)^\trans$；行空间的标准正交基：$\frac1{\sqrt3}(1,1,0,1)^\trans$，$\frac1{\sqrt{15}}(-2,1,3,1)^\trans$. 两组之间两两内积为零，维数 $2+2=4$，故互为正交补.}

\ans{ex:ip:A14}{$\vv Q=\begin{pmatrix}\frac1{\sqrt2}&\frac1{\sqrt6}\\ \frac1{\sqrt2}&-\frac1{\sqrt6}\\0&\frac2{\sqrt6}\end{pmatrix}$，$\vv R=\begin{pmatrix}\sqrt2&\frac1{\sqrt2}\\0&\sqrt{\frac32}\end{pmatrix}$.}

\ans{ex:ip:A15}{$\vv H=\frac19\begin{pmatrix}7&-4&-4\\-4&1&-8\\-4&-8&1\end{pmatrix}$；$\vv H\vv u=-\vv u$；$(2,-1,0)^\trans$ 与 $\vv u$ 正交，$\vv H(2,-1,0)^\trans=(2,-1,0)^\trans$.}

\ans{ex:ip:A16}{第一个行列式为 $1$，是旋转，旋转角 $\arctan\frac43\approx53.1^\circ$；第二个行列式为 $-1$，是关于直线 $y=\frac x2$（方向 $(2,1)^\trans$）的反射.}

\ans{ex:ip:A17}{设计矩阵只有一列 $\vv x$，正规方程为 $(\vv x^\trans\vv x)b=\vv x^\trans\vv y$.}

\ans{ex:ip:A18}{投影为 $2(1,1,1)^\trans$，距离为 $\|(-1,0,1)^\trans\|=\sqrt2$.}

\ans{ex:ip:A19}{过渡矩阵的第 $j$ 列是新基第 $j$ 个向量在旧（标准正交）基下的坐标，坐标的标准内积等于向量的内积，故各列标准正交.}

\ans{ex:ip:A20}{内积打分：$6,1,4$，排序为甲（第一条）$>$ 丙 $>$ 乙. 余弦相似度：约 $0.926$，$0.140$，$0.943$，排序为丙 $>$ 甲 $>$ 乙. 第一条法条更长，内积占了便宜.}

\ans{ex:ip:A21}{$\vv P^\trans=(\vv u^\trans)^\trans\vv u^\trans=\vv P$，$\vv P^2=\vv u(\vv u^\trans\vv u)\vv u^\trans=\vv P$. $\vv{Px}=\langle\vv x,\vv u\rangle\vv u$ 正是定理~\ref{thm:projformula} 中 $k=1$ 的公式. 一般非零 $\vv u$ 时投影矩阵为 $\dfrac{\vv u\vv u^\trans}{\vv u^\trans\vv u}$.}

\subsection*{B 组提示}

\ans{ex:ip:B1}{展开两个平方相减. 取 $\vv x=(1,0)^\trans$，$\vv y=(0,1)^\trans$ 检验平行四边形法则 $\|\vv x+\vv y\|^2+\|\vv x-\vv y\|^2=2\|\vv x\|^2+2\|\vv y\|^2$，左边为 $8$，右边为 $4$.}
\ans{ex:ip:B2}{对 $\vv\alpha-\sum\langle\vv\alpha,\vv e_i\rangle\vv e_i$ 用勾股定理.}
\ans{ex:ip:B3}{证明 $\vv x-\vv{Px}$ 与 $\imm\vv P$ 正交：$\langle\vv x-\vv{Px},\vv{Py}\rangle=\vv y^\trans\vv P^\trans(\vv I-\vv P)\vv x=0$. 只有幂等时是沿 $\ker\vv P$ 向 $\imm\vv P$ 的``斜投影''，例如 $\begin{pmatrix}1&1\\0&0\end{pmatrix}$.}
\ans{ex:ip:B4}{目标函数是 $\|\vv y-\vv A\vv\beta\|_W^2$，其中 $\langle\vv x,\vv y\rangle_W=\vv x^\trans\vv W\vv y$ 是加权内积；残差在这个内积下与各列正交.}
\ans{ex:ip:B5}{对 $\ln y=\ln c+kt$ 作最小二乘，约得 $k\approx0.30$，$c\approx2.0$. 它最小化的是对数残差的平方和，相当于对相对误差加权.}
\ans{ex:ip:B6}{$\vv G=\vv A^\trans\vv A$（$\vv A$ 以 $\vv\alpha_i$ 为列），用第十一章命题~\ref{prop:gram} 或直接 $\vv x^\trans\vv G\vv x=\|\vv A\vv x\|^2$. $k=2$ 时 $\det\vv G=\|\vv\alpha\|^2\|\vv\beta\|^2\sin^2\theta$.}
\ans{ex:ip:B7}{距离是 $\vv x_0-\vv x_1$（$\vv x_1$ 为超平面上任一点）在法向 $\vv a$ 上投影的长度.}
\ans{ex:ip:B8}{设 $\vv{Qz}=\lambda\vv z$，考察 $\bar{\vv z}^\trans\vv Q^\trans\vv Q\vv z$.}
\ans{ex:ip:B9}{用积化和差逐对计算积分. 投影为 $2\sin x-\sin2x$（偶函数的系数为零）.}
\ans{ex:ip:B10}{全部最小二乘解为 $\hat{\vv x}+\ker\vv A$，把 $\hat{\vv x}$ 分解到 $\imm\vv A^\trans\oplus\ker\vv A$，用勾股定理.}
\ans{ex:ip:B11}{$\|\vv x-\vv y\|^2=2\|\vv x\|^2-2\langle\vv x,\vv y\rangle$，直接计算 $\vv H\vv x=\vv x-2\vv u\vv u^\trans\vv x$.}
\ans{ex:ip:B12}{对中心化的数据用 A17，再代回截距；正规方程的第一个方程就是 $\bar y=a+b\bar x$.}
\ans{ex:ip:B13}{保持夹角就保持正交. 把一组标准正交基的像归一化，证明 $\|\mathscr T\vv e_i\|$ 都相等（考察 $\vv e_i\pm\vv e_j$ 的像是否正交）.}
\ans{ex:ip:B14}{截距为第一组均值 $4$，虚拟变量系数为两组均值之差 $8-4=4$.}
\ans{ex:ip:B15}{$\vv R(\alpha)\vv R(\beta)=\vv R(\alpha+\beta)$；$\vv F(\alpha)\vv F(\beta)=\vv R(\alpha-\beta)$；计算 $\vv R(\alpha)\vv F(\beta)$ 与 $\vv F(\beta)\vv R(\alpha)$ 比较.}

\section{第十章　矩阵分解}

\subsection*{A 组答案}

\ans{ex:dec:A1}{（$\vv Q$ 的列依次对应所列特征值，答案不唯一.）
\begin{enumerate}
	\item $\vv Q=\frac1{\sqrt2}\begin{pmatrix}1&1\\1&-1\end{pmatrix}$，对角矩阵 $\diag(3,-1)$.
	\item $\vv Q=\frac1{\sqrt2}\begin{pmatrix}1&1\\1&-1\end{pmatrix}$，对角矩阵 $\diag(2,4)$.
	\item 特征值 $4,0,1$，特征向量 $\frac1{\sqrt2}(1,1,0)^\trans$，$\frac1{\sqrt2}(1,-1,0)^\trans$，$(0,0,1)^\trans$.
	\item 特征值 $3,0,0$，特征向量 $\frac1{\sqrt3}(1,1,1)^\trans$，$\frac1{\sqrt2}(1,-1,0)^\trans$，$\frac1{\sqrt6}(1,1,-2)^\trans$（后两个须在特征子空间内正交化）.
\end{enumerate}}

\ans{ex:dec:A2}{$\begin{pmatrix}3&-1\\-1&3\end{pmatrix}=2\cdot\frac12\begin{pmatrix}1&1\\1&1\end{pmatrix}+4\cdot\frac12\begin{pmatrix}1&-1\\-1&1\end{pmatrix}$. 两个投影矩阵之和为 $\vv I$，乘积为 $\vv O$.}

\ans{ex:dec:A3}{$\vv A=1\cdot\frac12\begin{pmatrix}1&1\\1&1\end{pmatrix}+2\cdot\frac12\begin{pmatrix}1&-1\\-1&1\end{pmatrix}=\frac12\begin{pmatrix}3&-1\\-1&3\end{pmatrix}$.}

\ans{ex:dec:A4}{属于 $1$ 的特征子空间是 $(0,1,1)^\trans$ 的正交补，故 $\vv A=\vv I-2\vv P$，$\vv P$ 是向 $\spn\big((0,1,1)^\trans\big)$ 的投影：$\vv A=\begin{pmatrix}1&0&0\\0&0&-1\\0&-1&0\end{pmatrix}$.}

\ans{ex:dec:A5}{特征值 $4,1,0$. 最大值 $4$，在 $\pm\frac1{\sqrt2}(1,1,0)^\trans$ 处取到；最小值 $0$，在 $\pm\frac1{\sqrt2}(1,-1,0)^\trans$ 处取到.}

\ans{ex:dec:A6}{矩阵为 $\begin{pmatrix}1&2\\2&1\end{pmatrix}$，最大值 $3$（在 $\pm\frac1{\sqrt2}(1,1)$ 处），最小值 $-1$（在 $\pm\frac1{\sqrt2}(1,-1)$ 处）.}

\ans{ex:dec:A7}{(1) $\sigma_1=3$，$\sigma_2=2$；$\vv U=\begin{pmatrix}0&1\\-1&0\end{pmatrix}$，$\vv\Sigma=\diag(3,2)$，$\vv V=\begin{pmatrix}0&1\\1&0\end{pmatrix}$. (2) $\vv U=\vv V=\frac1{\sqrt2}\begin{pmatrix}1&1\\1&-1\end{pmatrix}$，$\vv\Sigma=\diag(2,0)$. (3) $\vv U=(1)$，$\vv\Sigma=(5,\ 0)$，$\vv V=\frac15\begin{pmatrix}3&-4\\4&3\end{pmatrix}$.}

\ans{ex:dec:A8}{$\vv A^\trans\vv A=\begin{pmatrix}1&1\\1&2\end{pmatrix}$，特征值 $\frac{3\pm\sqrt5}2$，奇异值为 $\frac{\sqrt5+1}2$ 与 $\frac{\sqrt5-1}2$（黄金分割比及其倒数）；特征值是 $1,1$.}

\ans{ex:dec:A9}{对称正定矩阵的奇异值分解就是谱分解. $\vv A_1=4\cdot\frac12\begin{pmatrix}1&1\\1&1\end{pmatrix}=\begin{pmatrix}2&2\\2&2\end{pmatrix}$，误差 $\|\vv A-\vv A_1\|_F=\sigma_2=2$.}

\ans{ex:dec:A10}{$\vv A=9\vv P_1+\vv P_2$，$\vv P_1=\frac12\begin{pmatrix}1&1\\1&1\end{pmatrix}$，$\vv P_2=\frac12\begin{pmatrix}1&-1\\-1&1\end{pmatrix}$. $\vv A^{10}=9^{10}\vv P_1+\vv P_2$；$\vv B=3\vv P_1+\vv P_2=\begin{pmatrix}2&1\\1&2\end{pmatrix}$.}

\ans{ex:dec:A11}{$\vv A=\vv Q\vv O\vv Q^\trans=\vv O$. 反例：$\begin{pmatrix}0&1\\0&0\end{pmatrix}$.}

\ans{ex:dec:A12}{$\vv A^2=\vv A^\trans\vv A=\vv I$，特征值满足 $\lambda^2=1$. 它是关于特征子空间 $V_1$ 的``镜面反射''：$V_1$ 中的向量不动，$V_{-1}=V_1^\perp$ 中的向量反向.}

\ans{ex:dec:A13}{由第八章章末习题~\ref{ex:eig:C1}（取 $\vv B=\vv A^\trans$），或直接：若 $\vv A^\trans\vv A\vv x=\lambda\vv x$，$\lambda\ne0$，则 $\vv{AA}^\trans(\vv{Ax})=\lambda\vv{Ax}$ 且 $\vv{Ax}\ne\zero$；两边的特征子空间维数相同.}

\ans{ex:dec:A14}{$\tr(\vv A^\trans\vv A)=\sum_{i,j}a_{ij}^2$；它也等于 $\vv A^\trans\vv A$ 的特征值之和 $\sum\sigma_i^2$.}

\ans{ex:dec:A15}{$\det\vv A=\det\vv U\det\vv\Sigma\det\vv V^\trans=\pm\sigma_1\cdots\sigma_n$；$\vv A^{-1}=\vv V\vv\Sigma^{-1}\vv U^\trans$.}

\ans{ex:dec:A16}{只有第二个可以：能被正交矩阵对角化的实矩阵必须对称（谱定理的逆命题）.}

\ans{ex:dec:A17}{$\vv S=\frac13\begin{pmatrix}10&8\\8&10\end{pmatrix}$，特征值 $6$ 与 $\frac23$. 第一主成分方向 $\frac1{\sqrt2}(1,1)^\trans$，解释方差比例 $\frac6{6+2/3}=90\%$.}

\ans{ex:dec:A18}{$\kappa\big(\diag(10,1)\big)=10$；$\begin{pmatrix}1&1\\1&1.01\end{pmatrix}$ 的特征值约为 $2.005$ 与 $0.004988$，$\kappa\approx402$. 二者行列式都等于 $0.01$ 的 $\diag(10,0.001)$ 条件数为 $10^4$，与后者不同：行列式不决定条件数.}

\ans{ex:dec:A19}{$40\times(800+600+1)=56040$ 个数，原图 $480000$ 个，约为原来的 $11.7\%$. $k\ge343$ 时存储量超过原图.}

\ans{ex:dec:A20}{由谱定理 $\vv A$ 相似于 $\diag(\lambda_1,\ldots,\lambda_n)$，秩不变. 奇异值的平方是 $\vv A^\trans\vv A$ 的非零特征值，个数为 $\rk(\vv A^\trans\vv A)=\rk\vv A$.}

\ans{ex:dec:A21}{取正交矩阵 $\vv Q$ 使 $\vv Q^\trans\vv{AQ}=\diag(\lambda_1,\ldots,\lambda_n)$，令 $\vv y=\vv Q^\trans\vv x$，则 $\vv x^\trans\vv A\vv x=\sum\lambda_iy_i^2$，且 $\vv x\ne\zero\iff\vv y\ne\zero$. 若某个 $\lambda_i\le0$，取 $\vv y=\vv e_i$；若全为正，则 $\vv y\ne\zero$ 时和为正.}

\subsection*{B 组提示}

\ans{ex:dec:B1}{$\vv A$ 的每个特征子空间都是 $\vv B$-不变的，$\vv B$ 限制在其上仍是对称的，在其中取 $\vv B$ 的标准正交特征向量.}
\ans{ex:dec:B2}{存在性：$\vv B=\vv Q\diag(\sqrt{\lambda_i})\vv Q^\trans$. 唯一性：$\vv B$ 的每个特征子空间（属于 $\mu$）都含于 $\vv A$ 的属于 $\mu^2$ 的特征子空间，而 $\mu\ge0$ 由 $\mu^2$ 唯一确定.}
\ans{ex:dec:B3}{在 $\vv q_1^\perp=\spn(\vv q_2,\ldots,\vv q_n)$ 中重复定理~\ref{thm:rayleigh} 的证明.}
\ans{ex:dec:B4}{$\|\vv{Ax}\|^2=\vv x^\trans\vv A^\trans\vv A\vv x\le\lambda_{\max}(\vv A^\trans\vv A)\|\vv x\|^2$，在 $\vv v_1$ 处取等.}
\ans{ex:dec:B5}{$\vv{Ax}=\lambda\vv x$ 时 $|\lambda|\,\|\vv x\|=\|\vv{Ax}\|$，再用 B4 及其``最小值''版本.}
\ans{ex:dec:B6}{$\vv A^\trans\vv A$ 的非零特征值为 $3,1$（$\vv A\vv A^\trans=\begin{pmatrix}2&1\\1&2\end{pmatrix}$ 更易算），$\sigma_1=\sqrt3$，$\sigma_2=1$. 先求 $\vv U$ 的列：$\frac1{\sqrt2}(1,1)^\trans$，$\frac1{\sqrt2}(1,-1)^\trans$，再用 $\vv v_i=\vv A^\trans\vv u_i/\sigma_i$，第三个 $\vv v_3=\frac1{\sqrt3}(1,1,-1)^\trans$.}
\ans{ex:dec:B7}{$\vv A^+\vv b$ 满足正规方程（在 $\vv V$、$\vv U$ 的坐标下逐分量验证），且属于行空间 $\imm\vv A^\trans$，由第九章章末习题~\ref{ex:ip:B10} 它是范数最小的一个.}
\ans{ex:dec:B8}{$\vv J$ 的特征值为 $p,0,\ldots,0$，故矩阵的特征值为 $1+(p-1)\rho$ 与 $1-\rho$（$p-1$ 重）. 第一主成分 $\frac1{\sqrt p}(1,\ldots,1)^\trans$，比例 $\frac{1+(p-1)\rho}p$，$\rho\to1$ 时趋于 $1$，$\rho\to0$ 时趋于 $\frac1p$.}
\ans{ex:dec:B9}{协方差矩阵的第一主成分几乎就是第一个指标的方向（方差 $100$ 远大于 $1$）；相关系数为 $0.6$，相关系数矩阵的第一主成分是 $\frac1{\sqrt2}(1,1)^\trans$. 量纲改变方差，从而改变主成分.}
\ans{ex:dec:B10}{$\vv A-\vv A_k=\sum_{i>k}\sigma_i\vv u_i\vv v_i^\trans$，它的最大奇异值是 $\sigma_{k+1}$.}
\ans{ex:dec:B11}{$\vv v_i$ 是 $\vv A^\trans\vv A$ 的单重特征值的单位特征向量，只能差一个符号；$\vv u_i=\vv A\vv v_i/\sigma_i$ 随之确定.}
\ans{ex:dec:B12}{设 $\vv A=\begin{pmatrix}a&b\\c&d\end{pmatrix}$，比较 $\vv A^\trans\vv A$ 与 $\vv{AA}^\trans$ 的元素，得 $b^2=c^2$ 与 $(a-d)(b-c)=0$.}
\ans{ex:dec:B13}{第一句由第八章章末习题~\ref{ex:eig:B6}（特征值为实数）. 第二句取 $\vv x=\vv e_i$ 用定理~\ref{thm:rayleigh}.}
\ans{ex:dec:B14}{$k=1$ 时每份文书只剩一个数，查询与四份文书中三份合同类文书的相似度都是 $1$，与第四份为 $0$——同 $k=2$ 时的结果，但``合同''内部的差别完全消失，信息损失更多.}
\ans{ex:dec:B15}{$\vv x^\trans(\vv A-a\vv I)\vv x\ge0$ 对一切 $\vv x$ 成立当且仅当 $\vv A-a\vv I$ 的特征值 $\lambda_i-a\ge0$.}

\section{第十一章　二次型}

\subsection*{A 组答案}

\ans{ex:quad:A1}{(1) $\begin{pmatrix}1&2&0\\2&-3&-1\\0&-1&0\end{pmatrix}$，秩 $3$. (2) 对角元为 $0$、非对角元为 $\frac12$ 的三阶矩阵，秩 $3$. (3) 元素全为 $1$ 的三阶矩阵，秩 $1$. (4) 对角元为 $1$、非对角元为 $\frac12$ 的 $n$ 阶矩阵，即 $\frac12(\vv I+\vv J)$，秩 $n$（特征值 $\frac{n+1}2$ 与 $\frac12$，都不为零）.}

\ans{ex:quad:A2}{$x_1^2+2x_2^2-2x_1x_2+6x_2x_3$.}

\ans{ex:quad:A3}{(1) $(x_1+x_2)^2+(x_2+3x_3)^2-4x_3^2$，令 $y_1=x_1+x_2$，$y_2=x_2+3x_3$，$y_3=x_3$，得 $y_1^2+y_2^2-4y_3^2$. (2) $(x_1-2x_2)^2-x_3^2$，令 $y_1=x_1-2x_2$，$y_2=x_3$，$y_3=x_2$，得 $y_1^2-y_2^2$. (3) $x_2(x_1-x_3)$，令 $y_1=\frac12(x_1+x_2-x_3)$，$y_2=\frac12(x_2-x_1+x_3)$，$y_3=x_3$，得 $y_1^2-y_2^2$. 替换都不唯一.}

\ans{ex:quad:A4}{特征值 $2,5,1$，$\vv Q=\begin{pmatrix}1&0&0\\0&\frac1{\sqrt2}&\frac1{\sqrt2}\\0&\frac1{\sqrt2}&-\frac1{\sqrt2}\end{pmatrix}$，标准形 $2y_1^2+5y_2^2+y_3^2$. 曲面是椭球面.}

\ans{ex:quad:A5}{(1) $p=2$，$q=1$，符号差 $1$，规范形 $z_1^2+z_2^2-z_3^2$. (2)(3) 都是 $p=q=1$，符号差 $0$，规范形 $z_1^2-z_2^2$.}

\ans{ex:quad:A6}{(1) 正定（顺序主子式 $2,1,3$）. (2) 不正定（$\Delta_2=-3$）. (3) 正定（顺序主子式 $5,1,1$）.}

\ans{ex:quad:A7}{矩阵 $\begin{pmatrix}2&1&0\\1&1&\frac t2\\0&\frac t2&1\end{pmatrix}$，$\Delta_3=1-\frac{t^2}2>0$，即 $-\sqrt2<t<\sqrt2$.}

\ans{ex:quad:A8}{(1) 特征值 $3,-1$，标准方程 $3u^2-v^2=1$，双曲线；$u$ 轴沿 $\frac1{\sqrt2}(1,-1)$，$v$ 轴沿 $\frac1{\sqrt2}(1,1)$. (2) 特征值 $4,2$，标准方程 $\frac{u^2}2+\frac{v^2}4=1$，椭圆；沿 $\frac1{\sqrt2}(1,1)$ 的半轴为 $\sqrt2$，沿 $\frac1{\sqrt2}(1,-1)$ 的半轴为 $2$.}

\ans{ex:quad:A9}{(1) 合同（秩都为 $3$，$p$ 都为 $2$）：取 $\vv C$ 使 $\vv C\vv e_1=\frac1{\sqrt3}\vv e_3$，$\vv C\vv e_2=\sqrt5\,\vv e_1$，$\vv C\vv e_3=\sqrt3\,\vv e_2$. (2) 不合同：$p=1$ 与 $p=2$. (3) 合同：$\vv C=\frac1{\sqrt2}\begin{pmatrix}1&1\\1&-1\end{pmatrix}$.}

\ans{ex:quad:A10}{$\vv A^*=(\det\vv A)\vv A^{-1}$ 是对称矩阵，特征值为 $\frac{\det\vv A}{\lambda_i}>0$.}

\ans{ex:quad:A11}{驻点 $(0,0)$ 与 $(1,1)$. Hessian $\begin{pmatrix}6x&-3\\-3&6y\end{pmatrix}$：在 $(0,0)$ 处不定，是鞍点；在 $(1,1)$ 处正定，取极小值 $-1$.}

\ans{ex:quad:A12}{顺序主子式 $2,3,4$，正定. $\vv L=\begin{pmatrix}\sqrt2&0&0\\-\frac1{\sqrt2}&\sqrt{\frac32}&0\\0&-\sqrt{\frac23}&\frac2{\sqrt3}\end{pmatrix}$.}

\ans{ex:quad:A13}{$\vv L=\begin{pmatrix}2&0&0\\1&3&0\\-1&1&2\end{pmatrix}$.}

\ans{ex:quad:A14}{方差 $\vv w^\trans\vv\Sigma\vv w=0.0215$，标准差约 $14.7\%$.}

\ans{ex:quad:A15}{$\vv w^*\propto\vv\Sigma^{-1}\vv 1\propto\big(\frac1{0.09},\frac1{0.04}\big)$，即 $\vv w^*=\big(\frac4{13},\frac9{13}\big)$，方差 $\frac9{325}$，标准差约 $16.6\%$.}

\ans{ex:quad:A16}{$(\vv C^\trans\vv{AC})^\trans=\vv C^\trans\vv A^\trans\vv C$；秩不变见推论~\ref{cor:rankinvariant}. $\det(\vv C^\trans\vv{AC})=(\det\vv C)^2\det\vv A$：保持符号，不保持数值.}

\ans{ex:quad:A17}{在 $\mathbb{R}$ 上为 $z_1^2+z_2^2-z_3^2$；在 $\mathbb{C}$ 上为 $z_1^2+z_2^2+z_3^2$.}

\ans{ex:quad:A18}{依次为单叶双曲面（$p=2,q=1$）、双叶双曲面（$p=1,q=2$）、圆锥面、圆柱面（秩为 $2$）.}

\ans{ex:quad:A19}{特征值 $3$ 与 $-1$. 取 $\vv x=(1,-1)^\trans$，二次型的值为 $-2$.}

\ans{ex:quad:A20}{对主子矩阵 $\vv A_I$（行、列指标集都为 $I$），把 $\vv x_I\ne\zero$ 在其余位置补零，$\vv x_I^\trans\vv A_I\vv x_I=\vv x^\trans\vv A\vv x>0$. 正定矩阵的行列式为正（特征值之积）.}

\subsection*{B 组提示}

\ans{ex:quad:B1}{可交换的对称矩阵可同时正交对角化（第十章章末习题~\ref{ex:dec:B1}），对角元对应相乘. 反例：$\vv A=\begin{pmatrix}2&1\\1&1\end{pmatrix}$，$\vv B=\diag(1,2)$.}
\ans{ex:quad:B2}{取 $\vv C$ 使 $\vv C^\trans\vv{AC}=\vv I$，$\vv C^\trans\vv{BC}=\diag(\mu_i)$，$\mu_i\ge0$，则 $\det(\vv A+\vv B)=\det\vv A\prod(1+\mu_i)$.}
\ans{ex:quad:B3}{$\vv x^\trans(\vv A+t\vv B)\vv x\ge(\lambda_{\min}(\vv A)-|t|\max_i|\lambda_i(\vv B)|)\|\vv x\|^2$，用 B5.}
\ans{ex:quad:B4}{第 $i,j$ 行列组成的二阶主子式为正.}
\ans{ex:quad:B5}{由谱定理写成 $\sum\lambda_iy_i^2$，$\|\vv y\|=\|\vv x\|$.}
\ans{ex:quad:B6}{特征值 $2,-1,-1$，$p=1$，$q=2$. 配方时先作 $x_1=y_1+y_2$，$x_2=y_1-y_2$.}
\ans{ex:quad:B7}{$\Delta_1=\lambda$，$\Delta_2=\lambda^2-1$，$\Delta_3=(\lambda-2)(\lambda+1)^2$，故 $\lambda>2$.}
\ans{ex:quad:B8}{驻点条件 $2\vv{Ax}=\mu\vv c$，代入约束定出 $\mu$. $\vv A$ 正定保证驻点是最小值点（第十一章最小二乘例题中的配方法同样适用）.}
\ans{ex:quad:B9}{驻点条件给出 $\vv w=a\vv 1+b\vv\mu$，由两个约束解出 $a,b$（关于 $m$ 线性），方差 $\sigma^2\|\vv w\|^2$ 是 $m$ 的二次函数.}
\ans{ex:quad:B10}{解 $\det(\vv B-\lambda\vv A)=3\lambda^2-1=0$，得 $\lambda=\pm\frac1{\sqrt3}$；对每个 $\lambda$ 求 $(\vv B-\lambda\vv A)\vv x=\zero$ 的解并使 $\vv x^\trans\vv A\vv x=1$，以它们为列作 $\vv C$，则 $\vv C^\trans\vv{BC}=\diag(\frac1{\sqrt3},-\frac1{\sqrt3})$.}
\ans{ex:quad:B11}{配方：$f=(\vv x+\frac12\vv A^{-1}\vv b)^\trans\vv A(\vv x+\frac12\vv A^{-1}\vv b)+c-\frac14\vv b^\trans\vv A^{-1}\vv b$. 半正定时不一定有最小值，例如 $\vv A=\vv O$，$\vv b\ne\zero$.}
\ans{ex:quad:B12}{写 $\vv A=\vv B^\trans\vv B$（$\vv B=\vv D^{1/2}\vv Q^\trans$），则 $\|\vv B\vv x\|^2=0$，$\vv{Ax}=\vv B^\trans\vv B\vv x=\zero$.}
\ans{ex:quad:B13}{用定理~\ref{thm:posdef} 证明中的分块合同变换（把 $\vv a$ 换成矩阵块）.}
\ans{ex:quad:B14}{主轴定理把椭圆化为 $\lambda_1u^2+\lambda_2v^2\le1$，面积 $\frac\pi{\sqrt{\lambda_1\lambda_2}}$，而正交变换保持面积.}
\ans{ex:quad:B15}{$\vv x^\trans(\vv A^\trans\vv A+\mu\vv I)\vv x=\|\vv{Ax}\|^2+\mu\|\vv x\|^2$. 目标函数对 $\vv x$ 的驻点条件给出正规方程的``岭''版本.}
\ans{ex:quad:B16}{用谱分解，$\vv B=\vv Q\diag(\sqrt{\lambda_i})\vv Q^\trans$；$\vv A^{-1}$ 的特征值为 $\frac1{\lambda_i}$.}

\section{第十二章　*对偶空间与双线性型}

\subsection*{A 组答案}

\ans{ex:bil:A1}{(1) 是. (2) 不是：$2\vv x\mapsto4x_1x_2$. (3) 是. (4) 不是：$f\mapsto2f$ 时值变为 $4$ 倍. (5) 是.}

\ans{ex:bil:A2}{以三个基向量为列作 $\vv P$，对偶基是 $\vv P^{-1}$ 的各行：$\vv\varepsilon_1^*(\vv x)=x_1-x_2$，$\vv\varepsilon_2^*(\vv x)=x_2-x_3$，$\vv\varepsilon_3^*(\vv x)=x_3$.}

\ans{ex:bil:A3}{$f=f(\vv\varepsilon_1)\vv\varepsilon_1^*+f(\vv\varepsilon_2)\vv\varepsilon_2^*+f(\vv\varepsilon_3)\vv\varepsilon_3^*=\vv\varepsilon_1^*+3\vv\varepsilon_2^*+6\vv\varepsilon_3^*$.}

\ans{ex:bil:A4}{$W^0$ 是一维的，基为 $\vv x\mapsto x_1-x_2+x_3$.}

\ans{ex:bil:A5}{$\mathscr T^*$ 的矩阵为 $\begin{pmatrix}1&0&3\\2&1&-1\end{pmatrix}$；$\mathscr T^*$ 把 $\vv y\mapsto y_1+y_2+y_3$ 变为 $\vv x\mapsto4x_1+2x_2$.}

\ans{ex:bil:A6}{标准基下 $\vv A=\begin{pmatrix}1&2\\-1&3\end{pmatrix}$；新基下 $\vv B=\begin{pmatrix}5&-5\\1&3\end{pmatrix}$，可直接验证 $\vv B=\vv P^\trans\vv{AP}$，$\vv P=\begin{pmatrix}1&1\\1&-1\end{pmatrix}$.}

\ans{ex:bil:A7}{对称部分的度量矩阵 $\begin{pmatrix}1&\frac12\\ \frac12&3\end{pmatrix}$，交错部分 $\begin{pmatrix}0&\frac32\\-\frac32&0\end{pmatrix}$.}

\ans{ex:bil:A8}{(1)(2) 非退化（度量矩阵行列式 $-1$）；(3) 退化（度量矩阵全为 $1$，不可逆）.}

\ans{ex:bil:A9}{对应二次型 $x_1^2+2x_1x_2+3x_2^2$. 由二次型得 $f(\vv x,\vv y)=x_1y_1-2x_1y_2-2x_2y_1+2x_2y_2$.}

\ans{ex:bil:A10}{$f(\vv x,\vv x)=0$，$f(\vv x,\vv y)=-f(\vv y,\vv x)$. 度量矩阵 $\begin{pmatrix}0&1\\-1&0\end{pmatrix}$，可逆，非退化.}

\ans{ex:bil:A11}{Lagrange 基：$\frac12(x-1)(x-2)$，$-x(x-2)$，$\frac12x(x-1)$.}

\ans{ex:bil:A12}{$\dim W=2$，$\dim W^0=2$，基为 $\vv x\mapsto x_1+x_2$ 与 $\vv x\mapsto x_3-x_4$.}

\ans{ex:bil:A13}{双线性显然，对称性由 $\tr(\vv{AB})=\tr(\vv{BA})$；非退化：$\tr(\vv A\vv E_{ji})=a_{ij}$. 不是内积：$\vv A=\begin{pmatrix}0&1\\-1&0\end{pmatrix}$ 时 $\tr(\vv A^2)=-2<0$.}

\ans{ex:bil:A14}{坐标为 $\big(\varphi(1),\varphi(x),\varphi(x^2)\big)=\big(1,\frac12,\frac13\big)$.}

\ans{ex:bil:A15}{对任意 $\vv\alpha=\sum a_i\vv e_i$，$\langle\vv\alpha,\vv\beta\rangle=\sum a_if(\vv e_i)=f(\vv\alpha)$. 唯一性：若 $\langle\vv\alpha,\vv\beta-\vv\beta'\rangle=0$ 对一切 $\vv\alpha$ 成立，取 $\vv\alpha=\vv\beta-\vv\beta'$.}

\subsection*{B 组提示}

\ans{ex:bil:B1}{取 $V$ 的一组基，映射的矩阵的第 $i$ 行恰是 $f_i$ 在对偶基下的坐标. 满射当且仅当这个 $k\times n$ 矩阵行满秩，即各行线性无关.}
\ans{ex:bil:B2}{取 $\vv\alpha_0$ 使 $f(\vv\alpha_0)=1$，则 $V=\ker f\oplus\spn(\vv\alpha_0)$，$c=g(\vv\alpha_0)$. 推广：对映射 $\vv\alpha\mapsto(f_1(\vv\alpha),\ldots,f_k(\vv\alpha))$ 用同样的思路，或用零化子的维数.}
\ans{ex:bil:B3}{记 $U=\spn(f_1,\ldots,f_k)\subset V^*$. 把 $f_i$ 写成坐标后，$\ker f_1\cap\cdots\cap\ker f_k$ 就是以这些坐标为行的齐次方程组的解空间，维数为 $\dim V-\dim U$. 方程组语言：齐次方程组只有零解，当且仅当各方程的系数行向量张成整个行向量空间.}
\ans{ex:bil:B4}{由定理~\ref{thm:dualrank}：$\rk\mathscr T^*=\rk\mathscr T$.}
\ans{ex:bil:B5}{$a=c=\frac13$，$b=\frac43$. 积分是 $\mathbb{R}[x]_3$ 上的线性函数，三个求值函数构成 $\mathbb{R}[x]_3^*$ 的基（A11）. 三次时验证 $\int_0^2(x-1)^3\,\mathrm dx=0$ 与右边相等.}
\ans{ex:bil:B6}{双线性型由度量矩阵唯一确定，反之任一矩阵给出双线性型. 后半句即第三章章末习题~\ref{ex:mat:A22}.}
\ans{ex:bil:B7}{由定理~\ref{thm:alternating}，非退化时没有零块，$\dim V=2s$.}
\ans{ex:bil:B8}{度量矩阵 $\begin{pmatrix}1&1&0\\1&1&0\\0&0&1\end{pmatrix}$ 的核为 $\spn\big((1,-1,0)^\trans\big)$.}
\ans{ex:bil:B9}{迷向向量满足 $x_1^2=x_2^2$，即两条直线 $x_2=\pm x_1$. 后半句：设 $\vv P=\begin{pmatrix}a&b\\c&d\end{pmatrix}$，由条件得 $a^2-c^2=1$，$d^2-b^2=1$，$ab=cd$，再令 $a=\cosh t$.}
\ans{ex:bil:B10}{对 $\vv\alpha\in W$，$f\in W^0$，有 $(\mathscr T^*f)(\vv\alpha)=f(\mathscr T\vv\alpha)=0$.}
（须在 %% =====================================================================
%% exercises.tex　《高等代数讲义》章末习题（A/B/C 三组）
%%
%% 用法
%%   1. 在导言区（\begin{document} 之前、全书统一记号之后）写 %% =====================================================================
%% exercises.tex　《高等代数讲义》章末习题（A/B/C 三组）
%%
%% 用法
%%   1. 在导言区（\begin{document} 之前、全书统一记号之后）写 \input{exercises}.
	%%   2. 在每章末尾（原 \begin{problemset}...\end{problemset} 之后，或替换它）
	%%      写对应的命令：
	%%        第一章 \ExElimination      第二章 \ExVectorSpace   第三章 \ExMatrix
	%%        第四章 \ExLinearMap        第五章 \ExDeterminant   第六章 \ExSolutions
	%%        第七章 \ExPolynomial       第八章 \ExEigen         第九章 \ExInnerProduct
	%%        第十章 \ExDecomposition    第十一章 \ExQuadratic   第十二章 \ExBilinear
	%%   3. 分组标准
	%%        A 组：直接运用本章定义与定理的计算和证明；
	%%        B 组：需要综合多个结论、或有一定技巧的计算和证明；
	%%        C 组：有挑战性的题目，或超出本课程大纲、为继续深造准备的内容.
	%%   4. 题目与正文中已有的节末练习、原 problemset 不重复；个别题目注明``接正文某题''.
	%%      第六章按扩充版（ch06-solutions.tex）的内容编写.
	%%   5. 每道题都带标签 ex:<章键>:<组号>，正文、习题与附录答案一律用 \ref 引用，
	%%      渲染为``A4''``B12''这样的题号. 章键依次为
	%%        elim vs mat lmap det sol poly eig ip dec quad bil（第一至十二章）.
	%%      例：第一章章末习题~\ref{ex:elim:B12}. 新增题目请加在组末，并补上标签.
	%%      A 组答案、B 组提示在 answers.tex 中，按同一标签对应.
	%% 依赖：导言区的 \vv \zero \trans \spn \rk \imm \diag，以及 enumitem（elegantbook 已加载）.
	%% 引用的前文标签：thm:sameSolution thm:solvability prop:linspc eg:pricedecomp cor:homogeneous thm:rouche thm:sol-determinable
	%%   thm:spectral thm:posdef thm:svd thm:orthcanon
	%% =====================================================================
	
	\providecommand{\tr}{\operatorname{tr}}
	\newcommand{\exsetheading}{%
		\section*{章末习题}%
		\markright{章末习题}%
	}
\newcommand{\exgroup}[2]{%
	\par\medskip{\centering\sffamily\bfseries #1\quad #2\par}\smallskip}
	\newenvironment{exA}{\begin{enumerate}[label=\textbf{A\arabic*.},ref=A\arabic*,leftmargin=2.8em,itemsep=4pt]}{\end{enumerate}}
	\newenvironment{exB}{\begin{enumerate}[label=\textbf{B\arabic*.},ref=B\arabic*,leftmargin=2.8em,itemsep=4pt]}{\end{enumerate}}
	\newenvironment{exC}{\begin{enumerate}[label=\textbf{C\arabic*.},ref=C\arabic*,leftmargin=2.8em,itemsep=4pt]}{\end{enumerate}}
	\newcommand{\Agroup}{\exgroup{A 组}{基础巩固}}
	\newcommand{\Bgroup}{\exgroup{B 组}{综合提升}}
	\newcommand{\Cgroup}{\exgroup{C 组}{拓展探索}}
	
	% =====================================================================
	% 第一章　线性方程组与消元法
	% =====================================================================
	\newcommand{\ExElimination}{%
		\exsetheading
		\Agroup
		\begin{exA}
			\item\label{ex:elim:A1} 判断下列矩阵哪些是阶梯形，哪些是简化阶梯形：
			\[
			\begin{pmatrix}1&2&0\\0&0&1\end{pmatrix},\quad
			\begin{pmatrix}1&0&3\\0&2&1\\0&0&0\end{pmatrix},\quad
			\begin{pmatrix}0&1&0\\1&0&0\end{pmatrix},\quad
			\begin{pmatrix}1&0&2&0\\0&1&-1&0\\0&0&0&1\end{pmatrix},\quad
			\begin{pmatrix}0&0\\0&1\end{pmatrix}.
			\]
			\item\label{ex:elim:A2} 写出方程组 $3x_1-x_3=2$，$x_2+4x_3=0$，$x_1-x_2=5$ 的增广矩阵；反过来，写出以 $\begin{pmatrix}1&-2&0&3\\0&1&5&-1\end{pmatrix}$ 为增广矩阵的方程组.
			\item\label{ex:elim:A3} 用消元法解方程组 $\left\{\begin{array}{l}x_1+2x_2+3x_3=14,\\2x_1-x_2+x_3=3,\\3x_1+x_2-x_3=2.\end{array}\right.$
			\item\label{ex:elim:A4} 用消元法解方程组 $\left\{\begin{array}{l}x_1+x_2+x_3+x_4=2,\\x_1-x_2+x_3-x_4=4,\\2x_1+x_2-x_3+3x_4=-1,\\x_1+2x_2+3x_3-2x_4=5.\end{array}\right.$
			\item\label{ex:elim:A5} 判断方程组 $\left\{\begin{array}{l}x_1-x_2+x_3=1,\\2x_1+x_2-x_3=2,\\x_1+2x_2-2x_3=3\end{array}\right.$ 是否有解，并说明消元过程中哪一行``暴露''了结论.
			\item\label{ex:elim:A6} 求方程组 $\left\{\begin{array}{l}x_1+2x_2-x_3+x_4=3,\\2x_1+4x_2-x_3+3x_4=7,\\-x_1-2x_2+2x_3=-2\end{array}\right.$ 的全部解，并指出主变量与自由变量.
			\item\label{ex:elim:A7} 求齐次方程组 $\left\{\begin{array}{l}x_1+x_2-x_3-x_4=0,\\2x_1-5x_2+3x_3+2x_4=0,\\7x_1-7x_2+3x_3+x_4=0\end{array}\right.$ 的全部解.
			\item\label{ex:elim:A8} 把矩阵 $\begin{pmatrix}2&4&-2&6\\1&2&1&1\\3&6&1&5\end{pmatrix}$ 化为简化阶梯形，并写出每一步所用的初等行变换.
			\item\label{ex:elim:A9} 写出全部 $2\times2$ 简化阶梯形矩阵的``形状''（可以含参数），并说明共有几类.
			\item\label{ex:elim:A10} $k$ 取何值时，方程组 $x+y=1$，$x-y=3$，$2x+ky=5$ 有解？
			\item\label{ex:elim:A11} 设 $a,b$ 为常数. 方程组 $x+2y=a$，$3x+6y=b$ 何时有解？有解时写出全部解.
			\item\label{ex:elim:A12} 求过三点 $(1,2)$，$(2,3)$，$(3,6)$ 的抛物线 $y=ax^2+bx+c$.
			\item\label{ex:elim:A13} 求过四点 $(-1,2)$，$(0,1)$，$(1,0)$，$(2,5)$ 的三次多项式函数 $y=a+bx+cx^2+dx^3$.
			\item\label{ex:elim:A14} 某剧场三种票价分别为 $30$，$50$，$80$ 元. 某场演出共售出 $100$ 张票，票房 $4800$ 元，且 $80$ 元票的张数恰为 $30$ 元票的一半. 三种票各售出多少张？
			\item\label{ex:elim:A15} 配平化学方程式 $x_1\,\mathrm{C_3H_8}+x_2\,\mathrm{O_2}\to x_3\,\mathrm{CO_2}+x_4\,\mathrm{H_2O}$：对 C，H，O 三种原子分别写出守恒方程，解这个齐次方程组，并取最小的正整数解.
			\item\label{ex:elim:A16} 某城区四个路口 $A,B,C,D$ 由单行道依次连成环路 $A\to B\to C\to D\to A$，环路上的车流量依次记为 $x_1$（$A\to B$），$x_2$（$B\to C$），$x_3$（$C\to D$），$x_4$（$D\to A$）（单位：辆/小时）. 各路口与外部道路的进出量为：$A$ 驶入 $400$、驶出 $300$；$B$ 驶入 $200$、驶出 $350$；$C$ 驶入 $250$、驶出 $100$；$D$ 驶入 $150$、驶出 $250$. 设每个路口驶入总量等于驶出总量.
			\begin{enumerate}
				\item 列出方程组并求全部解；
				\item 若各路段车流量都不能为负，$x_1$ 最少是多少？
				\item 若把 $D$ 的驶出量改为 $350$，方程组会怎样？这反映了什么守恒律？
			\end{enumerate}
			\item\label{ex:elim:A17} 证明：二元一次方程组 $ax+by=e$，$cx+dy=f$ 对任意的 $e,f$ 都有唯一解，当且仅当 $ad-bc\ne0$.（只用消元法，不用行列式.）
			\item\label{ex:elim:A18} 证明：若线性方程组有两个不同的解 $\vv x_1,\vv x_2$，则对任意数 $t$，$\vv x_1+t(\vv x_2-\vv x_1)$ 都是解. 由此说明：线性方程组的解要么不存在，要么唯一，要么有无穷多个.
			\item\label{ex:elim:A19} 证明方程组 $\{x+y+z=3,\ x-y=0\}$ 与 $\{2x+z=3,\ x+y+z=3\}$ 同解.
			\item\label{ex:elim:A20} 举例说明：若\textbf{同时}做``第一个方程加到第二个方程上''与``第二个方程加到第一个方程上''（两步都用原来的方程），所得方程组可能与原方程组不同解. 这与定理~\ref{thm:sameSolution}（初等行变换不改变解集）矛盾吗？
			\item\label{ex:elim:A21} 已知一个方程组的增广矩阵化为阶梯形后为 $\begin{pmatrix}1&3&-1&2&0\\0&0&2&1&4\\0&0&0&0&a\end{pmatrix}$. 讨论 $a$ 取何值时方程组有解；有解时它有几个自由变量？
		\end{exA}
		
		\Bgroup
		\begin{exB}
			\item\label{ex:elim:B1} 讨论 $a,b$ 取何值时，方程组
			\[\left\{\begin{array}{l}x_1+x_2+x_3+x_4=0,\\x_2+2x_3+2x_4=1,\\-x_2+(a-3)x_3-2x_4=b,\\3x_1+2x_2+x_3+ax_4=-1\end{array}\right.\]
			有唯一解、无解、有无穷多解；有无穷多解时求出全部解.
			\item\label{ex:elim:B2} 讨论 $\lambda$ 取何值时，方程组
			$(1+\lambda)x_1+x_2+x_3=0$，$x_1+(1+\lambda)x_2+x_3=3$，$x_1+x_2+(1+\lambda)x_3=\lambda$
			有唯一解、无解、有无穷多解.
			\item\label{ex:elim:B3} 设 $n$ 个方程、$n$ 个未知量的方程组对应的齐次方程组只有零解. 证明：无论常数项取什么值，原方程组都有唯一解.（提示：考察简化阶梯形中主元的个数.）
			\item\label{ex:elim:B4} 证明：线性方程组无解，当且仅当可以把若干个方程分别乘以适当的数再相加，得到矛盾方程 $0=1$.
			\item\label{ex:elim:B5} 证明：交换两行这一初等变换，可以由``倍加''与``某行乘以 $-1$''这两类变换的组合实现. 由此说明三类初等行变换并不是``独立''的.
			\item\label{ex:elim:B6} 证明：若线性方程组的系数和常数项都是有理数，且它有实数解，则它一定有有理数解；若它的实数解唯一，则这个解就是有理数解.
			\item\label{ex:elim:B7} 讨论方程组
			\[x_1-x_2=a_1,\quad x_2-x_3=a_2,\quad\ldots,\quad x_{n-1}-x_n=a_{n-1},\quad x_n-x_1=a_n\]
			有解的充要条件，并在有解时求全部解.
			\item\label{ex:elim:B8} 讨论方程组
			\[x_1+x_2=a_1,\quad x_2+x_3=a_2,\quad\ldots,\quad x_{n-1}+x_n=a_{n-1},\quad x_n+x_1=a_n\]
			的解. 证明：$n$ 为奇数时解唯一；$n$ 为偶数时有解当且仅当 $a_1-a_2+a_3-\cdots-a_n=0$.
			\item\label{ex:elim:B9} 某经济体有三个部门，直接消耗系数矩阵与最终需求分别为
			\[\vv A=\begin{pmatrix}0.1&0.2&0.2\\0.2&0.1&0.2\\0.1&0.3&0.1\end{pmatrix},\qquad \vv d=\begin{pmatrix}20\\130\\65\end{pmatrix}.\]
			按第一章``投入产出模型''的方式列出方程组，求各部门的总产出.
			\item\label{ex:elim:B10} 估计用消元法（化为阶梯形后回代）解 $n$ 个方程、$n$ 个未知量的方程组所需的乘除法次数，证明它约为 $\frac13n^3$（即与 $\frac13n^3$ 之比在 $n\to\infty$ 时趋于 $1$）.
			\item\label{ex:elim:B11} 设三个平面 $a_ix+b_iy+c_iz=d_i$（$i=1,2,3$）两两相交. 用消元法的语言说明：它们恰好交于一点、交于一条直线、或没有公共点（三条交线互相平行），并各举一例.
			\item\label{ex:elim:B12} 证明：$n$ 元齐次方程组的解集对加法和数乘封闭；而非齐次方程组的两个解之差是对应齐次方程组的解. 由此说明：非齐次方程组的全部解，等于它的任一个解加上齐次方程组的全部解.
			\item\label{ex:elim:B13} 一家律所有甲、乙、丙三位律师，本月分别承办了若干件合同、劳动、知识产权纠纷. 已知：三类案件合计 $30$ 件；甲承办的数量 $=$ 乙 $+$ 丙；每件合同、劳动、知识产权案件分别需要 $10$，$6$，$15$ 个工时，三人本月共用 $300$ 个工时……
			这些条件够不够确定每位律师承办了多少案件？请自己补全一个合理的设定（明确未知量是什么），列出方程组，用自由变量的个数说明还缺多少条独立的信息.
			\item\label{ex:elim:B14} 设方程组的系数矩阵每一行的元素之和都等于 $0$. 证明 $x_1=x_2=\cdots=x_n=1$ 是对应齐次方程组的解；若非齐次方程组有解 $\vv x_0$，则对任意数 $t$，$\vv x_0+t(1,\ldots,1)^\trans$ 也是解.
			\item\label{ex:elim:B15} 证明：如果方程组 (I) 的每个方程都是方程组 (II) 各方程的``线性组合''（各乘一个数再相加），反之亦然，那么 (I) 与 (II) 同解. 逆命题对于有解的方程组成立吗？先在两个未知量的情形下思考.
			\item\label{ex:elim:B16} 讨论 $\lambda$ 取何值时，方程组 $\lambda x+y+z=1$，$x+\lambda y+z=\lambda$，$x+y+\lambda z=\lambda^2$ 有唯一解、无解、有无穷多解；有唯一解时求出解.（第六章例~\ref{eg:lambda} 将用秩与行列式重新组织这个讨论.）
		\end{exB}
		
		\Cgroup
		\begin{exC}
			\item\label{ex:elim:C1} （百鸡问题）《张丘建算经》（约公元 5 世纪）：``今有鸡翁一，直钱五；鸡母一，直钱三；鸡雏三，直钱一. 凡百钱买鸡百只，问鸡翁、母、雏各几何？''列出线性方程组，求它的全部解，再从中找出所有\textbf{非负整数}解. 线性代数在哪一步之后就无能为力了？
			\item\label{ex:elim:C2} （幻方）若一个 $3\times3$ 数表的每行、每列和两条对角线上的三个数之和都等于同一个数 $s$，称它为幻方（不要求数互不相同）. 把 $9$ 个元素作为未知量、$s$ 也作为未知量，写出方程组，证明：
			\begin{enumerate}
				\item 中心元素一定等于 $s/3$；
				\item 全部幻方可以用 $3$ 个自由参数表示，写出通式；
				\item 用 $1,2,\ldots,9$ 各一次填成的幻方本质上只有一个（其余由旋转、翻转得到）.
			\end{enumerate}
			\item\label{ex:elim:C3} （点灯游戏）$3\times3$ 方格中每格有一盏灯，按下某格会使该格及其上下左右相邻的格子开关状态翻转. 在``$1+1=0$''的算术（即模 $2$ 运算）中，把``按哪些格''作为未知量，把``从某一状态变为全灭''写成线性方程组. 证明：从任意初始状态出发，总能把灯全部熄灭.（模 $2$ 运算的集合 $\{0,1\}$ 不是本书意义下的数域，但消元法照样适用——想一想为什么.）
			\item\label{ex:elim:C4} （舍入误差与主元）在``只保留三位有效数字''的计算机上解方程组 $0.0001x+y=1$，$x+y=2$：
			\begin{enumerate}
				\item 直接用第一个方程消去第二个方程中的 $x$；
				\item 先交换两个方程再消元.
			\end{enumerate}
			把两种结果与精确解比较，说明数值软件为什么要``选绝对值最大的元素作主元''.
			\item\label{ex:elim:C5} 设两个都有解的方程组 (I)(II) 都有 $m$ 个方程、$n$ 个未知量，且同解. 证明：它们的增广矩阵行等价.（提示：先证明二者的简化阶梯形的非零行完全相同.）
			\item\label{ex:elim:C6} （Jacobi 迭代）对方程组 $4x+y=9$，$x+3y=5$，把它改写为 $x=(9-y)/4$，$y=(5-x)/3$，从 $(x_0,y_0)=(0,0)$ 出发反复代入. 计算前五步，观察是否收敛到真解. 再对 $x+4y=9$，$3x+y=5$ 做同样的事，会发生什么？试找出一个保证收敛的充分条件.
			\item\label{ex:elim:C7} （电路）一个电路中有三条支路电流 $I_1,I_2,I_3$. 节点 $A$ 处 $I_1=I_2+I_3$；两个回路的电压方程为 $2I_1+I_2=8$，$I_2-2I_3=0$（单位从略）. 求各电流. 再说明：若把第三个节点 $B$ 的电流方程也写进来，方程组为什么不会因此``多一个条件''？
			\item\label{ex:elim:C8} 证明：一个 $m\times n$ 矩阵（$m,n\ge1$）通过初等行变换所能得到的不同的简化阶梯形``形状''（即主元所在列的位置组合）共有 $\sum_{r=0}^{\min(m,n)}\binom nr$ 种.
			\item\label{ex:elim:C9} 证明：一个矩阵的简化阶梯形是唯一的.（提示：设 $\vv R_1,\vv R_2$ 都是 $\vv M$ 的简化阶梯形，则以它们为系数矩阵的齐次方程组同解. 对列数归纳：删去最后一列后由归纳假设前面各列相同，再比较最后一列，考察它在两者中是否为主元列.）
		\end{exC}
	}
	
	% =====================================================================
	% 第二章　向量空间
	% =====================================================================
	\newcommand{\ExVectorSpace}{%
		\exsetheading
		\Agroup
		\begin{exA}
			\item\label{ex:vs:A1} 下列集合关于通常的四则运算是否构成数域？
			(1) $\{a+b\sqrt3\mid a,b\in\mathbb{Q}\}$；(2) $\{a+b\sqrt2\mid a\in\mathbb{Q},\ b\in\mathbb{Z}\}$；(3) 全体偶数；(4) $\{a+b\sqrt[3]{2}\mid a,b\in\mathbb{Q}\}$；(5) 全体纯虚数与 $0$.
			\item\label{ex:vs:A2} 下列集合关于通常的加法与数乘是否构成实向量空间？
			(1) 全体 $n$ 阶实对称矩阵；(2) $\{(x,y)^\trans\in\mathbb{R}^2\mid x\ge0\}$；(3) 次数恰为 $n$ 的实系数多项式全体；(4) $\{f\in\mathbb{R}[x]\mid f(1)=0\}$；(5) 满足 $a_{n+2}=a_{n+1}+a_n$ 的全体实数列；(6) $\{(x,y,z)^\trans\mid x+y+z=1\}$.
			\item\label{ex:vs:A3} 判断下列向量组的线性相关性：
			\begin{enumerate}
				\item $(1,2,-1)^\trans$，$(2,-3,1)^\trans$，$(4,1,-1)^\trans$；
				\item 在 $\mathbb{R}[x]_3$ 中：$1+x$，$1-x$，$x^2$；
				\item 在 $\mathbb{R}[x]_3$ 中：$1+x+x^2$，$1-x$，$2+x^2$，$x+x^2$；
				\item 在 $M_2(\mathbb{R})$ 中：$\begin{pmatrix}1&0\\0&1\end{pmatrix}$，$\begin{pmatrix}0&1\\1&0\end{pmatrix}$，$\begin{pmatrix}1&1\\1&1\end{pmatrix}$，$\begin{pmatrix}1&0\\0&0\end{pmatrix}$.
			\end{enumerate}
			\item\label{ex:vs:A4} 把 $\vv\beta=(1,2,1,1)^\trans$ 表示为 $\vv\alpha_1=(1,1,1,1)^\trans$，$\vv\alpha_2=(1,1,-1,-1)^\trans$，$\vv\alpha_3=(1,-1,1,-1)^\trans$，$\vv\alpha_4=(1,-1,-1,1)^\trans$ 的线性组合.
			\item\label{ex:vs:A5} $t$ 取何值时，$(1,1,0)^\trans$，$(1,3,-1)^\trans$，$(5,3,t)^\trans$ 线性相关？
			\item\label{ex:vs:A6} 求向量组 $\vv\alpha_1=(1,-1,2,4)^\trans$，$\vv\alpha_2=(0,3,1,2)^\trans$，$\vv\alpha_3=(3,0,7,14)^\trans$，$\vv\alpha_4=(1,-1,2,0)^\trans$，$\vv\alpha_5=(2,1,5,6)^\trans$ 的秩和一个极大无关组，并把其余向量用这个极大无关组线性表出.
			\item\label{ex:vs:A7} 设 $\vv\alpha_1,\vv\alpha_2,\vv\alpha_3$ 线性无关. 判断下列向量组的线性相关性：
			(1) $\vv\alpha_1-\vv\alpha_2$，$\vv\alpha_2-\vv\alpha_3$，$\vv\alpha_3-\vv\alpha_1$；
			(2) $\vv\alpha_1$，$\vv\alpha_1+\vv\alpha_2$，$\vv\alpha_1+\vv\alpha_2+\vv\alpha_3$；
			(3) $\vv\alpha_1+2\vv\alpha_2$，$\vv\alpha_2+2\vv\alpha_3$，$\vv\alpha_3+2\vv\alpha_1$.
			\item\label{ex:vs:A8} 证明：含有零向量的向量组线性相关；含有两个相同向量的向量组线性相关；单个向量 $\vv\alpha$ 线性无关当且仅当 $\vv\alpha\ne\zero$.
			\item\label{ex:vs:A9} 证明：两个向量线性相关，当且仅当其中一个是另一个的倍数. 三个向量的情形能这样说吗？
			\item\label{ex:vs:A10} 只用向量空间的八条公理与命题~\ref{prop:linspc}，证明：(1) $\vv\alpha+\vv\alpha=2\vv\alpha$；(2) 若 $k\ne0$ 且 $k\vv\alpha=k\vv\beta$，则 $\vv\alpha=\vv\beta$；(3) $-(-\vv\alpha)=\vv\alpha$.
			\item\label{ex:vs:A11} 证明下列集合是相应空间的子空间，并求一组基和维数：
			\begin{enumerate}
				\item $\{\vv x\in\mathbb{R}^4\mid x_1+x_2=x_3+x_4\}$；
				\item $n$ 阶上三角矩阵全体；
				\item 迹为零的 $n$ 阶方阵全体 $\{\vv X\mid x_{11}+\cdots+x_{nn}=0\}$；
				\item $\{f\in\mathbb{R}[x]_4\mid f(1)=0\}$.
			\end{enumerate}
			\item\label{ex:vs:A12} 判断下列子集是否为子空间：(1) $\{(a,b,c)^\trans\mid a=2b\}$；(2) $\{(a,b,c)^\trans\mid ab=0\}$；(3) $\{(a,b,c)^\trans\mid a,b,c\in\mathbb{Q}\}\subset\mathbb{R}^3$；(4) $\{f\in C[0,1]\mid f(0)=2f(1)\}$.
			\item\label{ex:vs:A13} 求 $\vv\alpha_1=(1,2,1,0)^\trans$，$\vv\alpha_2=(-1,1,1,1)^\trans$，$\vv\alpha_3=(0,3,2,1)^\trans$，$\vv\alpha_4=(1,0,1,2)^\trans$ 张成的子空间的维数和一组基.
			\item\label{ex:vs:A14} 求多项式 $2x^2-x+3$ 在基 $1,\ x-1,\ (x-1)^2$ 下的坐标.
			\item\label{ex:vs:A15} 求矩阵 $\begin{pmatrix}2&3\\-1&4\end{pmatrix}$ 在 $M_2(\mathbb{R})$ 的基 $\begin{pmatrix}1&0\\0&0\end{pmatrix},\begin{pmatrix}1&1\\0&0\end{pmatrix},\begin{pmatrix}1&1\\1&0\end{pmatrix},\begin{pmatrix}1&1\\1&1\end{pmatrix}$ 下的坐标.
			\item\label{ex:vs:A16} (1) $\mathbb{R}$ 是否是 $\mathbb{Q}$ 上的向量空间？(2) $\mathbb{Q}$ 是否是 $\mathbb{R}$ 上的向量空间？(3) $\mathbb{C}$ 作为 $\mathbb{R}$ 上的向量空间和作为 $\mathbb{C}$ 上的向量空间，维数各是多少？
			\item\label{ex:vs:A17} 写出 $\mathbb{F}[x]_n$ 与 $\mathbb{F}^n$ 之间、$M_2(\mathbb{F})$ 与 $\mathbb{F}^4$ 之间的一个同构映射，并验证它保持线性运算.
			\item\label{ex:vs:A18} 设 $W_1=\spn\big((1,0,0)^\trans,(0,1,0)^\trans\big)$，$W_2=\spn\big((1,1,1)^\trans\big)$，$W_3=\spn\big((1,1,0)^\trans\big)$. 分别判断 $W_1+W_2$，$W_1+W_3$ 是否为直和，并求和空间的维数.
			\item\label{ex:vs:A19} 设 $W_1=\spn\big((1,1,0,0)^\trans,(1,0,1,1)^\trans\big)$，$W_2=\spn\big((0,-1,1,1)^\trans,(1,2,0,1)^\trans\big)$. 求 $W_1+W_2$ 与 $W_1\cap W_2$ 的维数，并验证维数公式.
			\item\label{ex:vs:A20} 设 $\vv\alpha_1,\ldots,\vv\alpha_s$ 的秩为 $r$. 证明：其中任意 $r+1$ 个向量线性相关；任意 $r$ 个线性无关的向量都构成极大无关组.
		\end{exA}
		
		\Bgroup
		\begin{exB}
			\item\label{ex:vs:B1} 证明：两个子空间的并 $W_1\cup W_2$ 是子空间，当且仅当 $W_1\subset W_2$ 或 $W_2\subset W_1$.
			\item\label{ex:vs:B2} 证明：若向量组 $\vv\beta_1,\ldots,\vv\beta_t$ 可由 $\vv\alpha_1,\ldots,\vv\alpha_s$ 线性表出，且两组的秩相等，则两组等价.
			\item\label{ex:vs:B3} 设向量组 (I) 与 (II) 的秩分别为 $r_1,r_2$，把它们合并为 (III). 证明 $\max\{r_1,r_2\}\le\rk(\mathrm{III})\le r_1+r_2$. 再证明：$\rk(\vv\alpha_1+\vv\beta_1,\ldots,\vv\alpha_s+\vv\beta_s)\le\rk(\vv\alpha_1,\ldots,\vv\alpha_s)+\rk(\vv\beta_1,\ldots,\vv\beta_s)$.
			\item\label{ex:vs:B4} 设 $\vv\alpha_1,\ldots,\vv\alpha_n$ 是 $V$ 的基，$\vv\beta=a_1\vv\alpha_1+\cdots+a_n\vv\alpha_n$. 证明：把 $\vv\alpha_j$ 换成 $\vv\beta$ 后所得的向量组仍是基，当且仅当 $a_j\ne0$.
			\item\label{ex:vs:B5} 设 $p_0,p_1,\ldots,p_{n-1}\in\mathbb{F}[x]$ 满足 $\deg p_k=k$. 证明它们构成 $\mathbb{F}[x]_n$ 的一组基.
			\item\label{ex:vs:B6} 记 $V$ 为满足 $a_{n+2}=a_{n+1}+2a_n$（$n\ge1$）的全体实数列. 证明 $V$ 是向量空间，$\dim V=2$，并说明 $\{2^n\}$ 与 $\{(-1)^n\}$ 构成 $V$ 的一组基. 由此写出 $a_1=1$，$a_2=1$ 时的通项.
			\item\label{ex:vs:B7} 在 $C(\mathbb{R})$ 中：(1) 证明 $\mathrm e^x,\mathrm e^{2x},\mathrm e^{3x}$ 线性无关；(2) 证明 $\sin x,\cos x,\sin(x+1)$ 线性相关；(3) $1,\cos x,\cos2x,\cos^2x$ 的秩是多少？
			\item\label{ex:vs:B8} 证明 $\{a+b\sqrt2+c\sqrt3+d\sqrt6\mid a,b,c,d\in\mathbb{Q}\}$ 是数域，并证明 $1,\sqrt2,\sqrt3,\sqrt6$ 在 $\mathbb{Q}$ 上线性无关. 这个数域作为 $\mathbb{Q}$ 上的向量空间是几维的？
			\item\label{ex:vs:B9} 写出 $\mathbb{R}^2$ 与 $\mathbb{R}^3$ 的全部子空间（按维数分类，并说明几何意义）.
			\item\label{ex:vs:B10} 证明：$\mathbb{R}^5$ 中任意两个三维子空间的交至少是一维的. 一般地，$n$ 维空间中维数分别为 $k_1,k_2$ 的两个子空间，若 $k_1+k_2>n$，则交不为 $\{\zero\}$.
			\item\label{ex:vs:B11} 设 $a\in\mathbb{F}$. 证明 $\mathbb{F}[x]_n=\mathbb{F}\cdot1\oplus\{f\in\mathbb{F}[x]_n\mid f(a)=0\}$，并写出 $f$ 按这个直和的分解.
			\item\label{ex:vs:B12} 设 $V=C[-1,1]$，$E$，$O$ 分别为其中的偶函数、奇函数全体. 证明 $V=E\oplus O$，并把 $\mathrm e^x$ 按此分解.
			\item\label{ex:vs:B13} 设 $\dim V=n$，$W_1,\ldots,W_k$ 是 $V$ 的子空间，且 $\dim W_1+\cdots+\dim W_k>(k-1)n$. 证明 $W_1\cap\cdots\cap W_k\ne\{\zero\}$.
			\item\label{ex:vs:B14} 设 $\vv\alpha_1,\ldots,\vv\alpha_n$ 是 $V$ 的基，$W=\spn(\vv\alpha_1+\vv\alpha_2,\ \vv\alpha_2+\vv\alpha_3,\ \ldots,\ \vv\alpha_{n-1}+\vv\alpha_n)$. 求 $\dim W$，并求 $W$ 的一个补空间.
			\item\label{ex:vs:B15} 设 $V_1=\{\vv x\in\mathbb{F}^n\mid x_1=x_2=\cdots=x_n\}$，$V_2=\{\vv x\in\mathbb{F}^n\mid x_1=0\}$. 证明 $\mathbb{F}^n=V_1\oplus V_2$. 与第二章例~\ref{eg:pricedecomp} 的分解相比，哪一种分解更``自然''？（提示：思考补空间的不唯一性.）
			\item\label{ex:vs:B16} 在 $M_n(\mathbb{F})$ 中，设 $W_1$ 为对角线以下元素全为零的矩阵（上三角矩阵）全体，$W_2$ 为对角线及以上元素全为零的矩阵全体. 证明 $M_n(\mathbb{F})=W_1\oplus W_2$，并求 $\dim W_1$.
			\item\label{ex:vs:B17} （覆盖定理的另一个证明）设 $W_1,\ldots,W_s$ 是 $\mathbb{R}^n$ 的真子空间，令 $\vv v(t)=(1,t,t^2,\ldots,t^{n-1})^\trans$.
			\begin{enumerate}
				\item 证明：每个真子空间都包含于某个超平面 $\{\vv x\mid a_1x_1+\cdots+a_nx_n=0\}$（$a_i$ 不全为零）中.（提示：取 $W$ 的基并扩充为 $\mathbb{R}^n$ 的基，看``最后一个坐标''.）
				\item 证明：对上述超平面，满足 $\vv v(t)$ 落在其中的实数 $t$ 至多有 $n-1$ 个.
				\item 由此推出：$\vv v(t)$ 落在 $W_1\cup\cdots\cup W_s$ 中的 $t$ 至多有 $s(n-1)$ 个. 这给出了覆盖定理在 $\mathbb{R}^n$ 中的一个``构造性''证明. 对 $n=3$ 以及三个子空间 $x_1=x_2$，$x_2=x_3$，$x_1+x_2+x_3=0$，写出一个不在它们之中的 $\vv v(t)$.
			\end{enumerate}
			\item\label{ex:vs:B18} 设 $\dim V=n$，$W_1,\ldots,W_s$ 是 $V$ 的真子空间. 证明：存在 $V$ 的一组基 $\vv\varepsilon_1,\ldots,\vv\varepsilon_n$，其中\textbf{每个}向量都不属于 $W_1\cup\cdots\cup W_s$.（提示：逐个选取. 选 $\vv\varepsilon_k$ 时，要躲开的除了 $W_1,\ldots,W_s$，还有哪个真子空间？）
		\end{exB}
		
		\Cgroup
		\begin{exC}
			\item\label{ex:vs:C1} （覆盖定理在有限``数系''上的失效）设 $p$ 是素数，$\mathbb{F}_p=\{0,1,\ldots,p-1\}$ 按模 $p$ 作加法和乘法（可以承认它对加、减、乘以及除以非零元素都封闭. 它不包含 $\mathbb{Q}$，所以不是本书意义下的数域；但第二章中除定理~\ref{thm:Qsmallest} 与覆盖定理以外的定义和结论，在它上面照样成立）.
			\begin{enumerate}
				\item 证明 $\mathbb{F}_p^2$ 恰有 $p+1$ 个一维子空间，并且它们的并是整个 $\mathbb{F}_p^2$.
				\item 证明 $\mathbb{F}_p^2$ 不能写成 $p$ 个真子空间的并.（提示：每个一维子空间恰含 $p-1$ 个非零向量，数一数.）
				\item 指出定理~\ref{thm:cover} 的证明在 $\mathbb{F}_p$ 上从哪一步开始失效. 由第 2 问，``个数不超过 $p$ 时盖不住''，你能把正文的证明改写成这个``定量''的版本吗？
			\end{enumerate}
			\item\label{ex:vs:C2} 证明 $\ln2,\ln3,\ln5$ 在 $\mathbb{Q}$ 上线性无关. 由此说明 $\mathbb{R}$ 作为 $\mathbb{Q}$ 上的向量空间不是有限维的.（提示：算术基本定理.）
			\item\label{ex:vs:C3} 证明：全体实数列构成的向量空间不是有限维的；而``只有有限项非零''的实数列全体是它的子空间，且 $\vv e_1,\vv e_2,\ldots$（第 $i$ 项为 $1$，其余为 $0$）是它的一组``基''——每个元素都能唯一地写成其中\textbf{有限个}向量的线性组合. 这组向量为什么不是全体实数列的基？
			\item\label{ex:vs:C4} （模律）设 $W_1,W_2,W_3$ 是 $V$ 的子空间. 证明 $(W_1\cap W_3)+(W_2\cap W_3)\subset(W_1+W_2)\cap W_3$，在 $\mathbb{R}^2$ 中举例说明等号可以不成立；再证明：若 $W_1\subset W_3$，则等号一定成立.
			\item\label{ex:vs:C5} （有限域上的向量空间）在只含 $0,1$ 两个元素、按模 $2$ 运算的``数系'' $\mathbb{F}_2$ 上考虑 $\mathbb{F}_2^3$. (1) 它有多少个向量？(2) 有多少个一维子空间、多少个二维子空间？(3) 它有多少组（有序）基？推广到 $\mathbb{F}_2^n$.
			\item\label{ex:vs:C6} （塔律）设 $\mathbb{F}\subset\mathbb{K}\subset\mathbb{L}$ 是三个数域，$\mathbb{L}$ 作为 $\mathbb{K}$ 上的向量空间是 $m$ 维的，$\mathbb{K}$ 作为 $\mathbb{F}$ 上的向量空间是 $n$ 维的. 证明 $\mathbb{L}$ 作为 $\mathbb{F}$ 上的向量空间是 $mn$ 维的. 用它重新计算 B8 中数域的维数.
			\item\label{ex:vs:C7} 证明 $\mathbb{Q}(\sqrt2+\sqrt3)=\mathbb{Q}(\sqrt2,\sqrt3)$，即包含 $\sqrt2+\sqrt3$ 的最小数域与包含 $\sqrt2,\sqrt3$ 的最小数域相同.（提示：计算 $(\sqrt2+\sqrt3)^3$.）
			\item\label{ex:vs:C8} 设 $V$ 是 $n$ 维空间，$W$ 是 $k$ 维子空间（$0<k<n$）. 证明 $W$ 有无穷多个补空间. 在 $\mathbb{F}_2^2$ 中，一维子空间有几个补空间？
			\item\label{ex:vs:C9} （Fibonacci 空间的另一组基）在 B6 的记号下，改为研究满足 $a_{n+2}=a_{n+1}+a_n$ 的数列空间. 求出两个公比 $q$ 使 $\{q^n\}$ 属于这个空间，证明它们构成一组基，并推出 Fibonacci 数列的通项（Binet 公式）. 第八章会从另一个角度再次得到它.
		\end{exC}
	}
	
	% =====================================================================
	% 第三章　矩阵
	% =====================================================================
	\newcommand{\ExMatrix}{%
		\exsetheading
		\Agroup
		\begin{exA}
			\item\label{ex:mat:A1} 设 $\vv A=\begin{pmatrix}1&-1&2\\0&3&1\end{pmatrix}$，$\vv B=\begin{pmatrix}2&1\\-1&0\\1&4\end{pmatrix}$. 计算 $\vv{AB}$，$\vv{BA}$，$\vv A^\trans\vv B^\trans$，并验证 $(\vv{AB})^\trans=\vv B^\trans\vv A^\trans$.
			\item\label{ex:mat:A2} 按``列的线性组合''计算 $\begin{pmatrix}1&2&0\\3&-1&4\end{pmatrix}\begin{pmatrix}2\\-1\\1\end{pmatrix}$；再说明 $\vv e_i^\trans\vv A\vv e_j$ 是 $\vv A$ 的哪个元素，$\vv A\vv e_j$、$\vv e_i^\trans\vv A$ 又分别是什么.
			\item\label{ex:mat:A3} 计算 $\begin{pmatrix}1&1\\0&1\end{pmatrix}^n$，$\begin{pmatrix}\lambda&1\\0&\lambda\end{pmatrix}^n$，$\begin{pmatrix}\cos\theta&-\sin\theta\\ \sin\theta&\cos\theta\end{pmatrix}^n$，$\begin{pmatrix}0&1&0\\0&0&1\\0&0&0\end{pmatrix}^n$.
			\item\label{ex:mat:A4} 设 $\vv A=\begin{pmatrix}1&2\\0&1\end{pmatrix}$，$\vv B=\begin{pmatrix}1&0\\3&1\end{pmatrix}$. 计算 $(\vv A+\vv B)^2$ 与 $\vv A^2+2\vv{AB}+\vv B^2$，二者是否相等？$(\vv A+\vv B)(\vv A-\vv B)$ 与 $\vv A^2-\vv B^2$ 呢？
			\item\label{ex:mat:A5} 用初等行变换求下列矩阵的逆：
			\[\begin{pmatrix}1&2&3\\2&2&1\\3&4&3\end{pmatrix},\qquad\begin{pmatrix}1&1&1\\0&1&1\\0&0&1\end{pmatrix},\qquad \begin{pmatrix}2&1&0\\1&2&1\\0&1&2\end{pmatrix},\qquad\begin{pmatrix}0&0&0&1\\0&0&2&0\\0&3&0&0\\4&0&0&0\end{pmatrix}.\]
			\item\label{ex:mat:A6} 解矩阵方程：(1) $\begin{pmatrix}2&5\\1&3\end{pmatrix}\vv X=\begin{pmatrix}4&-6\\2&1\end{pmatrix}$；(2) $\vv X\begin{pmatrix}2&1\\1&1\end{pmatrix}=\begin{pmatrix}1&0\\0&2\end{pmatrix}$.
			\item\label{ex:mat:A7} 求矩阵 $\begin{pmatrix}1&1&a\\1&a&1\\a&1&1\end{pmatrix}$ 的秩（对 $a$ 分情况讨论）.
			\item\label{ex:mat:A8} 求矩阵 $\begin{pmatrix}1&2&-1&0&3\\2&4&1&3&0\\-1&-2&2&1&-3\\1&2&2&3&0\end{pmatrix}$ 的秩，以及列向量组的一个极大无关组，并用它表出其余各列.
			\item\label{ex:mat:A9} 写出 $3$ 阶初等矩阵 $\vv E(1,3)$，$\vv E(2(5))$，$\vv E(3,1(-2))$，并分别写出它们的逆. 计算 $\vv E(3,1(-2))\vv A$ 与 $\vv A\vv E(3,1(-2))$，其中 $\vv A$ 是任一 $3$ 阶方阵，说明效果.
			\item\label{ex:mat:A10} 把 $\begin{pmatrix}2&1\\4&3\end{pmatrix}$ 写成若干个初等矩阵的乘积.
			\item\label{ex:mat:A11} 证明：两个上三角矩阵之积仍是上三角矩阵，且积的对角元等于两者对角元之积. 可逆上三角矩阵的逆也是上三角矩阵吗？
			\item\label{ex:mat:A12} 设 $\vv A,\vv B$ 都是 $n$ 阶对称矩阵. 证明：(1) $\vv A^k$，$\vv A+\vv B$，$\vv P^\trans\vv A\vv P$ 都对称；(2) $\vv{AB}$ 对称当且仅当 $\vv{AB}=\vv{BA}$；(3) 对任意 $m\times n$ 矩阵 $\vv C$，$\vv C^\trans\vv C$ 与 $\vv C\vv C^\trans$ 都对称.
			\item\label{ex:mat:A13} (1) 设 $\vv x,\vv y\in\mathbb{F}^n$，证明 $\tr(\vv x\vv y^\trans)=\vv y^\trans\vv x$；(2) 设 $n$ 阶方阵 $\vv A$ 满足：对一切 $n$ 阶方阵 $\vv X$ 都有 $\tr(\vv{AX})=0$. 证明 $\vv A=\vv O$（提示：取 $\vv X=\vv E_{ji}$）；(3) 证明：不存在 $n$ 阶方阵 $\vv A,\vv B$ 使 $\vv{AB}-\vv{BA}=\vv I$.
			\item\label{ex:mat:A14} 设方阵 $\vv A$ 满足 $\vv A^2-\vv A-2\vv I=\vv O$. 证明 $\vv A$ 与 $\vv A+2\vv I$ 都可逆，并用 $\vv A$ 表示它们的逆. $\vv A+\vv I$ 一定可逆吗？
			\item\label{ex:mat:A15} 证明：若 $\vv A$ 可逆且 $\vv{AB}=\vv{AC}$，则 $\vv B=\vv C$；举例说明 $\vv A$ 不可逆时结论可能不成立.
			\item\label{ex:mat:A16} 设 $\vv A,\vv B$ 可逆. 求 $\begin{pmatrix}\vv A&\vv O\\ \vv O&\vv B\end{pmatrix}^{-1}$ 与 $\begin{pmatrix}\vv O&\vv A\\ \vv B&\vv O\end{pmatrix}^{-1}$，并用后者求 $\begin{pmatrix}0&0&1&2\\0&0&1&3\\2&0&0&0\\0&5&0&0\end{pmatrix}^{-1}$.
			\item\label{ex:mat:A17} 用分块乘法计算 $\begin{pmatrix}1&0&2&1\\0&1&1&3\\0&0&1&0\\0&0&0&1\end{pmatrix}^2$，并猜出 $n$ 次方的一般形式.
			\item\label{ex:mat:A18} 某公司在两个工厂生产三种产品，各工厂每天的产量（件）与三种产品的单位成本、单位售价（元）分别为
			\[\vv Q=\begin{pmatrix}100&80&60\\120&50&90\end{pmatrix}\ (\text{行：工厂；列：产品}),\qquad \vv C=\begin{pmatrix}20&30\\15&25\\40&55\end{pmatrix}\ (\text{列：成本、售价}).\]
			计算 $\vv Q\vv C$，说明它每个元素的含义，并求两个工厂每天的利润.
			\item\label{ex:mat:A19} 四座城市之间的直达航线为：$1\leftrightarrow2$，$1\leftrightarrow3$，$2\leftrightarrow3$，$3\leftrightarrow4$（双向）. 写出邻接矩阵 $\vv A$，计算 $\vv A^2$ 与 $\vv A^3$. $(\vv A^2)_{ii}$ 与城市 $i$ 的航线数有什么关系？从城市 $1$ 到城市 $4$ 恰好乘三次航班有几种走法？
			\item\label{ex:mat:A20} 设 $\vv A$ 是实反对称矩阵. 证明：对一切 $\vv x\in\mathbb{R}^n$ 有 $\vv x^\trans\vv A\vv x=0$；$\vv A^2$ 是对称矩阵.
			\item\label{ex:mat:A21} 证明：若 $\vv A$ 是 $m\times n$ 矩阵，$\vv P$ 是 $m$ 阶可逆矩阵，则 $\vv{PA}$ 与 $\vv A$ 的简化阶梯形相同.
			\item\label{ex:mat:A22} 证明：任一 $n$ 阶方阵都可以唯一地写成一个对称矩阵与一个反对称矩阵之和.（提示：$\vv A=\frac12(\vv A+\vv A^\trans)+\frac12(\vv A-\vv A^\trans)$.）
		\end{exA}
		
		\Bgroup
		\begin{exB}
			\item\label{ex:mat:B1} 分别求出全部满足 (1) $\vv A^2=\vv O$；(2) $\vv A^2=\vv I$；(3) $\vv A^2=\vv A$ 的二阶实方阵.
			\item\label{ex:mat:B2} 设 $a_1,\ldots,a_n$ 互不相同. 证明：与 $\diag(a_1,\ldots,a_n)$ 可交换的矩阵一定是对角矩阵.
			\item\label{ex:mat:B3} 证明：与\textbf{一切} $n$ 阶方阵都可交换的矩阵必为数量矩阵 $c\vv I$.（提示：考虑与 $\vv E_{ij}$ 交换.）
			\item\label{ex:mat:B4} 设 $\vv A,\vv B$ 都是幂零矩阵且 $\vv{AB}=\vv{BA}$. 证明 $\vv A+\vv B$ 与 $\vv{AB}$ 都是幂零矩阵. 举例说明``可交换''不能去掉.
			\item\label{ex:mat:B5} 证明：$\rk\begin{pmatrix}\vv A&\vv O\\ \vv O&\vv B\end{pmatrix}=\rk\vv A+\rk\vv B$，$\rk\begin{pmatrix}\vv A&\vv C\\ \vv O&\vv B\end{pmatrix}\ge\rk\vv A+\rk\vv B$，并举例说明后者可以取严格不等号.
			\item\label{ex:mat:B6} （Sherman–Morrison 公式）设 $\vv u,\vv v\in\mathbb{F}^n$，$1+\vv v^\trans\vv u\ne0$. 证明
			\[(\vv I+\vv u\vv v^\trans)^{-1}=\vv I-\frac{\vv u\vv v^\trans}{1+\vv v^\trans\vv u}.\]
			\item\label{ex:mat:B7} 设 $\vv A$ 为 $m\times n$ 矩阵，$\vv B$ 为 $n\times m$ 矩阵. 证明：若 $\vv I_m+\vv{AB}$ 可逆，则 $\vv I_n+\vv{BA}$ 也可逆，并且 $(\vv I_n+\vv{BA})^{-1}=\vv I_n-\vv B(\vv I_m+\vv{AB})^{-1}\vv A$.
			\item\label{ex:mat:B8} 证明：秩为 $1$ 的 $m\times n$ 矩阵都可以写成 $\vv u\vv v^\trans$（$\vv u\in\mathbb{F}^m$，$\vv v\in\mathbb{F}^n$ 非零）. 对这样的矩阵证明 $\vv A^2=(\tr\vv A)\vv A$（$m=n$ 时），并由此求 $\vv A^k$.
			\item\label{ex:mat:B9} 证明：秩为 $r$ 的矩阵可以写成 $r$ 个秩为 $1$ 的矩阵之和，但不能写成少于 $r$ 个秩为 $1$ 的矩阵之和.
			\item\label{ex:mat:B10} （满秩分解）证明：秩为 $r$ 的 $m\times n$ 矩阵 $\vv A$ 可以写成 $\vv A=\vv{BC}$，其中 $\vv B$ 是列满秩的 $m\times r$ 矩阵，$\vv C$ 是行满秩的 $r\times n$ 矩阵. 对 A8 中的矩阵写出一个满秩分解.
			\item\label{ex:mat:B11} 计算 $\begin{pmatrix}1&1&1\\1&1&1\\1&1&1\end{pmatrix}^n$ 与 $\begin{pmatrix}2&1&0\\0&2&1\\0&0&2\end{pmatrix}^n$.（提示：后者写成 $2\vv I+\vv N$.）
			\item\label{ex:mat:B12} 证明：$n$ 阶严格上三角矩阵（对角线及以下全为零）$\vv N$ 满足 $\vv N^n=\vv O$. 更精确地，$\vv N^k$ 的非零元只可能出现在第 $k$ 条``上对角线''及其右上方.
			\item\label{ex:mat:B13} 设 $\vv A$ 为实矩阵. 证明：若 $\vv A^\trans\vv A=\vv O$，则 $\vv A=\vv O$；若 $\tr(\vv A^\trans\vv A)=0$，结论同样成立. 复矩阵呢？
			\item\label{ex:mat:B14} 每行、每列恰有一个 $1$、其余元素为 $0$ 的方阵称为\textbf{置换矩阵}. 证明：置换矩阵之积仍是置换矩阵，置换矩阵 $\vv P$ 满足 $\vv P^\trans\vv P=\vv I$；$\vv{PA}$ 与 $\vv{AP}$ 分别对 $\vv A$ 做了什么？
			\item\label{ex:mat:B15} 对 $\vv A=\begin{pmatrix}1&2&1\\2&4&3\\1&2&2\end{pmatrix}$，求可逆矩阵 $\vv P,\vv Q$ 使 $\vv{PAQ}$ 为相抵标准形.
			\item\label{ex:mat:B16} 设 $\vv A$ 是 $n$ 阶方阵，$\vv A^k=\vv O$. 证明 $\vv I+\vv A$ 可逆，并求其逆. 若 $\vv A^2=2\vv A$，$\vv I-\vv A$ 是否可逆？
		\end{exB}
		
		\Cgroup
		\begin{exC}
			\item\label{ex:mat:C1} 设 $f:M_n(\mathbb{F})\to\mathbb{F}$ 满足 $f(\vv A+\vv B)=f(\vv A)+f(\vv B)$，$f(k\vv A)=kf(\vv A)$，且 $f(\vv{AB})=f(\vv{BA})$ 对一切 $\vv A,\vv B$ 成立. 证明 $f=c\cdot\tr$ 对某个常数 $c$ 成立.（提示：计算 $f(\vv E_{ij})$.）
			\item\label{ex:mat:C2} 设 $n$ 阶方阵 $\vv A,\vv B$ 满足 $\vv A+\vv B=\vv{AB}$. 证明 $\vv{AB}=\vv{BA}$.（提示：考察 $(\vv I-\vv A)(\vv I-\vv B)$.）
			\item\label{ex:mat:C3} （Strassen 算法）对二阶分块矩阵的乘法，直接计算需要 $8$ 次块乘法. 验证下面的公式只需 $7$ 次：令
			$M_1=(A_{11}+A_{22})(B_{11}+B_{22})$，$M_2=(A_{21}+A_{22})B_{11}$，$M_3=A_{11}(B_{12}-B_{22})$，$M_4=A_{22}(B_{21}-B_{11})$，$M_5=(A_{11}+A_{12})B_{22}$，$M_6=(A_{21}-A_{11})(B_{11}+B_{12})$，$M_7=(A_{12}-A_{22})(B_{21}+B_{22})$，则
			$C_{11}=M_1+M_4-M_5+M_7$，$C_{12}=M_3+M_5$，$C_{21}=M_2+M_4$，$C_{22}=M_1-M_2+M_3+M_6$. 由此说明 $n=2^k$ 阶矩阵乘法可以用约 $n^{\log_27}\approx n^{2.81}$ 次数的乘法完成.
			\item\label{ex:mat:C4} 设 $\vv A$ 是幂等矩阵. 用满秩分解 $\vv A=\vv{BC}$ 证明 $\vv{CB}=\vv I_r$，进而证明 $\rk\vv A=\tr\vv A$.
			\item\label{ex:mat:C5} （Kronecker 积）对 $m\times n$ 矩阵 $\vv A$ 与 $p\times q$ 矩阵 $\vv B$，定义 $\vv A\otimes\vv B=(a_{ij}\vv B)$（$mp\times nq$ 分块矩阵）. 证明 $(\vv A\otimes\vv B)(\vv C\otimes\vv D)=(\vv{AC})\otimes(\vv{BD})$，并由此说明 $\vv A,\vv B$ 可逆时 $(\vv A\otimes\vv B)^{-1}=\vv A^{-1}\otimes\vv B^{-1}$.
			\item\label{ex:mat:C6} （Hadamard 矩阵）元素全为 $\pm1$ 且满足 $\vv H\vv H^\trans=n\vv I$ 的 $n$ 阶矩阵称为 Hadamard 矩阵. 验证 $\vv H_2=\begin{pmatrix}1&1\\1&-1\end{pmatrix}$ 是，并证明 $\vv H_2\otimes\vv H_2$，$\vv H_2\otimes\vv H_2\otimes\vv H_2$ 也是. 证明：$n>2$ 时 Hadamard 矩阵的阶数必是 $4$ 的倍数.（这类矩阵用于实验设计与纠错编码.）
			\item\label{ex:mat:C7} 证明：全体形如 $\begin{pmatrix}a&-b\\b&a\end{pmatrix}$（$a,b\in\mathbb{R}$）的矩阵关于矩阵的加法和乘法``像复数一样''运算：加法、乘法封闭，乘法可交换，非零元可逆. 写出它与 $\mathbb{C}$ 之间的对应，并用它验证 $(1+\mathrm i)^8=16$.
			\item\label{ex:mat:C8} 设 $\vv A$ 是元素非负、各列之和都小于 $1$ 的方阵. 证明 $\vv A^k$ 的每个元素都随 $k\to\infty$ 趋于 $0$，并且级数 $\vv I+\vv A+\vv A^2+\cdots$（按元素）收敛到 $(\vv I-\vv A)^{-1}$. 解释它在投入产出模型中的意义：``直接消耗 $+$ 一次间接消耗 $+$ 二次间接消耗 $+\cdots$''.
			\item\label{ex:mat:C9} 设 $\vv A$ 为 $m\times n$ 矩阵. 证明：$\vv A$ 有右逆（存在 $\vv C$ 使 $\vv{AC}=\vv I_m$）当且仅当 $\rk\vv A=m$；有左逆当且仅当 $\rk\vv A=n$. 非方阵的单侧逆唯一吗？
		\end{exC}
	}
	
	% =====================================================================
	% 第四章　线性映射
	% =====================================================================
	\newcommand{\ExLinearMap}{%
		\exsetheading
		\Agroup
		\begin{exA}
			\item\label{ex:lmap:A1} 判断下列映射是否线性，是线性映射的求出核：
			\begin{enumerate}
				\item $\mathbb{R}^2\to\mathbb{R}^2$，$(x,y)^\trans\mapsto(x+y,\,2x)^\trans$；
				\item $\mathbb{R}^2\to\mathbb{R}^2$，$(x,y)^\trans\mapsto(x+1,\,y)^\trans$；
				\item $\mathbb{R}^2\to\mathbb{R}$，$(x,y)^\trans\mapsto|x|+y$；
				\item $\mathbb{R}[x]\to\mathbb{R}[x]$，$f(x)\mapsto f(x+1)$；
				\item $\mathbb{R}[x]\to\mathbb{R}[x]$，$f(x)\mapsto xf'(x)$；
				\item $M_n(\mathbb{R})\to M_n(\mathbb{R})$，$\vv X\mapsto\vv X^\trans$；
				\item $M_2(\mathbb{R})\to\mathbb{R}$，$\vv X\mapsto x_{11}x_{22}-x_{12}x_{21}$.
			\end{enumerate}
			\item\label{ex:lmap:A2} 求下列矩阵所定义的线性映射 $\vv x\mapsto\vv{Ax}$ 的核与像（各给出一组基），并验证秩—零化度定理：
			\[\begin{pmatrix}1&2&-1\\2&4&-2\end{pmatrix},\qquad\begin{pmatrix}1&0&2&1\\0&1&-1&1\\1&1&1&2\end{pmatrix}.\]
			\item\label{ex:lmap:A3} 写出线性映射 $\mathscr T:\mathbb{R}^3\to\mathbb{R}^2$，$\mathscr T(x_1,x_2,x_3)^\trans=(x_1-2x_2+x_3,\ 3x_2+4x_3)^\trans$ 在标准基下的矩阵.
			\item\label{ex:lmap:A4} 求下列平面变换在标准基下的矩阵：(1) 逆时针旋转 $90^\circ$；(2) 关于直线 $y=x$ 的反射；(3) 向直线 $y=x$ 的正投影；(4) 沿 $x$ 方向的剪切 $(x,y)\mapsto(x+2y,y)$. 计算 (2)(3) 的矩阵的平方，并解释结果.
			\item\label{ex:lmap:A5} 求求导变换 $\mathscr D:\mathbb{R}[x]_3\to\mathbb{R}[x]_3$ 在基 $1,x,x^2$ 下的矩阵，以及在基 $1,\ x-1,\ (x-1)^2$ 下的矩阵.
			\item\label{ex:lmap:A6} 求转置变换 $\vv X\mapsto\vv X^\trans$ 在 $M_2(\mathbb{F})$ 的基 $\vv E_{11},\vv E_{12},\vv E_{21},\vv E_{22}$ 下的矩阵. 它的平方是什么？
			\item\label{ex:lmap:A7} 在 $\mathbb{R}^3$ 中，求从基 $\vv\varepsilon_1=(1,0,0)^\trans$，$\vv\varepsilon_2=(1,1,0)^\trans$，$\vv\varepsilon_3=(1,1,1)^\trans$ 到基 $\vv\eta_1=(1,2,1)^\trans$，$\vv\eta_2=(2,3,3)^\trans$，$\vv\eta_3=(3,7,1)^\trans$ 的过渡矩阵，并求 $\vv\alpha=(3,4,2)^\trans$ 在两组基下的坐标.
			\item\label{ex:lmap:A8} 设线性变换 $\mathscr T$ 在基 $\vv\varepsilon_1,\vv\varepsilon_2$ 下的矩阵为 $\begin{pmatrix}2&1\\1&2\end{pmatrix}$. 求它在基 $\vv\eta_1=\vv\varepsilon_1+\vv\varepsilon_2$，$\vv\eta_2=\vv\varepsilon_1-\vv\varepsilon_2$ 下的矩阵.
			\item\label{ex:lmap:A9} 证明：线性映射的复合仍是线性映射；可逆（双射）线性映射的逆映射仍是线性映射.
			\item\label{ex:lmap:A10} 证明：线性映射 $\mathscr T$ 是单射，当且仅当它把每个线性无关组都映为线性无关组.
			\item\label{ex:lmap:A11} 是否存在线性变换 $\mathscr T:\mathbb{R}^2\to\mathbb{R}^2$，使
			(1) $\mathscr T(1,1)^\trans=(1,0)^\trans$，$\mathscr T(1,-1)^\trans=(0,1)^\trans$，$\mathscr T(2,0)^\trans=(1,1)^\trans$？
			(2) 同上，但最后一个条件改为 $\mathscr T(2,0)^\trans=(1,2)^\trans$？
			存在时求出 $\mathscr T(x,y)^\trans$.
			\item\label{ex:lmap:A12} (1) 线性映射 $\mathscr T:\mathbb{R}^5\to\mathbb{R}^3$ 是满射，$\dim\ker\mathscr T$ 是多少？(2) 是否存在从 $\mathbb{R}^3$ 到 $\mathbb{R}^5$ 的满射线性映射？是否存在从 $\mathbb{R}^5$ 到 $\mathbb{R}^3$ 的单射线性映射？
			\item\label{ex:lmap:A13} 设 $\mathscr T:U\to V$，$\mathscr S:V\to W$ 线性. 证明 $\ker\mathscr T\subset\ker(\mathscr{ST})$，$\imm(\mathscr{ST})\subset\imm\mathscr S$，从而 $\rk(\mathscr{ST})\le\min\{\rk\mathscr S,\rk\mathscr T\}$.
			\item\label{ex:lmap:A14} 求迹映射 $\tr:M_n(\mathbb{F})\to\mathbb{F}$ 的核的维数.
			\item\label{ex:lmap:A15} 证明：相似矩阵有相同的迹（利用命题~\ref{prop:trace}(3)）. 由此判断 $\begin{pmatrix}1&2\\0&3\end{pmatrix}$ 与 $\begin{pmatrix}1&0\\0&2\end{pmatrix}$ 是否相似.
			\item\label{ex:lmap:A16} 证明：与单位矩阵相似的矩阵只有单位矩阵本身. 由此说明 $\begin{pmatrix}1&1\\0&1\end{pmatrix}$ 与 $\vv I$ 不相似.
			\item\label{ex:lmap:A17} 定义求值映射 $\mathscr E:\mathbb{R}[x]_3\to\mathbb{R}^3$，$f\mapsto(f(0),f(1),f(2))^\trans$. 求它在基 $1,x,x^2$ 与标准基下的矩阵，并证明 $\mathscr E$ 是同构.
			\item\label{ex:lmap:A18} 设 $\mathscr T$ 是 $\mathbb{R}^3$ 上的线性变换，$\mathscr T\vv e_1=(1,1,0)^\trans$，$\mathscr T\vv e_2=(0,1,1)^\trans$，$\mathscr T\vv e_3=(1,2,1)^\trans$. 求 $\mathscr T(x_1,x_2,x_3)^\trans$，$\ker\mathscr T$ 与 $\imm\mathscr T$.
			\item\label{ex:lmap:A19} 称 $\vv x\mapsto\vv{Ax}+\vv b$ 为\textbf{仿射变换}. 证明两个仿射变换的复合仍是仿射变换；仿射变换是线性变换当且仅当 $\vv b=\zero$.
			\item\label{ex:lmap:A20} 设 $\vv R(\theta)=\begin{pmatrix}\cos\theta&-\sin\theta\\ \sin\theta&\cos\theta\end{pmatrix}$ 是平面上绕原点旋转角 $\theta$ 的变换在标准基下的矩阵. 从几何上说明``先转 $\beta$ 再转 $\alpha$''就是转 $\alpha+\beta$，从而 $\vv R(\alpha)\vv R(\beta)=\vv R(\alpha+\beta)$；比较两边矩阵的元素，推出两角和的正弦、余弦公式.
		\end{exA}
		
		\Bgroup
		\begin{exB}
			\item\label{ex:lmap:B1} 记 $L(V,W)$ 为 $V$ 到 $W$ 的全体线性映射. 证明它关于映射的加法与数乘构成向量空间，且 $\dim L(V,W)=\dim V\cdot\dim W$.（提示：取定基后，线性映射与矩阵一一对应.）
			\item\label{ex:lmap:B2} 设 $\mathscr T$ 是 $n$ 维空间 $V$ 上的线性变换，$\mathscr T^2=\mathscr O$. 证明 $\imm\mathscr T\subset\ker\mathscr T$，从而 $\rk\mathscr T\le n/2$；并证明存在一组基，使 $\mathscr T$ 的矩阵只在 $(1,2),(3,4),\ldots,(2r-1,2r)$ 这 $r$ 个位置上为 $1$，其余元素全为 $0$，其中 $r=\rk\mathscr T$.
			\item\label{ex:lmap:B3} 设 $\mathscr T,\mathscr S$ 是有限维空间上的线性变换. 证明 $\dim\ker(\mathscr{ST})\le\dim\ker\mathscr S+\dim\ker\mathscr T$.
			\item\label{ex:lmap:B4} 设 $V$ 有限维，$\mathscr S,\mathscr T$ 是 $V$ 上的线性变换且 $\mathscr{ST}=\mathscr I$. 证明 $\mathscr{TS}=\mathscr I$.
			\item\label{ex:lmap:B5} 在无限维空间 $\mathbb{R}[x]$ 上，令 $\mathscr D$ 为求导，$\mathscr J$ 为 $f\mapsto\int_0^xf(t)\,\mathrm dt$. 证明 $\mathscr{DJ}=\mathscr I$，但 $\mathscr{JD}\ne\mathscr I$；$\mathscr D$ 是满射但不是单射，$\mathscr J$ 是单射但不是满射. 这与 B4 是否矛盾？
			\item\label{ex:lmap:B6} 固定 $\vv A\in M_n(\mathbb{F})$，令 $\mathscr T(\vv X)=\vv{AX}-\vv{XA}$. 证明 $\imm\mathscr T$ 含于迹为零的矩阵全体，从而 $\mathscr T$ 不是满射；并证明 $\ker\mathscr T$ 包含 $\vv I,\vv A,\vv A^2,\ldots$.
			\item\label{ex:lmap:B7} 设 $\vv A=\begin{pmatrix}a&b\\c&d\end{pmatrix}$. 求线性变换 $\vv X\mapsto\vv{AX}$ 在 $M_2(\mathbb{F})$ 的基 $\vv E_{11},\vv E_{21},\vv E_{12},\vv E_{22}$ 下的矩阵. 若换成 $\vv X\mapsto\vv{XA}$ 呢？
			\item\label{ex:lmap:B8} 设线性变换 $\mathscr P$ 满足 $\mathscr P^2=\mathscr P$. 证明存在一组基使 $\mathscr P$ 的矩阵为 $\begin{pmatrix}\vv I_r&\vv O\\ \vv O&\vv O\end{pmatrix}$，其中 $r=\rk\mathscr P$.
			\item\label{ex:lmap:B9} 设线性变换 $\mathscr T$ 满足 $\mathscr T^2=\mathscr I$. 证明 $V=\ker(\mathscr T-\mathscr I)\oplus\ker(\mathscr T+\mathscr I)$，从而存在一组基使 $\mathscr T$ 的矩阵为 $\diag(\vv I_p,-\vv I_q)$.（提示：$\vv\alpha=\frac12(\vv\alpha+\mathscr T\vv\alpha)+\frac12(\vv\alpha-\mathscr T\vv\alpha)$.）
			\item\label{ex:lmap:B10} 设 $\mathscr T:V\to W$ 线性，$U$ 是 $V$ 的子空间. 证明 $\dim\mathscr T(U)=\dim U-\dim(U\cap\ker\mathscr T)$.
			\item\label{ex:lmap:B11} 在由 $\mathrm e^x\cos x$，$\mathrm e^x\sin x$ 张成的二维函数空间中，求求导变换的矩阵，证明它可逆，并利用逆矩阵求出 $\int\mathrm e^x\sin x\,\mathrm dx$（不必用分部积分）.
			\item\label{ex:lmap:B12} （齐次坐标）平面上的平移 $(x,y)\mapsto(x+a,y+b)$ 不是线性的. 把点 $(x,y)$ 写成三维向量 $(x,y,1)^\trans$，证明平移、旋转、伸缩都可以写成 $3\times3$ 矩阵的乘法，且仿射变换的复合对应矩阵的乘积. 写出``先绕原点旋转 $90^\circ$，再平移 $(1,2)$''与``先平移再旋转''的矩阵，它们相等吗？（计算机图形学正是这样处理变换的.）
			\item\label{ex:lmap:B13} 证明：若 $\vv A$ 与 $\vv B$ 相似，则 $\vv A^k$ 与 $\vv B^k$ 相似，$\vv A^\trans$ 与 $\vv B^\trans$ 相似，$f(\vv A)$ 与 $f(\vv B)$ 相似（$f$ 为多项式）.
			\item\label{ex:lmap:B14} 证明 $\begin{pmatrix}1&0\\0&0\end{pmatrix}$ 与 $\begin{pmatrix}0&1\\0&0\end{pmatrix}$ 相抵但不相似. 从``线性映射 $V\to W$''与``线性变换 $V\to V$''的区别解释：为什么相抵只需看秩，而相似要复杂得多？
			\item\label{ex:lmap:B15} 设 $\mathscr T$ 是 $V$ 上的线性变换，$\vv\alpha\in V$ 满足 $\mathscr T^{k-1}\vv\alpha\ne\zero$，$\mathscr T^k\vv\alpha=\zero$. 证明 $\vv\alpha,\mathscr T\vv\alpha,\ldots,\mathscr T^{k-1}\vv\alpha$ 线性无关. 若 $k=n=\dim V$，求 $\mathscr T$ 在基 $\mathscr T^{n-1}\vv\alpha,\ldots,\mathscr T\vv\alpha,\vv\alpha$ 下的矩阵.
		\end{exB}
		
		\Cgroup
		\begin{exC}
			\item\label{ex:lmap:C1} （Fitting 分解）设 $\mathscr T$ 是 $n$ 维空间 $V$ 上的线性变换. 证明：
			\begin{enumerate}
				\item $\ker\mathscr T\subset\ker\mathscr T^2\subset\cdots$，$\imm\mathscr T\supset\imm\mathscr T^2\supset\cdots$，且两条链都在某一步之后不再变化；
				\item $V=\ker\mathscr T^n\oplus\imm\mathscr T^n$，两者都是 $\mathscr T$ 的不变子空间，$\mathscr T$ 在前者上幂零，在后者上可逆.
			\end{enumerate}
			\item\label{ex:lmap:C2} （商空间）设 $W$ 是 $V$ 的子空间. 对 $\vv\alpha\in V$，称集合 $\vv\alpha+W$ 为 $\vv\alpha$ 所在的陪集，全体陪集记为 $V/W$. 定义 $(\vv\alpha+W)+(\vv\beta+W)=(\vv\alpha+\vv\beta)+W$，$k(\vv\alpha+W)=k\vv\alpha+W$.
			\begin{enumerate}
				\item 证明这两个运算是良定义的，且 $V/W$ 构成向量空间；
				\item 证明 $\dim V/W=\dim V-\dim W$；
				\item （第一同构定理/同态基本定理）对线性映射 $\mathscr T:V\to U$，证明 $V/\ker\mathscr T\cong\imm\mathscr T$. 这给出秩—零化度定理的另一种理解.
			\end{enumerate}
			\item\label{ex:lmap:C3} 把 $\mathbb{C}$ 看作 $\mathbb{R}$ 上的二维空间，基为 $1,\mathrm i$. (1) 求``乘以 $a+b\mathrm i$''这个变换的矩阵；(2) 求``取共轭''这个变换的矩阵；(3) 证明 $\mathbb{R}$-线性变换 $\begin{pmatrix}p&q\\r&s\end{pmatrix}$ 是``乘以某个复数''，当且仅当 $p=s$，$q=-r$.
			\item\label{ex:lmap:C4} 证明：$n$ 维空间上的任一线性变换 $\mathscr T$，都存在次数不超过 $n^2$ 的非零多项式 $f$ 使 $f(\mathscr T)=\mathscr O$.（提示：$\mathscr I,\mathscr T,\ldots,\mathscr T^{n^2}$ 在 $L(V,V)$ 中线性相关. 第八章的 Cayley–Hamilton 定理会把 $n^2$ 改进为 $n$.）
			\item\label{ex:lmap:C5} （Haar 小波）在 $\mathbb{R}^4$ 中取基 $\vv h_1=(1,1,1,1)^\trans$，$\vv h_2=(1,1,-1,-1)^\trans$，$\vv h_3=(1,-1,0,0)^\trans$，$\vv h_4=(0,0,1,-1)^\trans$. 求信号 $\vv x=(8,6,3,4)^\trans$ 在这组基下的坐标，并解释每个坐标的含义. 若把绝对值最小的坐标改为 $0$ 再变换回来，误差是多少？这是图像与音频压缩的思想之一.
			\item\label{ex:lmap:C6} 设 $\mathscr T$ 是 $V$ 上的线性变换. 若 $\mathscr T$ 与 $V$ 上的一切线性变换都可交换，证明 $\mathscr T=c\mathscr I$. 若只要求与一切\textbf{可逆}线性变换可交换呢？
			\item\label{ex:lmap:C7} 设 $\mathscr T:V\to V$ 线性，$\dim V=n$. 证明：$\rk\mathscr T^k-\rk\mathscr T^{k+1}$ 关于 $k$ 单调不增. 由此说明：若 $\rk\mathscr T=\rk\mathscr T^2$，则对一切 $k$，$\rk\mathscr T^k=\rk\mathscr T$.
			\item\label{ex:lmap:C8} （差分算子）在全体实数列构成的空间上，定义左移 $\mathscr S(a_1,a_2,a_3,\ldots)=(a_2,a_3,\ldots)$. 证明 $\mathscr S$ 是线性的、满射但不是单射；满足 $a_{n+2}=a_{n+1}+a_n$ 的数列全体恰好是 $\ker(\mathscr S^2-\mathscr S-\mathscr I)$. 由第二章章末习题~\ref{ex:vs:C9}，这个核是二维的.
		\end{exC}
	}
	
	% =====================================================================
	% 第五章　行列式
	% =====================================================================
	\newcommand{\ExDeterminant}{%
		\exsetheading
		\Agroup
		\begin{exA}
			\item\label{ex:det:A1} 求下列排列的逆序数并判断奇偶：$315426$，$42531$，$13\cdots(2n-1)\,24\cdots(2n)$.
			\item\label{ex:det:A2} (1) 在 $4$ 阶行列式中，项 $a_{13}a_{24}a_{32}a_{41}$ 带什么符号？(2) 选择 $i,j$，使 $a_{1i}a_{23}a_{35}a_{5j}a_{44}$ 成为 $5$ 阶行列式中带正号的一项.
			\item\label{ex:det:A3} 计算：
			\[\begin{vmatrix}2&1&-1\\3&0&2\\1&4&5\end{vmatrix},\qquad
			\begin{vmatrix}1&2&0&1\\2&1&1&0\\0&1&2&1\\1&0&1&2\end{vmatrix},\qquad
			\begin{vmatrix}3&1&1&1\\1&3&1&1\\1&1&3&1\\1&1&1&3\end{vmatrix}.\]
			\item\label{ex:det:A4} 计算 $\begin{vmatrix}1&-1&2&3\\2&0&1&-1\\0&1&-1&2\\3&1&0&1\end{vmatrix}$，分别用``打洞''化三角形与按第一列展开两种方法，比较哪种更快.
			\item\label{ex:det:A5} 计算 $n$ 阶行列式 $\begin{vmatrix}0&\cdots&0&a_1\\0&\cdots&a_2&0\\ \vdots&&&\vdots\\a_n&\cdots&0&0\end{vmatrix}$（只有副对角线上有元素）.
			\item\label{ex:det:A6} 不展开，证明：
			\[\begin{vmatrix}a+b&b+c&c+a\\a_1+b_1&b_1+c_1&c_1+a_1\\a_2+b_2&b_2+c_2&c_2+a_2\end{vmatrix}=2\begin{vmatrix}a&b&c\\a_1&b_1&c_1\\a_2&b_2&c_2\end{vmatrix}.\]
			\item\label{ex:det:A7} 解方程 $\begin{vmatrix}x&1&1\\1&x&1\\1&1&x\end{vmatrix}=0$ 与 $\begin{vmatrix}1&1&1\\1&2&x\\1&4&x^2\end{vmatrix}=0$.
			\item\label{ex:det:A8} 设 $\vv A$ 为 $3$ 阶方阵，$\det\vv A=3$. 求 $\det(2\vv A)$，$\det(-\vv A)$，$\det(\vv A^{-1})$，$\det(\vv A^*)$，$\det(\vv A^\trans\vv A)$，$\det(3\vv A^{-1}-2\vv A^*)$.
			\item\label{ex:det:A9} 证明：(1) 若 $n$ 阶方阵满足 $\vv A^\trans\vv A=\vv I$，则 $\det\vv A=\pm1$；(2) 若 $\vv A^2=\vv I$，则 $\det\vv A=\pm1$；(3) 若 $\vv A^2=\vv A$ 且 $\vv A\ne\vv I$，则 $\det\vv A=0$.
			\item\label{ex:det:A10} 求下列矩阵的伴随矩阵与逆矩阵：\[\begin{pmatrix}a&b\\c&d\end{pmatrix},\qquad\begin{pmatrix}1&2&3\\0&1&4\\5&6&0\end{pmatrix}.\]
			\item\label{ex:det:A11} 用 Cramer 法则解方程组 $2x_1+x_2+x_3=3$，$x_1+3x_2+x_3=0$，$x_1+x_2+4x_3=8$.
			\item\label{ex:det:A12} 求以 $(1,1)$，$(4,2)$，$(2,5)$ 为顶点的三角形的面积；判断 $(1,2)$，$(3,5)$，$(7,11)$ 三点是否共线.
			\item\label{ex:det:A13} 求以 $(1,0,1)^\trans$，$(2,1,0)^\trans$，$(0,3,1)^\trans$ 为棱的平行六面体的体积，以及以原点和这三个向量的终点为顶点的四面体的体积.
			\item\label{ex:det:A14} $\lambda$ 取何值时，齐次方程组 $(5-\lambda)x+2y+2z=0$，$2x+(6-\lambda)y=0$，$2x+(4-\lambda)z=0$ 有非零解？
			\item\label{ex:det:A15} 设 \[\vv A=\begin{pmatrix}1&2&3&4\\2&3&4&1\\3&4&1&2\\4&1&2&3\end{pmatrix}.\] 求 $A_{11}+A_{12}+A_{13}+A_{14}$（$A_{1j}$ 是代数余子式）.（提示：它是把第一行换成全 $1$ 后的行列式.）
			\item\label{ex:det:A16} 证明：过平面上两个不同点 $(x_1,y_1)$，$(x_2,y_2)$ 的直线方程可以写成 \[\begin{vmatrix}x&y&1\\x_1&y_1&1\\x_2&y_2&1\end{vmatrix}=0.\]
			\item\label{ex:det:A17} 设 $\vv A,\vv B$ 分别是 $m$ 阶、$n$ 阶方阵，证明 $\begin{vmatrix}\vv A&\vv O\\ \vv C&\vv B\end{vmatrix}=\det\vv A\det\vv B$，并计算 $\begin{vmatrix}\vv O&\vv A\\ \vv B&\vv O\end{vmatrix}$.
			\item\label{ex:det:A18} 计算 $4$ 阶三对角行列式 $\begin{vmatrix}2&-1&0&0\\-1&2&-1&0\\0&-1&2&-1\\0&0&-1&2\end{vmatrix}$.
			\item\label{ex:det:A19} 设 $\vv A$ 是 $n$ 阶可逆矩阵，$n\ge2$. 证明 $\vv A^*$ 可逆，且 $(\vv A^*)^{-1}=\frac1{\det\vv A}\vv A=(\vv A^{-1})^*$.
			\item\label{ex:det:A20} 用行列式判断：向量 $(1,2,3)^\trans$，$(4,5,6)^\trans$，$(7,8,9)^\trans$ 是否线性无关？把 $9$ 改成 $10$ 呢？
			\item\label{ex:det:A21} 设 $\vv A$ 是 $n$ 阶反对称矩阵（$\vv A^\trans=-\vv A$），$n$ 为奇数. 证明 $\det\vv A=0$.（提示：对 $\vv A^\trans=-\vv A$ 两边取行列式.）
		\end{exA}
		
		\Bgroup
		\begin{exB}
			\item\label{ex:det:B1} 计算 $n$ 阶三对角行列式 $D_n=\begin{vmatrix}2&-1&&\\-1&2&\ddots&\\&\ddots&\ddots&-1\\&&-1&2\end{vmatrix}$.（提示：按第一行展开得 $D_n=2D_{n-1}-D_{n-2}$.）
			\item\label{ex:det:B2} 计算 $n$ 阶行列式 $\begin{vmatrix}a&b&\cdots&b\\b&a&\cdots&b\\ \vdots&&\ddots&\vdots\\b&b&\cdots&a\end{vmatrix}$ 与 $\begin{vmatrix}x&a&\cdots&a\\-a&x&\cdots&a\\ \vdots&&\ddots&\vdots\\-a&-a&\cdots&x\end{vmatrix}$.
			\item\label{ex:det:B3} （爪形行列式）设 $a_1\cdots a_n\ne0$，计算 $\begin{vmatrix}a_0&1&1&\cdots&1\\1&a_1&0&\cdots&0\\1&0&a_2&\cdots&0\\ \vdots&&&\ddots&\\1&0&0&\cdots&a_n\end{vmatrix}$.
			\item\label{ex:det:B4} （矩阵行列式引理）证明 $\det(\vv I_n+\vv u\vv v^\trans)=1+\vv v^\trans\vv u$，更一般地，$\vv A$ 可逆时 $\det(\vv A+\vv u\vv v^\trans)=(1+\vv v^\trans\vv A^{-1}\vv u)\det\vv A$. 用它重做 B2 的第一个行列式.
			\item\label{ex:det:B5} 设 $\vv A$ 为 $2\times3$ 矩阵，$\vv B$ 为 $3\times2$ 矩阵. 证明 $\det(\vv{AB})$ 等于 $\vv A$ 的三个 $2$ 阶子式与 $\vv B$ 的对应 $2$ 阶子式乘积之和（这是 Cauchy–Binet 公式最简单的情形）. 由此说明 $\vv{BA}$（$3$ 阶方阵）的行列式为零.
			\item\label{ex:det:B6} 设 $\vv A$ 是整数矩阵. 证明：$\vv A^{-1}$ 存在且也是整数矩阵，当且仅当 $\det\vv A=\pm1$.
			\item\label{ex:det:B7} 证明：$n\ge2$ 时 $(\vv A^*)^*=(\det\vv A)^{n-2}\vv A$.（先设 $\vv A$ 可逆；不可逆的情形另行讨论.）
			\item\label{ex:det:B8} 设 $\vv A,\vv B$ 都是 $n$ 阶方阵. 证明 $\begin{vmatrix}\vv A&\vv B\\ \vv B&\vv A\end{vmatrix}=\det(\vv A+\vv B)\det(\vv A-\vv B)$.（提示：分块初等变换.）
			\item\label{ex:det:B9} 设 $\vv A$ 可逆. 证明 $\begin{vmatrix}\vv A&\vv B\\ \vv C&\vv D\end{vmatrix}=\det\vv A\cdot\det(\vv D-\vv{CA}^{-1}\vv B)$；若还有 $\vv{AC}=\vv{CA}$，则它等于 $\det(\vv{AD}-\vv{CB})$.
			\item\label{ex:det:B10} 用 Fibonacci 矩阵 $\vv F=\begin{pmatrix}1&1\\1&0\end{pmatrix}$ 满足 $\vv F^n=\begin{pmatrix}F_{n+1}&F_n\\F_n&F_{n-1}\end{pmatrix}$（归纳法证明），推出 Cassini 恒等式 $F_{n+1}F_{n-1}-F_n^2=(-1)^n$.
			\item\label{ex:det:B11} 计算 $n$ 阶行列式 $D_n=\begin{vmatrix}1+x^2&x&&\\x&1+x^2&\ddots&\\&\ddots&\ddots&x\\&&x&1+x^2\end{vmatrix}$.
			\item\label{ex:det:B12} 设 $\omega=\mathrm e^{2\pi\mathrm i/3}$. 证明
			\[\begin{vmatrix}a&b&c\\c&a&b\\b&c&a\end{vmatrix}=(a+b+c)(a+\omega b+\omega^2c)(a+\omega^2b+\omega c)=a^3+b^3+c^3-3abc.\]
			\item\label{ex:det:B13} 证明：元素为 $a_{ij}=\min(i,j)$ 的 $n$ 阶行列式等于 $1$.
			\item\label{ex:det:B14} 证明：平面上四点 $(x_i,y_i)$（$i=1,2,3,4$）共圆（或共线），当且仅当 $\begin{vmatrix}x_i^2+y_i^2&x_i&y_i&1\end{vmatrix}_{4\times4}=0$（每行对应一个点）.
						\item\label{ex:det:B15} 设 $\vv A$ 为 $n$ 阶方阵（$n\ge2$）. 证明
			\[\rk\vv A^*=\begin{cases}n,&\rk\vv A=n,\\1,&\rk\vv A=n-1,\\0,&\rk\vv A\le n-2.\end{cases}\]
			（提示：$\vv A^*$ 的元素都是 $\vv A$ 的 $n-1$ 阶子式，用定理~\ref{thm:minorrank}；$\rk\vv A=n-1$ 时考察 $\vv A\vv A^*=(\det\vv A)\vv I$.）
			\item\label{ex:det:B16} 设 $f(x)=\begin{vmatrix}x&1&2&3\\1&x&2&3\\1&2&x&3\\1&2&3&x\end{vmatrix}$. 不展开，求 $f(x)=0$ 的全部根.
		\end{exB}
		
		\Cgroup
		\begin{exC}
			\item\label{ex:det:C1} （行列式的唯一性）设 $D:M_n(\mathbb{F})\to\mathbb{F}$ 对每一行都是线性的、交换两行变号，且 $D(\vv I)=1$. 证明 $D=\det$.（提示：先说明 $D$ 在初等矩阵上的取值，再对任意矩阵用初等变换.）
			\item\label{ex:det:C2} （Cauchy 行列式）证明
			\[\det\Big(\frac1{a_i+b_j}\Big)_{n\times n}=\frac{\prod_{i<j}(a_j-a_i)(b_j-b_i)}{\prod_{i,j}(a_i+b_j)}.\]
			先对 $n=2$ 直接验证，再用``各行减去最后一行''归纳. 取 $a_i=i$，$b_j=j-1$，得到 Hilbert 矩阵 $(1/(i+j-1))$ 的行列式，它为什么在数值计算中臭名昭著？
			\item\label{ex:det:C3} （排列的符号是同态）对排列 $\sigma$，令 $\vv P_\sigma$ 为把 $\vv e_j$ 映到 $\vv e_{\sigma(j)}$ 的置换矩阵. 证明 $\vv P_\sigma\vv P_\tau=\vv P_{\sigma\tau}$，$\det\vv P_\sigma=(-1)^{\tau(\sigma)}$，并由行列式的乘积公式推出 $\operatorname{sgn}(\sigma\tau)=\operatorname{sgn}\sigma\cdot\operatorname{sgn}\tau$.
			\item\label{ex:det:C4} （矩阵—树定理）对一个连通无向图，令 $\vv D$ 为各顶点度数构成的对角矩阵，$\vv A$ 为邻接矩阵，$\vv L=\vv D-\vv A$ 称为 Laplace 矩阵. 定理说：图的生成树个数等于 $\vv L$ 的任一个 $(n-1)$ 阶主子式. 对 $4$ 个顶点的完全图验证它有 $16$ 棵生成树，对一个 $4$ 个顶点的环验证它有 $4$ 棵生成树.
			\item\label{ex:det:C5} 证明：若 $n$ 阶矩阵的元素全为 $\pm1$，则它的行列式能被 $2^{n-1}$ 整除. 求 $3$ 阶此类行列式的最大值.
			\item\label{ex:det:C6} （Jacobi 公式的雏形）证明 $\det(\vv I+t\vv B)=1+t\tr\vv B+(\text{$t^2$ 及更高次项})$. 由此说明：若 $\vv A(t)$ 的元素都是 $t$ 的可导函数，且 $\vv A(0)=\vv I$，则 $\frac{\mathrm d}{\mathrm dt}\det\vv A(t)\big|_{t=0}=\tr\vv A'(0)$.
			\item\label{ex:det:C7} （Sylvester 结式）设 $f=a_0x^2+a_1x+a_2$，$g=b_0x^2+b_1x+b_2$（$a_0b_0\ne0$）. 证明 $f,g$ 有公共根当且仅当
			\[\begin{vmatrix}a_0&a_1&a_2&0\\0&a_0&a_1&a_2\\b_0&b_1&b_2&0\\0&b_0&b_1&b_2\end{vmatrix}=0.\]
			（提示：有公共根当且仅当存在次数 $\le1$ 的不全为零的 $u,v$ 使 $uf+vg=0$，与第七章的 Bézout 等式对照.）
			\item\label{ex:det:C8} （Dodgson 凝聚法）对 $3$ 阶矩阵 $\vv A$，记 $M_{ij}$ 为余子式，$c=a_{22}$. 验证 $\det\vv A\cdot c=M_{11}M_{33}-M_{13}M_{31}$. 它允许只计算 $2$ 阶行列式就得到高阶行列式. 《爱丽丝漫游奇境》的作者刘易斯·卡罗尔（即 Dodgson）于 1866 年发表了这一方法.
			\item\label{ex:det:C9} 设 $\vv A$ 是 $n$ 阶实反对称矩阵. 证明 $\det(\vv I+\vv A)\ge1$.（提示：$\det(\vv I+\vv A)=\det(\vv I-\vv A)$，考察 $\det(\vv I+t\vv A)$ 作为 $t$ 的多项式只含偶次项，且 $\vv I+t\vv A$ 对一切实数 $t$ 都可逆.）
		\end{exC}
	}
	
	% =====================================================================
	% 第六章　线性方程组解的结构
	% =====================================================================
	\newcommand{\ExSolutions}{%
		\exsetheading
		\Agroup
		\begin{exA}
			\item\label{ex:sol:A1} 求下列齐次方程组的基础解系与通解：
			\[\text{(1)}\ \left\{\begin{array}{l}x_1+x_2-x_3-x_4=0,\\2x_1-5x_2+3x_3+2x_4=0,\\7x_1-7x_2+3x_3+x_4=0;\end{array}\right.\qquad
			\text{(2)}\ \left\{\begin{array}{l}x_1+2x_2+x_3-x_4=0,\\3x_1+6x_2-x_3-3x_4=0,\\5x_1+10x_2+x_3-5x_4=0.\end{array}\right.\]
			\item\label{ex:sol:A2} 求下列方程组的通解，写成``特解 $+$ 基础解系的线性组合''：
			\[\text{(1)}\ \left\{\begin{array}{l}x_1+x_2-3x_3-x_4=1,\\3x_1-x_2-3x_3+4x_4=4,\\x_1+5x_2-9x_3-8x_4=0;\end{array}\right.\qquad
			\text{(2)}\ \left\{\begin{array}{l}2x_1+x_2-x_3+x_4=1,\\4x_1+2x_2-2x_3+x_4=2,\\2x_1+x_2-x_3-x_4=1.\end{array}\right.\]
			\item\label{ex:sol:A3} 设 $\vv A$ 是 $4\times5$ 矩阵，$\rk\vv A=3$. (1) $\vv{Ax}=\zero$ 的基础解系含几个向量？(2) $\vv{Ax}=\vv b$ 是否一定有解？(3) 若 $\rk\vv A=4$ 呢？
			\item\label{ex:sol:A4} 不解方程，用秩判断下列方程组是否有解：(1) $x_1+x_2=1$，$x_1-x_2=1$，$2x_1+3x_2=2$；(2) $x_1+x_2+x_3=1$，$2x_1+2x_2+2x_3=3$.
			\item\label{ex:sol:A5} 设 $\vv\eta_1,\vv\eta_2$ 是 $\vv{Ax}=\vv b$（$\vv b\ne\zero$）的两个解. 下列向量中哪些是 $\vv{Ax}=\vv b$ 的解，哪些是 $\vv{Ax}=\zero$ 的解？
			$\vv\eta_1+\vv\eta_2$，$2\vv\eta_1-\vv\eta_2$，$\frac12(\vv\eta_1+\vv\eta_2)$，$\vv\eta_1-\vv\eta_2$，$3\vv\eta_1-2\vv\eta_2$.
			\item\label{ex:sol:A6} 设 $3$ 元方程组 $\vv{Ax}=\vv b$ 的系数矩阵秩为 $2$，$\vv\eta_1=(1,2,3)^\trans$ 与 $\vv\eta_2=(1,-2,-1)^\trans$ 是它的两个解. 求通解.
			\item\label{ex:sol:A7} 写出一个以 $\vv x=(1,0,2)^\trans+k(1,1,-1)^\trans$（$k$ 任意）为全部解的线性方程组.
			\item\label{ex:sol:A8} 判断三个平面 $x+y+z=1$，$x+2y+3z=2$，$2x+3y+4z=4$ 的位置关系；若把最后一个常数改为 $3$ 呢？
			\item\label{ex:sol:A9} 举例说明：$\vv{Ax}=\zero$ 只有零解时，$\vv{Ax}=\vv b$ 可能无解；但若有解，解必唯一.
			\item\label{ex:sol:A10} 解矩阵方程 $\begin{pmatrix}1&2&3\\2&2&1\\3&4&3\end{pmatrix}\vv X=\begin{pmatrix}2&5\\3&1\\4&3\end{pmatrix}$.
			\item\label{ex:sol:A11} 设 $\vv A=\begin{pmatrix}1&1\\2&2\end{pmatrix}$. 求全部满足 $\vv{AX}=\vv O$ 的二阶矩阵 $\vv X$，它们构成的空间是几维的？
			\item\label{ex:sol:A12} 设 $\vv A=\begin{pmatrix}1&2&-2\\4&t&3\\3&-1&1\end{pmatrix}$，且存在三阶非零矩阵 $\vv B$ 使 $\vv{AB}=\vv O$. 求 $t$，并求 $\rk\vv B$ 的可能值.
			\item\label{ex:sol:A13} 求方程 $x_1+x_2+x_3=1$ 的通解，并写出它的三个解，使得通解中的每个解都是这三个解的仿射组合.
			\item\label{ex:sol:A14} 把空间直线 $\{x+y+z=1,\ x-y=0\}$ 写成参数式；反过来，把直线 $\vv x=(1,2,0)^\trans+t(1,0,1)^\trans$ 写成两个平面的交线.
			\item\label{ex:sol:A15} 判断空间直线 $\vv x=(1,0,0)^\trans+s(1,1,0)^\trans$ 与 $\vv x=(0,1,1)^\trans+t(0,1,-1)^\trans$ 是否相交.（列出关于 $s,t$ 的方程组.）
			\item\label{ex:sol:A16} 物流网络：三个仓库 $A,B,C$ 之间有线路 $A\to B$（$x_1$），$B\to C$（$x_2$），$A\to C$（$x_3$）. $A$ 每天净发出 $30$ 吨，$C$ 每天净收到 $30$ 吨，$B$ 不囤货. 求全部方案. 哪些量（如 $x_1+x_3$，$x_1-x_2$，$x_3$）能被方程组唯一确定？用可确定性判据说明.
			\item\label{ex:sol:A17} 求 $\vv A=\begin{pmatrix}2&1&0\\4&3&1\\2&3&3\end{pmatrix}$ 的 LU 分解.
			\item\label{ex:sol:A18} 设 $\vv A$ 是 $m\times n$ 矩阵，$\rk\vv A=m<n$. 证明 $\vv{Ax}=\vv b$ 对每个 $\vv b$ 都有无穷多解.
			\item\label{ex:sol:A19} 求经过点 $(1,1,1)$，$(2,0,1)$，$(0,2,3)$ 的平面方程（用齐次方程组求法向量）.
			\item\label{ex:sol:A20} 设 $\vv A=\begin{pmatrix}1&1\\1&1.0001\end{pmatrix}$. 分别求 $\vv{Ax}=(2,2)^\trans$ 与 $\vv{Ax}=(2,2.0001)^\trans$ 的解，计算右端的相对变化与解的相对变化之比.
		\end{exA}
		
		\Bgroup
		\begin{exB}
			\item\label{ex:sol:B1} 求两个齐次方程组 $\{x_1+x_2=0,\ x_2-x_4=0\}$ 与 $\{x_1-x_2+x_3=0,\ x_2-x_3+x_4=0\}$ 的公共非零解.
			\item\label{ex:sol:B2} 设 $\vv A$ 是实矩阵. 证明 $\vv{Ax}=\zero$ 与 $\vv A^\trans\vv{Ax}=\zero$ 同解，从而 $\rk(\vv A^\trans\vv A)=\rk\vv A$.（提示：$\vv x^\trans\vv A^\trans\vv{Ax}=\|\vv{Ax}\|^2$ 在第九章才有名字，但这里只需``实数的平方和为零则每个都为零''.）
			\item\label{ex:sol:B3} 证明 $\rk(\vv A+\vv B)\le\rk(\vv A\mid\vv B)\le\rk\vv A+\rk\vv B$.
			\item\label{ex:sol:B4} 设 $\vv A$ 为 $n$ 阶方阵，$\vv A^2=\vv I$. 证明 $\rk(\vv A+\vv I)+\rk(\vv A-\vv I)=n$.
			\item\label{ex:sol:B5} 证明 $\rk(\vv{AB})=\rk\vv B$ 当且仅当 $\ker\vv A\cap\imm\vv B=\{\zero\}$.
			\item\label{ex:sol:B6} 设 $\vv\eta^*$ 是 $\vv{Ax}=\vv b$（$\vv b\ne\zero$）的一个解，$\vv\xi_1,\ldots,\vv\xi_k$ 是对应齐次方程组的基础解系. 证明 $\vv\eta^*,\vv\eta^*+\vv\xi_1,\ldots,\vv\eta^*+\vv\xi_k$ 线性无关.
			\item\label{ex:sol:B7} 一张 $3\times3$ 的统计表只公开了三个行和与三个列和. 至少还要再公开几个单元格，才能确定全表？给出一种``公开哪几个格''的选法，并说明为什么某些同样个数的选法却不行.
			\item\label{ex:sol:B8} 对 $\vv A=\begin{pmatrix}1&2&1\\2&4&3\\3&6&4\end{pmatrix}$，求 $\vv b=(b_1,b_2,b_3)^\trans$ 应满足的条件，使 $\vv{Ax}=\vv b$ 有解.（提示：找出把各方程组合成 $0=\cdots$ 的系数.）
			\item\label{ex:sol:B9} 设 $\vv A$ 为 $n$ 阶方阵，$\rk\vv A=n-1$，且 $\vv A$ 的每行元素之和为零. 求 $\vv{Ax}=\zero$ 的通解.
			\item\label{ex:sol:B10} 设 $\vv A,\vv B$ 为 $n$ 阶方阵，$\vv{AB}=\vv O$，$\vv A+\vv B$ 可逆. 证明 $\rk\vv A+\rk\vv B=n$.
			\item\label{ex:sol:B11} 证明：两个 $m\times n$ 齐次方程组 $\vv{Ax}=\zero$ 与 $\vv{Bx}=\zero$ 同解，当且仅当 $\vv A$ 与 $\vv B$ 的行向量组等价.
			\item\label{ex:sol:B12} 设 $\vv A$ 为 $m\times n$ 矩阵. 证明：$\vv{Ax}=\vv b$ 对一切 $\vv b$ 有解，当且仅当存在 $n\times m$ 矩阵 $\vv C$ 使 $\vv{AC}=\vv I_m$；这时 $\vv C\vv b$ 就是一个解.
			\item\label{ex:sol:B13} 设 $\vv A$ 为 $n$ 阶方阵（$n\ge2$），$\rk\vv A=n-1$，$\vv\xi$ 是 $\vv{Ax}=\zero$ 的非零解. 证明 $\vv A^*$ 的每一列都是 $\vv\xi$ 的倍数.
			\item\label{ex:sol:B14} 用可确定性判据回答：方程组 $x_1+x_2+x_3=6$，$x_2+x_3+x_4=9$，$x_3+x_4+x_5=12$（例如三段相邻时间窗口内的观测总量）能确定 $x_1-x_4$ 吗？$x_1+x_2+x_3+x_4+x_5$ 呢？$x_1-x_4+x_5$ 呢？
			\item\label{ex:sol:B15} 设 $\vv A$ 为 $m\times n$ 矩阵，$\vv B$ 为 $n\times p$ 矩阵. 证明 $\dim\ker(\vv{AB})\le\dim\ker\vv A+\dim\ker\vv B$，并由此重新推出 Sylvester 不等式.
		\end{exB}
		
		\Cgroup
		\begin{exC}
			\item\label{ex:sol:C1} （Farkas 引理的二维形式）设 $\vv a_1,\vv a_2\in\mathbb{R}^2$ 线性无关，$\vv b\in\mathbb{R}^2$. 证明：下面两件事恰有一件成立：(1) 存在 $x_1,x_2\ge0$ 使 $x_1\vv a_1+x_2\vv a_2=\vv b$；(2) 存在 $\vv y$ 使 $\vv y^\trans\vv a_1\ge0$，$\vv y^\trans\vv a_2\ge0$，$\vv y^\trans\vv b<0$. 画图解释. 这是``无套利定价''与线性规划对偶理论的出发点.
			\item\label{ex:sol:C2} 设 $\vv A$ 是 $n$ 阶整数矩阵. 证明：若对任意整数向量 $\vv b$，$\vv{Ax}=\vv b$ 都有整数解，则 $\det\vv A=\pm1$.
			\item\label{ex:sol:C3} 证明：方阵 $\vv A$ 不用行交换就能做 LU 分解（$\vv L$ 单位下三角，$\vv U$ 上三角且可逆），当且仅当 $\vv A$ 的各阶顺序主子式都不为零；这时分解唯一.
			\item\label{ex:sol:C4} 设 $\vv A$ 是三对角矩阵（只有主对角线和相邻两条斜线上可能有非零元）. 证明：若它有 LU 分解，则 $\vv L,\vv U$ 都只有两条非零斜线；因此解三对角方程组只需与 $n$ 成正比的运算量. 这种方法称为``追赶法''.
			\item\label{ex:sol:C5} 用分块初等变换
			\[\begin{pmatrix}\vv I_n&\vv O\\ -\vv A&\vv I_m\end{pmatrix}\begin{pmatrix}\vv I_n&\vv B\\ \vv A&\vv O\end{pmatrix}\begin{pmatrix}\vv I_n&-\vv B\\ \vv O&\vv I_p\end{pmatrix}=\begin{pmatrix}\vv I_n&\vv O\\ \vv O&-\vv{AB}\end{pmatrix}\]
			给出 Sylvester 不等式 $\rk(\vv{AB})\ge\rk\vv A+\rk\vv B-n$ 的另一个证明.（提示：用第三章章末习题~\ref{ex:mat:B5} 的分块秩不等式.）
			\item\label{ex:sol:C6} 证明：$\mathbb{F}^n$ 中任意 $k$ 个线性流形（不必过原点）的交，要么是空集，要么仍是线性流形. 给出两个二维线性流形在 $\mathbb{R}^3$ 中交为空集的例子，并用秩的语言刻画这种情形.
			\item\label{ex:sol:C7} （列昂惕夫价格模型）设 $\vv A$ 为直接消耗系数矩阵，$\vv v$ 为各部门单位产品的增加值（工资、利润等）. 价格 $\vv p$ 满足``单位价格 $=$ 单位投入成本 $+$ 单位增加值''：$\vv p=\vv A^\trans\vv p+\vv v$. 说明它是第一章产量模型 $\vv x=\vv{Ax}+\vv d$ 的``转置''；证明在两个模型都可解时有 $\vv p^\trans\vv d=\vv v^\trans\vv x$，并解释其经济含义（国民收入的两种核算方法：支出法与收入法）.
			\item\label{ex:sol:C8} 设 $\vv A$ 为 $n$ 阶方阵. 证明存在正整数 $k$ 使 $\rk\vv A^k=\rk\vv A^{k+1}=\rk\vv A^{k+2}=\cdots$，且这样的最小的 $k$ 不超过 $n$.
		\end{exC}
	}
	
	% =====================================================================
	% 第七章　多项式
	% =====================================================================
	\newcommand{\ExPolynomial}{%
		\exsetheading
		\Agroup
		\begin{exA}
			\item\label{ex:poly:A1} 求商与余式：(1) $x^5-1$ 除以 $x^2+x+1$；(2) $2x^4-3x^3+4x^2-5x+6$ 除以 $x-2$（用综合除法）.
			\item\label{ex:poly:A2} 求 $x^{100}$ 分别除以 $x^2-1$，$x^2+1$，$(x-1)^2$ 所得的余式.
			\item\label{ex:poly:A3} 求 $a,b$，使 $x^2+x+1$ 整除 $x^4+ax^2+b$.（提示：设 $\omega$ 是 $x^2+x+1$ 的根，则 $\omega^3=1$.）
			\item\label{ex:poly:A4} 用辗转相除法求最大公因式：(1) $x^4+3x^3-x^2-4x-3$ 与 $3x^3+10x^2+2x-3$；(2) $x^3-2x^2+x$ 与 $x^2-1$.
			\item\label{ex:poly:A5} 证明：对任意多项式 $h$，$\gcd(f,g)=\gcd(f,\ g-hf)$. 这正是辗转相除法每一步的依据.
			\item\label{ex:poly:A6} 求 $u,v$ 使 $u(x)(x^2+1)+v(x)(x+1)=1$.
			\item\label{ex:poly:A7} 求 $2x^3-5x^2+x+2$ 的全部有理根，并在 $\mathbb{Q}$ 上把它分解为一次因式之积.
			\item\label{ex:poly:A8} 分别在 $\mathbb{Q}$，$\mathbb{R}$，$\mathbb{C}$ 上分解：(1) $x^4+x^2+1$；(2) $x^6-1$；(3) $x^3-2$.
			\item\label{ex:poly:A9} 判断 $1$ 是 $x^5-5x+4$ 的几重根.
			\item\label{ex:poly:A10} $a$ 取何值时，$x^3-3x+a$ 有重根？求出重根.
			\item\label{ex:poly:A11} 设 $\alpha_1,\alpha_2,\alpha_3$ 是 $x^3-2x^2+3x-4$ 的根. 不解方程，求 $\alpha_1+\alpha_2+\alpha_3$，$\alpha_1^2+\alpha_2^2+\alpha_3^2$，$\frac1{\alpha_1}+\frac1{\alpha_2}+\frac1{\alpha_3}$.
			\item\label{ex:poly:A12} 求次数最低的首一有理系数多项式，使 $1+\sqrt2$ 与 $\mathrm i$ 都是它的根.
			\item\label{ex:poly:A13} 求一个三次实系数首一多项式，以 $2$ 和 $1+\mathrm i$ 为根.
			\item\label{ex:poly:A14} 设 $f(x)=x^2-2x+1$，$\vv A=\begin{pmatrix}1&2\\0&1\end{pmatrix}$，计算 $f(\vv A)$；再设 $g(x)=x^3-x$，计算 $g(\vv A)$.
			\item\label{ex:poly:A15} 求过点 $(0,1)$，$(1,0)$，$(2,3)$，$(3,10)$ 的三次以下的插值多项式（用 Lagrange 公式）.
			\item\label{ex:poly:A16} 证明 $(x-1)\mid x^n-1$，并写出商式. 一般地，对任意 $a$，$(x-a)\mid x^n-a^n$.
			\item\label{ex:poly:A17} 把 $f(x)=x^4-2x^3+3$ 按 $x-1$ 的幂展开，即写成 $\sum c_k(x-1)^k$.（提示：反复用 $x-1$ 作综合除法，或用 $c_k=f^{(k)}(1)/k!$.）
			\item\label{ex:poly:A18} 多项式 $f$ 除以 $x-1$ 余 $3$，除以 $x-2$ 余 $5$. 求 $f$ 除以 $(x-1)(x-2)$ 的余式.
			\item\label{ex:poly:A19} 一个项目现在投入 $100$ 万元，一年后回收 $55$ 万元，两年后回收 $60.5$ 万元. 写出内部收益率满足的方程并精确求解.
			\item\label{ex:poly:A20} 证明：若 $f\mid g$ 且 $g\mid f$，则 $f=cg$（$c$ 为非零常数）.
			\item\label{ex:poly:A21} 设 $\gcd(f,g)=d$，$f=df_1$，$g=dg_1$. 证明 $\gcd(f_1,g_1)=1$.
		\end{exA}
		
		\Bgroup
		\begin{exB}
			\item\label{ex:poly:B1} 证明：若 $\gcd(f,g)=1$ 且 $\gcd(f,h)=1$，则 $\gcd(f,gh)=1$.
			\item\label{ex:poly:B2} 证明：$\gcd(x^m-1,\ x^n-1)=x^{\gcd(m,n)}-1$.（提示：模仿整数的辗转相除，先证 $x^m-1=x^{m-n}(x^n-1)+(x^{m-n}-1)$.）
			\item\label{ex:poly:B3} 设 $f=(x-1)^3(x+2)^2$. 计算 $f/\gcd(f,f')$，说明它恰好是``去掉重复因式''的结果. 一般地证明：若 $f=c\,p_1^{k_1}\cdots p_m^{k_m}$ 是标准分解，则 $f/\gcd(f,f')=c\,p_1\cdots p_m$.
			\item\label{ex:poly:B4} 证明：若 $p$ 不可约且 $p\mid f^2$，则 $p\mid f$. 举例说明``不可约''不能去掉.
			\item\label{ex:poly:B5} 设 $\ell_1,\ldots,\ell_n$ 是关于互异节点 $x_1,\ldots,x_n$ 的 Lagrange 基函数. 证明 $\sum_i\ell_i(x)=1$，$\sum_ix_i\ell_i(x)=x$，一般地 $\sum_ix_i^k\ell_i(x)=x^k$（$0\le k\le n-1$）.
			\item\label{ex:poly:B6} 证明：若复系数多项式 $f$ 满足 $f(x+1)=f(x)$ 对一切 $x$ 成立，则 $f$ 是常数.
			\item\label{ex:poly:B7} 证明 $x^3+x+1$ 在 $\mathbb{Q}$ 上不可约. 为什么对三次多项式只需检验有无有理根？四次多项式可以这样检验吗？
			\item\label{ex:poly:B8} 求 $x^n$ 除以 $x^2-3x+2$ 的余式.
			\item\label{ex:poly:B9} （Hermite 插值）求次数最低的多项式 $f$，使 $f(1)=1$，$f'(1)=2$，$f(2)=4$.（提示：设 $f=1+2(x-1)+c(x-1)^2$.）
			\item\label{ex:poly:B10} 设 $f$ 是有理系数多项式，$a,b\in\mathbb{Q}$，$\sqrt2\notin\mathbb{Q}$. 证明：若 $a+b\sqrt2$ 是 $f$ 的根，则 $a-b\sqrt2$ 也是 $f$ 的根.（提示：用 $(x-a)^2-2b^2$ 除 $f$.）
			\item\label{ex:poly:B11} 设 $f$ 是整系数多项式，$f(0)$ 与 $f(1)$ 都是奇数. 证明 $f$ 没有整数根.
			\item\label{ex:poly:B12} 求最小多项式意义下的``零化''：分别求出次数最低的首一多项式 $g$，使 $g(\vv A)=\vv O$，其中 $\vv A=\begin{pmatrix}0&-1\\1&0\end{pmatrix}$，$\vv A=\begin{pmatrix}2&0\\0&2\end{pmatrix}$，$\vv A=\begin{pmatrix}2&1\\0&2\end{pmatrix}$.
			\item\label{ex:poly:B13} 证明：现金流 $-C_0,\ C_1,\ldots,C_n$（$C_0>0$，$C_i\ge0$ 且不全为零）的净现值 $N(v)=-C_0+C_1v+\cdots+C_nv^n$ 在 $v>0$ 时严格递增，从而 $N(v)=0$ 至多一个正根；当 $C_1+\cdots+C_n>C_0$ 时这个正根落在 $(0,1)$ 中，即 IRR 唯一且为正.
			\item\label{ex:poly:B14} 设 $f$ 是 $n$ 次多项式，在 $n+1$ 个整数点 $0,1,\ldots,n$ 上都取整数值. 证明 $f$ 在一切整数点上都取整数值.（提示：用基 $\binom x0,\binom x1,\ldots,\binom xn$，其中 $\binom xk=\frac{x(x-1)\cdots(x-k+1)}{k!}$.）
			\item\label{ex:poly:B15} 证明 Lagrange 插值的``唯一性''也可以由线性代数推出：求值映射 $\mathbb{F}[x]_n\to\mathbb{F}^n$，$f\mapsto(f(x_1),\ldots,f(x_n))^\trans$ 是单射，从而是同构.
			\item\label{ex:poly:B16} 设 $p$ 是 $\mathbb{F}$ 上的不可约多项式，$f\in\mathbb{F}[x]$. 证明：$\gcd(p,f)=1$ 或 $p\mid f$，二者必居其一.
		\end{exB}
		
		\Cgroup
		\begin{exC}
			\item\label{ex:poly:C1} （Eisenstein 判别法）设 $f=a_nx^n+\cdots+a_0$ 是整系数多项式，$p$ 是素数，$p\nmid a_n$，$p\mid a_i$（$i<n$），$p^2\nmid a_0$. 证明 $f$ 在 $\mathbb{Q}$ 上不可约（可以先承认 Gauss 引理：整系数多项式若能分解为两个低次有理系数多项式之积，则也能分解为两个低次整系数多项式之积）. 由此说明 $x^n-2$ 对一切 $n$ 都不可约，所以 $\mathbb{Q}[x]$ 中有任意高次的不可约多项式.
			\item\label{ex:poly:C2} 设 $p$ 是素数. 证明分圆多项式 $\Phi_p(x)=x^{p-1}+x^{p-2}+\cdots+1$ 在 $\mathbb{Q}$ 上不可约.（提示：对 $\Phi_p(x+1)$ 用 Eisenstein 判别法.）
			\item\label{ex:poly:C3} （Newton 恒等式）记 $s_k=\alpha_1^k+\alpha_2^k+\alpha_3^k$，$\sigma_1,\sigma_2,\sigma_3$ 为初等对称多项式. 证明 $s_1=\sigma_1$，$s_2=\sigma_1s_1-2\sigma_2$，$s_3=\sigma_1s_2-\sigma_2s_1+3\sigma_3$，并求 $x^3-2x^2+3x-4$ 的根的立方和.
			\item\label{ex:poly:C4} （三次方程的判别式）设 $x^3+px+q$ 的三个复根为 $\alpha_1,\alpha_2,\alpha_3$，$\Delta=\prod_{i<j}(\alpha_i-\alpha_j)^2$. 证明 $\Delta=-4p^3-27q^2$，并说明（$p,q$ 为实数时）：$\Delta>0$ 时三个不同实根，$\Delta<0$ 时一实二虚，$\Delta=0$ 时有重根.
			\item\label{ex:poly:C5} （Chebyshev 多项式）证明存在 $n$ 次多项式 $T_n$ 使 $\cos n\theta=T_n(\cos\theta)$，且 $T_{n+1}=2xT_n-T_{n-1}$. 求 $T_n$ 的全部根. 由此证明：$\cos20^\circ$ 是 $8x^3-6x-1$ 的根，并验证这个三次多项式在 $\mathbb{Q}$ 上不可约. （结合``尺规作图只能得到经若干次开平方得到的数''，这说明 $60^\circ$ 角不能用尺规三等分.）
			\item\label{ex:poly:C6} 证明：在 $\mathbb{R}$ 上处处非负的实系数多项式，都能写成两个实系数多项式的平方和.（提示：利用实根的重数必为偶数，以及 $|a+b\mathrm i|^2=a^2+b^2$ 与 $(a^2+b^2)(c^2+d^2)=(ac-bd)^2+(ad+bc)^2$.）
			\item\label{ex:poly:C7} 设 $f\in\mathbb{Q}[x]$ 不可约，$g\in\mathbb{Q}[x]$ 与 $f$ 有一个公共复根. 证明 $f\mid g$. 由此说明：一个复数若是某个有理系数多项式的根，则以它为根的次数最低的首一有理系数多项式（称为它的\textbf{极小多项式}）唯一，且不可约.
			\item\label{ex:poly:C8} （$\mathbb{F}_2$ 上的多项式）在系数按模 $2$ 运算的多项式中，证明 $x^2+x+1$ 是唯一的二次不可约多项式，并找出全部三次不可约多项式. 这类多项式是循环冗余校验（CRC）与密码学的基础.
			\item\label{ex:poly:C9} 设 $f$ 是实系数多项式，它的系数序列的``变号次数''为 $V$（忽略零系数）. Descartes 符号法则说：$f$ 的正根个数（按重数计）不超过 $V$，且与 $V$ 的差是偶数. 用它重新解释第七章``内部收益率''建模实验中两个项目的结论.
		\end{exC}
	}
	
	% =====================================================================
	% 第八章　特征值与特征向量
	% =====================================================================
	\newcommand{\ExEigen}{%
		\exsetheading
		\Agroup
		\begin{exA}
			\item\label{ex:eig:A1} 求下列矩阵的特征值与特征向量（在 $\mathbb{C}$ 上考虑）：
			\[\begin{pmatrix}3&1\\0&2\end{pmatrix},\quad\begin{pmatrix}1&2\\2&1\end{pmatrix},\quad\begin{pmatrix}0&1\\-1&0\end{pmatrix},\quad\begin{pmatrix}4&-2\\1&1\end{pmatrix},\]\[\begin{pmatrix}2&0&0\\0&1&1\\0&1&1\end{pmatrix},\quad\begin{pmatrix}1&2&2\\2&1&2\\2&2&1\end{pmatrix}.\]
			\item\label{ex:eig:A2} 判断 A1 中各矩阵能否在 $\mathbb{R}$ 上对角化；能的，求可逆矩阵 $\vv P$ 使 $\vv P^{-1}\vv{AP}$ 为对角矩阵.
			\item\label{ex:eig:A3} 判断下列矩阵能否对角化：$\begin{pmatrix}1&1\\0&2\end{pmatrix}$，$\begin{pmatrix}2&1\\0&2\end{pmatrix}$，$\begin{pmatrix}2&0&0\\0&2&0\\1&0&2\end{pmatrix}$，$\begin{pmatrix}1&1&0\\0&2&0\\0&0&1\end{pmatrix}$.
			\item\label{ex:eig:A4} 设 $3$ 阶方阵 $\vv A$ 的特征值为 $1,2,3$. 求 $\det\vv A$，$\tr\vv A$，$\vv A^{-1}$ 的特征值，$\vv A^2-2\vv A+\vv I$ 的特征值，$\det(\vv A^*)$，$\det(\vv A^2+\vv I)$.
			\item\label{ex:eig:A5} 已知二阶方阵的迹为 $5$，行列式为 $6$，求它的特征值. 若迹为 $2$、行列式为 $5$ 呢？
			\item\label{ex:eig:A6} 证明：$\vv A$ 与 $\vv A^\trans$ 有相同的特征多项式. 它们的特征向量是否相同？举例说明.
			\item\label{ex:eig:A7} 证明：三角矩阵的特征值恰为其对角元.
			\item\label{ex:eig:A8} 证明：若 $\vv A^k=\vv O$，则 $\vv A$ 的特征值全为 $0$；若 $\vv A^2=\vv A$，则特征值只能是 $0$ 或 $1$；若 $\vv A^2=\vv I$，则特征值只能是 $\pm1$.
			\item\label{ex:eig:A9} 设 $\vv A=\begin{pmatrix}4&-2\\1&1\end{pmatrix}$. 利用对角化计算 $\vv A^n$.
			\item\label{ex:eig:A10} 设品牌转移矩阵为 $\vv P=\begin{pmatrix}0.9&0.2\\0.1&0.8\end{pmatrix}$. 求稳态分布，并利用对角化求 $\lim_{n\to\infty}\vv P^n$. 解释极限矩阵两列相同的含义.
			\item\label{ex:eig:A11} 已知 $\vv\xi=(1,1,-1)^\trans$ 是 $\vv A=\begin{pmatrix}2&-1&2\\5&a&3\\-1&b&-2\end{pmatrix}$ 的特征向量. 求 $a,b$ 与对应的特征值.
			\item\label{ex:eig:A12} 设 $\vv A=\begin{pmatrix}0&0&1\\1&1&x\\1&0&0\end{pmatrix}$. $x$ 取何值时 $\vv A$ 可对角化？
			\item\label{ex:eig:A13} 对二阶方阵直接验证 Cayley–Hamilton 定理：$\vv A^2-(\tr\vv A)\vv A+(\det\vv A)\vv I=\vv O$.
			\item\label{ex:eig:A14} 求下列矩阵的最小多项式：$\diag(1,1,2)$，$\begin{pmatrix}1&1&0\\0&1&0\\0&0&2\end{pmatrix}$，$\begin{pmatrix}0&1&0\\0&0&1\\0&0&0\end{pmatrix}$，$\begin{pmatrix}1&1&1\\1&1&1\\1&1&1\end{pmatrix}$.
			\item\label{ex:eig:A15} 判断 $\begin{pmatrix}1&1\\0&1\end{pmatrix}$ 与 $\begin{pmatrix}1&0\\1&1\end{pmatrix}$ 是否相似；若相似，求出 $\vv P$.
			\item\label{ex:eig:A16} 设 $\vv A$ 是上三角矩阵. 证明 $\spn(\vv e_1)$，$\spn(\vv e_1,\vv e_2)$，……都是 $\vv A$ 的不变子空间.
			\item\label{ex:eig:A17} 解差分方程 $a_{k+1}=a_k+2a_{k-1}$，$a_0=0$，$a_1=1$.
			\item\label{ex:eig:A18} （种群模型）某种群分为幼年与成年两组，第 $k$ 年的数量记为 $\vv x_k=(y_k,z_k)^\trans$. 每只成年个体每年平均产下 $1.1$ 只幼体，幼体有 $60\%$ 存活并长为成年，成年个体有一半活到下一年，即
			$y_{k+1}=1.1z_k$，$z_{k+1}=0.6y_k+0.5z_k$. 求矩阵的特征值与特征向量，说明种群的长期增长率和幼体与成年个体的稳定比例.
			\item\label{ex:eig:A19} 证明：实二阶方阵 $\vv A$ 没有实特征值，当且仅当 $(\tr\vv A)^2<4\det\vv A$.
			\item\label{ex:eig:A20} 求平面上向直线 $y=2x$ 的正投影、关于该直线的反射这两个变换的特征值与特征向量（不必写出矩阵，从几何上直接回答），再写出矩阵加以验证.
		\end{exA}
		
		\Bgroup
		\begin{exB}
			\item\label{ex:eig:B1} 设 $\vv u,\vv v\in\mathbb{R}^n$ 非零，$\vv A=\vv u\vv v^\trans$. 求 $\vv A$ 的全部特征值及其重数，并证明：$\vv A$ 可对角化当且仅当 $\vv v^\trans\vv u\ne0$.
			\item\label{ex:eig:B2} 设 $\vv A$ 有 $n$ 个互异的特征值，$\vv{AB}=\vv{BA}$. 证明 $\vv A$ 的特征向量都是 $\vv B$ 的特征向量，从而 $\vv B$ 也可对角化；进一步证明 $\vv B$ 是 $\vv A$ 的多项式.
			\item\label{ex:eig:B3} 证明：若 $\vv A$ 可对角化，则 $\vv A^k$ 与 $f(\vv A)$ 都可对角化. 逆命题成立吗？
			\item\label{ex:eig:B4} 设 $\vv A$ 可对角化，$W$ 是 $\vv A$ 的不变子空间. 证明 $\vv A|_W$ 也可对角化.（提示：最小多项式.）
			\item\label{ex:eig:B5} 证明：满足 $\vv A^2=-\vv I$ 的二阶实方阵都与 $\begin{pmatrix}0&-1\\1&0\end{pmatrix}$ 相似.
			\item\label{ex:eig:B6} （Gershgorin 圆盘定理）证明：方阵 $\vv A$ 的每个特征值都落在某个圆盘 $\{z\in\mathbb{C}:|z-a_{ii}|\le\sum_{j\ne i}|a_{ij}|\}$ 中. 由此证明：严格对角占优的矩阵（$|a_{ii}|>\sum_{j\ne i}|a_{ij}|$）可逆.
			\item\label{ex:eig:B7} 设 $\vv A$ 的每行元素之和都是 $s$. 证明 $s$ 是 $\vv A$ 的特征值. 若 $\vv A$ 可逆，求 $\vv A^{-1}$ 每行元素之和.
			\item\label{ex:eig:B8} 设 $\vv A=\begin{pmatrix}a&b&c\\c&a&b\\b&c&a\end{pmatrix}$（循环矩阵），$\omega=\mathrm e^{2\pi\mathrm i/3}$. 证明 $(1,\omega^k,\omega^{2k})^\trans$（$k=0,1,2$）都是 $\vv A$ 的特征向量，并求对应的特征值. 由此重新得到第五章章末习题~\ref{ex:det:B12} 的行列式公式.
			\item\label{ex:eig:B9} 解差分方程 $a_{k+1}=4a_k-4a_{k-1}$，$a_0=1$，$a_1=4$. 解释通项中为什么出现 $k\cdot2^k$.
			\item\label{ex:eig:B10} 证明 $\tr(\vv A^k)=\lambda_1^k+\cdots+\lambda_n^k$. 若三阶方阵满足 $\tr\vv A=0$，$\tr\vv A^2=6$，$\det\vv A=0$，求它的特征值.
			\item\label{ex:eig:B11} 证明：若 $\vv A^3=\vv A$，则 $\vv A$ 可对角化，且特征值在 $\{0,1,-1\}$ 中.
			\item\label{ex:eig:B12} $a,b,c$ 满足什么条件时，$\begin{pmatrix}2&0&0\\a&2&0\\b&c&1\end{pmatrix}$ 可对角化？
			\item\label{ex:eig:B13} 三个网页之间的链接为：网页 $1$ 链接到 $2,3$，网页 $2$ 链接到 $3$，网页 $3$ 链接到 $1$.（与判例不同，网页之间可以``互相引用''，形成环路.） 取 $d=1$ 写出转移矩阵，手算它的稳态分布.
			\item\label{ex:eig:B14} （矩阵指数）对可对角化矩阵 $\vv A=\vv P\vv D\vv P^{-1}$，定义 $\mathrm e^{t\vv A}=\vv P\diag(\mathrm e^{t\lambda_1},\ldots,\mathrm e^{t\lambda_n})\vv P^{-1}$. 对 $\vv A=\begin{pmatrix}0&1\\1&0\end{pmatrix}$ 计算 $\mathrm e^{t\vv A}$，并验证 $\vv x(t)=\mathrm e^{t\vv A}\vv x_0$ 满足微分方程组 $x_1'=x_2$，$x_2'=x_1$.
			\item\label{ex:eig:B15} 三状态马尔可夫链的转移矩阵为 $\begin{pmatrix}0.5&0.25&0.25\\0.25&0.5&0.25\\0.25&0.25&0.5\end{pmatrix}$. 求稳态分布与全部特征值，说明偏离稳态的部分每一步缩小为原来的多少.
			\item\label{ex:eig:B16} 证明：若 $\lambda$ 是 $\vv A$ 的 $k$ 重特征值（代数重数），则 $\rk(\lambda\vv I-\vv A)\ge n-k$.
			\item\label{ex:eig:B17} 对第三章``判例引证网络''中的四份判决，分别取阻尼系数 $d=0.5,\ 0.85,\ 0.95,\ 1$，按第八章建模实验的方法求 PageRank 得分（可用计算机），列表比较各判决的得分与排序. $d$ 越接近 $1$，得分越``集中''在哪些判决上？为什么？
		\end{exB}
		
		\Cgroup
		\begin{exC}
			\item\label{ex:eig:C1} 设 $\vv A$ 为 $m\times n$ 矩阵，$\vv B$ 为 $n\times m$ 矩阵. 证明 $\lambda^n\det(\lambda\vv I_m-\vv{AB})=\lambda^m\det(\lambda\vv I_n-\vv{BA})$. 由此说明 $\vv{AB}$ 与 $\vv{BA}$ 的非零特征值（连同重数）相同.（提示：计算 $\begin{pmatrix}\vv I_m&-\vv A\\ \vv O&\lambda\vv I_n\end{pmatrix}\begin{pmatrix}\lambda\vv I_m&\vv A\\ \vv B&\vv I_n\end{pmatrix}$ 等两种乘积的行列式.）
			\item\label{ex:eig:C2} （Schur 三角化）证明：每个复方阵都相似于一个上三角矩阵. 由此给出 Cayley–Hamilton 定理的另一个证明：对上三角矩阵 $\vv T$，证明 $(\vv T-t_{11}\vv I)\cdots(\vv T-t_{nn}\vv I)=\vv O$.
			\item\label{ex:eig:C3} 求下列矩阵的 Jordan 标准形：
			\[\begin{pmatrix}2&1&0\\0&2&0\\0&0&2\end{pmatrix},\quad\begin{pmatrix}1&1&1\\0&1&1\\0&0&1\end{pmatrix},\quad\begin{pmatrix}0&1&0&0\\0&0&0&0\\0&0&0&1\\0&0&0&0\end{pmatrix}.\]
			特征多项式为 $(\lambda-2)^4$ 的 $4$ 阶复矩阵，在相似的意义下共有几类？
			\item\label{ex:eig:C4} 证明：若一组可对角化的矩阵两两可交换，则它们可以\textbf{同时}对角化（存在同一个 $\vv P$ 使它们都化为对角矩阵）.
			\item\label{ex:eig:C5} 设 $\vv A$ 可对角化. 证明 $\vv A^k\to\vv O$（按元素）当且仅当 $\vv A$ 的所有特征值的模都小于 $1$. 用它说明：投入产出模型中，若直接消耗系数矩阵 $\vv A$ 可对角化且特征值的模都小于 $1$，则 $(\vv I-\vv A)^{-1}=\vv I+\vv A+\vv A^2+\cdots$.
			\item\label{ex:eig:C6} （友矩阵）证明 $\begin{pmatrix}0&0&\cdots&0&-a_0\\1&0&\cdots&0&-a_1\\0&1&\cdots&0&-a_2\\ \vdots&&\ddots&&\vdots\\0&0&\cdots&1&-a_{n-1}\end{pmatrix}$ 的特征多项式为 $\lambda^n+a_{n-1}\lambda^{n-1}+\cdots+a_0$. 由此说明：求多项式的根等价于求矩阵的特征值（数值软件正是这样求多项式的根的）.
			\item\label{ex:eig:C7} （赌徒破产）一位赌徒有 $k$ 元（$0\le k\le N$），每局以概率 $\frac12$ 赢或输 $1$ 元，输光或达到 $N$ 元时停止. 设 $p_k$ 为最终达到 $N$ 元的概率，证明 $p_k=\frac12(p_{k+1}+p_{k-1})$，并用特征根方法解出 $p_k=k/N$. 赢的概率改为 $p\ne\frac12$ 时呢？
			\item\label{ex:eig:C8} 证明：$n$ 阶三对角矩阵 $\vv T=\begin{pmatrix}2&-1&&\\-1&2&\ddots&\\&\ddots&\ddots&-1\\&&-1&2\end{pmatrix}$ 的特征值为 $2-2\cos\frac{k\pi}{n+1}$（$k=1,\ldots,n$），对应特征向量的第 $j$ 个分量为 $\sin\frac{jk\pi}{n+1}$. 对 $n=2,3$ 直接验证.
			\item\label{ex:eig:C9} （Perron–Frobenius 的二阶情形）设 $\vv A$ 是元素全为正的二阶矩阵. 证明 $\vv A$ 有两个互异实特征值，较大者 $r$ 为正且大于另一个的绝对值，属于 $r$ 的特征向量可取为各分量都为正. 由此解释：为什么 PageRank 引入阻尼系数后，``全部元素为正''能保证排序唯一.
			\item\label{ex:eig:C10} 证明：满足 $\vv A^2=-\vv I$ 的 $n$ 阶实方阵存在当且仅当 $n$ 为偶数，且这时 $\vv A$ 相似于 $\diag(\vv J,\ldots,\vv J)$，$\vv J=\begin{pmatrix}0&-1\\1&0\end{pmatrix}$.
		\end{exC}
	}
	
	% =====================================================================
	% 第九章　内积空间
	% =====================================================================
	\newcommand{\ExInnerProduct}{%
		\exsetheading
		\Agroup
		\begin{exA}
			\item\label{ex:ip:A1} 在 $\mathbb{R}^4$ 中，设 $\vv\alpha=(1,2,2,0)^\trans$，$\vv\beta=(2,0,1,2)^\trans$. 求 $\|\vv\alpha\|$，$\|\vv\beta\|$，$\langle\vv\alpha,\vv\beta\rangle$，它们的夹角，以及 $\|\vv\alpha-\vv\beta\|$.
			\item\label{ex:ip:A2} 判断下列二元函数是否为 $\mathbb{R}^2$ 上的内积：
			(1) $\langle\vv x,\vv y\rangle=x_1y_1+2x_2y_2$；(2) $\langle\vv x,\vv y\rangle=x_1y_2+x_2y_1$；(3) $\langle\vv x,\vv y\rangle=2x_1y_1+x_1y_2+x_2y_1+x_2y_2$；(4) $\langle\vv x,\vv y\rangle=x_1y_1-x_2y_2$.
			\item\label{ex:ip:A3} 用 Gram–Schmidt 方法把下列向量组正交化并单位化：
			(1) $(1,1,1)^\trans$，$(1,2,0)^\trans$，$(0,1,2)^\trans$；(2) $(1,1,0,0)^\trans$，$(1,0,1,0)^\trans$，$(-1,0,0,1)^\trans$.
			\item\label{ex:ip:A4} 求与 $(1,1,1)^\trans$ 和 $(1,0,-1)^\trans$ 都正交的单位向量.
			\item\label{ex:ip:A5} 已知 $\vv e_1=\frac13(1,2,2)^\trans$，$\vv e_2=\frac13(2,1,-2)^\trans$，$\vv e_3=\frac13(2,-2,1)^\trans$. 验证它们构成 $\mathbb{R}^3$ 的标准正交基，并只用内积求 $(3,6,0)^\trans$ 在这组基下的坐标.
			\item\label{ex:ip:A6} 求 $\vv\alpha=(1,2,3,4)^\trans$ 在 $W=\spn\big((1,1,0,0)^\trans,(0,0,1,1)^\trans\big)$ 上的正交投影，以及 $\vv\alpha$ 到 $W$ 的距离.
			\item\label{ex:ip:A7} 用最小二乘法求过数据点 $(0,1)$，$(1,2)$，$(2,2)$，$(3,4)$，$(4,6)$ 的最佳直线 $y=a+bx$，并计算残差平方和.
			\item\label{ex:ip:A8} 求 $a,b,c$ 使 $\begin{pmatrix}\frac1{\sqrt2}&a&\frac1{\sqrt6}\\0&b&-\frac2{\sqrt6}\\-\frac1{\sqrt2}&c&\frac1{\sqrt6}\end{pmatrix}$ 为正交矩阵.
			\item\label{ex:ip:A9} 证明：若 $\vv\alpha_1,\ldots,\vv\alpha_k$ 两两正交，则 $\|\vv\alpha_1+\cdots+\vv\alpha_k\|^2=\|\vv\alpha_1\|^2+\cdots+\|\vv\alpha_k\|^2$.
			\item\label{ex:ip:A10} 证明：$\big|\|\vv\alpha\|-\|\vv\beta\|\big|\le\|\vv\alpha-\vv\beta\|$.
			\item\label{ex:ip:A11} 用 Cauchy–Schwarz 不等式证明：(1) $(a_1+\cdots+a_n)^2\le n(a_1^2+\cdots+a_n^2)$；(2) 若 $a_i>0$，则 $(a_1+\cdots+a_n)\big(\frac1{a_1}+\cdots+\frac1{a_n}\big)\ge n^2$.
			\item\label{ex:ip:A12} 在 $C[0,1]$ 中取 $\langle f,g\rangle=\int_0^1f(x)g(x)\,\mathrm dx$. 求 $\|1\|$，$\|x\|$，$1$ 与 $x$ 的夹角，并把 $1,x$ 正交化.
			\item\label{ex:ip:A13} 设 $\vv A=\begin{pmatrix}1&1&0&1\\0&1&1&1\end{pmatrix}$. 求 $\ker\vv A$ 与行空间 $\imm\vv A^\trans$ 的标准正交基，并验证它们互为正交补.
			\item\label{ex:ip:A14} 求 $\vv A=\begin{pmatrix}1&1\\1&0\\0&1\end{pmatrix}$ 的 QR 分解.
			\item\label{ex:ip:A15} 设 $\vv u=\frac13(1,2,2)^\trans$. 写出 Householder 矩阵 $\vv H=\vv I-2\vv u\vv u^\trans$，验证它是对称的正交矩阵，并求 $\vv H\vv u$ 与 $\vv H(2,-1,0)^\trans$.
			\item\label{ex:ip:A16} 判断下列正交矩阵是旋转还是反射；若是旋转求旋转角，若是反射求反射轴：$\frac15\begin{pmatrix}3&-4\\4&3\end{pmatrix}$，$\frac15\begin{pmatrix}3&4\\4&-3\end{pmatrix}$.
			\item\label{ex:ip:A17} 对``过原点的直线''模型 $y=bx$，推导最小二乘解 $\hat b=\sum x_iy_i\big/\sum x_i^2$.
			\item\label{ex:ip:A18} 在 $\mathbb{R}^3$ 中，求点 $(1,2,3)$ 到直线 $\vv x=t(1,1,1)^\trans$ 的距离.
			\item\label{ex:ip:A19} 证明：两组标准正交基之间的过渡矩阵是正交矩阵.
			\item\label{ex:ip:A20} 三个法条的词频向量（关键词：违约、赔偿、合同、离婚）分别为 $(2,1,3,0)^\trans$，$(0,1,0,4)^\trans$，$(1,1,2,0)^\trans$，查询为 $(1,1,1,0)^\trans$. 分别用内积和余弦相似度给三个法条排序，结果有何不同？
			\item\label{ex:ip:A21} 设 $\vv u\in\mathbb{R}^n$ 是单位向量. 证明 $\vv P=\vv u\vv u^\trans$ 满足 $\vv P^\trans=\vv P$，$\vv P^2=\vv P$，且对任意 $\vv x$，$\vv P\vv x$ 是 $\vv x$ 在 $\spn(\vv u)$ 上的正交投影. 若 $\vv u$ 不是单位向量，投影矩阵应当怎样修改？
		\end{exA}
		
		\Bgroup
		\begin{exB}
			\item\label{ex:ip:B1} （极化恒等式）证明实内积空间中 $\langle\vv\alpha,\vv\beta\rangle=\frac14\big(\|\vv\alpha+\vv\beta\|^2-\|\vv\alpha-\vv\beta\|^2\big)$. 再证明：$\mathbb{R}^2$ 上的``范数'' $\|\vv x\|_1=|x_1|+|x_2|$ 不满足平行四边形法则，因而不能来自任何内积.
			\item\label{ex:ip:B2} （Bessel 不等式）设 $\vv e_1,\ldots,\vv e_k$ 是标准正交组. 证明 $\sum_{i=1}^k\langle\vv\alpha,\vv e_i\rangle^2\le\|\vv\alpha\|^2$，等号成立当且仅当 $\vv\alpha\in\spn(\vv e_1,\ldots,\vv e_k)$.
			\item\label{ex:ip:B3} 设实矩阵 $\vv P$ 满足 $\vv P^\trans=\vv P$，$\vv P^2=\vv P$. 证明 $\vv P$ 是向 $\imm\vv P$ 的正交投影的矩阵. 若只有 $\vv P^2=\vv P$，它是什么样的投影？举例说明.
			\item\label{ex:ip:B4} （加权最小二乘）若各观测值的可靠程度不同，给第 $i$ 个残差权重 $w_i>0$，最小化 $\sum w_i(y_i-a-bx_i)^2$. 证明正规方程变为 $\vv A^\trans\vv W\vv A\hat{\vv\beta}=\vv A^\trans\vv W\vv y$，$\vv W=\diag(w_1,\ldots,w_n)$. 它对应于哪个内积下的正交投影？
			\item\label{ex:ip:B5} 某商品的销售额 $y$ 随时间 $t$ 近似按指数增长 $y=c\,\mathrm e^{kt}$. 已知 $t=0,1,2,3$ 时 $y$ 约为 $2.0,2.7,3.7,5.0$. 两边取对数化为线性模型，用最小二乘估计 $c,k$. 这样做最小化的是哪个量？与直接最小化 $\sum(y_i-c\mathrm e^{kt_i})^2$ 有何不同？
			\item\label{ex:ip:B6} （Gram 行列式）设 $\vv\alpha_1,\ldots,\vv\alpha_k\in\mathbb{R}^n$，$\vv G=(\langle\vv\alpha_i,\vv\alpha_j\rangle)$. 证明 $\det\vv G\ge0$，等号成立当且仅当它们线性相关. 在 $k=2$ 时说明 $\sqrt{\det\vv G}$ 是两向量张成的平行四边形的面积.
			\item\label{ex:ip:B7} 证明：$\mathbb{R}^n$ 中点 $\vv x_0$ 到超平面 $\vv a^\trans\vv x=b$（$\vv a\ne\zero$）的距离为 $|\vv a^\trans\vv x_0-b|/\|\vv a\|$.
			\item\label{ex:ip:B8} 证明：实正交矩阵的复特征值的模都等于 $1$.
			\item\label{ex:ip:B9} （Fourier 逼近）在 $C[-\pi,\pi]$ 中取 $\langle f,g\rangle=\int_{-\pi}^\pi fg\,\mathrm dx$. 证明 $1,\cos x,\sin x,\cos2x,\sin2x$ 两两正交. 求 $f(x)=x$ 在它们张成的子空间上的正交投影.
			\item\label{ex:ip:B10} 设 $\vv A$ 为实矩阵. 证明：在 $\vv{Ax}=\vv y$ 的全部最小二乘解中，范数最小的那一个唯一，并且它属于行空间 $\imm\vv A^\trans$.
			\item\label{ex:ip:B11} 设 $\vv x,\vv y\in\mathbb{R}^n$，$\|\vv x\|=\|\vv y\|$，$\vv x\ne\vv y$. 令 $\vv u=(\vv x-\vv y)/\|\vv x-\vv y\|$. 证明 Householder 矩阵 $\vv I-2\vv u\vv u^\trans$ 把 $\vv x$ 变为 $\vv y$. 这是用 Householder 反射做 QR 分解的基础.
			\item\label{ex:ip:B12} 设两个变量的样本标准差分别为 $s_x,s_y$，相关系数为 $r$. 证明回归直线 $y=a+bx$ 的斜率为 $b=r\,s_y/s_x$，并且直线经过点 $(\bar x,\bar y)$.
			\item\label{ex:ip:B13} 证明：保持\textbf{夹角}的可逆线性变换（对一切非零向量，$\mathscr T\vv\alpha$ 与 $\mathscr T\vv\beta$ 的夹角等于 $\vv\alpha$ 与 $\vv\beta$ 的夹角）必是正交变换的非零倍数.
			\item\label{ex:ip:B14} 在回归中加入一个``虚拟变量''：设有两组数据，第一组 $y$ 值为 $3,5,4$，第二组为 $7,8,9$. 以 $\vv A=(\vv 1,\vv d)$ 为设计矩阵（$\vv d$ 在第二组为 $1$、第一组为 $0$），求最小二乘解，并解释两个系数的含义.
			\item\label{ex:ip:B15} 证明：在 $\mathbb{R}^2$ 中，任意两个旋转可交换，而一个旋转与一个反射一般不可交换；两个反射的乘积是旋转.
		\end{exB}
		
		\Cgroup
		\begin{exC}
			\item\label{ex:ip:C1} （Hadamard 不等式）用 QR 分解证明：对 $n$ 阶实方阵 $\vv A=(\vv a_1,\ldots,\vv a_n)$，$|\det\vv A|\le\|\vv a_1\|\cdots\|\vv a_n\|$，并说明它的几何意义（平行多面体的体积不超过棱长之积）.
			\item\label{ex:ip:C2} （Gram–Schmidt 的数值不稳定）设 $\varepsilon$ 很小，对 $\vv a_1=(1,\varepsilon,0,0)^\trans$，$\vv a_2=(1,0,\varepsilon,0)^\trans$，$\vv a_3=(1,0,0,\varepsilon)^\trans$ 作 Gram–Schmidt，在计算中把 $1+\varepsilon^2$ 舍入为 $1$. 证明所得的 $\vv e_2,\vv e_3$ 之间的夹角是 $60^\circ$ 而不是 $90^\circ$. 再用``修正的 Gram–Schmidt''（每得到一个 $\vv e_k$，立即从\textbf{剩余的}所有向量中减去它的分量）重做，比较结果.
			\item\label{ex:ip:C3} （Legendre 多项式）令 $P_n(x)=\frac1{2^nn!}\frac{\mathrm d^n}{\mathrm dx^n}(x^2-1)^n$. 证明 $P_0,P_1,P_2,\ldots$ 在 $\langle f,g\rangle=\int_{-1}^1fg\,\mathrm dx$ 下两两正交，并求 $\|P_n\|$.
			\item\label{ex:ip:C4} （Gauss–Markov 定理的线性代数部分）设 $\vv A$ 列满秩. 若 $\vv C\vv y$ 是参数 $\vv\beta$ 的``线性无偏估计''，即 $\vv{CA}=\vv I$. 令 $\vv C_0=(\vv A^\trans\vv A)^{-1}\vv A^\trans$（最小二乘）. 证明 $\vv C\vv C^\trans-\vv C_0\vv C_0^\trans$ 是半正定矩阵（第十一章）. 学过概率论后可以知道，这说明最小二乘估计的方差在一切线性无偏估计中最小.
			\item\label{ex:ip:C5} （复内积空间）在 $\mathbb{C}^n$ 上定义 $\langle\vv x,\vv y\rangle=\sum x_i\bar y_i$. 证明它满足共轭对称性、对第一个变量线性、正定性；证明 Cauchy–Schwarz 不等式仍成立. 满足 $\bar{\vv U}^\trans\vv U=\vv I$ 的复矩阵称为\textbf{酉矩阵}，证明它保持这个内积.
			\item\label{ex:ip:C6} 设 $f:\mathbb{R}^n\to\mathbb{R}^n$ 保持距离（$\|f(\vv x)-f(\vv y)\|=\|\vv x-\vv y\|$）且 $f(\zero)=\zero$. 证明 $f$ 保持内积，进而证明 $f$ 是线性的正交变换. 由此说明：保持距离的变换都可以分解为正交变换与平移的复合.
			\item\label{ex:ip:C7} （伪逆）设 $\vv A$ 列满秩，$\vv A^+=(\vv A^\trans\vv A)^{-1}\vv A^\trans$. 验证 Penrose 四条件：$\vv{AA}^+\vv A=\vv A$，$\vv A^+\vv{AA}^+=\vv A^+$，$(\vv{AA}^+)^\trans=\vv{AA}^+$，$(\vv A^+\vv A)^\trans=\vv A^+\vv A$. 说明 $\vv{AA}^+$ 是向哪个子空间的正交投影.
			\item\label{ex:ip:C8} 求空间中两条异面直线 $\vv x=\vv p+s\vv u$ 与 $\vv x=\vv q+t\vv v$ 之间的最短距离：把问题写成关于 $(s,t)$ 的最小二乘问题，并对 $\vv p=\zero$，$\vv u=(1,0,0)^\trans$，$\vv q=(0,1,1)^\trans$，$\vv v=(0,1,-1)^\trans$ 计算.
		\end{exC}
	}
	
	% =====================================================================
	% 第十章　矩阵分解
	% =====================================================================
	\newcommand{\ExDecomposition}{%
		\exsetheading
		\Agroup
		\begin{exA}
			\item\label{ex:dec:A1} 求正交矩阵 $\vv Q$ 使 $\vv Q^\trans\vv{AQ}$ 为对角矩阵：
			\[\begin{pmatrix}1&2\\2&1\end{pmatrix},\qquad\begin{pmatrix}3&-1\\-1&3\end{pmatrix},\]\[\begin{pmatrix}2&2&0\\2&2&0\\0&0&1\end{pmatrix},\qquad\begin{pmatrix}1&1&1\\1&1&1\\1&1&1\end{pmatrix}.\]
			\item\label{ex:dec:A2} 写出 $\begin{pmatrix}3&-1\\-1&3\end{pmatrix}$ 的谱分解 $\lambda_1\vv q_1\vv q_1^\trans+\lambda_2\vv q_2\vv q_2^\trans$，并验证两个投影矩阵之和为 $\vv I$、之积为 $\vv O$.
			\item\label{ex:dec:A3} 求二阶实对称矩阵 $\vv A$，使它的特征值为 $1,2$，且 $(1,1)^\trans$ 是属于 $1$ 的特征向量.
			\item\label{ex:dec:A4} 设三阶实对称矩阵 $\vv A$ 的特征值为 $1,1,-1$，属于 $-1$ 的特征向量为 $(0,1,1)^\trans$. 求 $\vv A$.
			\item\label{ex:dec:A5} 求 $\vv x^\trans\vv{Ax}$ 在单位球面 $\|\vv x\|=1$ 上的最大值与最小值，以及取到它们的 $\vv x$：$\vv A=\begin{pmatrix}2&2&0\\2&2&0\\0&0&1\end{pmatrix}$.
			\item\label{ex:dec:A6} 求 $x_1^2+4x_1x_2+x_2^2$ 在单位圆 $x_1^2+x_2^2=1$ 上的最大值与最小值.
			\item\label{ex:dec:A7} 求下列矩阵的奇异值分解：$\begin{pmatrix}2&0\\0&-3\end{pmatrix}$，$\begin{pmatrix}1&1\\1&1\end{pmatrix}$，$(3,4)$（$1\times2$ 矩阵）.
			\item\label{ex:dec:A8} 求 $\begin{pmatrix}1&1\\0&1\end{pmatrix}$ 的奇异值，并与它的特征值比较.（答案与黄金分割比有关.）
			\item\label{ex:dec:A9} 求 $\begin{pmatrix}3&1\\1&3\end{pmatrix}$ 的最佳秩一近似，并计算误差 $\|\vv A-\vv A_1\|_F$.
			\item\label{ex:dec:A10} 设 $\vv A=\begin{pmatrix}5&4\\4&5\end{pmatrix}$. 利用谱分解求 $\vv A^{10}$（写成两个矩阵之和的形式即可），并求对称正定矩阵 $\vv B$ 使 $\vv B^2=\vv A$.
			\item\label{ex:dec:A11} 证明：特征值全为 $0$ 的实对称矩阵是零矩阵. 对非对称矩阵举出反例.
			\item\label{ex:dec:A12} 证明：既对称又正交的实矩阵 $\vv A$ 满足 $\vv A^2=\vv I$，其特征值为 $\pm1$；并说明它在几何上是什么变换.
			\item\label{ex:dec:A13} 证明：$\vv A^\trans\vv A$ 与 $\vv{AA}^\trans$ 的非零特征值相同（重数也相同）.
			\item\label{ex:dec:A14} 证明：$\|\vv A\|_F^2=\tr(\vv A^\trans\vv A)=\sigma_1^2+\cdots+\sigma_r^2$.
			\item\label{ex:dec:A15} 设 $\vv A$ 是可逆实方阵. 证明 $|\det\vv A|=\sigma_1\sigma_2\cdots\sigma_n$，$\vv A^{-1}$ 的奇异值为 $1/\sigma_i$.
			\item\label{ex:dec:A16} 判断下列矩阵能否用正交矩阵对角化，说明理由：$\begin{pmatrix}1&2\\0&3\end{pmatrix}$，$\begin{pmatrix}1&2\\2&-2\end{pmatrix}$，$\begin{pmatrix}0&-1\\1&0\end{pmatrix}$（在实数范围内）.
			\item\label{ex:dec:A17} 四个已中心化的二维数据点为 $(2,1)$，$(-2,-1)$，$(1,2)$，$(-1,-2)$. 求样本协方差矩阵、第一主成分方向及其解释的方差比例.
			\item\label{ex:dec:A18} 计算 $\diag(10,1)$ 与 $\begin{pmatrix}1&1\\1&1.01\end{pmatrix}$ 的条件数（后者用近似值）. 后者的行列式与 $\diag(10,0.001)$ 相同，二者的条件数相同吗？
			\item\label{ex:dec:A19} 一张 $800\times600$ 像素的灰度图片，保留前 $k=40$ 个奇异值时需存储多少个数？压缩比是多少？$k$ 取多少时存储量反而超过原图？
			\item\label{ex:dec:A20} 证明：实对称矩阵的秩等于非零特征值的个数；奇异值的个数（非零）等于秩.
			\item\label{ex:dec:A21} 证明：实对称矩阵 $\vv A$ 满足``对一切 $\vv x\ne\zero$，$\vv x^\trans\vv A\vv x>0$''，当且仅当 $\vv A$ 的特征值全为正.（提示：用谱定理把 $\vv x^\trans\vv A\vv x$ 写成 $\sum\lambda_iy_i^2$.）
		\end{exA}
		
		\Bgroup
		\begin{exB}
			\item\label{ex:dec:B1} 设 $\vv A,\vv B$ 都是 $n$ 阶实对称矩阵且 $\vv{AB}=\vv{BA}$. 证明存在正交矩阵 $\vv Q$ 使 $\vv Q^\trans\vv{AQ}$ 与 $\vv Q^\trans\vv{BQ}$ 同时为对角矩阵.
			\item\label{ex:dec:B2} 证明：半正定的实对称矩阵 $\vv A$（特征值全 $\ge0$）有唯一的半正定平方根 $\vv B$（$\vv B^2=\vv A$，$\vv B$ 对称且特征值全 $\ge0$）.
			\item\label{ex:dec:B3} 设实对称矩阵 $\vv A$ 的特征值为 $\lambda_1\ge\lambda_2\ge\cdots\ge\lambda_n$，$\vv q_1$ 是属于 $\lambda_1$ 的单位特征向量. 证明 $\lambda_2=\max\{\vv x^\trans\vv{Ax}:\|\vv x\|=1,\ \vv x\perp\vv q_1\}$. 这说明第二主成分是``与第一主成分垂直的方向中方差最大的''.
			\item\label{ex:dec:B4} 证明：对任意 $\vv x$，$\|\vv{Ax}\|\le\sigma_1\|\vv x\|$，且等号可以取到. 因此 $\sigma_1=\max_{\vv x\ne\zero}\|\vv{Ax}\|/\|\vv x\|$，称为 $\vv A$ 的\textbf{谱范数}\index{pu3 fan4shu4@谱范数} $\|\vv A\|_2$.
			\item\label{ex:dec:B5} 证明：方阵 $\vv A$ 的任一特征值 $\lambda$ 满足 $\sigma_{\min}\le|\lambda|\le\sigma_{\max}$.
			\item\label{ex:dec:B6} 求 $\vv A=\begin{pmatrix}1&0&1\\0&1&1\end{pmatrix}$ 的奇异值分解.
			\item\label{ex:dec:B7} （伪逆）设 $\vv A=\vv U\vv\Sigma\vv V^\trans$，定义 $\vv A^+=\vv V\vv\Sigma^+\vv U^\trans$，其中 $\vv\Sigma^+$ 是把 $\vv\Sigma$ 转置并把非零元取倒数. 证明 $\vv A^+\vv b$ 是 $\vv{Ax}=\vv b$ 的范数最小的最小二乘解. 当 $\vv A$ 列满秩时，它与第九章的 $(\vv A^\trans\vv A)^{-1}\vv A^\trans$ 一致.
			\item\label{ex:dec:B8} 设 $p$ 个标准化指标两两相关系数都为 $\rho$（$0<\rho<1$），相关系数矩阵为 $(1-\rho)\vv I+\rho\vv J$（$\vv J$ 为全 $1$ 矩阵）. 求它的特征值与第一主成分，并求第一主成分解释的方差比例. 当 $\rho\to1$ 与 $\rho\to0$ 时比例分别趋于多少？
			\item\label{ex:dec:B9} 设两个指标的协方差矩阵为 $\begin{pmatrix}100&6\\6&1\end{pmatrix}$（第一个指标以``元''计，第二个以``万元''计）. 对协方差矩阵与相关系数矩阵分别作主成分分析，比较第一主成分. 由此说明：主成分分析对量纲敏感，为什么通常先标准化.
			\item\label{ex:dec:B10} 证明 $\|\vv A-\vv A_k\|_2=\sigma_{k+1}$.
			\item\label{ex:dec:B11} 证明：若 $\vv A$ 的奇异值互不相同且都不为零，则 $\vv A$ 的奇异值分解中的 $\vv u_i,\vv v_i$ 除了同时改变符号之外是唯一的.
			\item\label{ex:dec:B12} （正规矩阵）满足 $\vv A^\trans\vv A=\vv{AA}^\trans$ 的实方阵称为正规矩阵. 证明二阶实正规矩阵只有两类：对称矩阵，以及形如 $\begin{pmatrix}a&-b\\b&a\end{pmatrix}$ 的矩阵（旋转与伸缩的复合）.
			\item\label{ex:dec:B13} 设 $\vv A$ 是实对称矩阵. 证明 $\vv A$ 的特征值都落在 $[\min_i(a_{ii}-r_i),\ \max_i(a_{ii}+r_i)]$ 中，$r_i=\sum_{j\ne i}|a_{ij}|$. 并证明 $\lambda_{\max}\ge\max_ia_{ii}$.
			\item\label{ex:dec:B14} 对第十章``潜在语义分析''的例子，取 $k=1$，说明此时四份文书分别被表示成什么，查询``解约''与哪几份文书相似. 与 $k=2$ 比较，``压缩过头''的现象是否更严重？
			\item\label{ex:dec:B15} 证明：实对称矩阵 $\vv A$ 的特征值都在 $[a,b]$ 中，当且仅当 $\vv A-a\vv I$ 与 $b\vv I-\vv A$ 的二次型都非负.
		\end{exB}
		
		\Cgroup
		\begin{exC}
			\item\label{ex:dec:C1} （Eckart–Young 定理的证明）设 $\rk\vv B\le k$. 证明 $\ker\vv B$ 与 $\spn(\vv v_1,\ldots,\vv v_{k+1})$ 有非零公共向量 $\vv x$，由此证明 $\|\vv A-\vv B\|_2\ge\sigma_{k+1}$. 再借助 Weyl 不等式 $\sigma_{i+j-1}(\vv X+\vv Y)\le\sigma_i(\vv X)+\sigma_j(\vv Y)$（可承认）证明 Frobenius 范数的情形.
			\item\label{ex:dec:C2} （Cauchy 交错定理）设实对称矩阵 $\vv A$ 的特征值为 $\lambda_1\ge\cdots\ge\lambda_n$，删去最后一行一列所得矩阵的特征值为 $\mu_1\ge\cdots\ge\mu_{n-1}$. 证明 $\lambda_1\ge\mu_1\ge\lambda_2\ge\mu_2\ge\cdots\ge\mu_{n-1}\ge\lambda_n$.（提示：Courant–Fischer 极小极大刻画.）
			\item\label{ex:dec:C3} （复谱定理）证明：复方阵 $\vv A$ 可以被酉矩阵对角化（存在 $\vv U$ 使 $\bar{\vv U}^\trans\vv{AU}$ 为对角矩阵），当且仅当 $\vv A\bar{\vv A}^\trans=\bar{\vv A}^\trans\vv A$. 实对称矩阵、实正交矩阵、实反对称矩阵都满足这个条件.
			\item\label{ex:dec:C4} （正交 Procrustes 问题）给定 $n\times k$ 矩阵 $\vv A,\vv B$，求正交矩阵 $\vv Q$ 使 $\|\vv{QA}-\vv B\|_F$ 最小. 证明：若 $\vv{BA}^\trans=\vv U\vv\Sigma\vv V^\trans$，则最优解为 $\vv Q=\vv{UV}^\trans$. 这个问题出现在比较两组测量数据、对齐两份词向量等场合.
			\item\label{ex:dec:C5} （Ky Fan 定理）设实对称矩阵的特征值为 $\lambda_1\ge\cdots\ge\lambda_n$. 证明对任意列标准正交的 $n\times k$ 矩阵 $\vv X$，$\tr(\vv X^\trans\vv{AX})\le\lambda_1+\cdots+\lambda_k$. 由此说明前 $k$ 个主成分``一起''解释的方差也是最大的.
			\item\label{ex:dec:C6} （幂迭代）设 $\vv A$ 可对角化，特征值满足 $|\lambda_1|>|\lambda_2|\ge\cdots$. 证明：对``一般的''初始向量 $\vv x_0$，$\vv x_k=\vv A^k\vv x_0/\|\vv A^k\vv x_0\|$ 收敛到属于 $\lambda_1$ 的单位特征向量（可能差一个符号），收敛速度由 $|\lambda_2/\lambda_1|$ 决定. 对 $\begin{pmatrix}2&1\\1&2\end{pmatrix}$ 从 $(1,0)^\trans$ 出发迭代五步验证.
			\item\label{ex:dec:C7} （图的 Laplace 矩阵）设 $\vv L=\vv D-\vv A$ 是无向图的 Laplace 矩阵（见第五章章末习题~\ref{ex:det:C4}）. 证明 $\vv x^\trans\vv L\vv x=\sum_{\text{边 }ij}(x_i-x_j)^2$，因而 $\vv L$ 半正定；特征值 $0$ 的重数等于图的连通分支数. 第二小的特征值对应的特征向量（Fiedler 向量）的正负号常被用来把图``切''成两半——这就是谱聚类的思想.
			\item\label{ex:dec:C8} （正交回归）对中心化数据点 $\vv z_1,\ldots,\vv z_n\in\mathbb{R}^2$，求过原点的直线使各点到直线的\textbf{垂直}距离平方和最小. 证明最优方向就是第一主成分方向. 与普通最小二乘（只看竖直方向的残差）比较，两条直线何时相同？
			\item\label{ex:dec:C9} （扰动分析）设 $\vv A$ 可逆，$\vv{Ax}=\vv b$，$\vv A(\vv x+\delta\vv x)=\vv b+\delta\vv b$. 证明 $\dfrac{\|\delta\vv x\|}{\|\vv x\|}\le\kappa(\vv A)\dfrac{\|\delta\vv b\|}{\|\vv b\|}$，$\kappa(\vv A)=\sigma_1/\sigma_n$，并说明等号可以取到. 用它解释第六章病态方程组的例子.
		\end{exC}
	}
	
	% =====================================================================
	% 第十一章　二次型
	% =====================================================================
	\newcommand{\ExQuadratic}{%
		\exsetheading
		\Agroup
		\begin{exA}
			\item\label{ex:quad:A1} 写出下列二次型的矩阵与秩：
			(1) $x_1^2-3x_2^2+4x_1x_2-2x_2x_3$；(2) $x_1x_2+x_2x_3+x_3x_1$；(3) $(x_1+x_2+x_3)^2$；(4) $\sum_{i=1}^n x_i^2+\sum_{i<j}x_ix_j$.
			\item\label{ex:quad:A2} 写出以 $\begin{pmatrix}1&-1&0\\-1&2&3\\0&3&0\end{pmatrix}$ 为矩阵的二次型.
			\item\label{ex:quad:A3} 用配方法化为标准形，并写出所用的可逆变量替换：
			(1) $x_1^2+2x_2^2+5x_3^2+2x_1x_2+6x_2x_3$；(2) $x_1^2-4x_1x_2+4x_2^2-x_3^2$；(3) $x_1x_2-x_2x_3$.
			\item\label{ex:quad:A4} 用正交变换法把 $2x_1^2+3x_2^2+3x_3^2+4x_2x_3$ 化为标准形，写出正交矩阵. 曲面 $2x_1^2+3x_2^2+3x_3^2+4x_2x_3=1$ 是什么形状？
			\item\label{ex:quad:A5} 求 A3 中各二次型的正、负惯性指数与符号差，并写出规范形.
			\item\label{ex:quad:A6} 判断下列二次型是否正定：
			(1) $2x_1^2+x_2^2+3x_3^2+2x_1x_2$；(2) $x_1^2+x_2^2+x_3^2+4x_1x_2$；(3) $5x_1^2+x_2^2+5x_3^2+4x_1x_2-8x_1x_3-4x_2x_3$.
			\item\label{ex:quad:A7} $t$ 取何值时，$2x_1^2+x_2^2+x_3^2+2x_1x_2+tx_2x_3$ 正定？
			\item\label{ex:quad:A8} 判断曲线类型并求主轴方向：(1) $x^2-4xy+y^2=1$；(2) $3x^2+2xy+3y^2=8$.
			\item\label{ex:quad:A9} 判断下列矩阵是否合同：(1) $\diag(1,2,-3)$ 与 $\diag(-1,5,6)$；(2) $\begin{pmatrix}0&1\\1&0\end{pmatrix}$ 与 $\vv I_2$；(3) $\begin{pmatrix}0&1\\1&0\end{pmatrix}$ 与 $\diag(1,-1)$. 合同时求出 $\vv C$.
			\item\label{ex:quad:A10} 证明：若 $\vv A$ 正定，则 $\vv A^*$ 也正定.
			\item\label{ex:quad:A11} 求 $f(x,y)=x^3+y^3-3xy$ 的全部驻点，并用 Hessian 矩阵判断其类型.
			\item\label{ex:quad:A12} 判断 $\begin{pmatrix}2&-1&0\\-1&2&-1\\0&-1&2\end{pmatrix}$ 是否正定，并求它的 Cholesky 分解.
			\item\label{ex:quad:A13} 求 $\begin{pmatrix}4&2&-2\\2&10&2\\-2&2&6\end{pmatrix}$ 的 Cholesky 分解.
			\item\label{ex:quad:A14} 三种资产收益率的协方差矩阵为 $\vv\Sigma=\begin{pmatrix}0.04&0.01&0\\0.01&0.09&0\\0&0&0.01\end{pmatrix}$. 求组合 $\vv w=(0.5,0.3,0.2)^\trans$ 的方差与标准差.
			\item\label{ex:quad:A15} 两种资产收益率的标准差为 $30\%$ 与 $20\%$，互不相关. 求最小方差组合及其标准差.
			\item\label{ex:quad:A16} 证明：合同变换保持矩阵的对称性与秩. 合同变换保持行列式的符号吗？保持行列式的值吗？
			\item\label{ex:quad:A17} 分别在 $\mathbb{R}$ 上与 $\mathbb{C}$ 上写出二次型 $x_1^2-x_2^2+2x_3^2$ 的规范形.
			\item\label{ex:quad:A18} 判断下列方程表示的二次曲面类型：$x^2+y^2-z^2=1$，$x^2-y^2-z^2=1$，$x^2+y^2-z^2=0$，$x^2+y^2=1$.（提示：先看正、负惯性指数.）
			\item\label{ex:quad:A19} 说明 $\begin{pmatrix}1&2\\2&1\end{pmatrix}$ 不正定，并找出一个使二次型取负值的向量.
			\item\label{ex:quad:A20} 设 $\vv A$ 正定. 证明：$\vv A$ 的任一主子矩阵（取相同的若干行和列）都正定，从而所有主子式都为正.
		\end{exA}
		
		\Bgroup
		\begin{exB}
			\item\label{ex:quad:B1} 设 $\vv A,\vv B$ 正定且 $\vv{AB}=\vv{BA}$. 证明 $\vv{AB}$ 正定. 举例说明``可交换''不能去掉（此时 $\vv{AB}$ 甚至不对称）.
			\item\label{ex:quad:B2} 设 $\vv A$ 正定，$\vv B$ 半正定. 证明 $\det(\vv A+\vv B)\ge\det\vv A$，等号成立当且仅当 $\vv B=\vv O$.（提示：同时合同对角化.）
			\item\label{ex:quad:B3} 设 $\vv A$ 正定，$\vv B$ 对称. 证明：当 $|t|$ 充分小时 $\vv A+t\vv B$ 正定.
			\item\label{ex:quad:B4} 设 $\vv A$ 正定. 证明对 $i\ne j$ 有 $|a_{ij}|<\sqrt{a_{ii}a_{jj}}$，从而 $\vv A$ 的绝对值最大的元素在对角线上.
			\item\label{ex:quad:B5} 证明：实对称矩阵 $\vv A$ 满足 $\lambda_{\min}\|\vv x\|^2\le\vv x^\trans\vv{Ax}\le\lambda_{\max}\|\vv x\|^2$.
			\item\label{ex:quad:B6} 求 $\begin{pmatrix}0&1&1\\1&0&1\\1&1&0\end{pmatrix}$ 的正、负惯性指数（用两种方法：特征值；配方法）.
			\item\label{ex:quad:B7} $\lambda$ 取何值时，$\lambda(x_1^2+x_2^2+x_3^2)+2x_1x_2+2x_1x_3-2x_2x_3$ 正定？
			\item\label{ex:quad:B8} 求二次型 $\vv x^\trans\vv{Ax}$ 在约束 $\vv c^\trans\vv x=1$ 下的最小值（$\vv A$ 正定）：用拉格朗日乘数法证明最优解为 $\vv x^*=\dfrac{\vv A^{-1}\vv c}{\vv c^\trans\vv A^{-1}\vv c}$. 最小方差组合是它的特例.
			\item\label{ex:quad:B9} 设三种资产互不相关、方差都为 $\sigma^2$，期望收益率为 $\mu_1<\mu_2<\mu_3$. 求期望收益为 $m$ 的最小方差组合的权重，并说明有效前沿是一条抛物线（在``方差—期望收益''平面上）.
			\item\label{ex:quad:B10} 设 $\vv A=\begin{pmatrix}2&1\\1&2\end{pmatrix}$，$\vv B=\begin{pmatrix}1&0\\0&-1\end{pmatrix}$. 求可逆矩阵 $\vv C$ 使 $\vv C^\trans\vv{AC}=\vv I$ 且 $\vv C^\trans\vv{BC}$ 为对角矩阵.
			\item\label{ex:quad:B11} 设 $\vv A$ 正定. 求 $f(\vv x)=\vv x^\trans\vv{Ax}+\vv b^\trans\vv x+c$ 的最小值点与最小值. 若 $\vv A$ 只是半正定，最小值一定存在吗？
			\item\label{ex:quad:B12} 设 $\vv A$ 半正定，$\vv x^\trans\vv{Ax}=0$. 证明 $\vv{Ax}=\zero$.
			\item\label{ex:quad:B13} （Schur 补判据）证明：对称分块矩阵 $\begin{pmatrix}\vv A&\vv B\\ \vv B^\trans&\vv C\end{pmatrix}$ 正定，当且仅当 $\vv A$ 正定且 $\vv C-\vv B^\trans\vv A^{-1}\vv B$ 正定.
			\item\label{ex:quad:B14} 设 $\vv A$ 是二阶正定矩阵. 证明椭圆区域 $\{\vv x:\vv x^\trans\vv{Ax}\le1\}$ 的面积为 $\pi/\sqrt{\det\vv A}$.（提示：用主轴定理化为标准椭圆.）
			\item\label{ex:quad:B15} 设 $\vv A$ 为 $m\times n$ 实矩阵，$\mu>0$. 证明 $\vv A^\trans\vv A+\mu\vv I$ 正定，并说明岭回归 $\min\|\vv y-\vv{Ax}\|^2+\mu\|\vv x\|^2$ 的解为 $(\vv A^\trans\vv A+\mu\vv I)^{-1}\vv A^\trans\vv y$.
			\item\label{ex:quad:B16} 证明：实对称矩阵 $\vv A$ 正定，当且仅当存在正定矩阵 $\vv B$ 使 $\vv A=\vv B^2$；当且仅当 $\vv A^{-1}$ 存在且正定.
		\end{exB}
		
		\Cgroup
		\begin{exC}
			\item\label{ex:quad:C1} （Schur 乘积定理）对同阶矩阵定义逐元素乘积 $(\vv A\circ\vv B)_{ij}=a_{ij}b_{ij}$. 证明：若 $\vv A,\vv B$ 都半正定，则 $\vv A\circ\vv B$ 也半正定.（提示：把 $\vv B$ 写成 $\sum\lambda_k\vv q_k\vv q_k^\trans$.）
			\item\label{ex:quad:C2} （Minkowski 行列式不等式）设 $\vv A,\vv B$ 都是 $n$ 阶正定矩阵. 证明 $\det(\vv A+\vv B)^{1/n}\ge\det(\vv A)^{1/n}+\det(\vv B)^{1/n}$.（提示：同时合同对角化后用 $\prod(1+\mu_i)^{1/n}\ge1+\prod\mu_i^{1/n}$.）
			\item\label{ex:quad:C3} （有理数域上的二次型）证明 $x^2+y^2$ 与 $3x^2+3y^2$ 在 $\mathbb{R}$ 上合同，但在 $\mathbb{Q}$ 上不合同（即不存在有理系数的可逆替换把一个变为另一个）.（提示：若合同，则存在有理数 $a,b$ 使 $a^2+b^2=3$.）这说明惯性定理依赖于实数可以随意开平方.
			\item\label{ex:quad:C4} 函数 $f=x^2+y^4$ 与 $g=x^2-y^4$ 在原点处的 Hessian 矩阵相同且只是半正定，但前者在原点取极小值，后者不取极值. 由此说明 Hessian 半正定时二阶判别法失效.
			\item\label{ex:quad:C5} （广义特征值）设 $\vv A$ 对称，$\vv B$ 正定. 证明 $\max_{\vv x\ne\zero}\dfrac{\vv x^\trans\vv{Ax}}{\vv x^\trans\vv{Bx}}$ 等于方程 $\det(\vv A-\lambda\vv B)=0$ 的最大根. 统计学中的判别分析（Fisher LDA）正是这样求最佳投影方向的.
			\item\label{ex:quad:C6} （Mahalanobis 距离）设 $\vv\Sigma$ 正定，定义 $d(\vv x,\vv y)=\sqrt{(\vv x-\vv y)^\trans\vv\Sigma^{-1}(\vv x-\vv y)}$. 证明它满足三角不等式，且在仿射变换 $\vv x\mapsto\vv{Px}+\vv b$、$\vv\Sigma\mapsto\vv P\vv\Sigma\vv P^\trans$ 下不变. 它与欧氏距离有何不同？
			\item\label{ex:quad:C7} （Hilbert 矩阵）证明 $n$ 阶矩阵 $\vv H=\big(\frac1{i+j-1}\big)$ 正定.（提示：$\frac1{i+j-1}=\int_0^1x^{i-1}x^{j-1}\,\mathrm dx$，它是某个内积下的 Gram 矩阵.）
			\item\label{ex:quad:C8} （切点组合）设有无风险资产，收益率为 $r_f$. 证明使``夏普比率''$\dfrac{\vv\mu^\trans\vv w-r_f}{\sqrt{\vv w^\trans\vv\Sigma\vv w}}$（约束 $\vv 1^\trans\vv w=1$）最大的风险资产组合为 $\vv w\propto\vv\Sigma^{-1}(\vv\mu-r_f\vv 1)$.（提示：先去掉约束，比值关于 $\vv w$ 是零次齐次的；再用 Cauchy–Schwarz 不等式.）
			\item\label{ex:quad:C9} 证明：秩为 $r$、正惯性指数为 $p$ 的实二次型可以写成 $p$ 个线性型的平方和减去 $r-p$ 个线性型的平方和，并且不能写成少于 $r$ 个线性型的平方（带 $\pm$ 号）的组合.
		\end{exC}
	}
	
	% =====================================================================
	% 第十二章　*对偶空间与双线性型（选读）
	% =====================================================================
	\newcommand{\ExBilinear}{%
		\exsetheading
		\Agroup
		\begin{exA}
			\item\label{ex:bil:A1} 判断下列映射是否为相应空间上的线性函数：(1) $\mathbb{R}^3\to\mathbb{R}$，$\vv x\mapsto x_1-2x_3$；(2) $\vv x\mapsto x_1x_2$；(3) $\mathbb{R}[x]\to\mathbb{R}$，$f\mapsto f'(0)$；(4) $f\mapsto\int_0^1f(x)^2\,\mathrm dx$；(5) $M_n(\mathbb{R})\to\mathbb{R}$，$\vv X\mapsto\tr(\vv{AX})$（$\vv A$ 固定）.
			\item\label{ex:bil:A2} 求 $\mathbb{R}^3$ 的基 $(1,0,0)^\trans$，$(1,1,0)^\trans$，$(1,1,1)^\trans$ 的对偶基（把每个对偶基向量写成 $\vv x\mapsto a_1x_1+a_2x_2+a_3x_3$ 的形式）.
			\item\label{ex:bil:A3} 把线性函数 $f(\vv x)=x_1+2x_2+3x_3$ 用 A2 中的对偶基线性表出.
			\item\label{ex:bil:A4} 求 $W=\spn\big((1,1,0)^\trans,(0,1,1)^\trans\big)$ 的零化子 $W^0\subset(\mathbb{R}^3)^*$ 的一组基.
			\item\label{ex:bil:A5} 设线性映射 $\mathscr T:\mathbb{R}^2\to\mathbb{R}^3$ 在标准基下的矩阵为 $\begin{pmatrix}1&2\\0&1\\3&-1\end{pmatrix}$. 求对偶映射 $\mathscr T^*$ 在对偶基下的矩阵，并计算 $\mathscr T^*$ 作用在线性函数 $\vv y\mapsto y_1+y_2+y_3$ 上的结果.
			\item\label{ex:bil:A6} 求双线性型 $f(\vv x,\vv y)=x_1y_1+2x_1y_2-x_2y_1+3x_2y_2$ 在标准基下的度量矩阵，以及在基 $(1,1)^\trans$，$(1,-1)^\trans$ 下的度量矩阵，并验证 $\vv B=\vv P^\trans\vv{AP}$.
			\item\label{ex:bil:A7} 把 A6 中的 $f$ 分解为对称部分与交错部分之和.
			\item\label{ex:bil:A8} 判断下列双线性型是否非退化：(1) $x_1y_1-x_2y_2$；(2) $x_1y_2+x_2y_1$；(3) $(x_1+x_2)(y_1+y_2)$.
			\item\label{ex:bil:A9} 写出对称双线性型 $f(\vv x,\vv y)=x_1y_1+x_1y_2+x_2y_1+3x_2y_2$ 对应的二次型；反过来，由二次型 $x_1^2-4x_1x_2+2x_2^2$ 用极化恒等式求对应的对称双线性型.
			\item\label{ex:bil:A10} 验证 $\mathbb{R}^2$ 上的 $f(\vv x,\vv y)=\det(\vv x,\vv y)$ 是交错双线性型，写出它的度量矩阵. 它非退化吗？
			\item\label{ex:bil:A11} 设 $a_1,a_2,a_3$ 互异，$f_i(p)=p(a_i)$ 是 $\mathbb{R}[x]_3$ 上的线性函数. 对 $a_i=0,1,2$，求以 $f_1,f_2,f_3$ 为对偶基的 $\mathbb{R}[x]_3$ 的基.
			\item\label{ex:bil:A12} 在 $\mathbb{R}^4$ 中，设 $W$ 由 $x_1+x_2=0$，$x_3-x_4=0$ 定义. 求 $\dim W$ 与 $\dim W^0$，并写出 $W^0$ 的一组基.
			\item\label{ex:bil:A13} 证明 $\langle\vv A,\vv B\rangle=\tr(\vv{AB})$ 是 $M_n(\mathbb{R})$ 上非退化的对称双线性型. 它是内积吗？
			\item\label{ex:bil:A14} 对 $\mathbb{R}[x]_3$ 上的线性函数 $\varphi(p)=\int_0^1p(x)\,\mathrm dx$，求它在 $1,x,x^2$ 的对偶基下的坐标.
			\item\label{ex:bil:A15} （Riesz 表示定理的有限维情形）设 $V$ 是有限维实内积空间. 证明：对每个 $f\in V^*$，存在唯一的 $\vv\beta\in V$ 使 $f(\vv\alpha)=\langle\vv\alpha,\vv\beta\rangle$ 对一切 $\vv\alpha$ 成立.（提示：取标准正交基 $\vv e_1,\ldots,\vv e_n$，令 $\vv\beta=\sum f(\vv e_i)\vv e_i$.）
		\end{exA}
		
		\Bgroup
		\begin{exB}
			\item\label{ex:bil:B1} 证明：$f_1,\ldots,f_k\in V^*$ 线性无关，当且仅当映射 $V\to\mathbb{F}^k$，$\vv\alpha\mapsto(f_1(\vv\alpha),\ldots,f_k(\vv\alpha))^\trans$ 是满射.
			\item\label{ex:bil:B2} 设 $f,g\in V^*$ 且 $\ker f\subset\ker g$. 证明存在常数 $c$ 使 $g=cf$. 推广：若 $\ker f_1\cap\cdots\cap\ker f_k\subset\ker g$，则 $g$ 是 $f_1,\ldots,f_k$ 的线性组合.
			\item\label{ex:bil:B3} 证明：$\ker f_1\cap\cdots\cap\ker f_k=\{\zero\}$ 当且仅当 $f_1,\ldots,f_k$ 张成 $V^*$. 用方程组的语言解释这个结论.
			\item\label{ex:bil:B4} 证明：$\mathscr T$ 是单射当且仅当 $\mathscr T^*$ 是满射；$\mathscr T$ 是满射当且仅当 $\mathscr T^*$ 是单射.
			\item\label{ex:bil:B5} （Simpson 公式）求常数 $a,b,c$，使 $\int_0^2p(x)\,\mathrm dx=a\,p(0)+b\,p(1)+c\,p(2)$ 对一切 $p\in\mathbb{R}[x]_3$ 成立. 用对偶基的语言说明为什么这样的 $a,b,c$ 存在且唯一. 验证这个公式对三次多项式也成立.
			\item\label{ex:bil:B6} 证明：$V$ 上全体双线性型构成的空间与 $M_n(\mathbb{F})$ 同构，其中对称型、交错型分别构成维数为 $\frac{n(n+1)}2$、$\frac{n(n-1)}2$ 的子空间，且整个空间是二者的直和.
			\item\label{ex:bil:B7} 证明：若 $V$ 上存在非退化的交错双线性型，则 $\dim V$ 是偶数.
			\item\label{ex:bil:B8} 求双线性型 $f(\vv x,\vv y)=x_1y_1+x_1y_2+x_2y_1+x_2y_2+x_3y_3$ 的``根'' $\{\vv\alpha:f(\vv\alpha,\cdot)=0\}$.
			\item\label{ex:bil:B9} 在 $\mathbb{R}^2$ 上取 $f(\vv x,\vv y)=x_1y_1-x_2y_2$. 求全部``迷向向量''（$f(\vv\alpha,\vv\alpha)=0$），并说明它们构成两条直线. 再证明保持 $f$ 的线性变换 $\vv P$（$\vv P^\trans\diag(1,-1)\vv P=\diag(1,-1)$，$\det\vv P=1$，且保持 $x_1>|x_2|$ 的区域）都形如 $\begin{pmatrix}\cosh t&\sinh t\\ \sinh t&\cosh t\end{pmatrix}$.
			\item\label{ex:bil:B10} 设 $W$ 是 $V$ 的子空间，$\mathscr T:V\to V$ 线性且 $\mathscr T(W)\subset W$. 证明 $\mathscr T^*(W^0)\subset W^0$.
		\end{exB}
		
		\Cgroup
		\begin{exC}
			\item\label{ex:bil:C1} 证明：对无限维空间 $\mathbb{F}[x]$，自然映射 $\vv\alpha\mapsto\hat{\vv\alpha}$ 是单射但不是满射.（提示：$\mathbb{F}[x]^*$ 中的元素由它在 $1,x,x^2,\ldots$ 上的取值任意决定，可以``无限多个非零''.）
			\item\label{ex:bil:C2} 证明：$V$ 上的双线性型与线性映射 $V\to V^*$ 一一对应：$f\mapsto(\vv\alpha\mapsto f(\vv\alpha,\cdot))$. 用这个对应和对偶基的换基规则，重新推出度量矩阵的换基公式 $\vv B=\vv P^\trans\vv{AP}$.
			\item\label{ex:bil:C3} 证明：二阶矩阵 $\vv P$ 满足 $\vv P^\trans\vv J\vv P=\vv J$（$\vv J=\begin{pmatrix}0&1\\-1&0\end{pmatrix}$）当且仅当 $\det\vv P=1$. 四阶时，满足 $\vv P^\trans\diag(\vv J,\vv J)\vv P=\diag(\vv J,\vv J)$ 的矩阵（辛矩阵）的行列式是多少？
			\item\label{ex:bil:C4} 求 $M_n(\mathbb{R})$ 上对称双线性型 $\tr(\vv{AB})$ 的正、负惯性指数.（提示：对称矩阵与反对称矩阵在这个型下``互相正交''.）
			\item\label{ex:bil:C5} 设 $f$ 是 $V$ 上的对称双线性型，$W$ 是子空间，$W^\perp=\{\vv\alpha:f(\vv\alpha,\vv\beta)=0\ \forall\vv\beta\in W\}$. 证明 $\dim W^\perp\ge\dim V-\dim W$，$f$ 非退化时等号成立；举例说明即使 $f$ 非退化，$W\cap W^\perp$ 也可能不为零.
			\item\label{ex:bil:C6} 证明：$\mathbb{F}^n$ 上全体交错 $n$ 重线性函数构成的空间是一维的，由行列式张成. 这给出第五章章末习题~\ref{ex:det:C1} 的``对偶空间版本''.
			\item\label{ex:bil:C7} （线性规划对偶）考虑问题：在 $\vv{Ax}=\vv b$，$\vv x\ge\zero$ 下最小化 $\vv c^\trans\vv x$. 证明对任意满足 $\vv A^\trans\vv y\le\vv c$ 的 $\vv y$ 与任意可行的 $\vv x$，都有 $\vv b^\trans\vv y\le\vv c^\trans\vv x$（弱对偶）. 用``价格是数量空间上的线性函数''解释 $\vv y$ 的经济含义（影子价格）.
		\end{exC}
	}.
	%%   2. 在每章末尾（原 \begin{problemset}...\end{problemset} 之后，或替换它）
	%%      写对应的命令：
	%%        第一章 \ExElimination      第二章 \ExVectorSpace   第三章 \ExMatrix
	%%        第四章 \ExLinearMap        第五章 \ExDeterminant   第六章 \ExSolutions
	%%        第七章 \ExPolynomial       第八章 \ExEigen         第九章 \ExInnerProduct
	%%        第十章 \ExDecomposition    第十一章 \ExQuadratic   第十二章 \ExBilinear
	%%   3. 分组标准
	%%        A 组：直接运用本章定义与定理的计算和证明；
	%%        B 组：需要综合多个结论、或有一定技巧的计算和证明；
	%%        C 组：有挑战性的题目，或超出本课程大纲、为继续深造准备的内容.
	%%   4. 题目与正文中已有的节末练习、原 problemset 不重复；个别题目注明``接正文某题''.
	%%      第六章按扩充版（ch06-solutions.tex）的内容编写.
	%%   5. 每道题都带标签 ex:<章键>:<组号>，正文、习题与附录答案一律用 \ref 引用，
	%%      渲染为``A4''``B12''这样的题号. 章键依次为
	%%        elim vs mat lmap det sol poly eig ip dec quad bil（第一至十二章）.
	%%      例：第一章章末习题~\ref{ex:elim:B12}. 新增题目请加在组末，并补上标签.
	%%      A 组答案、B 组提示在 answers.tex 中，按同一标签对应.
	%% 依赖：导言区的 \vv \zero \trans \spn \rk \imm \diag，以及 enumitem（elegantbook 已加载）.
	%% 引用的前文标签：thm:sameSolution thm:solvability prop:linspc eg:pricedecomp cor:homogeneous thm:rouche thm:sol-determinable
	%%   thm:spectral thm:posdef thm:svd thm:orthcanon
	%% =====================================================================
	
	\providecommand{\tr}{\operatorname{tr}}
	\newcommand{\exsetheading}{%
		\section*{章末习题}%
		\markright{章末习题}%
	}
\newcommand{\exgroup}[2]{%
	\par\medskip{\centering\sffamily\bfseries #1\quad #2\par}\smallskip}
	\newenvironment{exA}{\begin{enumerate}[label=\textbf{A\arabic*.},ref=A\arabic*,leftmargin=2.8em,itemsep=4pt]}{\end{enumerate}}
	\newenvironment{exB}{\begin{enumerate}[label=\textbf{B\arabic*.},ref=B\arabic*,leftmargin=2.8em,itemsep=4pt]}{\end{enumerate}}
	\newenvironment{exC}{\begin{enumerate}[label=\textbf{C\arabic*.},ref=C\arabic*,leftmargin=2.8em,itemsep=4pt]}{\end{enumerate}}
	\newcommand{\Agroup}{\exgroup{A 组}{基础巩固}}
	\newcommand{\Bgroup}{\exgroup{B 组}{综合提升}}
	\newcommand{\Cgroup}{\exgroup{C 组}{拓展探索}}
	
	% =====================================================================
	% 第一章　线性方程组与消元法
	% =====================================================================
	\newcommand{\ExElimination}{%
		\exsetheading
		\Agroup
		\begin{exA}
			\item\label{ex:elim:A1} 判断下列矩阵哪些是阶梯形，哪些是简化阶梯形：
			\[
			\begin{pmatrix}1&2&0\\0&0&1\end{pmatrix},\quad
			\begin{pmatrix}1&0&3\\0&2&1\\0&0&0\end{pmatrix},\quad
			\begin{pmatrix}0&1&0\\1&0&0\end{pmatrix},\quad
			\begin{pmatrix}1&0&2&0\\0&1&-1&0\\0&0&0&1\end{pmatrix},\quad
			\begin{pmatrix}0&0\\0&1\end{pmatrix}.
			\]
			\item\label{ex:elim:A2} 写出方程组 $3x_1-x_3=2$，$x_2+4x_3=0$，$x_1-x_2=5$ 的增广矩阵；反过来，写出以 $\begin{pmatrix}1&-2&0&3\\0&1&5&-1\end{pmatrix}$ 为增广矩阵的方程组.
			\item\label{ex:elim:A3} 用消元法解方程组 $\left\{\begin{array}{l}x_1+2x_2+3x_3=14,\\2x_1-x_2+x_3=3,\\3x_1+x_2-x_3=2.\end{array}\right.$
			\item\label{ex:elim:A4} 用消元法解方程组 $\left\{\begin{array}{l}x_1+x_2+x_3+x_4=2,\\x_1-x_2+x_3-x_4=4,\\2x_1+x_2-x_3+3x_4=-1,\\x_1+2x_2+3x_3-2x_4=5.\end{array}\right.$
			\item\label{ex:elim:A5} 判断方程组 $\left\{\begin{array}{l}x_1-x_2+x_3=1,\\2x_1+x_2-x_3=2,\\x_1+2x_2-2x_3=3\end{array}\right.$ 是否有解，并说明消元过程中哪一行``暴露''了结论.
			\item\label{ex:elim:A6} 求方程组 $\left\{\begin{array}{l}x_1+2x_2-x_3+x_4=3,\\2x_1+4x_2-x_3+3x_4=7,\\-x_1-2x_2+2x_3=-2\end{array}\right.$ 的全部解，并指出主变量与自由变量.
			\item\label{ex:elim:A7} 求齐次方程组 $\left\{\begin{array}{l}x_1+x_2-x_3-x_4=0,\\2x_1-5x_2+3x_3+2x_4=0,\\7x_1-7x_2+3x_3+x_4=0\end{array}\right.$ 的全部解.
			\item\label{ex:elim:A8} 把矩阵 $\begin{pmatrix}2&4&-2&6\\1&2&1&1\\3&6&1&5\end{pmatrix}$ 化为简化阶梯形，并写出每一步所用的初等行变换.
			\item\label{ex:elim:A9} 写出全部 $2\times2$ 简化阶梯形矩阵的``形状''（可以含参数），并说明共有几类.
			\item\label{ex:elim:A10} $k$ 取何值时，方程组 $x+y=1$，$x-y=3$，$2x+ky=5$ 有解？
			\item\label{ex:elim:A11} 设 $a,b$ 为常数. 方程组 $x+2y=a$，$3x+6y=b$ 何时有解？有解时写出全部解.
			\item\label{ex:elim:A12} 求过三点 $(1,2)$，$(2,3)$，$(3,6)$ 的抛物线 $y=ax^2+bx+c$.
			\item\label{ex:elim:A13} 求过四点 $(-1,2)$，$(0,1)$，$(1,0)$，$(2,5)$ 的三次多项式函数 $y=a+bx+cx^2+dx^3$.
			\item\label{ex:elim:A14} 某剧场三种票价分别为 $30$，$50$，$80$ 元. 某场演出共售出 $100$ 张票，票房 $4800$ 元，且 $80$ 元票的张数恰为 $30$ 元票的一半. 三种票各售出多少张？
			\item\label{ex:elim:A15} 配平化学方程式 $x_1\,\mathrm{C_3H_8}+x_2\,\mathrm{O_2}\to x_3\,\mathrm{CO_2}+x_4\,\mathrm{H_2O}$：对 C，H，O 三种原子分别写出守恒方程，解这个齐次方程组，并取最小的正整数解.
			\item\label{ex:elim:A16} 某城区四个路口 $A,B,C,D$ 由单行道依次连成环路 $A\to B\to C\to D\to A$，环路上的车流量依次记为 $x_1$（$A\to B$），$x_2$（$B\to C$），$x_3$（$C\to D$），$x_4$（$D\to A$）（单位：辆/小时）. 各路口与外部道路的进出量为：$A$ 驶入 $400$、驶出 $300$；$B$ 驶入 $200$、驶出 $350$；$C$ 驶入 $250$、驶出 $100$；$D$ 驶入 $150$、驶出 $250$. 设每个路口驶入总量等于驶出总量.
			\begin{enumerate}
				\item 列出方程组并求全部解；
				\item 若各路段车流量都不能为负，$x_1$ 最少是多少？
				\item 若把 $D$ 的驶出量改为 $350$，方程组会怎样？这反映了什么守恒律？
			\end{enumerate}
			\item\label{ex:elim:A17} 证明：二元一次方程组 $ax+by=e$，$cx+dy=f$ 对任意的 $e,f$ 都有唯一解，当且仅当 $ad-bc\ne0$.（只用消元法，不用行列式.）
			\item\label{ex:elim:A18} 证明：若线性方程组有两个不同的解 $\vv x_1,\vv x_2$，则对任意数 $t$，$\vv x_1+t(\vv x_2-\vv x_1)$ 都是解. 由此说明：线性方程组的解要么不存在，要么唯一，要么有无穷多个.
			\item\label{ex:elim:A19} 证明方程组 $\{x+y+z=3,\ x-y=0\}$ 与 $\{2x+z=3,\ x+y+z=3\}$ 同解.
			\item\label{ex:elim:A20} 举例说明：若\textbf{同时}做``第一个方程加到第二个方程上''与``第二个方程加到第一个方程上''（两步都用原来的方程），所得方程组可能与原方程组不同解. 这与定理~\ref{thm:sameSolution}（初等行变换不改变解集）矛盾吗？
			\item\label{ex:elim:A21} 已知一个方程组的增广矩阵化为阶梯形后为 $\begin{pmatrix}1&3&-1&2&0\\0&0&2&1&4\\0&0&0&0&a\end{pmatrix}$. 讨论 $a$ 取何值时方程组有解；有解时它有几个自由变量？
		\end{exA}
		
		\Bgroup
		\begin{exB}
			\item\label{ex:elim:B1} 讨论 $a,b$ 取何值时，方程组
			\[\left\{\begin{array}{l}x_1+x_2+x_3+x_4=0,\\x_2+2x_3+2x_4=1,\\-x_2+(a-3)x_3-2x_4=b,\\3x_1+2x_2+x_3+ax_4=-1\end{array}\right.\]
			有唯一解、无解、有无穷多解；有无穷多解时求出全部解.
			\item\label{ex:elim:B2} 讨论 $\lambda$ 取何值时，方程组
			$(1+\lambda)x_1+x_2+x_3=0$，$x_1+(1+\lambda)x_2+x_3=3$，$x_1+x_2+(1+\lambda)x_3=\lambda$
			有唯一解、无解、有无穷多解.
			\item\label{ex:elim:B3} 设 $n$ 个方程、$n$ 个未知量的方程组对应的齐次方程组只有零解. 证明：无论常数项取什么值，原方程组都有唯一解.（提示：考察简化阶梯形中主元的个数.）
			\item\label{ex:elim:B4} 证明：线性方程组无解，当且仅当可以把若干个方程分别乘以适当的数再相加，得到矛盾方程 $0=1$.
			\item\label{ex:elim:B5} 证明：交换两行这一初等变换，可以由``倍加''与``某行乘以 $-1$''这两类变换的组合实现. 由此说明三类初等行变换并不是``独立''的.
			\item\label{ex:elim:B6} 证明：若线性方程组的系数和常数项都是有理数，且它有实数解，则它一定有有理数解；若它的实数解唯一，则这个解就是有理数解.
			\item\label{ex:elim:B7} 讨论方程组
			\[x_1-x_2=a_1,\quad x_2-x_3=a_2,\quad\ldots,\quad x_{n-1}-x_n=a_{n-1},\quad x_n-x_1=a_n\]
			有解的充要条件，并在有解时求全部解.
			\item\label{ex:elim:B8} 讨论方程组
			\[x_1+x_2=a_1,\quad x_2+x_3=a_2,\quad\ldots,\quad x_{n-1}+x_n=a_{n-1},\quad x_n+x_1=a_n\]
			的解. 证明：$n$ 为奇数时解唯一；$n$ 为偶数时有解当且仅当 $a_1-a_2+a_3-\cdots-a_n=0$.
			\item\label{ex:elim:B9} 某经济体有三个部门，直接消耗系数矩阵与最终需求分别为
			\[\vv A=\begin{pmatrix}0.1&0.2&0.2\\0.2&0.1&0.2\\0.1&0.3&0.1\end{pmatrix},\qquad \vv d=\begin{pmatrix}20\\130\\65\end{pmatrix}.\]
			按第一章``投入产出模型''的方式列出方程组，求各部门的总产出.
			\item\label{ex:elim:B10} 估计用消元法（化为阶梯形后回代）解 $n$ 个方程、$n$ 个未知量的方程组所需的乘除法次数，证明它约为 $\frac13n^3$（即与 $\frac13n^3$ 之比在 $n\to\infty$ 时趋于 $1$）.
			\item\label{ex:elim:B11} 设三个平面 $a_ix+b_iy+c_iz=d_i$（$i=1,2,3$）两两相交. 用消元法的语言说明：它们恰好交于一点、交于一条直线、或没有公共点（三条交线互相平行），并各举一例.
			\item\label{ex:elim:B12} 证明：$n$ 元齐次方程组的解集对加法和数乘封闭；而非齐次方程组的两个解之差是对应齐次方程组的解. 由此说明：非齐次方程组的全部解，等于它的任一个解加上齐次方程组的全部解.
			\item\label{ex:elim:B13} 一家律所有甲、乙、丙三位律师，本月分别承办了若干件合同、劳动、知识产权纠纷. 已知：三类案件合计 $30$ 件；甲承办的数量 $=$ 乙 $+$ 丙；每件合同、劳动、知识产权案件分别需要 $10$，$6$，$15$ 个工时，三人本月共用 $300$ 个工时……
			这些条件够不够确定每位律师承办了多少案件？请自己补全一个合理的设定（明确未知量是什么），列出方程组，用自由变量的个数说明还缺多少条独立的信息.
			\item\label{ex:elim:B14} 设方程组的系数矩阵每一行的元素之和都等于 $0$. 证明 $x_1=x_2=\cdots=x_n=1$ 是对应齐次方程组的解；若非齐次方程组有解 $\vv x_0$，则对任意数 $t$，$\vv x_0+t(1,\ldots,1)^\trans$ 也是解.
			\item\label{ex:elim:B15} 证明：如果方程组 (I) 的每个方程都是方程组 (II) 各方程的``线性组合''（各乘一个数再相加），反之亦然，那么 (I) 与 (II) 同解. 逆命题对于有解的方程组成立吗？先在两个未知量的情形下思考.
			\item\label{ex:elim:B16} 讨论 $\lambda$ 取何值时，方程组 $\lambda x+y+z=1$，$x+\lambda y+z=\lambda$，$x+y+\lambda z=\lambda^2$ 有唯一解、无解、有无穷多解；有唯一解时求出解.（第六章例~\ref{eg:lambda} 将用秩与行列式重新组织这个讨论.）
		\end{exB}
		
		\Cgroup
		\begin{exC}
			\item\label{ex:elim:C1} （百鸡问题）《张丘建算经》（约公元 5 世纪）：``今有鸡翁一，直钱五；鸡母一，直钱三；鸡雏三，直钱一. 凡百钱买鸡百只，问鸡翁、母、雏各几何？''列出线性方程组，求它的全部解，再从中找出所有\textbf{非负整数}解. 线性代数在哪一步之后就无能为力了？
			\item\label{ex:elim:C2} （幻方）若一个 $3\times3$ 数表的每行、每列和两条对角线上的三个数之和都等于同一个数 $s$，称它为幻方（不要求数互不相同）. 把 $9$ 个元素作为未知量、$s$ 也作为未知量，写出方程组，证明：
			\begin{enumerate}
				\item 中心元素一定等于 $s/3$；
				\item 全部幻方可以用 $3$ 个自由参数表示，写出通式；
				\item 用 $1,2,\ldots,9$ 各一次填成的幻方本质上只有一个（其余由旋转、翻转得到）.
			\end{enumerate}
			\item\label{ex:elim:C3} （点灯游戏）$3\times3$ 方格中每格有一盏灯，按下某格会使该格及其上下左右相邻的格子开关状态翻转. 在``$1+1=0$''的算术（即模 $2$ 运算）中，把``按哪些格''作为未知量，把``从某一状态变为全灭''写成线性方程组. 证明：从任意初始状态出发，总能把灯全部熄灭.（模 $2$ 运算的集合 $\{0,1\}$ 不是本书意义下的数域，但消元法照样适用——想一想为什么.）
			\item\label{ex:elim:C4} （舍入误差与主元）在``只保留三位有效数字''的计算机上解方程组 $0.0001x+y=1$，$x+y=2$：
			\begin{enumerate}
				\item 直接用第一个方程消去第二个方程中的 $x$；
				\item 先交换两个方程再消元.
			\end{enumerate}
			把两种结果与精确解比较，说明数值软件为什么要``选绝对值最大的元素作主元''.
			\item\label{ex:elim:C5} 设两个都有解的方程组 (I)(II) 都有 $m$ 个方程、$n$ 个未知量，且同解. 证明：它们的增广矩阵行等价.（提示：先证明二者的简化阶梯形的非零行完全相同.）
			\item\label{ex:elim:C6} （Jacobi 迭代）对方程组 $4x+y=9$，$x+3y=5$，把它改写为 $x=(9-y)/4$，$y=(5-x)/3$，从 $(x_0,y_0)=(0,0)$ 出发反复代入. 计算前五步，观察是否收敛到真解. 再对 $x+4y=9$，$3x+y=5$ 做同样的事，会发生什么？试找出一个保证收敛的充分条件.
			\item\label{ex:elim:C7} （电路）一个电路中有三条支路电流 $I_1,I_2,I_3$. 节点 $A$ 处 $I_1=I_2+I_3$；两个回路的电压方程为 $2I_1+I_2=8$，$I_2-2I_3=0$（单位从略）. 求各电流. 再说明：若把第三个节点 $B$ 的电流方程也写进来，方程组为什么不会因此``多一个条件''？
			\item\label{ex:elim:C8} 证明：一个 $m\times n$ 矩阵（$m,n\ge1$）通过初等行变换所能得到的不同的简化阶梯形``形状''（即主元所在列的位置组合）共有 $\sum_{r=0}^{\min(m,n)}\binom nr$ 种.
			\item\label{ex:elim:C9} 证明：一个矩阵的简化阶梯形是唯一的.（提示：设 $\vv R_1,\vv R_2$ 都是 $\vv M$ 的简化阶梯形，则以它们为系数矩阵的齐次方程组同解. 对列数归纳：删去最后一列后由归纳假设前面各列相同，再比较最后一列，考察它在两者中是否为主元列.）
		\end{exC}
	}
	
	% =====================================================================
	% 第二章　向量空间
	% =====================================================================
	\newcommand{\ExVectorSpace}{%
		\exsetheading
		\Agroup
		\begin{exA}
			\item\label{ex:vs:A1} 下列集合关于通常的四则运算是否构成数域？
			(1) $\{a+b\sqrt3\mid a,b\in\mathbb{Q}\}$；(2) $\{a+b\sqrt2\mid a\in\mathbb{Q},\ b\in\mathbb{Z}\}$；(3) 全体偶数；(4) $\{a+b\sqrt[3]{2}\mid a,b\in\mathbb{Q}\}$；(5) 全体纯虚数与 $0$.
			\item\label{ex:vs:A2} 下列集合关于通常的加法与数乘是否构成实向量空间？
			(1) 全体 $n$ 阶实对称矩阵；(2) $\{(x,y)^\trans\in\mathbb{R}^2\mid x\ge0\}$；(3) 次数恰为 $n$ 的实系数多项式全体；(4) $\{f\in\mathbb{R}[x]\mid f(1)=0\}$；(5) 满足 $a_{n+2}=a_{n+1}+a_n$ 的全体实数列；(6) $\{(x,y,z)^\trans\mid x+y+z=1\}$.
			\item\label{ex:vs:A3} 判断下列向量组的线性相关性：
			\begin{enumerate}
				\item $(1,2,-1)^\trans$，$(2,-3,1)^\trans$，$(4,1,-1)^\trans$；
				\item 在 $\mathbb{R}[x]_3$ 中：$1+x$，$1-x$，$x^2$；
				\item 在 $\mathbb{R}[x]_3$ 中：$1+x+x^2$，$1-x$，$2+x^2$，$x+x^2$；
				\item 在 $M_2(\mathbb{R})$ 中：$\begin{pmatrix}1&0\\0&1\end{pmatrix}$，$\begin{pmatrix}0&1\\1&0\end{pmatrix}$，$\begin{pmatrix}1&1\\1&1\end{pmatrix}$，$\begin{pmatrix}1&0\\0&0\end{pmatrix}$.
			\end{enumerate}
			\item\label{ex:vs:A4} 把 $\vv\beta=(1,2,1,1)^\trans$ 表示为 $\vv\alpha_1=(1,1,1,1)^\trans$，$\vv\alpha_2=(1,1,-1,-1)^\trans$，$\vv\alpha_3=(1,-1,1,-1)^\trans$，$\vv\alpha_4=(1,-1,-1,1)^\trans$ 的线性组合.
			\item\label{ex:vs:A5} $t$ 取何值时，$(1,1,0)^\trans$，$(1,3,-1)^\trans$，$(5,3,t)^\trans$ 线性相关？
			\item\label{ex:vs:A6} 求向量组 $\vv\alpha_1=(1,-1,2,4)^\trans$，$\vv\alpha_2=(0,3,1,2)^\trans$，$\vv\alpha_3=(3,0,7,14)^\trans$，$\vv\alpha_4=(1,-1,2,0)^\trans$，$\vv\alpha_5=(2,1,5,6)^\trans$ 的秩和一个极大无关组，并把其余向量用这个极大无关组线性表出.
			\item\label{ex:vs:A7} 设 $\vv\alpha_1,\vv\alpha_2,\vv\alpha_3$ 线性无关. 判断下列向量组的线性相关性：
			(1) $\vv\alpha_1-\vv\alpha_2$，$\vv\alpha_2-\vv\alpha_3$，$\vv\alpha_3-\vv\alpha_1$；
			(2) $\vv\alpha_1$，$\vv\alpha_1+\vv\alpha_2$，$\vv\alpha_1+\vv\alpha_2+\vv\alpha_3$；
			(3) $\vv\alpha_1+2\vv\alpha_2$，$\vv\alpha_2+2\vv\alpha_3$，$\vv\alpha_3+2\vv\alpha_1$.
			\item\label{ex:vs:A8} 证明：含有零向量的向量组线性相关；含有两个相同向量的向量组线性相关；单个向量 $\vv\alpha$ 线性无关当且仅当 $\vv\alpha\ne\zero$.
			\item\label{ex:vs:A9} 证明：两个向量线性相关，当且仅当其中一个是另一个的倍数. 三个向量的情形能这样说吗？
			\item\label{ex:vs:A10} 只用向量空间的八条公理与命题~\ref{prop:linspc}，证明：(1) $\vv\alpha+\vv\alpha=2\vv\alpha$；(2) 若 $k\ne0$ 且 $k\vv\alpha=k\vv\beta$，则 $\vv\alpha=\vv\beta$；(3) $-(-\vv\alpha)=\vv\alpha$.
			\item\label{ex:vs:A11} 证明下列集合是相应空间的子空间，并求一组基和维数：
			\begin{enumerate}
				\item $\{\vv x\in\mathbb{R}^4\mid x_1+x_2=x_3+x_4\}$；
				\item $n$ 阶上三角矩阵全体；
				\item 迹为零的 $n$ 阶方阵全体 $\{\vv X\mid x_{11}+\cdots+x_{nn}=0\}$；
				\item $\{f\in\mathbb{R}[x]_4\mid f(1)=0\}$.
			\end{enumerate}
			\item\label{ex:vs:A12} 判断下列子集是否为子空间：(1) $\{(a,b,c)^\trans\mid a=2b\}$；(2) $\{(a,b,c)^\trans\mid ab=0\}$；(3) $\{(a,b,c)^\trans\mid a,b,c\in\mathbb{Q}\}\subset\mathbb{R}^3$；(4) $\{f\in C[0,1]\mid f(0)=2f(1)\}$.
			\item\label{ex:vs:A13} 求 $\vv\alpha_1=(1,2,1,0)^\trans$，$\vv\alpha_2=(-1,1,1,1)^\trans$，$\vv\alpha_3=(0,3,2,1)^\trans$，$\vv\alpha_4=(1,0,1,2)^\trans$ 张成的子空间的维数和一组基.
			\item\label{ex:vs:A14} 求多项式 $2x^2-x+3$ 在基 $1,\ x-1,\ (x-1)^2$ 下的坐标.
			\item\label{ex:vs:A15} 求矩阵 $\begin{pmatrix}2&3\\-1&4\end{pmatrix}$ 在 $M_2(\mathbb{R})$ 的基 $\begin{pmatrix}1&0\\0&0\end{pmatrix},\begin{pmatrix}1&1\\0&0\end{pmatrix},\begin{pmatrix}1&1\\1&0\end{pmatrix},\begin{pmatrix}1&1\\1&1\end{pmatrix}$ 下的坐标.
			\item\label{ex:vs:A16} (1) $\mathbb{R}$ 是否是 $\mathbb{Q}$ 上的向量空间？(2) $\mathbb{Q}$ 是否是 $\mathbb{R}$ 上的向量空间？(3) $\mathbb{C}$ 作为 $\mathbb{R}$ 上的向量空间和作为 $\mathbb{C}$ 上的向量空间，维数各是多少？
			\item\label{ex:vs:A17} 写出 $\mathbb{F}[x]_n$ 与 $\mathbb{F}^n$ 之间、$M_2(\mathbb{F})$ 与 $\mathbb{F}^4$ 之间的一个同构映射，并验证它保持线性运算.
			\item\label{ex:vs:A18} 设 $W_1=\spn\big((1,0,0)^\trans,(0,1,0)^\trans\big)$，$W_2=\spn\big((1,1,1)^\trans\big)$，$W_3=\spn\big((1,1,0)^\trans\big)$. 分别判断 $W_1+W_2$，$W_1+W_3$ 是否为直和，并求和空间的维数.
			\item\label{ex:vs:A19} 设 $W_1=\spn\big((1,1,0,0)^\trans,(1,0,1,1)^\trans\big)$，$W_2=\spn\big((0,-1,1,1)^\trans,(1,2,0,1)^\trans\big)$. 求 $W_1+W_2$ 与 $W_1\cap W_2$ 的维数，并验证维数公式.
			\item\label{ex:vs:A20} 设 $\vv\alpha_1,\ldots,\vv\alpha_s$ 的秩为 $r$. 证明：其中任意 $r+1$ 个向量线性相关；任意 $r$ 个线性无关的向量都构成极大无关组.
		\end{exA}
		
		\Bgroup
		\begin{exB}
			\item\label{ex:vs:B1} 证明：两个子空间的并 $W_1\cup W_2$ 是子空间，当且仅当 $W_1\subset W_2$ 或 $W_2\subset W_1$.
			\item\label{ex:vs:B2} 证明：若向量组 $\vv\beta_1,\ldots,\vv\beta_t$ 可由 $\vv\alpha_1,\ldots,\vv\alpha_s$ 线性表出，且两组的秩相等，则两组等价.
			\item\label{ex:vs:B3} 设向量组 (I) 与 (II) 的秩分别为 $r_1,r_2$，把它们合并为 (III). 证明 $\max\{r_1,r_2\}\le\rk(\mathrm{III})\le r_1+r_2$. 再证明：$\rk(\vv\alpha_1+\vv\beta_1,\ldots,\vv\alpha_s+\vv\beta_s)\le\rk(\vv\alpha_1,\ldots,\vv\alpha_s)+\rk(\vv\beta_1,\ldots,\vv\beta_s)$.
			\item\label{ex:vs:B4} 设 $\vv\alpha_1,\ldots,\vv\alpha_n$ 是 $V$ 的基，$\vv\beta=a_1\vv\alpha_1+\cdots+a_n\vv\alpha_n$. 证明：把 $\vv\alpha_j$ 换成 $\vv\beta$ 后所得的向量组仍是基，当且仅当 $a_j\ne0$.
			\item\label{ex:vs:B5} 设 $p_0,p_1,\ldots,p_{n-1}\in\mathbb{F}[x]$ 满足 $\deg p_k=k$. 证明它们构成 $\mathbb{F}[x]_n$ 的一组基.
			\item\label{ex:vs:B6} 记 $V$ 为满足 $a_{n+2}=a_{n+1}+2a_n$（$n\ge1$）的全体实数列. 证明 $V$ 是向量空间，$\dim V=2$，并说明 $\{2^n\}$ 与 $\{(-1)^n\}$ 构成 $V$ 的一组基. 由此写出 $a_1=1$，$a_2=1$ 时的通项.
			\item\label{ex:vs:B7} 在 $C(\mathbb{R})$ 中：(1) 证明 $\mathrm e^x,\mathrm e^{2x},\mathrm e^{3x}$ 线性无关；(2) 证明 $\sin x,\cos x,\sin(x+1)$ 线性相关；(3) $1,\cos x,\cos2x,\cos^2x$ 的秩是多少？
			\item\label{ex:vs:B8} 证明 $\{a+b\sqrt2+c\sqrt3+d\sqrt6\mid a,b,c,d\in\mathbb{Q}\}$ 是数域，并证明 $1,\sqrt2,\sqrt3,\sqrt6$ 在 $\mathbb{Q}$ 上线性无关. 这个数域作为 $\mathbb{Q}$ 上的向量空间是几维的？
			\item\label{ex:vs:B9} 写出 $\mathbb{R}^2$ 与 $\mathbb{R}^3$ 的全部子空间（按维数分类，并说明几何意义）.
			\item\label{ex:vs:B10} 证明：$\mathbb{R}^5$ 中任意两个三维子空间的交至少是一维的. 一般地，$n$ 维空间中维数分别为 $k_1,k_2$ 的两个子空间，若 $k_1+k_2>n$，则交不为 $\{\zero\}$.
			\item\label{ex:vs:B11} 设 $a\in\mathbb{F}$. 证明 $\mathbb{F}[x]_n=\mathbb{F}\cdot1\oplus\{f\in\mathbb{F}[x]_n\mid f(a)=0\}$，并写出 $f$ 按这个直和的分解.
			\item\label{ex:vs:B12} 设 $V=C[-1,1]$，$E$，$O$ 分别为其中的偶函数、奇函数全体. 证明 $V=E\oplus O$，并把 $\mathrm e^x$ 按此分解.
			\item\label{ex:vs:B13} 设 $\dim V=n$，$W_1,\ldots,W_k$ 是 $V$ 的子空间，且 $\dim W_1+\cdots+\dim W_k>(k-1)n$. 证明 $W_1\cap\cdots\cap W_k\ne\{\zero\}$.
			\item\label{ex:vs:B14} 设 $\vv\alpha_1,\ldots,\vv\alpha_n$ 是 $V$ 的基，$W=\spn(\vv\alpha_1+\vv\alpha_2,\ \vv\alpha_2+\vv\alpha_3,\ \ldots,\ \vv\alpha_{n-1}+\vv\alpha_n)$. 求 $\dim W$，并求 $W$ 的一个补空间.
			\item\label{ex:vs:B15} 设 $V_1=\{\vv x\in\mathbb{F}^n\mid x_1=x_2=\cdots=x_n\}$，$V_2=\{\vv x\in\mathbb{F}^n\mid x_1=0\}$. 证明 $\mathbb{F}^n=V_1\oplus V_2$. 与第二章例~\ref{eg:pricedecomp} 的分解相比，哪一种分解更``自然''？（提示：思考补空间的不唯一性.）
			\item\label{ex:vs:B16} 在 $M_n(\mathbb{F})$ 中，设 $W_1$ 为对角线以下元素全为零的矩阵（上三角矩阵）全体，$W_2$ 为对角线及以上元素全为零的矩阵全体. 证明 $M_n(\mathbb{F})=W_1\oplus W_2$，并求 $\dim W_1$.
			\item\label{ex:vs:B17} （覆盖定理的另一个证明）设 $W_1,\ldots,W_s$ 是 $\mathbb{R}^n$ 的真子空间，令 $\vv v(t)=(1,t,t^2,\ldots,t^{n-1})^\trans$.
			\begin{enumerate}
				\item 证明：每个真子空间都包含于某个超平面 $\{\vv x\mid a_1x_1+\cdots+a_nx_n=0\}$（$a_i$ 不全为零）中.（提示：取 $W$ 的基并扩充为 $\mathbb{R}^n$ 的基，看``最后一个坐标''.）
				\item 证明：对上述超平面，满足 $\vv v(t)$ 落在其中的实数 $t$ 至多有 $n-1$ 个.
				\item 由此推出：$\vv v(t)$ 落在 $W_1\cup\cdots\cup W_s$ 中的 $t$ 至多有 $s(n-1)$ 个. 这给出了覆盖定理在 $\mathbb{R}^n$ 中的一个``构造性''证明. 对 $n=3$ 以及三个子空间 $x_1=x_2$，$x_2=x_3$，$x_1+x_2+x_3=0$，写出一个不在它们之中的 $\vv v(t)$.
			\end{enumerate}
			\item\label{ex:vs:B18} 设 $\dim V=n$，$W_1,\ldots,W_s$ 是 $V$ 的真子空间. 证明：存在 $V$ 的一组基 $\vv\varepsilon_1,\ldots,\vv\varepsilon_n$，其中\textbf{每个}向量都不属于 $W_1\cup\cdots\cup W_s$.（提示：逐个选取. 选 $\vv\varepsilon_k$ 时，要躲开的除了 $W_1,\ldots,W_s$，还有哪个真子空间？）
		\end{exB}
		
		\Cgroup
		\begin{exC}
			\item\label{ex:vs:C1} （覆盖定理在有限``数系''上的失效）设 $p$ 是素数，$\mathbb{F}_p=\{0,1,\ldots,p-1\}$ 按模 $p$ 作加法和乘法（可以承认它对加、减、乘以及除以非零元素都封闭. 它不包含 $\mathbb{Q}$，所以不是本书意义下的数域；但第二章中除定理~\ref{thm:Qsmallest} 与覆盖定理以外的定义和结论，在它上面照样成立）.
			\begin{enumerate}
				\item 证明 $\mathbb{F}_p^2$ 恰有 $p+1$ 个一维子空间，并且它们的并是整个 $\mathbb{F}_p^2$.
				\item 证明 $\mathbb{F}_p^2$ 不能写成 $p$ 个真子空间的并.（提示：每个一维子空间恰含 $p-1$ 个非零向量，数一数.）
				\item 指出定理~\ref{thm:cover} 的证明在 $\mathbb{F}_p$ 上从哪一步开始失效. 由第 2 问，``个数不超过 $p$ 时盖不住''，你能把正文的证明改写成这个``定量''的版本吗？
			\end{enumerate}
			\item\label{ex:vs:C2} 证明 $\ln2,\ln3,\ln5$ 在 $\mathbb{Q}$ 上线性无关. 由此说明 $\mathbb{R}$ 作为 $\mathbb{Q}$ 上的向量空间不是有限维的.（提示：算术基本定理.）
			\item\label{ex:vs:C3} 证明：全体实数列构成的向量空间不是有限维的；而``只有有限项非零''的实数列全体是它的子空间，且 $\vv e_1,\vv e_2,\ldots$（第 $i$ 项为 $1$，其余为 $0$）是它的一组``基''——每个元素都能唯一地写成其中\textbf{有限个}向量的线性组合. 这组向量为什么不是全体实数列的基？
			\item\label{ex:vs:C4} （模律）设 $W_1,W_2,W_3$ 是 $V$ 的子空间. 证明 $(W_1\cap W_3)+(W_2\cap W_3)\subset(W_1+W_2)\cap W_3$，在 $\mathbb{R}^2$ 中举例说明等号可以不成立；再证明：若 $W_1\subset W_3$，则等号一定成立.
			\item\label{ex:vs:C5} （有限域上的向量空间）在只含 $0,1$ 两个元素、按模 $2$ 运算的``数系'' $\mathbb{F}_2$ 上考虑 $\mathbb{F}_2^3$. (1) 它有多少个向量？(2) 有多少个一维子空间、多少个二维子空间？(3) 它有多少组（有序）基？推广到 $\mathbb{F}_2^n$.
			\item\label{ex:vs:C6} （塔律）设 $\mathbb{F}\subset\mathbb{K}\subset\mathbb{L}$ 是三个数域，$\mathbb{L}$ 作为 $\mathbb{K}$ 上的向量空间是 $m$ 维的，$\mathbb{K}$ 作为 $\mathbb{F}$ 上的向量空间是 $n$ 维的. 证明 $\mathbb{L}$ 作为 $\mathbb{F}$ 上的向量空间是 $mn$ 维的. 用它重新计算 B8 中数域的维数.
			\item\label{ex:vs:C7} 证明 $\mathbb{Q}(\sqrt2+\sqrt3)=\mathbb{Q}(\sqrt2,\sqrt3)$，即包含 $\sqrt2+\sqrt3$ 的最小数域与包含 $\sqrt2,\sqrt3$ 的最小数域相同.（提示：计算 $(\sqrt2+\sqrt3)^3$.）
			\item\label{ex:vs:C8} 设 $V$ 是 $n$ 维空间，$W$ 是 $k$ 维子空间（$0<k<n$）. 证明 $W$ 有无穷多个补空间. 在 $\mathbb{F}_2^2$ 中，一维子空间有几个补空间？
			\item\label{ex:vs:C9} （Fibonacci 空间的另一组基）在 B6 的记号下，改为研究满足 $a_{n+2}=a_{n+1}+a_n$ 的数列空间. 求出两个公比 $q$ 使 $\{q^n\}$ 属于这个空间，证明它们构成一组基，并推出 Fibonacci 数列的通项（Binet 公式）. 第八章会从另一个角度再次得到它.
		\end{exC}
	}
	
	% =====================================================================
	% 第三章　矩阵
	% =====================================================================
	\newcommand{\ExMatrix}{%
		\exsetheading
		\Agroup
		\begin{exA}
			\item\label{ex:mat:A1} 设 $\vv A=\begin{pmatrix}1&-1&2\\0&3&1\end{pmatrix}$，$\vv B=\begin{pmatrix}2&1\\-1&0\\1&4\end{pmatrix}$. 计算 $\vv{AB}$，$\vv{BA}$，$\vv A^\trans\vv B^\trans$，并验证 $(\vv{AB})^\trans=\vv B^\trans\vv A^\trans$.
			\item\label{ex:mat:A2} 按``列的线性组合''计算 $\begin{pmatrix}1&2&0\\3&-1&4\end{pmatrix}\begin{pmatrix}2\\-1\\1\end{pmatrix}$；再说明 $\vv e_i^\trans\vv A\vv e_j$ 是 $\vv A$ 的哪个元素，$\vv A\vv e_j$、$\vv e_i^\trans\vv A$ 又分别是什么.
			\item\label{ex:mat:A3} 计算 $\begin{pmatrix}1&1\\0&1\end{pmatrix}^n$，$\begin{pmatrix}\lambda&1\\0&\lambda\end{pmatrix}^n$，$\begin{pmatrix}\cos\theta&-\sin\theta\\ \sin\theta&\cos\theta\end{pmatrix}^n$，$\begin{pmatrix}0&1&0\\0&0&1\\0&0&0\end{pmatrix}^n$.
			\item\label{ex:mat:A4} 设 $\vv A=\begin{pmatrix}1&2\\0&1\end{pmatrix}$，$\vv B=\begin{pmatrix}1&0\\3&1\end{pmatrix}$. 计算 $(\vv A+\vv B)^2$ 与 $\vv A^2+2\vv{AB}+\vv B^2$，二者是否相等？$(\vv A+\vv B)(\vv A-\vv B)$ 与 $\vv A^2-\vv B^2$ 呢？
			\item\label{ex:mat:A5} 用初等行变换求下列矩阵的逆：
			\[\begin{pmatrix}1&2&3\\2&2&1\\3&4&3\end{pmatrix},\qquad\begin{pmatrix}1&1&1\\0&1&1\\0&0&1\end{pmatrix},\qquad \begin{pmatrix}2&1&0\\1&2&1\\0&1&2\end{pmatrix},\qquad\begin{pmatrix}0&0&0&1\\0&0&2&0\\0&3&0&0\\4&0&0&0\end{pmatrix}.\]
			\item\label{ex:mat:A6} 解矩阵方程：(1) $\begin{pmatrix}2&5\\1&3\end{pmatrix}\vv X=\begin{pmatrix}4&-6\\2&1\end{pmatrix}$；(2) $\vv X\begin{pmatrix}2&1\\1&1\end{pmatrix}=\begin{pmatrix}1&0\\0&2\end{pmatrix}$.
			\item\label{ex:mat:A7} 求矩阵 $\begin{pmatrix}1&1&a\\1&a&1\\a&1&1\end{pmatrix}$ 的秩（对 $a$ 分情况讨论）.
			\item\label{ex:mat:A8} 求矩阵 $\begin{pmatrix}1&2&-1&0&3\\2&4&1&3&0\\-1&-2&2&1&-3\\1&2&2&3&0\end{pmatrix}$ 的秩，以及列向量组的一个极大无关组，并用它表出其余各列.
			\item\label{ex:mat:A9} 写出 $3$ 阶初等矩阵 $\vv E(1,3)$，$\vv E(2(5))$，$\vv E(3,1(-2))$，并分别写出它们的逆. 计算 $\vv E(3,1(-2))\vv A$ 与 $\vv A\vv E(3,1(-2))$，其中 $\vv A$ 是任一 $3$ 阶方阵，说明效果.
			\item\label{ex:mat:A10} 把 $\begin{pmatrix}2&1\\4&3\end{pmatrix}$ 写成若干个初等矩阵的乘积.
			\item\label{ex:mat:A11} 证明：两个上三角矩阵之积仍是上三角矩阵，且积的对角元等于两者对角元之积. 可逆上三角矩阵的逆也是上三角矩阵吗？
			\item\label{ex:mat:A12} 设 $\vv A,\vv B$ 都是 $n$ 阶对称矩阵. 证明：(1) $\vv A^k$，$\vv A+\vv B$，$\vv P^\trans\vv A\vv P$ 都对称；(2) $\vv{AB}$ 对称当且仅当 $\vv{AB}=\vv{BA}$；(3) 对任意 $m\times n$ 矩阵 $\vv C$，$\vv C^\trans\vv C$ 与 $\vv C\vv C^\trans$ 都对称.
			\item\label{ex:mat:A13} (1) 设 $\vv x,\vv y\in\mathbb{F}^n$，证明 $\tr(\vv x\vv y^\trans)=\vv y^\trans\vv x$；(2) 设 $n$ 阶方阵 $\vv A$ 满足：对一切 $n$ 阶方阵 $\vv X$ 都有 $\tr(\vv{AX})=0$. 证明 $\vv A=\vv O$（提示：取 $\vv X=\vv E_{ji}$）；(3) 证明：不存在 $n$ 阶方阵 $\vv A,\vv B$ 使 $\vv{AB}-\vv{BA}=\vv I$.
			\item\label{ex:mat:A14} 设方阵 $\vv A$ 满足 $\vv A^2-\vv A-2\vv I=\vv O$. 证明 $\vv A$ 与 $\vv A+2\vv I$ 都可逆，并用 $\vv A$ 表示它们的逆. $\vv A+\vv I$ 一定可逆吗？
			\item\label{ex:mat:A15} 证明：若 $\vv A$ 可逆且 $\vv{AB}=\vv{AC}$，则 $\vv B=\vv C$；举例说明 $\vv A$ 不可逆时结论可能不成立.
			\item\label{ex:mat:A16} 设 $\vv A,\vv B$ 可逆. 求 $\begin{pmatrix}\vv A&\vv O\\ \vv O&\vv B\end{pmatrix}^{-1}$ 与 $\begin{pmatrix}\vv O&\vv A\\ \vv B&\vv O\end{pmatrix}^{-1}$，并用后者求 $\begin{pmatrix}0&0&1&2\\0&0&1&3\\2&0&0&0\\0&5&0&0\end{pmatrix}^{-1}$.
			\item\label{ex:mat:A17} 用分块乘法计算 $\begin{pmatrix}1&0&2&1\\0&1&1&3\\0&0&1&0\\0&0&0&1\end{pmatrix}^2$，并猜出 $n$ 次方的一般形式.
			\item\label{ex:mat:A18} 某公司在两个工厂生产三种产品，各工厂每天的产量（件）与三种产品的单位成本、单位售价（元）分别为
			\[\vv Q=\begin{pmatrix}100&80&60\\120&50&90\end{pmatrix}\ (\text{行：工厂；列：产品}),\qquad \vv C=\begin{pmatrix}20&30\\15&25\\40&55\end{pmatrix}\ (\text{列：成本、售价}).\]
			计算 $\vv Q\vv C$，说明它每个元素的含义，并求两个工厂每天的利润.
			\item\label{ex:mat:A19} 四座城市之间的直达航线为：$1\leftrightarrow2$，$1\leftrightarrow3$，$2\leftrightarrow3$，$3\leftrightarrow4$（双向）. 写出邻接矩阵 $\vv A$，计算 $\vv A^2$ 与 $\vv A^3$. $(\vv A^2)_{ii}$ 与城市 $i$ 的航线数有什么关系？从城市 $1$ 到城市 $4$ 恰好乘三次航班有几种走法？
			\item\label{ex:mat:A20} 设 $\vv A$ 是实反对称矩阵. 证明：对一切 $\vv x\in\mathbb{R}^n$ 有 $\vv x^\trans\vv A\vv x=0$；$\vv A^2$ 是对称矩阵.
			\item\label{ex:mat:A21} 证明：若 $\vv A$ 是 $m\times n$ 矩阵，$\vv P$ 是 $m$ 阶可逆矩阵，则 $\vv{PA}$ 与 $\vv A$ 的简化阶梯形相同.
			\item\label{ex:mat:A22} 证明：任一 $n$ 阶方阵都可以唯一地写成一个对称矩阵与一个反对称矩阵之和.（提示：$\vv A=\frac12(\vv A+\vv A^\trans)+\frac12(\vv A-\vv A^\trans)$.）
		\end{exA}
		
		\Bgroup
		\begin{exB}
			\item\label{ex:mat:B1} 分别求出全部满足 (1) $\vv A^2=\vv O$；(2) $\vv A^2=\vv I$；(3) $\vv A^2=\vv A$ 的二阶实方阵.
			\item\label{ex:mat:B2} 设 $a_1,\ldots,a_n$ 互不相同. 证明：与 $\diag(a_1,\ldots,a_n)$ 可交换的矩阵一定是对角矩阵.
			\item\label{ex:mat:B3} 证明：与\textbf{一切} $n$ 阶方阵都可交换的矩阵必为数量矩阵 $c\vv I$.（提示：考虑与 $\vv E_{ij}$ 交换.）
			\item\label{ex:mat:B4} 设 $\vv A,\vv B$ 都是幂零矩阵且 $\vv{AB}=\vv{BA}$. 证明 $\vv A+\vv B$ 与 $\vv{AB}$ 都是幂零矩阵. 举例说明``可交换''不能去掉.
			\item\label{ex:mat:B5} 证明：$\rk\begin{pmatrix}\vv A&\vv O\\ \vv O&\vv B\end{pmatrix}=\rk\vv A+\rk\vv B$，$\rk\begin{pmatrix}\vv A&\vv C\\ \vv O&\vv B\end{pmatrix}\ge\rk\vv A+\rk\vv B$，并举例说明后者可以取严格不等号.
			\item\label{ex:mat:B6} （Sherman–Morrison 公式）设 $\vv u,\vv v\in\mathbb{F}^n$，$1+\vv v^\trans\vv u\ne0$. 证明
			\[(\vv I+\vv u\vv v^\trans)^{-1}=\vv I-\frac{\vv u\vv v^\trans}{1+\vv v^\trans\vv u}.\]
			\item\label{ex:mat:B7} 设 $\vv A$ 为 $m\times n$ 矩阵，$\vv B$ 为 $n\times m$ 矩阵. 证明：若 $\vv I_m+\vv{AB}$ 可逆，则 $\vv I_n+\vv{BA}$ 也可逆，并且 $(\vv I_n+\vv{BA})^{-1}=\vv I_n-\vv B(\vv I_m+\vv{AB})^{-1}\vv A$.
			\item\label{ex:mat:B8} 证明：秩为 $1$ 的 $m\times n$ 矩阵都可以写成 $\vv u\vv v^\trans$（$\vv u\in\mathbb{F}^m$，$\vv v\in\mathbb{F}^n$ 非零）. 对这样的矩阵证明 $\vv A^2=(\tr\vv A)\vv A$（$m=n$ 时），并由此求 $\vv A^k$.
			\item\label{ex:mat:B9} 证明：秩为 $r$ 的矩阵可以写成 $r$ 个秩为 $1$ 的矩阵之和，但不能写成少于 $r$ 个秩为 $1$ 的矩阵之和.
			\item\label{ex:mat:B10} （满秩分解）证明：秩为 $r$ 的 $m\times n$ 矩阵 $\vv A$ 可以写成 $\vv A=\vv{BC}$，其中 $\vv B$ 是列满秩的 $m\times r$ 矩阵，$\vv C$ 是行满秩的 $r\times n$ 矩阵. 对 A8 中的矩阵写出一个满秩分解.
			\item\label{ex:mat:B11} 计算 $\begin{pmatrix}1&1&1\\1&1&1\\1&1&1\end{pmatrix}^n$ 与 $\begin{pmatrix}2&1&0\\0&2&1\\0&0&2\end{pmatrix}^n$.（提示：后者写成 $2\vv I+\vv N$.）
			\item\label{ex:mat:B12} 证明：$n$ 阶严格上三角矩阵（对角线及以下全为零）$\vv N$ 满足 $\vv N^n=\vv O$. 更精确地，$\vv N^k$ 的非零元只可能出现在第 $k$ 条``上对角线''及其右上方.
			\item\label{ex:mat:B13} 设 $\vv A$ 为实矩阵. 证明：若 $\vv A^\trans\vv A=\vv O$，则 $\vv A=\vv O$；若 $\tr(\vv A^\trans\vv A)=0$，结论同样成立. 复矩阵呢？
			\item\label{ex:mat:B14} 每行、每列恰有一个 $1$、其余元素为 $0$ 的方阵称为\textbf{置换矩阵}. 证明：置换矩阵之积仍是置换矩阵，置换矩阵 $\vv P$ 满足 $\vv P^\trans\vv P=\vv I$；$\vv{PA}$ 与 $\vv{AP}$ 分别对 $\vv A$ 做了什么？
			\item\label{ex:mat:B15} 对 $\vv A=\begin{pmatrix}1&2&1\\2&4&3\\1&2&2\end{pmatrix}$，求可逆矩阵 $\vv P,\vv Q$ 使 $\vv{PAQ}$ 为相抵标准形.
			\item\label{ex:mat:B16} 设 $\vv A$ 是 $n$ 阶方阵，$\vv A^k=\vv O$. 证明 $\vv I+\vv A$ 可逆，并求其逆. 若 $\vv A^2=2\vv A$，$\vv I-\vv A$ 是否可逆？
		\end{exB}
		
		\Cgroup
		\begin{exC}
			\item\label{ex:mat:C1} 设 $f:M_n(\mathbb{F})\to\mathbb{F}$ 满足 $f(\vv A+\vv B)=f(\vv A)+f(\vv B)$，$f(k\vv A)=kf(\vv A)$，且 $f(\vv{AB})=f(\vv{BA})$ 对一切 $\vv A,\vv B$ 成立. 证明 $f=c\cdot\tr$ 对某个常数 $c$ 成立.（提示：计算 $f(\vv E_{ij})$.）
			\item\label{ex:mat:C2} 设 $n$ 阶方阵 $\vv A,\vv B$ 满足 $\vv A+\vv B=\vv{AB}$. 证明 $\vv{AB}=\vv{BA}$.（提示：考察 $(\vv I-\vv A)(\vv I-\vv B)$.）
			\item\label{ex:mat:C3} （Strassen 算法）对二阶分块矩阵的乘法，直接计算需要 $8$ 次块乘法. 验证下面的公式只需 $7$ 次：令
			$M_1=(A_{11}+A_{22})(B_{11}+B_{22})$，$M_2=(A_{21}+A_{22})B_{11}$，$M_3=A_{11}(B_{12}-B_{22})$，$M_4=A_{22}(B_{21}-B_{11})$，$M_5=(A_{11}+A_{12})B_{22}$，$M_6=(A_{21}-A_{11})(B_{11}+B_{12})$，$M_7=(A_{12}-A_{22})(B_{21}+B_{22})$，则
			$C_{11}=M_1+M_4-M_5+M_7$，$C_{12}=M_3+M_5$，$C_{21}=M_2+M_4$，$C_{22}=M_1-M_2+M_3+M_6$. 由此说明 $n=2^k$ 阶矩阵乘法可以用约 $n^{\log_27}\approx n^{2.81}$ 次数的乘法完成.
			\item\label{ex:mat:C4} 设 $\vv A$ 是幂等矩阵. 用满秩分解 $\vv A=\vv{BC}$ 证明 $\vv{CB}=\vv I_r$，进而证明 $\rk\vv A=\tr\vv A$.
			\item\label{ex:mat:C5} （Kronecker 积）对 $m\times n$ 矩阵 $\vv A$ 与 $p\times q$ 矩阵 $\vv B$，定义 $\vv A\otimes\vv B=(a_{ij}\vv B)$（$mp\times nq$ 分块矩阵）. 证明 $(\vv A\otimes\vv B)(\vv C\otimes\vv D)=(\vv{AC})\otimes(\vv{BD})$，并由此说明 $\vv A,\vv B$ 可逆时 $(\vv A\otimes\vv B)^{-1}=\vv A^{-1}\otimes\vv B^{-1}$.
			\item\label{ex:mat:C6} （Hadamard 矩阵）元素全为 $\pm1$ 且满足 $\vv H\vv H^\trans=n\vv I$ 的 $n$ 阶矩阵称为 Hadamard 矩阵. 验证 $\vv H_2=\begin{pmatrix}1&1\\1&-1\end{pmatrix}$ 是，并证明 $\vv H_2\otimes\vv H_2$，$\vv H_2\otimes\vv H_2\otimes\vv H_2$ 也是. 证明：$n>2$ 时 Hadamard 矩阵的阶数必是 $4$ 的倍数.（这类矩阵用于实验设计与纠错编码.）
			\item\label{ex:mat:C7} 证明：全体形如 $\begin{pmatrix}a&-b\\b&a\end{pmatrix}$（$a,b\in\mathbb{R}$）的矩阵关于矩阵的加法和乘法``像复数一样''运算：加法、乘法封闭，乘法可交换，非零元可逆. 写出它与 $\mathbb{C}$ 之间的对应，并用它验证 $(1+\mathrm i)^8=16$.
			\item\label{ex:mat:C8} 设 $\vv A$ 是元素非负、各列之和都小于 $1$ 的方阵. 证明 $\vv A^k$ 的每个元素都随 $k\to\infty$ 趋于 $0$，并且级数 $\vv I+\vv A+\vv A^2+\cdots$（按元素）收敛到 $(\vv I-\vv A)^{-1}$. 解释它在投入产出模型中的意义：``直接消耗 $+$ 一次间接消耗 $+$ 二次间接消耗 $+\cdots$''.
			\item\label{ex:mat:C9} 设 $\vv A$ 为 $m\times n$ 矩阵. 证明：$\vv A$ 有右逆（存在 $\vv C$ 使 $\vv{AC}=\vv I_m$）当且仅当 $\rk\vv A=m$；有左逆当且仅当 $\rk\vv A=n$. 非方阵的单侧逆唯一吗？
		\end{exC}
	}
	
	% =====================================================================
	% 第四章　线性映射
	% =====================================================================
	\newcommand{\ExLinearMap}{%
		\exsetheading
		\Agroup
		\begin{exA}
			\item\label{ex:lmap:A1} 判断下列映射是否线性，是线性映射的求出核：
			\begin{enumerate}
				\item $\mathbb{R}^2\to\mathbb{R}^2$，$(x,y)^\trans\mapsto(x+y,\,2x)^\trans$；
				\item $\mathbb{R}^2\to\mathbb{R}^2$，$(x,y)^\trans\mapsto(x+1,\,y)^\trans$；
				\item $\mathbb{R}^2\to\mathbb{R}$，$(x,y)^\trans\mapsto|x|+y$；
				\item $\mathbb{R}[x]\to\mathbb{R}[x]$，$f(x)\mapsto f(x+1)$；
				\item $\mathbb{R}[x]\to\mathbb{R}[x]$，$f(x)\mapsto xf'(x)$；
				\item $M_n(\mathbb{R})\to M_n(\mathbb{R})$，$\vv X\mapsto\vv X^\trans$；
				\item $M_2(\mathbb{R})\to\mathbb{R}$，$\vv X\mapsto x_{11}x_{22}-x_{12}x_{21}$.
			\end{enumerate}
			\item\label{ex:lmap:A2} 求下列矩阵所定义的线性映射 $\vv x\mapsto\vv{Ax}$ 的核与像（各给出一组基），并验证秩—零化度定理：
			\[\begin{pmatrix}1&2&-1\\2&4&-2\end{pmatrix},\qquad\begin{pmatrix}1&0&2&1\\0&1&-1&1\\1&1&1&2\end{pmatrix}.\]
			\item\label{ex:lmap:A3} 写出线性映射 $\mathscr T:\mathbb{R}^3\to\mathbb{R}^2$，$\mathscr T(x_1,x_2,x_3)^\trans=(x_1-2x_2+x_3,\ 3x_2+4x_3)^\trans$ 在标准基下的矩阵.
			\item\label{ex:lmap:A4} 求下列平面变换在标准基下的矩阵：(1) 逆时针旋转 $90^\circ$；(2) 关于直线 $y=x$ 的反射；(3) 向直线 $y=x$ 的正投影；(4) 沿 $x$ 方向的剪切 $(x,y)\mapsto(x+2y,y)$. 计算 (2)(3) 的矩阵的平方，并解释结果.
			\item\label{ex:lmap:A5} 求求导变换 $\mathscr D:\mathbb{R}[x]_3\to\mathbb{R}[x]_3$ 在基 $1,x,x^2$ 下的矩阵，以及在基 $1,\ x-1,\ (x-1)^2$ 下的矩阵.
			\item\label{ex:lmap:A6} 求转置变换 $\vv X\mapsto\vv X^\trans$ 在 $M_2(\mathbb{F})$ 的基 $\vv E_{11},\vv E_{12},\vv E_{21},\vv E_{22}$ 下的矩阵. 它的平方是什么？
			\item\label{ex:lmap:A7} 在 $\mathbb{R}^3$ 中，求从基 $\vv\varepsilon_1=(1,0,0)^\trans$，$\vv\varepsilon_2=(1,1,0)^\trans$，$\vv\varepsilon_3=(1,1,1)^\trans$ 到基 $\vv\eta_1=(1,2,1)^\trans$，$\vv\eta_2=(2,3,3)^\trans$，$\vv\eta_3=(3,7,1)^\trans$ 的过渡矩阵，并求 $\vv\alpha=(3,4,2)^\trans$ 在两组基下的坐标.
			\item\label{ex:lmap:A8} 设线性变换 $\mathscr T$ 在基 $\vv\varepsilon_1,\vv\varepsilon_2$ 下的矩阵为 $\begin{pmatrix}2&1\\1&2\end{pmatrix}$. 求它在基 $\vv\eta_1=\vv\varepsilon_1+\vv\varepsilon_2$，$\vv\eta_2=\vv\varepsilon_1-\vv\varepsilon_2$ 下的矩阵.
			\item\label{ex:lmap:A9} 证明：线性映射的复合仍是线性映射；可逆（双射）线性映射的逆映射仍是线性映射.
			\item\label{ex:lmap:A10} 证明：线性映射 $\mathscr T$ 是单射，当且仅当它把每个线性无关组都映为线性无关组.
			\item\label{ex:lmap:A11} 是否存在线性变换 $\mathscr T:\mathbb{R}^2\to\mathbb{R}^2$，使
			(1) $\mathscr T(1,1)^\trans=(1,0)^\trans$，$\mathscr T(1,-1)^\trans=(0,1)^\trans$，$\mathscr T(2,0)^\trans=(1,1)^\trans$？
			(2) 同上，但最后一个条件改为 $\mathscr T(2,0)^\trans=(1,2)^\trans$？
			存在时求出 $\mathscr T(x,y)^\trans$.
			\item\label{ex:lmap:A12} (1) 线性映射 $\mathscr T:\mathbb{R}^5\to\mathbb{R}^3$ 是满射，$\dim\ker\mathscr T$ 是多少？(2) 是否存在从 $\mathbb{R}^3$ 到 $\mathbb{R}^5$ 的满射线性映射？是否存在从 $\mathbb{R}^5$ 到 $\mathbb{R}^3$ 的单射线性映射？
			\item\label{ex:lmap:A13} 设 $\mathscr T:U\to V$，$\mathscr S:V\to W$ 线性. 证明 $\ker\mathscr T\subset\ker(\mathscr{ST})$，$\imm(\mathscr{ST})\subset\imm\mathscr S$，从而 $\rk(\mathscr{ST})\le\min\{\rk\mathscr S,\rk\mathscr T\}$.
			\item\label{ex:lmap:A14} 求迹映射 $\tr:M_n(\mathbb{F})\to\mathbb{F}$ 的核的维数.
			\item\label{ex:lmap:A15} 证明：相似矩阵有相同的迹（利用命题~\ref{prop:trace}(3)）. 由此判断 $\begin{pmatrix}1&2\\0&3\end{pmatrix}$ 与 $\begin{pmatrix}1&0\\0&2\end{pmatrix}$ 是否相似.
			\item\label{ex:lmap:A16} 证明：与单位矩阵相似的矩阵只有单位矩阵本身. 由此说明 $\begin{pmatrix}1&1\\0&1\end{pmatrix}$ 与 $\vv I$ 不相似.
			\item\label{ex:lmap:A17} 定义求值映射 $\mathscr E:\mathbb{R}[x]_3\to\mathbb{R}^3$，$f\mapsto(f(0),f(1),f(2))^\trans$. 求它在基 $1,x,x^2$ 与标准基下的矩阵，并证明 $\mathscr E$ 是同构.
			\item\label{ex:lmap:A18} 设 $\mathscr T$ 是 $\mathbb{R}^3$ 上的线性变换，$\mathscr T\vv e_1=(1,1,0)^\trans$，$\mathscr T\vv e_2=(0,1,1)^\trans$，$\mathscr T\vv e_3=(1,2,1)^\trans$. 求 $\mathscr T(x_1,x_2,x_3)^\trans$，$\ker\mathscr T$ 与 $\imm\mathscr T$.
			\item\label{ex:lmap:A19} 称 $\vv x\mapsto\vv{Ax}+\vv b$ 为\textbf{仿射变换}. 证明两个仿射变换的复合仍是仿射变换；仿射变换是线性变换当且仅当 $\vv b=\zero$.
			\item\label{ex:lmap:A20} 设 $\vv R(\theta)=\begin{pmatrix}\cos\theta&-\sin\theta\\ \sin\theta&\cos\theta\end{pmatrix}$ 是平面上绕原点旋转角 $\theta$ 的变换在标准基下的矩阵. 从几何上说明``先转 $\beta$ 再转 $\alpha$''就是转 $\alpha+\beta$，从而 $\vv R(\alpha)\vv R(\beta)=\vv R(\alpha+\beta)$；比较两边矩阵的元素，推出两角和的正弦、余弦公式.
		\end{exA}
		
		\Bgroup
		\begin{exB}
			\item\label{ex:lmap:B1} 记 $L(V,W)$ 为 $V$ 到 $W$ 的全体线性映射. 证明它关于映射的加法与数乘构成向量空间，且 $\dim L(V,W)=\dim V\cdot\dim W$.（提示：取定基后，线性映射与矩阵一一对应.）
			\item\label{ex:lmap:B2} 设 $\mathscr T$ 是 $n$ 维空间 $V$ 上的线性变换，$\mathscr T^2=\mathscr O$. 证明 $\imm\mathscr T\subset\ker\mathscr T$，从而 $\rk\mathscr T\le n/2$；并证明存在一组基，使 $\mathscr T$ 的矩阵只在 $(1,2),(3,4),\ldots,(2r-1,2r)$ 这 $r$ 个位置上为 $1$，其余元素全为 $0$，其中 $r=\rk\mathscr T$.
			\item\label{ex:lmap:B3} 设 $\mathscr T,\mathscr S$ 是有限维空间上的线性变换. 证明 $\dim\ker(\mathscr{ST})\le\dim\ker\mathscr S+\dim\ker\mathscr T$.
			\item\label{ex:lmap:B4} 设 $V$ 有限维，$\mathscr S,\mathscr T$ 是 $V$ 上的线性变换且 $\mathscr{ST}=\mathscr I$. 证明 $\mathscr{TS}=\mathscr I$.
			\item\label{ex:lmap:B5} 在无限维空间 $\mathbb{R}[x]$ 上，令 $\mathscr D$ 为求导，$\mathscr J$ 为 $f\mapsto\int_0^xf(t)\,\mathrm dt$. 证明 $\mathscr{DJ}=\mathscr I$，但 $\mathscr{JD}\ne\mathscr I$；$\mathscr D$ 是满射但不是单射，$\mathscr J$ 是单射但不是满射. 这与 B4 是否矛盾？
			\item\label{ex:lmap:B6} 固定 $\vv A\in M_n(\mathbb{F})$，令 $\mathscr T(\vv X)=\vv{AX}-\vv{XA}$. 证明 $\imm\mathscr T$ 含于迹为零的矩阵全体，从而 $\mathscr T$ 不是满射；并证明 $\ker\mathscr T$ 包含 $\vv I,\vv A,\vv A^2,\ldots$.
			\item\label{ex:lmap:B7} 设 $\vv A=\begin{pmatrix}a&b\\c&d\end{pmatrix}$. 求线性变换 $\vv X\mapsto\vv{AX}$ 在 $M_2(\mathbb{F})$ 的基 $\vv E_{11},\vv E_{21},\vv E_{12},\vv E_{22}$ 下的矩阵. 若换成 $\vv X\mapsto\vv{XA}$ 呢？
			\item\label{ex:lmap:B8} 设线性变换 $\mathscr P$ 满足 $\mathscr P^2=\mathscr P$. 证明存在一组基使 $\mathscr P$ 的矩阵为 $\begin{pmatrix}\vv I_r&\vv O\\ \vv O&\vv O\end{pmatrix}$，其中 $r=\rk\mathscr P$.
			\item\label{ex:lmap:B9} 设线性变换 $\mathscr T$ 满足 $\mathscr T^2=\mathscr I$. 证明 $V=\ker(\mathscr T-\mathscr I)\oplus\ker(\mathscr T+\mathscr I)$，从而存在一组基使 $\mathscr T$ 的矩阵为 $\diag(\vv I_p,-\vv I_q)$.（提示：$\vv\alpha=\frac12(\vv\alpha+\mathscr T\vv\alpha)+\frac12(\vv\alpha-\mathscr T\vv\alpha)$.）
			\item\label{ex:lmap:B10} 设 $\mathscr T:V\to W$ 线性，$U$ 是 $V$ 的子空间. 证明 $\dim\mathscr T(U)=\dim U-\dim(U\cap\ker\mathscr T)$.
			\item\label{ex:lmap:B11} 在由 $\mathrm e^x\cos x$，$\mathrm e^x\sin x$ 张成的二维函数空间中，求求导变换的矩阵，证明它可逆，并利用逆矩阵求出 $\int\mathrm e^x\sin x\,\mathrm dx$（不必用分部积分）.
			\item\label{ex:lmap:B12} （齐次坐标）平面上的平移 $(x,y)\mapsto(x+a,y+b)$ 不是线性的. 把点 $(x,y)$ 写成三维向量 $(x,y,1)^\trans$，证明平移、旋转、伸缩都可以写成 $3\times3$ 矩阵的乘法，且仿射变换的复合对应矩阵的乘积. 写出``先绕原点旋转 $90^\circ$，再平移 $(1,2)$''与``先平移再旋转''的矩阵，它们相等吗？（计算机图形学正是这样处理变换的.）
			\item\label{ex:lmap:B13} 证明：若 $\vv A$ 与 $\vv B$ 相似，则 $\vv A^k$ 与 $\vv B^k$ 相似，$\vv A^\trans$ 与 $\vv B^\trans$ 相似，$f(\vv A)$ 与 $f(\vv B)$ 相似（$f$ 为多项式）.
			\item\label{ex:lmap:B14} 证明 $\begin{pmatrix}1&0\\0&0\end{pmatrix}$ 与 $\begin{pmatrix}0&1\\0&0\end{pmatrix}$ 相抵但不相似. 从``线性映射 $V\to W$''与``线性变换 $V\to V$''的区别解释：为什么相抵只需看秩，而相似要复杂得多？
			\item\label{ex:lmap:B15} 设 $\mathscr T$ 是 $V$ 上的线性变换，$\vv\alpha\in V$ 满足 $\mathscr T^{k-1}\vv\alpha\ne\zero$，$\mathscr T^k\vv\alpha=\zero$. 证明 $\vv\alpha,\mathscr T\vv\alpha,\ldots,\mathscr T^{k-1}\vv\alpha$ 线性无关. 若 $k=n=\dim V$，求 $\mathscr T$ 在基 $\mathscr T^{n-1}\vv\alpha,\ldots,\mathscr T\vv\alpha,\vv\alpha$ 下的矩阵.
		\end{exB}
		
		\Cgroup
		\begin{exC}
			\item\label{ex:lmap:C1} （Fitting 分解）设 $\mathscr T$ 是 $n$ 维空间 $V$ 上的线性变换. 证明：
			\begin{enumerate}
				\item $\ker\mathscr T\subset\ker\mathscr T^2\subset\cdots$，$\imm\mathscr T\supset\imm\mathscr T^2\supset\cdots$，且两条链都在某一步之后不再变化；
				\item $V=\ker\mathscr T^n\oplus\imm\mathscr T^n$，两者都是 $\mathscr T$ 的不变子空间，$\mathscr T$ 在前者上幂零，在后者上可逆.
			\end{enumerate}
			\item\label{ex:lmap:C2} （商空间）设 $W$ 是 $V$ 的子空间. 对 $\vv\alpha\in V$，称集合 $\vv\alpha+W$ 为 $\vv\alpha$ 所在的陪集，全体陪集记为 $V/W$. 定义 $(\vv\alpha+W)+(\vv\beta+W)=(\vv\alpha+\vv\beta)+W$，$k(\vv\alpha+W)=k\vv\alpha+W$.
			\begin{enumerate}
				\item 证明这两个运算是良定义的，且 $V/W$ 构成向量空间；
				\item 证明 $\dim V/W=\dim V-\dim W$；
				\item （第一同构定理/同态基本定理）对线性映射 $\mathscr T:V\to U$，证明 $V/\ker\mathscr T\cong\imm\mathscr T$. 这给出秩—零化度定理的另一种理解.
			\end{enumerate}
			\item\label{ex:lmap:C3} 把 $\mathbb{C}$ 看作 $\mathbb{R}$ 上的二维空间，基为 $1,\mathrm i$. (1) 求``乘以 $a+b\mathrm i$''这个变换的矩阵；(2) 求``取共轭''这个变换的矩阵；(3) 证明 $\mathbb{R}$-线性变换 $\begin{pmatrix}p&q\\r&s\end{pmatrix}$ 是``乘以某个复数''，当且仅当 $p=s$，$q=-r$.
			\item\label{ex:lmap:C4} 证明：$n$ 维空间上的任一线性变换 $\mathscr T$，都存在次数不超过 $n^2$ 的非零多项式 $f$ 使 $f(\mathscr T)=\mathscr O$.（提示：$\mathscr I,\mathscr T,\ldots,\mathscr T^{n^2}$ 在 $L(V,V)$ 中线性相关. 第八章的 Cayley–Hamilton 定理会把 $n^2$ 改进为 $n$.）
			\item\label{ex:lmap:C5} （Haar 小波）在 $\mathbb{R}^4$ 中取基 $\vv h_1=(1,1,1,1)^\trans$，$\vv h_2=(1,1,-1,-1)^\trans$，$\vv h_3=(1,-1,0,0)^\trans$，$\vv h_4=(0,0,1,-1)^\trans$. 求信号 $\vv x=(8,6,3,4)^\trans$ 在这组基下的坐标，并解释每个坐标的含义. 若把绝对值最小的坐标改为 $0$ 再变换回来，误差是多少？这是图像与音频压缩的思想之一.
			\item\label{ex:lmap:C6} 设 $\mathscr T$ 是 $V$ 上的线性变换. 若 $\mathscr T$ 与 $V$ 上的一切线性变换都可交换，证明 $\mathscr T=c\mathscr I$. 若只要求与一切\textbf{可逆}线性变换可交换呢？
			\item\label{ex:lmap:C7} 设 $\mathscr T:V\to V$ 线性，$\dim V=n$. 证明：$\rk\mathscr T^k-\rk\mathscr T^{k+1}$ 关于 $k$ 单调不增. 由此说明：若 $\rk\mathscr T=\rk\mathscr T^2$，则对一切 $k$，$\rk\mathscr T^k=\rk\mathscr T$.
			\item\label{ex:lmap:C8} （差分算子）在全体实数列构成的空间上，定义左移 $\mathscr S(a_1,a_2,a_3,\ldots)=(a_2,a_3,\ldots)$. 证明 $\mathscr S$ 是线性的、满射但不是单射；满足 $a_{n+2}=a_{n+1}+a_n$ 的数列全体恰好是 $\ker(\mathscr S^2-\mathscr S-\mathscr I)$. 由第二章章末习题~\ref{ex:vs:C9}，这个核是二维的.
		\end{exC}
	}
	
	% =====================================================================
	% 第五章　行列式
	% =====================================================================
	\newcommand{\ExDeterminant}{%
		\exsetheading
		\Agroup
		\begin{exA}
			\item\label{ex:det:A1} 求下列排列的逆序数并判断奇偶：$315426$，$42531$，$13\cdots(2n-1)\,24\cdots(2n)$.
			\item\label{ex:det:A2} (1) 在 $4$ 阶行列式中，项 $a_{13}a_{24}a_{32}a_{41}$ 带什么符号？(2) 选择 $i,j$，使 $a_{1i}a_{23}a_{35}a_{5j}a_{44}$ 成为 $5$ 阶行列式中带正号的一项.
			\item\label{ex:det:A3} 计算：
			\[\begin{vmatrix}2&1&-1\\3&0&2\\1&4&5\end{vmatrix},\qquad
			\begin{vmatrix}1&2&0&1\\2&1&1&0\\0&1&2&1\\1&0&1&2\end{vmatrix},\qquad
			\begin{vmatrix}3&1&1&1\\1&3&1&1\\1&1&3&1\\1&1&1&3\end{vmatrix}.\]
			\item\label{ex:det:A4} 计算 $\begin{vmatrix}1&-1&2&3\\2&0&1&-1\\0&1&-1&2\\3&1&0&1\end{vmatrix}$，分别用``打洞''化三角形与按第一列展开两种方法，比较哪种更快.
			\item\label{ex:det:A5} 计算 $n$ 阶行列式 $\begin{vmatrix}0&\cdots&0&a_1\\0&\cdots&a_2&0\\ \vdots&&&\vdots\\a_n&\cdots&0&0\end{vmatrix}$（只有副对角线上有元素）.
			\item\label{ex:det:A6} 不展开，证明：
			\[\begin{vmatrix}a+b&b+c&c+a\\a_1+b_1&b_1+c_1&c_1+a_1\\a_2+b_2&b_2+c_2&c_2+a_2\end{vmatrix}=2\begin{vmatrix}a&b&c\\a_1&b_1&c_1\\a_2&b_2&c_2\end{vmatrix}.\]
			\item\label{ex:det:A7} 解方程 $\begin{vmatrix}x&1&1\\1&x&1\\1&1&x\end{vmatrix}=0$ 与 $\begin{vmatrix}1&1&1\\1&2&x\\1&4&x^2\end{vmatrix}=0$.
			\item\label{ex:det:A8} 设 $\vv A$ 为 $3$ 阶方阵，$\det\vv A=3$. 求 $\det(2\vv A)$，$\det(-\vv A)$，$\det(\vv A^{-1})$，$\det(\vv A^*)$，$\det(\vv A^\trans\vv A)$，$\det(3\vv A^{-1}-2\vv A^*)$.
			\item\label{ex:det:A9} 证明：(1) 若 $n$ 阶方阵满足 $\vv A^\trans\vv A=\vv I$，则 $\det\vv A=\pm1$；(2) 若 $\vv A^2=\vv I$，则 $\det\vv A=\pm1$；(3) 若 $\vv A^2=\vv A$ 且 $\vv A\ne\vv I$，则 $\det\vv A=0$.
			\item\label{ex:det:A10} 求下列矩阵的伴随矩阵与逆矩阵：\[\begin{pmatrix}a&b\\c&d\end{pmatrix},\qquad\begin{pmatrix}1&2&3\\0&1&4\\5&6&0\end{pmatrix}.\]
			\item\label{ex:det:A11} 用 Cramer 法则解方程组 $2x_1+x_2+x_3=3$，$x_1+3x_2+x_3=0$，$x_1+x_2+4x_3=8$.
			\item\label{ex:det:A12} 求以 $(1,1)$，$(4,2)$，$(2,5)$ 为顶点的三角形的面积；判断 $(1,2)$，$(3,5)$，$(7,11)$ 三点是否共线.
			\item\label{ex:det:A13} 求以 $(1,0,1)^\trans$，$(2,1,0)^\trans$，$(0,3,1)^\trans$ 为棱的平行六面体的体积，以及以原点和这三个向量的终点为顶点的四面体的体积.
			\item\label{ex:det:A14} $\lambda$ 取何值时，齐次方程组 $(5-\lambda)x+2y+2z=0$，$2x+(6-\lambda)y=0$，$2x+(4-\lambda)z=0$ 有非零解？
			\item\label{ex:det:A15} 设 \[\vv A=\begin{pmatrix}1&2&3&4\\2&3&4&1\\3&4&1&2\\4&1&2&3\end{pmatrix}.\] 求 $A_{11}+A_{12}+A_{13}+A_{14}$（$A_{1j}$ 是代数余子式）.（提示：它是把第一行换成全 $1$ 后的行列式.）
			\item\label{ex:det:A16} 证明：过平面上两个不同点 $(x_1,y_1)$，$(x_2,y_2)$ 的直线方程可以写成 \[\begin{vmatrix}x&y&1\\x_1&y_1&1\\x_2&y_2&1\end{vmatrix}=0.\]
			\item\label{ex:det:A17} 设 $\vv A,\vv B$ 分别是 $m$ 阶、$n$ 阶方阵，证明 $\begin{vmatrix}\vv A&\vv O\\ \vv C&\vv B\end{vmatrix}=\det\vv A\det\vv B$，并计算 $\begin{vmatrix}\vv O&\vv A\\ \vv B&\vv O\end{vmatrix}$.
			\item\label{ex:det:A18} 计算 $4$ 阶三对角行列式 $\begin{vmatrix}2&-1&0&0\\-1&2&-1&0\\0&-1&2&-1\\0&0&-1&2\end{vmatrix}$.
			\item\label{ex:det:A19} 设 $\vv A$ 是 $n$ 阶可逆矩阵，$n\ge2$. 证明 $\vv A^*$ 可逆，且 $(\vv A^*)^{-1}=\frac1{\det\vv A}\vv A=(\vv A^{-1})^*$.
			\item\label{ex:det:A20} 用行列式判断：向量 $(1,2,3)^\trans$，$(4,5,6)^\trans$，$(7,8,9)^\trans$ 是否线性无关？把 $9$ 改成 $10$ 呢？
			\item\label{ex:det:A21} 设 $\vv A$ 是 $n$ 阶反对称矩阵（$\vv A^\trans=-\vv A$），$n$ 为奇数. 证明 $\det\vv A=0$.（提示：对 $\vv A^\trans=-\vv A$ 两边取行列式.）
		\end{exA}
		
		\Bgroup
		\begin{exB}
			\item\label{ex:det:B1} 计算 $n$ 阶三对角行列式 $D_n=\begin{vmatrix}2&-1&&\\-1&2&\ddots&\\&\ddots&\ddots&-1\\&&-1&2\end{vmatrix}$.（提示：按第一行展开得 $D_n=2D_{n-1}-D_{n-2}$.）
			\item\label{ex:det:B2} 计算 $n$ 阶行列式 $\begin{vmatrix}a&b&\cdots&b\\b&a&\cdots&b\\ \vdots&&\ddots&\vdots\\b&b&\cdots&a\end{vmatrix}$ 与 $\begin{vmatrix}x&a&\cdots&a\\-a&x&\cdots&a\\ \vdots&&\ddots&\vdots\\-a&-a&\cdots&x\end{vmatrix}$.
			\item\label{ex:det:B3} （爪形行列式）设 $a_1\cdots a_n\ne0$，计算 $\begin{vmatrix}a_0&1&1&\cdots&1\\1&a_1&0&\cdots&0\\1&0&a_2&\cdots&0\\ \vdots&&&\ddots&\\1&0&0&\cdots&a_n\end{vmatrix}$.
			\item\label{ex:det:B4} （矩阵行列式引理）证明 $\det(\vv I_n+\vv u\vv v^\trans)=1+\vv v^\trans\vv u$，更一般地，$\vv A$ 可逆时 $\det(\vv A+\vv u\vv v^\trans)=(1+\vv v^\trans\vv A^{-1}\vv u)\det\vv A$. 用它重做 B2 的第一个行列式.
			\item\label{ex:det:B5} 设 $\vv A$ 为 $2\times3$ 矩阵，$\vv B$ 为 $3\times2$ 矩阵. 证明 $\det(\vv{AB})$ 等于 $\vv A$ 的三个 $2$ 阶子式与 $\vv B$ 的对应 $2$ 阶子式乘积之和（这是 Cauchy–Binet 公式最简单的情形）. 由此说明 $\vv{BA}$（$3$ 阶方阵）的行列式为零.
			\item\label{ex:det:B6} 设 $\vv A$ 是整数矩阵. 证明：$\vv A^{-1}$ 存在且也是整数矩阵，当且仅当 $\det\vv A=\pm1$.
			\item\label{ex:det:B7} 证明：$n\ge2$ 时 $(\vv A^*)^*=(\det\vv A)^{n-2}\vv A$.（先设 $\vv A$ 可逆；不可逆的情形另行讨论.）
			\item\label{ex:det:B8} 设 $\vv A,\vv B$ 都是 $n$ 阶方阵. 证明 $\begin{vmatrix}\vv A&\vv B\\ \vv B&\vv A\end{vmatrix}=\det(\vv A+\vv B)\det(\vv A-\vv B)$.（提示：分块初等变换.）
			\item\label{ex:det:B9} 设 $\vv A$ 可逆. 证明 $\begin{vmatrix}\vv A&\vv B\\ \vv C&\vv D\end{vmatrix}=\det\vv A\cdot\det(\vv D-\vv{CA}^{-1}\vv B)$；若还有 $\vv{AC}=\vv{CA}$，则它等于 $\det(\vv{AD}-\vv{CB})$.
			\item\label{ex:det:B10} 用 Fibonacci 矩阵 $\vv F=\begin{pmatrix}1&1\\1&0\end{pmatrix}$ 满足 $\vv F^n=\begin{pmatrix}F_{n+1}&F_n\\F_n&F_{n-1}\end{pmatrix}$（归纳法证明），推出 Cassini 恒等式 $F_{n+1}F_{n-1}-F_n^2=(-1)^n$.
			\item\label{ex:det:B11} 计算 $n$ 阶行列式 $D_n=\begin{vmatrix}1+x^2&x&&\\x&1+x^2&\ddots&\\&\ddots&\ddots&x\\&&x&1+x^2\end{vmatrix}$.
			\item\label{ex:det:B12} 设 $\omega=\mathrm e^{2\pi\mathrm i/3}$. 证明
			\[\begin{vmatrix}a&b&c\\c&a&b\\b&c&a\end{vmatrix}=(a+b+c)(a+\omega b+\omega^2c)(a+\omega^2b+\omega c)=a^3+b^3+c^3-3abc.\]
			\item\label{ex:det:B13} 证明：元素为 $a_{ij}=\min(i,j)$ 的 $n$ 阶行列式等于 $1$.
			\item\label{ex:det:B14} 证明：平面上四点 $(x_i,y_i)$（$i=1,2,3,4$）共圆（或共线），当且仅当 $\begin{vmatrix}x_i^2+y_i^2&x_i&y_i&1\end{vmatrix}_{4\times4}=0$（每行对应一个点）.
						\item\label{ex:det:B15} 设 $\vv A$ 为 $n$ 阶方阵（$n\ge2$）. 证明
			\[\rk\vv A^*=\begin{cases}n,&\rk\vv A=n,\\1,&\rk\vv A=n-1,\\0,&\rk\vv A\le n-2.\end{cases}\]
			（提示：$\vv A^*$ 的元素都是 $\vv A$ 的 $n-1$ 阶子式，用定理~\ref{thm:minorrank}；$\rk\vv A=n-1$ 时考察 $\vv A\vv A^*=(\det\vv A)\vv I$.）
			\item\label{ex:det:B16} 设 $f(x)=\begin{vmatrix}x&1&2&3\\1&x&2&3\\1&2&x&3\\1&2&3&x\end{vmatrix}$. 不展开，求 $f(x)=0$ 的全部根.
		\end{exB}
		
		\Cgroup
		\begin{exC}
			\item\label{ex:det:C1} （行列式的唯一性）设 $D:M_n(\mathbb{F})\to\mathbb{F}$ 对每一行都是线性的、交换两行变号，且 $D(\vv I)=1$. 证明 $D=\det$.（提示：先说明 $D$ 在初等矩阵上的取值，再对任意矩阵用初等变换.）
			\item\label{ex:det:C2} （Cauchy 行列式）证明
			\[\det\Big(\frac1{a_i+b_j}\Big)_{n\times n}=\frac{\prod_{i<j}(a_j-a_i)(b_j-b_i)}{\prod_{i,j}(a_i+b_j)}.\]
			先对 $n=2$ 直接验证，再用``各行减去最后一行''归纳. 取 $a_i=i$，$b_j=j-1$，得到 Hilbert 矩阵 $(1/(i+j-1))$ 的行列式，它为什么在数值计算中臭名昭著？
			\item\label{ex:det:C3} （排列的符号是同态）对排列 $\sigma$，令 $\vv P_\sigma$ 为把 $\vv e_j$ 映到 $\vv e_{\sigma(j)}$ 的置换矩阵. 证明 $\vv P_\sigma\vv P_\tau=\vv P_{\sigma\tau}$，$\det\vv P_\sigma=(-1)^{\tau(\sigma)}$，并由行列式的乘积公式推出 $\operatorname{sgn}(\sigma\tau)=\operatorname{sgn}\sigma\cdot\operatorname{sgn}\tau$.
			\item\label{ex:det:C4} （矩阵—树定理）对一个连通无向图，令 $\vv D$ 为各顶点度数构成的对角矩阵，$\vv A$ 为邻接矩阵，$\vv L=\vv D-\vv A$ 称为 Laplace 矩阵. 定理说：图的生成树个数等于 $\vv L$ 的任一个 $(n-1)$ 阶主子式. 对 $4$ 个顶点的完全图验证它有 $16$ 棵生成树，对一个 $4$ 个顶点的环验证它有 $4$ 棵生成树.
			\item\label{ex:det:C5} 证明：若 $n$ 阶矩阵的元素全为 $\pm1$，则它的行列式能被 $2^{n-1}$ 整除. 求 $3$ 阶此类行列式的最大值.
			\item\label{ex:det:C6} （Jacobi 公式的雏形）证明 $\det(\vv I+t\vv B)=1+t\tr\vv B+(\text{$t^2$ 及更高次项})$. 由此说明：若 $\vv A(t)$ 的元素都是 $t$ 的可导函数，且 $\vv A(0)=\vv I$，则 $\frac{\mathrm d}{\mathrm dt}\det\vv A(t)\big|_{t=0}=\tr\vv A'(0)$.
			\item\label{ex:det:C7} （Sylvester 结式）设 $f=a_0x^2+a_1x+a_2$，$g=b_0x^2+b_1x+b_2$（$a_0b_0\ne0$）. 证明 $f,g$ 有公共根当且仅当
			\[\begin{vmatrix}a_0&a_1&a_2&0\\0&a_0&a_1&a_2\\b_0&b_1&b_2&0\\0&b_0&b_1&b_2\end{vmatrix}=0.\]
			（提示：有公共根当且仅当存在次数 $\le1$ 的不全为零的 $u,v$ 使 $uf+vg=0$，与第七章的 Bézout 等式对照.）
			\item\label{ex:det:C8} （Dodgson 凝聚法）对 $3$ 阶矩阵 $\vv A$，记 $M_{ij}$ 为余子式，$c=a_{22}$. 验证 $\det\vv A\cdot c=M_{11}M_{33}-M_{13}M_{31}$. 它允许只计算 $2$ 阶行列式就得到高阶行列式. 《爱丽丝漫游奇境》的作者刘易斯·卡罗尔（即 Dodgson）于 1866 年发表了这一方法.
			\item\label{ex:det:C9} 设 $\vv A$ 是 $n$ 阶实反对称矩阵. 证明 $\det(\vv I+\vv A)\ge1$.（提示：$\det(\vv I+\vv A)=\det(\vv I-\vv A)$，考察 $\det(\vv I+t\vv A)$ 作为 $t$ 的多项式只含偶次项，且 $\vv I+t\vv A$ 对一切实数 $t$ 都可逆.）
		\end{exC}
	}
	
	% =====================================================================
	% 第六章　线性方程组解的结构
	% =====================================================================
	\newcommand{\ExSolutions}{%
		\exsetheading
		\Agroup
		\begin{exA}
			\item\label{ex:sol:A1} 求下列齐次方程组的基础解系与通解：
			\[\text{(1)}\ \left\{\begin{array}{l}x_1+x_2-x_3-x_4=0,\\2x_1-5x_2+3x_3+2x_4=0,\\7x_1-7x_2+3x_3+x_4=0;\end{array}\right.\qquad
			\text{(2)}\ \left\{\begin{array}{l}x_1+2x_2+x_3-x_4=0,\\3x_1+6x_2-x_3-3x_4=0,\\5x_1+10x_2+x_3-5x_4=0.\end{array}\right.\]
			\item\label{ex:sol:A2} 求下列方程组的通解，写成``特解 $+$ 基础解系的线性组合''：
			\[\text{(1)}\ \left\{\begin{array}{l}x_1+x_2-3x_3-x_4=1,\\3x_1-x_2-3x_3+4x_4=4,\\x_1+5x_2-9x_3-8x_4=0;\end{array}\right.\qquad
			\text{(2)}\ \left\{\begin{array}{l}2x_1+x_2-x_3+x_4=1,\\4x_1+2x_2-2x_3+x_4=2,\\2x_1+x_2-x_3-x_4=1.\end{array}\right.\]
			\item\label{ex:sol:A3} 设 $\vv A$ 是 $4\times5$ 矩阵，$\rk\vv A=3$. (1) $\vv{Ax}=\zero$ 的基础解系含几个向量？(2) $\vv{Ax}=\vv b$ 是否一定有解？(3) 若 $\rk\vv A=4$ 呢？
			\item\label{ex:sol:A4} 不解方程，用秩判断下列方程组是否有解：(1) $x_1+x_2=1$，$x_1-x_2=1$，$2x_1+3x_2=2$；(2) $x_1+x_2+x_3=1$，$2x_1+2x_2+2x_3=3$.
			\item\label{ex:sol:A5} 设 $\vv\eta_1,\vv\eta_2$ 是 $\vv{Ax}=\vv b$（$\vv b\ne\zero$）的两个解. 下列向量中哪些是 $\vv{Ax}=\vv b$ 的解，哪些是 $\vv{Ax}=\zero$ 的解？
			$\vv\eta_1+\vv\eta_2$，$2\vv\eta_1-\vv\eta_2$，$\frac12(\vv\eta_1+\vv\eta_2)$，$\vv\eta_1-\vv\eta_2$，$3\vv\eta_1-2\vv\eta_2$.
			\item\label{ex:sol:A6} 设 $3$ 元方程组 $\vv{Ax}=\vv b$ 的系数矩阵秩为 $2$，$\vv\eta_1=(1,2,3)^\trans$ 与 $\vv\eta_2=(1,-2,-1)^\trans$ 是它的两个解. 求通解.
			\item\label{ex:sol:A7} 写出一个以 $\vv x=(1,0,2)^\trans+k(1,1,-1)^\trans$（$k$ 任意）为全部解的线性方程组.
			\item\label{ex:sol:A8} 判断三个平面 $x+y+z=1$，$x+2y+3z=2$，$2x+3y+4z=4$ 的位置关系；若把最后一个常数改为 $3$ 呢？
			\item\label{ex:sol:A9} 举例说明：$\vv{Ax}=\zero$ 只有零解时，$\vv{Ax}=\vv b$ 可能无解；但若有解，解必唯一.
			\item\label{ex:sol:A10} 解矩阵方程 $\begin{pmatrix}1&2&3\\2&2&1\\3&4&3\end{pmatrix}\vv X=\begin{pmatrix}2&5\\3&1\\4&3\end{pmatrix}$.
			\item\label{ex:sol:A11} 设 $\vv A=\begin{pmatrix}1&1\\2&2\end{pmatrix}$. 求全部满足 $\vv{AX}=\vv O$ 的二阶矩阵 $\vv X$，它们构成的空间是几维的？
			\item\label{ex:sol:A12} 设 $\vv A=\begin{pmatrix}1&2&-2\\4&t&3\\3&-1&1\end{pmatrix}$，且存在三阶非零矩阵 $\vv B$ 使 $\vv{AB}=\vv O$. 求 $t$，并求 $\rk\vv B$ 的可能值.
			\item\label{ex:sol:A13} 求方程 $x_1+x_2+x_3=1$ 的通解，并写出它的三个解，使得通解中的每个解都是这三个解的仿射组合.
			\item\label{ex:sol:A14} 把空间直线 $\{x+y+z=1,\ x-y=0\}$ 写成参数式；反过来，把直线 $\vv x=(1,2,0)^\trans+t(1,0,1)^\trans$ 写成两个平面的交线.
			\item\label{ex:sol:A15} 判断空间直线 $\vv x=(1,0,0)^\trans+s(1,1,0)^\trans$ 与 $\vv x=(0,1,1)^\trans+t(0,1,-1)^\trans$ 是否相交.（列出关于 $s,t$ 的方程组.）
			\item\label{ex:sol:A16} 物流网络：三个仓库 $A,B,C$ 之间有线路 $A\to B$（$x_1$），$B\to C$（$x_2$），$A\to C$（$x_3$）. $A$ 每天净发出 $30$ 吨，$C$ 每天净收到 $30$ 吨，$B$ 不囤货. 求全部方案. 哪些量（如 $x_1+x_3$，$x_1-x_2$，$x_3$）能被方程组唯一确定？用可确定性判据说明.
			\item\label{ex:sol:A17} 求 $\vv A=\begin{pmatrix}2&1&0\\4&3&1\\2&3&3\end{pmatrix}$ 的 LU 分解.
			\item\label{ex:sol:A18} 设 $\vv A$ 是 $m\times n$ 矩阵，$\rk\vv A=m<n$. 证明 $\vv{Ax}=\vv b$ 对每个 $\vv b$ 都有无穷多解.
			\item\label{ex:sol:A19} 求经过点 $(1,1,1)$，$(2,0,1)$，$(0,2,3)$ 的平面方程（用齐次方程组求法向量）.
			\item\label{ex:sol:A20} 设 $\vv A=\begin{pmatrix}1&1\\1&1.0001\end{pmatrix}$. 分别求 $\vv{Ax}=(2,2)^\trans$ 与 $\vv{Ax}=(2,2.0001)^\trans$ 的解，计算右端的相对变化与解的相对变化之比.
		\end{exA}
		
		\Bgroup
		\begin{exB}
			\item\label{ex:sol:B1} 求两个齐次方程组 $\{x_1+x_2=0,\ x_2-x_4=0\}$ 与 $\{x_1-x_2+x_3=0,\ x_2-x_3+x_4=0\}$ 的公共非零解.
			\item\label{ex:sol:B2} 设 $\vv A$ 是实矩阵. 证明 $\vv{Ax}=\zero$ 与 $\vv A^\trans\vv{Ax}=\zero$ 同解，从而 $\rk(\vv A^\trans\vv A)=\rk\vv A$.（提示：$\vv x^\trans\vv A^\trans\vv{Ax}=\|\vv{Ax}\|^2$ 在第九章才有名字，但这里只需``实数的平方和为零则每个都为零''.）
			\item\label{ex:sol:B3} 证明 $\rk(\vv A+\vv B)\le\rk(\vv A\mid\vv B)\le\rk\vv A+\rk\vv B$.
			\item\label{ex:sol:B4} 设 $\vv A$ 为 $n$ 阶方阵，$\vv A^2=\vv I$. 证明 $\rk(\vv A+\vv I)+\rk(\vv A-\vv I)=n$.
			\item\label{ex:sol:B5} 证明 $\rk(\vv{AB})=\rk\vv B$ 当且仅当 $\ker\vv A\cap\imm\vv B=\{\zero\}$.
			\item\label{ex:sol:B6} 设 $\vv\eta^*$ 是 $\vv{Ax}=\vv b$（$\vv b\ne\zero$）的一个解，$\vv\xi_1,\ldots,\vv\xi_k$ 是对应齐次方程组的基础解系. 证明 $\vv\eta^*,\vv\eta^*+\vv\xi_1,\ldots,\vv\eta^*+\vv\xi_k$ 线性无关.
			\item\label{ex:sol:B7} 一张 $3\times3$ 的统计表只公开了三个行和与三个列和. 至少还要再公开几个单元格，才能确定全表？给出一种``公开哪几个格''的选法，并说明为什么某些同样个数的选法却不行.
			\item\label{ex:sol:B8} 对 $\vv A=\begin{pmatrix}1&2&1\\2&4&3\\3&6&4\end{pmatrix}$，求 $\vv b=(b_1,b_2,b_3)^\trans$ 应满足的条件，使 $\vv{Ax}=\vv b$ 有解.（提示：找出把各方程组合成 $0=\cdots$ 的系数.）
			\item\label{ex:sol:B9} 设 $\vv A$ 为 $n$ 阶方阵，$\rk\vv A=n-1$，且 $\vv A$ 的每行元素之和为零. 求 $\vv{Ax}=\zero$ 的通解.
			\item\label{ex:sol:B10} 设 $\vv A,\vv B$ 为 $n$ 阶方阵，$\vv{AB}=\vv O$，$\vv A+\vv B$ 可逆. 证明 $\rk\vv A+\rk\vv B=n$.
			\item\label{ex:sol:B11} 证明：两个 $m\times n$ 齐次方程组 $\vv{Ax}=\zero$ 与 $\vv{Bx}=\zero$ 同解，当且仅当 $\vv A$ 与 $\vv B$ 的行向量组等价.
			\item\label{ex:sol:B12} 设 $\vv A$ 为 $m\times n$ 矩阵. 证明：$\vv{Ax}=\vv b$ 对一切 $\vv b$ 有解，当且仅当存在 $n\times m$ 矩阵 $\vv C$ 使 $\vv{AC}=\vv I_m$；这时 $\vv C\vv b$ 就是一个解.
			\item\label{ex:sol:B13} 设 $\vv A$ 为 $n$ 阶方阵（$n\ge2$），$\rk\vv A=n-1$，$\vv\xi$ 是 $\vv{Ax}=\zero$ 的非零解. 证明 $\vv A^*$ 的每一列都是 $\vv\xi$ 的倍数.
			\item\label{ex:sol:B14} 用可确定性判据回答：方程组 $x_1+x_2+x_3=6$，$x_2+x_3+x_4=9$，$x_3+x_4+x_5=12$（例如三段相邻时间窗口内的观测总量）能确定 $x_1-x_4$ 吗？$x_1+x_2+x_3+x_4+x_5$ 呢？$x_1-x_4+x_5$ 呢？
			\item\label{ex:sol:B15} 设 $\vv A$ 为 $m\times n$ 矩阵，$\vv B$ 为 $n\times p$ 矩阵. 证明 $\dim\ker(\vv{AB})\le\dim\ker\vv A+\dim\ker\vv B$，并由此重新推出 Sylvester 不等式.
		\end{exB}
		
		\Cgroup
		\begin{exC}
			\item\label{ex:sol:C1} （Farkas 引理的二维形式）设 $\vv a_1,\vv a_2\in\mathbb{R}^2$ 线性无关，$\vv b\in\mathbb{R}^2$. 证明：下面两件事恰有一件成立：(1) 存在 $x_1,x_2\ge0$ 使 $x_1\vv a_1+x_2\vv a_2=\vv b$；(2) 存在 $\vv y$ 使 $\vv y^\trans\vv a_1\ge0$，$\vv y^\trans\vv a_2\ge0$，$\vv y^\trans\vv b<0$. 画图解释. 这是``无套利定价''与线性规划对偶理论的出发点.
			\item\label{ex:sol:C2} 设 $\vv A$ 是 $n$ 阶整数矩阵. 证明：若对任意整数向量 $\vv b$，$\vv{Ax}=\vv b$ 都有整数解，则 $\det\vv A=\pm1$.
			\item\label{ex:sol:C3} 证明：方阵 $\vv A$ 不用行交换就能做 LU 分解（$\vv L$ 单位下三角，$\vv U$ 上三角且可逆），当且仅当 $\vv A$ 的各阶顺序主子式都不为零；这时分解唯一.
			\item\label{ex:sol:C4} 设 $\vv A$ 是三对角矩阵（只有主对角线和相邻两条斜线上可能有非零元）. 证明：若它有 LU 分解，则 $\vv L,\vv U$ 都只有两条非零斜线；因此解三对角方程组只需与 $n$ 成正比的运算量. 这种方法称为``追赶法''.
			\item\label{ex:sol:C5} 用分块初等变换
			\[\begin{pmatrix}\vv I_n&\vv O\\ -\vv A&\vv I_m\end{pmatrix}\begin{pmatrix}\vv I_n&\vv B\\ \vv A&\vv O\end{pmatrix}\begin{pmatrix}\vv I_n&-\vv B\\ \vv O&\vv I_p\end{pmatrix}=\begin{pmatrix}\vv I_n&\vv O\\ \vv O&-\vv{AB}\end{pmatrix}\]
			给出 Sylvester 不等式 $\rk(\vv{AB})\ge\rk\vv A+\rk\vv B-n$ 的另一个证明.（提示：用第三章章末习题~\ref{ex:mat:B5} 的分块秩不等式.）
			\item\label{ex:sol:C6} 证明：$\mathbb{F}^n$ 中任意 $k$ 个线性流形（不必过原点）的交，要么是空集，要么仍是线性流形. 给出两个二维线性流形在 $\mathbb{R}^3$ 中交为空集的例子，并用秩的语言刻画这种情形.
			\item\label{ex:sol:C7} （列昂惕夫价格模型）设 $\vv A$ 为直接消耗系数矩阵，$\vv v$ 为各部门单位产品的增加值（工资、利润等）. 价格 $\vv p$ 满足``单位价格 $=$ 单位投入成本 $+$ 单位增加值''：$\vv p=\vv A^\trans\vv p+\vv v$. 说明它是第一章产量模型 $\vv x=\vv{Ax}+\vv d$ 的``转置''；证明在两个模型都可解时有 $\vv p^\trans\vv d=\vv v^\trans\vv x$，并解释其经济含义（国民收入的两种核算方法：支出法与收入法）.
			\item\label{ex:sol:C8} 设 $\vv A$ 为 $n$ 阶方阵. 证明存在正整数 $k$ 使 $\rk\vv A^k=\rk\vv A^{k+1}=\rk\vv A^{k+2}=\cdots$，且这样的最小的 $k$ 不超过 $n$.
		\end{exC}
	}
	
	% =====================================================================
	% 第七章　多项式
	% =====================================================================
	\newcommand{\ExPolynomial}{%
		\exsetheading
		\Agroup
		\begin{exA}
			\item\label{ex:poly:A1} 求商与余式：(1) $x^5-1$ 除以 $x^2+x+1$；(2) $2x^4-3x^3+4x^2-5x+6$ 除以 $x-2$（用综合除法）.
			\item\label{ex:poly:A2} 求 $x^{100}$ 分别除以 $x^2-1$，$x^2+1$，$(x-1)^2$ 所得的余式.
			\item\label{ex:poly:A3} 求 $a,b$，使 $x^2+x+1$ 整除 $x^4+ax^2+b$.（提示：设 $\omega$ 是 $x^2+x+1$ 的根，则 $\omega^3=1$.）
			\item\label{ex:poly:A4} 用辗转相除法求最大公因式：(1) $x^4+3x^3-x^2-4x-3$ 与 $3x^3+10x^2+2x-3$；(2) $x^3-2x^2+x$ 与 $x^2-1$.
			\item\label{ex:poly:A5} 证明：对任意多项式 $h$，$\gcd(f,g)=\gcd(f,\ g-hf)$. 这正是辗转相除法每一步的依据.
			\item\label{ex:poly:A6} 求 $u,v$ 使 $u(x)(x^2+1)+v(x)(x+1)=1$.
			\item\label{ex:poly:A7} 求 $2x^3-5x^2+x+2$ 的全部有理根，并在 $\mathbb{Q}$ 上把它分解为一次因式之积.
			\item\label{ex:poly:A8} 分别在 $\mathbb{Q}$，$\mathbb{R}$，$\mathbb{C}$ 上分解：(1) $x^4+x^2+1$；(2) $x^6-1$；(3) $x^3-2$.
			\item\label{ex:poly:A9} 判断 $1$ 是 $x^5-5x+4$ 的几重根.
			\item\label{ex:poly:A10} $a$ 取何值时，$x^3-3x+a$ 有重根？求出重根.
			\item\label{ex:poly:A11} 设 $\alpha_1,\alpha_2,\alpha_3$ 是 $x^3-2x^2+3x-4$ 的根. 不解方程，求 $\alpha_1+\alpha_2+\alpha_3$，$\alpha_1^2+\alpha_2^2+\alpha_3^2$，$\frac1{\alpha_1}+\frac1{\alpha_2}+\frac1{\alpha_3}$.
			\item\label{ex:poly:A12} 求次数最低的首一有理系数多项式，使 $1+\sqrt2$ 与 $\mathrm i$ 都是它的根.
			\item\label{ex:poly:A13} 求一个三次实系数首一多项式，以 $2$ 和 $1+\mathrm i$ 为根.
			\item\label{ex:poly:A14} 设 $f(x)=x^2-2x+1$，$\vv A=\begin{pmatrix}1&2\\0&1\end{pmatrix}$，计算 $f(\vv A)$；再设 $g(x)=x^3-x$，计算 $g(\vv A)$.
			\item\label{ex:poly:A15} 求过点 $(0,1)$，$(1,0)$，$(2,3)$，$(3,10)$ 的三次以下的插值多项式（用 Lagrange 公式）.
			\item\label{ex:poly:A16} 证明 $(x-1)\mid x^n-1$，并写出商式. 一般地，对任意 $a$，$(x-a)\mid x^n-a^n$.
			\item\label{ex:poly:A17} 把 $f(x)=x^4-2x^3+3$ 按 $x-1$ 的幂展开，即写成 $\sum c_k(x-1)^k$.（提示：反复用 $x-1$ 作综合除法，或用 $c_k=f^{(k)}(1)/k!$.）
			\item\label{ex:poly:A18} 多项式 $f$ 除以 $x-1$ 余 $3$，除以 $x-2$ 余 $5$. 求 $f$ 除以 $(x-1)(x-2)$ 的余式.
			\item\label{ex:poly:A19} 一个项目现在投入 $100$ 万元，一年后回收 $55$ 万元，两年后回收 $60.5$ 万元. 写出内部收益率满足的方程并精确求解.
			\item\label{ex:poly:A20} 证明：若 $f\mid g$ 且 $g\mid f$，则 $f=cg$（$c$ 为非零常数）.
			\item\label{ex:poly:A21} 设 $\gcd(f,g)=d$，$f=df_1$，$g=dg_1$. 证明 $\gcd(f_1,g_1)=1$.
		\end{exA}
		
		\Bgroup
		\begin{exB}
			\item\label{ex:poly:B1} 证明：若 $\gcd(f,g)=1$ 且 $\gcd(f,h)=1$，则 $\gcd(f,gh)=1$.
			\item\label{ex:poly:B2} 证明：$\gcd(x^m-1,\ x^n-1)=x^{\gcd(m,n)}-1$.（提示：模仿整数的辗转相除，先证 $x^m-1=x^{m-n}(x^n-1)+(x^{m-n}-1)$.）
			\item\label{ex:poly:B3} 设 $f=(x-1)^3(x+2)^2$. 计算 $f/\gcd(f,f')$，说明它恰好是``去掉重复因式''的结果. 一般地证明：若 $f=c\,p_1^{k_1}\cdots p_m^{k_m}$ 是标准分解，则 $f/\gcd(f,f')=c\,p_1\cdots p_m$.
			\item\label{ex:poly:B4} 证明：若 $p$ 不可约且 $p\mid f^2$，则 $p\mid f$. 举例说明``不可约''不能去掉.
			\item\label{ex:poly:B5} 设 $\ell_1,\ldots,\ell_n$ 是关于互异节点 $x_1,\ldots,x_n$ 的 Lagrange 基函数. 证明 $\sum_i\ell_i(x)=1$，$\sum_ix_i\ell_i(x)=x$，一般地 $\sum_ix_i^k\ell_i(x)=x^k$（$0\le k\le n-1$）.
			\item\label{ex:poly:B6} 证明：若复系数多项式 $f$ 满足 $f(x+1)=f(x)$ 对一切 $x$ 成立，则 $f$ 是常数.
			\item\label{ex:poly:B7} 证明 $x^3+x+1$ 在 $\mathbb{Q}$ 上不可约. 为什么对三次多项式只需检验有无有理根？四次多项式可以这样检验吗？
			\item\label{ex:poly:B8} 求 $x^n$ 除以 $x^2-3x+2$ 的余式.
			\item\label{ex:poly:B9} （Hermite 插值）求次数最低的多项式 $f$，使 $f(1)=1$，$f'(1)=2$，$f(2)=4$.（提示：设 $f=1+2(x-1)+c(x-1)^2$.）
			\item\label{ex:poly:B10} 设 $f$ 是有理系数多项式，$a,b\in\mathbb{Q}$，$\sqrt2\notin\mathbb{Q}$. 证明：若 $a+b\sqrt2$ 是 $f$ 的根，则 $a-b\sqrt2$ 也是 $f$ 的根.（提示：用 $(x-a)^2-2b^2$ 除 $f$.）
			\item\label{ex:poly:B11} 设 $f$ 是整系数多项式，$f(0)$ 与 $f(1)$ 都是奇数. 证明 $f$ 没有整数根.
			\item\label{ex:poly:B12} 求最小多项式意义下的``零化''：分别求出次数最低的首一多项式 $g$，使 $g(\vv A)=\vv O$，其中 $\vv A=\begin{pmatrix}0&-1\\1&0\end{pmatrix}$，$\vv A=\begin{pmatrix}2&0\\0&2\end{pmatrix}$，$\vv A=\begin{pmatrix}2&1\\0&2\end{pmatrix}$.
			\item\label{ex:poly:B13} 证明：现金流 $-C_0,\ C_1,\ldots,C_n$（$C_0>0$，$C_i\ge0$ 且不全为零）的净现值 $N(v)=-C_0+C_1v+\cdots+C_nv^n$ 在 $v>0$ 时严格递增，从而 $N(v)=0$ 至多一个正根；当 $C_1+\cdots+C_n>C_0$ 时这个正根落在 $(0,1)$ 中，即 IRR 唯一且为正.
			\item\label{ex:poly:B14} 设 $f$ 是 $n$ 次多项式，在 $n+1$ 个整数点 $0,1,\ldots,n$ 上都取整数值. 证明 $f$ 在一切整数点上都取整数值.（提示：用基 $\binom x0,\binom x1,\ldots,\binom xn$，其中 $\binom xk=\frac{x(x-1)\cdots(x-k+1)}{k!}$.）
			\item\label{ex:poly:B15} 证明 Lagrange 插值的``唯一性''也可以由线性代数推出：求值映射 $\mathbb{F}[x]_n\to\mathbb{F}^n$，$f\mapsto(f(x_1),\ldots,f(x_n))^\trans$ 是单射，从而是同构.
			\item\label{ex:poly:B16} 设 $p$ 是 $\mathbb{F}$ 上的不可约多项式，$f\in\mathbb{F}[x]$. 证明：$\gcd(p,f)=1$ 或 $p\mid f$，二者必居其一.
		\end{exB}
		
		\Cgroup
		\begin{exC}
			\item\label{ex:poly:C1} （Eisenstein 判别法）设 $f=a_nx^n+\cdots+a_0$ 是整系数多项式，$p$ 是素数，$p\nmid a_n$，$p\mid a_i$（$i<n$），$p^2\nmid a_0$. 证明 $f$ 在 $\mathbb{Q}$ 上不可约（可以先承认 Gauss 引理：整系数多项式若能分解为两个低次有理系数多项式之积，则也能分解为两个低次整系数多项式之积）. 由此说明 $x^n-2$ 对一切 $n$ 都不可约，所以 $\mathbb{Q}[x]$ 中有任意高次的不可约多项式.
			\item\label{ex:poly:C2} 设 $p$ 是素数. 证明分圆多项式 $\Phi_p(x)=x^{p-1}+x^{p-2}+\cdots+1$ 在 $\mathbb{Q}$ 上不可约.（提示：对 $\Phi_p(x+1)$ 用 Eisenstein 判别法.）
			\item\label{ex:poly:C3} （Newton 恒等式）记 $s_k=\alpha_1^k+\alpha_2^k+\alpha_3^k$，$\sigma_1,\sigma_2,\sigma_3$ 为初等对称多项式. 证明 $s_1=\sigma_1$，$s_2=\sigma_1s_1-2\sigma_2$，$s_3=\sigma_1s_2-\sigma_2s_1+3\sigma_3$，并求 $x^3-2x^2+3x-4$ 的根的立方和.
			\item\label{ex:poly:C4} （三次方程的判别式）设 $x^3+px+q$ 的三个复根为 $\alpha_1,\alpha_2,\alpha_3$，$\Delta=\prod_{i<j}(\alpha_i-\alpha_j)^2$. 证明 $\Delta=-4p^3-27q^2$，并说明（$p,q$ 为实数时）：$\Delta>0$ 时三个不同实根，$\Delta<0$ 时一实二虚，$\Delta=0$ 时有重根.
			\item\label{ex:poly:C5} （Chebyshev 多项式）证明存在 $n$ 次多项式 $T_n$ 使 $\cos n\theta=T_n(\cos\theta)$，且 $T_{n+1}=2xT_n-T_{n-1}$. 求 $T_n$ 的全部根. 由此证明：$\cos20^\circ$ 是 $8x^3-6x-1$ 的根，并验证这个三次多项式在 $\mathbb{Q}$ 上不可约. （结合``尺规作图只能得到经若干次开平方得到的数''，这说明 $60^\circ$ 角不能用尺规三等分.）
			\item\label{ex:poly:C6} 证明：在 $\mathbb{R}$ 上处处非负的实系数多项式，都能写成两个实系数多项式的平方和.（提示：利用实根的重数必为偶数，以及 $|a+b\mathrm i|^2=a^2+b^2$ 与 $(a^2+b^2)(c^2+d^2)=(ac-bd)^2+(ad+bc)^2$.）
			\item\label{ex:poly:C7} 设 $f\in\mathbb{Q}[x]$ 不可约，$g\in\mathbb{Q}[x]$ 与 $f$ 有一个公共复根. 证明 $f\mid g$. 由此说明：一个复数若是某个有理系数多项式的根，则以它为根的次数最低的首一有理系数多项式（称为它的\textbf{极小多项式}）唯一，且不可约.
			\item\label{ex:poly:C8} （$\mathbb{F}_2$ 上的多项式）在系数按模 $2$ 运算的多项式中，证明 $x^2+x+1$ 是唯一的二次不可约多项式，并找出全部三次不可约多项式. 这类多项式是循环冗余校验（CRC）与密码学的基础.
			\item\label{ex:poly:C9} 设 $f$ 是实系数多项式，它的系数序列的``变号次数''为 $V$（忽略零系数）. Descartes 符号法则说：$f$ 的正根个数（按重数计）不超过 $V$，且与 $V$ 的差是偶数. 用它重新解释第七章``内部收益率''建模实验中两个项目的结论.
		\end{exC}
	}
	
	% =====================================================================
	% 第八章　特征值与特征向量
	% =====================================================================
	\newcommand{\ExEigen}{%
		\exsetheading
		\Agroup
		\begin{exA}
			\item\label{ex:eig:A1} 求下列矩阵的特征值与特征向量（在 $\mathbb{C}$ 上考虑）：
			\[\begin{pmatrix}3&1\\0&2\end{pmatrix},\quad\begin{pmatrix}1&2\\2&1\end{pmatrix},\quad\begin{pmatrix}0&1\\-1&0\end{pmatrix},\quad\begin{pmatrix}4&-2\\1&1\end{pmatrix},\]\[\begin{pmatrix}2&0&0\\0&1&1\\0&1&1\end{pmatrix},\quad\begin{pmatrix}1&2&2\\2&1&2\\2&2&1\end{pmatrix}.\]
			\item\label{ex:eig:A2} 判断 A1 中各矩阵能否在 $\mathbb{R}$ 上对角化；能的，求可逆矩阵 $\vv P$ 使 $\vv P^{-1}\vv{AP}$ 为对角矩阵.
			\item\label{ex:eig:A3} 判断下列矩阵能否对角化：$\begin{pmatrix}1&1\\0&2\end{pmatrix}$，$\begin{pmatrix}2&1\\0&2\end{pmatrix}$，$\begin{pmatrix}2&0&0\\0&2&0\\1&0&2\end{pmatrix}$，$\begin{pmatrix}1&1&0\\0&2&0\\0&0&1\end{pmatrix}$.
			\item\label{ex:eig:A4} 设 $3$ 阶方阵 $\vv A$ 的特征值为 $1,2,3$. 求 $\det\vv A$，$\tr\vv A$，$\vv A^{-1}$ 的特征值，$\vv A^2-2\vv A+\vv I$ 的特征值，$\det(\vv A^*)$，$\det(\vv A^2+\vv I)$.
			\item\label{ex:eig:A5} 已知二阶方阵的迹为 $5$，行列式为 $6$，求它的特征值. 若迹为 $2$、行列式为 $5$ 呢？
			\item\label{ex:eig:A6} 证明：$\vv A$ 与 $\vv A^\trans$ 有相同的特征多项式. 它们的特征向量是否相同？举例说明.
			\item\label{ex:eig:A7} 证明：三角矩阵的特征值恰为其对角元.
			\item\label{ex:eig:A8} 证明：若 $\vv A^k=\vv O$，则 $\vv A$ 的特征值全为 $0$；若 $\vv A^2=\vv A$，则特征值只能是 $0$ 或 $1$；若 $\vv A^2=\vv I$，则特征值只能是 $\pm1$.
			\item\label{ex:eig:A9} 设 $\vv A=\begin{pmatrix}4&-2\\1&1\end{pmatrix}$. 利用对角化计算 $\vv A^n$.
			\item\label{ex:eig:A10} 设品牌转移矩阵为 $\vv P=\begin{pmatrix}0.9&0.2\\0.1&0.8\end{pmatrix}$. 求稳态分布，并利用对角化求 $\lim_{n\to\infty}\vv P^n$. 解释极限矩阵两列相同的含义.
			\item\label{ex:eig:A11} 已知 $\vv\xi=(1,1,-1)^\trans$ 是 $\vv A=\begin{pmatrix}2&-1&2\\5&a&3\\-1&b&-2\end{pmatrix}$ 的特征向量. 求 $a,b$ 与对应的特征值.
			\item\label{ex:eig:A12} 设 $\vv A=\begin{pmatrix}0&0&1\\1&1&x\\1&0&0\end{pmatrix}$. $x$ 取何值时 $\vv A$ 可对角化？
			\item\label{ex:eig:A13} 对二阶方阵直接验证 Cayley–Hamilton 定理：$\vv A^2-(\tr\vv A)\vv A+(\det\vv A)\vv I=\vv O$.
			\item\label{ex:eig:A14} 求下列矩阵的最小多项式：$\diag(1,1,2)$，$\begin{pmatrix}1&1&0\\0&1&0\\0&0&2\end{pmatrix}$，$\begin{pmatrix}0&1&0\\0&0&1\\0&0&0\end{pmatrix}$，$\begin{pmatrix}1&1&1\\1&1&1\\1&1&1\end{pmatrix}$.
			\item\label{ex:eig:A15} 判断 $\begin{pmatrix}1&1\\0&1\end{pmatrix}$ 与 $\begin{pmatrix}1&0\\1&1\end{pmatrix}$ 是否相似；若相似，求出 $\vv P$.
			\item\label{ex:eig:A16} 设 $\vv A$ 是上三角矩阵. 证明 $\spn(\vv e_1)$，$\spn(\vv e_1,\vv e_2)$，……都是 $\vv A$ 的不变子空间.
			\item\label{ex:eig:A17} 解差分方程 $a_{k+1}=a_k+2a_{k-1}$，$a_0=0$，$a_1=1$.
			\item\label{ex:eig:A18} （种群模型）某种群分为幼年与成年两组，第 $k$ 年的数量记为 $\vv x_k=(y_k,z_k)^\trans$. 每只成年个体每年平均产下 $1.1$ 只幼体，幼体有 $60\%$ 存活并长为成年，成年个体有一半活到下一年，即
			$y_{k+1}=1.1z_k$，$z_{k+1}=0.6y_k+0.5z_k$. 求矩阵的特征值与特征向量，说明种群的长期增长率和幼体与成年个体的稳定比例.
			\item\label{ex:eig:A19} 证明：实二阶方阵 $\vv A$ 没有实特征值，当且仅当 $(\tr\vv A)^2<4\det\vv A$.
			\item\label{ex:eig:A20} 求平面上向直线 $y=2x$ 的正投影、关于该直线的反射这两个变换的特征值与特征向量（不必写出矩阵，从几何上直接回答），再写出矩阵加以验证.
		\end{exA}
		
		\Bgroup
		\begin{exB}
			\item\label{ex:eig:B1} 设 $\vv u,\vv v\in\mathbb{R}^n$ 非零，$\vv A=\vv u\vv v^\trans$. 求 $\vv A$ 的全部特征值及其重数，并证明：$\vv A$ 可对角化当且仅当 $\vv v^\trans\vv u\ne0$.
			\item\label{ex:eig:B2} 设 $\vv A$ 有 $n$ 个互异的特征值，$\vv{AB}=\vv{BA}$. 证明 $\vv A$ 的特征向量都是 $\vv B$ 的特征向量，从而 $\vv B$ 也可对角化；进一步证明 $\vv B$ 是 $\vv A$ 的多项式.
			\item\label{ex:eig:B3} 证明：若 $\vv A$ 可对角化，则 $\vv A^k$ 与 $f(\vv A)$ 都可对角化. 逆命题成立吗？
			\item\label{ex:eig:B4} 设 $\vv A$ 可对角化，$W$ 是 $\vv A$ 的不变子空间. 证明 $\vv A|_W$ 也可对角化.（提示：最小多项式.）
			\item\label{ex:eig:B5} 证明：满足 $\vv A^2=-\vv I$ 的二阶实方阵都与 $\begin{pmatrix}0&-1\\1&0\end{pmatrix}$ 相似.
			\item\label{ex:eig:B6} （Gershgorin 圆盘定理）证明：方阵 $\vv A$ 的每个特征值都落在某个圆盘 $\{z\in\mathbb{C}:|z-a_{ii}|\le\sum_{j\ne i}|a_{ij}|\}$ 中. 由此证明：严格对角占优的矩阵（$|a_{ii}|>\sum_{j\ne i}|a_{ij}|$）可逆.
			\item\label{ex:eig:B7} 设 $\vv A$ 的每行元素之和都是 $s$. 证明 $s$ 是 $\vv A$ 的特征值. 若 $\vv A$ 可逆，求 $\vv A^{-1}$ 每行元素之和.
			\item\label{ex:eig:B8} 设 $\vv A=\begin{pmatrix}a&b&c\\c&a&b\\b&c&a\end{pmatrix}$（循环矩阵），$\omega=\mathrm e^{2\pi\mathrm i/3}$. 证明 $(1,\omega^k,\omega^{2k})^\trans$（$k=0,1,2$）都是 $\vv A$ 的特征向量，并求对应的特征值. 由此重新得到第五章章末习题~\ref{ex:det:B12} 的行列式公式.
			\item\label{ex:eig:B9} 解差分方程 $a_{k+1}=4a_k-4a_{k-1}$，$a_0=1$，$a_1=4$. 解释通项中为什么出现 $k\cdot2^k$.
			\item\label{ex:eig:B10} 证明 $\tr(\vv A^k)=\lambda_1^k+\cdots+\lambda_n^k$. 若三阶方阵满足 $\tr\vv A=0$，$\tr\vv A^2=6$，$\det\vv A=0$，求它的特征值.
			\item\label{ex:eig:B11} 证明：若 $\vv A^3=\vv A$，则 $\vv A$ 可对角化，且特征值在 $\{0,1,-1\}$ 中.
			\item\label{ex:eig:B12} $a,b,c$ 满足什么条件时，$\begin{pmatrix}2&0&0\\a&2&0\\b&c&1\end{pmatrix}$ 可对角化？
			\item\label{ex:eig:B13} 三个网页之间的链接为：网页 $1$ 链接到 $2,3$，网页 $2$ 链接到 $3$，网页 $3$ 链接到 $1$.（与判例不同，网页之间可以``互相引用''，形成环路.） 取 $d=1$ 写出转移矩阵，手算它的稳态分布.
			\item\label{ex:eig:B14} （矩阵指数）对可对角化矩阵 $\vv A=\vv P\vv D\vv P^{-1}$，定义 $\mathrm e^{t\vv A}=\vv P\diag(\mathrm e^{t\lambda_1},\ldots,\mathrm e^{t\lambda_n})\vv P^{-1}$. 对 $\vv A=\begin{pmatrix}0&1\\1&0\end{pmatrix}$ 计算 $\mathrm e^{t\vv A}$，并验证 $\vv x(t)=\mathrm e^{t\vv A}\vv x_0$ 满足微分方程组 $x_1'=x_2$，$x_2'=x_1$.
			\item\label{ex:eig:B15} 三状态马尔可夫链的转移矩阵为 $\begin{pmatrix}0.5&0.25&0.25\\0.25&0.5&0.25\\0.25&0.25&0.5\end{pmatrix}$. 求稳态分布与全部特征值，说明偏离稳态的部分每一步缩小为原来的多少.
			\item\label{ex:eig:B16} 证明：若 $\lambda$ 是 $\vv A$ 的 $k$ 重特征值（代数重数），则 $\rk(\lambda\vv I-\vv A)\ge n-k$.
			\item\label{ex:eig:B17} 对第三章``判例引证网络''中的四份判决，分别取阻尼系数 $d=0.5,\ 0.85,\ 0.95,\ 1$，按第八章建模实验的方法求 PageRank 得分（可用计算机），列表比较各判决的得分与排序. $d$ 越接近 $1$，得分越``集中''在哪些判决上？为什么？
		\end{exB}
		
		\Cgroup
		\begin{exC}
			\item\label{ex:eig:C1} 设 $\vv A$ 为 $m\times n$ 矩阵，$\vv B$ 为 $n\times m$ 矩阵. 证明 $\lambda^n\det(\lambda\vv I_m-\vv{AB})=\lambda^m\det(\lambda\vv I_n-\vv{BA})$. 由此说明 $\vv{AB}$ 与 $\vv{BA}$ 的非零特征值（连同重数）相同.（提示：计算 $\begin{pmatrix}\vv I_m&-\vv A\\ \vv O&\lambda\vv I_n\end{pmatrix}\begin{pmatrix}\lambda\vv I_m&\vv A\\ \vv B&\vv I_n\end{pmatrix}$ 等两种乘积的行列式.）
			\item\label{ex:eig:C2} （Schur 三角化）证明：每个复方阵都相似于一个上三角矩阵. 由此给出 Cayley–Hamilton 定理的另一个证明：对上三角矩阵 $\vv T$，证明 $(\vv T-t_{11}\vv I)\cdots(\vv T-t_{nn}\vv I)=\vv O$.
			\item\label{ex:eig:C3} 求下列矩阵的 Jordan 标准形：
			\[\begin{pmatrix}2&1&0\\0&2&0\\0&0&2\end{pmatrix},\quad\begin{pmatrix}1&1&1\\0&1&1\\0&0&1\end{pmatrix},\quad\begin{pmatrix}0&1&0&0\\0&0&0&0\\0&0&0&1\\0&0&0&0\end{pmatrix}.\]
			特征多项式为 $(\lambda-2)^4$ 的 $4$ 阶复矩阵，在相似的意义下共有几类？
			\item\label{ex:eig:C4} 证明：若一组可对角化的矩阵两两可交换，则它们可以\textbf{同时}对角化（存在同一个 $\vv P$ 使它们都化为对角矩阵）.
			\item\label{ex:eig:C5} 设 $\vv A$ 可对角化. 证明 $\vv A^k\to\vv O$（按元素）当且仅当 $\vv A$ 的所有特征值的模都小于 $1$. 用它说明：投入产出模型中，若直接消耗系数矩阵 $\vv A$ 可对角化且特征值的模都小于 $1$，则 $(\vv I-\vv A)^{-1}=\vv I+\vv A+\vv A^2+\cdots$.
			\item\label{ex:eig:C6} （友矩阵）证明 $\begin{pmatrix}0&0&\cdots&0&-a_0\\1&0&\cdots&0&-a_1\\0&1&\cdots&0&-a_2\\ \vdots&&\ddots&&\vdots\\0&0&\cdots&1&-a_{n-1}\end{pmatrix}$ 的特征多项式为 $\lambda^n+a_{n-1}\lambda^{n-1}+\cdots+a_0$. 由此说明：求多项式的根等价于求矩阵的特征值（数值软件正是这样求多项式的根的）.
			\item\label{ex:eig:C7} （赌徒破产）一位赌徒有 $k$ 元（$0\le k\le N$），每局以概率 $\frac12$ 赢或输 $1$ 元，输光或达到 $N$ 元时停止. 设 $p_k$ 为最终达到 $N$ 元的概率，证明 $p_k=\frac12(p_{k+1}+p_{k-1})$，并用特征根方法解出 $p_k=k/N$. 赢的概率改为 $p\ne\frac12$ 时呢？
			\item\label{ex:eig:C8} 证明：$n$ 阶三对角矩阵 $\vv T=\begin{pmatrix}2&-1&&\\-1&2&\ddots&\\&\ddots&\ddots&-1\\&&-1&2\end{pmatrix}$ 的特征值为 $2-2\cos\frac{k\pi}{n+1}$（$k=1,\ldots,n$），对应特征向量的第 $j$ 个分量为 $\sin\frac{jk\pi}{n+1}$. 对 $n=2,3$ 直接验证.
			\item\label{ex:eig:C9} （Perron–Frobenius 的二阶情形）设 $\vv A$ 是元素全为正的二阶矩阵. 证明 $\vv A$ 有两个互异实特征值，较大者 $r$ 为正且大于另一个的绝对值，属于 $r$ 的特征向量可取为各分量都为正. 由此解释：为什么 PageRank 引入阻尼系数后，``全部元素为正''能保证排序唯一.
			\item\label{ex:eig:C10} 证明：满足 $\vv A^2=-\vv I$ 的 $n$ 阶实方阵存在当且仅当 $n$ 为偶数，且这时 $\vv A$ 相似于 $\diag(\vv J,\ldots,\vv J)$，$\vv J=\begin{pmatrix}0&-1\\1&0\end{pmatrix}$.
		\end{exC}
	}
	
	% =====================================================================
	% 第九章　内积空间
	% =====================================================================
	\newcommand{\ExInnerProduct}{%
		\exsetheading
		\Agroup
		\begin{exA}
			\item\label{ex:ip:A1} 在 $\mathbb{R}^4$ 中，设 $\vv\alpha=(1,2,2,0)^\trans$，$\vv\beta=(2,0,1,2)^\trans$. 求 $\|\vv\alpha\|$，$\|\vv\beta\|$，$\langle\vv\alpha,\vv\beta\rangle$，它们的夹角，以及 $\|\vv\alpha-\vv\beta\|$.
			\item\label{ex:ip:A2} 判断下列二元函数是否为 $\mathbb{R}^2$ 上的内积：
			(1) $\langle\vv x,\vv y\rangle=x_1y_1+2x_2y_2$；(2) $\langle\vv x,\vv y\rangle=x_1y_2+x_2y_1$；(3) $\langle\vv x,\vv y\rangle=2x_1y_1+x_1y_2+x_2y_1+x_2y_2$；(4) $\langle\vv x,\vv y\rangle=x_1y_1-x_2y_2$.
			\item\label{ex:ip:A3} 用 Gram–Schmidt 方法把下列向量组正交化并单位化：
			(1) $(1,1,1)^\trans$，$(1,2,0)^\trans$，$(0,1,2)^\trans$；(2) $(1,1,0,0)^\trans$，$(1,0,1,0)^\trans$，$(-1,0,0,1)^\trans$.
			\item\label{ex:ip:A4} 求与 $(1,1,1)^\trans$ 和 $(1,0,-1)^\trans$ 都正交的单位向量.
			\item\label{ex:ip:A5} 已知 $\vv e_1=\frac13(1,2,2)^\trans$，$\vv e_2=\frac13(2,1,-2)^\trans$，$\vv e_3=\frac13(2,-2,1)^\trans$. 验证它们构成 $\mathbb{R}^3$ 的标准正交基，并只用内积求 $(3,6,0)^\trans$ 在这组基下的坐标.
			\item\label{ex:ip:A6} 求 $\vv\alpha=(1,2,3,4)^\trans$ 在 $W=\spn\big((1,1,0,0)^\trans,(0,0,1,1)^\trans\big)$ 上的正交投影，以及 $\vv\alpha$ 到 $W$ 的距离.
			\item\label{ex:ip:A7} 用最小二乘法求过数据点 $(0,1)$，$(1,2)$，$(2,2)$，$(3,4)$，$(4,6)$ 的最佳直线 $y=a+bx$，并计算残差平方和.
			\item\label{ex:ip:A8} 求 $a,b,c$ 使 $\begin{pmatrix}\frac1{\sqrt2}&a&\frac1{\sqrt6}\\0&b&-\frac2{\sqrt6}\\-\frac1{\sqrt2}&c&\frac1{\sqrt6}\end{pmatrix}$ 为正交矩阵.
			\item\label{ex:ip:A9} 证明：若 $\vv\alpha_1,\ldots,\vv\alpha_k$ 两两正交，则 $\|\vv\alpha_1+\cdots+\vv\alpha_k\|^2=\|\vv\alpha_1\|^2+\cdots+\|\vv\alpha_k\|^2$.
			\item\label{ex:ip:A10} 证明：$\big|\|\vv\alpha\|-\|\vv\beta\|\big|\le\|\vv\alpha-\vv\beta\|$.
			\item\label{ex:ip:A11} 用 Cauchy–Schwarz 不等式证明：(1) $(a_1+\cdots+a_n)^2\le n(a_1^2+\cdots+a_n^2)$；(2) 若 $a_i>0$，则 $(a_1+\cdots+a_n)\big(\frac1{a_1}+\cdots+\frac1{a_n}\big)\ge n^2$.
			\item\label{ex:ip:A12} 在 $C[0,1]$ 中取 $\langle f,g\rangle=\int_0^1f(x)g(x)\,\mathrm dx$. 求 $\|1\|$，$\|x\|$，$1$ 与 $x$ 的夹角，并把 $1,x$ 正交化.
			\item\label{ex:ip:A13} 设 $\vv A=\begin{pmatrix}1&1&0&1\\0&1&1&1\end{pmatrix}$. 求 $\ker\vv A$ 与行空间 $\imm\vv A^\trans$ 的标准正交基，并验证它们互为正交补.
			\item\label{ex:ip:A14} 求 $\vv A=\begin{pmatrix}1&1\\1&0\\0&1\end{pmatrix}$ 的 QR 分解.
			\item\label{ex:ip:A15} 设 $\vv u=\frac13(1,2,2)^\trans$. 写出 Householder 矩阵 $\vv H=\vv I-2\vv u\vv u^\trans$，验证它是对称的正交矩阵，并求 $\vv H\vv u$ 与 $\vv H(2,-1,0)^\trans$.
			\item\label{ex:ip:A16} 判断下列正交矩阵是旋转还是反射；若是旋转求旋转角，若是反射求反射轴：$\frac15\begin{pmatrix}3&-4\\4&3\end{pmatrix}$，$\frac15\begin{pmatrix}3&4\\4&-3\end{pmatrix}$.
			\item\label{ex:ip:A17} 对``过原点的直线''模型 $y=bx$，推导最小二乘解 $\hat b=\sum x_iy_i\big/\sum x_i^2$.
			\item\label{ex:ip:A18} 在 $\mathbb{R}^3$ 中，求点 $(1,2,3)$ 到直线 $\vv x=t(1,1,1)^\trans$ 的距离.
			\item\label{ex:ip:A19} 证明：两组标准正交基之间的过渡矩阵是正交矩阵.
			\item\label{ex:ip:A20} 三个法条的词频向量（关键词：违约、赔偿、合同、离婚）分别为 $(2,1,3,0)^\trans$，$(0,1,0,4)^\trans$，$(1,1,2,0)^\trans$，查询为 $(1,1,1,0)^\trans$. 分别用内积和余弦相似度给三个法条排序，结果有何不同？
			\item\label{ex:ip:A21} 设 $\vv u\in\mathbb{R}^n$ 是单位向量. 证明 $\vv P=\vv u\vv u^\trans$ 满足 $\vv P^\trans=\vv P$，$\vv P^2=\vv P$，且对任意 $\vv x$，$\vv P\vv x$ 是 $\vv x$ 在 $\spn(\vv u)$ 上的正交投影. 若 $\vv u$ 不是单位向量，投影矩阵应当怎样修改？
		\end{exA}
		
		\Bgroup
		\begin{exB}
			\item\label{ex:ip:B1} （极化恒等式）证明实内积空间中 $\langle\vv\alpha,\vv\beta\rangle=\frac14\big(\|\vv\alpha+\vv\beta\|^2-\|\vv\alpha-\vv\beta\|^2\big)$. 再证明：$\mathbb{R}^2$ 上的``范数'' $\|\vv x\|_1=|x_1|+|x_2|$ 不满足平行四边形法则，因而不能来自任何内积.
			\item\label{ex:ip:B2} （Bessel 不等式）设 $\vv e_1,\ldots,\vv e_k$ 是标准正交组. 证明 $\sum_{i=1}^k\langle\vv\alpha,\vv e_i\rangle^2\le\|\vv\alpha\|^2$，等号成立当且仅当 $\vv\alpha\in\spn(\vv e_1,\ldots,\vv e_k)$.
			\item\label{ex:ip:B3} 设实矩阵 $\vv P$ 满足 $\vv P^\trans=\vv P$，$\vv P^2=\vv P$. 证明 $\vv P$ 是向 $\imm\vv P$ 的正交投影的矩阵. 若只有 $\vv P^2=\vv P$，它是什么样的投影？举例说明.
			\item\label{ex:ip:B4} （加权最小二乘）若各观测值的可靠程度不同，给第 $i$ 个残差权重 $w_i>0$，最小化 $\sum w_i(y_i-a-bx_i)^2$. 证明正规方程变为 $\vv A^\trans\vv W\vv A\hat{\vv\beta}=\vv A^\trans\vv W\vv y$，$\vv W=\diag(w_1,\ldots,w_n)$. 它对应于哪个内积下的正交投影？
			\item\label{ex:ip:B5} 某商品的销售额 $y$ 随时间 $t$ 近似按指数增长 $y=c\,\mathrm e^{kt}$. 已知 $t=0,1,2,3$ 时 $y$ 约为 $2.0,2.7,3.7,5.0$. 两边取对数化为线性模型，用最小二乘估计 $c,k$. 这样做最小化的是哪个量？与直接最小化 $\sum(y_i-c\mathrm e^{kt_i})^2$ 有何不同？
			\item\label{ex:ip:B6} （Gram 行列式）设 $\vv\alpha_1,\ldots,\vv\alpha_k\in\mathbb{R}^n$，$\vv G=(\langle\vv\alpha_i,\vv\alpha_j\rangle)$. 证明 $\det\vv G\ge0$，等号成立当且仅当它们线性相关. 在 $k=2$ 时说明 $\sqrt{\det\vv G}$ 是两向量张成的平行四边形的面积.
			\item\label{ex:ip:B7} 证明：$\mathbb{R}^n$ 中点 $\vv x_0$ 到超平面 $\vv a^\trans\vv x=b$（$\vv a\ne\zero$）的距离为 $|\vv a^\trans\vv x_0-b|/\|\vv a\|$.
			\item\label{ex:ip:B8} 证明：实正交矩阵的复特征值的模都等于 $1$.
			\item\label{ex:ip:B9} （Fourier 逼近）在 $C[-\pi,\pi]$ 中取 $\langle f,g\rangle=\int_{-\pi}^\pi fg\,\mathrm dx$. 证明 $1,\cos x,\sin x,\cos2x,\sin2x$ 两两正交. 求 $f(x)=x$ 在它们张成的子空间上的正交投影.
			\item\label{ex:ip:B10} 设 $\vv A$ 为实矩阵. 证明：在 $\vv{Ax}=\vv y$ 的全部最小二乘解中，范数最小的那一个唯一，并且它属于行空间 $\imm\vv A^\trans$.
			\item\label{ex:ip:B11} 设 $\vv x,\vv y\in\mathbb{R}^n$，$\|\vv x\|=\|\vv y\|$，$\vv x\ne\vv y$. 令 $\vv u=(\vv x-\vv y)/\|\vv x-\vv y\|$. 证明 Householder 矩阵 $\vv I-2\vv u\vv u^\trans$ 把 $\vv x$ 变为 $\vv y$. 这是用 Householder 反射做 QR 分解的基础.
			\item\label{ex:ip:B12} 设两个变量的样本标准差分别为 $s_x,s_y$，相关系数为 $r$. 证明回归直线 $y=a+bx$ 的斜率为 $b=r\,s_y/s_x$，并且直线经过点 $(\bar x,\bar y)$.
			\item\label{ex:ip:B13} 证明：保持\textbf{夹角}的可逆线性变换（对一切非零向量，$\mathscr T\vv\alpha$ 与 $\mathscr T\vv\beta$ 的夹角等于 $\vv\alpha$ 与 $\vv\beta$ 的夹角）必是正交变换的非零倍数.
			\item\label{ex:ip:B14} 在回归中加入一个``虚拟变量''：设有两组数据，第一组 $y$ 值为 $3,5,4$，第二组为 $7,8,9$. 以 $\vv A=(\vv 1,\vv d)$ 为设计矩阵（$\vv d$ 在第二组为 $1$、第一组为 $0$），求最小二乘解，并解释两个系数的含义.
			\item\label{ex:ip:B15} 证明：在 $\mathbb{R}^2$ 中，任意两个旋转可交换，而一个旋转与一个反射一般不可交换；两个反射的乘积是旋转.
		\end{exB}
		
		\Cgroup
		\begin{exC}
			\item\label{ex:ip:C1} （Hadamard 不等式）用 QR 分解证明：对 $n$ 阶实方阵 $\vv A=(\vv a_1,\ldots,\vv a_n)$，$|\det\vv A|\le\|\vv a_1\|\cdots\|\vv a_n\|$，并说明它的几何意义（平行多面体的体积不超过棱长之积）.
			\item\label{ex:ip:C2} （Gram–Schmidt 的数值不稳定）设 $\varepsilon$ 很小，对 $\vv a_1=(1,\varepsilon,0,0)^\trans$，$\vv a_2=(1,0,\varepsilon,0)^\trans$，$\vv a_3=(1,0,0,\varepsilon)^\trans$ 作 Gram–Schmidt，在计算中把 $1+\varepsilon^2$ 舍入为 $1$. 证明所得的 $\vv e_2,\vv e_3$ 之间的夹角是 $60^\circ$ 而不是 $90^\circ$. 再用``修正的 Gram–Schmidt''（每得到一个 $\vv e_k$，立即从\textbf{剩余的}所有向量中减去它的分量）重做，比较结果.
			\item\label{ex:ip:C3} （Legendre 多项式）令 $P_n(x)=\frac1{2^nn!}\frac{\mathrm d^n}{\mathrm dx^n}(x^2-1)^n$. 证明 $P_0,P_1,P_2,\ldots$ 在 $\langle f,g\rangle=\int_{-1}^1fg\,\mathrm dx$ 下两两正交，并求 $\|P_n\|$.
			\item\label{ex:ip:C4} （Gauss–Markov 定理的线性代数部分）设 $\vv A$ 列满秩. 若 $\vv C\vv y$ 是参数 $\vv\beta$ 的``线性无偏估计''，即 $\vv{CA}=\vv I$. 令 $\vv C_0=(\vv A^\trans\vv A)^{-1}\vv A^\trans$（最小二乘）. 证明 $\vv C\vv C^\trans-\vv C_0\vv C_0^\trans$ 是半正定矩阵（第十一章）. 学过概率论后可以知道，这说明最小二乘估计的方差在一切线性无偏估计中最小.
			\item\label{ex:ip:C5} （复内积空间）在 $\mathbb{C}^n$ 上定义 $\langle\vv x,\vv y\rangle=\sum x_i\bar y_i$. 证明它满足共轭对称性、对第一个变量线性、正定性；证明 Cauchy–Schwarz 不等式仍成立. 满足 $\bar{\vv U}^\trans\vv U=\vv I$ 的复矩阵称为\textbf{酉矩阵}，证明它保持这个内积.
			\item\label{ex:ip:C6} 设 $f:\mathbb{R}^n\to\mathbb{R}^n$ 保持距离（$\|f(\vv x)-f(\vv y)\|=\|\vv x-\vv y\|$）且 $f(\zero)=\zero$. 证明 $f$ 保持内积，进而证明 $f$ 是线性的正交变换. 由此说明：保持距离的变换都可以分解为正交变换与平移的复合.
			\item\label{ex:ip:C7} （伪逆）设 $\vv A$ 列满秩，$\vv A^+=(\vv A^\trans\vv A)^{-1}\vv A^\trans$. 验证 Penrose 四条件：$\vv{AA}^+\vv A=\vv A$，$\vv A^+\vv{AA}^+=\vv A^+$，$(\vv{AA}^+)^\trans=\vv{AA}^+$，$(\vv A^+\vv A)^\trans=\vv A^+\vv A$. 说明 $\vv{AA}^+$ 是向哪个子空间的正交投影.
			\item\label{ex:ip:C8} 求空间中两条异面直线 $\vv x=\vv p+s\vv u$ 与 $\vv x=\vv q+t\vv v$ 之间的最短距离：把问题写成关于 $(s,t)$ 的最小二乘问题，并对 $\vv p=\zero$，$\vv u=(1,0,0)^\trans$，$\vv q=(0,1,1)^\trans$，$\vv v=(0,1,-1)^\trans$ 计算.
		\end{exC}
	}
	
	% =====================================================================
	% 第十章　矩阵分解
	% =====================================================================
	\newcommand{\ExDecomposition}{%
		\exsetheading
		\Agroup
		\begin{exA}
			\item\label{ex:dec:A1} 求正交矩阵 $\vv Q$ 使 $\vv Q^\trans\vv{AQ}$ 为对角矩阵：
			\[\begin{pmatrix}1&2\\2&1\end{pmatrix},\qquad\begin{pmatrix}3&-1\\-1&3\end{pmatrix},\]\[\begin{pmatrix}2&2&0\\2&2&0\\0&0&1\end{pmatrix},\qquad\begin{pmatrix}1&1&1\\1&1&1\\1&1&1\end{pmatrix}.\]
			\item\label{ex:dec:A2} 写出 $\begin{pmatrix}3&-1\\-1&3\end{pmatrix}$ 的谱分解 $\lambda_1\vv q_1\vv q_1^\trans+\lambda_2\vv q_2\vv q_2^\trans$，并验证两个投影矩阵之和为 $\vv I$、之积为 $\vv O$.
			\item\label{ex:dec:A3} 求二阶实对称矩阵 $\vv A$，使它的特征值为 $1,2$，且 $(1,1)^\trans$ 是属于 $1$ 的特征向量.
			\item\label{ex:dec:A4} 设三阶实对称矩阵 $\vv A$ 的特征值为 $1,1,-1$，属于 $-1$ 的特征向量为 $(0,1,1)^\trans$. 求 $\vv A$.
			\item\label{ex:dec:A5} 求 $\vv x^\trans\vv{Ax}$ 在单位球面 $\|\vv x\|=1$ 上的最大值与最小值，以及取到它们的 $\vv x$：$\vv A=\begin{pmatrix}2&2&0\\2&2&0\\0&0&1\end{pmatrix}$.
			\item\label{ex:dec:A6} 求 $x_1^2+4x_1x_2+x_2^2$ 在单位圆 $x_1^2+x_2^2=1$ 上的最大值与最小值.
			\item\label{ex:dec:A7} 求下列矩阵的奇异值分解：$\begin{pmatrix}2&0\\0&-3\end{pmatrix}$，$\begin{pmatrix}1&1\\1&1\end{pmatrix}$，$(3,4)$（$1\times2$ 矩阵）.
			\item\label{ex:dec:A8} 求 $\begin{pmatrix}1&1\\0&1\end{pmatrix}$ 的奇异值，并与它的特征值比较.（答案与黄金分割比有关.）
			\item\label{ex:dec:A9} 求 $\begin{pmatrix}3&1\\1&3\end{pmatrix}$ 的最佳秩一近似，并计算误差 $\|\vv A-\vv A_1\|_F$.
			\item\label{ex:dec:A10} 设 $\vv A=\begin{pmatrix}5&4\\4&5\end{pmatrix}$. 利用谱分解求 $\vv A^{10}$（写成两个矩阵之和的形式即可），并求对称正定矩阵 $\vv B$ 使 $\vv B^2=\vv A$.
			\item\label{ex:dec:A11} 证明：特征值全为 $0$ 的实对称矩阵是零矩阵. 对非对称矩阵举出反例.
			\item\label{ex:dec:A12} 证明：既对称又正交的实矩阵 $\vv A$ 满足 $\vv A^2=\vv I$，其特征值为 $\pm1$；并说明它在几何上是什么变换.
			\item\label{ex:dec:A13} 证明：$\vv A^\trans\vv A$ 与 $\vv{AA}^\trans$ 的非零特征值相同（重数也相同）.
			\item\label{ex:dec:A14} 证明：$\|\vv A\|_F^2=\tr(\vv A^\trans\vv A)=\sigma_1^2+\cdots+\sigma_r^2$.
			\item\label{ex:dec:A15} 设 $\vv A$ 是可逆实方阵. 证明 $|\det\vv A|=\sigma_1\sigma_2\cdots\sigma_n$，$\vv A^{-1}$ 的奇异值为 $1/\sigma_i$.
			\item\label{ex:dec:A16} 判断下列矩阵能否用正交矩阵对角化，说明理由：$\begin{pmatrix}1&2\\0&3\end{pmatrix}$，$\begin{pmatrix}1&2\\2&-2\end{pmatrix}$，$\begin{pmatrix}0&-1\\1&0\end{pmatrix}$（在实数范围内）.
			\item\label{ex:dec:A17} 四个已中心化的二维数据点为 $(2,1)$，$(-2,-1)$，$(1,2)$，$(-1,-2)$. 求样本协方差矩阵、第一主成分方向及其解释的方差比例.
			\item\label{ex:dec:A18} 计算 $\diag(10,1)$ 与 $\begin{pmatrix}1&1\\1&1.01\end{pmatrix}$ 的条件数（后者用近似值）. 后者的行列式与 $\diag(10,0.001)$ 相同，二者的条件数相同吗？
			\item\label{ex:dec:A19} 一张 $800\times600$ 像素的灰度图片，保留前 $k=40$ 个奇异值时需存储多少个数？压缩比是多少？$k$ 取多少时存储量反而超过原图？
			\item\label{ex:dec:A20} 证明：实对称矩阵的秩等于非零特征值的个数；奇异值的个数（非零）等于秩.
			\item\label{ex:dec:A21} 证明：实对称矩阵 $\vv A$ 满足``对一切 $\vv x\ne\zero$，$\vv x^\trans\vv A\vv x>0$''，当且仅当 $\vv A$ 的特征值全为正.（提示：用谱定理把 $\vv x^\trans\vv A\vv x$ 写成 $\sum\lambda_iy_i^2$.）
		\end{exA}
		
		\Bgroup
		\begin{exB}
			\item\label{ex:dec:B1} 设 $\vv A,\vv B$ 都是 $n$ 阶实对称矩阵且 $\vv{AB}=\vv{BA}$. 证明存在正交矩阵 $\vv Q$ 使 $\vv Q^\trans\vv{AQ}$ 与 $\vv Q^\trans\vv{BQ}$ 同时为对角矩阵.
			\item\label{ex:dec:B2} 证明：半正定的实对称矩阵 $\vv A$（特征值全 $\ge0$）有唯一的半正定平方根 $\vv B$（$\vv B^2=\vv A$，$\vv B$ 对称且特征值全 $\ge0$）.
			\item\label{ex:dec:B3} 设实对称矩阵 $\vv A$ 的特征值为 $\lambda_1\ge\lambda_2\ge\cdots\ge\lambda_n$，$\vv q_1$ 是属于 $\lambda_1$ 的单位特征向量. 证明 $\lambda_2=\max\{\vv x^\trans\vv{Ax}:\|\vv x\|=1,\ \vv x\perp\vv q_1\}$. 这说明第二主成分是``与第一主成分垂直的方向中方差最大的''.
			\item\label{ex:dec:B4} 证明：对任意 $\vv x$，$\|\vv{Ax}\|\le\sigma_1\|\vv x\|$，且等号可以取到. 因此 $\sigma_1=\max_{\vv x\ne\zero}\|\vv{Ax}\|/\|\vv x\|$，称为 $\vv A$ 的\textbf{谱范数}\index{pu3 fan4shu4@谱范数} $\|\vv A\|_2$.
			\item\label{ex:dec:B5} 证明：方阵 $\vv A$ 的任一特征值 $\lambda$ 满足 $\sigma_{\min}\le|\lambda|\le\sigma_{\max}$.
			\item\label{ex:dec:B6} 求 $\vv A=\begin{pmatrix}1&0&1\\0&1&1\end{pmatrix}$ 的奇异值分解.
			\item\label{ex:dec:B7} （伪逆）设 $\vv A=\vv U\vv\Sigma\vv V^\trans$，定义 $\vv A^+=\vv V\vv\Sigma^+\vv U^\trans$，其中 $\vv\Sigma^+$ 是把 $\vv\Sigma$ 转置并把非零元取倒数. 证明 $\vv A^+\vv b$ 是 $\vv{Ax}=\vv b$ 的范数最小的最小二乘解. 当 $\vv A$ 列满秩时，它与第九章的 $(\vv A^\trans\vv A)^{-1}\vv A^\trans$ 一致.
			\item\label{ex:dec:B8} 设 $p$ 个标准化指标两两相关系数都为 $\rho$（$0<\rho<1$），相关系数矩阵为 $(1-\rho)\vv I+\rho\vv J$（$\vv J$ 为全 $1$ 矩阵）. 求它的特征值与第一主成分，并求第一主成分解释的方差比例. 当 $\rho\to1$ 与 $\rho\to0$ 时比例分别趋于多少？
			\item\label{ex:dec:B9} 设两个指标的协方差矩阵为 $\begin{pmatrix}100&6\\6&1\end{pmatrix}$（第一个指标以``元''计，第二个以``万元''计）. 对协方差矩阵与相关系数矩阵分别作主成分分析，比较第一主成分. 由此说明：主成分分析对量纲敏感，为什么通常先标准化.
			\item\label{ex:dec:B10} 证明 $\|\vv A-\vv A_k\|_2=\sigma_{k+1}$.
			\item\label{ex:dec:B11} 证明：若 $\vv A$ 的奇异值互不相同且都不为零，则 $\vv A$ 的奇异值分解中的 $\vv u_i,\vv v_i$ 除了同时改变符号之外是唯一的.
			\item\label{ex:dec:B12} （正规矩阵）满足 $\vv A^\trans\vv A=\vv{AA}^\trans$ 的实方阵称为正规矩阵. 证明二阶实正规矩阵只有两类：对称矩阵，以及形如 $\begin{pmatrix}a&-b\\b&a\end{pmatrix}$ 的矩阵（旋转与伸缩的复合）.
			\item\label{ex:dec:B13} 设 $\vv A$ 是实对称矩阵. 证明 $\vv A$ 的特征值都落在 $[\min_i(a_{ii}-r_i),\ \max_i(a_{ii}+r_i)]$ 中，$r_i=\sum_{j\ne i}|a_{ij}|$. 并证明 $\lambda_{\max}\ge\max_ia_{ii}$.
			\item\label{ex:dec:B14} 对第十章``潜在语义分析''的例子，取 $k=1$，说明此时四份文书分别被表示成什么，查询``解约''与哪几份文书相似. 与 $k=2$ 比较，``压缩过头''的现象是否更严重？
			\item\label{ex:dec:B15} 证明：实对称矩阵 $\vv A$ 的特征值都在 $[a,b]$ 中，当且仅当 $\vv A-a\vv I$ 与 $b\vv I-\vv A$ 的二次型都非负.
		\end{exB}
		
		\Cgroup
		\begin{exC}
			\item\label{ex:dec:C1} （Eckart–Young 定理的证明）设 $\rk\vv B\le k$. 证明 $\ker\vv B$ 与 $\spn(\vv v_1,\ldots,\vv v_{k+1})$ 有非零公共向量 $\vv x$，由此证明 $\|\vv A-\vv B\|_2\ge\sigma_{k+1}$. 再借助 Weyl 不等式 $\sigma_{i+j-1}(\vv X+\vv Y)\le\sigma_i(\vv X)+\sigma_j(\vv Y)$（可承认）证明 Frobenius 范数的情形.
			\item\label{ex:dec:C2} （Cauchy 交错定理）设实对称矩阵 $\vv A$ 的特征值为 $\lambda_1\ge\cdots\ge\lambda_n$，删去最后一行一列所得矩阵的特征值为 $\mu_1\ge\cdots\ge\mu_{n-1}$. 证明 $\lambda_1\ge\mu_1\ge\lambda_2\ge\mu_2\ge\cdots\ge\mu_{n-1}\ge\lambda_n$.（提示：Courant–Fischer 极小极大刻画.）
			\item\label{ex:dec:C3} （复谱定理）证明：复方阵 $\vv A$ 可以被酉矩阵对角化（存在 $\vv U$ 使 $\bar{\vv U}^\trans\vv{AU}$ 为对角矩阵），当且仅当 $\vv A\bar{\vv A}^\trans=\bar{\vv A}^\trans\vv A$. 实对称矩阵、实正交矩阵、实反对称矩阵都满足这个条件.
			\item\label{ex:dec:C4} （正交 Procrustes 问题）给定 $n\times k$ 矩阵 $\vv A,\vv B$，求正交矩阵 $\vv Q$ 使 $\|\vv{QA}-\vv B\|_F$ 最小. 证明：若 $\vv{BA}^\trans=\vv U\vv\Sigma\vv V^\trans$，则最优解为 $\vv Q=\vv{UV}^\trans$. 这个问题出现在比较两组测量数据、对齐两份词向量等场合.
			\item\label{ex:dec:C5} （Ky Fan 定理）设实对称矩阵的特征值为 $\lambda_1\ge\cdots\ge\lambda_n$. 证明对任意列标准正交的 $n\times k$ 矩阵 $\vv X$，$\tr(\vv X^\trans\vv{AX})\le\lambda_1+\cdots+\lambda_k$. 由此说明前 $k$ 个主成分``一起''解释的方差也是最大的.
			\item\label{ex:dec:C6} （幂迭代）设 $\vv A$ 可对角化，特征值满足 $|\lambda_1|>|\lambda_2|\ge\cdots$. 证明：对``一般的''初始向量 $\vv x_0$，$\vv x_k=\vv A^k\vv x_0/\|\vv A^k\vv x_0\|$ 收敛到属于 $\lambda_1$ 的单位特征向量（可能差一个符号），收敛速度由 $|\lambda_2/\lambda_1|$ 决定. 对 $\begin{pmatrix}2&1\\1&2\end{pmatrix}$ 从 $(1,0)^\trans$ 出发迭代五步验证.
			\item\label{ex:dec:C7} （图的 Laplace 矩阵）设 $\vv L=\vv D-\vv A$ 是无向图的 Laplace 矩阵（见第五章章末习题~\ref{ex:det:C4}）. 证明 $\vv x^\trans\vv L\vv x=\sum_{\text{边 }ij}(x_i-x_j)^2$，因而 $\vv L$ 半正定；特征值 $0$ 的重数等于图的连通分支数. 第二小的特征值对应的特征向量（Fiedler 向量）的正负号常被用来把图``切''成两半——这就是谱聚类的思想.
			\item\label{ex:dec:C8} （正交回归）对中心化数据点 $\vv z_1,\ldots,\vv z_n\in\mathbb{R}^2$，求过原点的直线使各点到直线的\textbf{垂直}距离平方和最小. 证明最优方向就是第一主成分方向. 与普通最小二乘（只看竖直方向的残差）比较，两条直线何时相同？
			\item\label{ex:dec:C9} （扰动分析）设 $\vv A$ 可逆，$\vv{Ax}=\vv b$，$\vv A(\vv x+\delta\vv x)=\vv b+\delta\vv b$. 证明 $\dfrac{\|\delta\vv x\|}{\|\vv x\|}\le\kappa(\vv A)\dfrac{\|\delta\vv b\|}{\|\vv b\|}$，$\kappa(\vv A)=\sigma_1/\sigma_n$，并说明等号可以取到. 用它解释第六章病态方程组的例子.
		\end{exC}
	}
	
	% =====================================================================
	% 第十一章　二次型
	% =====================================================================
	\newcommand{\ExQuadratic}{%
		\exsetheading
		\Agroup
		\begin{exA}
			\item\label{ex:quad:A1} 写出下列二次型的矩阵与秩：
			(1) $x_1^2-3x_2^2+4x_1x_2-2x_2x_3$；(2) $x_1x_2+x_2x_3+x_3x_1$；(3) $(x_1+x_2+x_3)^2$；(4) $\sum_{i=1}^n x_i^2+\sum_{i<j}x_ix_j$.
			\item\label{ex:quad:A2} 写出以 $\begin{pmatrix}1&-1&0\\-1&2&3\\0&3&0\end{pmatrix}$ 为矩阵的二次型.
			\item\label{ex:quad:A3} 用配方法化为标准形，并写出所用的可逆变量替换：
			(1) $x_1^2+2x_2^2+5x_3^2+2x_1x_2+6x_2x_3$；(2) $x_1^2-4x_1x_2+4x_2^2-x_3^2$；(3) $x_1x_2-x_2x_3$.
			\item\label{ex:quad:A4} 用正交变换法把 $2x_1^2+3x_2^2+3x_3^2+4x_2x_3$ 化为标准形，写出正交矩阵. 曲面 $2x_1^2+3x_2^2+3x_3^2+4x_2x_3=1$ 是什么形状？
			\item\label{ex:quad:A5} 求 A3 中各二次型的正、负惯性指数与符号差，并写出规范形.
			\item\label{ex:quad:A6} 判断下列二次型是否正定：
			(1) $2x_1^2+x_2^2+3x_3^2+2x_1x_2$；(2) $x_1^2+x_2^2+x_3^2+4x_1x_2$；(3) $5x_1^2+x_2^2+5x_3^2+4x_1x_2-8x_1x_3-4x_2x_3$.
			\item\label{ex:quad:A7} $t$ 取何值时，$2x_1^2+x_2^2+x_3^2+2x_1x_2+tx_2x_3$ 正定？
			\item\label{ex:quad:A8} 判断曲线类型并求主轴方向：(1) $x^2-4xy+y^2=1$；(2) $3x^2+2xy+3y^2=8$.
			\item\label{ex:quad:A9} 判断下列矩阵是否合同：(1) $\diag(1,2,-3)$ 与 $\diag(-1,5,6)$；(2) $\begin{pmatrix}0&1\\1&0\end{pmatrix}$ 与 $\vv I_2$；(3) $\begin{pmatrix}0&1\\1&0\end{pmatrix}$ 与 $\diag(1,-1)$. 合同时求出 $\vv C$.
			\item\label{ex:quad:A10} 证明：若 $\vv A$ 正定，则 $\vv A^*$ 也正定.
			\item\label{ex:quad:A11} 求 $f(x,y)=x^3+y^3-3xy$ 的全部驻点，并用 Hessian 矩阵判断其类型.
			\item\label{ex:quad:A12} 判断 $\begin{pmatrix}2&-1&0\\-1&2&-1\\0&-1&2\end{pmatrix}$ 是否正定，并求它的 Cholesky 分解.
			\item\label{ex:quad:A13} 求 $\begin{pmatrix}4&2&-2\\2&10&2\\-2&2&6\end{pmatrix}$ 的 Cholesky 分解.
			\item\label{ex:quad:A14} 三种资产收益率的协方差矩阵为 $\vv\Sigma=\begin{pmatrix}0.04&0.01&0\\0.01&0.09&0\\0&0&0.01\end{pmatrix}$. 求组合 $\vv w=(0.5,0.3,0.2)^\trans$ 的方差与标准差.
			\item\label{ex:quad:A15} 两种资产收益率的标准差为 $30\%$ 与 $20\%$，互不相关. 求最小方差组合及其标准差.
			\item\label{ex:quad:A16} 证明：合同变换保持矩阵的对称性与秩. 合同变换保持行列式的符号吗？保持行列式的值吗？
			\item\label{ex:quad:A17} 分别在 $\mathbb{R}$ 上与 $\mathbb{C}$ 上写出二次型 $x_1^2-x_2^2+2x_3^2$ 的规范形.
			\item\label{ex:quad:A18} 判断下列方程表示的二次曲面类型：$x^2+y^2-z^2=1$，$x^2-y^2-z^2=1$，$x^2+y^2-z^2=0$，$x^2+y^2=1$.（提示：先看正、负惯性指数.）
			\item\label{ex:quad:A19} 说明 $\begin{pmatrix}1&2\\2&1\end{pmatrix}$ 不正定，并找出一个使二次型取负值的向量.
			\item\label{ex:quad:A20} 设 $\vv A$ 正定. 证明：$\vv A$ 的任一主子矩阵（取相同的若干行和列）都正定，从而所有主子式都为正.
		\end{exA}
		
		\Bgroup
		\begin{exB}
			\item\label{ex:quad:B1} 设 $\vv A,\vv B$ 正定且 $\vv{AB}=\vv{BA}$. 证明 $\vv{AB}$ 正定. 举例说明``可交换''不能去掉（此时 $\vv{AB}$ 甚至不对称）.
			\item\label{ex:quad:B2} 设 $\vv A$ 正定，$\vv B$ 半正定. 证明 $\det(\vv A+\vv B)\ge\det\vv A$，等号成立当且仅当 $\vv B=\vv O$.（提示：同时合同对角化.）
			\item\label{ex:quad:B3} 设 $\vv A$ 正定，$\vv B$ 对称. 证明：当 $|t|$ 充分小时 $\vv A+t\vv B$ 正定.
			\item\label{ex:quad:B4} 设 $\vv A$ 正定. 证明对 $i\ne j$ 有 $|a_{ij}|<\sqrt{a_{ii}a_{jj}}$，从而 $\vv A$ 的绝对值最大的元素在对角线上.
			\item\label{ex:quad:B5} 证明：实对称矩阵 $\vv A$ 满足 $\lambda_{\min}\|\vv x\|^2\le\vv x^\trans\vv{Ax}\le\lambda_{\max}\|\vv x\|^2$.
			\item\label{ex:quad:B6} 求 $\begin{pmatrix}0&1&1\\1&0&1\\1&1&0\end{pmatrix}$ 的正、负惯性指数（用两种方法：特征值；配方法）.
			\item\label{ex:quad:B7} $\lambda$ 取何值时，$\lambda(x_1^2+x_2^2+x_3^2)+2x_1x_2+2x_1x_3-2x_2x_3$ 正定？
			\item\label{ex:quad:B8} 求二次型 $\vv x^\trans\vv{Ax}$ 在约束 $\vv c^\trans\vv x=1$ 下的最小值（$\vv A$ 正定）：用拉格朗日乘数法证明最优解为 $\vv x^*=\dfrac{\vv A^{-1}\vv c}{\vv c^\trans\vv A^{-1}\vv c}$. 最小方差组合是它的特例.
			\item\label{ex:quad:B9} 设三种资产互不相关、方差都为 $\sigma^2$，期望收益率为 $\mu_1<\mu_2<\mu_3$. 求期望收益为 $m$ 的最小方差组合的权重，并说明有效前沿是一条抛物线（在``方差—期望收益''平面上）.
			\item\label{ex:quad:B10} 设 $\vv A=\begin{pmatrix}2&1\\1&2\end{pmatrix}$，$\vv B=\begin{pmatrix}1&0\\0&-1\end{pmatrix}$. 求可逆矩阵 $\vv C$ 使 $\vv C^\trans\vv{AC}=\vv I$ 且 $\vv C^\trans\vv{BC}$ 为对角矩阵.
			\item\label{ex:quad:B11} 设 $\vv A$ 正定. 求 $f(\vv x)=\vv x^\trans\vv{Ax}+\vv b^\trans\vv x+c$ 的最小值点与最小值. 若 $\vv A$ 只是半正定，最小值一定存在吗？
			\item\label{ex:quad:B12} 设 $\vv A$ 半正定，$\vv x^\trans\vv{Ax}=0$. 证明 $\vv{Ax}=\zero$.
			\item\label{ex:quad:B13} （Schur 补判据）证明：对称分块矩阵 $\begin{pmatrix}\vv A&\vv B\\ \vv B^\trans&\vv C\end{pmatrix}$ 正定，当且仅当 $\vv A$ 正定且 $\vv C-\vv B^\trans\vv A^{-1}\vv B$ 正定.
			\item\label{ex:quad:B14} 设 $\vv A$ 是二阶正定矩阵. 证明椭圆区域 $\{\vv x:\vv x^\trans\vv{Ax}\le1\}$ 的面积为 $\pi/\sqrt{\det\vv A}$.（提示：用主轴定理化为标准椭圆.）
			\item\label{ex:quad:B15} 设 $\vv A$ 为 $m\times n$ 实矩阵，$\mu>0$. 证明 $\vv A^\trans\vv A+\mu\vv I$ 正定，并说明岭回归 $\min\|\vv y-\vv{Ax}\|^2+\mu\|\vv x\|^2$ 的解为 $(\vv A^\trans\vv A+\mu\vv I)^{-1}\vv A^\trans\vv y$.
			\item\label{ex:quad:B16} 证明：实对称矩阵 $\vv A$ 正定，当且仅当存在正定矩阵 $\vv B$ 使 $\vv A=\vv B^2$；当且仅当 $\vv A^{-1}$ 存在且正定.
		\end{exB}
		
		\Cgroup
		\begin{exC}
			\item\label{ex:quad:C1} （Schur 乘积定理）对同阶矩阵定义逐元素乘积 $(\vv A\circ\vv B)_{ij}=a_{ij}b_{ij}$. 证明：若 $\vv A,\vv B$ 都半正定，则 $\vv A\circ\vv B$ 也半正定.（提示：把 $\vv B$ 写成 $\sum\lambda_k\vv q_k\vv q_k^\trans$.）
			\item\label{ex:quad:C2} （Minkowski 行列式不等式）设 $\vv A,\vv B$ 都是 $n$ 阶正定矩阵. 证明 $\det(\vv A+\vv B)^{1/n}\ge\det(\vv A)^{1/n}+\det(\vv B)^{1/n}$.（提示：同时合同对角化后用 $\prod(1+\mu_i)^{1/n}\ge1+\prod\mu_i^{1/n}$.）
			\item\label{ex:quad:C3} （有理数域上的二次型）证明 $x^2+y^2$ 与 $3x^2+3y^2$ 在 $\mathbb{R}$ 上合同，但在 $\mathbb{Q}$ 上不合同（即不存在有理系数的可逆替换把一个变为另一个）.（提示：若合同，则存在有理数 $a,b$ 使 $a^2+b^2=3$.）这说明惯性定理依赖于实数可以随意开平方.
			\item\label{ex:quad:C4} 函数 $f=x^2+y^4$ 与 $g=x^2-y^4$ 在原点处的 Hessian 矩阵相同且只是半正定，但前者在原点取极小值，后者不取极值. 由此说明 Hessian 半正定时二阶判别法失效.
			\item\label{ex:quad:C5} （广义特征值）设 $\vv A$ 对称，$\vv B$ 正定. 证明 $\max_{\vv x\ne\zero}\dfrac{\vv x^\trans\vv{Ax}}{\vv x^\trans\vv{Bx}}$ 等于方程 $\det(\vv A-\lambda\vv B)=0$ 的最大根. 统计学中的判别分析（Fisher LDA）正是这样求最佳投影方向的.
			\item\label{ex:quad:C6} （Mahalanobis 距离）设 $\vv\Sigma$ 正定，定义 $d(\vv x,\vv y)=\sqrt{(\vv x-\vv y)^\trans\vv\Sigma^{-1}(\vv x-\vv y)}$. 证明它满足三角不等式，且在仿射变换 $\vv x\mapsto\vv{Px}+\vv b$、$\vv\Sigma\mapsto\vv P\vv\Sigma\vv P^\trans$ 下不变. 它与欧氏距离有何不同？
			\item\label{ex:quad:C7} （Hilbert 矩阵）证明 $n$ 阶矩阵 $\vv H=\big(\frac1{i+j-1}\big)$ 正定.（提示：$\frac1{i+j-1}=\int_0^1x^{i-1}x^{j-1}\,\mathrm dx$，它是某个内积下的 Gram 矩阵.）
			\item\label{ex:quad:C8} （切点组合）设有无风险资产，收益率为 $r_f$. 证明使``夏普比率''$\dfrac{\vv\mu^\trans\vv w-r_f}{\sqrt{\vv w^\trans\vv\Sigma\vv w}}$（约束 $\vv 1^\trans\vv w=1$）最大的风险资产组合为 $\vv w\propto\vv\Sigma^{-1}(\vv\mu-r_f\vv 1)$.（提示：先去掉约束，比值关于 $\vv w$ 是零次齐次的；再用 Cauchy–Schwarz 不等式.）
			\item\label{ex:quad:C9} 证明：秩为 $r$、正惯性指数为 $p$ 的实二次型可以写成 $p$ 个线性型的平方和减去 $r-p$ 个线性型的平方和，并且不能写成少于 $r$ 个线性型的平方（带 $\pm$ 号）的组合.
		\end{exC}
	}
	
	% =====================================================================
	% 第十二章　*对偶空间与双线性型（选读）
	% =====================================================================
	\newcommand{\ExBilinear}{%
		\exsetheading
		\Agroup
		\begin{exA}
			\item\label{ex:bil:A1} 判断下列映射是否为相应空间上的线性函数：(1) $\mathbb{R}^3\to\mathbb{R}$，$\vv x\mapsto x_1-2x_3$；(2) $\vv x\mapsto x_1x_2$；(3) $\mathbb{R}[x]\to\mathbb{R}$，$f\mapsto f'(0)$；(4) $f\mapsto\int_0^1f(x)^2\,\mathrm dx$；(5) $M_n(\mathbb{R})\to\mathbb{R}$，$\vv X\mapsto\tr(\vv{AX})$（$\vv A$ 固定）.
			\item\label{ex:bil:A2} 求 $\mathbb{R}^3$ 的基 $(1,0,0)^\trans$，$(1,1,0)^\trans$，$(1,1,1)^\trans$ 的对偶基（把每个对偶基向量写成 $\vv x\mapsto a_1x_1+a_2x_2+a_3x_3$ 的形式）.
			\item\label{ex:bil:A3} 把线性函数 $f(\vv x)=x_1+2x_2+3x_3$ 用 A2 中的对偶基线性表出.
			\item\label{ex:bil:A4} 求 $W=\spn\big((1,1,0)^\trans,(0,1,1)^\trans\big)$ 的零化子 $W^0\subset(\mathbb{R}^3)^*$ 的一组基.
			\item\label{ex:bil:A5} 设线性映射 $\mathscr T:\mathbb{R}^2\to\mathbb{R}^3$ 在标准基下的矩阵为 $\begin{pmatrix}1&2\\0&1\\3&-1\end{pmatrix}$. 求对偶映射 $\mathscr T^*$ 在对偶基下的矩阵，并计算 $\mathscr T^*$ 作用在线性函数 $\vv y\mapsto y_1+y_2+y_3$ 上的结果.
			\item\label{ex:bil:A6} 求双线性型 $f(\vv x,\vv y)=x_1y_1+2x_1y_2-x_2y_1+3x_2y_2$ 在标准基下的度量矩阵，以及在基 $(1,1)^\trans$，$(1,-1)^\trans$ 下的度量矩阵，并验证 $\vv B=\vv P^\trans\vv{AP}$.
			\item\label{ex:bil:A7} 把 A6 中的 $f$ 分解为对称部分与交错部分之和.
			\item\label{ex:bil:A8} 判断下列双线性型是否非退化：(1) $x_1y_1-x_2y_2$；(2) $x_1y_2+x_2y_1$；(3) $(x_1+x_2)(y_1+y_2)$.
			\item\label{ex:bil:A9} 写出对称双线性型 $f(\vv x,\vv y)=x_1y_1+x_1y_2+x_2y_1+3x_2y_2$ 对应的二次型；反过来，由二次型 $x_1^2-4x_1x_2+2x_2^2$ 用极化恒等式求对应的对称双线性型.
			\item\label{ex:bil:A10} 验证 $\mathbb{R}^2$ 上的 $f(\vv x,\vv y)=\det(\vv x,\vv y)$ 是交错双线性型，写出它的度量矩阵. 它非退化吗？
			\item\label{ex:bil:A11} 设 $a_1,a_2,a_3$ 互异，$f_i(p)=p(a_i)$ 是 $\mathbb{R}[x]_3$ 上的线性函数. 对 $a_i=0,1,2$，求以 $f_1,f_2,f_3$ 为对偶基的 $\mathbb{R}[x]_3$ 的基.
			\item\label{ex:bil:A12} 在 $\mathbb{R}^4$ 中，设 $W$ 由 $x_1+x_2=0$，$x_3-x_4=0$ 定义. 求 $\dim W$ 与 $\dim W^0$，并写出 $W^0$ 的一组基.
			\item\label{ex:bil:A13} 证明 $\langle\vv A,\vv B\rangle=\tr(\vv{AB})$ 是 $M_n(\mathbb{R})$ 上非退化的对称双线性型. 它是内积吗？
			\item\label{ex:bil:A14} 对 $\mathbb{R}[x]_3$ 上的线性函数 $\varphi(p)=\int_0^1p(x)\,\mathrm dx$，求它在 $1,x,x^2$ 的对偶基下的坐标.
			\item\label{ex:bil:A15} （Riesz 表示定理的有限维情形）设 $V$ 是有限维实内积空间. 证明：对每个 $f\in V^*$，存在唯一的 $\vv\beta\in V$ 使 $f(\vv\alpha)=\langle\vv\alpha,\vv\beta\rangle$ 对一切 $\vv\alpha$ 成立.（提示：取标准正交基 $\vv e_1,\ldots,\vv e_n$，令 $\vv\beta=\sum f(\vv e_i)\vv e_i$.）
		\end{exA}
		
		\Bgroup
		\begin{exB}
			\item\label{ex:bil:B1} 证明：$f_1,\ldots,f_k\in V^*$ 线性无关，当且仅当映射 $V\to\mathbb{F}^k$，$\vv\alpha\mapsto(f_1(\vv\alpha),\ldots,f_k(\vv\alpha))^\trans$ 是满射.
			\item\label{ex:bil:B2} 设 $f,g\in V^*$ 且 $\ker f\subset\ker g$. 证明存在常数 $c$ 使 $g=cf$. 推广：若 $\ker f_1\cap\cdots\cap\ker f_k\subset\ker g$，则 $g$ 是 $f_1,\ldots,f_k$ 的线性组合.
			\item\label{ex:bil:B3} 证明：$\ker f_1\cap\cdots\cap\ker f_k=\{\zero\}$ 当且仅当 $f_1,\ldots,f_k$ 张成 $V^*$. 用方程组的语言解释这个结论.
			\item\label{ex:bil:B4} 证明：$\mathscr T$ 是单射当且仅当 $\mathscr T^*$ 是满射；$\mathscr T$ 是满射当且仅当 $\mathscr T^*$ 是单射.
			\item\label{ex:bil:B5} （Simpson 公式）求常数 $a,b,c$，使 $\int_0^2p(x)\,\mathrm dx=a\,p(0)+b\,p(1)+c\,p(2)$ 对一切 $p\in\mathbb{R}[x]_3$ 成立. 用对偶基的语言说明为什么这样的 $a,b,c$ 存在且唯一. 验证这个公式对三次多项式也成立.
			\item\label{ex:bil:B6} 证明：$V$ 上全体双线性型构成的空间与 $M_n(\mathbb{F})$ 同构，其中对称型、交错型分别构成维数为 $\frac{n(n+1)}2$、$\frac{n(n-1)}2$ 的子空间，且整个空间是二者的直和.
			\item\label{ex:bil:B7} 证明：若 $V$ 上存在非退化的交错双线性型，则 $\dim V$ 是偶数.
			\item\label{ex:bil:B8} 求双线性型 $f(\vv x,\vv y)=x_1y_1+x_1y_2+x_2y_1+x_2y_2+x_3y_3$ 的``根'' $\{\vv\alpha:f(\vv\alpha,\cdot)=0\}$.
			\item\label{ex:bil:B9} 在 $\mathbb{R}^2$ 上取 $f(\vv x,\vv y)=x_1y_1-x_2y_2$. 求全部``迷向向量''（$f(\vv\alpha,\vv\alpha)=0$），并说明它们构成两条直线. 再证明保持 $f$ 的线性变换 $\vv P$（$\vv P^\trans\diag(1,-1)\vv P=\diag(1,-1)$，$\det\vv P=1$，且保持 $x_1>|x_2|$ 的区域）都形如 $\begin{pmatrix}\cosh t&\sinh t\\ \sinh t&\cosh t\end{pmatrix}$.
			\item\label{ex:bil:B10} 设 $W$ 是 $V$ 的子空间，$\mathscr T:V\to V$ 线性且 $\mathscr T(W)\subset W$. 证明 $\mathscr T^*(W^0)\subset W^0$.
		\end{exB}
		
		\Cgroup
		\begin{exC}
			\item\label{ex:bil:C1} 证明：对无限维空间 $\mathbb{F}[x]$，自然映射 $\vv\alpha\mapsto\hat{\vv\alpha}$ 是单射但不是满射.（提示：$\mathbb{F}[x]^*$ 中的元素由它在 $1,x,x^2,\ldots$ 上的取值任意决定，可以``无限多个非零''.）
			\item\label{ex:bil:C2} 证明：$V$ 上的双线性型与线性映射 $V\to V^*$ 一一对应：$f\mapsto(\vv\alpha\mapsto f(\vv\alpha,\cdot))$. 用这个对应和对偶基的换基规则，重新推出度量矩阵的换基公式 $\vv B=\vv P^\trans\vv{AP}$.
			\item\label{ex:bil:C3} 证明：二阶矩阵 $\vv P$ 满足 $\vv P^\trans\vv J\vv P=\vv J$（$\vv J=\begin{pmatrix}0&1\\-1&0\end{pmatrix}$）当且仅当 $\det\vv P=1$. 四阶时，满足 $\vv P^\trans\diag(\vv J,\vv J)\vv P=\diag(\vv J,\vv J)$ 的矩阵（辛矩阵）的行列式是多少？
			\item\label{ex:bil:C4} 求 $M_n(\mathbb{R})$ 上对称双线性型 $\tr(\vv{AB})$ 的正、负惯性指数.（提示：对称矩阵与反对称矩阵在这个型下``互相正交''.）
			\item\label{ex:bil:C5} 设 $f$ 是 $V$ 上的对称双线性型，$W$ 是子空间，$W^\perp=\{\vv\alpha:f(\vv\alpha,\vv\beta)=0\ \forall\vv\beta\in W\}$. 证明 $\dim W^\perp\ge\dim V-\dim W$，$f$ 非退化时等号成立；举例说明即使 $f$ 非退化，$W\cap W^\perp$ 也可能不为零.
			\item\label{ex:bil:C6} 证明：$\mathbb{F}^n$ 上全体交错 $n$ 重线性函数构成的空间是一维的，由行列式张成. 这给出第五章章末习题~\ref{ex:det:C1} 的``对偶空间版本''.
			\item\label{ex:bil:C7} （线性规划对偶）考虑问题：在 $\vv{Ax}=\vv b$，$\vv x\ge\zero$ 下最小化 $\vv c^\trans\vv x$. 证明对任意满足 $\vv A^\trans\vv y\le\vv c$ 的 $\vv y$ 与任意可行的 $\vv x$，都有 $\vv b^\trans\vv y\le\vv c^\trans\vv x$（弱对偶）. 用``价格是数量空间上的线性函数''解释 $\vv y$ 的经济含义（影子价格）.
		\end{exC}
	} 之后，
%%      exercises.tex 中每道题都有标签 ex:<章键>:<组号>，本文件按标签引用题号）.
%% 章键：elim vs mat lmap det sol poly eig ip dec quad bil（依次为第一至十二章）.
%% =====================================================================

\providecommand{\ans}[2]{\par\smallskip\noindent\textbf{\ref{#1}.}\ #2\par}
\setlength{\emergencystretch}{3em}
\chapter{习题答案与提示}\label{app:answers}

本附录给出各章章末习题 A 组的答案与 B 组的提示. A 组的计算题给出最终结果，必要时附关键步骤；证明题只给出思路，完整的书写请读者自行完成. B 组只给出关键的一步或方向，希望读者在看提示之前先独立思考至少一刻钟. C 组一般不提供提示；正文或其他题目的答案中引用到的 C 组题例外，在该章末尾给出证明要点.

计算题的答案往往不唯一（例如基础解系、过渡矩阵的选取、特征向量的倍数），只要满足题目的全部要求，与这里的答案不同也是正确的. 检验的方法通常比答案本身更重要：代回原方程、计算迹与行列式、验证正交性.


\section{第一章　线性方程组与消元法}

\subsection*{A 组答案}

\ans{ex:elim:A1}{第 1、4 个是简化阶梯形；第 2 个是阶梯形但不是简化阶梯形（主元 $2$ 不为 $1$）；第 3、5 个不是阶梯形（第 3 个第二行主元在第一行主元左侧；第 5 个零行在非零行上方）.}

\ans{ex:elim:A2}{$\begin{pmatrix}3&0&-1&2\\0&1&4&0\\1&-1&0&5\end{pmatrix}$；反过来为 $x_1-2x_2=3$，$x_2+5x_3=-1$（未知量为 $x_1,x_2,x_3$，最后一列为常数项；$x_3$ 在第一个方程中的系数为 $0$）.}

\ans{ex:elim:A3}{$(x_1,x_2,x_3)=(1,2,3)$.}

\ans{ex:elim:A4}{$(x_1,x_2,x_3,x_4)=(1,-1,2,0)$.}

\ans{ex:elim:A5}{无解. $(r_3-r_1)-(r_2-2r_1)$ 给出 $(0,0,0\mid 2)$，即 $0=2$.}

\ans{ex:elim:A6}{主变量 $x_1,x_3$，自由变量 $x_2,x_4$. 全部解为 $x_1=4-2s-2t$，$x_2=s$，$x_3=1-t$，$x_4=t$（$s,t$ 任意）.}

\ans{ex:elim:A7}{$x_1=\frac27x_3+\frac37x_4$，$x_2=\frac57x_3+\frac47x_4$. 令 $x_3=7s$，$x_4=7t$，全部解为 $s(2,5,7,0)^\trans+t(3,4,0,7)^\trans$.}

\ans{ex:elim:A8}{$\begin{pmatrix}1&2&0&2\\0&0&1&-1\\0&0&0&0\end{pmatrix}$. 一种做法：$r_1\leftrightarrow r_2$；$r_2-2r_1$，$r_3-3r_1$；$-\frac14r_2$；$r_3+2r_2$；$r_1-r_2$.}

\ans{ex:elim:A9}{共四类：$\vv O$；$\begin{pmatrix}1&a\\0&0\end{pmatrix}$（$a$ 任意）；$\begin{pmatrix}0&1\\0&0\end{pmatrix}$；$\vv I_2$.}

\ans{ex:elim:A10}{由前两个方程得 $x=2$，$y=-1$，代入第三个得 $4-k=5$，$k=-1$.}

\ans{ex:elim:A11}{有解当且仅当 $b=3a$. 这时全部解为 $x=a-2t$，$y=t$（$t$ 任意）.}

\ans{ex:elim:A12}{$y=x^2-2x+3$.}

\ans{ex:elim:A13}{$y=1-2x+x^3$，即 $a=1$，$b=-2$，$c=0$，$d=1$.}

\ans{ex:elim:A14}{$30$ 元、$50$ 元、$80$ 元票分别为 $40$，$40$，$20$ 张.}

\ans{ex:elim:A15}{守恒方程 $3x_1=x_3$，$8x_1=2x_4$，$2x_2=2x_3+x_4$. 解为 $(x_1,x_2,x_3,x_4)=t(1,5,3,4)$，最小正整数解：$\mathrm{C_3H_8}+5\,\mathrm{O_2}\to3\,\mathrm{CO_2}+4\,\mathrm{H_2O}$.}

\ans{ex:elim:A16}{(1) 四个路口的方程为 $x_1-x_4=100$，$x_1-x_2=150$，$x_3-x_2=150$，$x_3-x_4=100$，全部解为 $x_1=t+100$，$x_2=t-50$，$x_3=t+100$，$x_4=t$. (2) 由 $x_2\ge0$ 得 $t\ge50$，故 $x_1$ 最少为 $150$. (3) 方程组无解：外部驶入总量 $1000$，驶出总量变为 $1100$，第一、三个方程之和减去第二、四个方程，得 $0=-100$. 守恒律是``进入路网的车辆总数等于离开的总数''.}

\ans{ex:elim:A17}{若 $a\ne0$，用第一个方程消去第二个方程中的 $x$，得 $(ad-bc)y=af-ce$. 解唯一当且仅当 $y$ 的系数 $ad-bc\ne0$. $a=0$ 时交换两个方程（或对 $c$ 讨论），同样得到条件 $ad-bc\ne0$.}

\ans{ex:elim:A18}{设 $\vv A\vv x_1=\vv A\vv x_2=\vv b$（这里先借用第三章的记号，也可以逐个方程写出），则 $\vv A(\vv x_1+t(\vv x_2-\vv x_1))=(1-t)\vv b+t\vv b=\vv b$. 若有两个不同的解，$t$ 取不同值给出不同的解，于是有无穷多个.}

\ans{ex:elim:A19}{第二个方程组的第一个方程等于第一个方程组两个方程之和，第二个方程相同，所以 (I) 的解都是 (II) 的解；反之 (I) 的第二个方程是 (II) 的第一个方程减去第二个方程. 两者的解集都是 $\{(t,t,3-2t)\}$.}

\ans{ex:elim:A20}{例如方程组 $x=0$，$y=0$. 同时做两步后得到 $x+y=0$，$x+y=0$，解集变大. 不矛盾：``同时''做的这两步合起来不是一次初等行变换，也不能被撤销（它把两行变成了相同的两行，丢失了信息）.}

\ans{ex:elim:A21}{有解当且仅当 $a=0$. 这时系数部分有 $2$ 个主元，$4$ 个未知量，有 $2$ 个自由变量.}

\subsection*{B 组提示}

\ans{ex:elim:B1}{用 $r_4-3r_1$，再加上 $r_2$，得 $(0,0,0,a-1\mid0)$；$r_3+r_2$ 得 $(0,0,a-1,0\mid b+1)$. $a\ne1$ 时解唯一；$a=1$，$b\ne-1$ 时无解；$a=1$，$b=-1$ 时有无穷多解：$(-1+s+t,\ 1-2s-2t,\ s,\ t)$.}

\ans{ex:elim:B2}{把第三个方程移到最前面再消元，得到的阶梯形后两行为 $(0,\lambda,-\lambda\mid3-\lambda)$ 与 $\big(0,0,-\lambda(\lambda+3)\mid-(\lambda+3)(\lambda-1)\big)$. $\lambda\ne0,-3$ 时有三个主元，解唯一；$\lambda=0$ 时倒数第二行变为 $(0,0,0\mid3)$，无解；$\lambda=-3$ 时最后一行为零行，只有两个主元，有无穷多解.}

\ans{ex:elim:B3}{齐次方程组只有零解说明简化阶梯形有 $n$ 个主元，即每一列都有主元. 而系数部分只有 $n$ 行，不可能出现 $(0,\ldots,0\mid d)$ 形的行.}

\ans{ex:elim:B4}{初等行变换得到的每一行都是原方程的线性组合. 无解时简化阶梯形中出现 $(0,\ldots,0\mid d)$，$d\ne0$，把这一行除以 $d$ 即可. 反之若组合得 $0=1$，任何解代入都会矛盾.}

\ans{ex:elim:B5}{依次做 $r_i+r_j$，$r_j-r_i$，$r_i+r_j$，最后把第 $j$ 行乘以 $-1$，逐步写出两行的变化即可.}

\ans{ex:elim:B6}{消元只用到四则运算，简化阶梯形的元素仍是有理数. 有解时令全部自由变量取 $0$，主变量由有理数算出. 解唯一时没有自由变量，唯一的解就是这个有理数解.}

\ans{ex:elim:B7}{把 $n$ 个方程相加，左边为 $0$，故必要条件是 $a_1+\cdots+a_n=0$. 它也充分：令 $x_n=t$，则 $x_k=t+a_k+a_{k+1}+\cdots+a_{n-1}$（$k<n$）.}

\ans{ex:elim:B8}{把第一个方程减第二个、加第三个……，交替相加. $n$ 为奇数时得 $2x_1=a_1-a_2+a_3-\cdots+a_n$，$x_1$ 唯一确定，其余依次确定；$n$ 为偶数时左边相消为 $0$.}

\ans{ex:elim:B9}{方程组为 $(\vv I-\vv A)\vv x=\vv d$（逐个写出即可），解得各部门总产出为 $(100,\ 200,\ 150)^\trans$.}

\ans{ex:elim:B10}{处理第 $k$ 列时，需要对下面 $n-k$ 行各做一次倍加，每次约 $n-k+1$ 次乘法，求和 $\sum_{k}(n-k)(n-k+1)\approx\frac13n^3$；回代只需约 $\frac12n^2$ 次.}

\ans{ex:elim:B11}{考察三个方程的系数部分化为阶梯形后的主元个数：$3$ 个时交于一点；$2$ 个且相容时交于一条直线；$2$ 个且出现 $0=d$ 时，三条交线方向相同（都平行于系数矩阵的解空间方向），没有公共点. 例子可以取 $z=0$，$x=0$，$x+z=1$.}

\ans{ex:elim:B12}{直接验证 $\vv A(\vv x+\vv y)=\vv A\vv x+\vv A\vv y$，$\vv A(k\vv x)=k\vv A\vv x$（或逐个方程验证）. 后半句：设 $\vv x_0$ 是一个解，任一解 $\vv x$ 满足 $\vv x-\vv x_0$ 是齐次方程组的解.}

\ans{ex:elim:B13}{开放题. 一种设定：未知量为三位律师各承办三类案件的件数，共 $9$ 个；已知条件给出 $3$ 个独立方程，故至少还缺 $6$ 条独立的信息. 关键是先写清楚未知量，再数主元个数.}

\ans{ex:elim:B14}{每行之和为零即 $\vv A(1,\ldots,1)^\trans=\zero$. 后半句由 B12 的结论.}

\ans{ex:elim:B15}{(I) 的每个解满足 (II) 的每个方程，因为后者是前者的线性组合；反之亦然. 逆命题对有解的方程组成立，证明要用到第六章的结论（齐次方程组同解当且仅当系数矩阵行向量组等价，第六章章末习题~\ref{ex:sol:B11}）.}

\ans{ex:elim:B16}{$\lambda\ne1,-2$ 时解唯一：$x=-\frac{\lambda+1}{\lambda+2}$，$y=\frac1{\lambda+2}$，$z=\frac{(\lambda+1)^2}{\lambda+2}$；$\lambda=1$ 时三个方程相同，有无穷多解；$\lambda=-2$ 时三个方程相加得 $0=3$，无解.}

\subsection*{C 组证明要点（正文引用的题目）}

\ans{ex:elim:C9}{设 $\vv R_1,\vv R_2$ 都是 $\vv M$ 的简化阶梯形. 它们行等价，以它们为系数矩阵的齐次方程组同解（定理~\ref{thm:sameSolution}）. 对列数 $n$ 归纳. $n=1$ 时，简化阶梯形是零列或 $\vv e_1$，由 $\vv M$ 是否为零唯一决定.
	
	设 $n>1$. 删去最后一列，得到的 $\vv R_1',\vv R_2'$ 仍是简化阶梯形，且仍行等价（用同一串初等行变换），由归纳假设 $\vv R_1'=\vv R_2'$. 于是两者在前 $n-1$ 列中的主元位置相同，设共有 $r$ 个主元. 再比较最后一列：
	\begin{enumerate}
		\item 若最后一列在 $\vv R_1$ 中是主元列，则 $\vv R_1$ 的第 $r+1$ 行为 $(0,\ldots,0,1)$，$\vv R_1\vv x=\zero$ 的解都有 $x_n=0$. 若它在 $\vv R_2$ 中不是主元列，令 $x_n=1$、其余自由变量取 $0$，可得 $\vv R_2\vv x=\zero$ 的一个解，其 $x_n=1$，与同解矛盾. 所以它在两者中同为主元列，都等于 $\vv e_{r+1}$.
		\item 若最后一列在两者中都不是主元列，令 $x_n=1$、其余自由变量取 $0$，两个方程组各有唯一一个这样的解，其主变量的值分别是两矩阵最后一列前 $r$ 个分量的相反数. 两方程组同解，故这两个解相同，最后一列相同.
\end{enumerate}}

\section{第二章　向量空间}

\subsection*{A 组答案}

\ans{ex:vs:A1}{(1) 是. (2) 不是：$\frac12\cdot\sqrt2=\frac{\sqrt2}2$ 的 $\sqrt2$ 系数不是整数，乘法不封闭. (3) 不是：不含 $1$，也不对除法封闭. (4) 不是：$\sqrt[3]2\cdot\sqrt[3]2=\sqrt[3]4$ 不在其中. (5) 不是：$\mathrm i\cdot\mathrm i=-1$ 不是纯虚数.}

\ans{ex:vs:A2}{(1) 是. (2) 不是：$(1,0)^\trans$ 的负向量不在其中. (3) 不是：不含零多项式，且两个 $n$ 次多项式之和可能降次. (4) 是. (5) 是. (6) 不是：不含零向量.}

\ans{ex:vs:A3}{(1) 线性相关：$(4,1,-1)^\trans=2(1,2,-1)^\trans+(2,-3,1)^\trans$. (2) 线性无关. (3) 线性相关：$\mathbb{R}[x]_3$ 是三维的，四个向量必相关. (4) 线性相关：全 $1$ 矩阵等于前两个之和.}

\ans{ex:vs:A4}{$\vv\beta=\frac54\vv\alpha_1+\frac14\vv\alpha_2-\frac14\vv\alpha_3-\frac14\vv\alpha_4$. （这四个向量两两``对应分量乘积之和''为零，每个系数都可以直接算出.）}

\ans{ex:vs:A5}{以三个向量为列作矩阵并消元：$r_2-r_1$ 得 $(0,2,-2)$，再 $r_3+\frac12r_2$ 得 $(0,0,t-1)$. $t=1$ 时出现自由变量，线性相关；$t\ne1$ 时线性无关.（学完第五章后，可用行列式 $2t-2$ 验证.）}

\ans{ex:vs:A6}{秩为 $3$，一个极大无关组为 $\vv\alpha_1,\vv\alpha_2,\vv\alpha_4$；$\vv\alpha_3=3\vv\alpha_1+\vv\alpha_2$，$\vv\alpha_5=\vv\alpha_1+\vv\alpha_2+\vv\alpha_4$.}

\ans{ex:vs:A7}{(1) 线性相关：三者之和为 $\zero$. (2) 线性无关. (3) 线性无关：设 $k_1(\vv\alpha_1+2\vv\alpha_2)+k_2(\vv\alpha_2+2\vv\alpha_3)+k_3(\vv\alpha_3+2\vv\alpha_1)=\zero$，由 $\vv\alpha_i$ 线性无关得 $k_1+2k_3=0$，$2k_1+k_2=0$，$2k_2+k_3=0$. 依次代入得 $k_1=-2k_3$，$k_2=4k_3$，$k_3=-8k_3$，故 $k_3=0$，从而 $k_1=k_2=0$.}

\ans{ex:vs:A8}{$1\cdot\zero=\zero$ 是系数不全为零的组合；$\vv\alpha-\vv\alpha=\zero$ 同理；$k\vv\alpha=\zero$ 且 $\vv\alpha\ne\zero$ 时由命题~\ref{prop:linspc}(5) 得 $k=0$.}

\ans{ex:vs:A9}{由定理~\ref{thm:relvandexp}，两个向量相关当且仅当其中一个可由另一个表出，即是另一个的倍数. 三个向量时只能说``其中一个是另两个的线性组合''，不一定是某一个的倍数，例如 $\vv e_1,\vv e_2,\vv e_1+\vv e_2$.}

\ans{ex:vs:A10}{(1) $\vv\alpha+\vv\alpha=1\vv\alpha+1\vv\alpha=(1+1)\vv\alpha$. (2) 两边乘 $k^{-1}$. (3) $-\vv\alpha$ 与 $\vv\alpha$ 之和为零，由负向量的唯一性.}

\ans{ex:vs:A11}{(1) $\dim=3$，基 $(-1,1,0,0)^\trans$，$(1,0,1,0)^\trans$，$(1,0,0,1)^\trans$. (2) $\dim=\frac{n(n+1)}2$，基 $\vv E_{ij}$（$i\le j$）. (3) $\dim=n^2-1$，基 $\vv E_{ij}$（$i\ne j$）与 $\vv E_{11}-\vv E_{ii}$（$i=2,\ldots,n$）. (4) $\dim=3$，基 $x-1$，$(x-1)^2$，$(x-1)^3$.}

\ans{ex:vs:A12}{(1) 是. (2) 不是：$(1,0,0)^\trans+(0,1,0)^\trans$ 不在其中. (3) 不是：对实数数乘不封闭. (4) 是.}

\ans{ex:vs:A13}{维数为 $3$，一组基为 $\vv\alpha_1,\vv\alpha_2,\vv\alpha_4$（$\vv\alpha_3=\vv\alpha_1+\vv\alpha_2$）.}

\ans{ex:vs:A14}{$2x^2-x+3=4+3(x-1)+2(x-1)^2$，坐标为 $(4,3,2)^\trans$.}

\ans{ex:vs:A15}{坐标为 $(-1,4,-5,4)^\trans$（从右下角元素开始依次确定）.}

\ans{ex:vs:A16}{(1) 是. (2) 不是：$\mathbb{Q}$ 对实数的数乘不封闭. (3) 作为 $\mathbb{R}$ 上的空间是 $2$ 维，作为 $\mathbb{C}$ 上的空间是 $1$ 维.}

\ans{ex:vs:A17}{例如 $a_0+a_1x+\cdots+a_{n-1}x^{n-1}\mapsto(a_0,\ldots,a_{n-1})^\trans$；$\begin{pmatrix}a&b\\c&d\end{pmatrix}\mapsto(a,b,c,d)^\trans$. 两者都是取坐标的映射，保持线性运算.}

\ans{ex:vs:A18}{$W_1+W_2$ 是直和（$(1,1,1)^\trans\notin W_1$，交为零），等于 $\mathbb{R}^3$，维数 $3$；$W_3\subset W_1$，$W_1+W_3=W_1$ 不是直和，维数 $2$.}

\ans{ex:vs:A19}{$\dim(W_1+W_2)=3$，$\dim(W_1\cap W_2)=1$：$(0,-1,1,1)^\trans=-(1,1,0,0)^\trans+(1,0,1,1)^\trans$ 同属两者. $2+2-1=3$，与维数公式一致.}

\ans{ex:vs:A20}{它们都可由一个含 $r$ 个向量的极大无关组线性表出，由基本引理任意 $r+1$ 个相关. 任意 $r$ 个无关的向量再添一个就相关，由引理~\ref{lem:addone} 的逆否形式，组中每个向量都可由它们线性表出.}

\subsection*{B 组提示}

\ans{ex:vs:B1}{设两者互不包含，取 $\vv\alpha\in W_1\setminus W_2$，$\vv\beta\in W_2\setminus W_1$，证明 $\vv\alpha+\vv\beta$ 不属于并集.}
\ans{ex:vs:B2}{取 (II) 的一个极大无关组，它含 $r$ 个向量且可由 (I) 的极大无关组表出，二者秩相同，用定理~\ref{thm:ndimbasis}(2) 在 (I) 张成的空间中论证.}
\ans{ex:vs:B3}{(III) 可由两者的极大无关组合并表出. 第二个不等式：左边的每个向量可由两组的极大无关组合并表出.}
\ans{ex:vs:B4}{新组含 $n$ 个向量，只需证明它能表出 $\vv\alpha_j$：$a_j\ne0$ 时 $\vv\alpha_j=\frac1{a_j}(\vv\beta-\sum_{i\ne j}a_i\vv\alpha_i)$；$a_j=0$ 时新组线性相关.}
\ans{ex:vs:B5}{按次数从高到低比较系数，证明线性无关；个数等于维数 $n$.}
\ans{ex:vs:B6}{$a_1,a_2$ 确定整个数列，映射 $\{a_n\}\mapsto(a_1,a_2)$ 是同构. 代入验证 $2^n$ 与 $(-1)^n$ 满足递推且线性无关. $a_1=a_2=1$ 时 $a_n=\frac13\big(2^n-(-1)^n\big)$.}
\ans{ex:vs:B7}{(1) 设组合为零，令 $x=0$ 并逐次求导，或除以 $\mathrm e^{3x}$ 后令 $x\to+\infty$. (2) $\sin(x+1)=\cos1\cdot\sin x+\sin1\cdot\cos x$. (3) 秩为 $3$：$\cos2x=2\cos^2x-1$.}
\ans{ex:vs:B8}{证明对加、减、乘封闭，再对除法用两次``分母有理化''. 线性无关可设 $a+b\sqrt2=-(c+d\sqrt2)\sqrt3$ 后平方，利用 $\sqrt2,\sqrt3,\sqrt6$ 都是无理数. 维数为 $4$.}
\ans{ex:vs:B9}{按维数分：$\{\zero\}$；过原点的直线；（$\mathbb{R}^3$ 中）过原点的平面；整个空间.}
\ans{ex:vs:B10}{用维数公式：$\dim(W_1\cap W_2)=k_1+k_2-\dim(W_1+W_2)\ge k_1+k_2-n>0$.}
\ans{ex:vs:B11}{$f=f(a)+(f-f(a))$. 交为零：常数多项式在 $a$ 处为零只能是零.}
\ans{ex:vs:B12}{$f(x)=\frac{f(x)+f(-x)}2+\frac{f(x)-f(-x)}2$；$\mathrm e^x=\cosh x+\sinh x$.}
\ans{ex:vs:B13}{对 $k$ 归纳，反复使用维数公式：$\dim(U\cap W)\ge\dim U+\dim W-n$.}
\ans{ex:vs:B14}{这 $n-1$ 个向量线性无关（用坐标验证），$\dim W=n-1$. 补空间可取 $\spn(\vv\alpha_n)$（需验证 $\vv\alpha_n\notin W$：$W$ 中向量的坐标交错和为零）.}
\ans{ex:vs:B15}{$\vv x=x_1(1,\ldots,1)^\trans+(\vv x-x_1(1,\ldots,1)^\trans)$. 两种分解都是直和；与 $V_1$ 互补的子空间有无穷多个，第九章会说明例~\ref{eg:pricedecomp} 中的 $R$ 是其中``垂直''的那一个.}
\ans{ex:vs:B16}{每个矩阵唯一地拆成``对角线及以上''与``对角线以下''两部分. $\dim W_1=\frac{n(n+1)}2$.}
\ans{ex:vs:B17}{(1) 设 $W$ 的基为 $\vv\varepsilon_1,\ldots,\vv\varepsilon_k$（$k<n$），扩充为 $\mathbb{R}^n$ 的基 $\vv\varepsilon_1,\ldots,\vv\varepsilon_n$. ``取 $\vv x$ 在这组基下的第 $n$ 个坐标''是 $x_1,\ldots,x_n$ 的一次齐次式（解一个线性方程组即可写出），它在 $W$ 上恒为零，又在 $\vv\varepsilon_n$ 处为 $1$，系数不全为零. (2) 条件是 $a_1+a_2t+\cdots+a_nt^{n-1}=0$，这是一个次数不超过 $n-1$ 的非零多项式，至多 $n-1$ 个根. (3) 例如 $t=2$：$\vv v(2)=(1,2,4)^\trans$ 不满足三个方程中的任何一个.}

\ans{ex:vs:B18}{对 $k$ 归纳. 已选好线性无关的 $\vv\varepsilon_1,\ldots,\vv\varepsilon_{k-1}$（$k\le n$）后，$U=\spn(\vv\varepsilon_1,\ldots,\vv\varepsilon_{k-1})$ 的维数 $k-1<n$，也是真子空间. 由覆盖定理取 $\vv\varepsilon_k\notin W_1\cup\cdots\cup W_s\cup U$；由引理~\ref{lem:addone}，$\vv\varepsilon_1,\ldots,\vv\varepsilon_k$ 仍线性无关. 选满 $n$ 个即得一组基.}

\section{第三章　矩阵}

\subsection*{A 组答案}

\ans{ex:mat:A1}{$\vv{AB}=\begin{pmatrix}5&9\\-2&4\end{pmatrix}$，$\vv{BA}=\begin{pmatrix}2&1&5\\-1&1&-2\\1&11&6\end{pmatrix}$，$\vv A^\trans\vv B^\trans=(\vv{BA})^\trans$. 由此也看出 $(\vv{AB})^\trans=\vv B^\trans\vv A^\trans\ne\vv A^\trans\vv B^\trans$（后者甚至与前者不同型）.}

\ans{ex:mat:A2}{$2(1,3)^\trans-(2,-1)^\trans+(0,4)^\trans=(0,11)^\trans$. $\vv e_i^\trans\vv A\vv e_j=a_{ij}$；$\vv A\vv e_j$ 是 $\vv A$ 的第 $j$ 列，$\vv e_i^\trans\vv A$ 是第 $i$ 行.}

\ans{ex:mat:A3}{$\begin{pmatrix}1&n\\0&1\end{pmatrix}$；$\begin{pmatrix}\lambda^n&n\lambda^{n-1}\\0&\lambda^n\end{pmatrix}$；$\begin{pmatrix}\cos n\theta&-\sin n\theta\\ \sin n\theta&\cos n\theta\end{pmatrix}$；第四个矩阵记为 $\vv N$，$\vv N^2=\vv E_{13}$，$n\ge3$ 时 $\vv N^n=\vv O$. 均可用归纳法证明.}

\ans{ex:mat:A4}{$(\vv A+\vv B)^2=\begin{pmatrix}10&8\\12&10\end{pmatrix}$，$\vv A^2+2\vv{AB}+\vv B^2=\begin{pmatrix}16&8\\12&4\end{pmatrix}$，不相等；$(\vv A+\vv B)(\vv A-\vv B)=\begin{pmatrix}-6&4\\-6&6\end{pmatrix}\ne\vv A^2-\vv B^2=\begin{pmatrix}0&4\\-6&0\end{pmatrix}$. 原因是 $\vv{AB}\ne\vv{BA}$.}

\ans{ex:mat:A5}{依次为
$\begin{pmatrix}1&3&-2\\-\frac32&-3&\frac52\\1&1&-1\end{pmatrix}$，
$\begin{pmatrix}1&-1&0\\0&1&-1\\0&0&1\end{pmatrix}$，
$\frac14\begin{pmatrix}3&-2&1\\-2&4&-2\\1&-2&3\end{pmatrix}$，
$\begin{pmatrix}0&0&0&\frac14\\0&0&\frac13&0\\0&\frac12&0&0\\1&0&0&0\end{pmatrix}$.}

\ans{ex:mat:A6}{(1) $\vv X=\begin{pmatrix}2&-23\\0&8\end{pmatrix}$. (2) $\vv X=\begin{pmatrix}1&-1\\-2&4\end{pmatrix}$（注意 $\vv X=\vv B\vv A^{-1}$，逆矩阵乘在右边）.}

\ans{ex:mat:A7}{$r_2-r_1$，$r_3-ar_1$，再 $r_3+r_2$，得阶梯形的后两行为 $(0,a-1,1-a)$ 与 $\big(0,0,(1-a)(a+2)\big)$. $a\ne1$ 且 $a\ne-2$ 时秩为 $3$；$a=1$ 时后两行全为零，秩为 $1$；$a=-2$ 时第二行非零、第三行为零，秩为 $2$.（学完第五章后，可用行列式 $-(a-1)^2(a+2)$ 验证.）}

\ans{ex:mat:A8}{秩为 $3$. 简化阶梯形的主元在第 $1,3,5$ 列，故第 $1,3,5$ 列构成极大无关组；第 $2$ 列 $=2\times$第 $1$ 列，第 $4$ 列 $=$ 第 $1$ 列 $+$ 第 $3$ 列.}

\ans{ex:mat:A9}{$\vv E(1,3)=\begin{pmatrix}0&0&1\\0&1&0\\1&0&0\end{pmatrix}$，逆为自身；$\vv E(2(5))=\diag(1,5,1)$，逆为 $\diag(1,\frac15,1)$；$\vv E(3,1(-2))=\begin{pmatrix}1&0&0\\0&1&0\\-2&0&1\end{pmatrix}$，逆为 $\vv E(3,1(2))$. $\vv E(3,1(-2))\vv A$ 把 $\vv A$ 第 $1$ 行的 $-2$ 倍加到第 $3$ 行；$\vv A\vv E(3,1(-2))$ 把 $\vv A$ 第 $3$ 列的 $-2$ 倍加到第 $1$ 列.}

\ans{ex:mat:A10}{$\begin{pmatrix}2&1\\4&3\end{pmatrix}=\vv E(2,1(2))\,\vv E(1,2(1))\,\vv E(1(2))=\begin{pmatrix}1&0\\2&1\end{pmatrix}\begin{pmatrix}1&1\\0&1\end{pmatrix}\begin{pmatrix}2&0\\0&1\end{pmatrix}$. 写法不唯一.}

\ans{ex:mat:A11}{$(\vv{AB})_{ij}=\sum_ka_{ik}b_{kj}$，$i>j$ 时每一项中 $k<i$ 使 $a_{ik}=0$ 或 $k\ge i>j$ 使 $b_{kj}=0$；$i=j$ 时只剩 $a_{ii}b_{ii}$. 可逆上三角矩阵的逆也是上三角矩阵：对 $(\vv A\mid\vv I)$ 作消元时只需用下面的行消去上面的行.}

\ans{ex:mat:A12}{(1) 用 $(\vv{AB})^\trans=\vv B^\trans\vv A^\trans$ 逐一验证. (2) $(\vv{AB})^\trans=\vv B^\trans\vv A^\trans=\vv{BA}$. (3) $(\vv C^\trans\vv C)^\trans=\vv C^\trans\vv C$.}

\ans{ex:mat:A13}{(1) $\vv x\vv y^\trans$ 的 $(i,i)$ 元为 $x_iy_i$，求和即得. (2) $\tr(\vv A\vv E_{ji})=\vv e_i^\trans\vv A\vv e_j=a_{ij}$，故每个 $a_{ij}=0$. (3) 两边取迹，由命题~\ref{prop:trace}，左边为 $0$，右边为 $n$.}

\ans{ex:mat:A14}{$\vv A(\vv A-\vv I)=2\vv I$，故 $\vv A^{-1}=\frac12(\vv A-\vv I)$；$(\vv A+2\vv I)(\vv A-3\vv I)=\vv A^2-\vv A-6\vv I=-4\vv I$，故 $(\vv A+2\vv I)^{-1}=-\frac14(\vv A-3\vv I)$. $\vv A+\vv I$ 不一定可逆：$\vv A=-\vv I$ 满足条件，$\vv A+\vv I=\vv O$.}

\ans{ex:mat:A15}{左乘 $\vv A^{-1}$. 反例：$\vv A=\begin{pmatrix}1&0\\0&0\end{pmatrix}$，$\vv B=\vv O$，$\vv C=\begin{pmatrix}0&0\\0&1\end{pmatrix}$.}

\ans{ex:mat:A16}{$\begin{pmatrix}\vv A^{-1}&\vv O\\ \vv O&\vv B^{-1}\end{pmatrix}$，$\begin{pmatrix}\vv O&\vv B^{-1}\\ \vv A^{-1}&\vv O\end{pmatrix}$. 所求逆矩阵为 $\begin{pmatrix}0&0&\frac12&0\\0&0&0&\frac15\\3&-2&0&0\\-1&1&0&0\end{pmatrix}$.}

\ans{ex:mat:A17}{记 $\vv B=\begin{pmatrix}2&1\\1&3\end{pmatrix}$，矩阵为 $\begin{pmatrix}\vv I&\vv B\\ \vv O&\vv I\end{pmatrix}$，其平方为 $\begin{pmatrix}\vv I&2\vv B\\ \vv O&\vv I\end{pmatrix}=\begin{pmatrix}1&0&4&2\\0&1&2&6\\0&0&1&0\\0&0&0&1\end{pmatrix}$，$n$ 次方为 $\begin{pmatrix}\vv I&n\vv B\\ \vv O&\vv I\end{pmatrix}$.}

\ans{ex:mat:A18}{$\vv Q\vv C=\begin{pmatrix}5600&8300\\6750&9800\end{pmatrix}$：第 $i$ 行是第 $i$ 个工厂每天的总成本与总销售额. 每天利润分别为 $2700$ 元与 $3050$ 元.}

\ans{ex:mat:A19}{$\vv A=\begin{pmatrix}0&1&1&0\\1&0&1&0\\1&1&0&1\\0&0&1&0\end{pmatrix}$，$\vv A^2=\begin{pmatrix}2&1&1&1\\1&2&1&1\\1&1&3&0\\1&1&0&1\end{pmatrix}$，$\vv A^3=\begin{pmatrix}2&3&4&1\\3&2&4&1\\4&4&2&3\\1&1&3&0\end{pmatrix}$. $(\vv A^2)_{ii}$ 等于城市 $i$ 的航线数；从 $1$ 到 $4$ 恰乘三次航班有 $(\vv A^3)_{14}=1$ 种走法（$1\to2\to3\to4$）.}

\ans{ex:mat:A20}{$\vv x^\trans\vv A\vv x$ 是一个数，等于它的转置 $\vv x^\trans\vv A^\trans\vv x=-\vv x^\trans\vv A\vv x$，故为 $0$. $(\vv A^2)^\trans=(\vv A^\trans)^2=(-\vv A)^2=\vv A^2$.}

\ans{ex:mat:A21}{$\vv P$ 是初等矩阵之积，$\vv{PA}$ 与 $\vv A$ 行等价，因而有相同的简化阶梯形（简化阶梯形的唯一性见第一章章末习题~\ref{ex:elim:C9}）.}

\ans{ex:mat:A22}{存在性由提示中的等式. 唯一性：若 $\vv S_1+\vv K_1=\vv S_2+\vv K_2$（$\vv S_i$ 对称，$\vv K_i$ 反对称），则 $\vv S_1-\vv S_2=\vv K_2-\vv K_1$ 既对称又反对称，只能是零矩阵.}

\subsection*{B 组提示}

\ans{ex:mat:B1}{设 $\vv A=\begin{pmatrix}a&b\\c&d\end{pmatrix}$，用 Cayley–Hamilton 型恒等式 $\vv A^2=(a+d)\vv A-(ad-bc)\vv I$（直接验证）. (1) $\vv A^2=\vv O$ 当且仅当 $a+d=0$ 且 $ad-bc=0$. (2) $\vv A=\pm\vv I$，或 $a+d=0$ 且 $ad-bc=-1$. (3) $\vv A=\vv O$，$\vv A=\vv I$，或 $a+d=1$ 且 $ad-bc=0$.}
\ans{ex:mat:B2}{比较 $\vv{AD}$ 与 $\vv{DA}$ 的 $(i,j)$ 元：$a_ix_{ij}=x_{ij}a_j$.}
\ans{ex:mat:B3}{由 $\vv{AE}_{ij}=\vv E_{ij}\vv A$ 比较元素，得 $\vv A$ 的非对角元为零，且 $a_{ii}=a_{jj}$.}
\ans{ex:mat:B4}{用二项式定理（可交换时成立）：$(\vv A+\vv B)^{p+q}$ 的每一项中 $\vv A$ 的次数 $\ge p$ 或 $\vv B$ 的次数 $\ge q$. 反例：$\vv A=\vv E_{12}$，$\vv B=\vv E_{21}$.}
\ans{ex:mat:B5}{等式：分块对角矩阵的行（列）向量组的极大无关组由两块各自的极大无关组拼成. 不等式：用初等变换把 $\vv A$、$\vv B$ 分别化为标准形，看非零子矩阵. 例子：$\vv A=\vv B=(0)$，$\vv C=(1)$.}
\ans{ex:mat:B6}{把右边与 $\vv I+\vv u\vv v^\trans$ 相乘，利用 $\vv v^\trans\vv u$ 是一个数.}
\ans{ex:mat:B7}{直接验证 $(\vv I_n+\vv{BA})\big(\vv I_n-\vv B(\vv I_m+\vv{AB})^{-1}\vv A\big)=\vv I_n$，展开后把 $\vv B\vv A\vv B$ 写成 $\vv B(\vv{AB})$.}
\ans{ex:mat:B8}{各列都是某个非零列 $\vv u$ 的倍数，倍数组成 $\vv v$. $\vv A^2=\vv u(\vv v^\trans\vv u)\vv v^\trans$，而 $\vv v^\trans\vv u=\tr\vv A$. 于是 $\vv A^k=(\tr\vv A)^{k-1}\vv A$.}
\ans{ex:mat:B9}{用相抵标准形 $\vv A=\vv P\diag(\vv I_r,\vv O)\vv Q$，把 $\vv I_r$ 拆成 $r$ 个 $\vv E_{ii}$. 后半句：秩一矩阵之和的秩不超过个数.}
\ans{ex:mat:B10}{取 $\vv B$ 为 $\vv A$ 列向量组的一个极大无关组排成的矩阵，$\vv C$ 的各列是 $\vv A$ 各列在这组向量下的坐标. 对第 A8 题，$\vv B$ 取第 $1,3,5$ 列，$\vv C$ 是简化阶梯形的前三行.}
\ans{ex:mat:B11}{全 $1$ 矩阵 $\vv J$ 满足 $\vv J^2=3\vv J$，故 $\vv J^n=3^{n-1}\vv J$. 第二个：$\vv N^3=\vv O$，$(2\vv I+\vv N)^n=2^n\vv I+n2^{n-1}\vv N+\binom n22^{n-2}\vv N^2$.}
\ans{ex:mat:B12}{对 $k$ 归纳，计算 $\vv N^k$ 的 $(i,j)$ 元 $\sum_l(\vv N^{k-1})_{il}n_{lj}$ 中哪些项可能非零.}
\ans{ex:mat:B13}{$\tr(\vv A^\trans\vv A)=\sum_{i,j}a_{ij}^2$. 复矩阵不成立：$\vv A=\begin{pmatrix}1&0\\\mathrm i&0\end{pmatrix}$ 满足 $\vv A^\trans\vv A=\vv O$；复矩阵应改用 $\bar{\vv A}^\trans\vv A$.}
\ans{ex:mat:B14}{置换矩阵就是对单位矩阵的行作排列得到的矩阵. $\vv{PA}$ 对 $\vv A$ 的行作相应排列，$\vv{AP}$ 对列作排列.}
\ans{ex:mat:B15}{秩为 $2$. 先用行变换化为阶梯形（记录 $\vv P$），再用列变换（记录 $\vv Q$）. 答案不唯一，可用 $\vv{PAQ}$ 直接检验.}
\ans{ex:mat:B16}{$(\vv I+\vv A)(\vv I-\vv A+\vv A^2-\cdots+(-1)^{k-1}\vv A^{k-1})=\vv I$. 若 $\vv A^2=2\vv A$，则 $(\vv I-\vv A)^2=\vv I$，可逆且逆为自身.}

\section{第四章　线性映射}

\subsection*{A 组答案}

\ans{ex:lmap:A1}{(1) 线性，核为 $\{\zero\}$. (2) 不线性：$\zero\mapsto(1,0)^\trans$. (3) 不线性：$(-1,0)^\trans$ 与 $(1,0)^\trans$ 之和的像不等于像之和. (4) 线性，核为 $\{0\}$. (5) 线性，核为常数多项式全体. (6) 线性，核为 $\{\vv O\}$. (7) 不线性：$\det(2\vv X)=4\det\vv X$.}

\ans{ex:lmap:A2}{第一个：核的基 $(-2,1,0)^\trans$，$(1,0,1)^\trans$；像的基 $(1,2)^\trans$；$3=2+1$. 第二个：核的基 $(-2,1,1,0)^\trans$，$(-1,-1,0,1)^\trans$；像的基 $(1,0,1)^\trans$，$(0,1,1)^\trans$；$4=2+2$.}

\ans{ex:lmap:A3}{$\begin{pmatrix}1&-2&1\\0&3&4\end{pmatrix}$.}

\ans{ex:lmap:A4}{(1) $\begin{pmatrix}0&-1\\1&0\end{pmatrix}$. (2) $\begin{pmatrix}0&1\\1&0\end{pmatrix}$，平方为 $\vv I$：反射两次回到原处. (3) $\frac12\begin{pmatrix}1&1\\1&1\end{pmatrix}$，平方等于自身：投影两次与投影一次相同. (4) $\begin{pmatrix}1&2\\0&1\end{pmatrix}$.}

\ans{ex:lmap:A5}{两组基下的矩阵都是 $\begin{pmatrix}0&1&0\\0&0&2\\0&0&0\end{pmatrix}$：$\mathscr D(x-1)=1$，$\mathscr D(x-1)^2=2(x-1)$.}

\ans{ex:lmap:A6}{$\begin{pmatrix}1&0&0&0\\0&0&1&0\\0&1&0&0\\0&0&0&1\end{pmatrix}$，它的平方是 $\vv I_4$.}

\ans{ex:lmap:A7}{过渡矩阵 $\vv P=\begin{pmatrix}-1&-1&-4\\1&0&6\\1&3&1\end{pmatrix}$. $\vv\alpha$ 在 $\vv\varepsilon$ 下的坐标为 $(-1,2,2)^\trans$，在 $\vv\eta$ 下的坐标为 $(-16,5,3)^\trans$. 可用 $\vv x=\vv P\vv x'$ 检验.}

\ans{ex:lmap:A8}{$\diag(3,1)$. 过渡矩阵 $\vv P=\begin{pmatrix}1&1\\1&-1\end{pmatrix}$，也可以直接看出 $\mathscr T\vv\eta_1=3\vv\eta_1$，$\mathscr T\vv\eta_2=\vv\eta_2$.}

\ans{ex:lmap:A9}{$\mathscr S\mathscr T(a\vv\alpha+b\vv\beta)=\mathscr S(a\mathscr T\vv\alpha+b\mathscr T\vv\beta)=a\mathscr{ST}\vv\alpha+b\mathscr{ST}\vv\beta$. 逆映射：设 $\vv\alpha=\mathscr T\vv\alpha'$，$\vv\beta=\mathscr T\vv\beta'$，则 $\mathscr T^{-1}(a\vv\alpha+b\vv\beta)=\mathscr T^{-1}\mathscr T(a\vv\alpha'+b\vv\beta')=a\mathscr T^{-1}\vv\alpha+b\mathscr T^{-1}\vv\beta$.}

\ans{ex:lmap:A10}{若 $\sum k_i\mathscr T\vv\alpha_i=\zero$，则 $\sum k_i\vv\alpha_i\in\ker\mathscr T=\{\zero\}$. 反之，若 $\mathscr T\vv\alpha=\zero$ 而 $\vv\alpha\ne\zero$，则线性无关组 $\{\vv\alpha\}$ 被映为相关组 $\{\zero\}$.}

\ans{ex:lmap:A11}{(1) 存在：前两个向量是基，$\mathscr T(x,y)^\trans=\big(\frac{x+y}2,\frac{x-y}2\big)^\trans$，它恰好满足第三个条件. (2) 不存在：$(2,0)^\trans=(1,1)^\trans+(1,-1)^\trans$，线性性要求 $\mathscr T(2,0)^\trans=(1,1)^\trans$.}

\ans{ex:lmap:A12}{(1) $5-3=2$. (2) 都不存在：像的维数 $\le3<5$；核的维数 $\ge5-3=2>0$.}

\ans{ex:lmap:A13}{$\mathscr T\vv\alpha=\zero\Rightarrow\mathscr{ST}\vv\alpha=\zero$；$\mathscr{ST}\vv\alpha=\mathscr S(\mathscr T\vv\alpha)$. 由前者 $\rk(\mathscr{ST})=\dim U-\dim\ker(\mathscr{ST})\le\dim U-\dim\ker\mathscr T=\rk\mathscr T$，由后者 $\rk(\mathscr{ST})\le\rk\mathscr S$.}

\ans{ex:lmap:A14}{$\tr$ 是满射（$\tr(c\vv E_{11})=c$），由秩—零化度定理 $\dim\ker\tr=n^2-1$.}

\ans{ex:lmap:A15}{$\tr(\vv P^{-1}\vv{AP})=\tr(\vv{AP}\vv P^{-1})=\tr\vv A$. 两矩阵迹分别为 $4$ 与 $3$，不相似.}

\ans{ex:lmap:A16}{$\vv P^{-1}\vv I\vv P=\vv I$. 由于 $\begin{pmatrix}1&1\\0&1\end{pmatrix}\ne\vv I$，它与 $\vv I$ 不相似.}

\ans{ex:lmap:A17}{矩阵为 $\begin{pmatrix}1&0&0\\1&1&1\\1&2&4\end{pmatrix}$. 依次做 $r_2-r_1$，$r_3-r_1$，$r_3-2r_2$，得 $\begin{pmatrix}1&0&0\\0&1&1\\0&0&2\end{pmatrix}$，有 $3$ 个主元，由定理~\ref{thm:invertible} 矩阵可逆，故 $\mathscr E$ 是同构.（第五章会看到，它的行列式是 Vandermonde 型的.）}

\ans{ex:lmap:A18}{$\mathscr T(x_1,x_2,x_3)^\trans=(x_1+x_3,\ x_1+x_2+2x_3,\ x_2+x_3)^\trans$. $\ker\mathscr T=\spn\big((1,1,-1)^\trans\big)$，$\imm\mathscr T=\spn\big((1,1,0)^\trans,(0,1,1)^\trans\big)$.}

\ans{ex:lmap:A19}{$\vv A_2(\vv A_1\vv x+\vv b_1)+\vv b_2=(\vv A_2\vv A_1)\vv x+(\vv A_2\vv b_1+\vv b_2)$. 仿射变换线性当且仅当把 $\zero$ 映为 $\zero$，即 $\vv b=\zero$.}

\ans{ex:lmap:A20}{第一列是 $\vv e_1$ 转 $\theta$ 后的位置 $(\cos\theta,\sin\theta)^\trans$，第二列同理. 复合``先转 $\beta$ 再转 $\alpha$''的矩阵是 $\vv R(\alpha)\vv R(\beta)$（定理~\ref{thm:composition}），比较 $(1,1)$ 元与 $(2,1)$ 元，得 $\cos(\alpha+\beta)=\cos\alpha\cos\beta-\sin\alpha\sin\beta$，$\sin(\alpha+\beta)=\sin\alpha\cos\beta+\cos\alpha\sin\beta$.}

\subsection*{B 组提示}

\ans{ex:lmap:B1}{取定基后 $\mathscr T\mapsto$ 矩阵是 $L(V,W)$ 到 $M_{m\times n}(\mathbb{F})$ 的同构.}
\ans{ex:lmap:B2}{$\mathscr T^2=\mathscr O$ 即 $\mathscr T(\mathscr T\vv\alpha)=\zero$. 取 $\imm\mathscr T$ 的基 $\vv\beta_1,\ldots,\vv\beta_r$，$\vv\beta_i=\mathscr T\vv\alpha_i$，证明 $\vv\beta_1,\vv\alpha_1,\ldots,\vv\beta_r,\vv\alpha_r$ 线性无关，再用 $\ker\mathscr T$ 中的向量补成基.}
\ans{ex:lmap:B3}{对 $\mathscr T|_{\ker(\mathscr{ST})}$ 用秩—零化度定理：它的像含于 $\ker\mathscr S$，核含于 $\ker\mathscr T$.}
\ans{ex:lmap:B4}{$\mathscr{ST}=\mathscr I$ 说明 $\mathscr T$ 是单射，有限维时由推论~\ref{cor:nullity} 它是双射，于是 $\mathscr S=\mathscr T^{-1}$.}
\ans{ex:lmap:B5}{$\mathscr{JD}(1)=0$. 不矛盾：B4 依赖有限维，$\mathbb{R}[x]$ 是无限维的.}
\ans{ex:lmap:B6}{$\tr(\vv{AX}-\vv{XA})=0$. 后半句：$\vv A^k$ 与 $\vv A$ 可交换.}
\ans{ex:lmap:B7}{逐个计算 $\vv A\vv E_{11}$ 等并取坐标. $\vv X\mapsto\vv{AX}$ 的矩阵为 $\begin{pmatrix}\vv A&\vv O\\ \vv O&\vv A\end{pmatrix}$（在给定的基序下）；$\vv X\mapsto\vv{XA}$ 的矩阵为 $\begin{pmatrix}a\vv I&c\vv I\\ b\vv I&d\vv I\end{pmatrix}$.}
\ans{ex:lmap:B8}{证明 $V=\imm\mathscr P\oplus\ker\mathscr P$：$\vv\alpha=\mathscr P\vv\alpha+(\vv\alpha-\mathscr P\vv\alpha)$，且 $\mathscr P$ 在 $\imm\mathscr P$ 上是恒等. 合并两者的基.}
\ans{ex:lmap:B9}{按提示分解，并验证两部分分别在 $\ker(\mathscr T\mp\mathscr I)$ 中；交为零是因为 $\mathscr T\vv\alpha=\vv\alpha=-\vv\alpha$ 推出 $\vv\alpha=\zero$.}
\ans{ex:lmap:B10}{对 $\mathscr T|_U$ 用秩—零化度定理.}
\ans{ex:lmap:B11}{在基 $\mathrm e^x\cos x$，$\mathrm e^x\sin x$ 下 $\mathscr D$ 的矩阵为 $\begin{pmatrix}1&1\\-1&1\end{pmatrix}$，逆为 $\frac12\begin{pmatrix}1&-1\\1&1\end{pmatrix}$，故 $\int\mathrm e^x\sin x\,\mathrm dx=\frac12\mathrm e^x(\sin x-\cos x)+C$.}
\ans{ex:lmap:B12}{平移的矩阵为 $\begin{pmatrix}1&0&a\\0&1&b\\0&0&1\end{pmatrix}$. 两种顺序的矩阵不相等：先转后移把原点送到 $(1,2)$，先移后转把原点送到 $(-2,1)$.}
\ans{ex:lmap:B13}{$\vv P^{-1}\vv A^k\vv P=(\vv P^{-1}\vv{AP})^k$；转置时用 $\vv Q=(\vv P^{-1})^\trans$.}
\ans{ex:lmap:B14}{两者秩都为 $1$，故相抵. 若相似，则 $\vv P^{-1}\vv E_{11}\vv P$ 是幂等的，而 $\vv E_{12}^2=\vv O\ne\vv E_{12}$. 相抵允许在定义域与陪域中各自换基，相似必须在同一个空间中用同一组基.}
\ans{ex:lmap:B15}{设 $\sum c_i\mathscr T^i\vv\alpha=\zero$，依次施以 $\mathscr T^{k-1},\mathscr T^{k-2},\ldots$. 所求矩阵是主对角线之上一条斜线全为 $1$、其余为 $0$ 的矩阵.}

\section{第五章　行列式}

\subsection*{A 组答案}

\ans{ex:det:A1}{$\tau(315426)=5$，奇排列；$\tau(42531)=7$，奇排列；第三个排列的逆序数为 $0+1+2+\cdots+(n-1)=\frac{n(n-1)}2$，$n\equiv0,1\pmod4$ 时为偶排列，$n\equiv2,3\pmod4$ 时为奇排列.}

\ans{ex:det:A2}{(1) $\tau(3421)=5$，带负号. (2) 按行号排好后列号排列为 $i\,3\,5\,4\,j$，$\{i,j\}=\{1,2\}$. $\tau(13542)=4$，故 $i=1$，$j=2$.}

\ans{ex:det:A3}{$-41$，$-16$，$48$. 第三个：把各列加到第一列后提出 $6$，再化为三角形得 $6\cdot2^3=48$.}

\ans{ex:det:A4}{$6$.}

\ans{ex:det:A5}{$(-1)^{\frac{n(n-1)}2}a_1a_2\cdots a_n$：只有一项非零，它的列号排列是 $n(n-1)\cdots1$.}

\ans{ex:det:A6}{第一列减去第二列、再加上第三列（或按多重线性拆成 $8$ 个行列式，其中有两列相同的都为零），得 $2\begin{vmatrix}a&b&c\\a_1&b_1&c_1\\a_2&b_2&c_2\end{vmatrix}$.}

\ans{ex:det:A7}{第一个行列式等于 $(x-1)^2(x+2)$，根为 $x=1$（二重）与 $x=-2$；第二个是 Vandermonde 行列式 $(2-1)(x-1)(x-2)$，根为 $x=1,2$.}

\ans{ex:det:A8}{$24$，$-3$，$\frac13$，$9$，$9$. 最后一个：$\vv A^*=3\vv A^{-1}$，故 $3\vv A^{-1}-2\vv A^*=-3\vv A^{-1}$，行列式为 $(-3)^3\cdot\frac13=-9$.}

\ans{ex:det:A9}{(1) $(\det\vv A)^2=\det(\vv A^\trans\vv A)=1$. (2) $(\det\vv A)^2=1$. (3) 若 $\det\vv A\ne0$，则 $\vv A$ 可逆，由 $\vv A^2=\vv A$ 得 $\vv A=\vv I$，矛盾.}

\ans{ex:det:A10}{$\begin{pmatrix}d&-b\\-c&a\end{pmatrix}$，$ad-bc\ne0$ 时逆为 $\frac1{ad-bc}\begin{pmatrix}d&-b\\-c&a\end{pmatrix}$. 第二个矩阵行列式为 $1$，伴随矩阵就是逆矩阵：$\begin{pmatrix}-24&18&5\\20&-15&-4\\-5&4&1\end{pmatrix}$.}

\ans{ex:det:A11}{$D=17$，$D_1=17$，$D_2=-17$，$D_3=34$，解为 $(1,-1,2)$.}

\ans{ex:det:A12}{面积为 $\frac12\left|\det\begin{pmatrix}3&1\\1&4\end{pmatrix}\right|=\frac{11}2$. 三点共线：$(2,3)^\trans$ 与 $(6,9)^\trans$ 张成的平行四边形面积为 $0$.}

\ans{ex:det:A13}{平行六面体体积为 $7$，四面体体积为 $\frac76$.}

\ans{ex:det:A14}{系数行列式为 $(5-\lambda)(6-\lambda)(4-\lambda)-4(4-\lambda)-4(6-\lambda)=-(\lambda-2)(\lambda-5)(\lambda-8)$，故 $\lambda=2,5,8$.}

\ans{ex:det:A15}{把第一行换成 $(1,1,1,1)$ 后按第一行展开即得所求之和，计算得 $16$.}

\ans{ex:det:A16}{点 $(x,y)$ 在直线上当且仅当第一行是后两行的``仿射组合''，即三行线性相关. 也可以按第一行展开，得到一个关于 $x,y$ 的一次方程，并验证两点满足它.}

\ans{ex:det:A17}{按定义，非零项必须在左上角 $\vv A$ 中取 $m$ 个元素、在右下角 $\vv B$ 中取 $n$ 个，符号相乘. 第二个等于 $(-1)^{mn}\det\vv A\det\vv B$：把后 $n$ 列逐一移到前面，共 $mn$ 次相邻列交换.}

\ans{ex:det:A18}{$5$（一般地，这类 $n$ 阶行列式等于 $n+1$，见 B1）.}

\ans{ex:det:A19}{$\vv A\vv A^*=(\det\vv A)\vv I$ 两边取行列式得 $\det\vv A^*=(\det\vv A)^{n-1}\ne0$；再由 $\vv A^*=(\det\vv A)\vv A^{-1}$ 得 $(\vv A^*)^{-1}=\frac1{\det\vv A}\vv A$；对 $\vv A^{-1}$ 用同一公式得 $(\vv A^{-1})^*=\det(\vv A^{-1})\vv A=\frac1{\det\vv A}\vv A$.}

\ans{ex:det:A20}{行列式为 $0$，线性相关（$(7,8,9)^\trans=2(4,5,6)^\trans-(1,2,3)^\trans$）；改成 $10$ 后行列式为 $-3$，线性无关.}

\ans{ex:det:A21}{$\det\vv A=\det\vv A^\trans=\det(-\vv A)=(-1)^n\det\vv A=-\det\vv A$，故 $\det\vv A=0$.}

\subsection*{B 组提示}

\ans{ex:det:B1}{由 $D_n=2D_{n-1}-D_{n-2}$，$D_1=2$，$D_2=3$，得 $D_n-D_{n-1}=1$，$D_n=n+1$.}
\ans{ex:det:B2}{第一个：各列加到第一列，提出 $a+(n-1)b$，各行减去第一行，得 $[a+(n-1)b](a-b)^{n-1}$. 第二个：先证递推 $D_n=(x-a)D_{n-1}+a(x+a)^{n-1}$（或利用转置后 $a$ 变 $-a$），得 $\frac12\big[(x+a)^n+(x-a)^n\big]$.}
\ans{ex:det:B3}{第 $i$ 列乘 $\frac1{a_i}$ 后从第一列减去，化为上三角形：$a_1\cdots a_n\big(a_0-\sum_{i=1}^n\frac1{a_i}\big)$.}
\ans{ex:det:B4}{计算 $\begin{pmatrix}\vv I&\zero\\ -\vv v^\trans&1\end{pmatrix}\begin{pmatrix}\vv I+\vv u\vv v^\trans&\vv u\\ \zero^\trans&1\end{pmatrix}\begin{pmatrix}\vv I&\zero\\ \vv v^\trans&1\end{pmatrix}$. B2 的第一个行列式即 $\det\big((a-b)\vv I+b\vv 1\vv 1^\trans\big)$.}
\ans{ex:det:B5}{把 $\vv{AB}$ 的两列按第三章``列是线性组合''的看法展开，用多重线性拆成若干个行列式，其中列号重复的为零. $\vv{BA}$ 的秩 $\le2<3$.}
\ans{ex:det:B6}{必要性：$\det\vv A\det\vv A^{-1}=1$，两者都是整数. 充分性：$\vv A^{-1}=\frac1{\det\vv A}\vv A^*$，$\vv A^*$ 是整数矩阵.}
\ans{ex:det:B7}{可逆时用 $\vv A^*=(\det\vv A)\vv A^{-1}$ 两次. 不可逆时 $\rk\vv A^*\le1$，$n\ge3$ 时 $(\vv A^*)^*=\vv O$；$n=2$ 时直接验证.}
\ans{ex:det:B8}{对分块矩阵先把第二块行加到第一块行，再把第一块列从第二块列减去，化为 $\begin{pmatrix}\vv A+\vv B&\vv O\\ \vv B&\vv A-\vv B\end{pmatrix}$.}
\ans{ex:det:B9}{左乘 $\begin{pmatrix}\vv I&\vv O\\ -\vv{CA}^{-1}&\vv I\end{pmatrix}$. 可交换时 $\det\vv A\det(\vv D-\vv{CA}^{-1}\vv B)=\det(\vv{AD}-\vv{ACA}^{-1}\vv B)=\det(\vv{AD}-\vv{CB})$.}
\ans{ex:det:B10}{两边取行列式：$\det\vv F^n=(\det\vv F)^n=(-1)^n$.}
\ans{ex:det:B11}{按第一行展开得 $D_n=(1+x^2)D_{n-1}-x^2D_{n-2}$，即 $D_n-D_{n-1}=x^2(D_{n-1}-D_{n-2})$，故 $D_n=1+x^2+\cdots+x^{2n}$.}
\ans{ex:det:B12}{把第二、三列分别乘 $\omega^k$，$\omega^{2k}$ 加到第一列，提出公因式；也可直接展开验证第二个等号.}
\ans{ex:det:B13}{从最后一行起，每一行减去上一行，化为对角线全为 $1$ 的上三角形.}
\ans{ex:det:B14}{圆（或直线）的方程可写成 $a(x^2+y^2)+bx+cy+d=0$，系数不全为零. 四点都满足它，即一个以 $(a,b,c,d)$ 为未知量的齐次方程组有非零解.}
\ans{ex:det:B15}{$\rk\vv A=n$ 时 $\vv A^*=(\det\vv A)\vv A^{-1}$ 可逆. $\rk\vv A\le n-2$ 时，由定理~\ref{thm:minorrank} 所有 $n-1$ 阶子式为零，$\vv A^*=\vv O$. $\rk\vv A=n-1$ 时有非零的 $n-1$ 阶子式，故 $\vv A^*\ne\vv O$；又 $\vv A\vv A^*=(\det\vv A)\vv I=\vv O$，$\vv A^*$ 的每一列都属于 $\ker\vv A$，而 $\dim\ker\vv A=n-(n-1)=1$（推论~\ref{cor:nullity}），故 $\rk\vv A^*=1$.}
\ans{ex:det:B16}{$x=1,2,3$ 时有两行相同；每行之和都是 $x+6$，把各列加到第一列可提出因子 $x+6$. $f$ 是四次多项式，首项系数为 $1$，故 $f(x)=(x+6)(x-1)(x-2)(x-3)$，根为 $1,2,3,-6$.}

\section{第六章　线性方程组解的结构}

\subsection*{A 组答案}

\ans{ex:sol:A1}{(1) 基础解系 $(2,5,7,0)^\trans$，$(3,4,0,7)^\trans$，通解为二者的任意线性组合. (2) 基础解系 $(-2,1,0,0)^\trans$，$(1,0,0,1)^\trans$.}

\ans{ex:sol:A2}{(1) $\vv x=\big(\tfrac54,-\tfrac14,0,0\big)^\trans+k_1(3,3,2,0)^\trans+k_2(-3,7,0,4)^\trans$. (2) $\vv x=\big(\tfrac12,0,0,0\big)^\trans+k_1(-1,2,0,0)^\trans+k_2(1,0,2,0)^\trans$.}

\ans{ex:sol:A3}{(1) $5-3=2$ 个. (2) 不一定，要看 $\rk(\vv A\mid\vv b)$ 是否为 $3$. (3) 一定有解（行满秩，推论~\ref{cor:sol-fullrank}），且有无穷多解.}

\ans{ex:sol:A4}{(1) $\rk\vv A=\rk(\vv A\mid\vv b)=2$，有解（且唯一）. (2) $\rk\vv A=1$，$\rk(\vv A\mid\vv b)=2$，无解.}

\ans{ex:sol:A5}{系数和为 $1$ 的 $2\vv\eta_1-\vv\eta_2$，$\frac12(\vv\eta_1+\vv\eta_2)$，$3\vv\eta_1-2\vv\eta_2$ 是 $\vv{Ax}=\vv b$ 的解；系数和为 $0$ 的 $\vv\eta_1-\vv\eta_2$ 是齐次方程组的解；$\vv\eta_1+\vv\eta_2$ 两者都不是（命题~\ref{prop:sol-affine}）.}

\ans{ex:sol:A6}{基础解系含 $3-2=1$ 个向量，$\vv\eta_1-\vv\eta_2=(0,4,4)^\trans$ 是非零解. 通解为 $(1,2,3)^\trans+k(0,1,1)^\trans$.}

\ans{ex:sol:A7}{例如 $x_1-x_2=1$，$x_1+x_3=3$. 做法：求与 $(1,1,-1)^\trans$ 垂直（即满足 $y_1+y_2-y_3=0$）的两个线性无关的系数向量，再代入特解求常数项. 答案不唯一.}

\ans{ex:sol:A8}{$\rk\vv A=2$（第三个法向量是前两个之和，且任意两个不平行），$\rk(\vv A\mid\vv b)=3$：三个平面两两相交成三条平行直线，没有公共点. 常数改为 $3$ 后 $\rk(\vv A\mid\vv b)=2$，三个平面交于一条直线.}

\ans{ex:sol:A9}{例如 $\vv A=\begin{pmatrix}1\\0\end{pmatrix}$，$\vv b=(0,1)^\trans$：$\vv{Ax}=\zero$ 只有零解，$\vv{Ax}=\vv b$ 无解. 若有解，两个解之差是齐次方程组的解，只能为零.}

\ans{ex:sol:A10}{$\vv X=\begin{pmatrix}3&2\\-2&-3\\1&3\end{pmatrix}$.}

\ans{ex:sol:A11}{$\vv X=\begin{pmatrix}a&b\\-a&-b\end{pmatrix}$（$a,b$ 任意），构成二维空间：$\vv X$ 的每一列都属于 $\ker\vv A=\spn\big((1,-1)^\trans\big)$.}

\ans{ex:sol:A12}{$\vv B$ 的非零列是 $\vv{Ax}=\zero$ 的非零解，故 $\det\vv A=7t+21=0$，$t=-3$. 这时 $\rk\vv A=2$，由 $\rk\vv A+\rk\vv B\le3$ 且 $\vv B\ne\vv O$ 得 $\rk\vv B=1$.}

\ans{ex:sol:A13}{通解 $(1,0,0)^\trans+s(-1,1,0)^\trans+t(-1,0,1)^\trans$. 取 $\vv e_1,\vv e_2,\vv e_3$：任一解 $(x_1,x_2,x_3)^\trans=x_1\vv e_1+x_2\vv e_2+x_3\vv e_3$，而 $x_1+x_2+x_3=1$.}

\ans{ex:sol:A14}{第一条直线：$(0,0,1)^\trans+t(1,1,-2)^\trans$. 第二条直线：$x_1-x_3=1$，$x_2=2$（答案不唯一）.}

\ans{ex:sol:A15}{方程组 $1+s=0$，$s=1+t$，$0=1-t$ 无解，两直线不相交（它们的方向也不平行，是异面直线）.}

\ans{ex:sol:A16}{方程 $x_1+x_3=30$，$x_1-x_2=0$，$x_2+x_3=30$（秩为 $2$）. 全部方案为 $x_1=x_2=s$，$x_3=30-s$. $x_1+x_3=30$ 与 $x_1-x_2=0$ 被唯一确定（它们是方程左边的线性组合）；$x_3$ 不被确定：它在基础解系 $(1,1,-1)^\trans$ 上的值不为零.}

\ans{ex:sol:A17}{$\vv L=\begin{pmatrix}1&0&0\\2&1&0\\1&2&1\end{pmatrix}$，$\vv U=\begin{pmatrix}2&1&0\\0&1&1\\0&0&1\end{pmatrix}$.}

\ans{ex:sol:A18}{$\rk(\vv A\mid\vv b)$ 不超过行数 $m=\rk\vv A$，故相容；解空间维数 $n-m>0$，有无穷多解.}

\ans{ex:sol:A19}{两个方向 $(1,-1,0)^\trans$，$(-1,1,2)^\trans$；求与二者都垂直的系数向量，即解齐次方程组，得 $(1,1,0)^\trans$. 平面方程为 $x+y=2$.}

\ans{ex:sol:A20}{解分别为 $(2,0)$ 与 $(1,1)$. 右端的相对变化约为 $\frac{0.0001}{2\sqrt2}\approx3.5\times10^{-5}$，解的相对变化为 $\frac{\sqrt2}2\approx0.71$，二者之比约为 $2\times10^4$.}

\subsection*{B 组提示}

\ans{ex:sol:B1}{把四个方程联立求解，得公共解 $k(-1,1,2,1)^\trans$.}
\ans{ex:sol:B2}{$\vv A^\trans\vv{Ax}=\zero\Rightarrow\vv x^\trans\vv A^\trans\vv A\vv x=0$，即 $\vv{Ax}$ 各分量的平方和为零. 两个齐次方程组同解，解空间维数相同，再用推论~\ref{cor:nullity}.}
\ans{ex:sol:B3}{$\vv A+\vv B$ 的列是 $(\vv A\mid\vv B)$ 的列的线性组合；后者的列由两个极大无关组合并表出.}
\ans{ex:sol:B4}{$(\vv A+\vv I)(\vv A-\vv I)=\vv O$ 给出 $\le n$（Sylvester 不等式的推论）；$(\vv A+\vv I)-(\vv A-\vv I)=2\vv I$ 给出 $\ge n$（B3）.}
\ans{ex:sol:B5}{由定理~\ref{thm:sylvester} 证明中的等式 $\rk\vv B=\rk(\vv{AB})+\dim(\ker\vv A\cap\imm\vv B)$.}
\ans{ex:sol:B6}{设 $k_0\vv\eta^*+\sum k_i(\vv\eta^*+\vv\xi_i)=\zero$，左乘 $\vv A$ 得 $(k_0+\sum k_i)\vv b=\zero$，再用 $\vv\xi_i$ 线性无关.}
\ans{ex:sol:B7}{$9$ 个未知数，$6$ 个方程的秩为 $5$，自由度为 $4$，至少再公开 $4$ 个单元格. 公开左上角的 $2\times2$ 块即可；但若公开的 $4$ 格中有一行的全部 $3$ 格，那一行有一格是多余的，全表仍不能确定.}
\ans{ex:sol:B8}{三行满足 $r_1+r_2=r_3$，故条件为 $b_1+b_2=b_3$.}
\ans{ex:sol:B9}{解空间是一维的，$(1,\ldots,1)^\trans$ 是非零解. 通解为 $k(1,\ldots,1)^\trans$.}
\ans{ex:sol:B10}{$\rk\vv A+\rk\vv B\le n$（B4 中的推论）且 $\ge\rk(\vv A+\vv B)=n$（B3）.}
\ans{ex:sol:B11}{同解时，$\vv B$ 的每一行在 $\vv A$ 的解空间上恒为零，由可确定性判据（定理~\ref{thm:sol-determinable} 的 (2)$\Rightarrow$(3)）它是 $\vv A$ 各行的线性组合.}
\ans{ex:sol:B12}{必要性：对 $\vv b=\vv e_1,\ldots,\vv e_m$ 分别取解，排成 $\vv C$. 充分性：$\vv{A}(\vv C\vv b)=\vv b$.}
\ans{ex:sol:B13}{$\vv A\vv A^*=(\det\vv A)\vv I=\vv O$，$\vv A^*$ 的每一列都在一维解空间中.}
\ans{ex:sol:B14}{$x_1-x_4$ 是第一、二个方程之差，被确定为 $-3$；后两个量都不被确定：找不到系数 $a,b,c$ 使 $a,b,c$ 倍的三个方程之和的左边恰为所求的量.}
\ans{ex:sol:B15}{对 $\vv B|_{\ker(\vv{AB})}$ 用秩—零化度定理：它的核是 $\ker\vv B$，像含于 $\ker\vv A$.}

\section{第七章　多项式}

\subsection*{A 组答案}

\ans{ex:poly:A1}{(1) 商 $x^3-x^2+1$，余式 $-x-2$. (2) 商 $2x^3+x^2+6x+7$，余式 $20$（$=f(2)$）.}

\ans{ex:poly:A2}{依次为 $1$，$1$，$100x-99$. 前两个用 $x^2\equiv1$，$x^2\equiv-1$；第三个设余式为 $ax+b$，由 $x^{100}-(ax+b)$ 以 $1$ 为二重根，在 $x=1$ 处函数值与导数都为零.}

\ans{ex:poly:A3}{$\omega^4=\omega$，$\omega^2=-1-\omega$，代入得 $(1-a)\omega+(b-a)=0$，而 $\omega$ 不是实数，故 $a=b=1$. 验证：$x^4+x^2+1=(x^2+x+1)(x^2-x+1)$.}

\ans{ex:poly:A4}{(1) $x+3$. (2) $x-1$（$x^3-2x^2+x=x(x-1)^2$，$x^2-1=(x-1)(x+1)$）.}

\ans{ex:poly:A5}{$f$ 与 $g$ 的公因式整除 $g-hf$；$f$ 与 $g-hf$ 的公因式整除 $(g-hf)+hf=g$. 两组公因式相同，最大公因式也相同.}

\ans{ex:poly:A6}{$x^2+1=(x-1)(x+1)+2$，故 $u=\frac12$，$v=-\frac12(x-1)$.}

\ans{ex:poly:A7}{有理根只可能是 $\pm1,\pm2,\pm\frac12$. 得 $2x^3-5x^2+x+2=(x-1)(x-2)(2x+1)$，有理根为 $1,2,-\frac12$.}

\ans{ex:poly:A8}{(1) 在 $\mathbb{Q}$ 与 $\mathbb{R}$ 上为 $(x^2+x+1)(x^2-x+1)$；在 $\mathbb{C}$ 上为 $(x-\omega)(x-\bar\omega)(x+\omega)(x+\bar\omega)$，$\omega=\mathrm e^{2\pi\mathrm i/3}$. (2) 在 $\mathbb{Q}$ 与 $\mathbb{R}$ 上为 $(x-1)(x+1)(x^2+x+1)(x^2-x+1)$；在 $\mathbb{C}$ 上为 $\prod_{k=0}^5(x-\mathrm e^{k\pi\mathrm i/3})$. (3) 在 $\mathbb{Q}$ 上不可约；在 $\mathbb{R}$ 上为 $(x-\sqrt[3]2)(x^2+\sqrt[3]2x+\sqrt[3]4)$；在 $\mathbb{C}$ 上为 $(x-\sqrt[3]2)(x-\sqrt[3]2\,\omega)(x-\sqrt[3]2\,\omega^2)$.}

\ans{ex:poly:A9}{$f(1)=f'(1)=0$，$f''(1)=20\ne0$，是二重根. 也可以连续两次用 $x-1$ 作综合除法.}

\ans{ex:poly:A10}{有重根当且仅当与导数 $3x^2-3$ 有公共根，即 $f(1)=0$ 或 $f(-1)=0$. $a=2$ 时重根为 $1$；$a=-2$ 时重根为 $-1$.}

\ans{ex:poly:A11}{$\sigma_1=2$，$\sigma_2=3$，$\sigma_3=4$. 所求依次为 $2$；$\sigma_1^2-2\sigma_2=-2$；$\sigma_2/\sigma_3=\frac34$.}

\ans{ex:poly:A12}{$(x^2-2x-1)(x^2+1)=x^4-2x^3-2x-1$. （$1+\sqrt2$ 的共轭 $1-\sqrt2$、$\mathrm i$ 的共轭 $-\mathrm i$ 也必须是根，见第七章章末习题~\ref{ex:poly:B10}.）}

\ans{ex:poly:A13}{$(x-2)(x^2-2x+2)=x^3-4x^2+6x-4$.}

\ans{ex:poly:A14}{$f(\vv A)=(\vv A-\vv I)^2=\vv O$. 于是 $\vv A^2=2\vv A-\vv I$，$\vv A^3=3\vv A-2\vv I$，$g(\vv A)=2\vv A-2\vv I=\begin{pmatrix}0&4\\0&0\end{pmatrix}$.}

\ans{ex:poly:A15}{$f(x)=2x^2-3x+1$（四个点恰好落在一条抛物线上，插值多项式的三次项系数为 $0$）.}

\ans{ex:poly:A16}{$x^n-1=(x-1)(x^{n-1}+x^{n-2}+\cdots+1)$. 一般地 $x^n-a^n=(x-a)(x^{n-1}+ax^{n-2}+\cdots+a^{n-1})$；也可由余数定理：$x=a$ 时 $x^n-a^n=0$.}

\ans{ex:poly:A17}{$f(x)=2-2(x-1)+2(x-1)^3+(x-1)^4$.}

\ans{ex:poly:A18}{设余式为 $ax+b$，由 $a+b=3$，$2a+b=5$ 得余式为 $2x+1$.}

\ans{ex:poly:A19}{$60.5v^2+55v-100=0$，即 $121v^2+110v-200=0$，正根 $v=\frac{10}{11}$，内部收益率 $r=10\%$.}

\ans{ex:poly:A20}{比较次数得 $\deg f=\deg g$，$f=hg$ 中 $h$ 为常数，且 $f\ne0$ 时 $h\ne0$.}

\ans{ex:poly:A21}{设 $h$ 是 $f_1,g_1$ 的公因式，则 $dh$ 是 $f,g$ 的公因式，故 $dh\mid d$，$h$ 是常数.}

\subsection*{B 组提示}

\ans{ex:poly:B1}{由 Bézout 等式 $uf+vg=1$，$u'f+v'h=1$，两式相乘.}
\ans{ex:poly:B2}{用提示中的等式反复降低指数，指数的变化与整数的辗转相除 $\gcd(m,n)=\gcd(m-n,n)$ 完全平行.}
\ans{ex:poly:B3}{$f/\gcd(f,f')=(x-1)(x+2)$. 一般地，$p_i^{k_i}$ 在 $f'$ 中恰好以 $p_i^{k_i-1}$ 出现（需要用到 $p_i$ 与 $p_i'$ 互素）.}
\ans{ex:poly:B4}{$p$ 不可约，由引理~\ref{lem:primediv} 即得. 反例：$p=x^2$，$f=x$.}
\ans{ex:poly:B5}{两边都是次数小于 $n$ 的多项式，在 $x_1,\ldots,x_n$ 处取值相同（$\ell_i(x_j)$ 的性质），用推论~\ref{cor:rootcount}.}
\ans{ex:poly:B6}{$f(x)-f(0)$ 以 $0,1,2,\ldots$ 为根，有无穷多个根.}
\ans{ex:poly:B7}{三次多项式可约当且仅当有一次因式，即有有理根；可能的有理根只有 $\pm1$，都不是根. 四次多项式可能分解为两个二次因式之积而没有有理根，例如 $(x^2+1)^2$.}
\ans{ex:poly:B8}{设余式为 $ax+b$，代入 $x=1,2$：$(2^n-1)x+2-2^n$.}
\ans{ex:poly:B9}{由 $f(2)=4$ 得 $c=1$，$f(x)=x^2$.}
\ans{ex:poly:B10}{$(x-a)^2-2b^2$ 是有理系数的，用它除 $f$，余式 $rx+s$ 在 $a+b\sqrt2$ 处为零，由 $\sqrt2$ 无理得 $r=s=0$.}
\ans{ex:poly:B11}{对整数 $n$，$f(n)$ 与 $f(n\bmod2)$ 的奇偶性相同（$n-m\mid f(n)-f(m)$）.}
\ans{ex:poly:B12}{依次为 $x^2+1$，$x-2$，$(x-2)^2$.}
\ans{ex:poly:B13}{$N'(v)=C_1+2C_2v+\cdots>0$. 由 $N(0)=-C_0<0$ 与 $N(1)=\sum C_i-C_0>0$，用零点定理.}
\ans{ex:poly:B14}{$\binom xk$ 在整数点取整数值；把 $f$ 写成 $\sum c_k\binom xk$，依次代入 $x=0,1,\ldots,n$ 解出 $c_k$，它们都是整数.}
\ans{ex:poly:B15}{核中的多项式有 $n$ 个不同的根而次数小于 $n$，只能是零.}
\ans{ex:poly:B16}{$\gcd(p,f)$ 整除 $p$，只能是 $1$ 或 $p$ 的常数倍.}

\section{第八章　特征值与特征向量}

\subsection*{A 组答案}

\ans{ex:eig:A1}{以下``属于 $\lambda$''后给出特征子空间的一组基，特征向量是它们的非零线性组合.
\begin{enumerate}
	\item $\lambda=3$：$(1,0)^\trans$；$\lambda=2$：$(1,-1)^\trans$.
	\item $\lambda=3$：$(1,1)^\trans$；$\lambda=-1$：$(1,-1)^\trans$.
	\item $\lambda=\mathrm i$：$(1,\mathrm i)^\trans$；$\lambda=-\mathrm i$：$(1,-\mathrm i)^\trans$（在 $\mathbb{R}$ 上没有特征值）.
	\item $\lambda=2$：$(1,1)^\trans$；$\lambda=3$：$(2,1)^\trans$.
	\item $\lambda=2$（二重）：$(1,0,0)^\trans$，$(0,1,1)^\trans$；$\lambda=0$：$(0,1,-1)^\trans$.
	\item $\lambda=5$：$(1,1,1)^\trans$；$\lambda=-1$（二重）：$(1,-1,0)^\trans$，$(1,0,-1)^\trans$.
\end{enumerate}}

\ans{ex:eig:A2}{除第 3 个（在 $\mathbb{R}$ 上没有特征值）外都可对角化，$\vv P$ 取 A1 中各特征向量为列. 例如第 4 个：$\vv P=\begin{pmatrix}1&2\\1&1\end{pmatrix}$，$\vv P^{-1}\vv{AP}=\diag(2,3)$；第 6 个：$\vv P=\begin{pmatrix}1&1&1\\1&-1&0\\1&0&-1\end{pmatrix}$，$\vv P^{-1}\vv{AP}=\diag(5,-1,-1)$.}

\ans{ex:eig:A3}{第 1 个可以（两个不同的特征值）；第 2 个不可以（$2$ 的几何重数为 $1$）；第 3 个不可以（$2$ 是三重特征值，$2\vv I-\vv A$ 的秩为 $1$，几何重数为 $2$）；第 4 个可以（$1$ 的代数重数与几何重数都是 $2$）.}

\ans{ex:eig:A4}{$\det\vv A=6$，$\tr\vv A=6$；$\vv A^{-1}$ 的特征值为 $1,\frac12,\frac13$；$\vv A^2-2\vv A+\vv I=(\vv A-\vv I)^2$ 的特征值为 $0,1,4$；$\det\vv A^*=(\det\vv A)^2=36$；$\det(\vv A^2+\vv I)=2\cdot5\cdot10=100$.}

\ans{ex:eig:A5}{特征多项式为 $\lambda^2-5\lambda+6$，特征值 $2,3$；第二种情形为 $\lambda^2-2\lambda+5$，特征值 $1\pm2\mathrm i$.}

\ans{ex:eig:A6}{$\det(\lambda\vv I-\vv A^\trans)=\det\big((\lambda\vv I-\vv A)^\trans\big)$. 特征向量一般不同：$\vv A=\begin{pmatrix}1&1\\0&2\end{pmatrix}$ 属于 $1$ 的特征向量为 $(1,0)^\trans$，$\vv A^\trans$ 属于 $1$ 的特征向量为 $(1,-1)^\trans$.}

\ans{ex:eig:A7}{$\lambda\vv I-\vv A$ 仍是三角矩阵，行列式为 $\prod(\lambda-a_{ii})$.}

\ans{ex:eig:A8}{若 $\vv{Ax}=\lambda\vv x$，$\vv x\ne\zero$，则 $\vv A^k\vv x=\lambda^k\vv x$. 由 $\vv A^k=\vv O$ 得 $\lambda^k=0$；由 $\vv A^2=\vv A$ 得 $\lambda^2=\lambda$；由 $\vv A^2=\vv I$ 得 $\lambda^2=1$.}

\ans{ex:eig:A9}{$\vv A=\vv P\diag(2,3)\vv P^{-1}$，$\vv P=\begin{pmatrix}1&2\\1&1\end{pmatrix}$，
$\vv A^n=\begin{pmatrix}2\cdot3^n-2^n&2^{n+1}-2\cdot3^n\\3^n-2^n&2^{n+1}-3^n\end{pmatrix}$.}

\ans{ex:eig:A10}{特征值 $1$ 与 $0.7$. 稳态分布 $\vv\pi=(\frac23,\frac13)^\trans$. $\vv P^n\to\begin{pmatrix}\frac23&\frac23\\ \frac13&\frac13\end{pmatrix}$：两列相同，说明无论开始时顾客在哪一家，长期看都以同样的概率分布在两家，初始状态被``遗忘''了.}

\ans{ex:eig:A11}{$\vv A\vv\xi=(-1,a+2,b+1)^\trans=\lambda(1,1,-1)^\trans$，得 $\lambda=-1$，$a=-3$，$b=0$.}

\ans{ex:eig:A12}{特征多项式 $(\lambda-1)^2(\lambda+1)$. 可对角化当且仅当 $\rk(\vv I-\vv A)=1$，即 $x=-1$.}

\ans{ex:eig:A13}{$\vv A^2=\begin{pmatrix}a^2+bc&b(a+d)\\c(a+d)&d^2+bc\end{pmatrix}$，代入直接验证.}

\ans{ex:eig:A14}{依次为 $(x-1)(x-2)$，$(x-1)^2(x-2)$，$x^3$，$x(x-3)$.}

\ans{ex:eig:A15}{相似. 取 $\vv P=\begin{pmatrix}0&1\\1&0\end{pmatrix}$（交换两个基向量的次序），$\vv P^{-1}\begin{pmatrix}1&1\\0&1\end{pmatrix}\vv P=\begin{pmatrix}1&0\\1&1\end{pmatrix}$.}

\ans{ex:eig:A16}{$\vv A\vv e_1=a_{11}\vv e_1$，$\vv A\vv e_2=a_{12}\vv e_1+a_{22}\vv e_2$，……，$\vv A\vv e_k$ 只含 $\vv e_1,\ldots,\vv e_k$.}

\ans{ex:eig:A17}{特征根 $2,-1$，$a_k=\frac13\big(2^k-(-1)^k\big)$.}

\ans{ex:eig:A18}{矩阵 $\begin{pmatrix}0&1.1\\0.6&0.5\end{pmatrix}$ 的特征值为 $1.1$ 与 $-0.6$. 长期年增长率为 $10\%$；属于 $1.1$ 的特征向量为 $(1,1)^\trans$，幼体与成年个体之比稳定在 $1:1$.}

\ans{ex:eig:A19}{特征多项式 $\lambda^2-(\tr\vv A)\lambda+\det\vv A$ 没有实根，当且仅当判别式 $(\tr\vv A)^2-4\det\vv A<0$.}

\ans{ex:eig:A20}{投影：方向 $(1,2)^\trans$ 上的向量不动（特征值 $1$），垂直方向 $(2,-1)^\trans$ 上的向量变为零（特征值 $0$），矩阵 $\frac15\begin{pmatrix}1&2\\2&4\end{pmatrix}$. 反射：$(1,2)^\trans$ 属于 $1$，$(2,-1)^\trans$ 属于 $-1$，矩阵 $\frac15\begin{pmatrix}-3&4\\4&3\end{pmatrix}$.}

\subsection*{B 组提示}

\ans{ex:eig:B1}{$\vv A\vv u=(\vv v^\trans\vv u)\vv u$；$\rk\vv A=1$，$0$ 的几何重数为 $n-1$. 特征值为 $\vv v^\trans\vv u$ 与 $0$（$n-1$ 重）. $\vv v^\trans\vv u=0$ 时 $0$ 是 $n$ 重特征值，但几何重数只有 $n-1$.}
\ans{ex:eig:B2}{$\vv A\vv B\vv p=\vv B\vv A\vv p=\lambda\vv B\vv p$，而 $V_\lambda$ 是一维的. 后半句：用 Lagrange 插值构造多项式 $f$ 使 $f(\lambda_i)$ 等于 $\vv B$ 的对应特征值.}
\ans{ex:eig:B3}{$\vv P^{-1}f(\vv A)\vv P=f(\vv D)$. 逆命题不成立：$\vv A=\begin{pmatrix}0&1\\0&0\end{pmatrix}$ 不可对角化，$\vv A^2=\vv O$ 可以.}
\ans{ex:eig:B4}{$\vv A|_W$ 的最小多项式整除 $\vv A$ 的最小多项式，用定理~\ref{thm:minpolydiag}.}
\ans{ex:eig:B5}{取 $\vv v\ne\zero$，证明 $\vv v,\vv{Av}$ 线性无关，在这组基下 $\vv A$ 的矩阵为 $\begin{pmatrix}0&-1\\1&0\end{pmatrix}$.}
\ans{ex:eig:B6}{设 $\vv{Ax}=\lambda\vv x$，取 $|x_i|$ 最大的分量，看第 $i$ 个方程 $(\lambda-a_{ii})x_i=\sum_{j\ne i}a_{ij}x_j$.}
\ans{ex:eig:B7}{$\vv A(1,\ldots,1)^\trans=s(1,\ldots,1)^\trans$. 可逆时 $s\ne0$，$\vv A^{-1}$ 每行之和为 $\frac1s$.}
\ans{ex:eig:B8}{特征值为 $a+b\omega^k+c\omega^{2k}$（$k=0,1,2$），行列式是三者之积.}
\ans{ex:eig:B9}{特征根 $2$ 是二重的，通项形如 $(c_1+c_2k)2^k$，得 $a_k=(k+1)2^k$. $k2^k$ 来自 Jordan 块的幂（第八章选读节）.}
\ans{ex:eig:B10}{对三角化（或可对角化）的情形看对角元. 特征值为 $0,\sqrt3,-\sqrt3$.}
\ans{ex:eig:B11}{$x^3-x=x(x-1)(x+1)$，最小多项式整除它.}
\ans{ex:eig:B12}{$2$ 是二重特征值，需 $\rk(2\vv I-\vv A)=1$，条件为 $a=0$（$b,c$ 任意）.}
\ans{ex:eig:B13}{转移矩阵 $\begin{pmatrix}0&0&1\\ \frac12&0&0\\ \frac12&1&0\end{pmatrix}$，稳态分布 $\frac15(2,1,2)^\trans$.}
\ans{ex:eig:B14}{特征值 $\pm1$，$\mathrm e^{t\vv A}=\begin{pmatrix}\cosh t&\sinh t\\ \sinh t&\cosh t\end{pmatrix}$. 对 $t$ 求导验证.}
\ans{ex:eig:B15}{稳态分布 $(\frac13,\frac13,\frac13)^\trans$；特征值 $1,\frac14,\frac14$，偏离部分每一步缩小为原来的 $\frac14$.}
\ans{ex:eig:B16}{几何重数 $n-\rk(\lambda\vv I-\vv A)$ 不超过代数重数 $k$.}
\ans{ex:eig:B17}{得分（按判决 $1,2,3,4$）约为：$d=0.5$ 时 $(0.169,0.212,0.265,0.354)$；$d=0.85$ 时 $(0.128,0.183,0.261,0.428)$；$d=0.95$ 时 $(0.119,0.175,0.258,0.448)$；$d=1$ 时恰为 $\frac1{35}(4,6,9,16)$. 排序不变；$d$ 越大，随机跳转越少，得分越集中到引证链末端的判决 $4$ 上.}

\subsection*{C 组证明要点（正文引用的题目）}

\ans{ex:eig:C1}{直接计算两个分块乘积：
	\[\begin{pmatrix}\vv I_m&-\vv A\\ \vv O&\lambda\vv I_n\end{pmatrix}\begin{pmatrix}\lambda\vv I_m&\vv A\\ \vv B&\vv I_n\end{pmatrix}=\begin{pmatrix}\lambda\vv I_m-\vv{AB}&\vv O\\ \lambda\vv B&\lambda\vv I_n\end{pmatrix},\]\[
	\begin{pmatrix}\vv I_m&\vv O\\ -\vv B&\lambda\vv I_n\end{pmatrix}\begin{pmatrix}\lambda\vv I_m&\vv A\\ \vv B&\vv I_n\end{pmatrix}=\begin{pmatrix}\lambda\vv I_m&\vv A\\ \vv O&\lambda\vv I_n-\vv{BA}\end{pmatrix}.\]
	记中间矩阵的行列式为 $D$. 对两式取行列式（第五章章末习题~\ref{ex:det:A17}），得 $\lambda^nD=\lambda^n\det(\lambda\vv I_m-\vv{AB})$ 与 $\lambda^nD=\lambda^m\det(\lambda\vv I_n-\vv{BA})$，因此 $\lambda^n\det(\lambda\vv I_m-\vv{AB})=\lambda^m\det(\lambda\vv I_n-\vv{BA})$. 两边作为 $\lambda$ 的多项式相等，因子 $\lambda^n$、$\lambda^m$ 只影响根 $0$ 的重数，所以两个特征多项式的非零根（连同重数）相同.}

\section{第九章　内积空间}

\subsection*{A 组答案}

\ans{ex:ip:A1}{$\|\vv\alpha\|=\|\vv\beta\|=3$，$\langle\vv\alpha,\vv\beta\rangle=4$，夹角 $\arccos\frac49\approx63.6^\circ$，$\|\vv\alpha-\vv\beta\|=\sqrt{10}$.}

\ans{ex:ip:A2}{(1) 是. (2) 不是：$\vv x=(1,-1)^\trans$ 时 $\langle\vv x,\vv x\rangle=-2$. (3) 是：$\langle\vv x,\vv x\rangle=x_1^2+(x_1+x_2)^2$. (4) 不是：$\langle\vv e_2,\vv e_2\rangle=-1$.}

\ans{ex:ip:A3}{(1) $\frac1{\sqrt3}(1,1,1)^\trans$，$\frac1{\sqrt2}(0,1,-1)^\trans$，$\frac1{\sqrt6}(-2,1,1)^\trans$. (2) $\frac1{\sqrt2}(1,1,0,0)^\trans$，$\frac1{\sqrt6}(1,-1,2,0)^\trans$，$\frac1{2\sqrt3}(-1,1,1,3)^\trans$.}

\ans{ex:ip:A4}{解 $x_1+x_2+x_3=0$，$x_1-x_3=0$，得 $\pm\frac1{\sqrt6}(1,-2,1)^\trans$.}

\ans{ex:ip:A5}{逐对验证内积；坐标为 $\langle\vv\alpha,\vv e_i\rangle$，即 $(5,4,-2)^\trans$.}

\ans{ex:ip:A6}{两个生成向量已经正交，投影为 $\frac32(1,1,0,0)^\trans+\frac72(0,0,1,1)^\trans=\big(\frac32,\frac32,\frac72,\frac72\big)^\trans$，
	
	距离为 $$\|(-\frac12,\frac12,-\frac12,\frac12)^\trans\|=1$$.}

\ans{ex:ip:A7}{正规方程 $5a+10b=15$，$10a+30b=42$，得 $y=0.6+1.2x$，残差平方和为 $1.6$.}

\ans{ex:ip:A8}{第二列须与第一、三列正交且为单位向量：$a=c$，$b=a$，$3a^2=1$. 故 $a=b=c=\frac1{\sqrt3}$ 或 $a=b=c=-\frac1{\sqrt3}$.}

\ans{ex:ip:A9}{展开 $\big\langle\sum\vv\alpha_i,\sum\vv\alpha_j\big\rangle$，交叉项都为零.}

\ans{ex:ip:A10}{由三角不等式 $\|\vv\alpha\|\le\|\vv\alpha-\vv\beta\|+\|\vv\beta\|$，再交换 $\vv\alpha,\vv\beta$.}

\ans{ex:ip:A11}{(1) 对 $(a_1,\ldots,a_n)$ 与 $(1,\ldots,1)$ 用 Cauchy–Schwarz 不等式. (2) 对 $(\sqrt{a_1},\ldots,\sqrt{a_n})$ 与 $\big(\frac1{\sqrt{a_1}},\ldots,\frac1{\sqrt{a_n}}\big)$ 用之.}

\ans{ex:ip:A12}{$\|1\|=1$，$\|x\|=\frac1{\sqrt3}$，$\langle1,x\rangle=\frac12$，夹角 $\frac\pi6$. 正交化并单位化得 $1$，$\sqrt3(2x-1)$.}

\ans{ex:ip:A13}{$\ker\vv A$ 的标准正交基：$\frac1{\sqrt3}(1,-1,1,0)^\trans$，$\frac1{\sqrt{15}}(-1,-2,-1,3)^\trans$；行空间的标准正交基：$\frac1{\sqrt3}(1,1,0,1)^\trans$，$\frac1{\sqrt{15}}(-2,1,3,1)^\trans$. 两组之间两两内积为零，维数 $2+2=4$，故互为正交补.}

\ans{ex:ip:A14}{$\vv Q=\begin{pmatrix}\frac1{\sqrt2}&\frac1{\sqrt6}\\ \frac1{\sqrt2}&-\frac1{\sqrt6}\\0&\frac2{\sqrt6}\end{pmatrix}$，$\vv R=\begin{pmatrix}\sqrt2&\frac1{\sqrt2}\\0&\sqrt{\frac32}\end{pmatrix}$.}

\ans{ex:ip:A15}{$\vv H=\frac19\begin{pmatrix}7&-4&-4\\-4&1&-8\\-4&-8&1\end{pmatrix}$；$\vv H\vv u=-\vv u$；$(2,-1,0)^\trans$ 与 $\vv u$ 正交，$\vv H(2,-1,0)^\trans=(2,-1,0)^\trans$.}

\ans{ex:ip:A16}{第一个行列式为 $1$，是旋转，旋转角 $\arctan\frac43\approx53.1^\circ$；第二个行列式为 $-1$，是关于直线 $y=\frac x2$（方向 $(2,1)^\trans$）的反射.}

\ans{ex:ip:A17}{设计矩阵只有一列 $\vv x$，正规方程为 $(\vv x^\trans\vv x)b=\vv x^\trans\vv y$.}

\ans{ex:ip:A18}{投影为 $2(1,1,1)^\trans$，距离为 $\|(-1,0,1)^\trans\|=\sqrt2$.}

\ans{ex:ip:A19}{过渡矩阵的第 $j$ 列是新基第 $j$ 个向量在旧（标准正交）基下的坐标，坐标的标准内积等于向量的内积，故各列标准正交.}

\ans{ex:ip:A20}{内积打分：$6,1,4$，排序为甲（第一条）$>$ 丙 $>$ 乙. 余弦相似度：约 $0.926$，$0.140$，$0.943$，排序为丙 $>$ 甲 $>$ 乙. 第一条法条更长，内积占了便宜.}

\ans{ex:ip:A21}{$\vv P^\trans=(\vv u^\trans)^\trans\vv u^\trans=\vv P$，$\vv P^2=\vv u(\vv u^\trans\vv u)\vv u^\trans=\vv P$. $\vv{Px}=\langle\vv x,\vv u\rangle\vv u$ 正是定理~\ref{thm:projformula} 中 $k=1$ 的公式. 一般非零 $\vv u$ 时投影矩阵为 $\dfrac{\vv u\vv u^\trans}{\vv u^\trans\vv u}$.}

\subsection*{B 组提示}

\ans{ex:ip:B1}{展开两个平方相减. 取 $\vv x=(1,0)^\trans$，$\vv y=(0,1)^\trans$ 检验平行四边形法则 $\|\vv x+\vv y\|^2+\|\vv x-\vv y\|^2=2\|\vv x\|^2+2\|\vv y\|^2$，左边为 $8$，右边为 $4$.}
\ans{ex:ip:B2}{对 $\vv\alpha-\sum\langle\vv\alpha,\vv e_i\rangle\vv e_i$ 用勾股定理.}
\ans{ex:ip:B3}{证明 $\vv x-\vv{Px}$ 与 $\imm\vv P$ 正交：$\langle\vv x-\vv{Px},\vv{Py}\rangle=\vv y^\trans\vv P^\trans(\vv I-\vv P)\vv x=0$. 只有幂等时是沿 $\ker\vv P$ 向 $\imm\vv P$ 的``斜投影''，例如 $\begin{pmatrix}1&1\\0&0\end{pmatrix}$.}
\ans{ex:ip:B4}{目标函数是 $\|\vv y-\vv A\vv\beta\|_W^2$，其中 $\langle\vv x,\vv y\rangle_W=\vv x^\trans\vv W\vv y$ 是加权内积；残差在这个内积下与各列正交.}
\ans{ex:ip:B5}{对 $\ln y=\ln c+kt$ 作最小二乘，约得 $k\approx0.30$，$c\approx2.0$. 它最小化的是对数残差的平方和，相当于对相对误差加权.}
\ans{ex:ip:B6}{$\vv G=\vv A^\trans\vv A$（$\vv A$ 以 $\vv\alpha_i$ 为列），用第十一章命题~\ref{prop:gram} 或直接 $\vv x^\trans\vv G\vv x=\|\vv A\vv x\|^2$. $k=2$ 时 $\det\vv G=\|\vv\alpha\|^2\|\vv\beta\|^2\sin^2\theta$.}
\ans{ex:ip:B7}{距离是 $\vv x_0-\vv x_1$（$\vv x_1$ 为超平面上任一点）在法向 $\vv a$ 上投影的长度.}
\ans{ex:ip:B8}{设 $\vv{Qz}=\lambda\vv z$，考察 $\bar{\vv z}^\trans\vv Q^\trans\vv Q\vv z$.}
\ans{ex:ip:B9}{用积化和差逐对计算积分. 投影为 $2\sin x-\sin2x$（偶函数的系数为零）.}
\ans{ex:ip:B10}{全部最小二乘解为 $\hat{\vv x}+\ker\vv A$，把 $\hat{\vv x}$ 分解到 $\imm\vv A^\trans\oplus\ker\vv A$，用勾股定理.}
\ans{ex:ip:B11}{$\|\vv x-\vv y\|^2=2\|\vv x\|^2-2\langle\vv x,\vv y\rangle$，直接计算 $\vv H\vv x=\vv x-2\vv u\vv u^\trans\vv x$.}
\ans{ex:ip:B12}{对中心化的数据用 A17，再代回截距；正规方程的第一个方程就是 $\bar y=a+b\bar x$.}
\ans{ex:ip:B13}{保持夹角就保持正交. 把一组标准正交基的像归一化，证明 $\|\mathscr T\vv e_i\|$ 都相等（考察 $\vv e_i\pm\vv e_j$ 的像是否正交）.}
\ans{ex:ip:B14}{截距为第一组均值 $4$，虚拟变量系数为两组均值之差 $8-4=4$.}
\ans{ex:ip:B15}{$\vv R(\alpha)\vv R(\beta)=\vv R(\alpha+\beta)$；$\vv F(\alpha)\vv F(\beta)=\vv R(\alpha-\beta)$；计算 $\vv R(\alpha)\vv F(\beta)$ 与 $\vv F(\beta)\vv R(\alpha)$ 比较.}

\section{第十章　矩阵分解}

\subsection*{A 组答案}

\ans{ex:dec:A1}{（$\vv Q$ 的列依次对应所列特征值，答案不唯一.）
\begin{enumerate}
	\item $\vv Q=\frac1{\sqrt2}\begin{pmatrix}1&1\\1&-1\end{pmatrix}$，对角矩阵 $\diag(3,-1)$.
	\item $\vv Q=\frac1{\sqrt2}\begin{pmatrix}1&1\\1&-1\end{pmatrix}$，对角矩阵 $\diag(2,4)$.
	\item 特征值 $4,0,1$，特征向量 $\frac1{\sqrt2}(1,1,0)^\trans$，$\frac1{\sqrt2}(1,-1,0)^\trans$，$(0,0,1)^\trans$.
	\item 特征值 $3,0,0$，特征向量 $\frac1{\sqrt3}(1,1,1)^\trans$，$\frac1{\sqrt2}(1,-1,0)^\trans$，$\frac1{\sqrt6}(1,1,-2)^\trans$（后两个须在特征子空间内正交化）.
\end{enumerate}}

\ans{ex:dec:A2}{$\begin{pmatrix}3&-1\\-1&3\end{pmatrix}=2\cdot\frac12\begin{pmatrix}1&1\\1&1\end{pmatrix}+4\cdot\frac12\begin{pmatrix}1&-1\\-1&1\end{pmatrix}$. 两个投影矩阵之和为 $\vv I$，乘积为 $\vv O$.}

\ans{ex:dec:A3}{$\vv A=1\cdot\frac12\begin{pmatrix}1&1\\1&1\end{pmatrix}+2\cdot\frac12\begin{pmatrix}1&-1\\-1&1\end{pmatrix}=\frac12\begin{pmatrix}3&-1\\-1&3\end{pmatrix}$.}

\ans{ex:dec:A4}{属于 $1$ 的特征子空间是 $(0,1,1)^\trans$ 的正交补，故 $\vv A=\vv I-2\vv P$，$\vv P$ 是向 $\spn\big((0,1,1)^\trans\big)$ 的投影：$\vv A=\begin{pmatrix}1&0&0\\0&0&-1\\0&-1&0\end{pmatrix}$.}

\ans{ex:dec:A5}{特征值 $4,1,0$. 最大值 $4$，在 $\pm\frac1{\sqrt2}(1,1,0)^\trans$ 处取到；最小值 $0$，在 $\pm\frac1{\sqrt2}(1,-1,0)^\trans$ 处取到.}

\ans{ex:dec:A6}{矩阵为 $\begin{pmatrix}1&2\\2&1\end{pmatrix}$，最大值 $3$（在 $\pm\frac1{\sqrt2}(1,1)$ 处），最小值 $-1$（在 $\pm\frac1{\sqrt2}(1,-1)$ 处）.}

\ans{ex:dec:A7}{(1) $\sigma_1=3$，$\sigma_2=2$；$\vv U=\begin{pmatrix}0&1\\-1&0\end{pmatrix}$，$\vv\Sigma=\diag(3,2)$，$\vv V=\begin{pmatrix}0&1\\1&0\end{pmatrix}$. (2) $\vv U=\vv V=\frac1{\sqrt2}\begin{pmatrix}1&1\\1&-1\end{pmatrix}$，$\vv\Sigma=\diag(2,0)$. (3) $\vv U=(1)$，$\vv\Sigma=(5,\ 0)$，$\vv V=\frac15\begin{pmatrix}3&-4\\4&3\end{pmatrix}$.}

\ans{ex:dec:A8}{$\vv A^\trans\vv A=\begin{pmatrix}1&1\\1&2\end{pmatrix}$，特征值 $\frac{3\pm\sqrt5}2$，奇异值为 $\frac{\sqrt5+1}2$ 与 $\frac{\sqrt5-1}2$（黄金分割比及其倒数）；特征值是 $1,1$.}

\ans{ex:dec:A9}{对称正定矩阵的奇异值分解就是谱分解. $\vv A_1=4\cdot\frac12\begin{pmatrix}1&1\\1&1\end{pmatrix}=\begin{pmatrix}2&2\\2&2\end{pmatrix}$，误差 $\|\vv A-\vv A_1\|_F=\sigma_2=2$.}

\ans{ex:dec:A10}{$\vv A=9\vv P_1+\vv P_2$，$\vv P_1=\frac12\begin{pmatrix}1&1\\1&1\end{pmatrix}$，$\vv P_2=\frac12\begin{pmatrix}1&-1\\-1&1\end{pmatrix}$. $\vv A^{10}=9^{10}\vv P_1+\vv P_2$；$\vv B=3\vv P_1+\vv P_2=\begin{pmatrix}2&1\\1&2\end{pmatrix}$.}

\ans{ex:dec:A11}{$\vv A=\vv Q\vv O\vv Q^\trans=\vv O$. 反例：$\begin{pmatrix}0&1\\0&0\end{pmatrix}$.}

\ans{ex:dec:A12}{$\vv A^2=\vv A^\trans\vv A=\vv I$，特征值满足 $\lambda^2=1$. 它是关于特征子空间 $V_1$ 的``镜面反射''：$V_1$ 中的向量不动，$V_{-1}=V_1^\perp$ 中的向量反向.}

\ans{ex:dec:A13}{由第八章章末习题~\ref{ex:eig:C1}（取 $\vv B=\vv A^\trans$），或直接：若 $\vv A^\trans\vv A\vv x=\lambda\vv x$，$\lambda\ne0$，则 $\vv{AA}^\trans(\vv{Ax})=\lambda\vv{Ax}$ 且 $\vv{Ax}\ne\zero$；两边的特征子空间维数相同.}

\ans{ex:dec:A14}{$\tr(\vv A^\trans\vv A)=\sum_{i,j}a_{ij}^2$；它也等于 $\vv A^\trans\vv A$ 的特征值之和 $\sum\sigma_i^2$.}

\ans{ex:dec:A15}{$\det\vv A=\det\vv U\det\vv\Sigma\det\vv V^\trans=\pm\sigma_1\cdots\sigma_n$；$\vv A^{-1}=\vv V\vv\Sigma^{-1}\vv U^\trans$.}

\ans{ex:dec:A16}{只有第二个可以：能被正交矩阵对角化的实矩阵必须对称（谱定理的逆命题）.}

\ans{ex:dec:A17}{$\vv S=\frac13\begin{pmatrix}10&8\\8&10\end{pmatrix}$，特征值 $6$ 与 $\frac23$. 第一主成分方向 $\frac1{\sqrt2}(1,1)^\trans$，解释方差比例 $\frac6{6+2/3}=90\%$.}

\ans{ex:dec:A18}{$\kappa\big(\diag(10,1)\big)=10$；$\begin{pmatrix}1&1\\1&1.01\end{pmatrix}$ 的特征值约为 $2.005$ 与 $0.004988$，$\kappa\approx402$. 二者行列式都等于 $0.01$ 的 $\diag(10,0.001)$ 条件数为 $10^4$，与后者不同：行列式不决定条件数.}

\ans{ex:dec:A19}{$40\times(800+600+1)=56040$ 个数，原图 $480000$ 个，约为原来的 $11.7\%$. $k\ge343$ 时存储量超过原图.}

\ans{ex:dec:A20}{由谱定理 $\vv A$ 相似于 $\diag(\lambda_1,\ldots,\lambda_n)$，秩不变. 奇异值的平方是 $\vv A^\trans\vv A$ 的非零特征值，个数为 $\rk(\vv A^\trans\vv A)=\rk\vv A$.}

\ans{ex:dec:A21}{取正交矩阵 $\vv Q$ 使 $\vv Q^\trans\vv{AQ}=\diag(\lambda_1,\ldots,\lambda_n)$，令 $\vv y=\vv Q^\trans\vv x$，则 $\vv x^\trans\vv A\vv x=\sum\lambda_iy_i^2$，且 $\vv x\ne\zero\iff\vv y\ne\zero$. 若某个 $\lambda_i\le0$，取 $\vv y=\vv e_i$；若全为正，则 $\vv y\ne\zero$ 时和为正.}

\subsection*{B 组提示}

\ans{ex:dec:B1}{$\vv A$ 的每个特征子空间都是 $\vv B$-不变的，$\vv B$ 限制在其上仍是对称的，在其中取 $\vv B$ 的标准正交特征向量.}
\ans{ex:dec:B2}{存在性：$\vv B=\vv Q\diag(\sqrt{\lambda_i})\vv Q^\trans$. 唯一性：$\vv B$ 的每个特征子空间（属于 $\mu$）都含于 $\vv A$ 的属于 $\mu^2$ 的特征子空间，而 $\mu\ge0$ 由 $\mu^2$ 唯一确定.}
\ans{ex:dec:B3}{在 $\vv q_1^\perp=\spn(\vv q_2,\ldots,\vv q_n)$ 中重复定理~\ref{thm:rayleigh} 的证明.}
\ans{ex:dec:B4}{$\|\vv{Ax}\|^2=\vv x^\trans\vv A^\trans\vv A\vv x\le\lambda_{\max}(\vv A^\trans\vv A)\|\vv x\|^2$，在 $\vv v_1$ 处取等.}
\ans{ex:dec:B5}{$\vv{Ax}=\lambda\vv x$ 时 $|\lambda|\,\|\vv x\|=\|\vv{Ax}\|$，再用 B4 及其``最小值''版本.}
\ans{ex:dec:B6}{$\vv A^\trans\vv A$ 的非零特征值为 $3,1$（$\vv A\vv A^\trans=\begin{pmatrix}2&1\\1&2\end{pmatrix}$ 更易算），$\sigma_1=\sqrt3$，$\sigma_2=1$. 先求 $\vv U$ 的列：$\frac1{\sqrt2}(1,1)^\trans$，$\frac1{\sqrt2}(1,-1)^\trans$，再用 $\vv v_i=\vv A^\trans\vv u_i/\sigma_i$，第三个 $\vv v_3=\frac1{\sqrt3}(1,1,-1)^\trans$.}
\ans{ex:dec:B7}{$\vv A^+\vv b$ 满足正规方程（在 $\vv V$、$\vv U$ 的坐标下逐分量验证），且属于行空间 $\imm\vv A^\trans$，由第九章章末习题~\ref{ex:ip:B10} 它是范数最小的一个.}
\ans{ex:dec:B8}{$\vv J$ 的特征值为 $p,0,\ldots,0$，故矩阵的特征值为 $1+(p-1)\rho$ 与 $1-\rho$（$p-1$ 重）. 第一主成分 $\frac1{\sqrt p}(1,\ldots,1)^\trans$，比例 $\frac{1+(p-1)\rho}p$，$\rho\to1$ 时趋于 $1$，$\rho\to0$ 时趋于 $\frac1p$.}
\ans{ex:dec:B9}{协方差矩阵的第一主成分几乎就是第一个指标的方向（方差 $100$ 远大于 $1$）；相关系数为 $0.6$，相关系数矩阵的第一主成分是 $\frac1{\sqrt2}(1,1)^\trans$. 量纲改变方差，从而改变主成分.}
\ans{ex:dec:B10}{$\vv A-\vv A_k=\sum_{i>k}\sigma_i\vv u_i\vv v_i^\trans$，它的最大奇异值是 $\sigma_{k+1}$.}
\ans{ex:dec:B11}{$\vv v_i$ 是 $\vv A^\trans\vv A$ 的单重特征值的单位特征向量，只能差一个符号；$\vv u_i=\vv A\vv v_i/\sigma_i$ 随之确定.}
\ans{ex:dec:B12}{设 $\vv A=\begin{pmatrix}a&b\\c&d\end{pmatrix}$，比较 $\vv A^\trans\vv A$ 与 $\vv{AA}^\trans$ 的元素，得 $b^2=c^2$ 与 $(a-d)(b-c)=0$.}
\ans{ex:dec:B13}{第一句由第八章章末习题~\ref{ex:eig:B6}（特征值为实数）. 第二句取 $\vv x=\vv e_i$ 用定理~\ref{thm:rayleigh}.}
\ans{ex:dec:B14}{$k=1$ 时每份文书只剩一个数，查询与四份文书中三份合同类文书的相似度都是 $1$，与第四份为 $0$——同 $k=2$ 时的结果，但``合同''内部的差别完全消失，信息损失更多.}
\ans{ex:dec:B15}{$\vv x^\trans(\vv A-a\vv I)\vv x\ge0$ 对一切 $\vv x$ 成立当且仅当 $\vv A-a\vv I$ 的特征值 $\lambda_i-a\ge0$.}

\section{第十一章　二次型}

\subsection*{A 组答案}

\ans{ex:quad:A1}{(1) $\begin{pmatrix}1&2&0\\2&-3&-1\\0&-1&0\end{pmatrix}$，秩 $3$. (2) 对角元为 $0$、非对角元为 $\frac12$ 的三阶矩阵，秩 $3$. (3) 元素全为 $1$ 的三阶矩阵，秩 $1$. (4) 对角元为 $1$、非对角元为 $\frac12$ 的 $n$ 阶矩阵，即 $\frac12(\vv I+\vv J)$，秩 $n$（特征值 $\frac{n+1}2$ 与 $\frac12$，都不为零）.}

\ans{ex:quad:A2}{$x_1^2+2x_2^2-2x_1x_2+6x_2x_3$.}

\ans{ex:quad:A3}{(1) $(x_1+x_2)^2+(x_2+3x_3)^2-4x_3^2$，令 $y_1=x_1+x_2$，$y_2=x_2+3x_3$，$y_3=x_3$，得 $y_1^2+y_2^2-4y_3^2$. (2) $(x_1-2x_2)^2-x_3^2$，令 $y_1=x_1-2x_2$，$y_2=x_3$，$y_3=x_2$，得 $y_1^2-y_2^2$. (3) $x_2(x_1-x_3)$，令 $y_1=\frac12(x_1+x_2-x_3)$，$y_2=\frac12(x_2-x_1+x_3)$，$y_3=x_3$，得 $y_1^2-y_2^2$. 替换都不唯一.}

\ans{ex:quad:A4}{特征值 $2,5,1$，$\vv Q=\begin{pmatrix}1&0&0\\0&\frac1{\sqrt2}&\frac1{\sqrt2}\\0&\frac1{\sqrt2}&-\frac1{\sqrt2}\end{pmatrix}$，标准形 $2y_1^2+5y_2^2+y_3^2$. 曲面是椭球面.}

\ans{ex:quad:A5}{(1) $p=2$，$q=1$，符号差 $1$，规范形 $z_1^2+z_2^2-z_3^2$. (2)(3) 都是 $p=q=1$，符号差 $0$，规范形 $z_1^2-z_2^2$.}

\ans{ex:quad:A6}{(1) 正定（顺序主子式 $2,1,3$）. (2) 不正定（$\Delta_2=-3$）. (3) 正定（顺序主子式 $5,1,1$）.}

\ans{ex:quad:A7}{矩阵 $\begin{pmatrix}2&1&0\\1&1&\frac t2\\0&\frac t2&1\end{pmatrix}$，$\Delta_3=1-\frac{t^2}2>0$，即 $-\sqrt2<t<\sqrt2$.}

\ans{ex:quad:A8}{(1) 特征值 $3,-1$，标准方程 $3u^2-v^2=1$，双曲线；$u$ 轴沿 $\frac1{\sqrt2}(1,-1)$，$v$ 轴沿 $\frac1{\sqrt2}(1,1)$. (2) 特征值 $4,2$，标准方程 $\frac{u^2}2+\frac{v^2}4=1$，椭圆；沿 $\frac1{\sqrt2}(1,1)$ 的半轴为 $\sqrt2$，沿 $\frac1{\sqrt2}(1,-1)$ 的半轴为 $2$.}

\ans{ex:quad:A9}{(1) 合同（秩都为 $3$，$p$ 都为 $2$）：取 $\vv C$ 使 $\vv C\vv e_1=\frac1{\sqrt3}\vv e_3$，$\vv C\vv e_2=\sqrt5\,\vv e_1$，$\vv C\vv e_3=\sqrt3\,\vv e_2$. (2) 不合同：$p=1$ 与 $p=2$. (3) 合同：$\vv C=\frac1{\sqrt2}\begin{pmatrix}1&1\\1&-1\end{pmatrix}$.}

\ans{ex:quad:A10}{$\vv A^*=(\det\vv A)\vv A^{-1}$ 是对称矩阵，特征值为 $\frac{\det\vv A}{\lambda_i}>0$.}

\ans{ex:quad:A11}{驻点 $(0,0)$ 与 $(1,1)$. Hessian $\begin{pmatrix}6x&-3\\-3&6y\end{pmatrix}$：在 $(0,0)$ 处不定，是鞍点；在 $(1,1)$ 处正定，取极小值 $-1$.}

\ans{ex:quad:A12}{顺序主子式 $2,3,4$，正定. $\vv L=\begin{pmatrix}\sqrt2&0&0\\-\frac1{\sqrt2}&\sqrt{\frac32}&0\\0&-\sqrt{\frac23}&\frac2{\sqrt3}\end{pmatrix}$.}

\ans{ex:quad:A13}{$\vv L=\begin{pmatrix}2&0&0\\1&3&0\\-1&1&2\end{pmatrix}$.}

\ans{ex:quad:A14}{方差 $\vv w^\trans\vv\Sigma\vv w=0.0215$，标准差约 $14.7\%$.}

\ans{ex:quad:A15}{$\vv w^*\propto\vv\Sigma^{-1}\vv 1\propto\big(\frac1{0.09},\frac1{0.04}\big)$，即 $\vv w^*=\big(\frac4{13},\frac9{13}\big)$，方差 $\frac9{325}$，标准差约 $16.6\%$.}

\ans{ex:quad:A16}{$(\vv C^\trans\vv{AC})^\trans=\vv C^\trans\vv A^\trans\vv C$；秩不变见推论~\ref{cor:rankinvariant}. $\det(\vv C^\trans\vv{AC})=(\det\vv C)^2\det\vv A$：保持符号，不保持数值.}

\ans{ex:quad:A17}{在 $\mathbb{R}$ 上为 $z_1^2+z_2^2-z_3^2$；在 $\mathbb{C}$ 上为 $z_1^2+z_2^2+z_3^2$.}

\ans{ex:quad:A18}{依次为单叶双曲面（$p=2,q=1$）、双叶双曲面（$p=1,q=2$）、圆锥面、圆柱面（秩为 $2$）.}

\ans{ex:quad:A19}{特征值 $3$ 与 $-1$. 取 $\vv x=(1,-1)^\trans$，二次型的值为 $-2$.}

\ans{ex:quad:A20}{对主子矩阵 $\vv A_I$（行、列指标集都为 $I$），把 $\vv x_I\ne\zero$ 在其余位置补零，$\vv x_I^\trans\vv A_I\vv x_I=\vv x^\trans\vv A\vv x>0$. 正定矩阵的行列式为正（特征值之积）.}

\subsection*{B 组提示}

\ans{ex:quad:B1}{可交换的对称矩阵可同时正交对角化（第十章章末习题~\ref{ex:dec:B1}），对角元对应相乘. 反例：$\vv A=\begin{pmatrix}2&1\\1&1\end{pmatrix}$，$\vv B=\diag(1,2)$.}
\ans{ex:quad:B2}{取 $\vv C$ 使 $\vv C^\trans\vv{AC}=\vv I$，$\vv C^\trans\vv{BC}=\diag(\mu_i)$，$\mu_i\ge0$，则 $\det(\vv A+\vv B)=\det\vv A\prod(1+\mu_i)$.}
\ans{ex:quad:B3}{$\vv x^\trans(\vv A+t\vv B)\vv x\ge(\lambda_{\min}(\vv A)-|t|\max_i|\lambda_i(\vv B)|)\|\vv x\|^2$，用 B5.}
\ans{ex:quad:B4}{第 $i,j$ 行列组成的二阶主子式为正.}
\ans{ex:quad:B5}{由谱定理写成 $\sum\lambda_iy_i^2$，$\|\vv y\|=\|\vv x\|$.}
\ans{ex:quad:B6}{特征值 $2,-1,-1$，$p=1$，$q=2$. 配方时先作 $x_1=y_1+y_2$，$x_2=y_1-y_2$.}
\ans{ex:quad:B7}{$\Delta_1=\lambda$，$\Delta_2=\lambda^2-1$，$\Delta_3=(\lambda-2)(\lambda+1)^2$，故 $\lambda>2$.}
\ans{ex:quad:B8}{驻点条件 $2\vv{Ax}=\mu\vv c$，代入约束定出 $\mu$. $\vv A$ 正定保证驻点是最小值点（第十一章最小二乘例题中的配方法同样适用）.}
\ans{ex:quad:B9}{驻点条件给出 $\vv w=a\vv 1+b\vv\mu$，由两个约束解出 $a,b$（关于 $m$ 线性），方差 $\sigma^2\|\vv w\|^2$ 是 $m$ 的二次函数.}
\ans{ex:quad:B10}{解 $\det(\vv B-\lambda\vv A)=3\lambda^2-1=0$，得 $\lambda=\pm\frac1{\sqrt3}$；对每个 $\lambda$ 求 $(\vv B-\lambda\vv A)\vv x=\zero$ 的解并使 $\vv x^\trans\vv A\vv x=1$，以它们为列作 $\vv C$，则 $\vv C^\trans\vv{BC}=\diag(\frac1{\sqrt3},-\frac1{\sqrt3})$.}
\ans{ex:quad:B11}{配方：$f=(\vv x+\frac12\vv A^{-1}\vv b)^\trans\vv A(\vv x+\frac12\vv A^{-1}\vv b)+c-\frac14\vv b^\trans\vv A^{-1}\vv b$. 半正定时不一定有最小值，例如 $\vv A=\vv O$，$\vv b\ne\zero$.}
\ans{ex:quad:B12}{写 $\vv A=\vv B^\trans\vv B$（$\vv B=\vv D^{1/2}\vv Q^\trans$），则 $\|\vv B\vv x\|^2=0$，$\vv{Ax}=\vv B^\trans\vv B\vv x=\zero$.}
\ans{ex:quad:B13}{用定理~\ref{thm:posdef} 证明中的分块合同变换（把 $\vv a$ 换成矩阵块）.}
\ans{ex:quad:B14}{主轴定理把椭圆化为 $\lambda_1u^2+\lambda_2v^2\le1$，面积 $\frac\pi{\sqrt{\lambda_1\lambda_2}}$，而正交变换保持面积.}
\ans{ex:quad:B15}{$\vv x^\trans(\vv A^\trans\vv A+\mu\vv I)\vv x=\|\vv{Ax}\|^2+\mu\|\vv x\|^2$. 目标函数对 $\vv x$ 的驻点条件给出正规方程的``岭''版本.}
\ans{ex:quad:B16}{用谱分解，$\vv B=\vv Q\diag(\sqrt{\lambda_i})\vv Q^\trans$；$\vv A^{-1}$ 的特征值为 $\frac1{\lambda_i}$.}

\section{第十二章　*对偶空间与双线性型}

\subsection*{A 组答案}

\ans{ex:bil:A1}{(1) 是. (2) 不是：$2\vv x\mapsto4x_1x_2$. (3) 是. (4) 不是：$f\mapsto2f$ 时值变为 $4$ 倍. (5) 是.}

\ans{ex:bil:A2}{以三个基向量为列作 $\vv P$，对偶基是 $\vv P^{-1}$ 的各行：$\vv\varepsilon_1^*(\vv x)=x_1-x_2$，$\vv\varepsilon_2^*(\vv x)=x_2-x_3$，$\vv\varepsilon_3^*(\vv x)=x_3$.}

\ans{ex:bil:A3}{$f=f(\vv\varepsilon_1)\vv\varepsilon_1^*+f(\vv\varepsilon_2)\vv\varepsilon_2^*+f(\vv\varepsilon_3)\vv\varepsilon_3^*=\vv\varepsilon_1^*+3\vv\varepsilon_2^*+6\vv\varepsilon_3^*$.}

\ans{ex:bil:A4}{$W^0$ 是一维的，基为 $\vv x\mapsto x_1-x_2+x_3$.}

\ans{ex:bil:A5}{$\mathscr T^*$ 的矩阵为 $\begin{pmatrix}1&0&3\\2&1&-1\end{pmatrix}$；$\mathscr T^*$ 把 $\vv y\mapsto y_1+y_2+y_3$ 变为 $\vv x\mapsto4x_1+2x_2$.}

\ans{ex:bil:A6}{标准基下 $\vv A=\begin{pmatrix}1&2\\-1&3\end{pmatrix}$；新基下 $\vv B=\begin{pmatrix}5&-5\\1&3\end{pmatrix}$，可直接验证 $\vv B=\vv P^\trans\vv{AP}$，$\vv P=\begin{pmatrix}1&1\\1&-1\end{pmatrix}$.}

\ans{ex:bil:A7}{对称部分的度量矩阵 $\begin{pmatrix}1&\frac12\\ \frac12&3\end{pmatrix}$，交错部分 $\begin{pmatrix}0&\frac32\\-\frac32&0\end{pmatrix}$.}

\ans{ex:bil:A8}{(1)(2) 非退化（度量矩阵行列式 $-1$）；(3) 退化（度量矩阵全为 $1$，不可逆）.}

\ans{ex:bil:A9}{对应二次型 $x_1^2+2x_1x_2+3x_2^2$. 由二次型得 $f(\vv x,\vv y)=x_1y_1-2x_1y_2-2x_2y_1+2x_2y_2$.}

\ans{ex:bil:A10}{$f(\vv x,\vv x)=0$，$f(\vv x,\vv y)=-f(\vv y,\vv x)$. 度量矩阵 $\begin{pmatrix}0&1\\-1&0\end{pmatrix}$，可逆，非退化.}

\ans{ex:bil:A11}{Lagrange 基：$\frac12(x-1)(x-2)$，$-x(x-2)$，$\frac12x(x-1)$.}

\ans{ex:bil:A12}{$\dim W=2$，$\dim W^0=2$，基为 $\vv x\mapsto x_1+x_2$ 与 $\vv x\mapsto x_3-x_4$.}

\ans{ex:bil:A13}{双线性显然，对称性由 $\tr(\vv{AB})=\tr(\vv{BA})$；非退化：$\tr(\vv A\vv E_{ji})=a_{ij}$. 不是内积：$\vv A=\begin{pmatrix}0&1\\-1&0\end{pmatrix}$ 时 $\tr(\vv A^2)=-2<0$.}

\ans{ex:bil:A14}{坐标为 $\big(\varphi(1),\varphi(x),\varphi(x^2)\big)=\big(1,\frac12,\frac13\big)$.}

\ans{ex:bil:A15}{对任意 $\vv\alpha=\sum a_i\vv e_i$，$\langle\vv\alpha,\vv\beta\rangle=\sum a_if(\vv e_i)=f(\vv\alpha)$. 唯一性：若 $\langle\vv\alpha,\vv\beta-\vv\beta'\rangle=0$ 对一切 $\vv\alpha$ 成立，取 $\vv\alpha=\vv\beta-\vv\beta'$.}

\subsection*{B 组提示}

\ans{ex:bil:B1}{取 $V$ 的一组基，映射的矩阵的第 $i$ 行恰是 $f_i$ 在对偶基下的坐标. 满射当且仅当这个 $k\times n$ 矩阵行满秩，即各行线性无关.}
\ans{ex:bil:B2}{取 $\vv\alpha_0$ 使 $f(\vv\alpha_0)=1$，则 $V=\ker f\oplus\spn(\vv\alpha_0)$，$c=g(\vv\alpha_0)$. 推广：对映射 $\vv\alpha\mapsto(f_1(\vv\alpha),\ldots,f_k(\vv\alpha))$ 用同样的思路，或用零化子的维数.}
\ans{ex:bil:B3}{记 $U=\spn(f_1,\ldots,f_k)\subset V^*$. 把 $f_i$ 写成坐标后，$\ker f_1\cap\cdots\cap\ker f_k$ 就是以这些坐标为行的齐次方程组的解空间，维数为 $\dim V-\dim U$. 方程组语言：齐次方程组只有零解，当且仅当各方程的系数行向量张成整个行向量空间.}
\ans{ex:bil:B4}{由定理~\ref{thm:dualrank}：$\rk\mathscr T^*=\rk\mathscr T$.}
\ans{ex:bil:B5}{$a=c=\frac13$，$b=\frac43$. 积分是 $\mathbb{R}[x]_3$ 上的线性函数，三个求值函数构成 $\mathbb{R}[x]_3^*$ 的基（A11）. 三次时验证 $\int_0^2(x-1)^3\,\mathrm dx=0$ 与右边相等.}
\ans{ex:bil:B6}{双线性型由度量矩阵唯一确定，反之任一矩阵给出双线性型. 后半句即第三章章末习题~\ref{ex:mat:A22}.}
\ans{ex:bil:B7}{由定理~\ref{thm:alternating}，非退化时没有零块，$\dim V=2s$.}
\ans{ex:bil:B8}{度量矩阵 $\begin{pmatrix}1&1&0\\1&1&0\\0&0&1\end{pmatrix}$ 的核为 $\spn\big((1,-1,0)^\trans\big)$.}
\ans{ex:bil:B9}{迷向向量满足 $x_1^2=x_2^2$，即两条直线 $x_2=\pm x_1$. 后半句：设 $\vv P=\begin{pmatrix}a&b\\c&d\end{pmatrix}$，由条件得 $a^2-c^2=1$，$d^2-b^2=1$，$ab=cd$，再令 $a=\cosh t$.}
\ans{ex:bil:B10}{对 $\vv\alpha\in W$，$f\in W^0$，有 $(\mathscr T^*f)(\vv\alpha)=f(\mathscr T\vv\alpha)=0$.}
（须在 %% =====================================================================
%% exercises.tex　《高等代数讲义》章末习题（A/B/C 三组）
%%
%% 用法
%%   1. 在导言区（\begin{document} 之前、全书统一记号之后）写 %% =====================================================================
%% exercises.tex　《高等代数讲义》章末习题（A/B/C 三组）
%%
%% 用法
%%   1. 在导言区（\begin{document} 之前、全书统一记号之后）写 %% =====================================================================
%% exercises.tex　《高等代数讲义》章末习题（A/B/C 三组）
%%
%% 用法
%%   1. 在导言区（\begin{document} 之前、全书统一记号之后）写 \input{exercises}.
	%%   2. 在每章末尾（原 \begin{problemset}...\end{problemset} 之后，或替换它）
	%%      写对应的命令：
	%%        第一章 \ExElimination      第二章 \ExVectorSpace   第三章 \ExMatrix
	%%        第四章 \ExLinearMap        第五章 \ExDeterminant   第六章 \ExSolutions
	%%        第七章 \ExPolynomial       第八章 \ExEigen         第九章 \ExInnerProduct
	%%        第十章 \ExDecomposition    第十一章 \ExQuadratic   第十二章 \ExBilinear
	%%   3. 分组标准
	%%        A 组：直接运用本章定义与定理的计算和证明；
	%%        B 组：需要综合多个结论、或有一定技巧的计算和证明；
	%%        C 组：有挑战性的题目，或超出本课程大纲、为继续深造准备的内容.
	%%   4. 题目与正文中已有的节末练习、原 problemset 不重复；个别题目注明``接正文某题''.
	%%      第六章按扩充版（ch06-solutions.tex）的内容编写.
	%%   5. 每道题都带标签 ex:<章键>:<组号>，正文、习题与附录答案一律用 \ref 引用，
	%%      渲染为``A4''``B12''这样的题号. 章键依次为
	%%        elim vs mat lmap det sol poly eig ip dec quad bil（第一至十二章）.
	%%      例：第一章章末习题~\ref{ex:elim:B12}. 新增题目请加在组末，并补上标签.
	%%      A 组答案、B 组提示在 answers.tex 中，按同一标签对应.
	%% 依赖：导言区的 \vv \zero \trans \spn \rk \imm \diag，以及 enumitem（elegantbook 已加载）.
	%% 引用的前文标签：thm:sameSolution thm:solvability prop:linspc eg:pricedecomp cor:homogeneous thm:rouche thm:sol-determinable
	%%   thm:spectral thm:posdef thm:svd thm:orthcanon
	%% =====================================================================
	
	\providecommand{\tr}{\operatorname{tr}}
	\newcommand{\exsetheading}{%
		\section*{章末习题}%
		\markright{章末习题}%
	}
\newcommand{\exgroup}[2]{%
	\par\medskip{\centering\sffamily\bfseries #1\quad #2\par}\smallskip}
	\newenvironment{exA}{\begin{enumerate}[label=\textbf{A\arabic*.},ref=A\arabic*,leftmargin=2.8em,itemsep=4pt]}{\end{enumerate}}
	\newenvironment{exB}{\begin{enumerate}[label=\textbf{B\arabic*.},ref=B\arabic*,leftmargin=2.8em,itemsep=4pt]}{\end{enumerate}}
	\newenvironment{exC}{\begin{enumerate}[label=\textbf{C\arabic*.},ref=C\arabic*,leftmargin=2.8em,itemsep=4pt]}{\end{enumerate}}
	\newcommand{\Agroup}{\exgroup{A 组}{基础巩固}}
	\newcommand{\Bgroup}{\exgroup{B 组}{综合提升}}
	\newcommand{\Cgroup}{\exgroup{C 组}{拓展探索}}
	
	% =====================================================================
	% 第一章　线性方程组与消元法
	% =====================================================================
	\newcommand{\ExElimination}{%
		\exsetheading
		\Agroup
		\begin{exA}
			\item\label{ex:elim:A1} 判断下列矩阵哪些是阶梯形，哪些是简化阶梯形：
			\[
			\begin{pmatrix}1&2&0\\0&0&1\end{pmatrix},\quad
			\begin{pmatrix}1&0&3\\0&2&1\\0&0&0\end{pmatrix},\quad
			\begin{pmatrix}0&1&0\\1&0&0\end{pmatrix},\quad
			\begin{pmatrix}1&0&2&0\\0&1&-1&0\\0&0&0&1\end{pmatrix},\quad
			\begin{pmatrix}0&0\\0&1\end{pmatrix}.
			\]
			\item\label{ex:elim:A2} 写出方程组 $3x_1-x_3=2$，$x_2+4x_3=0$，$x_1-x_2=5$ 的增广矩阵；反过来，写出以 $\begin{pmatrix}1&-2&0&3\\0&1&5&-1\end{pmatrix}$ 为增广矩阵的方程组.
			\item\label{ex:elim:A3} 用消元法解方程组 $\left\{\begin{array}{l}x_1+2x_2+3x_3=14,\\2x_1-x_2+x_3=3,\\3x_1+x_2-x_3=2.\end{array}\right.$
			\item\label{ex:elim:A4} 用消元法解方程组 $\left\{\begin{array}{l}x_1+x_2+x_3+x_4=2,\\x_1-x_2+x_3-x_4=4,\\2x_1+x_2-x_3+3x_4=-1,\\x_1+2x_2+3x_3-2x_4=5.\end{array}\right.$
			\item\label{ex:elim:A5} 判断方程组 $\left\{\begin{array}{l}x_1-x_2+x_3=1,\\2x_1+x_2-x_3=2,\\x_1+2x_2-2x_3=3\end{array}\right.$ 是否有解，并说明消元过程中哪一行``暴露''了结论.
			\item\label{ex:elim:A6} 求方程组 $\left\{\begin{array}{l}x_1+2x_2-x_3+x_4=3,\\2x_1+4x_2-x_3+3x_4=7,\\-x_1-2x_2+2x_3=-2\end{array}\right.$ 的全部解，并指出主变量与自由变量.
			\item\label{ex:elim:A7} 求齐次方程组 $\left\{\begin{array}{l}x_1+x_2-x_3-x_4=0,\\2x_1-5x_2+3x_3+2x_4=0,\\7x_1-7x_2+3x_3+x_4=0\end{array}\right.$ 的全部解.
			\item\label{ex:elim:A8} 把矩阵 $\begin{pmatrix}2&4&-2&6\\1&2&1&1\\3&6&1&5\end{pmatrix}$ 化为简化阶梯形，并写出每一步所用的初等行变换.
			\item\label{ex:elim:A9} 写出全部 $2\times2$ 简化阶梯形矩阵的``形状''（可以含参数），并说明共有几类.
			\item\label{ex:elim:A10} $k$ 取何值时，方程组 $x+y=1$，$x-y=3$，$2x+ky=5$ 有解？
			\item\label{ex:elim:A11} 设 $a,b$ 为常数. 方程组 $x+2y=a$，$3x+6y=b$ 何时有解？有解时写出全部解.
			\item\label{ex:elim:A12} 求过三点 $(1,2)$，$(2,3)$，$(3,6)$ 的抛物线 $y=ax^2+bx+c$.
			\item\label{ex:elim:A13} 求过四点 $(-1,2)$，$(0,1)$，$(1,0)$，$(2,5)$ 的三次多项式函数 $y=a+bx+cx^2+dx^3$.
			\item\label{ex:elim:A14} 某剧场三种票价分别为 $30$，$50$，$80$ 元. 某场演出共售出 $100$ 张票，票房 $4800$ 元，且 $80$ 元票的张数恰为 $30$ 元票的一半. 三种票各售出多少张？
			\item\label{ex:elim:A15} 配平化学方程式 $x_1\,\mathrm{C_3H_8}+x_2\,\mathrm{O_2}\to x_3\,\mathrm{CO_2}+x_4\,\mathrm{H_2O}$：对 C，H，O 三种原子分别写出守恒方程，解这个齐次方程组，并取最小的正整数解.
			\item\label{ex:elim:A16} 某城区四个路口 $A,B,C,D$ 由单行道依次连成环路 $A\to B\to C\to D\to A$，环路上的车流量依次记为 $x_1$（$A\to B$），$x_2$（$B\to C$），$x_3$（$C\to D$），$x_4$（$D\to A$）（单位：辆/小时）. 各路口与外部道路的进出量为：$A$ 驶入 $400$、驶出 $300$；$B$ 驶入 $200$、驶出 $350$；$C$ 驶入 $250$、驶出 $100$；$D$ 驶入 $150$、驶出 $250$. 设每个路口驶入总量等于驶出总量.
			\begin{enumerate}
				\item 列出方程组并求全部解；
				\item 若各路段车流量都不能为负，$x_1$ 最少是多少？
				\item 若把 $D$ 的驶出量改为 $350$，方程组会怎样？这反映了什么守恒律？
			\end{enumerate}
			\item\label{ex:elim:A17} 证明：二元一次方程组 $ax+by=e$，$cx+dy=f$ 对任意的 $e,f$ 都有唯一解，当且仅当 $ad-bc\ne0$.（只用消元法，不用行列式.）
			\item\label{ex:elim:A18} 证明：若线性方程组有两个不同的解 $\vv x_1,\vv x_2$，则对任意数 $t$，$\vv x_1+t(\vv x_2-\vv x_1)$ 都是解. 由此说明：线性方程组的解要么不存在，要么唯一，要么有无穷多个.
			\item\label{ex:elim:A19} 证明方程组 $\{x+y+z=3,\ x-y=0\}$ 与 $\{2x+z=3,\ x+y+z=3\}$ 同解.
			\item\label{ex:elim:A20} 举例说明：若\textbf{同时}做``第一个方程加到第二个方程上''与``第二个方程加到第一个方程上''（两步都用原来的方程），所得方程组可能与原方程组不同解. 这与定理~\ref{thm:sameSolution}（初等行变换不改变解集）矛盾吗？
			\item\label{ex:elim:A21} 已知一个方程组的增广矩阵化为阶梯形后为 $\begin{pmatrix}1&3&-1&2&0\\0&0&2&1&4\\0&0&0&0&a\end{pmatrix}$. 讨论 $a$ 取何值时方程组有解；有解时它有几个自由变量？
		\end{exA}
		
		\Bgroup
		\begin{exB}
			\item\label{ex:elim:B1} 讨论 $a,b$ 取何值时，方程组
			\[\left\{\begin{array}{l}x_1+x_2+x_3+x_4=0,\\x_2+2x_3+2x_4=1,\\-x_2+(a-3)x_3-2x_4=b,\\3x_1+2x_2+x_3+ax_4=-1\end{array}\right.\]
			有唯一解、无解、有无穷多解；有无穷多解时求出全部解.
			\item\label{ex:elim:B2} 讨论 $\lambda$ 取何值时，方程组
			$(1+\lambda)x_1+x_2+x_3=0$，$x_1+(1+\lambda)x_2+x_3=3$，$x_1+x_2+(1+\lambda)x_3=\lambda$
			有唯一解、无解、有无穷多解.
			\item\label{ex:elim:B3} 设 $n$ 个方程、$n$ 个未知量的方程组对应的齐次方程组只有零解. 证明：无论常数项取什么值，原方程组都有唯一解.（提示：考察简化阶梯形中主元的个数.）
			\item\label{ex:elim:B4} 证明：线性方程组无解，当且仅当可以把若干个方程分别乘以适当的数再相加，得到矛盾方程 $0=1$.
			\item\label{ex:elim:B5} 证明：交换两行这一初等变换，可以由``倍加''与``某行乘以 $-1$''这两类变换的组合实现. 由此说明三类初等行变换并不是``独立''的.
			\item\label{ex:elim:B6} 证明：若线性方程组的系数和常数项都是有理数，且它有实数解，则它一定有有理数解；若它的实数解唯一，则这个解就是有理数解.
			\item\label{ex:elim:B7} 讨论方程组
			\[x_1-x_2=a_1,\quad x_2-x_3=a_2,\quad\ldots,\quad x_{n-1}-x_n=a_{n-1},\quad x_n-x_1=a_n\]
			有解的充要条件，并在有解时求全部解.
			\item\label{ex:elim:B8} 讨论方程组
			\[x_1+x_2=a_1,\quad x_2+x_3=a_2,\quad\ldots,\quad x_{n-1}+x_n=a_{n-1},\quad x_n+x_1=a_n\]
			的解. 证明：$n$ 为奇数时解唯一；$n$ 为偶数时有解当且仅当 $a_1-a_2+a_3-\cdots-a_n=0$.
			\item\label{ex:elim:B9} 某经济体有三个部门，直接消耗系数矩阵与最终需求分别为
			\[\vv A=\begin{pmatrix}0.1&0.2&0.2\\0.2&0.1&0.2\\0.1&0.3&0.1\end{pmatrix},\qquad \vv d=\begin{pmatrix}20\\130\\65\end{pmatrix}.\]
			按第一章``投入产出模型''的方式列出方程组，求各部门的总产出.
			\item\label{ex:elim:B10} 估计用消元法（化为阶梯形后回代）解 $n$ 个方程、$n$ 个未知量的方程组所需的乘除法次数，证明它约为 $\frac13n^3$（即与 $\frac13n^3$ 之比在 $n\to\infty$ 时趋于 $1$）.
			\item\label{ex:elim:B11} 设三个平面 $a_ix+b_iy+c_iz=d_i$（$i=1,2,3$）两两相交. 用消元法的语言说明：它们恰好交于一点、交于一条直线、或没有公共点（三条交线互相平行），并各举一例.
			\item\label{ex:elim:B12} 证明：$n$ 元齐次方程组的解集对加法和数乘封闭；而非齐次方程组的两个解之差是对应齐次方程组的解. 由此说明：非齐次方程组的全部解，等于它的任一个解加上齐次方程组的全部解.
			\item\label{ex:elim:B13} 一家律所有甲、乙、丙三位律师，本月分别承办了若干件合同、劳动、知识产权纠纷. 已知：三类案件合计 $30$ 件；甲承办的数量 $=$ 乙 $+$ 丙；每件合同、劳动、知识产权案件分别需要 $10$，$6$，$15$ 个工时，三人本月共用 $300$ 个工时……
			这些条件够不够确定每位律师承办了多少案件？请自己补全一个合理的设定（明确未知量是什么），列出方程组，用自由变量的个数说明还缺多少条独立的信息.
			\item\label{ex:elim:B14} 设方程组的系数矩阵每一行的元素之和都等于 $0$. 证明 $x_1=x_2=\cdots=x_n=1$ 是对应齐次方程组的解；若非齐次方程组有解 $\vv x_0$，则对任意数 $t$，$\vv x_0+t(1,\ldots,1)^\trans$ 也是解.
			\item\label{ex:elim:B15} 证明：如果方程组 (I) 的每个方程都是方程组 (II) 各方程的``线性组合''（各乘一个数再相加），反之亦然，那么 (I) 与 (II) 同解. 逆命题对于有解的方程组成立吗？先在两个未知量的情形下思考.
			\item\label{ex:elim:B16} 讨论 $\lambda$ 取何值时，方程组 $\lambda x+y+z=1$，$x+\lambda y+z=\lambda$，$x+y+\lambda z=\lambda^2$ 有唯一解、无解、有无穷多解；有唯一解时求出解.（第六章例~\ref{eg:lambda} 将用秩与行列式重新组织这个讨论.）
		\end{exB}
		
		\Cgroup
		\begin{exC}
			\item\label{ex:elim:C1} （百鸡问题）《张丘建算经》（约公元 5 世纪）：``今有鸡翁一，直钱五；鸡母一，直钱三；鸡雏三，直钱一. 凡百钱买鸡百只，问鸡翁、母、雏各几何？''列出线性方程组，求它的全部解，再从中找出所有\textbf{非负整数}解. 线性代数在哪一步之后就无能为力了？
			\item\label{ex:elim:C2} （幻方）若一个 $3\times3$ 数表的每行、每列和两条对角线上的三个数之和都等于同一个数 $s$，称它为幻方（不要求数互不相同）. 把 $9$ 个元素作为未知量、$s$ 也作为未知量，写出方程组，证明：
			\begin{enumerate}
				\item 中心元素一定等于 $s/3$；
				\item 全部幻方可以用 $3$ 个自由参数表示，写出通式；
				\item 用 $1,2,\ldots,9$ 各一次填成的幻方本质上只有一个（其余由旋转、翻转得到）.
			\end{enumerate}
			\item\label{ex:elim:C3} （点灯游戏）$3\times3$ 方格中每格有一盏灯，按下某格会使该格及其上下左右相邻的格子开关状态翻转. 在``$1+1=0$''的算术（即模 $2$ 运算）中，把``按哪些格''作为未知量，把``从某一状态变为全灭''写成线性方程组. 证明：从任意初始状态出发，总能把灯全部熄灭.（模 $2$ 运算的集合 $\{0,1\}$ 不是本书意义下的数域，但消元法照样适用——想一想为什么.）
			\item\label{ex:elim:C4} （舍入误差与主元）在``只保留三位有效数字''的计算机上解方程组 $0.0001x+y=1$，$x+y=2$：
			\begin{enumerate}
				\item 直接用第一个方程消去第二个方程中的 $x$；
				\item 先交换两个方程再消元.
			\end{enumerate}
			把两种结果与精确解比较，说明数值软件为什么要``选绝对值最大的元素作主元''.
			\item\label{ex:elim:C5} 设两个都有解的方程组 (I)(II) 都有 $m$ 个方程、$n$ 个未知量，且同解. 证明：它们的增广矩阵行等价.（提示：先证明二者的简化阶梯形的非零行完全相同.）
			\item\label{ex:elim:C6} （Jacobi 迭代）对方程组 $4x+y=9$，$x+3y=5$，把它改写为 $x=(9-y)/4$，$y=(5-x)/3$，从 $(x_0,y_0)=(0,0)$ 出发反复代入. 计算前五步，观察是否收敛到真解. 再对 $x+4y=9$，$3x+y=5$ 做同样的事，会发生什么？试找出一个保证收敛的充分条件.
			\item\label{ex:elim:C7} （电路）一个电路中有三条支路电流 $I_1,I_2,I_3$. 节点 $A$ 处 $I_1=I_2+I_3$；两个回路的电压方程为 $2I_1+I_2=8$，$I_2-2I_3=0$（单位从略）. 求各电流. 再说明：若把第三个节点 $B$ 的电流方程也写进来，方程组为什么不会因此``多一个条件''？
			\item\label{ex:elim:C8} 证明：一个 $m\times n$ 矩阵（$m,n\ge1$）通过初等行变换所能得到的不同的简化阶梯形``形状''（即主元所在列的位置组合）共有 $\sum_{r=0}^{\min(m,n)}\binom nr$ 种.
			\item\label{ex:elim:C9} 证明：一个矩阵的简化阶梯形是唯一的.（提示：设 $\vv R_1,\vv R_2$ 都是 $\vv M$ 的简化阶梯形，则以它们为系数矩阵的齐次方程组同解. 对列数归纳：删去最后一列后由归纳假设前面各列相同，再比较最后一列，考察它在两者中是否为主元列.）
		\end{exC}
	}
	
	% =====================================================================
	% 第二章　向量空间
	% =====================================================================
	\newcommand{\ExVectorSpace}{%
		\exsetheading
		\Agroup
		\begin{exA}
			\item\label{ex:vs:A1} 下列集合关于通常的四则运算是否构成数域？
			(1) $\{a+b\sqrt3\mid a,b\in\mathbb{Q}\}$；(2) $\{a+b\sqrt2\mid a\in\mathbb{Q},\ b\in\mathbb{Z}\}$；(3) 全体偶数；(4) $\{a+b\sqrt[3]{2}\mid a,b\in\mathbb{Q}\}$；(5) 全体纯虚数与 $0$.
			\item\label{ex:vs:A2} 下列集合关于通常的加法与数乘是否构成实向量空间？
			(1) 全体 $n$ 阶实对称矩阵；(2) $\{(x,y)^\trans\in\mathbb{R}^2\mid x\ge0\}$；(3) 次数恰为 $n$ 的实系数多项式全体；(4) $\{f\in\mathbb{R}[x]\mid f(1)=0\}$；(5) 满足 $a_{n+2}=a_{n+1}+a_n$ 的全体实数列；(6) $\{(x,y,z)^\trans\mid x+y+z=1\}$.
			\item\label{ex:vs:A3} 判断下列向量组的线性相关性：
			\begin{enumerate}
				\item $(1,2,-1)^\trans$，$(2,-3,1)^\trans$，$(4,1,-1)^\trans$；
				\item 在 $\mathbb{R}[x]_3$ 中：$1+x$，$1-x$，$x^2$；
				\item 在 $\mathbb{R}[x]_3$ 中：$1+x+x^2$，$1-x$，$2+x^2$，$x+x^2$；
				\item 在 $M_2(\mathbb{R})$ 中：$\begin{pmatrix}1&0\\0&1\end{pmatrix}$，$\begin{pmatrix}0&1\\1&0\end{pmatrix}$，$\begin{pmatrix}1&1\\1&1\end{pmatrix}$，$\begin{pmatrix}1&0\\0&0\end{pmatrix}$.
			\end{enumerate}
			\item\label{ex:vs:A4} 把 $\vv\beta=(1,2,1,1)^\trans$ 表示为 $\vv\alpha_1=(1,1,1,1)^\trans$，$\vv\alpha_2=(1,1,-1,-1)^\trans$，$\vv\alpha_3=(1,-1,1,-1)^\trans$，$\vv\alpha_4=(1,-1,-1,1)^\trans$ 的线性组合.
			\item\label{ex:vs:A5} $t$ 取何值时，$(1,1,0)^\trans$，$(1,3,-1)^\trans$，$(5,3,t)^\trans$ 线性相关？
			\item\label{ex:vs:A6} 求向量组 $\vv\alpha_1=(1,-1,2,4)^\trans$，$\vv\alpha_2=(0,3,1,2)^\trans$，$\vv\alpha_3=(3,0,7,14)^\trans$，$\vv\alpha_4=(1,-1,2,0)^\trans$，$\vv\alpha_5=(2,1,5,6)^\trans$ 的秩和一个极大无关组，并把其余向量用这个极大无关组线性表出.
			\item\label{ex:vs:A7} 设 $\vv\alpha_1,\vv\alpha_2,\vv\alpha_3$ 线性无关. 判断下列向量组的线性相关性：
			(1) $\vv\alpha_1-\vv\alpha_2$，$\vv\alpha_2-\vv\alpha_3$，$\vv\alpha_3-\vv\alpha_1$；
			(2) $\vv\alpha_1$，$\vv\alpha_1+\vv\alpha_2$，$\vv\alpha_1+\vv\alpha_2+\vv\alpha_3$；
			(3) $\vv\alpha_1+2\vv\alpha_2$，$\vv\alpha_2+2\vv\alpha_3$，$\vv\alpha_3+2\vv\alpha_1$.
			\item\label{ex:vs:A8} 证明：含有零向量的向量组线性相关；含有两个相同向量的向量组线性相关；单个向量 $\vv\alpha$ 线性无关当且仅当 $\vv\alpha\ne\zero$.
			\item\label{ex:vs:A9} 证明：两个向量线性相关，当且仅当其中一个是另一个的倍数. 三个向量的情形能这样说吗？
			\item\label{ex:vs:A10} 只用向量空间的八条公理与命题~\ref{prop:linspc}，证明：(1) $\vv\alpha+\vv\alpha=2\vv\alpha$；(2) 若 $k\ne0$ 且 $k\vv\alpha=k\vv\beta$，则 $\vv\alpha=\vv\beta$；(3) $-(-\vv\alpha)=\vv\alpha$.
			\item\label{ex:vs:A11} 证明下列集合是相应空间的子空间，并求一组基和维数：
			\begin{enumerate}
				\item $\{\vv x\in\mathbb{R}^4\mid x_1+x_2=x_3+x_4\}$；
				\item $n$ 阶上三角矩阵全体；
				\item 迹为零的 $n$ 阶方阵全体 $\{\vv X\mid x_{11}+\cdots+x_{nn}=0\}$；
				\item $\{f\in\mathbb{R}[x]_4\mid f(1)=0\}$.
			\end{enumerate}
			\item\label{ex:vs:A12} 判断下列子集是否为子空间：(1) $\{(a,b,c)^\trans\mid a=2b\}$；(2) $\{(a,b,c)^\trans\mid ab=0\}$；(3) $\{(a,b,c)^\trans\mid a,b,c\in\mathbb{Q}\}\subset\mathbb{R}^3$；(4) $\{f\in C[0,1]\mid f(0)=2f(1)\}$.
			\item\label{ex:vs:A13} 求 $\vv\alpha_1=(1,2,1,0)^\trans$，$\vv\alpha_2=(-1,1,1,1)^\trans$，$\vv\alpha_3=(0,3,2,1)^\trans$，$\vv\alpha_4=(1,0,1,2)^\trans$ 张成的子空间的维数和一组基.
			\item\label{ex:vs:A14} 求多项式 $2x^2-x+3$ 在基 $1,\ x-1,\ (x-1)^2$ 下的坐标.
			\item\label{ex:vs:A15} 求矩阵 $\begin{pmatrix}2&3\\-1&4\end{pmatrix}$ 在 $M_2(\mathbb{R})$ 的基 $\begin{pmatrix}1&0\\0&0\end{pmatrix},\begin{pmatrix}1&1\\0&0\end{pmatrix},\begin{pmatrix}1&1\\1&0\end{pmatrix},\begin{pmatrix}1&1\\1&1\end{pmatrix}$ 下的坐标.
			\item\label{ex:vs:A16} (1) $\mathbb{R}$ 是否是 $\mathbb{Q}$ 上的向量空间？(2) $\mathbb{Q}$ 是否是 $\mathbb{R}$ 上的向量空间？(3) $\mathbb{C}$ 作为 $\mathbb{R}$ 上的向量空间和作为 $\mathbb{C}$ 上的向量空间，维数各是多少？
			\item\label{ex:vs:A17} 写出 $\mathbb{F}[x]_n$ 与 $\mathbb{F}^n$ 之间、$M_2(\mathbb{F})$ 与 $\mathbb{F}^4$ 之间的一个同构映射，并验证它保持线性运算.
			\item\label{ex:vs:A18} 设 $W_1=\spn\big((1,0,0)^\trans,(0,1,0)^\trans\big)$，$W_2=\spn\big((1,1,1)^\trans\big)$，$W_3=\spn\big((1,1,0)^\trans\big)$. 分别判断 $W_1+W_2$，$W_1+W_3$ 是否为直和，并求和空间的维数.
			\item\label{ex:vs:A19} 设 $W_1=\spn\big((1,1,0,0)^\trans,(1,0,1,1)^\trans\big)$，$W_2=\spn\big((0,-1,1,1)^\trans,(1,2,0,1)^\trans\big)$. 求 $W_1+W_2$ 与 $W_1\cap W_2$ 的维数，并验证维数公式.
			\item\label{ex:vs:A20} 设 $\vv\alpha_1,\ldots,\vv\alpha_s$ 的秩为 $r$. 证明：其中任意 $r+1$ 个向量线性相关；任意 $r$ 个线性无关的向量都构成极大无关组.
		\end{exA}
		
		\Bgroup
		\begin{exB}
			\item\label{ex:vs:B1} 证明：两个子空间的并 $W_1\cup W_2$ 是子空间，当且仅当 $W_1\subset W_2$ 或 $W_2\subset W_1$.
			\item\label{ex:vs:B2} 证明：若向量组 $\vv\beta_1,\ldots,\vv\beta_t$ 可由 $\vv\alpha_1,\ldots,\vv\alpha_s$ 线性表出，且两组的秩相等，则两组等价.
			\item\label{ex:vs:B3} 设向量组 (I) 与 (II) 的秩分别为 $r_1,r_2$，把它们合并为 (III). 证明 $\max\{r_1,r_2\}\le\rk(\mathrm{III})\le r_1+r_2$. 再证明：$\rk(\vv\alpha_1+\vv\beta_1,\ldots,\vv\alpha_s+\vv\beta_s)\le\rk(\vv\alpha_1,\ldots,\vv\alpha_s)+\rk(\vv\beta_1,\ldots,\vv\beta_s)$.
			\item\label{ex:vs:B4} 设 $\vv\alpha_1,\ldots,\vv\alpha_n$ 是 $V$ 的基，$\vv\beta=a_1\vv\alpha_1+\cdots+a_n\vv\alpha_n$. 证明：把 $\vv\alpha_j$ 换成 $\vv\beta$ 后所得的向量组仍是基，当且仅当 $a_j\ne0$.
			\item\label{ex:vs:B5} 设 $p_0,p_1,\ldots,p_{n-1}\in\mathbb{F}[x]$ 满足 $\deg p_k=k$. 证明它们构成 $\mathbb{F}[x]_n$ 的一组基.
			\item\label{ex:vs:B6} 记 $V$ 为满足 $a_{n+2}=a_{n+1}+2a_n$（$n\ge1$）的全体实数列. 证明 $V$ 是向量空间，$\dim V=2$，并说明 $\{2^n\}$ 与 $\{(-1)^n\}$ 构成 $V$ 的一组基. 由此写出 $a_1=1$，$a_2=1$ 时的通项.
			\item\label{ex:vs:B7} 在 $C(\mathbb{R})$ 中：(1) 证明 $\mathrm e^x,\mathrm e^{2x},\mathrm e^{3x}$ 线性无关；(2) 证明 $\sin x,\cos x,\sin(x+1)$ 线性相关；(3) $1,\cos x,\cos2x,\cos^2x$ 的秩是多少？
			\item\label{ex:vs:B8} 证明 $\{a+b\sqrt2+c\sqrt3+d\sqrt6\mid a,b,c,d\in\mathbb{Q}\}$ 是数域，并证明 $1,\sqrt2,\sqrt3,\sqrt6$ 在 $\mathbb{Q}$ 上线性无关. 这个数域作为 $\mathbb{Q}$ 上的向量空间是几维的？
			\item\label{ex:vs:B9} 写出 $\mathbb{R}^2$ 与 $\mathbb{R}^3$ 的全部子空间（按维数分类，并说明几何意义）.
			\item\label{ex:vs:B10} 证明：$\mathbb{R}^5$ 中任意两个三维子空间的交至少是一维的. 一般地，$n$ 维空间中维数分别为 $k_1,k_2$ 的两个子空间，若 $k_1+k_2>n$，则交不为 $\{\zero\}$.
			\item\label{ex:vs:B11} 设 $a\in\mathbb{F}$. 证明 $\mathbb{F}[x]_n=\mathbb{F}\cdot1\oplus\{f\in\mathbb{F}[x]_n\mid f(a)=0\}$，并写出 $f$ 按这个直和的分解.
			\item\label{ex:vs:B12} 设 $V=C[-1,1]$，$E$，$O$ 分别为其中的偶函数、奇函数全体. 证明 $V=E\oplus O$，并把 $\mathrm e^x$ 按此分解.
			\item\label{ex:vs:B13} 设 $\dim V=n$，$W_1,\ldots,W_k$ 是 $V$ 的子空间，且 $\dim W_1+\cdots+\dim W_k>(k-1)n$. 证明 $W_1\cap\cdots\cap W_k\ne\{\zero\}$.
			\item\label{ex:vs:B14} 设 $\vv\alpha_1,\ldots,\vv\alpha_n$ 是 $V$ 的基，$W=\spn(\vv\alpha_1+\vv\alpha_2,\ \vv\alpha_2+\vv\alpha_3,\ \ldots,\ \vv\alpha_{n-1}+\vv\alpha_n)$. 求 $\dim W$，并求 $W$ 的一个补空间.
			\item\label{ex:vs:B15} 设 $V_1=\{\vv x\in\mathbb{F}^n\mid x_1=x_2=\cdots=x_n\}$，$V_2=\{\vv x\in\mathbb{F}^n\mid x_1=0\}$. 证明 $\mathbb{F}^n=V_1\oplus V_2$. 与第二章例~\ref{eg:pricedecomp} 的分解相比，哪一种分解更``自然''？（提示：思考补空间的不唯一性.）
			\item\label{ex:vs:B16} 在 $M_n(\mathbb{F})$ 中，设 $W_1$ 为对角线以下元素全为零的矩阵（上三角矩阵）全体，$W_2$ 为对角线及以上元素全为零的矩阵全体. 证明 $M_n(\mathbb{F})=W_1\oplus W_2$，并求 $\dim W_1$.
			\item\label{ex:vs:B17} （覆盖定理的另一个证明）设 $W_1,\ldots,W_s$ 是 $\mathbb{R}^n$ 的真子空间，令 $\vv v(t)=(1,t,t^2,\ldots,t^{n-1})^\trans$.
			\begin{enumerate}
				\item 证明：每个真子空间都包含于某个超平面 $\{\vv x\mid a_1x_1+\cdots+a_nx_n=0\}$（$a_i$ 不全为零）中.（提示：取 $W$ 的基并扩充为 $\mathbb{R}^n$ 的基，看``最后一个坐标''.）
				\item 证明：对上述超平面，满足 $\vv v(t)$ 落在其中的实数 $t$ 至多有 $n-1$ 个.
				\item 由此推出：$\vv v(t)$ 落在 $W_1\cup\cdots\cup W_s$ 中的 $t$ 至多有 $s(n-1)$ 个. 这给出了覆盖定理在 $\mathbb{R}^n$ 中的一个``构造性''证明. 对 $n=3$ 以及三个子空间 $x_1=x_2$，$x_2=x_3$，$x_1+x_2+x_3=0$，写出一个不在它们之中的 $\vv v(t)$.
			\end{enumerate}
			\item\label{ex:vs:B18} 设 $\dim V=n$，$W_1,\ldots,W_s$ 是 $V$ 的真子空间. 证明：存在 $V$ 的一组基 $\vv\varepsilon_1,\ldots,\vv\varepsilon_n$，其中\textbf{每个}向量都不属于 $W_1\cup\cdots\cup W_s$.（提示：逐个选取. 选 $\vv\varepsilon_k$ 时，要躲开的除了 $W_1,\ldots,W_s$，还有哪个真子空间？）
		\end{exB}
		
		\Cgroup
		\begin{exC}
			\item\label{ex:vs:C1} （覆盖定理在有限``数系''上的失效）设 $p$ 是素数，$\mathbb{F}_p=\{0,1,\ldots,p-1\}$ 按模 $p$ 作加法和乘法（可以承认它对加、减、乘以及除以非零元素都封闭. 它不包含 $\mathbb{Q}$，所以不是本书意义下的数域；但第二章中除定理~\ref{thm:Qsmallest} 与覆盖定理以外的定义和结论，在它上面照样成立）.
			\begin{enumerate}
				\item 证明 $\mathbb{F}_p^2$ 恰有 $p+1$ 个一维子空间，并且它们的并是整个 $\mathbb{F}_p^2$.
				\item 证明 $\mathbb{F}_p^2$ 不能写成 $p$ 个真子空间的并.（提示：每个一维子空间恰含 $p-1$ 个非零向量，数一数.）
				\item 指出定理~\ref{thm:cover} 的证明在 $\mathbb{F}_p$ 上从哪一步开始失效. 由第 2 问，``个数不超过 $p$ 时盖不住''，你能把正文的证明改写成这个``定量''的版本吗？
			\end{enumerate}
			\item\label{ex:vs:C2} 证明 $\ln2,\ln3,\ln5$ 在 $\mathbb{Q}$ 上线性无关. 由此说明 $\mathbb{R}$ 作为 $\mathbb{Q}$ 上的向量空间不是有限维的.（提示：算术基本定理.）
			\item\label{ex:vs:C3} 证明：全体实数列构成的向量空间不是有限维的；而``只有有限项非零''的实数列全体是它的子空间，且 $\vv e_1,\vv e_2,\ldots$（第 $i$ 项为 $1$，其余为 $0$）是它的一组``基''——每个元素都能唯一地写成其中\textbf{有限个}向量的线性组合. 这组向量为什么不是全体实数列的基？
			\item\label{ex:vs:C4} （模律）设 $W_1,W_2,W_3$ 是 $V$ 的子空间. 证明 $(W_1\cap W_3)+(W_2\cap W_3)\subset(W_1+W_2)\cap W_3$，在 $\mathbb{R}^2$ 中举例说明等号可以不成立；再证明：若 $W_1\subset W_3$，则等号一定成立.
			\item\label{ex:vs:C5} （有限域上的向量空间）在只含 $0,1$ 两个元素、按模 $2$ 运算的``数系'' $\mathbb{F}_2$ 上考虑 $\mathbb{F}_2^3$. (1) 它有多少个向量？(2) 有多少个一维子空间、多少个二维子空间？(3) 它有多少组（有序）基？推广到 $\mathbb{F}_2^n$.
			\item\label{ex:vs:C6} （塔律）设 $\mathbb{F}\subset\mathbb{K}\subset\mathbb{L}$ 是三个数域，$\mathbb{L}$ 作为 $\mathbb{K}$ 上的向量空间是 $m$ 维的，$\mathbb{K}$ 作为 $\mathbb{F}$ 上的向量空间是 $n$ 维的. 证明 $\mathbb{L}$ 作为 $\mathbb{F}$ 上的向量空间是 $mn$ 维的. 用它重新计算 B8 中数域的维数.
			\item\label{ex:vs:C7} 证明 $\mathbb{Q}(\sqrt2+\sqrt3)=\mathbb{Q}(\sqrt2,\sqrt3)$，即包含 $\sqrt2+\sqrt3$ 的最小数域与包含 $\sqrt2,\sqrt3$ 的最小数域相同.（提示：计算 $(\sqrt2+\sqrt3)^3$.）
			\item\label{ex:vs:C8} 设 $V$ 是 $n$ 维空间，$W$ 是 $k$ 维子空间（$0<k<n$）. 证明 $W$ 有无穷多个补空间. 在 $\mathbb{F}_2^2$ 中，一维子空间有几个补空间？
			\item\label{ex:vs:C9} （Fibonacci 空间的另一组基）在 B6 的记号下，改为研究满足 $a_{n+2}=a_{n+1}+a_n$ 的数列空间. 求出两个公比 $q$ 使 $\{q^n\}$ 属于这个空间，证明它们构成一组基，并推出 Fibonacci 数列的通项（Binet 公式）. 第八章会从另一个角度再次得到它.
		\end{exC}
	}
	
	% =====================================================================
	% 第三章　矩阵
	% =====================================================================
	\newcommand{\ExMatrix}{%
		\exsetheading
		\Agroup
		\begin{exA}
			\item\label{ex:mat:A1} 设 $\vv A=\begin{pmatrix}1&-1&2\\0&3&1\end{pmatrix}$，$\vv B=\begin{pmatrix}2&1\\-1&0\\1&4\end{pmatrix}$. 计算 $\vv{AB}$，$\vv{BA}$，$\vv A^\trans\vv B^\trans$，并验证 $(\vv{AB})^\trans=\vv B^\trans\vv A^\trans$.
			\item\label{ex:mat:A2} 按``列的线性组合''计算 $\begin{pmatrix}1&2&0\\3&-1&4\end{pmatrix}\begin{pmatrix}2\\-1\\1\end{pmatrix}$；再说明 $\vv e_i^\trans\vv A\vv e_j$ 是 $\vv A$ 的哪个元素，$\vv A\vv e_j$、$\vv e_i^\trans\vv A$ 又分别是什么.
			\item\label{ex:mat:A3} 计算 $\begin{pmatrix}1&1\\0&1\end{pmatrix}^n$，$\begin{pmatrix}\lambda&1\\0&\lambda\end{pmatrix}^n$，$\begin{pmatrix}\cos\theta&-\sin\theta\\ \sin\theta&\cos\theta\end{pmatrix}^n$，$\begin{pmatrix}0&1&0\\0&0&1\\0&0&0\end{pmatrix}^n$.
			\item\label{ex:mat:A4} 设 $\vv A=\begin{pmatrix}1&2\\0&1\end{pmatrix}$，$\vv B=\begin{pmatrix}1&0\\3&1\end{pmatrix}$. 计算 $(\vv A+\vv B)^2$ 与 $\vv A^2+2\vv{AB}+\vv B^2$，二者是否相等？$(\vv A+\vv B)(\vv A-\vv B)$ 与 $\vv A^2-\vv B^2$ 呢？
			\item\label{ex:mat:A5} 用初等行变换求下列矩阵的逆：
			\[\begin{pmatrix}1&2&3\\2&2&1\\3&4&3\end{pmatrix},\qquad\begin{pmatrix}1&1&1\\0&1&1\\0&0&1\end{pmatrix},\qquad \begin{pmatrix}2&1&0\\1&2&1\\0&1&2\end{pmatrix},\qquad\begin{pmatrix}0&0&0&1\\0&0&2&0\\0&3&0&0\\4&0&0&0\end{pmatrix}.\]
			\item\label{ex:mat:A6} 解矩阵方程：(1) $\begin{pmatrix}2&5\\1&3\end{pmatrix}\vv X=\begin{pmatrix}4&-6\\2&1\end{pmatrix}$；(2) $\vv X\begin{pmatrix}2&1\\1&1\end{pmatrix}=\begin{pmatrix}1&0\\0&2\end{pmatrix}$.
			\item\label{ex:mat:A7} 求矩阵 $\begin{pmatrix}1&1&a\\1&a&1\\a&1&1\end{pmatrix}$ 的秩（对 $a$ 分情况讨论）.
			\item\label{ex:mat:A8} 求矩阵 $\begin{pmatrix}1&2&-1&0&3\\2&4&1&3&0\\-1&-2&2&1&-3\\1&2&2&3&0\end{pmatrix}$ 的秩，以及列向量组的一个极大无关组，并用它表出其余各列.
			\item\label{ex:mat:A9} 写出 $3$ 阶初等矩阵 $\vv E(1,3)$，$\vv E(2(5))$，$\vv E(3,1(-2))$，并分别写出它们的逆. 计算 $\vv E(3,1(-2))\vv A$ 与 $\vv A\vv E(3,1(-2))$，其中 $\vv A$ 是任一 $3$ 阶方阵，说明效果.
			\item\label{ex:mat:A10} 把 $\begin{pmatrix}2&1\\4&3\end{pmatrix}$ 写成若干个初等矩阵的乘积.
			\item\label{ex:mat:A11} 证明：两个上三角矩阵之积仍是上三角矩阵，且积的对角元等于两者对角元之积. 可逆上三角矩阵的逆也是上三角矩阵吗？
			\item\label{ex:mat:A12} 设 $\vv A,\vv B$ 都是 $n$ 阶对称矩阵. 证明：(1) $\vv A^k$，$\vv A+\vv B$，$\vv P^\trans\vv A\vv P$ 都对称；(2) $\vv{AB}$ 对称当且仅当 $\vv{AB}=\vv{BA}$；(3) 对任意 $m\times n$ 矩阵 $\vv C$，$\vv C^\trans\vv C$ 与 $\vv C\vv C^\trans$ 都对称.
			\item\label{ex:mat:A13} (1) 设 $\vv x,\vv y\in\mathbb{F}^n$，证明 $\tr(\vv x\vv y^\trans)=\vv y^\trans\vv x$；(2) 设 $n$ 阶方阵 $\vv A$ 满足：对一切 $n$ 阶方阵 $\vv X$ 都有 $\tr(\vv{AX})=0$. 证明 $\vv A=\vv O$（提示：取 $\vv X=\vv E_{ji}$）；(3) 证明：不存在 $n$ 阶方阵 $\vv A,\vv B$ 使 $\vv{AB}-\vv{BA}=\vv I$.
			\item\label{ex:mat:A14} 设方阵 $\vv A$ 满足 $\vv A^2-\vv A-2\vv I=\vv O$. 证明 $\vv A$ 与 $\vv A+2\vv I$ 都可逆，并用 $\vv A$ 表示它们的逆. $\vv A+\vv I$ 一定可逆吗？
			\item\label{ex:mat:A15} 证明：若 $\vv A$ 可逆且 $\vv{AB}=\vv{AC}$，则 $\vv B=\vv C$；举例说明 $\vv A$ 不可逆时结论可能不成立.
			\item\label{ex:mat:A16} 设 $\vv A,\vv B$ 可逆. 求 $\begin{pmatrix}\vv A&\vv O\\ \vv O&\vv B\end{pmatrix}^{-1}$ 与 $\begin{pmatrix}\vv O&\vv A\\ \vv B&\vv O\end{pmatrix}^{-1}$，并用后者求 $\begin{pmatrix}0&0&1&2\\0&0&1&3\\2&0&0&0\\0&5&0&0\end{pmatrix}^{-1}$.
			\item\label{ex:mat:A17} 用分块乘法计算 $\begin{pmatrix}1&0&2&1\\0&1&1&3\\0&0&1&0\\0&0&0&1\end{pmatrix}^2$，并猜出 $n$ 次方的一般形式.
			\item\label{ex:mat:A18} 某公司在两个工厂生产三种产品，各工厂每天的产量（件）与三种产品的单位成本、单位售价（元）分别为
			\[\vv Q=\begin{pmatrix}100&80&60\\120&50&90\end{pmatrix}\ (\text{行：工厂；列：产品}),\qquad \vv C=\begin{pmatrix}20&30\\15&25\\40&55\end{pmatrix}\ (\text{列：成本、售价}).\]
			计算 $\vv Q\vv C$，说明它每个元素的含义，并求两个工厂每天的利润.
			\item\label{ex:mat:A19} 四座城市之间的直达航线为：$1\leftrightarrow2$，$1\leftrightarrow3$，$2\leftrightarrow3$，$3\leftrightarrow4$（双向）. 写出邻接矩阵 $\vv A$，计算 $\vv A^2$ 与 $\vv A^3$. $(\vv A^2)_{ii}$ 与城市 $i$ 的航线数有什么关系？从城市 $1$ 到城市 $4$ 恰好乘三次航班有几种走法？
			\item\label{ex:mat:A20} 设 $\vv A$ 是实反对称矩阵. 证明：对一切 $\vv x\in\mathbb{R}^n$ 有 $\vv x^\trans\vv A\vv x=0$；$\vv A^2$ 是对称矩阵.
			\item\label{ex:mat:A21} 证明：若 $\vv A$ 是 $m\times n$ 矩阵，$\vv P$ 是 $m$ 阶可逆矩阵，则 $\vv{PA}$ 与 $\vv A$ 的简化阶梯形相同.
			\item\label{ex:mat:A22} 证明：任一 $n$ 阶方阵都可以唯一地写成一个对称矩阵与一个反对称矩阵之和.（提示：$\vv A=\frac12(\vv A+\vv A^\trans)+\frac12(\vv A-\vv A^\trans)$.）
		\end{exA}
		
		\Bgroup
		\begin{exB}
			\item\label{ex:mat:B1} 分别求出全部满足 (1) $\vv A^2=\vv O$；(2) $\vv A^2=\vv I$；(3) $\vv A^2=\vv A$ 的二阶实方阵.
			\item\label{ex:mat:B2} 设 $a_1,\ldots,a_n$ 互不相同. 证明：与 $\diag(a_1,\ldots,a_n)$ 可交换的矩阵一定是对角矩阵.
			\item\label{ex:mat:B3} 证明：与\textbf{一切} $n$ 阶方阵都可交换的矩阵必为数量矩阵 $c\vv I$.（提示：考虑与 $\vv E_{ij}$ 交换.）
			\item\label{ex:mat:B4} 设 $\vv A,\vv B$ 都是幂零矩阵且 $\vv{AB}=\vv{BA}$. 证明 $\vv A+\vv B$ 与 $\vv{AB}$ 都是幂零矩阵. 举例说明``可交换''不能去掉.
			\item\label{ex:mat:B5} 证明：$\rk\begin{pmatrix}\vv A&\vv O\\ \vv O&\vv B\end{pmatrix}=\rk\vv A+\rk\vv B$，$\rk\begin{pmatrix}\vv A&\vv C\\ \vv O&\vv B\end{pmatrix}\ge\rk\vv A+\rk\vv B$，并举例说明后者可以取严格不等号.
			\item\label{ex:mat:B6} （Sherman–Morrison 公式）设 $\vv u,\vv v\in\mathbb{F}^n$，$1+\vv v^\trans\vv u\ne0$. 证明
			\[(\vv I+\vv u\vv v^\trans)^{-1}=\vv I-\frac{\vv u\vv v^\trans}{1+\vv v^\trans\vv u}.\]
			\item\label{ex:mat:B7} 设 $\vv A$ 为 $m\times n$ 矩阵，$\vv B$ 为 $n\times m$ 矩阵. 证明：若 $\vv I_m+\vv{AB}$ 可逆，则 $\vv I_n+\vv{BA}$ 也可逆，并且 $(\vv I_n+\vv{BA})^{-1}=\vv I_n-\vv B(\vv I_m+\vv{AB})^{-1}\vv A$.
			\item\label{ex:mat:B8} 证明：秩为 $1$ 的 $m\times n$ 矩阵都可以写成 $\vv u\vv v^\trans$（$\vv u\in\mathbb{F}^m$，$\vv v\in\mathbb{F}^n$ 非零）. 对这样的矩阵证明 $\vv A^2=(\tr\vv A)\vv A$（$m=n$ 时），并由此求 $\vv A^k$.
			\item\label{ex:mat:B9} 证明：秩为 $r$ 的矩阵可以写成 $r$ 个秩为 $1$ 的矩阵之和，但不能写成少于 $r$ 个秩为 $1$ 的矩阵之和.
			\item\label{ex:mat:B10} （满秩分解）证明：秩为 $r$ 的 $m\times n$ 矩阵 $\vv A$ 可以写成 $\vv A=\vv{BC}$，其中 $\vv B$ 是列满秩的 $m\times r$ 矩阵，$\vv C$ 是行满秩的 $r\times n$ 矩阵. 对 A8 中的矩阵写出一个满秩分解.
			\item\label{ex:mat:B11} 计算 $\begin{pmatrix}1&1&1\\1&1&1\\1&1&1\end{pmatrix}^n$ 与 $\begin{pmatrix}2&1&0\\0&2&1\\0&0&2\end{pmatrix}^n$.（提示：后者写成 $2\vv I+\vv N$.）
			\item\label{ex:mat:B12} 证明：$n$ 阶严格上三角矩阵（对角线及以下全为零）$\vv N$ 满足 $\vv N^n=\vv O$. 更精确地，$\vv N^k$ 的非零元只可能出现在第 $k$ 条``上对角线''及其右上方.
			\item\label{ex:mat:B13} 设 $\vv A$ 为实矩阵. 证明：若 $\vv A^\trans\vv A=\vv O$，则 $\vv A=\vv O$；若 $\tr(\vv A^\trans\vv A)=0$，结论同样成立. 复矩阵呢？
			\item\label{ex:mat:B14} 每行、每列恰有一个 $1$、其余元素为 $0$ 的方阵称为\textbf{置换矩阵}. 证明：置换矩阵之积仍是置换矩阵，置换矩阵 $\vv P$ 满足 $\vv P^\trans\vv P=\vv I$；$\vv{PA}$ 与 $\vv{AP}$ 分别对 $\vv A$ 做了什么？
			\item\label{ex:mat:B15} 对 $\vv A=\begin{pmatrix}1&2&1\\2&4&3\\1&2&2\end{pmatrix}$，求可逆矩阵 $\vv P,\vv Q$ 使 $\vv{PAQ}$ 为相抵标准形.
			\item\label{ex:mat:B16} 设 $\vv A$ 是 $n$ 阶方阵，$\vv A^k=\vv O$. 证明 $\vv I+\vv A$ 可逆，并求其逆. 若 $\vv A^2=2\vv A$，$\vv I-\vv A$ 是否可逆？
		\end{exB}
		
		\Cgroup
		\begin{exC}
			\item\label{ex:mat:C1} 设 $f:M_n(\mathbb{F})\to\mathbb{F}$ 满足 $f(\vv A+\vv B)=f(\vv A)+f(\vv B)$，$f(k\vv A)=kf(\vv A)$，且 $f(\vv{AB})=f(\vv{BA})$ 对一切 $\vv A,\vv B$ 成立. 证明 $f=c\cdot\tr$ 对某个常数 $c$ 成立.（提示：计算 $f(\vv E_{ij})$.）
			\item\label{ex:mat:C2} 设 $n$ 阶方阵 $\vv A,\vv B$ 满足 $\vv A+\vv B=\vv{AB}$. 证明 $\vv{AB}=\vv{BA}$.（提示：考察 $(\vv I-\vv A)(\vv I-\vv B)$.）
			\item\label{ex:mat:C3} （Strassen 算法）对二阶分块矩阵的乘法，直接计算需要 $8$ 次块乘法. 验证下面的公式只需 $7$ 次：令
			$M_1=(A_{11}+A_{22})(B_{11}+B_{22})$，$M_2=(A_{21}+A_{22})B_{11}$，$M_3=A_{11}(B_{12}-B_{22})$，$M_4=A_{22}(B_{21}-B_{11})$，$M_5=(A_{11}+A_{12})B_{22}$，$M_6=(A_{21}-A_{11})(B_{11}+B_{12})$，$M_7=(A_{12}-A_{22})(B_{21}+B_{22})$，则
			$C_{11}=M_1+M_4-M_5+M_7$，$C_{12}=M_3+M_5$，$C_{21}=M_2+M_4$，$C_{22}=M_1-M_2+M_3+M_6$. 由此说明 $n=2^k$ 阶矩阵乘法可以用约 $n^{\log_27}\approx n^{2.81}$ 次数的乘法完成.
			\item\label{ex:mat:C4} 设 $\vv A$ 是幂等矩阵. 用满秩分解 $\vv A=\vv{BC}$ 证明 $\vv{CB}=\vv I_r$，进而证明 $\rk\vv A=\tr\vv A$.
			\item\label{ex:mat:C5} （Kronecker 积）对 $m\times n$ 矩阵 $\vv A$ 与 $p\times q$ 矩阵 $\vv B$，定义 $\vv A\otimes\vv B=(a_{ij}\vv B)$（$mp\times nq$ 分块矩阵）. 证明 $(\vv A\otimes\vv B)(\vv C\otimes\vv D)=(\vv{AC})\otimes(\vv{BD})$，并由此说明 $\vv A,\vv B$ 可逆时 $(\vv A\otimes\vv B)^{-1}=\vv A^{-1}\otimes\vv B^{-1}$.
			\item\label{ex:mat:C6} （Hadamard 矩阵）元素全为 $\pm1$ 且满足 $\vv H\vv H^\trans=n\vv I$ 的 $n$ 阶矩阵称为 Hadamard 矩阵. 验证 $\vv H_2=\begin{pmatrix}1&1\\1&-1\end{pmatrix}$ 是，并证明 $\vv H_2\otimes\vv H_2$，$\vv H_2\otimes\vv H_2\otimes\vv H_2$ 也是. 证明：$n>2$ 时 Hadamard 矩阵的阶数必是 $4$ 的倍数.（这类矩阵用于实验设计与纠错编码.）
			\item\label{ex:mat:C7} 证明：全体形如 $\begin{pmatrix}a&-b\\b&a\end{pmatrix}$（$a,b\in\mathbb{R}$）的矩阵关于矩阵的加法和乘法``像复数一样''运算：加法、乘法封闭，乘法可交换，非零元可逆. 写出它与 $\mathbb{C}$ 之间的对应，并用它验证 $(1+\mathrm i)^8=16$.
			\item\label{ex:mat:C8} 设 $\vv A$ 是元素非负、各列之和都小于 $1$ 的方阵. 证明 $\vv A^k$ 的每个元素都随 $k\to\infty$ 趋于 $0$，并且级数 $\vv I+\vv A+\vv A^2+\cdots$（按元素）收敛到 $(\vv I-\vv A)^{-1}$. 解释它在投入产出模型中的意义：``直接消耗 $+$ 一次间接消耗 $+$ 二次间接消耗 $+\cdots$''.
			\item\label{ex:mat:C9} 设 $\vv A$ 为 $m\times n$ 矩阵. 证明：$\vv A$ 有右逆（存在 $\vv C$ 使 $\vv{AC}=\vv I_m$）当且仅当 $\rk\vv A=m$；有左逆当且仅当 $\rk\vv A=n$. 非方阵的单侧逆唯一吗？
		\end{exC}
	}
	
	% =====================================================================
	% 第四章　线性映射
	% =====================================================================
	\newcommand{\ExLinearMap}{%
		\exsetheading
		\Agroup
		\begin{exA}
			\item\label{ex:lmap:A1} 判断下列映射是否线性，是线性映射的求出核：
			\begin{enumerate}
				\item $\mathbb{R}^2\to\mathbb{R}^2$，$(x,y)^\trans\mapsto(x+y,\,2x)^\trans$；
				\item $\mathbb{R}^2\to\mathbb{R}^2$，$(x,y)^\trans\mapsto(x+1,\,y)^\trans$；
				\item $\mathbb{R}^2\to\mathbb{R}$，$(x,y)^\trans\mapsto|x|+y$；
				\item $\mathbb{R}[x]\to\mathbb{R}[x]$，$f(x)\mapsto f(x+1)$；
				\item $\mathbb{R}[x]\to\mathbb{R}[x]$，$f(x)\mapsto xf'(x)$；
				\item $M_n(\mathbb{R})\to M_n(\mathbb{R})$，$\vv X\mapsto\vv X^\trans$；
				\item $M_2(\mathbb{R})\to\mathbb{R}$，$\vv X\mapsto x_{11}x_{22}-x_{12}x_{21}$.
			\end{enumerate}
			\item\label{ex:lmap:A2} 求下列矩阵所定义的线性映射 $\vv x\mapsto\vv{Ax}$ 的核与像（各给出一组基），并验证秩—零化度定理：
			\[\begin{pmatrix}1&2&-1\\2&4&-2\end{pmatrix},\qquad\begin{pmatrix}1&0&2&1\\0&1&-1&1\\1&1&1&2\end{pmatrix}.\]
			\item\label{ex:lmap:A3} 写出线性映射 $\mathscr T:\mathbb{R}^3\to\mathbb{R}^2$，$\mathscr T(x_1,x_2,x_3)^\trans=(x_1-2x_2+x_3,\ 3x_2+4x_3)^\trans$ 在标准基下的矩阵.
			\item\label{ex:lmap:A4} 求下列平面变换在标准基下的矩阵：(1) 逆时针旋转 $90^\circ$；(2) 关于直线 $y=x$ 的反射；(3) 向直线 $y=x$ 的正投影；(4) 沿 $x$ 方向的剪切 $(x,y)\mapsto(x+2y,y)$. 计算 (2)(3) 的矩阵的平方，并解释结果.
			\item\label{ex:lmap:A5} 求求导变换 $\mathscr D:\mathbb{R}[x]_3\to\mathbb{R}[x]_3$ 在基 $1,x,x^2$ 下的矩阵，以及在基 $1,\ x-1,\ (x-1)^2$ 下的矩阵.
			\item\label{ex:lmap:A6} 求转置变换 $\vv X\mapsto\vv X^\trans$ 在 $M_2(\mathbb{F})$ 的基 $\vv E_{11},\vv E_{12},\vv E_{21},\vv E_{22}$ 下的矩阵. 它的平方是什么？
			\item\label{ex:lmap:A7} 在 $\mathbb{R}^3$ 中，求从基 $\vv\varepsilon_1=(1,0,0)^\trans$，$\vv\varepsilon_2=(1,1,0)^\trans$，$\vv\varepsilon_3=(1,1,1)^\trans$ 到基 $\vv\eta_1=(1,2,1)^\trans$，$\vv\eta_2=(2,3,3)^\trans$，$\vv\eta_3=(3,7,1)^\trans$ 的过渡矩阵，并求 $\vv\alpha=(3,4,2)^\trans$ 在两组基下的坐标.
			\item\label{ex:lmap:A8} 设线性变换 $\mathscr T$ 在基 $\vv\varepsilon_1,\vv\varepsilon_2$ 下的矩阵为 $\begin{pmatrix}2&1\\1&2\end{pmatrix}$. 求它在基 $\vv\eta_1=\vv\varepsilon_1+\vv\varepsilon_2$，$\vv\eta_2=\vv\varepsilon_1-\vv\varepsilon_2$ 下的矩阵.
			\item\label{ex:lmap:A9} 证明：线性映射的复合仍是线性映射；可逆（双射）线性映射的逆映射仍是线性映射.
			\item\label{ex:lmap:A10} 证明：线性映射 $\mathscr T$ 是单射，当且仅当它把每个线性无关组都映为线性无关组.
			\item\label{ex:lmap:A11} 是否存在线性变换 $\mathscr T:\mathbb{R}^2\to\mathbb{R}^2$，使
			(1) $\mathscr T(1,1)^\trans=(1,0)^\trans$，$\mathscr T(1,-1)^\trans=(0,1)^\trans$，$\mathscr T(2,0)^\trans=(1,1)^\trans$？
			(2) 同上，但最后一个条件改为 $\mathscr T(2,0)^\trans=(1,2)^\trans$？
			存在时求出 $\mathscr T(x,y)^\trans$.
			\item\label{ex:lmap:A12} (1) 线性映射 $\mathscr T:\mathbb{R}^5\to\mathbb{R}^3$ 是满射，$\dim\ker\mathscr T$ 是多少？(2) 是否存在从 $\mathbb{R}^3$ 到 $\mathbb{R}^5$ 的满射线性映射？是否存在从 $\mathbb{R}^5$ 到 $\mathbb{R}^3$ 的单射线性映射？
			\item\label{ex:lmap:A13} 设 $\mathscr T:U\to V$，$\mathscr S:V\to W$ 线性. 证明 $\ker\mathscr T\subset\ker(\mathscr{ST})$，$\imm(\mathscr{ST})\subset\imm\mathscr S$，从而 $\rk(\mathscr{ST})\le\min\{\rk\mathscr S,\rk\mathscr T\}$.
			\item\label{ex:lmap:A14} 求迹映射 $\tr:M_n(\mathbb{F})\to\mathbb{F}$ 的核的维数.
			\item\label{ex:lmap:A15} 证明：相似矩阵有相同的迹（利用命题~\ref{prop:trace}(3)）. 由此判断 $\begin{pmatrix}1&2\\0&3\end{pmatrix}$ 与 $\begin{pmatrix}1&0\\0&2\end{pmatrix}$ 是否相似.
			\item\label{ex:lmap:A16} 证明：与单位矩阵相似的矩阵只有单位矩阵本身. 由此说明 $\begin{pmatrix}1&1\\0&1\end{pmatrix}$ 与 $\vv I$ 不相似.
			\item\label{ex:lmap:A17} 定义求值映射 $\mathscr E:\mathbb{R}[x]_3\to\mathbb{R}^3$，$f\mapsto(f(0),f(1),f(2))^\trans$. 求它在基 $1,x,x^2$ 与标准基下的矩阵，并证明 $\mathscr E$ 是同构.
			\item\label{ex:lmap:A18} 设 $\mathscr T$ 是 $\mathbb{R}^3$ 上的线性变换，$\mathscr T\vv e_1=(1,1,0)^\trans$，$\mathscr T\vv e_2=(0,1,1)^\trans$，$\mathscr T\vv e_3=(1,2,1)^\trans$. 求 $\mathscr T(x_1,x_2,x_3)^\trans$，$\ker\mathscr T$ 与 $\imm\mathscr T$.
			\item\label{ex:lmap:A19} 称 $\vv x\mapsto\vv{Ax}+\vv b$ 为\textbf{仿射变换}. 证明两个仿射变换的复合仍是仿射变换；仿射变换是线性变换当且仅当 $\vv b=\zero$.
			\item\label{ex:lmap:A20} 设 $\vv R(\theta)=\begin{pmatrix}\cos\theta&-\sin\theta\\ \sin\theta&\cos\theta\end{pmatrix}$ 是平面上绕原点旋转角 $\theta$ 的变换在标准基下的矩阵. 从几何上说明``先转 $\beta$ 再转 $\alpha$''就是转 $\alpha+\beta$，从而 $\vv R(\alpha)\vv R(\beta)=\vv R(\alpha+\beta)$；比较两边矩阵的元素，推出两角和的正弦、余弦公式.
		\end{exA}
		
		\Bgroup
		\begin{exB}
			\item\label{ex:lmap:B1} 记 $L(V,W)$ 为 $V$ 到 $W$ 的全体线性映射. 证明它关于映射的加法与数乘构成向量空间，且 $\dim L(V,W)=\dim V\cdot\dim W$.（提示：取定基后，线性映射与矩阵一一对应.）
			\item\label{ex:lmap:B2} 设 $\mathscr T$ 是 $n$ 维空间 $V$ 上的线性变换，$\mathscr T^2=\mathscr O$. 证明 $\imm\mathscr T\subset\ker\mathscr T$，从而 $\rk\mathscr T\le n/2$；并证明存在一组基，使 $\mathscr T$ 的矩阵只在 $(1,2),(3,4),\ldots,(2r-1,2r)$ 这 $r$ 个位置上为 $1$，其余元素全为 $0$，其中 $r=\rk\mathscr T$.
			\item\label{ex:lmap:B3} 设 $\mathscr T,\mathscr S$ 是有限维空间上的线性变换. 证明 $\dim\ker(\mathscr{ST})\le\dim\ker\mathscr S+\dim\ker\mathscr T$.
			\item\label{ex:lmap:B4} 设 $V$ 有限维，$\mathscr S,\mathscr T$ 是 $V$ 上的线性变换且 $\mathscr{ST}=\mathscr I$. 证明 $\mathscr{TS}=\mathscr I$.
			\item\label{ex:lmap:B5} 在无限维空间 $\mathbb{R}[x]$ 上，令 $\mathscr D$ 为求导，$\mathscr J$ 为 $f\mapsto\int_0^xf(t)\,\mathrm dt$. 证明 $\mathscr{DJ}=\mathscr I$，但 $\mathscr{JD}\ne\mathscr I$；$\mathscr D$ 是满射但不是单射，$\mathscr J$ 是单射但不是满射. 这与 B4 是否矛盾？
			\item\label{ex:lmap:B6} 固定 $\vv A\in M_n(\mathbb{F})$，令 $\mathscr T(\vv X)=\vv{AX}-\vv{XA}$. 证明 $\imm\mathscr T$ 含于迹为零的矩阵全体，从而 $\mathscr T$ 不是满射；并证明 $\ker\mathscr T$ 包含 $\vv I,\vv A,\vv A^2,\ldots$.
			\item\label{ex:lmap:B7} 设 $\vv A=\begin{pmatrix}a&b\\c&d\end{pmatrix}$. 求线性变换 $\vv X\mapsto\vv{AX}$ 在 $M_2(\mathbb{F})$ 的基 $\vv E_{11},\vv E_{21},\vv E_{12},\vv E_{22}$ 下的矩阵. 若换成 $\vv X\mapsto\vv{XA}$ 呢？
			\item\label{ex:lmap:B8} 设线性变换 $\mathscr P$ 满足 $\mathscr P^2=\mathscr P$. 证明存在一组基使 $\mathscr P$ 的矩阵为 $\begin{pmatrix}\vv I_r&\vv O\\ \vv O&\vv O\end{pmatrix}$，其中 $r=\rk\mathscr P$.
			\item\label{ex:lmap:B9} 设线性变换 $\mathscr T$ 满足 $\mathscr T^2=\mathscr I$. 证明 $V=\ker(\mathscr T-\mathscr I)\oplus\ker(\mathscr T+\mathscr I)$，从而存在一组基使 $\mathscr T$ 的矩阵为 $\diag(\vv I_p,-\vv I_q)$.（提示：$\vv\alpha=\frac12(\vv\alpha+\mathscr T\vv\alpha)+\frac12(\vv\alpha-\mathscr T\vv\alpha)$.）
			\item\label{ex:lmap:B10} 设 $\mathscr T:V\to W$ 线性，$U$ 是 $V$ 的子空间. 证明 $\dim\mathscr T(U)=\dim U-\dim(U\cap\ker\mathscr T)$.
			\item\label{ex:lmap:B11} 在由 $\mathrm e^x\cos x$，$\mathrm e^x\sin x$ 张成的二维函数空间中，求求导变换的矩阵，证明它可逆，并利用逆矩阵求出 $\int\mathrm e^x\sin x\,\mathrm dx$（不必用分部积分）.
			\item\label{ex:lmap:B12} （齐次坐标）平面上的平移 $(x,y)\mapsto(x+a,y+b)$ 不是线性的. 把点 $(x,y)$ 写成三维向量 $(x,y,1)^\trans$，证明平移、旋转、伸缩都可以写成 $3\times3$ 矩阵的乘法，且仿射变换的复合对应矩阵的乘积. 写出``先绕原点旋转 $90^\circ$，再平移 $(1,2)$''与``先平移再旋转''的矩阵，它们相等吗？（计算机图形学正是这样处理变换的.）
			\item\label{ex:lmap:B13} 证明：若 $\vv A$ 与 $\vv B$ 相似，则 $\vv A^k$ 与 $\vv B^k$ 相似，$\vv A^\trans$ 与 $\vv B^\trans$ 相似，$f(\vv A)$ 与 $f(\vv B)$ 相似（$f$ 为多项式）.
			\item\label{ex:lmap:B14} 证明 $\begin{pmatrix}1&0\\0&0\end{pmatrix}$ 与 $\begin{pmatrix}0&1\\0&0\end{pmatrix}$ 相抵但不相似. 从``线性映射 $V\to W$''与``线性变换 $V\to V$''的区别解释：为什么相抵只需看秩，而相似要复杂得多？
			\item\label{ex:lmap:B15} 设 $\mathscr T$ 是 $V$ 上的线性变换，$\vv\alpha\in V$ 满足 $\mathscr T^{k-1}\vv\alpha\ne\zero$，$\mathscr T^k\vv\alpha=\zero$. 证明 $\vv\alpha,\mathscr T\vv\alpha,\ldots,\mathscr T^{k-1}\vv\alpha$ 线性无关. 若 $k=n=\dim V$，求 $\mathscr T$ 在基 $\mathscr T^{n-1}\vv\alpha,\ldots,\mathscr T\vv\alpha,\vv\alpha$ 下的矩阵.
		\end{exB}
		
		\Cgroup
		\begin{exC}
			\item\label{ex:lmap:C1} （Fitting 分解）设 $\mathscr T$ 是 $n$ 维空间 $V$ 上的线性变换. 证明：
			\begin{enumerate}
				\item $\ker\mathscr T\subset\ker\mathscr T^2\subset\cdots$，$\imm\mathscr T\supset\imm\mathscr T^2\supset\cdots$，且两条链都在某一步之后不再变化；
				\item $V=\ker\mathscr T^n\oplus\imm\mathscr T^n$，两者都是 $\mathscr T$ 的不变子空间，$\mathscr T$ 在前者上幂零，在后者上可逆.
			\end{enumerate}
			\item\label{ex:lmap:C2} （商空间）设 $W$ 是 $V$ 的子空间. 对 $\vv\alpha\in V$，称集合 $\vv\alpha+W$ 为 $\vv\alpha$ 所在的陪集，全体陪集记为 $V/W$. 定义 $(\vv\alpha+W)+(\vv\beta+W)=(\vv\alpha+\vv\beta)+W$，$k(\vv\alpha+W)=k\vv\alpha+W$.
			\begin{enumerate}
				\item 证明这两个运算是良定义的，且 $V/W$ 构成向量空间；
				\item 证明 $\dim V/W=\dim V-\dim W$；
				\item （第一同构定理/同态基本定理）对线性映射 $\mathscr T:V\to U$，证明 $V/\ker\mathscr T\cong\imm\mathscr T$. 这给出秩—零化度定理的另一种理解.
			\end{enumerate}
			\item\label{ex:lmap:C3} 把 $\mathbb{C}$ 看作 $\mathbb{R}$ 上的二维空间，基为 $1,\mathrm i$. (1) 求``乘以 $a+b\mathrm i$''这个变换的矩阵；(2) 求``取共轭''这个变换的矩阵；(3) 证明 $\mathbb{R}$-线性变换 $\begin{pmatrix}p&q\\r&s\end{pmatrix}$ 是``乘以某个复数''，当且仅当 $p=s$，$q=-r$.
			\item\label{ex:lmap:C4} 证明：$n$ 维空间上的任一线性变换 $\mathscr T$，都存在次数不超过 $n^2$ 的非零多项式 $f$ 使 $f(\mathscr T)=\mathscr O$.（提示：$\mathscr I,\mathscr T,\ldots,\mathscr T^{n^2}$ 在 $L(V,V)$ 中线性相关. 第八章的 Cayley–Hamilton 定理会把 $n^2$ 改进为 $n$.）
			\item\label{ex:lmap:C5} （Haar 小波）在 $\mathbb{R}^4$ 中取基 $\vv h_1=(1,1,1,1)^\trans$，$\vv h_2=(1,1,-1,-1)^\trans$，$\vv h_3=(1,-1,0,0)^\trans$，$\vv h_4=(0,0,1,-1)^\trans$. 求信号 $\vv x=(8,6,3,4)^\trans$ 在这组基下的坐标，并解释每个坐标的含义. 若把绝对值最小的坐标改为 $0$ 再变换回来，误差是多少？这是图像与音频压缩的思想之一.
			\item\label{ex:lmap:C6} 设 $\mathscr T$ 是 $V$ 上的线性变换. 若 $\mathscr T$ 与 $V$ 上的一切线性变换都可交换，证明 $\mathscr T=c\mathscr I$. 若只要求与一切\textbf{可逆}线性变换可交换呢？
			\item\label{ex:lmap:C7} 设 $\mathscr T:V\to V$ 线性，$\dim V=n$. 证明：$\rk\mathscr T^k-\rk\mathscr T^{k+1}$ 关于 $k$ 单调不增. 由此说明：若 $\rk\mathscr T=\rk\mathscr T^2$，则对一切 $k$，$\rk\mathscr T^k=\rk\mathscr T$.
			\item\label{ex:lmap:C8} （差分算子）在全体实数列构成的空间上，定义左移 $\mathscr S(a_1,a_2,a_3,\ldots)=(a_2,a_3,\ldots)$. 证明 $\mathscr S$ 是线性的、满射但不是单射；满足 $a_{n+2}=a_{n+1}+a_n$ 的数列全体恰好是 $\ker(\mathscr S^2-\mathscr S-\mathscr I)$. 由第二章章末习题~\ref{ex:vs:C9}，这个核是二维的.
		\end{exC}
	}
	
	% =====================================================================
	% 第五章　行列式
	% =====================================================================
	\newcommand{\ExDeterminant}{%
		\exsetheading
		\Agroup
		\begin{exA}
			\item\label{ex:det:A1} 求下列排列的逆序数并判断奇偶：$315426$，$42531$，$13\cdots(2n-1)\,24\cdots(2n)$.
			\item\label{ex:det:A2} (1) 在 $4$ 阶行列式中，项 $a_{13}a_{24}a_{32}a_{41}$ 带什么符号？(2) 选择 $i,j$，使 $a_{1i}a_{23}a_{35}a_{5j}a_{44}$ 成为 $5$ 阶行列式中带正号的一项.
			\item\label{ex:det:A3} 计算：
			\[\begin{vmatrix}2&1&-1\\3&0&2\\1&4&5\end{vmatrix},\qquad
			\begin{vmatrix}1&2&0&1\\2&1&1&0\\0&1&2&1\\1&0&1&2\end{vmatrix},\qquad
			\begin{vmatrix}3&1&1&1\\1&3&1&1\\1&1&3&1\\1&1&1&3\end{vmatrix}.\]
			\item\label{ex:det:A4} 计算 $\begin{vmatrix}1&-1&2&3\\2&0&1&-1\\0&1&-1&2\\3&1&0&1\end{vmatrix}$，分别用``打洞''化三角形与按第一列展开两种方法，比较哪种更快.
			\item\label{ex:det:A5} 计算 $n$ 阶行列式 $\begin{vmatrix}0&\cdots&0&a_1\\0&\cdots&a_2&0\\ \vdots&&&\vdots\\a_n&\cdots&0&0\end{vmatrix}$（只有副对角线上有元素）.
			\item\label{ex:det:A6} 不展开，证明：
			\[\begin{vmatrix}a+b&b+c&c+a\\a_1+b_1&b_1+c_1&c_1+a_1\\a_2+b_2&b_2+c_2&c_2+a_2\end{vmatrix}=2\begin{vmatrix}a&b&c\\a_1&b_1&c_1\\a_2&b_2&c_2\end{vmatrix}.\]
			\item\label{ex:det:A7} 解方程 $\begin{vmatrix}x&1&1\\1&x&1\\1&1&x\end{vmatrix}=0$ 与 $\begin{vmatrix}1&1&1\\1&2&x\\1&4&x^2\end{vmatrix}=0$.
			\item\label{ex:det:A8} 设 $\vv A$ 为 $3$ 阶方阵，$\det\vv A=3$. 求 $\det(2\vv A)$，$\det(-\vv A)$，$\det(\vv A^{-1})$，$\det(\vv A^*)$，$\det(\vv A^\trans\vv A)$，$\det(3\vv A^{-1}-2\vv A^*)$.
			\item\label{ex:det:A9} 证明：(1) 若 $n$ 阶方阵满足 $\vv A^\trans\vv A=\vv I$，则 $\det\vv A=\pm1$；(2) 若 $\vv A^2=\vv I$，则 $\det\vv A=\pm1$；(3) 若 $\vv A^2=\vv A$ 且 $\vv A\ne\vv I$，则 $\det\vv A=0$.
			\item\label{ex:det:A10} 求下列矩阵的伴随矩阵与逆矩阵：\[\begin{pmatrix}a&b\\c&d\end{pmatrix},\qquad\begin{pmatrix}1&2&3\\0&1&4\\5&6&0\end{pmatrix}.\]
			\item\label{ex:det:A11} 用 Cramer 法则解方程组 $2x_1+x_2+x_3=3$，$x_1+3x_2+x_3=0$，$x_1+x_2+4x_3=8$.
			\item\label{ex:det:A12} 求以 $(1,1)$，$(4,2)$，$(2,5)$ 为顶点的三角形的面积；判断 $(1,2)$，$(3,5)$，$(7,11)$ 三点是否共线.
			\item\label{ex:det:A13} 求以 $(1,0,1)^\trans$，$(2,1,0)^\trans$，$(0,3,1)^\trans$ 为棱的平行六面体的体积，以及以原点和这三个向量的终点为顶点的四面体的体积.
			\item\label{ex:det:A14} $\lambda$ 取何值时，齐次方程组 $(5-\lambda)x+2y+2z=0$，$2x+(6-\lambda)y=0$，$2x+(4-\lambda)z=0$ 有非零解？
			\item\label{ex:det:A15} 设 \[\vv A=\begin{pmatrix}1&2&3&4\\2&3&4&1\\3&4&1&2\\4&1&2&3\end{pmatrix}.\] 求 $A_{11}+A_{12}+A_{13}+A_{14}$（$A_{1j}$ 是代数余子式）.（提示：它是把第一行换成全 $1$ 后的行列式.）
			\item\label{ex:det:A16} 证明：过平面上两个不同点 $(x_1,y_1)$，$(x_2,y_2)$ 的直线方程可以写成 \[\begin{vmatrix}x&y&1\\x_1&y_1&1\\x_2&y_2&1\end{vmatrix}=0.\]
			\item\label{ex:det:A17} 设 $\vv A,\vv B$ 分别是 $m$ 阶、$n$ 阶方阵，证明 $\begin{vmatrix}\vv A&\vv O\\ \vv C&\vv B\end{vmatrix}=\det\vv A\det\vv B$，并计算 $\begin{vmatrix}\vv O&\vv A\\ \vv B&\vv O\end{vmatrix}$.
			\item\label{ex:det:A18} 计算 $4$ 阶三对角行列式 $\begin{vmatrix}2&-1&0&0\\-1&2&-1&0\\0&-1&2&-1\\0&0&-1&2\end{vmatrix}$.
			\item\label{ex:det:A19} 设 $\vv A$ 是 $n$ 阶可逆矩阵，$n\ge2$. 证明 $\vv A^*$ 可逆，且 $(\vv A^*)^{-1}=\frac1{\det\vv A}\vv A=(\vv A^{-1})^*$.
			\item\label{ex:det:A20} 用行列式判断：向量 $(1,2,3)^\trans$，$(4,5,6)^\trans$，$(7,8,9)^\trans$ 是否线性无关？把 $9$ 改成 $10$ 呢？
			\item\label{ex:det:A21} 设 $\vv A$ 是 $n$ 阶反对称矩阵（$\vv A^\trans=-\vv A$），$n$ 为奇数. 证明 $\det\vv A=0$.（提示：对 $\vv A^\trans=-\vv A$ 两边取行列式.）
		\end{exA}
		
		\Bgroup
		\begin{exB}
			\item\label{ex:det:B1} 计算 $n$ 阶三对角行列式 $D_n=\begin{vmatrix}2&-1&&\\-1&2&\ddots&\\&\ddots&\ddots&-1\\&&-1&2\end{vmatrix}$.（提示：按第一行展开得 $D_n=2D_{n-1}-D_{n-2}$.）
			\item\label{ex:det:B2} 计算 $n$ 阶行列式 $\begin{vmatrix}a&b&\cdots&b\\b&a&\cdots&b\\ \vdots&&\ddots&\vdots\\b&b&\cdots&a\end{vmatrix}$ 与 $\begin{vmatrix}x&a&\cdots&a\\-a&x&\cdots&a\\ \vdots&&\ddots&\vdots\\-a&-a&\cdots&x\end{vmatrix}$.
			\item\label{ex:det:B3} （爪形行列式）设 $a_1\cdots a_n\ne0$，计算 $\begin{vmatrix}a_0&1&1&\cdots&1\\1&a_1&0&\cdots&0\\1&0&a_2&\cdots&0\\ \vdots&&&\ddots&\\1&0&0&\cdots&a_n\end{vmatrix}$.
			\item\label{ex:det:B4} （矩阵行列式引理）证明 $\det(\vv I_n+\vv u\vv v^\trans)=1+\vv v^\trans\vv u$，更一般地，$\vv A$ 可逆时 $\det(\vv A+\vv u\vv v^\trans)=(1+\vv v^\trans\vv A^{-1}\vv u)\det\vv A$. 用它重做 B2 的第一个行列式.
			\item\label{ex:det:B5} 设 $\vv A$ 为 $2\times3$ 矩阵，$\vv B$ 为 $3\times2$ 矩阵. 证明 $\det(\vv{AB})$ 等于 $\vv A$ 的三个 $2$ 阶子式与 $\vv B$ 的对应 $2$ 阶子式乘积之和（这是 Cauchy–Binet 公式最简单的情形）. 由此说明 $\vv{BA}$（$3$ 阶方阵）的行列式为零.
			\item\label{ex:det:B6} 设 $\vv A$ 是整数矩阵. 证明：$\vv A^{-1}$ 存在且也是整数矩阵，当且仅当 $\det\vv A=\pm1$.
			\item\label{ex:det:B7} 证明：$n\ge2$ 时 $(\vv A^*)^*=(\det\vv A)^{n-2}\vv A$.（先设 $\vv A$ 可逆；不可逆的情形另行讨论.）
			\item\label{ex:det:B8} 设 $\vv A,\vv B$ 都是 $n$ 阶方阵. 证明 $\begin{vmatrix}\vv A&\vv B\\ \vv B&\vv A\end{vmatrix}=\det(\vv A+\vv B)\det(\vv A-\vv B)$.（提示：分块初等变换.）
			\item\label{ex:det:B9} 设 $\vv A$ 可逆. 证明 $\begin{vmatrix}\vv A&\vv B\\ \vv C&\vv D\end{vmatrix}=\det\vv A\cdot\det(\vv D-\vv{CA}^{-1}\vv B)$；若还有 $\vv{AC}=\vv{CA}$，则它等于 $\det(\vv{AD}-\vv{CB})$.
			\item\label{ex:det:B10} 用 Fibonacci 矩阵 $\vv F=\begin{pmatrix}1&1\\1&0\end{pmatrix}$ 满足 $\vv F^n=\begin{pmatrix}F_{n+1}&F_n\\F_n&F_{n-1}\end{pmatrix}$（归纳法证明），推出 Cassini 恒等式 $F_{n+1}F_{n-1}-F_n^2=(-1)^n$.
			\item\label{ex:det:B11} 计算 $n$ 阶行列式 $D_n=\begin{vmatrix}1+x^2&x&&\\x&1+x^2&\ddots&\\&\ddots&\ddots&x\\&&x&1+x^2\end{vmatrix}$.
			\item\label{ex:det:B12} 设 $\omega=\mathrm e^{2\pi\mathrm i/3}$. 证明
			\[\begin{vmatrix}a&b&c\\c&a&b\\b&c&a\end{vmatrix}=(a+b+c)(a+\omega b+\omega^2c)(a+\omega^2b+\omega c)=a^3+b^3+c^3-3abc.\]
			\item\label{ex:det:B13} 证明：元素为 $a_{ij}=\min(i,j)$ 的 $n$ 阶行列式等于 $1$.
			\item\label{ex:det:B14} 证明：平面上四点 $(x_i,y_i)$（$i=1,2,3,4$）共圆（或共线），当且仅当 $\begin{vmatrix}x_i^2+y_i^2&x_i&y_i&1\end{vmatrix}_{4\times4}=0$（每行对应一个点）.
						\item\label{ex:det:B15} 设 $\vv A$ 为 $n$ 阶方阵（$n\ge2$）. 证明
			\[\rk\vv A^*=\begin{cases}n,&\rk\vv A=n,\\1,&\rk\vv A=n-1,\\0,&\rk\vv A\le n-2.\end{cases}\]
			（提示：$\vv A^*$ 的元素都是 $\vv A$ 的 $n-1$ 阶子式，用定理~\ref{thm:minorrank}；$\rk\vv A=n-1$ 时考察 $\vv A\vv A^*=(\det\vv A)\vv I$.）
			\item\label{ex:det:B16} 设 $f(x)=\begin{vmatrix}x&1&2&3\\1&x&2&3\\1&2&x&3\\1&2&3&x\end{vmatrix}$. 不展开，求 $f(x)=0$ 的全部根.
		\end{exB}
		
		\Cgroup
		\begin{exC}
			\item\label{ex:det:C1} （行列式的唯一性）设 $D:M_n(\mathbb{F})\to\mathbb{F}$ 对每一行都是线性的、交换两行变号，且 $D(\vv I)=1$. 证明 $D=\det$.（提示：先说明 $D$ 在初等矩阵上的取值，再对任意矩阵用初等变换.）
			\item\label{ex:det:C2} （Cauchy 行列式）证明
			\[\det\Big(\frac1{a_i+b_j}\Big)_{n\times n}=\frac{\prod_{i<j}(a_j-a_i)(b_j-b_i)}{\prod_{i,j}(a_i+b_j)}.\]
			先对 $n=2$ 直接验证，再用``各行减去最后一行''归纳. 取 $a_i=i$，$b_j=j-1$，得到 Hilbert 矩阵 $(1/(i+j-1))$ 的行列式，它为什么在数值计算中臭名昭著？
			\item\label{ex:det:C3} （排列的符号是同态）对排列 $\sigma$，令 $\vv P_\sigma$ 为把 $\vv e_j$ 映到 $\vv e_{\sigma(j)}$ 的置换矩阵. 证明 $\vv P_\sigma\vv P_\tau=\vv P_{\sigma\tau}$，$\det\vv P_\sigma=(-1)^{\tau(\sigma)}$，并由行列式的乘积公式推出 $\operatorname{sgn}(\sigma\tau)=\operatorname{sgn}\sigma\cdot\operatorname{sgn}\tau$.
			\item\label{ex:det:C4} （矩阵—树定理）对一个连通无向图，令 $\vv D$ 为各顶点度数构成的对角矩阵，$\vv A$ 为邻接矩阵，$\vv L=\vv D-\vv A$ 称为 Laplace 矩阵. 定理说：图的生成树个数等于 $\vv L$ 的任一个 $(n-1)$ 阶主子式. 对 $4$ 个顶点的完全图验证它有 $16$ 棵生成树，对一个 $4$ 个顶点的环验证它有 $4$ 棵生成树.
			\item\label{ex:det:C5} 证明：若 $n$ 阶矩阵的元素全为 $\pm1$，则它的行列式能被 $2^{n-1}$ 整除. 求 $3$ 阶此类行列式的最大值.
			\item\label{ex:det:C6} （Jacobi 公式的雏形）证明 $\det(\vv I+t\vv B)=1+t\tr\vv B+(\text{$t^2$ 及更高次项})$. 由此说明：若 $\vv A(t)$ 的元素都是 $t$ 的可导函数，且 $\vv A(0)=\vv I$，则 $\frac{\mathrm d}{\mathrm dt}\det\vv A(t)\big|_{t=0}=\tr\vv A'(0)$.
			\item\label{ex:det:C7} （Sylvester 结式）设 $f=a_0x^2+a_1x+a_2$，$g=b_0x^2+b_1x+b_2$（$a_0b_0\ne0$）. 证明 $f,g$ 有公共根当且仅当
			\[\begin{vmatrix}a_0&a_1&a_2&0\\0&a_0&a_1&a_2\\b_0&b_1&b_2&0\\0&b_0&b_1&b_2\end{vmatrix}=0.\]
			（提示：有公共根当且仅当存在次数 $\le1$ 的不全为零的 $u,v$ 使 $uf+vg=0$，与第七章的 Bézout 等式对照.）
			\item\label{ex:det:C8} （Dodgson 凝聚法）对 $3$ 阶矩阵 $\vv A$，记 $M_{ij}$ 为余子式，$c=a_{22}$. 验证 $\det\vv A\cdot c=M_{11}M_{33}-M_{13}M_{31}$. 它允许只计算 $2$ 阶行列式就得到高阶行列式. 《爱丽丝漫游奇境》的作者刘易斯·卡罗尔（即 Dodgson）于 1866 年发表了这一方法.
			\item\label{ex:det:C9} 设 $\vv A$ 是 $n$ 阶实反对称矩阵. 证明 $\det(\vv I+\vv A)\ge1$.（提示：$\det(\vv I+\vv A)=\det(\vv I-\vv A)$，考察 $\det(\vv I+t\vv A)$ 作为 $t$ 的多项式只含偶次项，且 $\vv I+t\vv A$ 对一切实数 $t$ 都可逆.）
		\end{exC}
	}
	
	% =====================================================================
	% 第六章　线性方程组解的结构
	% =====================================================================
	\newcommand{\ExSolutions}{%
		\exsetheading
		\Agroup
		\begin{exA}
			\item\label{ex:sol:A1} 求下列齐次方程组的基础解系与通解：
			\[\text{(1)}\ \left\{\begin{array}{l}x_1+x_2-x_3-x_4=0,\\2x_1-5x_2+3x_3+2x_4=0,\\7x_1-7x_2+3x_3+x_4=0;\end{array}\right.\qquad
			\text{(2)}\ \left\{\begin{array}{l}x_1+2x_2+x_3-x_4=0,\\3x_1+6x_2-x_3-3x_4=0,\\5x_1+10x_2+x_3-5x_4=0.\end{array}\right.\]
			\item\label{ex:sol:A2} 求下列方程组的通解，写成``特解 $+$ 基础解系的线性组合''：
			\[\text{(1)}\ \left\{\begin{array}{l}x_1+x_2-3x_3-x_4=1,\\3x_1-x_2-3x_3+4x_4=4,\\x_1+5x_2-9x_3-8x_4=0;\end{array}\right.\qquad
			\text{(2)}\ \left\{\begin{array}{l}2x_1+x_2-x_3+x_4=1,\\4x_1+2x_2-2x_3+x_4=2,\\2x_1+x_2-x_3-x_4=1.\end{array}\right.\]
			\item\label{ex:sol:A3} 设 $\vv A$ 是 $4\times5$ 矩阵，$\rk\vv A=3$. (1) $\vv{Ax}=\zero$ 的基础解系含几个向量？(2) $\vv{Ax}=\vv b$ 是否一定有解？(3) 若 $\rk\vv A=4$ 呢？
			\item\label{ex:sol:A4} 不解方程，用秩判断下列方程组是否有解：(1) $x_1+x_2=1$，$x_1-x_2=1$，$2x_1+3x_2=2$；(2) $x_1+x_2+x_3=1$，$2x_1+2x_2+2x_3=3$.
			\item\label{ex:sol:A5} 设 $\vv\eta_1,\vv\eta_2$ 是 $\vv{Ax}=\vv b$（$\vv b\ne\zero$）的两个解. 下列向量中哪些是 $\vv{Ax}=\vv b$ 的解，哪些是 $\vv{Ax}=\zero$ 的解？
			$\vv\eta_1+\vv\eta_2$，$2\vv\eta_1-\vv\eta_2$，$\frac12(\vv\eta_1+\vv\eta_2)$，$\vv\eta_1-\vv\eta_2$，$3\vv\eta_1-2\vv\eta_2$.
			\item\label{ex:sol:A6} 设 $3$ 元方程组 $\vv{Ax}=\vv b$ 的系数矩阵秩为 $2$，$\vv\eta_1=(1,2,3)^\trans$ 与 $\vv\eta_2=(1,-2,-1)^\trans$ 是它的两个解. 求通解.
			\item\label{ex:sol:A7} 写出一个以 $\vv x=(1,0,2)^\trans+k(1,1,-1)^\trans$（$k$ 任意）为全部解的线性方程组.
			\item\label{ex:sol:A8} 判断三个平面 $x+y+z=1$，$x+2y+3z=2$，$2x+3y+4z=4$ 的位置关系；若把最后一个常数改为 $3$ 呢？
			\item\label{ex:sol:A9} 举例说明：$\vv{Ax}=\zero$ 只有零解时，$\vv{Ax}=\vv b$ 可能无解；但若有解，解必唯一.
			\item\label{ex:sol:A10} 解矩阵方程 $\begin{pmatrix}1&2&3\\2&2&1\\3&4&3\end{pmatrix}\vv X=\begin{pmatrix}2&5\\3&1\\4&3\end{pmatrix}$.
			\item\label{ex:sol:A11} 设 $\vv A=\begin{pmatrix}1&1\\2&2\end{pmatrix}$. 求全部满足 $\vv{AX}=\vv O$ 的二阶矩阵 $\vv X$，它们构成的空间是几维的？
			\item\label{ex:sol:A12} 设 $\vv A=\begin{pmatrix}1&2&-2\\4&t&3\\3&-1&1\end{pmatrix}$，且存在三阶非零矩阵 $\vv B$ 使 $\vv{AB}=\vv O$. 求 $t$，并求 $\rk\vv B$ 的可能值.
			\item\label{ex:sol:A13} 求方程 $x_1+x_2+x_3=1$ 的通解，并写出它的三个解，使得通解中的每个解都是这三个解的仿射组合.
			\item\label{ex:sol:A14} 把空间直线 $\{x+y+z=1,\ x-y=0\}$ 写成参数式；反过来，把直线 $\vv x=(1,2,0)^\trans+t(1,0,1)^\trans$ 写成两个平面的交线.
			\item\label{ex:sol:A15} 判断空间直线 $\vv x=(1,0,0)^\trans+s(1,1,0)^\trans$ 与 $\vv x=(0,1,1)^\trans+t(0,1,-1)^\trans$ 是否相交.（列出关于 $s,t$ 的方程组.）
			\item\label{ex:sol:A16} 物流网络：三个仓库 $A,B,C$ 之间有线路 $A\to B$（$x_1$），$B\to C$（$x_2$），$A\to C$（$x_3$）. $A$ 每天净发出 $30$ 吨，$C$ 每天净收到 $30$ 吨，$B$ 不囤货. 求全部方案. 哪些量（如 $x_1+x_3$，$x_1-x_2$，$x_3$）能被方程组唯一确定？用可确定性判据说明.
			\item\label{ex:sol:A17} 求 $\vv A=\begin{pmatrix}2&1&0\\4&3&1\\2&3&3\end{pmatrix}$ 的 LU 分解.
			\item\label{ex:sol:A18} 设 $\vv A$ 是 $m\times n$ 矩阵，$\rk\vv A=m<n$. 证明 $\vv{Ax}=\vv b$ 对每个 $\vv b$ 都有无穷多解.
			\item\label{ex:sol:A19} 求经过点 $(1,1,1)$，$(2,0,1)$，$(0,2,3)$ 的平面方程（用齐次方程组求法向量）.
			\item\label{ex:sol:A20} 设 $\vv A=\begin{pmatrix}1&1\\1&1.0001\end{pmatrix}$. 分别求 $\vv{Ax}=(2,2)^\trans$ 与 $\vv{Ax}=(2,2.0001)^\trans$ 的解，计算右端的相对变化与解的相对变化之比.
		\end{exA}
		
		\Bgroup
		\begin{exB}
			\item\label{ex:sol:B1} 求两个齐次方程组 $\{x_1+x_2=0,\ x_2-x_4=0\}$ 与 $\{x_1-x_2+x_3=0,\ x_2-x_3+x_4=0\}$ 的公共非零解.
			\item\label{ex:sol:B2} 设 $\vv A$ 是实矩阵. 证明 $\vv{Ax}=\zero$ 与 $\vv A^\trans\vv{Ax}=\zero$ 同解，从而 $\rk(\vv A^\trans\vv A)=\rk\vv A$.（提示：$\vv x^\trans\vv A^\trans\vv{Ax}=\|\vv{Ax}\|^2$ 在第九章才有名字，但这里只需``实数的平方和为零则每个都为零''.）
			\item\label{ex:sol:B3} 证明 $\rk(\vv A+\vv B)\le\rk(\vv A\mid\vv B)\le\rk\vv A+\rk\vv B$.
			\item\label{ex:sol:B4} 设 $\vv A$ 为 $n$ 阶方阵，$\vv A^2=\vv I$. 证明 $\rk(\vv A+\vv I)+\rk(\vv A-\vv I)=n$.
			\item\label{ex:sol:B5} 证明 $\rk(\vv{AB})=\rk\vv B$ 当且仅当 $\ker\vv A\cap\imm\vv B=\{\zero\}$.
			\item\label{ex:sol:B6} 设 $\vv\eta^*$ 是 $\vv{Ax}=\vv b$（$\vv b\ne\zero$）的一个解，$\vv\xi_1,\ldots,\vv\xi_k$ 是对应齐次方程组的基础解系. 证明 $\vv\eta^*,\vv\eta^*+\vv\xi_1,\ldots,\vv\eta^*+\vv\xi_k$ 线性无关.
			\item\label{ex:sol:B7} 一张 $3\times3$ 的统计表只公开了三个行和与三个列和. 至少还要再公开几个单元格，才能确定全表？给出一种``公开哪几个格''的选法，并说明为什么某些同样个数的选法却不行.
			\item\label{ex:sol:B8} 对 $\vv A=\begin{pmatrix}1&2&1\\2&4&3\\3&6&4\end{pmatrix}$，求 $\vv b=(b_1,b_2,b_3)^\trans$ 应满足的条件，使 $\vv{Ax}=\vv b$ 有解.（提示：找出把各方程组合成 $0=\cdots$ 的系数.）
			\item\label{ex:sol:B9} 设 $\vv A$ 为 $n$ 阶方阵，$\rk\vv A=n-1$，且 $\vv A$ 的每行元素之和为零. 求 $\vv{Ax}=\zero$ 的通解.
			\item\label{ex:sol:B10} 设 $\vv A,\vv B$ 为 $n$ 阶方阵，$\vv{AB}=\vv O$，$\vv A+\vv B$ 可逆. 证明 $\rk\vv A+\rk\vv B=n$.
			\item\label{ex:sol:B11} 证明：两个 $m\times n$ 齐次方程组 $\vv{Ax}=\zero$ 与 $\vv{Bx}=\zero$ 同解，当且仅当 $\vv A$ 与 $\vv B$ 的行向量组等价.
			\item\label{ex:sol:B12} 设 $\vv A$ 为 $m\times n$ 矩阵. 证明：$\vv{Ax}=\vv b$ 对一切 $\vv b$ 有解，当且仅当存在 $n\times m$ 矩阵 $\vv C$ 使 $\vv{AC}=\vv I_m$；这时 $\vv C\vv b$ 就是一个解.
			\item\label{ex:sol:B13} 设 $\vv A$ 为 $n$ 阶方阵（$n\ge2$），$\rk\vv A=n-1$，$\vv\xi$ 是 $\vv{Ax}=\zero$ 的非零解. 证明 $\vv A^*$ 的每一列都是 $\vv\xi$ 的倍数.
			\item\label{ex:sol:B14} 用可确定性判据回答：方程组 $x_1+x_2+x_3=6$，$x_2+x_3+x_4=9$，$x_3+x_4+x_5=12$（例如三段相邻时间窗口内的观测总量）能确定 $x_1-x_4$ 吗？$x_1+x_2+x_3+x_4+x_5$ 呢？$x_1-x_4+x_5$ 呢？
			\item\label{ex:sol:B15} 设 $\vv A$ 为 $m\times n$ 矩阵，$\vv B$ 为 $n\times p$ 矩阵. 证明 $\dim\ker(\vv{AB})\le\dim\ker\vv A+\dim\ker\vv B$，并由此重新推出 Sylvester 不等式.
		\end{exB}
		
		\Cgroup
		\begin{exC}
			\item\label{ex:sol:C1} （Farkas 引理的二维形式）设 $\vv a_1,\vv a_2\in\mathbb{R}^2$ 线性无关，$\vv b\in\mathbb{R}^2$. 证明：下面两件事恰有一件成立：(1) 存在 $x_1,x_2\ge0$ 使 $x_1\vv a_1+x_2\vv a_2=\vv b$；(2) 存在 $\vv y$ 使 $\vv y^\trans\vv a_1\ge0$，$\vv y^\trans\vv a_2\ge0$，$\vv y^\trans\vv b<0$. 画图解释. 这是``无套利定价''与线性规划对偶理论的出发点.
			\item\label{ex:sol:C2} 设 $\vv A$ 是 $n$ 阶整数矩阵. 证明：若对任意整数向量 $\vv b$，$\vv{Ax}=\vv b$ 都有整数解，则 $\det\vv A=\pm1$.
			\item\label{ex:sol:C3} 证明：方阵 $\vv A$ 不用行交换就能做 LU 分解（$\vv L$ 单位下三角，$\vv U$ 上三角且可逆），当且仅当 $\vv A$ 的各阶顺序主子式都不为零；这时分解唯一.
			\item\label{ex:sol:C4} 设 $\vv A$ 是三对角矩阵（只有主对角线和相邻两条斜线上可能有非零元）. 证明：若它有 LU 分解，则 $\vv L,\vv U$ 都只有两条非零斜线；因此解三对角方程组只需与 $n$ 成正比的运算量. 这种方法称为``追赶法''.
			\item\label{ex:sol:C5} 用分块初等变换
			\[\begin{pmatrix}\vv I_n&\vv O\\ -\vv A&\vv I_m\end{pmatrix}\begin{pmatrix}\vv I_n&\vv B\\ \vv A&\vv O\end{pmatrix}\begin{pmatrix}\vv I_n&-\vv B\\ \vv O&\vv I_p\end{pmatrix}=\begin{pmatrix}\vv I_n&\vv O\\ \vv O&-\vv{AB}\end{pmatrix}\]
			给出 Sylvester 不等式 $\rk(\vv{AB})\ge\rk\vv A+\rk\vv B-n$ 的另一个证明.（提示：用第三章章末习题~\ref{ex:mat:B5} 的分块秩不等式.）
			\item\label{ex:sol:C6} 证明：$\mathbb{F}^n$ 中任意 $k$ 个线性流形（不必过原点）的交，要么是空集，要么仍是线性流形. 给出两个二维线性流形在 $\mathbb{R}^3$ 中交为空集的例子，并用秩的语言刻画这种情形.
			\item\label{ex:sol:C7} （列昂惕夫价格模型）设 $\vv A$ 为直接消耗系数矩阵，$\vv v$ 为各部门单位产品的增加值（工资、利润等）. 价格 $\vv p$ 满足``单位价格 $=$ 单位投入成本 $+$ 单位增加值''：$\vv p=\vv A^\trans\vv p+\vv v$. 说明它是第一章产量模型 $\vv x=\vv{Ax}+\vv d$ 的``转置''；证明在两个模型都可解时有 $\vv p^\trans\vv d=\vv v^\trans\vv x$，并解释其经济含义（国民收入的两种核算方法：支出法与收入法）.
			\item\label{ex:sol:C8} 设 $\vv A$ 为 $n$ 阶方阵. 证明存在正整数 $k$ 使 $\rk\vv A^k=\rk\vv A^{k+1}=\rk\vv A^{k+2}=\cdots$，且这样的最小的 $k$ 不超过 $n$.
		\end{exC}
	}
	
	% =====================================================================
	% 第七章　多项式
	% =====================================================================
	\newcommand{\ExPolynomial}{%
		\exsetheading
		\Agroup
		\begin{exA}
			\item\label{ex:poly:A1} 求商与余式：(1) $x^5-1$ 除以 $x^2+x+1$；(2) $2x^4-3x^3+4x^2-5x+6$ 除以 $x-2$（用综合除法）.
			\item\label{ex:poly:A2} 求 $x^{100}$ 分别除以 $x^2-1$，$x^2+1$，$(x-1)^2$ 所得的余式.
			\item\label{ex:poly:A3} 求 $a,b$，使 $x^2+x+1$ 整除 $x^4+ax^2+b$.（提示：设 $\omega$ 是 $x^2+x+1$ 的根，则 $\omega^3=1$.）
			\item\label{ex:poly:A4} 用辗转相除法求最大公因式：(1) $x^4+3x^3-x^2-4x-3$ 与 $3x^3+10x^2+2x-3$；(2) $x^3-2x^2+x$ 与 $x^2-1$.
			\item\label{ex:poly:A5} 证明：对任意多项式 $h$，$\gcd(f,g)=\gcd(f,\ g-hf)$. 这正是辗转相除法每一步的依据.
			\item\label{ex:poly:A6} 求 $u,v$ 使 $u(x)(x^2+1)+v(x)(x+1)=1$.
			\item\label{ex:poly:A7} 求 $2x^3-5x^2+x+2$ 的全部有理根，并在 $\mathbb{Q}$ 上把它分解为一次因式之积.
			\item\label{ex:poly:A8} 分别在 $\mathbb{Q}$，$\mathbb{R}$，$\mathbb{C}$ 上分解：(1) $x^4+x^2+1$；(2) $x^6-1$；(3) $x^3-2$.
			\item\label{ex:poly:A9} 判断 $1$ 是 $x^5-5x+4$ 的几重根.
			\item\label{ex:poly:A10} $a$ 取何值时，$x^3-3x+a$ 有重根？求出重根.
			\item\label{ex:poly:A11} 设 $\alpha_1,\alpha_2,\alpha_3$ 是 $x^3-2x^2+3x-4$ 的根. 不解方程，求 $\alpha_1+\alpha_2+\alpha_3$，$\alpha_1^2+\alpha_2^2+\alpha_3^2$，$\frac1{\alpha_1}+\frac1{\alpha_2}+\frac1{\alpha_3}$.
			\item\label{ex:poly:A12} 求次数最低的首一有理系数多项式，使 $1+\sqrt2$ 与 $\mathrm i$ 都是它的根.
			\item\label{ex:poly:A13} 求一个三次实系数首一多项式，以 $2$ 和 $1+\mathrm i$ 为根.
			\item\label{ex:poly:A14} 设 $f(x)=x^2-2x+1$，$\vv A=\begin{pmatrix}1&2\\0&1\end{pmatrix}$，计算 $f(\vv A)$；再设 $g(x)=x^3-x$，计算 $g(\vv A)$.
			\item\label{ex:poly:A15} 求过点 $(0,1)$，$(1,0)$，$(2,3)$，$(3,10)$ 的三次以下的插值多项式（用 Lagrange 公式）.
			\item\label{ex:poly:A16} 证明 $(x-1)\mid x^n-1$，并写出商式. 一般地，对任意 $a$，$(x-a)\mid x^n-a^n$.
			\item\label{ex:poly:A17} 把 $f(x)=x^4-2x^3+3$ 按 $x-1$ 的幂展开，即写成 $\sum c_k(x-1)^k$.（提示：反复用 $x-1$ 作综合除法，或用 $c_k=f^{(k)}(1)/k!$.）
			\item\label{ex:poly:A18} 多项式 $f$ 除以 $x-1$ 余 $3$，除以 $x-2$ 余 $5$. 求 $f$ 除以 $(x-1)(x-2)$ 的余式.
			\item\label{ex:poly:A19} 一个项目现在投入 $100$ 万元，一年后回收 $55$ 万元，两年后回收 $60.5$ 万元. 写出内部收益率满足的方程并精确求解.
			\item\label{ex:poly:A20} 证明：若 $f\mid g$ 且 $g\mid f$，则 $f=cg$（$c$ 为非零常数）.
			\item\label{ex:poly:A21} 设 $\gcd(f,g)=d$，$f=df_1$，$g=dg_1$. 证明 $\gcd(f_1,g_1)=1$.
		\end{exA}
		
		\Bgroup
		\begin{exB}
			\item\label{ex:poly:B1} 证明：若 $\gcd(f,g)=1$ 且 $\gcd(f,h)=1$，则 $\gcd(f,gh)=1$.
			\item\label{ex:poly:B2} 证明：$\gcd(x^m-1,\ x^n-1)=x^{\gcd(m,n)}-1$.（提示：模仿整数的辗转相除，先证 $x^m-1=x^{m-n}(x^n-1)+(x^{m-n}-1)$.）
			\item\label{ex:poly:B3} 设 $f=(x-1)^3(x+2)^2$. 计算 $f/\gcd(f,f')$，说明它恰好是``去掉重复因式''的结果. 一般地证明：若 $f=c\,p_1^{k_1}\cdots p_m^{k_m}$ 是标准分解，则 $f/\gcd(f,f')=c\,p_1\cdots p_m$.
			\item\label{ex:poly:B4} 证明：若 $p$ 不可约且 $p\mid f^2$，则 $p\mid f$. 举例说明``不可约''不能去掉.
			\item\label{ex:poly:B5} 设 $\ell_1,\ldots,\ell_n$ 是关于互异节点 $x_1,\ldots,x_n$ 的 Lagrange 基函数. 证明 $\sum_i\ell_i(x)=1$，$\sum_ix_i\ell_i(x)=x$，一般地 $\sum_ix_i^k\ell_i(x)=x^k$（$0\le k\le n-1$）.
			\item\label{ex:poly:B6} 证明：若复系数多项式 $f$ 满足 $f(x+1)=f(x)$ 对一切 $x$ 成立，则 $f$ 是常数.
			\item\label{ex:poly:B7} 证明 $x^3+x+1$ 在 $\mathbb{Q}$ 上不可约. 为什么对三次多项式只需检验有无有理根？四次多项式可以这样检验吗？
			\item\label{ex:poly:B8} 求 $x^n$ 除以 $x^2-3x+2$ 的余式.
			\item\label{ex:poly:B9} （Hermite 插值）求次数最低的多项式 $f$，使 $f(1)=1$，$f'(1)=2$，$f(2)=4$.（提示：设 $f=1+2(x-1)+c(x-1)^2$.）
			\item\label{ex:poly:B10} 设 $f$ 是有理系数多项式，$a,b\in\mathbb{Q}$，$\sqrt2\notin\mathbb{Q}$. 证明：若 $a+b\sqrt2$ 是 $f$ 的根，则 $a-b\sqrt2$ 也是 $f$ 的根.（提示：用 $(x-a)^2-2b^2$ 除 $f$.）
			\item\label{ex:poly:B11} 设 $f$ 是整系数多项式，$f(0)$ 与 $f(1)$ 都是奇数. 证明 $f$ 没有整数根.
			\item\label{ex:poly:B12} 求最小多项式意义下的``零化''：分别求出次数最低的首一多项式 $g$，使 $g(\vv A)=\vv O$，其中 $\vv A=\begin{pmatrix}0&-1\\1&0\end{pmatrix}$，$\vv A=\begin{pmatrix}2&0\\0&2\end{pmatrix}$，$\vv A=\begin{pmatrix}2&1\\0&2\end{pmatrix}$.
			\item\label{ex:poly:B13} 证明：现金流 $-C_0,\ C_1,\ldots,C_n$（$C_0>0$，$C_i\ge0$ 且不全为零）的净现值 $N(v)=-C_0+C_1v+\cdots+C_nv^n$ 在 $v>0$ 时严格递增，从而 $N(v)=0$ 至多一个正根；当 $C_1+\cdots+C_n>C_0$ 时这个正根落在 $(0,1)$ 中，即 IRR 唯一且为正.
			\item\label{ex:poly:B14} 设 $f$ 是 $n$ 次多项式，在 $n+1$ 个整数点 $0,1,\ldots,n$ 上都取整数值. 证明 $f$ 在一切整数点上都取整数值.（提示：用基 $\binom x0,\binom x1,\ldots,\binom xn$，其中 $\binom xk=\frac{x(x-1)\cdots(x-k+1)}{k!}$.）
			\item\label{ex:poly:B15} 证明 Lagrange 插值的``唯一性''也可以由线性代数推出：求值映射 $\mathbb{F}[x]_n\to\mathbb{F}^n$，$f\mapsto(f(x_1),\ldots,f(x_n))^\trans$ 是单射，从而是同构.
			\item\label{ex:poly:B16} 设 $p$ 是 $\mathbb{F}$ 上的不可约多项式，$f\in\mathbb{F}[x]$. 证明：$\gcd(p,f)=1$ 或 $p\mid f$，二者必居其一.
		\end{exB}
		
		\Cgroup
		\begin{exC}
			\item\label{ex:poly:C1} （Eisenstein 判别法）设 $f=a_nx^n+\cdots+a_0$ 是整系数多项式，$p$ 是素数，$p\nmid a_n$，$p\mid a_i$（$i<n$），$p^2\nmid a_0$. 证明 $f$ 在 $\mathbb{Q}$ 上不可约（可以先承认 Gauss 引理：整系数多项式若能分解为两个低次有理系数多项式之积，则也能分解为两个低次整系数多项式之积）. 由此说明 $x^n-2$ 对一切 $n$ 都不可约，所以 $\mathbb{Q}[x]$ 中有任意高次的不可约多项式.
			\item\label{ex:poly:C2} 设 $p$ 是素数. 证明分圆多项式 $\Phi_p(x)=x^{p-1}+x^{p-2}+\cdots+1$ 在 $\mathbb{Q}$ 上不可约.（提示：对 $\Phi_p(x+1)$ 用 Eisenstein 判别法.）
			\item\label{ex:poly:C3} （Newton 恒等式）记 $s_k=\alpha_1^k+\alpha_2^k+\alpha_3^k$，$\sigma_1,\sigma_2,\sigma_3$ 为初等对称多项式. 证明 $s_1=\sigma_1$，$s_2=\sigma_1s_1-2\sigma_2$，$s_3=\sigma_1s_2-\sigma_2s_1+3\sigma_3$，并求 $x^3-2x^2+3x-4$ 的根的立方和.
			\item\label{ex:poly:C4} （三次方程的判别式）设 $x^3+px+q$ 的三个复根为 $\alpha_1,\alpha_2,\alpha_3$，$\Delta=\prod_{i<j}(\alpha_i-\alpha_j)^2$. 证明 $\Delta=-4p^3-27q^2$，并说明（$p,q$ 为实数时）：$\Delta>0$ 时三个不同实根，$\Delta<0$ 时一实二虚，$\Delta=0$ 时有重根.
			\item\label{ex:poly:C5} （Chebyshev 多项式）证明存在 $n$ 次多项式 $T_n$ 使 $\cos n\theta=T_n(\cos\theta)$，且 $T_{n+1}=2xT_n-T_{n-1}$. 求 $T_n$ 的全部根. 由此证明：$\cos20^\circ$ 是 $8x^3-6x-1$ 的根，并验证这个三次多项式在 $\mathbb{Q}$ 上不可约. （结合``尺规作图只能得到经若干次开平方得到的数''，这说明 $60^\circ$ 角不能用尺规三等分.）
			\item\label{ex:poly:C6} 证明：在 $\mathbb{R}$ 上处处非负的实系数多项式，都能写成两个实系数多项式的平方和.（提示：利用实根的重数必为偶数，以及 $|a+b\mathrm i|^2=a^2+b^2$ 与 $(a^2+b^2)(c^2+d^2)=(ac-bd)^2+(ad+bc)^2$.）
			\item\label{ex:poly:C7} 设 $f\in\mathbb{Q}[x]$ 不可约，$g\in\mathbb{Q}[x]$ 与 $f$ 有一个公共复根. 证明 $f\mid g$. 由此说明：一个复数若是某个有理系数多项式的根，则以它为根的次数最低的首一有理系数多项式（称为它的\textbf{极小多项式}）唯一，且不可约.
			\item\label{ex:poly:C8} （$\mathbb{F}_2$ 上的多项式）在系数按模 $2$ 运算的多项式中，证明 $x^2+x+1$ 是唯一的二次不可约多项式，并找出全部三次不可约多项式. 这类多项式是循环冗余校验（CRC）与密码学的基础.
			\item\label{ex:poly:C9} 设 $f$ 是实系数多项式，它的系数序列的``变号次数''为 $V$（忽略零系数）. Descartes 符号法则说：$f$ 的正根个数（按重数计）不超过 $V$，且与 $V$ 的差是偶数. 用它重新解释第七章``内部收益率''建模实验中两个项目的结论.
		\end{exC}
	}
	
	% =====================================================================
	% 第八章　特征值与特征向量
	% =====================================================================
	\newcommand{\ExEigen}{%
		\exsetheading
		\Agroup
		\begin{exA}
			\item\label{ex:eig:A1} 求下列矩阵的特征值与特征向量（在 $\mathbb{C}$ 上考虑）：
			\[\begin{pmatrix}3&1\\0&2\end{pmatrix},\quad\begin{pmatrix}1&2\\2&1\end{pmatrix},\quad\begin{pmatrix}0&1\\-1&0\end{pmatrix},\quad\begin{pmatrix}4&-2\\1&1\end{pmatrix},\]\[\begin{pmatrix}2&0&0\\0&1&1\\0&1&1\end{pmatrix},\quad\begin{pmatrix}1&2&2\\2&1&2\\2&2&1\end{pmatrix}.\]
			\item\label{ex:eig:A2} 判断 A1 中各矩阵能否在 $\mathbb{R}$ 上对角化；能的，求可逆矩阵 $\vv P$ 使 $\vv P^{-1}\vv{AP}$ 为对角矩阵.
			\item\label{ex:eig:A3} 判断下列矩阵能否对角化：$\begin{pmatrix}1&1\\0&2\end{pmatrix}$，$\begin{pmatrix}2&1\\0&2\end{pmatrix}$，$\begin{pmatrix}2&0&0\\0&2&0\\1&0&2\end{pmatrix}$，$\begin{pmatrix}1&1&0\\0&2&0\\0&0&1\end{pmatrix}$.
			\item\label{ex:eig:A4} 设 $3$ 阶方阵 $\vv A$ 的特征值为 $1,2,3$. 求 $\det\vv A$，$\tr\vv A$，$\vv A^{-1}$ 的特征值，$\vv A^2-2\vv A+\vv I$ 的特征值，$\det(\vv A^*)$，$\det(\vv A^2+\vv I)$.
			\item\label{ex:eig:A5} 已知二阶方阵的迹为 $5$，行列式为 $6$，求它的特征值. 若迹为 $2$、行列式为 $5$ 呢？
			\item\label{ex:eig:A6} 证明：$\vv A$ 与 $\vv A^\trans$ 有相同的特征多项式. 它们的特征向量是否相同？举例说明.
			\item\label{ex:eig:A7} 证明：三角矩阵的特征值恰为其对角元.
			\item\label{ex:eig:A8} 证明：若 $\vv A^k=\vv O$，则 $\vv A$ 的特征值全为 $0$；若 $\vv A^2=\vv A$，则特征值只能是 $0$ 或 $1$；若 $\vv A^2=\vv I$，则特征值只能是 $\pm1$.
			\item\label{ex:eig:A9} 设 $\vv A=\begin{pmatrix}4&-2\\1&1\end{pmatrix}$. 利用对角化计算 $\vv A^n$.
			\item\label{ex:eig:A10} 设品牌转移矩阵为 $\vv P=\begin{pmatrix}0.9&0.2\\0.1&0.8\end{pmatrix}$. 求稳态分布，并利用对角化求 $\lim_{n\to\infty}\vv P^n$. 解释极限矩阵两列相同的含义.
			\item\label{ex:eig:A11} 已知 $\vv\xi=(1,1,-1)^\trans$ 是 $\vv A=\begin{pmatrix}2&-1&2\\5&a&3\\-1&b&-2\end{pmatrix}$ 的特征向量. 求 $a,b$ 与对应的特征值.
			\item\label{ex:eig:A12} 设 $\vv A=\begin{pmatrix}0&0&1\\1&1&x\\1&0&0\end{pmatrix}$. $x$ 取何值时 $\vv A$ 可对角化？
			\item\label{ex:eig:A13} 对二阶方阵直接验证 Cayley–Hamilton 定理：$\vv A^2-(\tr\vv A)\vv A+(\det\vv A)\vv I=\vv O$.
			\item\label{ex:eig:A14} 求下列矩阵的最小多项式：$\diag(1,1,2)$，$\begin{pmatrix}1&1&0\\0&1&0\\0&0&2\end{pmatrix}$，$\begin{pmatrix}0&1&0\\0&0&1\\0&0&0\end{pmatrix}$，$\begin{pmatrix}1&1&1\\1&1&1\\1&1&1\end{pmatrix}$.
			\item\label{ex:eig:A15} 判断 $\begin{pmatrix}1&1\\0&1\end{pmatrix}$ 与 $\begin{pmatrix}1&0\\1&1\end{pmatrix}$ 是否相似；若相似，求出 $\vv P$.
			\item\label{ex:eig:A16} 设 $\vv A$ 是上三角矩阵. 证明 $\spn(\vv e_1)$，$\spn(\vv e_1,\vv e_2)$，……都是 $\vv A$ 的不变子空间.
			\item\label{ex:eig:A17} 解差分方程 $a_{k+1}=a_k+2a_{k-1}$，$a_0=0$，$a_1=1$.
			\item\label{ex:eig:A18} （种群模型）某种群分为幼年与成年两组，第 $k$ 年的数量记为 $\vv x_k=(y_k,z_k)^\trans$. 每只成年个体每年平均产下 $1.1$ 只幼体，幼体有 $60\%$ 存活并长为成年，成年个体有一半活到下一年，即
			$y_{k+1}=1.1z_k$，$z_{k+1}=0.6y_k+0.5z_k$. 求矩阵的特征值与特征向量，说明种群的长期增长率和幼体与成年个体的稳定比例.
			\item\label{ex:eig:A19} 证明：实二阶方阵 $\vv A$ 没有实特征值，当且仅当 $(\tr\vv A)^2<4\det\vv A$.
			\item\label{ex:eig:A20} 求平面上向直线 $y=2x$ 的正投影、关于该直线的反射这两个变换的特征值与特征向量（不必写出矩阵，从几何上直接回答），再写出矩阵加以验证.
		\end{exA}
		
		\Bgroup
		\begin{exB}
			\item\label{ex:eig:B1} 设 $\vv u,\vv v\in\mathbb{R}^n$ 非零，$\vv A=\vv u\vv v^\trans$. 求 $\vv A$ 的全部特征值及其重数，并证明：$\vv A$ 可对角化当且仅当 $\vv v^\trans\vv u\ne0$.
			\item\label{ex:eig:B2} 设 $\vv A$ 有 $n$ 个互异的特征值，$\vv{AB}=\vv{BA}$. 证明 $\vv A$ 的特征向量都是 $\vv B$ 的特征向量，从而 $\vv B$ 也可对角化；进一步证明 $\vv B$ 是 $\vv A$ 的多项式.
			\item\label{ex:eig:B3} 证明：若 $\vv A$ 可对角化，则 $\vv A^k$ 与 $f(\vv A)$ 都可对角化. 逆命题成立吗？
			\item\label{ex:eig:B4} 设 $\vv A$ 可对角化，$W$ 是 $\vv A$ 的不变子空间. 证明 $\vv A|_W$ 也可对角化.（提示：最小多项式.）
			\item\label{ex:eig:B5} 证明：满足 $\vv A^2=-\vv I$ 的二阶实方阵都与 $\begin{pmatrix}0&-1\\1&0\end{pmatrix}$ 相似.
			\item\label{ex:eig:B6} （Gershgorin 圆盘定理）证明：方阵 $\vv A$ 的每个特征值都落在某个圆盘 $\{z\in\mathbb{C}:|z-a_{ii}|\le\sum_{j\ne i}|a_{ij}|\}$ 中. 由此证明：严格对角占优的矩阵（$|a_{ii}|>\sum_{j\ne i}|a_{ij}|$）可逆.
			\item\label{ex:eig:B7} 设 $\vv A$ 的每行元素之和都是 $s$. 证明 $s$ 是 $\vv A$ 的特征值. 若 $\vv A$ 可逆，求 $\vv A^{-1}$ 每行元素之和.
			\item\label{ex:eig:B8} 设 $\vv A=\begin{pmatrix}a&b&c\\c&a&b\\b&c&a\end{pmatrix}$（循环矩阵），$\omega=\mathrm e^{2\pi\mathrm i/3}$. 证明 $(1,\omega^k,\omega^{2k})^\trans$（$k=0,1,2$）都是 $\vv A$ 的特征向量，并求对应的特征值. 由此重新得到第五章章末习题~\ref{ex:det:B12} 的行列式公式.
			\item\label{ex:eig:B9} 解差分方程 $a_{k+1}=4a_k-4a_{k-1}$，$a_0=1$，$a_1=4$. 解释通项中为什么出现 $k\cdot2^k$.
			\item\label{ex:eig:B10} 证明 $\tr(\vv A^k)=\lambda_1^k+\cdots+\lambda_n^k$. 若三阶方阵满足 $\tr\vv A=0$，$\tr\vv A^2=6$，$\det\vv A=0$，求它的特征值.
			\item\label{ex:eig:B11} 证明：若 $\vv A^3=\vv A$，则 $\vv A$ 可对角化，且特征值在 $\{0,1,-1\}$ 中.
			\item\label{ex:eig:B12} $a,b,c$ 满足什么条件时，$\begin{pmatrix}2&0&0\\a&2&0\\b&c&1\end{pmatrix}$ 可对角化？
			\item\label{ex:eig:B13} 三个网页之间的链接为：网页 $1$ 链接到 $2,3$，网页 $2$ 链接到 $3$，网页 $3$ 链接到 $1$.（与判例不同，网页之间可以``互相引用''，形成环路.） 取 $d=1$ 写出转移矩阵，手算它的稳态分布.
			\item\label{ex:eig:B14} （矩阵指数）对可对角化矩阵 $\vv A=\vv P\vv D\vv P^{-1}$，定义 $\mathrm e^{t\vv A}=\vv P\diag(\mathrm e^{t\lambda_1},\ldots,\mathrm e^{t\lambda_n})\vv P^{-1}$. 对 $\vv A=\begin{pmatrix}0&1\\1&0\end{pmatrix}$ 计算 $\mathrm e^{t\vv A}$，并验证 $\vv x(t)=\mathrm e^{t\vv A}\vv x_0$ 满足微分方程组 $x_1'=x_2$，$x_2'=x_1$.
			\item\label{ex:eig:B15} 三状态马尔可夫链的转移矩阵为 $\begin{pmatrix}0.5&0.25&0.25\\0.25&0.5&0.25\\0.25&0.25&0.5\end{pmatrix}$. 求稳态分布与全部特征值，说明偏离稳态的部分每一步缩小为原来的多少.
			\item\label{ex:eig:B16} 证明：若 $\lambda$ 是 $\vv A$ 的 $k$ 重特征值（代数重数），则 $\rk(\lambda\vv I-\vv A)\ge n-k$.
			\item\label{ex:eig:B17} 对第三章``判例引证网络''中的四份判决，分别取阻尼系数 $d=0.5,\ 0.85,\ 0.95,\ 1$，按第八章建模实验的方法求 PageRank 得分（可用计算机），列表比较各判决的得分与排序. $d$ 越接近 $1$，得分越``集中''在哪些判决上？为什么？
		\end{exB}
		
		\Cgroup
		\begin{exC}
			\item\label{ex:eig:C1} 设 $\vv A$ 为 $m\times n$ 矩阵，$\vv B$ 为 $n\times m$ 矩阵. 证明 $\lambda^n\det(\lambda\vv I_m-\vv{AB})=\lambda^m\det(\lambda\vv I_n-\vv{BA})$. 由此说明 $\vv{AB}$ 与 $\vv{BA}$ 的非零特征值（连同重数）相同.（提示：计算 $\begin{pmatrix}\vv I_m&-\vv A\\ \vv O&\lambda\vv I_n\end{pmatrix}\begin{pmatrix}\lambda\vv I_m&\vv A\\ \vv B&\vv I_n\end{pmatrix}$ 等两种乘积的行列式.）
			\item\label{ex:eig:C2} （Schur 三角化）证明：每个复方阵都相似于一个上三角矩阵. 由此给出 Cayley–Hamilton 定理的另一个证明：对上三角矩阵 $\vv T$，证明 $(\vv T-t_{11}\vv I)\cdots(\vv T-t_{nn}\vv I)=\vv O$.
			\item\label{ex:eig:C3} 求下列矩阵的 Jordan 标准形：
			\[\begin{pmatrix}2&1&0\\0&2&0\\0&0&2\end{pmatrix},\quad\begin{pmatrix}1&1&1\\0&1&1\\0&0&1\end{pmatrix},\quad\begin{pmatrix}0&1&0&0\\0&0&0&0\\0&0&0&1\\0&0&0&0\end{pmatrix}.\]
			特征多项式为 $(\lambda-2)^4$ 的 $4$ 阶复矩阵，在相似的意义下共有几类？
			\item\label{ex:eig:C4} 证明：若一组可对角化的矩阵两两可交换，则它们可以\textbf{同时}对角化（存在同一个 $\vv P$ 使它们都化为对角矩阵）.
			\item\label{ex:eig:C5} 设 $\vv A$ 可对角化. 证明 $\vv A^k\to\vv O$（按元素）当且仅当 $\vv A$ 的所有特征值的模都小于 $1$. 用它说明：投入产出模型中，若直接消耗系数矩阵 $\vv A$ 可对角化且特征值的模都小于 $1$，则 $(\vv I-\vv A)^{-1}=\vv I+\vv A+\vv A^2+\cdots$.
			\item\label{ex:eig:C6} （友矩阵）证明 $\begin{pmatrix}0&0&\cdots&0&-a_0\\1&0&\cdots&0&-a_1\\0&1&\cdots&0&-a_2\\ \vdots&&\ddots&&\vdots\\0&0&\cdots&1&-a_{n-1}\end{pmatrix}$ 的特征多项式为 $\lambda^n+a_{n-1}\lambda^{n-1}+\cdots+a_0$. 由此说明：求多项式的根等价于求矩阵的特征值（数值软件正是这样求多项式的根的）.
			\item\label{ex:eig:C7} （赌徒破产）一位赌徒有 $k$ 元（$0\le k\le N$），每局以概率 $\frac12$ 赢或输 $1$ 元，输光或达到 $N$ 元时停止. 设 $p_k$ 为最终达到 $N$ 元的概率，证明 $p_k=\frac12(p_{k+1}+p_{k-1})$，并用特征根方法解出 $p_k=k/N$. 赢的概率改为 $p\ne\frac12$ 时呢？
			\item\label{ex:eig:C8} 证明：$n$ 阶三对角矩阵 $\vv T=\begin{pmatrix}2&-1&&\\-1&2&\ddots&\\&\ddots&\ddots&-1\\&&-1&2\end{pmatrix}$ 的特征值为 $2-2\cos\frac{k\pi}{n+1}$（$k=1,\ldots,n$），对应特征向量的第 $j$ 个分量为 $\sin\frac{jk\pi}{n+1}$. 对 $n=2,3$ 直接验证.
			\item\label{ex:eig:C9} （Perron–Frobenius 的二阶情形）设 $\vv A$ 是元素全为正的二阶矩阵. 证明 $\vv A$ 有两个互异实特征值，较大者 $r$ 为正且大于另一个的绝对值，属于 $r$ 的特征向量可取为各分量都为正. 由此解释：为什么 PageRank 引入阻尼系数后，``全部元素为正''能保证排序唯一.
			\item\label{ex:eig:C10} 证明：满足 $\vv A^2=-\vv I$ 的 $n$ 阶实方阵存在当且仅当 $n$ 为偶数，且这时 $\vv A$ 相似于 $\diag(\vv J,\ldots,\vv J)$，$\vv J=\begin{pmatrix}0&-1\\1&0\end{pmatrix}$.
		\end{exC}
	}
	
	% =====================================================================
	% 第九章　内积空间
	% =====================================================================
	\newcommand{\ExInnerProduct}{%
		\exsetheading
		\Agroup
		\begin{exA}
			\item\label{ex:ip:A1} 在 $\mathbb{R}^4$ 中，设 $\vv\alpha=(1,2,2,0)^\trans$，$\vv\beta=(2,0,1,2)^\trans$. 求 $\|\vv\alpha\|$，$\|\vv\beta\|$，$\langle\vv\alpha,\vv\beta\rangle$，它们的夹角，以及 $\|\vv\alpha-\vv\beta\|$.
			\item\label{ex:ip:A2} 判断下列二元函数是否为 $\mathbb{R}^2$ 上的内积：
			(1) $\langle\vv x,\vv y\rangle=x_1y_1+2x_2y_2$；(2) $\langle\vv x,\vv y\rangle=x_1y_2+x_2y_1$；(3) $\langle\vv x,\vv y\rangle=2x_1y_1+x_1y_2+x_2y_1+x_2y_2$；(4) $\langle\vv x,\vv y\rangle=x_1y_1-x_2y_2$.
			\item\label{ex:ip:A3} 用 Gram–Schmidt 方法把下列向量组正交化并单位化：
			(1) $(1,1,1)^\trans$，$(1,2,0)^\trans$，$(0,1,2)^\trans$；(2) $(1,1,0,0)^\trans$，$(1,0,1,0)^\trans$，$(-1,0,0,1)^\trans$.
			\item\label{ex:ip:A4} 求与 $(1,1,1)^\trans$ 和 $(1,0,-1)^\trans$ 都正交的单位向量.
			\item\label{ex:ip:A5} 已知 $\vv e_1=\frac13(1,2,2)^\trans$，$\vv e_2=\frac13(2,1,-2)^\trans$，$\vv e_3=\frac13(2,-2,1)^\trans$. 验证它们构成 $\mathbb{R}^3$ 的标准正交基，并只用内积求 $(3,6,0)^\trans$ 在这组基下的坐标.
			\item\label{ex:ip:A6} 求 $\vv\alpha=(1,2,3,4)^\trans$ 在 $W=\spn\big((1,1,0,0)^\trans,(0,0,1,1)^\trans\big)$ 上的正交投影，以及 $\vv\alpha$ 到 $W$ 的距离.
			\item\label{ex:ip:A7} 用最小二乘法求过数据点 $(0,1)$，$(1,2)$，$(2,2)$，$(3,4)$，$(4,6)$ 的最佳直线 $y=a+bx$，并计算残差平方和.
			\item\label{ex:ip:A8} 求 $a,b,c$ 使 $\begin{pmatrix}\frac1{\sqrt2}&a&\frac1{\sqrt6}\\0&b&-\frac2{\sqrt6}\\-\frac1{\sqrt2}&c&\frac1{\sqrt6}\end{pmatrix}$ 为正交矩阵.
			\item\label{ex:ip:A9} 证明：若 $\vv\alpha_1,\ldots,\vv\alpha_k$ 两两正交，则 $\|\vv\alpha_1+\cdots+\vv\alpha_k\|^2=\|\vv\alpha_1\|^2+\cdots+\|\vv\alpha_k\|^2$.
			\item\label{ex:ip:A10} 证明：$\big|\|\vv\alpha\|-\|\vv\beta\|\big|\le\|\vv\alpha-\vv\beta\|$.
			\item\label{ex:ip:A11} 用 Cauchy–Schwarz 不等式证明：(1) $(a_1+\cdots+a_n)^2\le n(a_1^2+\cdots+a_n^2)$；(2) 若 $a_i>0$，则 $(a_1+\cdots+a_n)\big(\frac1{a_1}+\cdots+\frac1{a_n}\big)\ge n^2$.
			\item\label{ex:ip:A12} 在 $C[0,1]$ 中取 $\langle f,g\rangle=\int_0^1f(x)g(x)\,\mathrm dx$. 求 $\|1\|$，$\|x\|$，$1$ 与 $x$ 的夹角，并把 $1,x$ 正交化.
			\item\label{ex:ip:A13} 设 $\vv A=\begin{pmatrix}1&1&0&1\\0&1&1&1\end{pmatrix}$. 求 $\ker\vv A$ 与行空间 $\imm\vv A^\trans$ 的标准正交基，并验证它们互为正交补.
			\item\label{ex:ip:A14} 求 $\vv A=\begin{pmatrix}1&1\\1&0\\0&1\end{pmatrix}$ 的 QR 分解.
			\item\label{ex:ip:A15} 设 $\vv u=\frac13(1,2,2)^\trans$. 写出 Householder 矩阵 $\vv H=\vv I-2\vv u\vv u^\trans$，验证它是对称的正交矩阵，并求 $\vv H\vv u$ 与 $\vv H(2,-1,0)^\trans$.
			\item\label{ex:ip:A16} 判断下列正交矩阵是旋转还是反射；若是旋转求旋转角，若是反射求反射轴：$\frac15\begin{pmatrix}3&-4\\4&3\end{pmatrix}$，$\frac15\begin{pmatrix}3&4\\4&-3\end{pmatrix}$.
			\item\label{ex:ip:A17} 对``过原点的直线''模型 $y=bx$，推导最小二乘解 $\hat b=\sum x_iy_i\big/\sum x_i^2$.
			\item\label{ex:ip:A18} 在 $\mathbb{R}^3$ 中，求点 $(1,2,3)$ 到直线 $\vv x=t(1,1,1)^\trans$ 的距离.
			\item\label{ex:ip:A19} 证明：两组标准正交基之间的过渡矩阵是正交矩阵.
			\item\label{ex:ip:A20} 三个法条的词频向量（关键词：违约、赔偿、合同、离婚）分别为 $(2,1,3,0)^\trans$，$(0,1,0,4)^\trans$，$(1,1,2,0)^\trans$，查询为 $(1,1,1,0)^\trans$. 分别用内积和余弦相似度给三个法条排序，结果有何不同？
			\item\label{ex:ip:A21} 设 $\vv u\in\mathbb{R}^n$ 是单位向量. 证明 $\vv P=\vv u\vv u^\trans$ 满足 $\vv P^\trans=\vv P$，$\vv P^2=\vv P$，且对任意 $\vv x$，$\vv P\vv x$ 是 $\vv x$ 在 $\spn(\vv u)$ 上的正交投影. 若 $\vv u$ 不是单位向量，投影矩阵应当怎样修改？
		\end{exA}
		
		\Bgroup
		\begin{exB}
			\item\label{ex:ip:B1} （极化恒等式）证明实内积空间中 $\langle\vv\alpha,\vv\beta\rangle=\frac14\big(\|\vv\alpha+\vv\beta\|^2-\|\vv\alpha-\vv\beta\|^2\big)$. 再证明：$\mathbb{R}^2$ 上的``范数'' $\|\vv x\|_1=|x_1|+|x_2|$ 不满足平行四边形法则，因而不能来自任何内积.
			\item\label{ex:ip:B2} （Bessel 不等式）设 $\vv e_1,\ldots,\vv e_k$ 是标准正交组. 证明 $\sum_{i=1}^k\langle\vv\alpha,\vv e_i\rangle^2\le\|\vv\alpha\|^2$，等号成立当且仅当 $\vv\alpha\in\spn(\vv e_1,\ldots,\vv e_k)$.
			\item\label{ex:ip:B3} 设实矩阵 $\vv P$ 满足 $\vv P^\trans=\vv P$，$\vv P^2=\vv P$. 证明 $\vv P$ 是向 $\imm\vv P$ 的正交投影的矩阵. 若只有 $\vv P^2=\vv P$，它是什么样的投影？举例说明.
			\item\label{ex:ip:B4} （加权最小二乘）若各观测值的可靠程度不同，给第 $i$ 个残差权重 $w_i>0$，最小化 $\sum w_i(y_i-a-bx_i)^2$. 证明正规方程变为 $\vv A^\trans\vv W\vv A\hat{\vv\beta}=\vv A^\trans\vv W\vv y$，$\vv W=\diag(w_1,\ldots,w_n)$. 它对应于哪个内积下的正交投影？
			\item\label{ex:ip:B5} 某商品的销售额 $y$ 随时间 $t$ 近似按指数增长 $y=c\,\mathrm e^{kt}$. 已知 $t=0,1,2,3$ 时 $y$ 约为 $2.0,2.7,3.7,5.0$. 两边取对数化为线性模型，用最小二乘估计 $c,k$. 这样做最小化的是哪个量？与直接最小化 $\sum(y_i-c\mathrm e^{kt_i})^2$ 有何不同？
			\item\label{ex:ip:B6} （Gram 行列式）设 $\vv\alpha_1,\ldots,\vv\alpha_k\in\mathbb{R}^n$，$\vv G=(\langle\vv\alpha_i,\vv\alpha_j\rangle)$. 证明 $\det\vv G\ge0$，等号成立当且仅当它们线性相关. 在 $k=2$ 时说明 $\sqrt{\det\vv G}$ 是两向量张成的平行四边形的面积.
			\item\label{ex:ip:B7} 证明：$\mathbb{R}^n$ 中点 $\vv x_0$ 到超平面 $\vv a^\trans\vv x=b$（$\vv a\ne\zero$）的距离为 $|\vv a^\trans\vv x_0-b|/\|\vv a\|$.
			\item\label{ex:ip:B8} 证明：实正交矩阵的复特征值的模都等于 $1$.
			\item\label{ex:ip:B9} （Fourier 逼近）在 $C[-\pi,\pi]$ 中取 $\langle f,g\rangle=\int_{-\pi}^\pi fg\,\mathrm dx$. 证明 $1,\cos x,\sin x,\cos2x,\sin2x$ 两两正交. 求 $f(x)=x$ 在它们张成的子空间上的正交投影.
			\item\label{ex:ip:B10} 设 $\vv A$ 为实矩阵. 证明：在 $\vv{Ax}=\vv y$ 的全部最小二乘解中，范数最小的那一个唯一，并且它属于行空间 $\imm\vv A^\trans$.
			\item\label{ex:ip:B11} 设 $\vv x,\vv y\in\mathbb{R}^n$，$\|\vv x\|=\|\vv y\|$，$\vv x\ne\vv y$. 令 $\vv u=(\vv x-\vv y)/\|\vv x-\vv y\|$. 证明 Householder 矩阵 $\vv I-2\vv u\vv u^\trans$ 把 $\vv x$ 变为 $\vv y$. 这是用 Householder 反射做 QR 分解的基础.
			\item\label{ex:ip:B12} 设两个变量的样本标准差分别为 $s_x,s_y$，相关系数为 $r$. 证明回归直线 $y=a+bx$ 的斜率为 $b=r\,s_y/s_x$，并且直线经过点 $(\bar x,\bar y)$.
			\item\label{ex:ip:B13} 证明：保持\textbf{夹角}的可逆线性变换（对一切非零向量，$\mathscr T\vv\alpha$ 与 $\mathscr T\vv\beta$ 的夹角等于 $\vv\alpha$ 与 $\vv\beta$ 的夹角）必是正交变换的非零倍数.
			\item\label{ex:ip:B14} 在回归中加入一个``虚拟变量''：设有两组数据，第一组 $y$ 值为 $3,5,4$，第二组为 $7,8,9$. 以 $\vv A=(\vv 1,\vv d)$ 为设计矩阵（$\vv d$ 在第二组为 $1$、第一组为 $0$），求最小二乘解，并解释两个系数的含义.
			\item\label{ex:ip:B15} 证明：在 $\mathbb{R}^2$ 中，任意两个旋转可交换，而一个旋转与一个反射一般不可交换；两个反射的乘积是旋转.
		\end{exB}
		
		\Cgroup
		\begin{exC}
			\item\label{ex:ip:C1} （Hadamard 不等式）用 QR 分解证明：对 $n$ 阶实方阵 $\vv A=(\vv a_1,\ldots,\vv a_n)$，$|\det\vv A|\le\|\vv a_1\|\cdots\|\vv a_n\|$，并说明它的几何意义（平行多面体的体积不超过棱长之积）.
			\item\label{ex:ip:C2} （Gram–Schmidt 的数值不稳定）设 $\varepsilon$ 很小，对 $\vv a_1=(1,\varepsilon,0,0)^\trans$，$\vv a_2=(1,0,\varepsilon,0)^\trans$，$\vv a_3=(1,0,0,\varepsilon)^\trans$ 作 Gram–Schmidt，在计算中把 $1+\varepsilon^2$ 舍入为 $1$. 证明所得的 $\vv e_2,\vv e_3$ 之间的夹角是 $60^\circ$ 而不是 $90^\circ$. 再用``修正的 Gram–Schmidt''（每得到一个 $\vv e_k$，立即从\textbf{剩余的}所有向量中减去它的分量）重做，比较结果.
			\item\label{ex:ip:C3} （Legendre 多项式）令 $P_n(x)=\frac1{2^nn!}\frac{\mathrm d^n}{\mathrm dx^n}(x^2-1)^n$. 证明 $P_0,P_1,P_2,\ldots$ 在 $\langle f,g\rangle=\int_{-1}^1fg\,\mathrm dx$ 下两两正交，并求 $\|P_n\|$.
			\item\label{ex:ip:C4} （Gauss–Markov 定理的线性代数部分）设 $\vv A$ 列满秩. 若 $\vv C\vv y$ 是参数 $\vv\beta$ 的``线性无偏估计''，即 $\vv{CA}=\vv I$. 令 $\vv C_0=(\vv A^\trans\vv A)^{-1}\vv A^\trans$（最小二乘）. 证明 $\vv C\vv C^\trans-\vv C_0\vv C_0^\trans$ 是半正定矩阵（第十一章）. 学过概率论后可以知道，这说明最小二乘估计的方差在一切线性无偏估计中最小.
			\item\label{ex:ip:C5} （复内积空间）在 $\mathbb{C}^n$ 上定义 $\langle\vv x,\vv y\rangle=\sum x_i\bar y_i$. 证明它满足共轭对称性、对第一个变量线性、正定性；证明 Cauchy–Schwarz 不等式仍成立. 满足 $\bar{\vv U}^\trans\vv U=\vv I$ 的复矩阵称为\textbf{酉矩阵}，证明它保持这个内积.
			\item\label{ex:ip:C6} 设 $f:\mathbb{R}^n\to\mathbb{R}^n$ 保持距离（$\|f(\vv x)-f(\vv y)\|=\|\vv x-\vv y\|$）且 $f(\zero)=\zero$. 证明 $f$ 保持内积，进而证明 $f$ 是线性的正交变换. 由此说明：保持距离的变换都可以分解为正交变换与平移的复合.
			\item\label{ex:ip:C7} （伪逆）设 $\vv A$ 列满秩，$\vv A^+=(\vv A^\trans\vv A)^{-1}\vv A^\trans$. 验证 Penrose 四条件：$\vv{AA}^+\vv A=\vv A$，$\vv A^+\vv{AA}^+=\vv A^+$，$(\vv{AA}^+)^\trans=\vv{AA}^+$，$(\vv A^+\vv A)^\trans=\vv A^+\vv A$. 说明 $\vv{AA}^+$ 是向哪个子空间的正交投影.
			\item\label{ex:ip:C8} 求空间中两条异面直线 $\vv x=\vv p+s\vv u$ 与 $\vv x=\vv q+t\vv v$ 之间的最短距离：把问题写成关于 $(s,t)$ 的最小二乘问题，并对 $\vv p=\zero$，$\vv u=(1,0,0)^\trans$，$\vv q=(0,1,1)^\trans$，$\vv v=(0,1,-1)^\trans$ 计算.
		\end{exC}
	}
	
	% =====================================================================
	% 第十章　矩阵分解
	% =====================================================================
	\newcommand{\ExDecomposition}{%
		\exsetheading
		\Agroup
		\begin{exA}
			\item\label{ex:dec:A1} 求正交矩阵 $\vv Q$ 使 $\vv Q^\trans\vv{AQ}$ 为对角矩阵：
			\[\begin{pmatrix}1&2\\2&1\end{pmatrix},\qquad\begin{pmatrix}3&-1\\-1&3\end{pmatrix},\]\[\begin{pmatrix}2&2&0\\2&2&0\\0&0&1\end{pmatrix},\qquad\begin{pmatrix}1&1&1\\1&1&1\\1&1&1\end{pmatrix}.\]
			\item\label{ex:dec:A2} 写出 $\begin{pmatrix}3&-1\\-1&3\end{pmatrix}$ 的谱分解 $\lambda_1\vv q_1\vv q_1^\trans+\lambda_2\vv q_2\vv q_2^\trans$，并验证两个投影矩阵之和为 $\vv I$、之积为 $\vv O$.
			\item\label{ex:dec:A3} 求二阶实对称矩阵 $\vv A$，使它的特征值为 $1,2$，且 $(1,1)^\trans$ 是属于 $1$ 的特征向量.
			\item\label{ex:dec:A4} 设三阶实对称矩阵 $\vv A$ 的特征值为 $1,1,-1$，属于 $-1$ 的特征向量为 $(0,1,1)^\trans$. 求 $\vv A$.
			\item\label{ex:dec:A5} 求 $\vv x^\trans\vv{Ax}$ 在单位球面 $\|\vv x\|=1$ 上的最大值与最小值，以及取到它们的 $\vv x$：$\vv A=\begin{pmatrix}2&2&0\\2&2&0\\0&0&1\end{pmatrix}$.
			\item\label{ex:dec:A6} 求 $x_1^2+4x_1x_2+x_2^2$ 在单位圆 $x_1^2+x_2^2=1$ 上的最大值与最小值.
			\item\label{ex:dec:A7} 求下列矩阵的奇异值分解：$\begin{pmatrix}2&0\\0&-3\end{pmatrix}$，$\begin{pmatrix}1&1\\1&1\end{pmatrix}$，$(3,4)$（$1\times2$ 矩阵）.
			\item\label{ex:dec:A8} 求 $\begin{pmatrix}1&1\\0&1\end{pmatrix}$ 的奇异值，并与它的特征值比较.（答案与黄金分割比有关.）
			\item\label{ex:dec:A9} 求 $\begin{pmatrix}3&1\\1&3\end{pmatrix}$ 的最佳秩一近似，并计算误差 $\|\vv A-\vv A_1\|_F$.
			\item\label{ex:dec:A10} 设 $\vv A=\begin{pmatrix}5&4\\4&5\end{pmatrix}$. 利用谱分解求 $\vv A^{10}$（写成两个矩阵之和的形式即可），并求对称正定矩阵 $\vv B$ 使 $\vv B^2=\vv A$.
			\item\label{ex:dec:A11} 证明：特征值全为 $0$ 的实对称矩阵是零矩阵. 对非对称矩阵举出反例.
			\item\label{ex:dec:A12} 证明：既对称又正交的实矩阵 $\vv A$ 满足 $\vv A^2=\vv I$，其特征值为 $\pm1$；并说明它在几何上是什么变换.
			\item\label{ex:dec:A13} 证明：$\vv A^\trans\vv A$ 与 $\vv{AA}^\trans$ 的非零特征值相同（重数也相同）.
			\item\label{ex:dec:A14} 证明：$\|\vv A\|_F^2=\tr(\vv A^\trans\vv A)=\sigma_1^2+\cdots+\sigma_r^2$.
			\item\label{ex:dec:A15} 设 $\vv A$ 是可逆实方阵. 证明 $|\det\vv A|=\sigma_1\sigma_2\cdots\sigma_n$，$\vv A^{-1}$ 的奇异值为 $1/\sigma_i$.
			\item\label{ex:dec:A16} 判断下列矩阵能否用正交矩阵对角化，说明理由：$\begin{pmatrix}1&2\\0&3\end{pmatrix}$，$\begin{pmatrix}1&2\\2&-2\end{pmatrix}$，$\begin{pmatrix}0&-1\\1&0\end{pmatrix}$（在实数范围内）.
			\item\label{ex:dec:A17} 四个已中心化的二维数据点为 $(2,1)$，$(-2,-1)$，$(1,2)$，$(-1,-2)$. 求样本协方差矩阵、第一主成分方向及其解释的方差比例.
			\item\label{ex:dec:A18} 计算 $\diag(10,1)$ 与 $\begin{pmatrix}1&1\\1&1.01\end{pmatrix}$ 的条件数（后者用近似值）. 后者的行列式与 $\diag(10,0.001)$ 相同，二者的条件数相同吗？
			\item\label{ex:dec:A19} 一张 $800\times600$ 像素的灰度图片，保留前 $k=40$ 个奇异值时需存储多少个数？压缩比是多少？$k$ 取多少时存储量反而超过原图？
			\item\label{ex:dec:A20} 证明：实对称矩阵的秩等于非零特征值的个数；奇异值的个数（非零）等于秩.
			\item\label{ex:dec:A21} 证明：实对称矩阵 $\vv A$ 满足``对一切 $\vv x\ne\zero$，$\vv x^\trans\vv A\vv x>0$''，当且仅当 $\vv A$ 的特征值全为正.（提示：用谱定理把 $\vv x^\trans\vv A\vv x$ 写成 $\sum\lambda_iy_i^2$.）
		\end{exA}
		
		\Bgroup
		\begin{exB}
			\item\label{ex:dec:B1} 设 $\vv A,\vv B$ 都是 $n$ 阶实对称矩阵且 $\vv{AB}=\vv{BA}$. 证明存在正交矩阵 $\vv Q$ 使 $\vv Q^\trans\vv{AQ}$ 与 $\vv Q^\trans\vv{BQ}$ 同时为对角矩阵.
			\item\label{ex:dec:B2} 证明：半正定的实对称矩阵 $\vv A$（特征值全 $\ge0$）有唯一的半正定平方根 $\vv B$（$\vv B^2=\vv A$，$\vv B$ 对称且特征值全 $\ge0$）.
			\item\label{ex:dec:B3} 设实对称矩阵 $\vv A$ 的特征值为 $\lambda_1\ge\lambda_2\ge\cdots\ge\lambda_n$，$\vv q_1$ 是属于 $\lambda_1$ 的单位特征向量. 证明 $\lambda_2=\max\{\vv x^\trans\vv{Ax}:\|\vv x\|=1,\ \vv x\perp\vv q_1\}$. 这说明第二主成分是``与第一主成分垂直的方向中方差最大的''.
			\item\label{ex:dec:B4} 证明：对任意 $\vv x$，$\|\vv{Ax}\|\le\sigma_1\|\vv x\|$，且等号可以取到. 因此 $\sigma_1=\max_{\vv x\ne\zero}\|\vv{Ax}\|/\|\vv x\|$，称为 $\vv A$ 的\textbf{谱范数}\index{pu3 fan4shu4@谱范数} $\|\vv A\|_2$.
			\item\label{ex:dec:B5} 证明：方阵 $\vv A$ 的任一特征值 $\lambda$ 满足 $\sigma_{\min}\le|\lambda|\le\sigma_{\max}$.
			\item\label{ex:dec:B6} 求 $\vv A=\begin{pmatrix}1&0&1\\0&1&1\end{pmatrix}$ 的奇异值分解.
			\item\label{ex:dec:B7} （伪逆）设 $\vv A=\vv U\vv\Sigma\vv V^\trans$，定义 $\vv A^+=\vv V\vv\Sigma^+\vv U^\trans$，其中 $\vv\Sigma^+$ 是把 $\vv\Sigma$ 转置并把非零元取倒数. 证明 $\vv A^+\vv b$ 是 $\vv{Ax}=\vv b$ 的范数最小的最小二乘解. 当 $\vv A$ 列满秩时，它与第九章的 $(\vv A^\trans\vv A)^{-1}\vv A^\trans$ 一致.
			\item\label{ex:dec:B8} 设 $p$ 个标准化指标两两相关系数都为 $\rho$（$0<\rho<1$），相关系数矩阵为 $(1-\rho)\vv I+\rho\vv J$（$\vv J$ 为全 $1$ 矩阵）. 求它的特征值与第一主成分，并求第一主成分解释的方差比例. 当 $\rho\to1$ 与 $\rho\to0$ 时比例分别趋于多少？
			\item\label{ex:dec:B9} 设两个指标的协方差矩阵为 $\begin{pmatrix}100&6\\6&1\end{pmatrix}$（第一个指标以``元''计，第二个以``万元''计）. 对协方差矩阵与相关系数矩阵分别作主成分分析，比较第一主成分. 由此说明：主成分分析对量纲敏感，为什么通常先标准化.
			\item\label{ex:dec:B10} 证明 $\|\vv A-\vv A_k\|_2=\sigma_{k+1}$.
			\item\label{ex:dec:B11} 证明：若 $\vv A$ 的奇异值互不相同且都不为零，则 $\vv A$ 的奇异值分解中的 $\vv u_i,\vv v_i$ 除了同时改变符号之外是唯一的.
			\item\label{ex:dec:B12} （正规矩阵）满足 $\vv A^\trans\vv A=\vv{AA}^\trans$ 的实方阵称为正规矩阵. 证明二阶实正规矩阵只有两类：对称矩阵，以及形如 $\begin{pmatrix}a&-b\\b&a\end{pmatrix}$ 的矩阵（旋转与伸缩的复合）.
			\item\label{ex:dec:B13} 设 $\vv A$ 是实对称矩阵. 证明 $\vv A$ 的特征值都落在 $[\min_i(a_{ii}-r_i),\ \max_i(a_{ii}+r_i)]$ 中，$r_i=\sum_{j\ne i}|a_{ij}|$. 并证明 $\lambda_{\max}\ge\max_ia_{ii}$.
			\item\label{ex:dec:B14} 对第十章``潜在语义分析''的例子，取 $k=1$，说明此时四份文书分别被表示成什么，查询``解约''与哪几份文书相似. 与 $k=2$ 比较，``压缩过头''的现象是否更严重？
			\item\label{ex:dec:B15} 证明：实对称矩阵 $\vv A$ 的特征值都在 $[a,b]$ 中，当且仅当 $\vv A-a\vv I$ 与 $b\vv I-\vv A$ 的二次型都非负.
		\end{exB}
		
		\Cgroup
		\begin{exC}
			\item\label{ex:dec:C1} （Eckart–Young 定理的证明）设 $\rk\vv B\le k$. 证明 $\ker\vv B$ 与 $\spn(\vv v_1,\ldots,\vv v_{k+1})$ 有非零公共向量 $\vv x$，由此证明 $\|\vv A-\vv B\|_2\ge\sigma_{k+1}$. 再借助 Weyl 不等式 $\sigma_{i+j-1}(\vv X+\vv Y)\le\sigma_i(\vv X)+\sigma_j(\vv Y)$（可承认）证明 Frobenius 范数的情形.
			\item\label{ex:dec:C2} （Cauchy 交错定理）设实对称矩阵 $\vv A$ 的特征值为 $\lambda_1\ge\cdots\ge\lambda_n$，删去最后一行一列所得矩阵的特征值为 $\mu_1\ge\cdots\ge\mu_{n-1}$. 证明 $\lambda_1\ge\mu_1\ge\lambda_2\ge\mu_2\ge\cdots\ge\mu_{n-1}\ge\lambda_n$.（提示：Courant–Fischer 极小极大刻画.）
			\item\label{ex:dec:C3} （复谱定理）证明：复方阵 $\vv A$ 可以被酉矩阵对角化（存在 $\vv U$ 使 $\bar{\vv U}^\trans\vv{AU}$ 为对角矩阵），当且仅当 $\vv A\bar{\vv A}^\trans=\bar{\vv A}^\trans\vv A$. 实对称矩阵、实正交矩阵、实反对称矩阵都满足这个条件.
			\item\label{ex:dec:C4} （正交 Procrustes 问题）给定 $n\times k$ 矩阵 $\vv A,\vv B$，求正交矩阵 $\vv Q$ 使 $\|\vv{QA}-\vv B\|_F$ 最小. 证明：若 $\vv{BA}^\trans=\vv U\vv\Sigma\vv V^\trans$，则最优解为 $\vv Q=\vv{UV}^\trans$. 这个问题出现在比较两组测量数据、对齐两份词向量等场合.
			\item\label{ex:dec:C5} （Ky Fan 定理）设实对称矩阵的特征值为 $\lambda_1\ge\cdots\ge\lambda_n$. 证明对任意列标准正交的 $n\times k$ 矩阵 $\vv X$，$\tr(\vv X^\trans\vv{AX})\le\lambda_1+\cdots+\lambda_k$. 由此说明前 $k$ 个主成分``一起''解释的方差也是最大的.
			\item\label{ex:dec:C6} （幂迭代）设 $\vv A$ 可对角化，特征值满足 $|\lambda_1|>|\lambda_2|\ge\cdots$. 证明：对``一般的''初始向量 $\vv x_0$，$\vv x_k=\vv A^k\vv x_0/\|\vv A^k\vv x_0\|$ 收敛到属于 $\lambda_1$ 的单位特征向量（可能差一个符号），收敛速度由 $|\lambda_2/\lambda_1|$ 决定. 对 $\begin{pmatrix}2&1\\1&2\end{pmatrix}$ 从 $(1,0)^\trans$ 出发迭代五步验证.
			\item\label{ex:dec:C7} （图的 Laplace 矩阵）设 $\vv L=\vv D-\vv A$ 是无向图的 Laplace 矩阵（见第五章章末习题~\ref{ex:det:C4}）. 证明 $\vv x^\trans\vv L\vv x=\sum_{\text{边 }ij}(x_i-x_j)^2$，因而 $\vv L$ 半正定；特征值 $0$ 的重数等于图的连通分支数. 第二小的特征值对应的特征向量（Fiedler 向量）的正负号常被用来把图``切''成两半——这就是谱聚类的思想.
			\item\label{ex:dec:C8} （正交回归）对中心化数据点 $\vv z_1,\ldots,\vv z_n\in\mathbb{R}^2$，求过原点的直线使各点到直线的\textbf{垂直}距离平方和最小. 证明最优方向就是第一主成分方向. 与普通最小二乘（只看竖直方向的残差）比较，两条直线何时相同？
			\item\label{ex:dec:C9} （扰动分析）设 $\vv A$ 可逆，$\vv{Ax}=\vv b$，$\vv A(\vv x+\delta\vv x)=\vv b+\delta\vv b$. 证明 $\dfrac{\|\delta\vv x\|}{\|\vv x\|}\le\kappa(\vv A)\dfrac{\|\delta\vv b\|}{\|\vv b\|}$，$\kappa(\vv A)=\sigma_1/\sigma_n$，并说明等号可以取到. 用它解释第六章病态方程组的例子.
		\end{exC}
	}
	
	% =====================================================================
	% 第十一章　二次型
	% =====================================================================
	\newcommand{\ExQuadratic}{%
		\exsetheading
		\Agroup
		\begin{exA}
			\item\label{ex:quad:A1} 写出下列二次型的矩阵与秩：
			(1) $x_1^2-3x_2^2+4x_1x_2-2x_2x_3$；(2) $x_1x_2+x_2x_3+x_3x_1$；(3) $(x_1+x_2+x_3)^2$；(4) $\sum_{i=1}^n x_i^2+\sum_{i<j}x_ix_j$.
			\item\label{ex:quad:A2} 写出以 $\begin{pmatrix}1&-1&0\\-1&2&3\\0&3&0\end{pmatrix}$ 为矩阵的二次型.
			\item\label{ex:quad:A3} 用配方法化为标准形，并写出所用的可逆变量替换：
			(1) $x_1^2+2x_2^2+5x_3^2+2x_1x_2+6x_2x_3$；(2) $x_1^2-4x_1x_2+4x_2^2-x_3^2$；(3) $x_1x_2-x_2x_3$.
			\item\label{ex:quad:A4} 用正交变换法把 $2x_1^2+3x_2^2+3x_3^2+4x_2x_3$ 化为标准形，写出正交矩阵. 曲面 $2x_1^2+3x_2^2+3x_3^2+4x_2x_3=1$ 是什么形状？
			\item\label{ex:quad:A5} 求 A3 中各二次型的正、负惯性指数与符号差，并写出规范形.
			\item\label{ex:quad:A6} 判断下列二次型是否正定：
			(1) $2x_1^2+x_2^2+3x_3^2+2x_1x_2$；(2) $x_1^2+x_2^2+x_3^2+4x_1x_2$；(3) $5x_1^2+x_2^2+5x_3^2+4x_1x_2-8x_1x_3-4x_2x_3$.
			\item\label{ex:quad:A7} $t$ 取何值时，$2x_1^2+x_2^2+x_3^2+2x_1x_2+tx_2x_3$ 正定？
			\item\label{ex:quad:A8} 判断曲线类型并求主轴方向：(1) $x^2-4xy+y^2=1$；(2) $3x^2+2xy+3y^2=8$.
			\item\label{ex:quad:A9} 判断下列矩阵是否合同：(1) $\diag(1,2,-3)$ 与 $\diag(-1,5,6)$；(2) $\begin{pmatrix}0&1\\1&0\end{pmatrix}$ 与 $\vv I_2$；(3) $\begin{pmatrix}0&1\\1&0\end{pmatrix}$ 与 $\diag(1,-1)$. 合同时求出 $\vv C$.
			\item\label{ex:quad:A10} 证明：若 $\vv A$ 正定，则 $\vv A^*$ 也正定.
			\item\label{ex:quad:A11} 求 $f(x,y)=x^3+y^3-3xy$ 的全部驻点，并用 Hessian 矩阵判断其类型.
			\item\label{ex:quad:A12} 判断 $\begin{pmatrix}2&-1&0\\-1&2&-1\\0&-1&2\end{pmatrix}$ 是否正定，并求它的 Cholesky 分解.
			\item\label{ex:quad:A13} 求 $\begin{pmatrix}4&2&-2\\2&10&2\\-2&2&6\end{pmatrix}$ 的 Cholesky 分解.
			\item\label{ex:quad:A14} 三种资产收益率的协方差矩阵为 $\vv\Sigma=\begin{pmatrix}0.04&0.01&0\\0.01&0.09&0\\0&0&0.01\end{pmatrix}$. 求组合 $\vv w=(0.5,0.3,0.2)^\trans$ 的方差与标准差.
			\item\label{ex:quad:A15} 两种资产收益率的标准差为 $30\%$ 与 $20\%$，互不相关. 求最小方差组合及其标准差.
			\item\label{ex:quad:A16} 证明：合同变换保持矩阵的对称性与秩. 合同变换保持行列式的符号吗？保持行列式的值吗？
			\item\label{ex:quad:A17} 分别在 $\mathbb{R}$ 上与 $\mathbb{C}$ 上写出二次型 $x_1^2-x_2^2+2x_3^2$ 的规范形.
			\item\label{ex:quad:A18} 判断下列方程表示的二次曲面类型：$x^2+y^2-z^2=1$，$x^2-y^2-z^2=1$，$x^2+y^2-z^2=0$，$x^2+y^2=1$.（提示：先看正、负惯性指数.）
			\item\label{ex:quad:A19} 说明 $\begin{pmatrix}1&2\\2&1\end{pmatrix}$ 不正定，并找出一个使二次型取负值的向量.
			\item\label{ex:quad:A20} 设 $\vv A$ 正定. 证明：$\vv A$ 的任一主子矩阵（取相同的若干行和列）都正定，从而所有主子式都为正.
		\end{exA}
		
		\Bgroup
		\begin{exB}
			\item\label{ex:quad:B1} 设 $\vv A,\vv B$ 正定且 $\vv{AB}=\vv{BA}$. 证明 $\vv{AB}$ 正定. 举例说明``可交换''不能去掉（此时 $\vv{AB}$ 甚至不对称）.
			\item\label{ex:quad:B2} 设 $\vv A$ 正定，$\vv B$ 半正定. 证明 $\det(\vv A+\vv B)\ge\det\vv A$，等号成立当且仅当 $\vv B=\vv O$.（提示：同时合同对角化.）
			\item\label{ex:quad:B3} 设 $\vv A$ 正定，$\vv B$ 对称. 证明：当 $|t|$ 充分小时 $\vv A+t\vv B$ 正定.
			\item\label{ex:quad:B4} 设 $\vv A$ 正定. 证明对 $i\ne j$ 有 $|a_{ij}|<\sqrt{a_{ii}a_{jj}}$，从而 $\vv A$ 的绝对值最大的元素在对角线上.
			\item\label{ex:quad:B5} 证明：实对称矩阵 $\vv A$ 满足 $\lambda_{\min}\|\vv x\|^2\le\vv x^\trans\vv{Ax}\le\lambda_{\max}\|\vv x\|^2$.
			\item\label{ex:quad:B6} 求 $\begin{pmatrix}0&1&1\\1&0&1\\1&1&0\end{pmatrix}$ 的正、负惯性指数（用两种方法：特征值；配方法）.
			\item\label{ex:quad:B7} $\lambda$ 取何值时，$\lambda(x_1^2+x_2^2+x_3^2)+2x_1x_2+2x_1x_3-2x_2x_3$ 正定？
			\item\label{ex:quad:B8} 求二次型 $\vv x^\trans\vv{Ax}$ 在约束 $\vv c^\trans\vv x=1$ 下的最小值（$\vv A$ 正定）：用拉格朗日乘数法证明最优解为 $\vv x^*=\dfrac{\vv A^{-1}\vv c}{\vv c^\trans\vv A^{-1}\vv c}$. 最小方差组合是它的特例.
			\item\label{ex:quad:B9} 设三种资产互不相关、方差都为 $\sigma^2$，期望收益率为 $\mu_1<\mu_2<\mu_3$. 求期望收益为 $m$ 的最小方差组合的权重，并说明有效前沿是一条抛物线（在``方差—期望收益''平面上）.
			\item\label{ex:quad:B10} 设 $\vv A=\begin{pmatrix}2&1\\1&2\end{pmatrix}$，$\vv B=\begin{pmatrix}1&0\\0&-1\end{pmatrix}$. 求可逆矩阵 $\vv C$ 使 $\vv C^\trans\vv{AC}=\vv I$ 且 $\vv C^\trans\vv{BC}$ 为对角矩阵.
			\item\label{ex:quad:B11} 设 $\vv A$ 正定. 求 $f(\vv x)=\vv x^\trans\vv{Ax}+\vv b^\trans\vv x+c$ 的最小值点与最小值. 若 $\vv A$ 只是半正定，最小值一定存在吗？
			\item\label{ex:quad:B12} 设 $\vv A$ 半正定，$\vv x^\trans\vv{Ax}=0$. 证明 $\vv{Ax}=\zero$.
			\item\label{ex:quad:B13} （Schur 补判据）证明：对称分块矩阵 $\begin{pmatrix}\vv A&\vv B\\ \vv B^\trans&\vv C\end{pmatrix}$ 正定，当且仅当 $\vv A$ 正定且 $\vv C-\vv B^\trans\vv A^{-1}\vv B$ 正定.
			\item\label{ex:quad:B14} 设 $\vv A$ 是二阶正定矩阵. 证明椭圆区域 $\{\vv x:\vv x^\trans\vv{Ax}\le1\}$ 的面积为 $\pi/\sqrt{\det\vv A}$.（提示：用主轴定理化为标准椭圆.）
			\item\label{ex:quad:B15} 设 $\vv A$ 为 $m\times n$ 实矩阵，$\mu>0$. 证明 $\vv A^\trans\vv A+\mu\vv I$ 正定，并说明岭回归 $\min\|\vv y-\vv{Ax}\|^2+\mu\|\vv x\|^2$ 的解为 $(\vv A^\trans\vv A+\mu\vv I)^{-1}\vv A^\trans\vv y$.
			\item\label{ex:quad:B16} 证明：实对称矩阵 $\vv A$ 正定，当且仅当存在正定矩阵 $\vv B$ 使 $\vv A=\vv B^2$；当且仅当 $\vv A^{-1}$ 存在且正定.
		\end{exB}
		
		\Cgroup
		\begin{exC}
			\item\label{ex:quad:C1} （Schur 乘积定理）对同阶矩阵定义逐元素乘积 $(\vv A\circ\vv B)_{ij}=a_{ij}b_{ij}$. 证明：若 $\vv A,\vv B$ 都半正定，则 $\vv A\circ\vv B$ 也半正定.（提示：把 $\vv B$ 写成 $\sum\lambda_k\vv q_k\vv q_k^\trans$.）
			\item\label{ex:quad:C2} （Minkowski 行列式不等式）设 $\vv A,\vv B$ 都是 $n$ 阶正定矩阵. 证明 $\det(\vv A+\vv B)^{1/n}\ge\det(\vv A)^{1/n}+\det(\vv B)^{1/n}$.（提示：同时合同对角化后用 $\prod(1+\mu_i)^{1/n}\ge1+\prod\mu_i^{1/n}$.）
			\item\label{ex:quad:C3} （有理数域上的二次型）证明 $x^2+y^2$ 与 $3x^2+3y^2$ 在 $\mathbb{R}$ 上合同，但在 $\mathbb{Q}$ 上不合同（即不存在有理系数的可逆替换把一个变为另一个）.（提示：若合同，则存在有理数 $a,b$ 使 $a^2+b^2=3$.）这说明惯性定理依赖于实数可以随意开平方.
			\item\label{ex:quad:C4} 函数 $f=x^2+y^4$ 与 $g=x^2-y^4$ 在原点处的 Hessian 矩阵相同且只是半正定，但前者在原点取极小值，后者不取极值. 由此说明 Hessian 半正定时二阶判别法失效.
			\item\label{ex:quad:C5} （广义特征值）设 $\vv A$ 对称，$\vv B$ 正定. 证明 $\max_{\vv x\ne\zero}\dfrac{\vv x^\trans\vv{Ax}}{\vv x^\trans\vv{Bx}}$ 等于方程 $\det(\vv A-\lambda\vv B)=0$ 的最大根. 统计学中的判别分析（Fisher LDA）正是这样求最佳投影方向的.
			\item\label{ex:quad:C6} （Mahalanobis 距离）设 $\vv\Sigma$ 正定，定义 $d(\vv x,\vv y)=\sqrt{(\vv x-\vv y)^\trans\vv\Sigma^{-1}(\vv x-\vv y)}$. 证明它满足三角不等式，且在仿射变换 $\vv x\mapsto\vv{Px}+\vv b$、$\vv\Sigma\mapsto\vv P\vv\Sigma\vv P^\trans$ 下不变. 它与欧氏距离有何不同？
			\item\label{ex:quad:C7} （Hilbert 矩阵）证明 $n$ 阶矩阵 $\vv H=\big(\frac1{i+j-1}\big)$ 正定.（提示：$\frac1{i+j-1}=\int_0^1x^{i-1}x^{j-1}\,\mathrm dx$，它是某个内积下的 Gram 矩阵.）
			\item\label{ex:quad:C8} （切点组合）设有无风险资产，收益率为 $r_f$. 证明使``夏普比率''$\dfrac{\vv\mu^\trans\vv w-r_f}{\sqrt{\vv w^\trans\vv\Sigma\vv w}}$（约束 $\vv 1^\trans\vv w=1$）最大的风险资产组合为 $\vv w\propto\vv\Sigma^{-1}(\vv\mu-r_f\vv 1)$.（提示：先去掉约束，比值关于 $\vv w$ 是零次齐次的；再用 Cauchy–Schwarz 不等式.）
			\item\label{ex:quad:C9} 证明：秩为 $r$、正惯性指数为 $p$ 的实二次型可以写成 $p$ 个线性型的平方和减去 $r-p$ 个线性型的平方和，并且不能写成少于 $r$ 个线性型的平方（带 $\pm$ 号）的组合.
		\end{exC}
	}
	
	% =====================================================================
	% 第十二章　*对偶空间与双线性型（选读）
	% =====================================================================
	\newcommand{\ExBilinear}{%
		\exsetheading
		\Agroup
		\begin{exA}
			\item\label{ex:bil:A1} 判断下列映射是否为相应空间上的线性函数：(1) $\mathbb{R}^3\to\mathbb{R}$，$\vv x\mapsto x_1-2x_3$；(2) $\vv x\mapsto x_1x_2$；(3) $\mathbb{R}[x]\to\mathbb{R}$，$f\mapsto f'(0)$；(4) $f\mapsto\int_0^1f(x)^2\,\mathrm dx$；(5) $M_n(\mathbb{R})\to\mathbb{R}$，$\vv X\mapsto\tr(\vv{AX})$（$\vv A$ 固定）.
			\item\label{ex:bil:A2} 求 $\mathbb{R}^3$ 的基 $(1,0,0)^\trans$，$(1,1,0)^\trans$，$(1,1,1)^\trans$ 的对偶基（把每个对偶基向量写成 $\vv x\mapsto a_1x_1+a_2x_2+a_3x_3$ 的形式）.
			\item\label{ex:bil:A3} 把线性函数 $f(\vv x)=x_1+2x_2+3x_3$ 用 A2 中的对偶基线性表出.
			\item\label{ex:bil:A4} 求 $W=\spn\big((1,1,0)^\trans,(0,1,1)^\trans\big)$ 的零化子 $W^0\subset(\mathbb{R}^3)^*$ 的一组基.
			\item\label{ex:bil:A5} 设线性映射 $\mathscr T:\mathbb{R}^2\to\mathbb{R}^3$ 在标准基下的矩阵为 $\begin{pmatrix}1&2\\0&1\\3&-1\end{pmatrix}$. 求对偶映射 $\mathscr T^*$ 在对偶基下的矩阵，并计算 $\mathscr T^*$ 作用在线性函数 $\vv y\mapsto y_1+y_2+y_3$ 上的结果.
			\item\label{ex:bil:A6} 求双线性型 $f(\vv x,\vv y)=x_1y_1+2x_1y_2-x_2y_1+3x_2y_2$ 在标准基下的度量矩阵，以及在基 $(1,1)^\trans$，$(1,-1)^\trans$ 下的度量矩阵，并验证 $\vv B=\vv P^\trans\vv{AP}$.
			\item\label{ex:bil:A7} 把 A6 中的 $f$ 分解为对称部分与交错部分之和.
			\item\label{ex:bil:A8} 判断下列双线性型是否非退化：(1) $x_1y_1-x_2y_2$；(2) $x_1y_2+x_2y_1$；(3) $(x_1+x_2)(y_1+y_2)$.
			\item\label{ex:bil:A9} 写出对称双线性型 $f(\vv x,\vv y)=x_1y_1+x_1y_2+x_2y_1+3x_2y_2$ 对应的二次型；反过来，由二次型 $x_1^2-4x_1x_2+2x_2^2$ 用极化恒等式求对应的对称双线性型.
			\item\label{ex:bil:A10} 验证 $\mathbb{R}^2$ 上的 $f(\vv x,\vv y)=\det(\vv x,\vv y)$ 是交错双线性型，写出它的度量矩阵. 它非退化吗？
			\item\label{ex:bil:A11} 设 $a_1,a_2,a_3$ 互异，$f_i(p)=p(a_i)$ 是 $\mathbb{R}[x]_3$ 上的线性函数. 对 $a_i=0,1,2$，求以 $f_1,f_2,f_3$ 为对偶基的 $\mathbb{R}[x]_3$ 的基.
			\item\label{ex:bil:A12} 在 $\mathbb{R}^4$ 中，设 $W$ 由 $x_1+x_2=0$，$x_3-x_4=0$ 定义. 求 $\dim W$ 与 $\dim W^0$，并写出 $W^0$ 的一组基.
			\item\label{ex:bil:A13} 证明 $\langle\vv A,\vv B\rangle=\tr(\vv{AB})$ 是 $M_n(\mathbb{R})$ 上非退化的对称双线性型. 它是内积吗？
			\item\label{ex:bil:A14} 对 $\mathbb{R}[x]_3$ 上的线性函数 $\varphi(p)=\int_0^1p(x)\,\mathrm dx$，求它在 $1,x,x^2$ 的对偶基下的坐标.
			\item\label{ex:bil:A15} （Riesz 表示定理的有限维情形）设 $V$ 是有限维实内积空间. 证明：对每个 $f\in V^*$，存在唯一的 $\vv\beta\in V$ 使 $f(\vv\alpha)=\langle\vv\alpha,\vv\beta\rangle$ 对一切 $\vv\alpha$ 成立.（提示：取标准正交基 $\vv e_1,\ldots,\vv e_n$，令 $\vv\beta=\sum f(\vv e_i)\vv e_i$.）
		\end{exA}
		
		\Bgroup
		\begin{exB}
			\item\label{ex:bil:B1} 证明：$f_1,\ldots,f_k\in V^*$ 线性无关，当且仅当映射 $V\to\mathbb{F}^k$，$\vv\alpha\mapsto(f_1(\vv\alpha),\ldots,f_k(\vv\alpha))^\trans$ 是满射.
			\item\label{ex:bil:B2} 设 $f,g\in V^*$ 且 $\ker f\subset\ker g$. 证明存在常数 $c$ 使 $g=cf$. 推广：若 $\ker f_1\cap\cdots\cap\ker f_k\subset\ker g$，则 $g$ 是 $f_1,\ldots,f_k$ 的线性组合.
			\item\label{ex:bil:B3} 证明：$\ker f_1\cap\cdots\cap\ker f_k=\{\zero\}$ 当且仅当 $f_1,\ldots,f_k$ 张成 $V^*$. 用方程组的语言解释这个结论.
			\item\label{ex:bil:B4} 证明：$\mathscr T$ 是单射当且仅当 $\mathscr T^*$ 是满射；$\mathscr T$ 是满射当且仅当 $\mathscr T^*$ 是单射.
			\item\label{ex:bil:B5} （Simpson 公式）求常数 $a,b,c$，使 $\int_0^2p(x)\,\mathrm dx=a\,p(0)+b\,p(1)+c\,p(2)$ 对一切 $p\in\mathbb{R}[x]_3$ 成立. 用对偶基的语言说明为什么这样的 $a,b,c$ 存在且唯一. 验证这个公式对三次多项式也成立.
			\item\label{ex:bil:B6} 证明：$V$ 上全体双线性型构成的空间与 $M_n(\mathbb{F})$ 同构，其中对称型、交错型分别构成维数为 $\frac{n(n+1)}2$、$\frac{n(n-1)}2$ 的子空间，且整个空间是二者的直和.
			\item\label{ex:bil:B7} 证明：若 $V$ 上存在非退化的交错双线性型，则 $\dim V$ 是偶数.
			\item\label{ex:bil:B8} 求双线性型 $f(\vv x,\vv y)=x_1y_1+x_1y_2+x_2y_1+x_2y_2+x_3y_3$ 的``根'' $\{\vv\alpha:f(\vv\alpha,\cdot)=0\}$.
			\item\label{ex:bil:B9} 在 $\mathbb{R}^2$ 上取 $f(\vv x,\vv y)=x_1y_1-x_2y_2$. 求全部``迷向向量''（$f(\vv\alpha,\vv\alpha)=0$），并说明它们构成两条直线. 再证明保持 $f$ 的线性变换 $\vv P$（$\vv P^\trans\diag(1,-1)\vv P=\diag(1,-1)$，$\det\vv P=1$，且保持 $x_1>|x_2|$ 的区域）都形如 $\begin{pmatrix}\cosh t&\sinh t\\ \sinh t&\cosh t\end{pmatrix}$.
			\item\label{ex:bil:B10} 设 $W$ 是 $V$ 的子空间，$\mathscr T:V\to V$ 线性且 $\mathscr T(W)\subset W$. 证明 $\mathscr T^*(W^0)\subset W^0$.
		\end{exB}
		
		\Cgroup
		\begin{exC}
			\item\label{ex:bil:C1} 证明：对无限维空间 $\mathbb{F}[x]$，自然映射 $\vv\alpha\mapsto\hat{\vv\alpha}$ 是单射但不是满射.（提示：$\mathbb{F}[x]^*$ 中的元素由它在 $1,x,x^2,\ldots$ 上的取值任意决定，可以``无限多个非零''.）
			\item\label{ex:bil:C2} 证明：$V$ 上的双线性型与线性映射 $V\to V^*$ 一一对应：$f\mapsto(\vv\alpha\mapsto f(\vv\alpha,\cdot))$. 用这个对应和对偶基的换基规则，重新推出度量矩阵的换基公式 $\vv B=\vv P^\trans\vv{AP}$.
			\item\label{ex:bil:C3} 证明：二阶矩阵 $\vv P$ 满足 $\vv P^\trans\vv J\vv P=\vv J$（$\vv J=\begin{pmatrix}0&1\\-1&0\end{pmatrix}$）当且仅当 $\det\vv P=1$. 四阶时，满足 $\vv P^\trans\diag(\vv J,\vv J)\vv P=\diag(\vv J,\vv J)$ 的矩阵（辛矩阵）的行列式是多少？
			\item\label{ex:bil:C4} 求 $M_n(\mathbb{R})$ 上对称双线性型 $\tr(\vv{AB})$ 的正、负惯性指数.（提示：对称矩阵与反对称矩阵在这个型下``互相正交''.）
			\item\label{ex:bil:C5} 设 $f$ 是 $V$ 上的对称双线性型，$W$ 是子空间，$W^\perp=\{\vv\alpha:f(\vv\alpha,\vv\beta)=0\ \forall\vv\beta\in W\}$. 证明 $\dim W^\perp\ge\dim V-\dim W$，$f$ 非退化时等号成立；举例说明即使 $f$ 非退化，$W\cap W^\perp$ 也可能不为零.
			\item\label{ex:bil:C6} 证明：$\mathbb{F}^n$ 上全体交错 $n$ 重线性函数构成的空间是一维的，由行列式张成. 这给出第五章章末习题~\ref{ex:det:C1} 的``对偶空间版本''.
			\item\label{ex:bil:C7} （线性规划对偶）考虑问题：在 $\vv{Ax}=\vv b$，$\vv x\ge\zero$ 下最小化 $\vv c^\trans\vv x$. 证明对任意满足 $\vv A^\trans\vv y\le\vv c$ 的 $\vv y$ 与任意可行的 $\vv x$，都有 $\vv b^\trans\vv y\le\vv c^\trans\vv x$（弱对偶）. 用``价格是数量空间上的线性函数''解释 $\vv y$ 的经济含义（影子价格）.
		\end{exC}
	}.
	%%   2. 在每章末尾（原 \begin{problemset}...\end{problemset} 之后，或替换它）
	%%      写对应的命令：
	%%        第一章 \ExElimination      第二章 \ExVectorSpace   第三章 \ExMatrix
	%%        第四章 \ExLinearMap        第五章 \ExDeterminant   第六章 \ExSolutions
	%%        第七章 \ExPolynomial       第八章 \ExEigen         第九章 \ExInnerProduct
	%%        第十章 \ExDecomposition    第十一章 \ExQuadratic   第十二章 \ExBilinear
	%%   3. 分组标准
	%%        A 组：直接运用本章定义与定理的计算和证明；
	%%        B 组：需要综合多个结论、或有一定技巧的计算和证明；
	%%        C 组：有挑战性的题目，或超出本课程大纲、为继续深造准备的内容.
	%%   4. 题目与正文中已有的节末练习、原 problemset 不重复；个别题目注明``接正文某题''.
	%%      第六章按扩充版（ch06-solutions.tex）的内容编写.
	%%   5. 每道题都带标签 ex:<章键>:<组号>，正文、习题与附录答案一律用 \ref 引用，
	%%      渲染为``A4''``B12''这样的题号. 章键依次为
	%%        elim vs mat lmap det sol poly eig ip dec quad bil（第一至十二章）.
	%%      例：第一章章末习题~\ref{ex:elim:B12}. 新增题目请加在组末，并补上标签.
	%%      A 组答案、B 组提示在 answers.tex 中，按同一标签对应.
	%% 依赖：导言区的 \vv \zero \trans \spn \rk \imm \diag，以及 enumitem（elegantbook 已加载）.
	%% 引用的前文标签：thm:sameSolution thm:solvability prop:linspc eg:pricedecomp cor:homogeneous thm:rouche thm:sol-determinable
	%%   thm:spectral thm:posdef thm:svd thm:orthcanon
	%% =====================================================================
	
	\providecommand{\tr}{\operatorname{tr}}
	\newcommand{\exsetheading}{%
		\section*{章末习题}%
		\markright{章末习题}%
	}
\newcommand{\exgroup}[2]{%
	\par\medskip{\centering\sffamily\bfseries #1\quad #2\par}\smallskip}
	\newenvironment{exA}{\begin{enumerate}[label=\textbf{A\arabic*.},ref=A\arabic*,leftmargin=2.8em,itemsep=4pt]}{\end{enumerate}}
	\newenvironment{exB}{\begin{enumerate}[label=\textbf{B\arabic*.},ref=B\arabic*,leftmargin=2.8em,itemsep=4pt]}{\end{enumerate}}
	\newenvironment{exC}{\begin{enumerate}[label=\textbf{C\arabic*.},ref=C\arabic*,leftmargin=2.8em,itemsep=4pt]}{\end{enumerate}}
	\newcommand{\Agroup}{\exgroup{A 组}{基础巩固}}
	\newcommand{\Bgroup}{\exgroup{B 组}{综合提升}}
	\newcommand{\Cgroup}{\exgroup{C 组}{拓展探索}}
	
	% =====================================================================
	% 第一章　线性方程组与消元法
	% =====================================================================
	\newcommand{\ExElimination}{%
		\exsetheading
		\Agroup
		\begin{exA}
			\item\label{ex:elim:A1} 判断下列矩阵哪些是阶梯形，哪些是简化阶梯形：
			\[
			\begin{pmatrix}1&2&0\\0&0&1\end{pmatrix},\quad
			\begin{pmatrix}1&0&3\\0&2&1\\0&0&0\end{pmatrix},\quad
			\begin{pmatrix}0&1&0\\1&0&0\end{pmatrix},\quad
			\begin{pmatrix}1&0&2&0\\0&1&-1&0\\0&0&0&1\end{pmatrix},\quad
			\begin{pmatrix}0&0\\0&1\end{pmatrix}.
			\]
			\item\label{ex:elim:A2} 写出方程组 $3x_1-x_3=2$，$x_2+4x_3=0$，$x_1-x_2=5$ 的增广矩阵；反过来，写出以 $\begin{pmatrix}1&-2&0&3\\0&1&5&-1\end{pmatrix}$ 为增广矩阵的方程组.
			\item\label{ex:elim:A3} 用消元法解方程组 $\left\{\begin{array}{l}x_1+2x_2+3x_3=14,\\2x_1-x_2+x_3=3,\\3x_1+x_2-x_3=2.\end{array}\right.$
			\item\label{ex:elim:A4} 用消元法解方程组 $\left\{\begin{array}{l}x_1+x_2+x_3+x_4=2,\\x_1-x_2+x_3-x_4=4,\\2x_1+x_2-x_3+3x_4=-1,\\x_1+2x_2+3x_3-2x_4=5.\end{array}\right.$
			\item\label{ex:elim:A5} 判断方程组 $\left\{\begin{array}{l}x_1-x_2+x_3=1,\\2x_1+x_2-x_3=2,\\x_1+2x_2-2x_3=3\end{array}\right.$ 是否有解，并说明消元过程中哪一行``暴露''了结论.
			\item\label{ex:elim:A6} 求方程组 $\left\{\begin{array}{l}x_1+2x_2-x_3+x_4=3,\\2x_1+4x_2-x_3+3x_4=7,\\-x_1-2x_2+2x_3=-2\end{array}\right.$ 的全部解，并指出主变量与自由变量.
			\item\label{ex:elim:A7} 求齐次方程组 $\left\{\begin{array}{l}x_1+x_2-x_3-x_4=0,\\2x_1-5x_2+3x_3+2x_4=0,\\7x_1-7x_2+3x_3+x_4=0\end{array}\right.$ 的全部解.
			\item\label{ex:elim:A8} 把矩阵 $\begin{pmatrix}2&4&-2&6\\1&2&1&1\\3&6&1&5\end{pmatrix}$ 化为简化阶梯形，并写出每一步所用的初等行变换.
			\item\label{ex:elim:A9} 写出全部 $2\times2$ 简化阶梯形矩阵的``形状''（可以含参数），并说明共有几类.
			\item\label{ex:elim:A10} $k$ 取何值时，方程组 $x+y=1$，$x-y=3$，$2x+ky=5$ 有解？
			\item\label{ex:elim:A11} 设 $a,b$ 为常数. 方程组 $x+2y=a$，$3x+6y=b$ 何时有解？有解时写出全部解.
			\item\label{ex:elim:A12} 求过三点 $(1,2)$，$(2,3)$，$(3,6)$ 的抛物线 $y=ax^2+bx+c$.
			\item\label{ex:elim:A13} 求过四点 $(-1,2)$，$(0,1)$，$(1,0)$，$(2,5)$ 的三次多项式函数 $y=a+bx+cx^2+dx^3$.
			\item\label{ex:elim:A14} 某剧场三种票价分别为 $30$，$50$，$80$ 元. 某场演出共售出 $100$ 张票，票房 $4800$ 元，且 $80$ 元票的张数恰为 $30$ 元票的一半. 三种票各售出多少张？
			\item\label{ex:elim:A15} 配平化学方程式 $x_1\,\mathrm{C_3H_8}+x_2\,\mathrm{O_2}\to x_3\,\mathrm{CO_2}+x_4\,\mathrm{H_2O}$：对 C，H，O 三种原子分别写出守恒方程，解这个齐次方程组，并取最小的正整数解.
			\item\label{ex:elim:A16} 某城区四个路口 $A,B,C,D$ 由单行道依次连成环路 $A\to B\to C\to D\to A$，环路上的车流量依次记为 $x_1$（$A\to B$），$x_2$（$B\to C$），$x_3$（$C\to D$），$x_4$（$D\to A$）（单位：辆/小时）. 各路口与外部道路的进出量为：$A$ 驶入 $400$、驶出 $300$；$B$ 驶入 $200$、驶出 $350$；$C$ 驶入 $250$、驶出 $100$；$D$ 驶入 $150$、驶出 $250$. 设每个路口驶入总量等于驶出总量.
			\begin{enumerate}
				\item 列出方程组并求全部解；
				\item 若各路段车流量都不能为负，$x_1$ 最少是多少？
				\item 若把 $D$ 的驶出量改为 $350$，方程组会怎样？这反映了什么守恒律？
			\end{enumerate}
			\item\label{ex:elim:A17} 证明：二元一次方程组 $ax+by=e$，$cx+dy=f$ 对任意的 $e,f$ 都有唯一解，当且仅当 $ad-bc\ne0$.（只用消元法，不用行列式.）
			\item\label{ex:elim:A18} 证明：若线性方程组有两个不同的解 $\vv x_1,\vv x_2$，则对任意数 $t$，$\vv x_1+t(\vv x_2-\vv x_1)$ 都是解. 由此说明：线性方程组的解要么不存在，要么唯一，要么有无穷多个.
			\item\label{ex:elim:A19} 证明方程组 $\{x+y+z=3,\ x-y=0\}$ 与 $\{2x+z=3,\ x+y+z=3\}$ 同解.
			\item\label{ex:elim:A20} 举例说明：若\textbf{同时}做``第一个方程加到第二个方程上''与``第二个方程加到第一个方程上''（两步都用原来的方程），所得方程组可能与原方程组不同解. 这与定理~\ref{thm:sameSolution}（初等行变换不改变解集）矛盾吗？
			\item\label{ex:elim:A21} 已知一个方程组的增广矩阵化为阶梯形后为 $\begin{pmatrix}1&3&-1&2&0\\0&0&2&1&4\\0&0&0&0&a\end{pmatrix}$. 讨论 $a$ 取何值时方程组有解；有解时它有几个自由变量？
		\end{exA}
		
		\Bgroup
		\begin{exB}
			\item\label{ex:elim:B1} 讨论 $a,b$ 取何值时，方程组
			\[\left\{\begin{array}{l}x_1+x_2+x_3+x_4=0,\\x_2+2x_3+2x_4=1,\\-x_2+(a-3)x_3-2x_4=b,\\3x_1+2x_2+x_3+ax_4=-1\end{array}\right.\]
			有唯一解、无解、有无穷多解；有无穷多解时求出全部解.
			\item\label{ex:elim:B2} 讨论 $\lambda$ 取何值时，方程组
			$(1+\lambda)x_1+x_2+x_3=0$，$x_1+(1+\lambda)x_2+x_3=3$，$x_1+x_2+(1+\lambda)x_3=\lambda$
			有唯一解、无解、有无穷多解.
			\item\label{ex:elim:B3} 设 $n$ 个方程、$n$ 个未知量的方程组对应的齐次方程组只有零解. 证明：无论常数项取什么值，原方程组都有唯一解.（提示：考察简化阶梯形中主元的个数.）
			\item\label{ex:elim:B4} 证明：线性方程组无解，当且仅当可以把若干个方程分别乘以适当的数再相加，得到矛盾方程 $0=1$.
			\item\label{ex:elim:B5} 证明：交换两行这一初等变换，可以由``倍加''与``某行乘以 $-1$''这两类变换的组合实现. 由此说明三类初等行变换并不是``独立''的.
			\item\label{ex:elim:B6} 证明：若线性方程组的系数和常数项都是有理数，且它有实数解，则它一定有有理数解；若它的实数解唯一，则这个解就是有理数解.
			\item\label{ex:elim:B7} 讨论方程组
			\[x_1-x_2=a_1,\quad x_2-x_3=a_2,\quad\ldots,\quad x_{n-1}-x_n=a_{n-1},\quad x_n-x_1=a_n\]
			有解的充要条件，并在有解时求全部解.
			\item\label{ex:elim:B8} 讨论方程组
			\[x_1+x_2=a_1,\quad x_2+x_3=a_2,\quad\ldots,\quad x_{n-1}+x_n=a_{n-1},\quad x_n+x_1=a_n\]
			的解. 证明：$n$ 为奇数时解唯一；$n$ 为偶数时有解当且仅当 $a_1-a_2+a_3-\cdots-a_n=0$.
			\item\label{ex:elim:B9} 某经济体有三个部门，直接消耗系数矩阵与最终需求分别为
			\[\vv A=\begin{pmatrix}0.1&0.2&0.2\\0.2&0.1&0.2\\0.1&0.3&0.1\end{pmatrix},\qquad \vv d=\begin{pmatrix}20\\130\\65\end{pmatrix}.\]
			按第一章``投入产出模型''的方式列出方程组，求各部门的总产出.
			\item\label{ex:elim:B10} 估计用消元法（化为阶梯形后回代）解 $n$ 个方程、$n$ 个未知量的方程组所需的乘除法次数，证明它约为 $\frac13n^3$（即与 $\frac13n^3$ 之比在 $n\to\infty$ 时趋于 $1$）.
			\item\label{ex:elim:B11} 设三个平面 $a_ix+b_iy+c_iz=d_i$（$i=1,2,3$）两两相交. 用消元法的语言说明：它们恰好交于一点、交于一条直线、或没有公共点（三条交线互相平行），并各举一例.
			\item\label{ex:elim:B12} 证明：$n$ 元齐次方程组的解集对加法和数乘封闭；而非齐次方程组的两个解之差是对应齐次方程组的解. 由此说明：非齐次方程组的全部解，等于它的任一个解加上齐次方程组的全部解.
			\item\label{ex:elim:B13} 一家律所有甲、乙、丙三位律师，本月分别承办了若干件合同、劳动、知识产权纠纷. 已知：三类案件合计 $30$ 件；甲承办的数量 $=$ 乙 $+$ 丙；每件合同、劳动、知识产权案件分别需要 $10$，$6$，$15$ 个工时，三人本月共用 $300$ 个工时……
			这些条件够不够确定每位律师承办了多少案件？请自己补全一个合理的设定（明确未知量是什么），列出方程组，用自由变量的个数说明还缺多少条独立的信息.
			\item\label{ex:elim:B14} 设方程组的系数矩阵每一行的元素之和都等于 $0$. 证明 $x_1=x_2=\cdots=x_n=1$ 是对应齐次方程组的解；若非齐次方程组有解 $\vv x_0$，则对任意数 $t$，$\vv x_0+t(1,\ldots,1)^\trans$ 也是解.
			\item\label{ex:elim:B15} 证明：如果方程组 (I) 的每个方程都是方程组 (II) 各方程的``线性组合''（各乘一个数再相加），反之亦然，那么 (I) 与 (II) 同解. 逆命题对于有解的方程组成立吗？先在两个未知量的情形下思考.
			\item\label{ex:elim:B16} 讨论 $\lambda$ 取何值时，方程组 $\lambda x+y+z=1$，$x+\lambda y+z=\lambda$，$x+y+\lambda z=\lambda^2$ 有唯一解、无解、有无穷多解；有唯一解时求出解.（第六章例~\ref{eg:lambda} 将用秩与行列式重新组织这个讨论.）
		\end{exB}
		
		\Cgroup
		\begin{exC}
			\item\label{ex:elim:C1} （百鸡问题）《张丘建算经》（约公元 5 世纪）：``今有鸡翁一，直钱五；鸡母一，直钱三；鸡雏三，直钱一. 凡百钱买鸡百只，问鸡翁、母、雏各几何？''列出线性方程组，求它的全部解，再从中找出所有\textbf{非负整数}解. 线性代数在哪一步之后就无能为力了？
			\item\label{ex:elim:C2} （幻方）若一个 $3\times3$ 数表的每行、每列和两条对角线上的三个数之和都等于同一个数 $s$，称它为幻方（不要求数互不相同）. 把 $9$ 个元素作为未知量、$s$ 也作为未知量，写出方程组，证明：
			\begin{enumerate}
				\item 中心元素一定等于 $s/3$；
				\item 全部幻方可以用 $3$ 个自由参数表示，写出通式；
				\item 用 $1,2,\ldots,9$ 各一次填成的幻方本质上只有一个（其余由旋转、翻转得到）.
			\end{enumerate}
			\item\label{ex:elim:C3} （点灯游戏）$3\times3$ 方格中每格有一盏灯，按下某格会使该格及其上下左右相邻的格子开关状态翻转. 在``$1+1=0$''的算术（即模 $2$ 运算）中，把``按哪些格''作为未知量，把``从某一状态变为全灭''写成线性方程组. 证明：从任意初始状态出发，总能把灯全部熄灭.（模 $2$ 运算的集合 $\{0,1\}$ 不是本书意义下的数域，但消元法照样适用——想一想为什么.）
			\item\label{ex:elim:C4} （舍入误差与主元）在``只保留三位有效数字''的计算机上解方程组 $0.0001x+y=1$，$x+y=2$：
			\begin{enumerate}
				\item 直接用第一个方程消去第二个方程中的 $x$；
				\item 先交换两个方程再消元.
			\end{enumerate}
			把两种结果与精确解比较，说明数值软件为什么要``选绝对值最大的元素作主元''.
			\item\label{ex:elim:C5} 设两个都有解的方程组 (I)(II) 都有 $m$ 个方程、$n$ 个未知量，且同解. 证明：它们的增广矩阵行等价.（提示：先证明二者的简化阶梯形的非零行完全相同.）
			\item\label{ex:elim:C6} （Jacobi 迭代）对方程组 $4x+y=9$，$x+3y=5$，把它改写为 $x=(9-y)/4$，$y=(5-x)/3$，从 $(x_0,y_0)=(0,0)$ 出发反复代入. 计算前五步，观察是否收敛到真解. 再对 $x+4y=9$，$3x+y=5$ 做同样的事，会发生什么？试找出一个保证收敛的充分条件.
			\item\label{ex:elim:C7} （电路）一个电路中有三条支路电流 $I_1,I_2,I_3$. 节点 $A$ 处 $I_1=I_2+I_3$；两个回路的电压方程为 $2I_1+I_2=8$，$I_2-2I_3=0$（单位从略）. 求各电流. 再说明：若把第三个节点 $B$ 的电流方程也写进来，方程组为什么不会因此``多一个条件''？
			\item\label{ex:elim:C8} 证明：一个 $m\times n$ 矩阵（$m,n\ge1$）通过初等行变换所能得到的不同的简化阶梯形``形状''（即主元所在列的位置组合）共有 $\sum_{r=0}^{\min(m,n)}\binom nr$ 种.
			\item\label{ex:elim:C9} 证明：一个矩阵的简化阶梯形是唯一的.（提示：设 $\vv R_1,\vv R_2$ 都是 $\vv M$ 的简化阶梯形，则以它们为系数矩阵的齐次方程组同解. 对列数归纳：删去最后一列后由归纳假设前面各列相同，再比较最后一列，考察它在两者中是否为主元列.）
		\end{exC}
	}
	
	% =====================================================================
	% 第二章　向量空间
	% =====================================================================
	\newcommand{\ExVectorSpace}{%
		\exsetheading
		\Agroup
		\begin{exA}
			\item\label{ex:vs:A1} 下列集合关于通常的四则运算是否构成数域？
			(1) $\{a+b\sqrt3\mid a,b\in\mathbb{Q}\}$；(2) $\{a+b\sqrt2\mid a\in\mathbb{Q},\ b\in\mathbb{Z}\}$；(3) 全体偶数；(4) $\{a+b\sqrt[3]{2}\mid a,b\in\mathbb{Q}\}$；(5) 全体纯虚数与 $0$.
			\item\label{ex:vs:A2} 下列集合关于通常的加法与数乘是否构成实向量空间？
			(1) 全体 $n$ 阶实对称矩阵；(2) $\{(x,y)^\trans\in\mathbb{R}^2\mid x\ge0\}$；(3) 次数恰为 $n$ 的实系数多项式全体；(4) $\{f\in\mathbb{R}[x]\mid f(1)=0\}$；(5) 满足 $a_{n+2}=a_{n+1}+a_n$ 的全体实数列；(6) $\{(x,y,z)^\trans\mid x+y+z=1\}$.
			\item\label{ex:vs:A3} 判断下列向量组的线性相关性：
			\begin{enumerate}
				\item $(1,2,-1)^\trans$，$(2,-3,1)^\trans$，$(4,1,-1)^\trans$；
				\item 在 $\mathbb{R}[x]_3$ 中：$1+x$，$1-x$，$x^2$；
				\item 在 $\mathbb{R}[x]_3$ 中：$1+x+x^2$，$1-x$，$2+x^2$，$x+x^2$；
				\item 在 $M_2(\mathbb{R})$ 中：$\begin{pmatrix}1&0\\0&1\end{pmatrix}$，$\begin{pmatrix}0&1\\1&0\end{pmatrix}$，$\begin{pmatrix}1&1\\1&1\end{pmatrix}$，$\begin{pmatrix}1&0\\0&0\end{pmatrix}$.
			\end{enumerate}
			\item\label{ex:vs:A4} 把 $\vv\beta=(1,2,1,1)^\trans$ 表示为 $\vv\alpha_1=(1,1,1,1)^\trans$，$\vv\alpha_2=(1,1,-1,-1)^\trans$，$\vv\alpha_3=(1,-1,1,-1)^\trans$，$\vv\alpha_4=(1,-1,-1,1)^\trans$ 的线性组合.
			\item\label{ex:vs:A5} $t$ 取何值时，$(1,1,0)^\trans$，$(1,3,-1)^\trans$，$(5,3,t)^\trans$ 线性相关？
			\item\label{ex:vs:A6} 求向量组 $\vv\alpha_1=(1,-1,2,4)^\trans$，$\vv\alpha_2=(0,3,1,2)^\trans$，$\vv\alpha_3=(3,0,7,14)^\trans$，$\vv\alpha_4=(1,-1,2,0)^\trans$，$\vv\alpha_5=(2,1,5,6)^\trans$ 的秩和一个极大无关组，并把其余向量用这个极大无关组线性表出.
			\item\label{ex:vs:A7} 设 $\vv\alpha_1,\vv\alpha_2,\vv\alpha_3$ 线性无关. 判断下列向量组的线性相关性：
			(1) $\vv\alpha_1-\vv\alpha_2$，$\vv\alpha_2-\vv\alpha_3$，$\vv\alpha_3-\vv\alpha_1$；
			(2) $\vv\alpha_1$，$\vv\alpha_1+\vv\alpha_2$，$\vv\alpha_1+\vv\alpha_2+\vv\alpha_3$；
			(3) $\vv\alpha_1+2\vv\alpha_2$，$\vv\alpha_2+2\vv\alpha_3$，$\vv\alpha_3+2\vv\alpha_1$.
			\item\label{ex:vs:A8} 证明：含有零向量的向量组线性相关；含有两个相同向量的向量组线性相关；单个向量 $\vv\alpha$ 线性无关当且仅当 $\vv\alpha\ne\zero$.
			\item\label{ex:vs:A9} 证明：两个向量线性相关，当且仅当其中一个是另一个的倍数. 三个向量的情形能这样说吗？
			\item\label{ex:vs:A10} 只用向量空间的八条公理与命题~\ref{prop:linspc}，证明：(1) $\vv\alpha+\vv\alpha=2\vv\alpha$；(2) 若 $k\ne0$ 且 $k\vv\alpha=k\vv\beta$，则 $\vv\alpha=\vv\beta$；(3) $-(-\vv\alpha)=\vv\alpha$.
			\item\label{ex:vs:A11} 证明下列集合是相应空间的子空间，并求一组基和维数：
			\begin{enumerate}
				\item $\{\vv x\in\mathbb{R}^4\mid x_1+x_2=x_3+x_4\}$；
				\item $n$ 阶上三角矩阵全体；
				\item 迹为零的 $n$ 阶方阵全体 $\{\vv X\mid x_{11}+\cdots+x_{nn}=0\}$；
				\item $\{f\in\mathbb{R}[x]_4\mid f(1)=0\}$.
			\end{enumerate}
			\item\label{ex:vs:A12} 判断下列子集是否为子空间：(1) $\{(a,b,c)^\trans\mid a=2b\}$；(2) $\{(a,b,c)^\trans\mid ab=0\}$；(3) $\{(a,b,c)^\trans\mid a,b,c\in\mathbb{Q}\}\subset\mathbb{R}^3$；(4) $\{f\in C[0,1]\mid f(0)=2f(1)\}$.
			\item\label{ex:vs:A13} 求 $\vv\alpha_1=(1,2,1,0)^\trans$，$\vv\alpha_2=(-1,1,1,1)^\trans$，$\vv\alpha_3=(0,3,2,1)^\trans$，$\vv\alpha_4=(1,0,1,2)^\trans$ 张成的子空间的维数和一组基.
			\item\label{ex:vs:A14} 求多项式 $2x^2-x+3$ 在基 $1,\ x-1,\ (x-1)^2$ 下的坐标.
			\item\label{ex:vs:A15} 求矩阵 $\begin{pmatrix}2&3\\-1&4\end{pmatrix}$ 在 $M_2(\mathbb{R})$ 的基 $\begin{pmatrix}1&0\\0&0\end{pmatrix},\begin{pmatrix}1&1\\0&0\end{pmatrix},\begin{pmatrix}1&1\\1&0\end{pmatrix},\begin{pmatrix}1&1\\1&1\end{pmatrix}$ 下的坐标.
			\item\label{ex:vs:A16} (1) $\mathbb{R}$ 是否是 $\mathbb{Q}$ 上的向量空间？(2) $\mathbb{Q}$ 是否是 $\mathbb{R}$ 上的向量空间？(3) $\mathbb{C}$ 作为 $\mathbb{R}$ 上的向量空间和作为 $\mathbb{C}$ 上的向量空间，维数各是多少？
			\item\label{ex:vs:A17} 写出 $\mathbb{F}[x]_n$ 与 $\mathbb{F}^n$ 之间、$M_2(\mathbb{F})$ 与 $\mathbb{F}^4$ 之间的一个同构映射，并验证它保持线性运算.
			\item\label{ex:vs:A18} 设 $W_1=\spn\big((1,0,0)^\trans,(0,1,0)^\trans\big)$，$W_2=\spn\big((1,1,1)^\trans\big)$，$W_3=\spn\big((1,1,0)^\trans\big)$. 分别判断 $W_1+W_2$，$W_1+W_3$ 是否为直和，并求和空间的维数.
			\item\label{ex:vs:A19} 设 $W_1=\spn\big((1,1,0,0)^\trans,(1,0,1,1)^\trans\big)$，$W_2=\spn\big((0,-1,1,1)^\trans,(1,2,0,1)^\trans\big)$. 求 $W_1+W_2$ 与 $W_1\cap W_2$ 的维数，并验证维数公式.
			\item\label{ex:vs:A20} 设 $\vv\alpha_1,\ldots,\vv\alpha_s$ 的秩为 $r$. 证明：其中任意 $r+1$ 个向量线性相关；任意 $r$ 个线性无关的向量都构成极大无关组.
		\end{exA}
		
		\Bgroup
		\begin{exB}
			\item\label{ex:vs:B1} 证明：两个子空间的并 $W_1\cup W_2$ 是子空间，当且仅当 $W_1\subset W_2$ 或 $W_2\subset W_1$.
			\item\label{ex:vs:B2} 证明：若向量组 $\vv\beta_1,\ldots,\vv\beta_t$ 可由 $\vv\alpha_1,\ldots,\vv\alpha_s$ 线性表出，且两组的秩相等，则两组等价.
			\item\label{ex:vs:B3} 设向量组 (I) 与 (II) 的秩分别为 $r_1,r_2$，把它们合并为 (III). 证明 $\max\{r_1,r_2\}\le\rk(\mathrm{III})\le r_1+r_2$. 再证明：$\rk(\vv\alpha_1+\vv\beta_1,\ldots,\vv\alpha_s+\vv\beta_s)\le\rk(\vv\alpha_1,\ldots,\vv\alpha_s)+\rk(\vv\beta_1,\ldots,\vv\beta_s)$.
			\item\label{ex:vs:B4} 设 $\vv\alpha_1,\ldots,\vv\alpha_n$ 是 $V$ 的基，$\vv\beta=a_1\vv\alpha_1+\cdots+a_n\vv\alpha_n$. 证明：把 $\vv\alpha_j$ 换成 $\vv\beta$ 后所得的向量组仍是基，当且仅当 $a_j\ne0$.
			\item\label{ex:vs:B5} 设 $p_0,p_1,\ldots,p_{n-1}\in\mathbb{F}[x]$ 满足 $\deg p_k=k$. 证明它们构成 $\mathbb{F}[x]_n$ 的一组基.
			\item\label{ex:vs:B6} 记 $V$ 为满足 $a_{n+2}=a_{n+1}+2a_n$（$n\ge1$）的全体实数列. 证明 $V$ 是向量空间，$\dim V=2$，并说明 $\{2^n\}$ 与 $\{(-1)^n\}$ 构成 $V$ 的一组基. 由此写出 $a_1=1$，$a_2=1$ 时的通项.
			\item\label{ex:vs:B7} 在 $C(\mathbb{R})$ 中：(1) 证明 $\mathrm e^x,\mathrm e^{2x},\mathrm e^{3x}$ 线性无关；(2) 证明 $\sin x,\cos x,\sin(x+1)$ 线性相关；(3) $1,\cos x,\cos2x,\cos^2x$ 的秩是多少？
			\item\label{ex:vs:B8} 证明 $\{a+b\sqrt2+c\sqrt3+d\sqrt6\mid a,b,c,d\in\mathbb{Q}\}$ 是数域，并证明 $1,\sqrt2,\sqrt3,\sqrt6$ 在 $\mathbb{Q}$ 上线性无关. 这个数域作为 $\mathbb{Q}$ 上的向量空间是几维的？
			\item\label{ex:vs:B9} 写出 $\mathbb{R}^2$ 与 $\mathbb{R}^3$ 的全部子空间（按维数分类，并说明几何意义）.
			\item\label{ex:vs:B10} 证明：$\mathbb{R}^5$ 中任意两个三维子空间的交至少是一维的. 一般地，$n$ 维空间中维数分别为 $k_1,k_2$ 的两个子空间，若 $k_1+k_2>n$，则交不为 $\{\zero\}$.
			\item\label{ex:vs:B11} 设 $a\in\mathbb{F}$. 证明 $\mathbb{F}[x]_n=\mathbb{F}\cdot1\oplus\{f\in\mathbb{F}[x]_n\mid f(a)=0\}$，并写出 $f$ 按这个直和的分解.
			\item\label{ex:vs:B12} 设 $V=C[-1,1]$，$E$，$O$ 分别为其中的偶函数、奇函数全体. 证明 $V=E\oplus O$，并把 $\mathrm e^x$ 按此分解.
			\item\label{ex:vs:B13} 设 $\dim V=n$，$W_1,\ldots,W_k$ 是 $V$ 的子空间，且 $\dim W_1+\cdots+\dim W_k>(k-1)n$. 证明 $W_1\cap\cdots\cap W_k\ne\{\zero\}$.
			\item\label{ex:vs:B14} 设 $\vv\alpha_1,\ldots,\vv\alpha_n$ 是 $V$ 的基，$W=\spn(\vv\alpha_1+\vv\alpha_2,\ \vv\alpha_2+\vv\alpha_3,\ \ldots,\ \vv\alpha_{n-1}+\vv\alpha_n)$. 求 $\dim W$，并求 $W$ 的一个补空间.
			\item\label{ex:vs:B15} 设 $V_1=\{\vv x\in\mathbb{F}^n\mid x_1=x_2=\cdots=x_n\}$，$V_2=\{\vv x\in\mathbb{F}^n\mid x_1=0\}$. 证明 $\mathbb{F}^n=V_1\oplus V_2$. 与第二章例~\ref{eg:pricedecomp} 的分解相比，哪一种分解更``自然''？（提示：思考补空间的不唯一性.）
			\item\label{ex:vs:B16} 在 $M_n(\mathbb{F})$ 中，设 $W_1$ 为对角线以下元素全为零的矩阵（上三角矩阵）全体，$W_2$ 为对角线及以上元素全为零的矩阵全体. 证明 $M_n(\mathbb{F})=W_1\oplus W_2$，并求 $\dim W_1$.
			\item\label{ex:vs:B17} （覆盖定理的另一个证明）设 $W_1,\ldots,W_s$ 是 $\mathbb{R}^n$ 的真子空间，令 $\vv v(t)=(1,t,t^2,\ldots,t^{n-1})^\trans$.
			\begin{enumerate}
				\item 证明：每个真子空间都包含于某个超平面 $\{\vv x\mid a_1x_1+\cdots+a_nx_n=0\}$（$a_i$ 不全为零）中.（提示：取 $W$ 的基并扩充为 $\mathbb{R}^n$ 的基，看``最后一个坐标''.）
				\item 证明：对上述超平面，满足 $\vv v(t)$ 落在其中的实数 $t$ 至多有 $n-1$ 个.
				\item 由此推出：$\vv v(t)$ 落在 $W_1\cup\cdots\cup W_s$ 中的 $t$ 至多有 $s(n-1)$ 个. 这给出了覆盖定理在 $\mathbb{R}^n$ 中的一个``构造性''证明. 对 $n=3$ 以及三个子空间 $x_1=x_2$，$x_2=x_3$，$x_1+x_2+x_3=0$，写出一个不在它们之中的 $\vv v(t)$.
			\end{enumerate}
			\item\label{ex:vs:B18} 设 $\dim V=n$，$W_1,\ldots,W_s$ 是 $V$ 的真子空间. 证明：存在 $V$ 的一组基 $\vv\varepsilon_1,\ldots,\vv\varepsilon_n$，其中\textbf{每个}向量都不属于 $W_1\cup\cdots\cup W_s$.（提示：逐个选取. 选 $\vv\varepsilon_k$ 时，要躲开的除了 $W_1,\ldots,W_s$，还有哪个真子空间？）
		\end{exB}
		
		\Cgroup
		\begin{exC}
			\item\label{ex:vs:C1} （覆盖定理在有限``数系''上的失效）设 $p$ 是素数，$\mathbb{F}_p=\{0,1,\ldots,p-1\}$ 按模 $p$ 作加法和乘法（可以承认它对加、减、乘以及除以非零元素都封闭. 它不包含 $\mathbb{Q}$，所以不是本书意义下的数域；但第二章中除定理~\ref{thm:Qsmallest} 与覆盖定理以外的定义和结论，在它上面照样成立）.
			\begin{enumerate}
				\item 证明 $\mathbb{F}_p^2$ 恰有 $p+1$ 个一维子空间，并且它们的并是整个 $\mathbb{F}_p^2$.
				\item 证明 $\mathbb{F}_p^2$ 不能写成 $p$ 个真子空间的并.（提示：每个一维子空间恰含 $p-1$ 个非零向量，数一数.）
				\item 指出定理~\ref{thm:cover} 的证明在 $\mathbb{F}_p$ 上从哪一步开始失效. 由第 2 问，``个数不超过 $p$ 时盖不住''，你能把正文的证明改写成这个``定量''的版本吗？
			\end{enumerate}
			\item\label{ex:vs:C2} 证明 $\ln2,\ln3,\ln5$ 在 $\mathbb{Q}$ 上线性无关. 由此说明 $\mathbb{R}$ 作为 $\mathbb{Q}$ 上的向量空间不是有限维的.（提示：算术基本定理.）
			\item\label{ex:vs:C3} 证明：全体实数列构成的向量空间不是有限维的；而``只有有限项非零''的实数列全体是它的子空间，且 $\vv e_1,\vv e_2,\ldots$（第 $i$ 项为 $1$，其余为 $0$）是它的一组``基''——每个元素都能唯一地写成其中\textbf{有限个}向量的线性组合. 这组向量为什么不是全体实数列的基？
			\item\label{ex:vs:C4} （模律）设 $W_1,W_2,W_3$ 是 $V$ 的子空间. 证明 $(W_1\cap W_3)+(W_2\cap W_3)\subset(W_1+W_2)\cap W_3$，在 $\mathbb{R}^2$ 中举例说明等号可以不成立；再证明：若 $W_1\subset W_3$，则等号一定成立.
			\item\label{ex:vs:C5} （有限域上的向量空间）在只含 $0,1$ 两个元素、按模 $2$ 运算的``数系'' $\mathbb{F}_2$ 上考虑 $\mathbb{F}_2^3$. (1) 它有多少个向量？(2) 有多少个一维子空间、多少个二维子空间？(3) 它有多少组（有序）基？推广到 $\mathbb{F}_2^n$.
			\item\label{ex:vs:C6} （塔律）设 $\mathbb{F}\subset\mathbb{K}\subset\mathbb{L}$ 是三个数域，$\mathbb{L}$ 作为 $\mathbb{K}$ 上的向量空间是 $m$ 维的，$\mathbb{K}$ 作为 $\mathbb{F}$ 上的向量空间是 $n$ 维的. 证明 $\mathbb{L}$ 作为 $\mathbb{F}$ 上的向量空间是 $mn$ 维的. 用它重新计算 B8 中数域的维数.
			\item\label{ex:vs:C7} 证明 $\mathbb{Q}(\sqrt2+\sqrt3)=\mathbb{Q}(\sqrt2,\sqrt3)$，即包含 $\sqrt2+\sqrt3$ 的最小数域与包含 $\sqrt2,\sqrt3$ 的最小数域相同.（提示：计算 $(\sqrt2+\sqrt3)^3$.）
			\item\label{ex:vs:C8} 设 $V$ 是 $n$ 维空间，$W$ 是 $k$ 维子空间（$0<k<n$）. 证明 $W$ 有无穷多个补空间. 在 $\mathbb{F}_2^2$ 中，一维子空间有几个补空间？
			\item\label{ex:vs:C9} （Fibonacci 空间的另一组基）在 B6 的记号下，改为研究满足 $a_{n+2}=a_{n+1}+a_n$ 的数列空间. 求出两个公比 $q$ 使 $\{q^n\}$ 属于这个空间，证明它们构成一组基，并推出 Fibonacci 数列的通项（Binet 公式）. 第八章会从另一个角度再次得到它.
		\end{exC}
	}
	
	% =====================================================================
	% 第三章　矩阵
	% =====================================================================
	\newcommand{\ExMatrix}{%
		\exsetheading
		\Agroup
		\begin{exA}
			\item\label{ex:mat:A1} 设 $\vv A=\begin{pmatrix}1&-1&2\\0&3&1\end{pmatrix}$，$\vv B=\begin{pmatrix}2&1\\-1&0\\1&4\end{pmatrix}$. 计算 $\vv{AB}$，$\vv{BA}$，$\vv A^\trans\vv B^\trans$，并验证 $(\vv{AB})^\trans=\vv B^\trans\vv A^\trans$.
			\item\label{ex:mat:A2} 按``列的线性组合''计算 $\begin{pmatrix}1&2&0\\3&-1&4\end{pmatrix}\begin{pmatrix}2\\-1\\1\end{pmatrix}$；再说明 $\vv e_i^\trans\vv A\vv e_j$ 是 $\vv A$ 的哪个元素，$\vv A\vv e_j$、$\vv e_i^\trans\vv A$ 又分别是什么.
			\item\label{ex:mat:A3} 计算 $\begin{pmatrix}1&1\\0&1\end{pmatrix}^n$，$\begin{pmatrix}\lambda&1\\0&\lambda\end{pmatrix}^n$，$\begin{pmatrix}\cos\theta&-\sin\theta\\ \sin\theta&\cos\theta\end{pmatrix}^n$，$\begin{pmatrix}0&1&0\\0&0&1\\0&0&0\end{pmatrix}^n$.
			\item\label{ex:mat:A4} 设 $\vv A=\begin{pmatrix}1&2\\0&1\end{pmatrix}$，$\vv B=\begin{pmatrix}1&0\\3&1\end{pmatrix}$. 计算 $(\vv A+\vv B)^2$ 与 $\vv A^2+2\vv{AB}+\vv B^2$，二者是否相等？$(\vv A+\vv B)(\vv A-\vv B)$ 与 $\vv A^2-\vv B^2$ 呢？
			\item\label{ex:mat:A5} 用初等行变换求下列矩阵的逆：
			\[\begin{pmatrix}1&2&3\\2&2&1\\3&4&3\end{pmatrix},\qquad\begin{pmatrix}1&1&1\\0&1&1\\0&0&1\end{pmatrix},\qquad \begin{pmatrix}2&1&0\\1&2&1\\0&1&2\end{pmatrix},\qquad\begin{pmatrix}0&0&0&1\\0&0&2&0\\0&3&0&0\\4&0&0&0\end{pmatrix}.\]
			\item\label{ex:mat:A6} 解矩阵方程：(1) $\begin{pmatrix}2&5\\1&3\end{pmatrix}\vv X=\begin{pmatrix}4&-6\\2&1\end{pmatrix}$；(2) $\vv X\begin{pmatrix}2&1\\1&1\end{pmatrix}=\begin{pmatrix}1&0\\0&2\end{pmatrix}$.
			\item\label{ex:mat:A7} 求矩阵 $\begin{pmatrix}1&1&a\\1&a&1\\a&1&1\end{pmatrix}$ 的秩（对 $a$ 分情况讨论）.
			\item\label{ex:mat:A8} 求矩阵 $\begin{pmatrix}1&2&-1&0&3\\2&4&1&3&0\\-1&-2&2&1&-3\\1&2&2&3&0\end{pmatrix}$ 的秩，以及列向量组的一个极大无关组，并用它表出其余各列.
			\item\label{ex:mat:A9} 写出 $3$ 阶初等矩阵 $\vv E(1,3)$，$\vv E(2(5))$，$\vv E(3,1(-2))$，并分别写出它们的逆. 计算 $\vv E(3,1(-2))\vv A$ 与 $\vv A\vv E(3,1(-2))$，其中 $\vv A$ 是任一 $3$ 阶方阵，说明效果.
			\item\label{ex:mat:A10} 把 $\begin{pmatrix}2&1\\4&3\end{pmatrix}$ 写成若干个初等矩阵的乘积.
			\item\label{ex:mat:A11} 证明：两个上三角矩阵之积仍是上三角矩阵，且积的对角元等于两者对角元之积. 可逆上三角矩阵的逆也是上三角矩阵吗？
			\item\label{ex:mat:A12} 设 $\vv A,\vv B$ 都是 $n$ 阶对称矩阵. 证明：(1) $\vv A^k$，$\vv A+\vv B$，$\vv P^\trans\vv A\vv P$ 都对称；(2) $\vv{AB}$ 对称当且仅当 $\vv{AB}=\vv{BA}$；(3) 对任意 $m\times n$ 矩阵 $\vv C$，$\vv C^\trans\vv C$ 与 $\vv C\vv C^\trans$ 都对称.
			\item\label{ex:mat:A13} (1) 设 $\vv x,\vv y\in\mathbb{F}^n$，证明 $\tr(\vv x\vv y^\trans)=\vv y^\trans\vv x$；(2) 设 $n$ 阶方阵 $\vv A$ 满足：对一切 $n$ 阶方阵 $\vv X$ 都有 $\tr(\vv{AX})=0$. 证明 $\vv A=\vv O$（提示：取 $\vv X=\vv E_{ji}$）；(3) 证明：不存在 $n$ 阶方阵 $\vv A,\vv B$ 使 $\vv{AB}-\vv{BA}=\vv I$.
			\item\label{ex:mat:A14} 设方阵 $\vv A$ 满足 $\vv A^2-\vv A-2\vv I=\vv O$. 证明 $\vv A$ 与 $\vv A+2\vv I$ 都可逆，并用 $\vv A$ 表示它们的逆. $\vv A+\vv I$ 一定可逆吗？
			\item\label{ex:mat:A15} 证明：若 $\vv A$ 可逆且 $\vv{AB}=\vv{AC}$，则 $\vv B=\vv C$；举例说明 $\vv A$ 不可逆时结论可能不成立.
			\item\label{ex:mat:A16} 设 $\vv A,\vv B$ 可逆. 求 $\begin{pmatrix}\vv A&\vv O\\ \vv O&\vv B\end{pmatrix}^{-1}$ 与 $\begin{pmatrix}\vv O&\vv A\\ \vv B&\vv O\end{pmatrix}^{-1}$，并用后者求 $\begin{pmatrix}0&0&1&2\\0&0&1&3\\2&0&0&0\\0&5&0&0\end{pmatrix}^{-1}$.
			\item\label{ex:mat:A17} 用分块乘法计算 $\begin{pmatrix}1&0&2&1\\0&1&1&3\\0&0&1&0\\0&0&0&1\end{pmatrix}^2$，并猜出 $n$ 次方的一般形式.
			\item\label{ex:mat:A18} 某公司在两个工厂生产三种产品，各工厂每天的产量（件）与三种产品的单位成本、单位售价（元）分别为
			\[\vv Q=\begin{pmatrix}100&80&60\\120&50&90\end{pmatrix}\ (\text{行：工厂；列：产品}),\qquad \vv C=\begin{pmatrix}20&30\\15&25\\40&55\end{pmatrix}\ (\text{列：成本、售价}).\]
			计算 $\vv Q\vv C$，说明它每个元素的含义，并求两个工厂每天的利润.
			\item\label{ex:mat:A19} 四座城市之间的直达航线为：$1\leftrightarrow2$，$1\leftrightarrow3$，$2\leftrightarrow3$，$3\leftrightarrow4$（双向）. 写出邻接矩阵 $\vv A$，计算 $\vv A^2$ 与 $\vv A^3$. $(\vv A^2)_{ii}$ 与城市 $i$ 的航线数有什么关系？从城市 $1$ 到城市 $4$ 恰好乘三次航班有几种走法？
			\item\label{ex:mat:A20} 设 $\vv A$ 是实反对称矩阵. 证明：对一切 $\vv x\in\mathbb{R}^n$ 有 $\vv x^\trans\vv A\vv x=0$；$\vv A^2$ 是对称矩阵.
			\item\label{ex:mat:A21} 证明：若 $\vv A$ 是 $m\times n$ 矩阵，$\vv P$ 是 $m$ 阶可逆矩阵，则 $\vv{PA}$ 与 $\vv A$ 的简化阶梯形相同.
			\item\label{ex:mat:A22} 证明：任一 $n$ 阶方阵都可以唯一地写成一个对称矩阵与一个反对称矩阵之和.（提示：$\vv A=\frac12(\vv A+\vv A^\trans)+\frac12(\vv A-\vv A^\trans)$.）
		\end{exA}
		
		\Bgroup
		\begin{exB}
			\item\label{ex:mat:B1} 分别求出全部满足 (1) $\vv A^2=\vv O$；(2) $\vv A^2=\vv I$；(3) $\vv A^2=\vv A$ 的二阶实方阵.
			\item\label{ex:mat:B2} 设 $a_1,\ldots,a_n$ 互不相同. 证明：与 $\diag(a_1,\ldots,a_n)$ 可交换的矩阵一定是对角矩阵.
			\item\label{ex:mat:B3} 证明：与\textbf{一切} $n$ 阶方阵都可交换的矩阵必为数量矩阵 $c\vv I$.（提示：考虑与 $\vv E_{ij}$ 交换.）
			\item\label{ex:mat:B4} 设 $\vv A,\vv B$ 都是幂零矩阵且 $\vv{AB}=\vv{BA}$. 证明 $\vv A+\vv B$ 与 $\vv{AB}$ 都是幂零矩阵. 举例说明``可交换''不能去掉.
			\item\label{ex:mat:B5} 证明：$\rk\begin{pmatrix}\vv A&\vv O\\ \vv O&\vv B\end{pmatrix}=\rk\vv A+\rk\vv B$，$\rk\begin{pmatrix}\vv A&\vv C\\ \vv O&\vv B\end{pmatrix}\ge\rk\vv A+\rk\vv B$，并举例说明后者可以取严格不等号.
			\item\label{ex:mat:B6} （Sherman–Morrison 公式）设 $\vv u,\vv v\in\mathbb{F}^n$，$1+\vv v^\trans\vv u\ne0$. 证明
			\[(\vv I+\vv u\vv v^\trans)^{-1}=\vv I-\frac{\vv u\vv v^\trans}{1+\vv v^\trans\vv u}.\]
			\item\label{ex:mat:B7} 设 $\vv A$ 为 $m\times n$ 矩阵，$\vv B$ 为 $n\times m$ 矩阵. 证明：若 $\vv I_m+\vv{AB}$ 可逆，则 $\vv I_n+\vv{BA}$ 也可逆，并且 $(\vv I_n+\vv{BA})^{-1}=\vv I_n-\vv B(\vv I_m+\vv{AB})^{-1}\vv A$.
			\item\label{ex:mat:B8} 证明：秩为 $1$ 的 $m\times n$ 矩阵都可以写成 $\vv u\vv v^\trans$（$\vv u\in\mathbb{F}^m$，$\vv v\in\mathbb{F}^n$ 非零）. 对这样的矩阵证明 $\vv A^2=(\tr\vv A)\vv A$（$m=n$ 时），并由此求 $\vv A^k$.
			\item\label{ex:mat:B9} 证明：秩为 $r$ 的矩阵可以写成 $r$ 个秩为 $1$ 的矩阵之和，但不能写成少于 $r$ 个秩为 $1$ 的矩阵之和.
			\item\label{ex:mat:B10} （满秩分解）证明：秩为 $r$ 的 $m\times n$ 矩阵 $\vv A$ 可以写成 $\vv A=\vv{BC}$，其中 $\vv B$ 是列满秩的 $m\times r$ 矩阵，$\vv C$ 是行满秩的 $r\times n$ 矩阵. 对 A8 中的矩阵写出一个满秩分解.
			\item\label{ex:mat:B11} 计算 $\begin{pmatrix}1&1&1\\1&1&1\\1&1&1\end{pmatrix}^n$ 与 $\begin{pmatrix}2&1&0\\0&2&1\\0&0&2\end{pmatrix}^n$.（提示：后者写成 $2\vv I+\vv N$.）
			\item\label{ex:mat:B12} 证明：$n$ 阶严格上三角矩阵（对角线及以下全为零）$\vv N$ 满足 $\vv N^n=\vv O$. 更精确地，$\vv N^k$ 的非零元只可能出现在第 $k$ 条``上对角线''及其右上方.
			\item\label{ex:mat:B13} 设 $\vv A$ 为实矩阵. 证明：若 $\vv A^\trans\vv A=\vv O$，则 $\vv A=\vv O$；若 $\tr(\vv A^\trans\vv A)=0$，结论同样成立. 复矩阵呢？
			\item\label{ex:mat:B14} 每行、每列恰有一个 $1$、其余元素为 $0$ 的方阵称为\textbf{置换矩阵}. 证明：置换矩阵之积仍是置换矩阵，置换矩阵 $\vv P$ 满足 $\vv P^\trans\vv P=\vv I$；$\vv{PA}$ 与 $\vv{AP}$ 分别对 $\vv A$ 做了什么？
			\item\label{ex:mat:B15} 对 $\vv A=\begin{pmatrix}1&2&1\\2&4&3\\1&2&2\end{pmatrix}$，求可逆矩阵 $\vv P,\vv Q$ 使 $\vv{PAQ}$ 为相抵标准形.
			\item\label{ex:mat:B16} 设 $\vv A$ 是 $n$ 阶方阵，$\vv A^k=\vv O$. 证明 $\vv I+\vv A$ 可逆，并求其逆. 若 $\vv A^2=2\vv A$，$\vv I-\vv A$ 是否可逆？
		\end{exB}
		
		\Cgroup
		\begin{exC}
			\item\label{ex:mat:C1} 设 $f:M_n(\mathbb{F})\to\mathbb{F}$ 满足 $f(\vv A+\vv B)=f(\vv A)+f(\vv B)$，$f(k\vv A)=kf(\vv A)$，且 $f(\vv{AB})=f(\vv{BA})$ 对一切 $\vv A,\vv B$ 成立. 证明 $f=c\cdot\tr$ 对某个常数 $c$ 成立.（提示：计算 $f(\vv E_{ij})$.）
			\item\label{ex:mat:C2} 设 $n$ 阶方阵 $\vv A,\vv B$ 满足 $\vv A+\vv B=\vv{AB}$. 证明 $\vv{AB}=\vv{BA}$.（提示：考察 $(\vv I-\vv A)(\vv I-\vv B)$.）
			\item\label{ex:mat:C3} （Strassen 算法）对二阶分块矩阵的乘法，直接计算需要 $8$ 次块乘法. 验证下面的公式只需 $7$ 次：令
			$M_1=(A_{11}+A_{22})(B_{11}+B_{22})$，$M_2=(A_{21}+A_{22})B_{11}$，$M_3=A_{11}(B_{12}-B_{22})$，$M_4=A_{22}(B_{21}-B_{11})$，$M_5=(A_{11}+A_{12})B_{22}$，$M_6=(A_{21}-A_{11})(B_{11}+B_{12})$，$M_7=(A_{12}-A_{22})(B_{21}+B_{22})$，则
			$C_{11}=M_1+M_4-M_5+M_7$，$C_{12}=M_3+M_5$，$C_{21}=M_2+M_4$，$C_{22}=M_1-M_2+M_3+M_6$. 由此说明 $n=2^k$ 阶矩阵乘法可以用约 $n^{\log_27}\approx n^{2.81}$ 次数的乘法完成.
			\item\label{ex:mat:C4} 设 $\vv A$ 是幂等矩阵. 用满秩分解 $\vv A=\vv{BC}$ 证明 $\vv{CB}=\vv I_r$，进而证明 $\rk\vv A=\tr\vv A$.
			\item\label{ex:mat:C5} （Kronecker 积）对 $m\times n$ 矩阵 $\vv A$ 与 $p\times q$ 矩阵 $\vv B$，定义 $\vv A\otimes\vv B=(a_{ij}\vv B)$（$mp\times nq$ 分块矩阵）. 证明 $(\vv A\otimes\vv B)(\vv C\otimes\vv D)=(\vv{AC})\otimes(\vv{BD})$，并由此说明 $\vv A,\vv B$ 可逆时 $(\vv A\otimes\vv B)^{-1}=\vv A^{-1}\otimes\vv B^{-1}$.
			\item\label{ex:mat:C6} （Hadamard 矩阵）元素全为 $\pm1$ 且满足 $\vv H\vv H^\trans=n\vv I$ 的 $n$ 阶矩阵称为 Hadamard 矩阵. 验证 $\vv H_2=\begin{pmatrix}1&1\\1&-1\end{pmatrix}$ 是，并证明 $\vv H_2\otimes\vv H_2$，$\vv H_2\otimes\vv H_2\otimes\vv H_2$ 也是. 证明：$n>2$ 时 Hadamard 矩阵的阶数必是 $4$ 的倍数.（这类矩阵用于实验设计与纠错编码.）
			\item\label{ex:mat:C7} 证明：全体形如 $\begin{pmatrix}a&-b\\b&a\end{pmatrix}$（$a,b\in\mathbb{R}$）的矩阵关于矩阵的加法和乘法``像复数一样''运算：加法、乘法封闭，乘法可交换，非零元可逆. 写出它与 $\mathbb{C}$ 之间的对应，并用它验证 $(1+\mathrm i)^8=16$.
			\item\label{ex:mat:C8} 设 $\vv A$ 是元素非负、各列之和都小于 $1$ 的方阵. 证明 $\vv A^k$ 的每个元素都随 $k\to\infty$ 趋于 $0$，并且级数 $\vv I+\vv A+\vv A^2+\cdots$（按元素）收敛到 $(\vv I-\vv A)^{-1}$. 解释它在投入产出模型中的意义：``直接消耗 $+$ 一次间接消耗 $+$ 二次间接消耗 $+\cdots$''.
			\item\label{ex:mat:C9} 设 $\vv A$ 为 $m\times n$ 矩阵. 证明：$\vv A$ 有右逆（存在 $\vv C$ 使 $\vv{AC}=\vv I_m$）当且仅当 $\rk\vv A=m$；有左逆当且仅当 $\rk\vv A=n$. 非方阵的单侧逆唯一吗？
		\end{exC}
	}
	
	% =====================================================================
	% 第四章　线性映射
	% =====================================================================
	\newcommand{\ExLinearMap}{%
		\exsetheading
		\Agroup
		\begin{exA}
			\item\label{ex:lmap:A1} 判断下列映射是否线性，是线性映射的求出核：
			\begin{enumerate}
				\item $\mathbb{R}^2\to\mathbb{R}^2$，$(x,y)^\trans\mapsto(x+y,\,2x)^\trans$；
				\item $\mathbb{R}^2\to\mathbb{R}^2$，$(x,y)^\trans\mapsto(x+1,\,y)^\trans$；
				\item $\mathbb{R}^2\to\mathbb{R}$，$(x,y)^\trans\mapsto|x|+y$；
				\item $\mathbb{R}[x]\to\mathbb{R}[x]$，$f(x)\mapsto f(x+1)$；
				\item $\mathbb{R}[x]\to\mathbb{R}[x]$，$f(x)\mapsto xf'(x)$；
				\item $M_n(\mathbb{R})\to M_n(\mathbb{R})$，$\vv X\mapsto\vv X^\trans$；
				\item $M_2(\mathbb{R})\to\mathbb{R}$，$\vv X\mapsto x_{11}x_{22}-x_{12}x_{21}$.
			\end{enumerate}
			\item\label{ex:lmap:A2} 求下列矩阵所定义的线性映射 $\vv x\mapsto\vv{Ax}$ 的核与像（各给出一组基），并验证秩—零化度定理：
			\[\begin{pmatrix}1&2&-1\\2&4&-2\end{pmatrix},\qquad\begin{pmatrix}1&0&2&1\\0&1&-1&1\\1&1&1&2\end{pmatrix}.\]
			\item\label{ex:lmap:A3} 写出线性映射 $\mathscr T:\mathbb{R}^3\to\mathbb{R}^2$，$\mathscr T(x_1,x_2,x_3)^\trans=(x_1-2x_2+x_3,\ 3x_2+4x_3)^\trans$ 在标准基下的矩阵.
			\item\label{ex:lmap:A4} 求下列平面变换在标准基下的矩阵：(1) 逆时针旋转 $90^\circ$；(2) 关于直线 $y=x$ 的反射；(3) 向直线 $y=x$ 的正投影；(4) 沿 $x$ 方向的剪切 $(x,y)\mapsto(x+2y,y)$. 计算 (2)(3) 的矩阵的平方，并解释结果.
			\item\label{ex:lmap:A5} 求求导变换 $\mathscr D:\mathbb{R}[x]_3\to\mathbb{R}[x]_3$ 在基 $1,x,x^2$ 下的矩阵，以及在基 $1,\ x-1,\ (x-1)^2$ 下的矩阵.
			\item\label{ex:lmap:A6} 求转置变换 $\vv X\mapsto\vv X^\trans$ 在 $M_2(\mathbb{F})$ 的基 $\vv E_{11},\vv E_{12},\vv E_{21},\vv E_{22}$ 下的矩阵. 它的平方是什么？
			\item\label{ex:lmap:A7} 在 $\mathbb{R}^3$ 中，求从基 $\vv\varepsilon_1=(1,0,0)^\trans$，$\vv\varepsilon_2=(1,1,0)^\trans$，$\vv\varepsilon_3=(1,1,1)^\trans$ 到基 $\vv\eta_1=(1,2,1)^\trans$，$\vv\eta_2=(2,3,3)^\trans$，$\vv\eta_3=(3,7,1)^\trans$ 的过渡矩阵，并求 $\vv\alpha=(3,4,2)^\trans$ 在两组基下的坐标.
			\item\label{ex:lmap:A8} 设线性变换 $\mathscr T$ 在基 $\vv\varepsilon_1,\vv\varepsilon_2$ 下的矩阵为 $\begin{pmatrix}2&1\\1&2\end{pmatrix}$. 求它在基 $\vv\eta_1=\vv\varepsilon_1+\vv\varepsilon_2$，$\vv\eta_2=\vv\varepsilon_1-\vv\varepsilon_2$ 下的矩阵.
			\item\label{ex:lmap:A9} 证明：线性映射的复合仍是线性映射；可逆（双射）线性映射的逆映射仍是线性映射.
			\item\label{ex:lmap:A10} 证明：线性映射 $\mathscr T$ 是单射，当且仅当它把每个线性无关组都映为线性无关组.
			\item\label{ex:lmap:A11} 是否存在线性变换 $\mathscr T:\mathbb{R}^2\to\mathbb{R}^2$，使
			(1) $\mathscr T(1,1)^\trans=(1,0)^\trans$，$\mathscr T(1,-1)^\trans=(0,1)^\trans$，$\mathscr T(2,0)^\trans=(1,1)^\trans$？
			(2) 同上，但最后一个条件改为 $\mathscr T(2,0)^\trans=(1,2)^\trans$？
			存在时求出 $\mathscr T(x,y)^\trans$.
			\item\label{ex:lmap:A12} (1) 线性映射 $\mathscr T:\mathbb{R}^5\to\mathbb{R}^3$ 是满射，$\dim\ker\mathscr T$ 是多少？(2) 是否存在从 $\mathbb{R}^3$ 到 $\mathbb{R}^5$ 的满射线性映射？是否存在从 $\mathbb{R}^5$ 到 $\mathbb{R}^3$ 的单射线性映射？
			\item\label{ex:lmap:A13} 设 $\mathscr T:U\to V$，$\mathscr S:V\to W$ 线性. 证明 $\ker\mathscr T\subset\ker(\mathscr{ST})$，$\imm(\mathscr{ST})\subset\imm\mathscr S$，从而 $\rk(\mathscr{ST})\le\min\{\rk\mathscr S,\rk\mathscr T\}$.
			\item\label{ex:lmap:A14} 求迹映射 $\tr:M_n(\mathbb{F})\to\mathbb{F}$ 的核的维数.
			\item\label{ex:lmap:A15} 证明：相似矩阵有相同的迹（利用命题~\ref{prop:trace}(3)）. 由此判断 $\begin{pmatrix}1&2\\0&3\end{pmatrix}$ 与 $\begin{pmatrix}1&0\\0&2\end{pmatrix}$ 是否相似.
			\item\label{ex:lmap:A16} 证明：与单位矩阵相似的矩阵只有单位矩阵本身. 由此说明 $\begin{pmatrix}1&1\\0&1\end{pmatrix}$ 与 $\vv I$ 不相似.
			\item\label{ex:lmap:A17} 定义求值映射 $\mathscr E:\mathbb{R}[x]_3\to\mathbb{R}^3$，$f\mapsto(f(0),f(1),f(2))^\trans$. 求它在基 $1,x,x^2$ 与标准基下的矩阵，并证明 $\mathscr E$ 是同构.
			\item\label{ex:lmap:A18} 设 $\mathscr T$ 是 $\mathbb{R}^3$ 上的线性变换，$\mathscr T\vv e_1=(1,1,0)^\trans$，$\mathscr T\vv e_2=(0,1,1)^\trans$，$\mathscr T\vv e_3=(1,2,1)^\trans$. 求 $\mathscr T(x_1,x_2,x_3)^\trans$，$\ker\mathscr T$ 与 $\imm\mathscr T$.
			\item\label{ex:lmap:A19} 称 $\vv x\mapsto\vv{Ax}+\vv b$ 为\textbf{仿射变换}. 证明两个仿射变换的复合仍是仿射变换；仿射变换是线性变换当且仅当 $\vv b=\zero$.
			\item\label{ex:lmap:A20} 设 $\vv R(\theta)=\begin{pmatrix}\cos\theta&-\sin\theta\\ \sin\theta&\cos\theta\end{pmatrix}$ 是平面上绕原点旋转角 $\theta$ 的变换在标准基下的矩阵. 从几何上说明``先转 $\beta$ 再转 $\alpha$''就是转 $\alpha+\beta$，从而 $\vv R(\alpha)\vv R(\beta)=\vv R(\alpha+\beta)$；比较两边矩阵的元素，推出两角和的正弦、余弦公式.
		\end{exA}
		
		\Bgroup
		\begin{exB}
			\item\label{ex:lmap:B1} 记 $L(V,W)$ 为 $V$ 到 $W$ 的全体线性映射. 证明它关于映射的加法与数乘构成向量空间，且 $\dim L(V,W)=\dim V\cdot\dim W$.（提示：取定基后，线性映射与矩阵一一对应.）
			\item\label{ex:lmap:B2} 设 $\mathscr T$ 是 $n$ 维空间 $V$ 上的线性变换，$\mathscr T^2=\mathscr O$. 证明 $\imm\mathscr T\subset\ker\mathscr T$，从而 $\rk\mathscr T\le n/2$；并证明存在一组基，使 $\mathscr T$ 的矩阵只在 $(1,2),(3,4),\ldots,(2r-1,2r)$ 这 $r$ 个位置上为 $1$，其余元素全为 $0$，其中 $r=\rk\mathscr T$.
			\item\label{ex:lmap:B3} 设 $\mathscr T,\mathscr S$ 是有限维空间上的线性变换. 证明 $\dim\ker(\mathscr{ST})\le\dim\ker\mathscr S+\dim\ker\mathscr T$.
			\item\label{ex:lmap:B4} 设 $V$ 有限维，$\mathscr S,\mathscr T$ 是 $V$ 上的线性变换且 $\mathscr{ST}=\mathscr I$. 证明 $\mathscr{TS}=\mathscr I$.
			\item\label{ex:lmap:B5} 在无限维空间 $\mathbb{R}[x]$ 上，令 $\mathscr D$ 为求导，$\mathscr J$ 为 $f\mapsto\int_0^xf(t)\,\mathrm dt$. 证明 $\mathscr{DJ}=\mathscr I$，但 $\mathscr{JD}\ne\mathscr I$；$\mathscr D$ 是满射但不是单射，$\mathscr J$ 是单射但不是满射. 这与 B4 是否矛盾？
			\item\label{ex:lmap:B6} 固定 $\vv A\in M_n(\mathbb{F})$，令 $\mathscr T(\vv X)=\vv{AX}-\vv{XA}$. 证明 $\imm\mathscr T$ 含于迹为零的矩阵全体，从而 $\mathscr T$ 不是满射；并证明 $\ker\mathscr T$ 包含 $\vv I,\vv A,\vv A^2,\ldots$.
			\item\label{ex:lmap:B7} 设 $\vv A=\begin{pmatrix}a&b\\c&d\end{pmatrix}$. 求线性变换 $\vv X\mapsto\vv{AX}$ 在 $M_2(\mathbb{F})$ 的基 $\vv E_{11},\vv E_{21},\vv E_{12},\vv E_{22}$ 下的矩阵. 若换成 $\vv X\mapsto\vv{XA}$ 呢？
			\item\label{ex:lmap:B8} 设线性变换 $\mathscr P$ 满足 $\mathscr P^2=\mathscr P$. 证明存在一组基使 $\mathscr P$ 的矩阵为 $\begin{pmatrix}\vv I_r&\vv O\\ \vv O&\vv O\end{pmatrix}$，其中 $r=\rk\mathscr P$.
			\item\label{ex:lmap:B9} 设线性变换 $\mathscr T$ 满足 $\mathscr T^2=\mathscr I$. 证明 $V=\ker(\mathscr T-\mathscr I)\oplus\ker(\mathscr T+\mathscr I)$，从而存在一组基使 $\mathscr T$ 的矩阵为 $\diag(\vv I_p,-\vv I_q)$.（提示：$\vv\alpha=\frac12(\vv\alpha+\mathscr T\vv\alpha)+\frac12(\vv\alpha-\mathscr T\vv\alpha)$.）
			\item\label{ex:lmap:B10} 设 $\mathscr T:V\to W$ 线性，$U$ 是 $V$ 的子空间. 证明 $\dim\mathscr T(U)=\dim U-\dim(U\cap\ker\mathscr T)$.
			\item\label{ex:lmap:B11} 在由 $\mathrm e^x\cos x$，$\mathrm e^x\sin x$ 张成的二维函数空间中，求求导变换的矩阵，证明它可逆，并利用逆矩阵求出 $\int\mathrm e^x\sin x\,\mathrm dx$（不必用分部积分）.
			\item\label{ex:lmap:B12} （齐次坐标）平面上的平移 $(x,y)\mapsto(x+a,y+b)$ 不是线性的. 把点 $(x,y)$ 写成三维向量 $(x,y,1)^\trans$，证明平移、旋转、伸缩都可以写成 $3\times3$ 矩阵的乘法，且仿射变换的复合对应矩阵的乘积. 写出``先绕原点旋转 $90^\circ$，再平移 $(1,2)$''与``先平移再旋转''的矩阵，它们相等吗？（计算机图形学正是这样处理变换的.）
			\item\label{ex:lmap:B13} 证明：若 $\vv A$ 与 $\vv B$ 相似，则 $\vv A^k$ 与 $\vv B^k$ 相似，$\vv A^\trans$ 与 $\vv B^\trans$ 相似，$f(\vv A)$ 与 $f(\vv B)$ 相似（$f$ 为多项式）.
			\item\label{ex:lmap:B14} 证明 $\begin{pmatrix}1&0\\0&0\end{pmatrix}$ 与 $\begin{pmatrix}0&1\\0&0\end{pmatrix}$ 相抵但不相似. 从``线性映射 $V\to W$''与``线性变换 $V\to V$''的区别解释：为什么相抵只需看秩，而相似要复杂得多？
			\item\label{ex:lmap:B15} 设 $\mathscr T$ 是 $V$ 上的线性变换，$\vv\alpha\in V$ 满足 $\mathscr T^{k-1}\vv\alpha\ne\zero$，$\mathscr T^k\vv\alpha=\zero$. 证明 $\vv\alpha,\mathscr T\vv\alpha,\ldots,\mathscr T^{k-1}\vv\alpha$ 线性无关. 若 $k=n=\dim V$，求 $\mathscr T$ 在基 $\mathscr T^{n-1}\vv\alpha,\ldots,\mathscr T\vv\alpha,\vv\alpha$ 下的矩阵.
		\end{exB}
		
		\Cgroup
		\begin{exC}
			\item\label{ex:lmap:C1} （Fitting 分解）设 $\mathscr T$ 是 $n$ 维空间 $V$ 上的线性变换. 证明：
			\begin{enumerate}
				\item $\ker\mathscr T\subset\ker\mathscr T^2\subset\cdots$，$\imm\mathscr T\supset\imm\mathscr T^2\supset\cdots$，且两条链都在某一步之后不再变化；
				\item $V=\ker\mathscr T^n\oplus\imm\mathscr T^n$，两者都是 $\mathscr T$ 的不变子空间，$\mathscr T$ 在前者上幂零，在后者上可逆.
			\end{enumerate}
			\item\label{ex:lmap:C2} （商空间）设 $W$ 是 $V$ 的子空间. 对 $\vv\alpha\in V$，称集合 $\vv\alpha+W$ 为 $\vv\alpha$ 所在的陪集，全体陪集记为 $V/W$. 定义 $(\vv\alpha+W)+(\vv\beta+W)=(\vv\alpha+\vv\beta)+W$，$k(\vv\alpha+W)=k\vv\alpha+W$.
			\begin{enumerate}
				\item 证明这两个运算是良定义的，且 $V/W$ 构成向量空间；
				\item 证明 $\dim V/W=\dim V-\dim W$；
				\item （第一同构定理/同态基本定理）对线性映射 $\mathscr T:V\to U$，证明 $V/\ker\mathscr T\cong\imm\mathscr T$. 这给出秩—零化度定理的另一种理解.
			\end{enumerate}
			\item\label{ex:lmap:C3} 把 $\mathbb{C}$ 看作 $\mathbb{R}$ 上的二维空间，基为 $1,\mathrm i$. (1) 求``乘以 $a+b\mathrm i$''这个变换的矩阵；(2) 求``取共轭''这个变换的矩阵；(3) 证明 $\mathbb{R}$-线性变换 $\begin{pmatrix}p&q\\r&s\end{pmatrix}$ 是``乘以某个复数''，当且仅当 $p=s$，$q=-r$.
			\item\label{ex:lmap:C4} 证明：$n$ 维空间上的任一线性变换 $\mathscr T$，都存在次数不超过 $n^2$ 的非零多项式 $f$ 使 $f(\mathscr T)=\mathscr O$.（提示：$\mathscr I,\mathscr T,\ldots,\mathscr T^{n^2}$ 在 $L(V,V)$ 中线性相关. 第八章的 Cayley–Hamilton 定理会把 $n^2$ 改进为 $n$.）
			\item\label{ex:lmap:C5} （Haar 小波）在 $\mathbb{R}^4$ 中取基 $\vv h_1=(1,1,1,1)^\trans$，$\vv h_2=(1,1,-1,-1)^\trans$，$\vv h_3=(1,-1,0,0)^\trans$，$\vv h_4=(0,0,1,-1)^\trans$. 求信号 $\vv x=(8,6,3,4)^\trans$ 在这组基下的坐标，并解释每个坐标的含义. 若把绝对值最小的坐标改为 $0$ 再变换回来，误差是多少？这是图像与音频压缩的思想之一.
			\item\label{ex:lmap:C6} 设 $\mathscr T$ 是 $V$ 上的线性变换. 若 $\mathscr T$ 与 $V$ 上的一切线性变换都可交换，证明 $\mathscr T=c\mathscr I$. 若只要求与一切\textbf{可逆}线性变换可交换呢？
			\item\label{ex:lmap:C7} 设 $\mathscr T:V\to V$ 线性，$\dim V=n$. 证明：$\rk\mathscr T^k-\rk\mathscr T^{k+1}$ 关于 $k$ 单调不增. 由此说明：若 $\rk\mathscr T=\rk\mathscr T^2$，则对一切 $k$，$\rk\mathscr T^k=\rk\mathscr T$.
			\item\label{ex:lmap:C8} （差分算子）在全体实数列构成的空间上，定义左移 $\mathscr S(a_1,a_2,a_3,\ldots)=(a_2,a_3,\ldots)$. 证明 $\mathscr S$ 是线性的、满射但不是单射；满足 $a_{n+2}=a_{n+1}+a_n$ 的数列全体恰好是 $\ker(\mathscr S^2-\mathscr S-\mathscr I)$. 由第二章章末习题~\ref{ex:vs:C9}，这个核是二维的.
		\end{exC}
	}
	
	% =====================================================================
	% 第五章　行列式
	% =====================================================================
	\newcommand{\ExDeterminant}{%
		\exsetheading
		\Agroup
		\begin{exA}
			\item\label{ex:det:A1} 求下列排列的逆序数并判断奇偶：$315426$，$42531$，$13\cdots(2n-1)\,24\cdots(2n)$.
			\item\label{ex:det:A2} (1) 在 $4$ 阶行列式中，项 $a_{13}a_{24}a_{32}a_{41}$ 带什么符号？(2) 选择 $i,j$，使 $a_{1i}a_{23}a_{35}a_{5j}a_{44}$ 成为 $5$ 阶行列式中带正号的一项.
			\item\label{ex:det:A3} 计算：
			\[\begin{vmatrix}2&1&-1\\3&0&2\\1&4&5\end{vmatrix},\qquad
			\begin{vmatrix}1&2&0&1\\2&1&1&0\\0&1&2&1\\1&0&1&2\end{vmatrix},\qquad
			\begin{vmatrix}3&1&1&1\\1&3&1&1\\1&1&3&1\\1&1&1&3\end{vmatrix}.\]
			\item\label{ex:det:A4} 计算 $\begin{vmatrix}1&-1&2&3\\2&0&1&-1\\0&1&-1&2\\3&1&0&1\end{vmatrix}$，分别用``打洞''化三角形与按第一列展开两种方法，比较哪种更快.
			\item\label{ex:det:A5} 计算 $n$ 阶行列式 $\begin{vmatrix}0&\cdots&0&a_1\\0&\cdots&a_2&0\\ \vdots&&&\vdots\\a_n&\cdots&0&0\end{vmatrix}$（只有副对角线上有元素）.
			\item\label{ex:det:A6} 不展开，证明：
			\[\begin{vmatrix}a+b&b+c&c+a\\a_1+b_1&b_1+c_1&c_1+a_1\\a_2+b_2&b_2+c_2&c_2+a_2\end{vmatrix}=2\begin{vmatrix}a&b&c\\a_1&b_1&c_1\\a_2&b_2&c_2\end{vmatrix}.\]
			\item\label{ex:det:A7} 解方程 $\begin{vmatrix}x&1&1\\1&x&1\\1&1&x\end{vmatrix}=0$ 与 $\begin{vmatrix}1&1&1\\1&2&x\\1&4&x^2\end{vmatrix}=0$.
			\item\label{ex:det:A8} 设 $\vv A$ 为 $3$ 阶方阵，$\det\vv A=3$. 求 $\det(2\vv A)$，$\det(-\vv A)$，$\det(\vv A^{-1})$，$\det(\vv A^*)$，$\det(\vv A^\trans\vv A)$，$\det(3\vv A^{-1}-2\vv A^*)$.
			\item\label{ex:det:A9} 证明：(1) 若 $n$ 阶方阵满足 $\vv A^\trans\vv A=\vv I$，则 $\det\vv A=\pm1$；(2) 若 $\vv A^2=\vv I$，则 $\det\vv A=\pm1$；(3) 若 $\vv A^2=\vv A$ 且 $\vv A\ne\vv I$，则 $\det\vv A=0$.
			\item\label{ex:det:A10} 求下列矩阵的伴随矩阵与逆矩阵：\[\begin{pmatrix}a&b\\c&d\end{pmatrix},\qquad\begin{pmatrix}1&2&3\\0&1&4\\5&6&0\end{pmatrix}.\]
			\item\label{ex:det:A11} 用 Cramer 法则解方程组 $2x_1+x_2+x_3=3$，$x_1+3x_2+x_3=0$，$x_1+x_2+4x_3=8$.
			\item\label{ex:det:A12} 求以 $(1,1)$，$(4,2)$，$(2,5)$ 为顶点的三角形的面积；判断 $(1,2)$，$(3,5)$，$(7,11)$ 三点是否共线.
			\item\label{ex:det:A13} 求以 $(1,0,1)^\trans$，$(2,1,0)^\trans$，$(0,3,1)^\trans$ 为棱的平行六面体的体积，以及以原点和这三个向量的终点为顶点的四面体的体积.
			\item\label{ex:det:A14} $\lambda$ 取何值时，齐次方程组 $(5-\lambda)x+2y+2z=0$，$2x+(6-\lambda)y=0$，$2x+(4-\lambda)z=0$ 有非零解？
			\item\label{ex:det:A15} 设 \[\vv A=\begin{pmatrix}1&2&3&4\\2&3&4&1\\3&4&1&2\\4&1&2&3\end{pmatrix}.\] 求 $A_{11}+A_{12}+A_{13}+A_{14}$（$A_{1j}$ 是代数余子式）.（提示：它是把第一行换成全 $1$ 后的行列式.）
			\item\label{ex:det:A16} 证明：过平面上两个不同点 $(x_1,y_1)$，$(x_2,y_2)$ 的直线方程可以写成 \[\begin{vmatrix}x&y&1\\x_1&y_1&1\\x_2&y_2&1\end{vmatrix}=0.\]
			\item\label{ex:det:A17} 设 $\vv A,\vv B$ 分别是 $m$ 阶、$n$ 阶方阵，证明 $\begin{vmatrix}\vv A&\vv O\\ \vv C&\vv B\end{vmatrix}=\det\vv A\det\vv B$，并计算 $\begin{vmatrix}\vv O&\vv A\\ \vv B&\vv O\end{vmatrix}$.
			\item\label{ex:det:A18} 计算 $4$ 阶三对角行列式 $\begin{vmatrix}2&-1&0&0\\-1&2&-1&0\\0&-1&2&-1\\0&0&-1&2\end{vmatrix}$.
			\item\label{ex:det:A19} 设 $\vv A$ 是 $n$ 阶可逆矩阵，$n\ge2$. 证明 $\vv A^*$ 可逆，且 $(\vv A^*)^{-1}=\frac1{\det\vv A}\vv A=(\vv A^{-1})^*$.
			\item\label{ex:det:A20} 用行列式判断：向量 $(1,2,3)^\trans$，$(4,5,6)^\trans$，$(7,8,9)^\trans$ 是否线性无关？把 $9$ 改成 $10$ 呢？
			\item\label{ex:det:A21} 设 $\vv A$ 是 $n$ 阶反对称矩阵（$\vv A^\trans=-\vv A$），$n$ 为奇数. 证明 $\det\vv A=0$.（提示：对 $\vv A^\trans=-\vv A$ 两边取行列式.）
		\end{exA}
		
		\Bgroup
		\begin{exB}
			\item\label{ex:det:B1} 计算 $n$ 阶三对角行列式 $D_n=\begin{vmatrix}2&-1&&\\-1&2&\ddots&\\&\ddots&\ddots&-1\\&&-1&2\end{vmatrix}$.（提示：按第一行展开得 $D_n=2D_{n-1}-D_{n-2}$.）
			\item\label{ex:det:B2} 计算 $n$ 阶行列式 $\begin{vmatrix}a&b&\cdots&b\\b&a&\cdots&b\\ \vdots&&\ddots&\vdots\\b&b&\cdots&a\end{vmatrix}$ 与 $\begin{vmatrix}x&a&\cdots&a\\-a&x&\cdots&a\\ \vdots&&\ddots&\vdots\\-a&-a&\cdots&x\end{vmatrix}$.
			\item\label{ex:det:B3} （爪形行列式）设 $a_1\cdots a_n\ne0$，计算 $\begin{vmatrix}a_0&1&1&\cdots&1\\1&a_1&0&\cdots&0\\1&0&a_2&\cdots&0\\ \vdots&&&\ddots&\\1&0&0&\cdots&a_n\end{vmatrix}$.
			\item\label{ex:det:B4} （矩阵行列式引理）证明 $\det(\vv I_n+\vv u\vv v^\trans)=1+\vv v^\trans\vv u$，更一般地，$\vv A$ 可逆时 $\det(\vv A+\vv u\vv v^\trans)=(1+\vv v^\trans\vv A^{-1}\vv u)\det\vv A$. 用它重做 B2 的第一个行列式.
			\item\label{ex:det:B5} 设 $\vv A$ 为 $2\times3$ 矩阵，$\vv B$ 为 $3\times2$ 矩阵. 证明 $\det(\vv{AB})$ 等于 $\vv A$ 的三个 $2$ 阶子式与 $\vv B$ 的对应 $2$ 阶子式乘积之和（这是 Cauchy–Binet 公式最简单的情形）. 由此说明 $\vv{BA}$（$3$ 阶方阵）的行列式为零.
			\item\label{ex:det:B6} 设 $\vv A$ 是整数矩阵. 证明：$\vv A^{-1}$ 存在且也是整数矩阵，当且仅当 $\det\vv A=\pm1$.
			\item\label{ex:det:B7} 证明：$n\ge2$ 时 $(\vv A^*)^*=(\det\vv A)^{n-2}\vv A$.（先设 $\vv A$ 可逆；不可逆的情形另行讨论.）
			\item\label{ex:det:B8} 设 $\vv A,\vv B$ 都是 $n$ 阶方阵. 证明 $\begin{vmatrix}\vv A&\vv B\\ \vv B&\vv A\end{vmatrix}=\det(\vv A+\vv B)\det(\vv A-\vv B)$.（提示：分块初等变换.）
			\item\label{ex:det:B9} 设 $\vv A$ 可逆. 证明 $\begin{vmatrix}\vv A&\vv B\\ \vv C&\vv D\end{vmatrix}=\det\vv A\cdot\det(\vv D-\vv{CA}^{-1}\vv B)$；若还有 $\vv{AC}=\vv{CA}$，则它等于 $\det(\vv{AD}-\vv{CB})$.
			\item\label{ex:det:B10} 用 Fibonacci 矩阵 $\vv F=\begin{pmatrix}1&1\\1&0\end{pmatrix}$ 满足 $\vv F^n=\begin{pmatrix}F_{n+1}&F_n\\F_n&F_{n-1}\end{pmatrix}$（归纳法证明），推出 Cassini 恒等式 $F_{n+1}F_{n-1}-F_n^2=(-1)^n$.
			\item\label{ex:det:B11} 计算 $n$ 阶行列式 $D_n=\begin{vmatrix}1+x^2&x&&\\x&1+x^2&\ddots&\\&\ddots&\ddots&x\\&&x&1+x^2\end{vmatrix}$.
			\item\label{ex:det:B12} 设 $\omega=\mathrm e^{2\pi\mathrm i/3}$. 证明
			\[\begin{vmatrix}a&b&c\\c&a&b\\b&c&a\end{vmatrix}=(a+b+c)(a+\omega b+\omega^2c)(a+\omega^2b+\omega c)=a^3+b^3+c^3-3abc.\]
			\item\label{ex:det:B13} 证明：元素为 $a_{ij}=\min(i,j)$ 的 $n$ 阶行列式等于 $1$.
			\item\label{ex:det:B14} 证明：平面上四点 $(x_i,y_i)$（$i=1,2,3,4$）共圆（或共线），当且仅当 $\begin{vmatrix}x_i^2+y_i^2&x_i&y_i&1\end{vmatrix}_{4\times4}=0$（每行对应一个点）.
						\item\label{ex:det:B15} 设 $\vv A$ 为 $n$ 阶方阵（$n\ge2$）. 证明
			\[\rk\vv A^*=\begin{cases}n,&\rk\vv A=n,\\1,&\rk\vv A=n-1,\\0,&\rk\vv A\le n-2.\end{cases}\]
			（提示：$\vv A^*$ 的元素都是 $\vv A$ 的 $n-1$ 阶子式，用定理~\ref{thm:minorrank}；$\rk\vv A=n-1$ 时考察 $\vv A\vv A^*=(\det\vv A)\vv I$.）
			\item\label{ex:det:B16} 设 $f(x)=\begin{vmatrix}x&1&2&3\\1&x&2&3\\1&2&x&3\\1&2&3&x\end{vmatrix}$. 不展开，求 $f(x)=0$ 的全部根.
		\end{exB}
		
		\Cgroup
		\begin{exC}
			\item\label{ex:det:C1} （行列式的唯一性）设 $D:M_n(\mathbb{F})\to\mathbb{F}$ 对每一行都是线性的、交换两行变号，且 $D(\vv I)=1$. 证明 $D=\det$.（提示：先说明 $D$ 在初等矩阵上的取值，再对任意矩阵用初等变换.）
			\item\label{ex:det:C2} （Cauchy 行列式）证明
			\[\det\Big(\frac1{a_i+b_j}\Big)_{n\times n}=\frac{\prod_{i<j}(a_j-a_i)(b_j-b_i)}{\prod_{i,j}(a_i+b_j)}.\]
			先对 $n=2$ 直接验证，再用``各行减去最后一行''归纳. 取 $a_i=i$，$b_j=j-1$，得到 Hilbert 矩阵 $(1/(i+j-1))$ 的行列式，它为什么在数值计算中臭名昭著？
			\item\label{ex:det:C3} （排列的符号是同态）对排列 $\sigma$，令 $\vv P_\sigma$ 为把 $\vv e_j$ 映到 $\vv e_{\sigma(j)}$ 的置换矩阵. 证明 $\vv P_\sigma\vv P_\tau=\vv P_{\sigma\tau}$，$\det\vv P_\sigma=(-1)^{\tau(\sigma)}$，并由行列式的乘积公式推出 $\operatorname{sgn}(\sigma\tau)=\operatorname{sgn}\sigma\cdot\operatorname{sgn}\tau$.
			\item\label{ex:det:C4} （矩阵—树定理）对一个连通无向图，令 $\vv D$ 为各顶点度数构成的对角矩阵，$\vv A$ 为邻接矩阵，$\vv L=\vv D-\vv A$ 称为 Laplace 矩阵. 定理说：图的生成树个数等于 $\vv L$ 的任一个 $(n-1)$ 阶主子式. 对 $4$ 个顶点的完全图验证它有 $16$ 棵生成树，对一个 $4$ 个顶点的环验证它有 $4$ 棵生成树.
			\item\label{ex:det:C5} 证明：若 $n$ 阶矩阵的元素全为 $\pm1$，则它的行列式能被 $2^{n-1}$ 整除. 求 $3$ 阶此类行列式的最大值.
			\item\label{ex:det:C6} （Jacobi 公式的雏形）证明 $\det(\vv I+t\vv B)=1+t\tr\vv B+(\text{$t^2$ 及更高次项})$. 由此说明：若 $\vv A(t)$ 的元素都是 $t$ 的可导函数，且 $\vv A(0)=\vv I$，则 $\frac{\mathrm d}{\mathrm dt}\det\vv A(t)\big|_{t=0}=\tr\vv A'(0)$.
			\item\label{ex:det:C7} （Sylvester 结式）设 $f=a_0x^2+a_1x+a_2$，$g=b_0x^2+b_1x+b_2$（$a_0b_0\ne0$）. 证明 $f,g$ 有公共根当且仅当
			\[\begin{vmatrix}a_0&a_1&a_2&0\\0&a_0&a_1&a_2\\b_0&b_1&b_2&0\\0&b_0&b_1&b_2\end{vmatrix}=0.\]
			（提示：有公共根当且仅当存在次数 $\le1$ 的不全为零的 $u,v$ 使 $uf+vg=0$，与第七章的 Bézout 等式对照.）
			\item\label{ex:det:C8} （Dodgson 凝聚法）对 $3$ 阶矩阵 $\vv A$，记 $M_{ij}$ 为余子式，$c=a_{22}$. 验证 $\det\vv A\cdot c=M_{11}M_{33}-M_{13}M_{31}$. 它允许只计算 $2$ 阶行列式就得到高阶行列式. 《爱丽丝漫游奇境》的作者刘易斯·卡罗尔（即 Dodgson）于 1866 年发表了这一方法.
			\item\label{ex:det:C9} 设 $\vv A$ 是 $n$ 阶实反对称矩阵. 证明 $\det(\vv I+\vv A)\ge1$.（提示：$\det(\vv I+\vv A)=\det(\vv I-\vv A)$，考察 $\det(\vv I+t\vv A)$ 作为 $t$ 的多项式只含偶次项，且 $\vv I+t\vv A$ 对一切实数 $t$ 都可逆.）
		\end{exC}
	}
	
	% =====================================================================
	% 第六章　线性方程组解的结构
	% =====================================================================
	\newcommand{\ExSolutions}{%
		\exsetheading
		\Agroup
		\begin{exA}
			\item\label{ex:sol:A1} 求下列齐次方程组的基础解系与通解：
			\[\text{(1)}\ \left\{\begin{array}{l}x_1+x_2-x_3-x_4=0,\\2x_1-5x_2+3x_3+2x_4=0,\\7x_1-7x_2+3x_3+x_4=0;\end{array}\right.\qquad
			\text{(2)}\ \left\{\begin{array}{l}x_1+2x_2+x_3-x_4=0,\\3x_1+6x_2-x_3-3x_4=0,\\5x_1+10x_2+x_3-5x_4=0.\end{array}\right.\]
			\item\label{ex:sol:A2} 求下列方程组的通解，写成``特解 $+$ 基础解系的线性组合''：
			\[\text{(1)}\ \left\{\begin{array}{l}x_1+x_2-3x_3-x_4=1,\\3x_1-x_2-3x_3+4x_4=4,\\x_1+5x_2-9x_3-8x_4=0;\end{array}\right.\qquad
			\text{(2)}\ \left\{\begin{array}{l}2x_1+x_2-x_3+x_4=1,\\4x_1+2x_2-2x_3+x_4=2,\\2x_1+x_2-x_3-x_4=1.\end{array}\right.\]
			\item\label{ex:sol:A3} 设 $\vv A$ 是 $4\times5$ 矩阵，$\rk\vv A=3$. (1) $\vv{Ax}=\zero$ 的基础解系含几个向量？(2) $\vv{Ax}=\vv b$ 是否一定有解？(3) 若 $\rk\vv A=4$ 呢？
			\item\label{ex:sol:A4} 不解方程，用秩判断下列方程组是否有解：(1) $x_1+x_2=1$，$x_1-x_2=1$，$2x_1+3x_2=2$；(2) $x_1+x_2+x_3=1$，$2x_1+2x_2+2x_3=3$.
			\item\label{ex:sol:A5} 设 $\vv\eta_1,\vv\eta_2$ 是 $\vv{Ax}=\vv b$（$\vv b\ne\zero$）的两个解. 下列向量中哪些是 $\vv{Ax}=\vv b$ 的解，哪些是 $\vv{Ax}=\zero$ 的解？
			$\vv\eta_1+\vv\eta_2$，$2\vv\eta_1-\vv\eta_2$，$\frac12(\vv\eta_1+\vv\eta_2)$，$\vv\eta_1-\vv\eta_2$，$3\vv\eta_1-2\vv\eta_2$.
			\item\label{ex:sol:A6} 设 $3$ 元方程组 $\vv{Ax}=\vv b$ 的系数矩阵秩为 $2$，$\vv\eta_1=(1,2,3)^\trans$ 与 $\vv\eta_2=(1,-2,-1)^\trans$ 是它的两个解. 求通解.
			\item\label{ex:sol:A7} 写出一个以 $\vv x=(1,0,2)^\trans+k(1,1,-1)^\trans$（$k$ 任意）为全部解的线性方程组.
			\item\label{ex:sol:A8} 判断三个平面 $x+y+z=1$，$x+2y+3z=2$，$2x+3y+4z=4$ 的位置关系；若把最后一个常数改为 $3$ 呢？
			\item\label{ex:sol:A9} 举例说明：$\vv{Ax}=\zero$ 只有零解时，$\vv{Ax}=\vv b$ 可能无解；但若有解，解必唯一.
			\item\label{ex:sol:A10} 解矩阵方程 $\begin{pmatrix}1&2&3\\2&2&1\\3&4&3\end{pmatrix}\vv X=\begin{pmatrix}2&5\\3&1\\4&3\end{pmatrix}$.
			\item\label{ex:sol:A11} 设 $\vv A=\begin{pmatrix}1&1\\2&2\end{pmatrix}$. 求全部满足 $\vv{AX}=\vv O$ 的二阶矩阵 $\vv X$，它们构成的空间是几维的？
			\item\label{ex:sol:A12} 设 $\vv A=\begin{pmatrix}1&2&-2\\4&t&3\\3&-1&1\end{pmatrix}$，且存在三阶非零矩阵 $\vv B$ 使 $\vv{AB}=\vv O$. 求 $t$，并求 $\rk\vv B$ 的可能值.
			\item\label{ex:sol:A13} 求方程 $x_1+x_2+x_3=1$ 的通解，并写出它的三个解，使得通解中的每个解都是这三个解的仿射组合.
			\item\label{ex:sol:A14} 把空间直线 $\{x+y+z=1,\ x-y=0\}$ 写成参数式；反过来，把直线 $\vv x=(1,2,0)^\trans+t(1,0,1)^\trans$ 写成两个平面的交线.
			\item\label{ex:sol:A15} 判断空间直线 $\vv x=(1,0,0)^\trans+s(1,1,0)^\trans$ 与 $\vv x=(0,1,1)^\trans+t(0,1,-1)^\trans$ 是否相交.（列出关于 $s,t$ 的方程组.）
			\item\label{ex:sol:A16} 物流网络：三个仓库 $A,B,C$ 之间有线路 $A\to B$（$x_1$），$B\to C$（$x_2$），$A\to C$（$x_3$）. $A$ 每天净发出 $30$ 吨，$C$ 每天净收到 $30$ 吨，$B$ 不囤货. 求全部方案. 哪些量（如 $x_1+x_3$，$x_1-x_2$，$x_3$）能被方程组唯一确定？用可确定性判据说明.
			\item\label{ex:sol:A17} 求 $\vv A=\begin{pmatrix}2&1&0\\4&3&1\\2&3&3\end{pmatrix}$ 的 LU 分解.
			\item\label{ex:sol:A18} 设 $\vv A$ 是 $m\times n$ 矩阵，$\rk\vv A=m<n$. 证明 $\vv{Ax}=\vv b$ 对每个 $\vv b$ 都有无穷多解.
			\item\label{ex:sol:A19} 求经过点 $(1,1,1)$，$(2,0,1)$，$(0,2,3)$ 的平面方程（用齐次方程组求法向量）.
			\item\label{ex:sol:A20} 设 $\vv A=\begin{pmatrix}1&1\\1&1.0001\end{pmatrix}$. 分别求 $\vv{Ax}=(2,2)^\trans$ 与 $\vv{Ax}=(2,2.0001)^\trans$ 的解，计算右端的相对变化与解的相对变化之比.
		\end{exA}
		
		\Bgroup
		\begin{exB}
			\item\label{ex:sol:B1} 求两个齐次方程组 $\{x_1+x_2=0,\ x_2-x_4=0\}$ 与 $\{x_1-x_2+x_3=0,\ x_2-x_3+x_4=0\}$ 的公共非零解.
			\item\label{ex:sol:B2} 设 $\vv A$ 是实矩阵. 证明 $\vv{Ax}=\zero$ 与 $\vv A^\trans\vv{Ax}=\zero$ 同解，从而 $\rk(\vv A^\trans\vv A)=\rk\vv A$.（提示：$\vv x^\trans\vv A^\trans\vv{Ax}=\|\vv{Ax}\|^2$ 在第九章才有名字，但这里只需``实数的平方和为零则每个都为零''.）
			\item\label{ex:sol:B3} 证明 $\rk(\vv A+\vv B)\le\rk(\vv A\mid\vv B)\le\rk\vv A+\rk\vv B$.
			\item\label{ex:sol:B4} 设 $\vv A$ 为 $n$ 阶方阵，$\vv A^2=\vv I$. 证明 $\rk(\vv A+\vv I)+\rk(\vv A-\vv I)=n$.
			\item\label{ex:sol:B5} 证明 $\rk(\vv{AB})=\rk\vv B$ 当且仅当 $\ker\vv A\cap\imm\vv B=\{\zero\}$.
			\item\label{ex:sol:B6} 设 $\vv\eta^*$ 是 $\vv{Ax}=\vv b$（$\vv b\ne\zero$）的一个解，$\vv\xi_1,\ldots,\vv\xi_k$ 是对应齐次方程组的基础解系. 证明 $\vv\eta^*,\vv\eta^*+\vv\xi_1,\ldots,\vv\eta^*+\vv\xi_k$ 线性无关.
			\item\label{ex:sol:B7} 一张 $3\times3$ 的统计表只公开了三个行和与三个列和. 至少还要再公开几个单元格，才能确定全表？给出一种``公开哪几个格''的选法，并说明为什么某些同样个数的选法却不行.
			\item\label{ex:sol:B8} 对 $\vv A=\begin{pmatrix}1&2&1\\2&4&3\\3&6&4\end{pmatrix}$，求 $\vv b=(b_1,b_2,b_3)^\trans$ 应满足的条件，使 $\vv{Ax}=\vv b$ 有解.（提示：找出把各方程组合成 $0=\cdots$ 的系数.）
			\item\label{ex:sol:B9} 设 $\vv A$ 为 $n$ 阶方阵，$\rk\vv A=n-1$，且 $\vv A$ 的每行元素之和为零. 求 $\vv{Ax}=\zero$ 的通解.
			\item\label{ex:sol:B10} 设 $\vv A,\vv B$ 为 $n$ 阶方阵，$\vv{AB}=\vv O$，$\vv A+\vv B$ 可逆. 证明 $\rk\vv A+\rk\vv B=n$.
			\item\label{ex:sol:B11} 证明：两个 $m\times n$ 齐次方程组 $\vv{Ax}=\zero$ 与 $\vv{Bx}=\zero$ 同解，当且仅当 $\vv A$ 与 $\vv B$ 的行向量组等价.
			\item\label{ex:sol:B12} 设 $\vv A$ 为 $m\times n$ 矩阵. 证明：$\vv{Ax}=\vv b$ 对一切 $\vv b$ 有解，当且仅当存在 $n\times m$ 矩阵 $\vv C$ 使 $\vv{AC}=\vv I_m$；这时 $\vv C\vv b$ 就是一个解.
			\item\label{ex:sol:B13} 设 $\vv A$ 为 $n$ 阶方阵（$n\ge2$），$\rk\vv A=n-1$，$\vv\xi$ 是 $\vv{Ax}=\zero$ 的非零解. 证明 $\vv A^*$ 的每一列都是 $\vv\xi$ 的倍数.
			\item\label{ex:sol:B14} 用可确定性判据回答：方程组 $x_1+x_2+x_3=6$，$x_2+x_3+x_4=9$，$x_3+x_4+x_5=12$（例如三段相邻时间窗口内的观测总量）能确定 $x_1-x_4$ 吗？$x_1+x_2+x_3+x_4+x_5$ 呢？$x_1-x_4+x_5$ 呢？
			\item\label{ex:sol:B15} 设 $\vv A$ 为 $m\times n$ 矩阵，$\vv B$ 为 $n\times p$ 矩阵. 证明 $\dim\ker(\vv{AB})\le\dim\ker\vv A+\dim\ker\vv B$，并由此重新推出 Sylvester 不等式.
		\end{exB}
		
		\Cgroup
		\begin{exC}
			\item\label{ex:sol:C1} （Farkas 引理的二维形式）设 $\vv a_1,\vv a_2\in\mathbb{R}^2$ 线性无关，$\vv b\in\mathbb{R}^2$. 证明：下面两件事恰有一件成立：(1) 存在 $x_1,x_2\ge0$ 使 $x_1\vv a_1+x_2\vv a_2=\vv b$；(2) 存在 $\vv y$ 使 $\vv y^\trans\vv a_1\ge0$，$\vv y^\trans\vv a_2\ge0$，$\vv y^\trans\vv b<0$. 画图解释. 这是``无套利定价''与线性规划对偶理论的出发点.
			\item\label{ex:sol:C2} 设 $\vv A$ 是 $n$ 阶整数矩阵. 证明：若对任意整数向量 $\vv b$，$\vv{Ax}=\vv b$ 都有整数解，则 $\det\vv A=\pm1$.
			\item\label{ex:sol:C3} 证明：方阵 $\vv A$ 不用行交换就能做 LU 分解（$\vv L$ 单位下三角，$\vv U$ 上三角且可逆），当且仅当 $\vv A$ 的各阶顺序主子式都不为零；这时分解唯一.
			\item\label{ex:sol:C4} 设 $\vv A$ 是三对角矩阵（只有主对角线和相邻两条斜线上可能有非零元）. 证明：若它有 LU 分解，则 $\vv L,\vv U$ 都只有两条非零斜线；因此解三对角方程组只需与 $n$ 成正比的运算量. 这种方法称为``追赶法''.
			\item\label{ex:sol:C5} 用分块初等变换
			\[\begin{pmatrix}\vv I_n&\vv O\\ -\vv A&\vv I_m\end{pmatrix}\begin{pmatrix}\vv I_n&\vv B\\ \vv A&\vv O\end{pmatrix}\begin{pmatrix}\vv I_n&-\vv B\\ \vv O&\vv I_p\end{pmatrix}=\begin{pmatrix}\vv I_n&\vv O\\ \vv O&-\vv{AB}\end{pmatrix}\]
			给出 Sylvester 不等式 $\rk(\vv{AB})\ge\rk\vv A+\rk\vv B-n$ 的另一个证明.（提示：用第三章章末习题~\ref{ex:mat:B5} 的分块秩不等式.）
			\item\label{ex:sol:C6} 证明：$\mathbb{F}^n$ 中任意 $k$ 个线性流形（不必过原点）的交，要么是空集，要么仍是线性流形. 给出两个二维线性流形在 $\mathbb{R}^3$ 中交为空集的例子，并用秩的语言刻画这种情形.
			\item\label{ex:sol:C7} （列昂惕夫价格模型）设 $\vv A$ 为直接消耗系数矩阵，$\vv v$ 为各部门单位产品的增加值（工资、利润等）. 价格 $\vv p$ 满足``单位价格 $=$ 单位投入成本 $+$ 单位增加值''：$\vv p=\vv A^\trans\vv p+\vv v$. 说明它是第一章产量模型 $\vv x=\vv{Ax}+\vv d$ 的``转置''；证明在两个模型都可解时有 $\vv p^\trans\vv d=\vv v^\trans\vv x$，并解释其经济含义（国民收入的两种核算方法：支出法与收入法）.
			\item\label{ex:sol:C8} 设 $\vv A$ 为 $n$ 阶方阵. 证明存在正整数 $k$ 使 $\rk\vv A^k=\rk\vv A^{k+1}=\rk\vv A^{k+2}=\cdots$，且这样的最小的 $k$ 不超过 $n$.
		\end{exC}
	}
	
	% =====================================================================
	% 第七章　多项式
	% =====================================================================
	\newcommand{\ExPolynomial}{%
		\exsetheading
		\Agroup
		\begin{exA}
			\item\label{ex:poly:A1} 求商与余式：(1) $x^5-1$ 除以 $x^2+x+1$；(2) $2x^4-3x^3+4x^2-5x+6$ 除以 $x-2$（用综合除法）.
			\item\label{ex:poly:A2} 求 $x^{100}$ 分别除以 $x^2-1$，$x^2+1$，$(x-1)^2$ 所得的余式.
			\item\label{ex:poly:A3} 求 $a,b$，使 $x^2+x+1$ 整除 $x^4+ax^2+b$.（提示：设 $\omega$ 是 $x^2+x+1$ 的根，则 $\omega^3=1$.）
			\item\label{ex:poly:A4} 用辗转相除法求最大公因式：(1) $x^4+3x^3-x^2-4x-3$ 与 $3x^3+10x^2+2x-3$；(2) $x^3-2x^2+x$ 与 $x^2-1$.
			\item\label{ex:poly:A5} 证明：对任意多项式 $h$，$\gcd(f,g)=\gcd(f,\ g-hf)$. 这正是辗转相除法每一步的依据.
			\item\label{ex:poly:A6} 求 $u,v$ 使 $u(x)(x^2+1)+v(x)(x+1)=1$.
			\item\label{ex:poly:A7} 求 $2x^3-5x^2+x+2$ 的全部有理根，并在 $\mathbb{Q}$ 上把它分解为一次因式之积.
			\item\label{ex:poly:A8} 分别在 $\mathbb{Q}$，$\mathbb{R}$，$\mathbb{C}$ 上分解：(1) $x^4+x^2+1$；(2) $x^6-1$；(3) $x^3-2$.
			\item\label{ex:poly:A9} 判断 $1$ 是 $x^5-5x+4$ 的几重根.
			\item\label{ex:poly:A10} $a$ 取何值时，$x^3-3x+a$ 有重根？求出重根.
			\item\label{ex:poly:A11} 设 $\alpha_1,\alpha_2,\alpha_3$ 是 $x^3-2x^2+3x-4$ 的根. 不解方程，求 $\alpha_1+\alpha_2+\alpha_3$，$\alpha_1^2+\alpha_2^2+\alpha_3^2$，$\frac1{\alpha_1}+\frac1{\alpha_2}+\frac1{\alpha_3}$.
			\item\label{ex:poly:A12} 求次数最低的首一有理系数多项式，使 $1+\sqrt2$ 与 $\mathrm i$ 都是它的根.
			\item\label{ex:poly:A13} 求一个三次实系数首一多项式，以 $2$ 和 $1+\mathrm i$ 为根.
			\item\label{ex:poly:A14} 设 $f(x)=x^2-2x+1$，$\vv A=\begin{pmatrix}1&2\\0&1\end{pmatrix}$，计算 $f(\vv A)$；再设 $g(x)=x^3-x$，计算 $g(\vv A)$.
			\item\label{ex:poly:A15} 求过点 $(0,1)$，$(1,0)$，$(2,3)$，$(3,10)$ 的三次以下的插值多项式（用 Lagrange 公式）.
			\item\label{ex:poly:A16} 证明 $(x-1)\mid x^n-1$，并写出商式. 一般地，对任意 $a$，$(x-a)\mid x^n-a^n$.
			\item\label{ex:poly:A17} 把 $f(x)=x^4-2x^3+3$ 按 $x-1$ 的幂展开，即写成 $\sum c_k(x-1)^k$.（提示：反复用 $x-1$ 作综合除法，或用 $c_k=f^{(k)}(1)/k!$.）
			\item\label{ex:poly:A18} 多项式 $f$ 除以 $x-1$ 余 $3$，除以 $x-2$ 余 $5$. 求 $f$ 除以 $(x-1)(x-2)$ 的余式.
			\item\label{ex:poly:A19} 一个项目现在投入 $100$ 万元，一年后回收 $55$ 万元，两年后回收 $60.5$ 万元. 写出内部收益率满足的方程并精确求解.
			\item\label{ex:poly:A20} 证明：若 $f\mid g$ 且 $g\mid f$，则 $f=cg$（$c$ 为非零常数）.
			\item\label{ex:poly:A21} 设 $\gcd(f,g)=d$，$f=df_1$，$g=dg_1$. 证明 $\gcd(f_1,g_1)=1$.
		\end{exA}
		
		\Bgroup
		\begin{exB}
			\item\label{ex:poly:B1} 证明：若 $\gcd(f,g)=1$ 且 $\gcd(f,h)=1$，则 $\gcd(f,gh)=1$.
			\item\label{ex:poly:B2} 证明：$\gcd(x^m-1,\ x^n-1)=x^{\gcd(m,n)}-1$.（提示：模仿整数的辗转相除，先证 $x^m-1=x^{m-n}(x^n-1)+(x^{m-n}-1)$.）
			\item\label{ex:poly:B3} 设 $f=(x-1)^3(x+2)^2$. 计算 $f/\gcd(f,f')$，说明它恰好是``去掉重复因式''的结果. 一般地证明：若 $f=c\,p_1^{k_1}\cdots p_m^{k_m}$ 是标准分解，则 $f/\gcd(f,f')=c\,p_1\cdots p_m$.
			\item\label{ex:poly:B4} 证明：若 $p$ 不可约且 $p\mid f^2$，则 $p\mid f$. 举例说明``不可约''不能去掉.
			\item\label{ex:poly:B5} 设 $\ell_1,\ldots,\ell_n$ 是关于互异节点 $x_1,\ldots,x_n$ 的 Lagrange 基函数. 证明 $\sum_i\ell_i(x)=1$，$\sum_ix_i\ell_i(x)=x$，一般地 $\sum_ix_i^k\ell_i(x)=x^k$（$0\le k\le n-1$）.
			\item\label{ex:poly:B6} 证明：若复系数多项式 $f$ 满足 $f(x+1)=f(x)$ 对一切 $x$ 成立，则 $f$ 是常数.
			\item\label{ex:poly:B7} 证明 $x^3+x+1$ 在 $\mathbb{Q}$ 上不可约. 为什么对三次多项式只需检验有无有理根？四次多项式可以这样检验吗？
			\item\label{ex:poly:B8} 求 $x^n$ 除以 $x^2-3x+2$ 的余式.
			\item\label{ex:poly:B9} （Hermite 插值）求次数最低的多项式 $f$，使 $f(1)=1$，$f'(1)=2$，$f(2)=4$.（提示：设 $f=1+2(x-1)+c(x-1)^2$.）
			\item\label{ex:poly:B10} 设 $f$ 是有理系数多项式，$a,b\in\mathbb{Q}$，$\sqrt2\notin\mathbb{Q}$. 证明：若 $a+b\sqrt2$ 是 $f$ 的根，则 $a-b\sqrt2$ 也是 $f$ 的根.（提示：用 $(x-a)^2-2b^2$ 除 $f$.）
			\item\label{ex:poly:B11} 设 $f$ 是整系数多项式，$f(0)$ 与 $f(1)$ 都是奇数. 证明 $f$ 没有整数根.
			\item\label{ex:poly:B12} 求最小多项式意义下的``零化''：分别求出次数最低的首一多项式 $g$，使 $g(\vv A)=\vv O$，其中 $\vv A=\begin{pmatrix}0&-1\\1&0\end{pmatrix}$，$\vv A=\begin{pmatrix}2&0\\0&2\end{pmatrix}$，$\vv A=\begin{pmatrix}2&1\\0&2\end{pmatrix}$.
			\item\label{ex:poly:B13} 证明：现金流 $-C_0,\ C_1,\ldots,C_n$（$C_0>0$，$C_i\ge0$ 且不全为零）的净现值 $N(v)=-C_0+C_1v+\cdots+C_nv^n$ 在 $v>0$ 时严格递增，从而 $N(v)=0$ 至多一个正根；当 $C_1+\cdots+C_n>C_0$ 时这个正根落在 $(0,1)$ 中，即 IRR 唯一且为正.
			\item\label{ex:poly:B14} 设 $f$ 是 $n$ 次多项式，在 $n+1$ 个整数点 $0,1,\ldots,n$ 上都取整数值. 证明 $f$ 在一切整数点上都取整数值.（提示：用基 $\binom x0,\binom x1,\ldots,\binom xn$，其中 $\binom xk=\frac{x(x-1)\cdots(x-k+1)}{k!}$.）
			\item\label{ex:poly:B15} 证明 Lagrange 插值的``唯一性''也可以由线性代数推出：求值映射 $\mathbb{F}[x]_n\to\mathbb{F}^n$，$f\mapsto(f(x_1),\ldots,f(x_n))^\trans$ 是单射，从而是同构.
			\item\label{ex:poly:B16} 设 $p$ 是 $\mathbb{F}$ 上的不可约多项式，$f\in\mathbb{F}[x]$. 证明：$\gcd(p,f)=1$ 或 $p\mid f$，二者必居其一.
		\end{exB}
		
		\Cgroup
		\begin{exC}
			\item\label{ex:poly:C1} （Eisenstein 判别法）设 $f=a_nx^n+\cdots+a_0$ 是整系数多项式，$p$ 是素数，$p\nmid a_n$，$p\mid a_i$（$i<n$），$p^2\nmid a_0$. 证明 $f$ 在 $\mathbb{Q}$ 上不可约（可以先承认 Gauss 引理：整系数多项式若能分解为两个低次有理系数多项式之积，则也能分解为两个低次整系数多项式之积）. 由此说明 $x^n-2$ 对一切 $n$ 都不可约，所以 $\mathbb{Q}[x]$ 中有任意高次的不可约多项式.
			\item\label{ex:poly:C2} 设 $p$ 是素数. 证明分圆多项式 $\Phi_p(x)=x^{p-1}+x^{p-2}+\cdots+1$ 在 $\mathbb{Q}$ 上不可约.（提示：对 $\Phi_p(x+1)$ 用 Eisenstein 判别法.）
			\item\label{ex:poly:C3} （Newton 恒等式）记 $s_k=\alpha_1^k+\alpha_2^k+\alpha_3^k$，$\sigma_1,\sigma_2,\sigma_3$ 为初等对称多项式. 证明 $s_1=\sigma_1$，$s_2=\sigma_1s_1-2\sigma_2$，$s_3=\sigma_1s_2-\sigma_2s_1+3\sigma_3$，并求 $x^3-2x^2+3x-4$ 的根的立方和.
			\item\label{ex:poly:C4} （三次方程的判别式）设 $x^3+px+q$ 的三个复根为 $\alpha_1,\alpha_2,\alpha_3$，$\Delta=\prod_{i<j}(\alpha_i-\alpha_j)^2$. 证明 $\Delta=-4p^3-27q^2$，并说明（$p,q$ 为实数时）：$\Delta>0$ 时三个不同实根，$\Delta<0$ 时一实二虚，$\Delta=0$ 时有重根.
			\item\label{ex:poly:C5} （Chebyshev 多项式）证明存在 $n$ 次多项式 $T_n$ 使 $\cos n\theta=T_n(\cos\theta)$，且 $T_{n+1}=2xT_n-T_{n-1}$. 求 $T_n$ 的全部根. 由此证明：$\cos20^\circ$ 是 $8x^3-6x-1$ 的根，并验证这个三次多项式在 $\mathbb{Q}$ 上不可约. （结合``尺规作图只能得到经若干次开平方得到的数''，这说明 $60^\circ$ 角不能用尺规三等分.）
			\item\label{ex:poly:C6} 证明：在 $\mathbb{R}$ 上处处非负的实系数多项式，都能写成两个实系数多项式的平方和.（提示：利用实根的重数必为偶数，以及 $|a+b\mathrm i|^2=a^2+b^2$ 与 $(a^2+b^2)(c^2+d^2)=(ac-bd)^2+(ad+bc)^2$.）
			\item\label{ex:poly:C7} 设 $f\in\mathbb{Q}[x]$ 不可约，$g\in\mathbb{Q}[x]$ 与 $f$ 有一个公共复根. 证明 $f\mid g$. 由此说明：一个复数若是某个有理系数多项式的根，则以它为根的次数最低的首一有理系数多项式（称为它的\textbf{极小多项式}）唯一，且不可约.
			\item\label{ex:poly:C8} （$\mathbb{F}_2$ 上的多项式）在系数按模 $2$ 运算的多项式中，证明 $x^2+x+1$ 是唯一的二次不可约多项式，并找出全部三次不可约多项式. 这类多项式是循环冗余校验（CRC）与密码学的基础.
			\item\label{ex:poly:C9} 设 $f$ 是实系数多项式，它的系数序列的``变号次数''为 $V$（忽略零系数）. Descartes 符号法则说：$f$ 的正根个数（按重数计）不超过 $V$，且与 $V$ 的差是偶数. 用它重新解释第七章``内部收益率''建模实验中两个项目的结论.
		\end{exC}
	}
	
	% =====================================================================
	% 第八章　特征值与特征向量
	% =====================================================================
	\newcommand{\ExEigen}{%
		\exsetheading
		\Agroup
		\begin{exA}
			\item\label{ex:eig:A1} 求下列矩阵的特征值与特征向量（在 $\mathbb{C}$ 上考虑）：
			\[\begin{pmatrix}3&1\\0&2\end{pmatrix},\quad\begin{pmatrix}1&2\\2&1\end{pmatrix},\quad\begin{pmatrix}0&1\\-1&0\end{pmatrix},\quad\begin{pmatrix}4&-2\\1&1\end{pmatrix},\]\[\begin{pmatrix}2&0&0\\0&1&1\\0&1&1\end{pmatrix},\quad\begin{pmatrix}1&2&2\\2&1&2\\2&2&1\end{pmatrix}.\]
			\item\label{ex:eig:A2} 判断 A1 中各矩阵能否在 $\mathbb{R}$ 上对角化；能的，求可逆矩阵 $\vv P$ 使 $\vv P^{-1}\vv{AP}$ 为对角矩阵.
			\item\label{ex:eig:A3} 判断下列矩阵能否对角化：$\begin{pmatrix}1&1\\0&2\end{pmatrix}$，$\begin{pmatrix}2&1\\0&2\end{pmatrix}$，$\begin{pmatrix}2&0&0\\0&2&0\\1&0&2\end{pmatrix}$，$\begin{pmatrix}1&1&0\\0&2&0\\0&0&1\end{pmatrix}$.
			\item\label{ex:eig:A4} 设 $3$ 阶方阵 $\vv A$ 的特征值为 $1,2,3$. 求 $\det\vv A$，$\tr\vv A$，$\vv A^{-1}$ 的特征值，$\vv A^2-2\vv A+\vv I$ 的特征值，$\det(\vv A^*)$，$\det(\vv A^2+\vv I)$.
			\item\label{ex:eig:A5} 已知二阶方阵的迹为 $5$，行列式为 $6$，求它的特征值. 若迹为 $2$、行列式为 $5$ 呢？
			\item\label{ex:eig:A6} 证明：$\vv A$ 与 $\vv A^\trans$ 有相同的特征多项式. 它们的特征向量是否相同？举例说明.
			\item\label{ex:eig:A7} 证明：三角矩阵的特征值恰为其对角元.
			\item\label{ex:eig:A8} 证明：若 $\vv A^k=\vv O$，则 $\vv A$ 的特征值全为 $0$；若 $\vv A^2=\vv A$，则特征值只能是 $0$ 或 $1$；若 $\vv A^2=\vv I$，则特征值只能是 $\pm1$.
			\item\label{ex:eig:A9} 设 $\vv A=\begin{pmatrix}4&-2\\1&1\end{pmatrix}$. 利用对角化计算 $\vv A^n$.
			\item\label{ex:eig:A10} 设品牌转移矩阵为 $\vv P=\begin{pmatrix}0.9&0.2\\0.1&0.8\end{pmatrix}$. 求稳态分布，并利用对角化求 $\lim_{n\to\infty}\vv P^n$. 解释极限矩阵两列相同的含义.
			\item\label{ex:eig:A11} 已知 $\vv\xi=(1,1,-1)^\trans$ 是 $\vv A=\begin{pmatrix}2&-1&2\\5&a&3\\-1&b&-2\end{pmatrix}$ 的特征向量. 求 $a,b$ 与对应的特征值.
			\item\label{ex:eig:A12} 设 $\vv A=\begin{pmatrix}0&0&1\\1&1&x\\1&0&0\end{pmatrix}$. $x$ 取何值时 $\vv A$ 可对角化？
			\item\label{ex:eig:A13} 对二阶方阵直接验证 Cayley–Hamilton 定理：$\vv A^2-(\tr\vv A)\vv A+(\det\vv A)\vv I=\vv O$.
			\item\label{ex:eig:A14} 求下列矩阵的最小多项式：$\diag(1,1,2)$，$\begin{pmatrix}1&1&0\\0&1&0\\0&0&2\end{pmatrix}$，$\begin{pmatrix}0&1&0\\0&0&1\\0&0&0\end{pmatrix}$，$\begin{pmatrix}1&1&1\\1&1&1\\1&1&1\end{pmatrix}$.
			\item\label{ex:eig:A15} 判断 $\begin{pmatrix}1&1\\0&1\end{pmatrix}$ 与 $\begin{pmatrix}1&0\\1&1\end{pmatrix}$ 是否相似；若相似，求出 $\vv P$.
			\item\label{ex:eig:A16} 设 $\vv A$ 是上三角矩阵. 证明 $\spn(\vv e_1)$，$\spn(\vv e_1,\vv e_2)$，……都是 $\vv A$ 的不变子空间.
			\item\label{ex:eig:A17} 解差分方程 $a_{k+1}=a_k+2a_{k-1}$，$a_0=0$，$a_1=1$.
			\item\label{ex:eig:A18} （种群模型）某种群分为幼年与成年两组，第 $k$ 年的数量记为 $\vv x_k=(y_k,z_k)^\trans$. 每只成年个体每年平均产下 $1.1$ 只幼体，幼体有 $60\%$ 存活并长为成年，成年个体有一半活到下一年，即
			$y_{k+1}=1.1z_k$，$z_{k+1}=0.6y_k+0.5z_k$. 求矩阵的特征值与特征向量，说明种群的长期增长率和幼体与成年个体的稳定比例.
			\item\label{ex:eig:A19} 证明：实二阶方阵 $\vv A$ 没有实特征值，当且仅当 $(\tr\vv A)^2<4\det\vv A$.
			\item\label{ex:eig:A20} 求平面上向直线 $y=2x$ 的正投影、关于该直线的反射这两个变换的特征值与特征向量（不必写出矩阵，从几何上直接回答），再写出矩阵加以验证.
		\end{exA}
		
		\Bgroup
		\begin{exB}
			\item\label{ex:eig:B1} 设 $\vv u,\vv v\in\mathbb{R}^n$ 非零，$\vv A=\vv u\vv v^\trans$. 求 $\vv A$ 的全部特征值及其重数，并证明：$\vv A$ 可对角化当且仅当 $\vv v^\trans\vv u\ne0$.
			\item\label{ex:eig:B2} 设 $\vv A$ 有 $n$ 个互异的特征值，$\vv{AB}=\vv{BA}$. 证明 $\vv A$ 的特征向量都是 $\vv B$ 的特征向量，从而 $\vv B$ 也可对角化；进一步证明 $\vv B$ 是 $\vv A$ 的多项式.
			\item\label{ex:eig:B3} 证明：若 $\vv A$ 可对角化，则 $\vv A^k$ 与 $f(\vv A)$ 都可对角化. 逆命题成立吗？
			\item\label{ex:eig:B4} 设 $\vv A$ 可对角化，$W$ 是 $\vv A$ 的不变子空间. 证明 $\vv A|_W$ 也可对角化.（提示：最小多项式.）
			\item\label{ex:eig:B5} 证明：满足 $\vv A^2=-\vv I$ 的二阶实方阵都与 $\begin{pmatrix}0&-1\\1&0\end{pmatrix}$ 相似.
			\item\label{ex:eig:B6} （Gershgorin 圆盘定理）证明：方阵 $\vv A$ 的每个特征值都落在某个圆盘 $\{z\in\mathbb{C}:|z-a_{ii}|\le\sum_{j\ne i}|a_{ij}|\}$ 中. 由此证明：严格对角占优的矩阵（$|a_{ii}|>\sum_{j\ne i}|a_{ij}|$）可逆.
			\item\label{ex:eig:B7} 设 $\vv A$ 的每行元素之和都是 $s$. 证明 $s$ 是 $\vv A$ 的特征值. 若 $\vv A$ 可逆，求 $\vv A^{-1}$ 每行元素之和.
			\item\label{ex:eig:B8} 设 $\vv A=\begin{pmatrix}a&b&c\\c&a&b\\b&c&a\end{pmatrix}$（循环矩阵），$\omega=\mathrm e^{2\pi\mathrm i/3}$. 证明 $(1,\omega^k,\omega^{2k})^\trans$（$k=0,1,2$）都是 $\vv A$ 的特征向量，并求对应的特征值. 由此重新得到第五章章末习题~\ref{ex:det:B12} 的行列式公式.
			\item\label{ex:eig:B9} 解差分方程 $a_{k+1}=4a_k-4a_{k-1}$，$a_0=1$，$a_1=4$. 解释通项中为什么出现 $k\cdot2^k$.
			\item\label{ex:eig:B10} 证明 $\tr(\vv A^k)=\lambda_1^k+\cdots+\lambda_n^k$. 若三阶方阵满足 $\tr\vv A=0$，$\tr\vv A^2=6$，$\det\vv A=0$，求它的特征值.
			\item\label{ex:eig:B11} 证明：若 $\vv A^3=\vv A$，则 $\vv A$ 可对角化，且特征值在 $\{0,1,-1\}$ 中.
			\item\label{ex:eig:B12} $a,b,c$ 满足什么条件时，$\begin{pmatrix}2&0&0\\a&2&0\\b&c&1\end{pmatrix}$ 可对角化？
			\item\label{ex:eig:B13} 三个网页之间的链接为：网页 $1$ 链接到 $2,3$，网页 $2$ 链接到 $3$，网页 $3$ 链接到 $1$.（与判例不同，网页之间可以``互相引用''，形成环路.） 取 $d=1$ 写出转移矩阵，手算它的稳态分布.
			\item\label{ex:eig:B14} （矩阵指数）对可对角化矩阵 $\vv A=\vv P\vv D\vv P^{-1}$，定义 $\mathrm e^{t\vv A}=\vv P\diag(\mathrm e^{t\lambda_1},\ldots,\mathrm e^{t\lambda_n})\vv P^{-1}$. 对 $\vv A=\begin{pmatrix}0&1\\1&0\end{pmatrix}$ 计算 $\mathrm e^{t\vv A}$，并验证 $\vv x(t)=\mathrm e^{t\vv A}\vv x_0$ 满足微分方程组 $x_1'=x_2$，$x_2'=x_1$.
			\item\label{ex:eig:B15} 三状态马尔可夫链的转移矩阵为 $\begin{pmatrix}0.5&0.25&0.25\\0.25&0.5&0.25\\0.25&0.25&0.5\end{pmatrix}$. 求稳态分布与全部特征值，说明偏离稳态的部分每一步缩小为原来的多少.
			\item\label{ex:eig:B16} 证明：若 $\lambda$ 是 $\vv A$ 的 $k$ 重特征值（代数重数），则 $\rk(\lambda\vv I-\vv A)\ge n-k$.
			\item\label{ex:eig:B17} 对第三章``判例引证网络''中的四份判决，分别取阻尼系数 $d=0.5,\ 0.85,\ 0.95,\ 1$，按第八章建模实验的方法求 PageRank 得分（可用计算机），列表比较各判决的得分与排序. $d$ 越接近 $1$，得分越``集中''在哪些判决上？为什么？
		\end{exB}
		
		\Cgroup
		\begin{exC}
			\item\label{ex:eig:C1} 设 $\vv A$ 为 $m\times n$ 矩阵，$\vv B$ 为 $n\times m$ 矩阵. 证明 $\lambda^n\det(\lambda\vv I_m-\vv{AB})=\lambda^m\det(\lambda\vv I_n-\vv{BA})$. 由此说明 $\vv{AB}$ 与 $\vv{BA}$ 的非零特征值（连同重数）相同.（提示：计算 $\begin{pmatrix}\vv I_m&-\vv A\\ \vv O&\lambda\vv I_n\end{pmatrix}\begin{pmatrix}\lambda\vv I_m&\vv A\\ \vv B&\vv I_n\end{pmatrix}$ 等两种乘积的行列式.）
			\item\label{ex:eig:C2} （Schur 三角化）证明：每个复方阵都相似于一个上三角矩阵. 由此给出 Cayley–Hamilton 定理的另一个证明：对上三角矩阵 $\vv T$，证明 $(\vv T-t_{11}\vv I)\cdots(\vv T-t_{nn}\vv I)=\vv O$.
			\item\label{ex:eig:C3} 求下列矩阵的 Jordan 标准形：
			\[\begin{pmatrix}2&1&0\\0&2&0\\0&0&2\end{pmatrix},\quad\begin{pmatrix}1&1&1\\0&1&1\\0&0&1\end{pmatrix},\quad\begin{pmatrix}0&1&0&0\\0&0&0&0\\0&0&0&1\\0&0&0&0\end{pmatrix}.\]
			特征多项式为 $(\lambda-2)^4$ 的 $4$ 阶复矩阵，在相似的意义下共有几类？
			\item\label{ex:eig:C4} 证明：若一组可对角化的矩阵两两可交换，则它们可以\textbf{同时}对角化（存在同一个 $\vv P$ 使它们都化为对角矩阵）.
			\item\label{ex:eig:C5} 设 $\vv A$ 可对角化. 证明 $\vv A^k\to\vv O$（按元素）当且仅当 $\vv A$ 的所有特征值的模都小于 $1$. 用它说明：投入产出模型中，若直接消耗系数矩阵 $\vv A$ 可对角化且特征值的模都小于 $1$，则 $(\vv I-\vv A)^{-1}=\vv I+\vv A+\vv A^2+\cdots$.
			\item\label{ex:eig:C6} （友矩阵）证明 $\begin{pmatrix}0&0&\cdots&0&-a_0\\1&0&\cdots&0&-a_1\\0&1&\cdots&0&-a_2\\ \vdots&&\ddots&&\vdots\\0&0&\cdots&1&-a_{n-1}\end{pmatrix}$ 的特征多项式为 $\lambda^n+a_{n-1}\lambda^{n-1}+\cdots+a_0$. 由此说明：求多项式的根等价于求矩阵的特征值（数值软件正是这样求多项式的根的）.
			\item\label{ex:eig:C7} （赌徒破产）一位赌徒有 $k$ 元（$0\le k\le N$），每局以概率 $\frac12$ 赢或输 $1$ 元，输光或达到 $N$ 元时停止. 设 $p_k$ 为最终达到 $N$ 元的概率，证明 $p_k=\frac12(p_{k+1}+p_{k-1})$，并用特征根方法解出 $p_k=k/N$. 赢的概率改为 $p\ne\frac12$ 时呢？
			\item\label{ex:eig:C8} 证明：$n$ 阶三对角矩阵 $\vv T=\begin{pmatrix}2&-1&&\\-1&2&\ddots&\\&\ddots&\ddots&-1\\&&-1&2\end{pmatrix}$ 的特征值为 $2-2\cos\frac{k\pi}{n+1}$（$k=1,\ldots,n$），对应特征向量的第 $j$ 个分量为 $\sin\frac{jk\pi}{n+1}$. 对 $n=2,3$ 直接验证.
			\item\label{ex:eig:C9} （Perron–Frobenius 的二阶情形）设 $\vv A$ 是元素全为正的二阶矩阵. 证明 $\vv A$ 有两个互异实特征值，较大者 $r$ 为正且大于另一个的绝对值，属于 $r$ 的特征向量可取为各分量都为正. 由此解释：为什么 PageRank 引入阻尼系数后，``全部元素为正''能保证排序唯一.
			\item\label{ex:eig:C10} 证明：满足 $\vv A^2=-\vv I$ 的 $n$ 阶实方阵存在当且仅当 $n$ 为偶数，且这时 $\vv A$ 相似于 $\diag(\vv J,\ldots,\vv J)$，$\vv J=\begin{pmatrix}0&-1\\1&0\end{pmatrix}$.
		\end{exC}
	}
	
	% =====================================================================
	% 第九章　内积空间
	% =====================================================================
	\newcommand{\ExInnerProduct}{%
		\exsetheading
		\Agroup
		\begin{exA}
			\item\label{ex:ip:A1} 在 $\mathbb{R}^4$ 中，设 $\vv\alpha=(1,2,2,0)^\trans$，$\vv\beta=(2,0,1,2)^\trans$. 求 $\|\vv\alpha\|$，$\|\vv\beta\|$，$\langle\vv\alpha,\vv\beta\rangle$，它们的夹角，以及 $\|\vv\alpha-\vv\beta\|$.
			\item\label{ex:ip:A2} 判断下列二元函数是否为 $\mathbb{R}^2$ 上的内积：
			(1) $\langle\vv x,\vv y\rangle=x_1y_1+2x_2y_2$；(2) $\langle\vv x,\vv y\rangle=x_1y_2+x_2y_1$；(3) $\langle\vv x,\vv y\rangle=2x_1y_1+x_1y_2+x_2y_1+x_2y_2$；(4) $\langle\vv x,\vv y\rangle=x_1y_1-x_2y_2$.
			\item\label{ex:ip:A3} 用 Gram–Schmidt 方法把下列向量组正交化并单位化：
			(1) $(1,1,1)^\trans$，$(1,2,0)^\trans$，$(0,1,2)^\trans$；(2) $(1,1,0,0)^\trans$，$(1,0,1,0)^\trans$，$(-1,0,0,1)^\trans$.
			\item\label{ex:ip:A4} 求与 $(1,1,1)^\trans$ 和 $(1,0,-1)^\trans$ 都正交的单位向量.
			\item\label{ex:ip:A5} 已知 $\vv e_1=\frac13(1,2,2)^\trans$，$\vv e_2=\frac13(2,1,-2)^\trans$，$\vv e_3=\frac13(2,-2,1)^\trans$. 验证它们构成 $\mathbb{R}^3$ 的标准正交基，并只用内积求 $(3,6,0)^\trans$ 在这组基下的坐标.
			\item\label{ex:ip:A6} 求 $\vv\alpha=(1,2,3,4)^\trans$ 在 $W=\spn\big((1,1,0,0)^\trans,(0,0,1,1)^\trans\big)$ 上的正交投影，以及 $\vv\alpha$ 到 $W$ 的距离.
			\item\label{ex:ip:A7} 用最小二乘法求过数据点 $(0,1)$，$(1,2)$，$(2,2)$，$(3,4)$，$(4,6)$ 的最佳直线 $y=a+bx$，并计算残差平方和.
			\item\label{ex:ip:A8} 求 $a,b,c$ 使 $\begin{pmatrix}\frac1{\sqrt2}&a&\frac1{\sqrt6}\\0&b&-\frac2{\sqrt6}\\-\frac1{\sqrt2}&c&\frac1{\sqrt6}\end{pmatrix}$ 为正交矩阵.
			\item\label{ex:ip:A9} 证明：若 $\vv\alpha_1,\ldots,\vv\alpha_k$ 两两正交，则 $\|\vv\alpha_1+\cdots+\vv\alpha_k\|^2=\|\vv\alpha_1\|^2+\cdots+\|\vv\alpha_k\|^2$.
			\item\label{ex:ip:A10} 证明：$\big|\|\vv\alpha\|-\|\vv\beta\|\big|\le\|\vv\alpha-\vv\beta\|$.
			\item\label{ex:ip:A11} 用 Cauchy–Schwarz 不等式证明：(1) $(a_1+\cdots+a_n)^2\le n(a_1^2+\cdots+a_n^2)$；(2) 若 $a_i>0$，则 $(a_1+\cdots+a_n)\big(\frac1{a_1}+\cdots+\frac1{a_n}\big)\ge n^2$.
			\item\label{ex:ip:A12} 在 $C[0,1]$ 中取 $\langle f,g\rangle=\int_0^1f(x)g(x)\,\mathrm dx$. 求 $\|1\|$，$\|x\|$，$1$ 与 $x$ 的夹角，并把 $1,x$ 正交化.
			\item\label{ex:ip:A13} 设 $\vv A=\begin{pmatrix}1&1&0&1\\0&1&1&1\end{pmatrix}$. 求 $\ker\vv A$ 与行空间 $\imm\vv A^\trans$ 的标准正交基，并验证它们互为正交补.
			\item\label{ex:ip:A14} 求 $\vv A=\begin{pmatrix}1&1\\1&0\\0&1\end{pmatrix}$ 的 QR 分解.
			\item\label{ex:ip:A15} 设 $\vv u=\frac13(1,2,2)^\trans$. 写出 Householder 矩阵 $\vv H=\vv I-2\vv u\vv u^\trans$，验证它是对称的正交矩阵，并求 $\vv H\vv u$ 与 $\vv H(2,-1,0)^\trans$.
			\item\label{ex:ip:A16} 判断下列正交矩阵是旋转还是反射；若是旋转求旋转角，若是反射求反射轴：$\frac15\begin{pmatrix}3&-4\\4&3\end{pmatrix}$，$\frac15\begin{pmatrix}3&4\\4&-3\end{pmatrix}$.
			\item\label{ex:ip:A17} 对``过原点的直线''模型 $y=bx$，推导最小二乘解 $\hat b=\sum x_iy_i\big/\sum x_i^2$.
			\item\label{ex:ip:A18} 在 $\mathbb{R}^3$ 中，求点 $(1,2,3)$ 到直线 $\vv x=t(1,1,1)^\trans$ 的距离.
			\item\label{ex:ip:A19} 证明：两组标准正交基之间的过渡矩阵是正交矩阵.
			\item\label{ex:ip:A20} 三个法条的词频向量（关键词：违约、赔偿、合同、离婚）分别为 $(2,1,3,0)^\trans$，$(0,1,0,4)^\trans$，$(1,1,2,0)^\trans$，查询为 $(1,1,1,0)^\trans$. 分别用内积和余弦相似度给三个法条排序，结果有何不同？
			\item\label{ex:ip:A21} 设 $\vv u\in\mathbb{R}^n$ 是单位向量. 证明 $\vv P=\vv u\vv u^\trans$ 满足 $\vv P^\trans=\vv P$，$\vv P^2=\vv P$，且对任意 $\vv x$，$\vv P\vv x$ 是 $\vv x$ 在 $\spn(\vv u)$ 上的正交投影. 若 $\vv u$ 不是单位向量，投影矩阵应当怎样修改？
		\end{exA}
		
		\Bgroup
		\begin{exB}
			\item\label{ex:ip:B1} （极化恒等式）证明实内积空间中 $\langle\vv\alpha,\vv\beta\rangle=\frac14\big(\|\vv\alpha+\vv\beta\|^2-\|\vv\alpha-\vv\beta\|^2\big)$. 再证明：$\mathbb{R}^2$ 上的``范数'' $\|\vv x\|_1=|x_1|+|x_2|$ 不满足平行四边形法则，因而不能来自任何内积.
			\item\label{ex:ip:B2} （Bessel 不等式）设 $\vv e_1,\ldots,\vv e_k$ 是标准正交组. 证明 $\sum_{i=1}^k\langle\vv\alpha,\vv e_i\rangle^2\le\|\vv\alpha\|^2$，等号成立当且仅当 $\vv\alpha\in\spn(\vv e_1,\ldots,\vv e_k)$.
			\item\label{ex:ip:B3} 设实矩阵 $\vv P$ 满足 $\vv P^\trans=\vv P$，$\vv P^2=\vv P$. 证明 $\vv P$ 是向 $\imm\vv P$ 的正交投影的矩阵. 若只有 $\vv P^2=\vv P$，它是什么样的投影？举例说明.
			\item\label{ex:ip:B4} （加权最小二乘）若各观测值的可靠程度不同，给第 $i$ 个残差权重 $w_i>0$，最小化 $\sum w_i(y_i-a-bx_i)^2$. 证明正规方程变为 $\vv A^\trans\vv W\vv A\hat{\vv\beta}=\vv A^\trans\vv W\vv y$，$\vv W=\diag(w_1,\ldots,w_n)$. 它对应于哪个内积下的正交投影？
			\item\label{ex:ip:B5} 某商品的销售额 $y$ 随时间 $t$ 近似按指数增长 $y=c\,\mathrm e^{kt}$. 已知 $t=0,1,2,3$ 时 $y$ 约为 $2.0,2.7,3.7,5.0$. 两边取对数化为线性模型，用最小二乘估计 $c,k$. 这样做最小化的是哪个量？与直接最小化 $\sum(y_i-c\mathrm e^{kt_i})^2$ 有何不同？
			\item\label{ex:ip:B6} （Gram 行列式）设 $\vv\alpha_1,\ldots,\vv\alpha_k\in\mathbb{R}^n$，$\vv G=(\langle\vv\alpha_i,\vv\alpha_j\rangle)$. 证明 $\det\vv G\ge0$，等号成立当且仅当它们线性相关. 在 $k=2$ 时说明 $\sqrt{\det\vv G}$ 是两向量张成的平行四边形的面积.
			\item\label{ex:ip:B7} 证明：$\mathbb{R}^n$ 中点 $\vv x_0$ 到超平面 $\vv a^\trans\vv x=b$（$\vv a\ne\zero$）的距离为 $|\vv a^\trans\vv x_0-b|/\|\vv a\|$.
			\item\label{ex:ip:B8} 证明：实正交矩阵的复特征值的模都等于 $1$.
			\item\label{ex:ip:B9} （Fourier 逼近）在 $C[-\pi,\pi]$ 中取 $\langle f,g\rangle=\int_{-\pi}^\pi fg\,\mathrm dx$. 证明 $1,\cos x,\sin x,\cos2x,\sin2x$ 两两正交. 求 $f(x)=x$ 在它们张成的子空间上的正交投影.
			\item\label{ex:ip:B10} 设 $\vv A$ 为实矩阵. 证明：在 $\vv{Ax}=\vv y$ 的全部最小二乘解中，范数最小的那一个唯一，并且它属于行空间 $\imm\vv A^\trans$.
			\item\label{ex:ip:B11} 设 $\vv x,\vv y\in\mathbb{R}^n$，$\|\vv x\|=\|\vv y\|$，$\vv x\ne\vv y$. 令 $\vv u=(\vv x-\vv y)/\|\vv x-\vv y\|$. 证明 Householder 矩阵 $\vv I-2\vv u\vv u^\trans$ 把 $\vv x$ 变为 $\vv y$. 这是用 Householder 反射做 QR 分解的基础.
			\item\label{ex:ip:B12} 设两个变量的样本标准差分别为 $s_x,s_y$，相关系数为 $r$. 证明回归直线 $y=a+bx$ 的斜率为 $b=r\,s_y/s_x$，并且直线经过点 $(\bar x,\bar y)$.
			\item\label{ex:ip:B13} 证明：保持\textbf{夹角}的可逆线性变换（对一切非零向量，$\mathscr T\vv\alpha$ 与 $\mathscr T\vv\beta$ 的夹角等于 $\vv\alpha$ 与 $\vv\beta$ 的夹角）必是正交变换的非零倍数.
			\item\label{ex:ip:B14} 在回归中加入一个``虚拟变量''：设有两组数据，第一组 $y$ 值为 $3,5,4$，第二组为 $7,8,9$. 以 $\vv A=(\vv 1,\vv d)$ 为设计矩阵（$\vv d$ 在第二组为 $1$、第一组为 $0$），求最小二乘解，并解释两个系数的含义.
			\item\label{ex:ip:B15} 证明：在 $\mathbb{R}^2$ 中，任意两个旋转可交换，而一个旋转与一个反射一般不可交换；两个反射的乘积是旋转.
		\end{exB}
		
		\Cgroup
		\begin{exC}
			\item\label{ex:ip:C1} （Hadamard 不等式）用 QR 分解证明：对 $n$ 阶实方阵 $\vv A=(\vv a_1,\ldots,\vv a_n)$，$|\det\vv A|\le\|\vv a_1\|\cdots\|\vv a_n\|$，并说明它的几何意义（平行多面体的体积不超过棱长之积）.
			\item\label{ex:ip:C2} （Gram–Schmidt 的数值不稳定）设 $\varepsilon$ 很小，对 $\vv a_1=(1,\varepsilon,0,0)^\trans$，$\vv a_2=(1,0,\varepsilon,0)^\trans$，$\vv a_3=(1,0,0,\varepsilon)^\trans$ 作 Gram–Schmidt，在计算中把 $1+\varepsilon^2$ 舍入为 $1$. 证明所得的 $\vv e_2,\vv e_3$ 之间的夹角是 $60^\circ$ 而不是 $90^\circ$. 再用``修正的 Gram–Schmidt''（每得到一个 $\vv e_k$，立即从\textbf{剩余的}所有向量中减去它的分量）重做，比较结果.
			\item\label{ex:ip:C3} （Legendre 多项式）令 $P_n(x)=\frac1{2^nn!}\frac{\mathrm d^n}{\mathrm dx^n}(x^2-1)^n$. 证明 $P_0,P_1,P_2,\ldots$ 在 $\langle f,g\rangle=\int_{-1}^1fg\,\mathrm dx$ 下两两正交，并求 $\|P_n\|$.
			\item\label{ex:ip:C4} （Gauss–Markov 定理的线性代数部分）设 $\vv A$ 列满秩. 若 $\vv C\vv y$ 是参数 $\vv\beta$ 的``线性无偏估计''，即 $\vv{CA}=\vv I$. 令 $\vv C_0=(\vv A^\trans\vv A)^{-1}\vv A^\trans$（最小二乘）. 证明 $\vv C\vv C^\trans-\vv C_0\vv C_0^\trans$ 是半正定矩阵（第十一章）. 学过概率论后可以知道，这说明最小二乘估计的方差在一切线性无偏估计中最小.
			\item\label{ex:ip:C5} （复内积空间）在 $\mathbb{C}^n$ 上定义 $\langle\vv x,\vv y\rangle=\sum x_i\bar y_i$. 证明它满足共轭对称性、对第一个变量线性、正定性；证明 Cauchy–Schwarz 不等式仍成立. 满足 $\bar{\vv U}^\trans\vv U=\vv I$ 的复矩阵称为\textbf{酉矩阵}，证明它保持这个内积.
			\item\label{ex:ip:C6} 设 $f:\mathbb{R}^n\to\mathbb{R}^n$ 保持距离（$\|f(\vv x)-f(\vv y)\|=\|\vv x-\vv y\|$）且 $f(\zero)=\zero$. 证明 $f$ 保持内积，进而证明 $f$ 是线性的正交变换. 由此说明：保持距离的变换都可以分解为正交变换与平移的复合.
			\item\label{ex:ip:C7} （伪逆）设 $\vv A$ 列满秩，$\vv A^+=(\vv A^\trans\vv A)^{-1}\vv A^\trans$. 验证 Penrose 四条件：$\vv{AA}^+\vv A=\vv A$，$\vv A^+\vv{AA}^+=\vv A^+$，$(\vv{AA}^+)^\trans=\vv{AA}^+$，$(\vv A^+\vv A)^\trans=\vv A^+\vv A$. 说明 $\vv{AA}^+$ 是向哪个子空间的正交投影.
			\item\label{ex:ip:C8} 求空间中两条异面直线 $\vv x=\vv p+s\vv u$ 与 $\vv x=\vv q+t\vv v$ 之间的最短距离：把问题写成关于 $(s,t)$ 的最小二乘问题，并对 $\vv p=\zero$，$\vv u=(1,0,0)^\trans$，$\vv q=(0,1,1)^\trans$，$\vv v=(0,1,-1)^\trans$ 计算.
		\end{exC}
	}
	
	% =====================================================================
	% 第十章　矩阵分解
	% =====================================================================
	\newcommand{\ExDecomposition}{%
		\exsetheading
		\Agroup
		\begin{exA}
			\item\label{ex:dec:A1} 求正交矩阵 $\vv Q$ 使 $\vv Q^\trans\vv{AQ}$ 为对角矩阵：
			\[\begin{pmatrix}1&2\\2&1\end{pmatrix},\qquad\begin{pmatrix}3&-1\\-1&3\end{pmatrix},\]\[\begin{pmatrix}2&2&0\\2&2&0\\0&0&1\end{pmatrix},\qquad\begin{pmatrix}1&1&1\\1&1&1\\1&1&1\end{pmatrix}.\]
			\item\label{ex:dec:A2} 写出 $\begin{pmatrix}3&-1\\-1&3\end{pmatrix}$ 的谱分解 $\lambda_1\vv q_1\vv q_1^\trans+\lambda_2\vv q_2\vv q_2^\trans$，并验证两个投影矩阵之和为 $\vv I$、之积为 $\vv O$.
			\item\label{ex:dec:A3} 求二阶实对称矩阵 $\vv A$，使它的特征值为 $1,2$，且 $(1,1)^\trans$ 是属于 $1$ 的特征向量.
			\item\label{ex:dec:A4} 设三阶实对称矩阵 $\vv A$ 的特征值为 $1,1,-1$，属于 $-1$ 的特征向量为 $(0,1,1)^\trans$. 求 $\vv A$.
			\item\label{ex:dec:A5} 求 $\vv x^\trans\vv{Ax}$ 在单位球面 $\|\vv x\|=1$ 上的最大值与最小值，以及取到它们的 $\vv x$：$\vv A=\begin{pmatrix}2&2&0\\2&2&0\\0&0&1\end{pmatrix}$.
			\item\label{ex:dec:A6} 求 $x_1^2+4x_1x_2+x_2^2$ 在单位圆 $x_1^2+x_2^2=1$ 上的最大值与最小值.
			\item\label{ex:dec:A7} 求下列矩阵的奇异值分解：$\begin{pmatrix}2&0\\0&-3\end{pmatrix}$，$\begin{pmatrix}1&1\\1&1\end{pmatrix}$，$(3,4)$（$1\times2$ 矩阵）.
			\item\label{ex:dec:A8} 求 $\begin{pmatrix}1&1\\0&1\end{pmatrix}$ 的奇异值，并与它的特征值比较.（答案与黄金分割比有关.）
			\item\label{ex:dec:A9} 求 $\begin{pmatrix}3&1\\1&3\end{pmatrix}$ 的最佳秩一近似，并计算误差 $\|\vv A-\vv A_1\|_F$.
			\item\label{ex:dec:A10} 设 $\vv A=\begin{pmatrix}5&4\\4&5\end{pmatrix}$. 利用谱分解求 $\vv A^{10}$（写成两个矩阵之和的形式即可），并求对称正定矩阵 $\vv B$ 使 $\vv B^2=\vv A$.
			\item\label{ex:dec:A11} 证明：特征值全为 $0$ 的实对称矩阵是零矩阵. 对非对称矩阵举出反例.
			\item\label{ex:dec:A12} 证明：既对称又正交的实矩阵 $\vv A$ 满足 $\vv A^2=\vv I$，其特征值为 $\pm1$；并说明它在几何上是什么变换.
			\item\label{ex:dec:A13} 证明：$\vv A^\trans\vv A$ 与 $\vv{AA}^\trans$ 的非零特征值相同（重数也相同）.
			\item\label{ex:dec:A14} 证明：$\|\vv A\|_F^2=\tr(\vv A^\trans\vv A)=\sigma_1^2+\cdots+\sigma_r^2$.
			\item\label{ex:dec:A15} 设 $\vv A$ 是可逆实方阵. 证明 $|\det\vv A|=\sigma_1\sigma_2\cdots\sigma_n$，$\vv A^{-1}$ 的奇异值为 $1/\sigma_i$.
			\item\label{ex:dec:A16} 判断下列矩阵能否用正交矩阵对角化，说明理由：$\begin{pmatrix}1&2\\0&3\end{pmatrix}$，$\begin{pmatrix}1&2\\2&-2\end{pmatrix}$，$\begin{pmatrix}0&-1\\1&0\end{pmatrix}$（在实数范围内）.
			\item\label{ex:dec:A17} 四个已中心化的二维数据点为 $(2,1)$，$(-2,-1)$，$(1,2)$，$(-1,-2)$. 求样本协方差矩阵、第一主成分方向及其解释的方差比例.
			\item\label{ex:dec:A18} 计算 $\diag(10,1)$ 与 $\begin{pmatrix}1&1\\1&1.01\end{pmatrix}$ 的条件数（后者用近似值）. 后者的行列式与 $\diag(10,0.001)$ 相同，二者的条件数相同吗？
			\item\label{ex:dec:A19} 一张 $800\times600$ 像素的灰度图片，保留前 $k=40$ 个奇异值时需存储多少个数？压缩比是多少？$k$ 取多少时存储量反而超过原图？
			\item\label{ex:dec:A20} 证明：实对称矩阵的秩等于非零特征值的个数；奇异值的个数（非零）等于秩.
			\item\label{ex:dec:A21} 证明：实对称矩阵 $\vv A$ 满足``对一切 $\vv x\ne\zero$，$\vv x^\trans\vv A\vv x>0$''，当且仅当 $\vv A$ 的特征值全为正.（提示：用谱定理把 $\vv x^\trans\vv A\vv x$ 写成 $\sum\lambda_iy_i^2$.）
		\end{exA}
		
		\Bgroup
		\begin{exB}
			\item\label{ex:dec:B1} 设 $\vv A,\vv B$ 都是 $n$ 阶实对称矩阵且 $\vv{AB}=\vv{BA}$. 证明存在正交矩阵 $\vv Q$ 使 $\vv Q^\trans\vv{AQ}$ 与 $\vv Q^\trans\vv{BQ}$ 同时为对角矩阵.
			\item\label{ex:dec:B2} 证明：半正定的实对称矩阵 $\vv A$（特征值全 $\ge0$）有唯一的半正定平方根 $\vv B$（$\vv B^2=\vv A$，$\vv B$ 对称且特征值全 $\ge0$）.
			\item\label{ex:dec:B3} 设实对称矩阵 $\vv A$ 的特征值为 $\lambda_1\ge\lambda_2\ge\cdots\ge\lambda_n$，$\vv q_1$ 是属于 $\lambda_1$ 的单位特征向量. 证明 $\lambda_2=\max\{\vv x^\trans\vv{Ax}:\|\vv x\|=1,\ \vv x\perp\vv q_1\}$. 这说明第二主成分是``与第一主成分垂直的方向中方差最大的''.
			\item\label{ex:dec:B4} 证明：对任意 $\vv x$，$\|\vv{Ax}\|\le\sigma_1\|\vv x\|$，且等号可以取到. 因此 $\sigma_1=\max_{\vv x\ne\zero}\|\vv{Ax}\|/\|\vv x\|$，称为 $\vv A$ 的\textbf{谱范数}\index{pu3 fan4shu4@谱范数} $\|\vv A\|_2$.
			\item\label{ex:dec:B5} 证明：方阵 $\vv A$ 的任一特征值 $\lambda$ 满足 $\sigma_{\min}\le|\lambda|\le\sigma_{\max}$.
			\item\label{ex:dec:B6} 求 $\vv A=\begin{pmatrix}1&0&1\\0&1&1\end{pmatrix}$ 的奇异值分解.
			\item\label{ex:dec:B7} （伪逆）设 $\vv A=\vv U\vv\Sigma\vv V^\trans$，定义 $\vv A^+=\vv V\vv\Sigma^+\vv U^\trans$，其中 $\vv\Sigma^+$ 是把 $\vv\Sigma$ 转置并把非零元取倒数. 证明 $\vv A^+\vv b$ 是 $\vv{Ax}=\vv b$ 的范数最小的最小二乘解. 当 $\vv A$ 列满秩时，它与第九章的 $(\vv A^\trans\vv A)^{-1}\vv A^\trans$ 一致.
			\item\label{ex:dec:B8} 设 $p$ 个标准化指标两两相关系数都为 $\rho$（$0<\rho<1$），相关系数矩阵为 $(1-\rho)\vv I+\rho\vv J$（$\vv J$ 为全 $1$ 矩阵）. 求它的特征值与第一主成分，并求第一主成分解释的方差比例. 当 $\rho\to1$ 与 $\rho\to0$ 时比例分别趋于多少？
			\item\label{ex:dec:B9} 设两个指标的协方差矩阵为 $\begin{pmatrix}100&6\\6&1\end{pmatrix}$（第一个指标以``元''计，第二个以``万元''计）. 对协方差矩阵与相关系数矩阵分别作主成分分析，比较第一主成分. 由此说明：主成分分析对量纲敏感，为什么通常先标准化.
			\item\label{ex:dec:B10} 证明 $\|\vv A-\vv A_k\|_2=\sigma_{k+1}$.
			\item\label{ex:dec:B11} 证明：若 $\vv A$ 的奇异值互不相同且都不为零，则 $\vv A$ 的奇异值分解中的 $\vv u_i,\vv v_i$ 除了同时改变符号之外是唯一的.
			\item\label{ex:dec:B12} （正规矩阵）满足 $\vv A^\trans\vv A=\vv{AA}^\trans$ 的实方阵称为正规矩阵. 证明二阶实正规矩阵只有两类：对称矩阵，以及形如 $\begin{pmatrix}a&-b\\b&a\end{pmatrix}$ 的矩阵（旋转与伸缩的复合）.
			\item\label{ex:dec:B13} 设 $\vv A$ 是实对称矩阵. 证明 $\vv A$ 的特征值都落在 $[\min_i(a_{ii}-r_i),\ \max_i(a_{ii}+r_i)]$ 中，$r_i=\sum_{j\ne i}|a_{ij}|$. 并证明 $\lambda_{\max}\ge\max_ia_{ii}$.
			\item\label{ex:dec:B14} 对第十章``潜在语义分析''的例子，取 $k=1$，说明此时四份文书分别被表示成什么，查询``解约''与哪几份文书相似. 与 $k=2$ 比较，``压缩过头''的现象是否更严重？
			\item\label{ex:dec:B15} 证明：实对称矩阵 $\vv A$ 的特征值都在 $[a,b]$ 中，当且仅当 $\vv A-a\vv I$ 与 $b\vv I-\vv A$ 的二次型都非负.
		\end{exB}
		
		\Cgroup
		\begin{exC}
			\item\label{ex:dec:C1} （Eckart–Young 定理的证明）设 $\rk\vv B\le k$. 证明 $\ker\vv B$ 与 $\spn(\vv v_1,\ldots,\vv v_{k+1})$ 有非零公共向量 $\vv x$，由此证明 $\|\vv A-\vv B\|_2\ge\sigma_{k+1}$. 再借助 Weyl 不等式 $\sigma_{i+j-1}(\vv X+\vv Y)\le\sigma_i(\vv X)+\sigma_j(\vv Y)$（可承认）证明 Frobenius 范数的情形.
			\item\label{ex:dec:C2} （Cauchy 交错定理）设实对称矩阵 $\vv A$ 的特征值为 $\lambda_1\ge\cdots\ge\lambda_n$，删去最后一行一列所得矩阵的特征值为 $\mu_1\ge\cdots\ge\mu_{n-1}$. 证明 $\lambda_1\ge\mu_1\ge\lambda_2\ge\mu_2\ge\cdots\ge\mu_{n-1}\ge\lambda_n$.（提示：Courant–Fischer 极小极大刻画.）
			\item\label{ex:dec:C3} （复谱定理）证明：复方阵 $\vv A$ 可以被酉矩阵对角化（存在 $\vv U$ 使 $\bar{\vv U}^\trans\vv{AU}$ 为对角矩阵），当且仅当 $\vv A\bar{\vv A}^\trans=\bar{\vv A}^\trans\vv A$. 实对称矩阵、实正交矩阵、实反对称矩阵都满足这个条件.
			\item\label{ex:dec:C4} （正交 Procrustes 问题）给定 $n\times k$ 矩阵 $\vv A,\vv B$，求正交矩阵 $\vv Q$ 使 $\|\vv{QA}-\vv B\|_F$ 最小. 证明：若 $\vv{BA}^\trans=\vv U\vv\Sigma\vv V^\trans$，则最优解为 $\vv Q=\vv{UV}^\trans$. 这个问题出现在比较两组测量数据、对齐两份词向量等场合.
			\item\label{ex:dec:C5} （Ky Fan 定理）设实对称矩阵的特征值为 $\lambda_1\ge\cdots\ge\lambda_n$. 证明对任意列标准正交的 $n\times k$ 矩阵 $\vv X$，$\tr(\vv X^\trans\vv{AX})\le\lambda_1+\cdots+\lambda_k$. 由此说明前 $k$ 个主成分``一起''解释的方差也是最大的.
			\item\label{ex:dec:C6} （幂迭代）设 $\vv A$ 可对角化，特征值满足 $|\lambda_1|>|\lambda_2|\ge\cdots$. 证明：对``一般的''初始向量 $\vv x_0$，$\vv x_k=\vv A^k\vv x_0/\|\vv A^k\vv x_0\|$ 收敛到属于 $\lambda_1$ 的单位特征向量（可能差一个符号），收敛速度由 $|\lambda_2/\lambda_1|$ 决定. 对 $\begin{pmatrix}2&1\\1&2\end{pmatrix}$ 从 $(1,0)^\trans$ 出发迭代五步验证.
			\item\label{ex:dec:C7} （图的 Laplace 矩阵）设 $\vv L=\vv D-\vv A$ 是无向图的 Laplace 矩阵（见第五章章末习题~\ref{ex:det:C4}）. 证明 $\vv x^\trans\vv L\vv x=\sum_{\text{边 }ij}(x_i-x_j)^2$，因而 $\vv L$ 半正定；特征值 $0$ 的重数等于图的连通分支数. 第二小的特征值对应的特征向量（Fiedler 向量）的正负号常被用来把图``切''成两半——这就是谱聚类的思想.
			\item\label{ex:dec:C8} （正交回归）对中心化数据点 $\vv z_1,\ldots,\vv z_n\in\mathbb{R}^2$，求过原点的直线使各点到直线的\textbf{垂直}距离平方和最小. 证明最优方向就是第一主成分方向. 与普通最小二乘（只看竖直方向的残差）比较，两条直线何时相同？
			\item\label{ex:dec:C9} （扰动分析）设 $\vv A$ 可逆，$\vv{Ax}=\vv b$，$\vv A(\vv x+\delta\vv x)=\vv b+\delta\vv b$. 证明 $\dfrac{\|\delta\vv x\|}{\|\vv x\|}\le\kappa(\vv A)\dfrac{\|\delta\vv b\|}{\|\vv b\|}$，$\kappa(\vv A)=\sigma_1/\sigma_n$，并说明等号可以取到. 用它解释第六章病态方程组的例子.
		\end{exC}
	}
	
	% =====================================================================
	% 第十一章　二次型
	% =====================================================================
	\newcommand{\ExQuadratic}{%
		\exsetheading
		\Agroup
		\begin{exA}
			\item\label{ex:quad:A1} 写出下列二次型的矩阵与秩：
			(1) $x_1^2-3x_2^2+4x_1x_2-2x_2x_3$；(2) $x_1x_2+x_2x_3+x_3x_1$；(3) $(x_1+x_2+x_3)^2$；(4) $\sum_{i=1}^n x_i^2+\sum_{i<j}x_ix_j$.
			\item\label{ex:quad:A2} 写出以 $\begin{pmatrix}1&-1&0\\-1&2&3\\0&3&0\end{pmatrix}$ 为矩阵的二次型.
			\item\label{ex:quad:A3} 用配方法化为标准形，并写出所用的可逆变量替换：
			(1) $x_1^2+2x_2^2+5x_3^2+2x_1x_2+6x_2x_3$；(2) $x_1^2-4x_1x_2+4x_2^2-x_3^2$；(3) $x_1x_2-x_2x_3$.
			\item\label{ex:quad:A4} 用正交变换法把 $2x_1^2+3x_2^2+3x_3^2+4x_2x_3$ 化为标准形，写出正交矩阵. 曲面 $2x_1^2+3x_2^2+3x_3^2+4x_2x_3=1$ 是什么形状？
			\item\label{ex:quad:A5} 求 A3 中各二次型的正、负惯性指数与符号差，并写出规范形.
			\item\label{ex:quad:A6} 判断下列二次型是否正定：
			(1) $2x_1^2+x_2^2+3x_3^2+2x_1x_2$；(2) $x_1^2+x_2^2+x_3^2+4x_1x_2$；(3) $5x_1^2+x_2^2+5x_3^2+4x_1x_2-8x_1x_3-4x_2x_3$.
			\item\label{ex:quad:A7} $t$ 取何值时，$2x_1^2+x_2^2+x_3^2+2x_1x_2+tx_2x_3$ 正定？
			\item\label{ex:quad:A8} 判断曲线类型并求主轴方向：(1) $x^2-4xy+y^2=1$；(2) $3x^2+2xy+3y^2=8$.
			\item\label{ex:quad:A9} 判断下列矩阵是否合同：(1) $\diag(1,2,-3)$ 与 $\diag(-1,5,6)$；(2) $\begin{pmatrix}0&1\\1&0\end{pmatrix}$ 与 $\vv I_2$；(3) $\begin{pmatrix}0&1\\1&0\end{pmatrix}$ 与 $\diag(1,-1)$. 合同时求出 $\vv C$.
			\item\label{ex:quad:A10} 证明：若 $\vv A$ 正定，则 $\vv A^*$ 也正定.
			\item\label{ex:quad:A11} 求 $f(x,y)=x^3+y^3-3xy$ 的全部驻点，并用 Hessian 矩阵判断其类型.
			\item\label{ex:quad:A12} 判断 $\begin{pmatrix}2&-1&0\\-1&2&-1\\0&-1&2\end{pmatrix}$ 是否正定，并求它的 Cholesky 分解.
			\item\label{ex:quad:A13} 求 $\begin{pmatrix}4&2&-2\\2&10&2\\-2&2&6\end{pmatrix}$ 的 Cholesky 分解.
			\item\label{ex:quad:A14} 三种资产收益率的协方差矩阵为 $\vv\Sigma=\begin{pmatrix}0.04&0.01&0\\0.01&0.09&0\\0&0&0.01\end{pmatrix}$. 求组合 $\vv w=(0.5,0.3,0.2)^\trans$ 的方差与标准差.
			\item\label{ex:quad:A15} 两种资产收益率的标准差为 $30\%$ 与 $20\%$，互不相关. 求最小方差组合及其标准差.
			\item\label{ex:quad:A16} 证明：合同变换保持矩阵的对称性与秩. 合同变换保持行列式的符号吗？保持行列式的值吗？
			\item\label{ex:quad:A17} 分别在 $\mathbb{R}$ 上与 $\mathbb{C}$ 上写出二次型 $x_1^2-x_2^2+2x_3^2$ 的规范形.
			\item\label{ex:quad:A18} 判断下列方程表示的二次曲面类型：$x^2+y^2-z^2=1$，$x^2-y^2-z^2=1$，$x^2+y^2-z^2=0$，$x^2+y^2=1$.（提示：先看正、负惯性指数.）
			\item\label{ex:quad:A19} 说明 $\begin{pmatrix}1&2\\2&1\end{pmatrix}$ 不正定，并找出一个使二次型取负值的向量.
			\item\label{ex:quad:A20} 设 $\vv A$ 正定. 证明：$\vv A$ 的任一主子矩阵（取相同的若干行和列）都正定，从而所有主子式都为正.
		\end{exA}
		
		\Bgroup
		\begin{exB}
			\item\label{ex:quad:B1} 设 $\vv A,\vv B$ 正定且 $\vv{AB}=\vv{BA}$. 证明 $\vv{AB}$ 正定. 举例说明``可交换''不能去掉（此时 $\vv{AB}$ 甚至不对称）.
			\item\label{ex:quad:B2} 设 $\vv A$ 正定，$\vv B$ 半正定. 证明 $\det(\vv A+\vv B)\ge\det\vv A$，等号成立当且仅当 $\vv B=\vv O$.（提示：同时合同对角化.）
			\item\label{ex:quad:B3} 设 $\vv A$ 正定，$\vv B$ 对称. 证明：当 $|t|$ 充分小时 $\vv A+t\vv B$ 正定.
			\item\label{ex:quad:B4} 设 $\vv A$ 正定. 证明对 $i\ne j$ 有 $|a_{ij}|<\sqrt{a_{ii}a_{jj}}$，从而 $\vv A$ 的绝对值最大的元素在对角线上.
			\item\label{ex:quad:B5} 证明：实对称矩阵 $\vv A$ 满足 $\lambda_{\min}\|\vv x\|^2\le\vv x^\trans\vv{Ax}\le\lambda_{\max}\|\vv x\|^2$.
			\item\label{ex:quad:B6} 求 $\begin{pmatrix}0&1&1\\1&0&1\\1&1&0\end{pmatrix}$ 的正、负惯性指数（用两种方法：特征值；配方法）.
			\item\label{ex:quad:B7} $\lambda$ 取何值时，$\lambda(x_1^2+x_2^2+x_3^2)+2x_1x_2+2x_1x_3-2x_2x_3$ 正定？
			\item\label{ex:quad:B8} 求二次型 $\vv x^\trans\vv{Ax}$ 在约束 $\vv c^\trans\vv x=1$ 下的最小值（$\vv A$ 正定）：用拉格朗日乘数法证明最优解为 $\vv x^*=\dfrac{\vv A^{-1}\vv c}{\vv c^\trans\vv A^{-1}\vv c}$. 最小方差组合是它的特例.
			\item\label{ex:quad:B9} 设三种资产互不相关、方差都为 $\sigma^2$，期望收益率为 $\mu_1<\mu_2<\mu_3$. 求期望收益为 $m$ 的最小方差组合的权重，并说明有效前沿是一条抛物线（在``方差—期望收益''平面上）.
			\item\label{ex:quad:B10} 设 $\vv A=\begin{pmatrix}2&1\\1&2\end{pmatrix}$，$\vv B=\begin{pmatrix}1&0\\0&-1\end{pmatrix}$. 求可逆矩阵 $\vv C$ 使 $\vv C^\trans\vv{AC}=\vv I$ 且 $\vv C^\trans\vv{BC}$ 为对角矩阵.
			\item\label{ex:quad:B11} 设 $\vv A$ 正定. 求 $f(\vv x)=\vv x^\trans\vv{Ax}+\vv b^\trans\vv x+c$ 的最小值点与最小值. 若 $\vv A$ 只是半正定，最小值一定存在吗？
			\item\label{ex:quad:B12} 设 $\vv A$ 半正定，$\vv x^\trans\vv{Ax}=0$. 证明 $\vv{Ax}=\zero$.
			\item\label{ex:quad:B13} （Schur 补判据）证明：对称分块矩阵 $\begin{pmatrix}\vv A&\vv B\\ \vv B^\trans&\vv C\end{pmatrix}$ 正定，当且仅当 $\vv A$ 正定且 $\vv C-\vv B^\trans\vv A^{-1}\vv B$ 正定.
			\item\label{ex:quad:B14} 设 $\vv A$ 是二阶正定矩阵. 证明椭圆区域 $\{\vv x:\vv x^\trans\vv{Ax}\le1\}$ 的面积为 $\pi/\sqrt{\det\vv A}$.（提示：用主轴定理化为标准椭圆.）
			\item\label{ex:quad:B15} 设 $\vv A$ 为 $m\times n$ 实矩阵，$\mu>0$. 证明 $\vv A^\trans\vv A+\mu\vv I$ 正定，并说明岭回归 $\min\|\vv y-\vv{Ax}\|^2+\mu\|\vv x\|^2$ 的解为 $(\vv A^\trans\vv A+\mu\vv I)^{-1}\vv A^\trans\vv y$.
			\item\label{ex:quad:B16} 证明：实对称矩阵 $\vv A$ 正定，当且仅当存在正定矩阵 $\vv B$ 使 $\vv A=\vv B^2$；当且仅当 $\vv A^{-1}$ 存在且正定.
		\end{exB}
		
		\Cgroup
		\begin{exC}
			\item\label{ex:quad:C1} （Schur 乘积定理）对同阶矩阵定义逐元素乘积 $(\vv A\circ\vv B)_{ij}=a_{ij}b_{ij}$. 证明：若 $\vv A,\vv B$ 都半正定，则 $\vv A\circ\vv B$ 也半正定.（提示：把 $\vv B$ 写成 $\sum\lambda_k\vv q_k\vv q_k^\trans$.）
			\item\label{ex:quad:C2} （Minkowski 行列式不等式）设 $\vv A,\vv B$ 都是 $n$ 阶正定矩阵. 证明 $\det(\vv A+\vv B)^{1/n}\ge\det(\vv A)^{1/n}+\det(\vv B)^{1/n}$.（提示：同时合同对角化后用 $\prod(1+\mu_i)^{1/n}\ge1+\prod\mu_i^{1/n}$.）
			\item\label{ex:quad:C3} （有理数域上的二次型）证明 $x^2+y^2$ 与 $3x^2+3y^2$ 在 $\mathbb{R}$ 上合同，但在 $\mathbb{Q}$ 上不合同（即不存在有理系数的可逆替换把一个变为另一个）.（提示：若合同，则存在有理数 $a,b$ 使 $a^2+b^2=3$.）这说明惯性定理依赖于实数可以随意开平方.
			\item\label{ex:quad:C4} 函数 $f=x^2+y^4$ 与 $g=x^2-y^4$ 在原点处的 Hessian 矩阵相同且只是半正定，但前者在原点取极小值，后者不取极值. 由此说明 Hessian 半正定时二阶判别法失效.
			\item\label{ex:quad:C5} （广义特征值）设 $\vv A$ 对称，$\vv B$ 正定. 证明 $\max_{\vv x\ne\zero}\dfrac{\vv x^\trans\vv{Ax}}{\vv x^\trans\vv{Bx}}$ 等于方程 $\det(\vv A-\lambda\vv B)=0$ 的最大根. 统计学中的判别分析（Fisher LDA）正是这样求最佳投影方向的.
			\item\label{ex:quad:C6} （Mahalanobis 距离）设 $\vv\Sigma$ 正定，定义 $d(\vv x,\vv y)=\sqrt{(\vv x-\vv y)^\trans\vv\Sigma^{-1}(\vv x-\vv y)}$. 证明它满足三角不等式，且在仿射变换 $\vv x\mapsto\vv{Px}+\vv b$、$\vv\Sigma\mapsto\vv P\vv\Sigma\vv P^\trans$ 下不变. 它与欧氏距离有何不同？
			\item\label{ex:quad:C7} （Hilbert 矩阵）证明 $n$ 阶矩阵 $\vv H=\big(\frac1{i+j-1}\big)$ 正定.（提示：$\frac1{i+j-1}=\int_0^1x^{i-1}x^{j-1}\,\mathrm dx$，它是某个内积下的 Gram 矩阵.）
			\item\label{ex:quad:C8} （切点组合）设有无风险资产，收益率为 $r_f$. 证明使``夏普比率''$\dfrac{\vv\mu^\trans\vv w-r_f}{\sqrt{\vv w^\trans\vv\Sigma\vv w}}$（约束 $\vv 1^\trans\vv w=1$）最大的风险资产组合为 $\vv w\propto\vv\Sigma^{-1}(\vv\mu-r_f\vv 1)$.（提示：先去掉约束，比值关于 $\vv w$ 是零次齐次的；再用 Cauchy–Schwarz 不等式.）
			\item\label{ex:quad:C9} 证明：秩为 $r$、正惯性指数为 $p$ 的实二次型可以写成 $p$ 个线性型的平方和减去 $r-p$ 个线性型的平方和，并且不能写成少于 $r$ 个线性型的平方（带 $\pm$ 号）的组合.
		\end{exC}
	}
	
	% =====================================================================
	% 第十二章　*对偶空间与双线性型（选读）
	% =====================================================================
	\newcommand{\ExBilinear}{%
		\exsetheading
		\Agroup
		\begin{exA}
			\item\label{ex:bil:A1} 判断下列映射是否为相应空间上的线性函数：(1) $\mathbb{R}^3\to\mathbb{R}$，$\vv x\mapsto x_1-2x_3$；(2) $\vv x\mapsto x_1x_2$；(3) $\mathbb{R}[x]\to\mathbb{R}$，$f\mapsto f'(0)$；(4) $f\mapsto\int_0^1f(x)^2\,\mathrm dx$；(5) $M_n(\mathbb{R})\to\mathbb{R}$，$\vv X\mapsto\tr(\vv{AX})$（$\vv A$ 固定）.
			\item\label{ex:bil:A2} 求 $\mathbb{R}^3$ 的基 $(1,0,0)^\trans$，$(1,1,0)^\trans$，$(1,1,1)^\trans$ 的对偶基（把每个对偶基向量写成 $\vv x\mapsto a_1x_1+a_2x_2+a_3x_3$ 的形式）.
			\item\label{ex:bil:A3} 把线性函数 $f(\vv x)=x_1+2x_2+3x_3$ 用 A2 中的对偶基线性表出.
			\item\label{ex:bil:A4} 求 $W=\spn\big((1,1,0)^\trans,(0,1,1)^\trans\big)$ 的零化子 $W^0\subset(\mathbb{R}^3)^*$ 的一组基.
			\item\label{ex:bil:A5} 设线性映射 $\mathscr T:\mathbb{R}^2\to\mathbb{R}^3$ 在标准基下的矩阵为 $\begin{pmatrix}1&2\\0&1\\3&-1\end{pmatrix}$. 求对偶映射 $\mathscr T^*$ 在对偶基下的矩阵，并计算 $\mathscr T^*$ 作用在线性函数 $\vv y\mapsto y_1+y_2+y_3$ 上的结果.
			\item\label{ex:bil:A6} 求双线性型 $f(\vv x,\vv y)=x_1y_1+2x_1y_2-x_2y_1+3x_2y_2$ 在标准基下的度量矩阵，以及在基 $(1,1)^\trans$，$(1,-1)^\trans$ 下的度量矩阵，并验证 $\vv B=\vv P^\trans\vv{AP}$.
			\item\label{ex:bil:A7} 把 A6 中的 $f$ 分解为对称部分与交错部分之和.
			\item\label{ex:bil:A8} 判断下列双线性型是否非退化：(1) $x_1y_1-x_2y_2$；(2) $x_1y_2+x_2y_1$；(3) $(x_1+x_2)(y_1+y_2)$.
			\item\label{ex:bil:A9} 写出对称双线性型 $f(\vv x,\vv y)=x_1y_1+x_1y_2+x_2y_1+3x_2y_2$ 对应的二次型；反过来，由二次型 $x_1^2-4x_1x_2+2x_2^2$ 用极化恒等式求对应的对称双线性型.
			\item\label{ex:bil:A10} 验证 $\mathbb{R}^2$ 上的 $f(\vv x,\vv y)=\det(\vv x,\vv y)$ 是交错双线性型，写出它的度量矩阵. 它非退化吗？
			\item\label{ex:bil:A11} 设 $a_1,a_2,a_3$ 互异，$f_i(p)=p(a_i)$ 是 $\mathbb{R}[x]_3$ 上的线性函数. 对 $a_i=0,1,2$，求以 $f_1,f_2,f_3$ 为对偶基的 $\mathbb{R}[x]_3$ 的基.
			\item\label{ex:bil:A12} 在 $\mathbb{R}^4$ 中，设 $W$ 由 $x_1+x_2=0$，$x_3-x_4=0$ 定义. 求 $\dim W$ 与 $\dim W^0$，并写出 $W^0$ 的一组基.
			\item\label{ex:bil:A13} 证明 $\langle\vv A,\vv B\rangle=\tr(\vv{AB})$ 是 $M_n(\mathbb{R})$ 上非退化的对称双线性型. 它是内积吗？
			\item\label{ex:bil:A14} 对 $\mathbb{R}[x]_3$ 上的线性函数 $\varphi(p)=\int_0^1p(x)\,\mathrm dx$，求它在 $1,x,x^2$ 的对偶基下的坐标.
			\item\label{ex:bil:A15} （Riesz 表示定理的有限维情形）设 $V$ 是有限维实内积空间. 证明：对每个 $f\in V^*$，存在唯一的 $\vv\beta\in V$ 使 $f(\vv\alpha)=\langle\vv\alpha,\vv\beta\rangle$ 对一切 $\vv\alpha$ 成立.（提示：取标准正交基 $\vv e_1,\ldots,\vv e_n$，令 $\vv\beta=\sum f(\vv e_i)\vv e_i$.）
		\end{exA}
		
		\Bgroup
		\begin{exB}
			\item\label{ex:bil:B1} 证明：$f_1,\ldots,f_k\in V^*$ 线性无关，当且仅当映射 $V\to\mathbb{F}^k$，$\vv\alpha\mapsto(f_1(\vv\alpha),\ldots,f_k(\vv\alpha))^\trans$ 是满射.
			\item\label{ex:bil:B2} 设 $f,g\in V^*$ 且 $\ker f\subset\ker g$. 证明存在常数 $c$ 使 $g=cf$. 推广：若 $\ker f_1\cap\cdots\cap\ker f_k\subset\ker g$，则 $g$ 是 $f_1,\ldots,f_k$ 的线性组合.
			\item\label{ex:bil:B3} 证明：$\ker f_1\cap\cdots\cap\ker f_k=\{\zero\}$ 当且仅当 $f_1,\ldots,f_k$ 张成 $V^*$. 用方程组的语言解释这个结论.
			\item\label{ex:bil:B4} 证明：$\mathscr T$ 是单射当且仅当 $\mathscr T^*$ 是满射；$\mathscr T$ 是满射当且仅当 $\mathscr T^*$ 是单射.
			\item\label{ex:bil:B5} （Simpson 公式）求常数 $a,b,c$，使 $\int_0^2p(x)\,\mathrm dx=a\,p(0)+b\,p(1)+c\,p(2)$ 对一切 $p\in\mathbb{R}[x]_3$ 成立. 用对偶基的语言说明为什么这样的 $a,b,c$ 存在且唯一. 验证这个公式对三次多项式也成立.
			\item\label{ex:bil:B6} 证明：$V$ 上全体双线性型构成的空间与 $M_n(\mathbb{F})$ 同构，其中对称型、交错型分别构成维数为 $\frac{n(n+1)}2$、$\frac{n(n-1)}2$ 的子空间，且整个空间是二者的直和.
			\item\label{ex:bil:B7} 证明：若 $V$ 上存在非退化的交错双线性型，则 $\dim V$ 是偶数.
			\item\label{ex:bil:B8} 求双线性型 $f(\vv x,\vv y)=x_1y_1+x_1y_2+x_2y_1+x_2y_2+x_3y_3$ 的``根'' $\{\vv\alpha:f(\vv\alpha,\cdot)=0\}$.
			\item\label{ex:bil:B9} 在 $\mathbb{R}^2$ 上取 $f(\vv x,\vv y)=x_1y_1-x_2y_2$. 求全部``迷向向量''（$f(\vv\alpha,\vv\alpha)=0$），并说明它们构成两条直线. 再证明保持 $f$ 的线性变换 $\vv P$（$\vv P^\trans\diag(1,-1)\vv P=\diag(1,-1)$，$\det\vv P=1$，且保持 $x_1>|x_2|$ 的区域）都形如 $\begin{pmatrix}\cosh t&\sinh t\\ \sinh t&\cosh t\end{pmatrix}$.
			\item\label{ex:bil:B10} 设 $W$ 是 $V$ 的子空间，$\mathscr T:V\to V$ 线性且 $\mathscr T(W)\subset W$. 证明 $\mathscr T^*(W^0)\subset W^0$.
		\end{exB}
		
		\Cgroup
		\begin{exC}
			\item\label{ex:bil:C1} 证明：对无限维空间 $\mathbb{F}[x]$，自然映射 $\vv\alpha\mapsto\hat{\vv\alpha}$ 是单射但不是满射.（提示：$\mathbb{F}[x]^*$ 中的元素由它在 $1,x,x^2,\ldots$ 上的取值任意决定，可以``无限多个非零''.）
			\item\label{ex:bil:C2} 证明：$V$ 上的双线性型与线性映射 $V\to V^*$ 一一对应：$f\mapsto(\vv\alpha\mapsto f(\vv\alpha,\cdot))$. 用这个对应和对偶基的换基规则，重新推出度量矩阵的换基公式 $\vv B=\vv P^\trans\vv{AP}$.
			\item\label{ex:bil:C3} 证明：二阶矩阵 $\vv P$ 满足 $\vv P^\trans\vv J\vv P=\vv J$（$\vv J=\begin{pmatrix}0&1\\-1&0\end{pmatrix}$）当且仅当 $\det\vv P=1$. 四阶时，满足 $\vv P^\trans\diag(\vv J,\vv J)\vv P=\diag(\vv J,\vv J)$ 的矩阵（辛矩阵）的行列式是多少？
			\item\label{ex:bil:C4} 求 $M_n(\mathbb{R})$ 上对称双线性型 $\tr(\vv{AB})$ 的正、负惯性指数.（提示：对称矩阵与反对称矩阵在这个型下``互相正交''.）
			\item\label{ex:bil:C5} 设 $f$ 是 $V$ 上的对称双线性型，$W$ 是子空间，$W^\perp=\{\vv\alpha:f(\vv\alpha,\vv\beta)=0\ \forall\vv\beta\in W\}$. 证明 $\dim W^\perp\ge\dim V-\dim W$，$f$ 非退化时等号成立；举例说明即使 $f$ 非退化，$W\cap W^\perp$ 也可能不为零.
			\item\label{ex:bil:C6} 证明：$\mathbb{F}^n$ 上全体交错 $n$ 重线性函数构成的空间是一维的，由行列式张成. 这给出第五章章末习题~\ref{ex:det:C1} 的``对偶空间版本''.
			\item\label{ex:bil:C7} （线性规划对偶）考虑问题：在 $\vv{Ax}=\vv b$，$\vv x\ge\zero$ 下最小化 $\vv c^\trans\vv x$. 证明对任意满足 $\vv A^\trans\vv y\le\vv c$ 的 $\vv y$ 与任意可行的 $\vv x$，都有 $\vv b^\trans\vv y\le\vv c^\trans\vv x$（弱对偶）. 用``价格是数量空间上的线性函数''解释 $\vv y$ 的经济含义（影子价格）.
		\end{exC}
	}.
	%%   2. 在每章末尾（原 \begin{problemset}...\end{problemset} 之后，或替换它）
	%%      写对应的命令：
	%%        第一章 \ExElimination      第二章 \ExVectorSpace   第三章 \ExMatrix
	%%        第四章 \ExLinearMap        第五章 \ExDeterminant   第六章 \ExSolutions
	%%        第七章 \ExPolynomial       第八章 \ExEigen         第九章 \ExInnerProduct
	%%        第十章 \ExDecomposition    第十一章 \ExQuadratic   第十二章 \ExBilinear
	%%   3. 分组标准
	%%        A 组：直接运用本章定义与定理的计算和证明；
	%%        B 组：需要综合多个结论、或有一定技巧的计算和证明；
	%%        C 组：有挑战性的题目，或超出本课程大纲、为继续深造准备的内容.
	%%   4. 题目与正文中已有的节末练习、原 problemset 不重复；个别题目注明``接正文某题''.
	%%      第六章按扩充版（ch06-solutions.tex）的内容编写.
	%%   5. 每道题都带标签 ex:<章键>:<组号>，正文、习题与附录答案一律用 \ref 引用，
	%%      渲染为``A4''``B12''这样的题号. 章键依次为
	%%        elim vs mat lmap det sol poly eig ip dec quad bil（第一至十二章）.
	%%      例：第一章章末习题~\ref{ex:elim:B12}. 新增题目请加在组末，并补上标签.
	%%      A 组答案、B 组提示在 answers.tex 中，按同一标签对应.
	%% 依赖：导言区的 \vv \zero \trans \spn \rk \imm \diag，以及 enumitem（elegantbook 已加载）.
	%% 引用的前文标签：thm:sameSolution thm:solvability prop:linspc eg:pricedecomp cor:homogeneous thm:rouche thm:sol-determinable
	%%   thm:spectral thm:posdef thm:svd thm:orthcanon
	%% =====================================================================
	
	\providecommand{\tr}{\operatorname{tr}}
	\newcommand{\exsetheading}{%
		\section*{章末习题}%
		\markright{章末习题}%
	}
\newcommand{\exgroup}[2]{%
	\par\medskip{\centering\sffamily\bfseries #1\quad #2\par}\smallskip}
	\newenvironment{exA}{\begin{enumerate}[label=\textbf{A\arabic*.},ref=A\arabic*,leftmargin=2.8em,itemsep=4pt]}{\end{enumerate}}
	\newenvironment{exB}{\begin{enumerate}[label=\textbf{B\arabic*.},ref=B\arabic*,leftmargin=2.8em,itemsep=4pt]}{\end{enumerate}}
	\newenvironment{exC}{\begin{enumerate}[label=\textbf{C\arabic*.},ref=C\arabic*,leftmargin=2.8em,itemsep=4pt]}{\end{enumerate}}
	\newcommand{\Agroup}{\exgroup{A 组}{基础巩固}}
	\newcommand{\Bgroup}{\exgroup{B 组}{综合提升}}
	\newcommand{\Cgroup}{\exgroup{C 组}{拓展探索}}
	
	% =====================================================================
	% 第一章　线性方程组与消元法
	% =====================================================================
	\newcommand{\ExElimination}{%
		\exsetheading
		\Agroup
		\begin{exA}
			\item\label{ex:elim:A1} 判断下列矩阵哪些是阶梯形，哪些是简化阶梯形：
			\[
			\begin{pmatrix}1&2&0\\0&0&1\end{pmatrix},\quad
			\begin{pmatrix}1&0&3\\0&2&1\\0&0&0\end{pmatrix},\quad
			\begin{pmatrix}0&1&0\\1&0&0\end{pmatrix},\quad
			\begin{pmatrix}1&0&2&0\\0&1&-1&0\\0&0&0&1\end{pmatrix},\quad
			\begin{pmatrix}0&0\\0&1\end{pmatrix}.
			\]
			\item\label{ex:elim:A2} 写出方程组 $3x_1-x_3=2$，$x_2+4x_3=0$，$x_1-x_2=5$ 的增广矩阵；反过来，写出以 $\begin{pmatrix}1&-2&0&3\\0&1&5&-1\end{pmatrix}$ 为增广矩阵的方程组.
			\item\label{ex:elim:A3} 用消元法解方程组 $\left\{\begin{array}{l}x_1+2x_2+3x_3=14,\\2x_1-x_2+x_3=3,\\3x_1+x_2-x_3=2.\end{array}\right.$
			\item\label{ex:elim:A4} 用消元法解方程组 $\left\{\begin{array}{l}x_1+x_2+x_3+x_4=2,\\x_1-x_2+x_3-x_4=4,\\2x_1+x_2-x_3+3x_4=-1,\\x_1+2x_2+3x_3-2x_4=5.\end{array}\right.$
			\item\label{ex:elim:A5} 判断方程组 $\left\{\begin{array}{l}x_1-x_2+x_3=1,\\2x_1+x_2-x_3=2,\\x_1+2x_2-2x_3=3\end{array}\right.$ 是否有解，并说明消元过程中哪一行``暴露''了结论.
			\item\label{ex:elim:A6} 求方程组 $\left\{\begin{array}{l}x_1+2x_2-x_3+x_4=3,\\2x_1+4x_2-x_3+3x_4=7,\\-x_1-2x_2+2x_3=-2\end{array}\right.$ 的全部解，并指出主变量与自由变量.
			\item\label{ex:elim:A7} 求齐次方程组 $\left\{\begin{array}{l}x_1+x_2-x_3-x_4=0,\\2x_1-5x_2+3x_3+2x_4=0,\\7x_1-7x_2+3x_3+x_4=0\end{array}\right.$ 的全部解.
			\item\label{ex:elim:A8} 把矩阵 $\begin{pmatrix}2&4&-2&6\\1&2&1&1\\3&6&1&5\end{pmatrix}$ 化为简化阶梯形，并写出每一步所用的初等行变换.
			\item\label{ex:elim:A9} 写出全部 $2\times2$ 简化阶梯形矩阵的``形状''（可以含参数），并说明共有几类.
			\item\label{ex:elim:A10} $k$ 取何值时，方程组 $x+y=1$，$x-y=3$，$2x+ky=5$ 有解？
			\item\label{ex:elim:A11} 设 $a,b$ 为常数. 方程组 $x+2y=a$，$3x+6y=b$ 何时有解？有解时写出全部解.
			\item\label{ex:elim:A12} 求过三点 $(1,2)$，$(2,3)$，$(3,6)$ 的抛物线 $y=ax^2+bx+c$.
			\item\label{ex:elim:A13} 求过四点 $(-1,2)$，$(0,1)$，$(1,0)$，$(2,5)$ 的三次多项式函数 $y=a+bx+cx^2+dx^3$.
			\item\label{ex:elim:A14} 某剧场三种票价分别为 $30$，$50$，$80$ 元. 某场演出共售出 $100$ 张票，票房 $4800$ 元，且 $80$ 元票的张数恰为 $30$ 元票的一半. 三种票各售出多少张？
			\item\label{ex:elim:A15} 配平化学方程式 $x_1\,\mathrm{C_3H_8}+x_2\,\mathrm{O_2}\to x_3\,\mathrm{CO_2}+x_4\,\mathrm{H_2O}$：对 C，H，O 三种原子分别写出守恒方程，解这个齐次方程组，并取最小的正整数解.
			\item\label{ex:elim:A16} 某城区四个路口 $A,B,C,D$ 由单行道依次连成环路 $A\to B\to C\to D\to A$，环路上的车流量依次记为 $x_1$（$A\to B$），$x_2$（$B\to C$），$x_3$（$C\to D$），$x_4$（$D\to A$）（单位：辆/小时）. 各路口与外部道路的进出量为：$A$ 驶入 $400$、驶出 $300$；$B$ 驶入 $200$、驶出 $350$；$C$ 驶入 $250$、驶出 $100$；$D$ 驶入 $150$、驶出 $250$. 设每个路口驶入总量等于驶出总量.
			\begin{enumerate}
				\item 列出方程组并求全部解；
				\item 若各路段车流量都不能为负，$x_1$ 最少是多少？
				\item 若把 $D$ 的驶出量改为 $350$，方程组会怎样？这反映了什么守恒律？
			\end{enumerate}
			\item\label{ex:elim:A17} 证明：二元一次方程组 $ax+by=e$，$cx+dy=f$ 对任意的 $e,f$ 都有唯一解，当且仅当 $ad-bc\ne0$.（只用消元法，不用行列式.）
			\item\label{ex:elim:A18} 证明：若线性方程组有两个不同的解 $\vv x_1,\vv x_2$，则对任意数 $t$，$\vv x_1+t(\vv x_2-\vv x_1)$ 都是解. 由此说明：线性方程组的解要么不存在，要么唯一，要么有无穷多个.
			\item\label{ex:elim:A19} 证明方程组 $\{x+y+z=3,\ x-y=0\}$ 与 $\{2x+z=3,\ x+y+z=3\}$ 同解.
			\item\label{ex:elim:A20} 举例说明：若\textbf{同时}做``第一个方程加到第二个方程上''与``第二个方程加到第一个方程上''（两步都用原来的方程），所得方程组可能与原方程组不同解. 这与定理~\ref{thm:sameSolution}（初等行变换不改变解集）矛盾吗？
			\item\label{ex:elim:A21} 已知一个方程组的增广矩阵化为阶梯形后为 $\begin{pmatrix}1&3&-1&2&0\\0&0&2&1&4\\0&0&0&0&a\end{pmatrix}$. 讨论 $a$ 取何值时方程组有解；有解时它有几个自由变量？
		\end{exA}
		
		\Bgroup
		\begin{exB}
			\item\label{ex:elim:B1} 讨论 $a,b$ 取何值时，方程组
			\[\left\{\begin{array}{l}x_1+x_2+x_3+x_4=0,\\x_2+2x_3+2x_4=1,\\-x_2+(a-3)x_3-2x_4=b,\\3x_1+2x_2+x_3+ax_4=-1\end{array}\right.\]
			有唯一解、无解、有无穷多解；有无穷多解时求出全部解.
			\item\label{ex:elim:B2} 讨论 $\lambda$ 取何值时，方程组
			$(1+\lambda)x_1+x_2+x_3=0$，$x_1+(1+\lambda)x_2+x_3=3$，$x_1+x_2+(1+\lambda)x_3=\lambda$
			有唯一解、无解、有无穷多解.
			\item\label{ex:elim:B3} 设 $n$ 个方程、$n$ 个未知量的方程组对应的齐次方程组只有零解. 证明：无论常数项取什么值，原方程组都有唯一解.（提示：考察简化阶梯形中主元的个数.）
			\item\label{ex:elim:B4} 证明：线性方程组无解，当且仅当可以把若干个方程分别乘以适当的数再相加，得到矛盾方程 $0=1$.
			\item\label{ex:elim:B5} 证明：交换两行这一初等变换，可以由``倍加''与``某行乘以 $-1$''这两类变换的组合实现. 由此说明三类初等行变换并不是``独立''的.
			\item\label{ex:elim:B6} 证明：若线性方程组的系数和常数项都是有理数，且它有实数解，则它一定有有理数解；若它的实数解唯一，则这个解就是有理数解.
			\item\label{ex:elim:B7} 讨论方程组
			\[x_1-x_2=a_1,\quad x_2-x_3=a_2,\quad\ldots,\quad x_{n-1}-x_n=a_{n-1},\quad x_n-x_1=a_n\]
			有解的充要条件，并在有解时求全部解.
			\item\label{ex:elim:B8} 讨论方程组
			\[x_1+x_2=a_1,\quad x_2+x_3=a_2,\quad\ldots,\quad x_{n-1}+x_n=a_{n-1},\quad x_n+x_1=a_n\]
			的解. 证明：$n$ 为奇数时解唯一；$n$ 为偶数时有解当且仅当 $a_1-a_2+a_3-\cdots-a_n=0$.
			\item\label{ex:elim:B9} 某经济体有三个部门，直接消耗系数矩阵与最终需求分别为
			\[\vv A=\begin{pmatrix}0.1&0.2&0.2\\0.2&0.1&0.2\\0.1&0.3&0.1\end{pmatrix},\qquad \vv d=\begin{pmatrix}20\\130\\65\end{pmatrix}.\]
			按第一章``投入产出模型''的方式列出方程组，求各部门的总产出.
			\item\label{ex:elim:B10} 估计用消元法（化为阶梯形后回代）解 $n$ 个方程、$n$ 个未知量的方程组所需的乘除法次数，证明它约为 $\frac13n^3$（即与 $\frac13n^3$ 之比在 $n\to\infty$ 时趋于 $1$）.
			\item\label{ex:elim:B11} 设三个平面 $a_ix+b_iy+c_iz=d_i$（$i=1,2,3$）两两相交. 用消元法的语言说明：它们恰好交于一点、交于一条直线、或没有公共点（三条交线互相平行），并各举一例.
			\item\label{ex:elim:B12} 证明：$n$ 元齐次方程组的解集对加法和数乘封闭；而非齐次方程组的两个解之差是对应齐次方程组的解. 由此说明：非齐次方程组的全部解，等于它的任一个解加上齐次方程组的全部解.
			\item\label{ex:elim:B13} 一家律所有甲、乙、丙三位律师，本月分别承办了若干件合同、劳动、知识产权纠纷. 已知：三类案件合计 $30$ 件；甲承办的数量 $=$ 乙 $+$ 丙；每件合同、劳动、知识产权案件分别需要 $10$，$6$，$15$ 个工时，三人本月共用 $300$ 个工时……
			这些条件够不够确定每位律师承办了多少案件？请自己补全一个合理的设定（明确未知量是什么），列出方程组，用自由变量的个数说明还缺多少条独立的信息.
			\item\label{ex:elim:B14} 设方程组的系数矩阵每一行的元素之和都等于 $0$. 证明 $x_1=x_2=\cdots=x_n=1$ 是对应齐次方程组的解；若非齐次方程组有解 $\vv x_0$，则对任意数 $t$，$\vv x_0+t(1,\ldots,1)^\trans$ 也是解.
			\item\label{ex:elim:B15} 证明：如果方程组 (I) 的每个方程都是方程组 (II) 各方程的``线性组合''（各乘一个数再相加），反之亦然，那么 (I) 与 (II) 同解. 逆命题对于有解的方程组成立吗？先在两个未知量的情形下思考.
			\item\label{ex:elim:B16} 讨论 $\lambda$ 取何值时，方程组 $\lambda x+y+z=1$，$x+\lambda y+z=\lambda$，$x+y+\lambda z=\lambda^2$ 有唯一解、无解、有无穷多解；有唯一解时求出解.（第六章例~\ref{eg:lambda} 将用秩与行列式重新组织这个讨论.）
		\end{exB}
		
		\Cgroup
		\begin{exC}
			\item\label{ex:elim:C1} （百鸡问题）《张丘建算经》（约公元 5 世纪）：``今有鸡翁一，直钱五；鸡母一，直钱三；鸡雏三，直钱一. 凡百钱买鸡百只，问鸡翁、母、雏各几何？''列出线性方程组，求它的全部解，再从中找出所有\textbf{非负整数}解. 线性代数在哪一步之后就无能为力了？
			\item\label{ex:elim:C2} （幻方）若一个 $3\times3$ 数表的每行、每列和两条对角线上的三个数之和都等于同一个数 $s$，称它为幻方（不要求数互不相同）. 把 $9$ 个元素作为未知量、$s$ 也作为未知量，写出方程组，证明：
			\begin{enumerate}
				\item 中心元素一定等于 $s/3$；
				\item 全部幻方可以用 $3$ 个自由参数表示，写出通式；
				\item 用 $1,2,\ldots,9$ 各一次填成的幻方本质上只有一个（其余由旋转、翻转得到）.
			\end{enumerate}
			\item\label{ex:elim:C3} （点灯游戏）$3\times3$ 方格中每格有一盏灯，按下某格会使该格及其上下左右相邻的格子开关状态翻转. 在``$1+1=0$''的算术（即模 $2$ 运算）中，把``按哪些格''作为未知量，把``从某一状态变为全灭''写成线性方程组. 证明：从任意初始状态出发，总能把灯全部熄灭.（模 $2$ 运算的集合 $\{0,1\}$ 不是本书意义下的数域，但消元法照样适用——想一想为什么.）
			\item\label{ex:elim:C4} （舍入误差与主元）在``只保留三位有效数字''的计算机上解方程组 $0.0001x+y=1$，$x+y=2$：
			\begin{enumerate}
				\item 直接用第一个方程消去第二个方程中的 $x$；
				\item 先交换两个方程再消元.
			\end{enumerate}
			把两种结果与精确解比较，说明数值软件为什么要``选绝对值最大的元素作主元''.
			\item\label{ex:elim:C5} 设两个都有解的方程组 (I)(II) 都有 $m$ 个方程、$n$ 个未知量，且同解. 证明：它们的增广矩阵行等价.（提示：先证明二者的简化阶梯形的非零行完全相同.）
			\item\label{ex:elim:C6} （Jacobi 迭代）对方程组 $4x+y=9$，$x+3y=5$，把它改写为 $x=(9-y)/4$，$y=(5-x)/3$，从 $(x_0,y_0)=(0,0)$ 出发反复代入. 计算前五步，观察是否收敛到真解. 再对 $x+4y=9$，$3x+y=5$ 做同样的事，会发生什么？试找出一个保证收敛的充分条件.
			\item\label{ex:elim:C7} （电路）一个电路中有三条支路电流 $I_1,I_2,I_3$. 节点 $A$ 处 $I_1=I_2+I_3$；两个回路的电压方程为 $2I_1+I_2=8$，$I_2-2I_3=0$（单位从略）. 求各电流. 再说明：若把第三个节点 $B$ 的电流方程也写进来，方程组为什么不会因此``多一个条件''？
			\item\label{ex:elim:C8} 证明：一个 $m\times n$ 矩阵（$m,n\ge1$）通过初等行变换所能得到的不同的简化阶梯形``形状''（即主元所在列的位置组合）共有 $\sum_{r=0}^{\min(m,n)}\binom nr$ 种.
			\item\label{ex:elim:C9} 证明：一个矩阵的简化阶梯形是唯一的.（提示：设 $\vv R_1,\vv R_2$ 都是 $\vv M$ 的简化阶梯形，则以它们为系数矩阵的齐次方程组同解. 对列数归纳：删去最后一列后由归纳假设前面各列相同，再比较最后一列，考察它在两者中是否为主元列.）
		\end{exC}
	}
	
	% =====================================================================
	% 第二章　向量空间
	% =====================================================================
	\newcommand{\ExVectorSpace}{%
		\exsetheading
		\Agroup
		\begin{exA}
			\item\label{ex:vs:A1} 下列集合关于通常的四则运算是否构成数域？
			(1) $\{a+b\sqrt3\mid a,b\in\mathbb{Q}\}$；(2) $\{a+b\sqrt2\mid a\in\mathbb{Q},\ b\in\mathbb{Z}\}$；(3) 全体偶数；(4) $\{a+b\sqrt[3]{2}\mid a,b\in\mathbb{Q}\}$；(5) 全体纯虚数与 $0$.
			\item\label{ex:vs:A2} 下列集合关于通常的加法与数乘是否构成实向量空间？
			(1) 全体 $n$ 阶实对称矩阵；(2) $\{(x,y)^\trans\in\mathbb{R}^2\mid x\ge0\}$；(3) 次数恰为 $n$ 的实系数多项式全体；(4) $\{f\in\mathbb{R}[x]\mid f(1)=0\}$；(5) 满足 $a_{n+2}=a_{n+1}+a_n$ 的全体实数列；(6) $\{(x,y,z)^\trans\mid x+y+z=1\}$.
			\item\label{ex:vs:A3} 判断下列向量组的线性相关性：
			\begin{enumerate}
				\item $(1,2,-1)^\trans$，$(2,-3,1)^\trans$，$(4,1,-1)^\trans$；
				\item 在 $\mathbb{R}[x]_3$ 中：$1+x$，$1-x$，$x^2$；
				\item 在 $\mathbb{R}[x]_3$ 中：$1+x+x^2$，$1-x$，$2+x^2$，$x+x^2$；
				\item 在 $M_2(\mathbb{R})$ 中：$\begin{pmatrix}1&0\\0&1\end{pmatrix}$，$\begin{pmatrix}0&1\\1&0\end{pmatrix}$，$\begin{pmatrix}1&1\\1&1\end{pmatrix}$，$\begin{pmatrix}1&0\\0&0\end{pmatrix}$.
			\end{enumerate}
			\item\label{ex:vs:A4} 把 $\vv\beta=(1,2,1,1)^\trans$ 表示为 $\vv\alpha_1=(1,1,1,1)^\trans$，$\vv\alpha_2=(1,1,-1,-1)^\trans$，$\vv\alpha_3=(1,-1,1,-1)^\trans$，$\vv\alpha_4=(1,-1,-1,1)^\trans$ 的线性组合.
			\item\label{ex:vs:A5} $t$ 取何值时，$(1,1,0)^\trans$，$(1,3,-1)^\trans$，$(5,3,t)^\trans$ 线性相关？
			\item\label{ex:vs:A6} 求向量组 $\vv\alpha_1=(1,-1,2,4)^\trans$，$\vv\alpha_2=(0,3,1,2)^\trans$，$\vv\alpha_3=(3,0,7,14)^\trans$，$\vv\alpha_4=(1,-1,2,0)^\trans$，$\vv\alpha_5=(2,1,5,6)^\trans$ 的秩和一个极大无关组，并把其余向量用这个极大无关组线性表出.
			\item\label{ex:vs:A7} 设 $\vv\alpha_1,\vv\alpha_2,\vv\alpha_3$ 线性无关. 判断下列向量组的线性相关性：
			(1) $\vv\alpha_1-\vv\alpha_2$，$\vv\alpha_2-\vv\alpha_3$，$\vv\alpha_3-\vv\alpha_1$；
			(2) $\vv\alpha_1$，$\vv\alpha_1+\vv\alpha_2$，$\vv\alpha_1+\vv\alpha_2+\vv\alpha_3$；
			(3) $\vv\alpha_1+2\vv\alpha_2$，$\vv\alpha_2+2\vv\alpha_3$，$\vv\alpha_3+2\vv\alpha_1$.
			\item\label{ex:vs:A8} 证明：含有零向量的向量组线性相关；含有两个相同向量的向量组线性相关；单个向量 $\vv\alpha$ 线性无关当且仅当 $\vv\alpha\ne\zero$.
			\item\label{ex:vs:A9} 证明：两个向量线性相关，当且仅当其中一个是另一个的倍数. 三个向量的情形能这样说吗？
			\item\label{ex:vs:A10} 只用向量空间的八条公理与命题~\ref{prop:linspc}，证明：(1) $\vv\alpha+\vv\alpha=2\vv\alpha$；(2) 若 $k\ne0$ 且 $k\vv\alpha=k\vv\beta$，则 $\vv\alpha=\vv\beta$；(3) $-(-\vv\alpha)=\vv\alpha$.
			\item\label{ex:vs:A11} 证明下列集合是相应空间的子空间，并求一组基和维数：
			\begin{enumerate}
				\item $\{\vv x\in\mathbb{R}^4\mid x_1+x_2=x_3+x_4\}$；
				\item $n$ 阶上三角矩阵全体；
				\item 迹为零的 $n$ 阶方阵全体 $\{\vv X\mid x_{11}+\cdots+x_{nn}=0\}$；
				\item $\{f\in\mathbb{R}[x]_4\mid f(1)=0\}$.
			\end{enumerate}
			\item\label{ex:vs:A12} 判断下列子集是否为子空间：(1) $\{(a,b,c)^\trans\mid a=2b\}$；(2) $\{(a,b,c)^\trans\mid ab=0\}$；(3) $\{(a,b,c)^\trans\mid a,b,c\in\mathbb{Q}\}\subset\mathbb{R}^3$；(4) $\{f\in C[0,1]\mid f(0)=2f(1)\}$.
			\item\label{ex:vs:A13} 求 $\vv\alpha_1=(1,2,1,0)^\trans$，$\vv\alpha_2=(-1,1,1,1)^\trans$，$\vv\alpha_3=(0,3,2,1)^\trans$，$\vv\alpha_4=(1,0,1,2)^\trans$ 张成的子空间的维数和一组基.
			\item\label{ex:vs:A14} 求多项式 $2x^2-x+3$ 在基 $1,\ x-1,\ (x-1)^2$ 下的坐标.
			\item\label{ex:vs:A15} 求矩阵 $\begin{pmatrix}2&3\\-1&4\end{pmatrix}$ 在 $M_2(\mathbb{R})$ 的基 $\begin{pmatrix}1&0\\0&0\end{pmatrix},\begin{pmatrix}1&1\\0&0\end{pmatrix},\begin{pmatrix}1&1\\1&0\end{pmatrix},\begin{pmatrix}1&1\\1&1\end{pmatrix}$ 下的坐标.
			\item\label{ex:vs:A16} (1) $\mathbb{R}$ 是否是 $\mathbb{Q}$ 上的向量空间？(2) $\mathbb{Q}$ 是否是 $\mathbb{R}$ 上的向量空间？(3) $\mathbb{C}$ 作为 $\mathbb{R}$ 上的向量空间和作为 $\mathbb{C}$ 上的向量空间，维数各是多少？
			\item\label{ex:vs:A17} 写出 $\mathbb{F}[x]_n$ 与 $\mathbb{F}^n$ 之间、$M_2(\mathbb{F})$ 与 $\mathbb{F}^4$ 之间的一个同构映射，并验证它保持线性运算.
			\item\label{ex:vs:A18} 设 $W_1=\spn\big((1,0,0)^\trans,(0,1,0)^\trans\big)$，$W_2=\spn\big((1,1,1)^\trans\big)$，$W_3=\spn\big((1,1,0)^\trans\big)$. 分别判断 $W_1+W_2$，$W_1+W_3$ 是否为直和，并求和空间的维数.
			\item\label{ex:vs:A19} 设 $W_1=\spn\big((1,1,0,0)^\trans,(1,0,1,1)^\trans\big)$，$W_2=\spn\big((0,-1,1,1)^\trans,(1,2,0,1)^\trans\big)$. 求 $W_1+W_2$ 与 $W_1\cap W_2$ 的维数，并验证维数公式.
			\item\label{ex:vs:A20} 设 $\vv\alpha_1,\ldots,\vv\alpha_s$ 的秩为 $r$. 证明：其中任意 $r+1$ 个向量线性相关；任意 $r$ 个线性无关的向量都构成极大无关组.
		\end{exA}
		
		\Bgroup
		\begin{exB}
			\item\label{ex:vs:B1} 证明：两个子空间的并 $W_1\cup W_2$ 是子空间，当且仅当 $W_1\subset W_2$ 或 $W_2\subset W_1$.
			\item\label{ex:vs:B2} 证明：若向量组 $\vv\beta_1,\ldots,\vv\beta_t$ 可由 $\vv\alpha_1,\ldots,\vv\alpha_s$ 线性表出，且两组的秩相等，则两组等价.
			\item\label{ex:vs:B3} 设向量组 (I) 与 (II) 的秩分别为 $r_1,r_2$，把它们合并为 (III). 证明 $\max\{r_1,r_2\}\le\rk(\mathrm{III})\le r_1+r_2$. 再证明：$\rk(\vv\alpha_1+\vv\beta_1,\ldots,\vv\alpha_s+\vv\beta_s)\le\rk(\vv\alpha_1,\ldots,\vv\alpha_s)+\rk(\vv\beta_1,\ldots,\vv\beta_s)$.
			\item\label{ex:vs:B4} 设 $\vv\alpha_1,\ldots,\vv\alpha_n$ 是 $V$ 的基，$\vv\beta=a_1\vv\alpha_1+\cdots+a_n\vv\alpha_n$. 证明：把 $\vv\alpha_j$ 换成 $\vv\beta$ 后所得的向量组仍是基，当且仅当 $a_j\ne0$.
			\item\label{ex:vs:B5} 设 $p_0,p_1,\ldots,p_{n-1}\in\mathbb{F}[x]$ 满足 $\deg p_k=k$. 证明它们构成 $\mathbb{F}[x]_n$ 的一组基.
			\item\label{ex:vs:B6} 记 $V$ 为满足 $a_{n+2}=a_{n+1}+2a_n$（$n\ge1$）的全体实数列. 证明 $V$ 是向量空间，$\dim V=2$，并说明 $\{2^n\}$ 与 $\{(-1)^n\}$ 构成 $V$ 的一组基. 由此写出 $a_1=1$，$a_2=1$ 时的通项.
			\item\label{ex:vs:B7} 在 $C(\mathbb{R})$ 中：(1) 证明 $\mathrm e^x,\mathrm e^{2x},\mathrm e^{3x}$ 线性无关；(2) 证明 $\sin x,\cos x,\sin(x+1)$ 线性相关；(3) $1,\cos x,\cos2x,\cos^2x$ 的秩是多少？
			\item\label{ex:vs:B8} 证明 $\{a+b\sqrt2+c\sqrt3+d\sqrt6\mid a,b,c,d\in\mathbb{Q}\}$ 是数域，并证明 $1,\sqrt2,\sqrt3,\sqrt6$ 在 $\mathbb{Q}$ 上线性无关. 这个数域作为 $\mathbb{Q}$ 上的向量空间是几维的？
			\item\label{ex:vs:B9} 写出 $\mathbb{R}^2$ 与 $\mathbb{R}^3$ 的全部子空间（按维数分类，并说明几何意义）.
			\item\label{ex:vs:B10} 证明：$\mathbb{R}^5$ 中任意两个三维子空间的交至少是一维的. 一般地，$n$ 维空间中维数分别为 $k_1,k_2$ 的两个子空间，若 $k_1+k_2>n$，则交不为 $\{\zero\}$.
			\item\label{ex:vs:B11} 设 $a\in\mathbb{F}$. 证明 $\mathbb{F}[x]_n=\mathbb{F}\cdot1\oplus\{f\in\mathbb{F}[x]_n\mid f(a)=0\}$，并写出 $f$ 按这个直和的分解.
			\item\label{ex:vs:B12} 设 $V=C[-1,1]$，$E$，$O$ 分别为其中的偶函数、奇函数全体. 证明 $V=E\oplus O$，并把 $\mathrm e^x$ 按此分解.
			\item\label{ex:vs:B13} 设 $\dim V=n$，$W_1,\ldots,W_k$ 是 $V$ 的子空间，且 $\dim W_1+\cdots+\dim W_k>(k-1)n$. 证明 $W_1\cap\cdots\cap W_k\ne\{\zero\}$.
			\item\label{ex:vs:B14} 设 $\vv\alpha_1,\ldots,\vv\alpha_n$ 是 $V$ 的基，$W=\spn(\vv\alpha_1+\vv\alpha_2,\ \vv\alpha_2+\vv\alpha_3,\ \ldots,\ \vv\alpha_{n-1}+\vv\alpha_n)$. 求 $\dim W$，并求 $W$ 的一个补空间.
			\item\label{ex:vs:B15} 设 $V_1=\{\vv x\in\mathbb{F}^n\mid x_1=x_2=\cdots=x_n\}$，$V_2=\{\vv x\in\mathbb{F}^n\mid x_1=0\}$. 证明 $\mathbb{F}^n=V_1\oplus V_2$. 与第二章例~\ref{eg:pricedecomp} 的分解相比，哪一种分解更``自然''？（提示：思考补空间的不唯一性.）
			\item\label{ex:vs:B16} 在 $M_n(\mathbb{F})$ 中，设 $W_1$ 为对角线以下元素全为零的矩阵（上三角矩阵）全体，$W_2$ 为对角线及以上元素全为零的矩阵全体. 证明 $M_n(\mathbb{F})=W_1\oplus W_2$，并求 $\dim W_1$.
			\item\label{ex:vs:B17} （覆盖定理的另一个证明）设 $W_1,\ldots,W_s$ 是 $\mathbb{R}^n$ 的真子空间，令 $\vv v(t)=(1,t,t^2,\ldots,t^{n-1})^\trans$.
			\begin{enumerate}
				\item 证明：每个真子空间都包含于某个超平面 $\{\vv x\mid a_1x_1+\cdots+a_nx_n=0\}$（$a_i$ 不全为零）中.（提示：取 $W$ 的基并扩充为 $\mathbb{R}^n$ 的基，看``最后一个坐标''.）
				\item 证明：对上述超平面，满足 $\vv v(t)$ 落在其中的实数 $t$ 至多有 $n-1$ 个.
				\item 由此推出：$\vv v(t)$ 落在 $W_1\cup\cdots\cup W_s$ 中的 $t$ 至多有 $s(n-1)$ 个. 这给出了覆盖定理在 $\mathbb{R}^n$ 中的一个``构造性''证明. 对 $n=3$ 以及三个子空间 $x_1=x_2$，$x_2=x_3$，$x_1+x_2+x_3=0$，写出一个不在它们之中的 $\vv v(t)$.
			\end{enumerate}
			\item\label{ex:vs:B18} 设 $\dim V=n$，$W_1,\ldots,W_s$ 是 $V$ 的真子空间. 证明：存在 $V$ 的一组基 $\vv\varepsilon_1,\ldots,\vv\varepsilon_n$，其中\textbf{每个}向量都不属于 $W_1\cup\cdots\cup W_s$.（提示：逐个选取. 选 $\vv\varepsilon_k$ 时，要躲开的除了 $W_1,\ldots,W_s$，还有哪个真子空间？）
		\end{exB}
		
		\Cgroup
		\begin{exC}
			\item\label{ex:vs:C1} （覆盖定理在有限``数系''上的失效）设 $p$ 是素数，$\mathbb{F}_p=\{0,1,\ldots,p-1\}$ 按模 $p$ 作加法和乘法（可以承认它对加、减、乘以及除以非零元素都封闭. 它不包含 $\mathbb{Q}$，所以不是本书意义下的数域；但第二章中除定理~\ref{thm:Qsmallest} 与覆盖定理以外的定义和结论，在它上面照样成立）.
			\begin{enumerate}
				\item 证明 $\mathbb{F}_p^2$ 恰有 $p+1$ 个一维子空间，并且它们的并是整个 $\mathbb{F}_p^2$.
				\item 证明 $\mathbb{F}_p^2$ 不能写成 $p$ 个真子空间的并.（提示：每个一维子空间恰含 $p-1$ 个非零向量，数一数.）
				\item 指出定理~\ref{thm:cover} 的证明在 $\mathbb{F}_p$ 上从哪一步开始失效. 由第 2 问，``个数不超过 $p$ 时盖不住''，你能把正文的证明改写成这个``定量''的版本吗？
			\end{enumerate}
			\item\label{ex:vs:C2} 证明 $\ln2,\ln3,\ln5$ 在 $\mathbb{Q}$ 上线性无关. 由此说明 $\mathbb{R}$ 作为 $\mathbb{Q}$ 上的向量空间不是有限维的.（提示：算术基本定理.）
			\item\label{ex:vs:C3} 证明：全体实数列构成的向量空间不是有限维的；而``只有有限项非零''的实数列全体是它的子空间，且 $\vv e_1,\vv e_2,\ldots$（第 $i$ 项为 $1$，其余为 $0$）是它的一组``基''——每个元素都能唯一地写成其中\textbf{有限个}向量的线性组合. 这组向量为什么不是全体实数列的基？
			\item\label{ex:vs:C4} （模律）设 $W_1,W_2,W_3$ 是 $V$ 的子空间. 证明 $(W_1\cap W_3)+(W_2\cap W_3)\subset(W_1+W_2)\cap W_3$，在 $\mathbb{R}^2$ 中举例说明等号可以不成立；再证明：若 $W_1\subset W_3$，则等号一定成立.
			\item\label{ex:vs:C5} （有限域上的向量空间）在只含 $0,1$ 两个元素、按模 $2$ 运算的``数系'' $\mathbb{F}_2$ 上考虑 $\mathbb{F}_2^3$. (1) 它有多少个向量？(2) 有多少个一维子空间、多少个二维子空间？(3) 它有多少组（有序）基？推广到 $\mathbb{F}_2^n$.
			\item\label{ex:vs:C6} （塔律）设 $\mathbb{F}\subset\mathbb{K}\subset\mathbb{L}$ 是三个数域，$\mathbb{L}$ 作为 $\mathbb{K}$ 上的向量空间是 $m$ 维的，$\mathbb{K}$ 作为 $\mathbb{F}$ 上的向量空间是 $n$ 维的. 证明 $\mathbb{L}$ 作为 $\mathbb{F}$ 上的向量空间是 $mn$ 维的. 用它重新计算 B8 中数域的维数.
			\item\label{ex:vs:C7} 证明 $\mathbb{Q}(\sqrt2+\sqrt3)=\mathbb{Q}(\sqrt2,\sqrt3)$，即包含 $\sqrt2+\sqrt3$ 的最小数域与包含 $\sqrt2,\sqrt3$ 的最小数域相同.（提示：计算 $(\sqrt2+\sqrt3)^3$.）
			\item\label{ex:vs:C8} 设 $V$ 是 $n$ 维空间，$W$ 是 $k$ 维子空间（$0<k<n$）. 证明 $W$ 有无穷多个补空间. 在 $\mathbb{F}_2^2$ 中，一维子空间有几个补空间？
			\item\label{ex:vs:C9} （Fibonacci 空间的另一组基）在 B6 的记号下，改为研究满足 $a_{n+2}=a_{n+1}+a_n$ 的数列空间. 求出两个公比 $q$ 使 $\{q^n\}$ 属于这个空间，证明它们构成一组基，并推出 Fibonacci 数列的通项（Binet 公式）. 第八章会从另一个角度再次得到它.
		\end{exC}
	}
	
	% =====================================================================
	% 第三章　矩阵
	% =====================================================================
	\newcommand{\ExMatrix}{%
		\exsetheading
		\Agroup
		\begin{exA}
			\item\label{ex:mat:A1} 设 $\vv A=\begin{pmatrix}1&-1&2\\0&3&1\end{pmatrix}$，$\vv B=\begin{pmatrix}2&1\\-1&0\\1&4\end{pmatrix}$. 计算 $\vv{AB}$，$\vv{BA}$，$\vv A^\trans\vv B^\trans$，并验证 $(\vv{AB})^\trans=\vv B^\trans\vv A^\trans$.
			\item\label{ex:mat:A2} 按``列的线性组合''计算 $\begin{pmatrix}1&2&0\\3&-1&4\end{pmatrix}\begin{pmatrix}2\\-1\\1\end{pmatrix}$；再说明 $\vv e_i^\trans\vv A\vv e_j$ 是 $\vv A$ 的哪个元素，$\vv A\vv e_j$、$\vv e_i^\trans\vv A$ 又分别是什么.
			\item\label{ex:mat:A3} 计算 $\begin{pmatrix}1&1\\0&1\end{pmatrix}^n$，$\begin{pmatrix}\lambda&1\\0&\lambda\end{pmatrix}^n$，$\begin{pmatrix}\cos\theta&-\sin\theta\\ \sin\theta&\cos\theta\end{pmatrix}^n$，$\begin{pmatrix}0&1&0\\0&0&1\\0&0&0\end{pmatrix}^n$.
			\item\label{ex:mat:A4} 设 $\vv A=\begin{pmatrix}1&2\\0&1\end{pmatrix}$，$\vv B=\begin{pmatrix}1&0\\3&1\end{pmatrix}$. 计算 $(\vv A+\vv B)^2$ 与 $\vv A^2+2\vv{AB}+\vv B^2$，二者是否相等？$(\vv A+\vv B)(\vv A-\vv B)$ 与 $\vv A^2-\vv B^2$ 呢？
			\item\label{ex:mat:A5} 用初等行变换求下列矩阵的逆：
			\[\begin{pmatrix}1&2&3\\2&2&1\\3&4&3\end{pmatrix},\qquad\begin{pmatrix}1&1&1\\0&1&1\\0&0&1\end{pmatrix},\qquad \begin{pmatrix}2&1&0\\1&2&1\\0&1&2\end{pmatrix},\qquad\begin{pmatrix}0&0&0&1\\0&0&2&0\\0&3&0&0\\4&0&0&0\end{pmatrix}.\]
			\item\label{ex:mat:A6} 解矩阵方程：(1) $\begin{pmatrix}2&5\\1&3\end{pmatrix}\vv X=\begin{pmatrix}4&-6\\2&1\end{pmatrix}$；(2) $\vv X\begin{pmatrix}2&1\\1&1\end{pmatrix}=\begin{pmatrix}1&0\\0&2\end{pmatrix}$.
			\item\label{ex:mat:A7} 求矩阵 $\begin{pmatrix}1&1&a\\1&a&1\\a&1&1\end{pmatrix}$ 的秩（对 $a$ 分情况讨论）.
			\item\label{ex:mat:A8} 求矩阵 $\begin{pmatrix}1&2&-1&0&3\\2&4&1&3&0\\-1&-2&2&1&-3\\1&2&2&3&0\end{pmatrix}$ 的秩，以及列向量组的一个极大无关组，并用它表出其余各列.
			\item\label{ex:mat:A9} 写出 $3$ 阶初等矩阵 $\vv E(1,3)$，$\vv E(2(5))$，$\vv E(3,1(-2))$，并分别写出它们的逆. 计算 $\vv E(3,1(-2))\vv A$ 与 $\vv A\vv E(3,1(-2))$，其中 $\vv A$ 是任一 $3$ 阶方阵，说明效果.
			\item\label{ex:mat:A10} 把 $\begin{pmatrix}2&1\\4&3\end{pmatrix}$ 写成若干个初等矩阵的乘积.
			\item\label{ex:mat:A11} 证明：两个上三角矩阵之积仍是上三角矩阵，且积的对角元等于两者对角元之积. 可逆上三角矩阵的逆也是上三角矩阵吗？
			\item\label{ex:mat:A12} 设 $\vv A,\vv B$ 都是 $n$ 阶对称矩阵. 证明：(1) $\vv A^k$，$\vv A+\vv B$，$\vv P^\trans\vv A\vv P$ 都对称；(2) $\vv{AB}$ 对称当且仅当 $\vv{AB}=\vv{BA}$；(3) 对任意 $m\times n$ 矩阵 $\vv C$，$\vv C^\trans\vv C$ 与 $\vv C\vv C^\trans$ 都对称.
			\item\label{ex:mat:A13} (1) 设 $\vv x,\vv y\in\mathbb{F}^n$，证明 $\tr(\vv x\vv y^\trans)=\vv y^\trans\vv x$；(2) 设 $n$ 阶方阵 $\vv A$ 满足：对一切 $n$ 阶方阵 $\vv X$ 都有 $\tr(\vv{AX})=0$. 证明 $\vv A=\vv O$（提示：取 $\vv X=\vv E_{ji}$）；(3) 证明：不存在 $n$ 阶方阵 $\vv A,\vv B$ 使 $\vv{AB}-\vv{BA}=\vv I$.
			\item\label{ex:mat:A14} 设方阵 $\vv A$ 满足 $\vv A^2-\vv A-2\vv I=\vv O$. 证明 $\vv A$ 与 $\vv A+2\vv I$ 都可逆，并用 $\vv A$ 表示它们的逆. $\vv A+\vv I$ 一定可逆吗？
			\item\label{ex:mat:A15} 证明：若 $\vv A$ 可逆且 $\vv{AB}=\vv{AC}$，则 $\vv B=\vv C$；举例说明 $\vv A$ 不可逆时结论可能不成立.
			\item\label{ex:mat:A16} 设 $\vv A,\vv B$ 可逆. 求 $\begin{pmatrix}\vv A&\vv O\\ \vv O&\vv B\end{pmatrix}^{-1}$ 与 $\begin{pmatrix}\vv O&\vv A\\ \vv B&\vv O\end{pmatrix}^{-1}$，并用后者求 $\begin{pmatrix}0&0&1&2\\0&0&1&3\\2&0&0&0\\0&5&0&0\end{pmatrix}^{-1}$.
			\item\label{ex:mat:A17} 用分块乘法计算 $\begin{pmatrix}1&0&2&1\\0&1&1&3\\0&0&1&0\\0&0&0&1\end{pmatrix}^2$，并猜出 $n$ 次方的一般形式.
			\item\label{ex:mat:A18} 某公司在两个工厂生产三种产品，各工厂每天的产量（件）与三种产品的单位成本、单位售价（元）分别为
			\[\vv Q=\begin{pmatrix}100&80&60\\120&50&90\end{pmatrix}\ (\text{行：工厂；列：产品}),\qquad \vv C=\begin{pmatrix}20&30\\15&25\\40&55\end{pmatrix}\ (\text{列：成本、售价}).\]
			计算 $\vv Q\vv C$，说明它每个元素的含义，并求两个工厂每天的利润.
			\item\label{ex:mat:A19} 四座城市之间的直达航线为：$1\leftrightarrow2$，$1\leftrightarrow3$，$2\leftrightarrow3$，$3\leftrightarrow4$（双向）. 写出邻接矩阵 $\vv A$，计算 $\vv A^2$ 与 $\vv A^3$. $(\vv A^2)_{ii}$ 与城市 $i$ 的航线数有什么关系？从城市 $1$ 到城市 $4$ 恰好乘三次航班有几种走法？
			\item\label{ex:mat:A20} 设 $\vv A$ 是实反对称矩阵. 证明：对一切 $\vv x\in\mathbb{R}^n$ 有 $\vv x^\trans\vv A\vv x=0$；$\vv A^2$ 是对称矩阵.
			\item\label{ex:mat:A21} 证明：若 $\vv A$ 是 $m\times n$ 矩阵，$\vv P$ 是 $m$ 阶可逆矩阵，则 $\vv{PA}$ 与 $\vv A$ 的简化阶梯形相同.
			\item\label{ex:mat:A22} 证明：任一 $n$ 阶方阵都可以唯一地写成一个对称矩阵与一个反对称矩阵之和.（提示：$\vv A=\frac12(\vv A+\vv A^\trans)+\frac12(\vv A-\vv A^\trans)$.）
		\end{exA}
		
		\Bgroup
		\begin{exB}
			\item\label{ex:mat:B1} 分别求出全部满足 (1) $\vv A^2=\vv O$；(2) $\vv A^2=\vv I$；(3) $\vv A^2=\vv A$ 的二阶实方阵.
			\item\label{ex:mat:B2} 设 $a_1,\ldots,a_n$ 互不相同. 证明：与 $\diag(a_1,\ldots,a_n)$ 可交换的矩阵一定是对角矩阵.
			\item\label{ex:mat:B3} 证明：与\textbf{一切} $n$ 阶方阵都可交换的矩阵必为数量矩阵 $c\vv I$.（提示：考虑与 $\vv E_{ij}$ 交换.）
			\item\label{ex:mat:B4} 设 $\vv A,\vv B$ 都是幂零矩阵且 $\vv{AB}=\vv{BA}$. 证明 $\vv A+\vv B$ 与 $\vv{AB}$ 都是幂零矩阵. 举例说明``可交换''不能去掉.
			\item\label{ex:mat:B5} 证明：$\rk\begin{pmatrix}\vv A&\vv O\\ \vv O&\vv B\end{pmatrix}=\rk\vv A+\rk\vv B$，$\rk\begin{pmatrix}\vv A&\vv C\\ \vv O&\vv B\end{pmatrix}\ge\rk\vv A+\rk\vv B$，并举例说明后者可以取严格不等号.
			\item\label{ex:mat:B6} （Sherman–Morrison 公式）设 $\vv u,\vv v\in\mathbb{F}^n$，$1+\vv v^\trans\vv u\ne0$. 证明
			\[(\vv I+\vv u\vv v^\trans)^{-1}=\vv I-\frac{\vv u\vv v^\trans}{1+\vv v^\trans\vv u}.\]
			\item\label{ex:mat:B7} 设 $\vv A$ 为 $m\times n$ 矩阵，$\vv B$ 为 $n\times m$ 矩阵. 证明：若 $\vv I_m+\vv{AB}$ 可逆，则 $\vv I_n+\vv{BA}$ 也可逆，并且 $(\vv I_n+\vv{BA})^{-1}=\vv I_n-\vv B(\vv I_m+\vv{AB})^{-1}\vv A$.
			\item\label{ex:mat:B8} 证明：秩为 $1$ 的 $m\times n$ 矩阵都可以写成 $\vv u\vv v^\trans$（$\vv u\in\mathbb{F}^m$，$\vv v\in\mathbb{F}^n$ 非零）. 对这样的矩阵证明 $\vv A^2=(\tr\vv A)\vv A$（$m=n$ 时），并由此求 $\vv A^k$.
			\item\label{ex:mat:B9} 证明：秩为 $r$ 的矩阵可以写成 $r$ 个秩为 $1$ 的矩阵之和，但不能写成少于 $r$ 个秩为 $1$ 的矩阵之和.
			\item\label{ex:mat:B10} （满秩分解）证明：秩为 $r$ 的 $m\times n$ 矩阵 $\vv A$ 可以写成 $\vv A=\vv{BC}$，其中 $\vv B$ 是列满秩的 $m\times r$ 矩阵，$\vv C$ 是行满秩的 $r\times n$ 矩阵. 对 A8 中的矩阵写出一个满秩分解.
			\item\label{ex:mat:B11} 计算 $\begin{pmatrix}1&1&1\\1&1&1\\1&1&1\end{pmatrix}^n$ 与 $\begin{pmatrix}2&1&0\\0&2&1\\0&0&2\end{pmatrix}^n$.（提示：后者写成 $2\vv I+\vv N$.）
			\item\label{ex:mat:B12} 证明：$n$ 阶严格上三角矩阵（对角线及以下全为零）$\vv N$ 满足 $\vv N^n=\vv O$. 更精确地，$\vv N^k$ 的非零元只可能出现在第 $k$ 条``上对角线''及其右上方.
			\item\label{ex:mat:B13} 设 $\vv A$ 为实矩阵. 证明：若 $\vv A^\trans\vv A=\vv O$，则 $\vv A=\vv O$；若 $\tr(\vv A^\trans\vv A)=0$，结论同样成立. 复矩阵呢？
			\item\label{ex:mat:B14} 每行、每列恰有一个 $1$、其余元素为 $0$ 的方阵称为\textbf{置换矩阵}. 证明：置换矩阵之积仍是置换矩阵，置换矩阵 $\vv P$ 满足 $\vv P^\trans\vv P=\vv I$；$\vv{PA}$ 与 $\vv{AP}$ 分别对 $\vv A$ 做了什么？
			\item\label{ex:mat:B15} 对 $\vv A=\begin{pmatrix}1&2&1\\2&4&3\\1&2&2\end{pmatrix}$，求可逆矩阵 $\vv P,\vv Q$ 使 $\vv{PAQ}$ 为相抵标准形.
			\item\label{ex:mat:B16} 设 $\vv A$ 是 $n$ 阶方阵，$\vv A^k=\vv O$. 证明 $\vv I+\vv A$ 可逆，并求其逆. 若 $\vv A^2=2\vv A$，$\vv I-\vv A$ 是否可逆？
		\end{exB}
		
		\Cgroup
		\begin{exC}
			\item\label{ex:mat:C1} 设 $f:M_n(\mathbb{F})\to\mathbb{F}$ 满足 $f(\vv A+\vv B)=f(\vv A)+f(\vv B)$，$f(k\vv A)=kf(\vv A)$，且 $f(\vv{AB})=f(\vv{BA})$ 对一切 $\vv A,\vv B$ 成立. 证明 $f=c\cdot\tr$ 对某个常数 $c$ 成立.（提示：计算 $f(\vv E_{ij})$.）
			\item\label{ex:mat:C2} 设 $n$ 阶方阵 $\vv A,\vv B$ 满足 $\vv A+\vv B=\vv{AB}$. 证明 $\vv{AB}=\vv{BA}$.（提示：考察 $(\vv I-\vv A)(\vv I-\vv B)$.）
			\item\label{ex:mat:C3} （Strassen 算法）对二阶分块矩阵的乘法，直接计算需要 $8$ 次块乘法. 验证下面的公式只需 $7$ 次：令
			$M_1=(A_{11}+A_{22})(B_{11}+B_{22})$，$M_2=(A_{21}+A_{22})B_{11}$，$M_3=A_{11}(B_{12}-B_{22})$，$M_4=A_{22}(B_{21}-B_{11})$，$M_5=(A_{11}+A_{12})B_{22}$，$M_6=(A_{21}-A_{11})(B_{11}+B_{12})$，$M_7=(A_{12}-A_{22})(B_{21}+B_{22})$，则
			$C_{11}=M_1+M_4-M_5+M_7$，$C_{12}=M_3+M_5$，$C_{21}=M_2+M_4$，$C_{22}=M_1-M_2+M_3+M_6$. 由此说明 $n=2^k$ 阶矩阵乘法可以用约 $n^{\log_27}\approx n^{2.81}$ 次数的乘法完成.
			\item\label{ex:mat:C4} 设 $\vv A$ 是幂等矩阵. 用满秩分解 $\vv A=\vv{BC}$ 证明 $\vv{CB}=\vv I_r$，进而证明 $\rk\vv A=\tr\vv A$.
			\item\label{ex:mat:C5} （Kronecker 积）对 $m\times n$ 矩阵 $\vv A$ 与 $p\times q$ 矩阵 $\vv B$，定义 $\vv A\otimes\vv B=(a_{ij}\vv B)$（$mp\times nq$ 分块矩阵）. 证明 $(\vv A\otimes\vv B)(\vv C\otimes\vv D)=(\vv{AC})\otimes(\vv{BD})$，并由此说明 $\vv A,\vv B$ 可逆时 $(\vv A\otimes\vv B)^{-1}=\vv A^{-1}\otimes\vv B^{-1}$.
			\item\label{ex:mat:C6} （Hadamard 矩阵）元素全为 $\pm1$ 且满足 $\vv H\vv H^\trans=n\vv I$ 的 $n$ 阶矩阵称为 Hadamard 矩阵. 验证 $\vv H_2=\begin{pmatrix}1&1\\1&-1\end{pmatrix}$ 是，并证明 $\vv H_2\otimes\vv H_2$，$\vv H_2\otimes\vv H_2\otimes\vv H_2$ 也是. 证明：$n>2$ 时 Hadamard 矩阵的阶数必是 $4$ 的倍数.（这类矩阵用于实验设计与纠错编码.）
			\item\label{ex:mat:C7} 证明：全体形如 $\begin{pmatrix}a&-b\\b&a\end{pmatrix}$（$a,b\in\mathbb{R}$）的矩阵关于矩阵的加法和乘法``像复数一样''运算：加法、乘法封闭，乘法可交换，非零元可逆. 写出它与 $\mathbb{C}$ 之间的对应，并用它验证 $(1+\mathrm i)^8=16$.
			\item\label{ex:mat:C8} 设 $\vv A$ 是元素非负、各列之和都小于 $1$ 的方阵. 证明 $\vv A^k$ 的每个元素都随 $k\to\infty$ 趋于 $0$，并且级数 $\vv I+\vv A+\vv A^2+\cdots$（按元素）收敛到 $(\vv I-\vv A)^{-1}$. 解释它在投入产出模型中的意义：``直接消耗 $+$ 一次间接消耗 $+$ 二次间接消耗 $+\cdots$''.
			\item\label{ex:mat:C9} 设 $\vv A$ 为 $m\times n$ 矩阵. 证明：$\vv A$ 有右逆（存在 $\vv C$ 使 $\vv{AC}=\vv I_m$）当且仅当 $\rk\vv A=m$；有左逆当且仅当 $\rk\vv A=n$. 非方阵的单侧逆唯一吗？
		\end{exC}
	}
	
	% =====================================================================
	% 第四章　线性映射
	% =====================================================================
	\newcommand{\ExLinearMap}{%
		\exsetheading
		\Agroup
		\begin{exA}
			\item\label{ex:lmap:A1} 判断下列映射是否线性，是线性映射的求出核：
			\begin{enumerate}
				\item $\mathbb{R}^2\to\mathbb{R}^2$，$(x,y)^\trans\mapsto(x+y,\,2x)^\trans$；
				\item $\mathbb{R}^2\to\mathbb{R}^2$，$(x,y)^\trans\mapsto(x+1,\,y)^\trans$；
				\item $\mathbb{R}^2\to\mathbb{R}$，$(x,y)^\trans\mapsto|x|+y$；
				\item $\mathbb{R}[x]\to\mathbb{R}[x]$，$f(x)\mapsto f(x+1)$；
				\item $\mathbb{R}[x]\to\mathbb{R}[x]$，$f(x)\mapsto xf'(x)$；
				\item $M_n(\mathbb{R})\to M_n(\mathbb{R})$，$\vv X\mapsto\vv X^\trans$；
				\item $M_2(\mathbb{R})\to\mathbb{R}$，$\vv X\mapsto x_{11}x_{22}-x_{12}x_{21}$.
			\end{enumerate}
			\item\label{ex:lmap:A2} 求下列矩阵所定义的线性映射 $\vv x\mapsto\vv{Ax}$ 的核与像（各给出一组基），并验证秩—零化度定理：
			\[\begin{pmatrix}1&2&-1\\2&4&-2\end{pmatrix},\qquad\begin{pmatrix}1&0&2&1\\0&1&-1&1\\1&1&1&2\end{pmatrix}.\]
			\item\label{ex:lmap:A3} 写出线性映射 $\mathscr T:\mathbb{R}^3\to\mathbb{R}^2$，$\mathscr T(x_1,x_2,x_3)^\trans=(x_1-2x_2+x_3,\ 3x_2+4x_3)^\trans$ 在标准基下的矩阵.
			\item\label{ex:lmap:A4} 求下列平面变换在标准基下的矩阵：(1) 逆时针旋转 $90^\circ$；(2) 关于直线 $y=x$ 的反射；(3) 向直线 $y=x$ 的正投影；(4) 沿 $x$ 方向的剪切 $(x,y)\mapsto(x+2y,y)$. 计算 (2)(3) 的矩阵的平方，并解释结果.
			\item\label{ex:lmap:A5} 求求导变换 $\mathscr D:\mathbb{R}[x]_3\to\mathbb{R}[x]_3$ 在基 $1,x,x^2$ 下的矩阵，以及在基 $1,\ x-1,\ (x-1)^2$ 下的矩阵.
			\item\label{ex:lmap:A6} 求转置变换 $\vv X\mapsto\vv X^\trans$ 在 $M_2(\mathbb{F})$ 的基 $\vv E_{11},\vv E_{12},\vv E_{21},\vv E_{22}$ 下的矩阵. 它的平方是什么？
			\item\label{ex:lmap:A7} 在 $\mathbb{R}^3$ 中，求从基 $\vv\varepsilon_1=(1,0,0)^\trans$，$\vv\varepsilon_2=(1,1,0)^\trans$，$\vv\varepsilon_3=(1,1,1)^\trans$ 到基 $\vv\eta_1=(1,2,1)^\trans$，$\vv\eta_2=(2,3,3)^\trans$，$\vv\eta_3=(3,7,1)^\trans$ 的过渡矩阵，并求 $\vv\alpha=(3,4,2)^\trans$ 在两组基下的坐标.
			\item\label{ex:lmap:A8} 设线性变换 $\mathscr T$ 在基 $\vv\varepsilon_1,\vv\varepsilon_2$ 下的矩阵为 $\begin{pmatrix}2&1\\1&2\end{pmatrix}$. 求它在基 $\vv\eta_1=\vv\varepsilon_1+\vv\varepsilon_2$，$\vv\eta_2=\vv\varepsilon_1-\vv\varepsilon_2$ 下的矩阵.
			\item\label{ex:lmap:A9} 证明：线性映射的复合仍是线性映射；可逆（双射）线性映射的逆映射仍是线性映射.
			\item\label{ex:lmap:A10} 证明：线性映射 $\mathscr T$ 是单射，当且仅当它把每个线性无关组都映为线性无关组.
			\item\label{ex:lmap:A11} 是否存在线性变换 $\mathscr T:\mathbb{R}^2\to\mathbb{R}^2$，使
			(1) $\mathscr T(1,1)^\trans=(1,0)^\trans$，$\mathscr T(1,-1)^\trans=(0,1)^\trans$，$\mathscr T(2,0)^\trans=(1,1)^\trans$？
			(2) 同上，但最后一个条件改为 $\mathscr T(2,0)^\trans=(1,2)^\trans$？
			存在时求出 $\mathscr T(x,y)^\trans$.
			\item\label{ex:lmap:A12} (1) 线性映射 $\mathscr T:\mathbb{R}^5\to\mathbb{R}^3$ 是满射，$\dim\ker\mathscr T$ 是多少？(2) 是否存在从 $\mathbb{R}^3$ 到 $\mathbb{R}^5$ 的满射线性映射？是否存在从 $\mathbb{R}^5$ 到 $\mathbb{R}^3$ 的单射线性映射？
			\item\label{ex:lmap:A13} 设 $\mathscr T:U\to V$，$\mathscr S:V\to W$ 线性. 证明 $\ker\mathscr T\subset\ker(\mathscr{ST})$，$\imm(\mathscr{ST})\subset\imm\mathscr S$，从而 $\rk(\mathscr{ST})\le\min\{\rk\mathscr S,\rk\mathscr T\}$.
			\item\label{ex:lmap:A14} 求迹映射 $\tr:M_n(\mathbb{F})\to\mathbb{F}$ 的核的维数.
			\item\label{ex:lmap:A15} 证明：相似矩阵有相同的迹（利用命题~\ref{prop:trace}(3)）. 由此判断 $\begin{pmatrix}1&2\\0&3\end{pmatrix}$ 与 $\begin{pmatrix}1&0\\0&2\end{pmatrix}$ 是否相似.
			\item\label{ex:lmap:A16} 证明：与单位矩阵相似的矩阵只有单位矩阵本身. 由此说明 $\begin{pmatrix}1&1\\0&1\end{pmatrix}$ 与 $\vv I$ 不相似.
			\item\label{ex:lmap:A17} 定义求值映射 $\mathscr E:\mathbb{R}[x]_3\to\mathbb{R}^3$，$f\mapsto(f(0),f(1),f(2))^\trans$. 求它在基 $1,x,x^2$ 与标准基下的矩阵，并证明 $\mathscr E$ 是同构.
			\item\label{ex:lmap:A18} 设 $\mathscr T$ 是 $\mathbb{R}^3$ 上的线性变换，$\mathscr T\vv e_1=(1,1,0)^\trans$，$\mathscr T\vv e_2=(0,1,1)^\trans$，$\mathscr T\vv e_3=(1,2,1)^\trans$. 求 $\mathscr T(x_1,x_2,x_3)^\trans$，$\ker\mathscr T$ 与 $\imm\mathscr T$.
			\item\label{ex:lmap:A19} 称 $\vv x\mapsto\vv{Ax}+\vv b$ 为\textbf{仿射变换}. 证明两个仿射变换的复合仍是仿射变换；仿射变换是线性变换当且仅当 $\vv b=\zero$.
			\item\label{ex:lmap:A20} 设 $\vv R(\theta)=\begin{pmatrix}\cos\theta&-\sin\theta\\ \sin\theta&\cos\theta\end{pmatrix}$ 是平面上绕原点旋转角 $\theta$ 的变换在标准基下的矩阵. 从几何上说明``先转 $\beta$ 再转 $\alpha$''就是转 $\alpha+\beta$，从而 $\vv R(\alpha)\vv R(\beta)=\vv R(\alpha+\beta)$；比较两边矩阵的元素，推出两角和的正弦、余弦公式.
		\end{exA}
		
		\Bgroup
		\begin{exB}
			\item\label{ex:lmap:B1} 记 $L(V,W)$ 为 $V$ 到 $W$ 的全体线性映射. 证明它关于映射的加法与数乘构成向量空间，且 $\dim L(V,W)=\dim V\cdot\dim W$.（提示：取定基后，线性映射与矩阵一一对应.）
			\item\label{ex:lmap:B2} 设 $\mathscr T$ 是 $n$ 维空间 $V$ 上的线性变换，$\mathscr T^2=\mathscr O$. 证明 $\imm\mathscr T\subset\ker\mathscr T$，从而 $\rk\mathscr T\le n/2$；并证明存在一组基，使 $\mathscr T$ 的矩阵只在 $(1,2),(3,4),\ldots,(2r-1,2r)$ 这 $r$ 个位置上为 $1$，其余元素全为 $0$，其中 $r=\rk\mathscr T$.
			\item\label{ex:lmap:B3} 设 $\mathscr T,\mathscr S$ 是有限维空间上的线性变换. 证明 $\dim\ker(\mathscr{ST})\le\dim\ker\mathscr S+\dim\ker\mathscr T$.
			\item\label{ex:lmap:B4} 设 $V$ 有限维，$\mathscr S,\mathscr T$ 是 $V$ 上的线性变换且 $\mathscr{ST}=\mathscr I$. 证明 $\mathscr{TS}=\mathscr I$.
			\item\label{ex:lmap:B5} 在无限维空间 $\mathbb{R}[x]$ 上，令 $\mathscr D$ 为求导，$\mathscr J$ 为 $f\mapsto\int_0^xf(t)\,\mathrm dt$. 证明 $\mathscr{DJ}=\mathscr I$，但 $\mathscr{JD}\ne\mathscr I$；$\mathscr D$ 是满射但不是单射，$\mathscr J$ 是单射但不是满射. 这与 B4 是否矛盾？
			\item\label{ex:lmap:B6} 固定 $\vv A\in M_n(\mathbb{F})$，令 $\mathscr T(\vv X)=\vv{AX}-\vv{XA}$. 证明 $\imm\mathscr T$ 含于迹为零的矩阵全体，从而 $\mathscr T$ 不是满射；并证明 $\ker\mathscr T$ 包含 $\vv I,\vv A,\vv A^2,\ldots$.
			\item\label{ex:lmap:B7} 设 $\vv A=\begin{pmatrix}a&b\\c&d\end{pmatrix}$. 求线性变换 $\vv X\mapsto\vv{AX}$ 在 $M_2(\mathbb{F})$ 的基 $\vv E_{11},\vv E_{21},\vv E_{12},\vv E_{22}$ 下的矩阵. 若换成 $\vv X\mapsto\vv{XA}$ 呢？
			\item\label{ex:lmap:B8} 设线性变换 $\mathscr P$ 满足 $\mathscr P^2=\mathscr P$. 证明存在一组基使 $\mathscr P$ 的矩阵为 $\begin{pmatrix}\vv I_r&\vv O\\ \vv O&\vv O\end{pmatrix}$，其中 $r=\rk\mathscr P$.
			\item\label{ex:lmap:B9} 设线性变换 $\mathscr T$ 满足 $\mathscr T^2=\mathscr I$. 证明 $V=\ker(\mathscr T-\mathscr I)\oplus\ker(\mathscr T+\mathscr I)$，从而存在一组基使 $\mathscr T$ 的矩阵为 $\diag(\vv I_p,-\vv I_q)$.（提示：$\vv\alpha=\frac12(\vv\alpha+\mathscr T\vv\alpha)+\frac12(\vv\alpha-\mathscr T\vv\alpha)$.）
			\item\label{ex:lmap:B10} 设 $\mathscr T:V\to W$ 线性，$U$ 是 $V$ 的子空间. 证明 $\dim\mathscr T(U)=\dim U-\dim(U\cap\ker\mathscr T)$.
			\item\label{ex:lmap:B11} 在由 $\mathrm e^x\cos x$，$\mathrm e^x\sin x$ 张成的二维函数空间中，求求导变换的矩阵，证明它可逆，并利用逆矩阵求出 $\int\mathrm e^x\sin x\,\mathrm dx$（不必用分部积分）.
			\item\label{ex:lmap:B12} （齐次坐标）平面上的平移 $(x,y)\mapsto(x+a,y+b)$ 不是线性的. 把点 $(x,y)$ 写成三维向量 $(x,y,1)^\trans$，证明平移、旋转、伸缩都可以写成 $3\times3$ 矩阵的乘法，且仿射变换的复合对应矩阵的乘积. 写出``先绕原点旋转 $90^\circ$，再平移 $(1,2)$''与``先平移再旋转''的矩阵，它们相等吗？（计算机图形学正是这样处理变换的.）
			\item\label{ex:lmap:B13} 证明：若 $\vv A$ 与 $\vv B$ 相似，则 $\vv A^k$ 与 $\vv B^k$ 相似，$\vv A^\trans$ 与 $\vv B^\trans$ 相似，$f(\vv A)$ 与 $f(\vv B)$ 相似（$f$ 为多项式）.
			\item\label{ex:lmap:B14} 证明 $\begin{pmatrix}1&0\\0&0\end{pmatrix}$ 与 $\begin{pmatrix}0&1\\0&0\end{pmatrix}$ 相抵但不相似. 从``线性映射 $V\to W$''与``线性变换 $V\to V$''的区别解释：为什么相抵只需看秩，而相似要复杂得多？
			\item\label{ex:lmap:B15} 设 $\mathscr T$ 是 $V$ 上的线性变换，$\vv\alpha\in V$ 满足 $\mathscr T^{k-1}\vv\alpha\ne\zero$，$\mathscr T^k\vv\alpha=\zero$. 证明 $\vv\alpha,\mathscr T\vv\alpha,\ldots,\mathscr T^{k-1}\vv\alpha$ 线性无关. 若 $k=n=\dim V$，求 $\mathscr T$ 在基 $\mathscr T^{n-1}\vv\alpha,\ldots,\mathscr T\vv\alpha,\vv\alpha$ 下的矩阵.
		\end{exB}
		
		\Cgroup
		\begin{exC}
			\item\label{ex:lmap:C1} （Fitting 分解）设 $\mathscr T$ 是 $n$ 维空间 $V$ 上的线性变换. 证明：
			\begin{enumerate}
				\item $\ker\mathscr T\subset\ker\mathscr T^2\subset\cdots$，$\imm\mathscr T\supset\imm\mathscr T^2\supset\cdots$，且两条链都在某一步之后不再变化；
				\item $V=\ker\mathscr T^n\oplus\imm\mathscr T^n$，两者都是 $\mathscr T$ 的不变子空间，$\mathscr T$ 在前者上幂零，在后者上可逆.
			\end{enumerate}
			\item\label{ex:lmap:C2} （商空间）设 $W$ 是 $V$ 的子空间. 对 $\vv\alpha\in V$，称集合 $\vv\alpha+W$ 为 $\vv\alpha$ 所在的陪集，全体陪集记为 $V/W$. 定义 $(\vv\alpha+W)+(\vv\beta+W)=(\vv\alpha+\vv\beta)+W$，$k(\vv\alpha+W)=k\vv\alpha+W$.
			\begin{enumerate}
				\item 证明这两个运算是良定义的，且 $V/W$ 构成向量空间；
				\item 证明 $\dim V/W=\dim V-\dim W$；
				\item （第一同构定理/同态基本定理）对线性映射 $\mathscr T:V\to U$，证明 $V/\ker\mathscr T\cong\imm\mathscr T$. 这给出秩—零化度定理的另一种理解.
			\end{enumerate}
			\item\label{ex:lmap:C3} 把 $\mathbb{C}$ 看作 $\mathbb{R}$ 上的二维空间，基为 $1,\mathrm i$. (1) 求``乘以 $a+b\mathrm i$''这个变换的矩阵；(2) 求``取共轭''这个变换的矩阵；(3) 证明 $\mathbb{R}$-线性变换 $\begin{pmatrix}p&q\\r&s\end{pmatrix}$ 是``乘以某个复数''，当且仅当 $p=s$，$q=-r$.
			\item\label{ex:lmap:C4} 证明：$n$ 维空间上的任一线性变换 $\mathscr T$，都存在次数不超过 $n^2$ 的非零多项式 $f$ 使 $f(\mathscr T)=\mathscr O$.（提示：$\mathscr I,\mathscr T,\ldots,\mathscr T^{n^2}$ 在 $L(V,V)$ 中线性相关. 第八章的 Cayley–Hamilton 定理会把 $n^2$ 改进为 $n$.）
			\item\label{ex:lmap:C5} （Haar 小波）在 $\mathbb{R}^4$ 中取基 $\vv h_1=(1,1,1,1)^\trans$，$\vv h_2=(1,1,-1,-1)^\trans$，$\vv h_3=(1,-1,0,0)^\trans$，$\vv h_4=(0,0,1,-1)^\trans$. 求信号 $\vv x=(8,6,3,4)^\trans$ 在这组基下的坐标，并解释每个坐标的含义. 若把绝对值最小的坐标改为 $0$ 再变换回来，误差是多少？这是图像与音频压缩的思想之一.
			\item\label{ex:lmap:C6} 设 $\mathscr T$ 是 $V$ 上的线性变换. 若 $\mathscr T$ 与 $V$ 上的一切线性变换都可交换，证明 $\mathscr T=c\mathscr I$. 若只要求与一切\textbf{可逆}线性变换可交换呢？
			\item\label{ex:lmap:C7} 设 $\mathscr T:V\to V$ 线性，$\dim V=n$. 证明：$\rk\mathscr T^k-\rk\mathscr T^{k+1}$ 关于 $k$ 单调不增. 由此说明：若 $\rk\mathscr T=\rk\mathscr T^2$，则对一切 $k$，$\rk\mathscr T^k=\rk\mathscr T$.
			\item\label{ex:lmap:C8} （差分算子）在全体实数列构成的空间上，定义左移 $\mathscr S(a_1,a_2,a_3,\ldots)=(a_2,a_3,\ldots)$. 证明 $\mathscr S$ 是线性的、满射但不是单射；满足 $a_{n+2}=a_{n+1}+a_n$ 的数列全体恰好是 $\ker(\mathscr S^2-\mathscr S-\mathscr I)$. 由第二章章末习题~\ref{ex:vs:C9}，这个核是二维的.
		\end{exC}
	}
	
	% =====================================================================
	% 第五章　行列式
	% =====================================================================
	\newcommand{\ExDeterminant}{%
		\exsetheading
		\Agroup
		\begin{exA}
			\item\label{ex:det:A1} 求下列排列的逆序数并判断奇偶：$315426$，$42531$，$13\cdots(2n-1)\,24\cdots(2n)$.
			\item\label{ex:det:A2} (1) 在 $4$ 阶行列式中，项 $a_{13}a_{24}a_{32}a_{41}$ 带什么符号？(2) 选择 $i,j$，使 $a_{1i}a_{23}a_{35}a_{5j}a_{44}$ 成为 $5$ 阶行列式中带正号的一项.
			\item\label{ex:det:A3} 计算：
			\[\begin{vmatrix}2&1&-1\\3&0&2\\1&4&5\end{vmatrix},\qquad
			\begin{vmatrix}1&2&0&1\\2&1&1&0\\0&1&2&1\\1&0&1&2\end{vmatrix},\qquad
			\begin{vmatrix}3&1&1&1\\1&3&1&1\\1&1&3&1\\1&1&1&3\end{vmatrix}.\]
			\item\label{ex:det:A4} 计算 $\begin{vmatrix}1&-1&2&3\\2&0&1&-1\\0&1&-1&2\\3&1&0&1\end{vmatrix}$，分别用``打洞''化三角形与按第一列展开两种方法，比较哪种更快.
			\item\label{ex:det:A5} 计算 $n$ 阶行列式 $\begin{vmatrix}0&\cdots&0&a_1\\0&\cdots&a_2&0\\ \vdots&&&\vdots\\a_n&\cdots&0&0\end{vmatrix}$（只有副对角线上有元素）.
			\item\label{ex:det:A6} 不展开，证明：
			\[\begin{vmatrix}a+b&b+c&c+a\\a_1+b_1&b_1+c_1&c_1+a_1\\a_2+b_2&b_2+c_2&c_2+a_2\end{vmatrix}=2\begin{vmatrix}a&b&c\\a_1&b_1&c_1\\a_2&b_2&c_2\end{vmatrix}.\]
			\item\label{ex:det:A7} 解方程 $\begin{vmatrix}x&1&1\\1&x&1\\1&1&x\end{vmatrix}=0$ 与 $\begin{vmatrix}1&1&1\\1&2&x\\1&4&x^2\end{vmatrix}=0$.
			\item\label{ex:det:A8} 设 $\vv A$ 为 $3$ 阶方阵，$\det\vv A=3$. 求 $\det(2\vv A)$，$\det(-\vv A)$，$\det(\vv A^{-1})$，$\det(\vv A^*)$，$\det(\vv A^\trans\vv A)$，$\det(3\vv A^{-1}-2\vv A^*)$.
			\item\label{ex:det:A9} 证明：(1) 若 $n$ 阶方阵满足 $\vv A^\trans\vv A=\vv I$，则 $\det\vv A=\pm1$；(2) 若 $\vv A^2=\vv I$，则 $\det\vv A=\pm1$；(3) 若 $\vv A^2=\vv A$ 且 $\vv A\ne\vv I$，则 $\det\vv A=0$.
			\item\label{ex:det:A10} 求下列矩阵的伴随矩阵与逆矩阵：\[\begin{pmatrix}a&b\\c&d\end{pmatrix},\qquad\begin{pmatrix}1&2&3\\0&1&4\\5&6&0\end{pmatrix}.\]
			\item\label{ex:det:A11} 用 Cramer 法则解方程组 $2x_1+x_2+x_3=3$，$x_1+3x_2+x_3=0$，$x_1+x_2+4x_3=8$.
			\item\label{ex:det:A12} 求以 $(1,1)$，$(4,2)$，$(2,5)$ 为顶点的三角形的面积；判断 $(1,2)$，$(3,5)$，$(7,11)$ 三点是否共线.
			\item\label{ex:det:A13} 求以 $(1,0,1)^\trans$，$(2,1,0)^\trans$，$(0,3,1)^\trans$ 为棱的平行六面体的体积，以及以原点和这三个向量的终点为顶点的四面体的体积.
			\item\label{ex:det:A14} $\lambda$ 取何值时，齐次方程组 $(5-\lambda)x+2y+2z=0$，$2x+(6-\lambda)y=0$，$2x+(4-\lambda)z=0$ 有非零解？
			\item\label{ex:det:A15} 设 \[\vv A=\begin{pmatrix}1&2&3&4\\2&3&4&1\\3&4&1&2\\4&1&2&3\end{pmatrix}.\] 求 $A_{11}+A_{12}+A_{13}+A_{14}$（$A_{1j}$ 是代数余子式）.（提示：它是把第一行换成全 $1$ 后的行列式.）
			\item\label{ex:det:A16} 证明：过平面上两个不同点 $(x_1,y_1)$，$(x_2,y_2)$ 的直线方程可以写成 \[\begin{vmatrix}x&y&1\\x_1&y_1&1\\x_2&y_2&1\end{vmatrix}=0.\]
			\item\label{ex:det:A17} 设 $\vv A,\vv B$ 分别是 $m$ 阶、$n$ 阶方阵，证明 $\begin{vmatrix}\vv A&\vv O\\ \vv C&\vv B\end{vmatrix}=\det\vv A\det\vv B$，并计算 $\begin{vmatrix}\vv O&\vv A\\ \vv B&\vv O\end{vmatrix}$.
			\item\label{ex:det:A18} 计算 $4$ 阶三对角行列式 $\begin{vmatrix}2&-1&0&0\\-1&2&-1&0\\0&-1&2&-1\\0&0&-1&2\end{vmatrix}$.
			\item\label{ex:det:A19} 设 $\vv A$ 是 $n$ 阶可逆矩阵，$n\ge2$. 证明 $\vv A^*$ 可逆，且 $(\vv A^*)^{-1}=\frac1{\det\vv A}\vv A=(\vv A^{-1})^*$.
			\item\label{ex:det:A20} 用行列式判断：向量 $(1,2,3)^\trans$，$(4,5,6)^\trans$，$(7,8,9)^\trans$ 是否线性无关？把 $9$ 改成 $10$ 呢？
			\item\label{ex:det:A21} 设 $\vv A$ 是 $n$ 阶反对称矩阵（$\vv A^\trans=-\vv A$），$n$ 为奇数. 证明 $\det\vv A=0$.（提示：对 $\vv A^\trans=-\vv A$ 两边取行列式.）
		\end{exA}
		
		\Bgroup
		\begin{exB}
			\item\label{ex:det:B1} 计算 $n$ 阶三对角行列式 $D_n=\begin{vmatrix}2&-1&&\\-1&2&\ddots&\\&\ddots&\ddots&-1\\&&-1&2\end{vmatrix}$.（提示：按第一行展开得 $D_n=2D_{n-1}-D_{n-2}$.）
			\item\label{ex:det:B2} 计算 $n$ 阶行列式 $\begin{vmatrix}a&b&\cdots&b\\b&a&\cdots&b\\ \vdots&&\ddots&\vdots\\b&b&\cdots&a\end{vmatrix}$ 与 $\begin{vmatrix}x&a&\cdots&a\\-a&x&\cdots&a\\ \vdots&&\ddots&\vdots\\-a&-a&\cdots&x\end{vmatrix}$.
			\item\label{ex:det:B3} （爪形行列式）设 $a_1\cdots a_n\ne0$，计算 $\begin{vmatrix}a_0&1&1&\cdots&1\\1&a_1&0&\cdots&0\\1&0&a_2&\cdots&0\\ \vdots&&&\ddots&\\1&0&0&\cdots&a_n\end{vmatrix}$.
			\item\label{ex:det:B4} （矩阵行列式引理）证明 $\det(\vv I_n+\vv u\vv v^\trans)=1+\vv v^\trans\vv u$，更一般地，$\vv A$ 可逆时 $\det(\vv A+\vv u\vv v^\trans)=(1+\vv v^\trans\vv A^{-1}\vv u)\det\vv A$. 用它重做 B2 的第一个行列式.
			\item\label{ex:det:B5} 设 $\vv A$ 为 $2\times3$ 矩阵，$\vv B$ 为 $3\times2$ 矩阵. 证明 $\det(\vv{AB})$ 等于 $\vv A$ 的三个 $2$ 阶子式与 $\vv B$ 的对应 $2$ 阶子式乘积之和（这是 Cauchy–Binet 公式最简单的情形）. 由此说明 $\vv{BA}$（$3$ 阶方阵）的行列式为零.
			\item\label{ex:det:B6} 设 $\vv A$ 是整数矩阵. 证明：$\vv A^{-1}$ 存在且也是整数矩阵，当且仅当 $\det\vv A=\pm1$.
			\item\label{ex:det:B7} 证明：$n\ge2$ 时 $(\vv A^*)^*=(\det\vv A)^{n-2}\vv A$.（先设 $\vv A$ 可逆；不可逆的情形另行讨论.）
			\item\label{ex:det:B8} 设 $\vv A,\vv B$ 都是 $n$ 阶方阵. 证明 $\begin{vmatrix}\vv A&\vv B\\ \vv B&\vv A\end{vmatrix}=\det(\vv A+\vv B)\det(\vv A-\vv B)$.（提示：分块初等变换.）
			\item\label{ex:det:B9} 设 $\vv A$ 可逆. 证明 $\begin{vmatrix}\vv A&\vv B\\ \vv C&\vv D\end{vmatrix}=\det\vv A\cdot\det(\vv D-\vv{CA}^{-1}\vv B)$；若还有 $\vv{AC}=\vv{CA}$，则它等于 $\det(\vv{AD}-\vv{CB})$.
			\item\label{ex:det:B10} 用 Fibonacci 矩阵 $\vv F=\begin{pmatrix}1&1\\1&0\end{pmatrix}$ 满足 $\vv F^n=\begin{pmatrix}F_{n+1}&F_n\\F_n&F_{n-1}\end{pmatrix}$（归纳法证明），推出 Cassini 恒等式 $F_{n+1}F_{n-1}-F_n^2=(-1)^n$.
			\item\label{ex:det:B11} 计算 $n$ 阶行列式 $D_n=\begin{vmatrix}1+x^2&x&&\\x&1+x^2&\ddots&\\&\ddots&\ddots&x\\&&x&1+x^2\end{vmatrix}$.
			\item\label{ex:det:B12} 设 $\omega=\mathrm e^{2\pi\mathrm i/3}$. 证明
			\[\begin{vmatrix}a&b&c\\c&a&b\\b&c&a\end{vmatrix}=(a+b+c)(a+\omega b+\omega^2c)(a+\omega^2b+\omega c)=a^3+b^3+c^3-3abc.\]
			\item\label{ex:det:B13} 证明：元素为 $a_{ij}=\min(i,j)$ 的 $n$ 阶行列式等于 $1$.
			\item\label{ex:det:B14} 证明：平面上四点 $(x_i,y_i)$（$i=1,2,3,4$）共圆（或共线），当且仅当 $\begin{vmatrix}x_i^2+y_i^2&x_i&y_i&1\end{vmatrix}_{4\times4}=0$（每行对应一个点）.
						\item\label{ex:det:B15} 设 $\vv A$ 为 $n$ 阶方阵（$n\ge2$）. 证明
			\[\rk\vv A^*=\begin{cases}n,&\rk\vv A=n,\\1,&\rk\vv A=n-1,\\0,&\rk\vv A\le n-2.\end{cases}\]
			（提示：$\vv A^*$ 的元素都是 $\vv A$ 的 $n-1$ 阶子式，用定理~\ref{thm:minorrank}；$\rk\vv A=n-1$ 时考察 $\vv A\vv A^*=(\det\vv A)\vv I$.）
			\item\label{ex:det:B16} 设 $f(x)=\begin{vmatrix}x&1&2&3\\1&x&2&3\\1&2&x&3\\1&2&3&x\end{vmatrix}$. 不展开，求 $f(x)=0$ 的全部根.
		\end{exB}
		
		\Cgroup
		\begin{exC}
			\item\label{ex:det:C1} （行列式的唯一性）设 $D:M_n(\mathbb{F})\to\mathbb{F}$ 对每一行都是线性的、交换两行变号，且 $D(\vv I)=1$. 证明 $D=\det$.（提示：先说明 $D$ 在初等矩阵上的取值，再对任意矩阵用初等变换.）
			\item\label{ex:det:C2} （Cauchy 行列式）证明
			\[\det\Big(\frac1{a_i+b_j}\Big)_{n\times n}=\frac{\prod_{i<j}(a_j-a_i)(b_j-b_i)}{\prod_{i,j}(a_i+b_j)}.\]
			先对 $n=2$ 直接验证，再用``各行减去最后一行''归纳. 取 $a_i=i$，$b_j=j-1$，得到 Hilbert 矩阵 $(1/(i+j-1))$ 的行列式，它为什么在数值计算中臭名昭著？
			\item\label{ex:det:C3} （排列的符号是同态）对排列 $\sigma$，令 $\vv P_\sigma$ 为把 $\vv e_j$ 映到 $\vv e_{\sigma(j)}$ 的置换矩阵. 证明 $\vv P_\sigma\vv P_\tau=\vv P_{\sigma\tau}$，$\det\vv P_\sigma=(-1)^{\tau(\sigma)}$，并由行列式的乘积公式推出 $\operatorname{sgn}(\sigma\tau)=\operatorname{sgn}\sigma\cdot\operatorname{sgn}\tau$.
			\item\label{ex:det:C4} （矩阵—树定理）对一个连通无向图，令 $\vv D$ 为各顶点度数构成的对角矩阵，$\vv A$ 为邻接矩阵，$\vv L=\vv D-\vv A$ 称为 Laplace 矩阵. 定理说：图的生成树个数等于 $\vv L$ 的任一个 $(n-1)$ 阶主子式. 对 $4$ 个顶点的完全图验证它有 $16$ 棵生成树，对一个 $4$ 个顶点的环验证它有 $4$ 棵生成树.
			\item\label{ex:det:C5} 证明：若 $n$ 阶矩阵的元素全为 $\pm1$，则它的行列式能被 $2^{n-1}$ 整除. 求 $3$ 阶此类行列式的最大值.
			\item\label{ex:det:C6} （Jacobi 公式的雏形）证明 $\det(\vv I+t\vv B)=1+t\tr\vv B+(\text{$t^2$ 及更高次项})$. 由此说明：若 $\vv A(t)$ 的元素都是 $t$ 的可导函数，且 $\vv A(0)=\vv I$，则 $\frac{\mathrm d}{\mathrm dt}\det\vv A(t)\big|_{t=0}=\tr\vv A'(0)$.
			\item\label{ex:det:C7} （Sylvester 结式）设 $f=a_0x^2+a_1x+a_2$，$g=b_0x^2+b_1x+b_2$（$a_0b_0\ne0$）. 证明 $f,g$ 有公共根当且仅当
			\[\begin{vmatrix}a_0&a_1&a_2&0\\0&a_0&a_1&a_2\\b_0&b_1&b_2&0\\0&b_0&b_1&b_2\end{vmatrix}=0.\]
			（提示：有公共根当且仅当存在次数 $\le1$ 的不全为零的 $u,v$ 使 $uf+vg=0$，与第七章的 Bézout 等式对照.）
			\item\label{ex:det:C8} （Dodgson 凝聚法）对 $3$ 阶矩阵 $\vv A$，记 $M_{ij}$ 为余子式，$c=a_{22}$. 验证 $\det\vv A\cdot c=M_{11}M_{33}-M_{13}M_{31}$. 它允许只计算 $2$ 阶行列式就得到高阶行列式. 《爱丽丝漫游奇境》的作者刘易斯·卡罗尔（即 Dodgson）于 1866 年发表了这一方法.
			\item\label{ex:det:C9} 设 $\vv A$ 是 $n$ 阶实反对称矩阵. 证明 $\det(\vv I+\vv A)\ge1$.（提示：$\det(\vv I+\vv A)=\det(\vv I-\vv A)$，考察 $\det(\vv I+t\vv A)$ 作为 $t$ 的多项式只含偶次项，且 $\vv I+t\vv A$ 对一切实数 $t$ 都可逆.）
		\end{exC}
	}
	
	% =====================================================================
	% 第六章　线性方程组解的结构
	% =====================================================================
	\newcommand{\ExSolutions}{%
		\exsetheading
		\Agroup
		\begin{exA}
			\item\label{ex:sol:A1} 求下列齐次方程组的基础解系与通解：
			\[\text{(1)}\ \left\{\begin{array}{l}x_1+x_2-x_3-x_4=0,\\2x_1-5x_2+3x_3+2x_4=0,\\7x_1-7x_2+3x_3+x_4=0;\end{array}\right.\qquad
			\text{(2)}\ \left\{\begin{array}{l}x_1+2x_2+x_3-x_4=0,\\3x_1+6x_2-x_3-3x_4=0,\\5x_1+10x_2+x_3-5x_4=0.\end{array}\right.\]
			\item\label{ex:sol:A2} 求下列方程组的通解，写成``特解 $+$ 基础解系的线性组合''：
			\[\text{(1)}\ \left\{\begin{array}{l}x_1+x_2-3x_3-x_4=1,\\3x_1-x_2-3x_3+4x_4=4,\\x_1+5x_2-9x_3-8x_4=0;\end{array}\right.\qquad
			\text{(2)}\ \left\{\begin{array}{l}2x_1+x_2-x_3+x_4=1,\\4x_1+2x_2-2x_3+x_4=2,\\2x_1+x_2-x_3-x_4=1.\end{array}\right.\]
			\item\label{ex:sol:A3} 设 $\vv A$ 是 $4\times5$ 矩阵，$\rk\vv A=3$. (1) $\vv{Ax}=\zero$ 的基础解系含几个向量？(2) $\vv{Ax}=\vv b$ 是否一定有解？(3) 若 $\rk\vv A=4$ 呢？
			\item\label{ex:sol:A4} 不解方程，用秩判断下列方程组是否有解：(1) $x_1+x_2=1$，$x_1-x_2=1$，$2x_1+3x_2=2$；(2) $x_1+x_2+x_3=1$，$2x_1+2x_2+2x_3=3$.
			\item\label{ex:sol:A5} 设 $\vv\eta_1,\vv\eta_2$ 是 $\vv{Ax}=\vv b$（$\vv b\ne\zero$）的两个解. 下列向量中哪些是 $\vv{Ax}=\vv b$ 的解，哪些是 $\vv{Ax}=\zero$ 的解？
			$\vv\eta_1+\vv\eta_2$，$2\vv\eta_1-\vv\eta_2$，$\frac12(\vv\eta_1+\vv\eta_2)$，$\vv\eta_1-\vv\eta_2$，$3\vv\eta_1-2\vv\eta_2$.
			\item\label{ex:sol:A6} 设 $3$ 元方程组 $\vv{Ax}=\vv b$ 的系数矩阵秩为 $2$，$\vv\eta_1=(1,2,3)^\trans$ 与 $\vv\eta_2=(1,-2,-1)^\trans$ 是它的两个解. 求通解.
			\item\label{ex:sol:A7} 写出一个以 $\vv x=(1,0,2)^\trans+k(1,1,-1)^\trans$（$k$ 任意）为全部解的线性方程组.
			\item\label{ex:sol:A8} 判断三个平面 $x+y+z=1$，$x+2y+3z=2$，$2x+3y+4z=4$ 的位置关系；若把最后一个常数改为 $3$ 呢？
			\item\label{ex:sol:A9} 举例说明：$\vv{Ax}=\zero$ 只有零解时，$\vv{Ax}=\vv b$ 可能无解；但若有解，解必唯一.
			\item\label{ex:sol:A10} 解矩阵方程 $\begin{pmatrix}1&2&3\\2&2&1\\3&4&3\end{pmatrix}\vv X=\begin{pmatrix}2&5\\3&1\\4&3\end{pmatrix}$.
			\item\label{ex:sol:A11} 设 $\vv A=\begin{pmatrix}1&1\\2&2\end{pmatrix}$. 求全部满足 $\vv{AX}=\vv O$ 的二阶矩阵 $\vv X$，它们构成的空间是几维的？
			\item\label{ex:sol:A12} 设 $\vv A=\begin{pmatrix}1&2&-2\\4&t&3\\3&-1&1\end{pmatrix}$，且存在三阶非零矩阵 $\vv B$ 使 $\vv{AB}=\vv O$. 求 $t$，并求 $\rk\vv B$ 的可能值.
			\item\label{ex:sol:A13} 求方程 $x_1+x_2+x_3=1$ 的通解，并写出它的三个解，使得通解中的每个解都是这三个解的仿射组合.
			\item\label{ex:sol:A14} 把空间直线 $\{x+y+z=1,\ x-y=0\}$ 写成参数式；反过来，把直线 $\vv x=(1,2,0)^\trans+t(1,0,1)^\trans$ 写成两个平面的交线.
			\item\label{ex:sol:A15} 判断空间直线 $\vv x=(1,0,0)^\trans+s(1,1,0)^\trans$ 与 $\vv x=(0,1,1)^\trans+t(0,1,-1)^\trans$ 是否相交.（列出关于 $s,t$ 的方程组.）
			\item\label{ex:sol:A16} 物流网络：三个仓库 $A,B,C$ 之间有线路 $A\to B$（$x_1$），$B\to C$（$x_2$），$A\to C$（$x_3$）. $A$ 每天净发出 $30$ 吨，$C$ 每天净收到 $30$ 吨，$B$ 不囤货. 求全部方案. 哪些量（如 $x_1+x_3$，$x_1-x_2$，$x_3$）能被方程组唯一确定？用可确定性判据说明.
			\item\label{ex:sol:A17} 求 $\vv A=\begin{pmatrix}2&1&0\\4&3&1\\2&3&3\end{pmatrix}$ 的 LU 分解.
			\item\label{ex:sol:A18} 设 $\vv A$ 是 $m\times n$ 矩阵，$\rk\vv A=m<n$. 证明 $\vv{Ax}=\vv b$ 对每个 $\vv b$ 都有无穷多解.
			\item\label{ex:sol:A19} 求经过点 $(1,1,1)$，$(2,0,1)$，$(0,2,3)$ 的平面方程（用齐次方程组求法向量）.
			\item\label{ex:sol:A20} 设 $\vv A=\begin{pmatrix}1&1\\1&1.0001\end{pmatrix}$. 分别求 $\vv{Ax}=(2,2)^\trans$ 与 $\vv{Ax}=(2,2.0001)^\trans$ 的解，计算右端的相对变化与解的相对变化之比.
		\end{exA}
		
		\Bgroup
		\begin{exB}
			\item\label{ex:sol:B1} 求两个齐次方程组 $\{x_1+x_2=0,\ x_2-x_4=0\}$ 与 $\{x_1-x_2+x_3=0,\ x_2-x_3+x_4=0\}$ 的公共非零解.
			\item\label{ex:sol:B2} 设 $\vv A$ 是实矩阵. 证明 $\vv{Ax}=\zero$ 与 $\vv A^\trans\vv{Ax}=\zero$ 同解，从而 $\rk(\vv A^\trans\vv A)=\rk\vv A$.（提示：$\vv x^\trans\vv A^\trans\vv{Ax}=\|\vv{Ax}\|^2$ 在第九章才有名字，但这里只需``实数的平方和为零则每个都为零''.）
			\item\label{ex:sol:B3} 证明 $\rk(\vv A+\vv B)\le\rk(\vv A\mid\vv B)\le\rk\vv A+\rk\vv B$.
			\item\label{ex:sol:B4} 设 $\vv A$ 为 $n$ 阶方阵，$\vv A^2=\vv I$. 证明 $\rk(\vv A+\vv I)+\rk(\vv A-\vv I)=n$.
			\item\label{ex:sol:B5} 证明 $\rk(\vv{AB})=\rk\vv B$ 当且仅当 $\ker\vv A\cap\imm\vv B=\{\zero\}$.
			\item\label{ex:sol:B6} 设 $\vv\eta^*$ 是 $\vv{Ax}=\vv b$（$\vv b\ne\zero$）的一个解，$\vv\xi_1,\ldots,\vv\xi_k$ 是对应齐次方程组的基础解系. 证明 $\vv\eta^*,\vv\eta^*+\vv\xi_1,\ldots,\vv\eta^*+\vv\xi_k$ 线性无关.
			\item\label{ex:sol:B7} 一张 $3\times3$ 的统计表只公开了三个行和与三个列和. 至少还要再公开几个单元格，才能确定全表？给出一种``公开哪几个格''的选法，并说明为什么某些同样个数的选法却不行.
			\item\label{ex:sol:B8} 对 $\vv A=\begin{pmatrix}1&2&1\\2&4&3\\3&6&4\end{pmatrix}$，求 $\vv b=(b_1,b_2,b_3)^\trans$ 应满足的条件，使 $\vv{Ax}=\vv b$ 有解.（提示：找出把各方程组合成 $0=\cdots$ 的系数.）
			\item\label{ex:sol:B9} 设 $\vv A$ 为 $n$ 阶方阵，$\rk\vv A=n-1$，且 $\vv A$ 的每行元素之和为零. 求 $\vv{Ax}=\zero$ 的通解.
			\item\label{ex:sol:B10} 设 $\vv A,\vv B$ 为 $n$ 阶方阵，$\vv{AB}=\vv O$，$\vv A+\vv B$ 可逆. 证明 $\rk\vv A+\rk\vv B=n$.
			\item\label{ex:sol:B11} 证明：两个 $m\times n$ 齐次方程组 $\vv{Ax}=\zero$ 与 $\vv{Bx}=\zero$ 同解，当且仅当 $\vv A$ 与 $\vv B$ 的行向量组等价.
			\item\label{ex:sol:B12} 设 $\vv A$ 为 $m\times n$ 矩阵. 证明：$\vv{Ax}=\vv b$ 对一切 $\vv b$ 有解，当且仅当存在 $n\times m$ 矩阵 $\vv C$ 使 $\vv{AC}=\vv I_m$；这时 $\vv C\vv b$ 就是一个解.
			\item\label{ex:sol:B13} 设 $\vv A$ 为 $n$ 阶方阵（$n\ge2$），$\rk\vv A=n-1$，$\vv\xi$ 是 $\vv{Ax}=\zero$ 的非零解. 证明 $\vv A^*$ 的每一列都是 $\vv\xi$ 的倍数.
			\item\label{ex:sol:B14} 用可确定性判据回答：方程组 $x_1+x_2+x_3=6$，$x_2+x_3+x_4=9$，$x_3+x_4+x_5=12$（例如三段相邻时间窗口内的观测总量）能确定 $x_1-x_4$ 吗？$x_1+x_2+x_3+x_4+x_5$ 呢？$x_1-x_4+x_5$ 呢？
			\item\label{ex:sol:B15} 设 $\vv A$ 为 $m\times n$ 矩阵，$\vv B$ 为 $n\times p$ 矩阵. 证明 $\dim\ker(\vv{AB})\le\dim\ker\vv A+\dim\ker\vv B$，并由此重新推出 Sylvester 不等式.
		\end{exB}
		
		\Cgroup
		\begin{exC}
			\item\label{ex:sol:C1} （Farkas 引理的二维形式）设 $\vv a_1,\vv a_2\in\mathbb{R}^2$ 线性无关，$\vv b\in\mathbb{R}^2$. 证明：下面两件事恰有一件成立：(1) 存在 $x_1,x_2\ge0$ 使 $x_1\vv a_1+x_2\vv a_2=\vv b$；(2) 存在 $\vv y$ 使 $\vv y^\trans\vv a_1\ge0$，$\vv y^\trans\vv a_2\ge0$，$\vv y^\trans\vv b<0$. 画图解释. 这是``无套利定价''与线性规划对偶理论的出发点.
			\item\label{ex:sol:C2} 设 $\vv A$ 是 $n$ 阶整数矩阵. 证明：若对任意整数向量 $\vv b$，$\vv{Ax}=\vv b$ 都有整数解，则 $\det\vv A=\pm1$.
			\item\label{ex:sol:C3} 证明：方阵 $\vv A$ 不用行交换就能做 LU 分解（$\vv L$ 单位下三角，$\vv U$ 上三角且可逆），当且仅当 $\vv A$ 的各阶顺序主子式都不为零；这时分解唯一.
			\item\label{ex:sol:C4} 设 $\vv A$ 是三对角矩阵（只有主对角线和相邻两条斜线上可能有非零元）. 证明：若它有 LU 分解，则 $\vv L,\vv U$ 都只有两条非零斜线；因此解三对角方程组只需与 $n$ 成正比的运算量. 这种方法称为``追赶法''.
			\item\label{ex:sol:C5} 用分块初等变换
			\[\begin{pmatrix}\vv I_n&\vv O\\ -\vv A&\vv I_m\end{pmatrix}\begin{pmatrix}\vv I_n&\vv B\\ \vv A&\vv O\end{pmatrix}\begin{pmatrix}\vv I_n&-\vv B\\ \vv O&\vv I_p\end{pmatrix}=\begin{pmatrix}\vv I_n&\vv O\\ \vv O&-\vv{AB}\end{pmatrix}\]
			给出 Sylvester 不等式 $\rk(\vv{AB})\ge\rk\vv A+\rk\vv B-n$ 的另一个证明.（提示：用第三章章末习题~\ref{ex:mat:B5} 的分块秩不等式.）
			\item\label{ex:sol:C6} 证明：$\mathbb{F}^n$ 中任意 $k$ 个线性流形（不必过原点）的交，要么是空集，要么仍是线性流形. 给出两个二维线性流形在 $\mathbb{R}^3$ 中交为空集的例子，并用秩的语言刻画这种情形.
			\item\label{ex:sol:C7} （列昂惕夫价格模型）设 $\vv A$ 为直接消耗系数矩阵，$\vv v$ 为各部门单位产品的增加值（工资、利润等）. 价格 $\vv p$ 满足``单位价格 $=$ 单位投入成本 $+$ 单位增加值''：$\vv p=\vv A^\trans\vv p+\vv v$. 说明它是第一章产量模型 $\vv x=\vv{Ax}+\vv d$ 的``转置''；证明在两个模型都可解时有 $\vv p^\trans\vv d=\vv v^\trans\vv x$，并解释其经济含义（国民收入的两种核算方法：支出法与收入法）.
			\item\label{ex:sol:C8} 设 $\vv A$ 为 $n$ 阶方阵. 证明存在正整数 $k$ 使 $\rk\vv A^k=\rk\vv A^{k+1}=\rk\vv A^{k+2}=\cdots$，且这样的最小的 $k$ 不超过 $n$.
		\end{exC}
	}
	
	% =====================================================================
	% 第七章　多项式
	% =====================================================================
	\newcommand{\ExPolynomial}{%
		\exsetheading
		\Agroup
		\begin{exA}
			\item\label{ex:poly:A1} 求商与余式：(1) $x^5-1$ 除以 $x^2+x+1$；(2) $2x^4-3x^3+4x^2-5x+6$ 除以 $x-2$（用综合除法）.
			\item\label{ex:poly:A2} 求 $x^{100}$ 分别除以 $x^2-1$，$x^2+1$，$(x-1)^2$ 所得的余式.
			\item\label{ex:poly:A3} 求 $a,b$，使 $x^2+x+1$ 整除 $x^4+ax^2+b$.（提示：设 $\omega$ 是 $x^2+x+1$ 的根，则 $\omega^3=1$.）
			\item\label{ex:poly:A4} 用辗转相除法求最大公因式：(1) $x^4+3x^3-x^2-4x-3$ 与 $3x^3+10x^2+2x-3$；(2) $x^3-2x^2+x$ 与 $x^2-1$.
			\item\label{ex:poly:A5} 证明：对任意多项式 $h$，$\gcd(f,g)=\gcd(f,\ g-hf)$. 这正是辗转相除法每一步的依据.
			\item\label{ex:poly:A6} 求 $u,v$ 使 $u(x)(x^2+1)+v(x)(x+1)=1$.
			\item\label{ex:poly:A7} 求 $2x^3-5x^2+x+2$ 的全部有理根，并在 $\mathbb{Q}$ 上把它分解为一次因式之积.
			\item\label{ex:poly:A8} 分别在 $\mathbb{Q}$，$\mathbb{R}$，$\mathbb{C}$ 上分解：(1) $x^4+x^2+1$；(2) $x^6-1$；(3) $x^3-2$.
			\item\label{ex:poly:A9} 判断 $1$ 是 $x^5-5x+4$ 的几重根.
			\item\label{ex:poly:A10} $a$ 取何值时，$x^3-3x+a$ 有重根？求出重根.
			\item\label{ex:poly:A11} 设 $\alpha_1,\alpha_2,\alpha_3$ 是 $x^3-2x^2+3x-4$ 的根. 不解方程，求 $\alpha_1+\alpha_2+\alpha_3$，$\alpha_1^2+\alpha_2^2+\alpha_3^2$，$\frac1{\alpha_1}+\frac1{\alpha_2}+\frac1{\alpha_3}$.
			\item\label{ex:poly:A12} 求次数最低的首一有理系数多项式，使 $1+\sqrt2$ 与 $\mathrm i$ 都是它的根.
			\item\label{ex:poly:A13} 求一个三次实系数首一多项式，以 $2$ 和 $1+\mathrm i$ 为根.
			\item\label{ex:poly:A14} 设 $f(x)=x^2-2x+1$，$\vv A=\begin{pmatrix}1&2\\0&1\end{pmatrix}$，计算 $f(\vv A)$；再设 $g(x)=x^3-x$，计算 $g(\vv A)$.
			\item\label{ex:poly:A15} 求过点 $(0,1)$，$(1,0)$，$(2,3)$，$(3,10)$ 的三次以下的插值多项式（用 Lagrange 公式）.
			\item\label{ex:poly:A16} 证明 $(x-1)\mid x^n-1$，并写出商式. 一般地，对任意 $a$，$(x-a)\mid x^n-a^n$.
			\item\label{ex:poly:A17} 把 $f(x)=x^4-2x^3+3$ 按 $x-1$ 的幂展开，即写成 $\sum c_k(x-1)^k$.（提示：反复用 $x-1$ 作综合除法，或用 $c_k=f^{(k)}(1)/k!$.）
			\item\label{ex:poly:A18} 多项式 $f$ 除以 $x-1$ 余 $3$，除以 $x-2$ 余 $5$. 求 $f$ 除以 $(x-1)(x-2)$ 的余式.
			\item\label{ex:poly:A19} 一个项目现在投入 $100$ 万元，一年后回收 $55$ 万元，两年后回收 $60.5$ 万元. 写出内部收益率满足的方程并精确求解.
			\item\label{ex:poly:A20} 证明：若 $f\mid g$ 且 $g\mid f$，则 $f=cg$（$c$ 为非零常数）.
			\item\label{ex:poly:A21} 设 $\gcd(f,g)=d$，$f=df_1$，$g=dg_1$. 证明 $\gcd(f_1,g_1)=1$.
		\end{exA}
		
		\Bgroup
		\begin{exB}
			\item\label{ex:poly:B1} 证明：若 $\gcd(f,g)=1$ 且 $\gcd(f,h)=1$，则 $\gcd(f,gh)=1$.
			\item\label{ex:poly:B2} 证明：$\gcd(x^m-1,\ x^n-1)=x^{\gcd(m,n)}-1$.（提示：模仿整数的辗转相除，先证 $x^m-1=x^{m-n}(x^n-1)+(x^{m-n}-1)$.）
			\item\label{ex:poly:B3} 设 $f=(x-1)^3(x+2)^2$. 计算 $f/\gcd(f,f')$，说明它恰好是``去掉重复因式''的结果. 一般地证明：若 $f=c\,p_1^{k_1}\cdots p_m^{k_m}$ 是标准分解，则 $f/\gcd(f,f')=c\,p_1\cdots p_m$.
			\item\label{ex:poly:B4} 证明：若 $p$ 不可约且 $p\mid f^2$，则 $p\mid f$. 举例说明``不可约''不能去掉.
			\item\label{ex:poly:B5} 设 $\ell_1,\ldots,\ell_n$ 是关于互异节点 $x_1,\ldots,x_n$ 的 Lagrange 基函数. 证明 $\sum_i\ell_i(x)=1$，$\sum_ix_i\ell_i(x)=x$，一般地 $\sum_ix_i^k\ell_i(x)=x^k$（$0\le k\le n-1$）.
			\item\label{ex:poly:B6} 证明：若复系数多项式 $f$ 满足 $f(x+1)=f(x)$ 对一切 $x$ 成立，则 $f$ 是常数.
			\item\label{ex:poly:B7} 证明 $x^3+x+1$ 在 $\mathbb{Q}$ 上不可约. 为什么对三次多项式只需检验有无有理根？四次多项式可以这样检验吗？
			\item\label{ex:poly:B8} 求 $x^n$ 除以 $x^2-3x+2$ 的余式.
			\item\label{ex:poly:B9} （Hermite 插值）求次数最低的多项式 $f$，使 $f(1)=1$，$f'(1)=2$，$f(2)=4$.（提示：设 $f=1+2(x-1)+c(x-1)^2$.）
			\item\label{ex:poly:B10} 设 $f$ 是有理系数多项式，$a,b\in\mathbb{Q}$，$\sqrt2\notin\mathbb{Q}$. 证明：若 $a+b\sqrt2$ 是 $f$ 的根，则 $a-b\sqrt2$ 也是 $f$ 的根.（提示：用 $(x-a)^2-2b^2$ 除 $f$.）
			\item\label{ex:poly:B11} 设 $f$ 是整系数多项式，$f(0)$ 与 $f(1)$ 都是奇数. 证明 $f$ 没有整数根.
			\item\label{ex:poly:B12} 求最小多项式意义下的``零化''：分别求出次数最低的首一多项式 $g$，使 $g(\vv A)=\vv O$，其中 $\vv A=\begin{pmatrix}0&-1\\1&0\end{pmatrix}$，$\vv A=\begin{pmatrix}2&0\\0&2\end{pmatrix}$，$\vv A=\begin{pmatrix}2&1\\0&2\end{pmatrix}$.
			\item\label{ex:poly:B13} 证明：现金流 $-C_0,\ C_1,\ldots,C_n$（$C_0>0$，$C_i\ge0$ 且不全为零）的净现值 $N(v)=-C_0+C_1v+\cdots+C_nv^n$ 在 $v>0$ 时严格递增，从而 $N(v)=0$ 至多一个正根；当 $C_1+\cdots+C_n>C_0$ 时这个正根落在 $(0,1)$ 中，即 IRR 唯一且为正.
			\item\label{ex:poly:B14} 设 $f$ 是 $n$ 次多项式，在 $n+1$ 个整数点 $0,1,\ldots,n$ 上都取整数值. 证明 $f$ 在一切整数点上都取整数值.（提示：用基 $\binom x0,\binom x1,\ldots,\binom xn$，其中 $\binom xk=\frac{x(x-1)\cdots(x-k+1)}{k!}$.）
			\item\label{ex:poly:B15} 证明 Lagrange 插值的``唯一性''也可以由线性代数推出：求值映射 $\mathbb{F}[x]_n\to\mathbb{F}^n$，$f\mapsto(f(x_1),\ldots,f(x_n))^\trans$ 是单射，从而是同构.
			\item\label{ex:poly:B16} 设 $p$ 是 $\mathbb{F}$ 上的不可约多项式，$f\in\mathbb{F}[x]$. 证明：$\gcd(p,f)=1$ 或 $p\mid f$，二者必居其一.
		\end{exB}
		
		\Cgroup
		\begin{exC}
			\item\label{ex:poly:C1} （Eisenstein 判别法）设 $f=a_nx^n+\cdots+a_0$ 是整系数多项式，$p$ 是素数，$p\nmid a_n$，$p\mid a_i$（$i<n$），$p^2\nmid a_0$. 证明 $f$ 在 $\mathbb{Q}$ 上不可约（可以先承认 Gauss 引理：整系数多项式若能分解为两个低次有理系数多项式之积，则也能分解为两个低次整系数多项式之积）. 由此说明 $x^n-2$ 对一切 $n$ 都不可约，所以 $\mathbb{Q}[x]$ 中有任意高次的不可约多项式.
			\item\label{ex:poly:C2} 设 $p$ 是素数. 证明分圆多项式 $\Phi_p(x)=x^{p-1}+x^{p-2}+\cdots+1$ 在 $\mathbb{Q}$ 上不可约.（提示：对 $\Phi_p(x+1)$ 用 Eisenstein 判别法.）
			\item\label{ex:poly:C3} （Newton 恒等式）记 $s_k=\alpha_1^k+\alpha_2^k+\alpha_3^k$，$\sigma_1,\sigma_2,\sigma_3$ 为初等对称多项式. 证明 $s_1=\sigma_1$，$s_2=\sigma_1s_1-2\sigma_2$，$s_3=\sigma_1s_2-\sigma_2s_1+3\sigma_3$，并求 $x^3-2x^2+3x-4$ 的根的立方和.
			\item\label{ex:poly:C4} （三次方程的判别式）设 $x^3+px+q$ 的三个复根为 $\alpha_1,\alpha_2,\alpha_3$，$\Delta=\prod_{i<j}(\alpha_i-\alpha_j)^2$. 证明 $\Delta=-4p^3-27q^2$，并说明（$p,q$ 为实数时）：$\Delta>0$ 时三个不同实根，$\Delta<0$ 时一实二虚，$\Delta=0$ 时有重根.
			\item\label{ex:poly:C5} （Chebyshev 多项式）证明存在 $n$ 次多项式 $T_n$ 使 $\cos n\theta=T_n(\cos\theta)$，且 $T_{n+1}=2xT_n-T_{n-1}$. 求 $T_n$ 的全部根. 由此证明：$\cos20^\circ$ 是 $8x^3-6x-1$ 的根，并验证这个三次多项式在 $\mathbb{Q}$ 上不可约. （结合``尺规作图只能得到经若干次开平方得到的数''，这说明 $60^\circ$ 角不能用尺规三等分.）
			\item\label{ex:poly:C6} 证明：在 $\mathbb{R}$ 上处处非负的实系数多项式，都能写成两个实系数多项式的平方和.（提示：利用实根的重数必为偶数，以及 $|a+b\mathrm i|^2=a^2+b^2$ 与 $(a^2+b^2)(c^2+d^2)=(ac-bd)^2+(ad+bc)^2$.）
			\item\label{ex:poly:C7} 设 $f\in\mathbb{Q}[x]$ 不可约，$g\in\mathbb{Q}[x]$ 与 $f$ 有一个公共复根. 证明 $f\mid g$. 由此说明：一个复数若是某个有理系数多项式的根，则以它为根的次数最低的首一有理系数多项式（称为它的\textbf{极小多项式}）唯一，且不可约.
			\item\label{ex:poly:C8} （$\mathbb{F}_2$ 上的多项式）在系数按模 $2$ 运算的多项式中，证明 $x^2+x+1$ 是唯一的二次不可约多项式，并找出全部三次不可约多项式. 这类多项式是循环冗余校验（CRC）与密码学的基础.
			\item\label{ex:poly:C9} 设 $f$ 是实系数多项式，它的系数序列的``变号次数''为 $V$（忽略零系数）. Descartes 符号法则说：$f$ 的正根个数（按重数计）不超过 $V$，且与 $V$ 的差是偶数. 用它重新解释第七章``内部收益率''建模实验中两个项目的结论.
		\end{exC}
	}
	
	% =====================================================================
	% 第八章　特征值与特征向量
	% =====================================================================
	\newcommand{\ExEigen}{%
		\exsetheading
		\Agroup
		\begin{exA}
			\item\label{ex:eig:A1} 求下列矩阵的特征值与特征向量（在 $\mathbb{C}$ 上考虑）：
			\[\begin{pmatrix}3&1\\0&2\end{pmatrix},\quad\begin{pmatrix}1&2\\2&1\end{pmatrix},\quad\begin{pmatrix}0&1\\-1&0\end{pmatrix},\quad\begin{pmatrix}4&-2\\1&1\end{pmatrix},\]\[\begin{pmatrix}2&0&0\\0&1&1\\0&1&1\end{pmatrix},\quad\begin{pmatrix}1&2&2\\2&1&2\\2&2&1\end{pmatrix}.\]
			\item\label{ex:eig:A2} 判断 A1 中各矩阵能否在 $\mathbb{R}$ 上对角化；能的，求可逆矩阵 $\vv P$ 使 $\vv P^{-1}\vv{AP}$ 为对角矩阵.
			\item\label{ex:eig:A3} 判断下列矩阵能否对角化：$\begin{pmatrix}1&1\\0&2\end{pmatrix}$，$\begin{pmatrix}2&1\\0&2\end{pmatrix}$，$\begin{pmatrix}2&0&0\\0&2&0\\1&0&2\end{pmatrix}$，$\begin{pmatrix}1&1&0\\0&2&0\\0&0&1\end{pmatrix}$.
			\item\label{ex:eig:A4} 设 $3$ 阶方阵 $\vv A$ 的特征值为 $1,2,3$. 求 $\det\vv A$，$\tr\vv A$，$\vv A^{-1}$ 的特征值，$\vv A^2-2\vv A+\vv I$ 的特征值，$\det(\vv A^*)$，$\det(\vv A^2+\vv I)$.
			\item\label{ex:eig:A5} 已知二阶方阵的迹为 $5$，行列式为 $6$，求它的特征值. 若迹为 $2$、行列式为 $5$ 呢？
			\item\label{ex:eig:A6} 证明：$\vv A$ 与 $\vv A^\trans$ 有相同的特征多项式. 它们的特征向量是否相同？举例说明.
			\item\label{ex:eig:A7} 证明：三角矩阵的特征值恰为其对角元.
			\item\label{ex:eig:A8} 证明：若 $\vv A^k=\vv O$，则 $\vv A$ 的特征值全为 $0$；若 $\vv A^2=\vv A$，则特征值只能是 $0$ 或 $1$；若 $\vv A^2=\vv I$，则特征值只能是 $\pm1$.
			\item\label{ex:eig:A9} 设 $\vv A=\begin{pmatrix}4&-2\\1&1\end{pmatrix}$. 利用对角化计算 $\vv A^n$.
			\item\label{ex:eig:A10} 设品牌转移矩阵为 $\vv P=\begin{pmatrix}0.9&0.2\\0.1&0.8\end{pmatrix}$. 求稳态分布，并利用对角化求 $\lim_{n\to\infty}\vv P^n$. 解释极限矩阵两列相同的含义.
			\item\label{ex:eig:A11} 已知 $\vv\xi=(1,1,-1)^\trans$ 是 $\vv A=\begin{pmatrix}2&-1&2\\5&a&3\\-1&b&-2\end{pmatrix}$ 的特征向量. 求 $a,b$ 与对应的特征值.
			\item\label{ex:eig:A12} 设 $\vv A=\begin{pmatrix}0&0&1\\1&1&x\\1&0&0\end{pmatrix}$. $x$ 取何值时 $\vv A$ 可对角化？
			\item\label{ex:eig:A13} 对二阶方阵直接验证 Cayley–Hamilton 定理：$\vv A^2-(\tr\vv A)\vv A+(\det\vv A)\vv I=\vv O$.
			\item\label{ex:eig:A14} 求下列矩阵的最小多项式：$\diag(1,1,2)$，$\begin{pmatrix}1&1&0\\0&1&0\\0&0&2\end{pmatrix}$，$\begin{pmatrix}0&1&0\\0&0&1\\0&0&0\end{pmatrix}$，$\begin{pmatrix}1&1&1\\1&1&1\\1&1&1\end{pmatrix}$.
			\item\label{ex:eig:A15} 判断 $\begin{pmatrix}1&1\\0&1\end{pmatrix}$ 与 $\begin{pmatrix}1&0\\1&1\end{pmatrix}$ 是否相似；若相似，求出 $\vv P$.
			\item\label{ex:eig:A16} 设 $\vv A$ 是上三角矩阵. 证明 $\spn(\vv e_1)$，$\spn(\vv e_1,\vv e_2)$，……都是 $\vv A$ 的不变子空间.
			\item\label{ex:eig:A17} 解差分方程 $a_{k+1}=a_k+2a_{k-1}$，$a_0=0$，$a_1=1$.
			\item\label{ex:eig:A18} （种群模型）某种群分为幼年与成年两组，第 $k$ 年的数量记为 $\vv x_k=(y_k,z_k)^\trans$. 每只成年个体每年平均产下 $1.1$ 只幼体，幼体有 $60\%$ 存活并长为成年，成年个体有一半活到下一年，即
			$y_{k+1}=1.1z_k$，$z_{k+1}=0.6y_k+0.5z_k$. 求矩阵的特征值与特征向量，说明种群的长期增长率和幼体与成年个体的稳定比例.
			\item\label{ex:eig:A19} 证明：实二阶方阵 $\vv A$ 没有实特征值，当且仅当 $(\tr\vv A)^2<4\det\vv A$.
			\item\label{ex:eig:A20} 求平面上向直线 $y=2x$ 的正投影、关于该直线的反射这两个变换的特征值与特征向量（不必写出矩阵，从几何上直接回答），再写出矩阵加以验证.
		\end{exA}
		
		\Bgroup
		\begin{exB}
			\item\label{ex:eig:B1} 设 $\vv u,\vv v\in\mathbb{R}^n$ 非零，$\vv A=\vv u\vv v^\trans$. 求 $\vv A$ 的全部特征值及其重数，并证明：$\vv A$ 可对角化当且仅当 $\vv v^\trans\vv u\ne0$.
			\item\label{ex:eig:B2} 设 $\vv A$ 有 $n$ 个互异的特征值，$\vv{AB}=\vv{BA}$. 证明 $\vv A$ 的特征向量都是 $\vv B$ 的特征向量，从而 $\vv B$ 也可对角化；进一步证明 $\vv B$ 是 $\vv A$ 的多项式.
			\item\label{ex:eig:B3} 证明：若 $\vv A$ 可对角化，则 $\vv A^k$ 与 $f(\vv A)$ 都可对角化. 逆命题成立吗？
			\item\label{ex:eig:B4} 设 $\vv A$ 可对角化，$W$ 是 $\vv A$ 的不变子空间. 证明 $\vv A|_W$ 也可对角化.（提示：最小多项式.）
			\item\label{ex:eig:B5} 证明：满足 $\vv A^2=-\vv I$ 的二阶实方阵都与 $\begin{pmatrix}0&-1\\1&0\end{pmatrix}$ 相似.
			\item\label{ex:eig:B6} （Gershgorin 圆盘定理）证明：方阵 $\vv A$ 的每个特征值都落在某个圆盘 $\{z\in\mathbb{C}:|z-a_{ii}|\le\sum_{j\ne i}|a_{ij}|\}$ 中. 由此证明：严格对角占优的矩阵（$|a_{ii}|>\sum_{j\ne i}|a_{ij}|$）可逆.
			\item\label{ex:eig:B7} 设 $\vv A$ 的每行元素之和都是 $s$. 证明 $s$ 是 $\vv A$ 的特征值. 若 $\vv A$ 可逆，求 $\vv A^{-1}$ 每行元素之和.
			\item\label{ex:eig:B8} 设 $\vv A=\begin{pmatrix}a&b&c\\c&a&b\\b&c&a\end{pmatrix}$（循环矩阵），$\omega=\mathrm e^{2\pi\mathrm i/3}$. 证明 $(1,\omega^k,\omega^{2k})^\trans$（$k=0,1,2$）都是 $\vv A$ 的特征向量，并求对应的特征值. 由此重新得到第五章章末习题~\ref{ex:det:B12} 的行列式公式.
			\item\label{ex:eig:B9} 解差分方程 $a_{k+1}=4a_k-4a_{k-1}$，$a_0=1$，$a_1=4$. 解释通项中为什么出现 $k\cdot2^k$.
			\item\label{ex:eig:B10} 证明 $\tr(\vv A^k)=\lambda_1^k+\cdots+\lambda_n^k$. 若三阶方阵满足 $\tr\vv A=0$，$\tr\vv A^2=6$，$\det\vv A=0$，求它的特征值.
			\item\label{ex:eig:B11} 证明：若 $\vv A^3=\vv A$，则 $\vv A$ 可对角化，且特征值在 $\{0,1,-1\}$ 中.
			\item\label{ex:eig:B12} $a,b,c$ 满足什么条件时，$\begin{pmatrix}2&0&0\\a&2&0\\b&c&1\end{pmatrix}$ 可对角化？
			\item\label{ex:eig:B13} 三个网页之间的链接为：网页 $1$ 链接到 $2,3$，网页 $2$ 链接到 $3$，网页 $3$ 链接到 $1$.（与判例不同，网页之间可以``互相引用''，形成环路.） 取 $d=1$ 写出转移矩阵，手算它的稳态分布.
			\item\label{ex:eig:B14} （矩阵指数）对可对角化矩阵 $\vv A=\vv P\vv D\vv P^{-1}$，定义 $\mathrm e^{t\vv A}=\vv P\diag(\mathrm e^{t\lambda_1},\ldots,\mathrm e^{t\lambda_n})\vv P^{-1}$. 对 $\vv A=\begin{pmatrix}0&1\\1&0\end{pmatrix}$ 计算 $\mathrm e^{t\vv A}$，并验证 $\vv x(t)=\mathrm e^{t\vv A}\vv x_0$ 满足微分方程组 $x_1'=x_2$，$x_2'=x_1$.
			\item\label{ex:eig:B15} 三状态马尔可夫链的转移矩阵为 $\begin{pmatrix}0.5&0.25&0.25\\0.25&0.5&0.25\\0.25&0.25&0.5\end{pmatrix}$. 求稳态分布与全部特征值，说明偏离稳态的部分每一步缩小为原来的多少.
			\item\label{ex:eig:B16} 证明：若 $\lambda$ 是 $\vv A$ 的 $k$ 重特征值（代数重数），则 $\rk(\lambda\vv I-\vv A)\ge n-k$.
			\item\label{ex:eig:B17} 对第三章``判例引证网络''中的四份判决，分别取阻尼系数 $d=0.5,\ 0.85,\ 0.95,\ 1$，按第八章建模实验的方法求 PageRank 得分（可用计算机），列表比较各判决的得分与排序. $d$ 越接近 $1$，得分越``集中''在哪些判决上？为什么？
		\end{exB}
		
		\Cgroup
		\begin{exC}
			\item\label{ex:eig:C1} 设 $\vv A$ 为 $m\times n$ 矩阵，$\vv B$ 为 $n\times m$ 矩阵. 证明 $\lambda^n\det(\lambda\vv I_m-\vv{AB})=\lambda^m\det(\lambda\vv I_n-\vv{BA})$. 由此说明 $\vv{AB}$ 与 $\vv{BA}$ 的非零特征值（连同重数）相同.（提示：计算 $\begin{pmatrix}\vv I_m&-\vv A\\ \vv O&\lambda\vv I_n\end{pmatrix}\begin{pmatrix}\lambda\vv I_m&\vv A\\ \vv B&\vv I_n\end{pmatrix}$ 等两种乘积的行列式.）
			\item\label{ex:eig:C2} （Schur 三角化）证明：每个复方阵都相似于一个上三角矩阵. 由此给出 Cayley–Hamilton 定理的另一个证明：对上三角矩阵 $\vv T$，证明 $(\vv T-t_{11}\vv I)\cdots(\vv T-t_{nn}\vv I)=\vv O$.
			\item\label{ex:eig:C3} 求下列矩阵的 Jordan 标准形：
			\[\begin{pmatrix}2&1&0\\0&2&0\\0&0&2\end{pmatrix},\quad\begin{pmatrix}1&1&1\\0&1&1\\0&0&1\end{pmatrix},\quad\begin{pmatrix}0&1&0&0\\0&0&0&0\\0&0&0&1\\0&0&0&0\end{pmatrix}.\]
			特征多项式为 $(\lambda-2)^4$ 的 $4$ 阶复矩阵，在相似的意义下共有几类？
			\item\label{ex:eig:C4} 证明：若一组可对角化的矩阵两两可交换，则它们可以\textbf{同时}对角化（存在同一个 $\vv P$ 使它们都化为对角矩阵）.
			\item\label{ex:eig:C5} 设 $\vv A$ 可对角化. 证明 $\vv A^k\to\vv O$（按元素）当且仅当 $\vv A$ 的所有特征值的模都小于 $1$. 用它说明：投入产出模型中，若直接消耗系数矩阵 $\vv A$ 可对角化且特征值的模都小于 $1$，则 $(\vv I-\vv A)^{-1}=\vv I+\vv A+\vv A^2+\cdots$.
			\item\label{ex:eig:C6} （友矩阵）证明 $\begin{pmatrix}0&0&\cdots&0&-a_0\\1&0&\cdots&0&-a_1\\0&1&\cdots&0&-a_2\\ \vdots&&\ddots&&\vdots\\0&0&\cdots&1&-a_{n-1}\end{pmatrix}$ 的特征多项式为 $\lambda^n+a_{n-1}\lambda^{n-1}+\cdots+a_0$. 由此说明：求多项式的根等价于求矩阵的特征值（数值软件正是这样求多项式的根的）.
			\item\label{ex:eig:C7} （赌徒破产）一位赌徒有 $k$ 元（$0\le k\le N$），每局以概率 $\frac12$ 赢或输 $1$ 元，输光或达到 $N$ 元时停止. 设 $p_k$ 为最终达到 $N$ 元的概率，证明 $p_k=\frac12(p_{k+1}+p_{k-1})$，并用特征根方法解出 $p_k=k/N$. 赢的概率改为 $p\ne\frac12$ 时呢？
			\item\label{ex:eig:C8} 证明：$n$ 阶三对角矩阵 $\vv T=\begin{pmatrix}2&-1&&\\-1&2&\ddots&\\&\ddots&\ddots&-1\\&&-1&2\end{pmatrix}$ 的特征值为 $2-2\cos\frac{k\pi}{n+1}$（$k=1,\ldots,n$），对应特征向量的第 $j$ 个分量为 $\sin\frac{jk\pi}{n+1}$. 对 $n=2,3$ 直接验证.
			\item\label{ex:eig:C9} （Perron–Frobenius 的二阶情形）设 $\vv A$ 是元素全为正的二阶矩阵. 证明 $\vv A$ 有两个互异实特征值，较大者 $r$ 为正且大于另一个的绝对值，属于 $r$ 的特征向量可取为各分量都为正. 由此解释：为什么 PageRank 引入阻尼系数后，``全部元素为正''能保证排序唯一.
			\item\label{ex:eig:C10} 证明：满足 $\vv A^2=-\vv I$ 的 $n$ 阶实方阵存在当且仅当 $n$ 为偶数，且这时 $\vv A$ 相似于 $\diag(\vv J,\ldots,\vv J)$，$\vv J=\begin{pmatrix}0&-1\\1&0\end{pmatrix}$.
		\end{exC}
	}
	
	% =====================================================================
	% 第九章　内积空间
	% =====================================================================
	\newcommand{\ExInnerProduct}{%
		\exsetheading
		\Agroup
		\begin{exA}
			\item\label{ex:ip:A1} 在 $\mathbb{R}^4$ 中，设 $\vv\alpha=(1,2,2,0)^\trans$，$\vv\beta=(2,0,1,2)^\trans$. 求 $\|\vv\alpha\|$，$\|\vv\beta\|$，$\langle\vv\alpha,\vv\beta\rangle$，它们的夹角，以及 $\|\vv\alpha-\vv\beta\|$.
			\item\label{ex:ip:A2} 判断下列二元函数是否为 $\mathbb{R}^2$ 上的内积：
			(1) $\langle\vv x,\vv y\rangle=x_1y_1+2x_2y_2$；(2) $\langle\vv x,\vv y\rangle=x_1y_2+x_2y_1$；(3) $\langle\vv x,\vv y\rangle=2x_1y_1+x_1y_2+x_2y_1+x_2y_2$；(4) $\langle\vv x,\vv y\rangle=x_1y_1-x_2y_2$.
			\item\label{ex:ip:A3} 用 Gram–Schmidt 方法把下列向量组正交化并单位化：
			(1) $(1,1,1)^\trans$，$(1,2,0)^\trans$，$(0,1,2)^\trans$；(2) $(1,1,0,0)^\trans$，$(1,0,1,0)^\trans$，$(-1,0,0,1)^\trans$.
			\item\label{ex:ip:A4} 求与 $(1,1,1)^\trans$ 和 $(1,0,-1)^\trans$ 都正交的单位向量.
			\item\label{ex:ip:A5} 已知 $\vv e_1=\frac13(1,2,2)^\trans$，$\vv e_2=\frac13(2,1,-2)^\trans$，$\vv e_3=\frac13(2,-2,1)^\trans$. 验证它们构成 $\mathbb{R}^3$ 的标准正交基，并只用内积求 $(3,6,0)^\trans$ 在这组基下的坐标.
			\item\label{ex:ip:A6} 求 $\vv\alpha=(1,2,3,4)^\trans$ 在 $W=\spn\big((1,1,0,0)^\trans,(0,0,1,1)^\trans\big)$ 上的正交投影，以及 $\vv\alpha$ 到 $W$ 的距离.
			\item\label{ex:ip:A7} 用最小二乘法求过数据点 $(0,1)$，$(1,2)$，$(2,2)$，$(3,4)$，$(4,6)$ 的最佳直线 $y=a+bx$，并计算残差平方和.
			\item\label{ex:ip:A8} 求 $a,b,c$ 使 $\begin{pmatrix}\frac1{\sqrt2}&a&\frac1{\sqrt6}\\0&b&-\frac2{\sqrt6}\\-\frac1{\sqrt2}&c&\frac1{\sqrt6}\end{pmatrix}$ 为正交矩阵.
			\item\label{ex:ip:A9} 证明：若 $\vv\alpha_1,\ldots,\vv\alpha_k$ 两两正交，则 $\|\vv\alpha_1+\cdots+\vv\alpha_k\|^2=\|\vv\alpha_1\|^2+\cdots+\|\vv\alpha_k\|^2$.
			\item\label{ex:ip:A10} 证明：$\big|\|\vv\alpha\|-\|\vv\beta\|\big|\le\|\vv\alpha-\vv\beta\|$.
			\item\label{ex:ip:A11} 用 Cauchy–Schwarz 不等式证明：(1) $(a_1+\cdots+a_n)^2\le n(a_1^2+\cdots+a_n^2)$；(2) 若 $a_i>0$，则 $(a_1+\cdots+a_n)\big(\frac1{a_1}+\cdots+\frac1{a_n}\big)\ge n^2$.
			\item\label{ex:ip:A12} 在 $C[0,1]$ 中取 $\langle f,g\rangle=\int_0^1f(x)g(x)\,\mathrm dx$. 求 $\|1\|$，$\|x\|$，$1$ 与 $x$ 的夹角，并把 $1,x$ 正交化.
			\item\label{ex:ip:A13} 设 $\vv A=\begin{pmatrix}1&1&0&1\\0&1&1&1\end{pmatrix}$. 求 $\ker\vv A$ 与行空间 $\imm\vv A^\trans$ 的标准正交基，并验证它们互为正交补.
			\item\label{ex:ip:A14} 求 $\vv A=\begin{pmatrix}1&1\\1&0\\0&1\end{pmatrix}$ 的 QR 分解.
			\item\label{ex:ip:A15} 设 $\vv u=\frac13(1,2,2)^\trans$. 写出 Householder 矩阵 $\vv H=\vv I-2\vv u\vv u^\trans$，验证它是对称的正交矩阵，并求 $\vv H\vv u$ 与 $\vv H(2,-1,0)^\trans$.
			\item\label{ex:ip:A16} 判断下列正交矩阵是旋转还是反射；若是旋转求旋转角，若是反射求反射轴：$\frac15\begin{pmatrix}3&-4\\4&3\end{pmatrix}$，$\frac15\begin{pmatrix}3&4\\4&-3\end{pmatrix}$.
			\item\label{ex:ip:A17} 对``过原点的直线''模型 $y=bx$，推导最小二乘解 $\hat b=\sum x_iy_i\big/\sum x_i^2$.
			\item\label{ex:ip:A18} 在 $\mathbb{R}^3$ 中，求点 $(1,2,3)$ 到直线 $\vv x=t(1,1,1)^\trans$ 的距离.
			\item\label{ex:ip:A19} 证明：两组标准正交基之间的过渡矩阵是正交矩阵.
			\item\label{ex:ip:A20} 三个法条的词频向量（关键词：违约、赔偿、合同、离婚）分别为 $(2,1,3,0)^\trans$，$(0,1,0,4)^\trans$，$(1,1,2,0)^\trans$，查询为 $(1,1,1,0)^\trans$. 分别用内积和余弦相似度给三个法条排序，结果有何不同？
			\item\label{ex:ip:A21} 设 $\vv u\in\mathbb{R}^n$ 是单位向量. 证明 $\vv P=\vv u\vv u^\trans$ 满足 $\vv P^\trans=\vv P$，$\vv P^2=\vv P$，且对任意 $\vv x$，$\vv P\vv x$ 是 $\vv x$ 在 $\spn(\vv u)$ 上的正交投影. 若 $\vv u$ 不是单位向量，投影矩阵应当怎样修改？
		\end{exA}
		
		\Bgroup
		\begin{exB}
			\item\label{ex:ip:B1} （极化恒等式）证明实内积空间中 $\langle\vv\alpha,\vv\beta\rangle=\frac14\big(\|\vv\alpha+\vv\beta\|^2-\|\vv\alpha-\vv\beta\|^2\big)$. 再证明：$\mathbb{R}^2$ 上的``范数'' $\|\vv x\|_1=|x_1|+|x_2|$ 不满足平行四边形法则，因而不能来自任何内积.
			\item\label{ex:ip:B2} （Bessel 不等式）设 $\vv e_1,\ldots,\vv e_k$ 是标准正交组. 证明 $\sum_{i=1}^k\langle\vv\alpha,\vv e_i\rangle^2\le\|\vv\alpha\|^2$，等号成立当且仅当 $\vv\alpha\in\spn(\vv e_1,\ldots,\vv e_k)$.
			\item\label{ex:ip:B3} 设实矩阵 $\vv P$ 满足 $\vv P^\trans=\vv P$，$\vv P^2=\vv P$. 证明 $\vv P$ 是向 $\imm\vv P$ 的正交投影的矩阵. 若只有 $\vv P^2=\vv P$，它是什么样的投影？举例说明.
			\item\label{ex:ip:B4} （加权最小二乘）若各观测值的可靠程度不同，给第 $i$ 个残差权重 $w_i>0$，最小化 $\sum w_i(y_i-a-bx_i)^2$. 证明正规方程变为 $\vv A^\trans\vv W\vv A\hat{\vv\beta}=\vv A^\trans\vv W\vv y$，$\vv W=\diag(w_1,\ldots,w_n)$. 它对应于哪个内积下的正交投影？
			\item\label{ex:ip:B5} 某商品的销售额 $y$ 随时间 $t$ 近似按指数增长 $y=c\,\mathrm e^{kt}$. 已知 $t=0,1,2,3$ 时 $y$ 约为 $2.0,2.7,3.7,5.0$. 两边取对数化为线性模型，用最小二乘估计 $c,k$. 这样做最小化的是哪个量？与直接最小化 $\sum(y_i-c\mathrm e^{kt_i})^2$ 有何不同？
			\item\label{ex:ip:B6} （Gram 行列式）设 $\vv\alpha_1,\ldots,\vv\alpha_k\in\mathbb{R}^n$，$\vv G=(\langle\vv\alpha_i,\vv\alpha_j\rangle)$. 证明 $\det\vv G\ge0$，等号成立当且仅当它们线性相关. 在 $k=2$ 时说明 $\sqrt{\det\vv G}$ 是两向量张成的平行四边形的面积.
			\item\label{ex:ip:B7} 证明：$\mathbb{R}^n$ 中点 $\vv x_0$ 到超平面 $\vv a^\trans\vv x=b$（$\vv a\ne\zero$）的距离为 $|\vv a^\trans\vv x_0-b|/\|\vv a\|$.
			\item\label{ex:ip:B8} 证明：实正交矩阵的复特征值的模都等于 $1$.
			\item\label{ex:ip:B9} （Fourier 逼近）在 $C[-\pi,\pi]$ 中取 $\langle f,g\rangle=\int_{-\pi}^\pi fg\,\mathrm dx$. 证明 $1,\cos x,\sin x,\cos2x,\sin2x$ 两两正交. 求 $f(x)=x$ 在它们张成的子空间上的正交投影.
			\item\label{ex:ip:B10} 设 $\vv A$ 为实矩阵. 证明：在 $\vv{Ax}=\vv y$ 的全部最小二乘解中，范数最小的那一个唯一，并且它属于行空间 $\imm\vv A^\trans$.
			\item\label{ex:ip:B11} 设 $\vv x,\vv y\in\mathbb{R}^n$，$\|\vv x\|=\|\vv y\|$，$\vv x\ne\vv y$. 令 $\vv u=(\vv x-\vv y)/\|\vv x-\vv y\|$. 证明 Householder 矩阵 $\vv I-2\vv u\vv u^\trans$ 把 $\vv x$ 变为 $\vv y$. 这是用 Householder 反射做 QR 分解的基础.
			\item\label{ex:ip:B12} 设两个变量的样本标准差分别为 $s_x,s_y$，相关系数为 $r$. 证明回归直线 $y=a+bx$ 的斜率为 $b=r\,s_y/s_x$，并且直线经过点 $(\bar x,\bar y)$.
			\item\label{ex:ip:B13} 证明：保持\textbf{夹角}的可逆线性变换（对一切非零向量，$\mathscr T\vv\alpha$ 与 $\mathscr T\vv\beta$ 的夹角等于 $\vv\alpha$ 与 $\vv\beta$ 的夹角）必是正交变换的非零倍数.
			\item\label{ex:ip:B14} 在回归中加入一个``虚拟变量''：设有两组数据，第一组 $y$ 值为 $3,5,4$，第二组为 $7,8,9$. 以 $\vv A=(\vv 1,\vv d)$ 为设计矩阵（$\vv d$ 在第二组为 $1$、第一组为 $0$），求最小二乘解，并解释两个系数的含义.
			\item\label{ex:ip:B15} 证明：在 $\mathbb{R}^2$ 中，任意两个旋转可交换，而一个旋转与一个反射一般不可交换；两个反射的乘积是旋转.
		\end{exB}
		
		\Cgroup
		\begin{exC}
			\item\label{ex:ip:C1} （Hadamard 不等式）用 QR 分解证明：对 $n$ 阶实方阵 $\vv A=(\vv a_1,\ldots,\vv a_n)$，$|\det\vv A|\le\|\vv a_1\|\cdots\|\vv a_n\|$，并说明它的几何意义（平行多面体的体积不超过棱长之积）.
			\item\label{ex:ip:C2} （Gram–Schmidt 的数值不稳定）设 $\varepsilon$ 很小，对 $\vv a_1=(1,\varepsilon,0,0)^\trans$，$\vv a_2=(1,0,\varepsilon,0)^\trans$，$\vv a_3=(1,0,0,\varepsilon)^\trans$ 作 Gram–Schmidt，在计算中把 $1+\varepsilon^2$ 舍入为 $1$. 证明所得的 $\vv e_2,\vv e_3$ 之间的夹角是 $60^\circ$ 而不是 $90^\circ$. 再用``修正的 Gram–Schmidt''（每得到一个 $\vv e_k$，立即从\textbf{剩余的}所有向量中减去它的分量）重做，比较结果.
			\item\label{ex:ip:C3} （Legendre 多项式）令 $P_n(x)=\frac1{2^nn!}\frac{\mathrm d^n}{\mathrm dx^n}(x^2-1)^n$. 证明 $P_0,P_1,P_2,\ldots$ 在 $\langle f,g\rangle=\int_{-1}^1fg\,\mathrm dx$ 下两两正交，并求 $\|P_n\|$.
			\item\label{ex:ip:C4} （Gauss–Markov 定理的线性代数部分）设 $\vv A$ 列满秩. 若 $\vv C\vv y$ 是参数 $\vv\beta$ 的``线性无偏估计''，即 $\vv{CA}=\vv I$. 令 $\vv C_0=(\vv A^\trans\vv A)^{-1}\vv A^\trans$（最小二乘）. 证明 $\vv C\vv C^\trans-\vv C_0\vv C_0^\trans$ 是半正定矩阵（第十一章）. 学过概率论后可以知道，这说明最小二乘估计的方差在一切线性无偏估计中最小.
			\item\label{ex:ip:C5} （复内积空间）在 $\mathbb{C}^n$ 上定义 $\langle\vv x,\vv y\rangle=\sum x_i\bar y_i$. 证明它满足共轭对称性、对第一个变量线性、正定性；证明 Cauchy–Schwarz 不等式仍成立. 满足 $\bar{\vv U}^\trans\vv U=\vv I$ 的复矩阵称为\textbf{酉矩阵}，证明它保持这个内积.
			\item\label{ex:ip:C6} 设 $f:\mathbb{R}^n\to\mathbb{R}^n$ 保持距离（$\|f(\vv x)-f(\vv y)\|=\|\vv x-\vv y\|$）且 $f(\zero)=\zero$. 证明 $f$ 保持内积，进而证明 $f$ 是线性的正交变换. 由此说明：保持距离的变换都可以分解为正交变换与平移的复合.
			\item\label{ex:ip:C7} （伪逆）设 $\vv A$ 列满秩，$\vv A^+=(\vv A^\trans\vv A)^{-1}\vv A^\trans$. 验证 Penrose 四条件：$\vv{AA}^+\vv A=\vv A$，$\vv A^+\vv{AA}^+=\vv A^+$，$(\vv{AA}^+)^\trans=\vv{AA}^+$，$(\vv A^+\vv A)^\trans=\vv A^+\vv A$. 说明 $\vv{AA}^+$ 是向哪个子空间的正交投影.
			\item\label{ex:ip:C8} 求空间中两条异面直线 $\vv x=\vv p+s\vv u$ 与 $\vv x=\vv q+t\vv v$ 之间的最短距离：把问题写成关于 $(s,t)$ 的最小二乘问题，并对 $\vv p=\zero$，$\vv u=(1,0,0)^\trans$，$\vv q=(0,1,1)^\trans$，$\vv v=(0,1,-1)^\trans$ 计算.
		\end{exC}
	}
	
	% =====================================================================
	% 第十章　矩阵分解
	% =====================================================================
	\newcommand{\ExDecomposition}{%
		\exsetheading
		\Agroup
		\begin{exA}
			\item\label{ex:dec:A1} 求正交矩阵 $\vv Q$ 使 $\vv Q^\trans\vv{AQ}$ 为对角矩阵：
			\[\begin{pmatrix}1&2\\2&1\end{pmatrix},\qquad\begin{pmatrix}3&-1\\-1&3\end{pmatrix},\]\[\begin{pmatrix}2&2&0\\2&2&0\\0&0&1\end{pmatrix},\qquad\begin{pmatrix}1&1&1\\1&1&1\\1&1&1\end{pmatrix}.\]
			\item\label{ex:dec:A2} 写出 $\begin{pmatrix}3&-1\\-1&3\end{pmatrix}$ 的谱分解 $\lambda_1\vv q_1\vv q_1^\trans+\lambda_2\vv q_2\vv q_2^\trans$，并验证两个投影矩阵之和为 $\vv I$、之积为 $\vv O$.
			\item\label{ex:dec:A3} 求二阶实对称矩阵 $\vv A$，使它的特征值为 $1,2$，且 $(1,1)^\trans$ 是属于 $1$ 的特征向量.
			\item\label{ex:dec:A4} 设三阶实对称矩阵 $\vv A$ 的特征值为 $1,1,-1$，属于 $-1$ 的特征向量为 $(0,1,1)^\trans$. 求 $\vv A$.
			\item\label{ex:dec:A5} 求 $\vv x^\trans\vv{Ax}$ 在单位球面 $\|\vv x\|=1$ 上的最大值与最小值，以及取到它们的 $\vv x$：$\vv A=\begin{pmatrix}2&2&0\\2&2&0\\0&0&1\end{pmatrix}$.
			\item\label{ex:dec:A6} 求 $x_1^2+4x_1x_2+x_2^2$ 在单位圆 $x_1^2+x_2^2=1$ 上的最大值与最小值.
			\item\label{ex:dec:A7} 求下列矩阵的奇异值分解：$\begin{pmatrix}2&0\\0&-3\end{pmatrix}$，$\begin{pmatrix}1&1\\1&1\end{pmatrix}$，$(3,4)$（$1\times2$ 矩阵）.
			\item\label{ex:dec:A8} 求 $\begin{pmatrix}1&1\\0&1\end{pmatrix}$ 的奇异值，并与它的特征值比较.（答案与黄金分割比有关.）
			\item\label{ex:dec:A9} 求 $\begin{pmatrix}3&1\\1&3\end{pmatrix}$ 的最佳秩一近似，并计算误差 $\|\vv A-\vv A_1\|_F$.
			\item\label{ex:dec:A10} 设 $\vv A=\begin{pmatrix}5&4\\4&5\end{pmatrix}$. 利用谱分解求 $\vv A^{10}$（写成两个矩阵之和的形式即可），并求对称正定矩阵 $\vv B$ 使 $\vv B^2=\vv A$.
			\item\label{ex:dec:A11} 证明：特征值全为 $0$ 的实对称矩阵是零矩阵. 对非对称矩阵举出反例.
			\item\label{ex:dec:A12} 证明：既对称又正交的实矩阵 $\vv A$ 满足 $\vv A^2=\vv I$，其特征值为 $\pm1$；并说明它在几何上是什么变换.
			\item\label{ex:dec:A13} 证明：$\vv A^\trans\vv A$ 与 $\vv{AA}^\trans$ 的非零特征值相同（重数也相同）.
			\item\label{ex:dec:A14} 证明：$\|\vv A\|_F^2=\tr(\vv A^\trans\vv A)=\sigma_1^2+\cdots+\sigma_r^2$.
			\item\label{ex:dec:A15} 设 $\vv A$ 是可逆实方阵. 证明 $|\det\vv A|=\sigma_1\sigma_2\cdots\sigma_n$，$\vv A^{-1}$ 的奇异值为 $1/\sigma_i$.
			\item\label{ex:dec:A16} 判断下列矩阵能否用正交矩阵对角化，说明理由：$\begin{pmatrix}1&2\\0&3\end{pmatrix}$，$\begin{pmatrix}1&2\\2&-2\end{pmatrix}$，$\begin{pmatrix}0&-1\\1&0\end{pmatrix}$（在实数范围内）.
			\item\label{ex:dec:A17} 四个已中心化的二维数据点为 $(2,1)$，$(-2,-1)$，$(1,2)$，$(-1,-2)$. 求样本协方差矩阵、第一主成分方向及其解释的方差比例.
			\item\label{ex:dec:A18} 计算 $\diag(10,1)$ 与 $\begin{pmatrix}1&1\\1&1.01\end{pmatrix}$ 的条件数（后者用近似值）. 后者的行列式与 $\diag(10,0.001)$ 相同，二者的条件数相同吗？
			\item\label{ex:dec:A19} 一张 $800\times600$ 像素的灰度图片，保留前 $k=40$ 个奇异值时需存储多少个数？压缩比是多少？$k$ 取多少时存储量反而超过原图？
			\item\label{ex:dec:A20} 证明：实对称矩阵的秩等于非零特征值的个数；奇异值的个数（非零）等于秩.
			\item\label{ex:dec:A21} 证明：实对称矩阵 $\vv A$ 满足``对一切 $\vv x\ne\zero$，$\vv x^\trans\vv A\vv x>0$''，当且仅当 $\vv A$ 的特征值全为正.（提示：用谱定理把 $\vv x^\trans\vv A\vv x$ 写成 $\sum\lambda_iy_i^2$.）
		\end{exA}
		
		\Bgroup
		\begin{exB}
			\item\label{ex:dec:B1} 设 $\vv A,\vv B$ 都是 $n$ 阶实对称矩阵且 $\vv{AB}=\vv{BA}$. 证明存在正交矩阵 $\vv Q$ 使 $\vv Q^\trans\vv{AQ}$ 与 $\vv Q^\trans\vv{BQ}$ 同时为对角矩阵.
			\item\label{ex:dec:B2} 证明：半正定的实对称矩阵 $\vv A$（特征值全 $\ge0$）有唯一的半正定平方根 $\vv B$（$\vv B^2=\vv A$，$\vv B$ 对称且特征值全 $\ge0$）.
			\item\label{ex:dec:B3} 设实对称矩阵 $\vv A$ 的特征值为 $\lambda_1\ge\lambda_2\ge\cdots\ge\lambda_n$，$\vv q_1$ 是属于 $\lambda_1$ 的单位特征向量. 证明 $\lambda_2=\max\{\vv x^\trans\vv{Ax}:\|\vv x\|=1,\ \vv x\perp\vv q_1\}$. 这说明第二主成分是``与第一主成分垂直的方向中方差最大的''.
			\item\label{ex:dec:B4} 证明：对任意 $\vv x$，$\|\vv{Ax}\|\le\sigma_1\|\vv x\|$，且等号可以取到. 因此 $\sigma_1=\max_{\vv x\ne\zero}\|\vv{Ax}\|/\|\vv x\|$，称为 $\vv A$ 的\textbf{谱范数}\index{pu3 fan4shu4@谱范数} $\|\vv A\|_2$.
			\item\label{ex:dec:B5} 证明：方阵 $\vv A$ 的任一特征值 $\lambda$ 满足 $\sigma_{\min}\le|\lambda|\le\sigma_{\max}$.
			\item\label{ex:dec:B6} 求 $\vv A=\begin{pmatrix}1&0&1\\0&1&1\end{pmatrix}$ 的奇异值分解.
			\item\label{ex:dec:B7} （伪逆）设 $\vv A=\vv U\vv\Sigma\vv V^\trans$，定义 $\vv A^+=\vv V\vv\Sigma^+\vv U^\trans$，其中 $\vv\Sigma^+$ 是把 $\vv\Sigma$ 转置并把非零元取倒数. 证明 $\vv A^+\vv b$ 是 $\vv{Ax}=\vv b$ 的范数最小的最小二乘解. 当 $\vv A$ 列满秩时，它与第九章的 $(\vv A^\trans\vv A)^{-1}\vv A^\trans$ 一致.
			\item\label{ex:dec:B8} 设 $p$ 个标准化指标两两相关系数都为 $\rho$（$0<\rho<1$），相关系数矩阵为 $(1-\rho)\vv I+\rho\vv J$（$\vv J$ 为全 $1$ 矩阵）. 求它的特征值与第一主成分，并求第一主成分解释的方差比例. 当 $\rho\to1$ 与 $\rho\to0$ 时比例分别趋于多少？
			\item\label{ex:dec:B9} 设两个指标的协方差矩阵为 $\begin{pmatrix}100&6\\6&1\end{pmatrix}$（第一个指标以``元''计，第二个以``万元''计）. 对协方差矩阵与相关系数矩阵分别作主成分分析，比较第一主成分. 由此说明：主成分分析对量纲敏感，为什么通常先标准化.
			\item\label{ex:dec:B10} 证明 $\|\vv A-\vv A_k\|_2=\sigma_{k+1}$.
			\item\label{ex:dec:B11} 证明：若 $\vv A$ 的奇异值互不相同且都不为零，则 $\vv A$ 的奇异值分解中的 $\vv u_i,\vv v_i$ 除了同时改变符号之外是唯一的.
			\item\label{ex:dec:B12} （正规矩阵）满足 $\vv A^\trans\vv A=\vv{AA}^\trans$ 的实方阵称为正规矩阵. 证明二阶实正规矩阵只有两类：对称矩阵，以及形如 $\begin{pmatrix}a&-b\\b&a\end{pmatrix}$ 的矩阵（旋转与伸缩的复合）.
			\item\label{ex:dec:B13} 设 $\vv A$ 是实对称矩阵. 证明 $\vv A$ 的特征值都落在 $[\min_i(a_{ii}-r_i),\ \max_i(a_{ii}+r_i)]$ 中，$r_i=\sum_{j\ne i}|a_{ij}|$. 并证明 $\lambda_{\max}\ge\max_ia_{ii}$.
			\item\label{ex:dec:B14} 对第十章``潜在语义分析''的例子，取 $k=1$，说明此时四份文书分别被表示成什么，查询``解约''与哪几份文书相似. 与 $k=2$ 比较，``压缩过头''的现象是否更严重？
			\item\label{ex:dec:B15} 证明：实对称矩阵 $\vv A$ 的特征值都在 $[a,b]$ 中，当且仅当 $\vv A-a\vv I$ 与 $b\vv I-\vv A$ 的二次型都非负.
		\end{exB}
		
		\Cgroup
		\begin{exC}
			\item\label{ex:dec:C1} （Eckart–Young 定理的证明）设 $\rk\vv B\le k$. 证明 $\ker\vv B$ 与 $\spn(\vv v_1,\ldots,\vv v_{k+1})$ 有非零公共向量 $\vv x$，由此证明 $\|\vv A-\vv B\|_2\ge\sigma_{k+1}$. 再借助 Weyl 不等式 $\sigma_{i+j-1}(\vv X+\vv Y)\le\sigma_i(\vv X)+\sigma_j(\vv Y)$（可承认）证明 Frobenius 范数的情形.
			\item\label{ex:dec:C2} （Cauchy 交错定理）设实对称矩阵 $\vv A$ 的特征值为 $\lambda_1\ge\cdots\ge\lambda_n$，删去最后一行一列所得矩阵的特征值为 $\mu_1\ge\cdots\ge\mu_{n-1}$. 证明 $\lambda_1\ge\mu_1\ge\lambda_2\ge\mu_2\ge\cdots\ge\mu_{n-1}\ge\lambda_n$.（提示：Courant–Fischer 极小极大刻画.）
			\item\label{ex:dec:C3} （复谱定理）证明：复方阵 $\vv A$ 可以被酉矩阵对角化（存在 $\vv U$ 使 $\bar{\vv U}^\trans\vv{AU}$ 为对角矩阵），当且仅当 $\vv A\bar{\vv A}^\trans=\bar{\vv A}^\trans\vv A$. 实对称矩阵、实正交矩阵、实反对称矩阵都满足这个条件.
			\item\label{ex:dec:C4} （正交 Procrustes 问题）给定 $n\times k$ 矩阵 $\vv A,\vv B$，求正交矩阵 $\vv Q$ 使 $\|\vv{QA}-\vv B\|_F$ 最小. 证明：若 $\vv{BA}^\trans=\vv U\vv\Sigma\vv V^\trans$，则最优解为 $\vv Q=\vv{UV}^\trans$. 这个问题出现在比较两组测量数据、对齐两份词向量等场合.
			\item\label{ex:dec:C5} （Ky Fan 定理）设实对称矩阵的特征值为 $\lambda_1\ge\cdots\ge\lambda_n$. 证明对任意列标准正交的 $n\times k$ 矩阵 $\vv X$，$\tr(\vv X^\trans\vv{AX})\le\lambda_1+\cdots+\lambda_k$. 由此说明前 $k$ 个主成分``一起''解释的方差也是最大的.
			\item\label{ex:dec:C6} （幂迭代）设 $\vv A$ 可对角化，特征值满足 $|\lambda_1|>|\lambda_2|\ge\cdots$. 证明：对``一般的''初始向量 $\vv x_0$，$\vv x_k=\vv A^k\vv x_0/\|\vv A^k\vv x_0\|$ 收敛到属于 $\lambda_1$ 的单位特征向量（可能差一个符号），收敛速度由 $|\lambda_2/\lambda_1|$ 决定. 对 $\begin{pmatrix}2&1\\1&2\end{pmatrix}$ 从 $(1,0)^\trans$ 出发迭代五步验证.
			\item\label{ex:dec:C7} （图的 Laplace 矩阵）设 $\vv L=\vv D-\vv A$ 是无向图的 Laplace 矩阵（见第五章章末习题~\ref{ex:det:C4}）. 证明 $\vv x^\trans\vv L\vv x=\sum_{\text{边 }ij}(x_i-x_j)^2$，因而 $\vv L$ 半正定；特征值 $0$ 的重数等于图的连通分支数. 第二小的特征值对应的特征向量（Fiedler 向量）的正负号常被用来把图``切''成两半——这就是谱聚类的思想.
			\item\label{ex:dec:C8} （正交回归）对中心化数据点 $\vv z_1,\ldots,\vv z_n\in\mathbb{R}^2$，求过原点的直线使各点到直线的\textbf{垂直}距离平方和最小. 证明最优方向就是第一主成分方向. 与普通最小二乘（只看竖直方向的残差）比较，两条直线何时相同？
			\item\label{ex:dec:C9} （扰动分析）设 $\vv A$ 可逆，$\vv{Ax}=\vv b$，$\vv A(\vv x+\delta\vv x)=\vv b+\delta\vv b$. 证明 $\dfrac{\|\delta\vv x\|}{\|\vv x\|}\le\kappa(\vv A)\dfrac{\|\delta\vv b\|}{\|\vv b\|}$，$\kappa(\vv A)=\sigma_1/\sigma_n$，并说明等号可以取到. 用它解释第六章病态方程组的例子.
		\end{exC}
	}
	
	% =====================================================================
	% 第十一章　二次型
	% =====================================================================
	\newcommand{\ExQuadratic}{%
		\exsetheading
		\Agroup
		\begin{exA}
			\item\label{ex:quad:A1} 写出下列二次型的矩阵与秩：
			(1) $x_1^2-3x_2^2+4x_1x_2-2x_2x_3$；(2) $x_1x_2+x_2x_3+x_3x_1$；(3) $(x_1+x_2+x_3)^2$；(4) $\sum_{i=1}^n x_i^2+\sum_{i<j}x_ix_j$.
			\item\label{ex:quad:A2} 写出以 $\begin{pmatrix}1&-1&0\\-1&2&3\\0&3&0\end{pmatrix}$ 为矩阵的二次型.
			\item\label{ex:quad:A3} 用配方法化为标准形，并写出所用的可逆变量替换：
			(1) $x_1^2+2x_2^2+5x_3^2+2x_1x_2+6x_2x_3$；(2) $x_1^2-4x_1x_2+4x_2^2-x_3^2$；(3) $x_1x_2-x_2x_3$.
			\item\label{ex:quad:A4} 用正交变换法把 $2x_1^2+3x_2^2+3x_3^2+4x_2x_3$ 化为标准形，写出正交矩阵. 曲面 $2x_1^2+3x_2^2+3x_3^2+4x_2x_3=1$ 是什么形状？
			\item\label{ex:quad:A5} 求 A3 中各二次型的正、负惯性指数与符号差，并写出规范形.
			\item\label{ex:quad:A6} 判断下列二次型是否正定：
			(1) $2x_1^2+x_2^2+3x_3^2+2x_1x_2$；(2) $x_1^2+x_2^2+x_3^2+4x_1x_2$；(3) $5x_1^2+x_2^2+5x_3^2+4x_1x_2-8x_1x_3-4x_2x_3$.
			\item\label{ex:quad:A7} $t$ 取何值时，$2x_1^2+x_2^2+x_3^2+2x_1x_2+tx_2x_3$ 正定？
			\item\label{ex:quad:A8} 判断曲线类型并求主轴方向：(1) $x^2-4xy+y^2=1$；(2) $3x^2+2xy+3y^2=8$.
			\item\label{ex:quad:A9} 判断下列矩阵是否合同：(1) $\diag(1,2,-3)$ 与 $\diag(-1,5,6)$；(2) $\begin{pmatrix}0&1\\1&0\end{pmatrix}$ 与 $\vv I_2$；(3) $\begin{pmatrix}0&1\\1&0\end{pmatrix}$ 与 $\diag(1,-1)$. 合同时求出 $\vv C$.
			\item\label{ex:quad:A10} 证明：若 $\vv A$ 正定，则 $\vv A^*$ 也正定.
			\item\label{ex:quad:A11} 求 $f(x,y)=x^3+y^3-3xy$ 的全部驻点，并用 Hessian 矩阵判断其类型.
			\item\label{ex:quad:A12} 判断 $\begin{pmatrix}2&-1&0\\-1&2&-1\\0&-1&2\end{pmatrix}$ 是否正定，并求它的 Cholesky 分解.
			\item\label{ex:quad:A13} 求 $\begin{pmatrix}4&2&-2\\2&10&2\\-2&2&6\end{pmatrix}$ 的 Cholesky 分解.
			\item\label{ex:quad:A14} 三种资产收益率的协方差矩阵为 $\vv\Sigma=\begin{pmatrix}0.04&0.01&0\\0.01&0.09&0\\0&0&0.01\end{pmatrix}$. 求组合 $\vv w=(0.5,0.3,0.2)^\trans$ 的方差与标准差.
			\item\label{ex:quad:A15} 两种资产收益率的标准差为 $30\%$ 与 $20\%$，互不相关. 求最小方差组合及其标准差.
			\item\label{ex:quad:A16} 证明：合同变换保持矩阵的对称性与秩. 合同变换保持行列式的符号吗？保持行列式的值吗？
			\item\label{ex:quad:A17} 分别在 $\mathbb{R}$ 上与 $\mathbb{C}$ 上写出二次型 $x_1^2-x_2^2+2x_3^2$ 的规范形.
			\item\label{ex:quad:A18} 判断下列方程表示的二次曲面类型：$x^2+y^2-z^2=1$，$x^2-y^2-z^2=1$，$x^2+y^2-z^2=0$，$x^2+y^2=1$.（提示：先看正、负惯性指数.）
			\item\label{ex:quad:A19} 说明 $\begin{pmatrix}1&2\\2&1\end{pmatrix}$ 不正定，并找出一个使二次型取负值的向量.
			\item\label{ex:quad:A20} 设 $\vv A$ 正定. 证明：$\vv A$ 的任一主子矩阵（取相同的若干行和列）都正定，从而所有主子式都为正.
		\end{exA}
		
		\Bgroup
		\begin{exB}
			\item\label{ex:quad:B1} 设 $\vv A,\vv B$ 正定且 $\vv{AB}=\vv{BA}$. 证明 $\vv{AB}$ 正定. 举例说明``可交换''不能去掉（此时 $\vv{AB}$ 甚至不对称）.
			\item\label{ex:quad:B2} 设 $\vv A$ 正定，$\vv B$ 半正定. 证明 $\det(\vv A+\vv B)\ge\det\vv A$，等号成立当且仅当 $\vv B=\vv O$.（提示：同时合同对角化.）
			\item\label{ex:quad:B3} 设 $\vv A$ 正定，$\vv B$ 对称. 证明：当 $|t|$ 充分小时 $\vv A+t\vv B$ 正定.
			\item\label{ex:quad:B4} 设 $\vv A$ 正定. 证明对 $i\ne j$ 有 $|a_{ij}|<\sqrt{a_{ii}a_{jj}}$，从而 $\vv A$ 的绝对值最大的元素在对角线上.
			\item\label{ex:quad:B5} 证明：实对称矩阵 $\vv A$ 满足 $\lambda_{\min}\|\vv x\|^2\le\vv x^\trans\vv{Ax}\le\lambda_{\max}\|\vv x\|^2$.
			\item\label{ex:quad:B6} 求 $\begin{pmatrix}0&1&1\\1&0&1\\1&1&0\end{pmatrix}$ 的正、负惯性指数（用两种方法：特征值；配方法）.
			\item\label{ex:quad:B7} $\lambda$ 取何值时，$\lambda(x_1^2+x_2^2+x_3^2)+2x_1x_2+2x_1x_3-2x_2x_3$ 正定？
			\item\label{ex:quad:B8} 求二次型 $\vv x^\trans\vv{Ax}$ 在约束 $\vv c^\trans\vv x=1$ 下的最小值（$\vv A$ 正定）：用拉格朗日乘数法证明最优解为 $\vv x^*=\dfrac{\vv A^{-1}\vv c}{\vv c^\trans\vv A^{-1}\vv c}$. 最小方差组合是它的特例.
			\item\label{ex:quad:B9} 设三种资产互不相关、方差都为 $\sigma^2$，期望收益率为 $\mu_1<\mu_2<\mu_3$. 求期望收益为 $m$ 的最小方差组合的权重，并说明有效前沿是一条抛物线（在``方差—期望收益''平面上）.
			\item\label{ex:quad:B10} 设 $\vv A=\begin{pmatrix}2&1\\1&2\end{pmatrix}$，$\vv B=\begin{pmatrix}1&0\\0&-1\end{pmatrix}$. 求可逆矩阵 $\vv C$ 使 $\vv C^\trans\vv{AC}=\vv I$ 且 $\vv C^\trans\vv{BC}$ 为对角矩阵.
			\item\label{ex:quad:B11} 设 $\vv A$ 正定. 求 $f(\vv x)=\vv x^\trans\vv{Ax}+\vv b^\trans\vv x+c$ 的最小值点与最小值. 若 $\vv A$ 只是半正定，最小值一定存在吗？
			\item\label{ex:quad:B12} 设 $\vv A$ 半正定，$\vv x^\trans\vv{Ax}=0$. 证明 $\vv{Ax}=\zero$.
			\item\label{ex:quad:B13} （Schur 补判据）证明：对称分块矩阵 $\begin{pmatrix}\vv A&\vv B\\ \vv B^\trans&\vv C\end{pmatrix}$ 正定，当且仅当 $\vv A$ 正定且 $\vv C-\vv B^\trans\vv A^{-1}\vv B$ 正定.
			\item\label{ex:quad:B14} 设 $\vv A$ 是二阶正定矩阵. 证明椭圆区域 $\{\vv x:\vv x^\trans\vv{Ax}\le1\}$ 的面积为 $\pi/\sqrt{\det\vv A}$.（提示：用主轴定理化为标准椭圆.）
			\item\label{ex:quad:B15} 设 $\vv A$ 为 $m\times n$ 实矩阵，$\mu>0$. 证明 $\vv A^\trans\vv A+\mu\vv I$ 正定，并说明岭回归 $\min\|\vv y-\vv{Ax}\|^2+\mu\|\vv x\|^2$ 的解为 $(\vv A^\trans\vv A+\mu\vv I)^{-1}\vv A^\trans\vv y$.
			\item\label{ex:quad:B16} 证明：实对称矩阵 $\vv A$ 正定，当且仅当存在正定矩阵 $\vv B$ 使 $\vv A=\vv B^2$；当且仅当 $\vv A^{-1}$ 存在且正定.
		\end{exB}
		
		\Cgroup
		\begin{exC}
			\item\label{ex:quad:C1} （Schur 乘积定理）对同阶矩阵定义逐元素乘积 $(\vv A\circ\vv B)_{ij}=a_{ij}b_{ij}$. 证明：若 $\vv A,\vv B$ 都半正定，则 $\vv A\circ\vv B$ 也半正定.（提示：把 $\vv B$ 写成 $\sum\lambda_k\vv q_k\vv q_k^\trans$.）
			\item\label{ex:quad:C2} （Minkowski 行列式不等式）设 $\vv A,\vv B$ 都是 $n$ 阶正定矩阵. 证明 $\det(\vv A+\vv B)^{1/n}\ge\det(\vv A)^{1/n}+\det(\vv B)^{1/n}$.（提示：同时合同对角化后用 $\prod(1+\mu_i)^{1/n}\ge1+\prod\mu_i^{1/n}$.）
			\item\label{ex:quad:C3} （有理数域上的二次型）证明 $x^2+y^2$ 与 $3x^2+3y^2$ 在 $\mathbb{R}$ 上合同，但在 $\mathbb{Q}$ 上不合同（即不存在有理系数的可逆替换把一个变为另一个）.（提示：若合同，则存在有理数 $a,b$ 使 $a^2+b^2=3$.）这说明惯性定理依赖于实数可以随意开平方.
			\item\label{ex:quad:C4} 函数 $f=x^2+y^4$ 与 $g=x^2-y^4$ 在原点处的 Hessian 矩阵相同且只是半正定，但前者在原点取极小值，后者不取极值. 由此说明 Hessian 半正定时二阶判别法失效.
			\item\label{ex:quad:C5} （广义特征值）设 $\vv A$ 对称，$\vv B$ 正定. 证明 $\max_{\vv x\ne\zero}\dfrac{\vv x^\trans\vv{Ax}}{\vv x^\trans\vv{Bx}}$ 等于方程 $\det(\vv A-\lambda\vv B)=0$ 的最大根. 统计学中的判别分析（Fisher LDA）正是这样求最佳投影方向的.
			\item\label{ex:quad:C6} （Mahalanobis 距离）设 $\vv\Sigma$ 正定，定义 $d(\vv x,\vv y)=\sqrt{(\vv x-\vv y)^\trans\vv\Sigma^{-1}(\vv x-\vv y)}$. 证明它满足三角不等式，且在仿射变换 $\vv x\mapsto\vv{Px}+\vv b$、$\vv\Sigma\mapsto\vv P\vv\Sigma\vv P^\trans$ 下不变. 它与欧氏距离有何不同？
			\item\label{ex:quad:C7} （Hilbert 矩阵）证明 $n$ 阶矩阵 $\vv H=\big(\frac1{i+j-1}\big)$ 正定.（提示：$\frac1{i+j-1}=\int_0^1x^{i-1}x^{j-1}\,\mathrm dx$，它是某个内积下的 Gram 矩阵.）
			\item\label{ex:quad:C8} （切点组合）设有无风险资产，收益率为 $r_f$. 证明使``夏普比率''$\dfrac{\vv\mu^\trans\vv w-r_f}{\sqrt{\vv w^\trans\vv\Sigma\vv w}}$（约束 $\vv 1^\trans\vv w=1$）最大的风险资产组合为 $\vv w\propto\vv\Sigma^{-1}(\vv\mu-r_f\vv 1)$.（提示：先去掉约束，比值关于 $\vv w$ 是零次齐次的；再用 Cauchy–Schwarz 不等式.）
			\item\label{ex:quad:C9} 证明：秩为 $r$、正惯性指数为 $p$ 的实二次型可以写成 $p$ 个线性型的平方和减去 $r-p$ 个线性型的平方和，并且不能写成少于 $r$ 个线性型的平方（带 $\pm$ 号）的组合.
		\end{exC}
	}
	
	% =====================================================================
	% 第十二章　*对偶空间与双线性型（选读）
	% =====================================================================
	\newcommand{\ExBilinear}{%
		\exsetheading
		\Agroup
		\begin{exA}
			\item\label{ex:bil:A1} 判断下列映射是否为相应空间上的线性函数：(1) $\mathbb{R}^3\to\mathbb{R}$，$\vv x\mapsto x_1-2x_3$；(2) $\vv x\mapsto x_1x_2$；(3) $\mathbb{R}[x]\to\mathbb{R}$，$f\mapsto f'(0)$；(4) $f\mapsto\int_0^1f(x)^2\,\mathrm dx$；(5) $M_n(\mathbb{R})\to\mathbb{R}$，$\vv X\mapsto\tr(\vv{AX})$（$\vv A$ 固定）.
			\item\label{ex:bil:A2} 求 $\mathbb{R}^3$ 的基 $(1,0,0)^\trans$，$(1,1,0)^\trans$，$(1,1,1)^\trans$ 的对偶基（把每个对偶基向量写成 $\vv x\mapsto a_1x_1+a_2x_2+a_3x_3$ 的形式）.
			\item\label{ex:bil:A3} 把线性函数 $f(\vv x)=x_1+2x_2+3x_3$ 用 A2 中的对偶基线性表出.
			\item\label{ex:bil:A4} 求 $W=\spn\big((1,1,0)^\trans,(0,1,1)^\trans\big)$ 的零化子 $W^0\subset(\mathbb{R}^3)^*$ 的一组基.
			\item\label{ex:bil:A5} 设线性映射 $\mathscr T:\mathbb{R}^2\to\mathbb{R}^3$ 在标准基下的矩阵为 $\begin{pmatrix}1&2\\0&1\\3&-1\end{pmatrix}$. 求对偶映射 $\mathscr T^*$ 在对偶基下的矩阵，并计算 $\mathscr T^*$ 作用在线性函数 $\vv y\mapsto y_1+y_2+y_3$ 上的结果.
			\item\label{ex:bil:A6} 求双线性型 $f(\vv x,\vv y)=x_1y_1+2x_1y_2-x_2y_1+3x_2y_2$ 在标准基下的度量矩阵，以及在基 $(1,1)^\trans$，$(1,-1)^\trans$ 下的度量矩阵，并验证 $\vv B=\vv P^\trans\vv{AP}$.
			\item\label{ex:bil:A7} 把 A6 中的 $f$ 分解为对称部分与交错部分之和.
			\item\label{ex:bil:A8} 判断下列双线性型是否非退化：(1) $x_1y_1-x_2y_2$；(2) $x_1y_2+x_2y_1$；(3) $(x_1+x_2)(y_1+y_2)$.
			\item\label{ex:bil:A9} 写出对称双线性型 $f(\vv x,\vv y)=x_1y_1+x_1y_2+x_2y_1+3x_2y_2$ 对应的二次型；反过来，由二次型 $x_1^2-4x_1x_2+2x_2^2$ 用极化恒等式求对应的对称双线性型.
			\item\label{ex:bil:A10} 验证 $\mathbb{R}^2$ 上的 $f(\vv x,\vv y)=\det(\vv x,\vv y)$ 是交错双线性型，写出它的度量矩阵. 它非退化吗？
			\item\label{ex:bil:A11} 设 $a_1,a_2,a_3$ 互异，$f_i(p)=p(a_i)$ 是 $\mathbb{R}[x]_3$ 上的线性函数. 对 $a_i=0,1,2$，求以 $f_1,f_2,f_3$ 为对偶基的 $\mathbb{R}[x]_3$ 的基.
			\item\label{ex:bil:A12} 在 $\mathbb{R}^4$ 中，设 $W$ 由 $x_1+x_2=0$，$x_3-x_4=0$ 定义. 求 $\dim W$ 与 $\dim W^0$，并写出 $W^0$ 的一组基.
			\item\label{ex:bil:A13} 证明 $\langle\vv A,\vv B\rangle=\tr(\vv{AB})$ 是 $M_n(\mathbb{R})$ 上非退化的对称双线性型. 它是内积吗？
			\item\label{ex:bil:A14} 对 $\mathbb{R}[x]_3$ 上的线性函数 $\varphi(p)=\int_0^1p(x)\,\mathrm dx$，求它在 $1,x,x^2$ 的对偶基下的坐标.
			\item\label{ex:bil:A15} （Riesz 表示定理的有限维情形）设 $V$ 是有限维实内积空间. 证明：对每个 $f\in V^*$，存在唯一的 $\vv\beta\in V$ 使 $f(\vv\alpha)=\langle\vv\alpha,\vv\beta\rangle$ 对一切 $\vv\alpha$ 成立.（提示：取标准正交基 $\vv e_1,\ldots,\vv e_n$，令 $\vv\beta=\sum f(\vv e_i)\vv e_i$.）
		\end{exA}
		
		\Bgroup
		\begin{exB}
			\item\label{ex:bil:B1} 证明：$f_1,\ldots,f_k\in V^*$ 线性无关，当且仅当映射 $V\to\mathbb{F}^k$，$\vv\alpha\mapsto(f_1(\vv\alpha),\ldots,f_k(\vv\alpha))^\trans$ 是满射.
			\item\label{ex:bil:B2} 设 $f,g\in V^*$ 且 $\ker f\subset\ker g$. 证明存在常数 $c$ 使 $g=cf$. 推广：若 $\ker f_1\cap\cdots\cap\ker f_k\subset\ker g$，则 $g$ 是 $f_1,\ldots,f_k$ 的线性组合.
			\item\label{ex:bil:B3} 证明：$\ker f_1\cap\cdots\cap\ker f_k=\{\zero\}$ 当且仅当 $f_1,\ldots,f_k$ 张成 $V^*$. 用方程组的语言解释这个结论.
			\item\label{ex:bil:B4} 证明：$\mathscr T$ 是单射当且仅当 $\mathscr T^*$ 是满射；$\mathscr T$ 是满射当且仅当 $\mathscr T^*$ 是单射.
			\item\label{ex:bil:B5} （Simpson 公式）求常数 $a,b,c$，使 $\int_0^2p(x)\,\mathrm dx=a\,p(0)+b\,p(1)+c\,p(2)$ 对一切 $p\in\mathbb{R}[x]_3$ 成立. 用对偶基的语言说明为什么这样的 $a,b,c$ 存在且唯一. 验证这个公式对三次多项式也成立.
			\item\label{ex:bil:B6} 证明：$V$ 上全体双线性型构成的空间与 $M_n(\mathbb{F})$ 同构，其中对称型、交错型分别构成维数为 $\frac{n(n+1)}2$、$\frac{n(n-1)}2$ 的子空间，且整个空间是二者的直和.
			\item\label{ex:bil:B7} 证明：若 $V$ 上存在非退化的交错双线性型，则 $\dim V$ 是偶数.
			\item\label{ex:bil:B8} 求双线性型 $f(\vv x,\vv y)=x_1y_1+x_1y_2+x_2y_1+x_2y_2+x_3y_3$ 的``根'' $\{\vv\alpha:f(\vv\alpha,\cdot)=0\}$.
			\item\label{ex:bil:B9} 在 $\mathbb{R}^2$ 上取 $f(\vv x,\vv y)=x_1y_1-x_2y_2$. 求全部``迷向向量''（$f(\vv\alpha,\vv\alpha)=0$），并说明它们构成两条直线. 再证明保持 $f$ 的线性变换 $\vv P$（$\vv P^\trans\diag(1,-1)\vv P=\diag(1,-1)$，$\det\vv P=1$，且保持 $x_1>|x_2|$ 的区域）都形如 $\begin{pmatrix}\cosh t&\sinh t\\ \sinh t&\cosh t\end{pmatrix}$.
			\item\label{ex:bil:B10} 设 $W$ 是 $V$ 的子空间，$\mathscr T:V\to V$ 线性且 $\mathscr T(W)\subset W$. 证明 $\mathscr T^*(W^0)\subset W^0$.
		\end{exB}
		
		\Cgroup
		\begin{exC}
			\item\label{ex:bil:C1} 证明：对无限维空间 $\mathbb{F}[x]$，自然映射 $\vv\alpha\mapsto\hat{\vv\alpha}$ 是单射但不是满射.（提示：$\mathbb{F}[x]^*$ 中的元素由它在 $1,x,x^2,\ldots$ 上的取值任意决定，可以``无限多个非零''.）
			\item\label{ex:bil:C2} 证明：$V$ 上的双线性型与线性映射 $V\to V^*$ 一一对应：$f\mapsto(\vv\alpha\mapsto f(\vv\alpha,\cdot))$. 用这个对应和对偶基的换基规则，重新推出度量矩阵的换基公式 $\vv B=\vv P^\trans\vv{AP}$.
			\item\label{ex:bil:C3} 证明：二阶矩阵 $\vv P$ 满足 $\vv P^\trans\vv J\vv P=\vv J$（$\vv J=\begin{pmatrix}0&1\\-1&0\end{pmatrix}$）当且仅当 $\det\vv P=1$. 四阶时，满足 $\vv P^\trans\diag(\vv J,\vv J)\vv P=\diag(\vv J,\vv J)$ 的矩阵（辛矩阵）的行列式是多少？
			\item\label{ex:bil:C4} 求 $M_n(\mathbb{R})$ 上对称双线性型 $\tr(\vv{AB})$ 的正、负惯性指数.（提示：对称矩阵与反对称矩阵在这个型下``互相正交''.）
			\item\label{ex:bil:C5} 设 $f$ 是 $V$ 上的对称双线性型，$W$ 是子空间，$W^\perp=\{\vv\alpha:f(\vv\alpha,\vv\beta)=0\ \forall\vv\beta\in W\}$. 证明 $\dim W^\perp\ge\dim V-\dim W$，$f$ 非退化时等号成立；举例说明即使 $f$ 非退化，$W\cap W^\perp$ 也可能不为零.
			\item\label{ex:bil:C6} 证明：$\mathbb{F}^n$ 上全体交错 $n$ 重线性函数构成的空间是一维的，由行列式张成. 这给出第五章章末习题~\ref{ex:det:C1} 的``对偶空间版本''.
			\item\label{ex:bil:C7} （线性规划对偶）考虑问题：在 $\vv{Ax}=\vv b$，$\vv x\ge\zero$ 下最小化 $\vv c^\trans\vv x$. 证明对任意满足 $\vv A^\trans\vv y\le\vv c$ 的 $\vv y$ 与任意可行的 $\vv x$，都有 $\vv b^\trans\vv y\le\vv c^\trans\vv x$（弱对偶）. 用``价格是数量空间上的线性函数''解释 $\vv y$ 的经济含义（影子价格）.
		\end{exC}
	} 之后，
%%      exercises.tex 中每道题都有标签 ex:<章键>:<组号>，本文件按标签引用题号）.
%% 章键：elim vs mat lmap det sol poly eig ip dec quad bil（依次为第一至十二章）.
%% =====================================================================

\providecommand{\ans}[2]{\par\smallskip\noindent\textbf{\ref{#1}.}\ #2\par}
\setlength{\emergencystretch}{3em}
\chapter{习题答案与提示}\label{app:answers}

本附录给出各章章末习题 A 组的答案与 B 组的提示. A 组的计算题给出最终结果，必要时附关键步骤；证明题只给出思路，完整的书写请读者自行完成. B 组只给出关键的一步或方向，希望读者在看提示之前先独立思考至少一刻钟. C 组一般不提供提示；正文或其他题目的答案中引用到的 C 组题例外，在该章末尾给出证明要点.

计算题的答案往往不唯一（例如基础解系、过渡矩阵的选取、特征向量的倍数），只要满足题目的全部要求，与这里的答案不同也是正确的. 检验的方法通常比答案本身更重要：代回原方程、计算迹与行列式、验证正交性.


\section{第一章　线性方程组与消元法}

\subsection*{A 组答案}

\ans{ex:elim:A1}{第 1、4 个是简化阶梯形；第 2 个是阶梯形但不是简化阶梯形（主元 $2$ 不为 $1$）；第 3、5 个不是阶梯形（第 3 个第二行主元在第一行主元左侧；第 5 个零行在非零行上方）.}

\ans{ex:elim:A2}{$\begin{pmatrix}3&0&-1&2\\0&1&4&0\\1&-1&0&5\end{pmatrix}$；反过来为 $x_1-2x_2=3$，$x_2+5x_3=-1$（未知量为 $x_1,x_2,x_3$，最后一列为常数项；$x_3$ 在第一个方程中的系数为 $0$）.}

\ans{ex:elim:A3}{$(x_1,x_2,x_3)=(1,2,3)$.}

\ans{ex:elim:A4}{$(x_1,x_2,x_3,x_4)=(1,-1,2,0)$.}

\ans{ex:elim:A5}{无解. $(r_3-r_1)-(r_2-2r_1)$ 给出 $(0,0,0\mid 2)$，即 $0=2$.}

\ans{ex:elim:A6}{主变量 $x_1,x_3$，自由变量 $x_2,x_4$. 全部解为 $x_1=4-2s-2t$，$x_2=s$，$x_3=1-t$，$x_4=t$（$s,t$ 任意）.}

\ans{ex:elim:A7}{$x_1=\frac27x_3+\frac37x_4$，$x_2=\frac57x_3+\frac47x_4$. 令 $x_3=7s$，$x_4=7t$，全部解为 $s(2,5,7,0)^\trans+t(3,4,0,7)^\trans$.}

\ans{ex:elim:A8}{$\begin{pmatrix}1&2&0&2\\0&0&1&-1\\0&0&0&0\end{pmatrix}$. 一种做法：$r_1\leftrightarrow r_2$；$r_2-2r_1$，$r_3-3r_1$；$-\frac14r_2$；$r_3+2r_2$；$r_1-r_2$.}

\ans{ex:elim:A9}{共四类：$\vv O$；$\begin{pmatrix}1&a\\0&0\end{pmatrix}$（$a$ 任意）；$\begin{pmatrix}0&1\\0&0\end{pmatrix}$；$\vv I_2$.}

\ans{ex:elim:A10}{由前两个方程得 $x=2$，$y=-1$，代入第三个得 $4-k=5$，$k=-1$.}

\ans{ex:elim:A11}{有解当且仅当 $b=3a$. 这时全部解为 $x=a-2t$，$y=t$（$t$ 任意）.}

\ans{ex:elim:A12}{$y=x^2-2x+3$.}

\ans{ex:elim:A13}{$y=1-2x+x^3$，即 $a=1$，$b=-2$，$c=0$，$d=1$.}

\ans{ex:elim:A14}{$30$ 元、$50$ 元、$80$ 元票分别为 $40$，$40$，$20$ 张.}

\ans{ex:elim:A15}{守恒方程 $3x_1=x_3$，$8x_1=2x_4$，$2x_2=2x_3+x_4$. 解为 $(x_1,x_2,x_3,x_4)=t(1,5,3,4)$，最小正整数解：$\mathrm{C_3H_8}+5\,\mathrm{O_2}\to3\,\mathrm{CO_2}+4\,\mathrm{H_2O}$.}

\ans{ex:elim:A16}{(1) 四个路口的方程为 $x_1-x_4=100$，$x_1-x_2=150$，$x_3-x_2=150$，$x_3-x_4=100$，全部解为 $x_1=t+100$，$x_2=t-50$，$x_3=t+100$，$x_4=t$. (2) 由 $x_2\ge0$ 得 $t\ge50$，故 $x_1$ 最少为 $150$. (3) 方程组无解：外部驶入总量 $1000$，驶出总量变为 $1100$，第一、三个方程之和减去第二、四个方程，得 $0=-100$. 守恒律是``进入路网的车辆总数等于离开的总数''.}

\ans{ex:elim:A17}{若 $a\ne0$，用第一个方程消去第二个方程中的 $x$，得 $(ad-bc)y=af-ce$. 解唯一当且仅当 $y$ 的系数 $ad-bc\ne0$. $a=0$ 时交换两个方程（或对 $c$ 讨论），同样得到条件 $ad-bc\ne0$.}

\ans{ex:elim:A18}{设 $\vv A\vv x_1=\vv A\vv x_2=\vv b$（这里先借用第三章的记号，也可以逐个方程写出），则 $\vv A(\vv x_1+t(\vv x_2-\vv x_1))=(1-t)\vv b+t\vv b=\vv b$. 若有两个不同的解，$t$ 取不同值给出不同的解，于是有无穷多个.}

\ans{ex:elim:A19}{第二个方程组的第一个方程等于第一个方程组两个方程之和，第二个方程相同，所以 (I) 的解都是 (II) 的解；反之 (I) 的第二个方程是 (II) 的第一个方程减去第二个方程. 两者的解集都是 $\{(t,t,3-2t)\}$.}

\ans{ex:elim:A20}{例如方程组 $x=0$，$y=0$. 同时做两步后得到 $x+y=0$，$x+y=0$，解集变大. 不矛盾：``同时''做的这两步合起来不是一次初等行变换，也不能被撤销（它把两行变成了相同的两行，丢失了信息）.}

\ans{ex:elim:A21}{有解当且仅当 $a=0$. 这时系数部分有 $2$ 个主元，$4$ 个未知量，有 $2$ 个自由变量.}

\subsection*{B 组提示}

\ans{ex:elim:B1}{用 $r_4-3r_1$，再加上 $r_2$，得 $(0,0,0,a-1\mid0)$；$r_3+r_2$ 得 $(0,0,a-1,0\mid b+1)$. $a\ne1$ 时解唯一；$a=1$，$b\ne-1$ 时无解；$a=1$，$b=-1$ 时有无穷多解：$(-1+s+t,\ 1-2s-2t,\ s,\ t)$.}

\ans{ex:elim:B2}{把第三个方程移到最前面再消元，得到的阶梯形后两行为 $(0,\lambda,-\lambda\mid3-\lambda)$ 与 $\big(0,0,-\lambda(\lambda+3)\mid-(\lambda+3)(\lambda-1)\big)$. $\lambda\ne0,-3$ 时有三个主元，解唯一；$\lambda=0$ 时倒数第二行变为 $(0,0,0\mid3)$，无解；$\lambda=-3$ 时最后一行为零行，只有两个主元，有无穷多解.}

\ans{ex:elim:B3}{齐次方程组只有零解说明简化阶梯形有 $n$ 个主元，即每一列都有主元. 而系数部分只有 $n$ 行，不可能出现 $(0,\ldots,0\mid d)$ 形的行.}

\ans{ex:elim:B4}{初等行变换得到的每一行都是原方程的线性组合. 无解时简化阶梯形中出现 $(0,\ldots,0\mid d)$，$d\ne0$，把这一行除以 $d$ 即可. 反之若组合得 $0=1$，任何解代入都会矛盾.}

\ans{ex:elim:B5}{依次做 $r_i+r_j$，$r_j-r_i$，$r_i+r_j$，最后把第 $j$ 行乘以 $-1$，逐步写出两行的变化即可.}

\ans{ex:elim:B6}{消元只用到四则运算，简化阶梯形的元素仍是有理数. 有解时令全部自由变量取 $0$，主变量由有理数算出. 解唯一时没有自由变量，唯一的解就是这个有理数解.}

\ans{ex:elim:B7}{把 $n$ 个方程相加，左边为 $0$，故必要条件是 $a_1+\cdots+a_n=0$. 它也充分：令 $x_n=t$，则 $x_k=t+a_k+a_{k+1}+\cdots+a_{n-1}$（$k<n$）.}

\ans{ex:elim:B8}{把第一个方程减第二个、加第三个……，交替相加. $n$ 为奇数时得 $2x_1=a_1-a_2+a_3-\cdots+a_n$，$x_1$ 唯一确定，其余依次确定；$n$ 为偶数时左边相消为 $0$.}

\ans{ex:elim:B9}{方程组为 $(\vv I-\vv A)\vv x=\vv d$（逐个写出即可），解得各部门总产出为 $(100,\ 200,\ 150)^\trans$.}

\ans{ex:elim:B10}{处理第 $k$ 列时，需要对下面 $n-k$ 行各做一次倍加，每次约 $n-k+1$ 次乘法，求和 $\sum_{k}(n-k)(n-k+1)\approx\frac13n^3$；回代只需约 $\frac12n^2$ 次.}

\ans{ex:elim:B11}{考察三个方程的系数部分化为阶梯形后的主元个数：$3$ 个时交于一点；$2$ 个且相容时交于一条直线；$2$ 个且出现 $0=d$ 时，三条交线方向相同（都平行于系数矩阵的解空间方向），没有公共点. 例子可以取 $z=0$，$x=0$，$x+z=1$.}

\ans{ex:elim:B12}{直接验证 $\vv A(\vv x+\vv y)=\vv A\vv x+\vv A\vv y$，$\vv A(k\vv x)=k\vv A\vv x$（或逐个方程验证）. 后半句：设 $\vv x_0$ 是一个解，任一解 $\vv x$ 满足 $\vv x-\vv x_0$ 是齐次方程组的解.}

\ans{ex:elim:B13}{开放题. 一种设定：未知量为三位律师各承办三类案件的件数，共 $9$ 个；已知条件给出 $3$ 个独立方程，故至少还缺 $6$ 条独立的信息. 关键是先写清楚未知量，再数主元个数.}

\ans{ex:elim:B14}{每行之和为零即 $\vv A(1,\ldots,1)^\trans=\zero$. 后半句由 B12 的结论.}

\ans{ex:elim:B15}{(I) 的每个解满足 (II) 的每个方程，因为后者是前者的线性组合；反之亦然. 逆命题对有解的方程组成立，证明要用到第六章的结论（齐次方程组同解当且仅当系数矩阵行向量组等价，第六章章末习题~\ref{ex:sol:B11}）.}

\ans{ex:elim:B16}{$\lambda\ne1,-2$ 时解唯一：$x=-\frac{\lambda+1}{\lambda+2}$，$y=\frac1{\lambda+2}$，$z=\frac{(\lambda+1)^2}{\lambda+2}$；$\lambda=1$ 时三个方程相同，有无穷多解；$\lambda=-2$ 时三个方程相加得 $0=3$，无解.}

\subsection*{C 组证明要点（正文引用的题目）}

\ans{ex:elim:C9}{设 $\vv R_1,\vv R_2$ 都是 $\vv M$ 的简化阶梯形. 它们行等价，以它们为系数矩阵的齐次方程组同解（定理~\ref{thm:sameSolution}）. 对列数 $n$ 归纳. $n=1$ 时，简化阶梯形是零列或 $\vv e_1$，由 $\vv M$ 是否为零唯一决定.
	
	设 $n>1$. 删去最后一列，得到的 $\vv R_1',\vv R_2'$ 仍是简化阶梯形，且仍行等价（用同一串初等行变换），由归纳假设 $\vv R_1'=\vv R_2'$. 于是两者在前 $n-1$ 列中的主元位置相同，设共有 $r$ 个主元. 再比较最后一列：
	\begin{enumerate}
		\item 若最后一列在 $\vv R_1$ 中是主元列，则 $\vv R_1$ 的第 $r+1$ 行为 $(0,\ldots,0,1)$，$\vv R_1\vv x=\zero$ 的解都有 $x_n=0$. 若它在 $\vv R_2$ 中不是主元列，令 $x_n=1$、其余自由变量取 $0$，可得 $\vv R_2\vv x=\zero$ 的一个解，其 $x_n=1$，与同解矛盾. 所以它在两者中同为主元列，都等于 $\vv e_{r+1}$.
		\item 若最后一列在两者中都不是主元列，令 $x_n=1$、其余自由变量取 $0$，两个方程组各有唯一一个这样的解，其主变量的值分别是两矩阵最后一列前 $r$ 个分量的相反数. 两方程组同解，故这两个解相同，最后一列相同.
\end{enumerate}}

\section{第二章　向量空间}

\subsection*{A 组答案}

\ans{ex:vs:A1}{(1) 是. (2) 不是：$\frac12\cdot\sqrt2=\frac{\sqrt2}2$ 的 $\sqrt2$ 系数不是整数，乘法不封闭. (3) 不是：不含 $1$，也不对除法封闭. (4) 不是：$\sqrt[3]2\cdot\sqrt[3]2=\sqrt[3]4$ 不在其中. (5) 不是：$\mathrm i\cdot\mathrm i=-1$ 不是纯虚数.}

\ans{ex:vs:A2}{(1) 是. (2) 不是：$(1,0)^\trans$ 的负向量不在其中. (3) 不是：不含零多项式，且两个 $n$ 次多项式之和可能降次. (4) 是. (5) 是. (6) 不是：不含零向量.}

\ans{ex:vs:A3}{(1) 线性相关：$(4,1,-1)^\trans=2(1,2,-1)^\trans+(2,-3,1)^\trans$. (2) 线性无关. (3) 线性相关：$\mathbb{R}[x]_3$ 是三维的，四个向量必相关. (4) 线性相关：全 $1$ 矩阵等于前两个之和.}

\ans{ex:vs:A4}{$\vv\beta=\frac54\vv\alpha_1+\frac14\vv\alpha_2-\frac14\vv\alpha_3-\frac14\vv\alpha_4$. （这四个向量两两``对应分量乘积之和''为零，每个系数都可以直接算出.）}

\ans{ex:vs:A5}{以三个向量为列作矩阵并消元：$r_2-r_1$ 得 $(0,2,-2)$，再 $r_3+\frac12r_2$ 得 $(0,0,t-1)$. $t=1$ 时出现自由变量，线性相关；$t\ne1$ 时线性无关.（学完第五章后，可用行列式 $2t-2$ 验证.）}

\ans{ex:vs:A6}{秩为 $3$，一个极大无关组为 $\vv\alpha_1,\vv\alpha_2,\vv\alpha_4$；$\vv\alpha_3=3\vv\alpha_1+\vv\alpha_2$，$\vv\alpha_5=\vv\alpha_1+\vv\alpha_2+\vv\alpha_4$.}

\ans{ex:vs:A7}{(1) 线性相关：三者之和为 $\zero$. (2) 线性无关. (3) 线性无关：设 $k_1(\vv\alpha_1+2\vv\alpha_2)+k_2(\vv\alpha_2+2\vv\alpha_3)+k_3(\vv\alpha_3+2\vv\alpha_1)=\zero$，由 $\vv\alpha_i$ 线性无关得 $k_1+2k_3=0$，$2k_1+k_2=0$，$2k_2+k_3=0$. 依次代入得 $k_1=-2k_3$，$k_2=4k_3$，$k_3=-8k_3$，故 $k_3=0$，从而 $k_1=k_2=0$.}

\ans{ex:vs:A8}{$1\cdot\zero=\zero$ 是系数不全为零的组合；$\vv\alpha-\vv\alpha=\zero$ 同理；$k\vv\alpha=\zero$ 且 $\vv\alpha\ne\zero$ 时由命题~\ref{prop:linspc}(5) 得 $k=0$.}

\ans{ex:vs:A9}{由定理~\ref{thm:relvandexp}，两个向量相关当且仅当其中一个可由另一个表出，即是另一个的倍数. 三个向量时只能说``其中一个是另两个的线性组合''，不一定是某一个的倍数，例如 $\vv e_1,\vv e_2,\vv e_1+\vv e_2$.}

\ans{ex:vs:A10}{(1) $\vv\alpha+\vv\alpha=1\vv\alpha+1\vv\alpha=(1+1)\vv\alpha$. (2) 两边乘 $k^{-1}$. (3) $-\vv\alpha$ 与 $\vv\alpha$ 之和为零，由负向量的唯一性.}

\ans{ex:vs:A11}{(1) $\dim=3$，基 $(-1,1,0,0)^\trans$，$(1,0,1,0)^\trans$，$(1,0,0,1)^\trans$. (2) $\dim=\frac{n(n+1)}2$，基 $\vv E_{ij}$（$i\le j$）. (3) $\dim=n^2-1$，基 $\vv E_{ij}$（$i\ne j$）与 $\vv E_{11}-\vv E_{ii}$（$i=2,\ldots,n$）. (4) $\dim=3$，基 $x-1$，$(x-1)^2$，$(x-1)^3$.}

\ans{ex:vs:A12}{(1) 是. (2) 不是：$(1,0,0)^\trans+(0,1,0)^\trans$ 不在其中. (3) 不是：对实数数乘不封闭. (4) 是.}

\ans{ex:vs:A13}{维数为 $3$，一组基为 $\vv\alpha_1,\vv\alpha_2,\vv\alpha_4$（$\vv\alpha_3=\vv\alpha_1+\vv\alpha_2$）.}

\ans{ex:vs:A14}{$2x^2-x+3=4+3(x-1)+2(x-1)^2$，坐标为 $(4,3,2)^\trans$.}

\ans{ex:vs:A15}{坐标为 $(-1,4,-5,4)^\trans$（从右下角元素开始依次确定）.}

\ans{ex:vs:A16}{(1) 是. (2) 不是：$\mathbb{Q}$ 对实数的数乘不封闭. (3) 作为 $\mathbb{R}$ 上的空间是 $2$ 维，作为 $\mathbb{C}$ 上的空间是 $1$ 维.}

\ans{ex:vs:A17}{例如 $a_0+a_1x+\cdots+a_{n-1}x^{n-1}\mapsto(a_0,\ldots,a_{n-1})^\trans$；$\begin{pmatrix}a&b\\c&d\end{pmatrix}\mapsto(a,b,c,d)^\trans$. 两者都是取坐标的映射，保持线性运算.}

\ans{ex:vs:A18}{$W_1+W_2$ 是直和（$(1,1,1)^\trans\notin W_1$，交为零），等于 $\mathbb{R}^3$，维数 $3$；$W_3\subset W_1$，$W_1+W_3=W_1$ 不是直和，维数 $2$.}

\ans{ex:vs:A19}{$\dim(W_1+W_2)=3$，$\dim(W_1\cap W_2)=1$：$(0,-1,1,1)^\trans=-(1,1,0,0)^\trans+(1,0,1,1)^\trans$ 同属两者. $2+2-1=3$，与维数公式一致.}

\ans{ex:vs:A20}{它们都可由一个含 $r$ 个向量的极大无关组线性表出，由基本引理任意 $r+1$ 个相关. 任意 $r$ 个无关的向量再添一个就相关，由引理~\ref{lem:addone} 的逆否形式，组中每个向量都可由它们线性表出.}

\subsection*{B 组提示}

\ans{ex:vs:B1}{设两者互不包含，取 $\vv\alpha\in W_1\setminus W_2$，$\vv\beta\in W_2\setminus W_1$，证明 $\vv\alpha+\vv\beta$ 不属于并集.}
\ans{ex:vs:B2}{取 (II) 的一个极大无关组，它含 $r$ 个向量且可由 (I) 的极大无关组表出，二者秩相同，用定理~\ref{thm:ndimbasis}(2) 在 (I) 张成的空间中论证.}
\ans{ex:vs:B3}{(III) 可由两者的极大无关组合并表出. 第二个不等式：左边的每个向量可由两组的极大无关组合并表出.}
\ans{ex:vs:B4}{新组含 $n$ 个向量，只需证明它能表出 $\vv\alpha_j$：$a_j\ne0$ 时 $\vv\alpha_j=\frac1{a_j}(\vv\beta-\sum_{i\ne j}a_i\vv\alpha_i)$；$a_j=0$ 时新组线性相关.}
\ans{ex:vs:B5}{按次数从高到低比较系数，证明线性无关；个数等于维数 $n$.}
\ans{ex:vs:B6}{$a_1,a_2$ 确定整个数列，映射 $\{a_n\}\mapsto(a_1,a_2)$ 是同构. 代入验证 $2^n$ 与 $(-1)^n$ 满足递推且线性无关. $a_1=a_2=1$ 时 $a_n=\frac13\big(2^n-(-1)^n\big)$.}
\ans{ex:vs:B7}{(1) 设组合为零，令 $x=0$ 并逐次求导，或除以 $\mathrm e^{3x}$ 后令 $x\to+\infty$. (2) $\sin(x+1)=\cos1\cdot\sin x+\sin1\cdot\cos x$. (3) 秩为 $3$：$\cos2x=2\cos^2x-1$.}
\ans{ex:vs:B8}{证明对加、减、乘封闭，再对除法用两次``分母有理化''. 线性无关可设 $a+b\sqrt2=-(c+d\sqrt2)\sqrt3$ 后平方，利用 $\sqrt2,\sqrt3,\sqrt6$ 都是无理数. 维数为 $4$.}
\ans{ex:vs:B9}{按维数分：$\{\zero\}$；过原点的直线；（$\mathbb{R}^3$ 中）过原点的平面；整个空间.}
\ans{ex:vs:B10}{用维数公式：$\dim(W_1\cap W_2)=k_1+k_2-\dim(W_1+W_2)\ge k_1+k_2-n>0$.}
\ans{ex:vs:B11}{$f=f(a)+(f-f(a))$. 交为零：常数多项式在 $a$ 处为零只能是零.}
\ans{ex:vs:B12}{$f(x)=\frac{f(x)+f(-x)}2+\frac{f(x)-f(-x)}2$；$\mathrm e^x=\cosh x+\sinh x$.}
\ans{ex:vs:B13}{对 $k$ 归纳，反复使用维数公式：$\dim(U\cap W)\ge\dim U+\dim W-n$.}
\ans{ex:vs:B14}{这 $n-1$ 个向量线性无关（用坐标验证），$\dim W=n-1$. 补空间可取 $\spn(\vv\alpha_n)$（需验证 $\vv\alpha_n\notin W$：$W$ 中向量的坐标交错和为零）.}
\ans{ex:vs:B15}{$\vv x=x_1(1,\ldots,1)^\trans+(\vv x-x_1(1,\ldots,1)^\trans)$. 两种分解都是直和；与 $V_1$ 互补的子空间有无穷多个，第九章会说明例~\ref{eg:pricedecomp} 中的 $R$ 是其中``垂直''的那一个.}
\ans{ex:vs:B16}{每个矩阵唯一地拆成``对角线及以上''与``对角线以下''两部分. $\dim W_1=\frac{n(n+1)}2$.}
\ans{ex:vs:B17}{(1) 设 $W$ 的基为 $\vv\varepsilon_1,\ldots,\vv\varepsilon_k$（$k<n$），扩充为 $\mathbb{R}^n$ 的基 $\vv\varepsilon_1,\ldots,\vv\varepsilon_n$. ``取 $\vv x$ 在这组基下的第 $n$ 个坐标''是 $x_1,\ldots,x_n$ 的一次齐次式（解一个线性方程组即可写出），它在 $W$ 上恒为零，又在 $\vv\varepsilon_n$ 处为 $1$，系数不全为零. (2) 条件是 $a_1+a_2t+\cdots+a_nt^{n-1}=0$，这是一个次数不超过 $n-1$ 的非零多项式，至多 $n-1$ 个根. (3) 例如 $t=2$：$\vv v(2)=(1,2,4)^\trans$ 不满足三个方程中的任何一个.}

\ans{ex:vs:B18}{对 $k$ 归纳. 已选好线性无关的 $\vv\varepsilon_1,\ldots,\vv\varepsilon_{k-1}$（$k\le n$）后，$U=\spn(\vv\varepsilon_1,\ldots,\vv\varepsilon_{k-1})$ 的维数 $k-1<n$，也是真子空间. 由覆盖定理取 $\vv\varepsilon_k\notin W_1\cup\cdots\cup W_s\cup U$；由引理~\ref{lem:addone}，$\vv\varepsilon_1,\ldots,\vv\varepsilon_k$ 仍线性无关. 选满 $n$ 个即得一组基.}

\section{第三章　矩阵}

\subsection*{A 组答案}

\ans{ex:mat:A1}{$\vv{AB}=\begin{pmatrix}5&9\\-2&4\end{pmatrix}$，$\vv{BA}=\begin{pmatrix}2&1&5\\-1&1&-2\\1&11&6\end{pmatrix}$，$\vv A^\trans\vv B^\trans=(\vv{BA})^\trans$. 由此也看出 $(\vv{AB})^\trans=\vv B^\trans\vv A^\trans\ne\vv A^\trans\vv B^\trans$（后者甚至与前者不同型）.}

\ans{ex:mat:A2}{$2(1,3)^\trans-(2,-1)^\trans+(0,4)^\trans=(0,11)^\trans$. $\vv e_i^\trans\vv A\vv e_j=a_{ij}$；$\vv A\vv e_j$ 是 $\vv A$ 的第 $j$ 列，$\vv e_i^\trans\vv A$ 是第 $i$ 行.}

\ans{ex:mat:A3}{$\begin{pmatrix}1&n\\0&1\end{pmatrix}$；$\begin{pmatrix}\lambda^n&n\lambda^{n-1}\\0&\lambda^n\end{pmatrix}$；$\begin{pmatrix}\cos n\theta&-\sin n\theta\\ \sin n\theta&\cos n\theta\end{pmatrix}$；第四个矩阵记为 $\vv N$，$\vv N^2=\vv E_{13}$，$n\ge3$ 时 $\vv N^n=\vv O$. 均可用归纳法证明.}

\ans{ex:mat:A4}{$(\vv A+\vv B)^2=\begin{pmatrix}10&8\\12&10\end{pmatrix}$，$\vv A^2+2\vv{AB}+\vv B^2=\begin{pmatrix}16&8\\12&4\end{pmatrix}$，不相等；$(\vv A+\vv B)(\vv A-\vv B)=\begin{pmatrix}-6&4\\-6&6\end{pmatrix}\ne\vv A^2-\vv B^2=\begin{pmatrix}0&4\\-6&0\end{pmatrix}$. 原因是 $\vv{AB}\ne\vv{BA}$.}

\ans{ex:mat:A5}{依次为
$\begin{pmatrix}1&3&-2\\-\frac32&-3&\frac52\\1&1&-1\end{pmatrix}$，
$\begin{pmatrix}1&-1&0\\0&1&-1\\0&0&1\end{pmatrix}$，
$\frac14\begin{pmatrix}3&-2&1\\-2&4&-2\\1&-2&3\end{pmatrix}$，
$\begin{pmatrix}0&0&0&\frac14\\0&0&\frac13&0\\0&\frac12&0&0\\1&0&0&0\end{pmatrix}$.}

\ans{ex:mat:A6}{(1) $\vv X=\begin{pmatrix}2&-23\\0&8\end{pmatrix}$. (2) $\vv X=\begin{pmatrix}1&-1\\-2&4\end{pmatrix}$（注意 $\vv X=\vv B\vv A^{-1}$，逆矩阵乘在右边）.}

\ans{ex:mat:A7}{$r_2-r_1$，$r_3-ar_1$，再 $r_3+r_2$，得阶梯形的后两行为 $(0,a-1,1-a)$ 与 $\big(0,0,(1-a)(a+2)\big)$. $a\ne1$ 且 $a\ne-2$ 时秩为 $3$；$a=1$ 时后两行全为零，秩为 $1$；$a=-2$ 时第二行非零、第三行为零，秩为 $2$.（学完第五章后，可用行列式 $-(a-1)^2(a+2)$ 验证.）}

\ans{ex:mat:A8}{秩为 $3$. 简化阶梯形的主元在第 $1,3,5$ 列，故第 $1,3,5$ 列构成极大无关组；第 $2$ 列 $=2\times$第 $1$ 列，第 $4$ 列 $=$ 第 $1$ 列 $+$ 第 $3$ 列.}

\ans{ex:mat:A9}{$\vv E(1,3)=\begin{pmatrix}0&0&1\\0&1&0\\1&0&0\end{pmatrix}$，逆为自身；$\vv E(2(5))=\diag(1,5,1)$，逆为 $\diag(1,\frac15,1)$；$\vv E(3,1(-2))=\begin{pmatrix}1&0&0\\0&1&0\\-2&0&1\end{pmatrix}$，逆为 $\vv E(3,1(2))$. $\vv E(3,1(-2))\vv A$ 把 $\vv A$ 第 $1$ 行的 $-2$ 倍加到第 $3$ 行；$\vv A\vv E(3,1(-2))$ 把 $\vv A$ 第 $3$ 列的 $-2$ 倍加到第 $1$ 列.}

\ans{ex:mat:A10}{$\begin{pmatrix}2&1\\4&3\end{pmatrix}=\vv E(2,1(2))\,\vv E(1,2(1))\,\vv E(1(2))=\begin{pmatrix}1&0\\2&1\end{pmatrix}\begin{pmatrix}1&1\\0&1\end{pmatrix}\begin{pmatrix}2&0\\0&1\end{pmatrix}$. 写法不唯一.}

\ans{ex:mat:A11}{$(\vv{AB})_{ij}=\sum_ka_{ik}b_{kj}$，$i>j$ 时每一项中 $k<i$ 使 $a_{ik}=0$ 或 $k\ge i>j$ 使 $b_{kj}=0$；$i=j$ 时只剩 $a_{ii}b_{ii}$. 可逆上三角矩阵的逆也是上三角矩阵：对 $(\vv A\mid\vv I)$ 作消元时只需用下面的行消去上面的行.}

\ans{ex:mat:A12}{(1) 用 $(\vv{AB})^\trans=\vv B^\trans\vv A^\trans$ 逐一验证. (2) $(\vv{AB})^\trans=\vv B^\trans\vv A^\trans=\vv{BA}$. (3) $(\vv C^\trans\vv C)^\trans=\vv C^\trans\vv C$.}

\ans{ex:mat:A13}{(1) $\vv x\vv y^\trans$ 的 $(i,i)$ 元为 $x_iy_i$，求和即得. (2) $\tr(\vv A\vv E_{ji})=\vv e_i^\trans\vv A\vv e_j=a_{ij}$，故每个 $a_{ij}=0$. (3) 两边取迹，由命题~\ref{prop:trace}，左边为 $0$，右边为 $n$.}

\ans{ex:mat:A14}{$\vv A(\vv A-\vv I)=2\vv I$，故 $\vv A^{-1}=\frac12(\vv A-\vv I)$；$(\vv A+2\vv I)(\vv A-3\vv I)=\vv A^2-\vv A-6\vv I=-4\vv I$，故 $(\vv A+2\vv I)^{-1}=-\frac14(\vv A-3\vv I)$. $\vv A+\vv I$ 不一定可逆：$\vv A=-\vv I$ 满足条件，$\vv A+\vv I=\vv O$.}

\ans{ex:mat:A15}{左乘 $\vv A^{-1}$. 反例：$\vv A=\begin{pmatrix}1&0\\0&0\end{pmatrix}$，$\vv B=\vv O$，$\vv C=\begin{pmatrix}0&0\\0&1\end{pmatrix}$.}

\ans{ex:mat:A16}{$\begin{pmatrix}\vv A^{-1}&\vv O\\ \vv O&\vv B^{-1}\end{pmatrix}$，$\begin{pmatrix}\vv O&\vv B^{-1}\\ \vv A^{-1}&\vv O\end{pmatrix}$. 所求逆矩阵为 $\begin{pmatrix}0&0&\frac12&0\\0&0&0&\frac15\\3&-2&0&0\\-1&1&0&0\end{pmatrix}$.}

\ans{ex:mat:A17}{记 $\vv B=\begin{pmatrix}2&1\\1&3\end{pmatrix}$，矩阵为 $\begin{pmatrix}\vv I&\vv B\\ \vv O&\vv I\end{pmatrix}$，其平方为 $\begin{pmatrix}\vv I&2\vv B\\ \vv O&\vv I\end{pmatrix}=\begin{pmatrix}1&0&4&2\\0&1&2&6\\0&0&1&0\\0&0&0&1\end{pmatrix}$，$n$ 次方为 $\begin{pmatrix}\vv I&n\vv B\\ \vv O&\vv I\end{pmatrix}$.}

\ans{ex:mat:A18}{$\vv Q\vv C=\begin{pmatrix}5600&8300\\6750&9800\end{pmatrix}$：第 $i$ 行是第 $i$ 个工厂每天的总成本与总销售额. 每天利润分别为 $2700$ 元与 $3050$ 元.}

\ans{ex:mat:A19}{$\vv A=\begin{pmatrix}0&1&1&0\\1&0&1&0\\1&1&0&1\\0&0&1&0\end{pmatrix}$，$\vv A^2=\begin{pmatrix}2&1&1&1\\1&2&1&1\\1&1&3&0\\1&1&0&1\end{pmatrix}$，$\vv A^3=\begin{pmatrix}2&3&4&1\\3&2&4&1\\4&4&2&3\\1&1&3&0\end{pmatrix}$. $(\vv A^2)_{ii}$ 等于城市 $i$ 的航线数；从 $1$ 到 $4$ 恰乘三次航班有 $(\vv A^3)_{14}=1$ 种走法（$1\to2\to3\to4$）.}

\ans{ex:mat:A20}{$\vv x^\trans\vv A\vv x$ 是一个数，等于它的转置 $\vv x^\trans\vv A^\trans\vv x=-\vv x^\trans\vv A\vv x$，故为 $0$. $(\vv A^2)^\trans=(\vv A^\trans)^2=(-\vv A)^2=\vv A^2$.}

\ans{ex:mat:A21}{$\vv P$ 是初等矩阵之积，$\vv{PA}$ 与 $\vv A$ 行等价，因而有相同的简化阶梯形（简化阶梯形的唯一性见第一章章末习题~\ref{ex:elim:C9}）.}

\ans{ex:mat:A22}{存在性由提示中的等式. 唯一性：若 $\vv S_1+\vv K_1=\vv S_2+\vv K_2$（$\vv S_i$ 对称，$\vv K_i$ 反对称），则 $\vv S_1-\vv S_2=\vv K_2-\vv K_1$ 既对称又反对称，只能是零矩阵.}

\subsection*{B 组提示}

\ans{ex:mat:B1}{设 $\vv A=\begin{pmatrix}a&b\\c&d\end{pmatrix}$，用 Cayley–Hamilton 型恒等式 $\vv A^2=(a+d)\vv A-(ad-bc)\vv I$（直接验证）. (1) $\vv A^2=\vv O$ 当且仅当 $a+d=0$ 且 $ad-bc=0$. (2) $\vv A=\pm\vv I$，或 $a+d=0$ 且 $ad-bc=-1$. (3) $\vv A=\vv O$，$\vv A=\vv I$，或 $a+d=1$ 且 $ad-bc=0$.}
\ans{ex:mat:B2}{比较 $\vv{AD}$ 与 $\vv{DA}$ 的 $(i,j)$ 元：$a_ix_{ij}=x_{ij}a_j$.}
\ans{ex:mat:B3}{由 $\vv{AE}_{ij}=\vv E_{ij}\vv A$ 比较元素，得 $\vv A$ 的非对角元为零，且 $a_{ii}=a_{jj}$.}
\ans{ex:mat:B4}{用二项式定理（可交换时成立）：$(\vv A+\vv B)^{p+q}$ 的每一项中 $\vv A$ 的次数 $\ge p$ 或 $\vv B$ 的次数 $\ge q$. 反例：$\vv A=\vv E_{12}$，$\vv B=\vv E_{21}$.}
\ans{ex:mat:B5}{等式：分块对角矩阵的行（列）向量组的极大无关组由两块各自的极大无关组拼成. 不等式：用初等变换把 $\vv A$、$\vv B$ 分别化为标准形，看非零子矩阵. 例子：$\vv A=\vv B=(0)$，$\vv C=(1)$.}
\ans{ex:mat:B6}{把右边与 $\vv I+\vv u\vv v^\trans$ 相乘，利用 $\vv v^\trans\vv u$ 是一个数.}
\ans{ex:mat:B7}{直接验证 $(\vv I_n+\vv{BA})\big(\vv I_n-\vv B(\vv I_m+\vv{AB})^{-1}\vv A\big)=\vv I_n$，展开后把 $\vv B\vv A\vv B$ 写成 $\vv B(\vv{AB})$.}
\ans{ex:mat:B8}{各列都是某个非零列 $\vv u$ 的倍数，倍数组成 $\vv v$. $\vv A^2=\vv u(\vv v^\trans\vv u)\vv v^\trans$，而 $\vv v^\trans\vv u=\tr\vv A$. 于是 $\vv A^k=(\tr\vv A)^{k-1}\vv A$.}
\ans{ex:mat:B9}{用相抵标准形 $\vv A=\vv P\diag(\vv I_r,\vv O)\vv Q$，把 $\vv I_r$ 拆成 $r$ 个 $\vv E_{ii}$. 后半句：秩一矩阵之和的秩不超过个数.}
\ans{ex:mat:B10}{取 $\vv B$ 为 $\vv A$ 列向量组的一个极大无关组排成的矩阵，$\vv C$ 的各列是 $\vv A$ 各列在这组向量下的坐标. 对第 A8 题，$\vv B$ 取第 $1,3,5$ 列，$\vv C$ 是简化阶梯形的前三行.}
\ans{ex:mat:B11}{全 $1$ 矩阵 $\vv J$ 满足 $\vv J^2=3\vv J$，故 $\vv J^n=3^{n-1}\vv J$. 第二个：$\vv N^3=\vv O$，$(2\vv I+\vv N)^n=2^n\vv I+n2^{n-1}\vv N+\binom n22^{n-2}\vv N^2$.}
\ans{ex:mat:B12}{对 $k$ 归纳，计算 $\vv N^k$ 的 $(i,j)$ 元 $\sum_l(\vv N^{k-1})_{il}n_{lj}$ 中哪些项可能非零.}
\ans{ex:mat:B13}{$\tr(\vv A^\trans\vv A)=\sum_{i,j}a_{ij}^2$. 复矩阵不成立：$\vv A=\begin{pmatrix}1&0\\\mathrm i&0\end{pmatrix}$ 满足 $\vv A^\trans\vv A=\vv O$；复矩阵应改用 $\bar{\vv A}^\trans\vv A$.}
\ans{ex:mat:B14}{置换矩阵就是对单位矩阵的行作排列得到的矩阵. $\vv{PA}$ 对 $\vv A$ 的行作相应排列，$\vv{AP}$ 对列作排列.}
\ans{ex:mat:B15}{秩为 $2$. 先用行变换化为阶梯形（记录 $\vv P$），再用列变换（记录 $\vv Q$）. 答案不唯一，可用 $\vv{PAQ}$ 直接检验.}
\ans{ex:mat:B16}{$(\vv I+\vv A)(\vv I-\vv A+\vv A^2-\cdots+(-1)^{k-1}\vv A^{k-1})=\vv I$. 若 $\vv A^2=2\vv A$，则 $(\vv I-\vv A)^2=\vv I$，可逆且逆为自身.}

\section{第四章　线性映射}

\subsection*{A 组答案}

\ans{ex:lmap:A1}{(1) 线性，核为 $\{\zero\}$. (2) 不线性：$\zero\mapsto(1,0)^\trans$. (3) 不线性：$(-1,0)^\trans$ 与 $(1,0)^\trans$ 之和的像不等于像之和. (4) 线性，核为 $\{0\}$. (5) 线性，核为常数多项式全体. (6) 线性，核为 $\{\vv O\}$. (7) 不线性：$\det(2\vv X)=4\det\vv X$.}

\ans{ex:lmap:A2}{第一个：核的基 $(-2,1,0)^\trans$，$(1,0,1)^\trans$；像的基 $(1,2)^\trans$；$3=2+1$. 第二个：核的基 $(-2,1,1,0)^\trans$，$(-1,-1,0,1)^\trans$；像的基 $(1,0,1)^\trans$，$(0,1,1)^\trans$；$4=2+2$.}

\ans{ex:lmap:A3}{$\begin{pmatrix}1&-2&1\\0&3&4\end{pmatrix}$.}

\ans{ex:lmap:A4}{(1) $\begin{pmatrix}0&-1\\1&0\end{pmatrix}$. (2) $\begin{pmatrix}0&1\\1&0\end{pmatrix}$，平方为 $\vv I$：反射两次回到原处. (3) $\frac12\begin{pmatrix}1&1\\1&1\end{pmatrix}$，平方等于自身：投影两次与投影一次相同. (4) $\begin{pmatrix}1&2\\0&1\end{pmatrix}$.}

\ans{ex:lmap:A5}{两组基下的矩阵都是 $\begin{pmatrix}0&1&0\\0&0&2\\0&0&0\end{pmatrix}$：$\mathscr D(x-1)=1$，$\mathscr D(x-1)^2=2(x-1)$.}

\ans{ex:lmap:A6}{$\begin{pmatrix}1&0&0&0\\0&0&1&0\\0&1&0&0\\0&0&0&1\end{pmatrix}$，它的平方是 $\vv I_4$.}

\ans{ex:lmap:A7}{过渡矩阵 $\vv P=\begin{pmatrix}-1&-1&-4\\1&0&6\\1&3&1\end{pmatrix}$. $\vv\alpha$ 在 $\vv\varepsilon$ 下的坐标为 $(-1,2,2)^\trans$，在 $\vv\eta$ 下的坐标为 $(-16,5,3)^\trans$. 可用 $\vv x=\vv P\vv x'$ 检验.}

\ans{ex:lmap:A8}{$\diag(3,1)$. 过渡矩阵 $\vv P=\begin{pmatrix}1&1\\1&-1\end{pmatrix}$，也可以直接看出 $\mathscr T\vv\eta_1=3\vv\eta_1$，$\mathscr T\vv\eta_2=\vv\eta_2$.}

\ans{ex:lmap:A9}{$\mathscr S\mathscr T(a\vv\alpha+b\vv\beta)=\mathscr S(a\mathscr T\vv\alpha+b\mathscr T\vv\beta)=a\mathscr{ST}\vv\alpha+b\mathscr{ST}\vv\beta$. 逆映射：设 $\vv\alpha=\mathscr T\vv\alpha'$，$\vv\beta=\mathscr T\vv\beta'$，则 $\mathscr T^{-1}(a\vv\alpha+b\vv\beta)=\mathscr T^{-1}\mathscr T(a\vv\alpha'+b\vv\beta')=a\mathscr T^{-1}\vv\alpha+b\mathscr T^{-1}\vv\beta$.}

\ans{ex:lmap:A10}{若 $\sum k_i\mathscr T\vv\alpha_i=\zero$，则 $\sum k_i\vv\alpha_i\in\ker\mathscr T=\{\zero\}$. 反之，若 $\mathscr T\vv\alpha=\zero$ 而 $\vv\alpha\ne\zero$，则线性无关组 $\{\vv\alpha\}$ 被映为相关组 $\{\zero\}$.}

\ans{ex:lmap:A11}{(1) 存在：前两个向量是基，$\mathscr T(x,y)^\trans=\big(\frac{x+y}2,\frac{x-y}2\big)^\trans$，它恰好满足第三个条件. (2) 不存在：$(2,0)^\trans=(1,1)^\trans+(1,-1)^\trans$，线性性要求 $\mathscr T(2,0)^\trans=(1,1)^\trans$.}

\ans{ex:lmap:A12}{(1) $5-3=2$. (2) 都不存在：像的维数 $\le3<5$；核的维数 $\ge5-3=2>0$.}

\ans{ex:lmap:A13}{$\mathscr T\vv\alpha=\zero\Rightarrow\mathscr{ST}\vv\alpha=\zero$；$\mathscr{ST}\vv\alpha=\mathscr S(\mathscr T\vv\alpha)$. 由前者 $\rk(\mathscr{ST})=\dim U-\dim\ker(\mathscr{ST})\le\dim U-\dim\ker\mathscr T=\rk\mathscr T$，由后者 $\rk(\mathscr{ST})\le\rk\mathscr S$.}

\ans{ex:lmap:A14}{$\tr$ 是满射（$\tr(c\vv E_{11})=c$），由秩—零化度定理 $\dim\ker\tr=n^2-1$.}

\ans{ex:lmap:A15}{$\tr(\vv P^{-1}\vv{AP})=\tr(\vv{AP}\vv P^{-1})=\tr\vv A$. 两矩阵迹分别为 $4$ 与 $3$，不相似.}

\ans{ex:lmap:A16}{$\vv P^{-1}\vv I\vv P=\vv I$. 由于 $\begin{pmatrix}1&1\\0&1\end{pmatrix}\ne\vv I$，它与 $\vv I$ 不相似.}

\ans{ex:lmap:A17}{矩阵为 $\begin{pmatrix}1&0&0\\1&1&1\\1&2&4\end{pmatrix}$. 依次做 $r_2-r_1$，$r_3-r_1$，$r_3-2r_2$，得 $\begin{pmatrix}1&0&0\\0&1&1\\0&0&2\end{pmatrix}$，有 $3$ 个主元，由定理~\ref{thm:invertible} 矩阵可逆，故 $\mathscr E$ 是同构.（第五章会看到，它的行列式是 Vandermonde 型的.）}

\ans{ex:lmap:A18}{$\mathscr T(x_1,x_2,x_3)^\trans=(x_1+x_3,\ x_1+x_2+2x_3,\ x_2+x_3)^\trans$. $\ker\mathscr T=\spn\big((1,1,-1)^\trans\big)$，$\imm\mathscr T=\spn\big((1,1,0)^\trans,(0,1,1)^\trans\big)$.}

\ans{ex:lmap:A19}{$\vv A_2(\vv A_1\vv x+\vv b_1)+\vv b_2=(\vv A_2\vv A_1)\vv x+(\vv A_2\vv b_1+\vv b_2)$. 仿射变换线性当且仅当把 $\zero$ 映为 $\zero$，即 $\vv b=\zero$.}

\ans{ex:lmap:A20}{第一列是 $\vv e_1$ 转 $\theta$ 后的位置 $(\cos\theta,\sin\theta)^\trans$，第二列同理. 复合``先转 $\beta$ 再转 $\alpha$''的矩阵是 $\vv R(\alpha)\vv R(\beta)$（定理~\ref{thm:composition}），比较 $(1,1)$ 元与 $(2,1)$ 元，得 $\cos(\alpha+\beta)=\cos\alpha\cos\beta-\sin\alpha\sin\beta$，$\sin(\alpha+\beta)=\sin\alpha\cos\beta+\cos\alpha\sin\beta$.}

\subsection*{B 组提示}

\ans{ex:lmap:B1}{取定基后 $\mathscr T\mapsto$ 矩阵是 $L(V,W)$ 到 $M_{m\times n}(\mathbb{F})$ 的同构.}
\ans{ex:lmap:B2}{$\mathscr T^2=\mathscr O$ 即 $\mathscr T(\mathscr T\vv\alpha)=\zero$. 取 $\imm\mathscr T$ 的基 $\vv\beta_1,\ldots,\vv\beta_r$，$\vv\beta_i=\mathscr T\vv\alpha_i$，证明 $\vv\beta_1,\vv\alpha_1,\ldots,\vv\beta_r,\vv\alpha_r$ 线性无关，再用 $\ker\mathscr T$ 中的向量补成基.}
\ans{ex:lmap:B3}{对 $\mathscr T|_{\ker(\mathscr{ST})}$ 用秩—零化度定理：它的像含于 $\ker\mathscr S$，核含于 $\ker\mathscr T$.}
\ans{ex:lmap:B4}{$\mathscr{ST}=\mathscr I$ 说明 $\mathscr T$ 是单射，有限维时由推论~\ref{cor:nullity} 它是双射，于是 $\mathscr S=\mathscr T^{-1}$.}
\ans{ex:lmap:B5}{$\mathscr{JD}(1)=0$. 不矛盾：B4 依赖有限维，$\mathbb{R}[x]$ 是无限维的.}
\ans{ex:lmap:B6}{$\tr(\vv{AX}-\vv{XA})=0$. 后半句：$\vv A^k$ 与 $\vv A$ 可交换.}
\ans{ex:lmap:B7}{逐个计算 $\vv A\vv E_{11}$ 等并取坐标. $\vv X\mapsto\vv{AX}$ 的矩阵为 $\begin{pmatrix}\vv A&\vv O\\ \vv O&\vv A\end{pmatrix}$（在给定的基序下）；$\vv X\mapsto\vv{XA}$ 的矩阵为 $\begin{pmatrix}a\vv I&c\vv I\\ b\vv I&d\vv I\end{pmatrix}$.}
\ans{ex:lmap:B8}{证明 $V=\imm\mathscr P\oplus\ker\mathscr P$：$\vv\alpha=\mathscr P\vv\alpha+(\vv\alpha-\mathscr P\vv\alpha)$，且 $\mathscr P$ 在 $\imm\mathscr P$ 上是恒等. 合并两者的基.}
\ans{ex:lmap:B9}{按提示分解，并验证两部分分别在 $\ker(\mathscr T\mp\mathscr I)$ 中；交为零是因为 $\mathscr T\vv\alpha=\vv\alpha=-\vv\alpha$ 推出 $\vv\alpha=\zero$.}
\ans{ex:lmap:B10}{对 $\mathscr T|_U$ 用秩—零化度定理.}
\ans{ex:lmap:B11}{在基 $\mathrm e^x\cos x$，$\mathrm e^x\sin x$ 下 $\mathscr D$ 的矩阵为 $\begin{pmatrix}1&1\\-1&1\end{pmatrix}$，逆为 $\frac12\begin{pmatrix}1&-1\\1&1\end{pmatrix}$，故 $\int\mathrm e^x\sin x\,\mathrm dx=\frac12\mathrm e^x(\sin x-\cos x)+C$.}
\ans{ex:lmap:B12}{平移的矩阵为 $\begin{pmatrix}1&0&a\\0&1&b\\0&0&1\end{pmatrix}$. 两种顺序的矩阵不相等：先转后移把原点送到 $(1,2)$，先移后转把原点送到 $(-2,1)$.}
\ans{ex:lmap:B13}{$\vv P^{-1}\vv A^k\vv P=(\vv P^{-1}\vv{AP})^k$；转置时用 $\vv Q=(\vv P^{-1})^\trans$.}
\ans{ex:lmap:B14}{两者秩都为 $1$，故相抵. 若相似，则 $\vv P^{-1}\vv E_{11}\vv P$ 是幂等的，而 $\vv E_{12}^2=\vv O\ne\vv E_{12}$. 相抵允许在定义域与陪域中各自换基，相似必须在同一个空间中用同一组基.}
\ans{ex:lmap:B15}{设 $\sum c_i\mathscr T^i\vv\alpha=\zero$，依次施以 $\mathscr T^{k-1},\mathscr T^{k-2},\ldots$. 所求矩阵是主对角线之上一条斜线全为 $1$、其余为 $0$ 的矩阵.}

\section{第五章　行列式}

\subsection*{A 组答案}

\ans{ex:det:A1}{$\tau(315426)=5$，奇排列；$\tau(42531)=7$，奇排列；第三个排列的逆序数为 $0+1+2+\cdots+(n-1)=\frac{n(n-1)}2$，$n\equiv0,1\pmod4$ 时为偶排列，$n\equiv2,3\pmod4$ 时为奇排列.}

\ans{ex:det:A2}{(1) $\tau(3421)=5$，带负号. (2) 按行号排好后列号排列为 $i\,3\,5\,4\,j$，$\{i,j\}=\{1,2\}$. $\tau(13542)=4$，故 $i=1$，$j=2$.}

\ans{ex:det:A3}{$-41$，$-16$，$48$. 第三个：把各列加到第一列后提出 $6$，再化为三角形得 $6\cdot2^3=48$.}

\ans{ex:det:A4}{$6$.}

\ans{ex:det:A5}{$(-1)^{\frac{n(n-1)}2}a_1a_2\cdots a_n$：只有一项非零，它的列号排列是 $n(n-1)\cdots1$.}

\ans{ex:det:A6}{第一列减去第二列、再加上第三列（或按多重线性拆成 $8$ 个行列式，其中有两列相同的都为零），得 $2\begin{vmatrix}a&b&c\\a_1&b_1&c_1\\a_2&b_2&c_2\end{vmatrix}$.}

\ans{ex:det:A7}{第一个行列式等于 $(x-1)^2(x+2)$，根为 $x=1$（二重）与 $x=-2$；第二个是 Vandermonde 行列式 $(2-1)(x-1)(x-2)$，根为 $x=1,2$.}

\ans{ex:det:A8}{$24$，$-3$，$\frac13$，$9$，$9$. 最后一个：$\vv A^*=3\vv A^{-1}$，故 $3\vv A^{-1}-2\vv A^*=-3\vv A^{-1}$，行列式为 $(-3)^3\cdot\frac13=-9$.}

\ans{ex:det:A9}{(1) $(\det\vv A)^2=\det(\vv A^\trans\vv A)=1$. (2) $(\det\vv A)^2=1$. (3) 若 $\det\vv A\ne0$，则 $\vv A$ 可逆，由 $\vv A^2=\vv A$ 得 $\vv A=\vv I$，矛盾.}

\ans{ex:det:A10}{$\begin{pmatrix}d&-b\\-c&a\end{pmatrix}$，$ad-bc\ne0$ 时逆为 $\frac1{ad-bc}\begin{pmatrix}d&-b\\-c&a\end{pmatrix}$. 第二个矩阵行列式为 $1$，伴随矩阵就是逆矩阵：$\begin{pmatrix}-24&18&5\\20&-15&-4\\-5&4&1\end{pmatrix}$.}

\ans{ex:det:A11}{$D=17$，$D_1=17$，$D_2=-17$，$D_3=34$，解为 $(1,-1,2)$.}

\ans{ex:det:A12}{面积为 $\frac12\left|\det\begin{pmatrix}3&1\\1&4\end{pmatrix}\right|=\frac{11}2$. 三点共线：$(2,3)^\trans$ 与 $(6,9)^\trans$ 张成的平行四边形面积为 $0$.}

\ans{ex:det:A13}{平行六面体体积为 $7$，四面体体积为 $\frac76$.}

\ans{ex:det:A14}{系数行列式为 $(5-\lambda)(6-\lambda)(4-\lambda)-4(4-\lambda)-4(6-\lambda)=-(\lambda-2)(\lambda-5)(\lambda-8)$，故 $\lambda=2,5,8$.}

\ans{ex:det:A15}{把第一行换成 $(1,1,1,1)$ 后按第一行展开即得所求之和，计算得 $16$.}

\ans{ex:det:A16}{点 $(x,y)$ 在直线上当且仅当第一行是后两行的``仿射组合''，即三行线性相关. 也可以按第一行展开，得到一个关于 $x,y$ 的一次方程，并验证两点满足它.}

\ans{ex:det:A17}{按定义，非零项必须在左上角 $\vv A$ 中取 $m$ 个元素、在右下角 $\vv B$ 中取 $n$ 个，符号相乘. 第二个等于 $(-1)^{mn}\det\vv A\det\vv B$：把后 $n$ 列逐一移到前面，共 $mn$ 次相邻列交换.}

\ans{ex:det:A18}{$5$（一般地，这类 $n$ 阶行列式等于 $n+1$，见 B1）.}

\ans{ex:det:A19}{$\vv A\vv A^*=(\det\vv A)\vv I$ 两边取行列式得 $\det\vv A^*=(\det\vv A)^{n-1}\ne0$；再由 $\vv A^*=(\det\vv A)\vv A^{-1}$ 得 $(\vv A^*)^{-1}=\frac1{\det\vv A}\vv A$；对 $\vv A^{-1}$ 用同一公式得 $(\vv A^{-1})^*=\det(\vv A^{-1})\vv A=\frac1{\det\vv A}\vv A$.}

\ans{ex:det:A20}{行列式为 $0$，线性相关（$(7,8,9)^\trans=2(4,5,6)^\trans-(1,2,3)^\trans$）；改成 $10$ 后行列式为 $-3$，线性无关.}

\ans{ex:det:A21}{$\det\vv A=\det\vv A^\trans=\det(-\vv A)=(-1)^n\det\vv A=-\det\vv A$，故 $\det\vv A=0$.}

\subsection*{B 组提示}

\ans{ex:det:B1}{由 $D_n=2D_{n-1}-D_{n-2}$，$D_1=2$，$D_2=3$，得 $D_n-D_{n-1}=1$，$D_n=n+1$.}
\ans{ex:det:B2}{第一个：各列加到第一列，提出 $a+(n-1)b$，各行减去第一行，得 $[a+(n-1)b](a-b)^{n-1}$. 第二个：先证递推 $D_n=(x-a)D_{n-1}+a(x+a)^{n-1}$（或利用转置后 $a$ 变 $-a$），得 $\frac12\big[(x+a)^n+(x-a)^n\big]$.}
\ans{ex:det:B3}{第 $i$ 列乘 $\frac1{a_i}$ 后从第一列减去，化为上三角形：$a_1\cdots a_n\big(a_0-\sum_{i=1}^n\frac1{a_i}\big)$.}
\ans{ex:det:B4}{计算 $\begin{pmatrix}\vv I&\zero\\ -\vv v^\trans&1\end{pmatrix}\begin{pmatrix}\vv I+\vv u\vv v^\trans&\vv u\\ \zero^\trans&1\end{pmatrix}\begin{pmatrix}\vv I&\zero\\ \vv v^\trans&1\end{pmatrix}$. B2 的第一个行列式即 $\det\big((a-b)\vv I+b\vv 1\vv 1^\trans\big)$.}
\ans{ex:det:B5}{把 $\vv{AB}$ 的两列按第三章``列是线性组合''的看法展开，用多重线性拆成若干个行列式，其中列号重复的为零. $\vv{BA}$ 的秩 $\le2<3$.}
\ans{ex:det:B6}{必要性：$\det\vv A\det\vv A^{-1}=1$，两者都是整数. 充分性：$\vv A^{-1}=\frac1{\det\vv A}\vv A^*$，$\vv A^*$ 是整数矩阵.}
\ans{ex:det:B7}{可逆时用 $\vv A^*=(\det\vv A)\vv A^{-1}$ 两次. 不可逆时 $\rk\vv A^*\le1$，$n\ge3$ 时 $(\vv A^*)^*=\vv O$；$n=2$ 时直接验证.}
\ans{ex:det:B8}{对分块矩阵先把第二块行加到第一块行，再把第一块列从第二块列减去，化为 $\begin{pmatrix}\vv A+\vv B&\vv O\\ \vv B&\vv A-\vv B\end{pmatrix}$.}
\ans{ex:det:B9}{左乘 $\begin{pmatrix}\vv I&\vv O\\ -\vv{CA}^{-1}&\vv I\end{pmatrix}$. 可交换时 $\det\vv A\det(\vv D-\vv{CA}^{-1}\vv B)=\det(\vv{AD}-\vv{ACA}^{-1}\vv B)=\det(\vv{AD}-\vv{CB})$.}
\ans{ex:det:B10}{两边取行列式：$\det\vv F^n=(\det\vv F)^n=(-1)^n$.}
\ans{ex:det:B11}{按第一行展开得 $D_n=(1+x^2)D_{n-1}-x^2D_{n-2}$，即 $D_n-D_{n-1}=x^2(D_{n-1}-D_{n-2})$，故 $D_n=1+x^2+\cdots+x^{2n}$.}
\ans{ex:det:B12}{把第二、三列分别乘 $\omega^k$，$\omega^{2k}$ 加到第一列，提出公因式；也可直接展开验证第二个等号.}
\ans{ex:det:B13}{从最后一行起，每一行减去上一行，化为对角线全为 $1$ 的上三角形.}
\ans{ex:det:B14}{圆（或直线）的方程可写成 $a(x^2+y^2)+bx+cy+d=0$，系数不全为零. 四点都满足它，即一个以 $(a,b,c,d)$ 为未知量的齐次方程组有非零解.}
\ans{ex:det:B15}{$\rk\vv A=n$ 时 $\vv A^*=(\det\vv A)\vv A^{-1}$ 可逆. $\rk\vv A\le n-2$ 时，由定理~\ref{thm:minorrank} 所有 $n-1$ 阶子式为零，$\vv A^*=\vv O$. $\rk\vv A=n-1$ 时有非零的 $n-1$ 阶子式，故 $\vv A^*\ne\vv O$；又 $\vv A\vv A^*=(\det\vv A)\vv I=\vv O$，$\vv A^*$ 的每一列都属于 $\ker\vv A$，而 $\dim\ker\vv A=n-(n-1)=1$（推论~\ref{cor:nullity}），故 $\rk\vv A^*=1$.}
\ans{ex:det:B16}{$x=1,2,3$ 时有两行相同；每行之和都是 $x+6$，把各列加到第一列可提出因子 $x+6$. $f$ 是四次多项式，首项系数为 $1$，故 $f(x)=(x+6)(x-1)(x-2)(x-3)$，根为 $1,2,3,-6$.}

\section{第六章　线性方程组解的结构}

\subsection*{A 组答案}

\ans{ex:sol:A1}{(1) 基础解系 $(2,5,7,0)^\trans$，$(3,4,0,7)^\trans$，通解为二者的任意线性组合. (2) 基础解系 $(-2,1,0,0)^\trans$，$(1,0,0,1)^\trans$.}

\ans{ex:sol:A2}{(1) $\vv x=\big(\tfrac54,-\tfrac14,0,0\big)^\trans+k_1(3,3,2,0)^\trans+k_2(-3,7,0,4)^\trans$. (2) $\vv x=\big(\tfrac12,0,0,0\big)^\trans+k_1(-1,2,0,0)^\trans+k_2(1,0,2,0)^\trans$.}

\ans{ex:sol:A3}{(1) $5-3=2$ 个. (2) 不一定，要看 $\rk(\vv A\mid\vv b)$ 是否为 $3$. (3) 一定有解（行满秩，推论~\ref{cor:sol-fullrank}），且有无穷多解.}

\ans{ex:sol:A4}{(1) $\rk\vv A=\rk(\vv A\mid\vv b)=2$，有解（且唯一）. (2) $\rk\vv A=1$，$\rk(\vv A\mid\vv b)=2$，无解.}

\ans{ex:sol:A5}{系数和为 $1$ 的 $2\vv\eta_1-\vv\eta_2$，$\frac12(\vv\eta_1+\vv\eta_2)$，$3\vv\eta_1-2\vv\eta_2$ 是 $\vv{Ax}=\vv b$ 的解；系数和为 $0$ 的 $\vv\eta_1-\vv\eta_2$ 是齐次方程组的解；$\vv\eta_1+\vv\eta_2$ 两者都不是（命题~\ref{prop:sol-affine}）.}

\ans{ex:sol:A6}{基础解系含 $3-2=1$ 个向量，$\vv\eta_1-\vv\eta_2=(0,4,4)^\trans$ 是非零解. 通解为 $(1,2,3)^\trans+k(0,1,1)^\trans$.}

\ans{ex:sol:A7}{例如 $x_1-x_2=1$，$x_1+x_3=3$. 做法：求与 $(1,1,-1)^\trans$ 垂直（即满足 $y_1+y_2-y_3=0$）的两个线性无关的系数向量，再代入特解求常数项. 答案不唯一.}

\ans{ex:sol:A8}{$\rk\vv A=2$（第三个法向量是前两个之和，且任意两个不平行），$\rk(\vv A\mid\vv b)=3$：三个平面两两相交成三条平行直线，没有公共点. 常数改为 $3$ 后 $\rk(\vv A\mid\vv b)=2$，三个平面交于一条直线.}

\ans{ex:sol:A9}{例如 $\vv A=\begin{pmatrix}1\\0\end{pmatrix}$，$\vv b=(0,1)^\trans$：$\vv{Ax}=\zero$ 只有零解，$\vv{Ax}=\vv b$ 无解. 若有解，两个解之差是齐次方程组的解，只能为零.}

\ans{ex:sol:A10}{$\vv X=\begin{pmatrix}3&2\\-2&-3\\1&3\end{pmatrix}$.}

\ans{ex:sol:A11}{$\vv X=\begin{pmatrix}a&b\\-a&-b\end{pmatrix}$（$a,b$ 任意），构成二维空间：$\vv X$ 的每一列都属于 $\ker\vv A=\spn\big((1,-1)^\trans\big)$.}

\ans{ex:sol:A12}{$\vv B$ 的非零列是 $\vv{Ax}=\zero$ 的非零解，故 $\det\vv A=7t+21=0$，$t=-3$. 这时 $\rk\vv A=2$，由 $\rk\vv A+\rk\vv B\le3$ 且 $\vv B\ne\vv O$ 得 $\rk\vv B=1$.}

\ans{ex:sol:A13}{通解 $(1,0,0)^\trans+s(-1,1,0)^\trans+t(-1,0,1)^\trans$. 取 $\vv e_1,\vv e_2,\vv e_3$：任一解 $(x_1,x_2,x_3)^\trans=x_1\vv e_1+x_2\vv e_2+x_3\vv e_3$，而 $x_1+x_2+x_3=1$.}

\ans{ex:sol:A14}{第一条直线：$(0,0,1)^\trans+t(1,1,-2)^\trans$. 第二条直线：$x_1-x_3=1$，$x_2=2$（答案不唯一）.}

\ans{ex:sol:A15}{方程组 $1+s=0$，$s=1+t$，$0=1-t$ 无解，两直线不相交（它们的方向也不平行，是异面直线）.}

\ans{ex:sol:A16}{方程 $x_1+x_3=30$，$x_1-x_2=0$，$x_2+x_3=30$（秩为 $2$）. 全部方案为 $x_1=x_2=s$，$x_3=30-s$. $x_1+x_3=30$ 与 $x_1-x_2=0$ 被唯一确定（它们是方程左边的线性组合）；$x_3$ 不被确定：它在基础解系 $(1,1,-1)^\trans$ 上的值不为零.}

\ans{ex:sol:A17}{$\vv L=\begin{pmatrix}1&0&0\\2&1&0\\1&2&1\end{pmatrix}$，$\vv U=\begin{pmatrix}2&1&0\\0&1&1\\0&0&1\end{pmatrix}$.}

\ans{ex:sol:A18}{$\rk(\vv A\mid\vv b)$ 不超过行数 $m=\rk\vv A$，故相容；解空间维数 $n-m>0$，有无穷多解.}

\ans{ex:sol:A19}{两个方向 $(1,-1,0)^\trans$，$(-1,1,2)^\trans$；求与二者都垂直的系数向量，即解齐次方程组，得 $(1,1,0)^\trans$. 平面方程为 $x+y=2$.}

\ans{ex:sol:A20}{解分别为 $(2,0)$ 与 $(1,1)$. 右端的相对变化约为 $\frac{0.0001}{2\sqrt2}\approx3.5\times10^{-5}$，解的相对变化为 $\frac{\sqrt2}2\approx0.71$，二者之比约为 $2\times10^4$.}

\subsection*{B 组提示}

\ans{ex:sol:B1}{把四个方程联立求解，得公共解 $k(-1,1,2,1)^\trans$.}
\ans{ex:sol:B2}{$\vv A^\trans\vv{Ax}=\zero\Rightarrow\vv x^\trans\vv A^\trans\vv A\vv x=0$，即 $\vv{Ax}$ 各分量的平方和为零. 两个齐次方程组同解，解空间维数相同，再用推论~\ref{cor:nullity}.}
\ans{ex:sol:B3}{$\vv A+\vv B$ 的列是 $(\vv A\mid\vv B)$ 的列的线性组合；后者的列由两个极大无关组合并表出.}
\ans{ex:sol:B4}{$(\vv A+\vv I)(\vv A-\vv I)=\vv O$ 给出 $\le n$（Sylvester 不等式的推论）；$(\vv A+\vv I)-(\vv A-\vv I)=2\vv I$ 给出 $\ge n$（B3）.}
\ans{ex:sol:B5}{由定理~\ref{thm:sylvester} 证明中的等式 $\rk\vv B=\rk(\vv{AB})+\dim(\ker\vv A\cap\imm\vv B)$.}
\ans{ex:sol:B6}{设 $k_0\vv\eta^*+\sum k_i(\vv\eta^*+\vv\xi_i)=\zero$，左乘 $\vv A$ 得 $(k_0+\sum k_i)\vv b=\zero$，再用 $\vv\xi_i$ 线性无关.}
\ans{ex:sol:B7}{$9$ 个未知数，$6$ 个方程的秩为 $5$，自由度为 $4$，至少再公开 $4$ 个单元格. 公开左上角的 $2\times2$ 块即可；但若公开的 $4$ 格中有一行的全部 $3$ 格，那一行有一格是多余的，全表仍不能确定.}
\ans{ex:sol:B8}{三行满足 $r_1+r_2=r_3$，故条件为 $b_1+b_2=b_3$.}
\ans{ex:sol:B9}{解空间是一维的，$(1,\ldots,1)^\trans$ 是非零解. 通解为 $k(1,\ldots,1)^\trans$.}
\ans{ex:sol:B10}{$\rk\vv A+\rk\vv B\le n$（B4 中的推论）且 $\ge\rk(\vv A+\vv B)=n$（B3）.}
\ans{ex:sol:B11}{同解时，$\vv B$ 的每一行在 $\vv A$ 的解空间上恒为零，由可确定性判据（定理~\ref{thm:sol-determinable} 的 (2)$\Rightarrow$(3)）它是 $\vv A$ 各行的线性组合.}
\ans{ex:sol:B12}{必要性：对 $\vv b=\vv e_1,\ldots,\vv e_m$ 分别取解，排成 $\vv C$. 充分性：$\vv{A}(\vv C\vv b)=\vv b$.}
\ans{ex:sol:B13}{$\vv A\vv A^*=(\det\vv A)\vv I=\vv O$，$\vv A^*$ 的每一列都在一维解空间中.}
\ans{ex:sol:B14}{$x_1-x_4$ 是第一、二个方程之差，被确定为 $-3$；后两个量都不被确定：找不到系数 $a,b,c$ 使 $a,b,c$ 倍的三个方程之和的左边恰为所求的量.}
\ans{ex:sol:B15}{对 $\vv B|_{\ker(\vv{AB})}$ 用秩—零化度定理：它的核是 $\ker\vv B$，像含于 $\ker\vv A$.}

\section{第七章　多项式}

\subsection*{A 组答案}

\ans{ex:poly:A1}{(1) 商 $x^3-x^2+1$，余式 $-x-2$. (2) 商 $2x^3+x^2+6x+7$，余式 $20$（$=f(2)$）.}

\ans{ex:poly:A2}{依次为 $1$，$1$，$100x-99$. 前两个用 $x^2\equiv1$，$x^2\equiv-1$；第三个设余式为 $ax+b$，由 $x^{100}-(ax+b)$ 以 $1$ 为二重根，在 $x=1$ 处函数值与导数都为零.}

\ans{ex:poly:A3}{$\omega^4=\omega$，$\omega^2=-1-\omega$，代入得 $(1-a)\omega+(b-a)=0$，而 $\omega$ 不是实数，故 $a=b=1$. 验证：$x^4+x^2+1=(x^2+x+1)(x^2-x+1)$.}

\ans{ex:poly:A4}{(1) $x+3$. (2) $x-1$（$x^3-2x^2+x=x(x-1)^2$，$x^2-1=(x-1)(x+1)$）.}

\ans{ex:poly:A5}{$f$ 与 $g$ 的公因式整除 $g-hf$；$f$ 与 $g-hf$ 的公因式整除 $(g-hf)+hf=g$. 两组公因式相同，最大公因式也相同.}

\ans{ex:poly:A6}{$x^2+1=(x-1)(x+1)+2$，故 $u=\frac12$，$v=-\frac12(x-1)$.}

\ans{ex:poly:A7}{有理根只可能是 $\pm1,\pm2,\pm\frac12$. 得 $2x^3-5x^2+x+2=(x-1)(x-2)(2x+1)$，有理根为 $1,2,-\frac12$.}

\ans{ex:poly:A8}{(1) 在 $\mathbb{Q}$ 与 $\mathbb{R}$ 上为 $(x^2+x+1)(x^2-x+1)$；在 $\mathbb{C}$ 上为 $(x-\omega)(x-\bar\omega)(x+\omega)(x+\bar\omega)$，$\omega=\mathrm e^{2\pi\mathrm i/3}$. (2) 在 $\mathbb{Q}$ 与 $\mathbb{R}$ 上为 $(x-1)(x+1)(x^2+x+1)(x^2-x+1)$；在 $\mathbb{C}$ 上为 $\prod_{k=0}^5(x-\mathrm e^{k\pi\mathrm i/3})$. (3) 在 $\mathbb{Q}$ 上不可约；在 $\mathbb{R}$ 上为 $(x-\sqrt[3]2)(x^2+\sqrt[3]2x+\sqrt[3]4)$；在 $\mathbb{C}$ 上为 $(x-\sqrt[3]2)(x-\sqrt[3]2\,\omega)(x-\sqrt[3]2\,\omega^2)$.}

\ans{ex:poly:A9}{$f(1)=f'(1)=0$，$f''(1)=20\ne0$，是二重根. 也可以连续两次用 $x-1$ 作综合除法.}

\ans{ex:poly:A10}{有重根当且仅当与导数 $3x^2-3$ 有公共根，即 $f(1)=0$ 或 $f(-1)=0$. $a=2$ 时重根为 $1$；$a=-2$ 时重根为 $-1$.}

\ans{ex:poly:A11}{$\sigma_1=2$，$\sigma_2=3$，$\sigma_3=4$. 所求依次为 $2$；$\sigma_1^2-2\sigma_2=-2$；$\sigma_2/\sigma_3=\frac34$.}

\ans{ex:poly:A12}{$(x^2-2x-1)(x^2+1)=x^4-2x^3-2x-1$. （$1+\sqrt2$ 的共轭 $1-\sqrt2$、$\mathrm i$ 的共轭 $-\mathrm i$ 也必须是根，见第七章章末习题~\ref{ex:poly:B10}.）}

\ans{ex:poly:A13}{$(x-2)(x^2-2x+2)=x^3-4x^2+6x-4$.}

\ans{ex:poly:A14}{$f(\vv A)=(\vv A-\vv I)^2=\vv O$. 于是 $\vv A^2=2\vv A-\vv I$，$\vv A^3=3\vv A-2\vv I$，$g(\vv A)=2\vv A-2\vv I=\begin{pmatrix}0&4\\0&0\end{pmatrix}$.}

\ans{ex:poly:A15}{$f(x)=2x^2-3x+1$（四个点恰好落在一条抛物线上，插值多项式的三次项系数为 $0$）.}

\ans{ex:poly:A16}{$x^n-1=(x-1)(x^{n-1}+x^{n-2}+\cdots+1)$. 一般地 $x^n-a^n=(x-a)(x^{n-1}+ax^{n-2}+\cdots+a^{n-1})$；也可由余数定理：$x=a$ 时 $x^n-a^n=0$.}

\ans{ex:poly:A17}{$f(x)=2-2(x-1)+2(x-1)^3+(x-1)^4$.}

\ans{ex:poly:A18}{设余式为 $ax+b$，由 $a+b=3$，$2a+b=5$ 得余式为 $2x+1$.}

\ans{ex:poly:A19}{$60.5v^2+55v-100=0$，即 $121v^2+110v-200=0$，正根 $v=\frac{10}{11}$，内部收益率 $r=10\%$.}

\ans{ex:poly:A20}{比较次数得 $\deg f=\deg g$，$f=hg$ 中 $h$ 为常数，且 $f\ne0$ 时 $h\ne0$.}

\ans{ex:poly:A21}{设 $h$ 是 $f_1,g_1$ 的公因式，则 $dh$ 是 $f,g$ 的公因式，故 $dh\mid d$，$h$ 是常数.}

\subsection*{B 组提示}

\ans{ex:poly:B1}{由 Bézout 等式 $uf+vg=1$，$u'f+v'h=1$，两式相乘.}
\ans{ex:poly:B2}{用提示中的等式反复降低指数，指数的变化与整数的辗转相除 $\gcd(m,n)=\gcd(m-n,n)$ 完全平行.}
\ans{ex:poly:B3}{$f/\gcd(f,f')=(x-1)(x+2)$. 一般地，$p_i^{k_i}$ 在 $f'$ 中恰好以 $p_i^{k_i-1}$ 出现（需要用到 $p_i$ 与 $p_i'$ 互素）.}
\ans{ex:poly:B4}{$p$ 不可约，由引理~\ref{lem:primediv} 即得. 反例：$p=x^2$，$f=x$.}
\ans{ex:poly:B5}{两边都是次数小于 $n$ 的多项式，在 $x_1,\ldots,x_n$ 处取值相同（$\ell_i(x_j)$ 的性质），用推论~\ref{cor:rootcount}.}
\ans{ex:poly:B6}{$f(x)-f(0)$ 以 $0,1,2,\ldots$ 为根，有无穷多个根.}
\ans{ex:poly:B7}{三次多项式可约当且仅当有一次因式，即有有理根；可能的有理根只有 $\pm1$，都不是根. 四次多项式可能分解为两个二次因式之积而没有有理根，例如 $(x^2+1)^2$.}
\ans{ex:poly:B8}{设余式为 $ax+b$，代入 $x=1,2$：$(2^n-1)x+2-2^n$.}
\ans{ex:poly:B9}{由 $f(2)=4$ 得 $c=1$，$f(x)=x^2$.}
\ans{ex:poly:B10}{$(x-a)^2-2b^2$ 是有理系数的，用它除 $f$，余式 $rx+s$ 在 $a+b\sqrt2$ 处为零，由 $\sqrt2$ 无理得 $r=s=0$.}
\ans{ex:poly:B11}{对整数 $n$，$f(n)$ 与 $f(n\bmod2)$ 的奇偶性相同（$n-m\mid f(n)-f(m)$）.}
\ans{ex:poly:B12}{依次为 $x^2+1$，$x-2$，$(x-2)^2$.}
\ans{ex:poly:B13}{$N'(v)=C_1+2C_2v+\cdots>0$. 由 $N(0)=-C_0<0$ 与 $N(1)=\sum C_i-C_0>0$，用零点定理.}
\ans{ex:poly:B14}{$\binom xk$ 在整数点取整数值；把 $f$ 写成 $\sum c_k\binom xk$，依次代入 $x=0,1,\ldots,n$ 解出 $c_k$，它们都是整数.}
\ans{ex:poly:B15}{核中的多项式有 $n$ 个不同的根而次数小于 $n$，只能是零.}
\ans{ex:poly:B16}{$\gcd(p,f)$ 整除 $p$，只能是 $1$ 或 $p$ 的常数倍.}

\section{第八章　特征值与特征向量}

\subsection*{A 组答案}

\ans{ex:eig:A1}{以下``属于 $\lambda$''后给出特征子空间的一组基，特征向量是它们的非零线性组合.
\begin{enumerate}
	\item $\lambda=3$：$(1,0)^\trans$；$\lambda=2$：$(1,-1)^\trans$.
	\item $\lambda=3$：$(1,1)^\trans$；$\lambda=-1$：$(1,-1)^\trans$.
	\item $\lambda=\mathrm i$：$(1,\mathrm i)^\trans$；$\lambda=-\mathrm i$：$(1,-\mathrm i)^\trans$（在 $\mathbb{R}$ 上没有特征值）.
	\item $\lambda=2$：$(1,1)^\trans$；$\lambda=3$：$(2,1)^\trans$.
	\item $\lambda=2$（二重）：$(1,0,0)^\trans$，$(0,1,1)^\trans$；$\lambda=0$：$(0,1,-1)^\trans$.
	\item $\lambda=5$：$(1,1,1)^\trans$；$\lambda=-1$（二重）：$(1,-1,0)^\trans$，$(1,0,-1)^\trans$.
\end{enumerate}}

\ans{ex:eig:A2}{除第 3 个（在 $\mathbb{R}$ 上没有特征值）外都可对角化，$\vv P$ 取 A1 中各特征向量为列. 例如第 4 个：$\vv P=\begin{pmatrix}1&2\\1&1\end{pmatrix}$，$\vv P^{-1}\vv{AP}=\diag(2,3)$；第 6 个：$\vv P=\begin{pmatrix}1&1&1\\1&-1&0\\1&0&-1\end{pmatrix}$，$\vv P^{-1}\vv{AP}=\diag(5,-1,-1)$.}

\ans{ex:eig:A3}{第 1 个可以（两个不同的特征值）；第 2 个不可以（$2$ 的几何重数为 $1$）；第 3 个不可以（$2$ 是三重特征值，$2\vv I-\vv A$ 的秩为 $1$，几何重数为 $2$）；第 4 个可以（$1$ 的代数重数与几何重数都是 $2$）.}

\ans{ex:eig:A4}{$\det\vv A=6$，$\tr\vv A=6$；$\vv A^{-1}$ 的特征值为 $1,\frac12,\frac13$；$\vv A^2-2\vv A+\vv I=(\vv A-\vv I)^2$ 的特征值为 $0,1,4$；$\det\vv A^*=(\det\vv A)^2=36$；$\det(\vv A^2+\vv I)=2\cdot5\cdot10=100$.}

\ans{ex:eig:A5}{特征多项式为 $\lambda^2-5\lambda+6$，特征值 $2,3$；第二种情形为 $\lambda^2-2\lambda+5$，特征值 $1\pm2\mathrm i$.}

\ans{ex:eig:A6}{$\det(\lambda\vv I-\vv A^\trans)=\det\big((\lambda\vv I-\vv A)^\trans\big)$. 特征向量一般不同：$\vv A=\begin{pmatrix}1&1\\0&2\end{pmatrix}$ 属于 $1$ 的特征向量为 $(1,0)^\trans$，$\vv A^\trans$ 属于 $1$ 的特征向量为 $(1,-1)^\trans$.}

\ans{ex:eig:A7}{$\lambda\vv I-\vv A$ 仍是三角矩阵，行列式为 $\prod(\lambda-a_{ii})$.}

\ans{ex:eig:A8}{若 $\vv{Ax}=\lambda\vv x$，$\vv x\ne\zero$，则 $\vv A^k\vv x=\lambda^k\vv x$. 由 $\vv A^k=\vv O$ 得 $\lambda^k=0$；由 $\vv A^2=\vv A$ 得 $\lambda^2=\lambda$；由 $\vv A^2=\vv I$ 得 $\lambda^2=1$.}

\ans{ex:eig:A9}{$\vv A=\vv P\diag(2,3)\vv P^{-1}$，$\vv P=\begin{pmatrix}1&2\\1&1\end{pmatrix}$，
$\vv A^n=\begin{pmatrix}2\cdot3^n-2^n&2^{n+1}-2\cdot3^n\\3^n-2^n&2^{n+1}-3^n\end{pmatrix}$.}

\ans{ex:eig:A10}{特征值 $1$ 与 $0.7$. 稳态分布 $\vv\pi=(\frac23,\frac13)^\trans$. $\vv P^n\to\begin{pmatrix}\frac23&\frac23\\ \frac13&\frac13\end{pmatrix}$：两列相同，说明无论开始时顾客在哪一家，长期看都以同样的概率分布在两家，初始状态被``遗忘''了.}

\ans{ex:eig:A11}{$\vv A\vv\xi=(-1,a+2,b+1)^\trans=\lambda(1,1,-1)^\trans$，得 $\lambda=-1$，$a=-3$，$b=0$.}

\ans{ex:eig:A12}{特征多项式 $(\lambda-1)^2(\lambda+1)$. 可对角化当且仅当 $\rk(\vv I-\vv A)=1$，即 $x=-1$.}

\ans{ex:eig:A13}{$\vv A^2=\begin{pmatrix}a^2+bc&b(a+d)\\c(a+d)&d^2+bc\end{pmatrix}$，代入直接验证.}

\ans{ex:eig:A14}{依次为 $(x-1)(x-2)$，$(x-1)^2(x-2)$，$x^3$，$x(x-3)$.}

\ans{ex:eig:A15}{相似. 取 $\vv P=\begin{pmatrix}0&1\\1&0\end{pmatrix}$（交换两个基向量的次序），$\vv P^{-1}\begin{pmatrix}1&1\\0&1\end{pmatrix}\vv P=\begin{pmatrix}1&0\\1&1\end{pmatrix}$.}

\ans{ex:eig:A16}{$\vv A\vv e_1=a_{11}\vv e_1$，$\vv A\vv e_2=a_{12}\vv e_1+a_{22}\vv e_2$，……，$\vv A\vv e_k$ 只含 $\vv e_1,\ldots,\vv e_k$.}

\ans{ex:eig:A17}{特征根 $2,-1$，$a_k=\frac13\big(2^k-(-1)^k\big)$.}

\ans{ex:eig:A18}{矩阵 $\begin{pmatrix}0&1.1\\0.6&0.5\end{pmatrix}$ 的特征值为 $1.1$ 与 $-0.6$. 长期年增长率为 $10\%$；属于 $1.1$ 的特征向量为 $(1,1)^\trans$，幼体与成年个体之比稳定在 $1:1$.}

\ans{ex:eig:A19}{特征多项式 $\lambda^2-(\tr\vv A)\lambda+\det\vv A$ 没有实根，当且仅当判别式 $(\tr\vv A)^2-4\det\vv A<0$.}

\ans{ex:eig:A20}{投影：方向 $(1,2)^\trans$ 上的向量不动（特征值 $1$），垂直方向 $(2,-1)^\trans$ 上的向量变为零（特征值 $0$），矩阵 $\frac15\begin{pmatrix}1&2\\2&4\end{pmatrix}$. 反射：$(1,2)^\trans$ 属于 $1$，$(2,-1)^\trans$ 属于 $-1$，矩阵 $\frac15\begin{pmatrix}-3&4\\4&3\end{pmatrix}$.}

\subsection*{B 组提示}

\ans{ex:eig:B1}{$\vv A\vv u=(\vv v^\trans\vv u)\vv u$；$\rk\vv A=1$，$0$ 的几何重数为 $n-1$. 特征值为 $\vv v^\trans\vv u$ 与 $0$（$n-1$ 重）. $\vv v^\trans\vv u=0$ 时 $0$ 是 $n$ 重特征值，但几何重数只有 $n-1$.}
\ans{ex:eig:B2}{$\vv A\vv B\vv p=\vv B\vv A\vv p=\lambda\vv B\vv p$，而 $V_\lambda$ 是一维的. 后半句：用 Lagrange 插值构造多项式 $f$ 使 $f(\lambda_i)$ 等于 $\vv B$ 的对应特征值.}
\ans{ex:eig:B3}{$\vv P^{-1}f(\vv A)\vv P=f(\vv D)$. 逆命题不成立：$\vv A=\begin{pmatrix}0&1\\0&0\end{pmatrix}$ 不可对角化，$\vv A^2=\vv O$ 可以.}
\ans{ex:eig:B4}{$\vv A|_W$ 的最小多项式整除 $\vv A$ 的最小多项式，用定理~\ref{thm:minpolydiag}.}
\ans{ex:eig:B5}{取 $\vv v\ne\zero$，证明 $\vv v,\vv{Av}$ 线性无关，在这组基下 $\vv A$ 的矩阵为 $\begin{pmatrix}0&-1\\1&0\end{pmatrix}$.}
\ans{ex:eig:B6}{设 $\vv{Ax}=\lambda\vv x$，取 $|x_i|$ 最大的分量，看第 $i$ 个方程 $(\lambda-a_{ii})x_i=\sum_{j\ne i}a_{ij}x_j$.}
\ans{ex:eig:B7}{$\vv A(1,\ldots,1)^\trans=s(1,\ldots,1)^\trans$. 可逆时 $s\ne0$，$\vv A^{-1}$ 每行之和为 $\frac1s$.}
\ans{ex:eig:B8}{特征值为 $a+b\omega^k+c\omega^{2k}$（$k=0,1,2$），行列式是三者之积.}
\ans{ex:eig:B9}{特征根 $2$ 是二重的，通项形如 $(c_1+c_2k)2^k$，得 $a_k=(k+1)2^k$. $k2^k$ 来自 Jordan 块的幂（第八章选读节）.}
\ans{ex:eig:B10}{对三角化（或可对角化）的情形看对角元. 特征值为 $0,\sqrt3,-\sqrt3$.}
\ans{ex:eig:B11}{$x^3-x=x(x-1)(x+1)$，最小多项式整除它.}
\ans{ex:eig:B12}{$2$ 是二重特征值，需 $\rk(2\vv I-\vv A)=1$，条件为 $a=0$（$b,c$ 任意）.}
\ans{ex:eig:B13}{转移矩阵 $\begin{pmatrix}0&0&1\\ \frac12&0&0\\ \frac12&1&0\end{pmatrix}$，稳态分布 $\frac15(2,1,2)^\trans$.}
\ans{ex:eig:B14}{特征值 $\pm1$，$\mathrm e^{t\vv A}=\begin{pmatrix}\cosh t&\sinh t\\ \sinh t&\cosh t\end{pmatrix}$. 对 $t$ 求导验证.}
\ans{ex:eig:B15}{稳态分布 $(\frac13,\frac13,\frac13)^\trans$；特征值 $1,\frac14,\frac14$，偏离部分每一步缩小为原来的 $\frac14$.}
\ans{ex:eig:B16}{几何重数 $n-\rk(\lambda\vv I-\vv A)$ 不超过代数重数 $k$.}
\ans{ex:eig:B17}{得分（按判决 $1,2,3,4$）约为：$d=0.5$ 时 $(0.169,0.212,0.265,0.354)$；$d=0.85$ 时 $(0.128,0.183,0.261,0.428)$；$d=0.95$ 时 $(0.119,0.175,0.258,0.448)$；$d=1$ 时恰为 $\frac1{35}(4,6,9,16)$. 排序不变；$d$ 越大，随机跳转越少，得分越集中到引证链末端的判决 $4$ 上.}

\subsection*{C 组证明要点（正文引用的题目）}

\ans{ex:eig:C1}{直接计算两个分块乘积：
	\[\begin{pmatrix}\vv I_m&-\vv A\\ \vv O&\lambda\vv I_n\end{pmatrix}\begin{pmatrix}\lambda\vv I_m&\vv A\\ \vv B&\vv I_n\end{pmatrix}=\begin{pmatrix}\lambda\vv I_m-\vv{AB}&\vv O\\ \lambda\vv B&\lambda\vv I_n\end{pmatrix},\]\[
	\begin{pmatrix}\vv I_m&\vv O\\ -\vv B&\lambda\vv I_n\end{pmatrix}\begin{pmatrix}\lambda\vv I_m&\vv A\\ \vv B&\vv I_n\end{pmatrix}=\begin{pmatrix}\lambda\vv I_m&\vv A\\ \vv O&\lambda\vv I_n-\vv{BA}\end{pmatrix}.\]
	记中间矩阵的行列式为 $D$. 对两式取行列式（第五章章末习题~\ref{ex:det:A17}），得 $\lambda^nD=\lambda^n\det(\lambda\vv I_m-\vv{AB})$ 与 $\lambda^nD=\lambda^m\det(\lambda\vv I_n-\vv{BA})$，因此 $\lambda^n\det(\lambda\vv I_m-\vv{AB})=\lambda^m\det(\lambda\vv I_n-\vv{BA})$. 两边作为 $\lambda$ 的多项式相等，因子 $\lambda^n$、$\lambda^m$ 只影响根 $0$ 的重数，所以两个特征多项式的非零根（连同重数）相同.}

\section{第九章　内积空间}

\subsection*{A 组答案}

\ans{ex:ip:A1}{$\|\vv\alpha\|=\|\vv\beta\|=3$，$\langle\vv\alpha,\vv\beta\rangle=4$，夹角 $\arccos\frac49\approx63.6^\circ$，$\|\vv\alpha-\vv\beta\|=\sqrt{10}$.}

\ans{ex:ip:A2}{(1) 是. (2) 不是：$\vv x=(1,-1)^\trans$ 时 $\langle\vv x,\vv x\rangle=-2$. (3) 是：$\langle\vv x,\vv x\rangle=x_1^2+(x_1+x_2)^2$. (4) 不是：$\langle\vv e_2,\vv e_2\rangle=-1$.}

\ans{ex:ip:A3}{(1) $\frac1{\sqrt3}(1,1,1)^\trans$，$\frac1{\sqrt2}(0,1,-1)^\trans$，$\frac1{\sqrt6}(-2,1,1)^\trans$. (2) $\frac1{\sqrt2}(1,1,0,0)^\trans$，$\frac1{\sqrt6}(1,-1,2,0)^\trans$，$\frac1{2\sqrt3}(-1,1,1,3)^\trans$.}

\ans{ex:ip:A4}{解 $x_1+x_2+x_3=0$，$x_1-x_3=0$，得 $\pm\frac1{\sqrt6}(1,-2,1)^\trans$.}

\ans{ex:ip:A5}{逐对验证内积；坐标为 $\langle\vv\alpha,\vv e_i\rangle$，即 $(5,4,-2)^\trans$.}

\ans{ex:ip:A6}{两个生成向量已经正交，投影为 $\frac32(1,1,0,0)^\trans+\frac72(0,0,1,1)^\trans=\big(\frac32,\frac32,\frac72,\frac72\big)^\trans$，
	
	距离为 $$\|(-\frac12,\frac12,-\frac12,\frac12)^\trans\|=1$$.}

\ans{ex:ip:A7}{正规方程 $5a+10b=15$，$10a+30b=42$，得 $y=0.6+1.2x$，残差平方和为 $1.6$.}

\ans{ex:ip:A8}{第二列须与第一、三列正交且为单位向量：$a=c$，$b=a$，$3a^2=1$. 故 $a=b=c=\frac1{\sqrt3}$ 或 $a=b=c=-\frac1{\sqrt3}$.}

\ans{ex:ip:A9}{展开 $\big\langle\sum\vv\alpha_i,\sum\vv\alpha_j\big\rangle$，交叉项都为零.}

\ans{ex:ip:A10}{由三角不等式 $\|\vv\alpha\|\le\|\vv\alpha-\vv\beta\|+\|\vv\beta\|$，再交换 $\vv\alpha,\vv\beta$.}

\ans{ex:ip:A11}{(1) 对 $(a_1,\ldots,a_n)$ 与 $(1,\ldots,1)$ 用 Cauchy–Schwarz 不等式. (2) 对 $(\sqrt{a_1},\ldots,\sqrt{a_n})$ 与 $\big(\frac1{\sqrt{a_1}},\ldots,\frac1{\sqrt{a_n}}\big)$ 用之.}

\ans{ex:ip:A12}{$\|1\|=1$，$\|x\|=\frac1{\sqrt3}$，$\langle1,x\rangle=\frac12$，夹角 $\frac\pi6$. 正交化并单位化得 $1$，$\sqrt3(2x-1)$.}

\ans{ex:ip:A13}{$\ker\vv A$ 的标准正交基：$\frac1{\sqrt3}(1,-1,1,0)^\trans$，$\frac1{\sqrt{15}}(-1,-2,-1,3)^\trans$；行空间的标准正交基：$\frac1{\sqrt3}(1,1,0,1)^\trans$，$\frac1{\sqrt{15}}(-2,1,3,1)^\trans$. 两组之间两两内积为零，维数 $2+2=4$，故互为正交补.}

\ans{ex:ip:A14}{$\vv Q=\begin{pmatrix}\frac1{\sqrt2}&\frac1{\sqrt6}\\ \frac1{\sqrt2}&-\frac1{\sqrt6}\\0&\frac2{\sqrt6}\end{pmatrix}$，$\vv R=\begin{pmatrix}\sqrt2&\frac1{\sqrt2}\\0&\sqrt{\frac32}\end{pmatrix}$.}

\ans{ex:ip:A15}{$\vv H=\frac19\begin{pmatrix}7&-4&-4\\-4&1&-8\\-4&-8&1\end{pmatrix}$；$\vv H\vv u=-\vv u$；$(2,-1,0)^\trans$ 与 $\vv u$ 正交，$\vv H(2,-1,0)^\trans=(2,-1,0)^\trans$.}

\ans{ex:ip:A16}{第一个行列式为 $1$，是旋转，旋转角 $\arctan\frac43\approx53.1^\circ$；第二个行列式为 $-1$，是关于直线 $y=\frac x2$（方向 $(2,1)^\trans$）的反射.}

\ans{ex:ip:A17}{设计矩阵只有一列 $\vv x$，正规方程为 $(\vv x^\trans\vv x)b=\vv x^\trans\vv y$.}

\ans{ex:ip:A18}{投影为 $2(1,1,1)^\trans$，距离为 $\|(-1,0,1)^\trans\|=\sqrt2$.}

\ans{ex:ip:A19}{过渡矩阵的第 $j$ 列是新基第 $j$ 个向量在旧（标准正交）基下的坐标，坐标的标准内积等于向量的内积，故各列标准正交.}

\ans{ex:ip:A20}{内积打分：$6,1,4$，排序为甲（第一条）$>$ 丙 $>$ 乙. 余弦相似度：约 $0.926$，$0.140$，$0.943$，排序为丙 $>$ 甲 $>$ 乙. 第一条法条更长，内积占了便宜.}

\ans{ex:ip:A21}{$\vv P^\trans=(\vv u^\trans)^\trans\vv u^\trans=\vv P$，$\vv P^2=\vv u(\vv u^\trans\vv u)\vv u^\trans=\vv P$. $\vv{Px}=\langle\vv x,\vv u\rangle\vv u$ 正是定理~\ref{thm:projformula} 中 $k=1$ 的公式. 一般非零 $\vv u$ 时投影矩阵为 $\dfrac{\vv u\vv u^\trans}{\vv u^\trans\vv u}$.}

\subsection*{B 组提示}

\ans{ex:ip:B1}{展开两个平方相减. 取 $\vv x=(1,0)^\trans$，$\vv y=(0,1)^\trans$ 检验平行四边形法则 $\|\vv x+\vv y\|^2+\|\vv x-\vv y\|^2=2\|\vv x\|^2+2\|\vv y\|^2$，左边为 $8$，右边为 $4$.}
\ans{ex:ip:B2}{对 $\vv\alpha-\sum\langle\vv\alpha,\vv e_i\rangle\vv e_i$ 用勾股定理.}
\ans{ex:ip:B3}{证明 $\vv x-\vv{Px}$ 与 $\imm\vv P$ 正交：$\langle\vv x-\vv{Px},\vv{Py}\rangle=\vv y^\trans\vv P^\trans(\vv I-\vv P)\vv x=0$. 只有幂等时是沿 $\ker\vv P$ 向 $\imm\vv P$ 的``斜投影''，例如 $\begin{pmatrix}1&1\\0&0\end{pmatrix}$.}
\ans{ex:ip:B4}{目标函数是 $\|\vv y-\vv A\vv\beta\|_W^2$，其中 $\langle\vv x,\vv y\rangle_W=\vv x^\trans\vv W\vv y$ 是加权内积；残差在这个内积下与各列正交.}
\ans{ex:ip:B5}{对 $\ln y=\ln c+kt$ 作最小二乘，约得 $k\approx0.30$，$c\approx2.0$. 它最小化的是对数残差的平方和，相当于对相对误差加权.}
\ans{ex:ip:B6}{$\vv G=\vv A^\trans\vv A$（$\vv A$ 以 $\vv\alpha_i$ 为列），用第十一章命题~\ref{prop:gram} 或直接 $\vv x^\trans\vv G\vv x=\|\vv A\vv x\|^2$. $k=2$ 时 $\det\vv G=\|\vv\alpha\|^2\|\vv\beta\|^2\sin^2\theta$.}
\ans{ex:ip:B7}{距离是 $\vv x_0-\vv x_1$（$\vv x_1$ 为超平面上任一点）在法向 $\vv a$ 上投影的长度.}
\ans{ex:ip:B8}{设 $\vv{Qz}=\lambda\vv z$，考察 $\bar{\vv z}^\trans\vv Q^\trans\vv Q\vv z$.}
\ans{ex:ip:B9}{用积化和差逐对计算积分. 投影为 $2\sin x-\sin2x$（偶函数的系数为零）.}
\ans{ex:ip:B10}{全部最小二乘解为 $\hat{\vv x}+\ker\vv A$，把 $\hat{\vv x}$ 分解到 $\imm\vv A^\trans\oplus\ker\vv A$，用勾股定理.}
\ans{ex:ip:B11}{$\|\vv x-\vv y\|^2=2\|\vv x\|^2-2\langle\vv x,\vv y\rangle$，直接计算 $\vv H\vv x=\vv x-2\vv u\vv u^\trans\vv x$.}
\ans{ex:ip:B12}{对中心化的数据用 A17，再代回截距；正规方程的第一个方程就是 $\bar y=a+b\bar x$.}
\ans{ex:ip:B13}{保持夹角就保持正交. 把一组标准正交基的像归一化，证明 $\|\mathscr T\vv e_i\|$ 都相等（考察 $\vv e_i\pm\vv e_j$ 的像是否正交）.}
\ans{ex:ip:B14}{截距为第一组均值 $4$，虚拟变量系数为两组均值之差 $8-4=4$.}
\ans{ex:ip:B15}{$\vv R(\alpha)\vv R(\beta)=\vv R(\alpha+\beta)$；$\vv F(\alpha)\vv F(\beta)=\vv R(\alpha-\beta)$；计算 $\vv R(\alpha)\vv F(\beta)$ 与 $\vv F(\beta)\vv R(\alpha)$ 比较.}

\section{第十章　矩阵分解}

\subsection*{A 组答案}

\ans{ex:dec:A1}{（$\vv Q$ 的列依次对应所列特征值，答案不唯一.）
\begin{enumerate}
	\item $\vv Q=\frac1{\sqrt2}\begin{pmatrix}1&1\\1&-1\end{pmatrix}$，对角矩阵 $\diag(3,-1)$.
	\item $\vv Q=\frac1{\sqrt2}\begin{pmatrix}1&1\\1&-1\end{pmatrix}$，对角矩阵 $\diag(2,4)$.
	\item 特征值 $4,0,1$，特征向量 $\frac1{\sqrt2}(1,1,0)^\trans$，$\frac1{\sqrt2}(1,-1,0)^\trans$，$(0,0,1)^\trans$.
	\item 特征值 $3,0,0$，特征向量 $\frac1{\sqrt3}(1,1,1)^\trans$，$\frac1{\sqrt2}(1,-1,0)^\trans$，$\frac1{\sqrt6}(1,1,-2)^\trans$（后两个须在特征子空间内正交化）.
\end{enumerate}}

\ans{ex:dec:A2}{$\begin{pmatrix}3&-1\\-1&3\end{pmatrix}=2\cdot\frac12\begin{pmatrix}1&1\\1&1\end{pmatrix}+4\cdot\frac12\begin{pmatrix}1&-1\\-1&1\end{pmatrix}$. 两个投影矩阵之和为 $\vv I$，乘积为 $\vv O$.}

\ans{ex:dec:A3}{$\vv A=1\cdot\frac12\begin{pmatrix}1&1\\1&1\end{pmatrix}+2\cdot\frac12\begin{pmatrix}1&-1\\-1&1\end{pmatrix}=\frac12\begin{pmatrix}3&-1\\-1&3\end{pmatrix}$.}

\ans{ex:dec:A4}{属于 $1$ 的特征子空间是 $(0,1,1)^\trans$ 的正交补，故 $\vv A=\vv I-2\vv P$，$\vv P$ 是向 $\spn\big((0,1,1)^\trans\big)$ 的投影：$\vv A=\begin{pmatrix}1&0&0\\0&0&-1\\0&-1&0\end{pmatrix}$.}

\ans{ex:dec:A5}{特征值 $4,1,0$. 最大值 $4$，在 $\pm\frac1{\sqrt2}(1,1,0)^\trans$ 处取到；最小值 $0$，在 $\pm\frac1{\sqrt2}(1,-1,0)^\trans$ 处取到.}

\ans{ex:dec:A6}{矩阵为 $\begin{pmatrix}1&2\\2&1\end{pmatrix}$，最大值 $3$（在 $\pm\frac1{\sqrt2}(1,1)$ 处），最小值 $-1$（在 $\pm\frac1{\sqrt2}(1,-1)$ 处）.}

\ans{ex:dec:A7}{(1) $\sigma_1=3$，$\sigma_2=2$；$\vv U=\begin{pmatrix}0&1\\-1&0\end{pmatrix}$，$\vv\Sigma=\diag(3,2)$，$\vv V=\begin{pmatrix}0&1\\1&0\end{pmatrix}$. (2) $\vv U=\vv V=\frac1{\sqrt2}\begin{pmatrix}1&1\\1&-1\end{pmatrix}$，$\vv\Sigma=\diag(2,0)$. (3) $\vv U=(1)$，$\vv\Sigma=(5,\ 0)$，$\vv V=\frac15\begin{pmatrix}3&-4\\4&3\end{pmatrix}$.}

\ans{ex:dec:A8}{$\vv A^\trans\vv A=\begin{pmatrix}1&1\\1&2\end{pmatrix}$，特征值 $\frac{3\pm\sqrt5}2$，奇异值为 $\frac{\sqrt5+1}2$ 与 $\frac{\sqrt5-1}2$（黄金分割比及其倒数）；特征值是 $1,1$.}

\ans{ex:dec:A9}{对称正定矩阵的奇异值分解就是谱分解. $\vv A_1=4\cdot\frac12\begin{pmatrix}1&1\\1&1\end{pmatrix}=\begin{pmatrix}2&2\\2&2\end{pmatrix}$，误差 $\|\vv A-\vv A_1\|_F=\sigma_2=2$.}

\ans{ex:dec:A10}{$\vv A=9\vv P_1+\vv P_2$，$\vv P_1=\frac12\begin{pmatrix}1&1\\1&1\end{pmatrix}$，$\vv P_2=\frac12\begin{pmatrix}1&-1\\-1&1\end{pmatrix}$. $\vv A^{10}=9^{10}\vv P_1+\vv P_2$；$\vv B=3\vv P_1+\vv P_2=\begin{pmatrix}2&1\\1&2\end{pmatrix}$.}

\ans{ex:dec:A11}{$\vv A=\vv Q\vv O\vv Q^\trans=\vv O$. 反例：$\begin{pmatrix}0&1\\0&0\end{pmatrix}$.}

\ans{ex:dec:A12}{$\vv A^2=\vv A^\trans\vv A=\vv I$，特征值满足 $\lambda^2=1$. 它是关于特征子空间 $V_1$ 的``镜面反射''：$V_1$ 中的向量不动，$V_{-1}=V_1^\perp$ 中的向量反向.}

\ans{ex:dec:A13}{由第八章章末习题~\ref{ex:eig:C1}（取 $\vv B=\vv A^\trans$），或直接：若 $\vv A^\trans\vv A\vv x=\lambda\vv x$，$\lambda\ne0$，则 $\vv{AA}^\trans(\vv{Ax})=\lambda\vv{Ax}$ 且 $\vv{Ax}\ne\zero$；两边的特征子空间维数相同.}

\ans{ex:dec:A14}{$\tr(\vv A^\trans\vv A)=\sum_{i,j}a_{ij}^2$；它也等于 $\vv A^\trans\vv A$ 的特征值之和 $\sum\sigma_i^2$.}

\ans{ex:dec:A15}{$\det\vv A=\det\vv U\det\vv\Sigma\det\vv V^\trans=\pm\sigma_1\cdots\sigma_n$；$\vv A^{-1}=\vv V\vv\Sigma^{-1}\vv U^\trans$.}

\ans{ex:dec:A16}{只有第二个可以：能被正交矩阵对角化的实矩阵必须对称（谱定理的逆命题）.}

\ans{ex:dec:A17}{$\vv S=\frac13\begin{pmatrix}10&8\\8&10\end{pmatrix}$，特征值 $6$ 与 $\frac23$. 第一主成分方向 $\frac1{\sqrt2}(1,1)^\trans$，解释方差比例 $\frac6{6+2/3}=90\%$.}

\ans{ex:dec:A18}{$\kappa\big(\diag(10,1)\big)=10$；$\begin{pmatrix}1&1\\1&1.01\end{pmatrix}$ 的特征值约为 $2.005$ 与 $0.004988$，$\kappa\approx402$. 二者行列式都等于 $0.01$ 的 $\diag(10,0.001)$ 条件数为 $10^4$，与后者不同：行列式不决定条件数.}

\ans{ex:dec:A19}{$40\times(800+600+1)=56040$ 个数，原图 $480000$ 个，约为原来的 $11.7\%$. $k\ge343$ 时存储量超过原图.}

\ans{ex:dec:A20}{由谱定理 $\vv A$ 相似于 $\diag(\lambda_1,\ldots,\lambda_n)$，秩不变. 奇异值的平方是 $\vv A^\trans\vv A$ 的非零特征值，个数为 $\rk(\vv A^\trans\vv A)=\rk\vv A$.}

\ans{ex:dec:A21}{取正交矩阵 $\vv Q$ 使 $\vv Q^\trans\vv{AQ}=\diag(\lambda_1,\ldots,\lambda_n)$，令 $\vv y=\vv Q^\trans\vv x$，则 $\vv x^\trans\vv A\vv x=\sum\lambda_iy_i^2$，且 $\vv x\ne\zero\iff\vv y\ne\zero$. 若某个 $\lambda_i\le0$，取 $\vv y=\vv e_i$；若全为正，则 $\vv y\ne\zero$ 时和为正.}

\subsection*{B 组提示}

\ans{ex:dec:B1}{$\vv A$ 的每个特征子空间都是 $\vv B$-不变的，$\vv B$ 限制在其上仍是对称的，在其中取 $\vv B$ 的标准正交特征向量.}
\ans{ex:dec:B2}{存在性：$\vv B=\vv Q\diag(\sqrt{\lambda_i})\vv Q^\trans$. 唯一性：$\vv B$ 的每个特征子空间（属于 $\mu$）都含于 $\vv A$ 的属于 $\mu^2$ 的特征子空间，而 $\mu\ge0$ 由 $\mu^2$ 唯一确定.}
\ans{ex:dec:B3}{在 $\vv q_1^\perp=\spn(\vv q_2,\ldots,\vv q_n)$ 中重复定理~\ref{thm:rayleigh} 的证明.}
\ans{ex:dec:B4}{$\|\vv{Ax}\|^2=\vv x^\trans\vv A^\trans\vv A\vv x\le\lambda_{\max}(\vv A^\trans\vv A)\|\vv x\|^2$，在 $\vv v_1$ 处取等.}
\ans{ex:dec:B5}{$\vv{Ax}=\lambda\vv x$ 时 $|\lambda|\,\|\vv x\|=\|\vv{Ax}\|$，再用 B4 及其``最小值''版本.}
\ans{ex:dec:B6}{$\vv A^\trans\vv A$ 的非零特征值为 $3,1$（$\vv A\vv A^\trans=\begin{pmatrix}2&1\\1&2\end{pmatrix}$ 更易算），$\sigma_1=\sqrt3$，$\sigma_2=1$. 先求 $\vv U$ 的列：$\frac1{\sqrt2}(1,1)^\trans$，$\frac1{\sqrt2}(1,-1)^\trans$，再用 $\vv v_i=\vv A^\trans\vv u_i/\sigma_i$，第三个 $\vv v_3=\frac1{\sqrt3}(1,1,-1)^\trans$.}
\ans{ex:dec:B7}{$\vv A^+\vv b$ 满足正规方程（在 $\vv V$、$\vv U$ 的坐标下逐分量验证），且属于行空间 $\imm\vv A^\trans$，由第九章章末习题~\ref{ex:ip:B10} 它是范数最小的一个.}
\ans{ex:dec:B8}{$\vv J$ 的特征值为 $p,0,\ldots,0$，故矩阵的特征值为 $1+(p-1)\rho$ 与 $1-\rho$（$p-1$ 重）. 第一主成分 $\frac1{\sqrt p}(1,\ldots,1)^\trans$，比例 $\frac{1+(p-1)\rho}p$，$\rho\to1$ 时趋于 $1$，$\rho\to0$ 时趋于 $\frac1p$.}
\ans{ex:dec:B9}{协方差矩阵的第一主成分几乎就是第一个指标的方向（方差 $100$ 远大于 $1$）；相关系数为 $0.6$，相关系数矩阵的第一主成分是 $\frac1{\sqrt2}(1,1)^\trans$. 量纲改变方差，从而改变主成分.}
\ans{ex:dec:B10}{$\vv A-\vv A_k=\sum_{i>k}\sigma_i\vv u_i\vv v_i^\trans$，它的最大奇异值是 $\sigma_{k+1}$.}
\ans{ex:dec:B11}{$\vv v_i$ 是 $\vv A^\trans\vv A$ 的单重特征值的单位特征向量，只能差一个符号；$\vv u_i=\vv A\vv v_i/\sigma_i$ 随之确定.}
\ans{ex:dec:B12}{设 $\vv A=\begin{pmatrix}a&b\\c&d\end{pmatrix}$，比较 $\vv A^\trans\vv A$ 与 $\vv{AA}^\trans$ 的元素，得 $b^2=c^2$ 与 $(a-d)(b-c)=0$.}
\ans{ex:dec:B13}{第一句由第八章章末习题~\ref{ex:eig:B6}（特征值为实数）. 第二句取 $\vv x=\vv e_i$ 用定理~\ref{thm:rayleigh}.}
\ans{ex:dec:B14}{$k=1$ 时每份文书只剩一个数，查询与四份文书中三份合同类文书的相似度都是 $1$，与第四份为 $0$——同 $k=2$ 时的结果，但``合同''内部的差别完全消失，信息损失更多.}
\ans{ex:dec:B15}{$\vv x^\trans(\vv A-a\vv I)\vv x\ge0$ 对一切 $\vv x$ 成立当且仅当 $\vv A-a\vv I$ 的特征值 $\lambda_i-a\ge0$.}

\section{第十一章　二次型}

\subsection*{A 组答案}

\ans{ex:quad:A1}{(1) $\begin{pmatrix}1&2&0\\2&-3&-1\\0&-1&0\end{pmatrix}$，秩 $3$. (2) 对角元为 $0$、非对角元为 $\frac12$ 的三阶矩阵，秩 $3$. (3) 元素全为 $1$ 的三阶矩阵，秩 $1$. (4) 对角元为 $1$、非对角元为 $\frac12$ 的 $n$ 阶矩阵，即 $\frac12(\vv I+\vv J)$，秩 $n$（特征值 $\frac{n+1}2$ 与 $\frac12$，都不为零）.}

\ans{ex:quad:A2}{$x_1^2+2x_2^2-2x_1x_2+6x_2x_3$.}

\ans{ex:quad:A3}{(1) $(x_1+x_2)^2+(x_2+3x_3)^2-4x_3^2$，令 $y_1=x_1+x_2$，$y_2=x_2+3x_3$，$y_3=x_3$，得 $y_1^2+y_2^2-4y_3^2$. (2) $(x_1-2x_2)^2-x_3^2$，令 $y_1=x_1-2x_2$，$y_2=x_3$，$y_3=x_2$，得 $y_1^2-y_2^2$. (3) $x_2(x_1-x_3)$，令 $y_1=\frac12(x_1+x_2-x_3)$，$y_2=\frac12(x_2-x_1+x_3)$，$y_3=x_3$，得 $y_1^2-y_2^2$. 替换都不唯一.}

\ans{ex:quad:A4}{特征值 $2,5,1$，$\vv Q=\begin{pmatrix}1&0&0\\0&\frac1{\sqrt2}&\frac1{\sqrt2}\\0&\frac1{\sqrt2}&-\frac1{\sqrt2}\end{pmatrix}$，标准形 $2y_1^2+5y_2^2+y_3^2$. 曲面是椭球面.}

\ans{ex:quad:A5}{(1) $p=2$，$q=1$，符号差 $1$，规范形 $z_1^2+z_2^2-z_3^2$. (2)(3) 都是 $p=q=1$，符号差 $0$，规范形 $z_1^2-z_2^2$.}

\ans{ex:quad:A6}{(1) 正定（顺序主子式 $2,1,3$）. (2) 不正定（$\Delta_2=-3$）. (3) 正定（顺序主子式 $5,1,1$）.}

\ans{ex:quad:A7}{矩阵 $\begin{pmatrix}2&1&0\\1&1&\frac t2\\0&\frac t2&1\end{pmatrix}$，$\Delta_3=1-\frac{t^2}2>0$，即 $-\sqrt2<t<\sqrt2$.}

\ans{ex:quad:A8}{(1) 特征值 $3,-1$，标准方程 $3u^2-v^2=1$，双曲线；$u$ 轴沿 $\frac1{\sqrt2}(1,-1)$，$v$ 轴沿 $\frac1{\sqrt2}(1,1)$. (2) 特征值 $4,2$，标准方程 $\frac{u^2}2+\frac{v^2}4=1$，椭圆；沿 $\frac1{\sqrt2}(1,1)$ 的半轴为 $\sqrt2$，沿 $\frac1{\sqrt2}(1,-1)$ 的半轴为 $2$.}

\ans{ex:quad:A9}{(1) 合同（秩都为 $3$，$p$ 都为 $2$）：取 $\vv C$ 使 $\vv C\vv e_1=\frac1{\sqrt3}\vv e_3$，$\vv C\vv e_2=\sqrt5\,\vv e_1$，$\vv C\vv e_3=\sqrt3\,\vv e_2$. (2) 不合同：$p=1$ 与 $p=2$. (3) 合同：$\vv C=\frac1{\sqrt2}\begin{pmatrix}1&1\\1&-1\end{pmatrix}$.}

\ans{ex:quad:A10}{$\vv A^*=(\det\vv A)\vv A^{-1}$ 是对称矩阵，特征值为 $\frac{\det\vv A}{\lambda_i}>0$.}

\ans{ex:quad:A11}{驻点 $(0,0)$ 与 $(1,1)$. Hessian $\begin{pmatrix}6x&-3\\-3&6y\end{pmatrix}$：在 $(0,0)$ 处不定，是鞍点；在 $(1,1)$ 处正定，取极小值 $-1$.}

\ans{ex:quad:A12}{顺序主子式 $2,3,4$，正定. $\vv L=\begin{pmatrix}\sqrt2&0&0\\-\frac1{\sqrt2}&\sqrt{\frac32}&0\\0&-\sqrt{\frac23}&\frac2{\sqrt3}\end{pmatrix}$.}

\ans{ex:quad:A13}{$\vv L=\begin{pmatrix}2&0&0\\1&3&0\\-1&1&2\end{pmatrix}$.}

\ans{ex:quad:A14}{方差 $\vv w^\trans\vv\Sigma\vv w=0.0215$，标准差约 $14.7\%$.}

\ans{ex:quad:A15}{$\vv w^*\propto\vv\Sigma^{-1}\vv 1\propto\big(\frac1{0.09},\frac1{0.04}\big)$，即 $\vv w^*=\big(\frac4{13},\frac9{13}\big)$，方差 $\frac9{325}$，标准差约 $16.6\%$.}

\ans{ex:quad:A16}{$(\vv C^\trans\vv{AC})^\trans=\vv C^\trans\vv A^\trans\vv C$；秩不变见推论~\ref{cor:rankinvariant}. $\det(\vv C^\trans\vv{AC})=(\det\vv C)^2\det\vv A$：保持符号，不保持数值.}

\ans{ex:quad:A17}{在 $\mathbb{R}$ 上为 $z_1^2+z_2^2-z_3^2$；在 $\mathbb{C}$ 上为 $z_1^2+z_2^2+z_3^2$.}

\ans{ex:quad:A18}{依次为单叶双曲面（$p=2,q=1$）、双叶双曲面（$p=1,q=2$）、圆锥面、圆柱面（秩为 $2$）.}

\ans{ex:quad:A19}{特征值 $3$ 与 $-1$. 取 $\vv x=(1,-1)^\trans$，二次型的值为 $-2$.}

\ans{ex:quad:A20}{对主子矩阵 $\vv A_I$（行、列指标集都为 $I$），把 $\vv x_I\ne\zero$ 在其余位置补零，$\vv x_I^\trans\vv A_I\vv x_I=\vv x^\trans\vv A\vv x>0$. 正定矩阵的行列式为正（特征值之积）.}

\subsection*{B 组提示}

\ans{ex:quad:B1}{可交换的对称矩阵可同时正交对角化（第十章章末习题~\ref{ex:dec:B1}），对角元对应相乘. 反例：$\vv A=\begin{pmatrix}2&1\\1&1\end{pmatrix}$，$\vv B=\diag(1,2)$.}
\ans{ex:quad:B2}{取 $\vv C$ 使 $\vv C^\trans\vv{AC}=\vv I$，$\vv C^\trans\vv{BC}=\diag(\mu_i)$，$\mu_i\ge0$，则 $\det(\vv A+\vv B)=\det\vv A\prod(1+\mu_i)$.}
\ans{ex:quad:B3}{$\vv x^\trans(\vv A+t\vv B)\vv x\ge(\lambda_{\min}(\vv A)-|t|\max_i|\lambda_i(\vv B)|)\|\vv x\|^2$，用 B5.}
\ans{ex:quad:B4}{第 $i,j$ 行列组成的二阶主子式为正.}
\ans{ex:quad:B5}{由谱定理写成 $\sum\lambda_iy_i^2$，$\|\vv y\|=\|\vv x\|$.}
\ans{ex:quad:B6}{特征值 $2,-1,-1$，$p=1$，$q=2$. 配方时先作 $x_1=y_1+y_2$，$x_2=y_1-y_2$.}
\ans{ex:quad:B7}{$\Delta_1=\lambda$，$\Delta_2=\lambda^2-1$，$\Delta_3=(\lambda-2)(\lambda+1)^2$，故 $\lambda>2$.}
\ans{ex:quad:B8}{驻点条件 $2\vv{Ax}=\mu\vv c$，代入约束定出 $\mu$. $\vv A$ 正定保证驻点是最小值点（第十一章最小二乘例题中的配方法同样适用）.}
\ans{ex:quad:B9}{驻点条件给出 $\vv w=a\vv 1+b\vv\mu$，由两个约束解出 $a,b$（关于 $m$ 线性），方差 $\sigma^2\|\vv w\|^2$ 是 $m$ 的二次函数.}
\ans{ex:quad:B10}{解 $\det(\vv B-\lambda\vv A)=3\lambda^2-1=0$，得 $\lambda=\pm\frac1{\sqrt3}$；对每个 $\lambda$ 求 $(\vv B-\lambda\vv A)\vv x=\zero$ 的解并使 $\vv x^\trans\vv A\vv x=1$，以它们为列作 $\vv C$，则 $\vv C^\trans\vv{BC}=\diag(\frac1{\sqrt3},-\frac1{\sqrt3})$.}
\ans{ex:quad:B11}{配方：$f=(\vv x+\frac12\vv A^{-1}\vv b)^\trans\vv A(\vv x+\frac12\vv A^{-1}\vv b)+c-\frac14\vv b^\trans\vv A^{-1}\vv b$. 半正定时不一定有最小值，例如 $\vv A=\vv O$，$\vv b\ne\zero$.}
\ans{ex:quad:B12}{写 $\vv A=\vv B^\trans\vv B$（$\vv B=\vv D^{1/2}\vv Q^\trans$），则 $\|\vv B\vv x\|^2=0$，$\vv{Ax}=\vv B^\trans\vv B\vv x=\zero$.}
\ans{ex:quad:B13}{用定理~\ref{thm:posdef} 证明中的分块合同变换（把 $\vv a$ 换成矩阵块）.}
\ans{ex:quad:B14}{主轴定理把椭圆化为 $\lambda_1u^2+\lambda_2v^2\le1$，面积 $\frac\pi{\sqrt{\lambda_1\lambda_2}}$，而正交变换保持面积.}
\ans{ex:quad:B15}{$\vv x^\trans(\vv A^\trans\vv A+\mu\vv I)\vv x=\|\vv{Ax}\|^2+\mu\|\vv x\|^2$. 目标函数对 $\vv x$ 的驻点条件给出正规方程的``岭''版本.}
\ans{ex:quad:B16}{用谱分解，$\vv B=\vv Q\diag(\sqrt{\lambda_i})\vv Q^\trans$；$\vv A^{-1}$ 的特征值为 $\frac1{\lambda_i}$.}

\section{第十二章　*对偶空间与双线性型}

\subsection*{A 组答案}

\ans{ex:bil:A1}{(1) 是. (2) 不是：$2\vv x\mapsto4x_1x_2$. (3) 是. (4) 不是：$f\mapsto2f$ 时值变为 $4$ 倍. (5) 是.}

\ans{ex:bil:A2}{以三个基向量为列作 $\vv P$，对偶基是 $\vv P^{-1}$ 的各行：$\vv\varepsilon_1^*(\vv x)=x_1-x_2$，$\vv\varepsilon_2^*(\vv x)=x_2-x_3$，$\vv\varepsilon_3^*(\vv x)=x_3$.}

\ans{ex:bil:A3}{$f=f(\vv\varepsilon_1)\vv\varepsilon_1^*+f(\vv\varepsilon_2)\vv\varepsilon_2^*+f(\vv\varepsilon_3)\vv\varepsilon_3^*=\vv\varepsilon_1^*+3\vv\varepsilon_2^*+6\vv\varepsilon_3^*$.}

\ans{ex:bil:A4}{$W^0$ 是一维的，基为 $\vv x\mapsto x_1-x_2+x_3$.}

\ans{ex:bil:A5}{$\mathscr T^*$ 的矩阵为 $\begin{pmatrix}1&0&3\\2&1&-1\end{pmatrix}$；$\mathscr T^*$ 把 $\vv y\mapsto y_1+y_2+y_3$ 变为 $\vv x\mapsto4x_1+2x_2$.}

\ans{ex:bil:A6}{标准基下 $\vv A=\begin{pmatrix}1&2\\-1&3\end{pmatrix}$；新基下 $\vv B=\begin{pmatrix}5&-5\\1&3\end{pmatrix}$，可直接验证 $\vv B=\vv P^\trans\vv{AP}$，$\vv P=\begin{pmatrix}1&1\\1&-1\end{pmatrix}$.}

\ans{ex:bil:A7}{对称部分的度量矩阵 $\begin{pmatrix}1&\frac12\\ \frac12&3\end{pmatrix}$，交错部分 $\begin{pmatrix}0&\frac32\\-\frac32&0\end{pmatrix}$.}

\ans{ex:bil:A8}{(1)(2) 非退化（度量矩阵行列式 $-1$）；(3) 退化（度量矩阵全为 $1$，不可逆）.}

\ans{ex:bil:A9}{对应二次型 $x_1^2+2x_1x_2+3x_2^2$. 由二次型得 $f(\vv x,\vv y)=x_1y_1-2x_1y_2-2x_2y_1+2x_2y_2$.}

\ans{ex:bil:A10}{$f(\vv x,\vv x)=0$，$f(\vv x,\vv y)=-f(\vv y,\vv x)$. 度量矩阵 $\begin{pmatrix}0&1\\-1&0\end{pmatrix}$，可逆，非退化.}

\ans{ex:bil:A11}{Lagrange 基：$\frac12(x-1)(x-2)$，$-x(x-2)$，$\frac12x(x-1)$.}

\ans{ex:bil:A12}{$\dim W=2$，$\dim W^0=2$，基为 $\vv x\mapsto x_1+x_2$ 与 $\vv x\mapsto x_3-x_4$.}

\ans{ex:bil:A13}{双线性显然，对称性由 $\tr(\vv{AB})=\tr(\vv{BA})$；非退化：$\tr(\vv A\vv E_{ji})=a_{ij}$. 不是内积：$\vv A=\begin{pmatrix}0&1\\-1&0\end{pmatrix}$ 时 $\tr(\vv A^2)=-2<0$.}

\ans{ex:bil:A14}{坐标为 $\big(\varphi(1),\varphi(x),\varphi(x^2)\big)=\big(1,\frac12,\frac13\big)$.}

\ans{ex:bil:A15}{对任意 $\vv\alpha=\sum a_i\vv e_i$，$\langle\vv\alpha,\vv\beta\rangle=\sum a_if(\vv e_i)=f(\vv\alpha)$. 唯一性：若 $\langle\vv\alpha,\vv\beta-\vv\beta'\rangle=0$ 对一切 $\vv\alpha$ 成立，取 $\vv\alpha=\vv\beta-\vv\beta'$.}

\subsection*{B 组提示}

\ans{ex:bil:B1}{取 $V$ 的一组基，映射的矩阵的第 $i$ 行恰是 $f_i$ 在对偶基下的坐标. 满射当且仅当这个 $k\times n$ 矩阵行满秩，即各行线性无关.}
\ans{ex:bil:B2}{取 $\vv\alpha_0$ 使 $f(\vv\alpha_0)=1$，则 $V=\ker f\oplus\spn(\vv\alpha_0)$，$c=g(\vv\alpha_0)$. 推广：对映射 $\vv\alpha\mapsto(f_1(\vv\alpha),\ldots,f_k(\vv\alpha))$ 用同样的思路，或用零化子的维数.}
\ans{ex:bil:B3}{记 $U=\spn(f_1,\ldots,f_k)\subset V^*$. 把 $f_i$ 写成坐标后，$\ker f_1\cap\cdots\cap\ker f_k$ 就是以这些坐标为行的齐次方程组的解空间，维数为 $\dim V-\dim U$. 方程组语言：齐次方程组只有零解，当且仅当各方程的系数行向量张成整个行向量空间.}
\ans{ex:bil:B4}{由定理~\ref{thm:dualrank}：$\rk\mathscr T^*=\rk\mathscr T$.}
\ans{ex:bil:B5}{$a=c=\frac13$，$b=\frac43$. 积分是 $\mathbb{R}[x]_3$ 上的线性函数，三个求值函数构成 $\mathbb{R}[x]_3^*$ 的基（A11）. 三次时验证 $\int_0^2(x-1)^3\,\mathrm dx=0$ 与右边相等.}
\ans{ex:bil:B6}{双线性型由度量矩阵唯一确定，反之任一矩阵给出双线性型. 后半句即第三章章末习题~\ref{ex:mat:A22}.}
\ans{ex:bil:B7}{由定理~\ref{thm:alternating}，非退化时没有零块，$\dim V=2s$.}
\ans{ex:bil:B8}{度量矩阵 $\begin{pmatrix}1&1&0\\1&1&0\\0&0&1\end{pmatrix}$ 的核为 $\spn\big((1,-1,0)^\trans\big)$.}
\ans{ex:bil:B9}{迷向向量满足 $x_1^2=x_2^2$，即两条直线 $x_2=\pm x_1$. 后半句：设 $\vv P=\begin{pmatrix}a&b\\c&d\end{pmatrix}$，由条件得 $a^2-c^2=1$，$d^2-b^2=1$，$ab=cd$，再令 $a=\cosh t$.}
\ans{ex:bil:B10}{对 $\vv\alpha\in W$，$f\in W^0$，有 $(\mathscr T^*f)(\vv\alpha)=f(\mathscr T\vv\alpha)=0$.}
